{
\tolerance=4000
\hbadness=10000
\vbadness=10000
\subsection{BCU Instructions}
\label{sec:BCU-instructions}

\subsubsection[{\small AWAIT: Wait for any event or interrupt.}]{AWAIT\hfill Wait for any event or interrupt.}
\label{instr:AWAIT}\index{await}
 
\paragraph{Encoding} Format \textsf{BCU\_IPC} (Section~\ref{sec:format-BCU-IPC}), \verb|bcucode3 = 01000|\\

\bitfields{ 1/P, 2/00, 2/01, 4/1111, 5/01000, 6/-- -- -- -- -- --, 6/-- -- -- -- -- --, 6/-- -- -- -- -- -- }


\paragraph{Syntax} \texttt{await}

\paragraph{Description} Wait until \textbf{WS.WU0} equals one (cf Idle Modes chapter). The main core clock is gated during the waiting period to reduce power consumption.


\paragraph{Execution} {Reservation table \textsf{ALL} (Section~\ref{sec:reservation-ALL})}
~
{\begin{lstlisting}
stage RR:
  idle(0);
\end{lstlisting}}
\newpage
\subsubsection[{\small BARRIER: Instruction Pipeline Barrier}]{BARRIER\hfill Instruction Pipeline Barrier}
\label{instr:BARRIER}\index{barrier}
 
\paragraph{Encoding} Format \textsf{BCU\_IPC} (Section~\ref{sec:format-BCU-IPC}), \verb|bcucode3 = 01011|\\

\bitfields{ 1/P, 2/00, 2/01, 4/1111, 5/01011, 6/-- -- -- -- -- --, 6/-- -- -- -- -- --, 6/-- -- -- -- -- -- }


\paragraph{Syntax} \texttt{barrier}

\paragraph{Description} Flush the core instruction pipeline and memory system.


\paragraph{Execution} {Reservation table \textsf{ALL} (Section~\ref{sec:reservation-ALL})}
~
{\begin{lstlisting}
stage RR:
  barrier();
\end{lstlisting}}
\newpage
\subsubsection[{\small BLEND: Blended execution prefix}]{BLEND\hfill Blended execution prefix}
\label{instr:BLEND}\index{blend}
 
\paragraph{Encoding} Format \textsf{BCU\_BLEND} (Section~\ref{sec:format-BCU-BLEND}), \verb|bcucode3 = 00001|\\

\bitfields{ 1/P, 2/00, 2/01, 4/1111, 5/00001, 2/lanetodo, 2/lanesize, 8/activate, 6/registerZ }


\paragraph{Syntax} \texttt{blend}\emph{lanetodo}\emph{lanesize} \emph{registerZ}? \emph{activate}

\paragraph{Description} The results of the execution units specified by the bit mask \emph{registerZ} are blended under the mask \emph{lanetodo} applied to the \emph{lanesize}.


\paragraph{Modifier(s)}
\noindent\begin{itemize}
\item lanetodo (cf. lanetodo in section \ref{par:modifier-lanetodo} on page \pageref{par:modifier-lanetodo})
\item lanesize (cf. lanesize in section \ref{par:modifier-lanesize} on page \pageref{par:modifier-lanesize})
\end{itemize}

\paragraph{Execution} {Reservation table \textsf{BCU\_BRRP} (Section~\ref{sec:reservation-BCU-BRRP})}
~
{\begin{lstlisting}
stage ID:
  new lanetodo = _ZX_2(lanetodo);
  new lanesize = _ZX_2(lanesize);
  new activate = _ZX_6(activate);
stage RR:
  new argument = GPR[registerZ];
  blend(lanetodo, lanesize, argument, activate);
\end{lstlisting}}
\newpage
\subsubsection[{\small BREAK: Debugger Breakpoint}]{BREAK\hfill Debugger Breakpoint}
\label{instr:BREAK}\index{break}
 
\paragraph{Encoding} Format \textsf{BCU\_BREAK} (Section~\ref{sec:format-BCU-BREAK}), \verb|bcucode2 = 1111|\\

\bitfields{ 1/P, 2/00, 2/01, 4/1111, 5/11010, 6/-- -- -- -- -- --, 6/-- -- -- -- -- --, 4/-- -- -- --, 2/brknumber }


\paragraph{Syntax} \texttt{break} \emph{brknumber}

\paragraph{Description} Breakpoint for the debugger.


\paragraph{Execution} {Reservation table \textsf{ALL} (Section~\ref{sec:reservation-ALL})}
~
{\begin{lstlisting}
stage ID:
  new argument1 = _ZX_2(brknumber);
stage RR:
  break(argument1);
\end{lstlisting}}
\newpage
\subsubsection[{\small CALL: Call Subroutine PC-Relative}]{CALL\hfill Call Subroutine PC-Relative}
\label{instr:CALL}\index{call}
 
\paragraph{Encoding} Format \textsf{BCU\_UB} (Section~\ref{sec:format-BCU-UB}), \verb|bcucode1 = 11|\\

\bitfields{ 1/P, 2/00, 2/11, 27/pcrel27s2 }


\paragraph{Syntax} \texttt{call} \emph{pcrel27s2}

\paragraph{Description} The sequential next bundle \textbf{PC} is copied into the \textbf{RA} register. The next bundle \textbf{PC} is computed by adding to this instruction \textbf{PC} the \emph{pcrel27s2} after 27-bit sign extension and scaling by 4.


\paragraph{Execution} {Reservation table \textsf{BCU\_XFER} (Section~\ref{sec:reservation-BCU-XFER})}
~
{\begin{lstlisting}
stage ID:
  new argument1 = (_SX_27(pcrel27s2)<<2);
  new next = NPC;
  new address = PC + argument1;
  NPC = address;
  branch_info(1, address);
stage RR:
  RA = next;
\end{lstlisting}}
\newpage
\subsubsection[{\small CALLX: Call Subroutine Extended PC-Relative}]{CALLX\hfill Call Subroutine Extended PC-Relative}
\label{instr:CALLX}\index{callx}
 
\paragraph{Encoding} Format \textsf{BCU\_UB.X} (Section~\ref{sec:format-BCU-UB.X}), \verb|bcucode1 = 11|\\

\bitfields{ 1/P, 2/00, 2/00, 27/upper27, 1/1, 2/00, 2/11, 27/lower27 }


\paragraph{Syntax} \texttt{callx} \emph{upper27\_\discretionary{}{}{}lower27}

\paragraph{Description} The sequential next bundle \textbf{PC} is copied into the \textbf{RA} register. The next bundle \textbf{PC} is computed by adding to this instruction \textbf{PC} the \emph{upper27\_\discretionary{}{}{}lower27} after 54-bit sign extension and scaling by 4.


\paragraph{Execution} {Reservation table \textsf{BCU2.X} (Section~\ref{sec:reservation-BCU2.X})}
~
{\begin{lstlisting}
stage ID:
  new argument1 = (_SX_54(upper27_lower27)<<2);
  new next = NPC;
  new address = PC + argument1;
  NPC = address;
  branch_info(1, address);
stage RR:
  RA = next;
\end{lstlisting}}
\newpage
\subsubsection[{\small CB: Conditional Branch PC-Relative}]{CB\hfill Conditional Branch PC-Relative}
\label{instr:CB}\index{cb}
 
\paragraph{Encoding} Format \textsf{BCU\_CB} (Section~\ref{sec:format-BCU-CB}), \verb|bcucode1 = 01|\\

\bitfields{ 1/P, 2/00, 2/01, 4/bcucond, 17/pcrel17s2, 6/registerZ }


\paragraph{Syntax} \texttt{cb}\emph{bcucond} \emph{registerZ}? \emph{pcrel17s2}

\paragraph{Description} If \emph{registerZ} satisfies condition \emph{bcucond}, the next bundle \textbf{PC} is computed by adding to this instruction \textbf{PC} the \emph{pcrel17s2} after 17-bit sign extension and scaling by 4. \textbf{SRHPC} register may be updated according to current privilege level and owners configuration.


\paragraph{Modifier(s)}
\noindent\begin{itemize}
\item bcucond (cf. bcucond in section \ref{par:modifier-bcucond} on page \pageref{par:modifier-bcucond})
\end{itemize}

\paragraph{Execution} {Reservation table \textsf{BCU\_BRRP} (Section~\ref{sec:reservation-BCU-BRRP})}
~
{\begin{lstlisting}
stage RR:
  new argument1 = _ZX_4(bcucond);
  new argument2 = GPR[registerZ];
  new argument3 = (_SX_17(pcrel17s2)<<2);
  if (bcucond(argument1, argument2)) {
    new address = PC + argument3;
    NPC = address;
    branch_info(1, address);
    srhpc_update();
  } else {
    branch_info(0, 0);
  }
\end{lstlisting}}
\newpage
\subsubsection[{\small CBX: Conditional Branch Extended PC-Relative}]{CBX\hfill Conditional Branch Extended PC-Relative}
\label{instr:CBX}\index{cbx}
 
\paragraph{Encoding} Format \textsf{BCU\_CB.X} (Section~\ref{sec:format-BCU-CB.X}), \verb|bcucode1 = 01|\\

\bitfields{ 1/P, 2/00, 2/00, 27/upper27, 1/1, 2/00, 2/01, 4/bcucond, 17/lower17, 6/registerZ }


\paragraph{Syntax} \texttt{cbx}\emph{bcucond} \emph{registerZ}? \emph{upper27\_\discretionary{}{}{}lower17}

\paragraph{Description} If \emph{registerZ} satisfies condition \emph{bcucond}, the next bundle \textbf{PC} is computed by adding to this instruction \textbf{PC} the \emph{upper27\_\discretionary{}{}{}lower17} after 44-bit sign extension and scaling by 4. \textbf{SRHPC} register may be updated according to current privilege level and owners configuration.


\paragraph{Modifier(s)}
\noindent\begin{itemize}
\item bcucond (cf. bcucond in section \ref{par:modifier-bcucond} on page \pageref{par:modifier-bcucond})
\end{itemize}

\paragraph{Execution} {Reservation table \textsf{BCU2.X} (Section~\ref{sec:reservation-BCU2.X})}
~
{\begin{lstlisting}
stage RR:
  new argument1 = _ZX_4(bcucond);
  new argument2 = GPR[registerZ];
  new argument3 = (_SX_44(upper27_lower17)<<2);
  if (bcucond(argument1, argument2)) {
    new address = PC + argument3;
    NPC = address;
    branch_info(1, address);
    srhpc_update();
  } else {
    branch_info(0, 0);
  }
\end{lstlisting}}
\newpage
\subsubsection[{\small CCB: Compare Conditional Branch PC-Relative}]{CCB\hfill Compare Conditional Branch PC-Relative}
\label{instr:CCB}\index{ccb}
 
\paragraph{Encoding} Format \textsf{BCU\_CCB} (Section~\ref{sec:format-BCU-CCB}), \verb|bcucode1 = 00|\\

\bitfields{ 1/P, 2/00, 2/00, 4/ccbcomp, 11/pcrel11s2, 6/registerY, 6/registerZ }


\paragraph{Syntax} \texttt{ccb}\emph{ccbcomp} \emph{registerZ}, \emph{registerY} ? \emph{pcrel11s2}

\paragraph{Description} If \emph{registerZ} and \emph{registerY} satisfy condition \emph{ccbcomp}, the next bundle \textbf{PC} is computed by adding to this instruction \textbf{PC} the \emph{pcrel11s2} after 11-bit sign extension and scaling by 4. \textbf{SRHPC} register may be updated according to current privilege level and owners configuration.


\paragraph{Modifier(s)}
\noindent\begin{itemize}
\item ccbcomp (cf. ccbcomp in section \ref{par:modifier-ccbcomp} on page \pageref{par:modifier-ccbcomp})
\end{itemize}

\paragraph{Execution} {Reservation table \textsf{BCU\_BRRP2} (Section~\ref{sec:reservation-BCU-BRRP2})}
~
{\begin{lstlisting}
stage ID:
  if (ccbcomp(argument1, argument2, argument3)) {
    new address = PC + argument4;
    NPC = address;
    branch_info(1, address);
    srhpc_update();
  } else {
    branch_info(0, 0);
  }
stage RR:
  new argument1 = _ZX_4(ccbcomp);
  new argument2 = GPR[registerZ];
  new argument3 = GPR[registerY];
  new argument4 = (_SX_11(pcrel11s2)<<2);
\end{lstlisting}}
\newpage
\subsubsection[{\small CCBX: Compare Conditional Extended Branch PC-Relative}]{CCBX\hfill Compare Conditional Extended Branch PC-Relative}
\label{instr:CCBX}\index{ccbx}
 
\paragraph{Encoding} Format \textsf{BCU\_CCB.X} (Section~\ref{sec:format-BCU-CCB.X}), \verb|bcucode1 = 00|\\

\bitfields{ 1/P, 2/00, 2/00, 27/upper27, 1/1, 2/00, 2/00, 4/ccbcomp, 11/lower11, 6/registerY, 6/registerZ }


\paragraph{Syntax} \texttt{ccbx}\emph{ccbcomp} \emph{registerZ}, \emph{registerY} ? \emph{upper27\_\discretionary{}{}{}lower11}

\paragraph{Description} If \emph{registerZ} and \emph{registerY} satisfy condition \emph{ccbcomp}, the next bundle \textbf{PC} is computed by adding to this instruction \textbf{PC} the \emph{upper27\_\discretionary{}{}{}lower11} after 38-bit sign extension and scaling by 4. \textbf{SRHPC} register may be updated according to current privilege level and owners configuration.


\paragraph{Modifier(s)}
\noindent\begin{itemize}
\item ccbcomp (cf. ccbcomp in section \ref{par:modifier-ccbcomp} on page \pageref{par:modifier-ccbcomp})
\end{itemize}

\paragraph{Execution} {Reservation table \textsf{BCU2.X} (Section~\ref{sec:reservation-BCU2.X})}
~
{\begin{lstlisting}
stage ID:
  if (ccbcomp(argument1, argument2, argument3)) {
    new address = PC + argument4;
    NPC = address;
    branch_info(1, address);
    srhpc_update();
  } else {
    branch_info(0, 0);
  }
stage RR:
  new argument1 = _ZX_4(ccbcomp);
  new argument2 = GPR[registerZ];
  new argument3 = GPR[registerY];
  new argument4 = (_SX_38(upper27_lower11)<<2);
\end{lstlisting}}
\newpage
\subsubsection[{\small GET: Get System Register}]{GET\hfill Get System Register}
\label{instr:GET}\index{get}
 
\paragraph{Encoding} Format \textsf{BCU\_GSR} (Section~\ref{sec:format-BCU-GSR}), \verb|bcucode3 = 10001|\\

\bitfields{ 1/P, 2/00, 2/01, 4/1111, 5/10001, 3/-- -- --, 9/systemS2, 6/registerZ }


\paragraph{Syntax} \texttt{get} \emph{registerZ} = \emph{systemS2}

\paragraph{Description} The \emph{registerZ} is copied from the \emph{systemS2} register.


\paragraph{Execution} {Reservation table \textsf{BCU2\_TINY\_LSU} (Section~\ref{sec:reservation-BCU2-TINY-LSU})}
~
{\begin{lstlisting}
stage RR:
  new argument2 = SFR[systemS2];
  if (get_check_access(systemS2, registerZ)) {
    new result1 = get(systemS2, _ZX_64(argument2));
  }
stage E1:
  GPR[registerZ] = result1;
\end{lstlisting}}
\newpage
\subsubsection[{\small GOTO: Branch Unconditional PC-Relative}]{GOTO\hfill Branch Unconditional PC-Relative}
\label{instr:GOTO}\index{goto}
 
\paragraph{Encoding} Format \textsf{BCU\_UB} (Section~\ref{sec:format-BCU-UB}), \verb|bcucode1 = 10|\\

\bitfields{ 1/P, 2/00, 2/10, 27/pcrel27s2 }


\paragraph{Syntax} \texttt{goto} \emph{pcrel27s2}

\paragraph{Description} The next bundle \textbf{PC} is computed by adding to this instruction \textbf{PC} the \emph{pcrel27s2} after 27-bit sign extension and scaling by 4. \textbf{SRHPC} register may be updated according to current privilege level and owners configuration.


\paragraph{Execution} {Reservation table \textsf{BCU\_XFER} (Section~\ref{sec:reservation-BCU-XFER})}
~
{\begin{lstlisting}
stage ID:
  new argument1 = (_SX_27(pcrel27s2)<<2);
  new address = PC + argument1;
  NPC = address;
  branch_info(1, address);
stage RR:
  srhpc_update();
\end{lstlisting}}
\newpage
\subsubsection[{\small GOTOX: Branch Unconditional Extended PC-Relative}]{GOTOX\hfill Branch Unconditional Extended PC-Relative}
\label{instr:GOTOX}\index{gotox}
 
\paragraph{Encoding} Format \textsf{BCU\_UB.X} (Section~\ref{sec:format-BCU-UB.X}), \verb|bcucode1 = 10|\\

\bitfields{ 1/P, 2/00, 2/00, 27/upper27, 1/1, 2/00, 2/10, 27/lower27 }


\paragraph{Syntax} \texttt{gotox} \emph{upper27\_\discretionary{}{}{}lower27}

\paragraph{Description} The next bundle \textbf{PC} is computed by adding to this instruction \textbf{PC} the \emph{upper27\_\discretionary{}{}{}lower27} after 54-bit sign extension and scaling by 4. \textbf{SRHPC} register may be updated according to current privilege level and owners configuration.


\paragraph{Execution} {Reservation table \textsf{BCU2.X} (Section~\ref{sec:reservation-BCU2.X})}
~
{\begin{lstlisting}
stage ID:
  new argument1 = (_SX_54(upper27_lower27)<<2);
  new address = PC + argument1;
  NPC = address;
  branch_info(1, address);
stage RR:
  srhpc_update();
\end{lstlisting}}
\newpage
\subsubsection[{\small GUARD: Guarded execution prefix.}]{GUARD\hfill Guarded execution prefix.}
\label{instr:GUARD}\index{guard}
 
\paragraph{Encoding} Format \textsf{BCU\_GUARD} (Section~\ref{sec:format-BCU-GUARD}), \verb|bcucode3 = 00000|\\

\bitfields{ 1/P, 2/00, 2/01, 4/1111, 5/00000, 4/execpred, 8/activate, 6/registerZ }


\paragraph{Syntax} \texttt{guard}\emph{execpred} \emph{registerZ}? \emph{activate}

\paragraph{Description} The execution units specified by the bit mask \emph{activate} are guarded by the condition \emph{execpred} applied to the \emph{registerZ}.


\paragraph{Modifier(s)}
\noindent\begin{itemize}
\item execpred (cf. bcucond in section \ref{par:modifier-bcucond} on page \pageref{par:modifier-bcucond})
\end{itemize}

\paragraph{Execution} {Reservation table \textsf{BCU\_BRRP} (Section~\ref{sec:reservation-BCU-BRRP})}
~
{\begin{lstlisting}
stage ID:
  new bcucond = _ZX_4(execpred);
  new activate = _ZX_6(activate);
stage RR:
  new argument = GPR[registerZ];
  guard(bcucond, argument, activate);
\end{lstlisting}}
\newpage
\subsubsection[{\small ICALL: Indirect Call from Register}]{ICALL\hfill Indirect Call from Register}
\label{instr:ICALL}\index{icall}
 
\paragraph{Encoding} Format \textsf{BCU\_IBC} (Section~\ref{sec:format-BCU-IBC}), \verb|bcucode3 = 10111|\\

\bitfields{ 1/P, 2/00, 2/01, 4/1111, 5/10111, 6/-- -- -- -- -- --, 6/-- -- -- -- -- --, 6/registerZ }


\paragraph{Syntax} \texttt{icall} \emph{registerZ}

\paragraph{Description} The sequential next bundle \textbf{PC} is copied into the \textbf{RA} register. The next bundle \textbf{PC} is computed by reading the \emph{registerZ}.


\paragraph{Execution} {Reservation table \textsf{BCU\_XFER\_BRRP} (Section~\ref{sec:reservation-BCU-XFER-BRRP})}
~
{\begin{lstlisting}
stage ID:
  new argument1 = GPR[registerZ];
  new next = NPC;
stage RR:
  NPC = argument1;
  RA = next;
  branch_info(1, argument1);
\end{lstlisting}}
\newpage
\subsubsection[{\small IGET: Indirect Get System Register}]{IGET\hfill Indirect Get System Register}
\label{instr:IGET}\index{iget}
 
\paragraph{Encoding} Format \textsf{BCU\_IGSR} (Section~\ref{sec:format-BCU-IGSR}), \verb|bcucode3 = 10011|\\

\bitfields{ 1/P, 2/00, 2/01, 4/1111, 5/10011, 6/-- -- -- -- -- --, 6/-- -- -- -- -- --, 6/registerZ }


\paragraph{Syntax} \texttt{iget} \emph{registerZ}

\paragraph{Description} The \emph{registerZ} is copied from the system register whose index is in the \emph{registerZ}.


\paragraph{Execution} {Reservation table \textsf{BCU2\_TINY\_LSU} (Section~\ref{sec:reservation-BCU2-TINY-LSU})}
~
{\begin{lstlisting}
stage ID:
  new index = _ZX_9(GPR[registerZ]);
stage RR:
  if (get_check_access(index, registerZ)) {
    new result1 = get(index, SFR[index]);
  }
stage E1:
  GPR[registerZ] = result1;
\end{lstlisting}}
\newpage
\subsubsection[{\small IGOTO: Indirect Branch from Register}]{IGOTO\hfill Indirect Branch from Register}
\label{instr:IGOTO}\index{igoto}
 
\paragraph{Encoding} Format \textsf{BCU\_IBC} (Section~\ref{sec:format-BCU-IBC}), \verb|bcucode3 = 10110|\\

\bitfields{ 1/P, 2/00, 2/01, 4/1111, 5/10110, 6/-- -- -- -- -- --, 6/-- -- -- -- -- --, 6/registerZ }


\paragraph{Syntax} \texttt{igoto} \emph{registerZ}

\paragraph{Description} The next bundle \textbf{PC} is computed by reading the \emph{registerZ}. \textbf{SRHPC} register may be updated according to current privilege level and owners configuration.


\paragraph{Execution} {Reservation table \textsf{BCU\_XFER\_BRRP} (Section~\ref{sec:reservation-BCU-XFER-BRRP})}
~
{\begin{lstlisting}
stage ID:
  new argument1 = GPR[registerZ];
stage RR:
  NPC = argument1;
  branch_info(1, argument1);
  srhpc_update();
\end{lstlisting}}
\newpage
\subsubsection[{\small LOOPDO: Hardware Loop setup and Do execute}]{LOOPDO\hfill Hardware Loop setup and Do execute}
\label{instr:LOOPDO}\index{loopdo}
 
\paragraph{Encoding} Format \textsf{BCU\_HLS} (Section~\ref{sec:format-BCU-HLS}), \verb|bcucode2 = 1110|\\

\bitfields{ 1/P, 2/00, 2/01, 4/1110, 17/pcrel17s2, 6/registerZ }


\paragraph{Syntax} \texttt{loopdo} \emph{registerZ}, \emph{pcrel17s2}

\paragraph{Description} If \textbf{PS.HLE} is cleared, this instruction throws an OPCODE trap. The \textbf{LS} register is set to the sequential next bundle \textbf{PC}. The \textbf{LE} register is set to the current bundle \textbf{PC} plus the \emph{pcrel17s2} after sign extension and scaling by 4. The \textbf{LC} register is set with the \emph{registerZ}. If the \textbf{PS.HLE} bit is set in the \textbf{PS} register, hardware loop execution proceeds by comparing the sequential next bundle \textbf{PC} to the value of the \textbf{LE} register: \begin{itemize} \item If different, the next bundle \textbf{PC} is computed as usual. \item If the sequential next bundle \textbf{PC} equals the value of the \textbf{LE} register: \begin{itemize}
  \item If the value of the \textbf{LC} register ${}-1$ is non-zero, the value of the next bundle \textbf{PC}
  is read from the \textbf{LS} register, unless a branch is taken in the current bundle (in which case the next bundle \textbf{PC} is the branch target bundle \textbf{PC}).
  \item If the value of the \textbf{LC} register ${}-1$ is zero, the value of the next bundle \textbf{PC} is
  the value of the \textbf{LE} register, unless a branch is taken in the current bundle (in which case the next bundle \textbf{PC} is the branch target bundle \textbf{PC}).
  \item The \textbf{LC} register is decremented if its value is not already zero.
  \end{itemize}
\end{itemize} According to these rules, the LOOPDO instruction loops forever whenever the value of the \emph{registerZ} is zero, as the the value of the \textbf{LC} register ${}-1$ is always non-zero, while the \textbf{LC} register is never decremented. This instruction is provided for counted loops which iterate at least once, and while loops by setting the counter to zero.


\paragraph{Execution} {Reservation table \textsf{ALL} (Section~\ref{sec:reservation-ALL})}
~
{\begin{lstlisting}
stage ID:
  new argument1 = GPR[registerZ];
  new argument2 = (_SX_17(pcrel17s2)<<2);
  if (!PS.HLE) {
    _THROW(OPCODE);
  }
  invalpfb();
  LC = argument1;
stage E2:
  LS = NPC;
  LE = PC + _SX_32(argument2);
\end{lstlisting}}
\newpage
\subsubsection[{\small RET: Return from Call}]{RET\hfill Return from Call}
\label{instr:RET}\index{ret}
 
\paragraph{Encoding} Format \textsf{BCU\_RTS} (Section~\ref{sec:format-BCU-RTS}), \verb|bcucode3 = 10100|\\

\bitfields{ 1/P, 2/00, 2/01, 4/1111, 5/10100, 6/-- -- -- -- -- --, 6/-- -- -- -- -- --, 6/-- -- -- -- -- -- }


\paragraph{Syntax} \texttt{ret}

\paragraph{Description} The next bundle \textbf{PC} is obtained by reading the contents of the \textbf{RA} register. \textbf{SRHPC} register may be updated according to current privilege level and owners configuration.


\paragraph{Execution} {Reservation table \textsf{BCU\_XFER} (Section~\ref{sec:reservation-BCU-XFER})}
~
{\begin{lstlisting}
stage ID:
  NPC = RA;
stage RR:
  branch_info(1, RA);
  srhpc_update();
\end{lstlisting}}
\newpage
\subsubsection[{\small RFE: Return from Exception}]{RFE\hfill Return from Exception}
\label{instr:RFE}\index{rfe}
 
\paragraph{Encoding} Format \textsf{BCU\_RTS} (Section~\ref{sec:format-BCU-RTS}), \verb|bcucode3 = 10101|\\

\bitfields{ 1/P, 2/00, 2/01, 4/1111, 5/10101, 6/-- -- -- -- -- --, 6/-- -- -- -- -- --, 6/-- -- -- -- -- -- }


\paragraph{Syntax} \texttt{rfe}

\paragraph{Description} The \textbf{PS} is restored from the \textbf{SPS}. The \textbf{SPS} is restored from the \textbf{SSPS}. The next bundle \textbf{PC} is obtained by reading the contents of the \textbf{SPC} register. The \textbf{SPC} is restored from the \textbf{SSPC}. \textbf{SRHPC} register may be updated according to current privilege level and owners configuration.


\paragraph{Execution} {Reservation table \textsf{ALL} (Section~\ref{sec:reservation-ALL})}
~
{\begin{lstlisting}
stage RR:
  if (rfe_owner()) {
    rfe();
    branch_info(1, SPC);
    srhpc_update();
  }
\end{lstlisting}}
\newpage
\subsubsection[{\small RSWAP: General and System Register Swap}]{RSWAP\hfill General and System Register Swap}
\label{instr:RSWAP}\index{rswap}
 
\paragraph{Encoding} Format \textsf{BCU\_RSWAP} (Section~\ref{sec:format-BCU-RSWAP}), \verb|bcucode3 = 10010|\\

\bitfields{ 1/P, 2/00, 2/01, 4/1111, 5/10010, 3/-- -- --, 9/systemS4, 6/registerZ }


\paragraph{Syntax} \texttt{rswap} \emph{registerZ} = \emph{systemS4}

\paragraph{Description} The contents of the \emph{registerZ} are swapped set with the contents of the \emph{systemS4}.


\paragraph{Execution} {Reservation table \textsf{BCU2\_TINY\_LSU} (Section~\ref{sec:reservation-BCU2-TINY-LSU})}
~
{\begin{lstlisting}
stage RR:
  new argument2 = SFR[systemS4];
  new argument1 = GPR[registerZ];
stage E1:
  if (get_check_access(systemS4, registerZ) && set_check_access(systemS4, _ZX_64(argument1), registerZ)) {
    new result2 = _ZX_64(argument1);
    new result1 = _ZX_64(argument2);
    GPR[registerZ] = result1;
    SFR[systemS4] = result2;
  }
\end{lstlisting}}
\newpage

\paragraph{Encoding} Format \textsf{BCU\_RSWAPA} (Section~\ref{sec:format-BCU-RSWAPA}), \verb|bcucode3 = 10010|\\

\bitfields{ 1/P, 2/00, 2/01, 4/1111, 5/10010, 3/-- -- --, 9/systemAlone, 6/registerZ }


\paragraph{Syntax} \texttt{rswap} \emph{registerZ} = \emph{systemAlone}

\paragraph{Description} The contents of the \emph{registerZ} are swapped set with the contents of the \emph{systemAlone}.


\paragraph{Execution} {Reservation table \textsf{ALL} (Section~\ref{sec:reservation-ALL})}
~
{\begin{lstlisting}
stage RR:
  new argument2 = SFR[systemAlone];
  new argument1 = GPR[registerZ];
stage E1:
  if (get_check_access(systemAlone, registerZ) && set_check_access(systemAlone, _ZX_64(argument1), registerZ)) {
    new result2 = _ZX_64(argument1);
    new result1 = _ZX_64(argument2);
    GPR[registerZ] = result1;
    SFR[systemAlone] = result2;
  }
\end{lstlisting}}
\newpage
\subsubsection[{\small SCALL: System Call}]{SCALL\hfill System Call}
\label{instr:SCALL}\index{scall}
 
\paragraph{Encoding} Format \textsf{BCU\_SCR} (Section~\ref{sec:format-BCU-SCR}), \verb|bcucode2 = 1111|\\

\bitfields{ 1/P, 2/00, 2/01, 4/1111, 5/11001, 6/-- -- -- -- -- --, 6/-- -- -- -- -- --, 6/registerZ }


\paragraph{Syntax} \texttt{scall} \emph{registerZ}

\paragraph{Description} System call with number \emph{}. Interrupts are disabled.


\paragraph{Execution} {Reservation table \textsf{ALL} (Section~\ref{sec:reservation-ALL})}
~
{\begin{lstlisting}
stage ID:
  new argument1 = GPR[registerZ];
  syscall(argument1);
  branch_info(1, argument1);
\end{lstlisting}}
\newpage

\paragraph{Encoding} Format \textsf{BCU\_SCI} (Section~\ref{sec:format-BCU-SCI}), \verb|bcucode2 = 1111|\\

\bitfields{ 1/P, 2/00, 2/01, 4/1111, 5/11000, 6/-- -- -- -- -- --, 12/sysnumber }


\paragraph{Syntax} \texttt{scall} \emph{sysnumber}

\paragraph{Description} System call with number \emph{}. Interrupts are disabled.


\paragraph{Execution} {Reservation table \textsf{ALL} (Section~\ref{sec:reservation-ALL})}
~
{\begin{lstlisting}
stage ID:
  new argument1 = _ZX_12(sysnumber);
  syscall(argument1);
  branch_info(1, argument1);
\end{lstlisting}}
\newpage
\subsubsection[{\small SET: Set System Register}]{SET\hfill Set System Register}
\label{instr:SET}\index{set}
 
\paragraph{Encoding} Format \textsf{BCU\_SET} (Section~\ref{sec:format-BCU-SET}), \verb|bcucode3 = 10000|\\

\bitfields{ 1/P, 2/00, 2/01, 4/1111, 5/10000, 3/-- -- --, 9/systemT3, 6/registerZ }


\paragraph{Syntax} \texttt{set} \emph{systemT3} = \emph{registerZ}

\paragraph{Description} The \emph{systemT3} register is set with the \emph{registerZ}.


\paragraph{Execution} {Reservation table \textsf{BCU2} (Section~\ref{sec:reservation-BCU2})}
~
{\begin{lstlisting}
stage ID:
  new argument2 = GPR[registerZ];
stage E3:
  if (set_check_access(systemT3, _ZX_64(argument2), registerZ)) {
    new result1 = _ZX_64(argument2);
    SFR[systemT3] = result1;
  }
\end{lstlisting}}
\newpage

\paragraph{Encoding} Format \textsf{BCU\_SETA} (Section~\ref{sec:format-BCU-SETA}), \verb|bcucode3 = 10000|\\

\bitfields{ 1/P, 2/00, 2/01, 4/1111, 5/10000, 3/-- -- --, 9/systemAlone, 6/registerZ }


\paragraph{Syntax} \texttt{set} \emph{systemAlone} = \emph{registerZ}

\paragraph{Description} The \emph{systemAlone} register is set with the \emph{registerZ}.


\paragraph{Execution} {Reservation table \textsf{ALL} (Section~\ref{sec:reservation-ALL})}
~
{\begin{lstlisting}
stage ID:
  new argument2 = GPR[registerZ];
stage E3:
  if (set_check_access(systemAlone, _ZX_64(argument2), registerZ)) {
    new result1 = _ZX_64(argument2);
    SFR[systemAlone] = result1;
  }
\end{lstlisting}}
\newpage

\paragraph{Encoding} Format \textsf{BCU\_SETRA} (Section~\ref{sec:format-BCU-SETRA}), \verb|bcucode3 = 10000|\\

\bitfields{ 1/P, 2/00, 2/01, 4/1111, 5/10000, 3/-- -- --, 9/systemRA, 6/registerZ }


\paragraph{Syntax} \texttt{set} \emph{systemRA} = \emph{registerZ}

\paragraph{Description} The \emph{systemRA} register is set with the \emph{registerZ}.


\paragraph{Execution} {Reservation table \textsf{BCU2} (Section~\ref{sec:reservation-BCU2})}
~
{\begin{lstlisting}
stage ID:
  new argument2 = GPR[registerZ];
stage E3:
  if (set_check_access(systemRA, _ZX_64(argument2), registerZ)) {
    new result1 = _ZX_64(argument2);
    SFR[systemRA] = result1;
  }
\end{lstlisting}}
\newpage
\subsubsection[{\small SLEEP: Wait for an eligible interrupt.}]{SLEEP\hfill Wait for an eligible interrupt.}
\label{instr:SLEEP}\index{sleep}
 
\paragraph{Encoding} Format \textsf{BCU\_IPC} (Section~\ref{sec:format-BCU-IPC}), \verb|bcucode3 = 01001|\\

\bitfields{ 1/P, 2/00, 2/01, 4/1111, 5/01001, 6/-- -- -- -- -- --, 6/-- -- -- -- -- --, 6/-- -- -- -- -- -- }


\paragraph{Syntax} \texttt{sleep}

\paragraph{Description} Wait until \textbf{WS.WU1} equals one (cf Idle Modes chapter). The main core clock is gated during the waiting period to reduce power consumption.


\paragraph{Execution} {Reservation table \textsf{ALL} (Section~\ref{sec:reservation-ALL})}
~
{\begin{lstlisting}
stage RR:
  idle(1);
\end{lstlisting}}
\newpage
\subsubsection[{\small STOP: Wait for power controller wake-up.}]{STOP\hfill Wait for power controller wake-up.}
\label{instr:STOP}\index{stop}
 
\paragraph{Encoding} Format \textsf{BCU\_IPC} (Section~\ref{sec:format-BCU-IPC}), \verb|bcucode3 = 01010|\\

\bitfields{ 1/P, 2/00, 2/01, 4/1111, 5/01010, 6/-- -- -- -- -- --, 6/-- -- -- -- -- --, 6/-- -- -- -- -- -- }


\paragraph{Syntax} \texttt{stop}

\paragraph{Description} Wait until \textbf{WS.WU2} equals one (cf Idle Modes chapter). The main core clock is gated during the waiting period to reduce power consumption.


\paragraph{Execution} {Reservation table \textsf{ALL} (Section~\ref{sec:reservation-ALL})}
~
{\begin{lstlisting}
stage RR:
  if (stop_owner()) {
    idle(2);
  } else {
    _THROW(PRIVILEGE);
  }
\end{lstlisting}}
\newpage
\subsubsection[{\small SYNCGROUP: Synchronize Group of Cores}]{SYNCGROUP\hfill Synchronize Group of Cores}
\label{instr:SYNCGROUP}\index{syncgroup}
 
\paragraph{Encoding} Format \textsf{BCU\_PGI} (Section~\ref{sec:format-BCU-PGI}), \verb|bcucode3 = 01101|\\

\bitfields{ 1/P, 2/00, 2/01, 4/1111, 5/01101, 6/-- -- -- -- -- --, 6/-- -- -- -- -- --, 6/registerZ }


\paragraph{Syntax} \texttt{syncgroup} \emph{registerZ}

\paragraph{Description} The \emph{registerZ} is interpreted as four 16-bit fields: \textbf{waitclrF}, \textbf{waitclrB}, \textbf{notifyF}, \textbf{notifyB}. Upon executing the instruction, the core waits until all the events designated by bits in \textbf{waitclrF} and \textbf{waitclrB} in the 32 least-significant bits of \textbf{IPE} are set. When all the designated events are set, the core clears them, then notifies the event lines corresponding to the bit sets in fields \textbf{notifyF} and and \textbf{notifyB}. This notification of events is not registered in the core, it is a signal sent to set bits into the \textbf{IPE.FE} and \textbf{IPE.BE} fields of the destination cores. The successor core always receives the \textbf{notifyF} events in its \textbf{IPE.FE} field; likewise, the predecessor core always receives the \textbf{notifyB} events in its \textbf{IPE.BE} field. Furthermore, other PEs may receive these events, depending on the setting of the \textbf{IPE.FM} and \textbf{IPE.BM} fields in these neighbor PE cores.


\paragraph{Execution} {Reservation table \textsf{BCU2} (Section~\ref{sec:reservation-BCU2})}
~
{\begin{lstlisting}
stage ID:
  new argument1 = GPR[registerZ];
stage RR:
  if (!syncgroup_owner()) {
    _THROW(PRIVILEGE);
  } else {
    new waitclrF = argument1 & 0xFFFF;
    new waitclrB = (argument1 >> 16) & 0xFFFF;
    new notifyF = (argument1 >> 32) & 0xFFFF;
    new notifyB = (argument1 >> 48) & 0xFFFF;
    syncgroup(waitclrF, waitclrB, notifyF, notifyB);
  }
\end{lstlisting}}
\newpage
\subsubsection[{\small TLBDINVAL: Invalidate Data Micro-TLB.}]{TLBDINVAL\hfill Invalidate Data Micro-TLB.}
\label{instr:TLBDINVAL}\index{tlbdinval}
 
\paragraph{Encoding} Format \textsf{BCU\_TLB} (Section~\ref{sec:format-BCU-TLB}), \verb|bcucode3 = 00110|\\

\bitfields{ 1/P, 2/00, 2/01, 4/1111, 5/00110, 6/-- -- -- -- -- --, 6/-- -- -- -- -- --, 6/-- -- -- -- -- -- }


\paragraph{Syntax} \texttt{tlbdinval}

\paragraph{Description} Invalidate Data Micro-TLB.


\paragraph{Execution} {Reservation table \textsf{ALL} (Section~\ref{sec:reservation-ALL})}
~
{\begin{lstlisting}
stage RR:
  invaldtlb();
\end{lstlisting}}
\newpage
\subsubsection[{\small TLBIINVAL: Invalidate Instruction Micro-TLB.}]{TLBIINVAL\hfill Invalidate Instruction Micro-TLB.}
\label{instr:TLBIINVAL}\index{tlbiinval}
 
\paragraph{Encoding} Format \textsf{BCU\_TLB} (Section~\ref{sec:format-BCU-TLB}), \verb|bcucode3 = 00111|\\

\bitfields{ 1/P, 2/00, 2/01, 4/1111, 5/00111, 6/-- -- -- -- -- --, 6/-- -- -- -- -- --, 6/-- -- -- -- -- -- }


\paragraph{Syntax} \texttt{tlbiinval}

\paragraph{Description} Invalidate Instruction Micro-TLB.


\paragraph{Execution} {Reservation table \textsf{ALL} (Section~\ref{sec:reservation-ALL})}
~
{\begin{lstlisting}
stage RR:
  invalitlb();
\end{lstlisting}}
\newpage
\subsubsection[{\small TLBPROBE: Probe TLB}]{TLBPROBE\hfill Probe TLB}
\label{instr:TLBPROBE}\index{tlbprobe}
 
\paragraph{Encoding} Format \textsf{BCU\_TLB} (Section~\ref{sec:format-BCU-TLB}), \verb|bcucode3 = 00011|\\

\bitfields{ 1/P, 2/00, 2/01, 4/1111, 5/00011, 6/-- -- -- -- -- --, 6/-- -- -- -- -- --, 6/-- -- -- -- -- -- }


\paragraph{Syntax} \texttt{tlbprobe}

\paragraph{Description} Probe the TLB for virtual address match in the TEL:TEH register pair and write to MMC.IDX.


\paragraph{Execution} {Reservation table \textsf{ALL} (Section~\ref{sec:reservation-ALL})}
~
{\begin{lstlisting}
stage ID:
  if (!mmi_owner()) {
    _THROW(PRIVILEGE);
  } else {
    probetlb();
  }
\end{lstlisting}}
\newpage
\subsubsection[{\small TLBREAD: Read TLB Entry}]{TLBREAD\hfill Read TLB Entry}
\label{instr:TLBREAD}\index{tlbread}
 
\paragraph{Encoding} Format \textsf{BCU\_TLB} (Section~\ref{sec:format-BCU-TLB}), \verb|bcucode3 = 00010|\\

\bitfields{ 1/P, 2/00, 2/01, 4/1111, 5/00010, 6/-- -- -- -- -- --, 6/-- -- -- -- -- --, 6/-- -- -- -- -- -- }


\paragraph{Syntax} \texttt{tlbread}

\paragraph{Description} Read the TLB entry into the TEL:TEH register pair using MMC.IDX.


\paragraph{Execution} {Reservation table \textsf{ALL} (Section~\ref{sec:reservation-ALL})}
~
{\begin{lstlisting}
stage ID:
  if (!mmi_owner()) {
    _THROW(PRIVILEGE);
  } else {
    readtlb();
  }
\end{lstlisting}}
\newpage
\subsubsection[{\small TLBWRITE: TLB Entry Write}]{TLBWRITE\hfill TLB Entry Write}
\label{instr:TLBWRITE}\index{tlbwrite}
 
\paragraph{Encoding} Format \textsf{BCU\_TLB} (Section~\ref{sec:format-BCU-TLB}), \verb|bcucode3 = 00100|\\

\bitfields{ 1/P, 2/00, 2/01, 4/1111, 5/00100, 6/-- -- -- -- -- --, 6/-- -- -- -- -- --, 6/-- -- -- -- -- -- }


\paragraph{Syntax} \texttt{tlbwrite}

\paragraph{Description} Write the TLB entry from the TEL:TEH register pair using MMC.IDX.


\paragraph{Execution} {Reservation table \textsf{ALL} (Section~\ref{sec:reservation-ALL})}
~
{\begin{lstlisting}
stage ID:
  if (!mmi_owner()) {
    _THROW(PRIVILEGE);
  } else {
    writetlb();
  }
\end{lstlisting}}
\newpage
\subsubsection[{\small WAITIT: Wait for Interrupt}]{WAITIT\hfill Wait for Interrupt}
\label{instr:WAITIT}\index{waitit}
 
\paragraph{Encoding} Format \textsf{BCU\_PGI} (Section~\ref{sec:format-BCU-PGI}), \verb|bcucode3 = 01100|\\

\bitfields{ 1/P, 2/00, 2/01, 4/1111, 5/01100, 6/-- -- -- -- -- --, 6/-- -- -- -- -- --, 6/registerZ }


\paragraph{Syntax} \texttt{waitit} \emph{registerZ}

\paragraph{Description} Execution waits for selected interrupts in the \textbf{ILR} register (masking non-owned interrupts, i.e. interrupts that belong to a more privileged PL. The \emph{registerZ} is interpreted as an OR-mask in the least significant word and a AND-mask in the most significant word. The wait ends if all the bits in \textbf{ILR} match the non-zero AND-mask, or if any bit in \textbf{ILR} matches the non-zero OR-mask. When the wait is over, the value of \textbf{ILR} (masked in the same way) is written to the \emph{registerZ}. When both the AND-mask and OR-mask are zero, no wait happens.


\paragraph{Execution} {Reservation table \textsf{BCU2\_TINY\_LSU} (Section~\ref{sec:reservation-BCU2-TINY-LSU})}
~
{\begin{lstlisting}
stage ID:
  new argument1 = GPR[registerZ];
stage RR:
  new ormask = argument1 & 0xFFFFFFFF;
  new andmask = argument1.32[1];
  new result1 = waitit(ormask, andmask);
  GPR[registerZ] = result1;
\end{lstlisting}}
\newpage
\subsubsection[{\small WFXL: Effects on Least Significant Word System Register}]{WFXL\hfill Effects on Least Significant Word System Register}
\label{instr:WFXL}\index{wfxl}
 
\paragraph{Encoding} Format \textsf{BCU\_WFX} (Section~\ref{sec:format-BCU-WFX}), \verb|bcucode3 = 01110|\\

\bitfields{ 1/P, 2/00, 2/01, 4/1111, 5/01110, 3/-- -- --, 9/systemT2, 6/registerZ }


\paragraph{Syntax} \texttt{wfxl} \emph{systemT2}, \emph{registerZ}

\paragraph{Description} The least significant word of the \emph{systemT2} is operated, with the most significant word of the \emph{registerZ} as a set mask and the least significant word of the \emph{registerZ} as a clear mask.


\paragraph{Execution} {Reservation table \textsf{BCU2} (Section~\ref{sec:reservation-BCU2})}
~
{\begin{lstlisting}
stage ID:
  new argument2 = GPR[registerZ];
stage E4:
  new argument1 = SFR[systemT2];
  if (wfxl_check_access(systemT2, registerZ)) {
    new result1 = wfxl(systemT2, _ZX_64(argument2));
  }
  SFR[systemT2] = result1;
\end{lstlisting}}
\newpage

\paragraph{Encoding} Format \textsf{BCU\_WFXA} (Section~\ref{sec:format-BCU-WFXA}), \verb|bcucode3 = 01110|\\

\bitfields{ 1/P, 2/00, 2/01, 4/1111, 5/01110, 3/-- -- --, 9/systemAlone, 6/registerZ }


\paragraph{Syntax} \texttt{wfxl} \emph{systemAlone}, \emph{registerZ}

\paragraph{Description} The least significant word of the \emph{systemAlone} is operated, with the most significant word of the \emph{registerZ} as a set mask and the least significant word of the \emph{registerZ} as a clear mask.


\paragraph{Execution} {Reservation table \textsf{ALL} (Section~\ref{sec:reservation-ALL})}
~
{\begin{lstlisting}
stage ID:
  new argument1 = SFR[systemAlone];
  new argument2 = GPR[registerZ];
stage RR:
  if (wfxl_check_access(systemAlone, registerZ)) {
    new result1 = wfxl(systemAlone, _ZX_64(argument2));
  }
stage E4:
  SFR[systemAlone] = result1;
\end{lstlisting}}
\newpage
\subsubsection[{\small WFXM: Word Effects on Most Significant Word System Register}]{WFXM\hfill Word Effects on Most Significant Word System Register}
\label{instr:WFXM}\index{wfxm}
 
\paragraph{Encoding} Format \textsf{BCU\_WFX} (Section~\ref{sec:format-BCU-WFX}), \verb|bcucode3 = 01111|\\

\bitfields{ 1/P, 2/00, 2/01, 4/1111, 5/01111, 3/-- -- --, 9/systemT2, 6/registerZ }


\paragraph{Syntax} \texttt{wfxm} \emph{systemT2}, \emph{registerZ}

\paragraph{Description} The most significant word of the \emph{systemT2} is operated, with the most significant word of the \emph{registerZ} as a set mask and the least significant word of the \emph{registerZ} as a clear mask.


\paragraph{Execution} {Reservation table \textsf{BCU2} (Section~\ref{sec:reservation-BCU2})}
~
{\begin{lstlisting}
stage ID:
  new argument2 = GPR[registerZ];
stage E4:
  new argument1 = SFR[systemT2];
  if (wfxm_check_access(systemT2, registerZ)) {
    new result1 = wfxm(systemT2, _ZX_64(argument2));
  }
  SFR[systemT2] = result1;
\end{lstlisting}}
\newpage

\paragraph{Encoding} Format \textsf{BCU\_WFXA} (Section~\ref{sec:format-BCU-WFXA}), \verb|bcucode3 = 01111|\\

\bitfields{ 1/P, 2/00, 2/01, 4/1111, 5/01111, 3/-- -- --, 9/systemAlone, 6/registerZ }


\paragraph{Syntax} \texttt{wfxm} \emph{systemAlone}, \emph{registerZ}

\paragraph{Description} The most significant word of the \emph{systemAlone} is operated, with the most significant word of the \emph{registerZ} as a set mask and the least significant word of the \emph{registerZ} as a clear mask.


\paragraph{Execution} {Reservation table \textsf{ALL} (Section~\ref{sec:reservation-ALL})}
~
{\begin{lstlisting}
stage ID:
  new argument1 = SFR[systemAlone];
  new argument2 = GPR[registerZ];
stage RR:
  if (wfxm_check_access(systemAlone, registerZ)) {
    new result1 = wfxm(systemAlone, _ZX_64(argument2));
  }
stage E4:
  SFR[systemAlone] = result1;
\end{lstlisting}}
\newpage
\subsection{LSU Instructions}
\label{sec:LSU-instructions}

\subsubsection[{\small ACSWAPB: Atomic Compare and Swap Byte}]{ACSWAPB\hfill Atomic Compare and Swap Byte}
\label{instr:ACSWAPB}\index{acswapb}
 
\paragraph{Encoding} Format \textsf{LSU\_ACSWAPSB} (Section~\ref{sec:format-LSU-ACSWAPSB}), \verb|lsucode5 = 000|\\

\bitfields{ 1/P, 2/01, 3/111, 2/coherency, 6/registerW, 2/11, 3/000, 1/boolcas, 5/registerO, 1/1, 6/registerZ }


\paragraph{Syntax} \texttt{acswapb}\emph{boolcas}\emph{coherency} \emph{registerW}, [\emph{registerZ}] = \emph{registerO}

\paragraph{Description} The effective address is the value of the \emph{registerZ}.
This instruction implements Compare-And-Swap (CAS) on memory byte (8 bits). The \emph{registerO} is interpreted as the concatenation of \texttt{update} (low 64 bits) and \texttt{expected} (high 64 bits). The memory byte at the effective address is read into \texttt{current}. If \texttt{current} equals \texttt{expected}, the swap is successful and \texttt{update} is stored into the memory byte at the effective address. Otherwise, the swap has failed and the memory is not updated. The instruction return value is stored into the \emph{registerW}. It is true/false in case \texttt{boolcas} is true, or the \texttt{current} value in case \texttt{boolcas} is false.


\paragraph{Modifier(s)}
\noindent\begin{itemize}
\item coherency (cf. coherency in section \ref{par:modifier-coherency} on page \pageref{par:modifier-coherency})
\item boolcas (cf. boolcas in section \ref{par:modifier-boolcas} on page \pageref{par:modifier-boolcas})
\end{itemize}

\paragraph{Execution} {Reservation table \textsf{LSU2\_MEMW\_AUXR\_AUXW} (Section~\ref{sec:reservation-LSU2-MEMW-AUXR-AUXW})}
~
{\begin{lstlisting}
stage ID:
  new boolcas = _ZX_1(boolcas);
  new coherency = _ZX_2(coherency);
stage RR:
  new address = GPR[registerZ];
stage E1:
  new arguments = PGR[registerO];
  new result1 = MEM.atomic_cas(address, 0x1, boolcas | (coherency << 32), arguments.8[0], arguments.8[8], registerW);
stage SR:
  GPR[registerW] = result1;
\end{lstlisting}}
\newpage

\paragraph{Encoding} Format \textsf{LSU\_ACSWAPSB.O} (Section~\ref{sec:format-LSU-ACSWAPSB.O}), \verb|lsucode5 = 000|\\

\bitfields{ 1/P, 2/00, 2/-- --, 27/offset27, 1/1, 2/01, 3/111, 2/coherency, 6/registerW, 2/11, 3/000, 1/boolcas, 5/registerO, 1/1, 6/registerZ }


\paragraph{Syntax} \texttt{acswapb}\emph{boolcas}\emph{coherency} \emph{registerW}, \emph{offset27}[\emph{registerZ}] = \emph{registerO}

\paragraph{Description} The effective address is computed by adding the \emph{offset27} to the \emph{registerZ}.
This instruction implements Compare-And-Swap (CAS) on memory byte (8 bits). The \emph{registerO} is interpreted as the concatenation of \texttt{update} (low 64 bits) and \texttt{expected} (high 64 bits). The memory byte at the effective address is read into \texttt{current}. If \texttt{current} equals \texttt{expected}, the swap is successful and \texttt{update} is stored into the memory byte at the effective address. Otherwise, the swap has failed and the memory is not updated. The instruction return value is stored into the \emph{registerW}. It is true/false in case \texttt{boolcas} is true, or the \texttt{current} value in case \texttt{boolcas} is false.


\paragraph{Modifier(s)}
\noindent\begin{itemize}
\item coherency (cf. coherency in section \ref{par:modifier-coherency} on page \pageref{par:modifier-coherency})
\item boolcas (cf. boolcas in section \ref{par:modifier-boolcas} on page \pageref{par:modifier-boolcas})
\end{itemize}

\paragraph{Execution} {Reservation table \textsf{LSU2\_MEMW\_AUXR\_AUXW.X} (Section~\ref{sec:reservation-LSU2-MEMW-AUXR-AUXW.X})}
~
{\begin{lstlisting}
stage ID:
  new offset = _SX_27(offset27);
  new boolcas = _ZX_1(boolcas);
  new coherency = _ZX_2(coherency);
stage RR:
  new base = GPR[registerZ];
  new address = base + offset;
stage E1:
  new arguments = PGR[registerO];
  new result1 = MEM.atomic_cas(address, 0x1, boolcas | (coherency << 32), arguments.8[0], arguments.8[8], registerW);
stage SR:
  GPR[registerW] = result1;
\end{lstlisting}}
\newpage

\paragraph{Encoding} Format \textsf{LSU\_ACSWAPSB.D} (Section~\ref{sec:format-LSU-ACSWAPSB.D}), \verb|lsucode5 = 000|\\

\bitfields{ 1/P, 2/00, 2/-- --, 27/extend27, 1/1, 2/00, 2/-- --, 27/offset27, 1/1, 2/01, 3/111, 2/coherency, 6/registerW, 2/11, 3/000, 1/boolcas, 5/registerO, 1/1, 6/registerZ }


\paragraph{Syntax} \texttt{acswapb}\emph{boolcas}\emph{coherency} \emph{registerW}, \emph{extend27\_\discretionary{}{}{}offset27}[\emph{registerZ}] = \emph{registerO}

\paragraph{Description} The effective address is computed by adding the \emph{extend27\_\discretionary{}{}{}offset27} to the \emph{registerZ}.
This instruction implements Compare-And-Swap (CAS) on memory byte (8 bits). The \emph{registerO} is interpreted as the concatenation of \texttt{update} (low 64 bits) and \texttt{expected} (high 64 bits). The memory byte at the effective address is read into \texttt{current}. If \texttt{current} equals \texttt{expected}, the swap is successful and \texttt{update} is stored into the memory byte at the effective address. Otherwise, the swap has failed and the memory is not updated. The instruction return value is stored into the \emph{registerW}. It is true/false in case \texttt{boolcas} is true, or the \texttt{current} value in case \texttt{boolcas} is false.


\paragraph{Modifier(s)}
\noindent\begin{itemize}
\item coherency (cf. coherency in section \ref{par:modifier-coherency} on page \pageref{par:modifier-coherency})
\item boolcas (cf. boolcas in section \ref{par:modifier-boolcas} on page \pageref{par:modifier-boolcas})
\end{itemize}

\paragraph{Execution} {Reservation table \textsf{LSU2\_MEMW\_AUXR\_AUXW.Y} (Section~\ref{sec:reservation-LSU2-MEMW-AUXR-AUXW.Y})}
~
{\begin{lstlisting}
stage ID:
  new offset = _SX_54(extend27_offset27);
  new boolcas = _ZX_1(boolcas);
  new coherency = _ZX_2(coherency);
stage RR:
  new base = GPR[registerZ];
  new address = base + offset;
stage E1:
  new arguments = PGR[registerO];
  new result1 = MEM.atomic_cas(address, 0x1, boolcas | (coherency << 32), arguments.8[0], arguments.8[8], registerW);
stage SR:
  GPR[registerW] = result1;
\end{lstlisting}}
\newpage
\subsubsection[{\small ACSWAPD: Atomic Compare and Swap Double Word}]{ACSWAPD\hfill Atomic Compare and Swap Double Word}
\label{instr:ACSWAPD}\index{acswapd}
 
\paragraph{Encoding} Format \textsf{LSU\_ACSWAPSB} (Section~\ref{sec:format-LSU-ACSWAPSB}), \verb|lsucode5 = 011|\\

\bitfields{ 1/P, 2/01, 3/111, 2/coherency, 6/registerW, 2/11, 3/011, 1/boolcas, 5/registerO, 1/1, 6/registerZ }


\paragraph{Syntax} \texttt{acswapd}\emph{boolcas}\emph{coherency} \emph{registerW}, [\emph{registerZ}] = \emph{registerO}

\paragraph{Description} The effective address is the value of the \emph{registerZ}.
This instruction implements Compare-And-Swap (CAS) on memory double word (64 bits). The \emph{registerO} is interpreted as the concatenation of \texttt{update} (low 64 bits) and \texttt{expected} (high 64 bits). The memory double word at the effective address is read into \texttt{current}. If \texttt{current} equals \texttt{expected}, the swap is successful and \texttt{update} is stored into the memory double word at the effective address. Otherwise, the swap has failed and the memory is not updated. The instruction return value is stored into the \emph{registerW}. It is true/false in case \texttt{boolcas} is true, or the \texttt{current} value in case \texttt{boolcas} is false. The effective address must be double word aligned (multiple of 8).


\paragraph{Modifier(s)}
\noindent\begin{itemize}
\item coherency (cf. coherency in section \ref{par:modifier-coherency} on page \pageref{par:modifier-coherency})
\item boolcas (cf. boolcas in section \ref{par:modifier-boolcas} on page \pageref{par:modifier-boolcas})
\end{itemize}

\paragraph{Execution} {Reservation table \textsf{LSU2\_MEMW\_AUXR\_AUXW} (Section~\ref{sec:reservation-LSU2-MEMW-AUXR-AUXW})}
~
{\begin{lstlisting}
stage ID:
  new boolcas = _ZX_1(boolcas);
  new coherency = _ZX_2(coherency);
stage RR:
  new address = GPR[registerZ];
stage E1:
  new arguments = PGR[registerO];
  new result1 = MEM.atomic_cas(address, 0xFF, boolcas | (coherency << 32), arguments.64[0], arguments.64[1], registerW);
stage SR:
  GPR[registerW] = result1;
\end{lstlisting}}
\newpage

\paragraph{Encoding} Format \textsf{LSU\_ACSWAPSB.O} (Section~\ref{sec:format-LSU-ACSWAPSB.O}), \verb|lsucode5 = 011|\\

\bitfields{ 1/P, 2/00, 2/-- --, 27/offset27, 1/1, 2/01, 3/111, 2/coherency, 6/registerW, 2/11, 3/011, 1/boolcas, 5/registerO, 1/1, 6/registerZ }


\paragraph{Syntax} \texttt{acswapd}\emph{boolcas}\emph{coherency} \emph{registerW}, \emph{offset27}[\emph{registerZ}] = \emph{registerO}

\paragraph{Description} The effective address is computed by adding the \emph{offset27} to the \emph{registerZ}.
This instruction implements Compare-And-Swap (CAS) on memory double word (64 bits). The \emph{registerO} is interpreted as the concatenation of \texttt{update} (low 64 bits) and \texttt{expected} (high 64 bits). The memory double word at the effective address is read into \texttt{current}. If \texttt{current} equals \texttt{expected}, the swap is successful and \texttt{update} is stored into the memory double word at the effective address. Otherwise, the swap has failed and the memory is not updated. The instruction return value is stored into the \emph{registerW}. It is true/false in case \texttt{boolcas} is true, or the \texttt{current} value in case \texttt{boolcas} is false. The effective address must be double word aligned (multiple of 8).


\paragraph{Modifier(s)}
\noindent\begin{itemize}
\item coherency (cf. coherency in section \ref{par:modifier-coherency} on page \pageref{par:modifier-coherency})
\item boolcas (cf. boolcas in section \ref{par:modifier-boolcas} on page \pageref{par:modifier-boolcas})
\end{itemize}

\paragraph{Execution} {Reservation table \textsf{LSU2\_MEMW\_AUXR\_AUXW.X} (Section~\ref{sec:reservation-LSU2-MEMW-AUXR-AUXW.X})}
~
{\begin{lstlisting}
stage ID:
  new offset = _SX_27(offset27);
  new boolcas = _ZX_1(boolcas);
  new coherency = _ZX_2(coherency);
stage RR:
  new base = GPR[registerZ];
  new address = base + offset;
stage E1:
  new arguments = PGR[registerO];
  new result1 = MEM.atomic_cas(address, 0xFF, boolcas | (coherency << 32), arguments.64[0], arguments.64[1], registerW);
stage SR:
  GPR[registerW] = result1;
\end{lstlisting}}
\newpage

\paragraph{Encoding} Format \textsf{LSU\_ACSWAPSB.D} (Section~\ref{sec:format-LSU-ACSWAPSB.D}), \verb|lsucode5 = 011|\\

\bitfields{ 1/P, 2/00, 2/-- --, 27/extend27, 1/1, 2/00, 2/-- --, 27/offset27, 1/1, 2/01, 3/111, 2/coherency, 6/registerW, 2/11, 3/011, 1/boolcas, 5/registerO, 1/1, 6/registerZ }


\paragraph{Syntax} \texttt{acswapd}\emph{boolcas}\emph{coherency} \emph{registerW}, \emph{extend27\_\discretionary{}{}{}offset27}[\emph{registerZ}] = \emph{registerO}

\paragraph{Description} The effective address is computed by adding the \emph{extend27\_\discretionary{}{}{}offset27} to the \emph{registerZ}.
This instruction implements Compare-And-Swap (CAS) on memory double word (64 bits). The \emph{registerO} is interpreted as the concatenation of \texttt{update} (low 64 bits) and \texttt{expected} (high 64 bits). The memory double word at the effective address is read into \texttt{current}. If \texttt{current} equals \texttt{expected}, the swap is successful and \texttt{update} is stored into the memory double word at the effective address. Otherwise, the swap has failed and the memory is not updated. The instruction return value is stored into the \emph{registerW}. It is true/false in case \texttt{boolcas} is true, or the \texttt{current} value in case \texttt{boolcas} is false. The effective address must be double word aligned (multiple of 8).


\paragraph{Modifier(s)}
\noindent\begin{itemize}
\item coherency (cf. coherency in section \ref{par:modifier-coherency} on page \pageref{par:modifier-coherency})
\item boolcas (cf. boolcas in section \ref{par:modifier-boolcas} on page \pageref{par:modifier-boolcas})
\end{itemize}

\paragraph{Execution} {Reservation table \textsf{LSU2\_MEMW\_AUXR\_AUXW.Y} (Section~\ref{sec:reservation-LSU2-MEMW-AUXR-AUXW.Y})}
~
{\begin{lstlisting}
stage ID:
  new offset = _SX_54(extend27_offset27);
  new boolcas = _ZX_1(boolcas);
  new coherency = _ZX_2(coherency);
stage RR:
  new base = GPR[registerZ];
  new address = base + offset;
stage E1:
  new arguments = PGR[registerO];
  new result1 = MEM.atomic_cas(address, 0xFF, boolcas | (coherency << 32), arguments.64[0], arguments.64[1], registerW);
stage SR:
  GPR[registerW] = result1;
\end{lstlisting}}
\newpage
\subsubsection[{\small ACSWAPH: Atomic Compare and Swap Half Word}]{ACSWAPH\hfill Atomic Compare and Swap Half Word}
\label{instr:ACSWAPH}\index{acswaph}
 
\paragraph{Encoding} Format \textsf{LSU\_ACSWAPSB} (Section~\ref{sec:format-LSU-ACSWAPSB}), \verb|lsucode5 = 001|\\

\bitfields{ 1/P, 2/01, 3/111, 2/coherency, 6/registerW, 2/11, 3/001, 1/boolcas, 5/registerO, 1/1, 6/registerZ }


\paragraph{Syntax} \texttt{acswaph}\emph{boolcas}\emph{coherency} \emph{registerW}, [\emph{registerZ}] = \emph{registerO}

\paragraph{Description} The effective address is the value of the \emph{registerZ}.
This instruction implements Compare-And-Swap (CAS) on memory half word (16 bits). The \emph{registerO} is interpreted as the concatenation of \texttt{update} (low 64 bits) and \texttt{expected} (high 64 bits). The memory half word at the effective address is read into \texttt{current}. If \texttt{current} equals \texttt{expected}, the swap is successful and \texttt{update} is stored into the memory half word at the effective address. Otherwise, the swap has failed and the memory is not updated. The instruction return value is stored into the \emph{registerW}. It is true/false in case \texttt{boolcas} is true, or the \texttt{current} value in case \texttt{boolcas} is false. The effective address must be half word aligned (multiple of 2).


\paragraph{Modifier(s)}
\noindent\begin{itemize}
\item coherency (cf. coherency in section \ref{par:modifier-coherency} on page \pageref{par:modifier-coherency})
\item boolcas (cf. boolcas in section \ref{par:modifier-boolcas} on page \pageref{par:modifier-boolcas})
\end{itemize}

\paragraph{Execution} {Reservation table \textsf{LSU2\_MEMW\_AUXR\_AUXW} (Section~\ref{sec:reservation-LSU2-MEMW-AUXR-AUXW})}
~
{\begin{lstlisting}
stage ID:
  new boolcas = _ZX_1(boolcas);
  new coherency = _ZX_2(coherency);
stage RR:
  new address = GPR[registerZ];
stage E1:
  new arguments = PGR[registerO];
  new result1 = MEM.atomic_cas(address, 0x3, boolcas | (coherency << 32), arguments.16[0], arguments.16[4], registerW);
stage SR:
  GPR[registerW] = result1;
\end{lstlisting}}
\newpage

\paragraph{Encoding} Format \textsf{LSU\_ACSWAPSB.O} (Section~\ref{sec:format-LSU-ACSWAPSB.O}), \verb|lsucode5 = 001|\\

\bitfields{ 1/P, 2/00, 2/-- --, 27/offset27, 1/1, 2/01, 3/111, 2/coherency, 6/registerW, 2/11, 3/001, 1/boolcas, 5/registerO, 1/1, 6/registerZ }


\paragraph{Syntax} \texttt{acswaph}\emph{boolcas}\emph{coherency} \emph{registerW}, \emph{offset27}[\emph{registerZ}] = \emph{registerO}

\paragraph{Description} The effective address is computed by adding the \emph{offset27} to the \emph{registerZ}.
This instruction implements Compare-And-Swap (CAS) on memory half word (16 bits). The \emph{registerO} is interpreted as the concatenation of \texttt{update} (low 64 bits) and \texttt{expected} (high 64 bits). The memory half word at the effective address is read into \texttt{current}. If \texttt{current} equals \texttt{expected}, the swap is successful and \texttt{update} is stored into the memory half word at the effective address. Otherwise, the swap has failed and the memory is not updated. The instruction return value is stored into the \emph{registerW}. It is true/false in case \texttt{boolcas} is true, or the \texttt{current} value in case \texttt{boolcas} is false. The effective address must be half word aligned (multiple of 2).


\paragraph{Modifier(s)}
\noindent\begin{itemize}
\item coherency (cf. coherency in section \ref{par:modifier-coherency} on page \pageref{par:modifier-coherency})
\item boolcas (cf. boolcas in section \ref{par:modifier-boolcas} on page \pageref{par:modifier-boolcas})
\end{itemize}

\paragraph{Execution} {Reservation table \textsf{LSU2\_MEMW\_AUXR\_AUXW.X} (Section~\ref{sec:reservation-LSU2-MEMW-AUXR-AUXW.X})}
~
{\begin{lstlisting}
stage ID:
  new offset = _SX_27(offset27);
  new boolcas = _ZX_1(boolcas);
  new coherency = _ZX_2(coherency);
stage RR:
  new base = GPR[registerZ];
  new address = base + offset;
stage E1:
  new arguments = PGR[registerO];
  new result1 = MEM.atomic_cas(address, 0x3, boolcas | (coherency << 32), arguments.16[0], arguments.16[4], registerW);
stage SR:
  GPR[registerW] = result1;
\end{lstlisting}}
\newpage

\paragraph{Encoding} Format \textsf{LSU\_ACSWAPSB.D} (Section~\ref{sec:format-LSU-ACSWAPSB.D}), \verb|lsucode5 = 001|\\

\bitfields{ 1/P, 2/00, 2/-- --, 27/extend27, 1/1, 2/00, 2/-- --, 27/offset27, 1/1, 2/01, 3/111, 2/coherency, 6/registerW, 2/11, 3/001, 1/boolcas, 5/registerO, 1/1, 6/registerZ }


\paragraph{Syntax} \texttt{acswaph}\emph{boolcas}\emph{coherency} \emph{registerW}, \emph{extend27\_\discretionary{}{}{}offset27}[\emph{registerZ}] = \emph{registerO}

\paragraph{Description} The effective address is computed by adding the \emph{extend27\_\discretionary{}{}{}offset27} to the \emph{registerZ}.
This instruction implements Compare-And-Swap (CAS) on memory half word (16 bits). The \emph{registerO} is interpreted as the concatenation of \texttt{update} (low 64 bits) and \texttt{expected} (high 64 bits). The memory half word at the effective address is read into \texttt{current}. If \texttt{current} equals \texttt{expected}, the swap is successful and \texttt{update} is stored into the memory half word at the effective address. Otherwise, the swap has failed and the memory is not updated. The instruction return value is stored into the \emph{registerW}. It is true/false in case \texttt{boolcas} is true, or the \texttt{current} value in case \texttt{boolcas} is false. The effective address must be half word aligned (multiple of 2).


\paragraph{Modifier(s)}
\noindent\begin{itemize}
\item coherency (cf. coherency in section \ref{par:modifier-coherency} on page \pageref{par:modifier-coherency})
\item boolcas (cf. boolcas in section \ref{par:modifier-boolcas} on page \pageref{par:modifier-boolcas})
\end{itemize}

\paragraph{Execution} {Reservation table \textsf{LSU2\_MEMW\_AUXR\_AUXW.Y} (Section~\ref{sec:reservation-LSU2-MEMW-AUXR-AUXW.Y})}
~
{\begin{lstlisting}
stage ID:
  new offset = _SX_54(extend27_offset27);
  new boolcas = _ZX_1(boolcas);
  new coherency = _ZX_2(coherency);
stage RR:
  new base = GPR[registerZ];
  new address = base + offset;
stage E1:
  new arguments = PGR[registerO];
  new result1 = MEM.atomic_cas(address, 0x3, boolcas | (coherency << 32), arguments.16[0], arguments.16[4], registerW);
stage SR:
  GPR[registerW] = result1;
\end{lstlisting}}
\newpage
\subsubsection[{\small ACSWAPQ: Atomic Compare and Swap Quadruple Word}]{ACSWAPQ\hfill Atomic Compare and Swap Quadruple Word}
\label{instr:ACSWAPQ}\index{acswapq}
 
\paragraph{Encoding} Format \textsf{LSU\_ACSWAPQB} (Section~\ref{sec:format-LSU-ACSWAPQB}), \verb|lsucode5 = 100|\\

\bitfields{ 1/P, 2/01, 3/111, 2/coherency, 5/registerM, 1/0, 2/11, 3/100, 1/boolcas, 4/registerQ, 1/--, 1/1, 6/registerZ }


\paragraph{Syntax} \texttt{acswapq}\emph{boolcas}\emph{coherency} \emph{registerM}, [\emph{registerZ}] = \emph{registerQ}

\paragraph{Description} The effective address is the value of the \emph{registerZ}.
This instruction implements Compare-And-Swap (CAS) on memory quadruple word (128 bits). The \emph{registerQ} is interpreted as the concatenation of \texttt{update} (low 128 bits) and \texttt{expected} (high 128 bits). The memory double word at the effective address is read into \texttt{current}. If \texttt{current} equals \texttt{expected}, the swap is successful and \texttt{update} is stored into the memory double word at the effective address. Otherwise, the swap has failed and the memory is not updated. The instruction return value is stored into the \emph{registerM}. It is true/false in case \texttt{boolcas} is true, or the \texttt{current} value in case \texttt{boolcas} is false. The effective address must be quadruple word aligned (multiple of 16).


\paragraph{Modifier(s)}
\noindent\begin{itemize}
\item coherency (cf. coherency in section \ref{par:modifier-coherency} on page \pageref{par:modifier-coherency})
\item boolcas (cf. boolcas in section \ref{par:modifier-boolcas} on page \pageref{par:modifier-boolcas})
\end{itemize}

\paragraph{Execution} {Reservation table \textsf{LSU2\_MEMW\_AUXR\_AUXW} (Section~\ref{sec:reservation-LSU2-MEMW-AUXR-AUXW})}
~
{\begin{lstlisting}
stage ID:
  new boolcas = _ZX_1(boolcas);
  new coherency = _ZX_2(coherency);
stage RR:
  new address = GPR[registerZ];
stage E1:
  new arguments = QGR[registerQ];
  new result1 = MEM.atomic_cas(address, 0xFFFF, boolcas | (coherency << 32), arguments.128[0], arguments.128[1], registerM << 1);
stage SR:
  PGR[registerM] = result1;
\end{lstlisting}}
\newpage

\paragraph{Encoding} Format \textsf{LSU\_ACSWAPQB.O} (Section~\ref{sec:format-LSU-ACSWAPQB.O}), \verb|lsucode5 = 100|\\

\bitfields{ 1/P, 2/00, 2/-- --, 27/offset27, 1/1, 2/01, 3/111, 2/coherency, 5/registerM, 1/0, 2/11, 3/100, 1/boolcas, 4/registerQ, 1/--, 1/1, 6/registerZ }


\paragraph{Syntax} \texttt{acswapq}\emph{boolcas}\emph{coherency} \emph{registerM}, \emph{offset27}[\emph{registerZ}] = \emph{registerQ}

\paragraph{Description} The effective address is computed by adding the \emph{offset27} to the \emph{registerZ}.
This instruction implements Compare-And-Swap (CAS) on memory quadruple word (128 bits). The \emph{registerQ} is interpreted as the concatenation of \texttt{update} (low 128 bits) and \texttt{expected} (high 128 bits). The memory double word at the effective address is read into \texttt{current}. If \texttt{current} equals \texttt{expected}, the swap is successful and \texttt{update} is stored into the memory double word at the effective address. Otherwise, the swap has failed and the memory is not updated. The instruction return value is stored into the \emph{registerM}. It is true/false in case \texttt{boolcas} is true, or the \texttt{current} value in case \texttt{boolcas} is false. The effective address must be quadruple word aligned (multiple of 16).


\paragraph{Modifier(s)}
\noindent\begin{itemize}
\item coherency (cf. coherency in section \ref{par:modifier-coherency} on page \pageref{par:modifier-coherency})
\item boolcas (cf. boolcas in section \ref{par:modifier-boolcas} on page \pageref{par:modifier-boolcas})
\end{itemize}

\paragraph{Execution} {Reservation table \textsf{LSU2\_MEMW\_AUXR\_AUXW.X} (Section~\ref{sec:reservation-LSU2-MEMW-AUXR-AUXW.X})}
~
{\begin{lstlisting}
stage ID:
  new offset = _SX_27(offset27);
  new boolcas = _ZX_1(boolcas);
  new coherency = _ZX_2(coherency);
stage RR:
  new base = GPR[registerZ];
  new address = base + offset;
stage E1:
  new arguments = QGR[registerQ];
  new result1 = MEM.atomic_cas(address, 0xFFFF, boolcas | (coherency << 32), arguments.128[0], arguments.128[1], registerM << 1);
stage SR:
  PGR[registerM] = result1;
\end{lstlisting}}
\newpage

\paragraph{Encoding} Format \textsf{LSU\_ACSWAPQB.D} (Section~\ref{sec:format-LSU-ACSWAPQB.D}), \verb|lsucode5 = 100|\\

\bitfields{ 1/P, 2/00, 2/-- --, 27/extend27, 1/1, 2/00, 2/-- --, 27/offset27, 1/1, 2/01, 3/111, 2/coherency, 5/registerM, 1/0, 2/11, 3/100, 1/boolcas, 4/registerQ, 1/--, 1/1, 6/registerZ }


\paragraph{Syntax} \texttt{acswapq}\emph{boolcas}\emph{coherency} \emph{registerM}, \emph{extend27\_\discretionary{}{}{}offset27}[\emph{registerZ}] = \emph{registerQ}

\paragraph{Description} The effective address is computed by adding the \emph{extend27\_\discretionary{}{}{}offset27} to the \emph{registerZ}.
This instruction implements Compare-And-Swap (CAS) on memory quadruple word (128 bits). The \emph{registerQ} is interpreted as the concatenation of \texttt{update} (low 128 bits) and \texttt{expected} (high 128 bits). The memory double word at the effective address is read into \texttt{current}. If \texttt{current} equals \texttt{expected}, the swap is successful and \texttt{update} is stored into the memory double word at the effective address. Otherwise, the swap has failed and the memory is not updated. The instruction return value is stored into the \emph{registerM}. It is true/false in case \texttt{boolcas} is true, or the \texttt{current} value in case \texttt{boolcas} is false. The effective address must be quadruple word aligned (multiple of 16).


\paragraph{Modifier(s)}
\noindent\begin{itemize}
\item coherency (cf. coherency in section \ref{par:modifier-coherency} on page \pageref{par:modifier-coherency})
\item boolcas (cf. boolcas in section \ref{par:modifier-boolcas} on page \pageref{par:modifier-boolcas})
\end{itemize}

\paragraph{Execution} {Reservation table \textsf{LSU2\_MEMW\_AUXR\_AUXW.Y} (Section~\ref{sec:reservation-LSU2-MEMW-AUXR-AUXW.Y})}
~
{\begin{lstlisting}
stage ID:
  new offset = _SX_54(extend27_offset27);
  new boolcas = _ZX_1(boolcas);
  new coherency = _ZX_2(coherency);
stage RR:
  new base = GPR[registerZ];
  new address = base + offset;
stage E1:
  new arguments = QGR[registerQ];
  new result1 = MEM.atomic_cas(address, 0xFFFF, boolcas | (coherency << 32), arguments.128[0], arguments.128[1], registerM << 1);
stage SR:
  PGR[registerM] = result1;
\end{lstlisting}}
\newpage
\subsubsection[{\small ACSWAPW: Atomic Compare and Swap Word}]{ACSWAPW\hfill Atomic Compare and Swap Word}
\label{instr:ACSWAPW}\index{acswapw}
 
\paragraph{Encoding} Format \textsf{LSU\_ACSWAPSB} (Section~\ref{sec:format-LSU-ACSWAPSB}), \verb|lsucode5 = 010|\\

\bitfields{ 1/P, 2/01, 3/111, 2/coherency, 6/registerW, 2/11, 3/010, 1/boolcas, 5/registerO, 1/1, 6/registerZ }


\paragraph{Syntax} \texttt{acswapw}\emph{boolcas}\emph{coherency} \emph{registerW}, [\emph{registerZ}] = \emph{registerO}

\paragraph{Description} The effective address is the value of the \emph{registerZ}.
This instruction implements Compare-And-Swap (CAS) on memory word (32 bits). The \emph{registerO} is interpreted as the concatenation of \texttt{update} (low 64 bits) and \texttt{expected} (high 64 bits). The memory word at the effective address is read into \texttt{current}. If \texttt{current} equals \texttt{expected}, the swap is successful and \texttt{update} is stored into the memory word at the effective address. Otherwise, the swap has failed and the memory is not updated. The instruction return value is stored into the \emph{registerW}. It is true/false in case \texttt{boolcas} is true, or the \texttt{current} value in case \texttt{boolcas} is false. The effective address must be word aligned (multiple of 4).


\paragraph{Modifier(s)}
\noindent\begin{itemize}
\item coherency (cf. coherency in section \ref{par:modifier-coherency} on page \pageref{par:modifier-coherency})
\item boolcas (cf. boolcas in section \ref{par:modifier-boolcas} on page \pageref{par:modifier-boolcas})
\end{itemize}

\paragraph{Execution} {Reservation table \textsf{LSU2\_MEMW\_AUXR\_AUXW} (Section~\ref{sec:reservation-LSU2-MEMW-AUXR-AUXW})}
~
{\begin{lstlisting}
stage ID:
  new boolcas = _ZX_1(boolcas);
  new coherency = _ZX_2(coherency);
stage RR:
  new address = GPR[registerZ];
stage E1:
  new arguments = PGR[registerO];
  new result1 = MEM.atomic_cas(address, 0xF, boolcas | (coherency << 32), arguments.32[0], arguments.32[2], registerW);
stage SR:
  GPR[registerW] = result1;
\end{lstlisting}}
\newpage

\paragraph{Encoding} Format \textsf{LSU\_ACSWAPSB.O} (Section~\ref{sec:format-LSU-ACSWAPSB.O}), \verb|lsucode5 = 010|\\

\bitfields{ 1/P, 2/00, 2/-- --, 27/offset27, 1/1, 2/01, 3/111, 2/coherency, 6/registerW, 2/11, 3/010, 1/boolcas, 5/registerO, 1/1, 6/registerZ }


\paragraph{Syntax} \texttt{acswapw}\emph{boolcas}\emph{coherency} \emph{registerW}, \emph{offset27}[\emph{registerZ}] = \emph{registerO}

\paragraph{Description} The effective address is computed by adding the \emph{offset27} to the \emph{registerZ}.
This instruction implements Compare-And-Swap (CAS) on memory word (32 bits). The \emph{registerO} is interpreted as the concatenation of \texttt{update} (low 64 bits) and \texttt{expected} (high 64 bits). The memory word at the effective address is read into \texttt{current}. If \texttt{current} equals \texttt{expected}, the swap is successful and \texttt{update} is stored into the memory word at the effective address. Otherwise, the swap has failed and the memory is not updated. The instruction return value is stored into the \emph{registerW}. It is true/false in case \texttt{boolcas} is true, or the \texttt{current} value in case \texttt{boolcas} is false. The effective address must be word aligned (multiple of 4).


\paragraph{Modifier(s)}
\noindent\begin{itemize}
\item coherency (cf. coherency in section \ref{par:modifier-coherency} on page \pageref{par:modifier-coherency})
\item boolcas (cf. boolcas in section \ref{par:modifier-boolcas} on page \pageref{par:modifier-boolcas})
\end{itemize}

\paragraph{Execution} {Reservation table \textsf{LSU2\_MEMW\_AUXR\_AUXW.X} (Section~\ref{sec:reservation-LSU2-MEMW-AUXR-AUXW.X})}
~
{\begin{lstlisting}
stage ID:
  new offset = _SX_27(offset27);
  new boolcas = _ZX_1(boolcas);
  new coherency = _ZX_2(coherency);
stage RR:
  new base = GPR[registerZ];
  new address = base + offset;
stage E1:
  new arguments = PGR[registerO];
  new result1 = MEM.atomic_cas(address, 0xF, boolcas | (coherency << 32), arguments.32[0], arguments.32[2], registerW);
stage SR:
  GPR[registerW] = result1;
\end{lstlisting}}
\newpage

\paragraph{Encoding} Format \textsf{LSU\_ACSWAPSB.D} (Section~\ref{sec:format-LSU-ACSWAPSB.D}), \verb|lsucode5 = 010|\\

\bitfields{ 1/P, 2/00, 2/-- --, 27/extend27, 1/1, 2/00, 2/-- --, 27/offset27, 1/1, 2/01, 3/111, 2/coherency, 6/registerW, 2/11, 3/010, 1/boolcas, 5/registerO, 1/1, 6/registerZ }


\paragraph{Syntax} \texttt{acswapw}\emph{boolcas}\emph{coherency} \emph{registerW}, \emph{extend27\_\discretionary{}{}{}offset27}[\emph{registerZ}] = \emph{registerO}

\paragraph{Description} The effective address is computed by adding the \emph{extend27\_\discretionary{}{}{}offset27} to the \emph{registerZ}.
This instruction implements Compare-And-Swap (CAS) on memory word (32 bits). The \emph{registerO} is interpreted as the concatenation of \texttt{update} (low 64 bits) and \texttt{expected} (high 64 bits). The memory word at the effective address is read into \texttt{current}. If \texttt{current} equals \texttt{expected}, the swap is successful and \texttt{update} is stored into the memory word at the effective address. Otherwise, the swap has failed and the memory is not updated. The instruction return value is stored into the \emph{registerW}. It is true/false in case \texttt{boolcas} is true, or the \texttt{current} value in case \texttt{boolcas} is false. The effective address must be word aligned (multiple of 4).


\paragraph{Modifier(s)}
\noindent\begin{itemize}
\item coherency (cf. coherency in section \ref{par:modifier-coherency} on page \pageref{par:modifier-coherency})
\item boolcas (cf. boolcas in section \ref{par:modifier-boolcas} on page \pageref{par:modifier-boolcas})
\end{itemize}

\paragraph{Execution} {Reservation table \textsf{LSU2\_MEMW\_AUXR\_AUXW.Y} (Section~\ref{sec:reservation-LSU2-MEMW-AUXR-AUXW.Y})}
~
{\begin{lstlisting}
stage ID:
  new offset = _SX_54(extend27_offset27);
  new boolcas = _ZX_1(boolcas);
  new coherency = _ZX_2(coherency);
stage RR:
  new base = GPR[registerZ];
  new address = base + offset;
stage E1:
  new arguments = PGR[registerO];
  new result1 = MEM.atomic_cas(address, 0xF, boolcas | (coherency << 32), arguments.32[0], arguments.32[2], registerW);
stage SR:
  GPR[registerW] = result1;
\end{lstlisting}}
\newpage
\subsubsection[{\small ALADDB: Atomic Load and Add Byte}]{ALADDB\hfill Atomic Load and Add Byte}
\label{instr:ALADDB}\index{aladdb}
 
\paragraph{Encoding} Format \textsf{LSU\_ALADDSB} (Section~\ref{sec:format-LSU-ALADDSB}), \verb|lsucode5 = 000|\\

\bitfields{ 1/P, 2/01, 3/111, 2/coherency, 6/registerT, 2/11, 3/000, 1/--, 5/00100, 1/0, 6/registerZ }


\paragraph{Syntax} \texttt{aladdb}\emph{coherency} [\emph{registerZ}] = \emph{registerT}

\paragraph{Description} The effective address is the value of the \emph{registerZ}.
This instruction implements atomic load-add-store on byte. The byte at effective address is added with the \emph{registerT}, and its previous value is returned into the \emph{registerT}.


\paragraph{Modifier(s)}
\noindent\begin{itemize}
\item coherency (cf. coherency in section \ref{par:modifier-coherency} on page \pageref{par:modifier-coherency})
\end{itemize}

\paragraph{Execution} {Reservation table \textsf{LSU2\_MEMW\_AUXR\_AUXW} (Section~\ref{sec:reservation-LSU2-MEMW-AUXR-AUXW})}
~
{\begin{lstlisting}
stage ID:
  new coherency = _ZX_2(coherency);
stage RR:
  new address = GPR[registerZ];
stage E1:
  new argument2 = GPR[registerT];
  new result2 = _ZX_8(MEM.atomic_add(address, 0x1, coherency << 32, argument2.8[0], registerT));
stage SR:
  GPR[registerT] = result2;
\end{lstlisting}}
\newpage

\paragraph{Encoding} Format \textsf{LSU\_ALADDSB.O} (Section~\ref{sec:format-LSU-ALADDSB.O}), \verb|lsucode5 = 000|\\

\bitfields{ 1/P, 2/00, 2/-- --, 27/offset27, 1/1, 2/01, 3/111, 2/coherency, 6/registerT, 2/11, 3/000, 1/--, 5/00100, 1/0, 6/registerZ }


\paragraph{Syntax} \texttt{aladdb}\emph{coherency} \emph{offset27}[\emph{registerZ}] = \emph{registerT}

\paragraph{Description} The effective address is computed by adding the \emph{offset27} to the \emph{registerZ}.
This instruction implements atomic load-add-store on byte. The byte at effective address is added with the \emph{registerT}, and its previous value is returned into the \emph{registerT}.


\paragraph{Modifier(s)}
\noindent\begin{itemize}
\item coherency (cf. coherency in section \ref{par:modifier-coherency} on page \pageref{par:modifier-coherency})
\end{itemize}

\paragraph{Execution} {Reservation table \textsf{LSU2\_MEMW\_AUXR\_AUXW.X} (Section~\ref{sec:reservation-LSU2-MEMW-AUXR-AUXW.X})}
~
{\begin{lstlisting}
stage ID:
  new offset = _SX_27(offset27);
  new coherency = _ZX_2(coherency);
stage RR:
  new base = GPR[registerZ];
  new address = base + offset;
stage E1:
  new argument2 = GPR[registerT];
  new result2 = _ZX_8(MEM.atomic_add(address, 0x1, coherency << 32, argument2.8[0], registerT));
stage SR:
  GPR[registerT] = result2;
\end{lstlisting}}
\newpage

\paragraph{Encoding} Format \textsf{LSU\_ALADDSB.D} (Section~\ref{sec:format-LSU-ALADDSB.D}), \verb|lsucode5 = 000|\\

\bitfields{ 1/P, 2/00, 2/-- --, 27/extend27, 1/1, 2/00, 2/-- --, 27/offset27, 1/1, 2/01, 3/111, 2/coherency, 6/registerT, 2/11, 3/000, 1/--, 5/00100, 1/0, 6/registerZ }


\paragraph{Syntax} \texttt{aladdb}\emph{coherency} \emph{extend27\_\discretionary{}{}{}offset27}[\emph{registerZ}] = \emph{registerT}

\paragraph{Description} The effective address is computed by adding the \emph{extend27\_\discretionary{}{}{}offset27} to the \emph{registerZ}.
This instruction implements atomic load-add-store on byte. The byte at effective address is added with the \emph{registerT}, and its previous value is returned into the \emph{registerT}.


\paragraph{Modifier(s)}
\noindent\begin{itemize}
\item coherency (cf. coherency in section \ref{par:modifier-coherency} on page \pageref{par:modifier-coherency})
\end{itemize}

\paragraph{Execution} {Reservation table \textsf{LSU2\_MEMW\_AUXR\_AUXW.Y} (Section~\ref{sec:reservation-LSU2-MEMW-AUXR-AUXW.Y})}
~
{\begin{lstlisting}
stage ID:
  new offset = _SX_54(extend27_offset27);
  new coherency = _ZX_2(coherency);
stage RR:
  new base = GPR[registerZ];
  new address = base + offset;
stage E1:
  new argument2 = GPR[registerT];
  new result2 = _ZX_8(MEM.atomic_add(address, 0x1, coherency << 32, argument2.8[0], registerT));
stage SR:
  GPR[registerT] = result2;
\end{lstlisting}}
\newpage
\subsubsection[{\small ALADDD: Atomic Load and Add Double Word}]{ALADDD\hfill Atomic Load and Add Double Word}
\label{instr:ALADDD}\index{aladdd}
 
\paragraph{Encoding} Format \textsf{LSU\_ALADDSB} (Section~\ref{sec:format-LSU-ALADDSB}), \verb|lsucode5 = 011|\\

\bitfields{ 1/P, 2/01, 3/111, 2/coherency, 6/registerT, 2/11, 3/011, 1/--, 5/00100, 1/0, 6/registerZ }


\paragraph{Syntax} \texttt{aladdd}\emph{coherency} [\emph{registerZ}] = \emph{registerT}

\paragraph{Description} The effective address is the value of the \emph{registerZ}.
This instruction implements atomic load-add-store on double word. The double word at effective address is added with the \emph{registerT}, and its previous value is returned into the \emph{registerT}. The effective address must be double word aligned (multiple of 8).


\paragraph{Modifier(s)}
\noindent\begin{itemize}
\item coherency (cf. coherency in section \ref{par:modifier-coherency} on page \pageref{par:modifier-coherency})
\end{itemize}

\paragraph{Execution} {Reservation table \textsf{LSU2\_MEMW\_AUXR\_AUXW} (Section~\ref{sec:reservation-LSU2-MEMW-AUXR-AUXW})}
~
{\begin{lstlisting}
stage ID:
  new coherency = _ZX_2(coherency);
stage RR:
  new address = GPR[registerZ];
stage E1:
  new argument2 = GPR[registerT];
  new result2 = MEM.atomic_add(address, 0xFF, coherency << 32, argument2.64[0], registerT);
stage SR:
  GPR[registerT] = result2;
\end{lstlisting}}
\newpage

\paragraph{Encoding} Format \textsf{LSU\_ALADDSB.O} (Section~\ref{sec:format-LSU-ALADDSB.O}), \verb|lsucode5 = 011|\\

\bitfields{ 1/P, 2/00, 2/-- --, 27/offset27, 1/1, 2/01, 3/111, 2/coherency, 6/registerT, 2/11, 3/011, 1/--, 5/00100, 1/0, 6/registerZ }


\paragraph{Syntax} \texttt{aladdd}\emph{coherency} \emph{offset27}[\emph{registerZ}] = \emph{registerT}

\paragraph{Description} The effective address is computed by adding the \emph{offset27} to the \emph{registerZ}.
This instruction implements atomic load-add-store on double word. The double word at effective address is added with the \emph{registerT}, and its previous value is returned into the \emph{registerT}. The effective address must be double word aligned (multiple of 8).


\paragraph{Modifier(s)}
\noindent\begin{itemize}
\item coherency (cf. coherency in section \ref{par:modifier-coherency} on page \pageref{par:modifier-coherency})
\end{itemize}

\paragraph{Execution} {Reservation table \textsf{LSU2\_MEMW\_AUXR\_AUXW.X} (Section~\ref{sec:reservation-LSU2-MEMW-AUXR-AUXW.X})}
~
{\begin{lstlisting}
stage ID:
  new offset = _SX_27(offset27);
  new coherency = _ZX_2(coherency);
stage RR:
  new base = GPR[registerZ];
  new address = base + offset;
stage E1:
  new argument2 = GPR[registerT];
  new result2 = MEM.atomic_add(address, 0xFF, coherency << 32, argument2.64[0], registerT);
stage SR:
  GPR[registerT] = result2;
\end{lstlisting}}
\newpage

\paragraph{Encoding} Format \textsf{LSU\_ALADDSB.D} (Section~\ref{sec:format-LSU-ALADDSB.D}), \verb|lsucode5 = 011|\\

\bitfields{ 1/P, 2/00, 2/-- --, 27/extend27, 1/1, 2/00, 2/-- --, 27/offset27, 1/1, 2/01, 3/111, 2/coherency, 6/registerT, 2/11, 3/011, 1/--, 5/00100, 1/0, 6/registerZ }


\paragraph{Syntax} \texttt{aladdd}\emph{coherency} \emph{extend27\_\discretionary{}{}{}offset27}[\emph{registerZ}] = \emph{registerT}

\paragraph{Description} The effective address is computed by adding the \emph{extend27\_\discretionary{}{}{}offset27} to the \emph{registerZ}.
This instruction implements atomic load-add-store on double word. The double word at effective address is added with the \emph{registerT}, and its previous value is returned into the \emph{registerT}. The effective address must be double word aligned (multiple of 8).


\paragraph{Modifier(s)}
\noindent\begin{itemize}
\item coherency (cf. coherency in section \ref{par:modifier-coherency} on page \pageref{par:modifier-coherency})
\end{itemize}

\paragraph{Execution} {Reservation table \textsf{LSU2\_MEMW\_AUXR\_AUXW.Y} (Section~\ref{sec:reservation-LSU2-MEMW-AUXR-AUXW.Y})}
~
{\begin{lstlisting}
stage ID:
  new offset = _SX_54(extend27_offset27);
  new coherency = _ZX_2(coherency);
stage RR:
  new base = GPR[registerZ];
  new address = base + offset;
stage E1:
  new argument2 = GPR[registerT];
  new result2 = MEM.atomic_add(address, 0xFF, coherency << 32, argument2.64[0], registerT);
stage SR:
  GPR[registerT] = result2;
\end{lstlisting}}
\newpage
\subsubsection[{\small ALADDH: Atomic Load and Add Half Word}]{ALADDH\hfill Atomic Load and Add Half Word}
\label{instr:ALADDH}\index{aladdh}
 
\paragraph{Encoding} Format \textsf{LSU\_ALADDSB} (Section~\ref{sec:format-LSU-ALADDSB}), \verb|lsucode5 = 001|\\

\bitfields{ 1/P, 2/01, 3/111, 2/coherency, 6/registerT, 2/11, 3/001, 1/--, 5/00100, 1/0, 6/registerZ }


\paragraph{Syntax} \texttt{aladdh}\emph{coherency} [\emph{registerZ}] = \emph{registerT}

\paragraph{Description} The effective address is the value of the \emph{registerZ}.
This instruction implements atomic load-add-store on half word. The half word at effective address is added with the \emph{registerT}, and its previous value is returned into the \emph{registerT}. The effective address must be half word aligned (multiple of 2).


\paragraph{Modifier(s)}
\noindent\begin{itemize}
\item coherency (cf. coherency in section \ref{par:modifier-coherency} on page \pageref{par:modifier-coherency})
\end{itemize}

\paragraph{Execution} {Reservation table \textsf{LSU2\_MEMW\_AUXR\_AUXW} (Section~\ref{sec:reservation-LSU2-MEMW-AUXR-AUXW})}
~
{\begin{lstlisting}
stage ID:
  new coherency = _ZX_2(coherency);
stage RR:
  new address = GPR[registerZ];
stage E1:
  new argument2 = GPR[registerT];
  new result2 = _ZX_16(MEM.atomic_add(address, 0x3, coherency << 32, argument2.16[0], registerT));
stage SR:
  GPR[registerT] = result2;
\end{lstlisting}}
\newpage

\paragraph{Encoding} Format \textsf{LSU\_ALADDSB.O} (Section~\ref{sec:format-LSU-ALADDSB.O}), \verb|lsucode5 = 001|\\

\bitfields{ 1/P, 2/00, 2/-- --, 27/offset27, 1/1, 2/01, 3/111, 2/coherency, 6/registerT, 2/11, 3/001, 1/--, 5/00100, 1/0, 6/registerZ }


\paragraph{Syntax} \texttt{aladdh}\emph{coherency} \emph{offset27}[\emph{registerZ}] = \emph{registerT}

\paragraph{Description} The effective address is computed by adding the \emph{offset27} to the \emph{registerZ}.
This instruction implements atomic load-add-store on half word. The half word at effective address is added with the \emph{registerT}, and its previous value is returned into the \emph{registerT}. The effective address must be half word aligned (multiple of 2).


\paragraph{Modifier(s)}
\noindent\begin{itemize}
\item coherency (cf. coherency in section \ref{par:modifier-coherency} on page \pageref{par:modifier-coherency})
\end{itemize}

\paragraph{Execution} {Reservation table \textsf{LSU2\_MEMW\_AUXR\_AUXW.X} (Section~\ref{sec:reservation-LSU2-MEMW-AUXR-AUXW.X})}
~
{\begin{lstlisting}
stage ID:
  new offset = _SX_27(offset27);
  new coherency = _ZX_2(coherency);
stage RR:
  new base = GPR[registerZ];
  new address = base + offset;
stage E1:
  new argument2 = GPR[registerT];
  new result2 = _ZX_16(MEM.atomic_add(address, 0x3, coherency << 32, argument2.16[0], registerT));
stage SR:
  GPR[registerT] = result2;
\end{lstlisting}}
\newpage

\paragraph{Encoding} Format \textsf{LSU\_ALADDSB.D} (Section~\ref{sec:format-LSU-ALADDSB.D}), \verb|lsucode5 = 001|\\

\bitfields{ 1/P, 2/00, 2/-- --, 27/extend27, 1/1, 2/00, 2/-- --, 27/offset27, 1/1, 2/01, 3/111, 2/coherency, 6/registerT, 2/11, 3/001, 1/--, 5/00100, 1/0, 6/registerZ }


\paragraph{Syntax} \texttt{aladdh}\emph{coherency} \emph{extend27\_\discretionary{}{}{}offset27}[\emph{registerZ}] = \emph{registerT}

\paragraph{Description} The effective address is computed by adding the \emph{extend27\_\discretionary{}{}{}offset27} to the \emph{registerZ}.
This instruction implements atomic load-add-store on half word. The half word at effective address is added with the \emph{registerT}, and its previous value is returned into the \emph{registerT}. The effective address must be half word aligned (multiple of 2).


\paragraph{Modifier(s)}
\noindent\begin{itemize}
\item coherency (cf. coherency in section \ref{par:modifier-coherency} on page \pageref{par:modifier-coherency})
\end{itemize}

\paragraph{Execution} {Reservation table \textsf{LSU2\_MEMW\_AUXR\_AUXW.Y} (Section~\ref{sec:reservation-LSU2-MEMW-AUXR-AUXW.Y})}
~
{\begin{lstlisting}
stage ID:
  new offset = _SX_54(extend27_offset27);
  new coherency = _ZX_2(coherency);
stage RR:
  new base = GPR[registerZ];
  new address = base + offset;
stage E1:
  new argument2 = GPR[registerT];
  new result2 = _ZX_16(MEM.atomic_add(address, 0x3, coherency << 32, argument2.16[0], registerT));
stage SR:
  GPR[registerT] = result2;
\end{lstlisting}}
\newpage
\subsubsection[{\small ALADDW: Atomic Load and Add Word}]{ALADDW\hfill Atomic Load and Add Word}
\label{instr:ALADDW}\index{aladdw}
 
\paragraph{Encoding} Format \textsf{LSU\_ALADDSB} (Section~\ref{sec:format-LSU-ALADDSB}), \verb|lsucode5 = 010|\\

\bitfields{ 1/P, 2/01, 3/111, 2/coherency, 6/registerT, 2/11, 3/010, 1/--, 5/00100, 1/0, 6/registerZ }


\paragraph{Syntax} \texttt{aladdw}\emph{coherency} [\emph{registerZ}] = \emph{registerT}

\paragraph{Description} The effective address is the value of the \emph{registerZ}.
This instruction implements atomic load-add-store on word. The word at effective address is added with the \emph{registerT}, and its previous value is returned into the \emph{registerT}. The effective address must be word aligned (multiple of 4).


\paragraph{Modifier(s)}
\noindent\begin{itemize}
\item coherency (cf. coherency in section \ref{par:modifier-coherency} on page \pageref{par:modifier-coherency})
\end{itemize}

\paragraph{Execution} {Reservation table \textsf{LSU2\_MEMW\_AUXR\_AUXW} (Section~\ref{sec:reservation-LSU2-MEMW-AUXR-AUXW})}
~
{\begin{lstlisting}
stage ID:
  new coherency = _ZX_2(coherency);
stage RR:
  new address = GPR[registerZ];
stage E1:
  new argument2 = GPR[registerT];
  new result2 = _ZX_32(MEM.atomic_add(address, 0xF, coherency << 32, argument2.32[0], registerT));
stage SR:
  GPR[registerT] = result2;
\end{lstlisting}}
\newpage

\paragraph{Encoding} Format \textsf{LSU\_ALADDSB.O} (Section~\ref{sec:format-LSU-ALADDSB.O}), \verb|lsucode5 = 010|\\

\bitfields{ 1/P, 2/00, 2/-- --, 27/offset27, 1/1, 2/01, 3/111, 2/coherency, 6/registerT, 2/11, 3/010, 1/--, 5/00100, 1/0, 6/registerZ }


\paragraph{Syntax} \texttt{aladdw}\emph{coherency} \emph{offset27}[\emph{registerZ}] = \emph{registerT}

\paragraph{Description} The effective address is computed by adding the \emph{offset27} to the \emph{registerZ}.
This instruction implements atomic load-add-store on word. The word at effective address is added with the \emph{registerT}, and its previous value is returned into the \emph{registerT}. The effective address must be word aligned (multiple of 4).


\paragraph{Modifier(s)}
\noindent\begin{itemize}
\item coherency (cf. coherency in section \ref{par:modifier-coherency} on page \pageref{par:modifier-coherency})
\end{itemize}

\paragraph{Execution} {Reservation table \textsf{LSU2\_MEMW\_AUXR\_AUXW.X} (Section~\ref{sec:reservation-LSU2-MEMW-AUXR-AUXW.X})}
~
{\begin{lstlisting}
stage ID:
  new offset = _SX_27(offset27);
  new coherency = _ZX_2(coherency);
stage RR:
  new base = GPR[registerZ];
  new address = base + offset;
stage E1:
  new argument2 = GPR[registerT];
  new result2 = _ZX_32(MEM.atomic_add(address, 0xF, coherency << 32, argument2.32[0], registerT));
stage SR:
  GPR[registerT] = result2;
\end{lstlisting}}
\newpage

\paragraph{Encoding} Format \textsf{LSU\_ALADDSB.D} (Section~\ref{sec:format-LSU-ALADDSB.D}), \verb|lsucode5 = 010|\\

\bitfields{ 1/P, 2/00, 2/-- --, 27/extend27, 1/1, 2/00, 2/-- --, 27/offset27, 1/1, 2/01, 3/111, 2/coherency, 6/registerT, 2/11, 3/010, 1/--, 5/00100, 1/0, 6/registerZ }


\paragraph{Syntax} \texttt{aladdw}\emph{coherency} \emph{extend27\_\discretionary{}{}{}offset27}[\emph{registerZ}] = \emph{registerT}

\paragraph{Description} The effective address is computed by adding the \emph{extend27\_\discretionary{}{}{}offset27} to the \emph{registerZ}.
This instruction implements atomic load-add-store on word. The word at effective address is added with the \emph{registerT}, and its previous value is returned into the \emph{registerT}. The effective address must be word aligned (multiple of 4).


\paragraph{Modifier(s)}
\noindent\begin{itemize}
\item coherency (cf. coherency in section \ref{par:modifier-coherency} on page \pageref{par:modifier-coherency})
\end{itemize}

\paragraph{Execution} {Reservation table \textsf{LSU2\_MEMW\_AUXR\_AUXW.Y} (Section~\ref{sec:reservation-LSU2-MEMW-AUXR-AUXW.Y})}
~
{\begin{lstlisting}
stage ID:
  new offset = _SX_54(extend27_offset27);
  new coherency = _ZX_2(coherency);
stage RR:
  new base = GPR[registerZ];
  new address = base + offset;
stage E1:
  new argument2 = GPR[registerT];
  new result2 = _ZX_32(MEM.atomic_add(address, 0xF, coherency << 32, argument2.32[0], registerT));
stage SR:
  GPR[registerT] = result2;
\end{lstlisting}}
\newpage
\subsubsection[{\small ALANDB: Atomic Load and AND Byte}]{ALANDB\hfill Atomic Load and AND Byte}
\label{instr:ALANDB}\index{alandb}
 
\paragraph{Encoding} Format \textsf{LSU\_ALANDSB} (Section~\ref{sec:format-LSU-ALANDSB}), \verb|lsucode5 = 000|\\

\bitfields{ 1/P, 2/01, 3/111, 2/coherency, 6/registerT, 2/11, 3/000, 1/--, 5/00101, 1/0, 6/registerZ }


\paragraph{Syntax} \texttt{alandb}\emph{coherency} [\emph{registerZ}] = \emph{registerT}

\paragraph{Description} The effective address is the value of the \emph{registerZ}.
This instruction implements atomic load-AND-store on byte. The byte at effective address is ANDed with the \emph{registerT}, and its previous value is returned into the \emph{registerT}.


\paragraph{Modifier(s)}
\noindent\begin{itemize}
\item coherency (cf. coherency in section \ref{par:modifier-coherency} on page \pageref{par:modifier-coherency})
\end{itemize}

\paragraph{Execution} {Reservation table \textsf{LSU2\_MEMW\_AUXR\_AUXW} (Section~\ref{sec:reservation-LSU2-MEMW-AUXR-AUXW})}
~
{\begin{lstlisting}
stage ID:
  new coherency = _ZX_2(coherency);
stage RR:
  new address = GPR[registerZ];
stage E1:
  new argument2 = GPR[registerT];
  new result2 = _ZX_8(MEM.atomic_and(address, 0x1, coherency << 32, argument2.8[0], registerT));
stage SR:
  GPR[registerT] = result2;
\end{lstlisting}}
\newpage

\paragraph{Encoding} Format \textsf{LSU\_ALANDSB.O} (Section~\ref{sec:format-LSU-ALANDSB.O}), \verb|lsucode5 = 000|\\

\bitfields{ 1/P, 2/00, 2/-- --, 27/offset27, 1/1, 2/01, 3/111, 2/coherency, 6/registerT, 2/11, 3/000, 1/--, 5/00101, 1/0, 6/registerZ }


\paragraph{Syntax} \texttt{alandb}\emph{coherency} \emph{offset27}[\emph{registerZ}] = \emph{registerT}

\paragraph{Description} The effective address is computed by adding the \emph{offset27} to the \emph{registerZ}.
This instruction implements atomic load-AND-store on byte. The byte at effective address is ANDed with the \emph{registerT}, and its previous value is returned into the \emph{registerT}.


\paragraph{Modifier(s)}
\noindent\begin{itemize}
\item coherency (cf. coherency in section \ref{par:modifier-coherency} on page \pageref{par:modifier-coherency})
\end{itemize}

\paragraph{Execution} {Reservation table \textsf{LSU2\_MEMW\_AUXR\_AUXW.X} (Section~\ref{sec:reservation-LSU2-MEMW-AUXR-AUXW.X})}
~
{\begin{lstlisting}
stage ID:
  new offset = _SX_27(offset27);
  new coherency = _ZX_2(coherency);
stage RR:
  new base = GPR[registerZ];
  new address = base + offset;
stage E1:
  new argument2 = GPR[registerT];
  new result2 = _ZX_8(MEM.atomic_and(address, 0x1, coherency << 32, argument2.8[0], registerT));
stage SR:
  GPR[registerT] = result2;
\end{lstlisting}}
\newpage

\paragraph{Encoding} Format \textsf{LSU\_ALANDSB.D} (Section~\ref{sec:format-LSU-ALANDSB.D}), \verb|lsucode5 = 000|\\

\bitfields{ 1/P, 2/00, 2/-- --, 27/extend27, 1/1, 2/00, 2/-- --, 27/offset27, 1/1, 2/01, 3/111, 2/coherency, 6/registerT, 2/11, 3/000, 1/--, 5/00101, 1/0, 6/registerZ }


\paragraph{Syntax} \texttt{alandb}\emph{coherency} \emph{extend27\_\discretionary{}{}{}offset27}[\emph{registerZ}] = \emph{registerT}

\paragraph{Description} The effective address is computed by adding the \emph{extend27\_\discretionary{}{}{}offset27} to the \emph{registerZ}.
This instruction implements atomic load-AND-store on byte. The byte at effective address is ANDed with the \emph{registerT}, and its previous value is returned into the \emph{registerT}.


\paragraph{Modifier(s)}
\noindent\begin{itemize}
\item coherency (cf. coherency in section \ref{par:modifier-coherency} on page \pageref{par:modifier-coherency})
\end{itemize}

\paragraph{Execution} {Reservation table \textsf{LSU2\_MEMW\_AUXR\_AUXW.Y} (Section~\ref{sec:reservation-LSU2-MEMW-AUXR-AUXW.Y})}
~
{\begin{lstlisting}
stage ID:
  new offset = _SX_54(extend27_offset27);
  new coherency = _ZX_2(coherency);
stage RR:
  new base = GPR[registerZ];
  new address = base + offset;
stage E1:
  new argument2 = GPR[registerT];
  new result2 = _ZX_8(MEM.atomic_and(address, 0x1, coherency << 32, argument2.8[0], registerT));
stage SR:
  GPR[registerT] = result2;
\end{lstlisting}}
\newpage
\subsubsection[{\small ALANDD: Atomic Load and AND Double Word}]{ALANDD\hfill Atomic Load and AND Double Word}
\label{instr:ALANDD}\index{alandd}
 
\paragraph{Encoding} Format \textsf{LSU\_ALANDSB} (Section~\ref{sec:format-LSU-ALANDSB}), \verb|lsucode5 = 011|\\

\bitfields{ 1/P, 2/01, 3/111, 2/coherency, 6/registerT, 2/11, 3/011, 1/--, 5/00101, 1/0, 6/registerZ }


\paragraph{Syntax} \texttt{alandd}\emph{coherency} [\emph{registerZ}] = \emph{registerT}

\paragraph{Description} The effective address is the value of the \emph{registerZ}.
This instruction implements atomic load-AND-store on double word. The double word at effective address is ANDed with the \emph{registerT}, and its previous value is returned into the \emph{registerT}. The effective address must be double word aligned (multiple of 8).


\paragraph{Modifier(s)}
\noindent\begin{itemize}
\item coherency (cf. coherency in section \ref{par:modifier-coherency} on page \pageref{par:modifier-coherency})
\end{itemize}

\paragraph{Execution} {Reservation table \textsf{LSU2\_MEMW\_AUXR\_AUXW} (Section~\ref{sec:reservation-LSU2-MEMW-AUXR-AUXW})}
~
{\begin{lstlisting}
stage ID:
  new coherency = _ZX_2(coherency);
stage RR:
  new address = GPR[registerZ];
stage E1:
  new argument2 = GPR[registerT];
  new result2 = MEM.atomic_and(address, 0xFF, coherency << 32, argument2.64[0], registerT);
stage SR:
  GPR[registerT] = result2;
\end{lstlisting}}
\newpage

\paragraph{Encoding} Format \textsf{LSU\_ALANDSB.O} (Section~\ref{sec:format-LSU-ALANDSB.O}), \verb|lsucode5 = 011|\\

\bitfields{ 1/P, 2/00, 2/-- --, 27/offset27, 1/1, 2/01, 3/111, 2/coherency, 6/registerT, 2/11, 3/011, 1/--, 5/00101, 1/0, 6/registerZ }


\paragraph{Syntax} \texttt{alandd}\emph{coherency} \emph{offset27}[\emph{registerZ}] = \emph{registerT}

\paragraph{Description} The effective address is computed by adding the \emph{offset27} to the \emph{registerZ}.
This instruction implements atomic load-AND-store on double word. The double word at effective address is ANDed with the \emph{registerT}, and its previous value is returned into the \emph{registerT}. The effective address must be double word aligned (multiple of 8).


\paragraph{Modifier(s)}
\noindent\begin{itemize}
\item coherency (cf. coherency in section \ref{par:modifier-coherency} on page \pageref{par:modifier-coherency})
\end{itemize}

\paragraph{Execution} {Reservation table \textsf{LSU2\_MEMW\_AUXR\_AUXW.X} (Section~\ref{sec:reservation-LSU2-MEMW-AUXR-AUXW.X})}
~
{\begin{lstlisting}
stage ID:
  new offset = _SX_27(offset27);
  new coherency = _ZX_2(coherency);
stage RR:
  new base = GPR[registerZ];
  new address = base + offset;
stage E1:
  new argument2 = GPR[registerT];
  new result2 = MEM.atomic_and(address, 0xFF, coherency << 32, argument2.64[0], registerT);
stage SR:
  GPR[registerT] = result2;
\end{lstlisting}}
\newpage

\paragraph{Encoding} Format \textsf{LSU\_ALANDSB.D} (Section~\ref{sec:format-LSU-ALANDSB.D}), \verb|lsucode5 = 011|\\

\bitfields{ 1/P, 2/00, 2/-- --, 27/extend27, 1/1, 2/00, 2/-- --, 27/offset27, 1/1, 2/01, 3/111, 2/coherency, 6/registerT, 2/11, 3/011, 1/--, 5/00101, 1/0, 6/registerZ }


\paragraph{Syntax} \texttt{alandd}\emph{coherency} \emph{extend27\_\discretionary{}{}{}offset27}[\emph{registerZ}] = \emph{registerT}

\paragraph{Description} The effective address is computed by adding the \emph{extend27\_\discretionary{}{}{}offset27} to the \emph{registerZ}.
This instruction implements atomic load-AND-store on double word. The double word at effective address is ANDed with the \emph{registerT}, and its previous value is returned into the \emph{registerT}. The effective address must be double word aligned (multiple of 8).


\paragraph{Modifier(s)}
\noindent\begin{itemize}
\item coherency (cf. coherency in section \ref{par:modifier-coherency} on page \pageref{par:modifier-coherency})
\end{itemize}

\paragraph{Execution} {Reservation table \textsf{LSU2\_MEMW\_AUXR\_AUXW.Y} (Section~\ref{sec:reservation-LSU2-MEMW-AUXR-AUXW.Y})}
~
{\begin{lstlisting}
stage ID:
  new offset = _SX_54(extend27_offset27);
  new coherency = _ZX_2(coherency);
stage RR:
  new base = GPR[registerZ];
  new address = base + offset;
stage E1:
  new argument2 = GPR[registerT];
  new result2 = MEM.atomic_and(address, 0xFF, coherency << 32, argument2.64[0], registerT);
stage SR:
  GPR[registerT] = result2;
\end{lstlisting}}
\newpage
\subsubsection[{\small ALANDH: Atomic Load and AND Half Word}]{ALANDH\hfill Atomic Load and AND Half Word}
\label{instr:ALANDH}\index{alandh}
 
\paragraph{Encoding} Format \textsf{LSU\_ALANDSB} (Section~\ref{sec:format-LSU-ALANDSB}), \verb|lsucode5 = 001|\\

\bitfields{ 1/P, 2/01, 3/111, 2/coherency, 6/registerT, 2/11, 3/001, 1/--, 5/00101, 1/0, 6/registerZ }


\paragraph{Syntax} \texttt{alandh}\emph{coherency} [\emph{registerZ}] = \emph{registerT}

\paragraph{Description} The effective address is the value of the \emph{registerZ}.
This instruction implements atomic load-AND-store on half word. The half word at effective address is ANDed with the \emph{registerT}, and its previous value is returned into the \emph{registerT}. The effective address must be half word aligned (multiple of 2).


\paragraph{Modifier(s)}
\noindent\begin{itemize}
\item coherency (cf. coherency in section \ref{par:modifier-coherency} on page \pageref{par:modifier-coherency})
\end{itemize}

\paragraph{Execution} {Reservation table \textsf{LSU2\_MEMW\_AUXR\_AUXW} (Section~\ref{sec:reservation-LSU2-MEMW-AUXR-AUXW})}
~
{\begin{lstlisting}
stage ID:
  new coherency = _ZX_2(coherency);
stage RR:
  new address = GPR[registerZ];
stage E1:
  new argument2 = GPR[registerT];
  new result2 = _ZX_16(MEM.atomic_and(address, 0x3, coherency << 32, argument2.16[0], registerT));
stage SR:
  GPR[registerT] = result2;
\end{lstlisting}}
\newpage

\paragraph{Encoding} Format \textsf{LSU\_ALANDSB.O} (Section~\ref{sec:format-LSU-ALANDSB.O}), \verb|lsucode5 = 001|\\

\bitfields{ 1/P, 2/00, 2/-- --, 27/offset27, 1/1, 2/01, 3/111, 2/coherency, 6/registerT, 2/11, 3/001, 1/--, 5/00101, 1/0, 6/registerZ }


\paragraph{Syntax} \texttt{alandh}\emph{coherency} \emph{offset27}[\emph{registerZ}] = \emph{registerT}

\paragraph{Description} The effective address is computed by adding the \emph{offset27} to the \emph{registerZ}.
This instruction implements atomic load-AND-store on half word. The half word at effective address is ANDed with the \emph{registerT}, and its previous value is returned into the \emph{registerT}. The effective address must be half word aligned (multiple of 2).


\paragraph{Modifier(s)}
\noindent\begin{itemize}
\item coherency (cf. coherency in section \ref{par:modifier-coherency} on page \pageref{par:modifier-coherency})
\end{itemize}

\paragraph{Execution} {Reservation table \textsf{LSU2\_MEMW\_AUXR\_AUXW.X} (Section~\ref{sec:reservation-LSU2-MEMW-AUXR-AUXW.X})}
~
{\begin{lstlisting}
stage ID:
  new offset = _SX_27(offset27);
  new coherency = _ZX_2(coherency);
stage RR:
  new base = GPR[registerZ];
  new address = base + offset;
stage E1:
  new argument2 = GPR[registerT];
  new result2 = _ZX_16(MEM.atomic_and(address, 0x3, coherency << 32, argument2.16[0], registerT));
stage SR:
  GPR[registerT] = result2;
\end{lstlisting}}
\newpage

\paragraph{Encoding} Format \textsf{LSU\_ALANDSB.D} (Section~\ref{sec:format-LSU-ALANDSB.D}), \verb|lsucode5 = 001|\\

\bitfields{ 1/P, 2/00, 2/-- --, 27/extend27, 1/1, 2/00, 2/-- --, 27/offset27, 1/1, 2/01, 3/111, 2/coherency, 6/registerT, 2/11, 3/001, 1/--, 5/00101, 1/0, 6/registerZ }


\paragraph{Syntax} \texttt{alandh}\emph{coherency} \emph{extend27\_\discretionary{}{}{}offset27}[\emph{registerZ}] = \emph{registerT}

\paragraph{Description} The effective address is computed by adding the \emph{extend27\_\discretionary{}{}{}offset27} to the \emph{registerZ}.
This instruction implements atomic load-AND-store on half word. The half word at effective address is ANDed with the \emph{registerT}, and its previous value is returned into the \emph{registerT}. The effective address must be half word aligned (multiple of 2).


\paragraph{Modifier(s)}
\noindent\begin{itemize}
\item coherency (cf. coherency in section \ref{par:modifier-coherency} on page \pageref{par:modifier-coherency})
\end{itemize}

\paragraph{Execution} {Reservation table \textsf{LSU2\_MEMW\_AUXR\_AUXW.Y} (Section~\ref{sec:reservation-LSU2-MEMW-AUXR-AUXW.Y})}
~
{\begin{lstlisting}
stage ID:
  new offset = _SX_54(extend27_offset27);
  new coherency = _ZX_2(coherency);
stage RR:
  new base = GPR[registerZ];
  new address = base + offset;
stage E1:
  new argument2 = GPR[registerT];
  new result2 = _ZX_16(MEM.atomic_and(address, 0x3, coherency << 32, argument2.16[0], registerT));
stage SR:
  GPR[registerT] = result2;
\end{lstlisting}}
\newpage
\subsubsection[{\small ALANDW: Atomic Load and AND Word}]{ALANDW\hfill Atomic Load and AND Word}
\label{instr:ALANDW}\index{alandw}
 
\paragraph{Encoding} Format \textsf{LSU\_ALANDSB} (Section~\ref{sec:format-LSU-ALANDSB}), \verb|lsucode5 = 010|\\

\bitfields{ 1/P, 2/01, 3/111, 2/coherency, 6/registerT, 2/11, 3/010, 1/--, 5/00101, 1/0, 6/registerZ }


\paragraph{Syntax} \texttt{alandw}\emph{coherency} [\emph{registerZ}] = \emph{registerT}

\paragraph{Description} The effective address is the value of the \emph{registerZ}.
This instruction implements atomic load-AND-store on word. The word at effective address is ANDed with the \emph{registerT}, and its previous value is returned into the \emph{registerT}. The effective address must be word aligned (multiple of 4).


\paragraph{Modifier(s)}
\noindent\begin{itemize}
\item coherency (cf. coherency in section \ref{par:modifier-coherency} on page \pageref{par:modifier-coherency})
\end{itemize}

\paragraph{Execution} {Reservation table \textsf{LSU2\_MEMW\_AUXR\_AUXW} (Section~\ref{sec:reservation-LSU2-MEMW-AUXR-AUXW})}
~
{\begin{lstlisting}
stage ID:
  new coherency = _ZX_2(coherency);
stage RR:
  new address = GPR[registerZ];
stage E1:
  new argument2 = GPR[registerT];
  new result2 = _ZX_32(MEM.atomic_and(address, 0xF, coherency << 32, argument2.32[0], registerT));
stage SR:
  GPR[registerT] = result2;
\end{lstlisting}}
\newpage

\paragraph{Encoding} Format \textsf{LSU\_ALANDSB.O} (Section~\ref{sec:format-LSU-ALANDSB.O}), \verb|lsucode5 = 010|\\

\bitfields{ 1/P, 2/00, 2/-- --, 27/offset27, 1/1, 2/01, 3/111, 2/coherency, 6/registerT, 2/11, 3/010, 1/--, 5/00101, 1/0, 6/registerZ }


\paragraph{Syntax} \texttt{alandw}\emph{coherency} \emph{offset27}[\emph{registerZ}] = \emph{registerT}

\paragraph{Description} The effective address is computed by adding the \emph{offset27} to the \emph{registerZ}.
This instruction implements atomic load-AND-store on word. The word at effective address is ANDed with the \emph{registerT}, and its previous value is returned into the \emph{registerT}. The effective address must be word aligned (multiple of 4).


\paragraph{Modifier(s)}
\noindent\begin{itemize}
\item coherency (cf. coherency in section \ref{par:modifier-coherency} on page \pageref{par:modifier-coherency})
\end{itemize}

\paragraph{Execution} {Reservation table \textsf{LSU2\_MEMW\_AUXR\_AUXW.X} (Section~\ref{sec:reservation-LSU2-MEMW-AUXR-AUXW.X})}
~
{\begin{lstlisting}
stage ID:
  new offset = _SX_27(offset27);
  new coherency = _ZX_2(coherency);
stage RR:
  new base = GPR[registerZ];
  new address = base + offset;
stage E1:
  new argument2 = GPR[registerT];
  new result2 = _ZX_32(MEM.atomic_and(address, 0xF, coherency << 32, argument2.32[0], registerT));
stage SR:
  GPR[registerT] = result2;
\end{lstlisting}}
\newpage

\paragraph{Encoding} Format \textsf{LSU\_ALANDSB.D} (Section~\ref{sec:format-LSU-ALANDSB.D}), \verb|lsucode5 = 010|\\

\bitfields{ 1/P, 2/00, 2/-- --, 27/extend27, 1/1, 2/00, 2/-- --, 27/offset27, 1/1, 2/01, 3/111, 2/coherency, 6/registerT, 2/11, 3/010, 1/--, 5/00101, 1/0, 6/registerZ }


\paragraph{Syntax} \texttt{alandw}\emph{coherency} \emph{extend27\_\discretionary{}{}{}offset27}[\emph{registerZ}] = \emph{registerT}

\paragraph{Description} The effective address is computed by adding the \emph{extend27\_\discretionary{}{}{}offset27} to the \emph{registerZ}.
This instruction implements atomic load-AND-store on word. The word at effective address is ANDed with the \emph{registerT}, and its previous value is returned into the \emph{registerT}. The effective address must be word aligned (multiple of 4).


\paragraph{Modifier(s)}
\noindent\begin{itemize}
\item coherency (cf. coherency in section \ref{par:modifier-coherency} on page \pageref{par:modifier-coherency})
\end{itemize}

\paragraph{Execution} {Reservation table \textsf{LSU2\_MEMW\_AUXR\_AUXW.Y} (Section~\ref{sec:reservation-LSU2-MEMW-AUXR-AUXW.Y})}
~
{\begin{lstlisting}
stage ID:
  new offset = _SX_54(extend27_offset27);
  new coherency = _ZX_2(coherency);
stage RR:
  new base = GPR[registerZ];
  new address = base + offset;
stage E1:
  new argument2 = GPR[registerT];
  new result2 = _ZX_32(MEM.atomic_and(address, 0xF, coherency << 32, argument2.32[0], registerT));
stage SR:
  GPR[registerT] = result2;
\end{lstlisting}}
\newpage
\subsubsection[{\small ALB: Atomic Load Byte}]{ALB\hfill Atomic Load Byte}
\label{instr:ALB}\index{alb}
 
\paragraph{Encoding} Format \textsf{LSU\_ALSB} (Section~\ref{sec:format-LSU-ALSB}), \verb|lsucode5 = 000|\\

\bitfields{ 1/P, 2/01, 3/111, 2/coherency, 6/registerW, 2/11, 3/000, 1/--, 5/00000, 1/0, 6/registerZ }


\paragraph{Syntax} \texttt{alb}\emph{coherency} \emph{registerW} = [\emph{registerZ}]

\paragraph{Description} The effective address is the value of the \emph{registerZ}.
The \emph{registerW} is atomically loaded with zero extension from the memory byte starting at the effective address. This instruction is provided to directly operate on the last level of the memory hierarchy, which is typically a processor DDR or a cluster TCM, irrespective of the MMU DCP (Data Cache Policy) settings.


\paragraph{Modifier(s)}
\noindent\begin{itemize}
\item coherency (cf. coherency in section \ref{par:modifier-coherency} on page \pageref{par:modifier-coherency})
\end{itemize}

\paragraph{Execution} {Reservation table \textsf{LSU\_MEMW\_AUXW} (Section~\ref{sec:reservation-LSU-MEMW-AUXW})}
~
{\begin{lstlisting}
stage ID:
  new coherency = _ZX_2(coherency);
stage RR:
  new address = GPR[registerZ];
  new result1 = _ZX_8(MEM.atomic_load(address, 0x1, coherency << 32, registerW));
stage CR:
  GPR[registerW] = result1;
\end{lstlisting}}
\newpage

\paragraph{Encoding} Format \textsf{LSU\_ALSB.O} (Section~\ref{sec:format-LSU-ALSB.O}), \verb|lsucode5 = 000|\\

\bitfields{ 1/P, 2/00, 2/-- --, 27/offset27, 1/1, 2/01, 3/111, 2/coherency, 6/registerW, 2/11, 3/000, 1/--, 5/00000, 1/0, 6/registerZ }


\paragraph{Syntax} \texttt{alb}\emph{coherency} \emph{registerW} = \emph{offset27}[\emph{registerZ}]

\paragraph{Description} The effective address is computed by adding the \emph{offset27} to the \emph{registerZ}.
The \emph{registerW} is atomically loaded with zero extension from the memory byte starting at the effective address. This instruction is provided to directly operate on the last level of the memory hierarchy, which is typically a processor DDR or a cluster TCM, irrespective of the MMU DCP (Data Cache Policy) settings.


\paragraph{Modifier(s)}
\noindent\begin{itemize}
\item coherency (cf. coherency in section \ref{par:modifier-coherency} on page \pageref{par:modifier-coherency})
\end{itemize}

\paragraph{Execution} {Reservation table \textsf{LSU\_MEMW\_AUXW.X} (Section~\ref{sec:reservation-LSU-MEMW-AUXW.X})}
~
{\begin{lstlisting}
stage ID:
  new offset = _SX_27(offset27);
  new coherency = _ZX_2(coherency);
stage RR:
  new base = GPR[registerZ];
  new address = base + offset;
  new result1 = _ZX_8(MEM.atomic_load(address, 0x1, coherency << 32, registerW));
stage CR:
  GPR[registerW] = result1;
\end{lstlisting}}
\newpage

\paragraph{Encoding} Format \textsf{LSU\_ALSB.D} (Section~\ref{sec:format-LSU-ALSB.D}), \verb|lsucode5 = 000|\\

\bitfields{ 1/P, 2/00, 2/-- --, 27/extend27, 1/1, 2/00, 2/-- --, 27/offset27, 1/1, 2/01, 3/111, 2/coherency, 6/registerW, 2/11, 3/000, 1/--, 5/00000, 1/0, 6/registerZ }


\paragraph{Syntax} \texttt{alb}\emph{coherency} \emph{registerW} = \emph{extend27\_\discretionary{}{}{}offset27}[\emph{registerZ}]

\paragraph{Description} The effective address is computed by adding the \emph{extend27\_\discretionary{}{}{}offset27} to the \emph{registerZ}.
The \emph{registerW} is atomically loaded with zero extension from the memory byte starting at the effective address. This instruction is provided to directly operate on the last level of the memory hierarchy, which is typically a processor DDR or a cluster TCM, irrespective of the MMU DCP (Data Cache Policy) settings.


\paragraph{Modifier(s)}
\noindent\begin{itemize}
\item coherency (cf. coherency in section \ref{par:modifier-coherency} on page \pageref{par:modifier-coherency})
\end{itemize}

\paragraph{Execution} {Reservation table \textsf{LSU\_MEMW\_AUXW.Y} (Section~\ref{sec:reservation-LSU-MEMW-AUXW.Y})}
~
{\begin{lstlisting}
stage ID:
  new offset = _SX_54(extend27_offset27);
  new coherency = _ZX_2(coherency);
stage RR:
  new base = GPR[registerZ];
  new address = base + offset;
  new result1 = _ZX_8(MEM.atomic_load(address, 0x1, coherency << 32, registerW));
stage CR:
  GPR[registerW] = result1;
\end{lstlisting}}
\newpage
\subsubsection[{\small ALCLRB: Atomic Load and Clear Byte}]{ALCLRB\hfill Atomic Load and Clear Byte}
\label{instr:ALCLRB}\index{alclrb}
 
\paragraph{Encoding} Format \textsf{LSU\_ALCLRSB} (Section~\ref{sec:format-LSU-ALCLRSB}), \verb|lsucode5 = 000|\\

\bitfields{ 1/P, 2/01, 3/111, 2/coherency, 6/registerW, 2/11, 3/000, 1/--, 5/00001, 1/0, 6/registerZ }


\paragraph{Syntax} \texttt{alclrb}\emph{coherency} \emph{registerW} = [\emph{registerZ}]

\paragraph{Description} The effective address is the value of the \emph{registerZ}.
The \emph{registerW} is loaded from the memory byte starting at the effective address. The value 0 is stored into the memory byte starting at the effective address. This atomic instruction implements Test-and-Clear in memory. A zero byte in memory means that the lock at the effective address is taken. Unlocking is achieved by storing a non-zero byte at the effective address.


\paragraph{Modifier(s)}
\noindent\begin{itemize}
\item coherency (cf. coherency in section \ref{par:modifier-coherency} on page \pageref{par:modifier-coherency})
\end{itemize}

\paragraph{Execution} {Reservation table \textsf{LSU2\_MEMW\_AUXW} (Section~\ref{sec:reservation-LSU2-MEMW-AUXW})}
~
{\begin{lstlisting}
stage ID:
  new coherency = _ZX_2(coherency);
stage RR:
  new address = GPR[registerZ];
  new result1 = _ZX_8(MEM.atomic_swap(address, 0x1, coherency << 32, 0, registerW));
stage CR:
  GPR[registerW] = result1;
\end{lstlisting}}
\newpage

\paragraph{Encoding} Format \textsf{LSU\_ALCLRSB.O} (Section~\ref{sec:format-LSU-ALCLRSB.O}), \verb|lsucode5 = 000|\\

\bitfields{ 1/P, 2/00, 2/-- --, 27/offset27, 1/1, 2/01, 3/111, 2/coherency, 6/registerW, 2/11, 3/000, 1/--, 5/00001, 1/0, 6/registerZ }


\paragraph{Syntax} \texttt{alclrb}\emph{coherency} \emph{registerW} = \emph{offset27}[\emph{registerZ}]

\paragraph{Description} The effective address is computed by adding the \emph{offset27} to the \emph{registerZ}.
The \emph{registerW} is loaded from the memory byte starting at the effective address. The value 0 is stored into the memory byte starting at the effective address. This atomic instruction implements Test-and-Clear in memory. A zero byte in memory means that the lock at the effective address is taken. Unlocking is achieved by storing a non-zero byte at the effective address.


\paragraph{Modifier(s)}
\noindent\begin{itemize}
\item coherency (cf. coherency in section \ref{par:modifier-coherency} on page \pageref{par:modifier-coherency})
\end{itemize}

\paragraph{Execution} {Reservation table \textsf{LSU2\_MEMW\_AUXW.X} (Section~\ref{sec:reservation-LSU2-MEMW-AUXW.X})}
~
{\begin{lstlisting}
stage ID:
  new offset = _SX_27(offset27);
  new coherency = _ZX_2(coherency);
stage RR:
  new base = GPR[registerZ];
  new address = base + offset;
  new result1 = _ZX_8(MEM.atomic_swap(address, 0x1, coherency << 32, 0, registerW));
stage CR:
  GPR[registerW] = result1;
\end{lstlisting}}
\newpage

\paragraph{Encoding} Format \textsf{LSU\_ALCLRSB.D} (Section~\ref{sec:format-LSU-ALCLRSB.D}), \verb|lsucode5 = 000|\\

\bitfields{ 1/P, 2/00, 2/-- --, 27/extend27, 1/1, 2/00, 2/-- --, 27/offset27, 1/1, 2/01, 3/111, 2/coherency, 6/registerW, 2/11, 3/000, 1/--, 5/00001, 1/0, 6/registerZ }


\paragraph{Syntax} \texttt{alclrb}\emph{coherency} \emph{registerW} = \emph{extend27\_\discretionary{}{}{}offset27}[\emph{registerZ}]

\paragraph{Description} The effective address is computed by adding the \emph{extend27\_\discretionary{}{}{}offset27} to the \emph{registerZ}.
The \emph{registerW} is loaded from the memory byte starting at the effective address. The value 0 is stored into the memory byte starting at the effective address. This atomic instruction implements Test-and-Clear in memory. A zero byte in memory means that the lock at the effective address is taken. Unlocking is achieved by storing a non-zero byte at the effective address.


\paragraph{Modifier(s)}
\noindent\begin{itemize}
\item coherency (cf. coherency in section \ref{par:modifier-coherency} on page \pageref{par:modifier-coherency})
\end{itemize}

\paragraph{Execution} {Reservation table \textsf{LSU2\_MEMW\_AUXW.Y} (Section~\ref{sec:reservation-LSU2-MEMW-AUXW.Y})}
~
{\begin{lstlisting}
stage ID:
  new offset = _SX_54(extend27_offset27);
  new coherency = _ZX_2(coherency);
stage RR:
  new base = GPR[registerZ];
  new address = base + offset;
  new result1 = _ZX_8(MEM.atomic_swap(address, 0x1, coherency << 32, 0, registerW));
stage CR:
  GPR[registerW] = result1;
\end{lstlisting}}
\newpage
\subsubsection[{\small ALCLRD: Atomic Load and Clear Double Word}]{ALCLRD\hfill Atomic Load and Clear Double Word}
\label{instr:ALCLRD}\index{alclrd}
 
\paragraph{Encoding} Format \textsf{LSU\_ALCLRSB} (Section~\ref{sec:format-LSU-ALCLRSB}), \verb|lsucode5 = 011|\\

\bitfields{ 1/P, 2/01, 3/111, 2/coherency, 6/registerW, 2/11, 3/011, 1/--, 5/00001, 1/0, 6/registerZ }


\paragraph{Syntax} \texttt{alclrd}\emph{coherency} \emph{registerW} = [\emph{registerZ}]

\paragraph{Description} The effective address is the value of the \emph{registerZ}.
The \emph{registerW} is loaded from the memory double word starting at the effective address. The value 0 is stored into the memory double word starting at the effective address. This atomic instruction implements Test-and-Clear in memory. A zero double word in memory means that the lock at the effective address is taken. Unlocking is achieved by storing a non-zero double word at the effective address.


\paragraph{Modifier(s)}
\noindent\begin{itemize}
\item coherency (cf. coherency in section \ref{par:modifier-coherency} on page \pageref{par:modifier-coherency})
\end{itemize}

\paragraph{Execution} {Reservation table \textsf{LSU2\_MEMW\_AUXW} (Section~\ref{sec:reservation-LSU2-MEMW-AUXW})}
~
{\begin{lstlisting}
stage ID:
  new coherency = _ZX_2(coherency);
stage RR:
  new address = GPR[registerZ];
  new result1 = MEM.atomic_swap(address, 0xFF, coherency << 32, 0, registerW);
stage CR:
  GPR[registerW] = result1;
\end{lstlisting}}
\newpage

\paragraph{Encoding} Format \textsf{LSU\_ALCLRSB.O} (Section~\ref{sec:format-LSU-ALCLRSB.O}), \verb|lsucode5 = 011|\\

\bitfields{ 1/P, 2/00, 2/-- --, 27/offset27, 1/1, 2/01, 3/111, 2/coherency, 6/registerW, 2/11, 3/011, 1/--, 5/00001, 1/0, 6/registerZ }


\paragraph{Syntax} \texttt{alclrd}\emph{coherency} \emph{registerW} = \emph{offset27}[\emph{registerZ}]

\paragraph{Description} The effective address is computed by adding the \emph{offset27} to the \emph{registerZ}.
The \emph{registerW} is loaded from the memory double word starting at the effective address. The value 0 is stored into the memory double word starting at the effective address. This atomic instruction implements Test-and-Clear in memory. A zero double word in memory means that the lock at the effective address is taken. Unlocking is achieved by storing a non-zero double word at the effective address.


\paragraph{Modifier(s)}
\noindent\begin{itemize}
\item coherency (cf. coherency in section \ref{par:modifier-coherency} on page \pageref{par:modifier-coherency})
\end{itemize}

\paragraph{Execution} {Reservation table \textsf{LSU2\_MEMW\_AUXW.X} (Section~\ref{sec:reservation-LSU2-MEMW-AUXW.X})}
~
{\begin{lstlisting}
stage ID:
  new offset = _SX_27(offset27);
  new coherency = _ZX_2(coherency);
stage RR:
  new base = GPR[registerZ];
  new address = base + offset;
  new result1 = MEM.atomic_swap(address, 0xFF, coherency << 32, 0, registerW);
stage CR:
  GPR[registerW] = result1;
\end{lstlisting}}
\newpage

\paragraph{Encoding} Format \textsf{LSU\_ALCLRSB.D} (Section~\ref{sec:format-LSU-ALCLRSB.D}), \verb|lsucode5 = 011|\\

\bitfields{ 1/P, 2/00, 2/-- --, 27/extend27, 1/1, 2/00, 2/-- --, 27/offset27, 1/1, 2/01, 3/111, 2/coherency, 6/registerW, 2/11, 3/011, 1/--, 5/00001, 1/0, 6/registerZ }


\paragraph{Syntax} \texttt{alclrd}\emph{coherency} \emph{registerW} = \emph{extend27\_\discretionary{}{}{}offset27}[\emph{registerZ}]

\paragraph{Description} The effective address is computed by adding the \emph{extend27\_\discretionary{}{}{}offset27} to the \emph{registerZ}.
The \emph{registerW} is loaded from the memory double word starting at the effective address. The value 0 is stored into the memory double word starting at the effective address. This atomic instruction implements Test-and-Clear in memory. A zero double word in memory means that the lock at the effective address is taken. Unlocking is achieved by storing a non-zero double word at the effective address.


\paragraph{Modifier(s)}
\noindent\begin{itemize}
\item coherency (cf. coherency in section \ref{par:modifier-coherency} on page \pageref{par:modifier-coherency})
\end{itemize}

\paragraph{Execution} {Reservation table \textsf{LSU2\_MEMW\_AUXW.Y} (Section~\ref{sec:reservation-LSU2-MEMW-AUXW.Y})}
~
{\begin{lstlisting}
stage ID:
  new offset = _SX_54(extend27_offset27);
  new coherency = _ZX_2(coherency);
stage RR:
  new base = GPR[registerZ];
  new address = base + offset;
  new result1 = MEM.atomic_swap(address, 0xFF, coherency << 32, 0, registerW);
stage CR:
  GPR[registerW] = result1;
\end{lstlisting}}
\newpage
\subsubsection[{\small ALCLRH: Atomic Load and Clear Half Word}]{ALCLRH\hfill Atomic Load and Clear Half Word}
\label{instr:ALCLRH}\index{alclrh}
 
\paragraph{Encoding} Format \textsf{LSU\_ALCLRSB} (Section~\ref{sec:format-LSU-ALCLRSB}), \verb|lsucode5 = 001|\\

\bitfields{ 1/P, 2/01, 3/111, 2/coherency, 6/registerW, 2/11, 3/001, 1/--, 5/00001, 1/0, 6/registerZ }


\paragraph{Syntax} \texttt{alclrh}\emph{coherency} \emph{registerW} = [\emph{registerZ}]

\paragraph{Description} The effective address is the value of the \emph{registerZ}.
The \emph{registerW} is loaded from the memory half word starting at the effective address. The value 0 is stored into the memory half word starting at the effective address. This atomic instruction implements Test-and-Clear in memory. A zero half word in memory means that the lock at the effective address is taken. Unlocking is achieved by storing a non-zero half word at the effective address.


\paragraph{Modifier(s)}
\noindent\begin{itemize}
\item coherency (cf. coherency in section \ref{par:modifier-coherency} on page \pageref{par:modifier-coherency})
\end{itemize}

\paragraph{Execution} {Reservation table \textsf{LSU2\_MEMW\_AUXW} (Section~\ref{sec:reservation-LSU2-MEMW-AUXW})}
~
{\begin{lstlisting}
stage ID:
  new coherency = _ZX_2(coherency);
stage RR:
  new address = GPR[registerZ];
  new result1 = _ZX_16(MEM.atomic_swap(address, 0x3, coherency << 32, 0, registerW));
stage CR:
  GPR[registerW] = result1;
\end{lstlisting}}
\newpage

\paragraph{Encoding} Format \textsf{LSU\_ALCLRSB.O} (Section~\ref{sec:format-LSU-ALCLRSB.O}), \verb|lsucode5 = 001|\\

\bitfields{ 1/P, 2/00, 2/-- --, 27/offset27, 1/1, 2/01, 3/111, 2/coherency, 6/registerW, 2/11, 3/001, 1/--, 5/00001, 1/0, 6/registerZ }


\paragraph{Syntax} \texttt{alclrh}\emph{coherency} \emph{registerW} = \emph{offset27}[\emph{registerZ}]

\paragraph{Description} The effective address is computed by adding the \emph{offset27} to the \emph{registerZ}.
The \emph{registerW} is loaded from the memory half word starting at the effective address. The value 0 is stored into the memory half word starting at the effective address. This atomic instruction implements Test-and-Clear in memory. A zero half word in memory means that the lock at the effective address is taken. Unlocking is achieved by storing a non-zero half word at the effective address.


\paragraph{Modifier(s)}
\noindent\begin{itemize}
\item coherency (cf. coherency in section \ref{par:modifier-coherency} on page \pageref{par:modifier-coherency})
\end{itemize}

\paragraph{Execution} {Reservation table \textsf{LSU2\_MEMW\_AUXW.X} (Section~\ref{sec:reservation-LSU2-MEMW-AUXW.X})}
~
{\begin{lstlisting}
stage ID:
  new offset = _SX_27(offset27);
  new coherency = _ZX_2(coherency);
stage RR:
  new base = GPR[registerZ];
  new address = base + offset;
  new result1 = _ZX_16(MEM.atomic_swap(address, 0x3, coherency << 32, 0, registerW));
stage CR:
  GPR[registerW] = result1;
\end{lstlisting}}
\newpage

\paragraph{Encoding} Format \textsf{LSU\_ALCLRSB.D} (Section~\ref{sec:format-LSU-ALCLRSB.D}), \verb|lsucode5 = 001|\\

\bitfields{ 1/P, 2/00, 2/-- --, 27/extend27, 1/1, 2/00, 2/-- --, 27/offset27, 1/1, 2/01, 3/111, 2/coherency, 6/registerW, 2/11, 3/001, 1/--, 5/00001, 1/0, 6/registerZ }


\paragraph{Syntax} \texttt{alclrh}\emph{coherency} \emph{registerW} = \emph{extend27\_\discretionary{}{}{}offset27}[\emph{registerZ}]

\paragraph{Description} The effective address is computed by adding the \emph{extend27\_\discretionary{}{}{}offset27} to the \emph{registerZ}.
The \emph{registerW} is loaded from the memory half word starting at the effective address. The value 0 is stored into the memory half word starting at the effective address. This atomic instruction implements Test-and-Clear in memory. A zero half word in memory means that the lock at the effective address is taken. Unlocking is achieved by storing a non-zero half word at the effective address.


\paragraph{Modifier(s)}
\noindent\begin{itemize}
\item coherency (cf. coherency in section \ref{par:modifier-coherency} on page \pageref{par:modifier-coherency})
\end{itemize}

\paragraph{Execution} {Reservation table \textsf{LSU2\_MEMW\_AUXW.Y} (Section~\ref{sec:reservation-LSU2-MEMW-AUXW.Y})}
~
{\begin{lstlisting}
stage ID:
  new offset = _SX_54(extend27_offset27);
  new coherency = _ZX_2(coherency);
stage RR:
  new base = GPR[registerZ];
  new address = base + offset;
  new result1 = _ZX_16(MEM.atomic_swap(address, 0x3, coherency << 32, 0, registerW));
stage CR:
  GPR[registerW] = result1;
\end{lstlisting}}
\newpage
\subsubsection[{\small ALCLRW: Atomic Load and Clear Word}]{ALCLRW\hfill Atomic Load and Clear Word}
\label{instr:ALCLRW}\index{alclrw}
 
\paragraph{Encoding} Format \textsf{LSU\_ALCLRSB} (Section~\ref{sec:format-LSU-ALCLRSB}), \verb|lsucode5 = 010|\\

\bitfields{ 1/P, 2/01, 3/111, 2/coherency, 6/registerW, 2/11, 3/010, 1/--, 5/00001, 1/0, 6/registerZ }


\paragraph{Syntax} \texttt{alclrw}\emph{coherency} \emph{registerW} = [\emph{registerZ}]

\paragraph{Description} The effective address is the value of the \emph{registerZ}.
The \emph{registerW} is loaded from the memory word starting at the effective address. The value 0 is stored into the memory word starting at the effective address. This atomic instruction implements Test-and-Clear in memory. A zero word in memory means that the lock at the effective address is taken. Unlocking is achieved by storing a non-zero word at the effective address.


\paragraph{Modifier(s)}
\noindent\begin{itemize}
\item coherency (cf. coherency in section \ref{par:modifier-coherency} on page \pageref{par:modifier-coherency})
\end{itemize}

\paragraph{Execution} {Reservation table \textsf{LSU2\_MEMW\_AUXW} (Section~\ref{sec:reservation-LSU2-MEMW-AUXW})}
~
{\begin{lstlisting}
stage ID:
  new coherency = _ZX_2(coherency);
stage RR:
  new address = GPR[registerZ];
  new result1 = _ZX_32(MEM.atomic_swap(address, 0xF, coherency << 32, 0, registerW));
stage CR:
  GPR[registerW] = result1;
\end{lstlisting}}
\newpage

\paragraph{Encoding} Format \textsf{LSU\_ALCLRSB.O} (Section~\ref{sec:format-LSU-ALCLRSB.O}), \verb|lsucode5 = 010|\\

\bitfields{ 1/P, 2/00, 2/-- --, 27/offset27, 1/1, 2/01, 3/111, 2/coherency, 6/registerW, 2/11, 3/010, 1/--, 5/00001, 1/0, 6/registerZ }


\paragraph{Syntax} \texttt{alclrw}\emph{coherency} \emph{registerW} = \emph{offset27}[\emph{registerZ}]

\paragraph{Description} The effective address is computed by adding the \emph{offset27} to the \emph{registerZ}.
The \emph{registerW} is loaded from the memory word starting at the effective address. The value 0 is stored into the memory word starting at the effective address. This atomic instruction implements Test-and-Clear in memory. A zero word in memory means that the lock at the effective address is taken. Unlocking is achieved by storing a non-zero word at the effective address.


\paragraph{Modifier(s)}
\noindent\begin{itemize}
\item coherency (cf. coherency in section \ref{par:modifier-coherency} on page \pageref{par:modifier-coherency})
\end{itemize}

\paragraph{Execution} {Reservation table \textsf{LSU2\_MEMW\_AUXW.X} (Section~\ref{sec:reservation-LSU2-MEMW-AUXW.X})}
~
{\begin{lstlisting}
stage ID:
  new offset = _SX_27(offset27);
  new coherency = _ZX_2(coherency);
stage RR:
  new base = GPR[registerZ];
  new address = base + offset;
  new result1 = _ZX_32(MEM.atomic_swap(address, 0xF, coherency << 32, 0, registerW));
stage CR:
  GPR[registerW] = result1;
\end{lstlisting}}
\newpage

\paragraph{Encoding} Format \textsf{LSU\_ALCLRSB.D} (Section~\ref{sec:format-LSU-ALCLRSB.D}), \verb|lsucode5 = 010|\\

\bitfields{ 1/P, 2/00, 2/-- --, 27/extend27, 1/1, 2/00, 2/-- --, 27/offset27, 1/1, 2/01, 3/111, 2/coherency, 6/registerW, 2/11, 3/010, 1/--, 5/00001, 1/0, 6/registerZ }


\paragraph{Syntax} \texttt{alclrw}\emph{coherency} \emph{registerW} = \emph{extend27\_\discretionary{}{}{}offset27}[\emph{registerZ}]

\paragraph{Description} The effective address is computed by adding the \emph{extend27\_\discretionary{}{}{}offset27} to the \emph{registerZ}.
The \emph{registerW} is loaded from the memory word starting at the effective address. The value 0 is stored into the memory word starting at the effective address. This atomic instruction implements Test-and-Clear in memory. A zero word in memory means that the lock at the effective address is taken. Unlocking is achieved by storing a non-zero word at the effective address.


\paragraph{Modifier(s)}
\noindent\begin{itemize}
\item coherency (cf. coherency in section \ref{par:modifier-coherency} on page \pageref{par:modifier-coherency})
\end{itemize}

\paragraph{Execution} {Reservation table \textsf{LSU2\_MEMW\_AUXW.Y} (Section~\ref{sec:reservation-LSU2-MEMW-AUXW.Y})}
~
{\begin{lstlisting}
stage ID:
  new offset = _SX_54(extend27_offset27);
  new coherency = _ZX_2(coherency);
stage RR:
  new base = GPR[registerZ];
  new address = base + offset;
  new result1 = _ZX_32(MEM.atomic_swap(address, 0xF, coherency << 32, 0, registerW));
stage CR:
  GPR[registerW] = result1;
\end{lstlisting}}
\newpage
\subsubsection[{\small ALD: Atomic Load Double Word}]{ALD\hfill Atomic Load Double Word}
\label{instr:ALD}\index{ald}
 
\paragraph{Encoding} Format \textsf{LSU\_ALSB} (Section~\ref{sec:format-LSU-ALSB}), \verb|lsucode5 = 011|\\

\bitfields{ 1/P, 2/01, 3/111, 2/coherency, 6/registerW, 2/11, 3/011, 1/--, 5/00000, 1/0, 6/registerZ }


\paragraph{Syntax} \texttt{ald}\emph{coherency} \emph{registerW} = [\emph{registerZ}]

\paragraph{Description} The effective address is the value of the \emph{registerZ}.
The \emph{registerW} is atomically loaded from the memory double word starting at the effective address. This instruction is provided to directly operate on the last level of the memory hierarchy, which is typically a processor DDR or a cluster TCM, irrespective of the MMU DCP (Data Cache Policy) settings. The effective address must be double word aligned (multiple of 8).


\paragraph{Modifier(s)}
\noindent\begin{itemize}
\item coherency (cf. coherency in section \ref{par:modifier-coherency} on page \pageref{par:modifier-coherency})
\end{itemize}

\paragraph{Execution} {Reservation table \textsf{LSU\_MEMW\_AUXW} (Section~\ref{sec:reservation-LSU-MEMW-AUXW})}
~
{\begin{lstlisting}
stage ID:
  new coherency = _ZX_2(coherency);
stage RR:
  new address = GPR[registerZ];
  new result1 = MEM.atomic_load(address, 0xFF, coherency << 32, registerW);
stage CR:
  GPR[registerW] = result1;
\end{lstlisting}}
\newpage

\paragraph{Encoding} Format \textsf{LSU\_ALSB.O} (Section~\ref{sec:format-LSU-ALSB.O}), \verb|lsucode5 = 011|\\

\bitfields{ 1/P, 2/00, 2/-- --, 27/offset27, 1/1, 2/01, 3/111, 2/coherency, 6/registerW, 2/11, 3/011, 1/--, 5/00000, 1/0, 6/registerZ }


\paragraph{Syntax} \texttt{ald}\emph{coherency} \emph{registerW} = \emph{offset27}[\emph{registerZ}]

\paragraph{Description} The effective address is computed by adding the \emph{offset27} to the \emph{registerZ}.
The \emph{registerW} is atomically loaded from the memory double word starting at the effective address. This instruction is provided to directly operate on the last level of the memory hierarchy, which is typically a processor DDR or a cluster TCM, irrespective of the MMU DCP (Data Cache Policy) settings. The effective address must be double word aligned (multiple of 8).


\paragraph{Modifier(s)}
\noindent\begin{itemize}
\item coherency (cf. coherency in section \ref{par:modifier-coherency} on page \pageref{par:modifier-coherency})
\end{itemize}

\paragraph{Execution} {Reservation table \textsf{LSU\_MEMW\_AUXW.X} (Section~\ref{sec:reservation-LSU-MEMW-AUXW.X})}
~
{\begin{lstlisting}
stage ID:
  new offset = _SX_27(offset27);
  new coherency = _ZX_2(coherency);
stage RR:
  new base = GPR[registerZ];
  new address = base + offset;
  new result1 = MEM.atomic_load(address, 0xFF, coherency << 32, registerW);
stage CR:
  GPR[registerW] = result1;
\end{lstlisting}}
\newpage

\paragraph{Encoding} Format \textsf{LSU\_ALSB.D} (Section~\ref{sec:format-LSU-ALSB.D}), \verb|lsucode5 = 011|\\

\bitfields{ 1/P, 2/00, 2/-- --, 27/extend27, 1/1, 2/00, 2/-- --, 27/offset27, 1/1, 2/01, 3/111, 2/coherency, 6/registerW, 2/11, 3/011, 1/--, 5/00000, 1/0, 6/registerZ }


\paragraph{Syntax} \texttt{ald}\emph{coherency} \emph{registerW} = \emph{extend27\_\discretionary{}{}{}offset27}[\emph{registerZ}]

\paragraph{Description} The effective address is computed by adding the \emph{extend27\_\discretionary{}{}{}offset27} to the \emph{registerZ}.
The \emph{registerW} is atomically loaded from the memory double word starting at the effective address. This instruction is provided to directly operate on the last level of the memory hierarchy, which is typically a processor DDR or a cluster TCM, irrespective of the MMU DCP (Data Cache Policy) settings. The effective address must be double word aligned (multiple of 8).


\paragraph{Modifier(s)}
\noindent\begin{itemize}
\item coherency (cf. coherency in section \ref{par:modifier-coherency} on page \pageref{par:modifier-coherency})
\end{itemize}

\paragraph{Execution} {Reservation table \textsf{LSU\_MEMW\_AUXW.Y} (Section~\ref{sec:reservation-LSU-MEMW-AUXW.Y})}
~
{\begin{lstlisting}
stage ID:
  new offset = _SX_54(extend27_offset27);
  new coherency = _ZX_2(coherency);
stage RR:
  new base = GPR[registerZ];
  new address = base + offset;
  new result1 = MEM.atomic_load(address, 0xFF, coherency << 32, registerW);
stage CR:
  GPR[registerW] = result1;
\end{lstlisting}}
\newpage
\subsubsection[{\small ALDUSB: Atomic Load and Decrement Unsigned Saturating Byte}]{ALDUSB\hfill Atomic Load and Decrement Unsigned Saturating Byte}
\label{instr:ALDUSB}\index{aldusb}
 
\paragraph{Encoding} Format \textsf{LSU\_ALDUSSB} (Section~\ref{sec:format-LSU-ALDUSSB}), \verb|lsucode5 = 000|\\

\bitfields{ 1/P, 2/01, 3/111, 2/coherency, 6/registerT, 2/11, 3/000, 1/--, 5/01110, 1/0, 6/registerZ }


\paragraph{Syntax} \texttt{aldusb}\emph{coherency} [\emph{registerZ}] = \emph{registerT}

\paragraph{Description} The effective address is the value of the \emph{registerZ}.
This instruction implements atomic load-DUS-store on signed byte (Decrement Unsigned Saturating). The byte at effective address is DUSed with the \emph{registerT}, and its previous value is returned into the \emph{registerT}.


\paragraph{Modifier(s)}
\noindent\begin{itemize}
\item coherency (cf. coherency in section \ref{par:modifier-coherency} on page \pageref{par:modifier-coherency})
\end{itemize}

\paragraph{Execution} {Reservation table \textsf{LSU2\_MEMW\_AUXR\_AUXW} (Section~\ref{sec:reservation-LSU2-MEMW-AUXR-AUXW})}
~
{\begin{lstlisting}
stage ID:
  new coherency = _ZX_2(coherency);
stage RR:
  new address = GPR[registerZ];
stage E1:
  new argument2 = GPR[registerT];
  new result2 = _ZX_8(MEM.atomic_dus(address, 0x1, coherency << 32, argument2.8[0], registerT));
stage SR:
  GPR[registerT] = result2;
\end{lstlisting}}
\newpage

\paragraph{Encoding} Format \textsf{LSU\_ALDUSSB.O} (Section~\ref{sec:format-LSU-ALDUSSB.O}), \verb|lsucode5 = 000|\\

\bitfields{ 1/P, 2/00, 2/-- --, 27/offset27, 1/1, 2/01, 3/111, 2/coherency, 6/registerT, 2/11, 3/000, 1/--, 5/01110, 1/0, 6/registerZ }


\paragraph{Syntax} \texttt{aldusb}\emph{coherency} \emph{offset27}[\emph{registerZ}] = \emph{registerT}

\paragraph{Description} The effective address is computed by adding the \emph{offset27} to the \emph{registerZ}.
This instruction implements atomic load-DUS-store on signed byte (Decrement Unsigned Saturating). The byte at effective address is DUSed with the \emph{registerT}, and its previous value is returned into the \emph{registerT}.


\paragraph{Modifier(s)}
\noindent\begin{itemize}
\item coherency (cf. coherency in section \ref{par:modifier-coherency} on page \pageref{par:modifier-coherency})
\end{itemize}

\paragraph{Execution} {Reservation table \textsf{LSU2\_MEMW\_AUXR\_AUXW.X} (Section~\ref{sec:reservation-LSU2-MEMW-AUXR-AUXW.X})}
~
{\begin{lstlisting}
stage ID:
  new offset = _SX_27(offset27);
  new coherency = _ZX_2(coherency);
stage RR:
  new base = GPR[registerZ];
  new address = base + offset;
stage E1:
  new argument2 = GPR[registerT];
  new result2 = _ZX_8(MEM.atomic_dus(address, 0x1, coherency << 32, argument2.8[0], registerT));
stage SR:
  GPR[registerT] = result2;
\end{lstlisting}}
\newpage

\paragraph{Encoding} Format \textsf{LSU\_ALDUSSB.D} (Section~\ref{sec:format-LSU-ALDUSSB.D}), \verb|lsucode5 = 000|\\

\bitfields{ 1/P, 2/00, 2/-- --, 27/extend27, 1/1, 2/00, 2/-- --, 27/offset27, 1/1, 2/01, 3/111, 2/coherency, 6/registerT, 2/11, 3/000, 1/--, 5/01110, 1/0, 6/registerZ }


\paragraph{Syntax} \texttt{aldusb}\emph{coherency} \emph{extend27\_\discretionary{}{}{}offset27}[\emph{registerZ}] = \emph{registerT}

\paragraph{Description} The effective address is computed by adding the \emph{extend27\_\discretionary{}{}{}offset27} to the \emph{registerZ}.
This instruction implements atomic load-DUS-store on signed byte (Decrement Unsigned Saturating). The byte at effective address is DUSed with the \emph{registerT}, and its previous value is returned into the \emph{registerT}.


\paragraph{Modifier(s)}
\noindent\begin{itemize}
\item coherency (cf. coherency in section \ref{par:modifier-coherency} on page \pageref{par:modifier-coherency})
\end{itemize}

\paragraph{Execution} {Reservation table \textsf{LSU2\_MEMW\_AUXR\_AUXW.Y} (Section~\ref{sec:reservation-LSU2-MEMW-AUXR-AUXW.Y})}
~
{\begin{lstlisting}
stage ID:
  new offset = _SX_54(extend27_offset27);
  new coherency = _ZX_2(coherency);
stage RR:
  new base = GPR[registerZ];
  new address = base + offset;
stage E1:
  new argument2 = GPR[registerT];
  new result2 = _ZX_8(MEM.atomic_dus(address, 0x1, coherency << 32, argument2.8[0], registerT));
stage SR:
  GPR[registerT] = result2;
\end{lstlisting}}
\newpage
\subsubsection[{\small ALDUSD: Atomic Load and DUS Double Word}]{ALDUSD\hfill Atomic Load and DUS Double Word}
\label{instr:ALDUSD}\index{aldusd}
 
\paragraph{Encoding} Format \textsf{LSU\_ALDUSSB} (Section~\ref{sec:format-LSU-ALDUSSB}), \verb|lsucode5 = 011|\\

\bitfields{ 1/P, 2/01, 3/111, 2/coherency, 6/registerT, 2/11, 3/011, 1/--, 5/01110, 1/0, 6/registerZ }


\paragraph{Syntax} \texttt{aldusd}\emph{coherency} [\emph{registerZ}] = \emph{registerT}

\paragraph{Description} The effective address is the value of the \emph{registerZ}.
This instruction implements atomic load-DUS-store on signed double word (Decrement Unsigned Saturating). The double word at effective address is DUSed with the \emph{registerT}, and its previous value is returned into the \emph{registerT}. The effective address must be double word aligned (multiple of 8).


\paragraph{Modifier(s)}
\noindent\begin{itemize}
\item coherency (cf. coherency in section \ref{par:modifier-coherency} on page \pageref{par:modifier-coherency})
\end{itemize}

\paragraph{Execution} {Reservation table \textsf{LSU2\_MEMW\_AUXR\_AUXW} (Section~\ref{sec:reservation-LSU2-MEMW-AUXR-AUXW})}
~
{\begin{lstlisting}
stage ID:
  new coherency = _ZX_2(coherency);
stage RR:
  new address = GPR[registerZ];
stage E1:
  new argument2 = GPR[registerT];
  new result2 = MEM.atomic_dus(address, 0xFF, coherency << 32, argument2.64[0], registerT);
stage SR:
  GPR[registerT] = result2;
\end{lstlisting}}
\newpage

\paragraph{Encoding} Format \textsf{LSU\_ALDUSSB.O} (Section~\ref{sec:format-LSU-ALDUSSB.O}), \verb|lsucode5 = 011|\\

\bitfields{ 1/P, 2/00, 2/-- --, 27/offset27, 1/1, 2/01, 3/111, 2/coherency, 6/registerT, 2/11, 3/011, 1/--, 5/01110, 1/0, 6/registerZ }


\paragraph{Syntax} \texttt{aldusd}\emph{coherency} \emph{offset27}[\emph{registerZ}] = \emph{registerT}

\paragraph{Description} The effective address is computed by adding the \emph{offset27} to the \emph{registerZ}.
This instruction implements atomic load-DUS-store on signed double word (Decrement Unsigned Saturating). The double word at effective address is DUSed with the \emph{registerT}, and its previous value is returned into the \emph{registerT}. The effective address must be double word aligned (multiple of 8).


\paragraph{Modifier(s)}
\noindent\begin{itemize}
\item coherency (cf. coherency in section \ref{par:modifier-coherency} on page \pageref{par:modifier-coherency})
\end{itemize}

\paragraph{Execution} {Reservation table \textsf{LSU2\_MEMW\_AUXR\_AUXW.X} (Section~\ref{sec:reservation-LSU2-MEMW-AUXR-AUXW.X})}
~
{\begin{lstlisting}
stage ID:
  new offset = _SX_27(offset27);
  new coherency = _ZX_2(coherency);
stage RR:
  new base = GPR[registerZ];
  new address = base + offset;
stage E1:
  new argument2 = GPR[registerT];
  new result2 = MEM.atomic_dus(address, 0xFF, coherency << 32, argument2.64[0], registerT);
stage SR:
  GPR[registerT] = result2;
\end{lstlisting}}
\newpage

\paragraph{Encoding} Format \textsf{LSU\_ALDUSSB.D} (Section~\ref{sec:format-LSU-ALDUSSB.D}), \verb|lsucode5 = 011|\\

\bitfields{ 1/P, 2/00, 2/-- --, 27/extend27, 1/1, 2/00, 2/-- --, 27/offset27, 1/1, 2/01, 3/111, 2/coherency, 6/registerT, 2/11, 3/011, 1/--, 5/01110, 1/0, 6/registerZ }


\paragraph{Syntax} \texttt{aldusd}\emph{coherency} \emph{extend27\_\discretionary{}{}{}offset27}[\emph{registerZ}] = \emph{registerT}

\paragraph{Description} The effective address is computed by adding the \emph{extend27\_\discretionary{}{}{}offset27} to the \emph{registerZ}.
This instruction implements atomic load-DUS-store on signed double word (Decrement Unsigned Saturating). The double word at effective address is DUSed with the \emph{registerT}, and its previous value is returned into the \emph{registerT}. The effective address must be double word aligned (multiple of 8).


\paragraph{Modifier(s)}
\noindent\begin{itemize}
\item coherency (cf. coherency in section \ref{par:modifier-coherency} on page \pageref{par:modifier-coherency})
\end{itemize}

\paragraph{Execution} {Reservation table \textsf{LSU2\_MEMW\_AUXR\_AUXW.Y} (Section~\ref{sec:reservation-LSU2-MEMW-AUXR-AUXW.Y})}
~
{\begin{lstlisting}
stage ID:
  new offset = _SX_54(extend27_offset27);
  new coherency = _ZX_2(coherency);
stage RR:
  new base = GPR[registerZ];
  new address = base + offset;
stage E1:
  new argument2 = GPR[registerT];
  new result2 = MEM.atomic_dus(address, 0xFF, coherency << 32, argument2.64[0], registerT);
stage SR:
  GPR[registerT] = result2;
\end{lstlisting}}
\newpage
\subsubsection[{\small ALDUSH: Atomic Load and DUS Half Word}]{ALDUSH\hfill Atomic Load and DUS Half Word}
\label{instr:ALDUSH}\index{aldush}
 
\paragraph{Encoding} Format \textsf{LSU\_ALDUSSB} (Section~\ref{sec:format-LSU-ALDUSSB}), \verb|lsucode5 = 001|\\

\bitfields{ 1/P, 2/01, 3/111, 2/coherency, 6/registerT, 2/11, 3/001, 1/--, 5/01110, 1/0, 6/registerZ }


\paragraph{Syntax} \texttt{aldush}\emph{coherency} [\emph{registerZ}] = \emph{registerT}

\paragraph{Description} The effective address is the value of the \emph{registerZ}.
This instruction implements atomic load-DUS-store on signed half word (Decrement Unsigned Saturating). The half word at effective address is DUSed with the \emph{registerT}, and its previous value is returned into the \emph{registerT}. The effective address must be half word aligned (multiple of 2).


\paragraph{Modifier(s)}
\noindent\begin{itemize}
\item coherency (cf. coherency in section \ref{par:modifier-coherency} on page \pageref{par:modifier-coherency})
\end{itemize}

\paragraph{Execution} {Reservation table \textsf{LSU2\_MEMW\_AUXR\_AUXW} (Section~\ref{sec:reservation-LSU2-MEMW-AUXR-AUXW})}
~
{\begin{lstlisting}
stage ID:
  new coherency = _ZX_2(coherency);
stage RR:
  new address = GPR[registerZ];
stage E1:
  new argument2 = GPR[registerT];
  new result2 = _ZX_16(MEM.atomic_dus(address, 0x3, coherency << 32, argument2.16[0], registerT));
stage SR:
  GPR[registerT] = result2;
\end{lstlisting}}
\newpage

\paragraph{Encoding} Format \textsf{LSU\_ALDUSSB.O} (Section~\ref{sec:format-LSU-ALDUSSB.O}), \verb|lsucode5 = 001|\\

\bitfields{ 1/P, 2/00, 2/-- --, 27/offset27, 1/1, 2/01, 3/111, 2/coherency, 6/registerT, 2/11, 3/001, 1/--, 5/01110, 1/0, 6/registerZ }


\paragraph{Syntax} \texttt{aldush}\emph{coherency} \emph{offset27}[\emph{registerZ}] = \emph{registerT}

\paragraph{Description} The effective address is computed by adding the \emph{offset27} to the \emph{registerZ}.
This instruction implements atomic load-DUS-store on signed half word (Decrement Unsigned Saturating). The half word at effective address is DUSed with the \emph{registerT}, and its previous value is returned into the \emph{registerT}. The effective address must be half word aligned (multiple of 2).


\paragraph{Modifier(s)}
\noindent\begin{itemize}
\item coherency (cf. coherency in section \ref{par:modifier-coherency} on page \pageref{par:modifier-coherency})
\end{itemize}

\paragraph{Execution} {Reservation table \textsf{LSU2\_MEMW\_AUXR\_AUXW.X} (Section~\ref{sec:reservation-LSU2-MEMW-AUXR-AUXW.X})}
~
{\begin{lstlisting}
stage ID:
  new offset = _SX_27(offset27);
  new coherency = _ZX_2(coherency);
stage RR:
  new base = GPR[registerZ];
  new address = base + offset;
stage E1:
  new argument2 = GPR[registerT];
  new result2 = _ZX_16(MEM.atomic_dus(address, 0x3, coherency << 32, argument2.16[0], registerT));
stage SR:
  GPR[registerT] = result2;
\end{lstlisting}}
\newpage

\paragraph{Encoding} Format \textsf{LSU\_ALDUSSB.D} (Section~\ref{sec:format-LSU-ALDUSSB.D}), \verb|lsucode5 = 001|\\

\bitfields{ 1/P, 2/00, 2/-- --, 27/extend27, 1/1, 2/00, 2/-- --, 27/offset27, 1/1, 2/01, 3/111, 2/coherency, 6/registerT, 2/11, 3/001, 1/--, 5/01110, 1/0, 6/registerZ }


\paragraph{Syntax} \texttt{aldush}\emph{coherency} \emph{extend27\_\discretionary{}{}{}offset27}[\emph{registerZ}] = \emph{registerT}

\paragraph{Description} The effective address is computed by adding the \emph{extend27\_\discretionary{}{}{}offset27} to the \emph{registerZ}.
This instruction implements atomic load-DUS-store on signed half word (Decrement Unsigned Saturating). The half word at effective address is DUSed with the \emph{registerT}, and its previous value is returned into the \emph{registerT}. The effective address must be half word aligned (multiple of 2).


\paragraph{Modifier(s)}
\noindent\begin{itemize}
\item coherency (cf. coherency in section \ref{par:modifier-coherency} on page \pageref{par:modifier-coherency})
\end{itemize}

\paragraph{Execution} {Reservation table \textsf{LSU2\_MEMW\_AUXR\_AUXW.Y} (Section~\ref{sec:reservation-LSU2-MEMW-AUXR-AUXW.Y})}
~
{\begin{lstlisting}
stage ID:
  new offset = _SX_54(extend27_offset27);
  new coherency = _ZX_2(coherency);
stage RR:
  new base = GPR[registerZ];
  new address = base + offset;
stage E1:
  new argument2 = GPR[registerT];
  new result2 = _ZX_16(MEM.atomic_dus(address, 0x3, coherency << 32, argument2.16[0], registerT));
stage SR:
  GPR[registerT] = result2;
\end{lstlisting}}
\newpage
\subsubsection[{\small ALDUSW: Atomic Load and DUS Word}]{ALDUSW\hfill Atomic Load and DUS Word}
\label{instr:ALDUSW}\index{aldusw}
 
\paragraph{Encoding} Format \textsf{LSU\_ALDUSSB} (Section~\ref{sec:format-LSU-ALDUSSB}), \verb|lsucode5 = 010|\\

\bitfields{ 1/P, 2/01, 3/111, 2/coherency, 6/registerT, 2/11, 3/010, 1/--, 5/01110, 1/0, 6/registerZ }


\paragraph{Syntax} \texttt{aldusw}\emph{coherency} [\emph{registerZ}] = \emph{registerT}

\paragraph{Description} The effective address is the value of the \emph{registerZ}.
This instruction implements atomic load-DUS-store on signed word (Decrement Unsigned Saturating). The word at effective address is DUSed with the \emph{registerT}, and its previous value is returned into the \emph{registerT}. The effective address must be word aligned (multiple of 4).


\paragraph{Modifier(s)}
\noindent\begin{itemize}
\item coherency (cf. coherency in section \ref{par:modifier-coherency} on page \pageref{par:modifier-coherency})
\end{itemize}

\paragraph{Execution} {Reservation table \textsf{LSU2\_MEMW\_AUXR\_AUXW} (Section~\ref{sec:reservation-LSU2-MEMW-AUXR-AUXW})}
~
{\begin{lstlisting}
stage ID:
  new coherency = _ZX_2(coherency);
stage RR:
  new address = GPR[registerZ];
stage E1:
  new argument2 = GPR[registerT];
  new result2 = _ZX_32(MEM.atomic_dus(address, 0xF, coherency << 32, argument2.32[0], registerT));
stage SR:
  GPR[registerT] = result2;
\end{lstlisting}}
\newpage

\paragraph{Encoding} Format \textsf{LSU\_ALDUSSB.O} (Section~\ref{sec:format-LSU-ALDUSSB.O}), \verb|lsucode5 = 010|\\

\bitfields{ 1/P, 2/00, 2/-- --, 27/offset27, 1/1, 2/01, 3/111, 2/coherency, 6/registerT, 2/11, 3/010, 1/--, 5/01110, 1/0, 6/registerZ }


\paragraph{Syntax} \texttt{aldusw}\emph{coherency} \emph{offset27}[\emph{registerZ}] = \emph{registerT}

\paragraph{Description} The effective address is computed by adding the \emph{offset27} to the \emph{registerZ}.
This instruction implements atomic load-DUS-store on signed word (Decrement Unsigned Saturating). The word at effective address is DUSed with the \emph{registerT}, and its previous value is returned into the \emph{registerT}. The effective address must be word aligned (multiple of 4).


\paragraph{Modifier(s)}
\noindent\begin{itemize}
\item coherency (cf. coherency in section \ref{par:modifier-coherency} on page \pageref{par:modifier-coherency})
\end{itemize}

\paragraph{Execution} {Reservation table \textsf{LSU2\_MEMW\_AUXR\_AUXW.X} (Section~\ref{sec:reservation-LSU2-MEMW-AUXR-AUXW.X})}
~
{\begin{lstlisting}
stage ID:
  new offset = _SX_27(offset27);
  new coherency = _ZX_2(coherency);
stage RR:
  new base = GPR[registerZ];
  new address = base + offset;
stage E1:
  new argument2 = GPR[registerT];
  new result2 = _ZX_32(MEM.atomic_dus(address, 0xF, coherency << 32, argument2.32[0], registerT));
stage SR:
  GPR[registerT] = result2;
\end{lstlisting}}
\newpage

\paragraph{Encoding} Format \textsf{LSU\_ALDUSSB.D} (Section~\ref{sec:format-LSU-ALDUSSB.D}), \verb|lsucode5 = 010|\\

\bitfields{ 1/P, 2/00, 2/-- --, 27/extend27, 1/1, 2/00, 2/-- --, 27/offset27, 1/1, 2/01, 3/111, 2/coherency, 6/registerT, 2/11, 3/010, 1/--, 5/01110, 1/0, 6/registerZ }


\paragraph{Syntax} \texttt{aldusw}\emph{coherency} \emph{extend27\_\discretionary{}{}{}offset27}[\emph{registerZ}] = \emph{registerT}

\paragraph{Description} The effective address is computed by adding the \emph{extend27\_\discretionary{}{}{}offset27} to the \emph{registerZ}.
This instruction implements atomic load-DUS-store on signed word (Decrement Unsigned Saturating). The word at effective address is DUSed with the \emph{registerT}, and its previous value is returned into the \emph{registerT}. The effective address must be word aligned (multiple of 4).


\paragraph{Modifier(s)}
\noindent\begin{itemize}
\item coherency (cf. coherency in section \ref{par:modifier-coherency} on page \pageref{par:modifier-coherency})
\end{itemize}

\paragraph{Execution} {Reservation table \textsf{LSU2\_MEMW\_AUXR\_AUXW.Y} (Section~\ref{sec:reservation-LSU2-MEMW-AUXR-AUXW.Y})}
~
{\begin{lstlisting}
stage ID:
  new offset = _SX_54(extend27_offset27);
  new coherency = _ZX_2(coherency);
stage RR:
  new base = GPR[registerZ];
  new address = base + offset;
stage E1:
  new argument2 = GPR[registerT];
  new result2 = _ZX_32(MEM.atomic_dus(address, 0xF, coherency << 32, argument2.32[0], registerT));
stage SR:
  GPR[registerT] = result2;
\end{lstlisting}}
\newpage
\subsubsection[{\small ALEORB: Atomic Load and EOR Byte}]{ALEORB\hfill Atomic Load and EOR Byte}
\label{instr:ALEORB}\index{aleorb}
 
\paragraph{Encoding} Format \textsf{LSU\_ALEORSB} (Section~\ref{sec:format-LSU-ALEORSB}), \verb|lsucode5 = 000|\\

\bitfields{ 1/P, 2/01, 3/111, 2/coherency, 6/registerT, 2/11, 3/000, 1/--, 5/00111, 1/0, 6/registerZ }


\paragraph{Syntax} \texttt{aleorb}\emph{coherency} [\emph{registerZ}] = \emph{registerT}

\paragraph{Description} The effective address is the value of the \emph{registerZ}.
This instruction implements atomic load-EOR-store on byte. The byte at effective address is EORed with the \emph{registerT}, and its previous value is returned into the \emph{registerT}.


\paragraph{Modifier(s)}
\noindent\begin{itemize}
\item coherency (cf. coherency in section \ref{par:modifier-coherency} on page \pageref{par:modifier-coherency})
\end{itemize}

\paragraph{Execution} {Reservation table \textsf{LSU2\_MEMW\_AUXR\_AUXW} (Section~\ref{sec:reservation-LSU2-MEMW-AUXR-AUXW})}
~
{\begin{lstlisting}
stage ID:
  new coherency = _ZX_2(coherency);
stage RR:
  new address = GPR[registerZ];
stage E1:
  new argument2 = GPR[registerT];
  new result2 = _ZX_8(MEM.atomic_eor(address, 0x1, coherency << 32, argument2.8[0], registerT));
stage SR:
  GPR[registerT] = result2;
\end{lstlisting}}
\newpage

\paragraph{Encoding} Format \textsf{LSU\_ALEORSB.O} (Section~\ref{sec:format-LSU-ALEORSB.O}), \verb|lsucode5 = 000|\\

\bitfields{ 1/P, 2/00, 2/-- --, 27/offset27, 1/1, 2/01, 3/111, 2/coherency, 6/registerT, 2/11, 3/000, 1/--, 5/00111, 1/0, 6/registerZ }


\paragraph{Syntax} \texttt{aleorb}\emph{coherency} \emph{offset27}[\emph{registerZ}] = \emph{registerT}

\paragraph{Description} The effective address is computed by adding the \emph{offset27} to the \emph{registerZ}.
This instruction implements atomic load-EOR-store on byte. The byte at effective address is EORed with the \emph{registerT}, and its previous value is returned into the \emph{registerT}.


\paragraph{Modifier(s)}
\noindent\begin{itemize}
\item coherency (cf. coherency in section \ref{par:modifier-coherency} on page \pageref{par:modifier-coherency})
\end{itemize}

\paragraph{Execution} {Reservation table \textsf{LSU2\_MEMW\_AUXR\_AUXW.X} (Section~\ref{sec:reservation-LSU2-MEMW-AUXR-AUXW.X})}
~
{\begin{lstlisting}
stage ID:
  new offset = _SX_27(offset27);
  new coherency = _ZX_2(coherency);
stage RR:
  new base = GPR[registerZ];
  new address = base + offset;
stage E1:
  new argument2 = GPR[registerT];
  new result2 = _ZX_8(MEM.atomic_eor(address, 0x1, coherency << 32, argument2.8[0], registerT));
stage SR:
  GPR[registerT] = result2;
\end{lstlisting}}
\newpage

\paragraph{Encoding} Format \textsf{LSU\_ALEORSB.D} (Section~\ref{sec:format-LSU-ALEORSB.D}), \verb|lsucode5 = 000|\\

\bitfields{ 1/P, 2/00, 2/-- --, 27/extend27, 1/1, 2/00, 2/-- --, 27/offset27, 1/1, 2/01, 3/111, 2/coherency, 6/registerT, 2/11, 3/000, 1/--, 5/00111, 1/0, 6/registerZ }


\paragraph{Syntax} \texttt{aleorb}\emph{coherency} \emph{extend27\_\discretionary{}{}{}offset27}[\emph{registerZ}] = \emph{registerT}

\paragraph{Description} The effective address is computed by adding the \emph{extend27\_\discretionary{}{}{}offset27} to the \emph{registerZ}.
This instruction implements atomic load-EOR-store on byte. The byte at effective address is EORed with the \emph{registerT}, and its previous value is returned into the \emph{registerT}.


\paragraph{Modifier(s)}
\noindent\begin{itemize}
\item coherency (cf. coherency in section \ref{par:modifier-coherency} on page \pageref{par:modifier-coherency})
\end{itemize}

\paragraph{Execution} {Reservation table \textsf{LSU2\_MEMW\_AUXR\_AUXW.Y} (Section~\ref{sec:reservation-LSU2-MEMW-AUXR-AUXW.Y})}
~
{\begin{lstlisting}
stage ID:
  new offset = _SX_54(extend27_offset27);
  new coherency = _ZX_2(coherency);
stage RR:
  new base = GPR[registerZ];
  new address = base + offset;
stage E1:
  new argument2 = GPR[registerT];
  new result2 = _ZX_8(MEM.atomic_eor(address, 0x1, coherency << 32, argument2.8[0], registerT));
stage SR:
  GPR[registerT] = result2;
\end{lstlisting}}
\newpage
\subsubsection[{\small ALEORD: Atomic Load and EOR Double Word}]{ALEORD\hfill Atomic Load and EOR Double Word}
\label{instr:ALEORD}\index{aleord}
 
\paragraph{Encoding} Format \textsf{LSU\_ALEORSB} (Section~\ref{sec:format-LSU-ALEORSB}), \verb|lsucode5 = 011|\\

\bitfields{ 1/P, 2/01, 3/111, 2/coherency, 6/registerT, 2/11, 3/011, 1/--, 5/00111, 1/0, 6/registerZ }


\paragraph{Syntax} \texttt{aleord}\emph{coherency} [\emph{registerZ}] = \emph{registerT}

\paragraph{Description} The effective address is the value of the \emph{registerZ}.
This instruction implements atomic load-EOR-store on double word. The double word at effective address is EORed with the \emph{registerT}, and its previous value is returned into the \emph{registerT}. The effective address must be double word aligned (multiple of 8).


\paragraph{Modifier(s)}
\noindent\begin{itemize}
\item coherency (cf. coherency in section \ref{par:modifier-coherency} on page \pageref{par:modifier-coherency})
\end{itemize}

\paragraph{Execution} {Reservation table \textsf{LSU2\_MEMW\_AUXR\_AUXW} (Section~\ref{sec:reservation-LSU2-MEMW-AUXR-AUXW})}
~
{\begin{lstlisting}
stage ID:
  new coherency = _ZX_2(coherency);
stage RR:
  new address = GPR[registerZ];
stage E1:
  new argument2 = GPR[registerT];
  new result2 = MEM.atomic_eor(address, 0xFF, coherency << 32, argument2.64[0], registerT);
stage SR:
  GPR[registerT] = result2;
\end{lstlisting}}
\newpage

\paragraph{Encoding} Format \textsf{LSU\_ALEORSB.O} (Section~\ref{sec:format-LSU-ALEORSB.O}), \verb|lsucode5 = 011|\\

\bitfields{ 1/P, 2/00, 2/-- --, 27/offset27, 1/1, 2/01, 3/111, 2/coherency, 6/registerT, 2/11, 3/011, 1/--, 5/00111, 1/0, 6/registerZ }


\paragraph{Syntax} \texttt{aleord}\emph{coherency} \emph{offset27}[\emph{registerZ}] = \emph{registerT}

\paragraph{Description} The effective address is computed by adding the \emph{offset27} to the \emph{registerZ}.
This instruction implements atomic load-EOR-store on double word. The double word at effective address is EORed with the \emph{registerT}, and its previous value is returned into the \emph{registerT}. The effective address must be double word aligned (multiple of 8).


\paragraph{Modifier(s)}
\noindent\begin{itemize}
\item coherency (cf. coherency in section \ref{par:modifier-coherency} on page \pageref{par:modifier-coherency})
\end{itemize}

\paragraph{Execution} {Reservation table \textsf{LSU2\_MEMW\_AUXR\_AUXW.X} (Section~\ref{sec:reservation-LSU2-MEMW-AUXR-AUXW.X})}
~
{\begin{lstlisting}
stage ID:
  new offset = _SX_27(offset27);
  new coherency = _ZX_2(coherency);
stage RR:
  new base = GPR[registerZ];
  new address = base + offset;
stage E1:
  new argument2 = GPR[registerT];
  new result2 = MEM.atomic_eor(address, 0xFF, coherency << 32, argument2.64[0], registerT);
stage SR:
  GPR[registerT] = result2;
\end{lstlisting}}
\newpage

\paragraph{Encoding} Format \textsf{LSU\_ALEORSB.D} (Section~\ref{sec:format-LSU-ALEORSB.D}), \verb|lsucode5 = 011|\\

\bitfields{ 1/P, 2/00, 2/-- --, 27/extend27, 1/1, 2/00, 2/-- --, 27/offset27, 1/1, 2/01, 3/111, 2/coherency, 6/registerT, 2/11, 3/011, 1/--, 5/00111, 1/0, 6/registerZ }


\paragraph{Syntax} \texttt{aleord}\emph{coherency} \emph{extend27\_\discretionary{}{}{}offset27}[\emph{registerZ}] = \emph{registerT}

\paragraph{Description} The effective address is computed by adding the \emph{extend27\_\discretionary{}{}{}offset27} to the \emph{registerZ}.
This instruction implements atomic load-EOR-store on double word. The double word at effective address is EORed with the \emph{registerT}, and its previous value is returned into the \emph{registerT}. The effective address must be double word aligned (multiple of 8).


\paragraph{Modifier(s)}
\noindent\begin{itemize}
\item coherency (cf. coherency in section \ref{par:modifier-coherency} on page \pageref{par:modifier-coherency})
\end{itemize}

\paragraph{Execution} {Reservation table \textsf{LSU2\_MEMW\_AUXR\_AUXW.Y} (Section~\ref{sec:reservation-LSU2-MEMW-AUXR-AUXW.Y})}
~
{\begin{lstlisting}
stage ID:
  new offset = _SX_54(extend27_offset27);
  new coherency = _ZX_2(coherency);
stage RR:
  new base = GPR[registerZ];
  new address = base + offset;
stage E1:
  new argument2 = GPR[registerT];
  new result2 = MEM.atomic_eor(address, 0xFF, coherency << 32, argument2.64[0], registerT);
stage SR:
  GPR[registerT] = result2;
\end{lstlisting}}
\newpage
\subsubsection[{\small ALEORH: Atomic Load and EOR Half Word}]{ALEORH\hfill Atomic Load and EOR Half Word}
\label{instr:ALEORH}\index{aleorh}
 
\paragraph{Encoding} Format \textsf{LSU\_ALEORSB} (Section~\ref{sec:format-LSU-ALEORSB}), \verb|lsucode5 = 001|\\

\bitfields{ 1/P, 2/01, 3/111, 2/coherency, 6/registerT, 2/11, 3/001, 1/--, 5/00111, 1/0, 6/registerZ }


\paragraph{Syntax} \texttt{aleorh}\emph{coherency} [\emph{registerZ}] = \emph{registerT}

\paragraph{Description} The effective address is the value of the \emph{registerZ}.
This instruction implements atomic load-EOR-store on half word. The half word at effective address is EORed with the \emph{registerT}, and its previous value is returned into the \emph{registerT}. The effective address must be half word aligned (multiple of 2).


\paragraph{Modifier(s)}
\noindent\begin{itemize}
\item coherency (cf. coherency in section \ref{par:modifier-coherency} on page \pageref{par:modifier-coherency})
\end{itemize}

\paragraph{Execution} {Reservation table \textsf{LSU2\_MEMW\_AUXR\_AUXW} (Section~\ref{sec:reservation-LSU2-MEMW-AUXR-AUXW})}
~
{\begin{lstlisting}
stage ID:
  new coherency = _ZX_2(coherency);
stage RR:
  new address = GPR[registerZ];
stage E1:
  new argument2 = GPR[registerT];
  new result2 = _ZX_16(MEM.atomic_eor(address, 0x3, coherency << 32, argument2.16[0], registerT));
stage SR:
  GPR[registerT] = result2;
\end{lstlisting}}
\newpage

\paragraph{Encoding} Format \textsf{LSU\_ALEORSB.O} (Section~\ref{sec:format-LSU-ALEORSB.O}), \verb|lsucode5 = 001|\\

\bitfields{ 1/P, 2/00, 2/-- --, 27/offset27, 1/1, 2/01, 3/111, 2/coherency, 6/registerT, 2/11, 3/001, 1/--, 5/00111, 1/0, 6/registerZ }


\paragraph{Syntax} \texttt{aleorh}\emph{coherency} \emph{offset27}[\emph{registerZ}] = \emph{registerT}

\paragraph{Description} The effective address is computed by adding the \emph{offset27} to the \emph{registerZ}.
This instruction implements atomic load-EOR-store on half word. The half word at effective address is EORed with the \emph{registerT}, and its previous value is returned into the \emph{registerT}. The effective address must be half word aligned (multiple of 2).


\paragraph{Modifier(s)}
\noindent\begin{itemize}
\item coherency (cf. coherency in section \ref{par:modifier-coherency} on page \pageref{par:modifier-coherency})
\end{itemize}

\paragraph{Execution} {Reservation table \textsf{LSU2\_MEMW\_AUXR\_AUXW.X} (Section~\ref{sec:reservation-LSU2-MEMW-AUXR-AUXW.X})}
~
{\begin{lstlisting}
stage ID:
  new offset = _SX_27(offset27);
  new coherency = _ZX_2(coherency);
stage RR:
  new base = GPR[registerZ];
  new address = base + offset;
stage E1:
  new argument2 = GPR[registerT];
  new result2 = _ZX_16(MEM.atomic_eor(address, 0x3, coherency << 32, argument2.16[0], registerT));
stage SR:
  GPR[registerT] = result2;
\end{lstlisting}}
\newpage

\paragraph{Encoding} Format \textsf{LSU\_ALEORSB.D} (Section~\ref{sec:format-LSU-ALEORSB.D}), \verb|lsucode5 = 001|\\

\bitfields{ 1/P, 2/00, 2/-- --, 27/extend27, 1/1, 2/00, 2/-- --, 27/offset27, 1/1, 2/01, 3/111, 2/coherency, 6/registerT, 2/11, 3/001, 1/--, 5/00111, 1/0, 6/registerZ }


\paragraph{Syntax} \texttt{aleorh}\emph{coherency} \emph{extend27\_\discretionary{}{}{}offset27}[\emph{registerZ}] = \emph{registerT}

\paragraph{Description} The effective address is computed by adding the \emph{extend27\_\discretionary{}{}{}offset27} to the \emph{registerZ}.
This instruction implements atomic load-EOR-store on half word. The half word at effective address is EORed with the \emph{registerT}, and its previous value is returned into the \emph{registerT}. The effective address must be half word aligned (multiple of 2).


\paragraph{Modifier(s)}
\noindent\begin{itemize}
\item coherency (cf. coherency in section \ref{par:modifier-coherency} on page \pageref{par:modifier-coherency})
\end{itemize}

\paragraph{Execution} {Reservation table \textsf{LSU2\_MEMW\_AUXR\_AUXW.Y} (Section~\ref{sec:reservation-LSU2-MEMW-AUXR-AUXW.Y})}
~
{\begin{lstlisting}
stage ID:
  new offset = _SX_54(extend27_offset27);
  new coherency = _ZX_2(coherency);
stage RR:
  new base = GPR[registerZ];
  new address = base + offset;
stage E1:
  new argument2 = GPR[registerT];
  new result2 = _ZX_16(MEM.atomic_eor(address, 0x3, coherency << 32, argument2.16[0], registerT));
stage SR:
  GPR[registerT] = result2;
\end{lstlisting}}
\newpage
\subsubsection[{\small ALEORW: Atomic Load and EOR Word}]{ALEORW\hfill Atomic Load and EOR Word}
\label{instr:ALEORW}\index{aleorw}
 
\paragraph{Encoding} Format \textsf{LSU\_ALEORSB} (Section~\ref{sec:format-LSU-ALEORSB}), \verb|lsucode5 = 010|\\

\bitfields{ 1/P, 2/01, 3/111, 2/coherency, 6/registerT, 2/11, 3/010, 1/--, 5/00111, 1/0, 6/registerZ }


\paragraph{Syntax} \texttt{aleorw}\emph{coherency} [\emph{registerZ}] = \emph{registerT}

\paragraph{Description} The effective address is the value of the \emph{registerZ}.
This instruction implements atomic load-EOR-store on word. The word at effective address is EORed with the \emph{registerT}, and its previous value is returned into the \emph{registerT}. The effective address must be word aligned (multiple of 4).


\paragraph{Modifier(s)}
\noindent\begin{itemize}
\item coherency (cf. coherency in section \ref{par:modifier-coherency} on page \pageref{par:modifier-coherency})
\end{itemize}

\paragraph{Execution} {Reservation table \textsf{LSU2\_MEMW\_AUXR\_AUXW} (Section~\ref{sec:reservation-LSU2-MEMW-AUXR-AUXW})}
~
{\begin{lstlisting}
stage ID:
  new coherency = _ZX_2(coherency);
stage RR:
  new address = GPR[registerZ];
stage E1:
  new argument2 = GPR[registerT];
  new result2 = _ZX_32(MEM.atomic_eor(address, 0xF, coherency << 32, argument2.32[0], registerT));
stage SR:
  GPR[registerT] = result2;
\end{lstlisting}}
\newpage

\paragraph{Encoding} Format \textsf{LSU\_ALEORSB.O} (Section~\ref{sec:format-LSU-ALEORSB.O}), \verb|lsucode5 = 010|\\

\bitfields{ 1/P, 2/00, 2/-- --, 27/offset27, 1/1, 2/01, 3/111, 2/coherency, 6/registerT, 2/11, 3/010, 1/--, 5/00111, 1/0, 6/registerZ }


\paragraph{Syntax} \texttt{aleorw}\emph{coherency} \emph{offset27}[\emph{registerZ}] = \emph{registerT}

\paragraph{Description} The effective address is computed by adding the \emph{offset27} to the \emph{registerZ}.
This instruction implements atomic load-EOR-store on word. The word at effective address is EORed with the \emph{registerT}, and its previous value is returned into the \emph{registerT}. The effective address must be word aligned (multiple of 4).


\paragraph{Modifier(s)}
\noindent\begin{itemize}
\item coherency (cf. coherency in section \ref{par:modifier-coherency} on page \pageref{par:modifier-coherency})
\end{itemize}

\paragraph{Execution} {Reservation table \textsf{LSU2\_MEMW\_AUXR\_AUXW.X} (Section~\ref{sec:reservation-LSU2-MEMW-AUXR-AUXW.X})}
~
{\begin{lstlisting}
stage ID:
  new offset = _SX_27(offset27);
  new coherency = _ZX_2(coherency);
stage RR:
  new base = GPR[registerZ];
  new address = base + offset;
stage E1:
  new argument2 = GPR[registerT];
  new result2 = _ZX_32(MEM.atomic_eor(address, 0xF, coherency << 32, argument2.32[0], registerT));
stage SR:
  GPR[registerT] = result2;
\end{lstlisting}}
\newpage

\paragraph{Encoding} Format \textsf{LSU\_ALEORSB.D} (Section~\ref{sec:format-LSU-ALEORSB.D}), \verb|lsucode5 = 010|\\

\bitfields{ 1/P, 2/00, 2/-- --, 27/extend27, 1/1, 2/00, 2/-- --, 27/offset27, 1/1, 2/01, 3/111, 2/coherency, 6/registerT, 2/11, 3/010, 1/--, 5/00111, 1/0, 6/registerZ }


\paragraph{Syntax} \texttt{aleorw}\emph{coherency} \emph{extend27\_\discretionary{}{}{}offset27}[\emph{registerZ}] = \emph{registerT}

\paragraph{Description} The effective address is computed by adding the \emph{extend27\_\discretionary{}{}{}offset27} to the \emph{registerZ}.
This instruction implements atomic load-EOR-store on word. The word at effective address is EORed with the \emph{registerT}, and its previous value is returned into the \emph{registerT}. The effective address must be word aligned (multiple of 4).


\paragraph{Modifier(s)}
\noindent\begin{itemize}
\item coherency (cf. coherency in section \ref{par:modifier-coherency} on page \pageref{par:modifier-coherency})
\end{itemize}

\paragraph{Execution} {Reservation table \textsf{LSU2\_MEMW\_AUXR\_AUXW.Y} (Section~\ref{sec:reservation-LSU2-MEMW-AUXR-AUXW.Y})}
~
{\begin{lstlisting}
stage ID:
  new offset = _SX_54(extend27_offset27);
  new coherency = _ZX_2(coherency);
stage RR:
  new base = GPR[registerZ];
  new address = base + offset;
stage E1:
  new argument2 = GPR[registerT];
  new result2 = _ZX_32(MEM.atomic_eor(address, 0xF, coherency << 32, argument2.32[0], registerT));
stage SR:
  GPR[registerT] = result2;
\end{lstlisting}}
\newpage
\subsubsection[{\small ALH: Atomic Load Half Word}]{ALH\hfill Atomic Load Half Word}
\label{instr:ALH}\index{alh}
 
\paragraph{Encoding} Format \textsf{LSU\_ALSB} (Section~\ref{sec:format-LSU-ALSB}), \verb|lsucode5 = 001|\\

\bitfields{ 1/P, 2/01, 3/111, 2/coherency, 6/registerW, 2/11, 3/001, 1/--, 5/00000, 1/0, 6/registerZ }


\paragraph{Syntax} \texttt{alh}\emph{coherency} \emph{registerW} = [\emph{registerZ}]

\paragraph{Description} The effective address is the value of the \emph{registerZ}.
The \emph{registerW} is atomically loaded with zero extension from the memory half word starting at the effective address. This instruction is provided to directly operate on the last level of the memory hierarchy, which is typically a processor DDR or a cluster TCM, irrespective of the MMU DCP (Data Cache Policy) settings. The effective address must be half word aligned (multiple of 2).


\paragraph{Modifier(s)}
\noindent\begin{itemize}
\item coherency (cf. coherency in section \ref{par:modifier-coherency} on page \pageref{par:modifier-coherency})
\end{itemize}

\paragraph{Execution} {Reservation table \textsf{LSU\_MEMW\_AUXW} (Section~\ref{sec:reservation-LSU-MEMW-AUXW})}
~
{\begin{lstlisting}
stage ID:
  new coherency = _ZX_2(coherency);
stage RR:
  new address = GPR[registerZ];
  new result1 = _ZX_16(MEM.atomic_load(address, 0x3, coherency << 32, registerW));
stage CR:
  GPR[registerW] = result1;
\end{lstlisting}}
\newpage

\paragraph{Encoding} Format \textsf{LSU\_ALSB.O} (Section~\ref{sec:format-LSU-ALSB.O}), \verb|lsucode5 = 001|\\

\bitfields{ 1/P, 2/00, 2/-- --, 27/offset27, 1/1, 2/01, 3/111, 2/coherency, 6/registerW, 2/11, 3/001, 1/--, 5/00000, 1/0, 6/registerZ }


\paragraph{Syntax} \texttt{alh}\emph{coherency} \emph{registerW} = \emph{offset27}[\emph{registerZ}]

\paragraph{Description} The effective address is computed by adding the \emph{offset27} to the \emph{registerZ}.
The \emph{registerW} is atomically loaded with zero extension from the memory half word starting at the effective address. This instruction is provided to directly operate on the last level of the memory hierarchy, which is typically a processor DDR or a cluster TCM, irrespective of the MMU DCP (Data Cache Policy) settings. The effective address must be half word aligned (multiple of 2).


\paragraph{Modifier(s)}
\noindent\begin{itemize}
\item coherency (cf. coherency in section \ref{par:modifier-coherency} on page \pageref{par:modifier-coherency})
\end{itemize}

\paragraph{Execution} {Reservation table \textsf{LSU\_MEMW\_AUXW.X} (Section~\ref{sec:reservation-LSU-MEMW-AUXW.X})}
~
{\begin{lstlisting}
stage ID:
  new offset = _SX_27(offset27);
  new coherency = _ZX_2(coherency);
stage RR:
  new base = GPR[registerZ];
  new address = base + offset;
  new result1 = _ZX_16(MEM.atomic_load(address, 0x3, coherency << 32, registerW));
stage CR:
  GPR[registerW] = result1;
\end{lstlisting}}
\newpage

\paragraph{Encoding} Format \textsf{LSU\_ALSB.D} (Section~\ref{sec:format-LSU-ALSB.D}), \verb|lsucode5 = 001|\\

\bitfields{ 1/P, 2/00, 2/-- --, 27/extend27, 1/1, 2/00, 2/-- --, 27/offset27, 1/1, 2/01, 3/111, 2/coherency, 6/registerW, 2/11, 3/001, 1/--, 5/00000, 1/0, 6/registerZ }


\paragraph{Syntax} \texttt{alh}\emph{coherency} \emph{registerW} = \emph{extend27\_\discretionary{}{}{}offset27}[\emph{registerZ}]

\paragraph{Description} The effective address is computed by adding the \emph{extend27\_\discretionary{}{}{}offset27} to the \emph{registerZ}.
The \emph{registerW} is atomically loaded with zero extension from the memory half word starting at the effective address. This instruction is provided to directly operate on the last level of the memory hierarchy, which is typically a processor DDR or a cluster TCM, irrespective of the MMU DCP (Data Cache Policy) settings. The effective address must be half word aligned (multiple of 2).


\paragraph{Modifier(s)}
\noindent\begin{itemize}
\item coherency (cf. coherency in section \ref{par:modifier-coherency} on page \pageref{par:modifier-coherency})
\end{itemize}

\paragraph{Execution} {Reservation table \textsf{LSU\_MEMW\_AUXW.Y} (Section~\ref{sec:reservation-LSU-MEMW-AUXW.Y})}
~
{\begin{lstlisting}
stage ID:
  new offset = _SX_54(extend27_offset27);
  new coherency = _ZX_2(coherency);
stage RR:
  new base = GPR[registerZ];
  new address = base + offset;
  new result1 = _ZX_16(MEM.atomic_load(address, 0x3, coherency << 32, registerW));
stage CR:
  GPR[registerW] = result1;
\end{lstlisting}}
\newpage
\subsubsection[{\small ALIORB: Atomic Load and IOR Byte}]{ALIORB\hfill Atomic Load and IOR Byte}
\label{instr:ALIORB}\index{aliorb}
 
\paragraph{Encoding} Format \textsf{LSU\_ALIORSB} (Section~\ref{sec:format-LSU-ALIORSB}), \verb|lsucode5 = 000|\\

\bitfields{ 1/P, 2/01, 3/111, 2/coherency, 6/registerT, 2/11, 3/000, 1/--, 5/00110, 1/0, 6/registerZ }


\paragraph{Syntax} \texttt{aliorb}\emph{coherency} [\emph{registerZ}] = \emph{registerT}

\paragraph{Description} The effective address is the value of the \emph{registerZ}.
This instruction implements atomic load-IOR-store on byte. The byte at effective address is IORed with the \emph{registerT}, and its previous value is returned into the \emph{registerT}.


\paragraph{Modifier(s)}
\noindent\begin{itemize}
\item coherency (cf. coherency in section \ref{par:modifier-coherency} on page \pageref{par:modifier-coherency})
\end{itemize}

\paragraph{Execution} {Reservation table \textsf{LSU2\_MEMW\_AUXR\_AUXW} (Section~\ref{sec:reservation-LSU2-MEMW-AUXR-AUXW})}
~
{\begin{lstlisting}
stage ID:
  new coherency = _ZX_2(coherency);
stage RR:
  new address = GPR[registerZ];
stage E1:
  new argument2 = GPR[registerT];
  new result2 = _ZX_8(MEM.atomic_ior(address, 0x1, coherency << 32, argument2.8[0], registerT));
stage SR:
  GPR[registerT] = result2;
\end{lstlisting}}
\newpage

\paragraph{Encoding} Format \textsf{LSU\_ALIORSB.O} (Section~\ref{sec:format-LSU-ALIORSB.O}), \verb|lsucode5 = 000|\\

\bitfields{ 1/P, 2/00, 2/-- --, 27/offset27, 1/1, 2/01, 3/111, 2/coherency, 6/registerT, 2/11, 3/000, 1/--, 5/00110, 1/0, 6/registerZ }


\paragraph{Syntax} \texttt{aliorb}\emph{coherency} \emph{offset27}[\emph{registerZ}] = \emph{registerT}

\paragraph{Description} The effective address is computed by adding the \emph{offset27} to the \emph{registerZ}.
This instruction implements atomic load-IOR-store on byte. The byte at effective address is IORed with the \emph{registerT}, and its previous value is returned into the \emph{registerT}.


\paragraph{Modifier(s)}
\noindent\begin{itemize}
\item coherency (cf. coherency in section \ref{par:modifier-coherency} on page \pageref{par:modifier-coherency})
\end{itemize}

\paragraph{Execution} {Reservation table \textsf{LSU2\_MEMW\_AUXR\_AUXW.X} (Section~\ref{sec:reservation-LSU2-MEMW-AUXR-AUXW.X})}
~
{\begin{lstlisting}
stage ID:
  new offset = _SX_27(offset27);
  new coherency = _ZX_2(coherency);
stage RR:
  new base = GPR[registerZ];
  new address = base + offset;
stage E1:
  new argument2 = GPR[registerT];
  new result2 = _ZX_8(MEM.atomic_ior(address, 0x1, coherency << 32, argument2.8[0], registerT));
stage SR:
  GPR[registerT] = result2;
\end{lstlisting}}
\newpage

\paragraph{Encoding} Format \textsf{LSU\_ALIORSB.D} (Section~\ref{sec:format-LSU-ALIORSB.D}), \verb|lsucode5 = 000|\\

\bitfields{ 1/P, 2/00, 2/-- --, 27/extend27, 1/1, 2/00, 2/-- --, 27/offset27, 1/1, 2/01, 3/111, 2/coherency, 6/registerT, 2/11, 3/000, 1/--, 5/00110, 1/0, 6/registerZ }


\paragraph{Syntax} \texttt{aliorb}\emph{coherency} \emph{extend27\_\discretionary{}{}{}offset27}[\emph{registerZ}] = \emph{registerT}

\paragraph{Description} The effective address is computed by adding the \emph{extend27\_\discretionary{}{}{}offset27} to the \emph{registerZ}.
This instruction implements atomic load-IOR-store on byte. The byte at effective address is IORed with the \emph{registerT}, and its previous value is returned into the \emph{registerT}.


\paragraph{Modifier(s)}
\noindent\begin{itemize}
\item coherency (cf. coherency in section \ref{par:modifier-coherency} on page \pageref{par:modifier-coherency})
\end{itemize}

\paragraph{Execution} {Reservation table \textsf{LSU2\_MEMW\_AUXR\_AUXW.Y} (Section~\ref{sec:reservation-LSU2-MEMW-AUXR-AUXW.Y})}
~
{\begin{lstlisting}
stage ID:
  new offset = _SX_54(extend27_offset27);
  new coherency = _ZX_2(coherency);
stage RR:
  new base = GPR[registerZ];
  new address = base + offset;
stage E1:
  new argument2 = GPR[registerT];
  new result2 = _ZX_8(MEM.atomic_ior(address, 0x1, coherency << 32, argument2.8[0], registerT));
stage SR:
  GPR[registerT] = result2;
\end{lstlisting}}
\newpage
\subsubsection[{\small ALIORD: Atomic Load and IOR Double Word}]{ALIORD\hfill Atomic Load and IOR Double Word}
\label{instr:ALIORD}\index{aliord}
 
\paragraph{Encoding} Format \textsf{LSU\_ALIORSB} (Section~\ref{sec:format-LSU-ALIORSB}), \verb|lsucode5 = 011|\\

\bitfields{ 1/P, 2/01, 3/111, 2/coherency, 6/registerT, 2/11, 3/011, 1/--, 5/00110, 1/0, 6/registerZ }


\paragraph{Syntax} \texttt{aliord}\emph{coherency} [\emph{registerZ}] = \emph{registerT}

\paragraph{Description} The effective address is the value of the \emph{registerZ}.
This instruction implements atomic load-IOR-store on double word. The double word at effective address is IORed with the \emph{registerT}, and its previous value is returned into the \emph{registerT}. The effective address must be double word aligned (multiple of 8).


\paragraph{Modifier(s)}
\noindent\begin{itemize}
\item coherency (cf. coherency in section \ref{par:modifier-coherency} on page \pageref{par:modifier-coherency})
\end{itemize}

\paragraph{Execution} {Reservation table \textsf{LSU2\_MEMW\_AUXR\_AUXW} (Section~\ref{sec:reservation-LSU2-MEMW-AUXR-AUXW})}
~
{\begin{lstlisting}
stage ID:
  new coherency = _ZX_2(coherency);
stage RR:
  new address = GPR[registerZ];
stage E1:
  new argument2 = GPR[registerT];
  new result2 = MEM.atomic_ior(address, 0xFF, coherency << 32, argument2.64[0], registerT);
stage SR:
  GPR[registerT] = result2;
\end{lstlisting}}
\newpage

\paragraph{Encoding} Format \textsf{LSU\_ALIORSB.O} (Section~\ref{sec:format-LSU-ALIORSB.O}), \verb|lsucode5 = 011|\\

\bitfields{ 1/P, 2/00, 2/-- --, 27/offset27, 1/1, 2/01, 3/111, 2/coherency, 6/registerT, 2/11, 3/011, 1/--, 5/00110, 1/0, 6/registerZ }


\paragraph{Syntax} \texttt{aliord}\emph{coherency} \emph{offset27}[\emph{registerZ}] = \emph{registerT}

\paragraph{Description} The effective address is computed by adding the \emph{offset27} to the \emph{registerZ}.
This instruction implements atomic load-IOR-store on double word. The double word at effective address is IORed with the \emph{registerT}, and its previous value is returned into the \emph{registerT}. The effective address must be double word aligned (multiple of 8).


\paragraph{Modifier(s)}
\noindent\begin{itemize}
\item coherency (cf. coherency in section \ref{par:modifier-coherency} on page \pageref{par:modifier-coherency})
\end{itemize}

\paragraph{Execution} {Reservation table \textsf{LSU2\_MEMW\_AUXR\_AUXW.X} (Section~\ref{sec:reservation-LSU2-MEMW-AUXR-AUXW.X})}
~
{\begin{lstlisting}
stage ID:
  new offset = _SX_27(offset27);
  new coherency = _ZX_2(coherency);
stage RR:
  new base = GPR[registerZ];
  new address = base + offset;
stage E1:
  new argument2 = GPR[registerT];
  new result2 = MEM.atomic_ior(address, 0xFF, coherency << 32, argument2.64[0], registerT);
stage SR:
  GPR[registerT] = result2;
\end{lstlisting}}
\newpage

\paragraph{Encoding} Format \textsf{LSU\_ALIORSB.D} (Section~\ref{sec:format-LSU-ALIORSB.D}), \verb|lsucode5 = 011|\\

\bitfields{ 1/P, 2/00, 2/-- --, 27/extend27, 1/1, 2/00, 2/-- --, 27/offset27, 1/1, 2/01, 3/111, 2/coherency, 6/registerT, 2/11, 3/011, 1/--, 5/00110, 1/0, 6/registerZ }


\paragraph{Syntax} \texttt{aliord}\emph{coherency} \emph{extend27\_\discretionary{}{}{}offset27}[\emph{registerZ}] = \emph{registerT}

\paragraph{Description} The effective address is computed by adding the \emph{extend27\_\discretionary{}{}{}offset27} to the \emph{registerZ}.
This instruction implements atomic load-IOR-store on double word. The double word at effective address is IORed with the \emph{registerT}, and its previous value is returned into the \emph{registerT}. The effective address must be double word aligned (multiple of 8).


\paragraph{Modifier(s)}
\noindent\begin{itemize}
\item coherency (cf. coherency in section \ref{par:modifier-coherency} on page \pageref{par:modifier-coherency})
\end{itemize}

\paragraph{Execution} {Reservation table \textsf{LSU2\_MEMW\_AUXR\_AUXW.Y} (Section~\ref{sec:reservation-LSU2-MEMW-AUXR-AUXW.Y})}
~
{\begin{lstlisting}
stage ID:
  new offset = _SX_54(extend27_offset27);
  new coherency = _ZX_2(coherency);
stage RR:
  new base = GPR[registerZ];
  new address = base + offset;
stage E1:
  new argument2 = GPR[registerT];
  new result2 = MEM.atomic_ior(address, 0xFF, coherency << 32, argument2.64[0], registerT);
stage SR:
  GPR[registerT] = result2;
\end{lstlisting}}
\newpage
\subsubsection[{\small ALIORH: Atomic Load and IOR Half Word}]{ALIORH\hfill Atomic Load and IOR Half Word}
\label{instr:ALIORH}\index{aliorh}
 
\paragraph{Encoding} Format \textsf{LSU\_ALIORSB} (Section~\ref{sec:format-LSU-ALIORSB}), \verb|lsucode5 = 001|\\

\bitfields{ 1/P, 2/01, 3/111, 2/coherency, 6/registerT, 2/11, 3/001, 1/--, 5/00110, 1/0, 6/registerZ }


\paragraph{Syntax} \texttt{aliorh}\emph{coherency} [\emph{registerZ}] = \emph{registerT}

\paragraph{Description} The effective address is the value of the \emph{registerZ}.
This instruction implements atomic load-IOR-store on half word. The half word at effective address is IORed with the \emph{registerT}, and its previous value is returned into the \emph{registerT}. The effective address must be half word aligned (multiple of 2).


\paragraph{Modifier(s)}
\noindent\begin{itemize}
\item coherency (cf. coherency in section \ref{par:modifier-coherency} on page \pageref{par:modifier-coherency})
\end{itemize}

\paragraph{Execution} {Reservation table \textsf{LSU2\_MEMW\_AUXR\_AUXW} (Section~\ref{sec:reservation-LSU2-MEMW-AUXR-AUXW})}
~
{\begin{lstlisting}
stage ID:
  new coherency = _ZX_2(coherency);
stage RR:
  new address = GPR[registerZ];
stage E1:
  new argument2 = GPR[registerT];
  new result2 = _ZX_16(MEM.atomic_ior(address, 0x3, coherency << 32, argument2.16[0], registerT));
stage SR:
  GPR[registerT] = result2;
\end{lstlisting}}
\newpage

\paragraph{Encoding} Format \textsf{LSU\_ALIORSB.O} (Section~\ref{sec:format-LSU-ALIORSB.O}), \verb|lsucode5 = 001|\\

\bitfields{ 1/P, 2/00, 2/-- --, 27/offset27, 1/1, 2/01, 3/111, 2/coherency, 6/registerT, 2/11, 3/001, 1/--, 5/00110, 1/0, 6/registerZ }


\paragraph{Syntax} \texttt{aliorh}\emph{coherency} \emph{offset27}[\emph{registerZ}] = \emph{registerT}

\paragraph{Description} The effective address is computed by adding the \emph{offset27} to the \emph{registerZ}.
This instruction implements atomic load-IOR-store on half word. The half word at effective address is IORed with the \emph{registerT}, and its previous value is returned into the \emph{registerT}. The effective address must be half word aligned (multiple of 2).


\paragraph{Modifier(s)}
\noindent\begin{itemize}
\item coherency (cf. coherency in section \ref{par:modifier-coherency} on page \pageref{par:modifier-coherency})
\end{itemize}

\paragraph{Execution} {Reservation table \textsf{LSU2\_MEMW\_AUXR\_AUXW.X} (Section~\ref{sec:reservation-LSU2-MEMW-AUXR-AUXW.X})}
~
{\begin{lstlisting}
stage ID:
  new offset = _SX_27(offset27);
  new coherency = _ZX_2(coherency);
stage RR:
  new base = GPR[registerZ];
  new address = base + offset;
stage E1:
  new argument2 = GPR[registerT];
  new result2 = _ZX_16(MEM.atomic_ior(address, 0x3, coherency << 32, argument2.16[0], registerT));
stage SR:
  GPR[registerT] = result2;
\end{lstlisting}}
\newpage

\paragraph{Encoding} Format \textsf{LSU\_ALIORSB.D} (Section~\ref{sec:format-LSU-ALIORSB.D}), \verb|lsucode5 = 001|\\

\bitfields{ 1/P, 2/00, 2/-- --, 27/extend27, 1/1, 2/00, 2/-- --, 27/offset27, 1/1, 2/01, 3/111, 2/coherency, 6/registerT, 2/11, 3/001, 1/--, 5/00110, 1/0, 6/registerZ }


\paragraph{Syntax} \texttt{aliorh}\emph{coherency} \emph{extend27\_\discretionary{}{}{}offset27}[\emph{registerZ}] = \emph{registerT}

\paragraph{Description} The effective address is computed by adding the \emph{extend27\_\discretionary{}{}{}offset27} to the \emph{registerZ}.
This instruction implements atomic load-IOR-store on half word. The half word at effective address is IORed with the \emph{registerT}, and its previous value is returned into the \emph{registerT}. The effective address must be half word aligned (multiple of 2).


\paragraph{Modifier(s)}
\noindent\begin{itemize}
\item coherency (cf. coherency in section \ref{par:modifier-coherency} on page \pageref{par:modifier-coherency})
\end{itemize}

\paragraph{Execution} {Reservation table \textsf{LSU2\_MEMW\_AUXR\_AUXW.Y} (Section~\ref{sec:reservation-LSU2-MEMW-AUXR-AUXW.Y})}
~
{\begin{lstlisting}
stage ID:
  new offset = _SX_54(extend27_offset27);
  new coherency = _ZX_2(coherency);
stage RR:
  new base = GPR[registerZ];
  new address = base + offset;
stage E1:
  new argument2 = GPR[registerT];
  new result2 = _ZX_16(MEM.atomic_ior(address, 0x3, coherency << 32, argument2.16[0], registerT));
stage SR:
  GPR[registerT] = result2;
\end{lstlisting}}
\newpage
\subsubsection[{\small ALIORW: Atomic Load and IOR Word}]{ALIORW\hfill Atomic Load and IOR Word}
\label{instr:ALIORW}\index{aliorw}
 
\paragraph{Encoding} Format \textsf{LSU\_ALIORSB} (Section~\ref{sec:format-LSU-ALIORSB}), \verb|lsucode5 = 010|\\

\bitfields{ 1/P, 2/01, 3/111, 2/coherency, 6/registerT, 2/11, 3/010, 1/--, 5/00110, 1/0, 6/registerZ }


\paragraph{Syntax} \texttt{aliorw}\emph{coherency} [\emph{registerZ}] = \emph{registerT}

\paragraph{Description} The effective address is the value of the \emph{registerZ}.
This instruction implements atomic load-IOR-store on word. The word at effective address is IORed with the \emph{registerT}, and its previous value is returned into the \emph{registerT}. The effective address must be word aligned (multiple of 4).


\paragraph{Modifier(s)}
\noindent\begin{itemize}
\item coherency (cf. coherency in section \ref{par:modifier-coherency} on page \pageref{par:modifier-coherency})
\end{itemize}

\paragraph{Execution} {Reservation table \textsf{LSU2\_MEMW\_AUXR\_AUXW} (Section~\ref{sec:reservation-LSU2-MEMW-AUXR-AUXW})}
~
{\begin{lstlisting}
stage ID:
  new coherency = _ZX_2(coherency);
stage RR:
  new address = GPR[registerZ];
stage E1:
  new argument2 = GPR[registerT];
  new result2 = _ZX_32(MEM.atomic_ior(address, 0xF, coherency << 32, argument2.32[0], registerT));
stage SR:
  GPR[registerT] = result2;
\end{lstlisting}}
\newpage

\paragraph{Encoding} Format \textsf{LSU\_ALIORSB.O} (Section~\ref{sec:format-LSU-ALIORSB.O}), \verb|lsucode5 = 010|\\

\bitfields{ 1/P, 2/00, 2/-- --, 27/offset27, 1/1, 2/01, 3/111, 2/coherency, 6/registerT, 2/11, 3/010, 1/--, 5/00110, 1/0, 6/registerZ }


\paragraph{Syntax} \texttt{aliorw}\emph{coherency} \emph{offset27}[\emph{registerZ}] = \emph{registerT}

\paragraph{Description} The effective address is computed by adding the \emph{offset27} to the \emph{registerZ}.
This instruction implements atomic load-IOR-store on word. The word at effective address is IORed with the \emph{registerT}, and its previous value is returned into the \emph{registerT}. The effective address must be word aligned (multiple of 4).


\paragraph{Modifier(s)}
\noindent\begin{itemize}
\item coherency (cf. coherency in section \ref{par:modifier-coherency} on page \pageref{par:modifier-coherency})
\end{itemize}

\paragraph{Execution} {Reservation table \textsf{LSU2\_MEMW\_AUXR\_AUXW.X} (Section~\ref{sec:reservation-LSU2-MEMW-AUXR-AUXW.X})}
~
{\begin{lstlisting}
stage ID:
  new offset = _SX_27(offset27);
  new coherency = _ZX_2(coherency);
stage RR:
  new base = GPR[registerZ];
  new address = base + offset;
stage E1:
  new argument2 = GPR[registerT];
  new result2 = _ZX_32(MEM.atomic_ior(address, 0xF, coherency << 32, argument2.32[0], registerT));
stage SR:
  GPR[registerT] = result2;
\end{lstlisting}}
\newpage

\paragraph{Encoding} Format \textsf{LSU\_ALIORSB.D} (Section~\ref{sec:format-LSU-ALIORSB.D}), \verb|lsucode5 = 010|\\

\bitfields{ 1/P, 2/00, 2/-- --, 27/extend27, 1/1, 2/00, 2/-- --, 27/offset27, 1/1, 2/01, 3/111, 2/coherency, 6/registerT, 2/11, 3/010, 1/--, 5/00110, 1/0, 6/registerZ }


\paragraph{Syntax} \texttt{aliorw}\emph{coherency} \emph{extend27\_\discretionary{}{}{}offset27}[\emph{registerZ}] = \emph{registerT}

\paragraph{Description} The effective address is computed by adding the \emph{extend27\_\discretionary{}{}{}offset27} to the \emph{registerZ}.
This instruction implements atomic load-IOR-store on word. The word at effective address is IORed with the \emph{registerT}, and its previous value is returned into the \emph{registerT}. The effective address must be word aligned (multiple of 4).


\paragraph{Modifier(s)}
\noindent\begin{itemize}
\item coherency (cf. coherency in section \ref{par:modifier-coherency} on page \pageref{par:modifier-coherency})
\end{itemize}

\paragraph{Execution} {Reservation table \textsf{LSU2\_MEMW\_AUXR\_AUXW.Y} (Section~\ref{sec:reservation-LSU2-MEMW-AUXR-AUXW.Y})}
~
{\begin{lstlisting}
stage ID:
  new offset = _SX_54(extend27_offset27);
  new coherency = _ZX_2(coherency);
stage RR:
  new base = GPR[registerZ];
  new address = base + offset;
stage E1:
  new argument2 = GPR[registerT];
  new result2 = _ZX_32(MEM.atomic_ior(address, 0xF, coherency << 32, argument2.32[0], registerT));
stage SR:
  GPR[registerT] = result2;
\end{lstlisting}}
\newpage
\subsubsection[{\small ALMAXB: Atomic Load and MAX Byte}]{ALMAXB\hfill Atomic Load and MAX Byte}
\label{instr:ALMAXB}\index{almaxb}
 
\paragraph{Encoding} Format \textsf{LSU\_ALMAXSB} (Section~\ref{sec:format-LSU-ALMAXSB}), \verb|lsucode5 = 000|\\

\bitfields{ 1/P, 2/01, 3/111, 2/coherency, 6/registerT, 2/11, 3/000, 1/--, 5/01001, 1/0, 6/registerZ }


\paragraph{Syntax} \texttt{almaxb}\emph{coherency} [\emph{registerZ}] = \emph{registerT}

\paragraph{Description} The effective address is the value of the \emph{registerZ}.
This instruction implements atomic load-MAX-store on signed byte. The byte at effective address is MAXed with the \emph{registerT}, and its previous value is returned into the \emph{registerT}.


\paragraph{Modifier(s)}
\noindent\begin{itemize}
\item coherency (cf. coherency in section \ref{par:modifier-coherency} on page \pageref{par:modifier-coherency})
\end{itemize}

\paragraph{Execution} {Reservation table \textsf{LSU2\_MEMW\_AUXR\_AUXW} (Section~\ref{sec:reservation-LSU2-MEMW-AUXR-AUXW})}
~
{\begin{lstlisting}
stage ID:
  new coherency = _ZX_2(coherency);
stage RR:
  new address = GPR[registerZ];
stage E1:
  new argument2 = GPR[registerT];
  new result2 = _ZX_8(MEM.atomic_max(address, 0x1, coherency << 32, argument2.8[0], registerT));
stage SR:
  GPR[registerT] = result2;
\end{lstlisting}}
\newpage

\paragraph{Encoding} Format \textsf{LSU\_ALMAXSB.O} (Section~\ref{sec:format-LSU-ALMAXSB.O}), \verb|lsucode5 = 000|\\

\bitfields{ 1/P, 2/00, 2/-- --, 27/offset27, 1/1, 2/01, 3/111, 2/coherency, 6/registerT, 2/11, 3/000, 1/--, 5/01001, 1/0, 6/registerZ }


\paragraph{Syntax} \texttt{almaxb}\emph{coherency} \emph{offset27}[\emph{registerZ}] = \emph{registerT}

\paragraph{Description} The effective address is computed by adding the \emph{offset27} to the \emph{registerZ}.
This instruction implements atomic load-MAX-store on signed byte. The byte at effective address is MAXed with the \emph{registerT}, and its previous value is returned into the \emph{registerT}.


\paragraph{Modifier(s)}
\noindent\begin{itemize}
\item coherency (cf. coherency in section \ref{par:modifier-coherency} on page \pageref{par:modifier-coherency})
\end{itemize}

\paragraph{Execution} {Reservation table \textsf{LSU2\_MEMW\_AUXR\_AUXW.X} (Section~\ref{sec:reservation-LSU2-MEMW-AUXR-AUXW.X})}
~
{\begin{lstlisting}
stage ID:
  new offset = _SX_27(offset27);
  new coherency = _ZX_2(coherency);
stage RR:
  new base = GPR[registerZ];
  new address = base + offset;
stage E1:
  new argument2 = GPR[registerT];
  new result2 = _ZX_8(MEM.atomic_max(address, 0x1, coherency << 32, argument2.8[0], registerT));
stage SR:
  GPR[registerT] = result2;
\end{lstlisting}}
\newpage

\paragraph{Encoding} Format \textsf{LSU\_ALMAXSB.D} (Section~\ref{sec:format-LSU-ALMAXSB.D}), \verb|lsucode5 = 000|\\

\bitfields{ 1/P, 2/00, 2/-- --, 27/extend27, 1/1, 2/00, 2/-- --, 27/offset27, 1/1, 2/01, 3/111, 2/coherency, 6/registerT, 2/11, 3/000, 1/--, 5/01001, 1/0, 6/registerZ }


\paragraph{Syntax} \texttt{almaxb}\emph{coherency} \emph{extend27\_\discretionary{}{}{}offset27}[\emph{registerZ}] = \emph{registerT}

\paragraph{Description} The effective address is computed by adding the \emph{extend27\_\discretionary{}{}{}offset27} to the \emph{registerZ}.
This instruction implements atomic load-MAX-store on signed byte. The byte at effective address is MAXed with the \emph{registerT}, and its previous value is returned into the \emph{registerT}.


\paragraph{Modifier(s)}
\noindent\begin{itemize}
\item coherency (cf. coherency in section \ref{par:modifier-coherency} on page \pageref{par:modifier-coherency})
\end{itemize}

\paragraph{Execution} {Reservation table \textsf{LSU2\_MEMW\_AUXR\_AUXW.Y} (Section~\ref{sec:reservation-LSU2-MEMW-AUXR-AUXW.Y})}
~
{\begin{lstlisting}
stage ID:
  new offset = _SX_54(extend27_offset27);
  new coherency = _ZX_2(coherency);
stage RR:
  new base = GPR[registerZ];
  new address = base + offset;
stage E1:
  new argument2 = GPR[registerT];
  new result2 = _ZX_8(MEM.atomic_max(address, 0x1, coherency << 32, argument2.8[0], registerT));
stage SR:
  GPR[registerT] = result2;
\end{lstlisting}}
\newpage
\subsubsection[{\small ALMAXD: Atomic Load and MAX Double Word}]{ALMAXD\hfill Atomic Load and MAX Double Word}
\label{instr:ALMAXD}\index{almaxd}
 
\paragraph{Encoding} Format \textsf{LSU\_ALMAXSB} (Section~\ref{sec:format-LSU-ALMAXSB}), \verb|lsucode5 = 011|\\

\bitfields{ 1/P, 2/01, 3/111, 2/coherency, 6/registerT, 2/11, 3/011, 1/--, 5/01001, 1/0, 6/registerZ }


\paragraph{Syntax} \texttt{almaxd}\emph{coherency} [\emph{registerZ}] = \emph{registerT}

\paragraph{Description} The effective address is the value of the \emph{registerZ}.
This instruction implements atomic load-MAX-store on signed double word. The double word at effective address is MAXed with the \emph{registerT}, and its previous value is returned into the \emph{registerT}. The effective address must be double word aligned (multiple of 8).


\paragraph{Modifier(s)}
\noindent\begin{itemize}
\item coherency (cf. coherency in section \ref{par:modifier-coherency} on page \pageref{par:modifier-coherency})
\end{itemize}

\paragraph{Execution} {Reservation table \textsf{LSU2\_MEMW\_AUXR\_AUXW} (Section~\ref{sec:reservation-LSU2-MEMW-AUXR-AUXW})}
~
{\begin{lstlisting}
stage ID:
  new coherency = _ZX_2(coherency);
stage RR:
  new address = GPR[registerZ];
stage E1:
  new argument2 = GPR[registerT];
  new result2 = MEM.atomic_max(address, 0xFF, coherency << 32, argument2.64[0], registerT);
stage SR:
  GPR[registerT] = result2;
\end{lstlisting}}
\newpage

\paragraph{Encoding} Format \textsf{LSU\_ALMAXSB.O} (Section~\ref{sec:format-LSU-ALMAXSB.O}), \verb|lsucode5 = 011|\\

\bitfields{ 1/P, 2/00, 2/-- --, 27/offset27, 1/1, 2/01, 3/111, 2/coherency, 6/registerT, 2/11, 3/011, 1/--, 5/01001, 1/0, 6/registerZ }


\paragraph{Syntax} \texttt{almaxd}\emph{coherency} \emph{offset27}[\emph{registerZ}] = \emph{registerT}

\paragraph{Description} The effective address is computed by adding the \emph{offset27} to the \emph{registerZ}.
This instruction implements atomic load-MAX-store on signed double word. The double word at effective address is MAXed with the \emph{registerT}, and its previous value is returned into the \emph{registerT}. The effective address must be double word aligned (multiple of 8).


\paragraph{Modifier(s)}
\noindent\begin{itemize}
\item coherency (cf. coherency in section \ref{par:modifier-coherency} on page \pageref{par:modifier-coherency})
\end{itemize}

\paragraph{Execution} {Reservation table \textsf{LSU2\_MEMW\_AUXR\_AUXW.X} (Section~\ref{sec:reservation-LSU2-MEMW-AUXR-AUXW.X})}
~
{\begin{lstlisting}
stage ID:
  new offset = _SX_27(offset27);
  new coherency = _ZX_2(coherency);
stage RR:
  new base = GPR[registerZ];
  new address = base + offset;
stage E1:
  new argument2 = GPR[registerT];
  new result2 = MEM.atomic_max(address, 0xFF, coherency << 32, argument2.64[0], registerT);
stage SR:
  GPR[registerT] = result2;
\end{lstlisting}}
\newpage

\paragraph{Encoding} Format \textsf{LSU\_ALMAXSB.D} (Section~\ref{sec:format-LSU-ALMAXSB.D}), \verb|lsucode5 = 011|\\

\bitfields{ 1/P, 2/00, 2/-- --, 27/extend27, 1/1, 2/00, 2/-- --, 27/offset27, 1/1, 2/01, 3/111, 2/coherency, 6/registerT, 2/11, 3/011, 1/--, 5/01001, 1/0, 6/registerZ }


\paragraph{Syntax} \texttt{almaxd}\emph{coherency} \emph{extend27\_\discretionary{}{}{}offset27}[\emph{registerZ}] = \emph{registerT}

\paragraph{Description} The effective address is computed by adding the \emph{extend27\_\discretionary{}{}{}offset27} to the \emph{registerZ}.
This instruction implements atomic load-MAX-store on signed double word. The double word at effective address is MAXed with the \emph{registerT}, and its previous value is returned into the \emph{registerT}. The effective address must be double word aligned (multiple of 8).


\paragraph{Modifier(s)}
\noindent\begin{itemize}
\item coherency (cf. coherency in section \ref{par:modifier-coherency} on page \pageref{par:modifier-coherency})
\end{itemize}

\paragraph{Execution} {Reservation table \textsf{LSU2\_MEMW\_AUXR\_AUXW.Y} (Section~\ref{sec:reservation-LSU2-MEMW-AUXR-AUXW.Y})}
~
{\begin{lstlisting}
stage ID:
  new offset = _SX_54(extend27_offset27);
  new coherency = _ZX_2(coherency);
stage RR:
  new base = GPR[registerZ];
  new address = base + offset;
stage E1:
  new argument2 = GPR[registerT];
  new result2 = MEM.atomic_max(address, 0xFF, coherency << 32, argument2.64[0], registerT);
stage SR:
  GPR[registerT] = result2;
\end{lstlisting}}
\newpage
\subsubsection[{\small ALMAXH: Atomic Load and MAX Half Word}]{ALMAXH\hfill Atomic Load and MAX Half Word}
\label{instr:ALMAXH}\index{almaxh}
 
\paragraph{Encoding} Format \textsf{LSU\_ALMAXSB} (Section~\ref{sec:format-LSU-ALMAXSB}), \verb|lsucode5 = 001|\\

\bitfields{ 1/P, 2/01, 3/111, 2/coherency, 6/registerT, 2/11, 3/001, 1/--, 5/01001, 1/0, 6/registerZ }


\paragraph{Syntax} \texttt{almaxh}\emph{coherency} [\emph{registerZ}] = \emph{registerT}

\paragraph{Description} The effective address is the value of the \emph{registerZ}.
This instruction implements atomic load-MAX-store on signed half word. The half word at effective address is MAXed with the \emph{registerT}, and its previous value is returned into the \emph{registerT}. The effective address must be half word aligned (multiple of 2).


\paragraph{Modifier(s)}
\noindent\begin{itemize}
\item coherency (cf. coherency in section \ref{par:modifier-coherency} on page \pageref{par:modifier-coherency})
\end{itemize}

\paragraph{Execution} {Reservation table \textsf{LSU2\_MEMW\_AUXR\_AUXW} (Section~\ref{sec:reservation-LSU2-MEMW-AUXR-AUXW})}
~
{\begin{lstlisting}
stage ID:
  new coherency = _ZX_2(coherency);
stage RR:
  new address = GPR[registerZ];
stage E1:
  new argument2 = GPR[registerT];
  new result2 = _ZX_16(MEM.atomic_max(address, 0x3, coherency << 32, argument2.16[0], registerT));
stage SR:
  GPR[registerT] = result2;
\end{lstlisting}}
\newpage

\paragraph{Encoding} Format \textsf{LSU\_ALMAXSB.O} (Section~\ref{sec:format-LSU-ALMAXSB.O}), \verb|lsucode5 = 001|\\

\bitfields{ 1/P, 2/00, 2/-- --, 27/offset27, 1/1, 2/01, 3/111, 2/coherency, 6/registerT, 2/11, 3/001, 1/--, 5/01001, 1/0, 6/registerZ }


\paragraph{Syntax} \texttt{almaxh}\emph{coherency} \emph{offset27}[\emph{registerZ}] = \emph{registerT}

\paragraph{Description} The effective address is computed by adding the \emph{offset27} to the \emph{registerZ}.
This instruction implements atomic load-MAX-store on signed half word. The half word at effective address is MAXed with the \emph{registerT}, and its previous value is returned into the \emph{registerT}. The effective address must be half word aligned (multiple of 2).


\paragraph{Modifier(s)}
\noindent\begin{itemize}
\item coherency (cf. coherency in section \ref{par:modifier-coherency} on page \pageref{par:modifier-coherency})
\end{itemize}

\paragraph{Execution} {Reservation table \textsf{LSU2\_MEMW\_AUXR\_AUXW.X} (Section~\ref{sec:reservation-LSU2-MEMW-AUXR-AUXW.X})}
~
{\begin{lstlisting}
stage ID:
  new offset = _SX_27(offset27);
  new coherency = _ZX_2(coherency);
stage RR:
  new base = GPR[registerZ];
  new address = base + offset;
stage E1:
  new argument2 = GPR[registerT];
  new result2 = _ZX_16(MEM.atomic_max(address, 0x3, coherency << 32, argument2.16[0], registerT));
stage SR:
  GPR[registerT] = result2;
\end{lstlisting}}
\newpage

\paragraph{Encoding} Format \textsf{LSU\_ALMAXSB.D} (Section~\ref{sec:format-LSU-ALMAXSB.D}), \verb|lsucode5 = 001|\\

\bitfields{ 1/P, 2/00, 2/-- --, 27/extend27, 1/1, 2/00, 2/-- --, 27/offset27, 1/1, 2/01, 3/111, 2/coherency, 6/registerT, 2/11, 3/001, 1/--, 5/01001, 1/0, 6/registerZ }


\paragraph{Syntax} \texttt{almaxh}\emph{coherency} \emph{extend27\_\discretionary{}{}{}offset27}[\emph{registerZ}] = \emph{registerT}

\paragraph{Description} The effective address is computed by adding the \emph{extend27\_\discretionary{}{}{}offset27} to the \emph{registerZ}.
This instruction implements atomic load-MAX-store on signed half word. The half word at effective address is MAXed with the \emph{registerT}, and its previous value is returned into the \emph{registerT}. The effective address must be half word aligned (multiple of 2).


\paragraph{Modifier(s)}
\noindent\begin{itemize}
\item coherency (cf. coherency in section \ref{par:modifier-coherency} on page \pageref{par:modifier-coherency})
\end{itemize}

\paragraph{Execution} {Reservation table \textsf{LSU2\_MEMW\_AUXR\_AUXW.Y} (Section~\ref{sec:reservation-LSU2-MEMW-AUXR-AUXW.Y})}
~
{\begin{lstlisting}
stage ID:
  new offset = _SX_54(extend27_offset27);
  new coherency = _ZX_2(coherency);
stage RR:
  new base = GPR[registerZ];
  new address = base + offset;
stage E1:
  new argument2 = GPR[registerT];
  new result2 = _ZX_16(MEM.atomic_max(address, 0x3, coherency << 32, argument2.16[0], registerT));
stage SR:
  GPR[registerT] = result2;
\end{lstlisting}}
\newpage
\subsubsection[{\small ALMAXUB: Atomic Load and MAX Unsigned Byte}]{ALMAXUB\hfill Atomic Load and MAX Unsigned Byte}
\label{instr:ALMAXUB}\index{almaxub}
 
\paragraph{Encoding} Format \textsf{LSU\_ALMAXUSB} (Section~\ref{sec:format-LSU-ALMAXUSB}), \verb|lsucode5 = 000|\\

\bitfields{ 1/P, 2/01, 3/111, 2/coherency, 6/registerT, 2/11, 3/000, 1/--, 5/01011, 1/0, 6/registerZ }


\paragraph{Syntax} \texttt{almaxub}\emph{coherency} [\emph{registerZ}] = \emph{registerT}

\paragraph{Description} The effective address is the value of the \emph{registerZ}.
This instruction implements atomic load-MAXU-store on signed byte. The byte at effective address is MAXUed with the \emph{registerT}, and its previous value is returned into the \emph{registerT}.


\paragraph{Modifier(s)}
\noindent\begin{itemize}
\item coherency (cf. coherency in section \ref{par:modifier-coherency} on page \pageref{par:modifier-coherency})
\end{itemize}

\paragraph{Execution} {Reservation table \textsf{LSU2\_MEMW\_AUXR\_AUXW} (Section~\ref{sec:reservation-LSU2-MEMW-AUXR-AUXW})}
~
{\begin{lstlisting}
stage ID:
  new coherency = _ZX_2(coherency);
stage RR:
  new address = GPR[registerZ];
stage E1:
  new argument2 = GPR[registerT];
  new result2 = _ZX_8(MEM.atomic_maxu(address, 0x1, coherency << 32, argument2.8[0], registerT));
stage SR:
  GPR[registerT] = result2;
\end{lstlisting}}
\newpage

\paragraph{Encoding} Format \textsf{LSU\_ALMAXUSB.O} (Section~\ref{sec:format-LSU-ALMAXUSB.O}), \verb|lsucode5 = 000|\\

\bitfields{ 1/P, 2/00, 2/-- --, 27/offset27, 1/1, 2/01, 3/111, 2/coherency, 6/registerT, 2/11, 3/000, 1/--, 5/01011, 1/0, 6/registerZ }


\paragraph{Syntax} \texttt{almaxub}\emph{coherency} \emph{offset27}[\emph{registerZ}] = \emph{registerT}

\paragraph{Description} The effective address is computed by adding the \emph{offset27} to the \emph{registerZ}.
This instruction implements atomic load-MAXU-store on signed byte. The byte at effective address is MAXUed with the \emph{registerT}, and its previous value is returned into the \emph{registerT}.


\paragraph{Modifier(s)}
\noindent\begin{itemize}
\item coherency (cf. coherency in section \ref{par:modifier-coherency} on page \pageref{par:modifier-coherency})
\end{itemize}

\paragraph{Execution} {Reservation table \textsf{LSU2\_MEMW\_AUXR\_AUXW.X} (Section~\ref{sec:reservation-LSU2-MEMW-AUXR-AUXW.X})}
~
{\begin{lstlisting}
stage ID:
  new offset = _SX_27(offset27);
  new coherency = _ZX_2(coherency);
stage RR:
  new base = GPR[registerZ];
  new address = base + offset;
stage E1:
  new argument2 = GPR[registerT];
  new result2 = _ZX_8(MEM.atomic_maxu(address, 0x1, coherency << 32, argument2.8[0], registerT));
stage SR:
  GPR[registerT] = result2;
\end{lstlisting}}
\newpage

\paragraph{Encoding} Format \textsf{LSU\_ALMAXUSB.D} (Section~\ref{sec:format-LSU-ALMAXUSB.D}), \verb|lsucode5 = 000|\\

\bitfields{ 1/P, 2/00, 2/-- --, 27/extend27, 1/1, 2/00, 2/-- --, 27/offset27, 1/1, 2/01, 3/111, 2/coherency, 6/registerT, 2/11, 3/000, 1/--, 5/01011, 1/0, 6/registerZ }


\paragraph{Syntax} \texttt{almaxub}\emph{coherency} \emph{extend27\_\discretionary{}{}{}offset27}[\emph{registerZ}] = \emph{registerT}

\paragraph{Description} The effective address is computed by adding the \emph{extend27\_\discretionary{}{}{}offset27} to the \emph{registerZ}.
This instruction implements atomic load-MAXU-store on signed byte. The byte at effective address is MAXUed with the \emph{registerT}, and its previous value is returned into the \emph{registerT}.


\paragraph{Modifier(s)}
\noindent\begin{itemize}
\item coherency (cf. coherency in section \ref{par:modifier-coherency} on page \pageref{par:modifier-coherency})
\end{itemize}

\paragraph{Execution} {Reservation table \textsf{LSU2\_MEMW\_AUXR\_AUXW.Y} (Section~\ref{sec:reservation-LSU2-MEMW-AUXR-AUXW.Y})}
~
{\begin{lstlisting}
stage ID:
  new offset = _SX_54(extend27_offset27);
  new coherency = _ZX_2(coherency);
stage RR:
  new base = GPR[registerZ];
  new address = base + offset;
stage E1:
  new argument2 = GPR[registerT];
  new result2 = _ZX_8(MEM.atomic_maxu(address, 0x1, coherency << 32, argument2.8[0], registerT));
stage SR:
  GPR[registerT] = result2;
\end{lstlisting}}
\newpage
\subsubsection[{\small ALMAXUD: Atomic Load and MAX Unsigned Double Word}]{ALMAXUD\hfill Atomic Load and MAX Unsigned Double Word}
\label{instr:ALMAXUD}\index{almaxud}
 
\paragraph{Encoding} Format \textsf{LSU\_ALMAXUSB} (Section~\ref{sec:format-LSU-ALMAXUSB}), \verb|lsucode5 = 011|\\

\bitfields{ 1/P, 2/01, 3/111, 2/coherency, 6/registerT, 2/11, 3/011, 1/--, 5/01011, 1/0, 6/registerZ }


\paragraph{Syntax} \texttt{almaxud}\emph{coherency} [\emph{registerZ}] = \emph{registerT}

\paragraph{Description} The effective address is the value of the \emph{registerZ}.
This instruction implements atomic load-MAXU-store on signed double word. The double word at effective address is MAXUed with the \emph{registerT}, and its previous value is returned into the \emph{registerT}. The effective address must be double word aligned (multiple of 8).


\paragraph{Modifier(s)}
\noindent\begin{itemize}
\item coherency (cf. coherency in section \ref{par:modifier-coherency} on page \pageref{par:modifier-coherency})
\end{itemize}

\paragraph{Execution} {Reservation table \textsf{LSU2\_MEMW\_AUXR\_AUXW} (Section~\ref{sec:reservation-LSU2-MEMW-AUXR-AUXW})}
~
{\begin{lstlisting}
stage ID:
  new coherency = _ZX_2(coherency);
stage RR:
  new address = GPR[registerZ];
stage E1:
  new argument2 = GPR[registerT];
  new result2 = MEM.atomic_maxu(address, 0xFF, coherency << 32, argument2.64[0], registerT);
stage SR:
  GPR[registerT] = result2;
\end{lstlisting}}
\newpage

\paragraph{Encoding} Format \textsf{LSU\_ALMAXUSB.O} (Section~\ref{sec:format-LSU-ALMAXUSB.O}), \verb|lsucode5 = 011|\\

\bitfields{ 1/P, 2/00, 2/-- --, 27/offset27, 1/1, 2/01, 3/111, 2/coherency, 6/registerT, 2/11, 3/011, 1/--, 5/01011, 1/0, 6/registerZ }


\paragraph{Syntax} \texttt{almaxud}\emph{coherency} \emph{offset27}[\emph{registerZ}] = \emph{registerT}

\paragraph{Description} The effective address is computed by adding the \emph{offset27} to the \emph{registerZ}.
This instruction implements atomic load-MAXU-store on signed double word. The double word at effective address is MAXUed with the \emph{registerT}, and its previous value is returned into the \emph{registerT}. The effective address must be double word aligned (multiple of 8).


\paragraph{Modifier(s)}
\noindent\begin{itemize}
\item coherency (cf. coherency in section \ref{par:modifier-coherency} on page \pageref{par:modifier-coherency})
\end{itemize}

\paragraph{Execution} {Reservation table \textsf{LSU2\_MEMW\_AUXR\_AUXW.X} (Section~\ref{sec:reservation-LSU2-MEMW-AUXR-AUXW.X})}
~
{\begin{lstlisting}
stage ID:
  new offset = _SX_27(offset27);
  new coherency = _ZX_2(coherency);
stage RR:
  new base = GPR[registerZ];
  new address = base + offset;
stage E1:
  new argument2 = GPR[registerT];
  new result2 = MEM.atomic_maxu(address, 0xFF, coherency << 32, argument2.64[0], registerT);
stage SR:
  GPR[registerT] = result2;
\end{lstlisting}}
\newpage

\paragraph{Encoding} Format \textsf{LSU\_ALMAXUSB.D} (Section~\ref{sec:format-LSU-ALMAXUSB.D}), \verb|lsucode5 = 011|\\

\bitfields{ 1/P, 2/00, 2/-- --, 27/extend27, 1/1, 2/00, 2/-- --, 27/offset27, 1/1, 2/01, 3/111, 2/coherency, 6/registerT, 2/11, 3/011, 1/--, 5/01011, 1/0, 6/registerZ }


\paragraph{Syntax} \texttt{almaxud}\emph{coherency} \emph{extend27\_\discretionary{}{}{}offset27}[\emph{registerZ}] = \emph{registerT}

\paragraph{Description} The effective address is computed by adding the \emph{extend27\_\discretionary{}{}{}offset27} to the \emph{registerZ}.
This instruction implements atomic load-MAXU-store on signed double word. The double word at effective address is MAXUed with the \emph{registerT}, and its previous value is returned into the \emph{registerT}. The effective address must be double word aligned (multiple of 8).


\paragraph{Modifier(s)}
\noindent\begin{itemize}
\item coherency (cf. coherency in section \ref{par:modifier-coherency} on page \pageref{par:modifier-coherency})
\end{itemize}

\paragraph{Execution} {Reservation table \textsf{LSU2\_MEMW\_AUXR\_AUXW.Y} (Section~\ref{sec:reservation-LSU2-MEMW-AUXR-AUXW.Y})}
~
{\begin{lstlisting}
stage ID:
  new offset = _SX_54(extend27_offset27);
  new coherency = _ZX_2(coherency);
stage RR:
  new base = GPR[registerZ];
  new address = base + offset;
stage E1:
  new argument2 = GPR[registerT];
  new result2 = MEM.atomic_maxu(address, 0xFF, coherency << 32, argument2.64[0], registerT);
stage SR:
  GPR[registerT] = result2;
\end{lstlisting}}
\newpage
\subsubsection[{\small ALMAXUH: Atomic Load and MAX Unsigned Half Word}]{ALMAXUH\hfill Atomic Load and MAX Unsigned Half Word}
\label{instr:ALMAXUH}\index{almaxuh}
 
\paragraph{Encoding} Format \textsf{LSU\_ALMAXUSB} (Section~\ref{sec:format-LSU-ALMAXUSB}), \verb|lsucode5 = 001|\\

\bitfields{ 1/P, 2/01, 3/111, 2/coherency, 6/registerT, 2/11, 3/001, 1/--, 5/01011, 1/0, 6/registerZ }


\paragraph{Syntax} \texttt{almaxuh}\emph{coherency} [\emph{registerZ}] = \emph{registerT}

\paragraph{Description} The effective address is the value of the \emph{registerZ}.
This instruction implements atomic load-MAXU-store on signed half word. The half word at effective address is MAXUed with the \emph{registerT}, and its previous value is returned into the \emph{registerT}. The effective address must be half word aligned (multiple of 2).


\paragraph{Modifier(s)}
\noindent\begin{itemize}
\item coherency (cf. coherency in section \ref{par:modifier-coherency} on page \pageref{par:modifier-coherency})
\end{itemize}

\paragraph{Execution} {Reservation table \textsf{LSU2\_MEMW\_AUXR\_AUXW} (Section~\ref{sec:reservation-LSU2-MEMW-AUXR-AUXW})}
~
{\begin{lstlisting}
stage ID:
  new coherency = _ZX_2(coherency);
stage RR:
  new address = GPR[registerZ];
stage E1:
  new argument2 = GPR[registerT];
  new result2 = _ZX_16(MEM.atomic_maxu(address, 0x3, coherency << 32, argument2.16[0], registerT));
stage SR:
  GPR[registerT] = result2;
\end{lstlisting}}
\newpage

\paragraph{Encoding} Format \textsf{LSU\_ALMAXUSB.O} (Section~\ref{sec:format-LSU-ALMAXUSB.O}), \verb|lsucode5 = 001|\\

\bitfields{ 1/P, 2/00, 2/-- --, 27/offset27, 1/1, 2/01, 3/111, 2/coherency, 6/registerT, 2/11, 3/001, 1/--, 5/01011, 1/0, 6/registerZ }


\paragraph{Syntax} \texttt{almaxuh}\emph{coherency} \emph{offset27}[\emph{registerZ}] = \emph{registerT}

\paragraph{Description} The effective address is computed by adding the \emph{offset27} to the \emph{registerZ}.
This instruction implements atomic load-MAXU-store on signed half word. The half word at effective address is MAXUed with the \emph{registerT}, and its previous value is returned into the \emph{registerT}. The effective address must be half word aligned (multiple of 2).


\paragraph{Modifier(s)}
\noindent\begin{itemize}
\item coherency (cf. coherency in section \ref{par:modifier-coherency} on page \pageref{par:modifier-coherency})
\end{itemize}

\paragraph{Execution} {Reservation table \textsf{LSU2\_MEMW\_AUXR\_AUXW.X} (Section~\ref{sec:reservation-LSU2-MEMW-AUXR-AUXW.X})}
~
{\begin{lstlisting}
stage ID:
  new offset = _SX_27(offset27);
  new coherency = _ZX_2(coherency);
stage RR:
  new base = GPR[registerZ];
  new address = base + offset;
stage E1:
  new argument2 = GPR[registerT];
  new result2 = _ZX_16(MEM.atomic_maxu(address, 0x3, coherency << 32, argument2.16[0], registerT));
stage SR:
  GPR[registerT] = result2;
\end{lstlisting}}
\newpage

\paragraph{Encoding} Format \textsf{LSU\_ALMAXUSB.D} (Section~\ref{sec:format-LSU-ALMAXUSB.D}), \verb|lsucode5 = 001|\\

\bitfields{ 1/P, 2/00, 2/-- --, 27/extend27, 1/1, 2/00, 2/-- --, 27/offset27, 1/1, 2/01, 3/111, 2/coherency, 6/registerT, 2/11, 3/001, 1/--, 5/01011, 1/0, 6/registerZ }


\paragraph{Syntax} \texttt{almaxuh}\emph{coherency} \emph{extend27\_\discretionary{}{}{}offset27}[\emph{registerZ}] = \emph{registerT}

\paragraph{Description} The effective address is computed by adding the \emph{extend27\_\discretionary{}{}{}offset27} to the \emph{registerZ}.
This instruction implements atomic load-MAXU-store on signed half word. The half word at effective address is MAXUed with the \emph{registerT}, and its previous value is returned into the \emph{registerT}. The effective address must be half word aligned (multiple of 2).


\paragraph{Modifier(s)}
\noindent\begin{itemize}
\item coherency (cf. coherency in section \ref{par:modifier-coherency} on page \pageref{par:modifier-coherency})
\end{itemize}

\paragraph{Execution} {Reservation table \textsf{LSU2\_MEMW\_AUXR\_AUXW.Y} (Section~\ref{sec:reservation-LSU2-MEMW-AUXR-AUXW.Y})}
~
{\begin{lstlisting}
stage ID:
  new offset = _SX_54(extend27_offset27);
  new coherency = _ZX_2(coherency);
stage RR:
  new base = GPR[registerZ];
  new address = base + offset;
stage E1:
  new argument2 = GPR[registerT];
  new result2 = _ZX_16(MEM.atomic_maxu(address, 0x3, coherency << 32, argument2.16[0], registerT));
stage SR:
  GPR[registerT] = result2;
\end{lstlisting}}
\newpage
\subsubsection[{\small ALMAXUW: Atomic Load and MAX Unsigned Word}]{ALMAXUW\hfill Atomic Load and MAX Unsigned Word}
\label{instr:ALMAXUW}\index{almaxuw}
 
\paragraph{Encoding} Format \textsf{LSU\_ALMAXUSB} (Section~\ref{sec:format-LSU-ALMAXUSB}), \verb|lsucode5 = 010|\\

\bitfields{ 1/P, 2/01, 3/111, 2/coherency, 6/registerT, 2/11, 3/010, 1/--, 5/01011, 1/0, 6/registerZ }


\paragraph{Syntax} \texttt{almaxuw}\emph{coherency} [\emph{registerZ}] = \emph{registerT}

\paragraph{Description} The effective address is the value of the \emph{registerZ}.
This instruction implements atomic load-MAXU-store on signed word. The word at effective address is MAXUed with the \emph{registerT}, and its previous value is returned into the \emph{registerT}. The effective address must be word aligned (multiple of 4).


\paragraph{Modifier(s)}
\noindent\begin{itemize}
\item coherency (cf. coherency in section \ref{par:modifier-coherency} on page \pageref{par:modifier-coherency})
\end{itemize}

\paragraph{Execution} {Reservation table \textsf{LSU2\_MEMW\_AUXR\_AUXW} (Section~\ref{sec:reservation-LSU2-MEMW-AUXR-AUXW})}
~
{\begin{lstlisting}
stage ID:
  new coherency = _ZX_2(coherency);
stage RR:
  new address = GPR[registerZ];
stage E1:
  new argument2 = GPR[registerT];
  new result2 = _ZX_32(MEM.atomic_maxu(address, 0xF, coherency << 32, argument2.32[0], registerT));
stage SR:
  GPR[registerT] = result2;
\end{lstlisting}}
\newpage

\paragraph{Encoding} Format \textsf{LSU\_ALMAXUSB.O} (Section~\ref{sec:format-LSU-ALMAXUSB.O}), \verb|lsucode5 = 010|\\

\bitfields{ 1/P, 2/00, 2/-- --, 27/offset27, 1/1, 2/01, 3/111, 2/coherency, 6/registerT, 2/11, 3/010, 1/--, 5/01011, 1/0, 6/registerZ }


\paragraph{Syntax} \texttt{almaxuw}\emph{coherency} \emph{offset27}[\emph{registerZ}] = \emph{registerT}

\paragraph{Description} The effective address is computed by adding the \emph{offset27} to the \emph{registerZ}.
This instruction implements atomic load-MAXU-store on signed word. The word at effective address is MAXUed with the \emph{registerT}, and its previous value is returned into the \emph{registerT}. The effective address must be word aligned (multiple of 4).


\paragraph{Modifier(s)}
\noindent\begin{itemize}
\item coherency (cf. coherency in section \ref{par:modifier-coherency} on page \pageref{par:modifier-coherency})
\end{itemize}

\paragraph{Execution} {Reservation table \textsf{LSU2\_MEMW\_AUXR\_AUXW.X} (Section~\ref{sec:reservation-LSU2-MEMW-AUXR-AUXW.X})}
~
{\begin{lstlisting}
stage ID:
  new offset = _SX_27(offset27);
  new coherency = _ZX_2(coherency);
stage RR:
  new base = GPR[registerZ];
  new address = base + offset;
stage E1:
  new argument2 = GPR[registerT];
  new result2 = _ZX_32(MEM.atomic_maxu(address, 0xF, coherency << 32, argument2.32[0], registerT));
stage SR:
  GPR[registerT] = result2;
\end{lstlisting}}
\newpage

\paragraph{Encoding} Format \textsf{LSU\_ALMAXUSB.D} (Section~\ref{sec:format-LSU-ALMAXUSB.D}), \verb|lsucode5 = 010|\\

\bitfields{ 1/P, 2/00, 2/-- --, 27/extend27, 1/1, 2/00, 2/-- --, 27/offset27, 1/1, 2/01, 3/111, 2/coherency, 6/registerT, 2/11, 3/010, 1/--, 5/01011, 1/0, 6/registerZ }


\paragraph{Syntax} \texttt{almaxuw}\emph{coherency} \emph{extend27\_\discretionary{}{}{}offset27}[\emph{registerZ}] = \emph{registerT}

\paragraph{Description} The effective address is computed by adding the \emph{extend27\_\discretionary{}{}{}offset27} to the \emph{registerZ}.
This instruction implements atomic load-MAXU-store on signed word. The word at effective address is MAXUed with the \emph{registerT}, and its previous value is returned into the \emph{registerT}. The effective address must be word aligned (multiple of 4).


\paragraph{Modifier(s)}
\noindent\begin{itemize}
\item coherency (cf. coherency in section \ref{par:modifier-coherency} on page \pageref{par:modifier-coherency})
\end{itemize}

\paragraph{Execution} {Reservation table \textsf{LSU2\_MEMW\_AUXR\_AUXW.Y} (Section~\ref{sec:reservation-LSU2-MEMW-AUXR-AUXW.Y})}
~
{\begin{lstlisting}
stage ID:
  new offset = _SX_54(extend27_offset27);
  new coherency = _ZX_2(coherency);
stage RR:
  new base = GPR[registerZ];
  new address = base + offset;
stage E1:
  new argument2 = GPR[registerT];
  new result2 = _ZX_32(MEM.atomic_maxu(address, 0xF, coherency << 32, argument2.32[0], registerT));
stage SR:
  GPR[registerT] = result2;
\end{lstlisting}}
\newpage
\subsubsection[{\small ALMAXW: Atomic Load and MAX Word}]{ALMAXW\hfill Atomic Load and MAX Word}
\label{instr:ALMAXW}\index{almaxw}
 
\paragraph{Encoding} Format \textsf{LSU\_ALMAXSB} (Section~\ref{sec:format-LSU-ALMAXSB}), \verb|lsucode5 = 010|\\

\bitfields{ 1/P, 2/01, 3/111, 2/coherency, 6/registerT, 2/11, 3/010, 1/--, 5/01001, 1/0, 6/registerZ }


\paragraph{Syntax} \texttt{almaxw}\emph{coherency} [\emph{registerZ}] = \emph{registerT}

\paragraph{Description} The effective address is the value of the \emph{registerZ}.
This instruction implements atomic load-MAX-store on signed word. The word at effective address is MAXed with the \emph{registerT}, and its previous value is returned into the \emph{registerT}. The effective address must be word aligned (multiple of 4).


\paragraph{Modifier(s)}
\noindent\begin{itemize}
\item coherency (cf. coherency in section \ref{par:modifier-coherency} on page \pageref{par:modifier-coherency})
\end{itemize}

\paragraph{Execution} {Reservation table \textsf{LSU2\_MEMW\_AUXR\_AUXW} (Section~\ref{sec:reservation-LSU2-MEMW-AUXR-AUXW})}
~
{\begin{lstlisting}
stage ID:
  new coherency = _ZX_2(coherency);
stage RR:
  new address = GPR[registerZ];
stage E1:
  new argument2 = GPR[registerT];
  new result2 = _ZX_32(MEM.atomic_max(address, 0xF, coherency << 32, argument2.32[0], registerT));
stage SR:
  GPR[registerT] = result2;
\end{lstlisting}}
\newpage

\paragraph{Encoding} Format \textsf{LSU\_ALMAXSB.O} (Section~\ref{sec:format-LSU-ALMAXSB.O}), \verb|lsucode5 = 010|\\

\bitfields{ 1/P, 2/00, 2/-- --, 27/offset27, 1/1, 2/01, 3/111, 2/coherency, 6/registerT, 2/11, 3/010, 1/--, 5/01001, 1/0, 6/registerZ }


\paragraph{Syntax} \texttt{almaxw}\emph{coherency} \emph{offset27}[\emph{registerZ}] = \emph{registerT}

\paragraph{Description} The effective address is computed by adding the \emph{offset27} to the \emph{registerZ}.
This instruction implements atomic load-MAX-store on signed word. The word at effective address is MAXed with the \emph{registerT}, and its previous value is returned into the \emph{registerT}. The effective address must be word aligned (multiple of 4).


\paragraph{Modifier(s)}
\noindent\begin{itemize}
\item coherency (cf. coherency in section \ref{par:modifier-coherency} on page \pageref{par:modifier-coherency})
\end{itemize}

\paragraph{Execution} {Reservation table \textsf{LSU2\_MEMW\_AUXR\_AUXW.X} (Section~\ref{sec:reservation-LSU2-MEMW-AUXR-AUXW.X})}
~
{\begin{lstlisting}
stage ID:
  new offset = _SX_27(offset27);
  new coherency = _ZX_2(coherency);
stage RR:
  new base = GPR[registerZ];
  new address = base + offset;
stage E1:
  new argument2 = GPR[registerT];
  new result2 = _ZX_32(MEM.atomic_max(address, 0xF, coherency << 32, argument2.32[0], registerT));
stage SR:
  GPR[registerT] = result2;
\end{lstlisting}}
\newpage

\paragraph{Encoding} Format \textsf{LSU\_ALMAXSB.D} (Section~\ref{sec:format-LSU-ALMAXSB.D}), \verb|lsucode5 = 010|\\

\bitfields{ 1/P, 2/00, 2/-- --, 27/extend27, 1/1, 2/00, 2/-- --, 27/offset27, 1/1, 2/01, 3/111, 2/coherency, 6/registerT, 2/11, 3/010, 1/--, 5/01001, 1/0, 6/registerZ }


\paragraph{Syntax} \texttt{almaxw}\emph{coherency} \emph{extend27\_\discretionary{}{}{}offset27}[\emph{registerZ}] = \emph{registerT}

\paragraph{Description} The effective address is computed by adding the \emph{extend27\_\discretionary{}{}{}offset27} to the \emph{registerZ}.
This instruction implements atomic load-MAX-store on signed word. The word at effective address is MAXed with the \emph{registerT}, and its previous value is returned into the \emph{registerT}. The effective address must be word aligned (multiple of 4).


\paragraph{Modifier(s)}
\noindent\begin{itemize}
\item coherency (cf. coherency in section \ref{par:modifier-coherency} on page \pageref{par:modifier-coherency})
\end{itemize}

\paragraph{Execution} {Reservation table \textsf{LSU2\_MEMW\_AUXR\_AUXW.Y} (Section~\ref{sec:reservation-LSU2-MEMW-AUXR-AUXW.Y})}
~
{\begin{lstlisting}
stage ID:
  new offset = _SX_54(extend27_offset27);
  new coherency = _ZX_2(coherency);
stage RR:
  new base = GPR[registerZ];
  new address = base + offset;
stage E1:
  new argument2 = GPR[registerT];
  new result2 = _ZX_32(MEM.atomic_max(address, 0xF, coherency << 32, argument2.32[0], registerT));
stage SR:
  GPR[registerT] = result2;
\end{lstlisting}}
\newpage
\subsubsection[{\small ALMINB: Atomic Load and MIN Byte}]{ALMINB\hfill Atomic Load and MIN Byte}
\label{instr:ALMINB}\index{alminb}
 
\paragraph{Encoding} Format \textsf{LSU\_ALMINSB} (Section~\ref{sec:format-LSU-ALMINSB}), \verb|lsucode5 = 000|\\

\bitfields{ 1/P, 2/01, 3/111, 2/coherency, 6/registerT, 2/11, 3/000, 1/--, 5/01000, 1/0, 6/registerZ }


\paragraph{Syntax} \texttt{alminb}\emph{coherency} [\emph{registerZ}] = \emph{registerT}

\paragraph{Description} The effective address is the value of the \emph{registerZ}.
This instruction implements atomic load-MIN-store on signed byte. The byte at effective address is MINed with the \emph{registerT}, and its previous value is returned into the \emph{registerT}.


\paragraph{Modifier(s)}
\noindent\begin{itemize}
\item coherency (cf. coherency in section \ref{par:modifier-coherency} on page \pageref{par:modifier-coherency})
\end{itemize}

\paragraph{Execution} {Reservation table \textsf{LSU2\_MEMW\_AUXR\_AUXW} (Section~\ref{sec:reservation-LSU2-MEMW-AUXR-AUXW})}
~
{\begin{lstlisting}
stage ID:
  new coherency = _ZX_2(coherency);
stage RR:
  new address = GPR[registerZ];
stage E1:
  new argument2 = GPR[registerT];
  new result2 = _ZX_8(MEM.atomic_min(address, 0x1, coherency << 32, argument2.8[0], registerT));
stage SR:
  GPR[registerT] = result2;
\end{lstlisting}}
\newpage

\paragraph{Encoding} Format \textsf{LSU\_ALMINSB.O} (Section~\ref{sec:format-LSU-ALMINSB.O}), \verb|lsucode5 = 000|\\

\bitfields{ 1/P, 2/00, 2/-- --, 27/offset27, 1/1, 2/01, 3/111, 2/coherency, 6/registerT, 2/11, 3/000, 1/--, 5/01000, 1/0, 6/registerZ }


\paragraph{Syntax} \texttt{alminb}\emph{coherency} \emph{offset27}[\emph{registerZ}] = \emph{registerT}

\paragraph{Description} The effective address is computed by adding the \emph{offset27} to the \emph{registerZ}.
This instruction implements atomic load-MIN-store on signed byte. The byte at effective address is MINed with the \emph{registerT}, and its previous value is returned into the \emph{registerT}.


\paragraph{Modifier(s)}
\noindent\begin{itemize}
\item coherency (cf. coherency in section \ref{par:modifier-coherency} on page \pageref{par:modifier-coherency})
\end{itemize}

\paragraph{Execution} {Reservation table \textsf{LSU2\_MEMW\_AUXR\_AUXW.X} (Section~\ref{sec:reservation-LSU2-MEMW-AUXR-AUXW.X})}
~
{\begin{lstlisting}
stage ID:
  new offset = _SX_27(offset27);
  new coherency = _ZX_2(coherency);
stage RR:
  new base = GPR[registerZ];
  new address = base + offset;
stage E1:
  new argument2 = GPR[registerT];
  new result2 = _ZX_8(MEM.atomic_min(address, 0x1, coherency << 32, argument2.8[0], registerT));
stage SR:
  GPR[registerT] = result2;
\end{lstlisting}}
\newpage

\paragraph{Encoding} Format \textsf{LSU\_ALMINSB.D} (Section~\ref{sec:format-LSU-ALMINSB.D}), \verb|lsucode5 = 000|\\

\bitfields{ 1/P, 2/00, 2/-- --, 27/extend27, 1/1, 2/00, 2/-- --, 27/offset27, 1/1, 2/01, 3/111, 2/coherency, 6/registerT, 2/11, 3/000, 1/--, 5/01000, 1/0, 6/registerZ }


\paragraph{Syntax} \texttt{alminb}\emph{coherency} \emph{extend27\_\discretionary{}{}{}offset27}[\emph{registerZ}] = \emph{registerT}

\paragraph{Description} The effective address is computed by adding the \emph{extend27\_\discretionary{}{}{}offset27} to the \emph{registerZ}.
This instruction implements atomic load-MIN-store on signed byte. The byte at effective address is MINed with the \emph{registerT}, and its previous value is returned into the \emph{registerT}.


\paragraph{Modifier(s)}
\noindent\begin{itemize}
\item coherency (cf. coherency in section \ref{par:modifier-coherency} on page \pageref{par:modifier-coherency})
\end{itemize}

\paragraph{Execution} {Reservation table \textsf{LSU2\_MEMW\_AUXR\_AUXW.Y} (Section~\ref{sec:reservation-LSU2-MEMW-AUXR-AUXW.Y})}
~
{\begin{lstlisting}
stage ID:
  new offset = _SX_54(extend27_offset27);
  new coherency = _ZX_2(coherency);
stage RR:
  new base = GPR[registerZ];
  new address = base + offset;
stage E1:
  new argument2 = GPR[registerT];
  new result2 = _ZX_8(MEM.atomic_min(address, 0x1, coherency << 32, argument2.8[0], registerT));
stage SR:
  GPR[registerT] = result2;
\end{lstlisting}}
\newpage
\subsubsection[{\small ALMIND: Atomic Load and MIN Double Word}]{ALMIND\hfill Atomic Load and MIN Double Word}
\label{instr:ALMIND}\index{almind}
 
\paragraph{Encoding} Format \textsf{LSU\_ALMINSB} (Section~\ref{sec:format-LSU-ALMINSB}), \verb|lsucode5 = 011|\\

\bitfields{ 1/P, 2/01, 3/111, 2/coherency, 6/registerT, 2/11, 3/011, 1/--, 5/01000, 1/0, 6/registerZ }


\paragraph{Syntax} \texttt{almind}\emph{coherency} [\emph{registerZ}] = \emph{registerT}

\paragraph{Description} The effective address is the value of the \emph{registerZ}.
This instruction implements atomic load-MIN-store on signed double word. The double word at effective address is MINed with the \emph{registerT}, and its previous value is returned into the \emph{registerT}. The effective address must be double word aligned (multiple of 8).


\paragraph{Modifier(s)}
\noindent\begin{itemize}
\item coherency (cf. coherency in section \ref{par:modifier-coherency} on page \pageref{par:modifier-coherency})
\end{itemize}

\paragraph{Execution} {Reservation table \textsf{LSU2\_MEMW\_AUXR\_AUXW} (Section~\ref{sec:reservation-LSU2-MEMW-AUXR-AUXW})}
~
{\begin{lstlisting}
stage ID:
  new coherency = _ZX_2(coherency);
stage RR:
  new address = GPR[registerZ];
stage E1:
  new argument2 = GPR[registerT];
  new result2 = MEM.atomic_min(address, 0xFF, coherency << 32, argument2.64[0], registerT);
stage SR:
  GPR[registerT] = result2;
\end{lstlisting}}
\newpage

\paragraph{Encoding} Format \textsf{LSU\_ALMINSB.O} (Section~\ref{sec:format-LSU-ALMINSB.O}), \verb|lsucode5 = 011|\\

\bitfields{ 1/P, 2/00, 2/-- --, 27/offset27, 1/1, 2/01, 3/111, 2/coherency, 6/registerT, 2/11, 3/011, 1/--, 5/01000, 1/0, 6/registerZ }


\paragraph{Syntax} \texttt{almind}\emph{coherency} \emph{offset27}[\emph{registerZ}] = \emph{registerT}

\paragraph{Description} The effective address is computed by adding the \emph{offset27} to the \emph{registerZ}.
This instruction implements atomic load-MIN-store on signed double word. The double word at effective address is MINed with the \emph{registerT}, and its previous value is returned into the \emph{registerT}. The effective address must be double word aligned (multiple of 8).


\paragraph{Modifier(s)}
\noindent\begin{itemize}
\item coherency (cf. coherency in section \ref{par:modifier-coherency} on page \pageref{par:modifier-coherency})
\end{itemize}

\paragraph{Execution} {Reservation table \textsf{LSU2\_MEMW\_AUXR\_AUXW.X} (Section~\ref{sec:reservation-LSU2-MEMW-AUXR-AUXW.X})}
~
{\begin{lstlisting}
stage ID:
  new offset = _SX_27(offset27);
  new coherency = _ZX_2(coherency);
stage RR:
  new base = GPR[registerZ];
  new address = base + offset;
stage E1:
  new argument2 = GPR[registerT];
  new result2 = MEM.atomic_min(address, 0xFF, coherency << 32, argument2.64[0], registerT);
stage SR:
  GPR[registerT] = result2;
\end{lstlisting}}
\newpage

\paragraph{Encoding} Format \textsf{LSU\_ALMINSB.D} (Section~\ref{sec:format-LSU-ALMINSB.D}), \verb|lsucode5 = 011|\\

\bitfields{ 1/P, 2/00, 2/-- --, 27/extend27, 1/1, 2/00, 2/-- --, 27/offset27, 1/1, 2/01, 3/111, 2/coherency, 6/registerT, 2/11, 3/011, 1/--, 5/01000, 1/0, 6/registerZ }


\paragraph{Syntax} \texttt{almind}\emph{coherency} \emph{extend27\_\discretionary{}{}{}offset27}[\emph{registerZ}] = \emph{registerT}

\paragraph{Description} The effective address is computed by adding the \emph{extend27\_\discretionary{}{}{}offset27} to the \emph{registerZ}.
This instruction implements atomic load-MIN-store on signed double word. The double word at effective address is MINed with the \emph{registerT}, and its previous value is returned into the \emph{registerT}. The effective address must be double word aligned (multiple of 8).


\paragraph{Modifier(s)}
\noindent\begin{itemize}
\item coherency (cf. coherency in section \ref{par:modifier-coherency} on page \pageref{par:modifier-coherency})
\end{itemize}

\paragraph{Execution} {Reservation table \textsf{LSU2\_MEMW\_AUXR\_AUXW.Y} (Section~\ref{sec:reservation-LSU2-MEMW-AUXR-AUXW.Y})}
~
{\begin{lstlisting}
stage ID:
  new offset = _SX_54(extend27_offset27);
  new coherency = _ZX_2(coherency);
stage RR:
  new base = GPR[registerZ];
  new address = base + offset;
stage E1:
  new argument2 = GPR[registerT];
  new result2 = MEM.atomic_min(address, 0xFF, coherency << 32, argument2.64[0], registerT);
stage SR:
  GPR[registerT] = result2;
\end{lstlisting}}
\newpage
\subsubsection[{\small ALMINH: Atomic Load and MIN Half Word}]{ALMINH\hfill Atomic Load and MIN Half Word}
\label{instr:ALMINH}\index{alminh}
 
\paragraph{Encoding} Format \textsf{LSU\_ALMINSB} (Section~\ref{sec:format-LSU-ALMINSB}), \verb|lsucode5 = 001|\\

\bitfields{ 1/P, 2/01, 3/111, 2/coherency, 6/registerT, 2/11, 3/001, 1/--, 5/01000, 1/0, 6/registerZ }


\paragraph{Syntax} \texttt{alminh}\emph{coherency} [\emph{registerZ}] = \emph{registerT}

\paragraph{Description} The effective address is the value of the \emph{registerZ}.
This instruction implements atomic load-MIN-store on signed half word. The half word at effective address is MINed with the \emph{registerT}, and its previous value is returned into the \emph{registerT}. The effective address must be half word aligned (multiple of 2).


\paragraph{Modifier(s)}
\noindent\begin{itemize}
\item coherency (cf. coherency in section \ref{par:modifier-coherency} on page \pageref{par:modifier-coherency})
\end{itemize}

\paragraph{Execution} {Reservation table \textsf{LSU2\_MEMW\_AUXR\_AUXW} (Section~\ref{sec:reservation-LSU2-MEMW-AUXR-AUXW})}
~
{\begin{lstlisting}
stage ID:
  new coherency = _ZX_2(coherency);
stage RR:
  new address = GPR[registerZ];
stage E1:
  new argument2 = GPR[registerT];
  new result2 = _ZX_16(MEM.atomic_min(address, 0x3, coherency << 32, argument2.16[0], registerT));
stage SR:
  GPR[registerT] = result2;
\end{lstlisting}}
\newpage

\paragraph{Encoding} Format \textsf{LSU\_ALMINSB.O} (Section~\ref{sec:format-LSU-ALMINSB.O}), \verb|lsucode5 = 001|\\

\bitfields{ 1/P, 2/00, 2/-- --, 27/offset27, 1/1, 2/01, 3/111, 2/coherency, 6/registerT, 2/11, 3/001, 1/--, 5/01000, 1/0, 6/registerZ }


\paragraph{Syntax} \texttt{alminh}\emph{coherency} \emph{offset27}[\emph{registerZ}] = \emph{registerT}

\paragraph{Description} The effective address is computed by adding the \emph{offset27} to the \emph{registerZ}.
This instruction implements atomic load-MIN-store on signed half word. The half word at effective address is MINed with the \emph{registerT}, and its previous value is returned into the \emph{registerT}. The effective address must be half word aligned (multiple of 2).


\paragraph{Modifier(s)}
\noindent\begin{itemize}
\item coherency (cf. coherency in section \ref{par:modifier-coherency} on page \pageref{par:modifier-coherency})
\end{itemize}

\paragraph{Execution} {Reservation table \textsf{LSU2\_MEMW\_AUXR\_AUXW.X} (Section~\ref{sec:reservation-LSU2-MEMW-AUXR-AUXW.X})}
~
{\begin{lstlisting}
stage ID:
  new offset = _SX_27(offset27);
  new coherency = _ZX_2(coherency);
stage RR:
  new base = GPR[registerZ];
  new address = base + offset;
stage E1:
  new argument2 = GPR[registerT];
  new result2 = _ZX_16(MEM.atomic_min(address, 0x3, coherency << 32, argument2.16[0], registerT));
stage SR:
  GPR[registerT] = result2;
\end{lstlisting}}
\newpage

\paragraph{Encoding} Format \textsf{LSU\_ALMINSB.D} (Section~\ref{sec:format-LSU-ALMINSB.D}), \verb|lsucode5 = 001|\\

\bitfields{ 1/P, 2/00, 2/-- --, 27/extend27, 1/1, 2/00, 2/-- --, 27/offset27, 1/1, 2/01, 3/111, 2/coherency, 6/registerT, 2/11, 3/001, 1/--, 5/01000, 1/0, 6/registerZ }


\paragraph{Syntax} \texttt{alminh}\emph{coherency} \emph{extend27\_\discretionary{}{}{}offset27}[\emph{registerZ}] = \emph{registerT}

\paragraph{Description} The effective address is computed by adding the \emph{extend27\_\discretionary{}{}{}offset27} to the \emph{registerZ}.
This instruction implements atomic load-MIN-store on signed half word. The half word at effective address is MINed with the \emph{registerT}, and its previous value is returned into the \emph{registerT}. The effective address must be half word aligned (multiple of 2).


\paragraph{Modifier(s)}
\noindent\begin{itemize}
\item coherency (cf. coherency in section \ref{par:modifier-coherency} on page \pageref{par:modifier-coherency})
\end{itemize}

\paragraph{Execution} {Reservation table \textsf{LSU2\_MEMW\_AUXR\_AUXW.Y} (Section~\ref{sec:reservation-LSU2-MEMW-AUXR-AUXW.Y})}
~
{\begin{lstlisting}
stage ID:
  new offset = _SX_54(extend27_offset27);
  new coherency = _ZX_2(coherency);
stage RR:
  new base = GPR[registerZ];
  new address = base + offset;
stage E1:
  new argument2 = GPR[registerT];
  new result2 = _ZX_16(MEM.atomic_min(address, 0x3, coherency << 32, argument2.16[0], registerT));
stage SR:
  GPR[registerT] = result2;
\end{lstlisting}}
\newpage
\subsubsection[{\small ALMINUB: Atomic Load and MIN Unsigned Byte}]{ALMINUB\hfill Atomic Load and MIN Unsigned Byte}
\label{instr:ALMINUB}\index{alminub}
 
\paragraph{Encoding} Format \textsf{LSU\_ALMINUSB} (Section~\ref{sec:format-LSU-ALMINUSB}), \verb|lsucode5 = 000|\\

\bitfields{ 1/P, 2/01, 3/111, 2/coherency, 6/registerT, 2/11, 3/000, 1/--, 5/01010, 1/0, 6/registerZ }


\paragraph{Syntax} \texttt{alminub}\emph{coherency} [\emph{registerZ}] = \emph{registerT}

\paragraph{Description} The effective address is the value of the \emph{registerZ}.
This instruction implements atomic load-MINU-store on signed byte. The byte at effective address is MINUed with the \emph{registerT}, and its previous value is returned into the \emph{registerT}.


\paragraph{Modifier(s)}
\noindent\begin{itemize}
\item coherency (cf. coherency in section \ref{par:modifier-coherency} on page \pageref{par:modifier-coherency})
\end{itemize}

\paragraph{Execution} {Reservation table \textsf{LSU2\_MEMW\_AUXR\_AUXW} (Section~\ref{sec:reservation-LSU2-MEMW-AUXR-AUXW})}
~
{\begin{lstlisting}
stage ID:
  new coherency = _ZX_2(coherency);
stage RR:
  new address = GPR[registerZ];
stage E1:
  new argument2 = GPR[registerT];
  new result2 = _ZX_8(MEM.atomic_minu(address, 0x1, coherency << 32, argument2.8[0], registerT));
stage SR:
  GPR[registerT] = result2;
\end{lstlisting}}
\newpage

\paragraph{Encoding} Format \textsf{LSU\_ALMINUSB.O} (Section~\ref{sec:format-LSU-ALMINUSB.O}), \verb|lsucode5 = 000|\\

\bitfields{ 1/P, 2/00, 2/-- --, 27/offset27, 1/1, 2/01, 3/111, 2/coherency, 6/registerT, 2/11, 3/000, 1/--, 5/01010, 1/0, 6/registerZ }


\paragraph{Syntax} \texttt{alminub}\emph{coherency} \emph{offset27}[\emph{registerZ}] = \emph{registerT}

\paragraph{Description} The effective address is computed by adding the \emph{offset27} to the \emph{registerZ}.
This instruction implements atomic load-MINU-store on signed byte. The byte at effective address is MINUed with the \emph{registerT}, and its previous value is returned into the \emph{registerT}.


\paragraph{Modifier(s)}
\noindent\begin{itemize}
\item coherency (cf. coherency in section \ref{par:modifier-coherency} on page \pageref{par:modifier-coherency})
\end{itemize}

\paragraph{Execution} {Reservation table \textsf{LSU2\_MEMW\_AUXR\_AUXW.X} (Section~\ref{sec:reservation-LSU2-MEMW-AUXR-AUXW.X})}
~
{\begin{lstlisting}
stage ID:
  new offset = _SX_27(offset27);
  new coherency = _ZX_2(coherency);
stage RR:
  new base = GPR[registerZ];
  new address = base + offset;
stage E1:
  new argument2 = GPR[registerT];
  new result2 = _ZX_8(MEM.atomic_minu(address, 0x1, coherency << 32, argument2.8[0], registerT));
stage SR:
  GPR[registerT] = result2;
\end{lstlisting}}
\newpage

\paragraph{Encoding} Format \textsf{LSU\_ALMINUSB.D} (Section~\ref{sec:format-LSU-ALMINUSB.D}), \verb|lsucode5 = 000|\\

\bitfields{ 1/P, 2/00, 2/-- --, 27/extend27, 1/1, 2/00, 2/-- --, 27/offset27, 1/1, 2/01, 3/111, 2/coherency, 6/registerT, 2/11, 3/000, 1/--, 5/01010, 1/0, 6/registerZ }


\paragraph{Syntax} \texttt{alminub}\emph{coherency} \emph{extend27\_\discretionary{}{}{}offset27}[\emph{registerZ}] = \emph{registerT}

\paragraph{Description} The effective address is computed by adding the \emph{extend27\_\discretionary{}{}{}offset27} to the \emph{registerZ}.
This instruction implements atomic load-MINU-store on signed byte. The byte at effective address is MINUed with the \emph{registerT}, and its previous value is returned into the \emph{registerT}.


\paragraph{Modifier(s)}
\noindent\begin{itemize}
\item coherency (cf. coherency in section \ref{par:modifier-coherency} on page \pageref{par:modifier-coherency})
\end{itemize}

\paragraph{Execution} {Reservation table \textsf{LSU2\_MEMW\_AUXR\_AUXW.Y} (Section~\ref{sec:reservation-LSU2-MEMW-AUXR-AUXW.Y})}
~
{\begin{lstlisting}
stage ID:
  new offset = _SX_54(extend27_offset27);
  new coherency = _ZX_2(coherency);
stage RR:
  new base = GPR[registerZ];
  new address = base + offset;
stage E1:
  new argument2 = GPR[registerT];
  new result2 = _ZX_8(MEM.atomic_minu(address, 0x1, coherency << 32, argument2.8[0], registerT));
stage SR:
  GPR[registerT] = result2;
\end{lstlisting}}
\newpage
\subsubsection[{\small ALMINUD: Atomic Load and MIN Unsigned Double Word}]{ALMINUD\hfill Atomic Load and MIN Unsigned Double Word}
\label{instr:ALMINUD}\index{alminud}
 
\paragraph{Encoding} Format \textsf{LSU\_ALMINUSB} (Section~\ref{sec:format-LSU-ALMINUSB}), \verb|lsucode5 = 011|\\

\bitfields{ 1/P, 2/01, 3/111, 2/coherency, 6/registerT, 2/11, 3/011, 1/--, 5/01010, 1/0, 6/registerZ }


\paragraph{Syntax} \texttt{alminud}\emph{coherency} [\emph{registerZ}] = \emph{registerT}

\paragraph{Description} The effective address is the value of the \emph{registerZ}.
This instruction implements atomic load-MINU-store on signed double word. The double word at effective address is MINUed with the \emph{registerT}, and its previous value is returned into the \emph{registerT}. The effective address must be double word aligned (multiple of 8).


\paragraph{Modifier(s)}
\noindent\begin{itemize}
\item coherency (cf. coherency in section \ref{par:modifier-coherency} on page \pageref{par:modifier-coherency})
\end{itemize}

\paragraph{Execution} {Reservation table \textsf{LSU2\_MEMW\_AUXR\_AUXW} (Section~\ref{sec:reservation-LSU2-MEMW-AUXR-AUXW})}
~
{\begin{lstlisting}
stage ID:
  new coherency = _ZX_2(coherency);
stage RR:
  new address = GPR[registerZ];
stage E1:
  new argument2 = GPR[registerT];
  new result2 = MEM.atomic_minu(address, 0xFF, coherency << 32, argument2.64[0], registerT);
stage SR:
  GPR[registerT] = result2;
\end{lstlisting}}
\newpage

\paragraph{Encoding} Format \textsf{LSU\_ALMINUSB.O} (Section~\ref{sec:format-LSU-ALMINUSB.O}), \verb|lsucode5 = 011|\\

\bitfields{ 1/P, 2/00, 2/-- --, 27/offset27, 1/1, 2/01, 3/111, 2/coherency, 6/registerT, 2/11, 3/011, 1/--, 5/01010, 1/0, 6/registerZ }


\paragraph{Syntax} \texttt{alminud}\emph{coherency} \emph{offset27}[\emph{registerZ}] = \emph{registerT}

\paragraph{Description} The effective address is computed by adding the \emph{offset27} to the \emph{registerZ}.
This instruction implements atomic load-MINU-store on signed double word. The double word at effective address is MINUed with the \emph{registerT}, and its previous value is returned into the \emph{registerT}. The effective address must be double word aligned (multiple of 8).


\paragraph{Modifier(s)}
\noindent\begin{itemize}
\item coherency (cf. coherency in section \ref{par:modifier-coherency} on page \pageref{par:modifier-coherency})
\end{itemize}

\paragraph{Execution} {Reservation table \textsf{LSU2\_MEMW\_AUXR\_AUXW.X} (Section~\ref{sec:reservation-LSU2-MEMW-AUXR-AUXW.X})}
~
{\begin{lstlisting}
stage ID:
  new offset = _SX_27(offset27);
  new coherency = _ZX_2(coherency);
stage RR:
  new base = GPR[registerZ];
  new address = base + offset;
stage E1:
  new argument2 = GPR[registerT];
  new result2 = MEM.atomic_minu(address, 0xFF, coherency << 32, argument2.64[0], registerT);
stage SR:
  GPR[registerT] = result2;
\end{lstlisting}}
\newpage

\paragraph{Encoding} Format \textsf{LSU\_ALMINUSB.D} (Section~\ref{sec:format-LSU-ALMINUSB.D}), \verb|lsucode5 = 011|\\

\bitfields{ 1/P, 2/00, 2/-- --, 27/extend27, 1/1, 2/00, 2/-- --, 27/offset27, 1/1, 2/01, 3/111, 2/coherency, 6/registerT, 2/11, 3/011, 1/--, 5/01010, 1/0, 6/registerZ }


\paragraph{Syntax} \texttt{alminud}\emph{coherency} \emph{extend27\_\discretionary{}{}{}offset27}[\emph{registerZ}] = \emph{registerT}

\paragraph{Description} The effective address is computed by adding the \emph{extend27\_\discretionary{}{}{}offset27} to the \emph{registerZ}.
This instruction implements atomic load-MINU-store on signed double word. The double word at effective address is MINUed with the \emph{registerT}, and its previous value is returned into the \emph{registerT}. The effective address must be double word aligned (multiple of 8).


\paragraph{Modifier(s)}
\noindent\begin{itemize}
\item coherency (cf. coherency in section \ref{par:modifier-coherency} on page \pageref{par:modifier-coherency})
\end{itemize}

\paragraph{Execution} {Reservation table \textsf{LSU2\_MEMW\_AUXR\_AUXW.Y} (Section~\ref{sec:reservation-LSU2-MEMW-AUXR-AUXW.Y})}
~
{\begin{lstlisting}
stage ID:
  new offset = _SX_54(extend27_offset27);
  new coherency = _ZX_2(coherency);
stage RR:
  new base = GPR[registerZ];
  new address = base + offset;
stage E1:
  new argument2 = GPR[registerT];
  new result2 = MEM.atomic_minu(address, 0xFF, coherency << 32, argument2.64[0], registerT);
stage SR:
  GPR[registerT] = result2;
\end{lstlisting}}
\newpage
\subsubsection[{\small ALMINUH: Atomic Load and MIN Unsigned Half Word}]{ALMINUH\hfill Atomic Load and MIN Unsigned Half Word}
\label{instr:ALMINUH}\index{alminuh}
 
\paragraph{Encoding} Format \textsf{LSU\_ALMINUSB} (Section~\ref{sec:format-LSU-ALMINUSB}), \verb|lsucode5 = 001|\\

\bitfields{ 1/P, 2/01, 3/111, 2/coherency, 6/registerT, 2/11, 3/001, 1/--, 5/01010, 1/0, 6/registerZ }


\paragraph{Syntax} \texttt{alminuh}\emph{coherency} [\emph{registerZ}] = \emph{registerT}

\paragraph{Description} The effective address is the value of the \emph{registerZ}.
This instruction implements atomic load-MINU-store on signed half word. The half word at effective address is MINUed with the \emph{registerT}, and its previous value is returned into the \emph{registerT}. The effective address must be half word aligned (multiple of 2).


\paragraph{Modifier(s)}
\noindent\begin{itemize}
\item coherency (cf. coherency in section \ref{par:modifier-coherency} on page \pageref{par:modifier-coherency})
\end{itemize}

\paragraph{Execution} {Reservation table \textsf{LSU2\_MEMW\_AUXR\_AUXW} (Section~\ref{sec:reservation-LSU2-MEMW-AUXR-AUXW})}
~
{\begin{lstlisting}
stage ID:
  new coherency = _ZX_2(coherency);
stage RR:
  new address = GPR[registerZ];
stage E1:
  new argument2 = GPR[registerT];
  new result2 = _ZX_16(MEM.atomic_minu(address, 0x3, coherency << 32, argument2.16[0], registerT));
stage SR:
  GPR[registerT] = result2;
\end{lstlisting}}
\newpage

\paragraph{Encoding} Format \textsf{LSU\_ALMINUSB.O} (Section~\ref{sec:format-LSU-ALMINUSB.O}), \verb|lsucode5 = 001|\\

\bitfields{ 1/P, 2/00, 2/-- --, 27/offset27, 1/1, 2/01, 3/111, 2/coherency, 6/registerT, 2/11, 3/001, 1/--, 5/01010, 1/0, 6/registerZ }


\paragraph{Syntax} \texttt{alminuh}\emph{coherency} \emph{offset27}[\emph{registerZ}] = \emph{registerT}

\paragraph{Description} The effective address is computed by adding the \emph{offset27} to the \emph{registerZ}.
This instruction implements atomic load-MINU-store on signed half word. The half word at effective address is MINUed with the \emph{registerT}, and its previous value is returned into the \emph{registerT}. The effective address must be half word aligned (multiple of 2).


\paragraph{Modifier(s)}
\noindent\begin{itemize}
\item coherency (cf. coherency in section \ref{par:modifier-coherency} on page \pageref{par:modifier-coherency})
\end{itemize}

\paragraph{Execution} {Reservation table \textsf{LSU2\_MEMW\_AUXR\_AUXW.X} (Section~\ref{sec:reservation-LSU2-MEMW-AUXR-AUXW.X})}
~
{\begin{lstlisting}
stage ID:
  new offset = _SX_27(offset27);
  new coherency = _ZX_2(coherency);
stage RR:
  new base = GPR[registerZ];
  new address = base + offset;
stage E1:
  new argument2 = GPR[registerT];
  new result2 = _ZX_16(MEM.atomic_minu(address, 0x3, coherency << 32, argument2.16[0], registerT));
stage SR:
  GPR[registerT] = result2;
\end{lstlisting}}
\newpage

\paragraph{Encoding} Format \textsf{LSU\_ALMINUSB.D} (Section~\ref{sec:format-LSU-ALMINUSB.D}), \verb|lsucode5 = 001|\\

\bitfields{ 1/P, 2/00, 2/-- --, 27/extend27, 1/1, 2/00, 2/-- --, 27/offset27, 1/1, 2/01, 3/111, 2/coherency, 6/registerT, 2/11, 3/001, 1/--, 5/01010, 1/0, 6/registerZ }


\paragraph{Syntax} \texttt{alminuh}\emph{coherency} \emph{extend27\_\discretionary{}{}{}offset27}[\emph{registerZ}] = \emph{registerT}

\paragraph{Description} The effective address is computed by adding the \emph{extend27\_\discretionary{}{}{}offset27} to the \emph{registerZ}.
This instruction implements atomic load-MINU-store on signed half word. The half word at effective address is MINUed with the \emph{registerT}, and its previous value is returned into the \emph{registerT}. The effective address must be half word aligned (multiple of 2).


\paragraph{Modifier(s)}
\noindent\begin{itemize}
\item coherency (cf. coherency in section \ref{par:modifier-coherency} on page \pageref{par:modifier-coherency})
\end{itemize}

\paragraph{Execution} {Reservation table \textsf{LSU2\_MEMW\_AUXR\_AUXW.Y} (Section~\ref{sec:reservation-LSU2-MEMW-AUXR-AUXW.Y})}
~
{\begin{lstlisting}
stage ID:
  new offset = _SX_54(extend27_offset27);
  new coherency = _ZX_2(coherency);
stage RR:
  new base = GPR[registerZ];
  new address = base + offset;
stage E1:
  new argument2 = GPR[registerT];
  new result2 = _ZX_16(MEM.atomic_minu(address, 0x3, coherency << 32, argument2.16[0], registerT));
stage SR:
  GPR[registerT] = result2;
\end{lstlisting}}
\newpage
\subsubsection[{\small ALMINUW: Atomic Load and MIN Unsigned Word}]{ALMINUW\hfill Atomic Load and MIN Unsigned Word}
\label{instr:ALMINUW}\index{alminuw}
 
\paragraph{Encoding} Format \textsf{LSU\_ALMINUSB} (Section~\ref{sec:format-LSU-ALMINUSB}), \verb|lsucode5 = 010|\\

\bitfields{ 1/P, 2/01, 3/111, 2/coherency, 6/registerT, 2/11, 3/010, 1/--, 5/01010, 1/0, 6/registerZ }


\paragraph{Syntax} \texttt{alminuw}\emph{coherency} [\emph{registerZ}] = \emph{registerT}

\paragraph{Description} The effective address is the value of the \emph{registerZ}.
This instruction implements atomic load-MINU-store on signed word. The word at effective address is MINUed with the \emph{registerT}, and its previous value is returned into the \emph{registerT}. The effective address must be word aligned (multiple of 4).


\paragraph{Modifier(s)}
\noindent\begin{itemize}
\item coherency (cf. coherency in section \ref{par:modifier-coherency} on page \pageref{par:modifier-coherency})
\end{itemize}

\paragraph{Execution} {Reservation table \textsf{LSU2\_MEMW\_AUXR\_AUXW} (Section~\ref{sec:reservation-LSU2-MEMW-AUXR-AUXW})}
~
{\begin{lstlisting}
stage ID:
  new coherency = _ZX_2(coherency);
stage RR:
  new address = GPR[registerZ];
stage E1:
  new argument2 = GPR[registerT];
  new result2 = _ZX_32(MEM.atomic_minu(address, 0xF, coherency << 32, argument2.32[0], registerT));
stage SR:
  GPR[registerT] = result2;
\end{lstlisting}}
\newpage

\paragraph{Encoding} Format \textsf{LSU\_ALMINUSB.O} (Section~\ref{sec:format-LSU-ALMINUSB.O}), \verb|lsucode5 = 010|\\

\bitfields{ 1/P, 2/00, 2/-- --, 27/offset27, 1/1, 2/01, 3/111, 2/coherency, 6/registerT, 2/11, 3/010, 1/--, 5/01010, 1/0, 6/registerZ }


\paragraph{Syntax} \texttt{alminuw}\emph{coherency} \emph{offset27}[\emph{registerZ}] = \emph{registerT}

\paragraph{Description} The effective address is computed by adding the \emph{offset27} to the \emph{registerZ}.
This instruction implements atomic load-MINU-store on signed word. The word at effective address is MINUed with the \emph{registerT}, and its previous value is returned into the \emph{registerT}. The effective address must be word aligned (multiple of 4).


\paragraph{Modifier(s)}
\noindent\begin{itemize}
\item coherency (cf. coherency in section \ref{par:modifier-coherency} on page \pageref{par:modifier-coherency})
\end{itemize}

\paragraph{Execution} {Reservation table \textsf{LSU2\_MEMW\_AUXR\_AUXW.X} (Section~\ref{sec:reservation-LSU2-MEMW-AUXR-AUXW.X})}
~
{\begin{lstlisting}
stage ID:
  new offset = _SX_27(offset27);
  new coherency = _ZX_2(coherency);
stage RR:
  new base = GPR[registerZ];
  new address = base + offset;
stage E1:
  new argument2 = GPR[registerT];
  new result2 = _ZX_32(MEM.atomic_minu(address, 0xF, coherency << 32, argument2.32[0], registerT));
stage SR:
  GPR[registerT] = result2;
\end{lstlisting}}
\newpage

\paragraph{Encoding} Format \textsf{LSU\_ALMINUSB.D} (Section~\ref{sec:format-LSU-ALMINUSB.D}), \verb|lsucode5 = 010|\\

\bitfields{ 1/P, 2/00, 2/-- --, 27/extend27, 1/1, 2/00, 2/-- --, 27/offset27, 1/1, 2/01, 3/111, 2/coherency, 6/registerT, 2/11, 3/010, 1/--, 5/01010, 1/0, 6/registerZ }


\paragraph{Syntax} \texttt{alminuw}\emph{coherency} \emph{extend27\_\discretionary{}{}{}offset27}[\emph{registerZ}] = \emph{registerT}

\paragraph{Description} The effective address is computed by adding the \emph{extend27\_\discretionary{}{}{}offset27} to the \emph{registerZ}.
This instruction implements atomic load-MINU-store on signed word. The word at effective address is MINUed with the \emph{registerT}, and its previous value is returned into the \emph{registerT}. The effective address must be word aligned (multiple of 4).


\paragraph{Modifier(s)}
\noindent\begin{itemize}
\item coherency (cf. coherency in section \ref{par:modifier-coherency} on page \pageref{par:modifier-coherency})
\end{itemize}

\paragraph{Execution} {Reservation table \textsf{LSU2\_MEMW\_AUXR\_AUXW.Y} (Section~\ref{sec:reservation-LSU2-MEMW-AUXR-AUXW.Y})}
~
{\begin{lstlisting}
stage ID:
  new offset = _SX_54(extend27_offset27);
  new coherency = _ZX_2(coherency);
stage RR:
  new base = GPR[registerZ];
  new address = base + offset;
stage E1:
  new argument2 = GPR[registerT];
  new result2 = _ZX_32(MEM.atomic_minu(address, 0xF, coherency << 32, argument2.32[0], registerT));
stage SR:
  GPR[registerT] = result2;
\end{lstlisting}}
\newpage
\subsubsection[{\small ALMINW: Atomic Load and MIN Word}]{ALMINW\hfill Atomic Load and MIN Word}
\label{instr:ALMINW}\index{alminw}
 
\paragraph{Encoding} Format \textsf{LSU\_ALMINSB} (Section~\ref{sec:format-LSU-ALMINSB}), \verb|lsucode5 = 010|\\

\bitfields{ 1/P, 2/01, 3/111, 2/coherency, 6/registerT, 2/11, 3/010, 1/--, 5/01000, 1/0, 6/registerZ }


\paragraph{Syntax} \texttt{alminw}\emph{coherency} [\emph{registerZ}] = \emph{registerT}

\paragraph{Description} The effective address is the value of the \emph{registerZ}.
This instruction implements atomic load-MIN-store on signed word. The word at effective address is MINed with the \emph{registerT}, and its previous value is returned into the \emph{registerT}. The effective address must be word aligned (multiple of 4).


\paragraph{Modifier(s)}
\noindent\begin{itemize}
\item coherency (cf. coherency in section \ref{par:modifier-coherency} on page \pageref{par:modifier-coherency})
\end{itemize}

\paragraph{Execution} {Reservation table \textsf{LSU2\_MEMW\_AUXR\_AUXW} (Section~\ref{sec:reservation-LSU2-MEMW-AUXR-AUXW})}
~
{\begin{lstlisting}
stage ID:
  new coherency = _ZX_2(coherency);
stage RR:
  new address = GPR[registerZ];
stage E1:
  new argument2 = GPR[registerT];
  new result2 = _ZX_32(MEM.atomic_min(address, 0xF, coherency << 32, argument2.32[0], registerT));
stage SR:
  GPR[registerT] = result2;
\end{lstlisting}}
\newpage

\paragraph{Encoding} Format \textsf{LSU\_ALMINSB.O} (Section~\ref{sec:format-LSU-ALMINSB.O}), \verb|lsucode5 = 010|\\

\bitfields{ 1/P, 2/00, 2/-- --, 27/offset27, 1/1, 2/01, 3/111, 2/coherency, 6/registerT, 2/11, 3/010, 1/--, 5/01000, 1/0, 6/registerZ }


\paragraph{Syntax} \texttt{alminw}\emph{coherency} \emph{offset27}[\emph{registerZ}] = \emph{registerT}

\paragraph{Description} The effective address is computed by adding the \emph{offset27} to the \emph{registerZ}.
This instruction implements atomic load-MIN-store on signed word. The word at effective address is MINed with the \emph{registerT}, and its previous value is returned into the \emph{registerT}. The effective address must be word aligned (multiple of 4).


\paragraph{Modifier(s)}
\noindent\begin{itemize}
\item coherency (cf. coherency in section \ref{par:modifier-coherency} on page \pageref{par:modifier-coherency})
\end{itemize}

\paragraph{Execution} {Reservation table \textsf{LSU2\_MEMW\_AUXR\_AUXW.X} (Section~\ref{sec:reservation-LSU2-MEMW-AUXR-AUXW.X})}
~
{\begin{lstlisting}
stage ID:
  new offset = _SX_27(offset27);
  new coherency = _ZX_2(coherency);
stage RR:
  new base = GPR[registerZ];
  new address = base + offset;
stage E1:
  new argument2 = GPR[registerT];
  new result2 = _ZX_32(MEM.atomic_min(address, 0xF, coherency << 32, argument2.32[0], registerT));
stage SR:
  GPR[registerT] = result2;
\end{lstlisting}}
\newpage

\paragraph{Encoding} Format \textsf{LSU\_ALMINSB.D} (Section~\ref{sec:format-LSU-ALMINSB.D}), \verb|lsucode5 = 010|\\

\bitfields{ 1/P, 2/00, 2/-- --, 27/extend27, 1/1, 2/00, 2/-- --, 27/offset27, 1/1, 2/01, 3/111, 2/coherency, 6/registerT, 2/11, 3/010, 1/--, 5/01000, 1/0, 6/registerZ }


\paragraph{Syntax} \texttt{alminw}\emph{coherency} \emph{extend27\_\discretionary{}{}{}offset27}[\emph{registerZ}] = \emph{registerT}

\paragraph{Description} The effective address is computed by adding the \emph{extend27\_\discretionary{}{}{}offset27} to the \emph{registerZ}.
This instruction implements atomic load-MIN-store on signed word. The word at effective address is MINed with the \emph{registerT}, and its previous value is returned into the \emph{registerT}. The effective address must be word aligned (multiple of 4).


\paragraph{Modifier(s)}
\noindent\begin{itemize}
\item coherency (cf. coherency in section \ref{par:modifier-coherency} on page \pageref{par:modifier-coherency})
\end{itemize}

\paragraph{Execution} {Reservation table \textsf{LSU2\_MEMW\_AUXR\_AUXW.Y} (Section~\ref{sec:reservation-LSU2-MEMW-AUXR-AUXW.Y})}
~
{\begin{lstlisting}
stage ID:
  new offset = _SX_54(extend27_offset27);
  new coherency = _ZX_2(coherency);
stage RR:
  new base = GPR[registerZ];
  new address = base + offset;
stage E1:
  new argument2 = GPR[registerT];
  new result2 = _ZX_32(MEM.atomic_min(address, 0xF, coherency << 32, argument2.32[0], registerT));
stage SR:
  GPR[registerT] = result2;
\end{lstlisting}}
\newpage
\subsubsection[{\small ALW: Atomic Load Word}]{ALW\hfill Atomic Load Word}
\label{instr:ALW}\index{alw}
 
\paragraph{Encoding} Format \textsf{LSU\_ALSB} (Section~\ref{sec:format-LSU-ALSB}), \verb|lsucode5 = 010|\\

\bitfields{ 1/P, 2/01, 3/111, 2/coherency, 6/registerW, 2/11, 3/010, 1/--, 5/00000, 1/0, 6/registerZ }


\paragraph{Syntax} \texttt{alw}\emph{coherency} \emph{registerW} = [\emph{registerZ}]

\paragraph{Description} The effective address is the value of the \emph{registerZ}.
The \emph{registerW} is atomically loaded with zero extension from the memory word starting at the effective address. This instruction is provided to directly operate on the last level of the memory hierarchy, which is typically a processor DDR or a cluster TCM, irrespective of the MMU DCP (Data Cache Policy) settings. The effective address must be word aligned (multiple of 4).


\paragraph{Modifier(s)}
\noindent\begin{itemize}
\item coherency (cf. coherency in section \ref{par:modifier-coherency} on page \pageref{par:modifier-coherency})
\end{itemize}

\paragraph{Execution} {Reservation table \textsf{LSU\_MEMW\_AUXW} (Section~\ref{sec:reservation-LSU-MEMW-AUXW})}
~
{\begin{lstlisting}
stage ID:
  new coherency = _ZX_2(coherency);
stage RR:
  new address = GPR[registerZ];
  new result1 = _ZX_32(MEM.atomic_load(address, 0xF, coherency << 32, registerW));
stage CR:
  GPR[registerW] = result1;
\end{lstlisting}}
\newpage

\paragraph{Encoding} Format \textsf{LSU\_ALSB.O} (Section~\ref{sec:format-LSU-ALSB.O}), \verb|lsucode5 = 010|\\

\bitfields{ 1/P, 2/00, 2/-- --, 27/offset27, 1/1, 2/01, 3/111, 2/coherency, 6/registerW, 2/11, 3/010, 1/--, 5/00000, 1/0, 6/registerZ }


\paragraph{Syntax} \texttt{alw}\emph{coherency} \emph{registerW} = \emph{offset27}[\emph{registerZ}]

\paragraph{Description} The effective address is computed by adding the \emph{offset27} to the \emph{registerZ}.
The \emph{registerW} is atomically loaded with zero extension from the memory word starting at the effective address. This instruction is provided to directly operate on the last level of the memory hierarchy, which is typically a processor DDR or a cluster TCM, irrespective of the MMU DCP (Data Cache Policy) settings. The effective address must be word aligned (multiple of 4).


\paragraph{Modifier(s)}
\noindent\begin{itemize}
\item coherency (cf. coherency in section \ref{par:modifier-coherency} on page \pageref{par:modifier-coherency})
\end{itemize}

\paragraph{Execution} {Reservation table \textsf{LSU\_MEMW\_AUXW.X} (Section~\ref{sec:reservation-LSU-MEMW-AUXW.X})}
~
{\begin{lstlisting}
stage ID:
  new offset = _SX_27(offset27);
  new coherency = _ZX_2(coherency);
stage RR:
  new base = GPR[registerZ];
  new address = base + offset;
  new result1 = _ZX_32(MEM.atomic_load(address, 0xF, coherency << 32, registerW));
stage CR:
  GPR[registerW] = result1;
\end{lstlisting}}
\newpage

\paragraph{Encoding} Format \textsf{LSU\_ALSB.D} (Section~\ref{sec:format-LSU-ALSB.D}), \verb|lsucode5 = 010|\\

\bitfields{ 1/P, 2/00, 2/-- --, 27/extend27, 1/1, 2/00, 2/-- --, 27/offset27, 1/1, 2/01, 3/111, 2/coherency, 6/registerW, 2/11, 3/010, 1/--, 5/00000, 1/0, 6/registerZ }


\paragraph{Syntax} \texttt{alw}\emph{coherency} \emph{registerW} = \emph{extend27\_\discretionary{}{}{}offset27}[\emph{registerZ}]

\paragraph{Description} The effective address is computed by adding the \emph{extend27\_\discretionary{}{}{}offset27} to the \emph{registerZ}.
The \emph{registerW} is atomically loaded with zero extension from the memory word starting at the effective address. This instruction is provided to directly operate on the last level of the memory hierarchy, which is typically a processor DDR or a cluster TCM, irrespective of the MMU DCP (Data Cache Policy) settings. The effective address must be word aligned (multiple of 4).


\paragraph{Modifier(s)}
\noindent\begin{itemize}
\item coherency (cf. coherency in section \ref{par:modifier-coherency} on page \pageref{par:modifier-coherency})
\end{itemize}

\paragraph{Execution} {Reservation table \textsf{LSU\_MEMW\_AUXW.Y} (Section~\ref{sec:reservation-LSU-MEMW-AUXW.Y})}
~
{\begin{lstlisting}
stage ID:
  new offset = _SX_54(extend27_offset27);
  new coherency = _ZX_2(coherency);
stage RR:
  new base = GPR[registerZ];
  new address = base + offset;
  new result1 = _ZX_32(MEM.atomic_load(address, 0xF, coherency << 32, registerW));
stage CR:
  GPR[registerW] = result1;
\end{lstlisting}}
\newpage
\subsubsection[{\small ASADDB: Atomic Store Add Byte}]{ASADDB\hfill Atomic Store Add Byte}
\label{instr:ASADDB}\index{asaddb}
 
\paragraph{Encoding} Format \textsf{LSU\_ASADDSB} (Section~\ref{sec:format-LSU-ASADDSB}), \verb|lsucode5 = 000|\\

\bitfields{ 1/P, 2/01, 3/111, 2/coherency, 6/registerT, 2/11, 3/000, 1/--, 5/10100, 1/0, 6/registerZ }


\paragraph{Syntax} \texttt{asaddb}\emph{coherency} [\emph{registerZ}] = \emph{registerT}

\paragraph{Description} The effective address is the value of the \emph{registerZ}.
This instruction implements atomic add-store on byte. The byte at effective address is added with the \emph{registerT}.


\paragraph{Modifier(s)}
\noindent\begin{itemize}
\item coherency (cf. coherency in section \ref{par:modifier-coherency} on page \pageref{par:modifier-coherency})
\end{itemize}

\paragraph{Execution} {Reservation table \textsf{LSU2\_MEMW\_AUXR} (Section~\ref{sec:reservation-LSU2-MEMW-AUXR})}
~
{\begin{lstlisting}
stage ID:
  new coherency = _ZX_2(coherency);
stage RR:
  new address = GPR[registerZ];
stage E1:
  new argument2 = GPR[registerT];
stage E3:
  MEM.atomic_add(address, 0x1, coherency << 32, argument2.8[0], registerT);
\end{lstlisting}}
\newpage

\paragraph{Encoding} Format \textsf{LSU\_ASADDSB.O} (Section~\ref{sec:format-LSU-ASADDSB.O}), \verb|lsucode5 = 000|\\

\bitfields{ 1/P, 2/00, 2/-- --, 27/offset27, 1/1, 2/01, 3/111, 2/coherency, 6/registerT, 2/11, 3/000, 1/--, 5/10100, 1/0, 6/registerZ }


\paragraph{Syntax} \texttt{asaddb}\emph{coherency} \emph{offset27}[\emph{registerZ}] = \emph{registerT}

\paragraph{Description} The effective address is computed by adding the \emph{offset27} to the \emph{registerZ}.
This instruction implements atomic add-store on byte. The byte at effective address is added with the \emph{registerT}.


\paragraph{Modifier(s)}
\noindent\begin{itemize}
\item coherency (cf. coherency in section \ref{par:modifier-coherency} on page \pageref{par:modifier-coherency})
\end{itemize}

\paragraph{Execution} {Reservation table \textsf{LSU2\_MEMW\_AUXR.X} (Section~\ref{sec:reservation-LSU2-MEMW-AUXR.X})}
~
{\begin{lstlisting}
stage ID:
  new offset = _SX_27(offset27);
  new coherency = _ZX_2(coherency);
stage RR:
  new base = GPR[registerZ];
  new address = base + offset;
stage E1:
  new argument2 = GPR[registerT];
stage E3:
  MEM.atomic_add(address, 0x1, coherency << 32, argument2.8[0], registerT);
\end{lstlisting}}
\newpage

\paragraph{Encoding} Format \textsf{LSU\_ASADDSB.D} (Section~\ref{sec:format-LSU-ASADDSB.D}), \verb|lsucode5 = 000|\\

\bitfields{ 1/P, 2/00, 2/-- --, 27/extend27, 1/1, 2/00, 2/-- --, 27/offset27, 1/1, 2/01, 3/111, 2/coherency, 6/registerT, 2/11, 3/000, 1/--, 5/10100, 1/0, 6/registerZ }


\paragraph{Syntax} \texttt{asaddb}\emph{coherency} \emph{extend27\_\discretionary{}{}{}offset27}[\emph{registerZ}] = \emph{registerT}

\paragraph{Description} The effective address is computed by adding the \emph{extend27\_\discretionary{}{}{}offset27} to the \emph{registerZ}.
This instruction implements atomic add-store on byte. The byte at effective address is added with the \emph{registerT}.


\paragraph{Modifier(s)}
\noindent\begin{itemize}
\item coherency (cf. coherency in section \ref{par:modifier-coherency} on page \pageref{par:modifier-coherency})
\end{itemize}

\paragraph{Execution} {Reservation table \textsf{LSU2\_MEMW\_AUXR.Y} (Section~\ref{sec:reservation-LSU2-MEMW-AUXR.Y})}
~
{\begin{lstlisting}
stage ID:
  new offset = _SX_54(extend27_offset27);
  new coherency = _ZX_2(coherency);
stage RR:
  new base = GPR[registerZ];
  new address = base + offset;
stage E1:
  new argument2 = GPR[registerT];
stage E3:
  MEM.atomic_add(address, 0x1, coherency << 32, argument2.8[0], registerT);
\end{lstlisting}}
\newpage
\subsubsection[{\small ASADDD: Atomic Store Add Double Word}]{ASADDD\hfill Atomic Store Add Double Word}
\label{instr:ASADDD}\index{asaddd}
 
\paragraph{Encoding} Format \textsf{LSU\_ASADDSB} (Section~\ref{sec:format-LSU-ASADDSB}), \verb|lsucode5 = 011|\\

\bitfields{ 1/P, 2/01, 3/111, 2/coherency, 6/registerT, 2/11, 3/011, 1/--, 5/10100, 1/0, 6/registerZ }


\paragraph{Syntax} \texttt{asaddd}\emph{coherency} [\emph{registerZ}] = \emph{registerT}

\paragraph{Description} The effective address is the value of the \emph{registerZ}.
This instruction implements atomic add-store on double word. The double word at effective address is added with the \emph{registerT}. The effective address must be double word aligned (multiple of 8).


\paragraph{Modifier(s)}
\noindent\begin{itemize}
\item coherency (cf. coherency in section \ref{par:modifier-coherency} on page \pageref{par:modifier-coherency})
\end{itemize}

\paragraph{Execution} {Reservation table \textsf{LSU2\_MEMW\_AUXR} (Section~\ref{sec:reservation-LSU2-MEMW-AUXR})}
~
{\begin{lstlisting}
stage ID:
  new coherency = _ZX_2(coherency);
stage RR:
  new address = GPR[registerZ];
stage E1:
  new argument2 = GPR[registerT];
stage E3:
  MEM.atomic_add(address, 0xFF, coherency << 32, argument2.64[0], registerT);
\end{lstlisting}}
\newpage

\paragraph{Encoding} Format \textsf{LSU\_ASADDSB.O} (Section~\ref{sec:format-LSU-ASADDSB.O}), \verb|lsucode5 = 011|\\

\bitfields{ 1/P, 2/00, 2/-- --, 27/offset27, 1/1, 2/01, 3/111, 2/coherency, 6/registerT, 2/11, 3/011, 1/--, 5/10100, 1/0, 6/registerZ }


\paragraph{Syntax} \texttt{asaddd}\emph{coherency} \emph{offset27}[\emph{registerZ}] = \emph{registerT}

\paragraph{Description} The effective address is computed by adding the \emph{offset27} to the \emph{registerZ}.
This instruction implements atomic add-store on double word. The double word at effective address is added with the \emph{registerT}. The effective address must be double word aligned (multiple of 8).


\paragraph{Modifier(s)}
\noindent\begin{itemize}
\item coherency (cf. coherency in section \ref{par:modifier-coherency} on page \pageref{par:modifier-coherency})
\end{itemize}

\paragraph{Execution} {Reservation table \textsf{LSU2\_MEMW\_AUXR.X} (Section~\ref{sec:reservation-LSU2-MEMW-AUXR.X})}
~
{\begin{lstlisting}
stage ID:
  new offset = _SX_27(offset27);
  new coherency = _ZX_2(coherency);
stage RR:
  new base = GPR[registerZ];
  new address = base + offset;
stage E1:
  new argument2 = GPR[registerT];
stage E3:
  MEM.atomic_add(address, 0xFF, coherency << 32, argument2.64[0], registerT);
\end{lstlisting}}
\newpage

\paragraph{Encoding} Format \textsf{LSU\_ASADDSB.D} (Section~\ref{sec:format-LSU-ASADDSB.D}), \verb|lsucode5 = 011|\\

\bitfields{ 1/P, 2/00, 2/-- --, 27/extend27, 1/1, 2/00, 2/-- --, 27/offset27, 1/1, 2/01, 3/111, 2/coherency, 6/registerT, 2/11, 3/011, 1/--, 5/10100, 1/0, 6/registerZ }


\paragraph{Syntax} \texttt{asaddd}\emph{coherency} \emph{extend27\_\discretionary{}{}{}offset27}[\emph{registerZ}] = \emph{registerT}

\paragraph{Description} The effective address is computed by adding the \emph{extend27\_\discretionary{}{}{}offset27} to the \emph{registerZ}.
This instruction implements atomic add-store on double word. The double word at effective address is added with the \emph{registerT}. The effective address must be double word aligned (multiple of 8).


\paragraph{Modifier(s)}
\noindent\begin{itemize}
\item coherency (cf. coherency in section \ref{par:modifier-coherency} on page \pageref{par:modifier-coherency})
\end{itemize}

\paragraph{Execution} {Reservation table \textsf{LSU2\_MEMW\_AUXR.Y} (Section~\ref{sec:reservation-LSU2-MEMW-AUXR.Y})}
~
{\begin{lstlisting}
stage ID:
  new offset = _SX_54(extend27_offset27);
  new coherency = _ZX_2(coherency);
stage RR:
  new base = GPR[registerZ];
  new address = base + offset;
stage E1:
  new argument2 = GPR[registerT];
stage E3:
  MEM.atomic_add(address, 0xFF, coherency << 32, argument2.64[0], registerT);
\end{lstlisting}}
\newpage
\subsubsection[{\small ASADDH: Atomic Store Add Half Word}]{ASADDH\hfill Atomic Store Add Half Word}
\label{instr:ASADDH}\index{asaddh}
 
\paragraph{Encoding} Format \textsf{LSU\_ASADDSB} (Section~\ref{sec:format-LSU-ASADDSB}), \verb|lsucode5 = 001|\\

\bitfields{ 1/P, 2/01, 3/111, 2/coherency, 6/registerT, 2/11, 3/001, 1/--, 5/10100, 1/0, 6/registerZ }


\paragraph{Syntax} \texttt{asaddh}\emph{coherency} [\emph{registerZ}] = \emph{registerT}

\paragraph{Description} The effective address is the value of the \emph{registerZ}.
This instruction implements atomic add-store on half word. The half word at effective address is added with the \emph{registerT}. The effective address must be half word aligned (multiple of 2).


\paragraph{Modifier(s)}
\noindent\begin{itemize}
\item coherency (cf. coherency in section \ref{par:modifier-coherency} on page \pageref{par:modifier-coherency})
\end{itemize}

\paragraph{Execution} {Reservation table \textsf{LSU2\_MEMW\_AUXR} (Section~\ref{sec:reservation-LSU2-MEMW-AUXR})}
~
{\begin{lstlisting}
stage ID:
  new coherency = _ZX_2(coherency);
stage RR:
  new address = GPR[registerZ];
stage E1:
  new argument2 = GPR[registerT];
stage E3:
  MEM.atomic_add(address, 0x3, coherency << 32, argument2.16[0], registerT);
\end{lstlisting}}
\newpage

\paragraph{Encoding} Format \textsf{LSU\_ASADDSB.O} (Section~\ref{sec:format-LSU-ASADDSB.O}), \verb|lsucode5 = 001|\\

\bitfields{ 1/P, 2/00, 2/-- --, 27/offset27, 1/1, 2/01, 3/111, 2/coherency, 6/registerT, 2/11, 3/001, 1/--, 5/10100, 1/0, 6/registerZ }


\paragraph{Syntax} \texttt{asaddh}\emph{coherency} \emph{offset27}[\emph{registerZ}] = \emph{registerT}

\paragraph{Description} The effective address is computed by adding the \emph{offset27} to the \emph{registerZ}.
This instruction implements atomic add-store on half word. The half word at effective address is added with the \emph{registerT}. The effective address must be half word aligned (multiple of 2).


\paragraph{Modifier(s)}
\noindent\begin{itemize}
\item coherency (cf. coherency in section \ref{par:modifier-coherency} on page \pageref{par:modifier-coherency})
\end{itemize}

\paragraph{Execution} {Reservation table \textsf{LSU2\_MEMW\_AUXR.X} (Section~\ref{sec:reservation-LSU2-MEMW-AUXR.X})}
~
{\begin{lstlisting}
stage ID:
  new offset = _SX_27(offset27);
  new coherency = _ZX_2(coherency);
stage RR:
  new base = GPR[registerZ];
  new address = base + offset;
stage E1:
  new argument2 = GPR[registerT];
stage E3:
  MEM.atomic_add(address, 0x3, coherency << 32, argument2.16[0], registerT);
\end{lstlisting}}
\newpage

\paragraph{Encoding} Format \textsf{LSU\_ASADDSB.D} (Section~\ref{sec:format-LSU-ASADDSB.D}), \verb|lsucode5 = 001|\\

\bitfields{ 1/P, 2/00, 2/-- --, 27/extend27, 1/1, 2/00, 2/-- --, 27/offset27, 1/1, 2/01, 3/111, 2/coherency, 6/registerT, 2/11, 3/001, 1/--, 5/10100, 1/0, 6/registerZ }


\paragraph{Syntax} \texttt{asaddh}\emph{coherency} \emph{extend27\_\discretionary{}{}{}offset27}[\emph{registerZ}] = \emph{registerT}

\paragraph{Description} The effective address is computed by adding the \emph{extend27\_\discretionary{}{}{}offset27} to the \emph{registerZ}.
This instruction implements atomic add-store on half word. The half word at effective address is added with the \emph{registerT}. The effective address must be half word aligned (multiple of 2).


\paragraph{Modifier(s)}
\noindent\begin{itemize}
\item coherency (cf. coherency in section \ref{par:modifier-coherency} on page \pageref{par:modifier-coherency})
\end{itemize}

\paragraph{Execution} {Reservation table \textsf{LSU2\_MEMW\_AUXR.Y} (Section~\ref{sec:reservation-LSU2-MEMW-AUXR.Y})}
~
{\begin{lstlisting}
stage ID:
  new offset = _SX_54(extend27_offset27);
  new coherency = _ZX_2(coherency);
stage RR:
  new base = GPR[registerZ];
  new address = base + offset;
stage E1:
  new argument2 = GPR[registerT];
stage E3:
  MEM.atomic_add(address, 0x3, coherency << 32, argument2.16[0], registerT);
\end{lstlisting}}
\newpage
\subsubsection[{\small ASADDW: Atomic Store Add Word}]{ASADDW\hfill Atomic Store Add Word}
\label{instr:ASADDW}\index{asaddw}
 
\paragraph{Encoding} Format \textsf{LSU\_ASADDSB} (Section~\ref{sec:format-LSU-ASADDSB}), \verb|lsucode5 = 010|\\

\bitfields{ 1/P, 2/01, 3/111, 2/coherency, 6/registerT, 2/11, 3/010, 1/--, 5/10100, 1/0, 6/registerZ }


\paragraph{Syntax} \texttt{asaddw}\emph{coherency} [\emph{registerZ}] = \emph{registerT}

\paragraph{Description} The effective address is the value of the \emph{registerZ}.
This instruction implements atomic add-store on word. The word at effective address is added with the \emph{registerT}. The effective address must be word aligned (multiple of 4).


\paragraph{Modifier(s)}
\noindent\begin{itemize}
\item coherency (cf. coherency in section \ref{par:modifier-coherency} on page \pageref{par:modifier-coherency})
\end{itemize}

\paragraph{Execution} {Reservation table \textsf{LSU2\_MEMW\_AUXR} (Section~\ref{sec:reservation-LSU2-MEMW-AUXR})}
~
{\begin{lstlisting}
stage ID:
  new coherency = _ZX_2(coherency);
stage RR:
  new address = GPR[registerZ];
stage E1:
  new argument2 = GPR[registerT];
stage E3:
  MEM.atomic_add(address, 0xF, coherency << 32, argument2.32[0], registerT);
\end{lstlisting}}
\newpage

\paragraph{Encoding} Format \textsf{LSU\_ASADDSB.O} (Section~\ref{sec:format-LSU-ASADDSB.O}), \verb|lsucode5 = 010|\\

\bitfields{ 1/P, 2/00, 2/-- --, 27/offset27, 1/1, 2/01, 3/111, 2/coherency, 6/registerT, 2/11, 3/010, 1/--, 5/10100, 1/0, 6/registerZ }


\paragraph{Syntax} \texttt{asaddw}\emph{coherency} \emph{offset27}[\emph{registerZ}] = \emph{registerT}

\paragraph{Description} The effective address is computed by adding the \emph{offset27} to the \emph{registerZ}.
This instruction implements atomic add-store on word. The word at effective address is added with the \emph{registerT}. The effective address must be word aligned (multiple of 4).


\paragraph{Modifier(s)}
\noindent\begin{itemize}
\item coherency (cf. coherency in section \ref{par:modifier-coherency} on page \pageref{par:modifier-coherency})
\end{itemize}

\paragraph{Execution} {Reservation table \textsf{LSU2\_MEMW\_AUXR.X} (Section~\ref{sec:reservation-LSU2-MEMW-AUXR.X})}
~
{\begin{lstlisting}
stage ID:
  new offset = _SX_27(offset27);
  new coherency = _ZX_2(coherency);
stage RR:
  new base = GPR[registerZ];
  new address = base + offset;
stage E1:
  new argument2 = GPR[registerT];
stage E3:
  MEM.atomic_add(address, 0xF, coherency << 32, argument2.32[0], registerT);
\end{lstlisting}}
\newpage

\paragraph{Encoding} Format \textsf{LSU\_ASADDSB.D} (Section~\ref{sec:format-LSU-ASADDSB.D}), \verb|lsucode5 = 010|\\

\bitfields{ 1/P, 2/00, 2/-- --, 27/extend27, 1/1, 2/00, 2/-- --, 27/offset27, 1/1, 2/01, 3/111, 2/coherency, 6/registerT, 2/11, 3/010, 1/--, 5/10100, 1/0, 6/registerZ }


\paragraph{Syntax} \texttt{asaddw}\emph{coherency} \emph{extend27\_\discretionary{}{}{}offset27}[\emph{registerZ}] = \emph{registerT}

\paragraph{Description} The effective address is computed by adding the \emph{extend27\_\discretionary{}{}{}offset27} to the \emph{registerZ}.
This instruction implements atomic add-store on word. The word at effective address is added with the \emph{registerT}. The effective address must be word aligned (multiple of 4).


\paragraph{Modifier(s)}
\noindent\begin{itemize}
\item coherency (cf. coherency in section \ref{par:modifier-coherency} on page \pageref{par:modifier-coherency})
\end{itemize}

\paragraph{Execution} {Reservation table \textsf{LSU2\_MEMW\_AUXR.Y} (Section~\ref{sec:reservation-LSU2-MEMW-AUXR.Y})}
~
{\begin{lstlisting}
stage ID:
  new offset = _SX_54(extend27_offset27);
  new coherency = _ZX_2(coherency);
stage RR:
  new base = GPR[registerZ];
  new address = base + offset;
stage E1:
  new argument2 = GPR[registerT];
stage E3:
  MEM.atomic_add(address, 0xF, coherency << 32, argument2.32[0], registerT);
\end{lstlisting}}
\newpage
\subsubsection[{\small ASANDB: Atomic Store AND Byte}]{ASANDB\hfill Atomic Store AND Byte}
\label{instr:ASANDB}\index{asandb}
 
\paragraph{Encoding} Format \textsf{LSU\_ASANDSB} (Section~\ref{sec:format-LSU-ASANDSB}), \verb|lsucode5 = 000|\\

\bitfields{ 1/P, 2/01, 3/111, 2/coherency, 6/registerT, 2/11, 3/000, 1/--, 5/10101, 1/0, 6/registerZ }


\paragraph{Syntax} \texttt{asandb}\emph{coherency} [\emph{registerZ}] = \emph{registerT}

\paragraph{Description} The effective address is the value of the \emph{registerZ}.
This instruction implements atomic AND-store on byte. The byte at effective address is ANDed with the \emph{registerT}.


\paragraph{Modifier(s)}
\noindent\begin{itemize}
\item coherency (cf. coherency in section \ref{par:modifier-coherency} on page \pageref{par:modifier-coherency})
\end{itemize}

\paragraph{Execution} {Reservation table \textsf{LSU2\_MEMW\_AUXR} (Section~\ref{sec:reservation-LSU2-MEMW-AUXR})}
~
{\begin{lstlisting}
stage ID:
  new coherency = _ZX_2(coherency);
stage RR:
  new address = GPR[registerZ];
stage E1:
  new argument2 = GPR[registerT];
stage E3:
  MEM.atomic_and(address, 0x1, coherency << 32, argument2.8[0], registerT);
\end{lstlisting}}
\newpage

\paragraph{Encoding} Format \textsf{LSU\_ASANDSB.O} (Section~\ref{sec:format-LSU-ASANDSB.O}), \verb|lsucode5 = 000|\\

\bitfields{ 1/P, 2/00, 2/-- --, 27/offset27, 1/1, 2/01, 3/111, 2/coherency, 6/registerT, 2/11, 3/000, 1/--, 5/10101, 1/0, 6/registerZ }


\paragraph{Syntax} \texttt{asandb}\emph{coherency} \emph{offset27}[\emph{registerZ}] = \emph{registerT}

\paragraph{Description} The effective address is computed by adding the \emph{offset27} to the \emph{registerZ}.
This instruction implements atomic AND-store on byte. The byte at effective address is ANDed with the \emph{registerT}.


\paragraph{Modifier(s)}
\noindent\begin{itemize}
\item coherency (cf. coherency in section \ref{par:modifier-coherency} on page \pageref{par:modifier-coherency})
\end{itemize}

\paragraph{Execution} {Reservation table \textsf{LSU2\_MEMW\_AUXR.X} (Section~\ref{sec:reservation-LSU2-MEMW-AUXR.X})}
~
{\begin{lstlisting}
stage ID:
  new offset = _SX_27(offset27);
  new coherency = _ZX_2(coherency);
stage RR:
  new base = GPR[registerZ];
  new address = base + offset;
stage E1:
  new argument2 = GPR[registerT];
stage E3:
  MEM.atomic_and(address, 0x1, coherency << 32, argument2.8[0], registerT);
\end{lstlisting}}
\newpage

\paragraph{Encoding} Format \textsf{LSU\_ASANDSB.D} (Section~\ref{sec:format-LSU-ASANDSB.D}), \verb|lsucode5 = 000|\\

\bitfields{ 1/P, 2/00, 2/-- --, 27/extend27, 1/1, 2/00, 2/-- --, 27/offset27, 1/1, 2/01, 3/111, 2/coherency, 6/registerT, 2/11, 3/000, 1/--, 5/10101, 1/0, 6/registerZ }


\paragraph{Syntax} \texttt{asandb}\emph{coherency} \emph{extend27\_\discretionary{}{}{}offset27}[\emph{registerZ}] = \emph{registerT}

\paragraph{Description} The effective address is computed by adding the \emph{extend27\_\discretionary{}{}{}offset27} to the \emph{registerZ}.
This instruction implements atomic AND-store on byte. The byte at effective address is ANDed with the \emph{registerT}.


\paragraph{Modifier(s)}
\noindent\begin{itemize}
\item coherency (cf. coherency in section \ref{par:modifier-coherency} on page \pageref{par:modifier-coherency})
\end{itemize}

\paragraph{Execution} {Reservation table \textsf{LSU2\_MEMW\_AUXR.Y} (Section~\ref{sec:reservation-LSU2-MEMW-AUXR.Y})}
~
{\begin{lstlisting}
stage ID:
  new offset = _SX_54(extend27_offset27);
  new coherency = _ZX_2(coherency);
stage RR:
  new base = GPR[registerZ];
  new address = base + offset;
stage E1:
  new argument2 = GPR[registerT];
stage E3:
  MEM.atomic_and(address, 0x1, coherency << 32, argument2.8[0], registerT);
\end{lstlisting}}
\newpage
\subsubsection[{\small ASANDD: Atomic Store AND Double Word}]{ASANDD\hfill Atomic Store AND Double Word}
\label{instr:ASANDD}\index{asandd}
 
\paragraph{Encoding} Format \textsf{LSU\_ASANDSB} (Section~\ref{sec:format-LSU-ASANDSB}), \verb|lsucode5 = 011|\\

\bitfields{ 1/P, 2/01, 3/111, 2/coherency, 6/registerT, 2/11, 3/011, 1/--, 5/10101, 1/0, 6/registerZ }


\paragraph{Syntax} \texttt{asandd}\emph{coherency} [\emph{registerZ}] = \emph{registerT}

\paragraph{Description} The effective address is the value of the \emph{registerZ}.
This instruction implements atomic AND-store on double word. The double word at effective address is ANDed with the \emph{registerT}. The effective address must be double word aligned (multiple of 8).


\paragraph{Modifier(s)}
\noindent\begin{itemize}
\item coherency (cf. coherency in section \ref{par:modifier-coherency} on page \pageref{par:modifier-coherency})
\end{itemize}

\paragraph{Execution} {Reservation table \textsf{LSU2\_MEMW\_AUXR} (Section~\ref{sec:reservation-LSU2-MEMW-AUXR})}
~
{\begin{lstlisting}
stage ID:
  new coherency = _ZX_2(coherency);
stage RR:
  new address = GPR[registerZ];
stage E1:
  new argument2 = GPR[registerT];
stage E3:
  MEM.atomic_and(address, 0xFF, coherency << 32, argument2.64[0], registerT);
\end{lstlisting}}
\newpage

\paragraph{Encoding} Format \textsf{LSU\_ASANDSB.O} (Section~\ref{sec:format-LSU-ASANDSB.O}), \verb|lsucode5 = 011|\\

\bitfields{ 1/P, 2/00, 2/-- --, 27/offset27, 1/1, 2/01, 3/111, 2/coherency, 6/registerT, 2/11, 3/011, 1/--, 5/10101, 1/0, 6/registerZ }


\paragraph{Syntax} \texttt{asandd}\emph{coherency} \emph{offset27}[\emph{registerZ}] = \emph{registerT}

\paragraph{Description} The effective address is computed by adding the \emph{offset27} to the \emph{registerZ}.
This instruction implements atomic AND-store on double word. The double word at effective address is ANDed with the \emph{registerT}. The effective address must be double word aligned (multiple of 8).


\paragraph{Modifier(s)}
\noindent\begin{itemize}
\item coherency (cf. coherency in section \ref{par:modifier-coherency} on page \pageref{par:modifier-coherency})
\end{itemize}

\paragraph{Execution} {Reservation table \textsf{LSU2\_MEMW\_AUXR.X} (Section~\ref{sec:reservation-LSU2-MEMW-AUXR.X})}
~
{\begin{lstlisting}
stage ID:
  new offset = _SX_27(offset27);
  new coherency = _ZX_2(coherency);
stage RR:
  new base = GPR[registerZ];
  new address = base + offset;
stage E1:
  new argument2 = GPR[registerT];
stage E3:
  MEM.atomic_and(address, 0xFF, coherency << 32, argument2.64[0], registerT);
\end{lstlisting}}
\newpage

\paragraph{Encoding} Format \textsf{LSU\_ASANDSB.D} (Section~\ref{sec:format-LSU-ASANDSB.D}), \verb|lsucode5 = 011|\\

\bitfields{ 1/P, 2/00, 2/-- --, 27/extend27, 1/1, 2/00, 2/-- --, 27/offset27, 1/1, 2/01, 3/111, 2/coherency, 6/registerT, 2/11, 3/011, 1/--, 5/10101, 1/0, 6/registerZ }


\paragraph{Syntax} \texttt{asandd}\emph{coherency} \emph{extend27\_\discretionary{}{}{}offset27}[\emph{registerZ}] = \emph{registerT}

\paragraph{Description} The effective address is computed by adding the \emph{extend27\_\discretionary{}{}{}offset27} to the \emph{registerZ}.
This instruction implements atomic AND-store on double word. The double word at effective address is ANDed with the \emph{registerT}. The effective address must be double word aligned (multiple of 8).


\paragraph{Modifier(s)}
\noindent\begin{itemize}
\item coherency (cf. coherency in section \ref{par:modifier-coherency} on page \pageref{par:modifier-coherency})
\end{itemize}

\paragraph{Execution} {Reservation table \textsf{LSU2\_MEMW\_AUXR.Y} (Section~\ref{sec:reservation-LSU2-MEMW-AUXR.Y})}
~
{\begin{lstlisting}
stage ID:
  new offset = _SX_54(extend27_offset27);
  new coherency = _ZX_2(coherency);
stage RR:
  new base = GPR[registerZ];
  new address = base + offset;
stage E1:
  new argument2 = GPR[registerT];
stage E3:
  MEM.atomic_and(address, 0xFF, coherency << 32, argument2.64[0], registerT);
\end{lstlisting}}
\newpage
\subsubsection[{\small ASANDH: Atomic Store AND Half Word}]{ASANDH\hfill Atomic Store AND Half Word}
\label{instr:ASANDH}\index{asandh}
 
\paragraph{Encoding} Format \textsf{LSU\_ASANDSB} (Section~\ref{sec:format-LSU-ASANDSB}), \verb|lsucode5 = 001|\\

\bitfields{ 1/P, 2/01, 3/111, 2/coherency, 6/registerT, 2/11, 3/001, 1/--, 5/10101, 1/0, 6/registerZ }


\paragraph{Syntax} \texttt{asandh}\emph{coherency} [\emph{registerZ}] = \emph{registerT}

\paragraph{Description} The effective address is the value of the \emph{registerZ}.
This instruction implements atomic AND-store on half word. The half word at effective address is ANDed with the \emph{registerT}. The effective address must be half word aligned (multiple of 2).


\paragraph{Modifier(s)}
\noindent\begin{itemize}
\item coherency (cf. coherency in section \ref{par:modifier-coherency} on page \pageref{par:modifier-coherency})
\end{itemize}

\paragraph{Execution} {Reservation table \textsf{LSU2\_MEMW\_AUXR} (Section~\ref{sec:reservation-LSU2-MEMW-AUXR})}
~
{\begin{lstlisting}
stage ID:
  new coherency = _ZX_2(coherency);
stage RR:
  new address = GPR[registerZ];
stage E1:
  new argument2 = GPR[registerT];
stage E3:
  MEM.atomic_and(address, 0x3, coherency << 32, argument2.16[0], registerT);
\end{lstlisting}}
\newpage

\paragraph{Encoding} Format \textsf{LSU\_ASANDSB.O} (Section~\ref{sec:format-LSU-ASANDSB.O}), \verb|lsucode5 = 001|\\

\bitfields{ 1/P, 2/00, 2/-- --, 27/offset27, 1/1, 2/01, 3/111, 2/coherency, 6/registerT, 2/11, 3/001, 1/--, 5/10101, 1/0, 6/registerZ }


\paragraph{Syntax} \texttt{asandh}\emph{coherency} \emph{offset27}[\emph{registerZ}] = \emph{registerT}

\paragraph{Description} The effective address is computed by adding the \emph{offset27} to the \emph{registerZ}.
This instruction implements atomic AND-store on half word. The half word at effective address is ANDed with the \emph{registerT}. The effective address must be half word aligned (multiple of 2).


\paragraph{Modifier(s)}
\noindent\begin{itemize}
\item coherency (cf. coherency in section \ref{par:modifier-coherency} on page \pageref{par:modifier-coherency})
\end{itemize}

\paragraph{Execution} {Reservation table \textsf{LSU2\_MEMW\_AUXR.X} (Section~\ref{sec:reservation-LSU2-MEMW-AUXR.X})}
~
{\begin{lstlisting}
stage ID:
  new offset = _SX_27(offset27);
  new coherency = _ZX_2(coherency);
stage RR:
  new base = GPR[registerZ];
  new address = base + offset;
stage E1:
  new argument2 = GPR[registerT];
stage E3:
  MEM.atomic_and(address, 0x3, coherency << 32, argument2.16[0], registerT);
\end{lstlisting}}
\newpage

\paragraph{Encoding} Format \textsf{LSU\_ASANDSB.D} (Section~\ref{sec:format-LSU-ASANDSB.D}), \verb|lsucode5 = 001|\\

\bitfields{ 1/P, 2/00, 2/-- --, 27/extend27, 1/1, 2/00, 2/-- --, 27/offset27, 1/1, 2/01, 3/111, 2/coherency, 6/registerT, 2/11, 3/001, 1/--, 5/10101, 1/0, 6/registerZ }


\paragraph{Syntax} \texttt{asandh}\emph{coherency} \emph{extend27\_\discretionary{}{}{}offset27}[\emph{registerZ}] = \emph{registerT}

\paragraph{Description} The effective address is computed by adding the \emph{extend27\_\discretionary{}{}{}offset27} to the \emph{registerZ}.
This instruction implements atomic AND-store on half word. The half word at effective address is ANDed with the \emph{registerT}. The effective address must be half word aligned (multiple of 2).


\paragraph{Modifier(s)}
\noindent\begin{itemize}
\item coherency (cf. coherency in section \ref{par:modifier-coherency} on page \pageref{par:modifier-coherency})
\end{itemize}

\paragraph{Execution} {Reservation table \textsf{LSU2\_MEMW\_AUXR.Y} (Section~\ref{sec:reservation-LSU2-MEMW-AUXR.Y})}
~
{\begin{lstlisting}
stage ID:
  new offset = _SX_54(extend27_offset27);
  new coherency = _ZX_2(coherency);
stage RR:
  new base = GPR[registerZ];
  new address = base + offset;
stage E1:
  new argument2 = GPR[registerT];
stage E3:
  MEM.atomic_and(address, 0x3, coherency << 32, argument2.16[0], registerT);
\end{lstlisting}}
\newpage
\subsubsection[{\small ASANDW: Atomic Store AND Word}]{ASANDW\hfill Atomic Store AND Word}
\label{instr:ASANDW}\index{asandw}
 
\paragraph{Encoding} Format \textsf{LSU\_ASANDSB} (Section~\ref{sec:format-LSU-ASANDSB}), \verb|lsucode5 = 010|\\

\bitfields{ 1/P, 2/01, 3/111, 2/coherency, 6/registerT, 2/11, 3/010, 1/--, 5/10101, 1/0, 6/registerZ }


\paragraph{Syntax} \texttt{asandw}\emph{coherency} [\emph{registerZ}] = \emph{registerT}

\paragraph{Description} The effective address is the value of the \emph{registerZ}.
This instruction implements atomic AND-store on word. The word at effective address is ANDed with the \emph{registerT}. The effective address must be word aligned (multiple of 4).


\paragraph{Modifier(s)}
\noindent\begin{itemize}
\item coherency (cf. coherency in section \ref{par:modifier-coherency} on page \pageref{par:modifier-coherency})
\end{itemize}

\paragraph{Execution} {Reservation table \textsf{LSU2\_MEMW\_AUXR} (Section~\ref{sec:reservation-LSU2-MEMW-AUXR})}
~
{\begin{lstlisting}
stage ID:
  new coherency = _ZX_2(coherency);
stage RR:
  new address = GPR[registerZ];
stage E1:
  new argument2 = GPR[registerT];
stage E3:
  MEM.atomic_and(address, 0xF, coherency << 32, argument2.32[0], registerT);
\end{lstlisting}}
\newpage

\paragraph{Encoding} Format \textsf{LSU\_ASANDSB.O} (Section~\ref{sec:format-LSU-ASANDSB.O}), \verb|lsucode5 = 010|\\

\bitfields{ 1/P, 2/00, 2/-- --, 27/offset27, 1/1, 2/01, 3/111, 2/coherency, 6/registerT, 2/11, 3/010, 1/--, 5/10101, 1/0, 6/registerZ }


\paragraph{Syntax} \texttt{asandw}\emph{coherency} \emph{offset27}[\emph{registerZ}] = \emph{registerT}

\paragraph{Description} The effective address is computed by adding the \emph{offset27} to the \emph{registerZ}.
This instruction implements atomic AND-store on word. The word at effective address is ANDed with the \emph{registerT}. The effective address must be word aligned (multiple of 4).


\paragraph{Modifier(s)}
\noindent\begin{itemize}
\item coherency (cf. coherency in section \ref{par:modifier-coherency} on page \pageref{par:modifier-coherency})
\end{itemize}

\paragraph{Execution} {Reservation table \textsf{LSU2\_MEMW\_AUXR.X} (Section~\ref{sec:reservation-LSU2-MEMW-AUXR.X})}
~
{\begin{lstlisting}
stage ID:
  new offset = _SX_27(offset27);
  new coherency = _ZX_2(coherency);
stage RR:
  new base = GPR[registerZ];
  new address = base + offset;
stage E1:
  new argument2 = GPR[registerT];
stage E3:
  MEM.atomic_and(address, 0xF, coherency << 32, argument2.32[0], registerT);
\end{lstlisting}}
\newpage

\paragraph{Encoding} Format \textsf{LSU\_ASANDSB.D} (Section~\ref{sec:format-LSU-ASANDSB.D}), \verb|lsucode5 = 010|\\

\bitfields{ 1/P, 2/00, 2/-- --, 27/extend27, 1/1, 2/00, 2/-- --, 27/offset27, 1/1, 2/01, 3/111, 2/coherency, 6/registerT, 2/11, 3/010, 1/--, 5/10101, 1/0, 6/registerZ }


\paragraph{Syntax} \texttt{asandw}\emph{coherency} \emph{extend27\_\discretionary{}{}{}offset27}[\emph{registerZ}] = \emph{registerT}

\paragraph{Description} The effective address is computed by adding the \emph{extend27\_\discretionary{}{}{}offset27} to the \emph{registerZ}.
This instruction implements atomic AND-store on word. The word at effective address is ANDed with the \emph{registerT}. The effective address must be word aligned (multiple of 4).


\paragraph{Modifier(s)}
\noindent\begin{itemize}
\item coherency (cf. coherency in section \ref{par:modifier-coherency} on page \pageref{par:modifier-coherency})
\end{itemize}

\paragraph{Execution} {Reservation table \textsf{LSU2\_MEMW\_AUXR.Y} (Section~\ref{sec:reservation-LSU2-MEMW-AUXR.Y})}
~
{\begin{lstlisting}
stage ID:
  new offset = _SX_54(extend27_offset27);
  new coherency = _ZX_2(coherency);
stage RR:
  new base = GPR[registerZ];
  new address = base + offset;
stage E1:
  new argument2 = GPR[registerT];
stage E3:
  MEM.atomic_and(address, 0xF, coherency << 32, argument2.32[0], registerT);
\end{lstlisting}}
\newpage
\subsubsection[{\small ASB: Atomic Store Byte}]{ASB\hfill Atomic Store Byte}
\label{instr:ASB}\index{asb}
 
\paragraph{Encoding} Format \textsf{LSU\_ASSB} (Section~\ref{sec:format-LSU-ASSB}), \verb|lsucode5 = 000|\\

\bitfields{ 1/P, 2/01, 3/111, 2/coherency, 6/registerT, 2/11, 3/000, 1/--, 5/10000, 1/0, 6/registerZ }


\paragraph{Syntax} \texttt{asb}\emph{coherency} [\emph{registerZ}] = \emph{registerT}

\paragraph{Description} The effective address is the value of the \emph{registerZ}.
The \emph{registerT} is atomically stored to the memory byte starting at the effective address. This instruction is provided to directly operate on the last level of the memory hierarchy, which is typically a processor DDR or a cluster TCM, irrespective of the MMU DCP (Data Cache Policy) settings.


\paragraph{Modifier(s)}
\noindent\begin{itemize}
\item coherency (cf. coherency in section \ref{par:modifier-coherency} on page \pageref{par:modifier-coherency})
\end{itemize}

\paragraph{Execution} {Reservation table \textsf{LSU\_MEMW\_AUXR} (Section~\ref{sec:reservation-LSU-MEMW-AUXR})}
~
{\begin{lstlisting}
stage ID:
  new coherency = _ZX_2(coherency);
stage RR:
  new address = GPR[registerZ];
stage E1:
  new argument2 = GPR[registerT];
stage E3:
  MEM.atomic_store(address, 0x1, coherency << 32, argument2, registerT);
\end{lstlisting}}
\newpage

\paragraph{Encoding} Format \textsf{LSU\_ASSB.O} (Section~\ref{sec:format-LSU-ASSB.O}), \verb|lsucode5 = 000|\\

\bitfields{ 1/P, 2/00, 2/-- --, 27/offset27, 1/1, 2/01, 3/111, 2/coherency, 6/registerT, 2/11, 3/000, 1/--, 5/10000, 1/0, 6/registerZ }


\paragraph{Syntax} \texttt{asb}\emph{coherency} \emph{offset27}[\emph{registerZ}] = \emph{registerT}

\paragraph{Description} The effective address is computed by adding the \emph{offset27} to the \emph{registerZ}.
The \emph{registerT} is atomically stored to the memory byte starting at the effective address. This instruction is provided to directly operate on the last level of the memory hierarchy, which is typically a processor DDR or a cluster TCM, irrespective of the MMU DCP (Data Cache Policy) settings.


\paragraph{Modifier(s)}
\noindent\begin{itemize}
\item coherency (cf. coherency in section \ref{par:modifier-coherency} on page \pageref{par:modifier-coherency})
\end{itemize}

\paragraph{Execution} {Reservation table \textsf{LSU\_MEMW\_AUXR.X} (Section~\ref{sec:reservation-LSU-MEMW-AUXR.X})}
~
{\begin{lstlisting}
stage ID:
  new offset = _SX_27(offset27);
  new coherency = _ZX_2(coherency);
stage RR:
  new base = GPR[registerZ];
  new address = base + offset;
stage E1:
  new argument2 = GPR[registerT];
stage E3:
  MEM.atomic_store(address, 0x1, coherency << 32, argument2, registerT);
\end{lstlisting}}
\newpage

\paragraph{Encoding} Format \textsf{LSU\_ASSB.D} (Section~\ref{sec:format-LSU-ASSB.D}), \verb|lsucode5 = 000|\\

\bitfields{ 1/P, 2/00, 2/-- --, 27/extend27, 1/1, 2/00, 2/-- --, 27/offset27, 1/1, 2/01, 3/111, 2/coherency, 6/registerT, 2/11, 3/000, 1/--, 5/10000, 1/0, 6/registerZ }


\paragraph{Syntax} \texttt{asb}\emph{coherency} \emph{extend27\_\discretionary{}{}{}offset27}[\emph{registerZ}] = \emph{registerT}

\paragraph{Description} The effective address is computed by adding the \emph{extend27\_\discretionary{}{}{}offset27} to the \emph{registerZ}.
The \emph{registerT} is atomically stored to the memory byte starting at the effective address. This instruction is provided to directly operate on the last level of the memory hierarchy, which is typically a processor DDR or a cluster TCM, irrespective of the MMU DCP (Data Cache Policy) settings.


\paragraph{Modifier(s)}
\noindent\begin{itemize}
\item coherency (cf. coherency in section \ref{par:modifier-coherency} on page \pageref{par:modifier-coherency})
\end{itemize}

\paragraph{Execution} {Reservation table \textsf{LSU\_MEMW\_AUXR.Y} (Section~\ref{sec:reservation-LSU-MEMW-AUXR.Y})}
~
{\begin{lstlisting}
stage ID:
  new offset = _SX_54(extend27_offset27);
  new coherency = _ZX_2(coherency);
stage RR:
  new base = GPR[registerZ];
  new address = base + offset;
stage E1:
  new argument2 = GPR[registerT];
stage E3:
  MEM.atomic_store(address, 0x1, coherency << 32, argument2, registerT);
\end{lstlisting}}
\newpage
\subsubsection[{\small ASD: Atomic Store Double Word}]{ASD\hfill Atomic Store Double Word}
\label{instr:ASD}\index{asd}
 
\paragraph{Encoding} Format \textsf{LSU\_ASSB} (Section~\ref{sec:format-LSU-ASSB}), \verb|lsucode5 = 011|\\

\bitfields{ 1/P, 2/01, 3/111, 2/coherency, 6/registerT, 2/11, 3/011, 1/--, 5/10000, 1/0, 6/registerZ }


\paragraph{Syntax} \texttt{asd}\emph{coherency} [\emph{registerZ}] = \emph{registerT}

\paragraph{Description} The effective address is the value of the \emph{registerZ}.
The \emph{registerT} is atomically stored to the memory double word starting at the effective address. This instruction is provided to directly operate on the last level of the memory hierarchy, which is typically a processor DDR or a cluster TCM, irrespective of the MMU DCP (Data Cache Policy) settings. The effective address must be double word aligned (multiple of 8).


\paragraph{Modifier(s)}
\noindent\begin{itemize}
\item coherency (cf. coherency in section \ref{par:modifier-coherency} on page \pageref{par:modifier-coherency})
\end{itemize}

\paragraph{Execution} {Reservation table \textsf{LSU\_MEMW\_AUXR} (Section~\ref{sec:reservation-LSU-MEMW-AUXR})}
~
{\begin{lstlisting}
stage ID:
  new coherency = _ZX_2(coherency);
stage RR:
  new address = GPR[registerZ];
stage E1:
  new argument2 = GPR[registerT];
stage E3:
  MEM.atomic_store(address, 0xFF, coherency << 32, argument2, registerT);
\end{lstlisting}}
\newpage

\paragraph{Encoding} Format \textsf{LSU\_ASSB.O} (Section~\ref{sec:format-LSU-ASSB.O}), \verb|lsucode5 = 011|\\

\bitfields{ 1/P, 2/00, 2/-- --, 27/offset27, 1/1, 2/01, 3/111, 2/coherency, 6/registerT, 2/11, 3/011, 1/--, 5/10000, 1/0, 6/registerZ }


\paragraph{Syntax} \texttt{asd}\emph{coherency} \emph{offset27}[\emph{registerZ}] = \emph{registerT}

\paragraph{Description} The effective address is computed by adding the \emph{offset27} to the \emph{registerZ}.
The \emph{registerT} is atomically stored to the memory double word starting at the effective address. This instruction is provided to directly operate on the last level of the memory hierarchy, which is typically a processor DDR or a cluster TCM, irrespective of the MMU DCP (Data Cache Policy) settings. The effective address must be double word aligned (multiple of 8).


\paragraph{Modifier(s)}
\noindent\begin{itemize}
\item coherency (cf. coherency in section \ref{par:modifier-coherency} on page \pageref{par:modifier-coherency})
\end{itemize}

\paragraph{Execution} {Reservation table \textsf{LSU\_MEMW\_AUXR.X} (Section~\ref{sec:reservation-LSU-MEMW-AUXR.X})}
~
{\begin{lstlisting}
stage ID:
  new offset = _SX_27(offset27);
  new coherency = _ZX_2(coherency);
stage RR:
  new base = GPR[registerZ];
  new address = base + offset;
stage E1:
  new argument2 = GPR[registerT];
stage E3:
  MEM.atomic_store(address, 0xFF, coherency << 32, argument2, registerT);
\end{lstlisting}}
\newpage

\paragraph{Encoding} Format \textsf{LSU\_ASSB.D} (Section~\ref{sec:format-LSU-ASSB.D}), \verb|lsucode5 = 011|\\

\bitfields{ 1/P, 2/00, 2/-- --, 27/extend27, 1/1, 2/00, 2/-- --, 27/offset27, 1/1, 2/01, 3/111, 2/coherency, 6/registerT, 2/11, 3/011, 1/--, 5/10000, 1/0, 6/registerZ }


\paragraph{Syntax} \texttt{asd}\emph{coherency} \emph{extend27\_\discretionary{}{}{}offset27}[\emph{registerZ}] = \emph{registerT}

\paragraph{Description} The effective address is computed by adding the \emph{extend27\_\discretionary{}{}{}offset27} to the \emph{registerZ}.
The \emph{registerT} is atomically stored to the memory double word starting at the effective address. This instruction is provided to directly operate on the last level of the memory hierarchy, which is typically a processor DDR or a cluster TCM, irrespective of the MMU DCP (Data Cache Policy) settings. The effective address must be double word aligned (multiple of 8).


\paragraph{Modifier(s)}
\noindent\begin{itemize}
\item coherency (cf. coherency in section \ref{par:modifier-coherency} on page \pageref{par:modifier-coherency})
\end{itemize}

\paragraph{Execution} {Reservation table \textsf{LSU\_MEMW\_AUXR.Y} (Section~\ref{sec:reservation-LSU-MEMW-AUXR.Y})}
~
{\begin{lstlisting}
stage ID:
  new offset = _SX_54(extend27_offset27);
  new coherency = _ZX_2(coherency);
stage RR:
  new base = GPR[registerZ];
  new address = base + offset;
stage E1:
  new argument2 = GPR[registerT];
stage E3:
  MEM.atomic_store(address, 0xFF, coherency << 32, argument2, registerT);
\end{lstlisting}}
\newpage
\subsubsection[{\small ASDUSB: Atomic Store Decrement Unsigned Saturating Byte}]{ASDUSB\hfill Atomic Store Decrement Unsigned Saturating Byte}
\label{instr:ASDUSB}\index{asdusb}
 
\paragraph{Encoding} Format \textsf{LSU\_ASDUSSB} (Section~\ref{sec:format-LSU-ASDUSSB}), \verb|lsucode5 = 000|\\

\bitfields{ 1/P, 2/01, 3/111, 2/coherency, 6/registerT, 2/11, 3/000, 1/--, 5/11110, 1/0, 6/registerZ }


\paragraph{Syntax} \texttt{asdusb}\emph{coherency} [\emph{registerZ}] = \emph{registerT}

\paragraph{Description} The effective address is the value of the \emph{registerZ}.
This instruction implements atomic DUS-store on signed byte (Decrement Unsigned Saturating). The byte at effective address is DUSed with the \emph{registerT}.


\paragraph{Modifier(s)}
\noindent\begin{itemize}
\item coherency (cf. coherency in section \ref{par:modifier-coherency} on page \pageref{par:modifier-coherency})
\end{itemize}

\paragraph{Execution} {Reservation table \textsf{LSU2\_MEMW\_AUXR} (Section~\ref{sec:reservation-LSU2-MEMW-AUXR})}
~
{\begin{lstlisting}
stage ID:
  new coherency = _ZX_2(coherency);
stage RR:
  new address = GPR[registerZ];
stage E1:
  new argument2 = GPR[registerT];
stage E3:
  MEM.atomic_dus(address, 0x1, coherency << 32, argument2.8[0], registerT);
\end{lstlisting}}
\newpage

\paragraph{Encoding} Format \textsf{LSU\_ASDUSSB.O} (Section~\ref{sec:format-LSU-ASDUSSB.O}), \verb|lsucode5 = 000|\\

\bitfields{ 1/P, 2/00, 2/-- --, 27/offset27, 1/1, 2/01, 3/111, 2/coherency, 6/registerT, 2/11, 3/000, 1/--, 5/11110, 1/0, 6/registerZ }


\paragraph{Syntax} \texttt{asdusb}\emph{coherency} \emph{offset27}[\emph{registerZ}] = \emph{registerT}

\paragraph{Description} The effective address is computed by adding the \emph{offset27} to the \emph{registerZ}.
This instruction implements atomic DUS-store on signed byte (Decrement Unsigned Saturating). The byte at effective address is DUSed with the \emph{registerT}.


\paragraph{Modifier(s)}
\noindent\begin{itemize}
\item coherency (cf. coherency in section \ref{par:modifier-coherency} on page \pageref{par:modifier-coherency})
\end{itemize}

\paragraph{Execution} {Reservation table \textsf{LSU2\_MEMW\_AUXR.X} (Section~\ref{sec:reservation-LSU2-MEMW-AUXR.X})}
~
{\begin{lstlisting}
stage ID:
  new offset = _SX_27(offset27);
  new coherency = _ZX_2(coherency);
stage RR:
  new base = GPR[registerZ];
  new address = base + offset;
stage E1:
  new argument2 = GPR[registerT];
stage E3:
  MEM.atomic_dus(address, 0x1, coherency << 32, argument2.8[0], registerT);
\end{lstlisting}}
\newpage

\paragraph{Encoding} Format \textsf{LSU\_ASDUSSB.D} (Section~\ref{sec:format-LSU-ASDUSSB.D}), \verb|lsucode5 = 000|\\

\bitfields{ 1/P, 2/00, 2/-- --, 27/extend27, 1/1, 2/00, 2/-- --, 27/offset27, 1/1, 2/01, 3/111, 2/coherency, 6/registerT, 2/11, 3/000, 1/--, 5/11110, 1/0, 6/registerZ }


\paragraph{Syntax} \texttt{asdusb}\emph{coherency} \emph{extend27\_\discretionary{}{}{}offset27}[\emph{registerZ}] = \emph{registerT}

\paragraph{Description} The effective address is computed by adding the \emph{extend27\_\discretionary{}{}{}offset27} to the \emph{registerZ}.
This instruction implements atomic DUS-store on signed byte (Decrement Unsigned Saturating). The byte at effective address is DUSed with the \emph{registerT}.


\paragraph{Modifier(s)}
\noindent\begin{itemize}
\item coherency (cf. coherency in section \ref{par:modifier-coherency} on page \pageref{par:modifier-coherency})
\end{itemize}

\paragraph{Execution} {Reservation table \textsf{LSU2\_MEMW\_AUXR.Y} (Section~\ref{sec:reservation-LSU2-MEMW-AUXR.Y})}
~
{\begin{lstlisting}
stage ID:
  new offset = _SX_54(extend27_offset27);
  new coherency = _ZX_2(coherency);
stage RR:
  new base = GPR[registerZ];
  new address = base + offset;
stage E1:
  new argument2 = GPR[registerT];
stage E3:
  MEM.atomic_dus(address, 0x1, coherency << 32, argument2.8[0], registerT);
\end{lstlisting}}
\newpage
\subsubsection[{\small ASDUSD: Atomic Store DUS Double Word}]{ASDUSD\hfill Atomic Store DUS Double Word}
\label{instr:ASDUSD}\index{asdusd}
 
\paragraph{Encoding} Format \textsf{LSU\_ASDUSSB} (Section~\ref{sec:format-LSU-ASDUSSB}), \verb|lsucode5 = 011|\\

\bitfields{ 1/P, 2/01, 3/111, 2/coherency, 6/registerT, 2/11, 3/011, 1/--, 5/11110, 1/0, 6/registerZ }


\paragraph{Syntax} \texttt{asdusd}\emph{coherency} [\emph{registerZ}] = \emph{registerT}

\paragraph{Description} The effective address is the value of the \emph{registerZ}.
This instruction implements atomic DUS-store on signed double word (Decrement Unsigned Saturating). The double word at effective address is DUSed with the \emph{registerT}. The effective address must be double word aligned (multiple of 8).


\paragraph{Modifier(s)}
\noindent\begin{itemize}
\item coherency (cf. coherency in section \ref{par:modifier-coherency} on page \pageref{par:modifier-coherency})
\end{itemize}

\paragraph{Execution} {Reservation table \textsf{LSU2\_MEMW\_AUXR} (Section~\ref{sec:reservation-LSU2-MEMW-AUXR})}
~
{\begin{lstlisting}
stage ID:
  new coherency = _ZX_2(coherency);
stage RR:
  new address = GPR[registerZ];
stage E1:
  new argument2 = GPR[registerT];
stage E3:
  MEM.atomic_dus(address, 0xFF, coherency << 32, argument2.64[0], registerT);
\end{lstlisting}}
\newpage

\paragraph{Encoding} Format \textsf{LSU\_ASDUSSB.O} (Section~\ref{sec:format-LSU-ASDUSSB.O}), \verb|lsucode5 = 011|\\

\bitfields{ 1/P, 2/00, 2/-- --, 27/offset27, 1/1, 2/01, 3/111, 2/coherency, 6/registerT, 2/11, 3/011, 1/--, 5/11110, 1/0, 6/registerZ }


\paragraph{Syntax} \texttt{asdusd}\emph{coherency} \emph{offset27}[\emph{registerZ}] = \emph{registerT}

\paragraph{Description} The effective address is computed by adding the \emph{offset27} to the \emph{registerZ}.
This instruction implements atomic DUS-store on signed double word (Decrement Unsigned Saturating). The double word at effective address is DUSed with the \emph{registerT}. The effective address must be double word aligned (multiple of 8).


\paragraph{Modifier(s)}
\noindent\begin{itemize}
\item coherency (cf. coherency in section \ref{par:modifier-coherency} on page \pageref{par:modifier-coherency})
\end{itemize}

\paragraph{Execution} {Reservation table \textsf{LSU2\_MEMW\_AUXR.X} (Section~\ref{sec:reservation-LSU2-MEMW-AUXR.X})}
~
{\begin{lstlisting}
stage ID:
  new offset = _SX_27(offset27);
  new coherency = _ZX_2(coherency);
stage RR:
  new base = GPR[registerZ];
  new address = base + offset;
stage E1:
  new argument2 = GPR[registerT];
stage E3:
  MEM.atomic_dus(address, 0xFF, coherency << 32, argument2.64[0], registerT);
\end{lstlisting}}
\newpage

\paragraph{Encoding} Format \textsf{LSU\_ASDUSSB.D} (Section~\ref{sec:format-LSU-ASDUSSB.D}), \verb|lsucode5 = 011|\\

\bitfields{ 1/P, 2/00, 2/-- --, 27/extend27, 1/1, 2/00, 2/-- --, 27/offset27, 1/1, 2/01, 3/111, 2/coherency, 6/registerT, 2/11, 3/011, 1/--, 5/11110, 1/0, 6/registerZ }


\paragraph{Syntax} \texttt{asdusd}\emph{coherency} \emph{extend27\_\discretionary{}{}{}offset27}[\emph{registerZ}] = \emph{registerT}

\paragraph{Description} The effective address is computed by adding the \emph{extend27\_\discretionary{}{}{}offset27} to the \emph{registerZ}.
This instruction implements atomic DUS-store on signed double word (Decrement Unsigned Saturating). The double word at effective address is DUSed with the \emph{registerT}. The effective address must be double word aligned (multiple of 8).


\paragraph{Modifier(s)}
\noindent\begin{itemize}
\item coherency (cf. coherency in section \ref{par:modifier-coherency} on page \pageref{par:modifier-coherency})
\end{itemize}

\paragraph{Execution} {Reservation table \textsf{LSU2\_MEMW\_AUXR.Y} (Section~\ref{sec:reservation-LSU2-MEMW-AUXR.Y})}
~
{\begin{lstlisting}
stage ID:
  new offset = _SX_54(extend27_offset27);
  new coherency = _ZX_2(coherency);
stage RR:
  new base = GPR[registerZ];
  new address = base + offset;
stage E1:
  new argument2 = GPR[registerT];
stage E3:
  MEM.atomic_dus(address, 0xFF, coherency << 32, argument2.64[0], registerT);
\end{lstlisting}}
\newpage
\subsubsection[{\small ASDUSH: Atomic Store DUS Half Word}]{ASDUSH\hfill Atomic Store DUS Half Word}
\label{instr:ASDUSH}\index{asdush}
 
\paragraph{Encoding} Format \textsf{LSU\_ASDUSSB} (Section~\ref{sec:format-LSU-ASDUSSB}), \verb|lsucode5 = 001|\\

\bitfields{ 1/P, 2/01, 3/111, 2/coherency, 6/registerT, 2/11, 3/001, 1/--, 5/11110, 1/0, 6/registerZ }


\paragraph{Syntax} \texttt{asdush}\emph{coherency} [\emph{registerZ}] = \emph{registerT}

\paragraph{Description} The effective address is the value of the \emph{registerZ}.
This instruction implements atomic DUS-store on signed half word (Decrement Unsigned Saturating). The half word at effective address is DUSed with the \emph{registerT}. The effective address must be half word aligned (multiple of 2).


\paragraph{Modifier(s)}
\noindent\begin{itemize}
\item coherency (cf. coherency in section \ref{par:modifier-coherency} on page \pageref{par:modifier-coherency})
\end{itemize}

\paragraph{Execution} {Reservation table \textsf{LSU2\_MEMW\_AUXR} (Section~\ref{sec:reservation-LSU2-MEMW-AUXR})}
~
{\begin{lstlisting}
stage ID:
  new coherency = _ZX_2(coherency);
stage RR:
  new address = GPR[registerZ];
stage E1:
  new argument2 = GPR[registerT];
stage E3:
  MEM.atomic_dus(address, 0x3, coherency << 32, argument2.16[0], registerT);
\end{lstlisting}}
\newpage

\paragraph{Encoding} Format \textsf{LSU\_ASDUSSB.O} (Section~\ref{sec:format-LSU-ASDUSSB.O}), \verb|lsucode5 = 001|\\

\bitfields{ 1/P, 2/00, 2/-- --, 27/offset27, 1/1, 2/01, 3/111, 2/coherency, 6/registerT, 2/11, 3/001, 1/--, 5/11110, 1/0, 6/registerZ }


\paragraph{Syntax} \texttt{asdush}\emph{coherency} \emph{offset27}[\emph{registerZ}] = \emph{registerT}

\paragraph{Description} The effective address is computed by adding the \emph{offset27} to the \emph{registerZ}.
This instruction implements atomic DUS-store on signed half word (Decrement Unsigned Saturating). The half word at effective address is DUSed with the \emph{registerT}. The effective address must be half word aligned (multiple of 2).


\paragraph{Modifier(s)}
\noindent\begin{itemize}
\item coherency (cf. coherency in section \ref{par:modifier-coherency} on page \pageref{par:modifier-coherency})
\end{itemize}

\paragraph{Execution} {Reservation table \textsf{LSU2\_MEMW\_AUXR.X} (Section~\ref{sec:reservation-LSU2-MEMW-AUXR.X})}
~
{\begin{lstlisting}
stage ID:
  new offset = _SX_27(offset27);
  new coherency = _ZX_2(coherency);
stage RR:
  new base = GPR[registerZ];
  new address = base + offset;
stage E1:
  new argument2 = GPR[registerT];
stage E3:
  MEM.atomic_dus(address, 0x3, coherency << 32, argument2.16[0], registerT);
\end{lstlisting}}
\newpage

\paragraph{Encoding} Format \textsf{LSU\_ASDUSSB.D} (Section~\ref{sec:format-LSU-ASDUSSB.D}), \verb|lsucode5 = 001|\\

\bitfields{ 1/P, 2/00, 2/-- --, 27/extend27, 1/1, 2/00, 2/-- --, 27/offset27, 1/1, 2/01, 3/111, 2/coherency, 6/registerT, 2/11, 3/001, 1/--, 5/11110, 1/0, 6/registerZ }


\paragraph{Syntax} \texttt{asdush}\emph{coherency} \emph{extend27\_\discretionary{}{}{}offset27}[\emph{registerZ}] = \emph{registerT}

\paragraph{Description} The effective address is computed by adding the \emph{extend27\_\discretionary{}{}{}offset27} to the \emph{registerZ}.
This instruction implements atomic DUS-store on signed half word (Decrement Unsigned Saturating). The half word at effective address is DUSed with the \emph{registerT}. The effective address must be half word aligned (multiple of 2).


\paragraph{Modifier(s)}
\noindent\begin{itemize}
\item coherency (cf. coherency in section \ref{par:modifier-coherency} on page \pageref{par:modifier-coherency})
\end{itemize}

\paragraph{Execution} {Reservation table \textsf{LSU2\_MEMW\_AUXR.Y} (Section~\ref{sec:reservation-LSU2-MEMW-AUXR.Y})}
~
{\begin{lstlisting}
stage ID:
  new offset = _SX_54(extend27_offset27);
  new coherency = _ZX_2(coherency);
stage RR:
  new base = GPR[registerZ];
  new address = base + offset;
stage E1:
  new argument2 = GPR[registerT];
stage E3:
  MEM.atomic_dus(address, 0x3, coherency << 32, argument2.16[0], registerT);
\end{lstlisting}}
\newpage
\subsubsection[{\small ASDUSW: Atomic Store DUS Word}]{ASDUSW\hfill Atomic Store DUS Word}
\label{instr:ASDUSW}\index{asdusw}
 
\paragraph{Encoding} Format \textsf{LSU\_ASDUSSB} (Section~\ref{sec:format-LSU-ASDUSSB}), \verb|lsucode5 = 010|\\

\bitfields{ 1/P, 2/01, 3/111, 2/coherency, 6/registerT, 2/11, 3/010, 1/--, 5/11110, 1/0, 6/registerZ }


\paragraph{Syntax} \texttt{asdusw}\emph{coherency} [\emph{registerZ}] = \emph{registerT}

\paragraph{Description} The effective address is the value of the \emph{registerZ}.
This instruction implements atomic DUS-store on signed word (Decrement Unsigned Saturating). The word at effective address is DUSed with the \emph{registerT}. The effective address must be word aligned (multiple of 4).


\paragraph{Modifier(s)}
\noindent\begin{itemize}
\item coherency (cf. coherency in section \ref{par:modifier-coherency} on page \pageref{par:modifier-coherency})
\end{itemize}

\paragraph{Execution} {Reservation table \textsf{LSU2\_MEMW\_AUXR} (Section~\ref{sec:reservation-LSU2-MEMW-AUXR})}
~
{\begin{lstlisting}
stage ID:
  new coherency = _ZX_2(coherency);
stage RR:
  new address = GPR[registerZ];
stage E1:
  new argument2 = GPR[registerT];
stage E3:
  MEM.atomic_dus(address, 0xF, coherency << 32, argument2.32[0], registerT);
\end{lstlisting}}
\newpage

\paragraph{Encoding} Format \textsf{LSU\_ASDUSSB.O} (Section~\ref{sec:format-LSU-ASDUSSB.O}), \verb|lsucode5 = 010|\\

\bitfields{ 1/P, 2/00, 2/-- --, 27/offset27, 1/1, 2/01, 3/111, 2/coherency, 6/registerT, 2/11, 3/010, 1/--, 5/11110, 1/0, 6/registerZ }


\paragraph{Syntax} \texttt{asdusw}\emph{coherency} \emph{offset27}[\emph{registerZ}] = \emph{registerT}

\paragraph{Description} The effective address is computed by adding the \emph{offset27} to the \emph{registerZ}.
This instruction implements atomic DUS-store on signed word (Decrement Unsigned Saturating). The word at effective address is DUSed with the \emph{registerT}. The effective address must be word aligned (multiple of 4).


\paragraph{Modifier(s)}
\noindent\begin{itemize}
\item coherency (cf. coherency in section \ref{par:modifier-coherency} on page \pageref{par:modifier-coherency})
\end{itemize}

\paragraph{Execution} {Reservation table \textsf{LSU2\_MEMW\_AUXR.X} (Section~\ref{sec:reservation-LSU2-MEMW-AUXR.X})}
~
{\begin{lstlisting}
stage ID:
  new offset = _SX_27(offset27);
  new coherency = _ZX_2(coherency);
stage RR:
  new base = GPR[registerZ];
  new address = base + offset;
stage E1:
  new argument2 = GPR[registerT];
stage E3:
  MEM.atomic_dus(address, 0xF, coherency << 32, argument2.32[0], registerT);
\end{lstlisting}}
\newpage

\paragraph{Encoding} Format \textsf{LSU\_ASDUSSB.D} (Section~\ref{sec:format-LSU-ASDUSSB.D}), \verb|lsucode5 = 010|\\

\bitfields{ 1/P, 2/00, 2/-- --, 27/extend27, 1/1, 2/00, 2/-- --, 27/offset27, 1/1, 2/01, 3/111, 2/coherency, 6/registerT, 2/11, 3/010, 1/--, 5/11110, 1/0, 6/registerZ }


\paragraph{Syntax} \texttt{asdusw}\emph{coherency} \emph{extend27\_\discretionary{}{}{}offset27}[\emph{registerZ}] = \emph{registerT}

\paragraph{Description} The effective address is computed by adding the \emph{extend27\_\discretionary{}{}{}offset27} to the \emph{registerZ}.
This instruction implements atomic DUS-store on signed word (Decrement Unsigned Saturating). The word at effective address is DUSed with the \emph{registerT}. The effective address must be word aligned (multiple of 4).


\paragraph{Modifier(s)}
\noindent\begin{itemize}
\item coherency (cf. coherency in section \ref{par:modifier-coherency} on page \pageref{par:modifier-coherency})
\end{itemize}

\paragraph{Execution} {Reservation table \textsf{LSU2\_MEMW\_AUXR.Y} (Section~\ref{sec:reservation-LSU2-MEMW-AUXR.Y})}
~
{\begin{lstlisting}
stage ID:
  new offset = _SX_54(extend27_offset27);
  new coherency = _ZX_2(coherency);
stage RR:
  new base = GPR[registerZ];
  new address = base + offset;
stage E1:
  new argument2 = GPR[registerT];
stage E3:
  MEM.atomic_dus(address, 0xF, coherency << 32, argument2.32[0], registerT);
\end{lstlisting}}
\newpage
\subsubsection[{\small ASEORB: Atomic Store EOR Byte}]{ASEORB\hfill Atomic Store EOR Byte}
\label{instr:ASEORB}\index{aseorb}
 
\paragraph{Encoding} Format \textsf{LSU\_ASEORSB} (Section~\ref{sec:format-LSU-ASEORSB}), \verb|lsucode5 = 000|\\

\bitfields{ 1/P, 2/01, 3/111, 2/coherency, 6/registerT, 2/11, 3/000, 1/--, 5/10111, 1/0, 6/registerZ }


\paragraph{Syntax} \texttt{aseorb}\emph{coherency} [\emph{registerZ}] = \emph{registerT}

\paragraph{Description} The effective address is the value of the \emph{registerZ}.
This instruction implements atomic EOR-store on byte. The byte at effective address is EORed with the \emph{registerT}.


\paragraph{Modifier(s)}
\noindent\begin{itemize}
\item coherency (cf. coherency in section \ref{par:modifier-coherency} on page \pageref{par:modifier-coherency})
\end{itemize}

\paragraph{Execution} {Reservation table \textsf{LSU2\_MEMW\_AUXR} (Section~\ref{sec:reservation-LSU2-MEMW-AUXR})}
~
{\begin{lstlisting}
stage ID:
  new coherency = _ZX_2(coherency);
stage RR:
  new address = GPR[registerZ];
stage E1:
  new argument2 = GPR[registerT];
stage E3:
  MEM.atomic_eor(address, 0x1, coherency << 32, argument2.8[0], registerT);
\end{lstlisting}}
\newpage

\paragraph{Encoding} Format \textsf{LSU\_ASEORSB.O} (Section~\ref{sec:format-LSU-ASEORSB.O}), \verb|lsucode5 = 000|\\

\bitfields{ 1/P, 2/00, 2/-- --, 27/offset27, 1/1, 2/01, 3/111, 2/coherency, 6/registerT, 2/11, 3/000, 1/--, 5/10111, 1/0, 6/registerZ }


\paragraph{Syntax} \texttt{aseorb}\emph{coherency} \emph{offset27}[\emph{registerZ}] = \emph{registerT}

\paragraph{Description} The effective address is computed by adding the \emph{offset27} to the \emph{registerZ}.
This instruction implements atomic EOR-store on byte. The byte at effective address is EORed with the \emph{registerT}.


\paragraph{Modifier(s)}
\noindent\begin{itemize}
\item coherency (cf. coherency in section \ref{par:modifier-coherency} on page \pageref{par:modifier-coherency})
\end{itemize}

\paragraph{Execution} {Reservation table \textsf{LSU2\_MEMW\_AUXR.X} (Section~\ref{sec:reservation-LSU2-MEMW-AUXR.X})}
~
{\begin{lstlisting}
stage ID:
  new offset = _SX_27(offset27);
  new coherency = _ZX_2(coherency);
stage RR:
  new base = GPR[registerZ];
  new address = base + offset;
stage E1:
  new argument2 = GPR[registerT];
stage E3:
  MEM.atomic_eor(address, 0x1, coherency << 32, argument2.8[0], registerT);
\end{lstlisting}}
\newpage

\paragraph{Encoding} Format \textsf{LSU\_ASEORSB.D} (Section~\ref{sec:format-LSU-ASEORSB.D}), \verb|lsucode5 = 000|\\

\bitfields{ 1/P, 2/00, 2/-- --, 27/extend27, 1/1, 2/00, 2/-- --, 27/offset27, 1/1, 2/01, 3/111, 2/coherency, 6/registerT, 2/11, 3/000, 1/--, 5/10111, 1/0, 6/registerZ }


\paragraph{Syntax} \texttt{aseorb}\emph{coherency} \emph{extend27\_\discretionary{}{}{}offset27}[\emph{registerZ}] = \emph{registerT}

\paragraph{Description} The effective address is computed by adding the \emph{extend27\_\discretionary{}{}{}offset27} to the \emph{registerZ}.
This instruction implements atomic EOR-store on byte. The byte at effective address is EORed with the \emph{registerT}.


\paragraph{Modifier(s)}
\noindent\begin{itemize}
\item coherency (cf. coherency in section \ref{par:modifier-coherency} on page \pageref{par:modifier-coherency})
\end{itemize}

\paragraph{Execution} {Reservation table \textsf{LSU2\_MEMW\_AUXR.Y} (Section~\ref{sec:reservation-LSU2-MEMW-AUXR.Y})}
~
{\begin{lstlisting}
stage ID:
  new offset = _SX_54(extend27_offset27);
  new coherency = _ZX_2(coherency);
stage RR:
  new base = GPR[registerZ];
  new address = base + offset;
stage E1:
  new argument2 = GPR[registerT];
stage E3:
  MEM.atomic_eor(address, 0x1, coherency << 32, argument2.8[0], registerT);
\end{lstlisting}}
\newpage
\subsubsection[{\small ASEORD: Atomic Store EOR Double Word}]{ASEORD\hfill Atomic Store EOR Double Word}
\label{instr:ASEORD}\index{aseord}
 
\paragraph{Encoding} Format \textsf{LSU\_ASEORSB} (Section~\ref{sec:format-LSU-ASEORSB}), \verb|lsucode5 = 011|\\

\bitfields{ 1/P, 2/01, 3/111, 2/coherency, 6/registerT, 2/11, 3/011, 1/--, 5/10111, 1/0, 6/registerZ }


\paragraph{Syntax} \texttt{aseord}\emph{coherency} [\emph{registerZ}] = \emph{registerT}

\paragraph{Description} The effective address is the value of the \emph{registerZ}.
This instruction implements atomic EOR-store on double word. The double word at effective address is EORed with the \emph{registerT}. The effective address must be double word aligned (multiple of 8).


\paragraph{Modifier(s)}
\noindent\begin{itemize}
\item coherency (cf. coherency in section \ref{par:modifier-coherency} on page \pageref{par:modifier-coherency})
\end{itemize}

\paragraph{Execution} {Reservation table \textsf{LSU2\_MEMW\_AUXR} (Section~\ref{sec:reservation-LSU2-MEMW-AUXR})}
~
{\begin{lstlisting}
stage ID:
  new coherency = _ZX_2(coherency);
stage RR:
  new address = GPR[registerZ];
stage E1:
  new argument2 = GPR[registerT];
stage E3:
  MEM.atomic_eor(address, 0xFF, coherency << 32, argument2.64[0], registerT);
\end{lstlisting}}
\newpage

\paragraph{Encoding} Format \textsf{LSU\_ASEORSB.O} (Section~\ref{sec:format-LSU-ASEORSB.O}), \verb|lsucode5 = 011|\\

\bitfields{ 1/P, 2/00, 2/-- --, 27/offset27, 1/1, 2/01, 3/111, 2/coherency, 6/registerT, 2/11, 3/011, 1/--, 5/10111, 1/0, 6/registerZ }


\paragraph{Syntax} \texttt{aseord}\emph{coherency} \emph{offset27}[\emph{registerZ}] = \emph{registerT}

\paragraph{Description} The effective address is computed by adding the \emph{offset27} to the \emph{registerZ}.
This instruction implements atomic EOR-store on double word. The double word at effective address is EORed with the \emph{registerT}. The effective address must be double word aligned (multiple of 8).


\paragraph{Modifier(s)}
\noindent\begin{itemize}
\item coherency (cf. coherency in section \ref{par:modifier-coherency} on page \pageref{par:modifier-coherency})
\end{itemize}

\paragraph{Execution} {Reservation table \textsf{LSU2\_MEMW\_AUXR.X} (Section~\ref{sec:reservation-LSU2-MEMW-AUXR.X})}
~
{\begin{lstlisting}
stage ID:
  new offset = _SX_27(offset27);
  new coherency = _ZX_2(coherency);
stage RR:
  new base = GPR[registerZ];
  new address = base + offset;
stage E1:
  new argument2 = GPR[registerT];
stage E3:
  MEM.atomic_eor(address, 0xFF, coherency << 32, argument2.64[0], registerT);
\end{lstlisting}}
\newpage

\paragraph{Encoding} Format \textsf{LSU\_ASEORSB.D} (Section~\ref{sec:format-LSU-ASEORSB.D}), \verb|lsucode5 = 011|\\

\bitfields{ 1/P, 2/00, 2/-- --, 27/extend27, 1/1, 2/00, 2/-- --, 27/offset27, 1/1, 2/01, 3/111, 2/coherency, 6/registerT, 2/11, 3/011, 1/--, 5/10111, 1/0, 6/registerZ }


\paragraph{Syntax} \texttt{aseord}\emph{coherency} \emph{extend27\_\discretionary{}{}{}offset27}[\emph{registerZ}] = \emph{registerT}

\paragraph{Description} The effective address is computed by adding the \emph{extend27\_\discretionary{}{}{}offset27} to the \emph{registerZ}.
This instruction implements atomic EOR-store on double word. The double word at effective address is EORed with the \emph{registerT}. The effective address must be double word aligned (multiple of 8).


\paragraph{Modifier(s)}
\noindent\begin{itemize}
\item coherency (cf. coherency in section \ref{par:modifier-coherency} on page \pageref{par:modifier-coherency})
\end{itemize}

\paragraph{Execution} {Reservation table \textsf{LSU2\_MEMW\_AUXR.Y} (Section~\ref{sec:reservation-LSU2-MEMW-AUXR.Y})}
~
{\begin{lstlisting}
stage ID:
  new offset = _SX_54(extend27_offset27);
  new coherency = _ZX_2(coherency);
stage RR:
  new base = GPR[registerZ];
  new address = base + offset;
stage E1:
  new argument2 = GPR[registerT];
stage E3:
  MEM.atomic_eor(address, 0xFF, coherency << 32, argument2.64[0], registerT);
\end{lstlisting}}
\newpage
\subsubsection[{\small ASEORH: Atomic Store EOR Half Word}]{ASEORH\hfill Atomic Store EOR Half Word}
\label{instr:ASEORH}\index{aseorh}
 
\paragraph{Encoding} Format \textsf{LSU\_ASEORSB} (Section~\ref{sec:format-LSU-ASEORSB}), \verb|lsucode5 = 001|\\

\bitfields{ 1/P, 2/01, 3/111, 2/coherency, 6/registerT, 2/11, 3/001, 1/--, 5/10111, 1/0, 6/registerZ }


\paragraph{Syntax} \texttt{aseorh}\emph{coherency} [\emph{registerZ}] = \emph{registerT}

\paragraph{Description} The effective address is the value of the \emph{registerZ}.
This instruction implements atomic EOR-store on half word. The half word at effective address is EORed with the \emph{registerT}. The effective address must be half word aligned (multiple of 2).


\paragraph{Modifier(s)}
\noindent\begin{itemize}
\item coherency (cf. coherency in section \ref{par:modifier-coherency} on page \pageref{par:modifier-coherency})
\end{itemize}

\paragraph{Execution} {Reservation table \textsf{LSU2\_MEMW\_AUXR} (Section~\ref{sec:reservation-LSU2-MEMW-AUXR})}
~
{\begin{lstlisting}
stage ID:
  new coherency = _ZX_2(coherency);
stage RR:
  new address = GPR[registerZ];
stage E1:
  new argument2 = GPR[registerT];
stage E3:
  MEM.atomic_eor(address, 0x3, coherency << 32, argument2.16[0], registerT);
\end{lstlisting}}
\newpage

\paragraph{Encoding} Format \textsf{LSU\_ASEORSB.O} (Section~\ref{sec:format-LSU-ASEORSB.O}), \verb|lsucode5 = 001|\\

\bitfields{ 1/P, 2/00, 2/-- --, 27/offset27, 1/1, 2/01, 3/111, 2/coherency, 6/registerT, 2/11, 3/001, 1/--, 5/10111, 1/0, 6/registerZ }


\paragraph{Syntax} \texttt{aseorh}\emph{coherency} \emph{offset27}[\emph{registerZ}] = \emph{registerT}

\paragraph{Description} The effective address is computed by adding the \emph{offset27} to the \emph{registerZ}.
This instruction implements atomic EOR-store on half word. The half word at effective address is EORed with the \emph{registerT}. The effective address must be half word aligned (multiple of 2).


\paragraph{Modifier(s)}
\noindent\begin{itemize}
\item coherency (cf. coherency in section \ref{par:modifier-coherency} on page \pageref{par:modifier-coherency})
\end{itemize}

\paragraph{Execution} {Reservation table \textsf{LSU2\_MEMW\_AUXR.X} (Section~\ref{sec:reservation-LSU2-MEMW-AUXR.X})}
~
{\begin{lstlisting}
stage ID:
  new offset = _SX_27(offset27);
  new coherency = _ZX_2(coherency);
stage RR:
  new base = GPR[registerZ];
  new address = base + offset;
stage E1:
  new argument2 = GPR[registerT];
stage E3:
  MEM.atomic_eor(address, 0x3, coherency << 32, argument2.16[0], registerT);
\end{lstlisting}}
\newpage

\paragraph{Encoding} Format \textsf{LSU\_ASEORSB.D} (Section~\ref{sec:format-LSU-ASEORSB.D}), \verb|lsucode5 = 001|\\

\bitfields{ 1/P, 2/00, 2/-- --, 27/extend27, 1/1, 2/00, 2/-- --, 27/offset27, 1/1, 2/01, 3/111, 2/coherency, 6/registerT, 2/11, 3/001, 1/--, 5/10111, 1/0, 6/registerZ }


\paragraph{Syntax} \texttt{aseorh}\emph{coherency} \emph{extend27\_\discretionary{}{}{}offset27}[\emph{registerZ}] = \emph{registerT}

\paragraph{Description} The effective address is computed by adding the \emph{extend27\_\discretionary{}{}{}offset27} to the \emph{registerZ}.
This instruction implements atomic EOR-store on half word. The half word at effective address is EORed with the \emph{registerT}. The effective address must be half word aligned (multiple of 2).


\paragraph{Modifier(s)}
\noindent\begin{itemize}
\item coherency (cf. coherency in section \ref{par:modifier-coherency} on page \pageref{par:modifier-coherency})
\end{itemize}

\paragraph{Execution} {Reservation table \textsf{LSU2\_MEMW\_AUXR.Y} (Section~\ref{sec:reservation-LSU2-MEMW-AUXR.Y})}
~
{\begin{lstlisting}
stage ID:
  new offset = _SX_54(extend27_offset27);
  new coherency = _ZX_2(coherency);
stage RR:
  new base = GPR[registerZ];
  new address = base + offset;
stage E1:
  new argument2 = GPR[registerT];
stage E3:
  MEM.atomic_eor(address, 0x3, coherency << 32, argument2.16[0], registerT);
\end{lstlisting}}
\newpage
\subsubsection[{\small ASEORW: Atomic Store EOR Word}]{ASEORW\hfill Atomic Store EOR Word}
\label{instr:ASEORW}\index{aseorw}
 
\paragraph{Encoding} Format \textsf{LSU\_ASEORSB} (Section~\ref{sec:format-LSU-ASEORSB}), \verb|lsucode5 = 010|\\

\bitfields{ 1/P, 2/01, 3/111, 2/coherency, 6/registerT, 2/11, 3/010, 1/--, 5/10111, 1/0, 6/registerZ }


\paragraph{Syntax} \texttt{aseorw}\emph{coherency} [\emph{registerZ}] = \emph{registerT}

\paragraph{Description} The effective address is the value of the \emph{registerZ}.
This instruction implements atomic EOR-store on word. The word at effective address is EORed with the \emph{registerT}. The effective address must be word aligned (multiple of 4).


\paragraph{Modifier(s)}
\noindent\begin{itemize}
\item coherency (cf. coherency in section \ref{par:modifier-coherency} on page \pageref{par:modifier-coherency})
\end{itemize}

\paragraph{Execution} {Reservation table \textsf{LSU2\_MEMW\_AUXR} (Section~\ref{sec:reservation-LSU2-MEMW-AUXR})}
~
{\begin{lstlisting}
stage ID:
  new coherency = _ZX_2(coherency);
stage RR:
  new address = GPR[registerZ];
stage E1:
  new argument2 = GPR[registerT];
stage E3:
  MEM.atomic_eor(address, 0xF, coherency << 32, argument2.32[0], registerT);
\end{lstlisting}}
\newpage

\paragraph{Encoding} Format \textsf{LSU\_ASEORSB.O} (Section~\ref{sec:format-LSU-ASEORSB.O}), \verb|lsucode5 = 010|\\

\bitfields{ 1/P, 2/00, 2/-- --, 27/offset27, 1/1, 2/01, 3/111, 2/coherency, 6/registerT, 2/11, 3/010, 1/--, 5/10111, 1/0, 6/registerZ }


\paragraph{Syntax} \texttt{aseorw}\emph{coherency} \emph{offset27}[\emph{registerZ}] = \emph{registerT}

\paragraph{Description} The effective address is computed by adding the \emph{offset27} to the \emph{registerZ}.
This instruction implements atomic EOR-store on word. The word at effective address is EORed with the \emph{registerT}. The effective address must be word aligned (multiple of 4).


\paragraph{Modifier(s)}
\noindent\begin{itemize}
\item coherency (cf. coherency in section \ref{par:modifier-coherency} on page \pageref{par:modifier-coherency})
\end{itemize}

\paragraph{Execution} {Reservation table \textsf{LSU2\_MEMW\_AUXR.X} (Section~\ref{sec:reservation-LSU2-MEMW-AUXR.X})}
~
{\begin{lstlisting}
stage ID:
  new offset = _SX_27(offset27);
  new coherency = _ZX_2(coherency);
stage RR:
  new base = GPR[registerZ];
  new address = base + offset;
stage E1:
  new argument2 = GPR[registerT];
stage E3:
  MEM.atomic_eor(address, 0xF, coherency << 32, argument2.32[0], registerT);
\end{lstlisting}}
\newpage

\paragraph{Encoding} Format \textsf{LSU\_ASEORSB.D} (Section~\ref{sec:format-LSU-ASEORSB.D}), \verb|lsucode5 = 010|\\

\bitfields{ 1/P, 2/00, 2/-- --, 27/extend27, 1/1, 2/00, 2/-- --, 27/offset27, 1/1, 2/01, 3/111, 2/coherency, 6/registerT, 2/11, 3/010, 1/--, 5/10111, 1/0, 6/registerZ }


\paragraph{Syntax} \texttt{aseorw}\emph{coherency} \emph{extend27\_\discretionary{}{}{}offset27}[\emph{registerZ}] = \emph{registerT}

\paragraph{Description} The effective address is computed by adding the \emph{extend27\_\discretionary{}{}{}offset27} to the \emph{registerZ}.
This instruction implements atomic EOR-store on word. The word at effective address is EORed with the \emph{registerT}. The effective address must be word aligned (multiple of 4).


\paragraph{Modifier(s)}
\noindent\begin{itemize}
\item coherency (cf. coherency in section \ref{par:modifier-coherency} on page \pageref{par:modifier-coherency})
\end{itemize}

\paragraph{Execution} {Reservation table \textsf{LSU2\_MEMW\_AUXR.Y} (Section~\ref{sec:reservation-LSU2-MEMW-AUXR.Y})}
~
{\begin{lstlisting}
stage ID:
  new offset = _SX_54(extend27_offset27);
  new coherency = _ZX_2(coherency);
stage RR:
  new base = GPR[registerZ];
  new address = base + offset;
stage E1:
  new argument2 = GPR[registerT];
stage E3:
  MEM.atomic_eor(address, 0xF, coherency << 32, argument2.32[0], registerT);
\end{lstlisting}}
\newpage
\subsubsection[{\small ASH: Atomic Store Half Word}]{ASH\hfill Atomic Store Half Word}
\label{instr:ASH}\index{ash}
 
\paragraph{Encoding} Format \textsf{LSU\_ASSB} (Section~\ref{sec:format-LSU-ASSB}), \verb|lsucode5 = 001|\\

\bitfields{ 1/P, 2/01, 3/111, 2/coherency, 6/registerT, 2/11, 3/001, 1/--, 5/10000, 1/0, 6/registerZ }


\paragraph{Syntax} \texttt{ash}\emph{coherency} [\emph{registerZ}] = \emph{registerT}

\paragraph{Description} The effective address is the value of the \emph{registerZ}.
The \emph{registerT} is atomically stored to the memory half word starting at the effective address. This instruction is provided to directly operate on the last level of the memory hierarchy, which is typically a processor DDR or a cluster TCM, irrespective of the MMU DCP (Data Cache Policy) settings. The effective address must be half word aligned (multiple of 2).


\paragraph{Modifier(s)}
\noindent\begin{itemize}
\item coherency (cf. coherency in section \ref{par:modifier-coherency} on page \pageref{par:modifier-coherency})
\end{itemize}

\paragraph{Execution} {Reservation table \textsf{LSU\_MEMW\_AUXR} (Section~\ref{sec:reservation-LSU-MEMW-AUXR})}
~
{\begin{lstlisting}
stage ID:
  new coherency = _ZX_2(coherency);
stage RR:
  new address = GPR[registerZ];
stage E1:
  new argument2 = GPR[registerT];
stage E3:
  MEM.atomic_store(address, 0x3, coherency << 32, argument2, registerT);
\end{lstlisting}}
\newpage

\paragraph{Encoding} Format \textsf{LSU\_ASSB.O} (Section~\ref{sec:format-LSU-ASSB.O}), \verb|lsucode5 = 001|\\

\bitfields{ 1/P, 2/00, 2/-- --, 27/offset27, 1/1, 2/01, 3/111, 2/coherency, 6/registerT, 2/11, 3/001, 1/--, 5/10000, 1/0, 6/registerZ }


\paragraph{Syntax} \texttt{ash}\emph{coherency} \emph{offset27}[\emph{registerZ}] = \emph{registerT}

\paragraph{Description} The effective address is computed by adding the \emph{offset27} to the \emph{registerZ}.
The \emph{registerT} is atomically stored to the memory half word starting at the effective address. This instruction is provided to directly operate on the last level of the memory hierarchy, which is typically a processor DDR or a cluster TCM, irrespective of the MMU DCP (Data Cache Policy) settings. The effective address must be half word aligned (multiple of 2).


\paragraph{Modifier(s)}
\noindent\begin{itemize}
\item coherency (cf. coherency in section \ref{par:modifier-coherency} on page \pageref{par:modifier-coherency})
\end{itemize}

\paragraph{Execution} {Reservation table \textsf{LSU\_MEMW\_AUXR.X} (Section~\ref{sec:reservation-LSU-MEMW-AUXR.X})}
~
{\begin{lstlisting}
stage ID:
  new offset = _SX_27(offset27);
  new coherency = _ZX_2(coherency);
stage RR:
  new base = GPR[registerZ];
  new address = base + offset;
stage E1:
  new argument2 = GPR[registerT];
stage E3:
  MEM.atomic_store(address, 0x3, coherency << 32, argument2, registerT);
\end{lstlisting}}
\newpage

\paragraph{Encoding} Format \textsf{LSU\_ASSB.D} (Section~\ref{sec:format-LSU-ASSB.D}), \verb|lsucode5 = 001|\\

\bitfields{ 1/P, 2/00, 2/-- --, 27/extend27, 1/1, 2/00, 2/-- --, 27/offset27, 1/1, 2/01, 3/111, 2/coherency, 6/registerT, 2/11, 3/001, 1/--, 5/10000, 1/0, 6/registerZ }


\paragraph{Syntax} \texttt{ash}\emph{coherency} \emph{extend27\_\discretionary{}{}{}offset27}[\emph{registerZ}] = \emph{registerT}

\paragraph{Description} The effective address is computed by adding the \emph{extend27\_\discretionary{}{}{}offset27} to the \emph{registerZ}.
The \emph{registerT} is atomically stored to the memory half word starting at the effective address. This instruction is provided to directly operate on the last level of the memory hierarchy, which is typically a processor DDR or a cluster TCM, irrespective of the MMU DCP (Data Cache Policy) settings. The effective address must be half word aligned (multiple of 2).


\paragraph{Modifier(s)}
\noindent\begin{itemize}
\item coherency (cf. coherency in section \ref{par:modifier-coherency} on page \pageref{par:modifier-coherency})
\end{itemize}

\paragraph{Execution} {Reservation table \textsf{LSU\_MEMW\_AUXR.Y} (Section~\ref{sec:reservation-LSU-MEMW-AUXR.Y})}
~
{\begin{lstlisting}
stage ID:
  new offset = _SX_54(extend27_offset27);
  new coherency = _ZX_2(coherency);
stage RR:
  new base = GPR[registerZ];
  new address = base + offset;
stage E1:
  new argument2 = GPR[registerT];
stage E3:
  MEM.atomic_store(address, 0x3, coherency << 32, argument2, registerT);
\end{lstlisting}}
\newpage
\subsubsection[{\small ASIORB: Atomic Store IOR Byte}]{ASIORB\hfill Atomic Store IOR Byte}
\label{instr:ASIORB}\index{asiorb}
 
\paragraph{Encoding} Format \textsf{LSU\_ASIORSB} (Section~\ref{sec:format-LSU-ASIORSB}), \verb|lsucode5 = 000|\\

\bitfields{ 1/P, 2/01, 3/111, 2/coherency, 6/registerT, 2/11, 3/000, 1/--, 5/10110, 1/0, 6/registerZ }


\paragraph{Syntax} \texttt{asiorb}\emph{coherency} [\emph{registerZ}] = \emph{registerT}

\paragraph{Description} The effective address is the value of the \emph{registerZ}.
This instruction implements atomic IOR-store on byte. The byte at effective address is IORed with the \emph{registerT}.


\paragraph{Modifier(s)}
\noindent\begin{itemize}
\item coherency (cf. coherency in section \ref{par:modifier-coherency} on page \pageref{par:modifier-coherency})
\end{itemize}

\paragraph{Execution} {Reservation table \textsf{LSU2\_MEMW\_AUXR} (Section~\ref{sec:reservation-LSU2-MEMW-AUXR})}
~
{\begin{lstlisting}
stage ID:
  new coherency = _ZX_2(coherency);
stage RR:
  new address = GPR[registerZ];
stage E1:
  new argument2 = GPR[registerT];
stage E3:
  MEM.atomic_ior(address, 0x1, coherency << 32, argument2.8[0], registerT);
\end{lstlisting}}
\newpage

\paragraph{Encoding} Format \textsf{LSU\_ASIORSB.O} (Section~\ref{sec:format-LSU-ASIORSB.O}), \verb|lsucode5 = 000|\\

\bitfields{ 1/P, 2/00, 2/-- --, 27/offset27, 1/1, 2/01, 3/111, 2/coherency, 6/registerT, 2/11, 3/000, 1/--, 5/10110, 1/0, 6/registerZ }


\paragraph{Syntax} \texttt{asiorb}\emph{coherency} \emph{offset27}[\emph{registerZ}] = \emph{registerT}

\paragraph{Description} The effective address is computed by adding the \emph{offset27} to the \emph{registerZ}.
This instruction implements atomic IOR-store on byte. The byte at effective address is IORed with the \emph{registerT}.


\paragraph{Modifier(s)}
\noindent\begin{itemize}
\item coherency (cf. coherency in section \ref{par:modifier-coherency} on page \pageref{par:modifier-coherency})
\end{itemize}

\paragraph{Execution} {Reservation table \textsf{LSU2\_MEMW\_AUXR.X} (Section~\ref{sec:reservation-LSU2-MEMW-AUXR.X})}
~
{\begin{lstlisting}
stage ID:
  new offset = _SX_27(offset27);
  new coherency = _ZX_2(coherency);
stage RR:
  new base = GPR[registerZ];
  new address = base + offset;
stage E1:
  new argument2 = GPR[registerT];
stage E3:
  MEM.atomic_ior(address, 0x1, coherency << 32, argument2.8[0], registerT);
\end{lstlisting}}
\newpage

\paragraph{Encoding} Format \textsf{LSU\_ASIORSB.D} (Section~\ref{sec:format-LSU-ASIORSB.D}), \verb|lsucode5 = 000|\\

\bitfields{ 1/P, 2/00, 2/-- --, 27/extend27, 1/1, 2/00, 2/-- --, 27/offset27, 1/1, 2/01, 3/111, 2/coherency, 6/registerT, 2/11, 3/000, 1/--, 5/10110, 1/0, 6/registerZ }


\paragraph{Syntax} \texttt{asiorb}\emph{coherency} \emph{extend27\_\discretionary{}{}{}offset27}[\emph{registerZ}] = \emph{registerT}

\paragraph{Description} The effective address is computed by adding the \emph{extend27\_\discretionary{}{}{}offset27} to the \emph{registerZ}.
This instruction implements atomic IOR-store on byte. The byte at effective address is IORed with the \emph{registerT}.


\paragraph{Modifier(s)}
\noindent\begin{itemize}
\item coherency (cf. coherency in section \ref{par:modifier-coherency} on page \pageref{par:modifier-coherency})
\end{itemize}

\paragraph{Execution} {Reservation table \textsf{LSU2\_MEMW\_AUXR.Y} (Section~\ref{sec:reservation-LSU2-MEMW-AUXR.Y})}
~
{\begin{lstlisting}
stage ID:
  new offset = _SX_54(extend27_offset27);
  new coherency = _ZX_2(coherency);
stage RR:
  new base = GPR[registerZ];
  new address = base + offset;
stage E1:
  new argument2 = GPR[registerT];
stage E3:
  MEM.atomic_ior(address, 0x1, coherency << 32, argument2.8[0], registerT);
\end{lstlisting}}
\newpage
\subsubsection[{\small ASIORD: Atomic Store IOR Double Word}]{ASIORD\hfill Atomic Store IOR Double Word}
\label{instr:ASIORD}\index{asiord}
 
\paragraph{Encoding} Format \textsf{LSU\_ASIORSB} (Section~\ref{sec:format-LSU-ASIORSB}), \verb|lsucode5 = 011|\\

\bitfields{ 1/P, 2/01, 3/111, 2/coherency, 6/registerT, 2/11, 3/011, 1/--, 5/10110, 1/0, 6/registerZ }


\paragraph{Syntax} \texttt{asiord}\emph{coherency} [\emph{registerZ}] = \emph{registerT}

\paragraph{Description} The effective address is the value of the \emph{registerZ}.
This instruction implements atomic IOR-store on double word. The double word at effective address is IORed with the \emph{registerT}. The effective address must be double word aligned (multiple of 8).


\paragraph{Modifier(s)}
\noindent\begin{itemize}
\item coherency (cf. coherency in section \ref{par:modifier-coherency} on page \pageref{par:modifier-coherency})
\end{itemize}

\paragraph{Execution} {Reservation table \textsf{LSU2\_MEMW\_AUXR} (Section~\ref{sec:reservation-LSU2-MEMW-AUXR})}
~
{\begin{lstlisting}
stage ID:
  new coherency = _ZX_2(coherency);
stage RR:
  new address = GPR[registerZ];
stage E1:
  new argument2 = GPR[registerT];
stage E3:
  MEM.atomic_ior(address, 0xFF, coherency << 32, argument2.64[0], registerT);
\end{lstlisting}}
\newpage

\paragraph{Encoding} Format \textsf{LSU\_ASIORSB.O} (Section~\ref{sec:format-LSU-ASIORSB.O}), \verb|lsucode5 = 011|\\

\bitfields{ 1/P, 2/00, 2/-- --, 27/offset27, 1/1, 2/01, 3/111, 2/coherency, 6/registerT, 2/11, 3/011, 1/--, 5/10110, 1/0, 6/registerZ }


\paragraph{Syntax} \texttt{asiord}\emph{coherency} \emph{offset27}[\emph{registerZ}] = \emph{registerT}

\paragraph{Description} The effective address is computed by adding the \emph{offset27} to the \emph{registerZ}.
This instruction implements atomic IOR-store on double word. The double word at effective address is IORed with the \emph{registerT}. The effective address must be double word aligned (multiple of 8).


\paragraph{Modifier(s)}
\noindent\begin{itemize}
\item coherency (cf. coherency in section \ref{par:modifier-coherency} on page \pageref{par:modifier-coherency})
\end{itemize}

\paragraph{Execution} {Reservation table \textsf{LSU2\_MEMW\_AUXR.X} (Section~\ref{sec:reservation-LSU2-MEMW-AUXR.X})}
~
{\begin{lstlisting}
stage ID:
  new offset = _SX_27(offset27);
  new coherency = _ZX_2(coherency);
stage RR:
  new base = GPR[registerZ];
  new address = base + offset;
stage E1:
  new argument2 = GPR[registerT];
stage E3:
  MEM.atomic_ior(address, 0xFF, coherency << 32, argument2.64[0], registerT);
\end{lstlisting}}
\newpage

\paragraph{Encoding} Format \textsf{LSU\_ASIORSB.D} (Section~\ref{sec:format-LSU-ASIORSB.D}), \verb|lsucode5 = 011|\\

\bitfields{ 1/P, 2/00, 2/-- --, 27/extend27, 1/1, 2/00, 2/-- --, 27/offset27, 1/1, 2/01, 3/111, 2/coherency, 6/registerT, 2/11, 3/011, 1/--, 5/10110, 1/0, 6/registerZ }


\paragraph{Syntax} \texttt{asiord}\emph{coherency} \emph{extend27\_\discretionary{}{}{}offset27}[\emph{registerZ}] = \emph{registerT}

\paragraph{Description} The effective address is computed by adding the \emph{extend27\_\discretionary{}{}{}offset27} to the \emph{registerZ}.
This instruction implements atomic IOR-store on double word. The double word at effective address is IORed with the \emph{registerT}. The effective address must be double word aligned (multiple of 8).


\paragraph{Modifier(s)}
\noindent\begin{itemize}
\item coherency (cf. coherency in section \ref{par:modifier-coherency} on page \pageref{par:modifier-coherency})
\end{itemize}

\paragraph{Execution} {Reservation table \textsf{LSU2\_MEMW\_AUXR.Y} (Section~\ref{sec:reservation-LSU2-MEMW-AUXR.Y})}
~
{\begin{lstlisting}
stage ID:
  new offset = _SX_54(extend27_offset27);
  new coherency = _ZX_2(coherency);
stage RR:
  new base = GPR[registerZ];
  new address = base + offset;
stage E1:
  new argument2 = GPR[registerT];
stage E3:
  MEM.atomic_ior(address, 0xFF, coherency << 32, argument2.64[0], registerT);
\end{lstlisting}}
\newpage
\subsubsection[{\small ASIORH: Atomic Store IOR Half Word}]{ASIORH\hfill Atomic Store IOR Half Word}
\label{instr:ASIORH}\index{asiorh}
 
\paragraph{Encoding} Format \textsf{LSU\_ASIORSB} (Section~\ref{sec:format-LSU-ASIORSB}), \verb|lsucode5 = 001|\\

\bitfields{ 1/P, 2/01, 3/111, 2/coherency, 6/registerT, 2/11, 3/001, 1/--, 5/10110, 1/0, 6/registerZ }


\paragraph{Syntax} \texttt{asiorh}\emph{coherency} [\emph{registerZ}] = \emph{registerT}

\paragraph{Description} The effective address is the value of the \emph{registerZ}.
This instruction implements atomic IOR-store on half word. The half word at effective address is IORed with the \emph{registerT}. The effective address must be half word aligned (multiple of 2).


\paragraph{Modifier(s)}
\noindent\begin{itemize}
\item coherency (cf. coherency in section \ref{par:modifier-coherency} on page \pageref{par:modifier-coherency})
\end{itemize}

\paragraph{Execution} {Reservation table \textsf{LSU2\_MEMW\_AUXR} (Section~\ref{sec:reservation-LSU2-MEMW-AUXR})}
~
{\begin{lstlisting}
stage ID:
  new coherency = _ZX_2(coherency);
stage RR:
  new address = GPR[registerZ];
stage E1:
  new argument2 = GPR[registerT];
stage E3:
  MEM.atomic_ior(address, 0x3, coherency << 32, argument2.16[0], registerT);
\end{lstlisting}}
\newpage

\paragraph{Encoding} Format \textsf{LSU\_ASIORSB.O} (Section~\ref{sec:format-LSU-ASIORSB.O}), \verb|lsucode5 = 001|\\

\bitfields{ 1/P, 2/00, 2/-- --, 27/offset27, 1/1, 2/01, 3/111, 2/coherency, 6/registerT, 2/11, 3/001, 1/--, 5/10110, 1/0, 6/registerZ }


\paragraph{Syntax} \texttt{asiorh}\emph{coherency} \emph{offset27}[\emph{registerZ}] = \emph{registerT}

\paragraph{Description} The effective address is computed by adding the \emph{offset27} to the \emph{registerZ}.
This instruction implements atomic IOR-store on half word. The half word at effective address is IORed with the \emph{registerT}. The effective address must be half word aligned (multiple of 2).


\paragraph{Modifier(s)}
\noindent\begin{itemize}
\item coherency (cf. coherency in section \ref{par:modifier-coherency} on page \pageref{par:modifier-coherency})
\end{itemize}

\paragraph{Execution} {Reservation table \textsf{LSU2\_MEMW\_AUXR.X} (Section~\ref{sec:reservation-LSU2-MEMW-AUXR.X})}
~
{\begin{lstlisting}
stage ID:
  new offset = _SX_27(offset27);
  new coherency = _ZX_2(coherency);
stage RR:
  new base = GPR[registerZ];
  new address = base + offset;
stage E1:
  new argument2 = GPR[registerT];
stage E3:
  MEM.atomic_ior(address, 0x3, coherency << 32, argument2.16[0], registerT);
\end{lstlisting}}
\newpage

\paragraph{Encoding} Format \textsf{LSU\_ASIORSB.D} (Section~\ref{sec:format-LSU-ASIORSB.D}), \verb|lsucode5 = 001|\\

\bitfields{ 1/P, 2/00, 2/-- --, 27/extend27, 1/1, 2/00, 2/-- --, 27/offset27, 1/1, 2/01, 3/111, 2/coherency, 6/registerT, 2/11, 3/001, 1/--, 5/10110, 1/0, 6/registerZ }


\paragraph{Syntax} \texttt{asiorh}\emph{coherency} \emph{extend27\_\discretionary{}{}{}offset27}[\emph{registerZ}] = \emph{registerT}

\paragraph{Description} The effective address is computed by adding the \emph{extend27\_\discretionary{}{}{}offset27} to the \emph{registerZ}.
This instruction implements atomic IOR-store on half word. The half word at effective address is IORed with the \emph{registerT}. The effective address must be half word aligned (multiple of 2).


\paragraph{Modifier(s)}
\noindent\begin{itemize}
\item coherency (cf. coherency in section \ref{par:modifier-coherency} on page \pageref{par:modifier-coherency})
\end{itemize}

\paragraph{Execution} {Reservation table \textsf{LSU2\_MEMW\_AUXR.Y} (Section~\ref{sec:reservation-LSU2-MEMW-AUXR.Y})}
~
{\begin{lstlisting}
stage ID:
  new offset = _SX_54(extend27_offset27);
  new coherency = _ZX_2(coherency);
stage RR:
  new base = GPR[registerZ];
  new address = base + offset;
stage E1:
  new argument2 = GPR[registerT];
stage E3:
  MEM.atomic_ior(address, 0x3, coherency << 32, argument2.16[0], registerT);
\end{lstlisting}}
\newpage
\subsubsection[{\small ASIORW: Atomic Store IOR Word}]{ASIORW\hfill Atomic Store IOR Word}
\label{instr:ASIORW}\index{asiorw}
 
\paragraph{Encoding} Format \textsf{LSU\_ASIORSB} (Section~\ref{sec:format-LSU-ASIORSB}), \verb|lsucode5 = 010|\\

\bitfields{ 1/P, 2/01, 3/111, 2/coherency, 6/registerT, 2/11, 3/010, 1/--, 5/10110, 1/0, 6/registerZ }


\paragraph{Syntax} \texttt{asiorw}\emph{coherency} [\emph{registerZ}] = \emph{registerT}

\paragraph{Description} The effective address is the value of the \emph{registerZ}.
This instruction implements atomic IOR-store on word. The word at effective address is IORed with the \emph{registerT}. The effective address must be word aligned (multiple of 4).


\paragraph{Modifier(s)}
\noindent\begin{itemize}
\item coherency (cf. coherency in section \ref{par:modifier-coherency} on page \pageref{par:modifier-coherency})
\end{itemize}

\paragraph{Execution} {Reservation table \textsf{LSU2\_MEMW\_AUXR} (Section~\ref{sec:reservation-LSU2-MEMW-AUXR})}
~
{\begin{lstlisting}
stage ID:
  new coherency = _ZX_2(coherency);
stage RR:
  new address = GPR[registerZ];
stage E1:
  new argument2 = GPR[registerT];
stage E3:
  MEM.atomic_ior(address, 0xF, coherency << 32, argument2.32[0], registerT);
\end{lstlisting}}
\newpage

\paragraph{Encoding} Format \textsf{LSU\_ASIORSB.O} (Section~\ref{sec:format-LSU-ASIORSB.O}), \verb|lsucode5 = 010|\\

\bitfields{ 1/P, 2/00, 2/-- --, 27/offset27, 1/1, 2/01, 3/111, 2/coherency, 6/registerT, 2/11, 3/010, 1/--, 5/10110, 1/0, 6/registerZ }


\paragraph{Syntax} \texttt{asiorw}\emph{coherency} \emph{offset27}[\emph{registerZ}] = \emph{registerT}

\paragraph{Description} The effective address is computed by adding the \emph{offset27} to the \emph{registerZ}.
This instruction implements atomic IOR-store on word. The word at effective address is IORed with the \emph{registerT}. The effective address must be word aligned (multiple of 4).


\paragraph{Modifier(s)}
\noindent\begin{itemize}
\item coherency (cf. coherency in section \ref{par:modifier-coherency} on page \pageref{par:modifier-coherency})
\end{itemize}

\paragraph{Execution} {Reservation table \textsf{LSU2\_MEMW\_AUXR.X} (Section~\ref{sec:reservation-LSU2-MEMW-AUXR.X})}
~
{\begin{lstlisting}
stage ID:
  new offset = _SX_27(offset27);
  new coherency = _ZX_2(coherency);
stage RR:
  new base = GPR[registerZ];
  new address = base + offset;
stage E1:
  new argument2 = GPR[registerT];
stage E3:
  MEM.atomic_ior(address, 0xF, coherency << 32, argument2.32[0], registerT);
\end{lstlisting}}
\newpage

\paragraph{Encoding} Format \textsf{LSU\_ASIORSB.D} (Section~\ref{sec:format-LSU-ASIORSB.D}), \verb|lsucode5 = 010|\\

\bitfields{ 1/P, 2/00, 2/-- --, 27/extend27, 1/1, 2/00, 2/-- --, 27/offset27, 1/1, 2/01, 3/111, 2/coherency, 6/registerT, 2/11, 3/010, 1/--, 5/10110, 1/0, 6/registerZ }


\paragraph{Syntax} \texttt{asiorw}\emph{coherency} \emph{extend27\_\discretionary{}{}{}offset27}[\emph{registerZ}] = \emph{registerT}

\paragraph{Description} The effective address is computed by adding the \emph{extend27\_\discretionary{}{}{}offset27} to the \emph{registerZ}.
This instruction implements atomic IOR-store on word. The word at effective address is IORed with the \emph{registerT}. The effective address must be word aligned (multiple of 4).


\paragraph{Modifier(s)}
\noindent\begin{itemize}
\item coherency (cf. coherency in section \ref{par:modifier-coherency} on page \pageref{par:modifier-coherency})
\end{itemize}

\paragraph{Execution} {Reservation table \textsf{LSU2\_MEMW\_AUXR.Y} (Section~\ref{sec:reservation-LSU2-MEMW-AUXR.Y})}
~
{\begin{lstlisting}
stage ID:
  new offset = _SX_54(extend27_offset27);
  new coherency = _ZX_2(coherency);
stage RR:
  new base = GPR[registerZ];
  new address = base + offset;
stage E1:
  new argument2 = GPR[registerT];
stage E3:
  MEM.atomic_ior(address, 0xF, coherency << 32, argument2.32[0], registerT);
\end{lstlisting}}
\newpage
\subsubsection[{\small ASMAXB: Atomic Store MAX Byte}]{ASMAXB\hfill Atomic Store MAX Byte}
\label{instr:ASMAXB}\index{asmaxb}
 
\paragraph{Encoding} Format \textsf{LSU\_ASMAXSB} (Section~\ref{sec:format-LSU-ASMAXSB}), \verb|lsucode5 = 000|\\

\bitfields{ 1/P, 2/01, 3/111, 2/coherency, 6/registerT, 2/11, 3/000, 1/--, 5/11001, 1/0, 6/registerZ }


\paragraph{Syntax} \texttt{asmaxb}\emph{coherency} [\emph{registerZ}] = \emph{registerT}

\paragraph{Description} The effective address is the value of the \emph{registerZ}.
This instruction implements atomic MAX-store on signed byte. The byte at effective address is MAXed with the \emph{registerT}.


\paragraph{Modifier(s)}
\noindent\begin{itemize}
\item coherency (cf. coherency in section \ref{par:modifier-coherency} on page \pageref{par:modifier-coherency})
\end{itemize}

\paragraph{Execution} {Reservation table \textsf{LSU2\_MEMW\_AUXR} (Section~\ref{sec:reservation-LSU2-MEMW-AUXR})}
~
{\begin{lstlisting}
stage ID:
  new coherency = _ZX_2(coherency);
stage RR:
  new address = GPR[registerZ];
stage E1:
  new argument2 = GPR[registerT];
stage E3:
  MEM.atomic_max(address, 0x1, coherency << 32, argument2.8[0], registerT);
\end{lstlisting}}
\newpage

\paragraph{Encoding} Format \textsf{LSU\_ASMAXSB.O} (Section~\ref{sec:format-LSU-ASMAXSB.O}), \verb|lsucode5 = 000|\\

\bitfields{ 1/P, 2/00, 2/-- --, 27/offset27, 1/1, 2/01, 3/111, 2/coherency, 6/registerT, 2/11, 3/000, 1/--, 5/11001, 1/0, 6/registerZ }


\paragraph{Syntax} \texttt{asmaxb}\emph{coherency} \emph{offset27}[\emph{registerZ}] = \emph{registerT}

\paragraph{Description} The effective address is computed by adding the \emph{offset27} to the \emph{registerZ}.
This instruction implements atomic MAX-store on signed byte. The byte at effective address is MAXed with the \emph{registerT}.


\paragraph{Modifier(s)}
\noindent\begin{itemize}
\item coherency (cf. coherency in section \ref{par:modifier-coherency} on page \pageref{par:modifier-coherency})
\end{itemize}

\paragraph{Execution} {Reservation table \textsf{LSU2\_MEMW\_AUXR.X} (Section~\ref{sec:reservation-LSU2-MEMW-AUXR.X})}
~
{\begin{lstlisting}
stage ID:
  new offset = _SX_27(offset27);
  new coherency = _ZX_2(coherency);
stage RR:
  new base = GPR[registerZ];
  new address = base + offset;
stage E1:
  new argument2 = GPR[registerT];
stage E3:
  MEM.atomic_max(address, 0x1, coherency << 32, argument2.8[0], registerT);
\end{lstlisting}}
\newpage

\paragraph{Encoding} Format \textsf{LSU\_ASMAXSB.D} (Section~\ref{sec:format-LSU-ASMAXSB.D}), \verb|lsucode5 = 000|\\

\bitfields{ 1/P, 2/00, 2/-- --, 27/extend27, 1/1, 2/00, 2/-- --, 27/offset27, 1/1, 2/01, 3/111, 2/coherency, 6/registerT, 2/11, 3/000, 1/--, 5/11001, 1/0, 6/registerZ }


\paragraph{Syntax} \texttt{asmaxb}\emph{coherency} \emph{extend27\_\discretionary{}{}{}offset27}[\emph{registerZ}] = \emph{registerT}

\paragraph{Description} The effective address is computed by adding the \emph{extend27\_\discretionary{}{}{}offset27} to the \emph{registerZ}.
This instruction implements atomic MAX-store on signed byte. The byte at effective address is MAXed with the \emph{registerT}.


\paragraph{Modifier(s)}
\noindent\begin{itemize}
\item coherency (cf. coherency in section \ref{par:modifier-coherency} on page \pageref{par:modifier-coherency})
\end{itemize}

\paragraph{Execution} {Reservation table \textsf{LSU2\_MEMW\_AUXR.Y} (Section~\ref{sec:reservation-LSU2-MEMW-AUXR.Y})}
~
{\begin{lstlisting}
stage ID:
  new offset = _SX_54(extend27_offset27);
  new coherency = _ZX_2(coherency);
stage RR:
  new base = GPR[registerZ];
  new address = base + offset;
stage E1:
  new argument2 = GPR[registerT];
stage E3:
  MEM.atomic_max(address, 0x1, coherency << 32, argument2.8[0], registerT);
\end{lstlisting}}
\newpage
\subsubsection[{\small ASMAXD: Atomic Store MAX Double Word}]{ASMAXD\hfill Atomic Store MAX Double Word}
\label{instr:ASMAXD}\index{asmaxd}
 
\paragraph{Encoding} Format \textsf{LSU\_ASMAXSB} (Section~\ref{sec:format-LSU-ASMAXSB}), \verb|lsucode5 = 011|\\

\bitfields{ 1/P, 2/01, 3/111, 2/coherency, 6/registerT, 2/11, 3/011, 1/--, 5/11001, 1/0, 6/registerZ }


\paragraph{Syntax} \texttt{asmaxd}\emph{coherency} [\emph{registerZ}] = \emph{registerT}

\paragraph{Description} The effective address is the value of the \emph{registerZ}.
This instruction implements atomic MAX-store on signed double word. The double word at effective address is MAXed with the \emph{registerT}. The effective address must be double word aligned (multiple of 8).


\paragraph{Modifier(s)}
\noindent\begin{itemize}
\item coherency (cf. coherency in section \ref{par:modifier-coherency} on page \pageref{par:modifier-coherency})
\end{itemize}

\paragraph{Execution} {Reservation table \textsf{LSU2\_MEMW\_AUXR} (Section~\ref{sec:reservation-LSU2-MEMW-AUXR})}
~
{\begin{lstlisting}
stage ID:
  new coherency = _ZX_2(coherency);
stage RR:
  new address = GPR[registerZ];
stage E1:
  new argument2 = GPR[registerT];
stage E3:
  MEM.atomic_max(address, 0xFF, coherency << 32, argument2.64[0], registerT);
\end{lstlisting}}
\newpage

\paragraph{Encoding} Format \textsf{LSU\_ASMAXSB.O} (Section~\ref{sec:format-LSU-ASMAXSB.O}), \verb|lsucode5 = 011|\\

\bitfields{ 1/P, 2/00, 2/-- --, 27/offset27, 1/1, 2/01, 3/111, 2/coherency, 6/registerT, 2/11, 3/011, 1/--, 5/11001, 1/0, 6/registerZ }


\paragraph{Syntax} \texttt{asmaxd}\emph{coherency} \emph{offset27}[\emph{registerZ}] = \emph{registerT}

\paragraph{Description} The effective address is computed by adding the \emph{offset27} to the \emph{registerZ}.
This instruction implements atomic MAX-store on signed double word. The double word at effective address is MAXed with the \emph{registerT}. The effective address must be double word aligned (multiple of 8).


\paragraph{Modifier(s)}
\noindent\begin{itemize}
\item coherency (cf. coherency in section \ref{par:modifier-coherency} on page \pageref{par:modifier-coherency})
\end{itemize}

\paragraph{Execution} {Reservation table \textsf{LSU2\_MEMW\_AUXR.X} (Section~\ref{sec:reservation-LSU2-MEMW-AUXR.X})}
~
{\begin{lstlisting}
stage ID:
  new offset = _SX_27(offset27);
  new coherency = _ZX_2(coherency);
stage RR:
  new base = GPR[registerZ];
  new address = base + offset;
stage E1:
  new argument2 = GPR[registerT];
stage E3:
  MEM.atomic_max(address, 0xFF, coherency << 32, argument2.64[0], registerT);
\end{lstlisting}}
\newpage

\paragraph{Encoding} Format \textsf{LSU\_ASMAXSB.D} (Section~\ref{sec:format-LSU-ASMAXSB.D}), \verb|lsucode5 = 011|\\

\bitfields{ 1/P, 2/00, 2/-- --, 27/extend27, 1/1, 2/00, 2/-- --, 27/offset27, 1/1, 2/01, 3/111, 2/coherency, 6/registerT, 2/11, 3/011, 1/--, 5/11001, 1/0, 6/registerZ }


\paragraph{Syntax} \texttt{asmaxd}\emph{coherency} \emph{extend27\_\discretionary{}{}{}offset27}[\emph{registerZ}] = \emph{registerT}

\paragraph{Description} The effective address is computed by adding the \emph{extend27\_\discretionary{}{}{}offset27} to the \emph{registerZ}.
This instruction implements atomic MAX-store on signed double word. The double word at effective address is MAXed with the \emph{registerT}. The effective address must be double word aligned (multiple of 8).


\paragraph{Modifier(s)}
\noindent\begin{itemize}
\item coherency (cf. coherency in section \ref{par:modifier-coherency} on page \pageref{par:modifier-coherency})
\end{itemize}

\paragraph{Execution} {Reservation table \textsf{LSU2\_MEMW\_AUXR.Y} (Section~\ref{sec:reservation-LSU2-MEMW-AUXR.Y})}
~
{\begin{lstlisting}
stage ID:
  new offset = _SX_54(extend27_offset27);
  new coherency = _ZX_2(coherency);
stage RR:
  new base = GPR[registerZ];
  new address = base + offset;
stage E1:
  new argument2 = GPR[registerT];
stage E3:
  MEM.atomic_max(address, 0xFF, coherency << 32, argument2.64[0], registerT);
\end{lstlisting}}
\newpage
\subsubsection[{\small ASMAXH: Atomic Store MAX Half Word}]{ASMAXH\hfill Atomic Store MAX Half Word}
\label{instr:ASMAXH}\index{asmaxh}
 
\paragraph{Encoding} Format \textsf{LSU\_ASMAXSB} (Section~\ref{sec:format-LSU-ASMAXSB}), \verb|lsucode5 = 001|\\

\bitfields{ 1/P, 2/01, 3/111, 2/coherency, 6/registerT, 2/11, 3/001, 1/--, 5/11001, 1/0, 6/registerZ }


\paragraph{Syntax} \texttt{asmaxh}\emph{coherency} [\emph{registerZ}] = \emph{registerT}

\paragraph{Description} The effective address is the value of the \emph{registerZ}.
This instruction implements atomic MAX-store on signed half word. The half word at effective address is MAXed with the \emph{registerT}. The effective address must be half word aligned (multiple of 2).


\paragraph{Modifier(s)}
\noindent\begin{itemize}
\item coherency (cf. coherency in section \ref{par:modifier-coherency} on page \pageref{par:modifier-coherency})
\end{itemize}

\paragraph{Execution} {Reservation table \textsf{LSU2\_MEMW\_AUXR} (Section~\ref{sec:reservation-LSU2-MEMW-AUXR})}
~
{\begin{lstlisting}
stage ID:
  new coherency = _ZX_2(coherency);
stage RR:
  new address = GPR[registerZ];
stage E1:
  new argument2 = GPR[registerT];
stage E3:
  MEM.atomic_max(address, 0x3, coherency << 32, argument2.16[0], registerT);
\end{lstlisting}}
\newpage

\paragraph{Encoding} Format \textsf{LSU\_ASMAXSB.O} (Section~\ref{sec:format-LSU-ASMAXSB.O}), \verb|lsucode5 = 001|\\

\bitfields{ 1/P, 2/00, 2/-- --, 27/offset27, 1/1, 2/01, 3/111, 2/coherency, 6/registerT, 2/11, 3/001, 1/--, 5/11001, 1/0, 6/registerZ }


\paragraph{Syntax} \texttt{asmaxh}\emph{coherency} \emph{offset27}[\emph{registerZ}] = \emph{registerT}

\paragraph{Description} The effective address is computed by adding the \emph{offset27} to the \emph{registerZ}.
This instruction implements atomic MAX-store on signed half word. The half word at effective address is MAXed with the \emph{registerT}. The effective address must be half word aligned (multiple of 2).


\paragraph{Modifier(s)}
\noindent\begin{itemize}
\item coherency (cf. coherency in section \ref{par:modifier-coherency} on page \pageref{par:modifier-coherency})
\end{itemize}

\paragraph{Execution} {Reservation table \textsf{LSU2\_MEMW\_AUXR.X} (Section~\ref{sec:reservation-LSU2-MEMW-AUXR.X})}
~
{\begin{lstlisting}
stage ID:
  new offset = _SX_27(offset27);
  new coherency = _ZX_2(coherency);
stage RR:
  new base = GPR[registerZ];
  new address = base + offset;
stage E1:
  new argument2 = GPR[registerT];
stage E3:
  MEM.atomic_max(address, 0x3, coherency << 32, argument2.16[0], registerT);
\end{lstlisting}}
\newpage

\paragraph{Encoding} Format \textsf{LSU\_ASMAXSB.D} (Section~\ref{sec:format-LSU-ASMAXSB.D}), \verb|lsucode5 = 001|\\

\bitfields{ 1/P, 2/00, 2/-- --, 27/extend27, 1/1, 2/00, 2/-- --, 27/offset27, 1/1, 2/01, 3/111, 2/coherency, 6/registerT, 2/11, 3/001, 1/--, 5/11001, 1/0, 6/registerZ }


\paragraph{Syntax} \texttt{asmaxh}\emph{coherency} \emph{extend27\_\discretionary{}{}{}offset27}[\emph{registerZ}] = \emph{registerT}

\paragraph{Description} The effective address is computed by adding the \emph{extend27\_\discretionary{}{}{}offset27} to the \emph{registerZ}.
This instruction implements atomic MAX-store on signed half word. The half word at effective address is MAXed with the \emph{registerT}. The effective address must be half word aligned (multiple of 2).


\paragraph{Modifier(s)}
\noindent\begin{itemize}
\item coherency (cf. coherency in section \ref{par:modifier-coherency} on page \pageref{par:modifier-coherency})
\end{itemize}

\paragraph{Execution} {Reservation table \textsf{LSU2\_MEMW\_AUXR.Y} (Section~\ref{sec:reservation-LSU2-MEMW-AUXR.Y})}
~
{\begin{lstlisting}
stage ID:
  new offset = _SX_54(extend27_offset27);
  new coherency = _ZX_2(coherency);
stage RR:
  new base = GPR[registerZ];
  new address = base + offset;
stage E1:
  new argument2 = GPR[registerT];
stage E3:
  MEM.atomic_max(address, 0x3, coherency << 32, argument2.16[0], registerT);
\end{lstlisting}}
\newpage
\subsubsection[{\small ASMAXUB: Atomic Store MAX Unsigned Byte}]{ASMAXUB\hfill Atomic Store MAX Unsigned Byte}
\label{instr:ASMAXUB}\index{asmaxub}
 
\paragraph{Encoding} Format \textsf{LSU\_ASMAXUSB} (Section~\ref{sec:format-LSU-ASMAXUSB}), \verb|lsucode5 = 000|\\

\bitfields{ 1/P, 2/01, 3/111, 2/coherency, 6/registerT, 2/11, 3/000, 1/--, 5/11011, 1/0, 6/registerZ }


\paragraph{Syntax} \texttt{asmaxub}\emph{coherency} [\emph{registerZ}] = \emph{registerT}

\paragraph{Description} The effective address is the value of the \emph{registerZ}.
This instruction implements atomic MAXU-store on signed byte. The byte at effective address is MAXUed with the \emph{registerT}.


\paragraph{Modifier(s)}
\noindent\begin{itemize}
\item coherency (cf. coherency in section \ref{par:modifier-coherency} on page \pageref{par:modifier-coherency})
\end{itemize}

\paragraph{Execution} {Reservation table \textsf{LSU2\_MEMW\_AUXR} (Section~\ref{sec:reservation-LSU2-MEMW-AUXR})}
~
{\begin{lstlisting}
stage ID:
  new coherency = _ZX_2(coherency);
stage RR:
  new address = GPR[registerZ];
stage E1:
  new argument2 = GPR[registerT];
stage E3:
  MEM.atomic_maxu(address, 0x1, coherency << 32, argument2.8[0], registerT);
\end{lstlisting}}
\newpage

\paragraph{Encoding} Format \textsf{LSU\_ASMAXUSB.O} (Section~\ref{sec:format-LSU-ASMAXUSB.O}), \verb|lsucode5 = 000|\\

\bitfields{ 1/P, 2/00, 2/-- --, 27/offset27, 1/1, 2/01, 3/111, 2/coherency, 6/registerT, 2/11, 3/000, 1/--, 5/11011, 1/0, 6/registerZ }


\paragraph{Syntax} \texttt{asmaxub}\emph{coherency} \emph{offset27}[\emph{registerZ}] = \emph{registerT}

\paragraph{Description} The effective address is computed by adding the \emph{offset27} to the \emph{registerZ}.
This instruction implements atomic MAXU-store on signed byte. The byte at effective address is MAXUed with the \emph{registerT}.


\paragraph{Modifier(s)}
\noindent\begin{itemize}
\item coherency (cf. coherency in section \ref{par:modifier-coherency} on page \pageref{par:modifier-coherency})
\end{itemize}

\paragraph{Execution} {Reservation table \textsf{LSU2\_MEMW\_AUXR.X} (Section~\ref{sec:reservation-LSU2-MEMW-AUXR.X})}
~
{\begin{lstlisting}
stage ID:
  new offset = _SX_27(offset27);
  new coherency = _ZX_2(coherency);
stage RR:
  new base = GPR[registerZ];
  new address = base + offset;
stage E1:
  new argument2 = GPR[registerT];
stage E3:
  MEM.atomic_maxu(address, 0x1, coherency << 32, argument2.8[0], registerT);
\end{lstlisting}}
\newpage

\paragraph{Encoding} Format \textsf{LSU\_ASMAXUSB.D} (Section~\ref{sec:format-LSU-ASMAXUSB.D}), \verb|lsucode5 = 000|\\

\bitfields{ 1/P, 2/00, 2/-- --, 27/extend27, 1/1, 2/00, 2/-- --, 27/offset27, 1/1, 2/01, 3/111, 2/coherency, 6/registerT, 2/11, 3/000, 1/--, 5/11011, 1/0, 6/registerZ }


\paragraph{Syntax} \texttt{asmaxub}\emph{coherency} \emph{extend27\_\discretionary{}{}{}offset27}[\emph{registerZ}] = \emph{registerT}

\paragraph{Description} The effective address is computed by adding the \emph{extend27\_\discretionary{}{}{}offset27} to the \emph{registerZ}.
This instruction implements atomic MAXU-store on signed byte. The byte at effective address is MAXUed with the \emph{registerT}.


\paragraph{Modifier(s)}
\noindent\begin{itemize}
\item coherency (cf. coherency in section \ref{par:modifier-coherency} on page \pageref{par:modifier-coherency})
\end{itemize}

\paragraph{Execution} {Reservation table \textsf{LSU2\_MEMW\_AUXR.Y} (Section~\ref{sec:reservation-LSU2-MEMW-AUXR.Y})}
~
{\begin{lstlisting}
stage ID:
  new offset = _SX_54(extend27_offset27);
  new coherency = _ZX_2(coherency);
stage RR:
  new base = GPR[registerZ];
  new address = base + offset;
stage E1:
  new argument2 = GPR[registerT];
stage E3:
  MEM.atomic_maxu(address, 0x1, coherency << 32, argument2.8[0], registerT);
\end{lstlisting}}
\newpage
\subsubsection[{\small ASMAXUD: Atomic Store MAX Unsigned Double Word}]{ASMAXUD\hfill Atomic Store MAX Unsigned Double Word}
\label{instr:ASMAXUD}\index{asmaxud}
 
\paragraph{Encoding} Format \textsf{LSU\_ASMAXUSB} (Section~\ref{sec:format-LSU-ASMAXUSB}), \verb|lsucode5 = 011|\\

\bitfields{ 1/P, 2/01, 3/111, 2/coherency, 6/registerT, 2/11, 3/011, 1/--, 5/11011, 1/0, 6/registerZ }


\paragraph{Syntax} \texttt{asmaxud}\emph{coherency} [\emph{registerZ}] = \emph{registerT}

\paragraph{Description} The effective address is the value of the \emph{registerZ}.
This instruction implements atomic MAXU-store on signed double word. The double word at effective address is MAXUed with the \emph{registerT}. The effective address must be double word aligned (multiple of 8).


\paragraph{Modifier(s)}
\noindent\begin{itemize}
\item coherency (cf. coherency in section \ref{par:modifier-coherency} on page \pageref{par:modifier-coherency})
\end{itemize}

\paragraph{Execution} {Reservation table \textsf{LSU2\_MEMW\_AUXR} (Section~\ref{sec:reservation-LSU2-MEMW-AUXR})}
~
{\begin{lstlisting}
stage ID:
  new coherency = _ZX_2(coherency);
stage RR:
  new address = GPR[registerZ];
stage E1:
  new argument2 = GPR[registerT];
stage E3:
  MEM.atomic_maxu(address, 0xFF, coherency << 32, argument2.64[0], registerT);
\end{lstlisting}}
\newpage

\paragraph{Encoding} Format \textsf{LSU\_ASMAXUSB.O} (Section~\ref{sec:format-LSU-ASMAXUSB.O}), \verb|lsucode5 = 011|\\

\bitfields{ 1/P, 2/00, 2/-- --, 27/offset27, 1/1, 2/01, 3/111, 2/coherency, 6/registerT, 2/11, 3/011, 1/--, 5/11011, 1/0, 6/registerZ }


\paragraph{Syntax} \texttt{asmaxud}\emph{coherency} \emph{offset27}[\emph{registerZ}] = \emph{registerT}

\paragraph{Description} The effective address is computed by adding the \emph{offset27} to the \emph{registerZ}.
This instruction implements atomic MAXU-store on signed double word. The double word at effective address is MAXUed with the \emph{registerT}. The effective address must be double word aligned (multiple of 8).


\paragraph{Modifier(s)}
\noindent\begin{itemize}
\item coherency (cf. coherency in section \ref{par:modifier-coherency} on page \pageref{par:modifier-coherency})
\end{itemize}

\paragraph{Execution} {Reservation table \textsf{LSU2\_MEMW\_AUXR.X} (Section~\ref{sec:reservation-LSU2-MEMW-AUXR.X})}
~
{\begin{lstlisting}
stage ID:
  new offset = _SX_27(offset27);
  new coherency = _ZX_2(coherency);
stage RR:
  new base = GPR[registerZ];
  new address = base + offset;
stage E1:
  new argument2 = GPR[registerT];
stage E3:
  MEM.atomic_maxu(address, 0xFF, coherency << 32, argument2.64[0], registerT);
\end{lstlisting}}
\newpage

\paragraph{Encoding} Format \textsf{LSU\_ASMAXUSB.D} (Section~\ref{sec:format-LSU-ASMAXUSB.D}), \verb|lsucode5 = 011|\\

\bitfields{ 1/P, 2/00, 2/-- --, 27/extend27, 1/1, 2/00, 2/-- --, 27/offset27, 1/1, 2/01, 3/111, 2/coherency, 6/registerT, 2/11, 3/011, 1/--, 5/11011, 1/0, 6/registerZ }


\paragraph{Syntax} \texttt{asmaxud}\emph{coherency} \emph{extend27\_\discretionary{}{}{}offset27}[\emph{registerZ}] = \emph{registerT}

\paragraph{Description} The effective address is computed by adding the \emph{extend27\_\discretionary{}{}{}offset27} to the \emph{registerZ}.
This instruction implements atomic MAXU-store on signed double word. The double word at effective address is MAXUed with the \emph{registerT}. The effective address must be double word aligned (multiple of 8).


\paragraph{Modifier(s)}
\noindent\begin{itemize}
\item coherency (cf. coherency in section \ref{par:modifier-coherency} on page \pageref{par:modifier-coherency})
\end{itemize}

\paragraph{Execution} {Reservation table \textsf{LSU2\_MEMW\_AUXR.Y} (Section~\ref{sec:reservation-LSU2-MEMW-AUXR.Y})}
~
{\begin{lstlisting}
stage ID:
  new offset = _SX_54(extend27_offset27);
  new coherency = _ZX_2(coherency);
stage RR:
  new base = GPR[registerZ];
  new address = base + offset;
stage E1:
  new argument2 = GPR[registerT];
stage E3:
  MEM.atomic_maxu(address, 0xFF, coherency << 32, argument2.64[0], registerT);
\end{lstlisting}}
\newpage
\subsubsection[{\small ASMAXUH: Atomic Store MAX Unsigned Half Word}]{ASMAXUH\hfill Atomic Store MAX Unsigned Half Word}
\label{instr:ASMAXUH}\index{asmaxuh}
 
\paragraph{Encoding} Format \textsf{LSU\_ASMAXUSB} (Section~\ref{sec:format-LSU-ASMAXUSB}), \verb|lsucode5 = 001|\\

\bitfields{ 1/P, 2/01, 3/111, 2/coherency, 6/registerT, 2/11, 3/001, 1/--, 5/11011, 1/0, 6/registerZ }


\paragraph{Syntax} \texttt{asmaxuh}\emph{coherency} [\emph{registerZ}] = \emph{registerT}

\paragraph{Description} The effective address is the value of the \emph{registerZ}.
This instruction implements atomic MAXU-store on signed half word. The half word at effective address is MAXUed with the \emph{registerT}. The effective address must be half word aligned (multiple of 2).


\paragraph{Modifier(s)}
\noindent\begin{itemize}
\item coherency (cf. coherency in section \ref{par:modifier-coherency} on page \pageref{par:modifier-coherency})
\end{itemize}

\paragraph{Execution} {Reservation table \textsf{LSU2\_MEMW\_AUXR} (Section~\ref{sec:reservation-LSU2-MEMW-AUXR})}
~
{\begin{lstlisting}
stage ID:
  new coherency = _ZX_2(coherency);
stage RR:
  new address = GPR[registerZ];
stage E1:
  new argument2 = GPR[registerT];
stage E3:
  MEM.atomic_maxu(address, 0x3, coherency << 32, argument2.16[0], registerT);
\end{lstlisting}}
\newpage

\paragraph{Encoding} Format \textsf{LSU\_ASMAXUSB.O} (Section~\ref{sec:format-LSU-ASMAXUSB.O}), \verb|lsucode5 = 001|\\

\bitfields{ 1/P, 2/00, 2/-- --, 27/offset27, 1/1, 2/01, 3/111, 2/coherency, 6/registerT, 2/11, 3/001, 1/--, 5/11011, 1/0, 6/registerZ }


\paragraph{Syntax} \texttt{asmaxuh}\emph{coherency} \emph{offset27}[\emph{registerZ}] = \emph{registerT}

\paragraph{Description} The effective address is computed by adding the \emph{offset27} to the \emph{registerZ}.
This instruction implements atomic MAXU-store on signed half word. The half word at effective address is MAXUed with the \emph{registerT}. The effective address must be half word aligned (multiple of 2).


\paragraph{Modifier(s)}
\noindent\begin{itemize}
\item coherency (cf. coherency in section \ref{par:modifier-coherency} on page \pageref{par:modifier-coherency})
\end{itemize}

\paragraph{Execution} {Reservation table \textsf{LSU2\_MEMW\_AUXR.X} (Section~\ref{sec:reservation-LSU2-MEMW-AUXR.X})}
~
{\begin{lstlisting}
stage ID:
  new offset = _SX_27(offset27);
  new coherency = _ZX_2(coherency);
stage RR:
  new base = GPR[registerZ];
  new address = base + offset;
stage E1:
  new argument2 = GPR[registerT];
stage E3:
  MEM.atomic_maxu(address, 0x3, coherency << 32, argument2.16[0], registerT);
\end{lstlisting}}
\newpage

\paragraph{Encoding} Format \textsf{LSU\_ASMAXUSB.D} (Section~\ref{sec:format-LSU-ASMAXUSB.D}), \verb|lsucode5 = 001|\\

\bitfields{ 1/P, 2/00, 2/-- --, 27/extend27, 1/1, 2/00, 2/-- --, 27/offset27, 1/1, 2/01, 3/111, 2/coherency, 6/registerT, 2/11, 3/001, 1/--, 5/11011, 1/0, 6/registerZ }


\paragraph{Syntax} \texttt{asmaxuh}\emph{coherency} \emph{extend27\_\discretionary{}{}{}offset27}[\emph{registerZ}] = \emph{registerT}

\paragraph{Description} The effective address is computed by adding the \emph{extend27\_\discretionary{}{}{}offset27} to the \emph{registerZ}.
This instruction implements atomic MAXU-store on signed half word. The half word at effective address is MAXUed with the \emph{registerT}. The effective address must be half word aligned (multiple of 2).


\paragraph{Modifier(s)}
\noindent\begin{itemize}
\item coherency (cf. coherency in section \ref{par:modifier-coherency} on page \pageref{par:modifier-coherency})
\end{itemize}

\paragraph{Execution} {Reservation table \textsf{LSU2\_MEMW\_AUXR.Y} (Section~\ref{sec:reservation-LSU2-MEMW-AUXR.Y})}
~
{\begin{lstlisting}
stage ID:
  new offset = _SX_54(extend27_offset27);
  new coherency = _ZX_2(coherency);
stage RR:
  new base = GPR[registerZ];
  new address = base + offset;
stage E1:
  new argument2 = GPR[registerT];
stage E3:
  MEM.atomic_maxu(address, 0x3, coherency << 32, argument2.16[0], registerT);
\end{lstlisting}}
\newpage
\subsubsection[{\small ASMAXUW: Atomic Store MAX Unsigned Word}]{ASMAXUW\hfill Atomic Store MAX Unsigned Word}
\label{instr:ASMAXUW}\index{asmaxuw}
 
\paragraph{Encoding} Format \textsf{LSU\_ASMAXUSB} (Section~\ref{sec:format-LSU-ASMAXUSB}), \verb|lsucode5 = 010|\\

\bitfields{ 1/P, 2/01, 3/111, 2/coherency, 6/registerT, 2/11, 3/010, 1/--, 5/11011, 1/0, 6/registerZ }


\paragraph{Syntax} \texttt{asmaxuw}\emph{coherency} [\emph{registerZ}] = \emph{registerT}

\paragraph{Description} The effective address is the value of the \emph{registerZ}.
This instruction implements atomic MAXU-store on signed word. The word at effective address is MAXUed with the \emph{registerT}. The effective address must be word aligned (multiple of 4).


\paragraph{Modifier(s)}
\noindent\begin{itemize}
\item coherency (cf. coherency in section \ref{par:modifier-coherency} on page \pageref{par:modifier-coherency})
\end{itemize}

\paragraph{Execution} {Reservation table \textsf{LSU2\_MEMW\_AUXR} (Section~\ref{sec:reservation-LSU2-MEMW-AUXR})}
~
{\begin{lstlisting}
stage ID:
  new coherency = _ZX_2(coherency);
stage RR:
  new address = GPR[registerZ];
stage E1:
  new argument2 = GPR[registerT];
stage E3:
  MEM.atomic_maxu(address, 0xF, coherency << 32, argument2.32[0], registerT);
\end{lstlisting}}
\newpage

\paragraph{Encoding} Format \textsf{LSU\_ASMAXUSB.O} (Section~\ref{sec:format-LSU-ASMAXUSB.O}), \verb|lsucode5 = 010|\\

\bitfields{ 1/P, 2/00, 2/-- --, 27/offset27, 1/1, 2/01, 3/111, 2/coherency, 6/registerT, 2/11, 3/010, 1/--, 5/11011, 1/0, 6/registerZ }


\paragraph{Syntax} \texttt{asmaxuw}\emph{coherency} \emph{offset27}[\emph{registerZ}] = \emph{registerT}

\paragraph{Description} The effective address is computed by adding the \emph{offset27} to the \emph{registerZ}.
This instruction implements atomic MAXU-store on signed word. The word at effective address is MAXUed with the \emph{registerT}. The effective address must be word aligned (multiple of 4).


\paragraph{Modifier(s)}
\noindent\begin{itemize}
\item coherency (cf. coherency in section \ref{par:modifier-coherency} on page \pageref{par:modifier-coherency})
\end{itemize}

\paragraph{Execution} {Reservation table \textsf{LSU2\_MEMW\_AUXR.X} (Section~\ref{sec:reservation-LSU2-MEMW-AUXR.X})}
~
{\begin{lstlisting}
stage ID:
  new offset = _SX_27(offset27);
  new coherency = _ZX_2(coherency);
stage RR:
  new base = GPR[registerZ];
  new address = base + offset;
stage E1:
  new argument2 = GPR[registerT];
stage E3:
  MEM.atomic_maxu(address, 0xF, coherency << 32, argument2.32[0], registerT);
\end{lstlisting}}
\newpage

\paragraph{Encoding} Format \textsf{LSU\_ASMAXUSB.D} (Section~\ref{sec:format-LSU-ASMAXUSB.D}), \verb|lsucode5 = 010|\\

\bitfields{ 1/P, 2/00, 2/-- --, 27/extend27, 1/1, 2/00, 2/-- --, 27/offset27, 1/1, 2/01, 3/111, 2/coherency, 6/registerT, 2/11, 3/010, 1/--, 5/11011, 1/0, 6/registerZ }


\paragraph{Syntax} \texttt{asmaxuw}\emph{coherency} \emph{extend27\_\discretionary{}{}{}offset27}[\emph{registerZ}] = \emph{registerT}

\paragraph{Description} The effective address is computed by adding the \emph{extend27\_\discretionary{}{}{}offset27} to the \emph{registerZ}.
This instruction implements atomic MAXU-store on signed word. The word at effective address is MAXUed with the \emph{registerT}. The effective address must be word aligned (multiple of 4).


\paragraph{Modifier(s)}
\noindent\begin{itemize}
\item coherency (cf. coherency in section \ref{par:modifier-coherency} on page \pageref{par:modifier-coherency})
\end{itemize}

\paragraph{Execution} {Reservation table \textsf{LSU2\_MEMW\_AUXR.Y} (Section~\ref{sec:reservation-LSU2-MEMW-AUXR.Y})}
~
{\begin{lstlisting}
stage ID:
  new offset = _SX_54(extend27_offset27);
  new coherency = _ZX_2(coherency);
stage RR:
  new base = GPR[registerZ];
  new address = base + offset;
stage E1:
  new argument2 = GPR[registerT];
stage E3:
  MEM.atomic_maxu(address, 0xF, coherency << 32, argument2.32[0], registerT);
\end{lstlisting}}
\newpage
\subsubsection[{\small ASMAXW: Atomic Store MAX Word}]{ASMAXW\hfill Atomic Store MAX Word}
\label{instr:ASMAXW}\index{asmaxw}
 
\paragraph{Encoding} Format \textsf{LSU\_ASMAXSB} (Section~\ref{sec:format-LSU-ASMAXSB}), \verb|lsucode5 = 010|\\

\bitfields{ 1/P, 2/01, 3/111, 2/coherency, 6/registerT, 2/11, 3/010, 1/--, 5/11001, 1/0, 6/registerZ }


\paragraph{Syntax} \texttt{asmaxw}\emph{coherency} [\emph{registerZ}] = \emph{registerT}

\paragraph{Description} The effective address is the value of the \emph{registerZ}.
This instruction implements atomic MAX-store on signed word. The word at effective address is MAXed with the \emph{registerT}. The effective address must be word aligned (multiple of 4).


\paragraph{Modifier(s)}
\noindent\begin{itemize}
\item coherency (cf. coherency in section \ref{par:modifier-coherency} on page \pageref{par:modifier-coherency})
\end{itemize}

\paragraph{Execution} {Reservation table \textsf{LSU2\_MEMW\_AUXR} (Section~\ref{sec:reservation-LSU2-MEMW-AUXR})}
~
{\begin{lstlisting}
stage ID:
  new coherency = _ZX_2(coherency);
stage RR:
  new address = GPR[registerZ];
stage E1:
  new argument2 = GPR[registerT];
stage E3:
  MEM.atomic_max(address, 0xF, coherency << 32, argument2.32[0], registerT);
\end{lstlisting}}
\newpage

\paragraph{Encoding} Format \textsf{LSU\_ASMAXSB.O} (Section~\ref{sec:format-LSU-ASMAXSB.O}), \verb|lsucode5 = 010|\\

\bitfields{ 1/P, 2/00, 2/-- --, 27/offset27, 1/1, 2/01, 3/111, 2/coherency, 6/registerT, 2/11, 3/010, 1/--, 5/11001, 1/0, 6/registerZ }


\paragraph{Syntax} \texttt{asmaxw}\emph{coherency} \emph{offset27}[\emph{registerZ}] = \emph{registerT}

\paragraph{Description} The effective address is computed by adding the \emph{offset27} to the \emph{registerZ}.
This instruction implements atomic MAX-store on signed word. The word at effective address is MAXed with the \emph{registerT}. The effective address must be word aligned (multiple of 4).


\paragraph{Modifier(s)}
\noindent\begin{itemize}
\item coherency (cf. coherency in section \ref{par:modifier-coherency} on page \pageref{par:modifier-coherency})
\end{itemize}

\paragraph{Execution} {Reservation table \textsf{LSU2\_MEMW\_AUXR.X} (Section~\ref{sec:reservation-LSU2-MEMW-AUXR.X})}
~
{\begin{lstlisting}
stage ID:
  new offset = _SX_27(offset27);
  new coherency = _ZX_2(coherency);
stage RR:
  new base = GPR[registerZ];
  new address = base + offset;
stage E1:
  new argument2 = GPR[registerT];
stage E3:
  MEM.atomic_max(address, 0xF, coherency << 32, argument2.32[0], registerT);
\end{lstlisting}}
\newpage

\paragraph{Encoding} Format \textsf{LSU\_ASMAXSB.D} (Section~\ref{sec:format-LSU-ASMAXSB.D}), \verb|lsucode5 = 010|\\

\bitfields{ 1/P, 2/00, 2/-- --, 27/extend27, 1/1, 2/00, 2/-- --, 27/offset27, 1/1, 2/01, 3/111, 2/coherency, 6/registerT, 2/11, 3/010, 1/--, 5/11001, 1/0, 6/registerZ }


\paragraph{Syntax} \texttt{asmaxw}\emph{coherency} \emph{extend27\_\discretionary{}{}{}offset27}[\emph{registerZ}] = \emph{registerT}

\paragraph{Description} The effective address is computed by adding the \emph{extend27\_\discretionary{}{}{}offset27} to the \emph{registerZ}.
This instruction implements atomic MAX-store on signed word. The word at effective address is MAXed with the \emph{registerT}. The effective address must be word aligned (multiple of 4).


\paragraph{Modifier(s)}
\noindent\begin{itemize}
\item coherency (cf. coherency in section \ref{par:modifier-coherency} on page \pageref{par:modifier-coherency})
\end{itemize}

\paragraph{Execution} {Reservation table \textsf{LSU2\_MEMW\_AUXR.Y} (Section~\ref{sec:reservation-LSU2-MEMW-AUXR.Y})}
~
{\begin{lstlisting}
stage ID:
  new offset = _SX_54(extend27_offset27);
  new coherency = _ZX_2(coherency);
stage RR:
  new base = GPR[registerZ];
  new address = base + offset;
stage E1:
  new argument2 = GPR[registerT];
stage E3:
  MEM.atomic_max(address, 0xF, coherency << 32, argument2.32[0], registerT);
\end{lstlisting}}
\newpage
\subsubsection[{\small ASMINB: Atomic Store MIN Byte}]{ASMINB\hfill Atomic Store MIN Byte}
\label{instr:ASMINB}\index{asminb}
 
\paragraph{Encoding} Format \textsf{LSU\_ASMINSB} (Section~\ref{sec:format-LSU-ASMINSB}), \verb|lsucode5 = 000|\\

\bitfields{ 1/P, 2/01, 3/111, 2/coherency, 6/registerT, 2/11, 3/000, 1/--, 5/11000, 1/0, 6/registerZ }


\paragraph{Syntax} \texttt{asminb}\emph{coherency} [\emph{registerZ}] = \emph{registerT}

\paragraph{Description} The effective address is the value of the \emph{registerZ}.
This instruction implements atomic MIN-store on signed byte. The byte at effective address is MINed with the \emph{registerT}.


\paragraph{Modifier(s)}
\noindent\begin{itemize}
\item coherency (cf. coherency in section \ref{par:modifier-coherency} on page \pageref{par:modifier-coherency})
\end{itemize}

\paragraph{Execution} {Reservation table \textsf{LSU2\_MEMW\_AUXR} (Section~\ref{sec:reservation-LSU2-MEMW-AUXR})}
~
{\begin{lstlisting}
stage ID:
  new coherency = _ZX_2(coherency);
stage RR:
  new address = GPR[registerZ];
stage E1:
  new argument2 = GPR[registerT];
stage E3:
  MEM.atomic_min(address, 0x1, coherency << 32, argument2.8[0], registerT);
\end{lstlisting}}
\newpage

\paragraph{Encoding} Format \textsf{LSU\_ASMINSB.O} (Section~\ref{sec:format-LSU-ASMINSB.O}), \verb|lsucode5 = 000|\\

\bitfields{ 1/P, 2/00, 2/-- --, 27/offset27, 1/1, 2/01, 3/111, 2/coherency, 6/registerT, 2/11, 3/000, 1/--, 5/11000, 1/0, 6/registerZ }


\paragraph{Syntax} \texttt{asminb}\emph{coherency} \emph{offset27}[\emph{registerZ}] = \emph{registerT}

\paragraph{Description} The effective address is computed by adding the \emph{offset27} to the \emph{registerZ}.
This instruction implements atomic MIN-store on signed byte. The byte at effective address is MINed with the \emph{registerT}.


\paragraph{Modifier(s)}
\noindent\begin{itemize}
\item coherency (cf. coherency in section \ref{par:modifier-coherency} on page \pageref{par:modifier-coherency})
\end{itemize}

\paragraph{Execution} {Reservation table \textsf{LSU2\_MEMW\_AUXR.X} (Section~\ref{sec:reservation-LSU2-MEMW-AUXR.X})}
~
{\begin{lstlisting}
stage ID:
  new offset = _SX_27(offset27);
  new coherency = _ZX_2(coherency);
stage RR:
  new base = GPR[registerZ];
  new address = base + offset;
stage E1:
  new argument2 = GPR[registerT];
stage E3:
  MEM.atomic_min(address, 0x1, coherency << 32, argument2.8[0], registerT);
\end{lstlisting}}
\newpage

\paragraph{Encoding} Format \textsf{LSU\_ASMINSB.D} (Section~\ref{sec:format-LSU-ASMINSB.D}), \verb|lsucode5 = 000|\\

\bitfields{ 1/P, 2/00, 2/-- --, 27/extend27, 1/1, 2/00, 2/-- --, 27/offset27, 1/1, 2/01, 3/111, 2/coherency, 6/registerT, 2/11, 3/000, 1/--, 5/11000, 1/0, 6/registerZ }


\paragraph{Syntax} \texttt{asminb}\emph{coherency} \emph{extend27\_\discretionary{}{}{}offset27}[\emph{registerZ}] = \emph{registerT}

\paragraph{Description} The effective address is computed by adding the \emph{extend27\_\discretionary{}{}{}offset27} to the \emph{registerZ}.
This instruction implements atomic MIN-store on signed byte. The byte at effective address is MINed with the \emph{registerT}.


\paragraph{Modifier(s)}
\noindent\begin{itemize}
\item coherency (cf. coherency in section \ref{par:modifier-coherency} on page \pageref{par:modifier-coherency})
\end{itemize}

\paragraph{Execution} {Reservation table \textsf{LSU2\_MEMW\_AUXR.Y} (Section~\ref{sec:reservation-LSU2-MEMW-AUXR.Y})}
~
{\begin{lstlisting}
stage ID:
  new offset = _SX_54(extend27_offset27);
  new coherency = _ZX_2(coherency);
stage RR:
  new base = GPR[registerZ];
  new address = base + offset;
stage E1:
  new argument2 = GPR[registerT];
stage E3:
  MEM.atomic_min(address, 0x1, coherency << 32, argument2.8[0], registerT);
\end{lstlisting}}
\newpage
\subsubsection[{\small ASMIND: Atomic Store MIN Double Word}]{ASMIND\hfill Atomic Store MIN Double Word}
\label{instr:ASMIND}\index{asmind}
 
\paragraph{Encoding} Format \textsf{LSU\_ASMINSB} (Section~\ref{sec:format-LSU-ASMINSB}), \verb|lsucode5 = 011|\\

\bitfields{ 1/P, 2/01, 3/111, 2/coherency, 6/registerT, 2/11, 3/011, 1/--, 5/11000, 1/0, 6/registerZ }


\paragraph{Syntax} \texttt{asmind}\emph{coherency} [\emph{registerZ}] = \emph{registerT}

\paragraph{Description} The effective address is the value of the \emph{registerZ}.
This instruction implements atomic MIN-store on signed double word. The double word at effective address is MINed with the \emph{registerT}. The effective address must be double word aligned (multiple of 8).


\paragraph{Modifier(s)}
\noindent\begin{itemize}
\item coherency (cf. coherency in section \ref{par:modifier-coherency} on page \pageref{par:modifier-coherency})
\end{itemize}

\paragraph{Execution} {Reservation table \textsf{LSU2\_MEMW\_AUXR} (Section~\ref{sec:reservation-LSU2-MEMW-AUXR})}
~
{\begin{lstlisting}
stage ID:
  new coherency = _ZX_2(coherency);
stage RR:
  new address = GPR[registerZ];
stage E1:
  new argument2 = GPR[registerT];
stage E3:
  MEM.atomic_min(address, 0xFF, coherency << 32, argument2.64[0], registerT);
\end{lstlisting}}
\newpage

\paragraph{Encoding} Format \textsf{LSU\_ASMINSB.O} (Section~\ref{sec:format-LSU-ASMINSB.O}), \verb|lsucode5 = 011|\\

\bitfields{ 1/P, 2/00, 2/-- --, 27/offset27, 1/1, 2/01, 3/111, 2/coherency, 6/registerT, 2/11, 3/011, 1/--, 5/11000, 1/0, 6/registerZ }


\paragraph{Syntax} \texttt{asmind}\emph{coherency} \emph{offset27}[\emph{registerZ}] = \emph{registerT}

\paragraph{Description} The effective address is computed by adding the \emph{offset27} to the \emph{registerZ}.
This instruction implements atomic MIN-store on signed double word. The double word at effective address is MINed with the \emph{registerT}. The effective address must be double word aligned (multiple of 8).


\paragraph{Modifier(s)}
\noindent\begin{itemize}
\item coherency (cf. coherency in section \ref{par:modifier-coherency} on page \pageref{par:modifier-coherency})
\end{itemize}

\paragraph{Execution} {Reservation table \textsf{LSU2\_MEMW\_AUXR.X} (Section~\ref{sec:reservation-LSU2-MEMW-AUXR.X})}
~
{\begin{lstlisting}
stage ID:
  new offset = _SX_27(offset27);
  new coherency = _ZX_2(coherency);
stage RR:
  new base = GPR[registerZ];
  new address = base + offset;
stage E1:
  new argument2 = GPR[registerT];
stage E3:
  MEM.atomic_min(address, 0xFF, coherency << 32, argument2.64[0], registerT);
\end{lstlisting}}
\newpage

\paragraph{Encoding} Format \textsf{LSU\_ASMINSB.D} (Section~\ref{sec:format-LSU-ASMINSB.D}), \verb|lsucode5 = 011|\\

\bitfields{ 1/P, 2/00, 2/-- --, 27/extend27, 1/1, 2/00, 2/-- --, 27/offset27, 1/1, 2/01, 3/111, 2/coherency, 6/registerT, 2/11, 3/011, 1/--, 5/11000, 1/0, 6/registerZ }


\paragraph{Syntax} \texttt{asmind}\emph{coherency} \emph{extend27\_\discretionary{}{}{}offset27}[\emph{registerZ}] = \emph{registerT}

\paragraph{Description} The effective address is computed by adding the \emph{extend27\_\discretionary{}{}{}offset27} to the \emph{registerZ}.
This instruction implements atomic MIN-store on signed double word. The double word at effective address is MINed with the \emph{registerT}. The effective address must be double word aligned (multiple of 8).


\paragraph{Modifier(s)}
\noindent\begin{itemize}
\item coherency (cf. coherency in section \ref{par:modifier-coherency} on page \pageref{par:modifier-coherency})
\end{itemize}

\paragraph{Execution} {Reservation table \textsf{LSU2\_MEMW\_AUXR.Y} (Section~\ref{sec:reservation-LSU2-MEMW-AUXR.Y})}
~
{\begin{lstlisting}
stage ID:
  new offset = _SX_54(extend27_offset27);
  new coherency = _ZX_2(coherency);
stage RR:
  new base = GPR[registerZ];
  new address = base + offset;
stage E1:
  new argument2 = GPR[registerT];
stage E3:
  MEM.atomic_min(address, 0xFF, coherency << 32, argument2.64[0], registerT);
\end{lstlisting}}
\newpage
\subsubsection[{\small ASMINH: Atomic Store MIN Half Word}]{ASMINH\hfill Atomic Store MIN Half Word}
\label{instr:ASMINH}\index{asminh}
 
\paragraph{Encoding} Format \textsf{LSU\_ASMINSB} (Section~\ref{sec:format-LSU-ASMINSB}), \verb|lsucode5 = 001|\\

\bitfields{ 1/P, 2/01, 3/111, 2/coherency, 6/registerT, 2/11, 3/001, 1/--, 5/11000, 1/0, 6/registerZ }


\paragraph{Syntax} \texttt{asminh}\emph{coherency} [\emph{registerZ}] = \emph{registerT}

\paragraph{Description} The effective address is the value of the \emph{registerZ}.
This instruction implements atomic MIN-store on signed half word. The half word at effective address is MINed with the \emph{registerT}. The effective address must be half word aligned (multiple of 2).


\paragraph{Modifier(s)}
\noindent\begin{itemize}
\item coherency (cf. coherency in section \ref{par:modifier-coherency} on page \pageref{par:modifier-coherency})
\end{itemize}

\paragraph{Execution} {Reservation table \textsf{LSU2\_MEMW\_AUXR} (Section~\ref{sec:reservation-LSU2-MEMW-AUXR})}
~
{\begin{lstlisting}
stage ID:
  new coherency = _ZX_2(coherency);
stage RR:
  new address = GPR[registerZ];
stage E1:
  new argument2 = GPR[registerT];
stage E3:
  MEM.atomic_min(address, 0x3, coherency << 32, argument2.16[0], registerT);
\end{lstlisting}}
\newpage

\paragraph{Encoding} Format \textsf{LSU\_ASMINSB.O} (Section~\ref{sec:format-LSU-ASMINSB.O}), \verb|lsucode5 = 001|\\

\bitfields{ 1/P, 2/00, 2/-- --, 27/offset27, 1/1, 2/01, 3/111, 2/coherency, 6/registerT, 2/11, 3/001, 1/--, 5/11000, 1/0, 6/registerZ }


\paragraph{Syntax} \texttt{asminh}\emph{coherency} \emph{offset27}[\emph{registerZ}] = \emph{registerT}

\paragraph{Description} The effective address is computed by adding the \emph{offset27} to the \emph{registerZ}.
This instruction implements atomic MIN-store on signed half word. The half word at effective address is MINed with the \emph{registerT}. The effective address must be half word aligned (multiple of 2).


\paragraph{Modifier(s)}
\noindent\begin{itemize}
\item coherency (cf. coherency in section \ref{par:modifier-coherency} on page \pageref{par:modifier-coherency})
\end{itemize}

\paragraph{Execution} {Reservation table \textsf{LSU2\_MEMW\_AUXR.X} (Section~\ref{sec:reservation-LSU2-MEMW-AUXR.X})}
~
{\begin{lstlisting}
stage ID:
  new offset = _SX_27(offset27);
  new coherency = _ZX_2(coherency);
stage RR:
  new base = GPR[registerZ];
  new address = base + offset;
stage E1:
  new argument2 = GPR[registerT];
stage E3:
  MEM.atomic_min(address, 0x3, coherency << 32, argument2.16[0], registerT);
\end{lstlisting}}
\newpage

\paragraph{Encoding} Format \textsf{LSU\_ASMINSB.D} (Section~\ref{sec:format-LSU-ASMINSB.D}), \verb|lsucode5 = 001|\\

\bitfields{ 1/P, 2/00, 2/-- --, 27/extend27, 1/1, 2/00, 2/-- --, 27/offset27, 1/1, 2/01, 3/111, 2/coherency, 6/registerT, 2/11, 3/001, 1/--, 5/11000, 1/0, 6/registerZ }


\paragraph{Syntax} \texttt{asminh}\emph{coherency} \emph{extend27\_\discretionary{}{}{}offset27}[\emph{registerZ}] = \emph{registerT}

\paragraph{Description} The effective address is computed by adding the \emph{extend27\_\discretionary{}{}{}offset27} to the \emph{registerZ}.
This instruction implements atomic MIN-store on signed half word. The half word at effective address is MINed with the \emph{registerT}. The effective address must be half word aligned (multiple of 2).


\paragraph{Modifier(s)}
\noindent\begin{itemize}
\item coherency (cf. coherency in section \ref{par:modifier-coherency} on page \pageref{par:modifier-coherency})
\end{itemize}

\paragraph{Execution} {Reservation table \textsf{LSU2\_MEMW\_AUXR.Y} (Section~\ref{sec:reservation-LSU2-MEMW-AUXR.Y})}
~
{\begin{lstlisting}
stage ID:
  new offset = _SX_54(extend27_offset27);
  new coherency = _ZX_2(coherency);
stage RR:
  new base = GPR[registerZ];
  new address = base + offset;
stage E1:
  new argument2 = GPR[registerT];
stage E3:
  MEM.atomic_min(address, 0x3, coherency << 32, argument2.16[0], registerT);
\end{lstlisting}}
\newpage
\subsubsection[{\small ASMINUB: Atomic Store MIN Unsigned Byte}]{ASMINUB\hfill Atomic Store MIN Unsigned Byte}
\label{instr:ASMINUB}\index{asminub}
 
\paragraph{Encoding} Format \textsf{LSU\_ASMINUSB} (Section~\ref{sec:format-LSU-ASMINUSB}), \verb|lsucode5 = 000|\\

\bitfields{ 1/P, 2/01, 3/111, 2/coherency, 6/registerT, 2/11, 3/000, 1/--, 5/11010, 1/0, 6/registerZ }


\paragraph{Syntax} \texttt{asminub}\emph{coherency} [\emph{registerZ}] = \emph{registerT}

\paragraph{Description} The effective address is the value of the \emph{registerZ}.
This instruction implements atomic MINU-store on signed byte. The byte at effective address is MINUed with the \emph{registerT}.


\paragraph{Modifier(s)}
\noindent\begin{itemize}
\item coherency (cf. coherency in section \ref{par:modifier-coherency} on page \pageref{par:modifier-coherency})
\end{itemize}

\paragraph{Execution} {Reservation table \textsf{LSU2\_MEMW\_AUXR} (Section~\ref{sec:reservation-LSU2-MEMW-AUXR})}
~
{\begin{lstlisting}
stage ID:
  new coherency = _ZX_2(coherency);
stage RR:
  new address = GPR[registerZ];
stage E1:
  new argument2 = GPR[registerT];
stage E3:
  MEM.atomic_minu(address, 0x1, coherency << 32, argument2.8[0], registerT);
\end{lstlisting}}
\newpage

\paragraph{Encoding} Format \textsf{LSU\_ASMINUSB.O} (Section~\ref{sec:format-LSU-ASMINUSB.O}), \verb|lsucode5 = 000|\\

\bitfields{ 1/P, 2/00, 2/-- --, 27/offset27, 1/1, 2/01, 3/111, 2/coherency, 6/registerT, 2/11, 3/000, 1/--, 5/11010, 1/0, 6/registerZ }


\paragraph{Syntax} \texttt{asminub}\emph{coherency} \emph{offset27}[\emph{registerZ}] = \emph{registerT}

\paragraph{Description} The effective address is computed by adding the \emph{offset27} to the \emph{registerZ}.
This instruction implements atomic MINU-store on signed byte. The byte at effective address is MINUed with the \emph{registerT}.


\paragraph{Modifier(s)}
\noindent\begin{itemize}
\item coherency (cf. coherency in section \ref{par:modifier-coherency} on page \pageref{par:modifier-coherency})
\end{itemize}

\paragraph{Execution} {Reservation table \textsf{LSU2\_MEMW\_AUXR.X} (Section~\ref{sec:reservation-LSU2-MEMW-AUXR.X})}
~
{\begin{lstlisting}
stage ID:
  new offset = _SX_27(offset27);
  new coherency = _ZX_2(coherency);
stage RR:
  new base = GPR[registerZ];
  new address = base + offset;
stage E1:
  new argument2 = GPR[registerT];
stage E3:
  MEM.atomic_minu(address, 0x1, coherency << 32, argument2.8[0], registerT);
\end{lstlisting}}
\newpage

\paragraph{Encoding} Format \textsf{LSU\_ASMINUSB.D} (Section~\ref{sec:format-LSU-ASMINUSB.D}), \verb|lsucode5 = 000|\\

\bitfields{ 1/P, 2/00, 2/-- --, 27/extend27, 1/1, 2/00, 2/-- --, 27/offset27, 1/1, 2/01, 3/111, 2/coherency, 6/registerT, 2/11, 3/000, 1/--, 5/11010, 1/0, 6/registerZ }


\paragraph{Syntax} \texttt{asminub}\emph{coherency} \emph{extend27\_\discretionary{}{}{}offset27}[\emph{registerZ}] = \emph{registerT}

\paragraph{Description} The effective address is computed by adding the \emph{extend27\_\discretionary{}{}{}offset27} to the \emph{registerZ}.
This instruction implements atomic MINU-store on signed byte. The byte at effective address is MINUed with the \emph{registerT}.


\paragraph{Modifier(s)}
\noindent\begin{itemize}
\item coherency (cf. coherency in section \ref{par:modifier-coherency} on page \pageref{par:modifier-coherency})
\end{itemize}

\paragraph{Execution} {Reservation table \textsf{LSU2\_MEMW\_AUXR.Y} (Section~\ref{sec:reservation-LSU2-MEMW-AUXR.Y})}
~
{\begin{lstlisting}
stage ID:
  new offset = _SX_54(extend27_offset27);
  new coherency = _ZX_2(coherency);
stage RR:
  new base = GPR[registerZ];
  new address = base + offset;
stage E1:
  new argument2 = GPR[registerT];
stage E3:
  MEM.atomic_minu(address, 0x1, coherency << 32, argument2.8[0], registerT);
\end{lstlisting}}
\newpage
\subsubsection[{\small ASMINUD: Atomic Store MIN Unsigned Double Word}]{ASMINUD\hfill Atomic Store MIN Unsigned Double Word}
\label{instr:ASMINUD}\index{asminud}
 
\paragraph{Encoding} Format \textsf{LSU\_ASMINUSB} (Section~\ref{sec:format-LSU-ASMINUSB}), \verb|lsucode5 = 011|\\

\bitfields{ 1/P, 2/01, 3/111, 2/coherency, 6/registerT, 2/11, 3/011, 1/--, 5/11010, 1/0, 6/registerZ }


\paragraph{Syntax} \texttt{asminud}\emph{coherency} [\emph{registerZ}] = \emph{registerT}

\paragraph{Description} The effective address is the value of the \emph{registerZ}.
This instruction implements atomic MINU-store on signed double word. The double word at effective address is MINUed with the \emph{registerT}. The effective address must be double word aligned (multiple of 8).


\paragraph{Modifier(s)}
\noindent\begin{itemize}
\item coherency (cf. coherency in section \ref{par:modifier-coherency} on page \pageref{par:modifier-coherency})
\end{itemize}

\paragraph{Execution} {Reservation table \textsf{LSU2\_MEMW\_AUXR} (Section~\ref{sec:reservation-LSU2-MEMW-AUXR})}
~
{\begin{lstlisting}
stage ID:
  new coherency = _ZX_2(coherency);
stage RR:
  new address = GPR[registerZ];
stage E1:
  new argument2 = GPR[registerT];
stage E3:
  MEM.atomic_minu(address, 0xFF, coherency << 32, argument2.64[0], registerT);
\end{lstlisting}}
\newpage

\paragraph{Encoding} Format \textsf{LSU\_ASMINUSB.O} (Section~\ref{sec:format-LSU-ASMINUSB.O}), \verb|lsucode5 = 011|\\

\bitfields{ 1/P, 2/00, 2/-- --, 27/offset27, 1/1, 2/01, 3/111, 2/coherency, 6/registerT, 2/11, 3/011, 1/--, 5/11010, 1/0, 6/registerZ }


\paragraph{Syntax} \texttt{asminud}\emph{coherency} \emph{offset27}[\emph{registerZ}] = \emph{registerT}

\paragraph{Description} The effective address is computed by adding the \emph{offset27} to the \emph{registerZ}.
This instruction implements atomic MINU-store on signed double word. The double word at effective address is MINUed with the \emph{registerT}. The effective address must be double word aligned (multiple of 8).


\paragraph{Modifier(s)}
\noindent\begin{itemize}
\item coherency (cf. coherency in section \ref{par:modifier-coherency} on page \pageref{par:modifier-coherency})
\end{itemize}

\paragraph{Execution} {Reservation table \textsf{LSU2\_MEMW\_AUXR.X} (Section~\ref{sec:reservation-LSU2-MEMW-AUXR.X})}
~
{\begin{lstlisting}
stage ID:
  new offset = _SX_27(offset27);
  new coherency = _ZX_2(coherency);
stage RR:
  new base = GPR[registerZ];
  new address = base + offset;
stage E1:
  new argument2 = GPR[registerT];
stage E3:
  MEM.atomic_minu(address, 0xFF, coherency << 32, argument2.64[0], registerT);
\end{lstlisting}}
\newpage

\paragraph{Encoding} Format \textsf{LSU\_ASMINUSB.D} (Section~\ref{sec:format-LSU-ASMINUSB.D}), \verb|lsucode5 = 011|\\

\bitfields{ 1/P, 2/00, 2/-- --, 27/extend27, 1/1, 2/00, 2/-- --, 27/offset27, 1/1, 2/01, 3/111, 2/coherency, 6/registerT, 2/11, 3/011, 1/--, 5/11010, 1/0, 6/registerZ }


\paragraph{Syntax} \texttt{asminud}\emph{coherency} \emph{extend27\_\discretionary{}{}{}offset27}[\emph{registerZ}] = \emph{registerT}

\paragraph{Description} The effective address is computed by adding the \emph{extend27\_\discretionary{}{}{}offset27} to the \emph{registerZ}.
This instruction implements atomic MINU-store on signed double word. The double word at effective address is MINUed with the \emph{registerT}. The effective address must be double word aligned (multiple of 8).


\paragraph{Modifier(s)}
\noindent\begin{itemize}
\item coherency (cf. coherency in section \ref{par:modifier-coherency} on page \pageref{par:modifier-coherency})
\end{itemize}

\paragraph{Execution} {Reservation table \textsf{LSU2\_MEMW\_AUXR.Y} (Section~\ref{sec:reservation-LSU2-MEMW-AUXR.Y})}
~
{\begin{lstlisting}
stage ID:
  new offset = _SX_54(extend27_offset27);
  new coherency = _ZX_2(coherency);
stage RR:
  new base = GPR[registerZ];
  new address = base + offset;
stage E1:
  new argument2 = GPR[registerT];
stage E3:
  MEM.atomic_minu(address, 0xFF, coherency << 32, argument2.64[0], registerT);
\end{lstlisting}}
\newpage
\subsubsection[{\small ASMINUH: Atomic Store MIN Unsigned Half Word}]{ASMINUH\hfill Atomic Store MIN Unsigned Half Word}
\label{instr:ASMINUH}\index{asminuh}
 
\paragraph{Encoding} Format \textsf{LSU\_ASMINUSB} (Section~\ref{sec:format-LSU-ASMINUSB}), \verb|lsucode5 = 001|\\

\bitfields{ 1/P, 2/01, 3/111, 2/coherency, 6/registerT, 2/11, 3/001, 1/--, 5/11010, 1/0, 6/registerZ }


\paragraph{Syntax} \texttt{asminuh}\emph{coherency} [\emph{registerZ}] = \emph{registerT}

\paragraph{Description} The effective address is the value of the \emph{registerZ}.
This instruction implements atomic MINU-store on signed half word. The half word at effective address is MINUed with the \emph{registerT}. The effective address must be half word aligned (multiple of 2).


\paragraph{Modifier(s)}
\noindent\begin{itemize}
\item coherency (cf. coherency in section \ref{par:modifier-coherency} on page \pageref{par:modifier-coherency})
\end{itemize}

\paragraph{Execution} {Reservation table \textsf{LSU2\_MEMW\_AUXR} (Section~\ref{sec:reservation-LSU2-MEMW-AUXR})}
~
{\begin{lstlisting}
stage ID:
  new coherency = _ZX_2(coherency);
stage RR:
  new address = GPR[registerZ];
stage E1:
  new argument2 = GPR[registerT];
stage E3:
  MEM.atomic_minu(address, 0x3, coherency << 32, argument2.16[0], registerT);
\end{lstlisting}}
\newpage

\paragraph{Encoding} Format \textsf{LSU\_ASMINUSB.O} (Section~\ref{sec:format-LSU-ASMINUSB.O}), \verb|lsucode5 = 001|\\

\bitfields{ 1/P, 2/00, 2/-- --, 27/offset27, 1/1, 2/01, 3/111, 2/coherency, 6/registerT, 2/11, 3/001, 1/--, 5/11010, 1/0, 6/registerZ }


\paragraph{Syntax} \texttt{asminuh}\emph{coherency} \emph{offset27}[\emph{registerZ}] = \emph{registerT}

\paragraph{Description} The effective address is computed by adding the \emph{offset27} to the \emph{registerZ}.
This instruction implements atomic MINU-store on signed half word. The half word at effective address is MINUed with the \emph{registerT}. The effective address must be half word aligned (multiple of 2).


\paragraph{Modifier(s)}
\noindent\begin{itemize}
\item coherency (cf. coherency in section \ref{par:modifier-coherency} on page \pageref{par:modifier-coherency})
\end{itemize}

\paragraph{Execution} {Reservation table \textsf{LSU2\_MEMW\_AUXR.X} (Section~\ref{sec:reservation-LSU2-MEMW-AUXR.X})}
~
{\begin{lstlisting}
stage ID:
  new offset = _SX_27(offset27);
  new coherency = _ZX_2(coherency);
stage RR:
  new base = GPR[registerZ];
  new address = base + offset;
stage E1:
  new argument2 = GPR[registerT];
stage E3:
  MEM.atomic_minu(address, 0x3, coherency << 32, argument2.16[0], registerT);
\end{lstlisting}}
\newpage

\paragraph{Encoding} Format \textsf{LSU\_ASMINUSB.D} (Section~\ref{sec:format-LSU-ASMINUSB.D}), \verb|lsucode5 = 001|\\

\bitfields{ 1/P, 2/00, 2/-- --, 27/extend27, 1/1, 2/00, 2/-- --, 27/offset27, 1/1, 2/01, 3/111, 2/coherency, 6/registerT, 2/11, 3/001, 1/--, 5/11010, 1/0, 6/registerZ }


\paragraph{Syntax} \texttt{asminuh}\emph{coherency} \emph{extend27\_\discretionary{}{}{}offset27}[\emph{registerZ}] = \emph{registerT}

\paragraph{Description} The effective address is computed by adding the \emph{extend27\_\discretionary{}{}{}offset27} to the \emph{registerZ}.
This instruction implements atomic MINU-store on signed half word. The half word at effective address is MINUed with the \emph{registerT}. The effective address must be half word aligned (multiple of 2).


\paragraph{Modifier(s)}
\noindent\begin{itemize}
\item coherency (cf. coherency in section \ref{par:modifier-coherency} on page \pageref{par:modifier-coherency})
\end{itemize}

\paragraph{Execution} {Reservation table \textsf{LSU2\_MEMW\_AUXR.Y} (Section~\ref{sec:reservation-LSU2-MEMW-AUXR.Y})}
~
{\begin{lstlisting}
stage ID:
  new offset = _SX_54(extend27_offset27);
  new coherency = _ZX_2(coherency);
stage RR:
  new base = GPR[registerZ];
  new address = base + offset;
stage E1:
  new argument2 = GPR[registerT];
stage E3:
  MEM.atomic_minu(address, 0x3, coherency << 32, argument2.16[0], registerT);
\end{lstlisting}}
\newpage
\subsubsection[{\small ASMINUW: Atomic Store MIN Unsigned Word}]{ASMINUW\hfill Atomic Store MIN Unsigned Word}
\label{instr:ASMINUW}\index{asminuw}
 
\paragraph{Encoding} Format \textsf{LSU\_ASMINUSB} (Section~\ref{sec:format-LSU-ASMINUSB}), \verb|lsucode5 = 010|\\

\bitfields{ 1/P, 2/01, 3/111, 2/coherency, 6/registerT, 2/11, 3/010, 1/--, 5/11010, 1/0, 6/registerZ }


\paragraph{Syntax} \texttt{asminuw}\emph{coherency} [\emph{registerZ}] = \emph{registerT}

\paragraph{Description} The effective address is the value of the \emph{registerZ}.
This instruction implements atomic MINU-store on signed word. The word at effective address is MINUed with the \emph{registerT}. The effective address must be word aligned (multiple of 4).


\paragraph{Modifier(s)}
\noindent\begin{itemize}
\item coherency (cf. coherency in section \ref{par:modifier-coherency} on page \pageref{par:modifier-coherency})
\end{itemize}

\paragraph{Execution} {Reservation table \textsf{LSU2\_MEMW\_AUXR} (Section~\ref{sec:reservation-LSU2-MEMW-AUXR})}
~
{\begin{lstlisting}
stage ID:
  new coherency = _ZX_2(coherency);
stage RR:
  new address = GPR[registerZ];
stage E1:
  new argument2 = GPR[registerT];
stage E3:
  MEM.atomic_minu(address, 0xF, coherency << 32, argument2.32[0], registerT);
\end{lstlisting}}
\newpage

\paragraph{Encoding} Format \textsf{LSU\_ASMINUSB.O} (Section~\ref{sec:format-LSU-ASMINUSB.O}), \verb|lsucode5 = 010|\\

\bitfields{ 1/P, 2/00, 2/-- --, 27/offset27, 1/1, 2/01, 3/111, 2/coherency, 6/registerT, 2/11, 3/010, 1/--, 5/11010, 1/0, 6/registerZ }


\paragraph{Syntax} \texttt{asminuw}\emph{coherency} \emph{offset27}[\emph{registerZ}] = \emph{registerT}

\paragraph{Description} The effective address is computed by adding the \emph{offset27} to the \emph{registerZ}.
This instruction implements atomic MINU-store on signed word. The word at effective address is MINUed with the \emph{registerT}. The effective address must be word aligned (multiple of 4).


\paragraph{Modifier(s)}
\noindent\begin{itemize}
\item coherency (cf. coherency in section \ref{par:modifier-coherency} on page \pageref{par:modifier-coherency})
\end{itemize}

\paragraph{Execution} {Reservation table \textsf{LSU2\_MEMW\_AUXR.X} (Section~\ref{sec:reservation-LSU2-MEMW-AUXR.X})}
~
{\begin{lstlisting}
stage ID:
  new offset = _SX_27(offset27);
  new coherency = _ZX_2(coherency);
stage RR:
  new base = GPR[registerZ];
  new address = base + offset;
stage E1:
  new argument2 = GPR[registerT];
stage E3:
  MEM.atomic_minu(address, 0xF, coherency << 32, argument2.32[0], registerT);
\end{lstlisting}}
\newpage

\paragraph{Encoding} Format \textsf{LSU\_ASMINUSB.D} (Section~\ref{sec:format-LSU-ASMINUSB.D}), \verb|lsucode5 = 010|\\

\bitfields{ 1/P, 2/00, 2/-- --, 27/extend27, 1/1, 2/00, 2/-- --, 27/offset27, 1/1, 2/01, 3/111, 2/coherency, 6/registerT, 2/11, 3/010, 1/--, 5/11010, 1/0, 6/registerZ }


\paragraph{Syntax} \texttt{asminuw}\emph{coherency} \emph{extend27\_\discretionary{}{}{}offset27}[\emph{registerZ}] = \emph{registerT}

\paragraph{Description} The effective address is computed by adding the \emph{extend27\_\discretionary{}{}{}offset27} to the \emph{registerZ}.
This instruction implements atomic MINU-store on signed word. The word at effective address is MINUed with the \emph{registerT}. The effective address must be word aligned (multiple of 4).


\paragraph{Modifier(s)}
\noindent\begin{itemize}
\item coherency (cf. coherency in section \ref{par:modifier-coherency} on page \pageref{par:modifier-coherency})
\end{itemize}

\paragraph{Execution} {Reservation table \textsf{LSU2\_MEMW\_AUXR.Y} (Section~\ref{sec:reservation-LSU2-MEMW-AUXR.Y})}
~
{\begin{lstlisting}
stage ID:
  new offset = _SX_54(extend27_offset27);
  new coherency = _ZX_2(coherency);
stage RR:
  new base = GPR[registerZ];
  new address = base + offset;
stage E1:
  new argument2 = GPR[registerT];
stage E3:
  MEM.atomic_minu(address, 0xF, coherency << 32, argument2.32[0], registerT);
\end{lstlisting}}
\newpage
\subsubsection[{\small ASMINW: Atomic Store MIN Word}]{ASMINW\hfill Atomic Store MIN Word}
\label{instr:ASMINW}\index{asminw}
 
\paragraph{Encoding} Format \textsf{LSU\_ASMINSB} (Section~\ref{sec:format-LSU-ASMINSB}), \verb|lsucode5 = 010|\\

\bitfields{ 1/P, 2/01, 3/111, 2/coherency, 6/registerT, 2/11, 3/010, 1/--, 5/11000, 1/0, 6/registerZ }


\paragraph{Syntax} \texttt{asminw}\emph{coherency} [\emph{registerZ}] = \emph{registerT}

\paragraph{Description} The effective address is the value of the \emph{registerZ}.
This instruction implements atomic MIN-store on signed word. The word at effective address is MINed with the \emph{registerT}. The effective address must be word aligned (multiple of 4).


\paragraph{Modifier(s)}
\noindent\begin{itemize}
\item coherency (cf. coherency in section \ref{par:modifier-coherency} on page \pageref{par:modifier-coherency})
\end{itemize}

\paragraph{Execution} {Reservation table \textsf{LSU2\_MEMW\_AUXR} (Section~\ref{sec:reservation-LSU2-MEMW-AUXR})}
~
{\begin{lstlisting}
stage ID:
  new coherency = _ZX_2(coherency);
stage RR:
  new address = GPR[registerZ];
stage E1:
  new argument2 = GPR[registerT];
stage E3:
  MEM.atomic_min(address, 0xF, coherency << 32, argument2.32[0], registerT);
\end{lstlisting}}
\newpage

\paragraph{Encoding} Format \textsf{LSU\_ASMINSB.O} (Section~\ref{sec:format-LSU-ASMINSB.O}), \verb|lsucode5 = 010|\\

\bitfields{ 1/P, 2/00, 2/-- --, 27/offset27, 1/1, 2/01, 3/111, 2/coherency, 6/registerT, 2/11, 3/010, 1/--, 5/11000, 1/0, 6/registerZ }


\paragraph{Syntax} \texttt{asminw}\emph{coherency} \emph{offset27}[\emph{registerZ}] = \emph{registerT}

\paragraph{Description} The effective address is computed by adding the \emph{offset27} to the \emph{registerZ}.
This instruction implements atomic MIN-store on signed word. The word at effective address is MINed with the \emph{registerT}. The effective address must be word aligned (multiple of 4).


\paragraph{Modifier(s)}
\noindent\begin{itemize}
\item coherency (cf. coherency in section \ref{par:modifier-coherency} on page \pageref{par:modifier-coherency})
\end{itemize}

\paragraph{Execution} {Reservation table \textsf{LSU2\_MEMW\_AUXR.X} (Section~\ref{sec:reservation-LSU2-MEMW-AUXR.X})}
~
{\begin{lstlisting}
stage ID:
  new offset = _SX_27(offset27);
  new coherency = _ZX_2(coherency);
stage RR:
  new base = GPR[registerZ];
  new address = base + offset;
stage E1:
  new argument2 = GPR[registerT];
stage E3:
  MEM.atomic_min(address, 0xF, coherency << 32, argument2.32[0], registerT);
\end{lstlisting}}
\newpage

\paragraph{Encoding} Format \textsf{LSU\_ASMINSB.D} (Section~\ref{sec:format-LSU-ASMINSB.D}), \verb|lsucode5 = 010|\\

\bitfields{ 1/P, 2/00, 2/-- --, 27/extend27, 1/1, 2/00, 2/-- --, 27/offset27, 1/1, 2/01, 3/111, 2/coherency, 6/registerT, 2/11, 3/010, 1/--, 5/11000, 1/0, 6/registerZ }


\paragraph{Syntax} \texttt{asminw}\emph{coherency} \emph{extend27\_\discretionary{}{}{}offset27}[\emph{registerZ}] = \emph{registerT}

\paragraph{Description} The effective address is computed by adding the \emph{extend27\_\discretionary{}{}{}offset27} to the \emph{registerZ}.
This instruction implements atomic MIN-store on signed word. The word at effective address is MINed with the \emph{registerT}. The effective address must be word aligned (multiple of 4).


\paragraph{Modifier(s)}
\noindent\begin{itemize}
\item coherency (cf. coherency in section \ref{par:modifier-coherency} on page \pageref{par:modifier-coherency})
\end{itemize}

\paragraph{Execution} {Reservation table \textsf{LSU2\_MEMW\_AUXR.Y} (Section~\ref{sec:reservation-LSU2-MEMW-AUXR.Y})}
~
{\begin{lstlisting}
stage ID:
  new offset = _SX_54(extend27_offset27);
  new coherency = _ZX_2(coherency);
stage RR:
  new base = GPR[registerZ];
  new address = base + offset;
stage E1:
  new argument2 = GPR[registerT];
stage E3:
  MEM.atomic_min(address, 0xF, coherency << 32, argument2.32[0], registerT);
\end{lstlisting}}
\newpage
\subsubsection[{\small ASW: Atomic Store Word}]{ASW\hfill Atomic Store Word}
\label{instr:ASW}\index{asw}
 
\paragraph{Encoding} Format \textsf{LSU\_ASSB} (Section~\ref{sec:format-LSU-ASSB}), \verb|lsucode5 = 010|\\

\bitfields{ 1/P, 2/01, 3/111, 2/coherency, 6/registerT, 2/11, 3/010, 1/--, 5/10000, 1/0, 6/registerZ }


\paragraph{Syntax} \texttt{asw}\emph{coherency} [\emph{registerZ}] = \emph{registerT}

\paragraph{Description} The effective address is the value of the \emph{registerZ}.
The \emph{registerT} is atomically stored to the memory word starting at the effective address. This instruction is provided to directly operate on the last level of the memory hierarchy, which is typically a processor DDR or a cluster TCM, irrespective of the MMU DCP (Data Cache Policy) settings. The effective address must be word aligned (multiple of 4).


\paragraph{Modifier(s)}
\noindent\begin{itemize}
\item coherency (cf. coherency in section \ref{par:modifier-coherency} on page \pageref{par:modifier-coherency})
\end{itemize}

\paragraph{Execution} {Reservation table \textsf{LSU\_MEMW\_AUXR} (Section~\ref{sec:reservation-LSU-MEMW-AUXR})}
~
{\begin{lstlisting}
stage ID:
  new coherency = _ZX_2(coherency);
stage RR:
  new address = GPR[registerZ];
stage E1:
  new argument2 = GPR[registerT];
stage E3:
  MEM.atomic_store(address, 0xF, coherency << 32, argument2, registerT);
\end{lstlisting}}
\newpage

\paragraph{Encoding} Format \textsf{LSU\_ASSB.O} (Section~\ref{sec:format-LSU-ASSB.O}), \verb|lsucode5 = 010|\\

\bitfields{ 1/P, 2/00, 2/-- --, 27/offset27, 1/1, 2/01, 3/111, 2/coherency, 6/registerT, 2/11, 3/010, 1/--, 5/10000, 1/0, 6/registerZ }


\paragraph{Syntax} \texttt{asw}\emph{coherency} \emph{offset27}[\emph{registerZ}] = \emph{registerT}

\paragraph{Description} The effective address is computed by adding the \emph{offset27} to the \emph{registerZ}.
The \emph{registerT} is atomically stored to the memory word starting at the effective address. This instruction is provided to directly operate on the last level of the memory hierarchy, which is typically a processor DDR or a cluster TCM, irrespective of the MMU DCP (Data Cache Policy) settings. The effective address must be word aligned (multiple of 4).


\paragraph{Modifier(s)}
\noindent\begin{itemize}
\item coherency (cf. coherency in section \ref{par:modifier-coherency} on page \pageref{par:modifier-coherency})
\end{itemize}

\paragraph{Execution} {Reservation table \textsf{LSU\_MEMW\_AUXR.X} (Section~\ref{sec:reservation-LSU-MEMW-AUXR.X})}
~
{\begin{lstlisting}
stage ID:
  new offset = _SX_27(offset27);
  new coherency = _ZX_2(coherency);
stage RR:
  new base = GPR[registerZ];
  new address = base + offset;
stage E1:
  new argument2 = GPR[registerT];
stage E3:
  MEM.atomic_store(address, 0xF, coherency << 32, argument2, registerT);
\end{lstlisting}}
\newpage

\paragraph{Encoding} Format \textsf{LSU\_ASSB.D} (Section~\ref{sec:format-LSU-ASSB.D}), \verb|lsucode5 = 010|\\

\bitfields{ 1/P, 2/00, 2/-- --, 27/extend27, 1/1, 2/00, 2/-- --, 27/offset27, 1/1, 2/01, 3/111, 2/coherency, 6/registerT, 2/11, 3/010, 1/--, 5/10000, 1/0, 6/registerZ }


\paragraph{Syntax} \texttt{asw}\emph{coherency} \emph{extend27\_\discretionary{}{}{}offset27}[\emph{registerZ}] = \emph{registerT}

\paragraph{Description} The effective address is computed by adding the \emph{extend27\_\discretionary{}{}{}offset27} to the \emph{registerZ}.
The \emph{registerT} is atomically stored to the memory word starting at the effective address. This instruction is provided to directly operate on the last level of the memory hierarchy, which is typically a processor DDR or a cluster TCM, irrespective of the MMU DCP (Data Cache Policy) settings. The effective address must be word aligned (multiple of 4).


\paragraph{Modifier(s)}
\noindent\begin{itemize}
\item coherency (cf. coherency in section \ref{par:modifier-coherency} on page \pageref{par:modifier-coherency})
\end{itemize}

\paragraph{Execution} {Reservation table \textsf{LSU\_MEMW\_AUXR.Y} (Section~\ref{sec:reservation-LSU-MEMW-AUXR.Y})}
~
{\begin{lstlisting}
stage ID:
  new offset = _SX_54(extend27_offset27);
  new coherency = _ZX_2(coherency);
stage RR:
  new base = GPR[registerZ];
  new address = base + offset;
stage E1:
  new argument2 = GPR[registerT];
stage E3:
  MEM.atomic_store(address, 0xF, coherency << 32, argument2, registerT);
\end{lstlisting}}
\newpage
\subsubsection[{\small ASWAPB: Atomic Swap Byte}]{ASWAPB\hfill Atomic Swap Byte}
\label{instr:ASWAPB}\index{aswapb}
 
\paragraph{Encoding} Format \textsf{LSU\_ASWAPSB} (Section~\ref{sec:format-LSU-ASWAPSB}), \verb|lsucode5 = 000|\\

\bitfields{ 1/P, 2/01, 3/111, 2/coherency, 6/registerT, 2/11, 3/000, 1/--, 5/00010, 1/0, 6/registerZ }


\paragraph{Syntax} \texttt{aswapb}\emph{coherency} [\emph{registerZ}] = \emph{registerT}

\paragraph{Description} The effective address is the value of the \emph{registerZ}.
This instruction implements an atomic memory swap of byte. The byte at effective address is written with the \emph{registerT} value, and its previous value is returned into the \emph{registerT}.


\paragraph{Modifier(s)}
\noindent\begin{itemize}
\item coherency (cf. coherency in section \ref{par:modifier-coherency} on page \pageref{par:modifier-coherency})
\end{itemize}

\paragraph{Execution} {Reservation table \textsf{LSU2\_MEMW\_AUXR\_AUXW} (Section~\ref{sec:reservation-LSU2-MEMW-AUXR-AUXW})}
~
{\begin{lstlisting}
stage ID:
  new coherency = _ZX_2(coherency);
stage RR:
  new address = GPR[registerZ];
stage E1:
  new argument2 = GPR[registerT];
  new result2 = _ZX_8(MEM.atomic_swap(address, 0x1, coherency << 32, argument2, registerT));
stage SR:
  GPR[registerT] = result2;
\end{lstlisting}}
\newpage

\paragraph{Encoding} Format \textsf{LSU\_ASWAPSB.O} (Section~\ref{sec:format-LSU-ASWAPSB.O}), \verb|lsucode5 = 000|\\

\bitfields{ 1/P, 2/00, 2/-- --, 27/offset27, 1/1, 2/01, 3/111, 2/coherency, 6/registerT, 2/11, 3/000, 1/--, 5/00010, 1/0, 6/registerZ }


\paragraph{Syntax} \texttt{aswapb}\emph{coherency} \emph{offset27}[\emph{registerZ}] = \emph{registerT}

\paragraph{Description} The effective address is computed by adding the \emph{offset27} to the \emph{registerZ}.
This instruction implements an atomic memory swap of byte. The byte at effective address is written with the \emph{registerT} value, and its previous value is returned into the \emph{registerT}.


\paragraph{Modifier(s)}
\noindent\begin{itemize}
\item coherency (cf. coherency in section \ref{par:modifier-coherency} on page \pageref{par:modifier-coherency})
\end{itemize}

\paragraph{Execution} {Reservation table \textsf{LSU2\_MEMW\_AUXR\_AUXW.X} (Section~\ref{sec:reservation-LSU2-MEMW-AUXR-AUXW.X})}
~
{\begin{lstlisting}
stage ID:
  new offset = _SX_27(offset27);
  new coherency = _ZX_2(coherency);
stage RR:
  new base = GPR[registerZ];
  new address = base + offset;
stage E1:
  new argument2 = GPR[registerT];
  new result2 = _ZX_8(MEM.atomic_swap(address, 0x1, coherency << 32, argument2, registerT));
stage SR:
  GPR[registerT] = result2;
\end{lstlisting}}
\newpage

\paragraph{Encoding} Format \textsf{LSU\_ASWAPSB.D} (Section~\ref{sec:format-LSU-ASWAPSB.D}), \verb|lsucode5 = 000|\\

\bitfields{ 1/P, 2/00, 2/-- --, 27/extend27, 1/1, 2/00, 2/-- --, 27/offset27, 1/1, 2/01, 3/111, 2/coherency, 6/registerT, 2/11, 3/000, 1/--, 5/00010, 1/0, 6/registerZ }


\paragraph{Syntax} \texttt{aswapb}\emph{coherency} \emph{extend27\_\discretionary{}{}{}offset27}[\emph{registerZ}] = \emph{registerT}

\paragraph{Description} The effective address is computed by adding the \emph{extend27\_\discretionary{}{}{}offset27} to the \emph{registerZ}.
This instruction implements an atomic memory swap of byte. The byte at effective address is written with the \emph{registerT} value, and its previous value is returned into the \emph{registerT}.


\paragraph{Modifier(s)}
\noindent\begin{itemize}
\item coherency (cf. coherency in section \ref{par:modifier-coherency} on page \pageref{par:modifier-coherency})
\end{itemize}

\paragraph{Execution} {Reservation table \textsf{LSU2\_MEMW\_AUXR\_AUXW.Y} (Section~\ref{sec:reservation-LSU2-MEMW-AUXR-AUXW.Y})}
~
{\begin{lstlisting}
stage ID:
  new offset = _SX_54(extend27_offset27);
  new coherency = _ZX_2(coherency);
stage RR:
  new base = GPR[registerZ];
  new address = base + offset;
stage E1:
  new argument2 = GPR[registerT];
  new result2 = _ZX_8(MEM.atomic_swap(address, 0x1, coherency << 32, argument2, registerT));
stage SR:
  GPR[registerT] = result2;
\end{lstlisting}}
\newpage
\subsubsection[{\small ASWAPD: Atomic Swap Double Word}]{ASWAPD\hfill Atomic Swap Double Word}
\label{instr:ASWAPD}\index{aswapd}
 
\paragraph{Encoding} Format \textsf{LSU\_ASWAPSB} (Section~\ref{sec:format-LSU-ASWAPSB}), \verb|lsucode5 = 011|\\

\bitfields{ 1/P, 2/01, 3/111, 2/coherency, 6/registerT, 2/11, 3/011, 1/--, 5/00010, 1/0, 6/registerZ }


\paragraph{Syntax} \texttt{aswapd}\emph{coherency} [\emph{registerZ}] = \emph{registerT}

\paragraph{Description} The effective address is the value of the \emph{registerZ}.
This instruction implements an atomic memory swap of double word. The double word at effective address is written with the \emph{registerT} value, and its previous value is returned into the \emph{registerT}.


\paragraph{Modifier(s)}
\noindent\begin{itemize}
\item coherency (cf. coherency in section \ref{par:modifier-coherency} on page \pageref{par:modifier-coherency})
\end{itemize}

\paragraph{Execution} {Reservation table \textsf{LSU2\_MEMW\_AUXR\_AUXW} (Section~\ref{sec:reservation-LSU2-MEMW-AUXR-AUXW})}
~
{\begin{lstlisting}
stage ID:
  new coherency = _ZX_2(coherency);
stage RR:
  new address = GPR[registerZ];
stage E1:
  new argument2 = GPR[registerT];
  new result2 = MEM.atomic_swap(address, 0xFF, coherency << 32, argument2, registerT);
stage SR:
  GPR[registerT] = result2;
\end{lstlisting}}
\newpage

\paragraph{Encoding} Format \textsf{LSU\_ASWAPSB.O} (Section~\ref{sec:format-LSU-ASWAPSB.O}), \verb|lsucode5 = 011|\\

\bitfields{ 1/P, 2/00, 2/-- --, 27/offset27, 1/1, 2/01, 3/111, 2/coherency, 6/registerT, 2/11, 3/011, 1/--, 5/00010, 1/0, 6/registerZ }


\paragraph{Syntax} \texttt{aswapd}\emph{coherency} \emph{offset27}[\emph{registerZ}] = \emph{registerT}

\paragraph{Description} The effective address is computed by adding the \emph{offset27} to the \emph{registerZ}.
This instruction implements an atomic memory swap of double word. The double word at effective address is written with the \emph{registerT} value, and its previous value is returned into the \emph{registerT}.


\paragraph{Modifier(s)}
\noindent\begin{itemize}
\item coherency (cf. coherency in section \ref{par:modifier-coherency} on page \pageref{par:modifier-coherency})
\end{itemize}

\paragraph{Execution} {Reservation table \textsf{LSU2\_MEMW\_AUXR\_AUXW.X} (Section~\ref{sec:reservation-LSU2-MEMW-AUXR-AUXW.X})}
~
{\begin{lstlisting}
stage ID:
  new offset = _SX_27(offset27);
  new coherency = _ZX_2(coherency);
stage RR:
  new base = GPR[registerZ];
  new address = base + offset;
stage E1:
  new argument2 = GPR[registerT];
  new result2 = MEM.atomic_swap(address, 0xFF, coherency << 32, argument2, registerT);
stage SR:
  GPR[registerT] = result2;
\end{lstlisting}}
\newpage

\paragraph{Encoding} Format \textsf{LSU\_ASWAPSB.D} (Section~\ref{sec:format-LSU-ASWAPSB.D}), \verb|lsucode5 = 011|\\

\bitfields{ 1/P, 2/00, 2/-- --, 27/extend27, 1/1, 2/00, 2/-- --, 27/offset27, 1/1, 2/01, 3/111, 2/coherency, 6/registerT, 2/11, 3/011, 1/--, 5/00010, 1/0, 6/registerZ }


\paragraph{Syntax} \texttt{aswapd}\emph{coherency} \emph{extend27\_\discretionary{}{}{}offset27}[\emph{registerZ}] = \emph{registerT}

\paragraph{Description} The effective address is computed by adding the \emph{extend27\_\discretionary{}{}{}offset27} to the \emph{registerZ}.
This instruction implements an atomic memory swap of double word. The double word at effective address is written with the \emph{registerT} value, and its previous value is returned into the \emph{registerT}.


\paragraph{Modifier(s)}
\noindent\begin{itemize}
\item coherency (cf. coherency in section \ref{par:modifier-coherency} on page \pageref{par:modifier-coherency})
\end{itemize}

\paragraph{Execution} {Reservation table \textsf{LSU2\_MEMW\_AUXR\_AUXW.Y} (Section~\ref{sec:reservation-LSU2-MEMW-AUXR-AUXW.Y})}
~
{\begin{lstlisting}
stage ID:
  new offset = _SX_54(extend27_offset27);
  new coherency = _ZX_2(coherency);
stage RR:
  new base = GPR[registerZ];
  new address = base + offset;
stage E1:
  new argument2 = GPR[registerT];
  new result2 = MEM.atomic_swap(address, 0xFF, coherency << 32, argument2, registerT);
stage SR:
  GPR[registerT] = result2;
\end{lstlisting}}
\newpage
\subsubsection[{\small ASWAPH: Atomic Swap Half Word}]{ASWAPH\hfill Atomic Swap Half Word}
\label{instr:ASWAPH}\index{aswaph}
 
\paragraph{Encoding} Format \textsf{LSU\_ASWAPSB} (Section~\ref{sec:format-LSU-ASWAPSB}), \verb|lsucode5 = 001|\\

\bitfields{ 1/P, 2/01, 3/111, 2/coherency, 6/registerT, 2/11, 3/001, 1/--, 5/00010, 1/0, 6/registerZ }


\paragraph{Syntax} \texttt{aswaph}\emph{coherency} [\emph{registerZ}] = \emph{registerT}

\paragraph{Description} The effective address is the value of the \emph{registerZ}.
This instruction implements an atomic memory swap of half word. The half word at effective address is written with the \emph{registerT} value, and its previous value is returned into the \emph{registerT}.


\paragraph{Modifier(s)}
\noindent\begin{itemize}
\item coherency (cf. coherency in section \ref{par:modifier-coherency} on page \pageref{par:modifier-coherency})
\end{itemize}

\paragraph{Execution} {Reservation table \textsf{LSU2\_MEMW\_AUXR\_AUXW} (Section~\ref{sec:reservation-LSU2-MEMW-AUXR-AUXW})}
~
{\begin{lstlisting}
stage ID:
  new coherency = _ZX_2(coherency);
stage RR:
  new address = GPR[registerZ];
stage E1:
  new argument2 = GPR[registerT];
  new result2 = _ZX_16(MEM.atomic_swap(address, 0x3, coherency << 32, argument2, registerT));
stage SR:
  GPR[registerT] = result2;
\end{lstlisting}}
\newpage

\paragraph{Encoding} Format \textsf{LSU\_ASWAPSB.O} (Section~\ref{sec:format-LSU-ASWAPSB.O}), \verb|lsucode5 = 001|\\

\bitfields{ 1/P, 2/00, 2/-- --, 27/offset27, 1/1, 2/01, 3/111, 2/coherency, 6/registerT, 2/11, 3/001, 1/--, 5/00010, 1/0, 6/registerZ }


\paragraph{Syntax} \texttt{aswaph}\emph{coherency} \emph{offset27}[\emph{registerZ}] = \emph{registerT}

\paragraph{Description} The effective address is computed by adding the \emph{offset27} to the \emph{registerZ}.
This instruction implements an atomic memory swap of half word. The half word at effective address is written with the \emph{registerT} value, and its previous value is returned into the \emph{registerT}.


\paragraph{Modifier(s)}
\noindent\begin{itemize}
\item coherency (cf. coherency in section \ref{par:modifier-coherency} on page \pageref{par:modifier-coherency})
\end{itemize}

\paragraph{Execution} {Reservation table \textsf{LSU2\_MEMW\_AUXR\_AUXW.X} (Section~\ref{sec:reservation-LSU2-MEMW-AUXR-AUXW.X})}
~
{\begin{lstlisting}
stage ID:
  new offset = _SX_27(offset27);
  new coherency = _ZX_2(coherency);
stage RR:
  new base = GPR[registerZ];
  new address = base + offset;
stage E1:
  new argument2 = GPR[registerT];
  new result2 = _ZX_16(MEM.atomic_swap(address, 0x3, coherency << 32, argument2, registerT));
stage SR:
  GPR[registerT] = result2;
\end{lstlisting}}
\newpage

\paragraph{Encoding} Format \textsf{LSU\_ASWAPSB.D} (Section~\ref{sec:format-LSU-ASWAPSB.D}), \verb|lsucode5 = 001|\\

\bitfields{ 1/P, 2/00, 2/-- --, 27/extend27, 1/1, 2/00, 2/-- --, 27/offset27, 1/1, 2/01, 3/111, 2/coherency, 6/registerT, 2/11, 3/001, 1/--, 5/00010, 1/0, 6/registerZ }


\paragraph{Syntax} \texttt{aswaph}\emph{coherency} \emph{extend27\_\discretionary{}{}{}offset27}[\emph{registerZ}] = \emph{registerT}

\paragraph{Description} The effective address is computed by adding the \emph{extend27\_\discretionary{}{}{}offset27} to the \emph{registerZ}.
This instruction implements an atomic memory swap of half word. The half word at effective address is written with the \emph{registerT} value, and its previous value is returned into the \emph{registerT}.


\paragraph{Modifier(s)}
\noindent\begin{itemize}
\item coherency (cf. coherency in section \ref{par:modifier-coherency} on page \pageref{par:modifier-coherency})
\end{itemize}

\paragraph{Execution} {Reservation table \textsf{LSU2\_MEMW\_AUXR\_AUXW.Y} (Section~\ref{sec:reservation-LSU2-MEMW-AUXR-AUXW.Y})}
~
{\begin{lstlisting}
stage ID:
  new offset = _SX_54(extend27_offset27);
  new coherency = _ZX_2(coherency);
stage RR:
  new base = GPR[registerZ];
  new address = base + offset;
stage E1:
  new argument2 = GPR[registerT];
  new result2 = _ZX_16(MEM.atomic_swap(address, 0x3, coherency << 32, argument2, registerT));
stage SR:
  GPR[registerT] = result2;
\end{lstlisting}}
\newpage
\subsubsection[{\small ASWAPW: Atomic Swap Word}]{ASWAPW\hfill Atomic Swap Word}
\label{instr:ASWAPW}\index{aswapw}
 
\paragraph{Encoding} Format \textsf{LSU\_ASWAPSB} (Section~\ref{sec:format-LSU-ASWAPSB}), \verb|lsucode5 = 010|\\

\bitfields{ 1/P, 2/01, 3/111, 2/coherency, 6/registerT, 2/11, 3/010, 1/--, 5/00010, 1/0, 6/registerZ }


\paragraph{Syntax} \texttt{aswapw}\emph{coherency} [\emph{registerZ}] = \emph{registerT}

\paragraph{Description} The effective address is the value of the \emph{registerZ}.
This instruction implements an atomic memory swap of word. The word at effective address is written with the \emph{registerT} value, and its previous value is returned into the \emph{registerT}.


\paragraph{Modifier(s)}
\noindent\begin{itemize}
\item coherency (cf. coherency in section \ref{par:modifier-coherency} on page \pageref{par:modifier-coherency})
\end{itemize}

\paragraph{Execution} {Reservation table \textsf{LSU2\_MEMW\_AUXR\_AUXW} (Section~\ref{sec:reservation-LSU2-MEMW-AUXR-AUXW})}
~
{\begin{lstlisting}
stage ID:
  new coherency = _ZX_2(coherency);
stage RR:
  new address = GPR[registerZ];
stage E1:
  new argument2 = GPR[registerT];
  new result2 = _ZX_32(MEM.atomic_swap(address, 0xF, coherency << 32, argument2, registerT));
stage SR:
  GPR[registerT] = result2;
\end{lstlisting}}
\newpage

\paragraph{Encoding} Format \textsf{LSU\_ASWAPSB.O} (Section~\ref{sec:format-LSU-ASWAPSB.O}), \verb|lsucode5 = 010|\\

\bitfields{ 1/P, 2/00, 2/-- --, 27/offset27, 1/1, 2/01, 3/111, 2/coherency, 6/registerT, 2/11, 3/010, 1/--, 5/00010, 1/0, 6/registerZ }


\paragraph{Syntax} \texttt{aswapw}\emph{coherency} \emph{offset27}[\emph{registerZ}] = \emph{registerT}

\paragraph{Description} The effective address is computed by adding the \emph{offset27} to the \emph{registerZ}.
This instruction implements an atomic memory swap of word. The word at effective address is written with the \emph{registerT} value, and its previous value is returned into the \emph{registerT}.


\paragraph{Modifier(s)}
\noindent\begin{itemize}
\item coherency (cf. coherency in section \ref{par:modifier-coherency} on page \pageref{par:modifier-coherency})
\end{itemize}

\paragraph{Execution} {Reservation table \textsf{LSU2\_MEMW\_AUXR\_AUXW.X} (Section~\ref{sec:reservation-LSU2-MEMW-AUXR-AUXW.X})}
~
{\begin{lstlisting}
stage ID:
  new offset = _SX_27(offset27);
  new coherency = _ZX_2(coherency);
stage RR:
  new base = GPR[registerZ];
  new address = base + offset;
stage E1:
  new argument2 = GPR[registerT];
  new result2 = _ZX_32(MEM.atomic_swap(address, 0xF, coherency << 32, argument2, registerT));
stage SR:
  GPR[registerT] = result2;
\end{lstlisting}}
\newpage

\paragraph{Encoding} Format \textsf{LSU\_ASWAPSB.D} (Section~\ref{sec:format-LSU-ASWAPSB.D}), \verb|lsucode5 = 010|\\

\bitfields{ 1/P, 2/00, 2/-- --, 27/extend27, 1/1, 2/00, 2/-- --, 27/offset27, 1/1, 2/01, 3/111, 2/coherency, 6/registerT, 2/11, 3/010, 1/--, 5/00010, 1/0, 6/registerZ }


\paragraph{Syntax} \texttt{aswapw}\emph{coherency} \emph{extend27\_\discretionary{}{}{}offset27}[\emph{registerZ}] = \emph{registerT}

\paragraph{Description} The effective address is computed by adding the \emph{extend27\_\discretionary{}{}{}offset27} to the \emph{registerZ}.
This instruction implements an atomic memory swap of word. The word at effective address is written with the \emph{registerT} value, and its previous value is returned into the \emph{registerT}.


\paragraph{Modifier(s)}
\noindent\begin{itemize}
\item coherency (cf. coherency in section \ref{par:modifier-coherency} on page \pageref{par:modifier-coherency})
\end{itemize}

\paragraph{Execution} {Reservation table \textsf{LSU2\_MEMW\_AUXR\_AUXW.Y} (Section~\ref{sec:reservation-LSU2-MEMW-AUXR-AUXW.Y})}
~
{\begin{lstlisting}
stage ID:
  new offset = _SX_54(extend27_offset27);
  new coherency = _ZX_2(coherency);
stage RR:
  new base = GPR[registerZ];
  new address = base + offset;
stage E1:
  new argument2 = GPR[registerT];
  new result2 = _ZX_32(MEM.atomic_swap(address, 0xF, coherency << 32, argument2, registerT));
stage SR:
  GPR[registerT] = result2;
\end{lstlisting}}
\newpage
\subsubsection[{\small COPYO: Copy Octuple Word}]{COPYO\hfill Copy Octuple Word}
\label{instr:COPYO}\index{copyo}
 
\paragraph{Encoding} Format \textsf{LSU\_COPYO} (Section~\ref{sec:format-LSU-COPYO}), \verb|lsucode4 = 01|\\

\bitfields{ 1/P, 2/01, 5/10100, 4/registerN, 1/1, 1/1, 2/01, 10/-- -- -- -- -- -- -- -- -- --, 4/registerR, 1/--, 1/-- }


\paragraph{Syntax} \texttt{copyo} \emph{registerN} = \emph{registerR}

\paragraph{Description} The \emph{registerN} operand receives the \emph{registerR} value.


\paragraph{Execution} {Reservation table \textsf{LSU\_AUXR\_AUXW} (Section~\ref{sec:reservation-LSU-AUXR-AUXW})}
~
{\begin{lstlisting}
stage RR:
  new argument2 = QGR[registerR];
  new result1 = argument2;
stage E3:
  QGR[registerN] = result1;
\end{lstlisting}}
\newpage
\subsubsection[{\small D1INVAL: Data Cache 1 Invalidate}]{D1INVAL\hfill Data Cache 1 Invalidate}
\label{instr:D1INVAL}\index{d1inval}
 
\paragraph{Encoding} Format \textsf{LSU\_MCC} (Section~\ref{sec:format-LSU-MCC}), \verb|lsucode2 = 1000|\\

\bitfields{ 1/P, 2/01, 3/111, 2/00, 4/1000, 1/1, 1/1, 2/00, 16/-- -- -- -- -- -- -- -- -- -- -- -- -- -- -- -- }


\paragraph{Syntax} \texttt{d1inval}

\paragraph{Description} Request to invalidate the data cache. Cache contents is invalidated even if it is modified.


\paragraph{Execution} {Reservation table \textsf{LSU2\_MEMW} (Section~\ref{sec:reservation-LSU2-MEMW})}
~
{\begin{lstlisting}
stage E3:
  MEM.d1inval();
\end{lstlisting}}
\newpage
\subsubsection[{\small DFLUSHL: Data Cache Flush Line}]{DFLUSHL\hfill Data Cache Flush Line}
\label{instr:DFLUSHL}\index{dflushl}
 
\paragraph{Encoding} Format \textsf{LSU\_FXBO} (Section~\ref{sec:format-LSU-FXBO}), \verb|lsucode2 = 0011|\\

\bitfields{ 1/P, 2/01, 3/111, 2/00, 4/0011, 1/1, 1/1, 2/00, 10/signed10, 6/registerZ }


\paragraph{Syntax} \texttt{dflushl} \emph{signed10}[\emph{registerZ}]

\paragraph{Description} The effective address is computed by adding the \emph{signed10} to the \emph{registerZ}.
The effective address is sent to the data cache, with request to evict the corresponding cache line. Has no effects if the line is not cached.


\paragraph{Execution} {Reservation table \textsf{LSU} (Section~\ref{sec:reservation-LSU})}
~
{\begin{lstlisting}
stage ID:
  new offset = _SX_10(signed10);
stage RR:
  new base = GPR[registerZ];
  new address = base + offset;
stage E3:
  MEM.dflushl(address);
\end{lstlisting}}
\newpage

\paragraph{Encoding} Format \textsf{LSU\_FXBO.X} (Section~\ref{sec:format-LSU-FXBO.X}), \verb|lsucode2 = 0011|\\

\bitfields{ 1/P, 2/00, 2/-- --, 27/upper27, 1/1, 2/01, 3/111, 2/00, 4/0011, 1/1, 1/1, 2/00, 10/lower10, 6/registerZ }


\paragraph{Syntax} \texttt{dflushl} \emph{upper27\_\discretionary{}{}{}lower10}[\emph{registerZ}]

\paragraph{Description} The effective address is computed by adding the \emph{upper27\_\discretionary{}{}{}lower10} to the \emph{registerZ}.
The effective address is sent to the data cache, with request to evict the corresponding cache line. Has no effects if the line is not cached.


\paragraph{Execution} {Reservation table \textsf{LSU.X} (Section~\ref{sec:reservation-LSU.X})}
~
{\begin{lstlisting}
stage ID:
  new offset = _SX_37(upper27_lower10);
stage RR:
  new base = GPR[registerZ];
  new address = base + offset;
stage E3:
  MEM.dflushl(address);
\end{lstlisting}}
\newpage

\paragraph{Encoding} Format \textsf{LSU\_FXBO.Y} (Section~\ref{sec:format-LSU-FXBO.Y}), \verb|lsucode2 = 0011|\\

\bitfields{ 1/P, 2/00, 2/-- --, 27/extend27, 1/1, 2/00, 2/-- --, 27/upper27, 1/1, 2/01, 3/111, 2/00, 4/0011, 1/1, 1/1, 2/00, 10/lower10, 6/registerZ }


\paragraph{Syntax} \texttt{dflushl} \emph{extend27\_\discretionary{}{}{}upper27\_\discretionary{}{}{}lower10}[\emph{registerZ}]

\paragraph{Description} The effective address is computed by adding the \emph{extend27\_\discretionary{}{}{}upper27\_\discretionary{}{}{}lower10} to the \emph{registerZ}.
The effective address is sent to the data cache, with request to evict the corresponding cache line. Has no effects if the line is not cached.


\paragraph{Execution} {Reservation table \textsf{LSU.Y} (Section~\ref{sec:reservation-LSU.Y})}
~
{\begin{lstlisting}
stage ID:
  new offset = _WX_64(extend27_upper27_lower10);
stage RR:
  new base = GPR[registerZ];
  new address = base + offset;
stage E3:
  MEM.dflushl(address);
\end{lstlisting}}
\newpage

\paragraph{Encoding} Format \textsf{LSU\_FXBI} (Section~\ref{sec:format-LSU-FXBI}), \verb|lsucode2 = 0011|\\

\bitfields{ 1/P, 2/01, 3/111, 2/00, 4/0011, 1/1, 1/1, 2/10, 3/111, 1/--, 6/registerY, 6/registerZ }


\paragraph{Syntax} \texttt{dflushl} \emph{registerY}[\emph{registerZ}]

\paragraph{Description} The effective address is computed by adding the \emph{registerZ} to the \emph{registerY}.
The effective address is sent to the data cache, with request to evict the corresponding cache line. Has no effects if the line is not cached.


\paragraph{Execution} {Reservation table \textsf{LSU} (Section~\ref{sec:reservation-LSU})}
~
{\begin{lstlisting}
stage RR:
  new doscale = 0;
  new base = GPR[registerZ];
  new index = GPR[registerY];
  new scaling = doscale ? 64 : 1;
  new address = base + (index * scaling);
stage E3:
  MEM.dflushl(address);
\end{lstlisting}}
\newpage
\subsubsection[{\small DFLUSHSW: Data Cache Flush Line by Set and Way}]{DFLUSHSW\hfill Data Cache Flush Line by Set and Way}
\label{instr:DFLUSHSW}\index{dflushsw}
 
\paragraph{Encoding} Format \textsf{LSU\_FXSW} (Section~\ref{sec:format-LSU-FXSW}), \verb|lsucode2 = 1011|\\

\bitfields{ 1/P, 2/01, 3/111, 2/cachelev, 4/1011, 1/1, 1/1, 2/10, 3/111, 1/--, 6/registerY, 6/registerZ }


\paragraph{Syntax} \texttt{dflushsw}\emph{cachelev} \emph{registerY}, \emph{registerZ}

\paragraph{Description} The set and way are sent to the data cache, with request to evict the corresponding cache line. Has no effects if the line is not cached.


\paragraph{Modifier(s)}
\noindent\begin{itemize}
\item cachelev (cf. cachelev in section \ref{par:modifier-cachelev} on page \pageref{par:modifier-cachelev})
\end{itemize}

\paragraph{Execution} {Reservation table \textsf{LSU2\_MEMW} (Section~\ref{sec:reservation-LSU2-MEMW})}
~
{\begin{lstlisting}
stage ID:
  new cachelev = _ZX_2(cachelev);
stage RR:
  new way = GPR[registerZ];
  new set = GPR[registerY];
stage E3:
  MEM.dflushlsw(set, way, cachelev);
\end{lstlisting}}
\newpage
\subsubsection[{\small DINVALL: Data Cache Invalidate Line}]{DINVALL\hfill Data Cache Invalidate Line}
\label{instr:DINVALL}\index{dinvall}
 
\paragraph{Encoding} Format \textsf{LSU\_FXBO} (Section~\ref{sec:format-LSU-FXBO}), \verb|lsucode2 = 0001|\\

\bitfields{ 1/P, 2/01, 3/111, 2/00, 4/0001, 1/1, 1/1, 2/00, 10/signed10, 6/registerZ }


\paragraph{Syntax} \texttt{dinvall} \emph{signed10}[\emph{registerZ}]

\paragraph{Description} The effective address is computed by adding the \emph{signed10} to the \emph{registerZ}.
The effective address is sent to the data cache, with request to invalidate the corresponding cache line. Has no effects if the line is not cached.


\paragraph{Execution} {Reservation table \textsf{LSU} (Section~\ref{sec:reservation-LSU})}
~
{\begin{lstlisting}
stage ID:
  new offset = _SX_10(signed10);
stage RR:
  new base = GPR[registerZ];
  new address = base + offset;
stage E3:
  MEM.dinvall(address);
\end{lstlisting}}
\newpage

\paragraph{Encoding} Format \textsf{LSU\_FXBO.X} (Section~\ref{sec:format-LSU-FXBO.X}), \verb|lsucode2 = 0001|\\

\bitfields{ 1/P, 2/00, 2/-- --, 27/upper27, 1/1, 2/01, 3/111, 2/00, 4/0001, 1/1, 1/1, 2/00, 10/lower10, 6/registerZ }


\paragraph{Syntax} \texttt{dinvall} \emph{upper27\_\discretionary{}{}{}lower10}[\emph{registerZ}]

\paragraph{Description} The effective address is computed by adding the \emph{upper27\_\discretionary{}{}{}lower10} to the \emph{registerZ}.
The effective address is sent to the data cache, with request to invalidate the corresponding cache line. Has no effects if the line is not cached.


\paragraph{Execution} {Reservation table \textsf{LSU.X} (Section~\ref{sec:reservation-LSU.X})}
~
{\begin{lstlisting}
stage ID:
  new offset = _SX_37(upper27_lower10);
stage RR:
  new base = GPR[registerZ];
  new address = base + offset;
stage E3:
  MEM.dinvall(address);
\end{lstlisting}}
\newpage

\paragraph{Encoding} Format \textsf{LSU\_FXBO.Y} (Section~\ref{sec:format-LSU-FXBO.Y}), \verb|lsucode2 = 0001|\\

\bitfields{ 1/P, 2/00, 2/-- --, 27/extend27, 1/1, 2/00, 2/-- --, 27/upper27, 1/1, 2/01, 3/111, 2/00, 4/0001, 1/1, 1/1, 2/00, 10/lower10, 6/registerZ }


\paragraph{Syntax} \texttt{dinvall} \emph{extend27\_\discretionary{}{}{}upper27\_\discretionary{}{}{}lower10}[\emph{registerZ}]

\paragraph{Description} The effective address is computed by adding the \emph{extend27\_\discretionary{}{}{}upper27\_\discretionary{}{}{}lower10} to the \emph{registerZ}.
The effective address is sent to the data cache, with request to invalidate the corresponding cache line. Has no effects if the line is not cached.


\paragraph{Execution} {Reservation table \textsf{LSU.Y} (Section~\ref{sec:reservation-LSU.Y})}
~
{\begin{lstlisting}
stage ID:
  new offset = _WX_64(extend27_upper27_lower10);
stage RR:
  new base = GPR[registerZ];
  new address = base + offset;
stage E3:
  MEM.dinvall(address);
\end{lstlisting}}
\newpage

\paragraph{Encoding} Format \textsf{LSU\_FXBI} (Section~\ref{sec:format-LSU-FXBI}), \verb|lsucode2 = 0001|\\

\bitfields{ 1/P, 2/01, 3/111, 2/00, 4/0001, 1/1, 1/1, 2/10, 3/111, 1/--, 6/registerY, 6/registerZ }


\paragraph{Syntax} \texttt{dinvall} \emph{registerY}[\emph{registerZ}]

\paragraph{Description} The effective address is computed by adding the \emph{registerZ} to the \emph{registerY}.
The effective address is sent to the data cache, with request to invalidate the corresponding cache line. Has no effects if the line is not cached.


\paragraph{Execution} {Reservation table \textsf{LSU} (Section~\ref{sec:reservation-LSU})}
~
{\begin{lstlisting}
stage RR:
  new doscale = 0;
  new base = GPR[registerZ];
  new index = GPR[registerY];
  new scaling = doscale ? 64 : 1;
  new address = base + (index * scaling);
stage E3:
  MEM.dinvall(address);
\end{lstlisting}}
\newpage
\subsubsection[{\small DINVALSW: Data Cache Invalidate Line by Set and Way}]{DINVALSW\hfill Data Cache Invalidate Line by Set and Way}
\label{instr:DINVALSW}\index{dinvalsw}
 
\paragraph{Encoding} Format \textsf{LSU\_FXSW} (Section~\ref{sec:format-LSU-FXSW}), \verb|lsucode2 = 1001|\\

\bitfields{ 1/P, 2/01, 3/111, 2/cachelev, 4/1001, 1/1, 1/1, 2/10, 3/111, 1/--, 6/registerY, 6/registerZ }


\paragraph{Syntax} \texttt{dinvalsw}\emph{cachelev} \emph{registerY}, \emph{registerZ}

\paragraph{Description} The set and way are sent to the data cache, with request to invalidate the corresponding cache line. Has no effects if the line is not cached.


\paragraph{Modifier(s)}
\noindent\begin{itemize}
\item cachelev (cf. cachelev in section \ref{par:modifier-cachelev} on page \pageref{par:modifier-cachelev})
\end{itemize}

\paragraph{Execution} {Reservation table \textsf{LSU2\_MEMW} (Section~\ref{sec:reservation-LSU2-MEMW})}
~
{\begin{lstlisting}
stage ID:
  new cachelev = _ZX_2(cachelev);
stage RR:
  new way = GPR[registerZ];
  new set = GPR[registerY];
stage E3:
  if (!dinvallsw_owner()) {
    _THROW(PRIVILEGE);
  } else {
    MEM.dinvallsw(set, way, cachelev);
  }
\end{lstlisting}}
\newpage
\subsubsection[{\small DPURGEL: Data Cache Purge Line}]{DPURGEL\hfill Data Cache Purge Line}
\label{instr:DPURGEL}\index{dpurgel}
 
\paragraph{Encoding} Format \textsf{LSU\_FXBO} (Section~\ref{sec:format-LSU-FXBO}), \verb|lsucode2 = 0010|\\

\bitfields{ 1/P, 2/01, 3/111, 2/00, 4/0010, 1/1, 1/1, 2/00, 10/signed10, 6/registerZ }


\paragraph{Syntax} \texttt{dpurgel} \emph{signed10}[\emph{registerZ}]

\paragraph{Description} The effective address is computed by adding the \emph{signed10} to the \emph{registerZ}.
The effective address is sent to the data cache, with request to purge the corresponding cache line. Has no effects if the line is not cached.


\paragraph{Execution} {Reservation table \textsf{LSU} (Section~\ref{sec:reservation-LSU})}
~
{\begin{lstlisting}
stage ID:
  new offset = _SX_10(signed10);
stage RR:
  new base = GPR[registerZ];
  new address = base + offset;
stage E3:
  MEM.dpurgel(address);
\end{lstlisting}}
\newpage

\paragraph{Encoding} Format \textsf{LSU\_FXBO.X} (Section~\ref{sec:format-LSU-FXBO.X}), \verb|lsucode2 = 0010|\\

\bitfields{ 1/P, 2/00, 2/-- --, 27/upper27, 1/1, 2/01, 3/111, 2/00, 4/0010, 1/1, 1/1, 2/00, 10/lower10, 6/registerZ }


\paragraph{Syntax} \texttt{dpurgel} \emph{upper27\_\discretionary{}{}{}lower10}[\emph{registerZ}]

\paragraph{Description} The effective address is computed by adding the \emph{upper27\_\discretionary{}{}{}lower10} to the \emph{registerZ}.
The effective address is sent to the data cache, with request to purge the corresponding cache line. Has no effects if the line is not cached.


\paragraph{Execution} {Reservation table \textsf{LSU.X} (Section~\ref{sec:reservation-LSU.X})}
~
{\begin{lstlisting}
stage ID:
  new offset = _SX_37(upper27_lower10);
stage RR:
  new base = GPR[registerZ];
  new address = base + offset;
stage E3:
  MEM.dpurgel(address);
\end{lstlisting}}
\newpage

\paragraph{Encoding} Format \textsf{LSU\_FXBO.Y} (Section~\ref{sec:format-LSU-FXBO.Y}), \verb|lsucode2 = 0010|\\

\bitfields{ 1/P, 2/00, 2/-- --, 27/extend27, 1/1, 2/00, 2/-- --, 27/upper27, 1/1, 2/01, 3/111, 2/00, 4/0010, 1/1, 1/1, 2/00, 10/lower10, 6/registerZ }


\paragraph{Syntax} \texttt{dpurgel} \emph{extend27\_\discretionary{}{}{}upper27\_\discretionary{}{}{}lower10}[\emph{registerZ}]

\paragraph{Description} The effective address is computed by adding the \emph{extend27\_\discretionary{}{}{}upper27\_\discretionary{}{}{}lower10} to the \emph{registerZ}.
The effective address is sent to the data cache, with request to purge the corresponding cache line. Has no effects if the line is not cached.


\paragraph{Execution} {Reservation table \textsf{LSU.Y} (Section~\ref{sec:reservation-LSU.Y})}
~
{\begin{lstlisting}
stage ID:
  new offset = _WX_64(extend27_upper27_lower10);
stage RR:
  new base = GPR[registerZ];
  new address = base + offset;
stage E3:
  MEM.dpurgel(address);
\end{lstlisting}}
\newpage

\paragraph{Encoding} Format \textsf{LSU\_FXBI} (Section~\ref{sec:format-LSU-FXBI}), \verb|lsucode2 = 0010|\\

\bitfields{ 1/P, 2/01, 3/111, 2/00, 4/0010, 1/1, 1/1, 2/10, 3/111, 1/--, 6/registerY, 6/registerZ }


\paragraph{Syntax} \texttt{dpurgel} \emph{registerY}[\emph{registerZ}]

\paragraph{Description} The effective address is computed by adding the \emph{registerZ} to the \emph{registerY}.
The effective address is sent to the data cache, with request to purge the corresponding cache line. Has no effects if the line is not cached.


\paragraph{Execution} {Reservation table \textsf{LSU} (Section~\ref{sec:reservation-LSU})}
~
{\begin{lstlisting}
stage RR:
  new doscale = 0;
  new base = GPR[registerZ];
  new index = GPR[registerY];
  new scaling = doscale ? 64 : 1;
  new address = base + (index * scaling);
stage E3:
  MEM.dpurgel(address);
\end{lstlisting}}
\newpage
\subsubsection[{\small DPURGESW: Data Cache Purge Line by Set and Way}]{DPURGESW\hfill Data Cache Purge Line by Set and Way}
\label{instr:DPURGESW}\index{dpurgesw}
 
\paragraph{Encoding} Format \textsf{LSU\_FXSW} (Section~\ref{sec:format-LSU-FXSW}), \verb|lsucode2 = 1010|\\

\bitfields{ 1/P, 2/01, 3/111, 2/cachelev, 4/1010, 1/1, 1/1, 2/10, 3/111, 1/--, 6/registerY, 6/registerZ }


\paragraph{Syntax} \texttt{dpurgesw}\emph{cachelev} \emph{registerY}, \emph{registerZ}

\paragraph{Description} The set and way are sent to the data cache, with request to purge the corresponding cache line. Has no effects if the line is not cached.


\paragraph{Modifier(s)}
\noindent\begin{itemize}
\item cachelev (cf. cachelev in section \ref{par:modifier-cachelev} on page \pageref{par:modifier-cachelev})
\end{itemize}

\paragraph{Execution} {Reservation table \textsf{LSU2\_MEMW} (Section~\ref{sec:reservation-LSU2-MEMW})}
~
{\begin{lstlisting}
stage ID:
  new cachelev = _ZX_2(cachelev);
stage RR:
  new way = GPR[registerZ];
  new set = GPR[registerY];
stage E3:
  MEM.dpurgelsw(set, way, cachelev);
\end{lstlisting}}
\newpage
\subsubsection[{\small DTOUCHL: Data Cache Touch Line}]{DTOUCHL\hfill Data Cache Touch Line}
\label{instr:DTOUCHL}\index{dtouchl}
 
\paragraph{Encoding} Format \textsf{LSU\_FXBO} (Section~\ref{sec:format-LSU-FXBO}), \verb|lsucode2 = 0000|\\

\bitfields{ 1/P, 2/01, 3/111, 2/00, 4/0000, 1/1, 1/1, 2/00, 10/signed10, 6/registerZ }


\paragraph{Syntax} \texttt{dtouchl} \emph{signed10}[\emph{registerZ}]

\paragraph{Description} The effective address is computed by adding the \emph{signed10} to the \emph{registerZ}.
The effective address is sent to the data cache, with request to prefetch the corresponding cache line.


\paragraph{Execution} {Reservation table \textsf{LSU} (Section~\ref{sec:reservation-LSU})}
~
{\begin{lstlisting}
stage ID:
  new offset = _SX_10(signed10);
stage RR:
  new base = GPR[registerZ];
  new address = base + offset;
stage E3:
  MEM.dtouchl(address);
\end{lstlisting}}
\newpage

\paragraph{Encoding} Format \textsf{LSU\_FXBO.X} (Section~\ref{sec:format-LSU-FXBO.X}), \verb|lsucode2 = 0000|\\

\bitfields{ 1/P, 2/00, 2/-- --, 27/upper27, 1/1, 2/01, 3/111, 2/00, 4/0000, 1/1, 1/1, 2/00, 10/lower10, 6/registerZ }


\paragraph{Syntax} \texttt{dtouchl} \emph{upper27\_\discretionary{}{}{}lower10}[\emph{registerZ}]

\paragraph{Description} The effective address is computed by adding the \emph{upper27\_\discretionary{}{}{}lower10} to the \emph{registerZ}.
The effective address is sent to the data cache, with request to prefetch the corresponding cache line.


\paragraph{Execution} {Reservation table \textsf{LSU.X} (Section~\ref{sec:reservation-LSU.X})}
~
{\begin{lstlisting}
stage ID:
  new offset = _SX_37(upper27_lower10);
stage RR:
  new base = GPR[registerZ];
  new address = base + offset;
stage E3:
  MEM.dtouchl(address);
\end{lstlisting}}
\newpage

\paragraph{Encoding} Format \textsf{LSU\_FXBO.Y} (Section~\ref{sec:format-LSU-FXBO.Y}), \verb|lsucode2 = 0000|\\

\bitfields{ 1/P, 2/00, 2/-- --, 27/extend27, 1/1, 2/00, 2/-- --, 27/upper27, 1/1, 2/01, 3/111, 2/00, 4/0000, 1/1, 1/1, 2/00, 10/lower10, 6/registerZ }


\paragraph{Syntax} \texttt{dtouchl} \emph{extend27\_\discretionary{}{}{}upper27\_\discretionary{}{}{}lower10}[\emph{registerZ}]

\paragraph{Description} The effective address is computed by adding the \emph{extend27\_\discretionary{}{}{}upper27\_\discretionary{}{}{}lower10} to the \emph{registerZ}.
The effective address is sent to the data cache, with request to prefetch the corresponding cache line.


\paragraph{Execution} {Reservation table \textsf{LSU.Y} (Section~\ref{sec:reservation-LSU.Y})}
~
{\begin{lstlisting}
stage ID:
  new offset = _WX_64(extend27_upper27_lower10);
stage RR:
  new base = GPR[registerZ];
  new address = base + offset;
stage E3:
  MEM.dtouchl(address);
\end{lstlisting}}
\newpage

\paragraph{Encoding} Format \textsf{LSU\_FXBI} (Section~\ref{sec:format-LSU-FXBI}), \verb|lsucode2 = 0000|\\

\bitfields{ 1/P, 2/01, 3/111, 2/00, 4/0000, 1/1, 1/1, 2/10, 3/111, 1/--, 6/registerY, 6/registerZ }


\paragraph{Syntax} \texttt{dtouchl} \emph{registerY}[\emph{registerZ}]

\paragraph{Description} The effective address is computed by adding the \emph{registerZ} to the \emph{registerY}.
The effective address is sent to the data cache, with request to prefetch the corresponding cache line.


\paragraph{Execution} {Reservation table \textsf{LSU} (Section~\ref{sec:reservation-LSU})}
~
{\begin{lstlisting}
stage RR:
  new doscale = 0;
  new base = GPR[registerZ];
  new index = GPR[registerY];
  new scaling = doscale ? 64 : 1;
  new address = base + (index * scaling);
stage E3:
  MEM.dtouchl(address);
\end{lstlisting}}
\newpage
\subsubsection[{\small FENCE: Memory Fence}]{FENCE\hfill Memory Fence}
\label{instr:FENCE}\index{fence}
 
\paragraph{Encoding} Format \textsf{LSU\_FENCE} (Section~\ref{sec:format-LSU-FENCE}), \verb|lsucode2 = 1111|\\

\bitfields{ 1/P, 2/01, 3/111, 2/accesses, 4/1111, 1/1, 1/1, 2/00, 16/-- -- -- -- -- -- -- -- -- -- -- -- -- -- -- -- }


\paragraph{Syntax} \texttt{fence}\emph{accesses}

\paragraph{Description} Ensures that all issued memory accesses are committed to memory. The \emph{accesses} specifies for what types of access. FENCE.R instructions cannot pass earlier load instructions. FENCE.W instructions cannot pass earlier store, atomic and cache instructions. FENCE instructions produces the effects of FENCE.R, FENCE.W and in addition wait for the completion of outstanding prefetch instructions.


\paragraph{Modifier(s)}
\noindent\begin{itemize}
\item accesses (cf. accesses in section \ref{par:modifier-accesses} on page \pageref{par:modifier-accesses})
\end{itemize}

\paragraph{Execution} {Reservation table \textsf{LSU2\_MEMW} (Section~\ref{sec:reservation-LSU2-MEMW})}
~
{\begin{lstlisting}
stage ID:
  new accesses = _ZX_2(accesses);
stage E3:
  MEM.fence(accesses);
\end{lstlisting}}
\newpage
\subsubsection[{\small I1INVAL: Instruction Cache 1 Invalidate}]{I1INVAL\hfill Instruction Cache 1 Invalidate}
\label{instr:I1INVAL}\index{i1inval}
 
\paragraph{Encoding} Format \textsf{LSU\_MCC} (Section~\ref{sec:format-LSU-MCC}), \verb|lsucode2 = 1100|\\

\bitfields{ 1/P, 2/01, 3/111, 2/00, 4/1100, 1/1, 1/1, 2/00, 16/-- -- -- -- -- -- -- -- -- -- -- -- -- -- -- -- }


\paragraph{Syntax} \texttt{i1inval}

\paragraph{Description} Request to invalidate the instruction cache.


\paragraph{Execution} {Reservation table \textsf{LSU2\_MEMW} (Section~\ref{sec:reservation-LSU2-MEMW})}
~
{\begin{lstlisting}
stage E3:
  MEM.i1inval();
\end{lstlisting}}
\newpage
\subsubsection[{\small I1INVALS: Instruction Cache 1 Invalidate Set}]{I1INVALS\hfill Instruction Cache 1 Invalidate Set}
\label{instr:I1INVALS}\index{i1invals}
 
\paragraph{Encoding} Format \textsf{LSU\_FXBO} (Section~\ref{sec:format-LSU-FXBO}), \verb|lsucode2 = 0101|\\

\bitfields{ 1/P, 2/01, 3/111, 2/00, 4/0101, 1/1, 1/1, 2/00, 10/signed10, 6/registerZ }


\paragraph{Syntax} \texttt{i1invals} \emph{signed10}[\emph{registerZ}]

\paragraph{Description} The effective address is computed by adding the \emph{signed10} to the \emph{registerZ}.
The effective address is sent to the instruction cache, with request to invalidate the corresponding L1 cache set. (All the cache lines with the same index.)


\paragraph{Execution} {Reservation table \textsf{LSU2\_MEMW} (Section~\ref{sec:reservation-LSU2-MEMW})}
~
{\begin{lstlisting}
stage ID:
  new offset = _SX_10(signed10);
stage RR:
  new base = GPR[registerZ];
  new address = base + offset;
stage E3:
  MEM.i1invals(address);
\end{lstlisting}}
\newpage

\paragraph{Encoding} Format \textsf{LSU\_FXBO.X} (Section~\ref{sec:format-LSU-FXBO.X}), \verb|lsucode2 = 0101|\\

\bitfields{ 1/P, 2/00, 2/-- --, 27/upper27, 1/1, 2/01, 3/111, 2/00, 4/0101, 1/1, 1/1, 2/00, 10/lower10, 6/registerZ }


\paragraph{Syntax} \texttt{i1invals} \emph{upper27\_\discretionary{}{}{}lower10}[\emph{registerZ}]

\paragraph{Description} The effective address is computed by adding the \emph{upper27\_\discretionary{}{}{}lower10} to the \emph{registerZ}.
The effective address is sent to the instruction cache, with request to invalidate the corresponding L1 cache set. (All the cache lines with the same index.)


\paragraph{Execution} {Reservation table \textsf{LSU2\_MEMW.X} (Section~\ref{sec:reservation-LSU2-MEMW.X})}
~
{\begin{lstlisting}
stage ID:
  new offset = _SX_37(upper27_lower10);
stage RR:
  new base = GPR[registerZ];
  new address = base + offset;
stage E3:
  MEM.i1invals(address);
\end{lstlisting}}
\newpage

\paragraph{Encoding} Format \textsf{LSU\_FXBO.Y} (Section~\ref{sec:format-LSU-FXBO.Y}), \verb|lsucode2 = 0101|\\

\bitfields{ 1/P, 2/00, 2/-- --, 27/extend27, 1/1, 2/00, 2/-- --, 27/upper27, 1/1, 2/01, 3/111, 2/00, 4/0101, 1/1, 1/1, 2/00, 10/lower10, 6/registerZ }


\paragraph{Syntax} \texttt{i1invals} \emph{extend27\_\discretionary{}{}{}upper27\_\discretionary{}{}{}lower10}[\emph{registerZ}]

\paragraph{Description} The effective address is computed by adding the \emph{extend27\_\discretionary{}{}{}upper27\_\discretionary{}{}{}lower10} to the \emph{registerZ}.
The effective address is sent to the instruction cache, with request to invalidate the corresponding L1 cache set. (All the cache lines with the same index.)


\paragraph{Execution} {Reservation table \textsf{LSU2\_MEMW.Y} (Section~\ref{sec:reservation-LSU2-MEMW.Y})}
~
{\begin{lstlisting}
stage ID:
  new offset = _WX_64(extend27_upper27_lower10);
stage RR:
  new base = GPR[registerZ];
  new address = base + offset;
stage E3:
  MEM.i1invals(address);
\end{lstlisting}}
\newpage

\paragraph{Encoding} Format \textsf{LSU\_FXBI} (Section~\ref{sec:format-LSU-FXBI}), \verb|lsucode2 = 0101|\\

\bitfields{ 1/P, 2/01, 3/111, 2/00, 4/0101, 1/1, 1/1, 2/10, 3/111, 1/--, 6/registerY, 6/registerZ }


\paragraph{Syntax} \texttt{i1invals} \emph{registerY}[\emph{registerZ}]

\paragraph{Description} The effective address is computed by adding the \emph{registerZ} to the \emph{registerY}.
The effective address is sent to the instruction cache, with request to invalidate the corresponding L1 cache set. (All the cache lines with the same index.)


\paragraph{Execution} {Reservation table \textsf{LSU2\_MEMW} (Section~\ref{sec:reservation-LSU2-MEMW})}
~
{\begin{lstlisting}
stage RR:
  new doscale = 0;
  new base = GPR[registerZ];
  new index = GPR[registerY];
  new scaling = doscale ? 64 : 1;
  new address = base + (index * scaling);
stage E3:
  MEM.i1invals(address);
\end{lstlisting}}
\newpage
\subsubsection[{\small LBS: Load Byte Sign Extended}]{LBS\hfill Load Byte Sign Extended}
\label{instr:LBS}\index{lbs}
 
\paragraph{Encoding} Format \textsf{LSU\_LSBO} (Section~\ref{sec:format-LSU-LSBO}), \verb|loadcode = 001|\\

\bitfields{ 1/P, 2/01, 3/001, 2/variant, 6/registerW, 2/00, 10/signed10, 6/registerZ }


\paragraph{Syntax} \texttt{lbs}\emph{variant} \emph{registerW} = \emph{signed10}[\emph{registerZ}]

\paragraph{Description} The effective address is computed by adding the \emph{signed10} to the \emph{registerZ}.
The \emph{registerW} is loaded from the memory byte at the effective address after sign extension.


\paragraph{Modifier(s)}
\noindent\begin{itemize}
\item variant (cf. variant in section \ref{par:modifier-variant} on page \pageref{par:modifier-variant})
\end{itemize}

\paragraph{Execution} {Reservation table \textsf{LSU\_AUXW} (Section~\ref{sec:reservation-LSU-AUXW})}
~
{\begin{lstlisting}
stage ID:
  new variant = _ZX_2(variant);
  new offset = _SX_10(signed10);
stage RR:
  new base = GPR[registerZ];
  new address = base + offset;
  new result1 = _SX_8(MEM.load(address, 0x1, variant, registerW));
stage CR:
  GPR[registerW] = result1;
\end{lstlisting}}
\newpage

\paragraph{Encoding} Format \textsf{LSU\_LSBO.X} (Section~\ref{sec:format-LSU-LSBO.X}), \verb|loadcode = 001|\\

\bitfields{ 1/P, 2/00, 2/-- --, 27/upper27, 1/1, 2/01, 3/001, 2/variant, 6/registerW, 2/00, 10/lower10, 6/registerZ }


\paragraph{Syntax} \texttt{lbs}\emph{variant} \emph{registerW} = \emph{upper27\_\discretionary{}{}{}lower10}[\emph{registerZ}]

\paragraph{Description} The effective address is computed by adding the \emph{upper27\_\discretionary{}{}{}lower10} to the \emph{registerZ}.
The \emph{registerW} is loaded from the memory byte at the effective address after sign extension.


\paragraph{Modifier(s)}
\noindent\begin{itemize}
\item variant (cf. variant in section \ref{par:modifier-variant} on page \pageref{par:modifier-variant})
\end{itemize}

\paragraph{Execution} {Reservation table \textsf{LSU\_AUXW.X} (Section~\ref{sec:reservation-LSU-AUXW.X})}
~
{\begin{lstlisting}
stage ID:
  new variant = _ZX_2(variant);
  new offset = _SX_37(upper27_lower10);
stage RR:
  new base = GPR[registerZ];
  new address = base + offset;
  new result1 = _SX_8(MEM.load(address, 0x1, variant, registerW));
stage CR:
  GPR[registerW] = result1;
\end{lstlisting}}
\newpage

\paragraph{Encoding} Format \textsf{LSU\_LSBO.Y} (Section~\ref{sec:format-LSU-LSBO.Y}), \verb|loadcode = 001|\\

\bitfields{ 1/P, 2/00, 2/-- --, 27/extend27, 1/1, 2/00, 2/-- --, 27/upper27, 1/1, 2/01, 3/001, 2/variant, 6/registerW, 2/00, 10/lower10, 6/registerZ }


\paragraph{Syntax} \texttt{lbs}\emph{variant} \emph{registerW} = \emph{extend27\_\discretionary{}{}{}upper27\_\discretionary{}{}{}lower10}[\emph{registerZ}]

\paragraph{Description} The effective address is computed by adding the \emph{extend27\_\discretionary{}{}{}upper27\_\discretionary{}{}{}lower10} to the \emph{registerZ}.
The \emph{registerW} is loaded from the memory byte at the effective address after sign extension.


\paragraph{Modifier(s)}
\noindent\begin{itemize}
\item variant (cf. variant in section \ref{par:modifier-variant} on page \pageref{par:modifier-variant})
\end{itemize}

\paragraph{Execution} {Reservation table \textsf{LSU\_AUXW.Y} (Section~\ref{sec:reservation-LSU-AUXW.Y})}
~
{\begin{lstlisting}
stage ID:
  new variant = _ZX_2(variant);
  new offset = _WX_64(extend27_upper27_lower10);
stage RR:
  new base = GPR[registerZ];
  new address = base + offset;
  new result1 = _SX_8(MEM.load(address, 0x1, variant, registerW));
stage CR:
  GPR[registerW] = result1;
\end{lstlisting}}
\newpage

\paragraph{Encoding} Format \textsf{LSU\_LSBI} (Section~\ref{sec:format-LSU-LSBI}), \verb|loadcode = 001|\\

\bitfields{ 1/P, 2/01, 3/001, 2/variant, 6/registerW, 2/10, 3/111, 1/--, 6/registerY, 6/registerZ }


\paragraph{Syntax} \texttt{lbs}\emph{variant} \emph{registerW} = \emph{registerY}[\emph{registerZ}]

\paragraph{Description} The effective address is computed by adding the \emph{registerZ} to the \emph{registerY} and scaled by the \emph{variant}.
The \emph{registerW} is loaded from the memory byte at the effective address after sign extension.


\paragraph{Modifier(s)}
\noindent\begin{itemize}
\item variant (cf. variant in section \ref{par:modifier-variant} on page \pageref{par:modifier-variant})
\end{itemize}

\paragraph{Execution} {Reservation table \textsf{LSU\_AUXW} (Section~\ref{sec:reservation-LSU-AUXW})}
~
{\begin{lstlisting}
stage ID:
  new variant = _ZX_2(variant);
stage RR:
  new doscale = 0;
  new base = GPR[registerZ];
  new index = GPR[registerY];
  new scaling = doscale ? 1 : 1;
  new address = base + (index * scaling);
  new result1 = _SX_8(MEM.load(address, 0x1, variant, registerW));
stage CR:
  GPR[registerW] = result1;
\end{lstlisting}}
\newpage
\subsubsection[{\small LBZ: Load Byte Zero Extended}]{LBZ\hfill Load Byte Zero Extended}
\label{instr:LBZ}\index{lbz}
 
\paragraph{Encoding} Format \textsf{LSU\_LSBO} (Section~\ref{sec:format-LSU-LSBO}), \verb|loadcode = 000|\\

\bitfields{ 1/P, 2/01, 3/000, 2/variant, 6/registerW, 2/00, 10/signed10, 6/registerZ }


\paragraph{Syntax} \texttt{lbz}\emph{variant} \emph{registerW} = \emph{signed10}[\emph{registerZ}]

\paragraph{Description} The effective address is computed by adding the \emph{signed10} to the \emph{registerZ}.
The \emph{registerW} is loaded from the memory byte at the effective address after zero extension.


\paragraph{Modifier(s)}
\noindent\begin{itemize}
\item variant (cf. variant in section \ref{par:modifier-variant} on page \pageref{par:modifier-variant})
\end{itemize}

\paragraph{Execution} {Reservation table \textsf{LSU\_AUXW} (Section~\ref{sec:reservation-LSU-AUXW})}
~
{\begin{lstlisting}
stage ID:
  new variant = _ZX_2(variant);
  new offset = _SX_10(signed10);
stage RR:
  new base = GPR[registerZ];
  new address = base + offset;
  new result1 = _ZX_8(MEM.load(address, 0x1, variant, registerW));
stage CR:
  GPR[registerW] = result1;
\end{lstlisting}}
\newpage

\paragraph{Encoding} Format \textsf{LSU\_LSBO.X} (Section~\ref{sec:format-LSU-LSBO.X}), \verb|loadcode = 000|\\

\bitfields{ 1/P, 2/00, 2/-- --, 27/upper27, 1/1, 2/01, 3/000, 2/variant, 6/registerW, 2/00, 10/lower10, 6/registerZ }


\paragraph{Syntax} \texttt{lbz}\emph{variant} \emph{registerW} = \emph{upper27\_\discretionary{}{}{}lower10}[\emph{registerZ}]

\paragraph{Description} The effective address is computed by adding the \emph{upper27\_\discretionary{}{}{}lower10} to the \emph{registerZ}.
The \emph{registerW} is loaded from the memory byte at the effective address after zero extension.


\paragraph{Modifier(s)}
\noindent\begin{itemize}
\item variant (cf. variant in section \ref{par:modifier-variant} on page \pageref{par:modifier-variant})
\end{itemize}

\paragraph{Execution} {Reservation table \textsf{LSU\_AUXW.X} (Section~\ref{sec:reservation-LSU-AUXW.X})}
~
{\begin{lstlisting}
stage ID:
  new variant = _ZX_2(variant);
  new offset = _SX_37(upper27_lower10);
stage RR:
  new base = GPR[registerZ];
  new address = base + offset;
  new result1 = _ZX_8(MEM.load(address, 0x1, variant, registerW));
stage CR:
  GPR[registerW] = result1;
\end{lstlisting}}
\newpage

\paragraph{Encoding} Format \textsf{LSU\_LSBO.Y} (Section~\ref{sec:format-LSU-LSBO.Y}), \verb|loadcode = 000|\\

\bitfields{ 1/P, 2/00, 2/-- --, 27/extend27, 1/1, 2/00, 2/-- --, 27/upper27, 1/1, 2/01, 3/000, 2/variant, 6/registerW, 2/00, 10/lower10, 6/registerZ }


\paragraph{Syntax} \texttt{lbz}\emph{variant} \emph{registerW} = \emph{extend27\_\discretionary{}{}{}upper27\_\discretionary{}{}{}lower10}[\emph{registerZ}]

\paragraph{Description} The effective address is computed by adding the \emph{extend27\_\discretionary{}{}{}upper27\_\discretionary{}{}{}lower10} to the \emph{registerZ}.
The \emph{registerW} is loaded from the memory byte at the effective address after zero extension.


\paragraph{Modifier(s)}
\noindent\begin{itemize}
\item variant (cf. variant in section \ref{par:modifier-variant} on page \pageref{par:modifier-variant})
\end{itemize}

\paragraph{Execution} {Reservation table \textsf{LSU\_AUXW.Y} (Section~\ref{sec:reservation-LSU-AUXW.Y})}
~
{\begin{lstlisting}
stage ID:
  new variant = _ZX_2(variant);
  new offset = _WX_64(extend27_upper27_lower10);
stage RR:
  new base = GPR[registerZ];
  new address = base + offset;
  new result1 = _ZX_8(MEM.load(address, 0x1, variant, registerW));
stage CR:
  GPR[registerW] = result1;
\end{lstlisting}}
\newpage

\paragraph{Encoding} Format \textsf{LSU\_LSBI} (Section~\ref{sec:format-LSU-LSBI}), \verb|loadcode = 000|\\

\bitfields{ 1/P, 2/01, 3/000, 2/variant, 6/registerW, 2/10, 3/111, 1/--, 6/registerY, 6/registerZ }


\paragraph{Syntax} \texttt{lbz}\emph{variant} \emph{registerW} = \emph{registerY}[\emph{registerZ}]

\paragraph{Description} The effective address is computed by adding the \emph{registerZ} to the \emph{registerY} and scaled by the \emph{variant}.
The \emph{registerW} is loaded from the memory byte at the effective address after zero extension.


\paragraph{Modifier(s)}
\noindent\begin{itemize}
\item variant (cf. variant in section \ref{par:modifier-variant} on page \pageref{par:modifier-variant})
\end{itemize}

\paragraph{Execution} {Reservation table \textsf{LSU\_AUXW} (Section~\ref{sec:reservation-LSU-AUXW})}
~
{\begin{lstlisting}
stage ID:
  new variant = _ZX_2(variant);
stage RR:
  new doscale = 0;
  new base = GPR[registerZ];
  new index = GPR[registerY];
  new scaling = doscale ? 1 : 1;
  new address = base + (index * scaling);
  new result1 = _ZX_8(MEM.load(address, 0x1, variant, registerW));
stage CR:
  GPR[registerW] = result1;
\end{lstlisting}}
\newpage
\subsubsection[{\small LD: Load Double Word}]{LD\hfill Load Double Word}
\label{instr:LD}\index{ld}
 
\paragraph{Encoding} Format \textsf{LSU\_LSBO} (Section~\ref{sec:format-LSU-LSBO}), \verb|loadcode = 110|\\

\bitfields{ 1/P, 2/01, 3/110, 2/variant, 6/registerW, 2/00, 10/signed10, 6/registerZ }


\paragraph{Syntax} \texttt{ld}\emph{variant} \emph{registerW} = \emph{signed10}[\emph{registerZ}]

\paragraph{Description} The effective address is computed by adding the \emph{signed10} to the \emph{registerZ}.
The \emph{registerW} is loaded from the memory double word at the effective address.


\paragraph{Modifier(s)}
\noindent\begin{itemize}
\item variant (cf. variant in section \ref{par:modifier-variant} on page \pageref{par:modifier-variant})
\end{itemize}

\paragraph{Execution} {Reservation table \textsf{LSU\_AUXW} (Section~\ref{sec:reservation-LSU-AUXW})}
~
{\begin{lstlisting}
stage ID:
  new variant = _ZX_2(variant);
  new offset = _SX_10(signed10);
stage RR:
  new base = GPR[registerZ];
  new address = base + offset;
  new result1 = MEM.load(address, 0xFF, variant, registerW);
stage CR:
  GPR[registerW] = result1;
\end{lstlisting}}
\newpage

\paragraph{Encoding} Format \textsf{LSU\_LSBO.X} (Section~\ref{sec:format-LSU-LSBO.X}), \verb|loadcode = 110|\\

\bitfields{ 1/P, 2/00, 2/-- --, 27/upper27, 1/1, 2/01, 3/110, 2/variant, 6/registerW, 2/00, 10/lower10, 6/registerZ }


\paragraph{Syntax} \texttt{ld}\emph{variant} \emph{registerW} = \emph{upper27\_\discretionary{}{}{}lower10}[\emph{registerZ}]

\paragraph{Description} The effective address is computed by adding the \emph{upper27\_\discretionary{}{}{}lower10} to the \emph{registerZ}.
The \emph{registerW} is loaded from the memory double word at the effective address.


\paragraph{Modifier(s)}
\noindent\begin{itemize}
\item variant (cf. variant in section \ref{par:modifier-variant} on page \pageref{par:modifier-variant})
\end{itemize}

\paragraph{Execution} {Reservation table \textsf{LSU\_AUXW.X} (Section~\ref{sec:reservation-LSU-AUXW.X})}
~
{\begin{lstlisting}
stage ID:
  new variant = _ZX_2(variant);
  new offset = _SX_37(upper27_lower10);
stage RR:
  new base = GPR[registerZ];
  new address = base + offset;
  new result1 = MEM.load(address, 0xFF, variant, registerW);
stage CR:
  GPR[registerW] = result1;
\end{lstlisting}}
\newpage

\paragraph{Encoding} Format \textsf{LSU\_LSBO.Y} (Section~\ref{sec:format-LSU-LSBO.Y}), \verb|loadcode = 110|\\

\bitfields{ 1/P, 2/00, 2/-- --, 27/extend27, 1/1, 2/00, 2/-- --, 27/upper27, 1/1, 2/01, 3/110, 2/variant, 6/registerW, 2/00, 10/lower10, 6/registerZ }


\paragraph{Syntax} \texttt{ld}\emph{variant} \emph{registerW} = \emph{extend27\_\discretionary{}{}{}upper27\_\discretionary{}{}{}lower10}[\emph{registerZ}]

\paragraph{Description} The effective address is computed by adding the \emph{extend27\_\discretionary{}{}{}upper27\_\discretionary{}{}{}lower10} to the \emph{registerZ}.
The \emph{registerW} is loaded from the memory double word at the effective address.


\paragraph{Modifier(s)}
\noindent\begin{itemize}
\item variant (cf. variant in section \ref{par:modifier-variant} on page \pageref{par:modifier-variant})
\end{itemize}

\paragraph{Execution} {Reservation table \textsf{LSU\_AUXW.Y} (Section~\ref{sec:reservation-LSU-AUXW.Y})}
~
{\begin{lstlisting}
stage ID:
  new variant = _ZX_2(variant);
  new offset = _WX_64(extend27_upper27_lower10);
stage RR:
  new base = GPR[registerZ];
  new address = base + offset;
  new result1 = MEM.load(address, 0xFF, variant, registerW);
stage CR:
  GPR[registerW] = result1;
\end{lstlisting}}
\newpage

\paragraph{Encoding} Format \textsf{LSU\_LSBI} (Section~\ref{sec:format-LSU-LSBI}), \verb|loadcode = 110|\\

\bitfields{ 1/P, 2/01, 3/110, 2/variant, 6/registerW, 2/10, 3/111, 1/--, 6/registerY, 6/registerZ }


\paragraph{Syntax} \texttt{ld}\emph{variant} \emph{registerW} = \emph{registerY}[\emph{registerZ}]

\paragraph{Description} The effective address is computed by adding the \emph{registerZ} to the \emph{registerY} and scaled by the \emph{variant}.
The \emph{registerW} is loaded from the memory double word at the effective address.


\paragraph{Modifier(s)}
\noindent\begin{itemize}
\item variant (cf. variant in section \ref{par:modifier-variant} on page \pageref{par:modifier-variant})
\end{itemize}

\paragraph{Execution} {Reservation table \textsf{LSU\_AUXW} (Section~\ref{sec:reservation-LSU-AUXW})}
~
{\begin{lstlisting}
stage ID:
  new variant = _ZX_2(variant);
stage RR:
  new doscale = 0;
  new base = GPR[registerZ];
  new index = GPR[registerY];
  new scaling = doscale ? 8 : 1;
  new address = base + (index * scaling);
  new result1 = MEM.load(address, 0xFF, variant, registerW);
stage CR:
  GPR[registerW] = result1;
\end{lstlisting}}
\newpage
\subsubsection[{\small LHS: Load Half Word Sign Extended}]{LHS\hfill Load Half Word Sign Extended}
\label{instr:LHS}\index{lhs}
 
\paragraph{Encoding} Format \textsf{LSU\_LSBO} (Section~\ref{sec:format-LSU-LSBO}), \verb|loadcode = 011|\\

\bitfields{ 1/P, 2/01, 3/011, 2/variant, 6/registerW, 2/00, 10/signed10, 6/registerZ }


\paragraph{Syntax} \texttt{lhs}\emph{variant} \emph{registerW} = \emph{signed10}[\emph{registerZ}]

\paragraph{Description} The effective address is computed by adding the \emph{signed10} to the \emph{registerZ}.
The \emph{registerW} is loaded from the memory half word at the effective address after sign extension.


\paragraph{Modifier(s)}
\noindent\begin{itemize}
\item variant (cf. variant in section \ref{par:modifier-variant} on page \pageref{par:modifier-variant})
\end{itemize}

\paragraph{Execution} {Reservation table \textsf{LSU\_AUXW} (Section~\ref{sec:reservation-LSU-AUXW})}
~
{\begin{lstlisting}
stage ID:
  new variant = _ZX_2(variant);
  new offset = _SX_10(signed10);
stage RR:
  new base = GPR[registerZ];
  new address = base + offset;
  new result1 = _SX_16(MEM.load(address, 0x3, variant, registerW));
stage CR:
  GPR[registerW] = result1;
\end{lstlisting}}
\newpage

\paragraph{Encoding} Format \textsf{LSU\_LSBO.X} (Section~\ref{sec:format-LSU-LSBO.X}), \verb|loadcode = 011|\\

\bitfields{ 1/P, 2/00, 2/-- --, 27/upper27, 1/1, 2/01, 3/011, 2/variant, 6/registerW, 2/00, 10/lower10, 6/registerZ }


\paragraph{Syntax} \texttt{lhs}\emph{variant} \emph{registerW} = \emph{upper27\_\discretionary{}{}{}lower10}[\emph{registerZ}]

\paragraph{Description} The effective address is computed by adding the \emph{upper27\_\discretionary{}{}{}lower10} to the \emph{registerZ}.
The \emph{registerW} is loaded from the memory half word at the effective address after sign extension.


\paragraph{Modifier(s)}
\noindent\begin{itemize}
\item variant (cf. variant in section \ref{par:modifier-variant} on page \pageref{par:modifier-variant})
\end{itemize}

\paragraph{Execution} {Reservation table \textsf{LSU\_AUXW.X} (Section~\ref{sec:reservation-LSU-AUXW.X})}
~
{\begin{lstlisting}
stage ID:
  new variant = _ZX_2(variant);
  new offset = _SX_37(upper27_lower10);
stage RR:
  new base = GPR[registerZ];
  new address = base + offset;
  new result1 = _SX_16(MEM.load(address, 0x3, variant, registerW));
stage CR:
  GPR[registerW] = result1;
\end{lstlisting}}
\newpage

\paragraph{Encoding} Format \textsf{LSU\_LSBO.Y} (Section~\ref{sec:format-LSU-LSBO.Y}), \verb|loadcode = 011|\\

\bitfields{ 1/P, 2/00, 2/-- --, 27/extend27, 1/1, 2/00, 2/-- --, 27/upper27, 1/1, 2/01, 3/011, 2/variant, 6/registerW, 2/00, 10/lower10, 6/registerZ }


\paragraph{Syntax} \texttt{lhs}\emph{variant} \emph{registerW} = \emph{extend27\_\discretionary{}{}{}upper27\_\discretionary{}{}{}lower10}[\emph{registerZ}]

\paragraph{Description} The effective address is computed by adding the \emph{extend27\_\discretionary{}{}{}upper27\_\discretionary{}{}{}lower10} to the \emph{registerZ}.
The \emph{registerW} is loaded from the memory half word at the effective address after sign extension.


\paragraph{Modifier(s)}
\noindent\begin{itemize}
\item variant (cf. variant in section \ref{par:modifier-variant} on page \pageref{par:modifier-variant})
\end{itemize}

\paragraph{Execution} {Reservation table \textsf{LSU\_AUXW.Y} (Section~\ref{sec:reservation-LSU-AUXW.Y})}
~
{\begin{lstlisting}
stage ID:
  new variant = _ZX_2(variant);
  new offset = _WX_64(extend27_upper27_lower10);
stage RR:
  new base = GPR[registerZ];
  new address = base + offset;
  new result1 = _SX_16(MEM.load(address, 0x3, variant, registerW));
stage CR:
  GPR[registerW] = result1;
\end{lstlisting}}
\newpage

\paragraph{Encoding} Format \textsf{LSU\_LSBI} (Section~\ref{sec:format-LSU-LSBI}), \verb|loadcode = 011|\\

\bitfields{ 1/P, 2/01, 3/011, 2/variant, 6/registerW, 2/10, 3/111, 1/--, 6/registerY, 6/registerZ }


\paragraph{Syntax} \texttt{lhs}\emph{variant} \emph{registerW} = \emph{registerY}[\emph{registerZ}]

\paragraph{Description} The effective address is computed by adding the \emph{registerZ} to the \emph{registerY} and scaled by the \emph{variant}.
The \emph{registerW} is loaded from the memory half word at the effective address after sign extension.


\paragraph{Modifier(s)}
\noindent\begin{itemize}
\item variant (cf. variant in section \ref{par:modifier-variant} on page \pageref{par:modifier-variant})
\end{itemize}

\paragraph{Execution} {Reservation table \textsf{LSU\_AUXW} (Section~\ref{sec:reservation-LSU-AUXW})}
~
{\begin{lstlisting}
stage ID:
  new variant = _ZX_2(variant);
stage RR:
  new doscale = 0;
  new base = GPR[registerZ];
  new index = GPR[registerY];
  new scaling = doscale ? 2 : 1;
  new address = base + (index * scaling);
  new result1 = _SX_16(MEM.load(address, 0x3, variant, registerW));
stage CR:
  GPR[registerW] = result1;
\end{lstlisting}}
\newpage
\subsubsection[{\small LHZ: Load Half Word Zero Extended}]{LHZ\hfill Load Half Word Zero Extended}
\label{instr:LHZ}\index{lhz}
 
\paragraph{Encoding} Format \textsf{LSU\_LSBO} (Section~\ref{sec:format-LSU-LSBO}), \verb|loadcode = 010|\\

\bitfields{ 1/P, 2/01, 3/010, 2/variant, 6/registerW, 2/00, 10/signed10, 6/registerZ }


\paragraph{Syntax} \texttt{lhz}\emph{variant} \emph{registerW} = \emph{signed10}[\emph{registerZ}]

\paragraph{Description} The effective address is computed by adding the \emph{signed10} to the \emph{registerZ}.
The \emph{registerW} is loaded from the memory half word at the effective address after zero extension.


\paragraph{Modifier(s)}
\noindent\begin{itemize}
\item variant (cf. variant in section \ref{par:modifier-variant} on page \pageref{par:modifier-variant})
\end{itemize}

\paragraph{Execution} {Reservation table \textsf{LSU\_AUXW} (Section~\ref{sec:reservation-LSU-AUXW})}
~
{\begin{lstlisting}
stage ID:
  new variant = _ZX_2(variant);
  new offset = _SX_10(signed10);
stage RR:
  new base = GPR[registerZ];
  new address = base + offset;
  new result1 = _ZX_16(MEM.load(address, 0x3, variant, registerW));
stage CR:
  GPR[registerW] = result1;
\end{lstlisting}}
\newpage

\paragraph{Encoding} Format \textsf{LSU\_LSBO.X} (Section~\ref{sec:format-LSU-LSBO.X}), \verb|loadcode = 010|\\

\bitfields{ 1/P, 2/00, 2/-- --, 27/upper27, 1/1, 2/01, 3/010, 2/variant, 6/registerW, 2/00, 10/lower10, 6/registerZ }


\paragraph{Syntax} \texttt{lhz}\emph{variant} \emph{registerW} = \emph{upper27\_\discretionary{}{}{}lower10}[\emph{registerZ}]

\paragraph{Description} The effective address is computed by adding the \emph{upper27\_\discretionary{}{}{}lower10} to the \emph{registerZ}.
The \emph{registerW} is loaded from the memory half word at the effective address after zero extension.


\paragraph{Modifier(s)}
\noindent\begin{itemize}
\item variant (cf. variant in section \ref{par:modifier-variant} on page \pageref{par:modifier-variant})
\end{itemize}

\paragraph{Execution} {Reservation table \textsf{LSU\_AUXW.X} (Section~\ref{sec:reservation-LSU-AUXW.X})}
~
{\begin{lstlisting}
stage ID:
  new variant = _ZX_2(variant);
  new offset = _SX_37(upper27_lower10);
stage RR:
  new base = GPR[registerZ];
  new address = base + offset;
  new result1 = _ZX_16(MEM.load(address, 0x3, variant, registerW));
stage CR:
  GPR[registerW] = result1;
\end{lstlisting}}
\newpage

\paragraph{Encoding} Format \textsf{LSU\_LSBO.Y} (Section~\ref{sec:format-LSU-LSBO.Y}), \verb|loadcode = 010|\\

\bitfields{ 1/P, 2/00, 2/-- --, 27/extend27, 1/1, 2/00, 2/-- --, 27/upper27, 1/1, 2/01, 3/010, 2/variant, 6/registerW, 2/00, 10/lower10, 6/registerZ }


\paragraph{Syntax} \texttt{lhz}\emph{variant} \emph{registerW} = \emph{extend27\_\discretionary{}{}{}upper27\_\discretionary{}{}{}lower10}[\emph{registerZ}]

\paragraph{Description} The effective address is computed by adding the \emph{extend27\_\discretionary{}{}{}upper27\_\discretionary{}{}{}lower10} to the \emph{registerZ}.
The \emph{registerW} is loaded from the memory half word at the effective address after zero extension.


\paragraph{Modifier(s)}
\noindent\begin{itemize}
\item variant (cf. variant in section \ref{par:modifier-variant} on page \pageref{par:modifier-variant})
\end{itemize}

\paragraph{Execution} {Reservation table \textsf{LSU\_AUXW.Y} (Section~\ref{sec:reservation-LSU-AUXW.Y})}
~
{\begin{lstlisting}
stage ID:
  new variant = _ZX_2(variant);
  new offset = _WX_64(extend27_upper27_lower10);
stage RR:
  new base = GPR[registerZ];
  new address = base + offset;
  new result1 = _ZX_16(MEM.load(address, 0x3, variant, registerW));
stage CR:
  GPR[registerW] = result1;
\end{lstlisting}}
\newpage

\paragraph{Encoding} Format \textsf{LSU\_LSBI} (Section~\ref{sec:format-LSU-LSBI}), \verb|loadcode = 010|\\

\bitfields{ 1/P, 2/01, 3/010, 2/variant, 6/registerW, 2/10, 3/111, 1/--, 6/registerY, 6/registerZ }


\paragraph{Syntax} \texttt{lhz}\emph{variant} \emph{registerW} = \emph{registerY}[\emph{registerZ}]

\paragraph{Description} The effective address is computed by adding the \emph{registerZ} to the \emph{registerY} and scaled by the \emph{variant}.
The \emph{registerW} is loaded from the memory half word at the effective address after zero extension.


\paragraph{Modifier(s)}
\noindent\begin{itemize}
\item variant (cf. variant in section \ref{par:modifier-variant} on page \pageref{par:modifier-variant})
\end{itemize}

\paragraph{Execution} {Reservation table \textsf{LSU\_AUXW} (Section~\ref{sec:reservation-LSU-AUXW})}
~
{\begin{lstlisting}
stage ID:
  new variant = _ZX_2(variant);
stage RR:
  new doscale = 0;
  new base = GPR[registerZ];
  new index = GPR[registerY];
  new scaling = doscale ? 2 : 1;
  new address = base + (index * scaling);
  new result1 = _ZX_16(MEM.load(address, 0x3, variant, registerW));
stage CR:
  GPR[registerW] = result1;
\end{lstlisting}}
\newpage
\subsubsection[{\small LO: Load Octuple Word}]{LO\hfill Load Octuple Word}
\label{instr:LO}\index{lo}
 
\paragraph{Encoding} Format \textsf{LSU\_LOBO} (Section~\ref{sec:format-LSU-LOBO}), \verb|loadcode = 111|\\

\bitfields{ 1/P, 2/01, 3/111, 2/variant, 4/registerN, 1/0, 1/1, 2/00, 10/signed10, 6/registerZ }


\paragraph{Syntax} \texttt{lo}\emph{variant} \emph{registerN} = \emph{signed10}[\emph{registerZ}]

\paragraph{Description} The effective address is computed by adding the \emph{signed10} to the \emph{registerZ}.
The \emph{registerN} is loaded from the memory octuple word at the effective address.


\paragraph{Modifier(s)}
\noindent\begin{itemize}
\item variant (cf. variant in section \ref{par:modifier-variant} on page \pageref{par:modifier-variant})
\end{itemize}

\paragraph{Execution} {Reservation table \textsf{LSU\_AUXW} (Section~\ref{sec:reservation-LSU-AUXW})}
~
{\begin{lstlisting}
stage ID:
  new variant = _ZX_2(variant);
  new offset = _SX_10(signed10);
stage RR:
  new bytemask = 0xFFFFFFFF;
  new base = GPR[registerZ];
  new address = base + offset;
  new result1 = MEM.load(address, 0xFFFFFFFF & bytemask, variant, registerN << 2);
stage CR:
  QGR[registerN] = result1;
\end{lstlisting}}
\newpage

\paragraph{Encoding} Format \textsf{LSU\_LOBO.X} (Section~\ref{sec:format-LSU-LOBO.X}), \verb|loadcode = 111|\\

\bitfields{ 1/P, 2/00, 2/-- --, 27/upper27, 1/1, 2/01, 3/111, 2/variant, 4/registerN, 1/0, 1/1, 2/00, 10/lower10, 6/registerZ }


\paragraph{Syntax} \texttt{lo}\emph{variant} \emph{registerN} = \emph{upper27\_\discretionary{}{}{}lower10}[\emph{registerZ}]

\paragraph{Description} The effective address is computed by adding the \emph{upper27\_\discretionary{}{}{}lower10} to the \emph{registerZ}.
The \emph{registerN} is loaded from the memory octuple word at the effective address.


\paragraph{Modifier(s)}
\noindent\begin{itemize}
\item variant (cf. variant in section \ref{par:modifier-variant} on page \pageref{par:modifier-variant})
\end{itemize}

\paragraph{Execution} {Reservation table \textsf{LSU\_AUXW.X} (Section~\ref{sec:reservation-LSU-AUXW.X})}
~
{\begin{lstlisting}
stage ID:
  new variant = _ZX_2(variant);
  new offset = _SX_37(upper27_lower10);
stage RR:
  new bytemask = 0xFFFFFFFF;
  new base = GPR[registerZ];
  new address = base + offset;
  new result1 = MEM.load(address, 0xFFFFFFFF & bytemask, variant, registerN << 2);
stage CR:
  QGR[registerN] = result1;
\end{lstlisting}}
\newpage

\paragraph{Encoding} Format \textsf{LSU\_LOBO.Y} (Section~\ref{sec:format-LSU-LOBO.Y}), \verb|loadcode = 111|\\

\bitfields{ 1/P, 2/00, 2/-- --, 27/extend27, 1/1, 2/00, 2/-- --, 27/upper27, 1/1, 2/01, 3/111, 2/variant, 4/registerN, 1/0, 1/1, 2/00, 10/lower10, 6/registerZ }


\paragraph{Syntax} \texttt{lo}\emph{variant} \emph{registerN} = \emph{extend27\_\discretionary{}{}{}upper27\_\discretionary{}{}{}lower10}[\emph{registerZ}]

\paragraph{Description} The effective address is computed by adding the \emph{extend27\_\discretionary{}{}{}upper27\_\discretionary{}{}{}lower10} to the \emph{registerZ}.
The \emph{registerN} is loaded from the memory octuple word at the effective address.


\paragraph{Modifier(s)}
\noindent\begin{itemize}
\item variant (cf. variant in section \ref{par:modifier-variant} on page \pageref{par:modifier-variant})
\end{itemize}

\paragraph{Execution} {Reservation table \textsf{LSU\_AUXW.Y} (Section~\ref{sec:reservation-LSU-AUXW.Y})}
~
{\begin{lstlisting}
stage ID:
  new variant = _ZX_2(variant);
  new offset = _WX_64(extend27_upper27_lower10);
stage RR:
  new bytemask = 0xFFFFFFFF;
  new base = GPR[registerZ];
  new address = base + offset;
  new result1 = MEM.load(address, 0xFFFFFFFF & bytemask, variant, registerN << 2);
stage CR:
  QGR[registerN] = result1;
\end{lstlisting}}
\newpage

\paragraph{Encoding} Format \textsf{LSU\_LOBI} (Section~\ref{sec:format-LSU-LOBI}), \verb|loadcode = 111|\\

\bitfields{ 1/P, 2/01, 3/111, 2/variant, 4/registerN, 1/0, 1/1, 2/10, 3/111, 1/--, 6/registerY, 6/registerZ }


\paragraph{Syntax} \texttt{lo}\emph{variant} \emph{registerN} = \emph{registerY}[\emph{registerZ}]

\paragraph{Description} The effective address is computed by adding the \emph{registerZ} to the \emph{registerY} and scaled by the \emph{variant}.
The \emph{registerN} is loaded from the memory octuple word at the effective address.


\paragraph{Modifier(s)}
\noindent\begin{itemize}
\item variant (cf. variant in section \ref{par:modifier-variant} on page \pageref{par:modifier-variant})
\end{itemize}

\paragraph{Execution} {Reservation table \textsf{LSU\_AUXW} (Section~\ref{sec:reservation-LSU-AUXW})}
~
{\begin{lstlisting}
stage ID:
  new variant = _ZX_2(variant);
stage RR:
  new bytemask = 0xFFFFFFFF;
  new doscale = 0;
  new base = GPR[registerZ];
  new index = GPR[registerY];
  new scaling = doscale ? 32 : 1;
  new address = base + (index * scaling);
  new result1 = MEM.load(address, 0xFFFFFFFF & bytemask, variant, registerN << 2);
stage CR:
  QGR[registerN] = result1;
\end{lstlisting}}
\newpage
\subsubsection[{\small LQ: Load Quadruple Word}]{LQ\hfill Load Quadruple Word}
\label{instr:LQ}\index{lq}
 
\paragraph{Encoding} Format \textsf{LSU\_LQBO} (Section~\ref{sec:format-LSU-LQBO}), \verb|loadcode = 111|\\

\bitfields{ 1/P, 2/01, 3/111, 2/variant, 5/registerM, 1/0, 2/00, 10/signed10, 6/registerZ }


\paragraph{Syntax} \texttt{lq}\emph{variant} \emph{registerM} = \emph{signed10}[\emph{registerZ}]

\paragraph{Description} The effective address is computed by adding the \emph{signed10} to the \emph{registerZ}.
The \emph{registerM} is loaded from the memory quadruple word at the effective address.


\paragraph{Modifier(s)}
\noindent\begin{itemize}
\item variant (cf. variant in section \ref{par:modifier-variant} on page \pageref{par:modifier-variant})
\end{itemize}

\paragraph{Execution} {Reservation table \textsf{LSU\_AUXW} (Section~\ref{sec:reservation-LSU-AUXW})}
~
{\begin{lstlisting}
stage ID:
  new variant = _ZX_2(variant);
  new offset = _SX_10(signed10);
stage RR:
  new base = GPR[registerZ];
  new address = base + offset;
  new result1 = MEM.load(address, 0xFFFF, variant, registerM << 1);
stage CR:
  PGR[registerM] = result1;
\end{lstlisting}}
\newpage

\paragraph{Encoding} Format \textsf{LSU\_LQBO.X} (Section~\ref{sec:format-LSU-LQBO.X}), \verb|loadcode = 111|\\

\bitfields{ 1/P, 2/00, 2/-- --, 27/upper27, 1/1, 2/01, 3/111, 2/variant, 5/registerM, 1/0, 2/00, 10/lower10, 6/registerZ }


\paragraph{Syntax} \texttt{lq}\emph{variant} \emph{registerM} = \emph{upper27\_\discretionary{}{}{}lower10}[\emph{registerZ}]

\paragraph{Description} The effective address is computed by adding the \emph{upper27\_\discretionary{}{}{}lower10} to the \emph{registerZ}.
The \emph{registerM} is loaded from the memory quadruple word at the effective address.


\paragraph{Modifier(s)}
\noindent\begin{itemize}
\item variant (cf. variant in section \ref{par:modifier-variant} on page \pageref{par:modifier-variant})
\end{itemize}

\paragraph{Execution} {Reservation table \textsf{LSU\_AUXW.X} (Section~\ref{sec:reservation-LSU-AUXW.X})}
~
{\begin{lstlisting}
stage ID:
  new variant = _ZX_2(variant);
  new offset = _SX_37(upper27_lower10);
stage RR:
  new base = GPR[registerZ];
  new address = base + offset;
  new result1 = MEM.load(address, 0xFFFF, variant, registerM << 1);
stage CR:
  PGR[registerM] = result1;
\end{lstlisting}}
\newpage

\paragraph{Encoding} Format \textsf{LSU\_LQBO.Y} (Section~\ref{sec:format-LSU-LQBO.Y}), \verb|loadcode = 111|\\

\bitfields{ 1/P, 2/00, 2/-- --, 27/extend27, 1/1, 2/00, 2/-- --, 27/upper27, 1/1, 2/01, 3/111, 2/variant, 5/registerM, 1/0, 2/00, 10/lower10, 6/registerZ }


\paragraph{Syntax} \texttt{lq}\emph{variant} \emph{registerM} = \emph{extend27\_\discretionary{}{}{}upper27\_\discretionary{}{}{}lower10}[\emph{registerZ}]

\paragraph{Description} The effective address is computed by adding the \emph{extend27\_\discretionary{}{}{}upper27\_\discretionary{}{}{}lower10} to the \emph{registerZ}.
The \emph{registerM} is loaded from the memory quadruple word at the effective address.


\paragraph{Modifier(s)}
\noindent\begin{itemize}
\item variant (cf. variant in section \ref{par:modifier-variant} on page \pageref{par:modifier-variant})
\end{itemize}

\paragraph{Execution} {Reservation table \textsf{LSU\_AUXW.Y} (Section~\ref{sec:reservation-LSU-AUXW.Y})}
~
{\begin{lstlisting}
stage ID:
  new variant = _ZX_2(variant);
  new offset = _WX_64(extend27_upper27_lower10);
stage RR:
  new base = GPR[registerZ];
  new address = base + offset;
  new result1 = MEM.load(address, 0xFFFF, variant, registerM << 1);
stage CR:
  PGR[registerM] = result1;
\end{lstlisting}}
\newpage

\paragraph{Encoding} Format \textsf{LSU\_LQBI} (Section~\ref{sec:format-LSU-LQBI}), \verb|loadcode = 111|\\

\bitfields{ 1/P, 2/01, 3/111, 2/variant, 5/registerM, 1/0, 2/10, 3/111, 1/--, 6/registerY, 6/registerZ }


\paragraph{Syntax} \texttt{lq}\emph{variant} \emph{registerM} = \emph{registerY}[\emph{registerZ}]

\paragraph{Description} The effective address is computed by adding the \emph{registerZ} to the \emph{registerY} and scaled by the \emph{variant}.
The \emph{registerM} is loaded from the memory quadruple word at the effective address.


\paragraph{Modifier(s)}
\noindent\begin{itemize}
\item variant (cf. variant in section \ref{par:modifier-variant} on page \pageref{par:modifier-variant})
\end{itemize}

\paragraph{Execution} {Reservation table \textsf{LSU\_AUXW} (Section~\ref{sec:reservation-LSU-AUXW})}
~
{\begin{lstlisting}
stage ID:
  new variant = _ZX_2(variant);
stage RR:
  new doscale = 0;
  new base = GPR[registerZ];
  new index = GPR[registerY];
  new scaling = doscale ? 16 : 1;
  new address = base + (index * scaling);
  new result1 = MEM.load(address, 0xFFFF, variant, registerM << 1);
stage CR:
  PGR[registerM] = result1;
\end{lstlisting}}
\newpage
\subsubsection[{\small LWS: Load Word Sign Extended}]{LWS\hfill Load Word Sign Extended}
\label{instr:LWS}\index{lws}
 
\paragraph{Encoding} Format \textsf{LSU\_LSBO} (Section~\ref{sec:format-LSU-LSBO}), \verb|loadcode = 101|\\

\bitfields{ 1/P, 2/01, 3/101, 2/variant, 6/registerW, 2/00, 10/signed10, 6/registerZ }


\paragraph{Syntax} \texttt{lws}\emph{variant} \emph{registerW} = \emph{signed10}[\emph{registerZ}]

\paragraph{Description} The effective address is computed by adding the \emph{signed10} to the \emph{registerZ}.
The \emph{registerW} is loaded from the memory word at the effective address after sign extension.


\paragraph{Modifier(s)}
\noindent\begin{itemize}
\item variant (cf. variant in section \ref{par:modifier-variant} on page \pageref{par:modifier-variant})
\end{itemize}

\paragraph{Execution} {Reservation table \textsf{LSU\_AUXW} (Section~\ref{sec:reservation-LSU-AUXW})}
~
{\begin{lstlisting}
stage ID:
  new variant = _ZX_2(variant);
  new offset = _SX_10(signed10);
stage RR:
  new base = GPR[registerZ];
  new address = base + offset;
  new result1 = _SX_32(MEM.load(address, 0xF, variant, registerW));
stage CR:
  GPR[registerW] = result1;
\end{lstlisting}}
\newpage

\paragraph{Encoding} Format \textsf{LSU\_LSBO.X} (Section~\ref{sec:format-LSU-LSBO.X}), \verb|loadcode = 101|\\

\bitfields{ 1/P, 2/00, 2/-- --, 27/upper27, 1/1, 2/01, 3/101, 2/variant, 6/registerW, 2/00, 10/lower10, 6/registerZ }


\paragraph{Syntax} \texttt{lws}\emph{variant} \emph{registerW} = \emph{upper27\_\discretionary{}{}{}lower10}[\emph{registerZ}]

\paragraph{Description} The effective address is computed by adding the \emph{upper27\_\discretionary{}{}{}lower10} to the \emph{registerZ}.
The \emph{registerW} is loaded from the memory word at the effective address after sign extension.


\paragraph{Modifier(s)}
\noindent\begin{itemize}
\item variant (cf. variant in section \ref{par:modifier-variant} on page \pageref{par:modifier-variant})
\end{itemize}

\paragraph{Execution} {Reservation table \textsf{LSU\_AUXW.X} (Section~\ref{sec:reservation-LSU-AUXW.X})}
~
{\begin{lstlisting}
stage ID:
  new variant = _ZX_2(variant);
  new offset = _SX_37(upper27_lower10);
stage RR:
  new base = GPR[registerZ];
  new address = base + offset;
  new result1 = _SX_32(MEM.load(address, 0xF, variant, registerW));
stage CR:
  GPR[registerW] = result1;
\end{lstlisting}}
\newpage

\paragraph{Encoding} Format \textsf{LSU\_LSBO.Y} (Section~\ref{sec:format-LSU-LSBO.Y}), \verb|loadcode = 101|\\

\bitfields{ 1/P, 2/00, 2/-- --, 27/extend27, 1/1, 2/00, 2/-- --, 27/upper27, 1/1, 2/01, 3/101, 2/variant, 6/registerW, 2/00, 10/lower10, 6/registerZ }


\paragraph{Syntax} \texttt{lws}\emph{variant} \emph{registerW} = \emph{extend27\_\discretionary{}{}{}upper27\_\discretionary{}{}{}lower10}[\emph{registerZ}]

\paragraph{Description} The effective address is computed by adding the \emph{extend27\_\discretionary{}{}{}upper27\_\discretionary{}{}{}lower10} to the \emph{registerZ}.
The \emph{registerW} is loaded from the memory word at the effective address after sign extension.


\paragraph{Modifier(s)}
\noindent\begin{itemize}
\item variant (cf. variant in section \ref{par:modifier-variant} on page \pageref{par:modifier-variant})
\end{itemize}

\paragraph{Execution} {Reservation table \textsf{LSU\_AUXW.Y} (Section~\ref{sec:reservation-LSU-AUXW.Y})}
~
{\begin{lstlisting}
stage ID:
  new variant = _ZX_2(variant);
  new offset = _WX_64(extend27_upper27_lower10);
stage RR:
  new base = GPR[registerZ];
  new address = base + offset;
  new result1 = _SX_32(MEM.load(address, 0xF, variant, registerW));
stage CR:
  GPR[registerW] = result1;
\end{lstlisting}}
\newpage

\paragraph{Encoding} Format \textsf{LSU\_LSBI} (Section~\ref{sec:format-LSU-LSBI}), \verb|loadcode = 101|\\

\bitfields{ 1/P, 2/01, 3/101, 2/variant, 6/registerW, 2/10, 3/111, 1/--, 6/registerY, 6/registerZ }


\paragraph{Syntax} \texttt{lws}\emph{variant} \emph{registerW} = \emph{registerY}[\emph{registerZ}]

\paragraph{Description} The effective address is computed by adding the \emph{registerZ} to the \emph{registerY} and scaled by the \emph{variant}.
The \emph{registerW} is loaded from the memory word at the effective address after sign extension.


\paragraph{Modifier(s)}
\noindent\begin{itemize}
\item variant (cf. variant in section \ref{par:modifier-variant} on page \pageref{par:modifier-variant})
\end{itemize}

\paragraph{Execution} {Reservation table \textsf{LSU\_AUXW} (Section~\ref{sec:reservation-LSU-AUXW})}
~
{\begin{lstlisting}
stage ID:
  new variant = _ZX_2(variant);
stage RR:
  new doscale = 0;
  new base = GPR[registerZ];
  new index = GPR[registerY];
  new scaling = doscale ? 4 : 1;
  new address = base + (index * scaling);
  new result1 = _SX_32(MEM.load(address, 0xF, variant, registerW));
stage CR:
  GPR[registerW] = result1;
\end{lstlisting}}
\newpage
\subsubsection[{\small LWZ: Load Word Zero Extended}]{LWZ\hfill Load Word Zero Extended}
\label{instr:LWZ}\index{lwz}
 
\paragraph{Encoding} Format \textsf{LSU\_LSBO} (Section~\ref{sec:format-LSU-LSBO}), \verb|loadcode = 100|\\

\bitfields{ 1/P, 2/01, 3/100, 2/variant, 6/registerW, 2/00, 10/signed10, 6/registerZ }


\paragraph{Syntax} \texttt{lwz}\emph{variant} \emph{registerW} = \emph{signed10}[\emph{registerZ}]

\paragraph{Description} The effective address is computed by adding the \emph{signed10} to the \emph{registerZ}.
The \emph{registerW} is loaded from the memory word at the effective address after zero extension.


\paragraph{Modifier(s)}
\noindent\begin{itemize}
\item variant (cf. variant in section \ref{par:modifier-variant} on page \pageref{par:modifier-variant})
\end{itemize}

\paragraph{Execution} {Reservation table \textsf{LSU\_AUXW} (Section~\ref{sec:reservation-LSU-AUXW})}
~
{\begin{lstlisting}
stage ID:
  new variant = _ZX_2(variant);
  new offset = _SX_10(signed10);
stage RR:
  new base = GPR[registerZ];
  new address = base + offset;
  new result1 = _ZX_32(MEM.load(address, 0xF, variant, registerW));
stage CR:
  GPR[registerW] = result1;
\end{lstlisting}}
\newpage

\paragraph{Encoding} Format \textsf{LSU\_LSBO.X} (Section~\ref{sec:format-LSU-LSBO.X}), \verb|loadcode = 100|\\

\bitfields{ 1/P, 2/00, 2/-- --, 27/upper27, 1/1, 2/01, 3/100, 2/variant, 6/registerW, 2/00, 10/lower10, 6/registerZ }


\paragraph{Syntax} \texttt{lwz}\emph{variant} \emph{registerW} = \emph{upper27\_\discretionary{}{}{}lower10}[\emph{registerZ}]

\paragraph{Description} The effective address is computed by adding the \emph{upper27\_\discretionary{}{}{}lower10} to the \emph{registerZ}.
The \emph{registerW} is loaded from the memory word at the effective address after zero extension.


\paragraph{Modifier(s)}
\noindent\begin{itemize}
\item variant (cf. variant in section \ref{par:modifier-variant} on page \pageref{par:modifier-variant})
\end{itemize}

\paragraph{Execution} {Reservation table \textsf{LSU\_AUXW.X} (Section~\ref{sec:reservation-LSU-AUXW.X})}
~
{\begin{lstlisting}
stage ID:
  new variant = _ZX_2(variant);
  new offset = _SX_37(upper27_lower10);
stage RR:
  new base = GPR[registerZ];
  new address = base + offset;
  new result1 = _ZX_32(MEM.load(address, 0xF, variant, registerW));
stage CR:
  GPR[registerW] = result1;
\end{lstlisting}}
\newpage

\paragraph{Encoding} Format \textsf{LSU\_LSBO.Y} (Section~\ref{sec:format-LSU-LSBO.Y}), \verb|loadcode = 100|\\

\bitfields{ 1/P, 2/00, 2/-- --, 27/extend27, 1/1, 2/00, 2/-- --, 27/upper27, 1/1, 2/01, 3/100, 2/variant, 6/registerW, 2/00, 10/lower10, 6/registerZ }


\paragraph{Syntax} \texttt{lwz}\emph{variant} \emph{registerW} = \emph{extend27\_\discretionary{}{}{}upper27\_\discretionary{}{}{}lower10}[\emph{registerZ}]

\paragraph{Description} The effective address is computed by adding the \emph{extend27\_\discretionary{}{}{}upper27\_\discretionary{}{}{}lower10} to the \emph{registerZ}.
The \emph{registerW} is loaded from the memory word at the effective address after zero extension.


\paragraph{Modifier(s)}
\noindent\begin{itemize}
\item variant (cf. variant in section \ref{par:modifier-variant} on page \pageref{par:modifier-variant})
\end{itemize}

\paragraph{Execution} {Reservation table \textsf{LSU\_AUXW.Y} (Section~\ref{sec:reservation-LSU-AUXW.Y})}
~
{\begin{lstlisting}
stage ID:
  new variant = _ZX_2(variant);
  new offset = _WX_64(extend27_upper27_lower10);
stage RR:
  new base = GPR[registerZ];
  new address = base + offset;
  new result1 = _ZX_32(MEM.load(address, 0xF, variant, registerW));
stage CR:
  GPR[registerW] = result1;
\end{lstlisting}}
\newpage

\paragraph{Encoding} Format \textsf{LSU\_LSBI} (Section~\ref{sec:format-LSU-LSBI}), \verb|loadcode = 100|\\

\bitfields{ 1/P, 2/01, 3/100, 2/variant, 6/registerW, 2/10, 3/111, 1/--, 6/registerY, 6/registerZ }


\paragraph{Syntax} \texttt{lwz}\emph{variant} \emph{registerW} = \emph{registerY}[\emph{registerZ}]

\paragraph{Description} The effective address is computed by adding the \emph{registerZ} to the \emph{registerY} and scaled by the \emph{variant}.
The \emph{registerW} is loaded from the memory word at the effective address after zero extension.


\paragraph{Modifier(s)}
\noindent\begin{itemize}
\item variant (cf. variant in section \ref{par:modifier-variant} on page \pageref{par:modifier-variant})
\end{itemize}

\paragraph{Execution} {Reservation table \textsf{LSU\_AUXW} (Section~\ref{sec:reservation-LSU-AUXW})}
~
{\begin{lstlisting}
stage ID:
  new variant = _ZX_2(variant);
stage RR:
  new doscale = 0;
  new base = GPR[registerZ];
  new index = GPR[registerY];
  new scaling = doscale ? 4 : 1;
  new address = base + (index * scaling);
  new result1 = _ZX_32(MEM.load(address, 0xF, variant, registerW));
stage CR:
  GPR[registerW] = result1;
\end{lstlisting}}
\newpage
\subsubsection[{\small SB: Store Byte}]{SB\hfill Store Byte}
\label{instr:SB}\index{sb}
 
\paragraph{Encoding} Format \textsf{LSU\_SSBO} (Section~\ref{sec:format-LSU-SSBO}), \verb|lsucode0 = 10000|\\

\bitfields{ 1/P, 2/01, 5/10000, 6/registerT, 2/01, 10/signed10, 6/registerZ }


\paragraph{Syntax} \texttt{sb} \emph{signed10}[\emph{registerZ}] = \emph{registerT}

\paragraph{Description} The effective address is computed by adding the \emph{signed10} to the \emph{registerZ}.
The \emph{registerT} is stored into the memory byte at the effective address.


\paragraph{Execution} {Reservation table \textsf{LSU\_MEMW\_AUXR} (Section~\ref{sec:reservation-LSU-MEMW-AUXR})}
~
{\begin{lstlisting}
stage ID:
  new offset = _SX_10(signed10);
stage RR:
  new base = GPR[registerZ];
  new address = base + offset;
stage E1:
  new argument3 = GPR[registerT];
stage E3:
  MEM.store(address, 0x1, 0, argument3, registerT);
\end{lstlisting}}
\newpage

\paragraph{Encoding} Format \textsf{LSU\_SSBO.X} (Section~\ref{sec:format-LSU-SSBO.X}), \verb|lsucode0 = 10000|\\

\bitfields{ 1/P, 2/00, 2/-- --, 27/upper27, 1/1, 2/01, 5/10000, 6/registerT, 2/01, 10/lower10, 6/registerZ }


\paragraph{Syntax} \texttt{sb} \emph{upper27\_\discretionary{}{}{}lower10}[\emph{registerZ}] = \emph{registerT}

\paragraph{Description} The effective address is computed by adding the \emph{upper27\_\discretionary{}{}{}lower10} to the \emph{registerZ}.
The \emph{registerT} is stored into the memory byte at the effective address.


\paragraph{Execution} {Reservation table \textsf{LSU\_MEMW\_AUXR.X} (Section~\ref{sec:reservation-LSU-MEMW-AUXR.X})}
~
{\begin{lstlisting}
stage ID:
  new offset = _SX_37(upper27_lower10);
stage RR:
  new base = GPR[registerZ];
  new address = base + offset;
stage E1:
  new argument3 = GPR[registerT];
stage E3:
  MEM.store(address, 0x1, 0, argument3, registerT);
\end{lstlisting}}
\newpage

\paragraph{Encoding} Format \textsf{LSU\_SSBO.Y} (Section~\ref{sec:format-LSU-SSBO.Y}), \verb|lsucode0 = 10000|\\

\bitfields{ 1/P, 2/00, 2/-- --, 27/extend27, 1/1, 2/00, 2/-- --, 27/upper27, 1/1, 2/01, 5/10000, 6/registerT, 2/01, 10/lower10, 6/registerZ }


\paragraph{Syntax} \texttt{sb} \emph{extend27\_\discretionary{}{}{}upper27\_\discretionary{}{}{}lower10}[\emph{registerZ}] = \emph{registerT}

\paragraph{Description} The effective address is computed by adding the \emph{extend27\_\discretionary{}{}{}upper27\_\discretionary{}{}{}lower10} to the \emph{registerZ}.
The \emph{registerT} is stored into the memory byte at the effective address.


\paragraph{Execution} {Reservation table \textsf{LSU\_MEMW\_AUXR.Y} (Section~\ref{sec:reservation-LSU-MEMW-AUXR.Y})}
~
{\begin{lstlisting}
stage ID:
  new offset = _WX_64(extend27_upper27_lower10);
stage RR:
  new base = GPR[registerZ];
  new address = base + offset;
stage E1:
  new argument3 = GPR[registerT];
stage E3:
  MEM.store(address, 0x1, 0, argument3, registerT);
\end{lstlisting}}
\newpage

\paragraph{Encoding} Format \textsf{LSU\_SSBI} (Section~\ref{sec:format-LSU-SSBI}), \verb|lsucode0 = 10000|\\

\bitfields{ 1/P, 2/01, 5/10000, 6/registerT, 2/11, 3/111, 1/--, 6/registerY, 6/registerZ }


\paragraph{Syntax} \texttt{sb} \emph{registerY}[\emph{registerZ}] = \emph{registerT}

\paragraph{Description} The effective address is computed by adding the \emph{registerZ} to the \emph{registerY} and scaled by the \emph{}.
The \emph{registerT} is stored into the memory byte at the effective address.


\paragraph{Execution} {Reservation table \textsf{LSU\_MEMW\_AUXR} (Section~\ref{sec:reservation-LSU-MEMW-AUXR})}
~
{\begin{lstlisting}
stage RR:
  new doscale = 0;
  new index = GPR[registerY];
  new base = GPR[registerZ];
  new scaling = doscale ? 1 : 1;
  new address = base + (index * scaling);
stage E1:
  new argument3 = GPR[registerT];
stage E3:
  MEM.store(address, 0x1, 0, argument3, registerT);
\end{lstlisting}}
\newpage
\subsubsection[{\small SD: Store Double Word}]{SD\hfill Store Double Word}
\label{instr:SD}\index{sd}
 
\paragraph{Encoding} Format \textsf{LSU\_SSBO} (Section~\ref{sec:format-LSU-SSBO}), \verb|lsucode0 = 10011|\\

\bitfields{ 1/P, 2/01, 5/10011, 6/registerT, 2/01, 10/signed10, 6/registerZ }


\paragraph{Syntax} \texttt{sd} \emph{signed10}[\emph{registerZ}] = \emph{registerT}

\paragraph{Description} The effective address is computed by adding the \emph{signed10} to the \emph{registerZ}.
The \emph{registerT} is stored into the memory double word at the effective address.


\paragraph{Execution} {Reservation table \textsf{LSU\_MEMW\_AUXR} (Section~\ref{sec:reservation-LSU-MEMW-AUXR})}
~
{\begin{lstlisting}
stage ID:
  new offset = _SX_10(signed10);
stage RR:
  new base = GPR[registerZ];
  new address = base + offset;
stage E1:
  new argument3 = GPR[registerT];
stage E3:
  MEM.store(address, 0xFF, 0, argument3, registerT);
\end{lstlisting}}
\newpage

\paragraph{Encoding} Format \textsf{LSU\_SSBO.X} (Section~\ref{sec:format-LSU-SSBO.X}), \verb|lsucode0 = 10011|\\

\bitfields{ 1/P, 2/00, 2/-- --, 27/upper27, 1/1, 2/01, 5/10011, 6/registerT, 2/01, 10/lower10, 6/registerZ }


\paragraph{Syntax} \texttt{sd} \emph{upper27\_\discretionary{}{}{}lower10}[\emph{registerZ}] = \emph{registerT}

\paragraph{Description} The effective address is computed by adding the \emph{upper27\_\discretionary{}{}{}lower10} to the \emph{registerZ}.
The \emph{registerT} is stored into the memory double word at the effective address.


\paragraph{Execution} {Reservation table \textsf{LSU\_MEMW\_AUXR.X} (Section~\ref{sec:reservation-LSU-MEMW-AUXR.X})}
~
{\begin{lstlisting}
stage ID:
  new offset = _SX_37(upper27_lower10);
stage RR:
  new base = GPR[registerZ];
  new address = base + offset;
stage E1:
  new argument3 = GPR[registerT];
stage E3:
  MEM.store(address, 0xFF, 0, argument3, registerT);
\end{lstlisting}}
\newpage

\paragraph{Encoding} Format \textsf{LSU\_SSBO.Y} (Section~\ref{sec:format-LSU-SSBO.Y}), \verb|lsucode0 = 10011|\\

\bitfields{ 1/P, 2/00, 2/-- --, 27/extend27, 1/1, 2/00, 2/-- --, 27/upper27, 1/1, 2/01, 5/10011, 6/registerT, 2/01, 10/lower10, 6/registerZ }


\paragraph{Syntax} \texttt{sd} \emph{extend27\_\discretionary{}{}{}upper27\_\discretionary{}{}{}lower10}[\emph{registerZ}] = \emph{registerT}

\paragraph{Description} The effective address is computed by adding the \emph{extend27\_\discretionary{}{}{}upper27\_\discretionary{}{}{}lower10} to the \emph{registerZ}.
The \emph{registerT} is stored into the memory double word at the effective address.


\paragraph{Execution} {Reservation table \textsf{LSU\_MEMW\_AUXR.Y} (Section~\ref{sec:reservation-LSU-MEMW-AUXR.Y})}
~
{\begin{lstlisting}
stage ID:
  new offset = _WX_64(extend27_upper27_lower10);
stage RR:
  new base = GPR[registerZ];
  new address = base + offset;
stage E1:
  new argument3 = GPR[registerT];
stage E3:
  MEM.store(address, 0xFF, 0, argument3, registerT);
\end{lstlisting}}
\newpage

\paragraph{Encoding} Format \textsf{LSU\_SSBI} (Section~\ref{sec:format-LSU-SSBI}), \verb|lsucode0 = 10011|\\

\bitfields{ 1/P, 2/01, 5/10011, 6/registerT, 2/11, 3/111, 1/--, 6/registerY, 6/registerZ }


\paragraph{Syntax} \texttt{sd} \emph{registerY}[\emph{registerZ}] = \emph{registerT}

\paragraph{Description} The effective address is computed by adding the \emph{registerZ} to the \emph{registerY} and scaled by the \emph{}.
The \emph{registerT} is stored into the memory double word at the effective address.


\paragraph{Execution} {Reservation table \textsf{LSU\_MEMW\_AUXR} (Section~\ref{sec:reservation-LSU-MEMW-AUXR})}
~
{\begin{lstlisting}
stage RR:
  new doscale = 0;
  new index = GPR[registerY];
  new base = GPR[registerZ];
  new scaling = doscale ? 8 : 1;
  new address = base + (index * scaling);
stage E1:
  new argument3 = GPR[registerT];
stage E3:
  MEM.store(address, 0xFF, 0, argument3, registerT);
\end{lstlisting}}
\newpage
\subsubsection[{\small SH: Store Half Word}]{SH\hfill Store Half Word}
\label{instr:SH}\index{sh}
 
\paragraph{Encoding} Format \textsf{LSU\_SSBO} (Section~\ref{sec:format-LSU-SSBO}), \verb|lsucode0 = 10001|\\

\bitfields{ 1/P, 2/01, 5/10001, 6/registerT, 2/01, 10/signed10, 6/registerZ }


\paragraph{Syntax} \texttt{sh} \emph{signed10}[\emph{registerZ}] = \emph{registerT}

\paragraph{Description} The effective address is computed by adding the \emph{signed10} to the \emph{registerZ}.
The \emph{registerT} is stored into the memory half word at the effective address.


\paragraph{Execution} {Reservation table \textsf{LSU\_MEMW\_AUXR} (Section~\ref{sec:reservation-LSU-MEMW-AUXR})}
~
{\begin{lstlisting}
stage ID:
  new offset = _SX_10(signed10);
stage RR:
  new base = GPR[registerZ];
  new address = base + offset;
stage E1:
  new argument3 = GPR[registerT];
stage E3:
  MEM.store(address, 0x3, 0, argument3, registerT);
\end{lstlisting}}
\newpage

\paragraph{Encoding} Format \textsf{LSU\_SSBO.X} (Section~\ref{sec:format-LSU-SSBO.X}), \verb|lsucode0 = 10001|\\

\bitfields{ 1/P, 2/00, 2/-- --, 27/upper27, 1/1, 2/01, 5/10001, 6/registerT, 2/01, 10/lower10, 6/registerZ }


\paragraph{Syntax} \texttt{sh} \emph{upper27\_\discretionary{}{}{}lower10}[\emph{registerZ}] = \emph{registerT}

\paragraph{Description} The effective address is computed by adding the \emph{upper27\_\discretionary{}{}{}lower10} to the \emph{registerZ}.
The \emph{registerT} is stored into the memory half word at the effective address.


\paragraph{Execution} {Reservation table \textsf{LSU\_MEMW\_AUXR.X} (Section~\ref{sec:reservation-LSU-MEMW-AUXR.X})}
~
{\begin{lstlisting}
stage ID:
  new offset = _SX_37(upper27_lower10);
stage RR:
  new base = GPR[registerZ];
  new address = base + offset;
stage E1:
  new argument3 = GPR[registerT];
stage E3:
  MEM.store(address, 0x3, 0, argument3, registerT);
\end{lstlisting}}
\newpage

\paragraph{Encoding} Format \textsf{LSU\_SSBO.Y} (Section~\ref{sec:format-LSU-SSBO.Y}), \verb|lsucode0 = 10001|\\

\bitfields{ 1/P, 2/00, 2/-- --, 27/extend27, 1/1, 2/00, 2/-- --, 27/upper27, 1/1, 2/01, 5/10001, 6/registerT, 2/01, 10/lower10, 6/registerZ }


\paragraph{Syntax} \texttt{sh} \emph{extend27\_\discretionary{}{}{}upper27\_\discretionary{}{}{}lower10}[\emph{registerZ}] = \emph{registerT}

\paragraph{Description} The effective address is computed by adding the \emph{extend27\_\discretionary{}{}{}upper27\_\discretionary{}{}{}lower10} to the \emph{registerZ}.
The \emph{registerT} is stored into the memory half word at the effective address.


\paragraph{Execution} {Reservation table \textsf{LSU\_MEMW\_AUXR.Y} (Section~\ref{sec:reservation-LSU-MEMW-AUXR.Y})}
~
{\begin{lstlisting}
stage ID:
  new offset = _WX_64(extend27_upper27_lower10);
stage RR:
  new base = GPR[registerZ];
  new address = base + offset;
stage E1:
  new argument3 = GPR[registerT];
stage E3:
  MEM.store(address, 0x3, 0, argument3, registerT);
\end{lstlisting}}
\newpage

\paragraph{Encoding} Format \textsf{LSU\_SSBI} (Section~\ref{sec:format-LSU-SSBI}), \verb|lsucode0 = 10001|\\

\bitfields{ 1/P, 2/01, 5/10001, 6/registerT, 2/11, 3/111, 1/--, 6/registerY, 6/registerZ }


\paragraph{Syntax} \texttt{sh} \emph{registerY}[\emph{registerZ}] = \emph{registerT}

\paragraph{Description} The effective address is computed by adding the \emph{registerZ} to the \emph{registerY} and scaled by the \emph{}.
The \emph{registerT} is stored into the memory half word at the effective address.


\paragraph{Execution} {Reservation table \textsf{LSU\_MEMW\_AUXR} (Section~\ref{sec:reservation-LSU-MEMW-AUXR})}
~
{\begin{lstlisting}
stage RR:
  new doscale = 0;
  new index = GPR[registerY];
  new base = GPR[registerZ];
  new scaling = doscale ? 2 : 1;
  new address = base + (index * scaling);
stage E1:
  new argument3 = GPR[registerT];
stage E3:
  MEM.store(address, 0x3, 0, argument3, registerT);
\end{lstlisting}}
\newpage
\subsubsection[{\small SO: Store Octuple Word}]{SO\hfill Store Octuple Word}
\label{instr:SO}\index{so}
 
\paragraph{Encoding} Format \textsf{LSU\_SOBO} (Section~\ref{sec:format-LSU-SOBO}), \verb|lsucode0 = 10100|\\

\bitfields{ 1/P, 2/01, 5/10100, 4/registerV, 1/0, 1/1, 2/01, 10/signed10, 6/registerZ }


\paragraph{Syntax} \texttt{so} \emph{signed10}[\emph{registerZ}] = \emph{registerV}

\paragraph{Description} The effective address is computed by adding the \emph{signed10} to the \emph{registerZ}.
The \emph{registerV} is stored into the memory octuple word at the effective address.


\paragraph{Execution} {Reservation table \textsf{LSU\_MEMW\_AUXR} (Section~\ref{sec:reservation-LSU-MEMW-AUXR})}
~
{\begin{lstlisting}
stage ID:
  new offset = _SX_10(signed10);
stage RR:
  new bytemask = 0xFFFFFFFF;
  new base = GPR[registerZ];
  new address = base + offset;
stage E1:
  new argument3 = QGR[registerV];
stage E3:
  MEM.store(address, 0xFFFFFFFF & bytemask, 0, argument3, registerV << 2);
\end{lstlisting}}
\newpage

\paragraph{Encoding} Format \textsf{LSU\_SOBO.X} (Section~\ref{sec:format-LSU-SOBO.X}), \verb|lsucode0 = 10100|\\

\bitfields{ 1/P, 2/00, 2/-- --, 27/upper27, 1/1, 2/01, 5/10100, 4/registerV, 1/0, 1/1, 2/01, 10/lower10, 6/registerZ }


\paragraph{Syntax} \texttt{so} \emph{upper27\_\discretionary{}{}{}lower10}[\emph{registerZ}] = \emph{registerV}

\paragraph{Description} The effective address is computed by adding the \emph{upper27\_\discretionary{}{}{}lower10} to the \emph{registerZ}.
The \emph{registerV} is stored into the memory octuple word at the effective address.


\paragraph{Execution} {Reservation table \textsf{LSU\_MEMW\_AUXR.X} (Section~\ref{sec:reservation-LSU-MEMW-AUXR.X})}
~
{\begin{lstlisting}
stage ID:
  new offset = _SX_37(upper27_lower10);
stage RR:
  new bytemask = 0xFFFFFFFF;
  new base = GPR[registerZ];
  new address = base + offset;
stage E1:
  new argument3 = QGR[registerV];
stage E3:
  MEM.store(address, 0xFFFFFFFF & bytemask, 0, argument3, registerV << 2);
\end{lstlisting}}
\newpage

\paragraph{Encoding} Format \textsf{LSU\_SOBO.Y} (Section~\ref{sec:format-LSU-SOBO.Y}), \verb|lsucode0 = 10100|\\

\bitfields{ 1/P, 2/00, 2/-- --, 27/extend27, 1/1, 2/00, 2/-- --, 27/upper27, 1/1, 2/01, 5/10100, 4/registerV, 1/0, 1/1, 2/01, 10/lower10, 6/registerZ }


\paragraph{Syntax} \texttt{so} \emph{extend27\_\discretionary{}{}{}upper27\_\discretionary{}{}{}lower10}[\emph{registerZ}] = \emph{registerV}

\paragraph{Description} The effective address is computed by adding the \emph{extend27\_\discretionary{}{}{}upper27\_\discretionary{}{}{}lower10} to the \emph{registerZ}.
The \emph{registerV} is stored into the memory octuple word at the effective address.


\paragraph{Execution} {Reservation table \textsf{LSU\_MEMW\_AUXR.Y} (Section~\ref{sec:reservation-LSU-MEMW-AUXR.Y})}
~
{\begin{lstlisting}
stage ID:
  new offset = _WX_64(extend27_upper27_lower10);
stage RR:
  new bytemask = 0xFFFFFFFF;
  new base = GPR[registerZ];
  new address = base + offset;
stage E1:
  new argument3 = QGR[registerV];
stage E3:
  MEM.store(address, 0xFFFFFFFF & bytemask, 0, argument3, registerV << 2);
\end{lstlisting}}
\newpage

\paragraph{Encoding} Format \textsf{LSU\_SOBI} (Section~\ref{sec:format-LSU-SOBI}), \verb|lsucode0 = 10100|\\

\bitfields{ 1/P, 2/01, 5/10100, 4/registerV, 1/0, 1/1, 2/11, 3/111, 1/--, 6/registerY, 6/registerZ }


\paragraph{Syntax} \texttt{so} \emph{registerY}[\emph{registerZ}] = \emph{registerV}

\paragraph{Description} The effective address is computed by adding the \emph{registerZ} to the \emph{registerY} and scaled by the \emph{}.
The \emph{registerV} is stored into the memory octuple word at the effective address.


\paragraph{Execution} {Reservation table \textsf{LSU\_MEMW\_AUXR} (Section~\ref{sec:reservation-LSU-MEMW-AUXR})}
~
{\begin{lstlisting}
stage RR:
  new bytemask = 0xFFFFFFFF;
  new doscale = 0;
  new index = GPR[registerY];
  new base = GPR[registerZ];
  new scaling = doscale ? 32 : 1;
  new address = base + (index * scaling);
stage E1:
  new argument3 = QGR[registerV];
stage E3:
  MEM.store(address, 0xFFFFFFFF & bytemask, 0, argument3, registerV << 2);
\end{lstlisting}}
\newpage
\subsubsection[{\small SQ: Store Quadruple Word}]{SQ\hfill Store Quadruple Word}
\label{instr:SQ}\index{sq}
 
\paragraph{Encoding} Format \textsf{LSU\_SQBO} (Section~\ref{sec:format-LSU-SQBO}), \verb|lsucode0 = 10100|\\

\bitfields{ 1/P, 2/01, 5/10100, 5/registerU, 1/0, 2/01, 10/signed10, 6/registerZ }


\paragraph{Syntax} \texttt{sq} \emph{signed10}[\emph{registerZ}] = \emph{registerU}

\paragraph{Description} The effective address is computed by adding the \emph{signed10} to the \emph{registerZ}.
The \emph{registerU} is stored into the memory quadruple word at the effective address.


\paragraph{Execution} {Reservation table \textsf{LSU\_MEMW\_AUXR} (Section~\ref{sec:reservation-LSU-MEMW-AUXR})}
~
{\begin{lstlisting}
stage ID:
  new offset = _SX_10(signed10);
stage RR:
  new base = GPR[registerZ];
  new address = base + offset;
stage E1:
  new argument3 = PGR[registerU];
stage E3:
  MEM.store(address, 0xFFFF, 0, argument3, registerU << 1);
\end{lstlisting}}
\newpage

\paragraph{Encoding} Format \textsf{LSU\_SQBO.X} (Section~\ref{sec:format-LSU-SQBO.X}), \verb|lsucode0 = 10100|\\

\bitfields{ 1/P, 2/00, 2/-- --, 27/upper27, 1/1, 2/01, 5/10100, 5/registerU, 1/0, 2/01, 10/lower10, 6/registerZ }


\paragraph{Syntax} \texttt{sq} \emph{upper27\_\discretionary{}{}{}lower10}[\emph{registerZ}] = \emph{registerU}

\paragraph{Description} The effective address is computed by adding the \emph{upper27\_\discretionary{}{}{}lower10} to the \emph{registerZ}.
The \emph{registerU} is stored into the memory quadruple word at the effective address.


\paragraph{Execution} {Reservation table \textsf{LSU\_MEMW\_AUXR.X} (Section~\ref{sec:reservation-LSU-MEMW-AUXR.X})}
~
{\begin{lstlisting}
stage ID:
  new offset = _SX_37(upper27_lower10);
stage RR:
  new base = GPR[registerZ];
  new address = base + offset;
stage E1:
  new argument3 = PGR[registerU];
stage E3:
  MEM.store(address, 0xFFFF, 0, argument3, registerU << 1);
\end{lstlisting}}
\newpage

\paragraph{Encoding} Format \textsf{LSU\_SQBO.Y} (Section~\ref{sec:format-LSU-SQBO.Y}), \verb|lsucode0 = 10100|\\

\bitfields{ 1/P, 2/00, 2/-- --, 27/extend27, 1/1, 2/00, 2/-- --, 27/upper27, 1/1, 2/01, 5/10100, 5/registerU, 1/0, 2/01, 10/lower10, 6/registerZ }


\paragraph{Syntax} \texttt{sq} \emph{extend27\_\discretionary{}{}{}upper27\_\discretionary{}{}{}lower10}[\emph{registerZ}] = \emph{registerU}

\paragraph{Description} The effective address is computed by adding the \emph{extend27\_\discretionary{}{}{}upper27\_\discretionary{}{}{}lower10} to the \emph{registerZ}.
The \emph{registerU} is stored into the memory quadruple word at the effective address.


\paragraph{Execution} {Reservation table \textsf{LSU\_MEMW\_AUXR.Y} (Section~\ref{sec:reservation-LSU-MEMW-AUXR.Y})}
~
{\begin{lstlisting}
stage ID:
  new offset = _WX_64(extend27_upper27_lower10);
stage RR:
  new base = GPR[registerZ];
  new address = base + offset;
stage E1:
  new argument3 = PGR[registerU];
stage E3:
  MEM.store(address, 0xFFFF, 0, argument3, registerU << 1);
\end{lstlisting}}
\newpage

\paragraph{Encoding} Format \textsf{LSU\_SQBI} (Section~\ref{sec:format-LSU-SQBI}), \verb|lsucode0 = 10100|\\

\bitfields{ 1/P, 2/01, 5/10100, 5/registerU, 1/0, 2/11, 3/111, 1/--, 6/registerY, 6/registerZ }


\paragraph{Syntax} \texttt{sq} \emph{registerY}[\emph{registerZ}] = \emph{registerU}

\paragraph{Description} The effective address is computed by adding the \emph{registerZ} to the \emph{registerY} and scaled by the \emph{}.
The \emph{registerU} is stored into the memory quadruple word at the effective address.


\paragraph{Execution} {Reservation table \textsf{LSU\_MEMW\_AUXR} (Section~\ref{sec:reservation-LSU-MEMW-AUXR})}
~
{\begin{lstlisting}
stage RR:
  new doscale = 0;
  new index = GPR[registerY];
  new base = GPR[registerZ];
  new scaling = doscale ? 16 : 1;
  new address = base + (index * scaling);
stage E1:
  new argument3 = PGR[registerU];
stage E3:
  MEM.store(address, 0xFFFF, 0, argument3, registerU << 1);
\end{lstlisting}}
\newpage
\subsubsection[{\small SW: Store Word}]{SW\hfill Store Word}
\label{instr:SW}\index{sw}
 
\paragraph{Encoding} Format \textsf{LSU\_SSBO} (Section~\ref{sec:format-LSU-SSBO}), \verb|lsucode0 = 10010|\\

\bitfields{ 1/P, 2/01, 5/10010, 6/registerT, 2/01, 10/signed10, 6/registerZ }


\paragraph{Syntax} \texttt{sw} \emph{signed10}[\emph{registerZ}] = \emph{registerT}

\paragraph{Description} The effective address is computed by adding the \emph{signed10} to the \emph{registerZ}.
The \emph{registerT} is stored into the memory word at the effective address.


\paragraph{Execution} {Reservation table \textsf{LSU\_MEMW\_AUXR} (Section~\ref{sec:reservation-LSU-MEMW-AUXR})}
~
{\begin{lstlisting}
stage ID:
  new offset = _SX_10(signed10);
stage RR:
  new base = GPR[registerZ];
  new address = base + offset;
stage E1:
  new argument3 = GPR[registerT];
stage E3:
  MEM.store(address, 0xF, 0, argument3, registerT);
\end{lstlisting}}
\newpage

\paragraph{Encoding} Format \textsf{LSU\_SSBO.X} (Section~\ref{sec:format-LSU-SSBO.X}), \verb|lsucode0 = 10010|\\

\bitfields{ 1/P, 2/00, 2/-- --, 27/upper27, 1/1, 2/01, 5/10010, 6/registerT, 2/01, 10/lower10, 6/registerZ }


\paragraph{Syntax} \texttt{sw} \emph{upper27\_\discretionary{}{}{}lower10}[\emph{registerZ}] = \emph{registerT}

\paragraph{Description} The effective address is computed by adding the \emph{upper27\_\discretionary{}{}{}lower10} to the \emph{registerZ}.
The \emph{registerT} is stored into the memory word at the effective address.


\paragraph{Execution} {Reservation table \textsf{LSU\_MEMW\_AUXR.X} (Section~\ref{sec:reservation-LSU-MEMW-AUXR.X})}
~
{\begin{lstlisting}
stage ID:
  new offset = _SX_37(upper27_lower10);
stage RR:
  new base = GPR[registerZ];
  new address = base + offset;
stage E1:
  new argument3 = GPR[registerT];
stage E3:
  MEM.store(address, 0xF, 0, argument3, registerT);
\end{lstlisting}}
\newpage

\paragraph{Encoding} Format \textsf{LSU\_SSBO.Y} (Section~\ref{sec:format-LSU-SSBO.Y}), \verb|lsucode0 = 10010|\\

\bitfields{ 1/P, 2/00, 2/-- --, 27/extend27, 1/1, 2/00, 2/-- --, 27/upper27, 1/1, 2/01, 5/10010, 6/registerT, 2/01, 10/lower10, 6/registerZ }


\paragraph{Syntax} \texttt{sw} \emph{extend27\_\discretionary{}{}{}upper27\_\discretionary{}{}{}lower10}[\emph{registerZ}] = \emph{registerT}

\paragraph{Description} The effective address is computed by adding the \emph{extend27\_\discretionary{}{}{}upper27\_\discretionary{}{}{}lower10} to the \emph{registerZ}.
The \emph{registerT} is stored into the memory word at the effective address.


\paragraph{Execution} {Reservation table \textsf{LSU\_MEMW\_AUXR.Y} (Section~\ref{sec:reservation-LSU-MEMW-AUXR.Y})}
~
{\begin{lstlisting}
stage ID:
  new offset = _WX_64(extend27_upper27_lower10);
stage RR:
  new base = GPR[registerZ];
  new address = base + offset;
stage E1:
  new argument3 = GPR[registerT];
stage E3:
  MEM.store(address, 0xF, 0, argument3, registerT);
\end{lstlisting}}
\newpage

\paragraph{Encoding} Format \textsf{LSU\_SSBI} (Section~\ref{sec:format-LSU-SSBI}), \verb|lsucode0 = 10010|\\

\bitfields{ 1/P, 2/01, 5/10010, 6/registerT, 2/11, 3/111, 1/--, 6/registerY, 6/registerZ }


\paragraph{Syntax} \texttt{sw} \emph{registerY}[\emph{registerZ}] = \emph{registerT}

\paragraph{Description} The effective address is computed by adding the \emph{registerZ} to the \emph{registerY} and scaled by the \emph{}.
The \emph{registerT} is stored into the memory word at the effective address.


\paragraph{Execution} {Reservation table \textsf{LSU\_MEMW\_AUXR} (Section~\ref{sec:reservation-LSU-MEMW-AUXR})}
~
{\begin{lstlisting}
stage RR:
  new doscale = 0;
  new index = GPR[registerY];
  new base = GPR[registerZ];
  new scaling = doscale ? 4 : 1;
  new address = base + (index * scaling);
stage E1:
  new argument3 = GPR[registerT];
stage E3:
  MEM.store(address, 0xF, 0, argument3, registerT);
\end{lstlisting}}
\newpage
\subsubsection[{\small XLO: Extension Load Octuple Word}]{XLO\hfill Extension Load Octuple Word}
\label{instr:XLO}\index{xlo}
 
\paragraph{Encoding} Format \textsf{LSU\_XLOBO} (Section~\ref{sec:format-LSU-XLOBO}), \verb|steering = 01|\\

\bitfields{ 1/P, 2/01, 3/000, 2/variant, 6/registerG, 2/01, 10/signed10, 6/registerZ }


\paragraph{Syntax} \texttt{xlo}\emph{variant} \emph{registerG} = \emph{signed10}[\emph{registerZ}]

\paragraph{Description} The effective address is computed by adding the \emph{signed10} to the \emph{registerZ}.
The \emph{registerG} is loaded from the memory octuple word at the effective address.


\paragraph{Modifier(s)}
\noindent\begin{itemize}
\item variant (cf. variant in section \ref{par:modifier-variant} on page \pageref{par:modifier-variant})
\end{itemize}

\paragraph{Execution} {Reservation table \textsf{LSU} (Section~\ref{sec:reservation-LSU})}
~
{\begin{lstlisting}
stage ID:
  new variant = _ZX_2(variant);
  new offset = _SX_10(signed10);
stage RR:
  new bytemask = 0xFFFFFFFF;
  new dri = registerG;
  new base = GPR[registerZ];
  new address = base + offset;
stage E1:
  CS.XMF = 1;
  new result1 = MEM.load(address, 0xFFFFFFFF & bytemask, variant, dri);
stage CR:
  XVR[registerG] = result1;
\end{lstlisting}}
\newpage

\paragraph{Encoding} Format \textsf{LSU\_XLOBO.X} (Section~\ref{sec:format-LSU-XLOBO.X}), \verb|steering = 01|\\

\bitfields{ 1/P, 2/00, 2/-- --, 27/upper27, 1/1, 2/01, 3/000, 2/variant, 6/registerG, 2/01, 10/lower10, 6/registerZ }


\paragraph{Syntax} \texttt{xlo}\emph{variant} \emph{registerG} = \emph{upper27\_\discretionary{}{}{}lower10}[\emph{registerZ}]

\paragraph{Description} The effective address is computed by adding the \emph{upper27\_\discretionary{}{}{}lower10} to the \emph{registerZ}.
The \emph{registerG} is loaded from the memory octuple word at the effective address.


\paragraph{Modifier(s)}
\noindent\begin{itemize}
\item variant (cf. variant in section \ref{par:modifier-variant} on page \pageref{par:modifier-variant})
\end{itemize}

\paragraph{Execution} {Reservation table \textsf{LSU.X} (Section~\ref{sec:reservation-LSU.X})}
~
{\begin{lstlisting}
stage ID:
  new variant = _ZX_2(variant);
  new offset = _SX_37(upper27_lower10);
stage RR:
  new bytemask = 0xFFFFFFFF;
  new dri = registerG;
  new base = GPR[registerZ];
  new address = base + offset;
stage E1:
  CS.XMF = 1;
  new result1 = MEM.load(address, 0xFFFFFFFF & bytemask, variant, dri);
stage CR:
  XVR[registerG] = result1;
\end{lstlisting}}
\newpage

\paragraph{Encoding} Format \textsf{LSU\_XLOBO.Y} (Section~\ref{sec:format-LSU-XLOBO.Y}), \verb|steering = 01|\\

\bitfields{ 1/P, 2/00, 2/-- --, 27/extend27, 1/1, 2/00, 2/-- --, 27/upper27, 1/1, 2/01, 3/000, 2/variant, 6/registerG, 2/01, 10/lower10, 6/registerZ }


\paragraph{Syntax} \texttt{xlo}\emph{variant} \emph{registerG} = \emph{extend27\_\discretionary{}{}{}upper27\_\discretionary{}{}{}lower10}[\emph{registerZ}]

\paragraph{Description} The effective address is computed by adding the \emph{extend27\_\discretionary{}{}{}upper27\_\discretionary{}{}{}lower10} to the \emph{registerZ}.
The \emph{registerG} is loaded from the memory octuple word at the effective address.


\paragraph{Modifier(s)}
\noindent\begin{itemize}
\item variant (cf. variant in section \ref{par:modifier-variant} on page \pageref{par:modifier-variant})
\end{itemize}

\paragraph{Execution} {Reservation table \textsf{LSU.Y} (Section~\ref{sec:reservation-LSU.Y})}
~
{\begin{lstlisting}
stage ID:
  new variant = _ZX_2(variant);
  new offset = _WX_64(extend27_upper27_lower10);
stage RR:
  new bytemask = 0xFFFFFFFF;
  new dri = registerG;
  new base = GPR[registerZ];
  new address = base + offset;
stage E1:
  CS.XMF = 1;
  new result1 = MEM.load(address, 0xFFFFFFFF & bytemask, variant, dri);
stage CR:
  XVR[registerG] = result1;
\end{lstlisting}}
\newpage

\paragraph{Encoding} Format \textsf{LSU\_XLOC4BO} (Section~\ref{sec:format-LSU-XLOC4BO}), \verb|steering = 01|\\

\bitfields{ 1/P, 2/01, 3/010, 2/variant, 4/registerGq, 2/qindex, 2/01, 10/signed10, 6/registerZ }


\paragraph{Syntax} \texttt{xlo}\emph{variant}\emph{qindex} \emph{registerGq} = \emph{signed10}[\emph{registerZ}]

\paragraph{Description} The effective address is computed by adding the \emph{signed10} to the \emph{registerZ}.
The \emph{registerGq} is loaded from the memory octuple word at the effective address.


\paragraph{Modifier(s)}
\noindent\begin{itemize}
\item qindex (cf. qindex in section \ref{par:modifier-qindex} on page \pageref{par:modifier-qindex})
\item variant (cf. variant in section \ref{par:modifier-variant} on page \pageref{par:modifier-variant})
\end{itemize}

\paragraph{Execution} {Reservation table \textsf{LSU} (Section~\ref{sec:reservation-LSU})}
~
{\begin{lstlisting}
stage ID:
  new variant = _ZX_2(variant);
  new qindex = _ZX_2(qindex);
  new offset = _SX_10(signed10);
stage RR:
  new bytemask = 0xFFFFFFFF;
  new dri = registerGq << 2;
  new base = GPR[registerZ];
  new address = base + offset;
stage E1:
  new argument1_0 = XMR[registerGq]:0;
  new argument1_1 = XMR[registerGq]:1;
  new argument1_2 = XMR[registerGq]:2;
  new argument1_3 = XMR[registerGq]:3;
  CS.XMF = 1;
  new result1 = MEM.load(address, 0xFFFFFFFF & bytemask, variant, dri);
  new result1_0 = _ZX_64(result1);
  new result1_1 = result1 >> 64;
  new result1_2 = result1 >> 128;
  new result1_3 = result1 >> 192;
stage E3:
  XMR[registerGq]:0 = insert_64(argument1_0, result1_0, qindex);
  XMR[registerGq]:1 = insert_64(argument1_1, result1_1, qindex);
  XMR[registerGq]:2 = insert_64(argument1_2, result1_2, qindex);
  XMR[registerGq]:3 = insert_64(argument1_3, result1_3, qindex);
\end{lstlisting}}
\newpage

\paragraph{Encoding} Format \textsf{LSU\_XLOC4BO.X} (Section~\ref{sec:format-LSU-XLOC4BO.X}), \verb|steering = 01|\\

\bitfields{ 1/P, 2/00, 2/-- --, 27/upper27, 1/1, 2/01, 3/010, 2/variant, 4/registerGq, 2/qindex, 2/01, 10/lower10, 6/registerZ }


\paragraph{Syntax} \texttt{xlo}\emph{variant}\emph{qindex} \emph{registerGq} = \emph{upper27\_\discretionary{}{}{}lower10}[\emph{registerZ}]

\paragraph{Description} The effective address is computed by adding the \emph{upper27\_\discretionary{}{}{}lower10} to the \emph{registerZ}.
The \emph{registerGq} is loaded from the memory octuple word at the effective address.


\paragraph{Modifier(s)}
\noindent\begin{itemize}
\item qindex (cf. qindex in section \ref{par:modifier-qindex} on page \pageref{par:modifier-qindex})
\item variant (cf. variant in section \ref{par:modifier-variant} on page \pageref{par:modifier-variant})
\end{itemize}

\paragraph{Execution} {Reservation table \textsf{LSU.X} (Section~\ref{sec:reservation-LSU.X})}
~
{\begin{lstlisting}
stage ID:
  new variant = _ZX_2(variant);
  new qindex = _ZX_2(qindex);
  new offset = _SX_37(upper27_lower10);
stage RR:
  new bytemask = 0xFFFFFFFF;
  new dri = registerGq << 2;
  new base = GPR[registerZ];
  new address = base + offset;
stage E1:
  new argument1_0 = XMR[registerGq]:0;
  new argument1_1 = XMR[registerGq]:1;
  new argument1_2 = XMR[registerGq]:2;
  new argument1_3 = XMR[registerGq]:3;
  CS.XMF = 1;
  new result1 = MEM.load(address, 0xFFFFFFFF & bytemask, variant, dri);
  new result1_0 = _ZX_64(result1);
  new result1_1 = result1 >> 64;
  new result1_2 = result1 >> 128;
  new result1_3 = result1 >> 192;
stage E3:
  XMR[registerGq]:0 = insert_64(argument1_0, result1_0, qindex);
  XMR[registerGq]:1 = insert_64(argument1_1, result1_1, qindex);
  XMR[registerGq]:2 = insert_64(argument1_2, result1_2, qindex);
  XMR[registerGq]:3 = insert_64(argument1_3, result1_3, qindex);
\end{lstlisting}}
\newpage

\paragraph{Encoding} Format \textsf{LSU\_XLOC4BO.Y} (Section~\ref{sec:format-LSU-XLOC4BO.Y}), \verb|steering = 01|\\

\bitfields{ 1/P, 2/00, 2/-- --, 27/extend27, 1/1, 2/00, 2/-- --, 27/upper27, 1/1, 2/01, 3/010, 2/variant, 4/registerGq, 2/qindex, 2/01, 10/lower10, 6/registerZ }


\paragraph{Syntax} \texttt{xlo}\emph{variant}\emph{qindex} \emph{registerGq} = \emph{extend27\_\discretionary{}{}{}upper27\_\discretionary{}{}{}lower10}[\emph{registerZ}]

\paragraph{Description} The effective address is computed by adding the \emph{extend27\_\discretionary{}{}{}upper27\_\discretionary{}{}{}lower10} to the \emph{registerZ}.
The \emph{registerGq} is loaded from the memory octuple word at the effective address.


\paragraph{Modifier(s)}
\noindent\begin{itemize}
\item qindex (cf. qindex in section \ref{par:modifier-qindex} on page \pageref{par:modifier-qindex})
\item variant (cf. variant in section \ref{par:modifier-variant} on page \pageref{par:modifier-variant})
\end{itemize}

\paragraph{Execution} {Reservation table \textsf{LSU.Y} (Section~\ref{sec:reservation-LSU.Y})}
~
{\begin{lstlisting}
stage ID:
  new variant = _ZX_2(variant);
  new qindex = _ZX_2(qindex);
  new offset = _WX_64(extend27_upper27_lower10);
stage RR:
  new bytemask = 0xFFFFFFFF;
  new dri = registerGq << 2;
  new base = GPR[registerZ];
  new address = base + offset;
stage E1:
  new argument1_0 = XMR[registerGq]:0;
  new argument1_1 = XMR[registerGq]:1;
  new argument1_2 = XMR[registerGq]:2;
  new argument1_3 = XMR[registerGq]:3;
  CS.XMF = 1;
  new result1 = MEM.load(address, 0xFFFFFFFF & bytemask, variant, dri);
  new result1_0 = _ZX_64(result1);
  new result1_1 = result1 >> 64;
  new result1_2 = result1 >> 128;
  new result1_3 = result1 >> 192;
stage E3:
  XMR[registerGq]:0 = insert_64(argument1_0, result1_0, qindex);
  XMR[registerGq]:1 = insert_64(argument1_1, result1_1, qindex);
  XMR[registerGq]:2 = insert_64(argument1_2, result1_2, qindex);
  XMR[registerGq]:3 = insert_64(argument1_3, result1_3, qindex);
\end{lstlisting}}
\newpage

\paragraph{Encoding} Format \textsf{LSU\_XLBI} (Section~\ref{sec:format-LSU-XLBI}), \verb|steering = 01|\\

\bitfields{ 1/P, 2/01, 3/000, 2/variant, 6/registerG, 2/11, 3/111, 1/--, 6/registerY, 6/registerZ }


\paragraph{Syntax} \texttt{xlo}\emph{variant} \emph{registerG} = \emph{registerY}[\emph{registerZ}]

\paragraph{Description} The effective address is computed by adding the \emph{registerZ} to the \emph{registerY} and scaled by the \emph{variant}.
The \emph{registerG} is loaded from the memory octuple word at the effective address.


\paragraph{Modifier(s)}
\noindent\begin{itemize}
\item variant (cf. variant in section \ref{par:modifier-variant} on page \pageref{par:modifier-variant})
\end{itemize}

\paragraph{Execution} {Reservation table \textsf{LSU} (Section~\ref{sec:reservation-LSU})}
~
{\begin{lstlisting}
stage ID:
  new variant = _ZX_2(variant);
stage RR:
  new bytemask = 0xFFFFFFFF;
  new doscale = 0;
  new dri = registerG;
  new base = GPR[registerZ];
  new index = GPR[registerY];
  new scaling = doscale ? 32 : 1;
  new address = base + (index * scaling);
stage E1:
  CS.XMF = 1;
  new result1 = MEM.load(address, 0xFFFFFFFF & bytemask, variant, dri);
stage CR:
  XVR[registerG] = result1;
\end{lstlisting}}
\newpage

\paragraph{Encoding} Format \textsf{LSU\_XLOC4BI} (Section~\ref{sec:format-LSU-XLOC4BI}), \verb|steering = 01|\\

\bitfields{ 1/P, 2/01, 3/010, 2/variant, 4/registerGq, 2/qindex, 2/11, 3/111, 1/--, 6/registerY, 6/registerZ }


\paragraph{Syntax} \texttt{xlo}\emph{variant}\emph{qindex} \emph{registerGq} = \emph{registerY}[\emph{registerZ}]

\paragraph{Description} The effective address is computed by adding the \emph{registerZ} to the \emph{registerY} and scaled by the \emph{qindex}.
The \emph{registerGq} is loaded from the memory octuple word at the effective address.


\paragraph{Modifier(s)}
\noindent\begin{itemize}
\item qindex (cf. qindex in section \ref{par:modifier-qindex} on page \pageref{par:modifier-qindex})
\item variant (cf. variant in section \ref{par:modifier-variant} on page \pageref{par:modifier-variant})
\end{itemize}

\paragraph{Execution} {Reservation table \textsf{LSU} (Section~\ref{sec:reservation-LSU})}
~
{\begin{lstlisting}
stage ID:
  new variant = _ZX_2(variant);
  new qindex = _ZX_2(qindex);
stage RR:
  new bytemask = 0xFFFFFFFF;
  new doscale = 0;
  new dri = registerGq << 2;
  new base = GPR[registerZ];
  new index = GPR[registerY];
  new scaling = doscale ? 32 : 1;
  new address = base + (index * scaling);
stage E1:
  new argument1_0 = XMR[registerGq]:0;
  new argument1_1 = XMR[registerGq]:1;
  new argument1_2 = XMR[registerGq]:2;
  new argument1_3 = XMR[registerGq]:3;
  CS.XMF = 1;
  new result1 = MEM.load(address, 0xFFFFFFFF & bytemask, variant, dri);
  new result1_0 = _ZX_64(result1);
  new result1_1 = result1 >> 64;
  new result1_2 = result1 >> 128;
  new result1_3 = result1 >> 192;
stage E3:
  XMR[registerGq]:0 = insert_64(argument1_0, result1_0, qindex);
  XMR[registerGq]:1 = insert_64(argument1_1, result1_1, qindex);
  XMR[registerGq]:2 = insert_64(argument1_2, result1_2, qindex);
  XMR[registerGq]:3 = insert_64(argument1_3, result1_3, qindex);
\end{lstlisting}}
\newpage
\subsubsection[{\small XPLB: Extension Preload Byte}]{XPLB\hfill Extension Preload Byte}
\label{instr:XPLB}\index{xplb}
 
\paragraph{Encoding} Format \textsf{LSU\_XPL2RBB} (Section~\ref{sec:format-LSU-XPL2RBB}), \verb|lsucode5 = 000|\\

\bitfields{ 1/P, 2/01, 3/011, 2/variant, 5/registerGg, 1/0, 2/11, 3/000, 1/--, 6/registerY, 6/registerZ }


\paragraph{Syntax} \texttt{xplb}\emph{variant} \emph{registerGg}, \emph{registerY} = [\emph{registerZ}]

\paragraph{Description} The \emph{registerGg} is preloaded from the memory byte at the effective address.


\paragraph{Modifier(s)}
\noindent\begin{itemize}
\item variant (cf. variant in section \ref{par:modifier-variant} on page \pageref{par:modifier-variant})
\end{itemize}

\paragraph{Execution} {Reservation table \textsf{LSU} (Section~\ref{sec:reservation-LSU})}
~
{\begin{lstlisting}
stage ID:
  new variant = _ZX_2(variant);
stage RR:
  new address = GPR[registerZ];
  new target = GPR[registerY];
  new shift = (target & 31) * 8;
  new index = (target >> 5) & 1;
  new where = (registerGg << 1) + index;
  new dri = where;
stage E1:
  new argument1 = XVR[where];
  new bytemask = 1;
  new targetbytes = bits2bytes(bytemask) << shift;
  CS.XMF = 1;
  new result1 = MEM.load(address, 0x1 & bytemask, variant, dri);
  result1 = ((result1 << shift) & targetbytes) | (argument1 & ~targetbytes);
stage CR:
  XVR[where] = result1;
\end{lstlisting}}
\newpage

\paragraph{Encoding} Format \textsf{LSU\_XPL2RBB.O} (Section~\ref{sec:format-LSU-XPL2RBB.O}), \verb|lsucode5 = 000|\\

\bitfields{ 1/P, 2/00, 2/-- --, 27/offset27, 1/1, 2/01, 3/011, 2/variant, 5/registerGg, 1/0, 2/11, 3/000, 1/--, 6/registerY, 6/registerZ }


\paragraph{Syntax} \texttt{xplb}\emph{variant} \emph{registerGg}, \emph{registerY} = \emph{offset27}[\emph{registerZ}]

\paragraph{Description} The \emph{registerGg} is preloaded from the memory byte at the effective address.


\paragraph{Modifier(s)}
\noindent\begin{itemize}
\item variant (cf. variant in section \ref{par:modifier-variant} on page \pageref{par:modifier-variant})
\end{itemize}

\paragraph{Execution} {Reservation table \textsf{LSU.X} (Section~\ref{sec:reservation-LSU.X})}
~
{\begin{lstlisting}
stage ID:
  new offset = _SX_27(offset27);
  new variant = _ZX_2(variant);
stage RR:
  new base = GPR[registerZ];
  new target = GPR[registerY];
  new address = base + offset;
  new shift = (target & 31) * 8;
  new index = (target >> 5) & 1;
  new where = (registerGg << 1) + index;
  new dri = where;
stage E1:
  new argument1 = XVR[where];
  new bytemask = 1;
  new targetbytes = bits2bytes(bytemask) << shift;
  CS.XMF = 1;
  new result1 = MEM.load(address, 0x1 & bytemask, variant, dri);
  result1 = ((result1 << shift) & targetbytes) | (argument1 & ~targetbytes);
stage CR:
  XVR[where] = result1;
\end{lstlisting}}
\newpage

\paragraph{Encoding} Format \textsf{LSU\_XPL2RBB.D} (Section~\ref{sec:format-LSU-XPL2RBB.D}), \verb|lsucode5 = 000|\\

\bitfields{ 1/P, 2/00, 2/-- --, 27/extend27, 1/1, 2/00, 2/-- --, 27/offset27, 1/1, 2/01, 3/011, 2/variant, 5/registerGg, 1/0, 2/11, 3/000, 1/--, 6/registerY, 6/registerZ }


\paragraph{Syntax} \texttt{xplb}\emph{variant} \emph{registerGg}, \emph{registerY} = \emph{extend27\_\discretionary{}{}{}offset27}[\emph{registerZ}]

\paragraph{Description} The \emph{registerGg} is preloaded from the memory byte at the effective address.


\paragraph{Modifier(s)}
\noindent\begin{itemize}
\item variant (cf. variant in section \ref{par:modifier-variant} on page \pageref{par:modifier-variant})
\end{itemize}

\paragraph{Execution} {Reservation table \textsf{LSU.Y} (Section~\ref{sec:reservation-LSU.Y})}
~
{\begin{lstlisting}
stage ID:
  new offset = _SX_54(extend27_offset27);
  new variant = _ZX_2(variant);
stage RR:
  new base = GPR[registerZ];
  new target = GPR[registerY];
  new address = base + offset;
  new shift = (target & 31) * 8;
  new index = (target >> 5) & 1;
  new where = (registerGg << 1) + index;
  new dri = where;
stage E1:
  new argument1 = XVR[where];
  new bytemask = 1;
  new targetbytes = bits2bytes(bytemask) << shift;
  CS.XMF = 1;
  new result1 = MEM.load(address, 0x1 & bytemask, variant, dri);
  result1 = ((result1 << shift) & targetbytes) | (argument1 & ~targetbytes);
stage CR:
  XVR[where] = result1;
\end{lstlisting}}
\newpage

\paragraph{Encoding} Format \textsf{LSU\_XPL4RBB} (Section~\ref{sec:format-LSU-XPL4RBB}), \verb|lsucode5 = 000|\\

\bitfields{ 1/P, 2/01, 3/011, 2/variant, 4/registerGh, 2/01, 2/11, 3/000, 1/--, 6/registerY, 6/registerZ }


\paragraph{Syntax} \texttt{xplb}\emph{variant} \emph{registerGh}, \emph{registerY} = [\emph{registerZ}]

\paragraph{Description} The \emph{registerGh} is preloaded from the memory byte at the effective address.


\paragraph{Modifier(s)}
\noindent\begin{itemize}
\item variant (cf. variant in section \ref{par:modifier-variant} on page \pageref{par:modifier-variant})
\end{itemize}

\paragraph{Execution} {Reservation table \textsf{LSU} (Section~\ref{sec:reservation-LSU})}
~
{\begin{lstlisting}
stage ID:
  new variant = _ZX_2(variant);
stage RR:
  new address = GPR[registerZ];
  new target = GPR[registerY];
  new shift = (target & 31) * 8;
  new index = (target >> 5) & 3;
  new where = (registerGh << 2) + index;
  new dri = where;
stage E1:
  new argument1 = XVR[where];
  new bytemask = 1;
  new targetbytes = bits2bytes(bytemask) << shift;
  CS.XMF = 1;
  new result1 = MEM.load(address, 0x1 & bytemask, variant, dri);
  result1 = ((result1 << shift) & targetbytes) | (argument1 & ~targetbytes);
stage CR:
  XVR[where] = result1;
\end{lstlisting}}
\newpage

\paragraph{Encoding} Format \textsf{LSU\_XPL4RBB.O} (Section~\ref{sec:format-LSU-XPL4RBB.O}), \verb|lsucode5 = 000|\\

\bitfields{ 1/P, 2/00, 2/-- --, 27/offset27, 1/1, 2/01, 3/011, 2/variant, 4/registerGh, 2/01, 2/11, 3/000, 1/--, 6/registerY, 6/registerZ }


\paragraph{Syntax} \texttt{xplb}\emph{variant} \emph{registerGh}, \emph{registerY} = \emph{offset27}[\emph{registerZ}]

\paragraph{Description} The \emph{registerGh} is preloaded from the memory byte at the effective address.


\paragraph{Modifier(s)}
\noindent\begin{itemize}
\item variant (cf. variant in section \ref{par:modifier-variant} on page \pageref{par:modifier-variant})
\end{itemize}

\paragraph{Execution} {Reservation table \textsf{LSU.X} (Section~\ref{sec:reservation-LSU.X})}
~
{\begin{lstlisting}
stage ID:
  new offset = _SX_27(offset27);
  new variant = _ZX_2(variant);
stage RR:
  new base = GPR[registerZ];
  new target = GPR[registerY];
  new address = base + offset;
  new shift = (target & 31) * 8;
  new index = (target >> 5) & 3;
  new where = (registerGh << 2) + index;
  new dri = where;
stage E1:
  new argument1 = XVR[where];
  new bytemask = 1;
  new targetbytes = bits2bytes(bytemask) << shift;
  CS.XMF = 1;
  new result1 = MEM.load(address, 0x1 & bytemask, variant, dri);
  result1 = ((result1 << shift) & targetbytes) | (argument1 & ~targetbytes);
stage CR:
  XVR[where] = result1;
\end{lstlisting}}
\newpage

\paragraph{Encoding} Format \textsf{LSU\_XPL4RBB.D} (Section~\ref{sec:format-LSU-XPL4RBB.D}), \verb|lsucode5 = 000|\\

\bitfields{ 1/P, 2/00, 2/-- --, 27/extend27, 1/1, 2/00, 2/-- --, 27/offset27, 1/1, 2/01, 3/011, 2/variant, 4/registerGh, 2/01, 2/11, 3/000, 1/--, 6/registerY, 6/registerZ }


\paragraph{Syntax} \texttt{xplb}\emph{variant} \emph{registerGh}, \emph{registerY} = \emph{extend27\_\discretionary{}{}{}offset27}[\emph{registerZ}]

\paragraph{Description} The \emph{registerGh} is preloaded from the memory byte at the effective address.


\paragraph{Modifier(s)}
\noindent\begin{itemize}
\item variant (cf. variant in section \ref{par:modifier-variant} on page \pageref{par:modifier-variant})
\end{itemize}

\paragraph{Execution} {Reservation table \textsf{LSU.Y} (Section~\ref{sec:reservation-LSU.Y})}
~
{\begin{lstlisting}
stage ID:
  new offset = _SX_54(extend27_offset27);
  new variant = _ZX_2(variant);
stage RR:
  new base = GPR[registerZ];
  new target = GPR[registerY];
  new address = base + offset;
  new shift = (target & 31) * 8;
  new index = (target >> 5) & 3;
  new where = (registerGh << 2) + index;
  new dri = where;
stage E1:
  new argument1 = XVR[where];
  new bytemask = 1;
  new targetbytes = bits2bytes(bytemask) << shift;
  CS.XMF = 1;
  new result1 = MEM.load(address, 0x1 & bytemask, variant, dri);
  result1 = ((result1 << shift) & targetbytes) | (argument1 & ~targetbytes);
stage CR:
  XVR[where] = result1;
\end{lstlisting}}
\newpage

\paragraph{Encoding} Format \textsf{LSU\_XPL8RBB} (Section~\ref{sec:format-LSU-XPL8RBB}), \verb|lsucode5 = 000|\\

\bitfields{ 1/P, 2/01, 3/011, 2/variant, 3/registerGi, 3/011, 2/11, 3/000, 1/--, 6/registerY, 6/registerZ }


\paragraph{Syntax} \texttt{xplb}\emph{variant} \emph{registerGi}, \emph{registerY} = [\emph{registerZ}]

\paragraph{Description} The \emph{registerGi} is preloaded from the memory byte at the effective address.


\paragraph{Modifier(s)}
\noindent\begin{itemize}
\item variant (cf. variant in section \ref{par:modifier-variant} on page \pageref{par:modifier-variant})
\end{itemize}

\paragraph{Execution} {Reservation table \textsf{LSU} (Section~\ref{sec:reservation-LSU})}
~
{\begin{lstlisting}
stage ID:
  new variant = _ZX_2(variant);
stage RR:
  new address = GPR[registerZ];
  new target = GPR[registerY];
  new shift = (target & 31) * 8;
  new index = (target >> 5) & 7;
  new where = (registerGi << 3) + index;
  new dri = where;
stage E1:
  new argument1 = XVR[where];
  new bytemask = 1;
  new targetbytes = bits2bytes(bytemask) << shift;
  CS.XMF = 1;
  new result1 = MEM.load(address, 0x1 & bytemask, variant, dri);
  result1 = ((result1 << shift) & targetbytes) | (argument1 & ~targetbytes);
stage CR:
  XVR[where] = result1;
\end{lstlisting}}
\newpage

\paragraph{Encoding} Format \textsf{LSU\_XPL8RBB.O} (Section~\ref{sec:format-LSU-XPL8RBB.O}), \verb|lsucode5 = 000|\\

\bitfields{ 1/P, 2/00, 2/-- --, 27/offset27, 1/1, 2/01, 3/011, 2/variant, 3/registerGi, 3/011, 2/11, 3/000, 1/--, 6/registerY, 6/registerZ }


\paragraph{Syntax} \texttt{xplb}\emph{variant} \emph{registerGi}, \emph{registerY} = \emph{offset27}[\emph{registerZ}]

\paragraph{Description} The \emph{registerGi} is preloaded from the memory byte at the effective address.


\paragraph{Modifier(s)}
\noindent\begin{itemize}
\item variant (cf. variant in section \ref{par:modifier-variant} on page \pageref{par:modifier-variant})
\end{itemize}

\paragraph{Execution} {Reservation table \textsf{LSU.X} (Section~\ref{sec:reservation-LSU.X})}
~
{\begin{lstlisting}
stage ID:
  new offset = _SX_27(offset27);
  new variant = _ZX_2(variant);
stage RR:
  new base = GPR[registerZ];
  new target = GPR[registerY];
  new address = base + offset;
  new shift = (target & 31) * 8;
  new index = (target >> 5) & 7;
  new where = (registerGi << 3) + index;
  new dri = where;
stage E1:
  new argument1 = XVR[where];
  new bytemask = 1;
  new targetbytes = bits2bytes(bytemask) << shift;
  CS.XMF = 1;
  new result1 = MEM.load(address, 0x1 & bytemask, variant, dri);
  result1 = ((result1 << shift) & targetbytes) | (argument1 & ~targetbytes);
stage CR:
  XVR[where] = result1;
\end{lstlisting}}
\newpage

\paragraph{Encoding} Format \textsf{LSU\_XPL8RBB.D} (Section~\ref{sec:format-LSU-XPL8RBB.D}), \verb|lsucode5 = 000|\\

\bitfields{ 1/P, 2/00, 2/-- --, 27/extend27, 1/1, 2/00, 2/-- --, 27/offset27, 1/1, 2/01, 3/011, 2/variant, 3/registerGi, 3/011, 2/11, 3/000, 1/--, 6/registerY, 6/registerZ }


\paragraph{Syntax} \texttt{xplb}\emph{variant} \emph{registerGi}, \emph{registerY} = \emph{extend27\_\discretionary{}{}{}offset27}[\emph{registerZ}]

\paragraph{Description} The \emph{registerGi} is preloaded from the memory byte at the effective address.


\paragraph{Modifier(s)}
\noindent\begin{itemize}
\item variant (cf. variant in section \ref{par:modifier-variant} on page \pageref{par:modifier-variant})
\end{itemize}

\paragraph{Execution} {Reservation table \textsf{LSU.Y} (Section~\ref{sec:reservation-LSU.Y})}
~
{\begin{lstlisting}
stage ID:
  new offset = _SX_54(extend27_offset27);
  new variant = _ZX_2(variant);
stage RR:
  new base = GPR[registerZ];
  new target = GPR[registerY];
  new address = base + offset;
  new shift = (target & 31) * 8;
  new index = (target >> 5) & 7;
  new where = (registerGi << 3) + index;
  new dri = where;
stage E1:
  new argument1 = XVR[where];
  new bytemask = 1;
  new targetbytes = bits2bytes(bytemask) << shift;
  CS.XMF = 1;
  new result1 = MEM.load(address, 0x1 & bytemask, variant, dri);
  result1 = ((result1 << shift) & targetbytes) | (argument1 & ~targetbytes);
stage CR:
  XVR[where] = result1;
\end{lstlisting}}
\newpage

\paragraph{Encoding} Format \textsf{LSU\_XPL16RBB} (Section~\ref{sec:format-LSU-XPL16RBB}), \verb|lsucode5 = 000|\\

\bitfields{ 1/P, 2/01, 3/011, 2/variant, 2/registerGj, 4/0111, 2/11, 3/000, 1/--, 6/registerY, 6/registerZ }


\paragraph{Syntax} \texttt{xplb}\emph{variant} \emph{registerGj}, \emph{registerY} = [\emph{registerZ}]

\paragraph{Description} The \emph{registerGj} is preloaded from the memory byte at the effective address.


\paragraph{Modifier(s)}
\noindent\begin{itemize}
\item variant (cf. variant in section \ref{par:modifier-variant} on page \pageref{par:modifier-variant})
\end{itemize}

\paragraph{Execution} {Reservation table \textsf{LSU} (Section~\ref{sec:reservation-LSU})}
~
{\begin{lstlisting}
stage ID:
  new variant = _ZX_2(variant);
stage RR:
  new address = GPR[registerZ];
  new target = GPR[registerY];
  new shift = (target & 31) * 8;
  new index = (target >> 5) & 15;
  new where = (registerGj << 4) + index;
  new dri = where;
stage E1:
  new argument1 = XVR[where];
  new bytemask = 1;
  new targetbytes = bits2bytes(bytemask) << shift;
  CS.XMF = 1;
  new result1 = MEM.load(address, 0x1 & bytemask, variant, dri);
  result1 = ((result1 << shift) & targetbytes) | (argument1 & ~targetbytes);
stage CR:
  XVR[where] = result1;
\end{lstlisting}}
\newpage

\paragraph{Encoding} Format \textsf{LSU\_XPL16RBB.O} (Section~\ref{sec:format-LSU-XPL16RBB.O}), \verb|lsucode5 = 000|\\

\bitfields{ 1/P, 2/00, 2/-- --, 27/offset27, 1/1, 2/01, 3/011, 2/variant, 2/registerGj, 4/0111, 2/11, 3/000, 1/--, 6/registerY, 6/registerZ }


\paragraph{Syntax} \texttt{xplb}\emph{variant} \emph{registerGj}, \emph{registerY} = \emph{offset27}[\emph{registerZ}]

\paragraph{Description} The \emph{registerGj} is preloaded from the memory byte at the effective address.


\paragraph{Modifier(s)}
\noindent\begin{itemize}
\item variant (cf. variant in section \ref{par:modifier-variant} on page \pageref{par:modifier-variant})
\end{itemize}

\paragraph{Execution} {Reservation table \textsf{LSU.X} (Section~\ref{sec:reservation-LSU.X})}
~
{\begin{lstlisting}
stage ID:
  new offset = _SX_27(offset27);
  new variant = _ZX_2(variant);
stage RR:
  new base = GPR[registerZ];
  new target = GPR[registerY];
  new address = base + offset;
  new shift = (target & 31) * 8;
  new index = (target >> 5) & 15;
  new where = (registerGj << 4) + index;
  new dri = where;
stage E1:
  new argument1 = XVR[where];
  new bytemask = 1;
  new targetbytes = bits2bytes(bytemask) << shift;
  CS.XMF = 1;
  new result1 = MEM.load(address, 0x1 & bytemask, variant, dri);
  result1 = ((result1 << shift) & targetbytes) | (argument1 & ~targetbytes);
stage CR:
  XVR[where] = result1;
\end{lstlisting}}
\newpage

\paragraph{Encoding} Format \textsf{LSU\_XPL16RBB.D} (Section~\ref{sec:format-LSU-XPL16RBB.D}), \verb|lsucode5 = 000|\\

\bitfields{ 1/P, 2/00, 2/-- --, 27/extend27, 1/1, 2/00, 2/-- --, 27/offset27, 1/1, 2/01, 3/011, 2/variant, 2/registerGj, 4/0111, 2/11, 3/000, 1/--, 6/registerY, 6/registerZ }


\paragraph{Syntax} \texttt{xplb}\emph{variant} \emph{registerGj}, \emph{registerY} = \emph{extend27\_\discretionary{}{}{}offset27}[\emph{registerZ}]

\paragraph{Description} The \emph{registerGj} is preloaded from the memory byte at the effective address.


\paragraph{Modifier(s)}
\noindent\begin{itemize}
\item variant (cf. variant in section \ref{par:modifier-variant} on page \pageref{par:modifier-variant})
\end{itemize}

\paragraph{Execution} {Reservation table \textsf{LSU.Y} (Section~\ref{sec:reservation-LSU.Y})}
~
{\begin{lstlisting}
stage ID:
  new offset = _SX_54(extend27_offset27);
  new variant = _ZX_2(variant);
stage RR:
  new base = GPR[registerZ];
  new target = GPR[registerY];
  new address = base + offset;
  new shift = (target & 31) * 8;
  new index = (target >> 5) & 15;
  new where = (registerGj << 4) + index;
  new dri = where;
stage E1:
  new argument1 = XVR[where];
  new bytemask = 1;
  new targetbytes = bits2bytes(bytemask) << shift;
  CS.XMF = 1;
  new result1 = MEM.load(address, 0x1 & bytemask, variant, dri);
  result1 = ((result1 << shift) & targetbytes) | (argument1 & ~targetbytes);
stage CR:
  XVR[where] = result1;
\end{lstlisting}}
\newpage

\paragraph{Encoding} Format \textsf{LSU\_XPL32RBB} (Section~\ref{sec:format-LSU-XPL32RBB}), \verb|lsucode5 = 000|\\

\bitfields{ 1/P, 2/01, 3/011, 2/variant, 1/registerGk, 5/01111, 2/11, 3/000, 1/--, 6/registerY, 6/registerZ }


\paragraph{Syntax} \texttt{xplb}\emph{variant} \emph{registerGk}, \emph{registerY} = [\emph{registerZ}]

\paragraph{Description} The \emph{registerGk} is preloaded from the memory byte at the effective address.


\paragraph{Modifier(s)}
\noindent\begin{itemize}
\item variant (cf. variant in section \ref{par:modifier-variant} on page \pageref{par:modifier-variant})
\end{itemize}

\paragraph{Execution} {Reservation table \textsf{LSU} (Section~\ref{sec:reservation-LSU})}
~
{\begin{lstlisting}
stage ID:
  new variant = _ZX_2(variant);
stage RR:
  new address = GPR[registerZ];
  new target = GPR[registerY];
  new shift = (target & 31) * 8;
  new index = (target >> 5) & 31;
  new where = (registerGk << 5) + index;
  new dri = where;
stage E1:
  new argument1 = XVR[where];
  new bytemask = 1;
  new targetbytes = bits2bytes(bytemask) << shift;
  CS.XMF = 1;
  new result1 = MEM.load(address, 0x1 & bytemask, variant, dri);
  result1 = ((result1 << shift) & targetbytes) | (argument1 & ~targetbytes);
stage CR:
  XVR[where] = result1;
\end{lstlisting}}
\newpage

\paragraph{Encoding} Format \textsf{LSU\_XPL32RBB.O} (Section~\ref{sec:format-LSU-XPL32RBB.O}), \verb|lsucode5 = 000|\\

\bitfields{ 1/P, 2/00, 2/-- --, 27/offset27, 1/1, 2/01, 3/011, 2/variant, 1/registerGk, 5/01111, 2/11, 3/000, 1/--, 6/registerY, 6/registerZ }


\paragraph{Syntax} \texttt{xplb}\emph{variant} \emph{registerGk}, \emph{registerY} = \emph{offset27}[\emph{registerZ}]

\paragraph{Description} The \emph{registerGk} is preloaded from the memory byte at the effective address.


\paragraph{Modifier(s)}
\noindent\begin{itemize}
\item variant (cf. variant in section \ref{par:modifier-variant} on page \pageref{par:modifier-variant})
\end{itemize}

\paragraph{Execution} {Reservation table \textsf{LSU.X} (Section~\ref{sec:reservation-LSU.X})}
~
{\begin{lstlisting}
stage ID:
  new offset = _SX_27(offset27);
  new variant = _ZX_2(variant);
stage RR:
  new base = GPR[registerZ];
  new target = GPR[registerY];
  new address = base + offset;
  new shift = (target & 31) * 8;
  new index = (target >> 5) & 31;
  new where = (registerGk << 5) + index;
  new dri = where;
stage E1:
  new argument1 = XVR[where];
  new bytemask = 1;
  new targetbytes = bits2bytes(bytemask) << shift;
  CS.XMF = 1;
  new result1 = MEM.load(address, 0x1 & bytemask, variant, dri);
  result1 = ((result1 << shift) & targetbytes) | (argument1 & ~targetbytes);
stage CR:
  XVR[where] = result1;
\end{lstlisting}}
\newpage

\paragraph{Encoding} Format \textsf{LSU\_XPL32RBB.D} (Section~\ref{sec:format-LSU-XPL32RBB.D}), \verb|lsucode5 = 000|\\

\bitfields{ 1/P, 2/00, 2/-- --, 27/extend27, 1/1, 2/00, 2/-- --, 27/offset27, 1/1, 2/01, 3/011, 2/variant, 1/registerGk, 5/01111, 2/11, 3/000, 1/--, 6/registerY, 6/registerZ }


\paragraph{Syntax} \texttt{xplb}\emph{variant} \emph{registerGk}, \emph{registerY} = \emph{extend27\_\discretionary{}{}{}offset27}[\emph{registerZ}]

\paragraph{Description} The \emph{registerGk} is preloaded from the memory byte at the effective address.


\paragraph{Modifier(s)}
\noindent\begin{itemize}
\item variant (cf. variant in section \ref{par:modifier-variant} on page \pageref{par:modifier-variant})
\end{itemize}

\paragraph{Execution} {Reservation table \textsf{LSU.Y} (Section~\ref{sec:reservation-LSU.Y})}
~
{\begin{lstlisting}
stage ID:
  new offset = _SX_54(extend27_offset27);
  new variant = _ZX_2(variant);
stage RR:
  new base = GPR[registerZ];
  new target = GPR[registerY];
  new address = base + offset;
  new shift = (target & 31) * 8;
  new index = (target >> 5) & 31;
  new where = (registerGk << 5) + index;
  new dri = where;
stage E1:
  new argument1 = XVR[where];
  new bytemask = 1;
  new targetbytes = bits2bytes(bytemask) << shift;
  CS.XMF = 1;
  new result1 = MEM.load(address, 0x1 & bytemask, variant, dri);
  result1 = ((result1 << shift) & targetbytes) | (argument1 & ~targetbytes);
stage CR:
  XVR[where] = result1;
\end{lstlisting}}
\newpage

\paragraph{Encoding} Format \textsf{LSU\_XPL64RBB} (Section~\ref{sec:format-LSU-XPL64RBB}), \verb|lsucode5 = 000|\\

\bitfields{ 1/P, 2/01, 3/011, 2/variant, 1/registerGl, 5/11111, 2/11, 3/000, 1/--, 6/registerY, 6/registerZ }


\paragraph{Syntax} \texttt{xplb}\emph{variant} \emph{registerGl}, \emph{registerY} = [\emph{registerZ}]

\paragraph{Description} The \emph{registerGl} is preloaded from the memory byte at the effective address.


\paragraph{Modifier(s)}
\noindent\begin{itemize}
\item variant (cf. variant in section \ref{par:modifier-variant} on page \pageref{par:modifier-variant})
\end{itemize}

\paragraph{Execution} {Reservation table \textsf{LSU} (Section~\ref{sec:reservation-LSU})}
~
{\begin{lstlisting}
stage ID:
  new variant = _ZX_2(variant);
stage RR:
  new address = GPR[registerZ];
  new target = GPR[registerY];
  new shift = (target & 31) * 8;
  new index = (target >> 5) & 63;
  new where = (registerGl << 6) + index;
  new dri = where;
stage E1:
  new argument1 = XVR[where];
  new bytemask = 1;
  new targetbytes = bits2bytes(bytemask) << shift;
  CS.XMF = 1;
  new result1 = MEM.load(address, 0x1 & bytemask, variant, dri);
  result1 = ((result1 << shift) & targetbytes) | (argument1 & ~targetbytes);
stage CR:
  XVR[where] = result1;
\end{lstlisting}}
\newpage

\paragraph{Encoding} Format \textsf{LSU\_XPL64RBB.O} (Section~\ref{sec:format-LSU-XPL64RBB.O}), \verb|lsucode5 = 000|\\

\bitfields{ 1/P, 2/00, 2/-- --, 27/offset27, 1/1, 2/01, 3/011, 2/variant, 1/registerGl, 5/11111, 2/11, 3/000, 1/--, 6/registerY, 6/registerZ }


\paragraph{Syntax} \texttt{xplb}\emph{variant} \emph{registerGl}, \emph{registerY} = \emph{offset27}[\emph{registerZ}]

\paragraph{Description} The \emph{registerGl} is preloaded from the memory byte at the effective address.


\paragraph{Modifier(s)}
\noindent\begin{itemize}
\item variant (cf. variant in section \ref{par:modifier-variant} on page \pageref{par:modifier-variant})
\end{itemize}

\paragraph{Execution} {Reservation table \textsf{LSU.X} (Section~\ref{sec:reservation-LSU.X})}
~
{\begin{lstlisting}
stage ID:
  new offset = _SX_27(offset27);
  new variant = _ZX_2(variant);
stage RR:
  new base = GPR[registerZ];
  new target = GPR[registerY];
  new address = base + offset;
  new shift = (target & 31) * 8;
  new index = (target >> 5) & 63;
  new where = (registerGl << 6) + index;
  new dri = where;
stage E1:
  new argument1 = XVR[where];
  new bytemask = 1;
  new targetbytes = bits2bytes(bytemask) << shift;
  CS.XMF = 1;
  new result1 = MEM.load(address, 0x1 & bytemask, variant, dri);
  result1 = ((result1 << shift) & targetbytes) | (argument1 & ~targetbytes);
stage CR:
  XVR[where] = result1;
\end{lstlisting}}
\newpage

\paragraph{Encoding} Format \textsf{LSU\_XPL64RBB.D} (Section~\ref{sec:format-LSU-XPL64RBB.D}), \verb|lsucode5 = 000|\\

\bitfields{ 1/P, 2/00, 2/-- --, 27/extend27, 1/1, 2/00, 2/-- --, 27/offset27, 1/1, 2/01, 3/011, 2/variant, 1/registerGl, 5/11111, 2/11, 3/000, 1/--, 6/registerY, 6/registerZ }


\paragraph{Syntax} \texttt{xplb}\emph{variant} \emph{registerGl}, \emph{registerY} = \emph{extend27\_\discretionary{}{}{}offset27}[\emph{registerZ}]

\paragraph{Description} The \emph{registerGl} is preloaded from the memory byte at the effective address.


\paragraph{Modifier(s)}
\noindent\begin{itemize}
\item variant (cf. variant in section \ref{par:modifier-variant} on page \pageref{par:modifier-variant})
\end{itemize}

\paragraph{Execution} {Reservation table \textsf{LSU.Y} (Section~\ref{sec:reservation-LSU.Y})}
~
{\begin{lstlisting}
stage ID:
  new offset = _SX_54(extend27_offset27);
  new variant = _ZX_2(variant);
stage RR:
  new base = GPR[registerZ];
  new target = GPR[registerY];
  new address = base + offset;
  new shift = (target & 31) * 8;
  new index = (target >> 5) & 63;
  new where = (registerGl << 6) + index;
  new dri = where;
stage E1:
  new argument1 = XVR[where];
  new bytemask = 1;
  new targetbytes = bits2bytes(bytemask) << shift;
  CS.XMF = 1;
  new result1 = MEM.load(address, 0x1 & bytemask, variant, dri);
  result1 = ((result1 << shift) & targetbytes) | (argument1 & ~targetbytes);
stage CR:
  XVR[where] = result1;
\end{lstlisting}}
\newpage
\subsubsection[{\small XPLD: Extension Preload Double Word}]{XPLD\hfill Extension Preload Double Word}
\label{instr:XPLD}\index{xpld}
 
\paragraph{Encoding} Format \textsf{LSU\_XPL2RBB} (Section~\ref{sec:format-LSU-XPL2RBB}), \verb|lsucode5 = 011|\\

\bitfields{ 1/P, 2/01, 3/011, 2/variant, 5/registerGg, 1/0, 2/11, 3/011, 1/--, 6/registerY, 6/registerZ }


\paragraph{Syntax} \texttt{xpld}\emph{variant} \emph{registerGg}, \emph{registerY} = [\emph{registerZ}]

\paragraph{Description} The \emph{registerGg} is preloaded from the memory double word at the effective address.


\paragraph{Modifier(s)}
\noindent\begin{itemize}
\item variant (cf. variant in section \ref{par:modifier-variant} on page \pageref{par:modifier-variant})
\end{itemize}

\paragraph{Execution} {Reservation table \textsf{LSU} (Section~\ref{sec:reservation-LSU})}
~
{\begin{lstlisting}
stage ID:
  new variant = _ZX_2(variant);
stage RR:
  new address = GPR[registerZ];
  new target = GPR[registerY];
  new shift = (target & 31) * 8;
  new index = (target >> 5) & 1;
  new where = (registerGg << 1) + index;
  new dri = where;
stage E1:
  new argument1 = XVR[where];
  new bytemask = 255;
  new targetbytes = bits2bytes(bytemask) << shift;
  CS.XMF = 1;
  new result1 = MEM.load(address, 0xFF & bytemask, variant, dri);
  result1 = ((result1 << shift) & targetbytes) | (argument1 & ~targetbytes);
stage CR:
  XVR[where] = result1;
\end{lstlisting}}
\newpage

\paragraph{Encoding} Format \textsf{LSU\_XPL2RBB.O} (Section~\ref{sec:format-LSU-XPL2RBB.O}), \verb|lsucode5 = 011|\\

\bitfields{ 1/P, 2/00, 2/-- --, 27/offset27, 1/1, 2/01, 3/011, 2/variant, 5/registerGg, 1/0, 2/11, 3/011, 1/--, 6/registerY, 6/registerZ }


\paragraph{Syntax} \texttt{xpld}\emph{variant} \emph{registerGg}, \emph{registerY} = \emph{offset27}[\emph{registerZ}]

\paragraph{Description} The \emph{registerGg} is preloaded from the memory double word at the effective address.


\paragraph{Modifier(s)}
\noindent\begin{itemize}
\item variant (cf. variant in section \ref{par:modifier-variant} on page \pageref{par:modifier-variant})
\end{itemize}

\paragraph{Execution} {Reservation table \textsf{LSU.X} (Section~\ref{sec:reservation-LSU.X})}
~
{\begin{lstlisting}
stage ID:
  new offset = _SX_27(offset27);
  new variant = _ZX_2(variant);
stage RR:
  new base = GPR[registerZ];
  new target = GPR[registerY];
  new address = base + offset;
  new shift = (target & 31) * 8;
  new index = (target >> 5) & 1;
  new where = (registerGg << 1) + index;
  new dri = where;
stage E1:
  new argument1 = XVR[where];
  new bytemask = 255;
  new targetbytes = bits2bytes(bytemask) << shift;
  CS.XMF = 1;
  new result1 = MEM.load(address, 0xFF & bytemask, variant, dri);
  result1 = ((result1 << shift) & targetbytes) | (argument1 & ~targetbytes);
stage CR:
  XVR[where] = result1;
\end{lstlisting}}
\newpage

\paragraph{Encoding} Format \textsf{LSU\_XPL2RBB.D} (Section~\ref{sec:format-LSU-XPL2RBB.D}), \verb|lsucode5 = 011|\\

\bitfields{ 1/P, 2/00, 2/-- --, 27/extend27, 1/1, 2/00, 2/-- --, 27/offset27, 1/1, 2/01, 3/011, 2/variant, 5/registerGg, 1/0, 2/11, 3/011, 1/--, 6/registerY, 6/registerZ }


\paragraph{Syntax} \texttt{xpld}\emph{variant} \emph{registerGg}, \emph{registerY} = \emph{extend27\_\discretionary{}{}{}offset27}[\emph{registerZ}]

\paragraph{Description} The \emph{registerGg} is preloaded from the memory double word at the effective address.


\paragraph{Modifier(s)}
\noindent\begin{itemize}
\item variant (cf. variant in section \ref{par:modifier-variant} on page \pageref{par:modifier-variant})
\end{itemize}

\paragraph{Execution} {Reservation table \textsf{LSU.Y} (Section~\ref{sec:reservation-LSU.Y})}
~
{\begin{lstlisting}
stage ID:
  new offset = _SX_54(extend27_offset27);
  new variant = _ZX_2(variant);
stage RR:
  new base = GPR[registerZ];
  new target = GPR[registerY];
  new address = base + offset;
  new shift = (target & 31) * 8;
  new index = (target >> 5) & 1;
  new where = (registerGg << 1) + index;
  new dri = where;
stage E1:
  new argument1 = XVR[where];
  new bytemask = 255;
  new targetbytes = bits2bytes(bytemask) << shift;
  CS.XMF = 1;
  new result1 = MEM.load(address, 0xFF & bytemask, variant, dri);
  result1 = ((result1 << shift) & targetbytes) | (argument1 & ~targetbytes);
stage CR:
  XVR[where] = result1;
\end{lstlisting}}
\newpage

\paragraph{Encoding} Format \textsf{LSU\_XPL4RBB} (Section~\ref{sec:format-LSU-XPL4RBB}), \verb|lsucode5 = 011|\\

\bitfields{ 1/P, 2/01, 3/011, 2/variant, 4/registerGh, 2/01, 2/11, 3/011, 1/--, 6/registerY, 6/registerZ }


\paragraph{Syntax} \texttt{xpld}\emph{variant} \emph{registerGh}, \emph{registerY} = [\emph{registerZ}]

\paragraph{Description} The \emph{registerGh} is preloaded from the memory double word at the effective address.


\paragraph{Modifier(s)}
\noindent\begin{itemize}
\item variant (cf. variant in section \ref{par:modifier-variant} on page \pageref{par:modifier-variant})
\end{itemize}

\paragraph{Execution} {Reservation table \textsf{LSU} (Section~\ref{sec:reservation-LSU})}
~
{\begin{lstlisting}
stage ID:
  new variant = _ZX_2(variant);
stage RR:
  new address = GPR[registerZ];
  new target = GPR[registerY];
  new shift = (target & 31) * 8;
  new index = (target >> 5) & 3;
  new where = (registerGh << 2) + index;
  new dri = where;
stage E1:
  new argument1 = XVR[where];
  new bytemask = 255;
  new targetbytes = bits2bytes(bytemask) << shift;
  CS.XMF = 1;
  new result1 = MEM.load(address, 0xFF & bytemask, variant, dri);
  result1 = ((result1 << shift) & targetbytes) | (argument1 & ~targetbytes);
stage CR:
  XVR[where] = result1;
\end{lstlisting}}
\newpage

\paragraph{Encoding} Format \textsf{LSU\_XPL4RBB.O} (Section~\ref{sec:format-LSU-XPL4RBB.O}), \verb|lsucode5 = 011|\\

\bitfields{ 1/P, 2/00, 2/-- --, 27/offset27, 1/1, 2/01, 3/011, 2/variant, 4/registerGh, 2/01, 2/11, 3/011, 1/--, 6/registerY, 6/registerZ }


\paragraph{Syntax} \texttt{xpld}\emph{variant} \emph{registerGh}, \emph{registerY} = \emph{offset27}[\emph{registerZ}]

\paragraph{Description} The \emph{registerGh} is preloaded from the memory double word at the effective address.


\paragraph{Modifier(s)}
\noindent\begin{itemize}
\item variant (cf. variant in section \ref{par:modifier-variant} on page \pageref{par:modifier-variant})
\end{itemize}

\paragraph{Execution} {Reservation table \textsf{LSU.X} (Section~\ref{sec:reservation-LSU.X})}
~
{\begin{lstlisting}
stage ID:
  new offset = _SX_27(offset27);
  new variant = _ZX_2(variant);
stage RR:
  new base = GPR[registerZ];
  new target = GPR[registerY];
  new address = base + offset;
  new shift = (target & 31) * 8;
  new index = (target >> 5) & 3;
  new where = (registerGh << 2) + index;
  new dri = where;
stage E1:
  new argument1 = XVR[where];
  new bytemask = 255;
  new targetbytes = bits2bytes(bytemask) << shift;
  CS.XMF = 1;
  new result1 = MEM.load(address, 0xFF & bytemask, variant, dri);
  result1 = ((result1 << shift) & targetbytes) | (argument1 & ~targetbytes);
stage CR:
  XVR[where] = result1;
\end{lstlisting}}
\newpage

\paragraph{Encoding} Format \textsf{LSU\_XPL4RBB.D} (Section~\ref{sec:format-LSU-XPL4RBB.D}), \verb|lsucode5 = 011|\\

\bitfields{ 1/P, 2/00, 2/-- --, 27/extend27, 1/1, 2/00, 2/-- --, 27/offset27, 1/1, 2/01, 3/011, 2/variant, 4/registerGh, 2/01, 2/11, 3/011, 1/--, 6/registerY, 6/registerZ }


\paragraph{Syntax} \texttt{xpld}\emph{variant} \emph{registerGh}, \emph{registerY} = \emph{extend27\_\discretionary{}{}{}offset27}[\emph{registerZ}]

\paragraph{Description} The \emph{registerGh} is preloaded from the memory double word at the effective address.


\paragraph{Modifier(s)}
\noindent\begin{itemize}
\item variant (cf. variant in section \ref{par:modifier-variant} on page \pageref{par:modifier-variant})
\end{itemize}

\paragraph{Execution} {Reservation table \textsf{LSU.Y} (Section~\ref{sec:reservation-LSU.Y})}
~
{\begin{lstlisting}
stage ID:
  new offset = _SX_54(extend27_offset27);
  new variant = _ZX_2(variant);
stage RR:
  new base = GPR[registerZ];
  new target = GPR[registerY];
  new address = base + offset;
  new shift = (target & 31) * 8;
  new index = (target >> 5) & 3;
  new where = (registerGh << 2) + index;
  new dri = where;
stage E1:
  new argument1 = XVR[where];
  new bytemask = 255;
  new targetbytes = bits2bytes(bytemask) << shift;
  CS.XMF = 1;
  new result1 = MEM.load(address, 0xFF & bytemask, variant, dri);
  result1 = ((result1 << shift) & targetbytes) | (argument1 & ~targetbytes);
stage CR:
  XVR[where] = result1;
\end{lstlisting}}
\newpage

\paragraph{Encoding} Format \textsf{LSU\_XPL8RBB} (Section~\ref{sec:format-LSU-XPL8RBB}), \verb|lsucode5 = 011|\\

\bitfields{ 1/P, 2/01, 3/011, 2/variant, 3/registerGi, 3/011, 2/11, 3/011, 1/--, 6/registerY, 6/registerZ }


\paragraph{Syntax} \texttt{xpld}\emph{variant} \emph{registerGi}, \emph{registerY} = [\emph{registerZ}]

\paragraph{Description} The \emph{registerGi} is preloaded from the memory double word at the effective address.


\paragraph{Modifier(s)}
\noindent\begin{itemize}
\item variant (cf. variant in section \ref{par:modifier-variant} on page \pageref{par:modifier-variant})
\end{itemize}

\paragraph{Execution} {Reservation table \textsf{LSU} (Section~\ref{sec:reservation-LSU})}
~
{\begin{lstlisting}
stage ID:
  new variant = _ZX_2(variant);
stage RR:
  new address = GPR[registerZ];
  new target = GPR[registerY];
  new shift = (target & 31) * 8;
  new index = (target >> 5) & 7;
  new where = (registerGi << 3) + index;
  new dri = where;
stage E1:
  new argument1 = XVR[where];
  new bytemask = 255;
  new targetbytes = bits2bytes(bytemask) << shift;
  CS.XMF = 1;
  new result1 = MEM.load(address, 0xFF & bytemask, variant, dri);
  result1 = ((result1 << shift) & targetbytes) | (argument1 & ~targetbytes);
stage CR:
  XVR[where] = result1;
\end{lstlisting}}
\newpage

\paragraph{Encoding} Format \textsf{LSU\_XPL8RBB.O} (Section~\ref{sec:format-LSU-XPL8RBB.O}), \verb|lsucode5 = 011|\\

\bitfields{ 1/P, 2/00, 2/-- --, 27/offset27, 1/1, 2/01, 3/011, 2/variant, 3/registerGi, 3/011, 2/11, 3/011, 1/--, 6/registerY, 6/registerZ }


\paragraph{Syntax} \texttt{xpld}\emph{variant} \emph{registerGi}, \emph{registerY} = \emph{offset27}[\emph{registerZ}]

\paragraph{Description} The \emph{registerGi} is preloaded from the memory double word at the effective address.


\paragraph{Modifier(s)}
\noindent\begin{itemize}
\item variant (cf. variant in section \ref{par:modifier-variant} on page \pageref{par:modifier-variant})
\end{itemize}

\paragraph{Execution} {Reservation table \textsf{LSU.X} (Section~\ref{sec:reservation-LSU.X})}
~
{\begin{lstlisting}
stage ID:
  new offset = _SX_27(offset27);
  new variant = _ZX_2(variant);
stage RR:
  new base = GPR[registerZ];
  new target = GPR[registerY];
  new address = base + offset;
  new shift = (target & 31) * 8;
  new index = (target >> 5) & 7;
  new where = (registerGi << 3) + index;
  new dri = where;
stage E1:
  new argument1 = XVR[where];
  new bytemask = 255;
  new targetbytes = bits2bytes(bytemask) << shift;
  CS.XMF = 1;
  new result1 = MEM.load(address, 0xFF & bytemask, variant, dri);
  result1 = ((result1 << shift) & targetbytes) | (argument1 & ~targetbytes);
stage CR:
  XVR[where] = result1;
\end{lstlisting}}
\newpage

\paragraph{Encoding} Format \textsf{LSU\_XPL8RBB.D} (Section~\ref{sec:format-LSU-XPL8RBB.D}), \verb|lsucode5 = 011|\\

\bitfields{ 1/P, 2/00, 2/-- --, 27/extend27, 1/1, 2/00, 2/-- --, 27/offset27, 1/1, 2/01, 3/011, 2/variant, 3/registerGi, 3/011, 2/11, 3/011, 1/--, 6/registerY, 6/registerZ }


\paragraph{Syntax} \texttt{xpld}\emph{variant} \emph{registerGi}, \emph{registerY} = \emph{extend27\_\discretionary{}{}{}offset27}[\emph{registerZ}]

\paragraph{Description} The \emph{registerGi} is preloaded from the memory double word at the effective address.


\paragraph{Modifier(s)}
\noindent\begin{itemize}
\item variant (cf. variant in section \ref{par:modifier-variant} on page \pageref{par:modifier-variant})
\end{itemize}

\paragraph{Execution} {Reservation table \textsf{LSU.Y} (Section~\ref{sec:reservation-LSU.Y})}
~
{\begin{lstlisting}
stage ID:
  new offset = _SX_54(extend27_offset27);
  new variant = _ZX_2(variant);
stage RR:
  new base = GPR[registerZ];
  new target = GPR[registerY];
  new address = base + offset;
  new shift = (target & 31) * 8;
  new index = (target >> 5) & 7;
  new where = (registerGi << 3) + index;
  new dri = where;
stage E1:
  new argument1 = XVR[where];
  new bytemask = 255;
  new targetbytes = bits2bytes(bytemask) << shift;
  CS.XMF = 1;
  new result1 = MEM.load(address, 0xFF & bytemask, variant, dri);
  result1 = ((result1 << shift) & targetbytes) | (argument1 & ~targetbytes);
stage CR:
  XVR[where] = result1;
\end{lstlisting}}
\newpage

\paragraph{Encoding} Format \textsf{LSU\_XPL16RBB} (Section~\ref{sec:format-LSU-XPL16RBB}), \verb|lsucode5 = 011|\\

\bitfields{ 1/P, 2/01, 3/011, 2/variant, 2/registerGj, 4/0111, 2/11, 3/011, 1/--, 6/registerY, 6/registerZ }


\paragraph{Syntax} \texttt{xpld}\emph{variant} \emph{registerGj}, \emph{registerY} = [\emph{registerZ}]

\paragraph{Description} The \emph{registerGj} is preloaded from the memory double word at the effective address.


\paragraph{Modifier(s)}
\noindent\begin{itemize}
\item variant (cf. variant in section \ref{par:modifier-variant} on page \pageref{par:modifier-variant})
\end{itemize}

\paragraph{Execution} {Reservation table \textsf{LSU} (Section~\ref{sec:reservation-LSU})}
~
{\begin{lstlisting}
stage ID:
  new variant = _ZX_2(variant);
stage RR:
  new address = GPR[registerZ];
  new target = GPR[registerY];
  new shift = (target & 31) * 8;
  new index = (target >> 5) & 15;
  new where = (registerGj << 4) + index;
  new dri = where;
stage E1:
  new argument1 = XVR[where];
  new bytemask = 255;
  new targetbytes = bits2bytes(bytemask) << shift;
  CS.XMF = 1;
  new result1 = MEM.load(address, 0xFF & bytemask, variant, dri);
  result1 = ((result1 << shift) & targetbytes) | (argument1 & ~targetbytes);
stage CR:
  XVR[where] = result1;
\end{lstlisting}}
\newpage

\paragraph{Encoding} Format \textsf{LSU\_XPL16RBB.O} (Section~\ref{sec:format-LSU-XPL16RBB.O}), \verb|lsucode5 = 011|\\

\bitfields{ 1/P, 2/00, 2/-- --, 27/offset27, 1/1, 2/01, 3/011, 2/variant, 2/registerGj, 4/0111, 2/11, 3/011, 1/--, 6/registerY, 6/registerZ }


\paragraph{Syntax} \texttt{xpld}\emph{variant} \emph{registerGj}, \emph{registerY} = \emph{offset27}[\emph{registerZ}]

\paragraph{Description} The \emph{registerGj} is preloaded from the memory double word at the effective address.


\paragraph{Modifier(s)}
\noindent\begin{itemize}
\item variant (cf. variant in section \ref{par:modifier-variant} on page \pageref{par:modifier-variant})
\end{itemize}

\paragraph{Execution} {Reservation table \textsf{LSU.X} (Section~\ref{sec:reservation-LSU.X})}
~
{\begin{lstlisting}
stage ID:
  new offset = _SX_27(offset27);
  new variant = _ZX_2(variant);
stage RR:
  new base = GPR[registerZ];
  new target = GPR[registerY];
  new address = base + offset;
  new shift = (target & 31) * 8;
  new index = (target >> 5) & 15;
  new where = (registerGj << 4) + index;
  new dri = where;
stage E1:
  new argument1 = XVR[where];
  new bytemask = 255;
  new targetbytes = bits2bytes(bytemask) << shift;
  CS.XMF = 1;
  new result1 = MEM.load(address, 0xFF & bytemask, variant, dri);
  result1 = ((result1 << shift) & targetbytes) | (argument1 & ~targetbytes);
stage CR:
  XVR[where] = result1;
\end{lstlisting}}
\newpage

\paragraph{Encoding} Format \textsf{LSU\_XPL16RBB.D} (Section~\ref{sec:format-LSU-XPL16RBB.D}), \verb|lsucode5 = 011|\\

\bitfields{ 1/P, 2/00, 2/-- --, 27/extend27, 1/1, 2/00, 2/-- --, 27/offset27, 1/1, 2/01, 3/011, 2/variant, 2/registerGj, 4/0111, 2/11, 3/011, 1/--, 6/registerY, 6/registerZ }


\paragraph{Syntax} \texttt{xpld}\emph{variant} \emph{registerGj}, \emph{registerY} = \emph{extend27\_\discretionary{}{}{}offset27}[\emph{registerZ}]

\paragraph{Description} The \emph{registerGj} is preloaded from the memory double word at the effective address.


\paragraph{Modifier(s)}
\noindent\begin{itemize}
\item variant (cf. variant in section \ref{par:modifier-variant} on page \pageref{par:modifier-variant})
\end{itemize}

\paragraph{Execution} {Reservation table \textsf{LSU.Y} (Section~\ref{sec:reservation-LSU.Y})}
~
{\begin{lstlisting}
stage ID:
  new offset = _SX_54(extend27_offset27);
  new variant = _ZX_2(variant);
stage RR:
  new base = GPR[registerZ];
  new target = GPR[registerY];
  new address = base + offset;
  new shift = (target & 31) * 8;
  new index = (target >> 5) & 15;
  new where = (registerGj << 4) + index;
  new dri = where;
stage E1:
  new argument1 = XVR[where];
  new bytemask = 255;
  new targetbytes = bits2bytes(bytemask) << shift;
  CS.XMF = 1;
  new result1 = MEM.load(address, 0xFF & bytemask, variant, dri);
  result1 = ((result1 << shift) & targetbytes) | (argument1 & ~targetbytes);
stage CR:
  XVR[where] = result1;
\end{lstlisting}}
\newpage

\paragraph{Encoding} Format \textsf{LSU\_XPL32RBB} (Section~\ref{sec:format-LSU-XPL32RBB}), \verb|lsucode5 = 011|\\

\bitfields{ 1/P, 2/01, 3/011, 2/variant, 1/registerGk, 5/01111, 2/11, 3/011, 1/--, 6/registerY, 6/registerZ }


\paragraph{Syntax} \texttt{xpld}\emph{variant} \emph{registerGk}, \emph{registerY} = [\emph{registerZ}]

\paragraph{Description} The \emph{registerGk} is preloaded from the memory double word at the effective address.


\paragraph{Modifier(s)}
\noindent\begin{itemize}
\item variant (cf. variant in section \ref{par:modifier-variant} on page \pageref{par:modifier-variant})
\end{itemize}

\paragraph{Execution} {Reservation table \textsf{LSU} (Section~\ref{sec:reservation-LSU})}
~
{\begin{lstlisting}
stage ID:
  new variant = _ZX_2(variant);
stage RR:
  new address = GPR[registerZ];
  new target = GPR[registerY];
  new shift = (target & 31) * 8;
  new index = (target >> 5) & 31;
  new where = (registerGk << 5) + index;
  new dri = where;
stage E1:
  new argument1 = XVR[where];
  new bytemask = 255;
  new targetbytes = bits2bytes(bytemask) << shift;
  CS.XMF = 1;
  new result1 = MEM.load(address, 0xFF & bytemask, variant, dri);
  result1 = ((result1 << shift) & targetbytes) | (argument1 & ~targetbytes);
stage CR:
  XVR[where] = result1;
\end{lstlisting}}
\newpage

\paragraph{Encoding} Format \textsf{LSU\_XPL32RBB.O} (Section~\ref{sec:format-LSU-XPL32RBB.O}), \verb|lsucode5 = 011|\\

\bitfields{ 1/P, 2/00, 2/-- --, 27/offset27, 1/1, 2/01, 3/011, 2/variant, 1/registerGk, 5/01111, 2/11, 3/011, 1/--, 6/registerY, 6/registerZ }


\paragraph{Syntax} \texttt{xpld}\emph{variant} \emph{registerGk}, \emph{registerY} = \emph{offset27}[\emph{registerZ}]

\paragraph{Description} The \emph{registerGk} is preloaded from the memory double word at the effective address.


\paragraph{Modifier(s)}
\noindent\begin{itemize}
\item variant (cf. variant in section \ref{par:modifier-variant} on page \pageref{par:modifier-variant})
\end{itemize}

\paragraph{Execution} {Reservation table \textsf{LSU.X} (Section~\ref{sec:reservation-LSU.X})}
~
{\begin{lstlisting}
stage ID:
  new offset = _SX_27(offset27);
  new variant = _ZX_2(variant);
stage RR:
  new base = GPR[registerZ];
  new target = GPR[registerY];
  new address = base + offset;
  new shift = (target & 31) * 8;
  new index = (target >> 5) & 31;
  new where = (registerGk << 5) + index;
  new dri = where;
stage E1:
  new argument1 = XVR[where];
  new bytemask = 255;
  new targetbytes = bits2bytes(bytemask) << shift;
  CS.XMF = 1;
  new result1 = MEM.load(address, 0xFF & bytemask, variant, dri);
  result1 = ((result1 << shift) & targetbytes) | (argument1 & ~targetbytes);
stage CR:
  XVR[where] = result1;
\end{lstlisting}}
\newpage

\paragraph{Encoding} Format \textsf{LSU\_XPL32RBB.D} (Section~\ref{sec:format-LSU-XPL32RBB.D}), \verb|lsucode5 = 011|\\

\bitfields{ 1/P, 2/00, 2/-- --, 27/extend27, 1/1, 2/00, 2/-- --, 27/offset27, 1/1, 2/01, 3/011, 2/variant, 1/registerGk, 5/01111, 2/11, 3/011, 1/--, 6/registerY, 6/registerZ }


\paragraph{Syntax} \texttt{xpld}\emph{variant} \emph{registerGk}, \emph{registerY} = \emph{extend27\_\discretionary{}{}{}offset27}[\emph{registerZ}]

\paragraph{Description} The \emph{registerGk} is preloaded from the memory double word at the effective address.


\paragraph{Modifier(s)}
\noindent\begin{itemize}
\item variant (cf. variant in section \ref{par:modifier-variant} on page \pageref{par:modifier-variant})
\end{itemize}

\paragraph{Execution} {Reservation table \textsf{LSU.Y} (Section~\ref{sec:reservation-LSU.Y})}
~
{\begin{lstlisting}
stage ID:
  new offset = _SX_54(extend27_offset27);
  new variant = _ZX_2(variant);
stage RR:
  new base = GPR[registerZ];
  new target = GPR[registerY];
  new address = base + offset;
  new shift = (target & 31) * 8;
  new index = (target >> 5) & 31;
  new where = (registerGk << 5) + index;
  new dri = where;
stage E1:
  new argument1 = XVR[where];
  new bytemask = 255;
  new targetbytes = bits2bytes(bytemask) << shift;
  CS.XMF = 1;
  new result1 = MEM.load(address, 0xFF & bytemask, variant, dri);
  result1 = ((result1 << shift) & targetbytes) | (argument1 & ~targetbytes);
stage CR:
  XVR[where] = result1;
\end{lstlisting}}
\newpage

\paragraph{Encoding} Format \textsf{LSU\_XPL64RBB} (Section~\ref{sec:format-LSU-XPL64RBB}), \verb|lsucode5 = 011|\\

\bitfields{ 1/P, 2/01, 3/011, 2/variant, 1/registerGl, 5/11111, 2/11, 3/011, 1/--, 6/registerY, 6/registerZ }


\paragraph{Syntax} \texttt{xpld}\emph{variant} \emph{registerGl}, \emph{registerY} = [\emph{registerZ}]

\paragraph{Description} The \emph{registerGl} is preloaded from the memory double word at the effective address.


\paragraph{Modifier(s)}
\noindent\begin{itemize}
\item variant (cf. variant in section \ref{par:modifier-variant} on page \pageref{par:modifier-variant})
\end{itemize}

\paragraph{Execution} {Reservation table \textsf{LSU} (Section~\ref{sec:reservation-LSU})}
~
{\begin{lstlisting}
stage ID:
  new variant = _ZX_2(variant);
stage RR:
  new address = GPR[registerZ];
  new target = GPR[registerY];
  new shift = (target & 31) * 8;
  new index = (target >> 5) & 63;
  new where = (registerGl << 6) + index;
  new dri = where;
stage E1:
  new argument1 = XVR[where];
  new bytemask = 255;
  new targetbytes = bits2bytes(bytemask) << shift;
  CS.XMF = 1;
  new result1 = MEM.load(address, 0xFF & bytemask, variant, dri);
  result1 = ((result1 << shift) & targetbytes) | (argument1 & ~targetbytes);
stage CR:
  XVR[where] = result1;
\end{lstlisting}}
\newpage

\paragraph{Encoding} Format \textsf{LSU\_XPL64RBB.O} (Section~\ref{sec:format-LSU-XPL64RBB.O}), \verb|lsucode5 = 011|\\

\bitfields{ 1/P, 2/00, 2/-- --, 27/offset27, 1/1, 2/01, 3/011, 2/variant, 1/registerGl, 5/11111, 2/11, 3/011, 1/--, 6/registerY, 6/registerZ }


\paragraph{Syntax} \texttt{xpld}\emph{variant} \emph{registerGl}, \emph{registerY} = \emph{offset27}[\emph{registerZ}]

\paragraph{Description} The \emph{registerGl} is preloaded from the memory double word at the effective address.


\paragraph{Modifier(s)}
\noindent\begin{itemize}
\item variant (cf. variant in section \ref{par:modifier-variant} on page \pageref{par:modifier-variant})
\end{itemize}

\paragraph{Execution} {Reservation table \textsf{LSU.X} (Section~\ref{sec:reservation-LSU.X})}
~
{\begin{lstlisting}
stage ID:
  new offset = _SX_27(offset27);
  new variant = _ZX_2(variant);
stage RR:
  new base = GPR[registerZ];
  new target = GPR[registerY];
  new address = base + offset;
  new shift = (target & 31) * 8;
  new index = (target >> 5) & 63;
  new where = (registerGl << 6) + index;
  new dri = where;
stage E1:
  new argument1 = XVR[where];
  new bytemask = 255;
  new targetbytes = bits2bytes(bytemask) << shift;
  CS.XMF = 1;
  new result1 = MEM.load(address, 0xFF & bytemask, variant, dri);
  result1 = ((result1 << shift) & targetbytes) | (argument1 & ~targetbytes);
stage CR:
  XVR[where] = result1;
\end{lstlisting}}
\newpage

\paragraph{Encoding} Format \textsf{LSU\_XPL64RBB.D} (Section~\ref{sec:format-LSU-XPL64RBB.D}), \verb|lsucode5 = 011|\\

\bitfields{ 1/P, 2/00, 2/-- --, 27/extend27, 1/1, 2/00, 2/-- --, 27/offset27, 1/1, 2/01, 3/011, 2/variant, 1/registerGl, 5/11111, 2/11, 3/011, 1/--, 6/registerY, 6/registerZ }


\paragraph{Syntax} \texttt{xpld}\emph{variant} \emph{registerGl}, \emph{registerY} = \emph{extend27\_\discretionary{}{}{}offset27}[\emph{registerZ}]

\paragraph{Description} The \emph{registerGl} is preloaded from the memory double word at the effective address.


\paragraph{Modifier(s)}
\noindent\begin{itemize}
\item variant (cf. variant in section \ref{par:modifier-variant} on page \pageref{par:modifier-variant})
\end{itemize}

\paragraph{Execution} {Reservation table \textsf{LSU.Y} (Section~\ref{sec:reservation-LSU.Y})}
~
{\begin{lstlisting}
stage ID:
  new offset = _SX_54(extend27_offset27);
  new variant = _ZX_2(variant);
stage RR:
  new base = GPR[registerZ];
  new target = GPR[registerY];
  new address = base + offset;
  new shift = (target & 31) * 8;
  new index = (target >> 5) & 63;
  new where = (registerGl << 6) + index;
  new dri = where;
stage E1:
  new argument1 = XVR[where];
  new bytemask = 255;
  new targetbytes = bits2bytes(bytemask) << shift;
  CS.XMF = 1;
  new result1 = MEM.load(address, 0xFF & bytemask, variant, dri);
  result1 = ((result1 << shift) & targetbytes) | (argument1 & ~targetbytes);
stage CR:
  XVR[where] = result1;
\end{lstlisting}}
\newpage
\subsubsection[{\small XPLH: Extension Preload Half}]{XPLH\hfill Extension Preload Half}
\label{instr:XPLH}\index{xplh}
 
\paragraph{Encoding} Format \textsf{LSU\_XPL2RBB} (Section~\ref{sec:format-LSU-XPL2RBB}), \verb|lsucode5 = 001|\\

\bitfields{ 1/P, 2/01, 3/011, 2/variant, 5/registerGg, 1/0, 2/11, 3/001, 1/--, 6/registerY, 6/registerZ }


\paragraph{Syntax} \texttt{xplh}\emph{variant} \emph{registerGg}, \emph{registerY} = [\emph{registerZ}]

\paragraph{Description} The \emph{registerGg} is preloaded from the memory half word at the effective address.


\paragraph{Modifier(s)}
\noindent\begin{itemize}
\item variant (cf. variant in section \ref{par:modifier-variant} on page \pageref{par:modifier-variant})
\end{itemize}

\paragraph{Execution} {Reservation table \textsf{LSU} (Section~\ref{sec:reservation-LSU})}
~
{\begin{lstlisting}
stage ID:
  new variant = _ZX_2(variant);
stage RR:
  new address = GPR[registerZ];
  new target = GPR[registerY];
  new shift = (target & 31) * 8;
  new index = (target >> 5) & 1;
  new where = (registerGg << 1) + index;
  new dri = where;
stage E1:
  new argument1 = XVR[where];
  new bytemask = 3;
  new targetbytes = bits2bytes(bytemask) << shift;
  CS.XMF = 1;
  new result1 = MEM.load(address, 0x3 & bytemask, variant, dri);
  result1 = ((result1 << shift) & targetbytes) | (argument1 & ~targetbytes);
stage CR:
  XVR[where] = result1;
\end{lstlisting}}
\newpage

\paragraph{Encoding} Format \textsf{LSU\_XPL2RBB.O} (Section~\ref{sec:format-LSU-XPL2RBB.O}), \verb|lsucode5 = 001|\\

\bitfields{ 1/P, 2/00, 2/-- --, 27/offset27, 1/1, 2/01, 3/011, 2/variant, 5/registerGg, 1/0, 2/11, 3/001, 1/--, 6/registerY, 6/registerZ }


\paragraph{Syntax} \texttt{xplh}\emph{variant} \emph{registerGg}, \emph{registerY} = \emph{offset27}[\emph{registerZ}]

\paragraph{Description} The \emph{registerGg} is preloaded from the memory half word at the effective address.


\paragraph{Modifier(s)}
\noindent\begin{itemize}
\item variant (cf. variant in section \ref{par:modifier-variant} on page \pageref{par:modifier-variant})
\end{itemize}

\paragraph{Execution} {Reservation table \textsf{LSU.X} (Section~\ref{sec:reservation-LSU.X})}
~
{\begin{lstlisting}
stage ID:
  new offset = _SX_27(offset27);
  new variant = _ZX_2(variant);
stage RR:
  new base = GPR[registerZ];
  new target = GPR[registerY];
  new address = base + offset;
  new shift = (target & 31) * 8;
  new index = (target >> 5) & 1;
  new where = (registerGg << 1) + index;
  new dri = where;
stage E1:
  new argument1 = XVR[where];
  new bytemask = 3;
  new targetbytes = bits2bytes(bytemask) << shift;
  CS.XMF = 1;
  new result1 = MEM.load(address, 0x3 & bytemask, variant, dri);
  result1 = ((result1 << shift) & targetbytes) | (argument1 & ~targetbytes);
stage CR:
  XVR[where] = result1;
\end{lstlisting}}
\newpage

\paragraph{Encoding} Format \textsf{LSU\_XPL2RBB.D} (Section~\ref{sec:format-LSU-XPL2RBB.D}), \verb|lsucode5 = 001|\\

\bitfields{ 1/P, 2/00, 2/-- --, 27/extend27, 1/1, 2/00, 2/-- --, 27/offset27, 1/1, 2/01, 3/011, 2/variant, 5/registerGg, 1/0, 2/11, 3/001, 1/--, 6/registerY, 6/registerZ }


\paragraph{Syntax} \texttt{xplh}\emph{variant} \emph{registerGg}, \emph{registerY} = \emph{extend27\_\discretionary{}{}{}offset27}[\emph{registerZ}]

\paragraph{Description} The \emph{registerGg} is preloaded from the memory half word at the effective address.


\paragraph{Modifier(s)}
\noindent\begin{itemize}
\item variant (cf. variant in section \ref{par:modifier-variant} on page \pageref{par:modifier-variant})
\end{itemize}

\paragraph{Execution} {Reservation table \textsf{LSU.Y} (Section~\ref{sec:reservation-LSU.Y})}
~
{\begin{lstlisting}
stage ID:
  new offset = _SX_54(extend27_offset27);
  new variant = _ZX_2(variant);
stage RR:
  new base = GPR[registerZ];
  new target = GPR[registerY];
  new address = base + offset;
  new shift = (target & 31) * 8;
  new index = (target >> 5) & 1;
  new where = (registerGg << 1) + index;
  new dri = where;
stage E1:
  new argument1 = XVR[where];
  new bytemask = 3;
  new targetbytes = bits2bytes(bytemask) << shift;
  CS.XMF = 1;
  new result1 = MEM.load(address, 0x3 & bytemask, variant, dri);
  result1 = ((result1 << shift) & targetbytes) | (argument1 & ~targetbytes);
stage CR:
  XVR[where] = result1;
\end{lstlisting}}
\newpage

\paragraph{Encoding} Format \textsf{LSU\_XPL4RBB} (Section~\ref{sec:format-LSU-XPL4RBB}), \verb|lsucode5 = 001|\\

\bitfields{ 1/P, 2/01, 3/011, 2/variant, 4/registerGh, 2/01, 2/11, 3/001, 1/--, 6/registerY, 6/registerZ }


\paragraph{Syntax} \texttt{xplh}\emph{variant} \emph{registerGh}, \emph{registerY} = [\emph{registerZ}]

\paragraph{Description} The \emph{registerGh} is preloaded from the memory half word at the effective address.


\paragraph{Modifier(s)}
\noindent\begin{itemize}
\item variant (cf. variant in section \ref{par:modifier-variant} on page \pageref{par:modifier-variant})
\end{itemize}

\paragraph{Execution} {Reservation table \textsf{LSU} (Section~\ref{sec:reservation-LSU})}
~
{\begin{lstlisting}
stage ID:
  new variant = _ZX_2(variant);
stage RR:
  new address = GPR[registerZ];
  new target = GPR[registerY];
  new shift = (target & 31) * 8;
  new index = (target >> 5) & 3;
  new where = (registerGh << 2) + index;
  new dri = where;
stage E1:
  new argument1 = XVR[where];
  new bytemask = 3;
  new targetbytes = bits2bytes(bytemask) << shift;
  CS.XMF = 1;
  new result1 = MEM.load(address, 0x3 & bytemask, variant, dri);
  result1 = ((result1 << shift) & targetbytes) | (argument1 & ~targetbytes);
stage CR:
  XVR[where] = result1;
\end{lstlisting}}
\newpage

\paragraph{Encoding} Format \textsf{LSU\_XPL4RBB.O} (Section~\ref{sec:format-LSU-XPL4RBB.O}), \verb|lsucode5 = 001|\\

\bitfields{ 1/P, 2/00, 2/-- --, 27/offset27, 1/1, 2/01, 3/011, 2/variant, 4/registerGh, 2/01, 2/11, 3/001, 1/--, 6/registerY, 6/registerZ }


\paragraph{Syntax} \texttt{xplh}\emph{variant} \emph{registerGh}, \emph{registerY} = \emph{offset27}[\emph{registerZ}]

\paragraph{Description} The \emph{registerGh} is preloaded from the memory half word at the effective address.


\paragraph{Modifier(s)}
\noindent\begin{itemize}
\item variant (cf. variant in section \ref{par:modifier-variant} on page \pageref{par:modifier-variant})
\end{itemize}

\paragraph{Execution} {Reservation table \textsf{LSU.X} (Section~\ref{sec:reservation-LSU.X})}
~
{\begin{lstlisting}
stage ID:
  new offset = _SX_27(offset27);
  new variant = _ZX_2(variant);
stage RR:
  new base = GPR[registerZ];
  new target = GPR[registerY];
  new address = base + offset;
  new shift = (target & 31) * 8;
  new index = (target >> 5) & 3;
  new where = (registerGh << 2) + index;
  new dri = where;
stage E1:
  new argument1 = XVR[where];
  new bytemask = 3;
  new targetbytes = bits2bytes(bytemask) << shift;
  CS.XMF = 1;
  new result1 = MEM.load(address, 0x3 & bytemask, variant, dri);
  result1 = ((result1 << shift) & targetbytes) | (argument1 & ~targetbytes);
stage CR:
  XVR[where] = result1;
\end{lstlisting}}
\newpage

\paragraph{Encoding} Format \textsf{LSU\_XPL4RBB.D} (Section~\ref{sec:format-LSU-XPL4RBB.D}), \verb|lsucode5 = 001|\\

\bitfields{ 1/P, 2/00, 2/-- --, 27/extend27, 1/1, 2/00, 2/-- --, 27/offset27, 1/1, 2/01, 3/011, 2/variant, 4/registerGh, 2/01, 2/11, 3/001, 1/--, 6/registerY, 6/registerZ }


\paragraph{Syntax} \texttt{xplh}\emph{variant} \emph{registerGh}, \emph{registerY} = \emph{extend27\_\discretionary{}{}{}offset27}[\emph{registerZ}]

\paragraph{Description} The \emph{registerGh} is preloaded from the memory half word at the effective address.


\paragraph{Modifier(s)}
\noindent\begin{itemize}
\item variant (cf. variant in section \ref{par:modifier-variant} on page \pageref{par:modifier-variant})
\end{itemize}

\paragraph{Execution} {Reservation table \textsf{LSU.Y} (Section~\ref{sec:reservation-LSU.Y})}
~
{\begin{lstlisting}
stage ID:
  new offset = _SX_54(extend27_offset27);
  new variant = _ZX_2(variant);
stage RR:
  new base = GPR[registerZ];
  new target = GPR[registerY];
  new address = base + offset;
  new shift = (target & 31) * 8;
  new index = (target >> 5) & 3;
  new where = (registerGh << 2) + index;
  new dri = where;
stage E1:
  new argument1 = XVR[where];
  new bytemask = 3;
  new targetbytes = bits2bytes(bytemask) << shift;
  CS.XMF = 1;
  new result1 = MEM.load(address, 0x3 & bytemask, variant, dri);
  result1 = ((result1 << shift) & targetbytes) | (argument1 & ~targetbytes);
stage CR:
  XVR[where] = result1;
\end{lstlisting}}
\newpage

\paragraph{Encoding} Format \textsf{LSU\_XPL8RBB} (Section~\ref{sec:format-LSU-XPL8RBB}), \verb|lsucode5 = 001|\\

\bitfields{ 1/P, 2/01, 3/011, 2/variant, 3/registerGi, 3/011, 2/11, 3/001, 1/--, 6/registerY, 6/registerZ }


\paragraph{Syntax} \texttt{xplh}\emph{variant} \emph{registerGi}, \emph{registerY} = [\emph{registerZ}]

\paragraph{Description} The \emph{registerGi} is preloaded from the memory half word at the effective address.


\paragraph{Modifier(s)}
\noindent\begin{itemize}
\item variant (cf. variant in section \ref{par:modifier-variant} on page \pageref{par:modifier-variant})
\end{itemize}

\paragraph{Execution} {Reservation table \textsf{LSU} (Section~\ref{sec:reservation-LSU})}
~
{\begin{lstlisting}
stage ID:
  new variant = _ZX_2(variant);
stage RR:
  new address = GPR[registerZ];
  new target = GPR[registerY];
  new shift = (target & 31) * 8;
  new index = (target >> 5) & 7;
  new where = (registerGi << 3) + index;
  new dri = where;
stage E1:
  new argument1 = XVR[where];
  new bytemask = 3;
  new targetbytes = bits2bytes(bytemask) << shift;
  CS.XMF = 1;
  new result1 = MEM.load(address, 0x3 & bytemask, variant, dri);
  result1 = ((result1 << shift) & targetbytes) | (argument1 & ~targetbytes);
stage CR:
  XVR[where] = result1;
\end{lstlisting}}
\newpage

\paragraph{Encoding} Format \textsf{LSU\_XPL8RBB.O} (Section~\ref{sec:format-LSU-XPL8RBB.O}), \verb|lsucode5 = 001|\\

\bitfields{ 1/P, 2/00, 2/-- --, 27/offset27, 1/1, 2/01, 3/011, 2/variant, 3/registerGi, 3/011, 2/11, 3/001, 1/--, 6/registerY, 6/registerZ }


\paragraph{Syntax} \texttt{xplh}\emph{variant} \emph{registerGi}, \emph{registerY} = \emph{offset27}[\emph{registerZ}]

\paragraph{Description} The \emph{registerGi} is preloaded from the memory half word at the effective address.


\paragraph{Modifier(s)}
\noindent\begin{itemize}
\item variant (cf. variant in section \ref{par:modifier-variant} on page \pageref{par:modifier-variant})
\end{itemize}

\paragraph{Execution} {Reservation table \textsf{LSU.X} (Section~\ref{sec:reservation-LSU.X})}
~
{\begin{lstlisting}
stage ID:
  new offset = _SX_27(offset27);
  new variant = _ZX_2(variant);
stage RR:
  new base = GPR[registerZ];
  new target = GPR[registerY];
  new address = base + offset;
  new shift = (target & 31) * 8;
  new index = (target >> 5) & 7;
  new where = (registerGi << 3) + index;
  new dri = where;
stage E1:
  new argument1 = XVR[where];
  new bytemask = 3;
  new targetbytes = bits2bytes(bytemask) << shift;
  CS.XMF = 1;
  new result1 = MEM.load(address, 0x3 & bytemask, variant, dri);
  result1 = ((result1 << shift) & targetbytes) | (argument1 & ~targetbytes);
stage CR:
  XVR[where] = result1;
\end{lstlisting}}
\newpage

\paragraph{Encoding} Format \textsf{LSU\_XPL8RBB.D} (Section~\ref{sec:format-LSU-XPL8RBB.D}), \verb|lsucode5 = 001|\\

\bitfields{ 1/P, 2/00, 2/-- --, 27/extend27, 1/1, 2/00, 2/-- --, 27/offset27, 1/1, 2/01, 3/011, 2/variant, 3/registerGi, 3/011, 2/11, 3/001, 1/--, 6/registerY, 6/registerZ }


\paragraph{Syntax} \texttt{xplh}\emph{variant} \emph{registerGi}, \emph{registerY} = \emph{extend27\_\discretionary{}{}{}offset27}[\emph{registerZ}]

\paragraph{Description} The \emph{registerGi} is preloaded from the memory half word at the effective address.


\paragraph{Modifier(s)}
\noindent\begin{itemize}
\item variant (cf. variant in section \ref{par:modifier-variant} on page \pageref{par:modifier-variant})
\end{itemize}

\paragraph{Execution} {Reservation table \textsf{LSU.Y} (Section~\ref{sec:reservation-LSU.Y})}
~
{\begin{lstlisting}
stage ID:
  new offset = _SX_54(extend27_offset27);
  new variant = _ZX_2(variant);
stage RR:
  new base = GPR[registerZ];
  new target = GPR[registerY];
  new address = base + offset;
  new shift = (target & 31) * 8;
  new index = (target >> 5) & 7;
  new where = (registerGi << 3) + index;
  new dri = where;
stage E1:
  new argument1 = XVR[where];
  new bytemask = 3;
  new targetbytes = bits2bytes(bytemask) << shift;
  CS.XMF = 1;
  new result1 = MEM.load(address, 0x3 & bytemask, variant, dri);
  result1 = ((result1 << shift) & targetbytes) | (argument1 & ~targetbytes);
stage CR:
  XVR[where] = result1;
\end{lstlisting}}
\newpage

\paragraph{Encoding} Format \textsf{LSU\_XPL16RBB} (Section~\ref{sec:format-LSU-XPL16RBB}), \verb|lsucode5 = 001|\\

\bitfields{ 1/P, 2/01, 3/011, 2/variant, 2/registerGj, 4/0111, 2/11, 3/001, 1/--, 6/registerY, 6/registerZ }


\paragraph{Syntax} \texttt{xplh}\emph{variant} \emph{registerGj}, \emph{registerY} = [\emph{registerZ}]

\paragraph{Description} The \emph{registerGj} is preloaded from the memory half word at the effective address.


\paragraph{Modifier(s)}
\noindent\begin{itemize}
\item variant (cf. variant in section \ref{par:modifier-variant} on page \pageref{par:modifier-variant})
\end{itemize}

\paragraph{Execution} {Reservation table \textsf{LSU} (Section~\ref{sec:reservation-LSU})}
~
{\begin{lstlisting}
stage ID:
  new variant = _ZX_2(variant);
stage RR:
  new address = GPR[registerZ];
  new target = GPR[registerY];
  new shift = (target & 31) * 8;
  new index = (target >> 5) & 15;
  new where = (registerGj << 4) + index;
  new dri = where;
stage E1:
  new argument1 = XVR[where];
  new bytemask = 3;
  new targetbytes = bits2bytes(bytemask) << shift;
  CS.XMF = 1;
  new result1 = MEM.load(address, 0x3 & bytemask, variant, dri);
  result1 = ((result1 << shift) & targetbytes) | (argument1 & ~targetbytes);
stage CR:
  XVR[where] = result1;
\end{lstlisting}}
\newpage

\paragraph{Encoding} Format \textsf{LSU\_XPL16RBB.O} (Section~\ref{sec:format-LSU-XPL16RBB.O}), \verb|lsucode5 = 001|\\

\bitfields{ 1/P, 2/00, 2/-- --, 27/offset27, 1/1, 2/01, 3/011, 2/variant, 2/registerGj, 4/0111, 2/11, 3/001, 1/--, 6/registerY, 6/registerZ }


\paragraph{Syntax} \texttt{xplh}\emph{variant} \emph{registerGj}, \emph{registerY} = \emph{offset27}[\emph{registerZ}]

\paragraph{Description} The \emph{registerGj} is preloaded from the memory half word at the effective address.


\paragraph{Modifier(s)}
\noindent\begin{itemize}
\item variant (cf. variant in section \ref{par:modifier-variant} on page \pageref{par:modifier-variant})
\end{itemize}

\paragraph{Execution} {Reservation table \textsf{LSU.X} (Section~\ref{sec:reservation-LSU.X})}
~
{\begin{lstlisting}
stage ID:
  new offset = _SX_27(offset27);
  new variant = _ZX_2(variant);
stage RR:
  new base = GPR[registerZ];
  new target = GPR[registerY];
  new address = base + offset;
  new shift = (target & 31) * 8;
  new index = (target >> 5) & 15;
  new where = (registerGj << 4) + index;
  new dri = where;
stage E1:
  new argument1 = XVR[where];
  new bytemask = 3;
  new targetbytes = bits2bytes(bytemask) << shift;
  CS.XMF = 1;
  new result1 = MEM.load(address, 0x3 & bytemask, variant, dri);
  result1 = ((result1 << shift) & targetbytes) | (argument1 & ~targetbytes);
stage CR:
  XVR[where] = result1;
\end{lstlisting}}
\newpage

\paragraph{Encoding} Format \textsf{LSU\_XPL16RBB.D} (Section~\ref{sec:format-LSU-XPL16RBB.D}), \verb|lsucode5 = 001|\\

\bitfields{ 1/P, 2/00, 2/-- --, 27/extend27, 1/1, 2/00, 2/-- --, 27/offset27, 1/1, 2/01, 3/011, 2/variant, 2/registerGj, 4/0111, 2/11, 3/001, 1/--, 6/registerY, 6/registerZ }


\paragraph{Syntax} \texttt{xplh}\emph{variant} \emph{registerGj}, \emph{registerY} = \emph{extend27\_\discretionary{}{}{}offset27}[\emph{registerZ}]

\paragraph{Description} The \emph{registerGj} is preloaded from the memory half word at the effective address.


\paragraph{Modifier(s)}
\noindent\begin{itemize}
\item variant (cf. variant in section \ref{par:modifier-variant} on page \pageref{par:modifier-variant})
\end{itemize}

\paragraph{Execution} {Reservation table \textsf{LSU.Y} (Section~\ref{sec:reservation-LSU.Y})}
~
{\begin{lstlisting}
stage ID:
  new offset = _SX_54(extend27_offset27);
  new variant = _ZX_2(variant);
stage RR:
  new base = GPR[registerZ];
  new target = GPR[registerY];
  new address = base + offset;
  new shift = (target & 31) * 8;
  new index = (target >> 5) & 15;
  new where = (registerGj << 4) + index;
  new dri = where;
stage E1:
  new argument1 = XVR[where];
  new bytemask = 3;
  new targetbytes = bits2bytes(bytemask) << shift;
  CS.XMF = 1;
  new result1 = MEM.load(address, 0x3 & bytemask, variant, dri);
  result1 = ((result1 << shift) & targetbytes) | (argument1 & ~targetbytes);
stage CR:
  XVR[where] = result1;
\end{lstlisting}}
\newpage

\paragraph{Encoding} Format \textsf{LSU\_XPL32RBB} (Section~\ref{sec:format-LSU-XPL32RBB}), \verb|lsucode5 = 001|\\

\bitfields{ 1/P, 2/01, 3/011, 2/variant, 1/registerGk, 5/01111, 2/11, 3/001, 1/--, 6/registerY, 6/registerZ }


\paragraph{Syntax} \texttt{xplh}\emph{variant} \emph{registerGk}, \emph{registerY} = [\emph{registerZ}]

\paragraph{Description} The \emph{registerGk} is preloaded from the memory half word at the effective address.


\paragraph{Modifier(s)}
\noindent\begin{itemize}
\item variant (cf. variant in section \ref{par:modifier-variant} on page \pageref{par:modifier-variant})
\end{itemize}

\paragraph{Execution} {Reservation table \textsf{LSU} (Section~\ref{sec:reservation-LSU})}
~
{\begin{lstlisting}
stage ID:
  new variant = _ZX_2(variant);
stage RR:
  new address = GPR[registerZ];
  new target = GPR[registerY];
  new shift = (target & 31) * 8;
  new index = (target >> 5) & 31;
  new where = (registerGk << 5) + index;
  new dri = where;
stage E1:
  new argument1 = XVR[where];
  new bytemask = 3;
  new targetbytes = bits2bytes(bytemask) << shift;
  CS.XMF = 1;
  new result1 = MEM.load(address, 0x3 & bytemask, variant, dri);
  result1 = ((result1 << shift) & targetbytes) | (argument1 & ~targetbytes);
stage CR:
  XVR[where] = result1;
\end{lstlisting}}
\newpage

\paragraph{Encoding} Format \textsf{LSU\_XPL32RBB.O} (Section~\ref{sec:format-LSU-XPL32RBB.O}), \verb|lsucode5 = 001|\\

\bitfields{ 1/P, 2/00, 2/-- --, 27/offset27, 1/1, 2/01, 3/011, 2/variant, 1/registerGk, 5/01111, 2/11, 3/001, 1/--, 6/registerY, 6/registerZ }


\paragraph{Syntax} \texttt{xplh}\emph{variant} \emph{registerGk}, \emph{registerY} = \emph{offset27}[\emph{registerZ}]

\paragraph{Description} The \emph{registerGk} is preloaded from the memory half word at the effective address.


\paragraph{Modifier(s)}
\noindent\begin{itemize}
\item variant (cf. variant in section \ref{par:modifier-variant} on page \pageref{par:modifier-variant})
\end{itemize}

\paragraph{Execution} {Reservation table \textsf{LSU.X} (Section~\ref{sec:reservation-LSU.X})}
~
{\begin{lstlisting}
stage ID:
  new offset = _SX_27(offset27);
  new variant = _ZX_2(variant);
stage RR:
  new base = GPR[registerZ];
  new target = GPR[registerY];
  new address = base + offset;
  new shift = (target & 31) * 8;
  new index = (target >> 5) & 31;
  new where = (registerGk << 5) + index;
  new dri = where;
stage E1:
  new argument1 = XVR[where];
  new bytemask = 3;
  new targetbytes = bits2bytes(bytemask) << shift;
  CS.XMF = 1;
  new result1 = MEM.load(address, 0x3 & bytemask, variant, dri);
  result1 = ((result1 << shift) & targetbytes) | (argument1 & ~targetbytes);
stage CR:
  XVR[where] = result1;
\end{lstlisting}}
\newpage

\paragraph{Encoding} Format \textsf{LSU\_XPL32RBB.D} (Section~\ref{sec:format-LSU-XPL32RBB.D}), \verb|lsucode5 = 001|\\

\bitfields{ 1/P, 2/00, 2/-- --, 27/extend27, 1/1, 2/00, 2/-- --, 27/offset27, 1/1, 2/01, 3/011, 2/variant, 1/registerGk, 5/01111, 2/11, 3/001, 1/--, 6/registerY, 6/registerZ }


\paragraph{Syntax} \texttt{xplh}\emph{variant} \emph{registerGk}, \emph{registerY} = \emph{extend27\_\discretionary{}{}{}offset27}[\emph{registerZ}]

\paragraph{Description} The \emph{registerGk} is preloaded from the memory half word at the effective address.


\paragraph{Modifier(s)}
\noindent\begin{itemize}
\item variant (cf. variant in section \ref{par:modifier-variant} on page \pageref{par:modifier-variant})
\end{itemize}

\paragraph{Execution} {Reservation table \textsf{LSU.Y} (Section~\ref{sec:reservation-LSU.Y})}
~
{\begin{lstlisting}
stage ID:
  new offset = _SX_54(extend27_offset27);
  new variant = _ZX_2(variant);
stage RR:
  new base = GPR[registerZ];
  new target = GPR[registerY];
  new address = base + offset;
  new shift = (target & 31) * 8;
  new index = (target >> 5) & 31;
  new where = (registerGk << 5) + index;
  new dri = where;
stage E1:
  new argument1 = XVR[where];
  new bytemask = 3;
  new targetbytes = bits2bytes(bytemask) << shift;
  CS.XMF = 1;
  new result1 = MEM.load(address, 0x3 & bytemask, variant, dri);
  result1 = ((result1 << shift) & targetbytes) | (argument1 & ~targetbytes);
stage CR:
  XVR[where] = result1;
\end{lstlisting}}
\newpage

\paragraph{Encoding} Format \textsf{LSU\_XPL64RBB} (Section~\ref{sec:format-LSU-XPL64RBB}), \verb|lsucode5 = 001|\\

\bitfields{ 1/P, 2/01, 3/011, 2/variant, 1/registerGl, 5/11111, 2/11, 3/001, 1/--, 6/registerY, 6/registerZ }


\paragraph{Syntax} \texttt{xplh}\emph{variant} \emph{registerGl}, \emph{registerY} = [\emph{registerZ}]

\paragraph{Description} The \emph{registerGl} is preloaded from the memory half word at the effective address.


\paragraph{Modifier(s)}
\noindent\begin{itemize}
\item variant (cf. variant in section \ref{par:modifier-variant} on page \pageref{par:modifier-variant})
\end{itemize}

\paragraph{Execution} {Reservation table \textsf{LSU} (Section~\ref{sec:reservation-LSU})}
~
{\begin{lstlisting}
stage ID:
  new variant = _ZX_2(variant);
stage RR:
  new address = GPR[registerZ];
  new target = GPR[registerY];
  new shift = (target & 31) * 8;
  new index = (target >> 5) & 63;
  new where = (registerGl << 6) + index;
  new dri = where;
stage E1:
  new argument1 = XVR[where];
  new bytemask = 3;
  new targetbytes = bits2bytes(bytemask) << shift;
  CS.XMF = 1;
  new result1 = MEM.load(address, 0x3 & bytemask, variant, dri);
  result1 = ((result1 << shift) & targetbytes) | (argument1 & ~targetbytes);
stage CR:
  XVR[where] = result1;
\end{lstlisting}}
\newpage

\paragraph{Encoding} Format \textsf{LSU\_XPL64RBB.O} (Section~\ref{sec:format-LSU-XPL64RBB.O}), \verb|lsucode5 = 001|\\

\bitfields{ 1/P, 2/00, 2/-- --, 27/offset27, 1/1, 2/01, 3/011, 2/variant, 1/registerGl, 5/11111, 2/11, 3/001, 1/--, 6/registerY, 6/registerZ }


\paragraph{Syntax} \texttt{xplh}\emph{variant} \emph{registerGl}, \emph{registerY} = \emph{offset27}[\emph{registerZ}]

\paragraph{Description} The \emph{registerGl} is preloaded from the memory half word at the effective address.


\paragraph{Modifier(s)}
\noindent\begin{itemize}
\item variant (cf. variant in section \ref{par:modifier-variant} on page \pageref{par:modifier-variant})
\end{itemize}

\paragraph{Execution} {Reservation table \textsf{LSU.X} (Section~\ref{sec:reservation-LSU.X})}
~
{\begin{lstlisting}
stage ID:
  new offset = _SX_27(offset27);
  new variant = _ZX_2(variant);
stage RR:
  new base = GPR[registerZ];
  new target = GPR[registerY];
  new address = base + offset;
  new shift = (target & 31) * 8;
  new index = (target >> 5) & 63;
  new where = (registerGl << 6) + index;
  new dri = where;
stage E1:
  new argument1 = XVR[where];
  new bytemask = 3;
  new targetbytes = bits2bytes(bytemask) << shift;
  CS.XMF = 1;
  new result1 = MEM.load(address, 0x3 & bytemask, variant, dri);
  result1 = ((result1 << shift) & targetbytes) | (argument1 & ~targetbytes);
stage CR:
  XVR[where] = result1;
\end{lstlisting}}
\newpage

\paragraph{Encoding} Format \textsf{LSU\_XPL64RBB.D} (Section~\ref{sec:format-LSU-XPL64RBB.D}), \verb|lsucode5 = 001|\\

\bitfields{ 1/P, 2/00, 2/-- --, 27/extend27, 1/1, 2/00, 2/-- --, 27/offset27, 1/1, 2/01, 3/011, 2/variant, 1/registerGl, 5/11111, 2/11, 3/001, 1/--, 6/registerY, 6/registerZ }


\paragraph{Syntax} \texttt{xplh}\emph{variant} \emph{registerGl}, \emph{registerY} = \emph{extend27\_\discretionary{}{}{}offset27}[\emph{registerZ}]

\paragraph{Description} The \emph{registerGl} is preloaded from the memory half word at the effective address.


\paragraph{Modifier(s)}
\noindent\begin{itemize}
\item variant (cf. variant in section \ref{par:modifier-variant} on page \pageref{par:modifier-variant})
\end{itemize}

\paragraph{Execution} {Reservation table \textsf{LSU.Y} (Section~\ref{sec:reservation-LSU.Y})}
~
{\begin{lstlisting}
stage ID:
  new offset = _SX_54(extend27_offset27);
  new variant = _ZX_2(variant);
stage RR:
  new base = GPR[registerZ];
  new target = GPR[registerY];
  new address = base + offset;
  new shift = (target & 31) * 8;
  new index = (target >> 5) & 63;
  new where = (registerGl << 6) + index;
  new dri = where;
stage E1:
  new argument1 = XVR[where];
  new bytemask = 3;
  new targetbytes = bits2bytes(bytemask) << shift;
  CS.XMF = 1;
  new result1 = MEM.load(address, 0x3 & bytemask, variant, dri);
  result1 = ((result1 << shift) & targetbytes) | (argument1 & ~targetbytes);
stage CR:
  XVR[where] = result1;
\end{lstlisting}}
\newpage
\subsubsection[{\small XPLO: Extension Preload Octuple Word}]{XPLO\hfill Extension Preload Octuple Word}
\label{instr:XPLO}\index{xplo}
 
\paragraph{Encoding} Format \textsf{LSU\_XPL2RBB} (Section~\ref{sec:format-LSU-XPL2RBB}), \verb|lsucode5 = 101|\\

\bitfields{ 1/P, 2/01, 3/011, 2/variant, 5/registerGg, 1/0, 2/11, 3/101, 1/--, 6/registerY, 6/registerZ }


\paragraph{Syntax} \texttt{xplo}\emph{variant} \emph{registerGg}, \emph{registerY} = [\emph{registerZ}]

\paragraph{Description} The \emph{registerGg} is preloaded from the memory octuple word at the effective address.


\paragraph{Modifier(s)}
\noindent\begin{itemize}
\item variant (cf. variant in section \ref{par:modifier-variant} on page \pageref{par:modifier-variant})
\end{itemize}

\paragraph{Execution} {Reservation table \textsf{LSU} (Section~\ref{sec:reservation-LSU})}
~
{\begin{lstlisting}
stage ID:
  new variant = _ZX_2(variant);
stage RR:
  new address = GPR[registerZ];
  new target = GPR[registerY];
  new shift = (target & 31) * 8;
  new index = (target >> 5) & 1;
  new where = (registerGg << 1) + index;
  new dri = where;
stage E1:
  new argument1 = XVR[where];
  new bytemask = 4294967295;
  new targetbytes = bits2bytes(bytemask) << shift;
  CS.XMF = 1;
  new result1 = MEM.load(address, 0xFFFFFFFF & bytemask, variant, dri);
  result1 = ((result1 << shift) & targetbytes) | (argument1 & ~targetbytes);
stage CR:
  XVR[where] = result1;
\end{lstlisting}}
\newpage

\paragraph{Encoding} Format \textsf{LSU\_XPL2RBB.O} (Section~\ref{sec:format-LSU-XPL2RBB.O}), \verb|lsucode5 = 101|\\

\bitfields{ 1/P, 2/00, 2/-- --, 27/offset27, 1/1, 2/01, 3/011, 2/variant, 5/registerGg, 1/0, 2/11, 3/101, 1/--, 6/registerY, 6/registerZ }


\paragraph{Syntax} \texttt{xplo}\emph{variant} \emph{registerGg}, \emph{registerY} = \emph{offset27}[\emph{registerZ}]

\paragraph{Description} The \emph{registerGg} is preloaded from the memory octuple word at the effective address.


\paragraph{Modifier(s)}
\noindent\begin{itemize}
\item variant (cf. variant in section \ref{par:modifier-variant} on page \pageref{par:modifier-variant})
\end{itemize}

\paragraph{Execution} {Reservation table \textsf{LSU.X} (Section~\ref{sec:reservation-LSU.X})}
~
{\begin{lstlisting}
stage ID:
  new offset = _SX_27(offset27);
  new variant = _ZX_2(variant);
stage RR:
  new base = GPR[registerZ];
  new target = GPR[registerY];
  new address = base + offset;
  new shift = (target & 31) * 8;
  new index = (target >> 5) & 1;
  new where = (registerGg << 1) + index;
  new dri = where;
stage E1:
  new argument1 = XVR[where];
  new bytemask = 4294967295;
  new targetbytes = bits2bytes(bytemask) << shift;
  CS.XMF = 1;
  new result1 = MEM.load(address, 0xFFFFFFFF & bytemask, variant, dri);
  result1 = ((result1 << shift) & targetbytes) | (argument1 & ~targetbytes);
stage CR:
  XVR[where] = result1;
\end{lstlisting}}
\newpage

\paragraph{Encoding} Format \textsf{LSU\_XPL2RBB.D} (Section~\ref{sec:format-LSU-XPL2RBB.D}), \verb|lsucode5 = 101|\\

\bitfields{ 1/P, 2/00, 2/-- --, 27/extend27, 1/1, 2/00, 2/-- --, 27/offset27, 1/1, 2/01, 3/011, 2/variant, 5/registerGg, 1/0, 2/11, 3/101, 1/--, 6/registerY, 6/registerZ }


\paragraph{Syntax} \texttt{xplo}\emph{variant} \emph{registerGg}, \emph{registerY} = \emph{extend27\_\discretionary{}{}{}offset27}[\emph{registerZ}]

\paragraph{Description} The \emph{registerGg} is preloaded from the memory octuple word at the effective address.


\paragraph{Modifier(s)}
\noindent\begin{itemize}
\item variant (cf. variant in section \ref{par:modifier-variant} on page \pageref{par:modifier-variant})
\end{itemize}

\paragraph{Execution} {Reservation table \textsf{LSU.Y} (Section~\ref{sec:reservation-LSU.Y})}
~
{\begin{lstlisting}
stage ID:
  new offset = _SX_54(extend27_offset27);
  new variant = _ZX_2(variant);
stage RR:
  new base = GPR[registerZ];
  new target = GPR[registerY];
  new address = base + offset;
  new shift = (target & 31) * 8;
  new index = (target >> 5) & 1;
  new where = (registerGg << 1) + index;
  new dri = where;
stage E1:
  new argument1 = XVR[where];
  new bytemask = 4294967295;
  new targetbytes = bits2bytes(bytemask) << shift;
  CS.XMF = 1;
  new result1 = MEM.load(address, 0xFFFFFFFF & bytemask, variant, dri);
  result1 = ((result1 << shift) & targetbytes) | (argument1 & ~targetbytes);
stage CR:
  XVR[where] = result1;
\end{lstlisting}}
\newpage

\paragraph{Encoding} Format \textsf{LSU\_XPL4RBB} (Section~\ref{sec:format-LSU-XPL4RBB}), \verb|lsucode5 = 101|\\

\bitfields{ 1/P, 2/01, 3/011, 2/variant, 4/registerGh, 2/01, 2/11, 3/101, 1/--, 6/registerY, 6/registerZ }


\paragraph{Syntax} \texttt{xplo}\emph{variant} \emph{registerGh}, \emph{registerY} = [\emph{registerZ}]

\paragraph{Description} The \emph{registerGh} is preloaded from the memory octuple word at the effective address.


\paragraph{Modifier(s)}
\noindent\begin{itemize}
\item variant (cf. variant in section \ref{par:modifier-variant} on page \pageref{par:modifier-variant})
\end{itemize}

\paragraph{Execution} {Reservation table \textsf{LSU} (Section~\ref{sec:reservation-LSU})}
~
{\begin{lstlisting}
stage ID:
  new variant = _ZX_2(variant);
stage RR:
  new address = GPR[registerZ];
  new target = GPR[registerY];
  new shift = (target & 31) * 8;
  new index = (target >> 5) & 3;
  new where = (registerGh << 2) + index;
  new dri = where;
stage E1:
  new argument1 = XVR[where];
  new bytemask = 4294967295;
  new targetbytes = bits2bytes(bytemask) << shift;
  CS.XMF = 1;
  new result1 = MEM.load(address, 0xFFFFFFFF & bytemask, variant, dri);
  result1 = ((result1 << shift) & targetbytes) | (argument1 & ~targetbytes);
stage CR:
  XVR[where] = result1;
\end{lstlisting}}
\newpage

\paragraph{Encoding} Format \textsf{LSU\_XPL4RBB.O} (Section~\ref{sec:format-LSU-XPL4RBB.O}), \verb|lsucode5 = 101|\\

\bitfields{ 1/P, 2/00, 2/-- --, 27/offset27, 1/1, 2/01, 3/011, 2/variant, 4/registerGh, 2/01, 2/11, 3/101, 1/--, 6/registerY, 6/registerZ }


\paragraph{Syntax} \texttt{xplo}\emph{variant} \emph{registerGh}, \emph{registerY} = \emph{offset27}[\emph{registerZ}]

\paragraph{Description} The \emph{registerGh} is preloaded from the memory octuple word at the effective address.


\paragraph{Modifier(s)}
\noindent\begin{itemize}
\item variant (cf. variant in section \ref{par:modifier-variant} on page \pageref{par:modifier-variant})
\end{itemize}

\paragraph{Execution} {Reservation table \textsf{LSU.X} (Section~\ref{sec:reservation-LSU.X})}
~
{\begin{lstlisting}
stage ID:
  new offset = _SX_27(offset27);
  new variant = _ZX_2(variant);
stage RR:
  new base = GPR[registerZ];
  new target = GPR[registerY];
  new address = base + offset;
  new shift = (target & 31) * 8;
  new index = (target >> 5) & 3;
  new where = (registerGh << 2) + index;
  new dri = where;
stage E1:
  new argument1 = XVR[where];
  new bytemask = 4294967295;
  new targetbytes = bits2bytes(bytemask) << shift;
  CS.XMF = 1;
  new result1 = MEM.load(address, 0xFFFFFFFF & bytemask, variant, dri);
  result1 = ((result1 << shift) & targetbytes) | (argument1 & ~targetbytes);
stage CR:
  XVR[where] = result1;
\end{lstlisting}}
\newpage

\paragraph{Encoding} Format \textsf{LSU\_XPL4RBB.D} (Section~\ref{sec:format-LSU-XPL4RBB.D}), \verb|lsucode5 = 101|\\

\bitfields{ 1/P, 2/00, 2/-- --, 27/extend27, 1/1, 2/00, 2/-- --, 27/offset27, 1/1, 2/01, 3/011, 2/variant, 4/registerGh, 2/01, 2/11, 3/101, 1/--, 6/registerY, 6/registerZ }


\paragraph{Syntax} \texttt{xplo}\emph{variant} \emph{registerGh}, \emph{registerY} = \emph{extend27\_\discretionary{}{}{}offset27}[\emph{registerZ}]

\paragraph{Description} The \emph{registerGh} is preloaded from the memory octuple word at the effective address.


\paragraph{Modifier(s)}
\noindent\begin{itemize}
\item variant (cf. variant in section \ref{par:modifier-variant} on page \pageref{par:modifier-variant})
\end{itemize}

\paragraph{Execution} {Reservation table \textsf{LSU.Y} (Section~\ref{sec:reservation-LSU.Y})}
~
{\begin{lstlisting}
stage ID:
  new offset = _SX_54(extend27_offset27);
  new variant = _ZX_2(variant);
stage RR:
  new base = GPR[registerZ];
  new target = GPR[registerY];
  new address = base + offset;
  new shift = (target & 31) * 8;
  new index = (target >> 5) & 3;
  new where = (registerGh << 2) + index;
  new dri = where;
stage E1:
  new argument1 = XVR[where];
  new bytemask = 4294967295;
  new targetbytes = bits2bytes(bytemask) << shift;
  CS.XMF = 1;
  new result1 = MEM.load(address, 0xFFFFFFFF & bytemask, variant, dri);
  result1 = ((result1 << shift) & targetbytes) | (argument1 & ~targetbytes);
stage CR:
  XVR[where] = result1;
\end{lstlisting}}
\newpage

\paragraph{Encoding} Format \textsf{LSU\_XPL8RBB} (Section~\ref{sec:format-LSU-XPL8RBB}), \verb|lsucode5 = 101|\\

\bitfields{ 1/P, 2/01, 3/011, 2/variant, 3/registerGi, 3/011, 2/11, 3/101, 1/--, 6/registerY, 6/registerZ }


\paragraph{Syntax} \texttt{xplo}\emph{variant} \emph{registerGi}, \emph{registerY} = [\emph{registerZ}]

\paragraph{Description} The \emph{registerGi} is preloaded from the memory octuple word at the effective address.


\paragraph{Modifier(s)}
\noindent\begin{itemize}
\item variant (cf. variant in section \ref{par:modifier-variant} on page \pageref{par:modifier-variant})
\end{itemize}

\paragraph{Execution} {Reservation table \textsf{LSU} (Section~\ref{sec:reservation-LSU})}
~
{\begin{lstlisting}
stage ID:
  new variant = _ZX_2(variant);
stage RR:
  new address = GPR[registerZ];
  new target = GPR[registerY];
  new shift = (target & 31) * 8;
  new index = (target >> 5) & 7;
  new where = (registerGi << 3) + index;
  new dri = where;
stage E1:
  new argument1 = XVR[where];
  new bytemask = 4294967295;
  new targetbytes = bits2bytes(bytemask) << shift;
  CS.XMF = 1;
  new result1 = MEM.load(address, 0xFFFFFFFF & bytemask, variant, dri);
  result1 = ((result1 << shift) & targetbytes) | (argument1 & ~targetbytes);
stage CR:
  XVR[where] = result1;
\end{lstlisting}}
\newpage

\paragraph{Encoding} Format \textsf{LSU\_XPL8RBB.O} (Section~\ref{sec:format-LSU-XPL8RBB.O}), \verb|lsucode5 = 101|\\

\bitfields{ 1/P, 2/00, 2/-- --, 27/offset27, 1/1, 2/01, 3/011, 2/variant, 3/registerGi, 3/011, 2/11, 3/101, 1/--, 6/registerY, 6/registerZ }


\paragraph{Syntax} \texttt{xplo}\emph{variant} \emph{registerGi}, \emph{registerY} = \emph{offset27}[\emph{registerZ}]

\paragraph{Description} The \emph{registerGi} is preloaded from the memory octuple word at the effective address.


\paragraph{Modifier(s)}
\noindent\begin{itemize}
\item variant (cf. variant in section \ref{par:modifier-variant} on page \pageref{par:modifier-variant})
\end{itemize}

\paragraph{Execution} {Reservation table \textsf{LSU.X} (Section~\ref{sec:reservation-LSU.X})}
~
{\begin{lstlisting}
stage ID:
  new offset = _SX_27(offset27);
  new variant = _ZX_2(variant);
stage RR:
  new base = GPR[registerZ];
  new target = GPR[registerY];
  new address = base + offset;
  new shift = (target & 31) * 8;
  new index = (target >> 5) & 7;
  new where = (registerGi << 3) + index;
  new dri = where;
stage E1:
  new argument1 = XVR[where];
  new bytemask = 4294967295;
  new targetbytes = bits2bytes(bytemask) << shift;
  CS.XMF = 1;
  new result1 = MEM.load(address, 0xFFFFFFFF & bytemask, variant, dri);
  result1 = ((result1 << shift) & targetbytes) | (argument1 & ~targetbytes);
stage CR:
  XVR[where] = result1;
\end{lstlisting}}
\newpage

\paragraph{Encoding} Format \textsf{LSU\_XPL8RBB.D} (Section~\ref{sec:format-LSU-XPL8RBB.D}), \verb|lsucode5 = 101|\\

\bitfields{ 1/P, 2/00, 2/-- --, 27/extend27, 1/1, 2/00, 2/-- --, 27/offset27, 1/1, 2/01, 3/011, 2/variant, 3/registerGi, 3/011, 2/11, 3/101, 1/--, 6/registerY, 6/registerZ }


\paragraph{Syntax} \texttt{xplo}\emph{variant} \emph{registerGi}, \emph{registerY} = \emph{extend27\_\discretionary{}{}{}offset27}[\emph{registerZ}]

\paragraph{Description} The \emph{registerGi} is preloaded from the memory octuple word at the effective address.


\paragraph{Modifier(s)}
\noindent\begin{itemize}
\item variant (cf. variant in section \ref{par:modifier-variant} on page \pageref{par:modifier-variant})
\end{itemize}

\paragraph{Execution} {Reservation table \textsf{LSU.Y} (Section~\ref{sec:reservation-LSU.Y})}
~
{\begin{lstlisting}
stage ID:
  new offset = _SX_54(extend27_offset27);
  new variant = _ZX_2(variant);
stage RR:
  new base = GPR[registerZ];
  new target = GPR[registerY];
  new address = base + offset;
  new shift = (target & 31) * 8;
  new index = (target >> 5) & 7;
  new where = (registerGi << 3) + index;
  new dri = where;
stage E1:
  new argument1 = XVR[where];
  new bytemask = 4294967295;
  new targetbytes = bits2bytes(bytemask) << shift;
  CS.XMF = 1;
  new result1 = MEM.load(address, 0xFFFFFFFF & bytemask, variant, dri);
  result1 = ((result1 << shift) & targetbytes) | (argument1 & ~targetbytes);
stage CR:
  XVR[where] = result1;
\end{lstlisting}}
\newpage

\paragraph{Encoding} Format \textsf{LSU\_XPL16RBB} (Section~\ref{sec:format-LSU-XPL16RBB}), \verb|lsucode5 = 101|\\

\bitfields{ 1/P, 2/01, 3/011, 2/variant, 2/registerGj, 4/0111, 2/11, 3/101, 1/--, 6/registerY, 6/registerZ }


\paragraph{Syntax} \texttt{xplo}\emph{variant} \emph{registerGj}, \emph{registerY} = [\emph{registerZ}]

\paragraph{Description} The \emph{registerGj} is preloaded from the memory octuple word at the effective address.


\paragraph{Modifier(s)}
\noindent\begin{itemize}
\item variant (cf. variant in section \ref{par:modifier-variant} on page \pageref{par:modifier-variant})
\end{itemize}

\paragraph{Execution} {Reservation table \textsf{LSU} (Section~\ref{sec:reservation-LSU})}
~
{\begin{lstlisting}
stage ID:
  new variant = _ZX_2(variant);
stage RR:
  new address = GPR[registerZ];
  new target = GPR[registerY];
  new shift = (target & 31) * 8;
  new index = (target >> 5) & 15;
  new where = (registerGj << 4) + index;
  new dri = where;
stage E1:
  new argument1 = XVR[where];
  new bytemask = 4294967295;
  new targetbytes = bits2bytes(bytemask) << shift;
  CS.XMF = 1;
  new result1 = MEM.load(address, 0xFFFFFFFF & bytemask, variant, dri);
  result1 = ((result1 << shift) & targetbytes) | (argument1 & ~targetbytes);
stage CR:
  XVR[where] = result1;
\end{lstlisting}}
\newpage

\paragraph{Encoding} Format \textsf{LSU\_XPL16RBB.O} (Section~\ref{sec:format-LSU-XPL16RBB.O}), \verb|lsucode5 = 101|\\

\bitfields{ 1/P, 2/00, 2/-- --, 27/offset27, 1/1, 2/01, 3/011, 2/variant, 2/registerGj, 4/0111, 2/11, 3/101, 1/--, 6/registerY, 6/registerZ }


\paragraph{Syntax} \texttt{xplo}\emph{variant} \emph{registerGj}, \emph{registerY} = \emph{offset27}[\emph{registerZ}]

\paragraph{Description} The \emph{registerGj} is preloaded from the memory octuple word at the effective address.


\paragraph{Modifier(s)}
\noindent\begin{itemize}
\item variant (cf. variant in section \ref{par:modifier-variant} on page \pageref{par:modifier-variant})
\end{itemize}

\paragraph{Execution} {Reservation table \textsf{LSU.X} (Section~\ref{sec:reservation-LSU.X})}
~
{\begin{lstlisting}
stage ID:
  new offset = _SX_27(offset27);
  new variant = _ZX_2(variant);
stage RR:
  new base = GPR[registerZ];
  new target = GPR[registerY];
  new address = base + offset;
  new shift = (target & 31) * 8;
  new index = (target >> 5) & 15;
  new where = (registerGj << 4) + index;
  new dri = where;
stage E1:
  new argument1 = XVR[where];
  new bytemask = 4294967295;
  new targetbytes = bits2bytes(bytemask) << shift;
  CS.XMF = 1;
  new result1 = MEM.load(address, 0xFFFFFFFF & bytemask, variant, dri);
  result1 = ((result1 << shift) & targetbytes) | (argument1 & ~targetbytes);
stage CR:
  XVR[where] = result1;
\end{lstlisting}}
\newpage

\paragraph{Encoding} Format \textsf{LSU\_XPL16RBB.D} (Section~\ref{sec:format-LSU-XPL16RBB.D}), \verb|lsucode5 = 101|\\

\bitfields{ 1/P, 2/00, 2/-- --, 27/extend27, 1/1, 2/00, 2/-- --, 27/offset27, 1/1, 2/01, 3/011, 2/variant, 2/registerGj, 4/0111, 2/11, 3/101, 1/--, 6/registerY, 6/registerZ }


\paragraph{Syntax} \texttt{xplo}\emph{variant} \emph{registerGj}, \emph{registerY} = \emph{extend27\_\discretionary{}{}{}offset27}[\emph{registerZ}]

\paragraph{Description} The \emph{registerGj} is preloaded from the memory octuple word at the effective address.


\paragraph{Modifier(s)}
\noindent\begin{itemize}
\item variant (cf. variant in section \ref{par:modifier-variant} on page \pageref{par:modifier-variant})
\end{itemize}

\paragraph{Execution} {Reservation table \textsf{LSU.Y} (Section~\ref{sec:reservation-LSU.Y})}
~
{\begin{lstlisting}
stage ID:
  new offset = _SX_54(extend27_offset27);
  new variant = _ZX_2(variant);
stage RR:
  new base = GPR[registerZ];
  new target = GPR[registerY];
  new address = base + offset;
  new shift = (target & 31) * 8;
  new index = (target >> 5) & 15;
  new where = (registerGj << 4) + index;
  new dri = where;
stage E1:
  new argument1 = XVR[where];
  new bytemask = 4294967295;
  new targetbytes = bits2bytes(bytemask) << shift;
  CS.XMF = 1;
  new result1 = MEM.load(address, 0xFFFFFFFF & bytemask, variant, dri);
  result1 = ((result1 << shift) & targetbytes) | (argument1 & ~targetbytes);
stage CR:
  XVR[where] = result1;
\end{lstlisting}}
\newpage

\paragraph{Encoding} Format \textsf{LSU\_XPL32RBB} (Section~\ref{sec:format-LSU-XPL32RBB}), \verb|lsucode5 = 101|\\

\bitfields{ 1/P, 2/01, 3/011, 2/variant, 1/registerGk, 5/01111, 2/11, 3/101, 1/--, 6/registerY, 6/registerZ }


\paragraph{Syntax} \texttt{xplo}\emph{variant} \emph{registerGk}, \emph{registerY} = [\emph{registerZ}]

\paragraph{Description} The \emph{registerGk} is preloaded from the memory octuple word at the effective address.


\paragraph{Modifier(s)}
\noindent\begin{itemize}
\item variant (cf. variant in section \ref{par:modifier-variant} on page \pageref{par:modifier-variant})
\end{itemize}

\paragraph{Execution} {Reservation table \textsf{LSU} (Section~\ref{sec:reservation-LSU})}
~
{\begin{lstlisting}
stage ID:
  new variant = _ZX_2(variant);
stage RR:
  new address = GPR[registerZ];
  new target = GPR[registerY];
  new shift = (target & 31) * 8;
  new index = (target >> 5) & 31;
  new where = (registerGk << 5) + index;
  new dri = where;
stage E1:
  new argument1 = XVR[where];
  new bytemask = 4294967295;
  new targetbytes = bits2bytes(bytemask) << shift;
  CS.XMF = 1;
  new result1 = MEM.load(address, 0xFFFFFFFF & bytemask, variant, dri);
  result1 = ((result1 << shift) & targetbytes) | (argument1 & ~targetbytes);
stage CR:
  XVR[where] = result1;
\end{lstlisting}}
\newpage

\paragraph{Encoding} Format \textsf{LSU\_XPL32RBB.O} (Section~\ref{sec:format-LSU-XPL32RBB.O}), \verb|lsucode5 = 101|\\

\bitfields{ 1/P, 2/00, 2/-- --, 27/offset27, 1/1, 2/01, 3/011, 2/variant, 1/registerGk, 5/01111, 2/11, 3/101, 1/--, 6/registerY, 6/registerZ }


\paragraph{Syntax} \texttt{xplo}\emph{variant} \emph{registerGk}, \emph{registerY} = \emph{offset27}[\emph{registerZ}]

\paragraph{Description} The \emph{registerGk} is preloaded from the memory octuple word at the effective address.


\paragraph{Modifier(s)}
\noindent\begin{itemize}
\item variant (cf. variant in section \ref{par:modifier-variant} on page \pageref{par:modifier-variant})
\end{itemize}

\paragraph{Execution} {Reservation table \textsf{LSU.X} (Section~\ref{sec:reservation-LSU.X})}
~
{\begin{lstlisting}
stage ID:
  new offset = _SX_27(offset27);
  new variant = _ZX_2(variant);
stage RR:
  new base = GPR[registerZ];
  new target = GPR[registerY];
  new address = base + offset;
  new shift = (target & 31) * 8;
  new index = (target >> 5) & 31;
  new where = (registerGk << 5) + index;
  new dri = where;
stage E1:
  new argument1 = XVR[where];
  new bytemask = 4294967295;
  new targetbytes = bits2bytes(bytemask) << shift;
  CS.XMF = 1;
  new result1 = MEM.load(address, 0xFFFFFFFF & bytemask, variant, dri);
  result1 = ((result1 << shift) & targetbytes) | (argument1 & ~targetbytes);
stage CR:
  XVR[where] = result1;
\end{lstlisting}}
\newpage

\paragraph{Encoding} Format \textsf{LSU\_XPL32RBB.D} (Section~\ref{sec:format-LSU-XPL32RBB.D}), \verb|lsucode5 = 101|\\

\bitfields{ 1/P, 2/00, 2/-- --, 27/extend27, 1/1, 2/00, 2/-- --, 27/offset27, 1/1, 2/01, 3/011, 2/variant, 1/registerGk, 5/01111, 2/11, 3/101, 1/--, 6/registerY, 6/registerZ }


\paragraph{Syntax} \texttt{xplo}\emph{variant} \emph{registerGk}, \emph{registerY} = \emph{extend27\_\discretionary{}{}{}offset27}[\emph{registerZ}]

\paragraph{Description} The \emph{registerGk} is preloaded from the memory octuple word at the effective address.


\paragraph{Modifier(s)}
\noindent\begin{itemize}
\item variant (cf. variant in section \ref{par:modifier-variant} on page \pageref{par:modifier-variant})
\end{itemize}

\paragraph{Execution} {Reservation table \textsf{LSU.Y} (Section~\ref{sec:reservation-LSU.Y})}
~
{\begin{lstlisting}
stage ID:
  new offset = _SX_54(extend27_offset27);
  new variant = _ZX_2(variant);
stage RR:
  new base = GPR[registerZ];
  new target = GPR[registerY];
  new address = base + offset;
  new shift = (target & 31) * 8;
  new index = (target >> 5) & 31;
  new where = (registerGk << 5) + index;
  new dri = where;
stage E1:
  new argument1 = XVR[where];
  new bytemask = 4294967295;
  new targetbytes = bits2bytes(bytemask) << shift;
  CS.XMF = 1;
  new result1 = MEM.load(address, 0xFFFFFFFF & bytemask, variant, dri);
  result1 = ((result1 << shift) & targetbytes) | (argument1 & ~targetbytes);
stage CR:
  XVR[where] = result1;
\end{lstlisting}}
\newpage

\paragraph{Encoding} Format \textsf{LSU\_XPL64RBB} (Section~\ref{sec:format-LSU-XPL64RBB}), \verb|lsucode5 = 101|\\

\bitfields{ 1/P, 2/01, 3/011, 2/variant, 1/registerGl, 5/11111, 2/11, 3/101, 1/--, 6/registerY, 6/registerZ }


\paragraph{Syntax} \texttt{xplo}\emph{variant} \emph{registerGl}, \emph{registerY} = [\emph{registerZ}]

\paragraph{Description} The \emph{registerGl} is preloaded from the memory octuple word at the effective address.


\paragraph{Modifier(s)}
\noindent\begin{itemize}
\item variant (cf. variant in section \ref{par:modifier-variant} on page \pageref{par:modifier-variant})
\end{itemize}

\paragraph{Execution} {Reservation table \textsf{LSU} (Section~\ref{sec:reservation-LSU})}
~
{\begin{lstlisting}
stage ID:
  new variant = _ZX_2(variant);
stage RR:
  new address = GPR[registerZ];
  new target = GPR[registerY];
  new shift = (target & 31) * 8;
  new index = (target >> 5) & 63;
  new where = (registerGl << 6) + index;
  new dri = where;
stage E1:
  new argument1 = XVR[where];
  new bytemask = 4294967295;
  new targetbytes = bits2bytes(bytemask) << shift;
  CS.XMF = 1;
  new result1 = MEM.load(address, 0xFFFFFFFF & bytemask, variant, dri);
  result1 = ((result1 << shift) & targetbytes) | (argument1 & ~targetbytes);
stage CR:
  XVR[where] = result1;
\end{lstlisting}}
\newpage

\paragraph{Encoding} Format \textsf{LSU\_XPL64RBB.O} (Section~\ref{sec:format-LSU-XPL64RBB.O}), \verb|lsucode5 = 101|\\

\bitfields{ 1/P, 2/00, 2/-- --, 27/offset27, 1/1, 2/01, 3/011, 2/variant, 1/registerGl, 5/11111, 2/11, 3/101, 1/--, 6/registerY, 6/registerZ }


\paragraph{Syntax} \texttt{xplo}\emph{variant} \emph{registerGl}, \emph{registerY} = \emph{offset27}[\emph{registerZ}]

\paragraph{Description} The \emph{registerGl} is preloaded from the memory octuple word at the effective address.


\paragraph{Modifier(s)}
\noindent\begin{itemize}
\item variant (cf. variant in section \ref{par:modifier-variant} on page \pageref{par:modifier-variant})
\end{itemize}

\paragraph{Execution} {Reservation table \textsf{LSU.X} (Section~\ref{sec:reservation-LSU.X})}
~
{\begin{lstlisting}
stage ID:
  new offset = _SX_27(offset27);
  new variant = _ZX_2(variant);
stage RR:
  new base = GPR[registerZ];
  new target = GPR[registerY];
  new address = base + offset;
  new shift = (target & 31) * 8;
  new index = (target >> 5) & 63;
  new where = (registerGl << 6) + index;
  new dri = where;
stage E1:
  new argument1 = XVR[where];
  new bytemask = 4294967295;
  new targetbytes = bits2bytes(bytemask) << shift;
  CS.XMF = 1;
  new result1 = MEM.load(address, 0xFFFFFFFF & bytemask, variant, dri);
  result1 = ((result1 << shift) & targetbytes) | (argument1 & ~targetbytes);
stage CR:
  XVR[where] = result1;
\end{lstlisting}}
\newpage

\paragraph{Encoding} Format \textsf{LSU\_XPL64RBB.D} (Section~\ref{sec:format-LSU-XPL64RBB.D}), \verb|lsucode5 = 101|\\

\bitfields{ 1/P, 2/00, 2/-- --, 27/extend27, 1/1, 2/00, 2/-- --, 27/offset27, 1/1, 2/01, 3/011, 2/variant, 1/registerGl, 5/11111, 2/11, 3/101, 1/--, 6/registerY, 6/registerZ }


\paragraph{Syntax} \texttt{xplo}\emph{variant} \emph{registerGl}, \emph{registerY} = \emph{extend27\_\discretionary{}{}{}offset27}[\emph{registerZ}]

\paragraph{Description} The \emph{registerGl} is preloaded from the memory octuple word at the effective address.


\paragraph{Modifier(s)}
\noindent\begin{itemize}
\item variant (cf. variant in section \ref{par:modifier-variant} on page \pageref{par:modifier-variant})
\end{itemize}

\paragraph{Execution} {Reservation table \textsf{LSU.Y} (Section~\ref{sec:reservation-LSU.Y})}
~
{\begin{lstlisting}
stage ID:
  new offset = _SX_54(extend27_offset27);
  new variant = _ZX_2(variant);
stage RR:
  new base = GPR[registerZ];
  new target = GPR[registerY];
  new address = base + offset;
  new shift = (target & 31) * 8;
  new index = (target >> 5) & 63;
  new where = (registerGl << 6) + index;
  new dri = where;
stage E1:
  new argument1 = XVR[where];
  new bytemask = 4294967295;
  new targetbytes = bits2bytes(bytemask) << shift;
  CS.XMF = 1;
  new result1 = MEM.load(address, 0xFFFFFFFF & bytemask, variant, dri);
  result1 = ((result1 << shift) & targetbytes) | (argument1 & ~targetbytes);
stage CR:
  XVR[where] = result1;
\end{lstlisting}}
\newpage
\subsubsection[{\small XPLQ: Extension Preload Quadruple Word}]{XPLQ\hfill Extension Preload Quadruple Word}
\label{instr:XPLQ}\index{xplq}
 
\paragraph{Encoding} Format \textsf{LSU\_XPL2RBB} (Section~\ref{sec:format-LSU-XPL2RBB}), \verb|lsucode5 = 100|\\

\bitfields{ 1/P, 2/01, 3/011, 2/variant, 5/registerGg, 1/0, 2/11, 3/100, 1/--, 6/registerY, 6/registerZ }


\paragraph{Syntax} \texttt{xplq}\emph{variant} \emph{registerGg}, \emph{registerY} = [\emph{registerZ}]

\paragraph{Description} The \emph{registerGg} is preloaded from the memory quadruple word at the effective address.


\paragraph{Modifier(s)}
\noindent\begin{itemize}
\item variant (cf. variant in section \ref{par:modifier-variant} on page \pageref{par:modifier-variant})
\end{itemize}

\paragraph{Execution} {Reservation table \textsf{LSU} (Section~\ref{sec:reservation-LSU})}
~
{\begin{lstlisting}
stage ID:
  new variant = _ZX_2(variant);
stage RR:
  new address = GPR[registerZ];
  new target = GPR[registerY];
  new shift = (target & 31) * 8;
  new index = (target >> 5) & 1;
  new where = (registerGg << 1) + index;
  new dri = where;
stage E1:
  new argument1 = XVR[where];
  new bytemask = 65535;
  new targetbytes = bits2bytes(bytemask) << shift;
  CS.XMF = 1;
  new result1 = MEM.load(address, 0xFFFF & bytemask, variant, dri);
  result1 = ((result1 << shift) & targetbytes) | (argument1 & ~targetbytes);
stage CR:
  XVR[where] = result1;
\end{lstlisting}}
\newpage

\paragraph{Encoding} Format \textsf{LSU\_XPL2RBB.O} (Section~\ref{sec:format-LSU-XPL2RBB.O}), \verb|lsucode5 = 100|\\

\bitfields{ 1/P, 2/00, 2/-- --, 27/offset27, 1/1, 2/01, 3/011, 2/variant, 5/registerGg, 1/0, 2/11, 3/100, 1/--, 6/registerY, 6/registerZ }


\paragraph{Syntax} \texttt{xplq}\emph{variant} \emph{registerGg}, \emph{registerY} = \emph{offset27}[\emph{registerZ}]

\paragraph{Description} The \emph{registerGg} is preloaded from the memory quadruple word at the effective address.


\paragraph{Modifier(s)}
\noindent\begin{itemize}
\item variant (cf. variant in section \ref{par:modifier-variant} on page \pageref{par:modifier-variant})
\end{itemize}

\paragraph{Execution} {Reservation table \textsf{LSU.X} (Section~\ref{sec:reservation-LSU.X})}
~
{\begin{lstlisting}
stage ID:
  new offset = _SX_27(offset27);
  new variant = _ZX_2(variant);
stage RR:
  new base = GPR[registerZ];
  new target = GPR[registerY];
  new address = base + offset;
  new shift = (target & 31) * 8;
  new index = (target >> 5) & 1;
  new where = (registerGg << 1) + index;
  new dri = where;
stage E1:
  new argument1 = XVR[where];
  new bytemask = 65535;
  new targetbytes = bits2bytes(bytemask) << shift;
  CS.XMF = 1;
  new result1 = MEM.load(address, 0xFFFF & bytemask, variant, dri);
  result1 = ((result1 << shift) & targetbytes) | (argument1 & ~targetbytes);
stage CR:
  XVR[where] = result1;
\end{lstlisting}}
\newpage

\paragraph{Encoding} Format \textsf{LSU\_XPL2RBB.D} (Section~\ref{sec:format-LSU-XPL2RBB.D}), \verb|lsucode5 = 100|\\

\bitfields{ 1/P, 2/00, 2/-- --, 27/extend27, 1/1, 2/00, 2/-- --, 27/offset27, 1/1, 2/01, 3/011, 2/variant, 5/registerGg, 1/0, 2/11, 3/100, 1/--, 6/registerY, 6/registerZ }


\paragraph{Syntax} \texttt{xplq}\emph{variant} \emph{registerGg}, \emph{registerY} = \emph{extend27\_\discretionary{}{}{}offset27}[\emph{registerZ}]

\paragraph{Description} The \emph{registerGg} is preloaded from the memory quadruple word at the effective address.


\paragraph{Modifier(s)}
\noindent\begin{itemize}
\item variant (cf. variant in section \ref{par:modifier-variant} on page \pageref{par:modifier-variant})
\end{itemize}

\paragraph{Execution} {Reservation table \textsf{LSU.Y} (Section~\ref{sec:reservation-LSU.Y})}
~
{\begin{lstlisting}
stage ID:
  new offset = _SX_54(extend27_offset27);
  new variant = _ZX_2(variant);
stage RR:
  new base = GPR[registerZ];
  new target = GPR[registerY];
  new address = base + offset;
  new shift = (target & 31) * 8;
  new index = (target >> 5) & 1;
  new where = (registerGg << 1) + index;
  new dri = where;
stage E1:
  new argument1 = XVR[where];
  new bytemask = 65535;
  new targetbytes = bits2bytes(bytemask) << shift;
  CS.XMF = 1;
  new result1 = MEM.load(address, 0xFFFF & bytemask, variant, dri);
  result1 = ((result1 << shift) & targetbytes) | (argument1 & ~targetbytes);
stage CR:
  XVR[where] = result1;
\end{lstlisting}}
\newpage

\paragraph{Encoding} Format \textsf{LSU\_XPL4RBB} (Section~\ref{sec:format-LSU-XPL4RBB}), \verb|lsucode5 = 100|\\

\bitfields{ 1/P, 2/01, 3/011, 2/variant, 4/registerGh, 2/01, 2/11, 3/100, 1/--, 6/registerY, 6/registerZ }


\paragraph{Syntax} \texttt{xplq}\emph{variant} \emph{registerGh}, \emph{registerY} = [\emph{registerZ}]

\paragraph{Description} The \emph{registerGh} is preloaded from the memory quadruple word at the effective address.


\paragraph{Modifier(s)}
\noindent\begin{itemize}
\item variant (cf. variant in section \ref{par:modifier-variant} on page \pageref{par:modifier-variant})
\end{itemize}

\paragraph{Execution} {Reservation table \textsf{LSU} (Section~\ref{sec:reservation-LSU})}
~
{\begin{lstlisting}
stage ID:
  new variant = _ZX_2(variant);
stage RR:
  new address = GPR[registerZ];
  new target = GPR[registerY];
  new shift = (target & 31) * 8;
  new index = (target >> 5) & 3;
  new where = (registerGh << 2) + index;
  new dri = where;
stage E1:
  new argument1 = XVR[where];
  new bytemask = 65535;
  new targetbytes = bits2bytes(bytemask) << shift;
  CS.XMF = 1;
  new result1 = MEM.load(address, 0xFFFF & bytemask, variant, dri);
  result1 = ((result1 << shift) & targetbytes) | (argument1 & ~targetbytes);
stage CR:
  XVR[where] = result1;
\end{lstlisting}}
\newpage

\paragraph{Encoding} Format \textsf{LSU\_XPL4RBB.O} (Section~\ref{sec:format-LSU-XPL4RBB.O}), \verb|lsucode5 = 100|\\

\bitfields{ 1/P, 2/00, 2/-- --, 27/offset27, 1/1, 2/01, 3/011, 2/variant, 4/registerGh, 2/01, 2/11, 3/100, 1/--, 6/registerY, 6/registerZ }


\paragraph{Syntax} \texttt{xplq}\emph{variant} \emph{registerGh}, \emph{registerY} = \emph{offset27}[\emph{registerZ}]

\paragraph{Description} The \emph{registerGh} is preloaded from the memory quadruple word at the effective address.


\paragraph{Modifier(s)}
\noindent\begin{itemize}
\item variant (cf. variant in section \ref{par:modifier-variant} on page \pageref{par:modifier-variant})
\end{itemize}

\paragraph{Execution} {Reservation table \textsf{LSU.X} (Section~\ref{sec:reservation-LSU.X})}
~
{\begin{lstlisting}
stage ID:
  new offset = _SX_27(offset27);
  new variant = _ZX_2(variant);
stage RR:
  new base = GPR[registerZ];
  new target = GPR[registerY];
  new address = base + offset;
  new shift = (target & 31) * 8;
  new index = (target >> 5) & 3;
  new where = (registerGh << 2) + index;
  new dri = where;
stage E1:
  new argument1 = XVR[where];
  new bytemask = 65535;
  new targetbytes = bits2bytes(bytemask) << shift;
  CS.XMF = 1;
  new result1 = MEM.load(address, 0xFFFF & bytemask, variant, dri);
  result1 = ((result1 << shift) & targetbytes) | (argument1 & ~targetbytes);
stage CR:
  XVR[where] = result1;
\end{lstlisting}}
\newpage

\paragraph{Encoding} Format \textsf{LSU\_XPL4RBB.D} (Section~\ref{sec:format-LSU-XPL4RBB.D}), \verb|lsucode5 = 100|\\

\bitfields{ 1/P, 2/00, 2/-- --, 27/extend27, 1/1, 2/00, 2/-- --, 27/offset27, 1/1, 2/01, 3/011, 2/variant, 4/registerGh, 2/01, 2/11, 3/100, 1/--, 6/registerY, 6/registerZ }


\paragraph{Syntax} \texttt{xplq}\emph{variant} \emph{registerGh}, \emph{registerY} = \emph{extend27\_\discretionary{}{}{}offset27}[\emph{registerZ}]

\paragraph{Description} The \emph{registerGh} is preloaded from the memory quadruple word at the effective address.


\paragraph{Modifier(s)}
\noindent\begin{itemize}
\item variant (cf. variant in section \ref{par:modifier-variant} on page \pageref{par:modifier-variant})
\end{itemize}

\paragraph{Execution} {Reservation table \textsf{LSU.Y} (Section~\ref{sec:reservation-LSU.Y})}
~
{\begin{lstlisting}
stage ID:
  new offset = _SX_54(extend27_offset27);
  new variant = _ZX_2(variant);
stage RR:
  new base = GPR[registerZ];
  new target = GPR[registerY];
  new address = base + offset;
  new shift = (target & 31) * 8;
  new index = (target >> 5) & 3;
  new where = (registerGh << 2) + index;
  new dri = where;
stage E1:
  new argument1 = XVR[where];
  new bytemask = 65535;
  new targetbytes = bits2bytes(bytemask) << shift;
  CS.XMF = 1;
  new result1 = MEM.load(address, 0xFFFF & bytemask, variant, dri);
  result1 = ((result1 << shift) & targetbytes) | (argument1 & ~targetbytes);
stage CR:
  XVR[where] = result1;
\end{lstlisting}}
\newpage

\paragraph{Encoding} Format \textsf{LSU\_XPL8RBB} (Section~\ref{sec:format-LSU-XPL8RBB}), \verb|lsucode5 = 100|\\

\bitfields{ 1/P, 2/01, 3/011, 2/variant, 3/registerGi, 3/011, 2/11, 3/100, 1/--, 6/registerY, 6/registerZ }


\paragraph{Syntax} \texttt{xplq}\emph{variant} \emph{registerGi}, \emph{registerY} = [\emph{registerZ}]

\paragraph{Description} The \emph{registerGi} is preloaded from the memory quadruple word at the effective address.


\paragraph{Modifier(s)}
\noindent\begin{itemize}
\item variant (cf. variant in section \ref{par:modifier-variant} on page \pageref{par:modifier-variant})
\end{itemize}

\paragraph{Execution} {Reservation table \textsf{LSU} (Section~\ref{sec:reservation-LSU})}
~
{\begin{lstlisting}
stage ID:
  new variant = _ZX_2(variant);
stage RR:
  new address = GPR[registerZ];
  new target = GPR[registerY];
  new shift = (target & 31) * 8;
  new index = (target >> 5) & 7;
  new where = (registerGi << 3) + index;
  new dri = where;
stage E1:
  new argument1 = XVR[where];
  new bytemask = 65535;
  new targetbytes = bits2bytes(bytemask) << shift;
  CS.XMF = 1;
  new result1 = MEM.load(address, 0xFFFF & bytemask, variant, dri);
  result1 = ((result1 << shift) & targetbytes) | (argument1 & ~targetbytes);
stage CR:
  XVR[where] = result1;
\end{lstlisting}}
\newpage

\paragraph{Encoding} Format \textsf{LSU\_XPL8RBB.O} (Section~\ref{sec:format-LSU-XPL8RBB.O}), \verb|lsucode5 = 100|\\

\bitfields{ 1/P, 2/00, 2/-- --, 27/offset27, 1/1, 2/01, 3/011, 2/variant, 3/registerGi, 3/011, 2/11, 3/100, 1/--, 6/registerY, 6/registerZ }


\paragraph{Syntax} \texttt{xplq}\emph{variant} \emph{registerGi}, \emph{registerY} = \emph{offset27}[\emph{registerZ}]

\paragraph{Description} The \emph{registerGi} is preloaded from the memory quadruple word at the effective address.


\paragraph{Modifier(s)}
\noindent\begin{itemize}
\item variant (cf. variant in section \ref{par:modifier-variant} on page \pageref{par:modifier-variant})
\end{itemize}

\paragraph{Execution} {Reservation table \textsf{LSU.X} (Section~\ref{sec:reservation-LSU.X})}
~
{\begin{lstlisting}
stage ID:
  new offset = _SX_27(offset27);
  new variant = _ZX_2(variant);
stage RR:
  new base = GPR[registerZ];
  new target = GPR[registerY];
  new address = base + offset;
  new shift = (target & 31) * 8;
  new index = (target >> 5) & 7;
  new where = (registerGi << 3) + index;
  new dri = where;
stage E1:
  new argument1 = XVR[where];
  new bytemask = 65535;
  new targetbytes = bits2bytes(bytemask) << shift;
  CS.XMF = 1;
  new result1 = MEM.load(address, 0xFFFF & bytemask, variant, dri);
  result1 = ((result1 << shift) & targetbytes) | (argument1 & ~targetbytes);
stage CR:
  XVR[where] = result1;
\end{lstlisting}}
\newpage

\paragraph{Encoding} Format \textsf{LSU\_XPL8RBB.D} (Section~\ref{sec:format-LSU-XPL8RBB.D}), \verb|lsucode5 = 100|\\

\bitfields{ 1/P, 2/00, 2/-- --, 27/extend27, 1/1, 2/00, 2/-- --, 27/offset27, 1/1, 2/01, 3/011, 2/variant, 3/registerGi, 3/011, 2/11, 3/100, 1/--, 6/registerY, 6/registerZ }


\paragraph{Syntax} \texttt{xplq}\emph{variant} \emph{registerGi}, \emph{registerY} = \emph{extend27\_\discretionary{}{}{}offset27}[\emph{registerZ}]

\paragraph{Description} The \emph{registerGi} is preloaded from the memory quadruple word at the effective address.


\paragraph{Modifier(s)}
\noindent\begin{itemize}
\item variant (cf. variant in section \ref{par:modifier-variant} on page \pageref{par:modifier-variant})
\end{itemize}

\paragraph{Execution} {Reservation table \textsf{LSU.Y} (Section~\ref{sec:reservation-LSU.Y})}
~
{\begin{lstlisting}
stage ID:
  new offset = _SX_54(extend27_offset27);
  new variant = _ZX_2(variant);
stage RR:
  new base = GPR[registerZ];
  new target = GPR[registerY];
  new address = base + offset;
  new shift = (target & 31) * 8;
  new index = (target >> 5) & 7;
  new where = (registerGi << 3) + index;
  new dri = where;
stage E1:
  new argument1 = XVR[where];
  new bytemask = 65535;
  new targetbytes = bits2bytes(bytemask) << shift;
  CS.XMF = 1;
  new result1 = MEM.load(address, 0xFFFF & bytemask, variant, dri);
  result1 = ((result1 << shift) & targetbytes) | (argument1 & ~targetbytes);
stage CR:
  XVR[where] = result1;
\end{lstlisting}}
\newpage

\paragraph{Encoding} Format \textsf{LSU\_XPL16RBB} (Section~\ref{sec:format-LSU-XPL16RBB}), \verb|lsucode5 = 100|\\

\bitfields{ 1/P, 2/01, 3/011, 2/variant, 2/registerGj, 4/0111, 2/11, 3/100, 1/--, 6/registerY, 6/registerZ }


\paragraph{Syntax} \texttt{xplq}\emph{variant} \emph{registerGj}, \emph{registerY} = [\emph{registerZ}]

\paragraph{Description} The \emph{registerGj} is preloaded from the memory quadruple word at the effective address.


\paragraph{Modifier(s)}
\noindent\begin{itemize}
\item variant (cf. variant in section \ref{par:modifier-variant} on page \pageref{par:modifier-variant})
\end{itemize}

\paragraph{Execution} {Reservation table \textsf{LSU} (Section~\ref{sec:reservation-LSU})}
~
{\begin{lstlisting}
stage ID:
  new variant = _ZX_2(variant);
stage RR:
  new address = GPR[registerZ];
  new target = GPR[registerY];
  new shift = (target & 31) * 8;
  new index = (target >> 5) & 15;
  new where = (registerGj << 4) + index;
  new dri = where;
stage E1:
  new argument1 = XVR[where];
  new bytemask = 65535;
  new targetbytes = bits2bytes(bytemask) << shift;
  CS.XMF = 1;
  new result1 = MEM.load(address, 0xFFFF & bytemask, variant, dri);
  result1 = ((result1 << shift) & targetbytes) | (argument1 & ~targetbytes);
stage CR:
  XVR[where] = result1;
\end{lstlisting}}
\newpage

\paragraph{Encoding} Format \textsf{LSU\_XPL16RBB.O} (Section~\ref{sec:format-LSU-XPL16RBB.O}), \verb|lsucode5 = 100|\\

\bitfields{ 1/P, 2/00, 2/-- --, 27/offset27, 1/1, 2/01, 3/011, 2/variant, 2/registerGj, 4/0111, 2/11, 3/100, 1/--, 6/registerY, 6/registerZ }


\paragraph{Syntax} \texttt{xplq}\emph{variant} \emph{registerGj}, \emph{registerY} = \emph{offset27}[\emph{registerZ}]

\paragraph{Description} The \emph{registerGj} is preloaded from the memory quadruple word at the effective address.


\paragraph{Modifier(s)}
\noindent\begin{itemize}
\item variant (cf. variant in section \ref{par:modifier-variant} on page \pageref{par:modifier-variant})
\end{itemize}

\paragraph{Execution} {Reservation table \textsf{LSU.X} (Section~\ref{sec:reservation-LSU.X})}
~
{\begin{lstlisting}
stage ID:
  new offset = _SX_27(offset27);
  new variant = _ZX_2(variant);
stage RR:
  new base = GPR[registerZ];
  new target = GPR[registerY];
  new address = base + offset;
  new shift = (target & 31) * 8;
  new index = (target >> 5) & 15;
  new where = (registerGj << 4) + index;
  new dri = where;
stage E1:
  new argument1 = XVR[where];
  new bytemask = 65535;
  new targetbytes = bits2bytes(bytemask) << shift;
  CS.XMF = 1;
  new result1 = MEM.load(address, 0xFFFF & bytemask, variant, dri);
  result1 = ((result1 << shift) & targetbytes) | (argument1 & ~targetbytes);
stage CR:
  XVR[where] = result1;
\end{lstlisting}}
\newpage

\paragraph{Encoding} Format \textsf{LSU\_XPL16RBB.D} (Section~\ref{sec:format-LSU-XPL16RBB.D}), \verb|lsucode5 = 100|\\

\bitfields{ 1/P, 2/00, 2/-- --, 27/extend27, 1/1, 2/00, 2/-- --, 27/offset27, 1/1, 2/01, 3/011, 2/variant, 2/registerGj, 4/0111, 2/11, 3/100, 1/--, 6/registerY, 6/registerZ }


\paragraph{Syntax} \texttt{xplq}\emph{variant} \emph{registerGj}, \emph{registerY} = \emph{extend27\_\discretionary{}{}{}offset27}[\emph{registerZ}]

\paragraph{Description} The \emph{registerGj} is preloaded from the memory quadruple word at the effective address.


\paragraph{Modifier(s)}
\noindent\begin{itemize}
\item variant (cf. variant in section \ref{par:modifier-variant} on page \pageref{par:modifier-variant})
\end{itemize}

\paragraph{Execution} {Reservation table \textsf{LSU.Y} (Section~\ref{sec:reservation-LSU.Y})}
~
{\begin{lstlisting}
stage ID:
  new offset = _SX_54(extend27_offset27);
  new variant = _ZX_2(variant);
stage RR:
  new base = GPR[registerZ];
  new target = GPR[registerY];
  new address = base + offset;
  new shift = (target & 31) * 8;
  new index = (target >> 5) & 15;
  new where = (registerGj << 4) + index;
  new dri = where;
stage E1:
  new argument1 = XVR[where];
  new bytemask = 65535;
  new targetbytes = bits2bytes(bytemask) << shift;
  CS.XMF = 1;
  new result1 = MEM.load(address, 0xFFFF & bytemask, variant, dri);
  result1 = ((result1 << shift) & targetbytes) | (argument1 & ~targetbytes);
stage CR:
  XVR[where] = result1;
\end{lstlisting}}
\newpage

\paragraph{Encoding} Format \textsf{LSU\_XPL32RBB} (Section~\ref{sec:format-LSU-XPL32RBB}), \verb|lsucode5 = 100|\\

\bitfields{ 1/P, 2/01, 3/011, 2/variant, 1/registerGk, 5/01111, 2/11, 3/100, 1/--, 6/registerY, 6/registerZ }


\paragraph{Syntax} \texttt{xplq}\emph{variant} \emph{registerGk}, \emph{registerY} = [\emph{registerZ}]

\paragraph{Description} The \emph{registerGk} is preloaded from the memory quadruple word at the effective address.


\paragraph{Modifier(s)}
\noindent\begin{itemize}
\item variant (cf. variant in section \ref{par:modifier-variant} on page \pageref{par:modifier-variant})
\end{itemize}

\paragraph{Execution} {Reservation table \textsf{LSU} (Section~\ref{sec:reservation-LSU})}
~
{\begin{lstlisting}
stage ID:
  new variant = _ZX_2(variant);
stage RR:
  new address = GPR[registerZ];
  new target = GPR[registerY];
  new shift = (target & 31) * 8;
  new index = (target >> 5) & 31;
  new where = (registerGk << 5) + index;
  new dri = where;
stage E1:
  new argument1 = XVR[where];
  new bytemask = 65535;
  new targetbytes = bits2bytes(bytemask) << shift;
  CS.XMF = 1;
  new result1 = MEM.load(address, 0xFFFF & bytemask, variant, dri);
  result1 = ((result1 << shift) & targetbytes) | (argument1 & ~targetbytes);
stage CR:
  XVR[where] = result1;
\end{lstlisting}}
\newpage

\paragraph{Encoding} Format \textsf{LSU\_XPL32RBB.O} (Section~\ref{sec:format-LSU-XPL32RBB.O}), \verb|lsucode5 = 100|\\

\bitfields{ 1/P, 2/00, 2/-- --, 27/offset27, 1/1, 2/01, 3/011, 2/variant, 1/registerGk, 5/01111, 2/11, 3/100, 1/--, 6/registerY, 6/registerZ }


\paragraph{Syntax} \texttt{xplq}\emph{variant} \emph{registerGk}, \emph{registerY} = \emph{offset27}[\emph{registerZ}]

\paragraph{Description} The \emph{registerGk} is preloaded from the memory quadruple word at the effective address.


\paragraph{Modifier(s)}
\noindent\begin{itemize}
\item variant (cf. variant in section \ref{par:modifier-variant} on page \pageref{par:modifier-variant})
\end{itemize}

\paragraph{Execution} {Reservation table \textsf{LSU.X} (Section~\ref{sec:reservation-LSU.X})}
~
{\begin{lstlisting}
stage ID:
  new offset = _SX_27(offset27);
  new variant = _ZX_2(variant);
stage RR:
  new base = GPR[registerZ];
  new target = GPR[registerY];
  new address = base + offset;
  new shift = (target & 31) * 8;
  new index = (target >> 5) & 31;
  new where = (registerGk << 5) + index;
  new dri = where;
stage E1:
  new argument1 = XVR[where];
  new bytemask = 65535;
  new targetbytes = bits2bytes(bytemask) << shift;
  CS.XMF = 1;
  new result1 = MEM.load(address, 0xFFFF & bytemask, variant, dri);
  result1 = ((result1 << shift) & targetbytes) | (argument1 & ~targetbytes);
stage CR:
  XVR[where] = result1;
\end{lstlisting}}
\newpage

\paragraph{Encoding} Format \textsf{LSU\_XPL32RBB.D} (Section~\ref{sec:format-LSU-XPL32RBB.D}), \verb|lsucode5 = 100|\\

\bitfields{ 1/P, 2/00, 2/-- --, 27/extend27, 1/1, 2/00, 2/-- --, 27/offset27, 1/1, 2/01, 3/011, 2/variant, 1/registerGk, 5/01111, 2/11, 3/100, 1/--, 6/registerY, 6/registerZ }


\paragraph{Syntax} \texttt{xplq}\emph{variant} \emph{registerGk}, \emph{registerY} = \emph{extend27\_\discretionary{}{}{}offset27}[\emph{registerZ}]

\paragraph{Description} The \emph{registerGk} is preloaded from the memory quadruple word at the effective address.


\paragraph{Modifier(s)}
\noindent\begin{itemize}
\item variant (cf. variant in section \ref{par:modifier-variant} on page \pageref{par:modifier-variant})
\end{itemize}

\paragraph{Execution} {Reservation table \textsf{LSU.Y} (Section~\ref{sec:reservation-LSU.Y})}
~
{\begin{lstlisting}
stage ID:
  new offset = _SX_54(extend27_offset27);
  new variant = _ZX_2(variant);
stage RR:
  new base = GPR[registerZ];
  new target = GPR[registerY];
  new address = base + offset;
  new shift = (target & 31) * 8;
  new index = (target >> 5) & 31;
  new where = (registerGk << 5) + index;
  new dri = where;
stage E1:
  new argument1 = XVR[where];
  new bytemask = 65535;
  new targetbytes = bits2bytes(bytemask) << shift;
  CS.XMF = 1;
  new result1 = MEM.load(address, 0xFFFF & bytemask, variant, dri);
  result1 = ((result1 << shift) & targetbytes) | (argument1 & ~targetbytes);
stage CR:
  XVR[where] = result1;
\end{lstlisting}}
\newpage

\paragraph{Encoding} Format \textsf{LSU\_XPL64RBB} (Section~\ref{sec:format-LSU-XPL64RBB}), \verb|lsucode5 = 100|\\

\bitfields{ 1/P, 2/01, 3/011, 2/variant, 1/registerGl, 5/11111, 2/11, 3/100, 1/--, 6/registerY, 6/registerZ }


\paragraph{Syntax} \texttt{xplq}\emph{variant} \emph{registerGl}, \emph{registerY} = [\emph{registerZ}]

\paragraph{Description} The \emph{registerGl} is preloaded from the memory quadruple word at the effective address.


\paragraph{Modifier(s)}
\noindent\begin{itemize}
\item variant (cf. variant in section \ref{par:modifier-variant} on page \pageref{par:modifier-variant})
\end{itemize}

\paragraph{Execution} {Reservation table \textsf{LSU} (Section~\ref{sec:reservation-LSU})}
~
{\begin{lstlisting}
stage ID:
  new variant = _ZX_2(variant);
stage RR:
  new address = GPR[registerZ];
  new target = GPR[registerY];
  new shift = (target & 31) * 8;
  new index = (target >> 5) & 63;
  new where = (registerGl << 6) + index;
  new dri = where;
stage E1:
  new argument1 = XVR[where];
  new bytemask = 65535;
  new targetbytes = bits2bytes(bytemask) << shift;
  CS.XMF = 1;
  new result1 = MEM.load(address, 0xFFFF & bytemask, variant, dri);
  result1 = ((result1 << shift) & targetbytes) | (argument1 & ~targetbytes);
stage CR:
  XVR[where] = result1;
\end{lstlisting}}
\newpage

\paragraph{Encoding} Format \textsf{LSU\_XPL64RBB.O} (Section~\ref{sec:format-LSU-XPL64RBB.O}), \verb|lsucode5 = 100|\\

\bitfields{ 1/P, 2/00, 2/-- --, 27/offset27, 1/1, 2/01, 3/011, 2/variant, 1/registerGl, 5/11111, 2/11, 3/100, 1/--, 6/registerY, 6/registerZ }


\paragraph{Syntax} \texttt{xplq}\emph{variant} \emph{registerGl}, \emph{registerY} = \emph{offset27}[\emph{registerZ}]

\paragraph{Description} The \emph{registerGl} is preloaded from the memory quadruple word at the effective address.


\paragraph{Modifier(s)}
\noindent\begin{itemize}
\item variant (cf. variant in section \ref{par:modifier-variant} on page \pageref{par:modifier-variant})
\end{itemize}

\paragraph{Execution} {Reservation table \textsf{LSU.X} (Section~\ref{sec:reservation-LSU.X})}
~
{\begin{lstlisting}
stage ID:
  new offset = _SX_27(offset27);
  new variant = _ZX_2(variant);
stage RR:
  new base = GPR[registerZ];
  new target = GPR[registerY];
  new address = base + offset;
  new shift = (target & 31) * 8;
  new index = (target >> 5) & 63;
  new where = (registerGl << 6) + index;
  new dri = where;
stage E1:
  new argument1 = XVR[where];
  new bytemask = 65535;
  new targetbytes = bits2bytes(bytemask) << shift;
  CS.XMF = 1;
  new result1 = MEM.load(address, 0xFFFF & bytemask, variant, dri);
  result1 = ((result1 << shift) & targetbytes) | (argument1 & ~targetbytes);
stage CR:
  XVR[where] = result1;
\end{lstlisting}}
\newpage

\paragraph{Encoding} Format \textsf{LSU\_XPL64RBB.D} (Section~\ref{sec:format-LSU-XPL64RBB.D}), \verb|lsucode5 = 100|\\

\bitfields{ 1/P, 2/00, 2/-- --, 27/extend27, 1/1, 2/00, 2/-- --, 27/offset27, 1/1, 2/01, 3/011, 2/variant, 1/registerGl, 5/11111, 2/11, 3/100, 1/--, 6/registerY, 6/registerZ }


\paragraph{Syntax} \texttt{xplq}\emph{variant} \emph{registerGl}, \emph{registerY} = \emph{extend27\_\discretionary{}{}{}offset27}[\emph{registerZ}]

\paragraph{Description} The \emph{registerGl} is preloaded from the memory quadruple word at the effective address.


\paragraph{Modifier(s)}
\noindent\begin{itemize}
\item variant (cf. variant in section \ref{par:modifier-variant} on page \pageref{par:modifier-variant})
\end{itemize}

\paragraph{Execution} {Reservation table \textsf{LSU.Y} (Section~\ref{sec:reservation-LSU.Y})}
~
{\begin{lstlisting}
stage ID:
  new offset = _SX_54(extend27_offset27);
  new variant = _ZX_2(variant);
stage RR:
  new base = GPR[registerZ];
  new target = GPR[registerY];
  new address = base + offset;
  new shift = (target & 31) * 8;
  new index = (target >> 5) & 63;
  new where = (registerGl << 6) + index;
  new dri = where;
stage E1:
  new argument1 = XVR[where];
  new bytemask = 65535;
  new targetbytes = bits2bytes(bytemask) << shift;
  CS.XMF = 1;
  new result1 = MEM.load(address, 0xFFFF & bytemask, variant, dri);
  result1 = ((result1 << shift) & targetbytes) | (argument1 & ~targetbytes);
stage CR:
  XVR[where] = result1;
\end{lstlisting}}
\newpage
\subsubsection[{\small XPLW: Extension Preload Word}]{XPLW\hfill Extension Preload Word}
\label{instr:XPLW}\index{xplw}
 
\paragraph{Encoding} Format \textsf{LSU\_XPL2RBB} (Section~\ref{sec:format-LSU-XPL2RBB}), \verb|lsucode5 = 010|\\

\bitfields{ 1/P, 2/01, 3/011, 2/variant, 5/registerGg, 1/0, 2/11, 3/010, 1/--, 6/registerY, 6/registerZ }


\paragraph{Syntax} \texttt{xplw}\emph{variant} \emph{registerGg}, \emph{registerY} = [\emph{registerZ}]

\paragraph{Description} The \emph{registerGg} is preloaded from the memory word at the effective address.


\paragraph{Modifier(s)}
\noindent\begin{itemize}
\item variant (cf. variant in section \ref{par:modifier-variant} on page \pageref{par:modifier-variant})
\end{itemize}

\paragraph{Execution} {Reservation table \textsf{LSU} (Section~\ref{sec:reservation-LSU})}
~
{\begin{lstlisting}
stage ID:
  new variant = _ZX_2(variant);
stage RR:
  new address = GPR[registerZ];
  new target = GPR[registerY];
  new shift = (target & 31) * 8;
  new index = (target >> 5) & 1;
  new where = (registerGg << 1) + index;
  new dri = where;
stage E1:
  new argument1 = XVR[where];
  new bytemask = 15;
  new targetbytes = bits2bytes(bytemask) << shift;
  CS.XMF = 1;
  new result1 = MEM.load(address, 0xF & bytemask, variant, dri);
  result1 = ((result1 << shift) & targetbytes) | (argument1 & ~targetbytes);
stage CR:
  XVR[where] = result1;
\end{lstlisting}}
\newpage

\paragraph{Encoding} Format \textsf{LSU\_XPL2RBB.O} (Section~\ref{sec:format-LSU-XPL2RBB.O}), \verb|lsucode5 = 010|\\

\bitfields{ 1/P, 2/00, 2/-- --, 27/offset27, 1/1, 2/01, 3/011, 2/variant, 5/registerGg, 1/0, 2/11, 3/010, 1/--, 6/registerY, 6/registerZ }


\paragraph{Syntax} \texttt{xplw}\emph{variant} \emph{registerGg}, \emph{registerY} = \emph{offset27}[\emph{registerZ}]

\paragraph{Description} The \emph{registerGg} is preloaded from the memory word at the effective address.


\paragraph{Modifier(s)}
\noindent\begin{itemize}
\item variant (cf. variant in section \ref{par:modifier-variant} on page \pageref{par:modifier-variant})
\end{itemize}

\paragraph{Execution} {Reservation table \textsf{LSU.X} (Section~\ref{sec:reservation-LSU.X})}
~
{\begin{lstlisting}
stage ID:
  new offset = _SX_27(offset27);
  new variant = _ZX_2(variant);
stage RR:
  new base = GPR[registerZ];
  new target = GPR[registerY];
  new address = base + offset;
  new shift = (target & 31) * 8;
  new index = (target >> 5) & 1;
  new where = (registerGg << 1) + index;
  new dri = where;
stage E1:
  new argument1 = XVR[where];
  new bytemask = 15;
  new targetbytes = bits2bytes(bytemask) << shift;
  CS.XMF = 1;
  new result1 = MEM.load(address, 0xF & bytemask, variant, dri);
  result1 = ((result1 << shift) & targetbytes) | (argument1 & ~targetbytes);
stage CR:
  XVR[where] = result1;
\end{lstlisting}}
\newpage

\paragraph{Encoding} Format \textsf{LSU\_XPL2RBB.D} (Section~\ref{sec:format-LSU-XPL2RBB.D}), \verb|lsucode5 = 010|\\

\bitfields{ 1/P, 2/00, 2/-- --, 27/extend27, 1/1, 2/00, 2/-- --, 27/offset27, 1/1, 2/01, 3/011, 2/variant, 5/registerGg, 1/0, 2/11, 3/010, 1/--, 6/registerY, 6/registerZ }


\paragraph{Syntax} \texttt{xplw}\emph{variant} \emph{registerGg}, \emph{registerY} = \emph{extend27\_\discretionary{}{}{}offset27}[\emph{registerZ}]

\paragraph{Description} The \emph{registerGg} is preloaded from the memory word at the effective address.


\paragraph{Modifier(s)}
\noindent\begin{itemize}
\item variant (cf. variant in section \ref{par:modifier-variant} on page \pageref{par:modifier-variant})
\end{itemize}

\paragraph{Execution} {Reservation table \textsf{LSU.Y} (Section~\ref{sec:reservation-LSU.Y})}
~
{\begin{lstlisting}
stage ID:
  new offset = _SX_54(extend27_offset27);
  new variant = _ZX_2(variant);
stage RR:
  new base = GPR[registerZ];
  new target = GPR[registerY];
  new address = base + offset;
  new shift = (target & 31) * 8;
  new index = (target >> 5) & 1;
  new where = (registerGg << 1) + index;
  new dri = where;
stage E1:
  new argument1 = XVR[where];
  new bytemask = 15;
  new targetbytes = bits2bytes(bytemask) << shift;
  CS.XMF = 1;
  new result1 = MEM.load(address, 0xF & bytemask, variant, dri);
  result1 = ((result1 << shift) & targetbytes) | (argument1 & ~targetbytes);
stage CR:
  XVR[where] = result1;
\end{lstlisting}}
\newpage

\paragraph{Encoding} Format \textsf{LSU\_XPL4RBB} (Section~\ref{sec:format-LSU-XPL4RBB}), \verb|lsucode5 = 010|\\

\bitfields{ 1/P, 2/01, 3/011, 2/variant, 4/registerGh, 2/01, 2/11, 3/010, 1/--, 6/registerY, 6/registerZ }


\paragraph{Syntax} \texttt{xplw}\emph{variant} \emph{registerGh}, \emph{registerY} = [\emph{registerZ}]

\paragraph{Description} The \emph{registerGh} is preloaded from the memory word at the effective address.


\paragraph{Modifier(s)}
\noindent\begin{itemize}
\item variant (cf. variant in section \ref{par:modifier-variant} on page \pageref{par:modifier-variant})
\end{itemize}

\paragraph{Execution} {Reservation table \textsf{LSU} (Section~\ref{sec:reservation-LSU})}
~
{\begin{lstlisting}
stage ID:
  new variant = _ZX_2(variant);
stage RR:
  new address = GPR[registerZ];
  new target = GPR[registerY];
  new shift = (target & 31) * 8;
  new index = (target >> 5) & 3;
  new where = (registerGh << 2) + index;
  new dri = where;
stage E1:
  new argument1 = XVR[where];
  new bytemask = 15;
  new targetbytes = bits2bytes(bytemask) << shift;
  CS.XMF = 1;
  new result1 = MEM.load(address, 0xF & bytemask, variant, dri);
  result1 = ((result1 << shift) & targetbytes) | (argument1 & ~targetbytes);
stage CR:
  XVR[where] = result1;
\end{lstlisting}}
\newpage

\paragraph{Encoding} Format \textsf{LSU\_XPL4RBB.O} (Section~\ref{sec:format-LSU-XPL4RBB.O}), \verb|lsucode5 = 010|\\

\bitfields{ 1/P, 2/00, 2/-- --, 27/offset27, 1/1, 2/01, 3/011, 2/variant, 4/registerGh, 2/01, 2/11, 3/010, 1/--, 6/registerY, 6/registerZ }


\paragraph{Syntax} \texttt{xplw}\emph{variant} \emph{registerGh}, \emph{registerY} = \emph{offset27}[\emph{registerZ}]

\paragraph{Description} The \emph{registerGh} is preloaded from the memory word at the effective address.


\paragraph{Modifier(s)}
\noindent\begin{itemize}
\item variant (cf. variant in section \ref{par:modifier-variant} on page \pageref{par:modifier-variant})
\end{itemize}

\paragraph{Execution} {Reservation table \textsf{LSU.X} (Section~\ref{sec:reservation-LSU.X})}
~
{\begin{lstlisting}
stage ID:
  new offset = _SX_27(offset27);
  new variant = _ZX_2(variant);
stage RR:
  new base = GPR[registerZ];
  new target = GPR[registerY];
  new address = base + offset;
  new shift = (target & 31) * 8;
  new index = (target >> 5) & 3;
  new where = (registerGh << 2) + index;
  new dri = where;
stage E1:
  new argument1 = XVR[where];
  new bytemask = 15;
  new targetbytes = bits2bytes(bytemask) << shift;
  CS.XMF = 1;
  new result1 = MEM.load(address, 0xF & bytemask, variant, dri);
  result1 = ((result1 << shift) & targetbytes) | (argument1 & ~targetbytes);
stage CR:
  XVR[where] = result1;
\end{lstlisting}}
\newpage

\paragraph{Encoding} Format \textsf{LSU\_XPL4RBB.D} (Section~\ref{sec:format-LSU-XPL4RBB.D}), \verb|lsucode5 = 010|\\

\bitfields{ 1/P, 2/00, 2/-- --, 27/extend27, 1/1, 2/00, 2/-- --, 27/offset27, 1/1, 2/01, 3/011, 2/variant, 4/registerGh, 2/01, 2/11, 3/010, 1/--, 6/registerY, 6/registerZ }


\paragraph{Syntax} \texttt{xplw}\emph{variant} \emph{registerGh}, \emph{registerY} = \emph{extend27\_\discretionary{}{}{}offset27}[\emph{registerZ}]

\paragraph{Description} The \emph{registerGh} is preloaded from the memory word at the effective address.


\paragraph{Modifier(s)}
\noindent\begin{itemize}
\item variant (cf. variant in section \ref{par:modifier-variant} on page \pageref{par:modifier-variant})
\end{itemize}

\paragraph{Execution} {Reservation table \textsf{LSU.Y} (Section~\ref{sec:reservation-LSU.Y})}
~
{\begin{lstlisting}
stage ID:
  new offset = _SX_54(extend27_offset27);
  new variant = _ZX_2(variant);
stage RR:
  new base = GPR[registerZ];
  new target = GPR[registerY];
  new address = base + offset;
  new shift = (target & 31) * 8;
  new index = (target >> 5) & 3;
  new where = (registerGh << 2) + index;
  new dri = where;
stage E1:
  new argument1 = XVR[where];
  new bytemask = 15;
  new targetbytes = bits2bytes(bytemask) << shift;
  CS.XMF = 1;
  new result1 = MEM.load(address, 0xF & bytemask, variant, dri);
  result1 = ((result1 << shift) & targetbytes) | (argument1 & ~targetbytes);
stage CR:
  XVR[where] = result1;
\end{lstlisting}}
\newpage

\paragraph{Encoding} Format \textsf{LSU\_XPL8RBB} (Section~\ref{sec:format-LSU-XPL8RBB}), \verb|lsucode5 = 010|\\

\bitfields{ 1/P, 2/01, 3/011, 2/variant, 3/registerGi, 3/011, 2/11, 3/010, 1/--, 6/registerY, 6/registerZ }


\paragraph{Syntax} \texttt{xplw}\emph{variant} \emph{registerGi}, \emph{registerY} = [\emph{registerZ}]

\paragraph{Description} The \emph{registerGi} is preloaded from the memory word at the effective address.


\paragraph{Modifier(s)}
\noindent\begin{itemize}
\item variant (cf. variant in section \ref{par:modifier-variant} on page \pageref{par:modifier-variant})
\end{itemize}

\paragraph{Execution} {Reservation table \textsf{LSU} (Section~\ref{sec:reservation-LSU})}
~
{\begin{lstlisting}
stage ID:
  new variant = _ZX_2(variant);
stage RR:
  new address = GPR[registerZ];
  new target = GPR[registerY];
  new shift = (target & 31) * 8;
  new index = (target >> 5) & 7;
  new where = (registerGi << 3) + index;
  new dri = where;
stage E1:
  new argument1 = XVR[where];
  new bytemask = 15;
  new targetbytes = bits2bytes(bytemask) << shift;
  CS.XMF = 1;
  new result1 = MEM.load(address, 0xF & bytemask, variant, dri);
  result1 = ((result1 << shift) & targetbytes) | (argument1 & ~targetbytes);
stage CR:
  XVR[where] = result1;
\end{lstlisting}}
\newpage

\paragraph{Encoding} Format \textsf{LSU\_XPL8RBB.O} (Section~\ref{sec:format-LSU-XPL8RBB.O}), \verb|lsucode5 = 010|\\

\bitfields{ 1/P, 2/00, 2/-- --, 27/offset27, 1/1, 2/01, 3/011, 2/variant, 3/registerGi, 3/011, 2/11, 3/010, 1/--, 6/registerY, 6/registerZ }


\paragraph{Syntax} \texttt{xplw}\emph{variant} \emph{registerGi}, \emph{registerY} = \emph{offset27}[\emph{registerZ}]

\paragraph{Description} The \emph{registerGi} is preloaded from the memory word at the effective address.


\paragraph{Modifier(s)}
\noindent\begin{itemize}
\item variant (cf. variant in section \ref{par:modifier-variant} on page \pageref{par:modifier-variant})
\end{itemize}

\paragraph{Execution} {Reservation table \textsf{LSU.X} (Section~\ref{sec:reservation-LSU.X})}
~
{\begin{lstlisting}
stage ID:
  new offset = _SX_27(offset27);
  new variant = _ZX_2(variant);
stage RR:
  new base = GPR[registerZ];
  new target = GPR[registerY];
  new address = base + offset;
  new shift = (target & 31) * 8;
  new index = (target >> 5) & 7;
  new where = (registerGi << 3) + index;
  new dri = where;
stage E1:
  new argument1 = XVR[where];
  new bytemask = 15;
  new targetbytes = bits2bytes(bytemask) << shift;
  CS.XMF = 1;
  new result1 = MEM.load(address, 0xF & bytemask, variant, dri);
  result1 = ((result1 << shift) & targetbytes) | (argument1 & ~targetbytes);
stage CR:
  XVR[where] = result1;
\end{lstlisting}}
\newpage

\paragraph{Encoding} Format \textsf{LSU\_XPL8RBB.D} (Section~\ref{sec:format-LSU-XPL8RBB.D}), \verb|lsucode5 = 010|\\

\bitfields{ 1/P, 2/00, 2/-- --, 27/extend27, 1/1, 2/00, 2/-- --, 27/offset27, 1/1, 2/01, 3/011, 2/variant, 3/registerGi, 3/011, 2/11, 3/010, 1/--, 6/registerY, 6/registerZ }


\paragraph{Syntax} \texttt{xplw}\emph{variant} \emph{registerGi}, \emph{registerY} = \emph{extend27\_\discretionary{}{}{}offset27}[\emph{registerZ}]

\paragraph{Description} The \emph{registerGi} is preloaded from the memory word at the effective address.


\paragraph{Modifier(s)}
\noindent\begin{itemize}
\item variant (cf. variant in section \ref{par:modifier-variant} on page \pageref{par:modifier-variant})
\end{itemize}

\paragraph{Execution} {Reservation table \textsf{LSU.Y} (Section~\ref{sec:reservation-LSU.Y})}
~
{\begin{lstlisting}
stage ID:
  new offset = _SX_54(extend27_offset27);
  new variant = _ZX_2(variant);
stage RR:
  new base = GPR[registerZ];
  new target = GPR[registerY];
  new address = base + offset;
  new shift = (target & 31) * 8;
  new index = (target >> 5) & 7;
  new where = (registerGi << 3) + index;
  new dri = where;
stage E1:
  new argument1 = XVR[where];
  new bytemask = 15;
  new targetbytes = bits2bytes(bytemask) << shift;
  CS.XMF = 1;
  new result1 = MEM.load(address, 0xF & bytemask, variant, dri);
  result1 = ((result1 << shift) & targetbytes) | (argument1 & ~targetbytes);
stage CR:
  XVR[where] = result1;
\end{lstlisting}}
\newpage

\paragraph{Encoding} Format \textsf{LSU\_XPL16RBB} (Section~\ref{sec:format-LSU-XPL16RBB}), \verb|lsucode5 = 010|\\

\bitfields{ 1/P, 2/01, 3/011, 2/variant, 2/registerGj, 4/0111, 2/11, 3/010, 1/--, 6/registerY, 6/registerZ }


\paragraph{Syntax} \texttt{xplw}\emph{variant} \emph{registerGj}, \emph{registerY} = [\emph{registerZ}]

\paragraph{Description} The \emph{registerGj} is preloaded from the memory word at the effective address.


\paragraph{Modifier(s)}
\noindent\begin{itemize}
\item variant (cf. variant in section \ref{par:modifier-variant} on page \pageref{par:modifier-variant})
\end{itemize}

\paragraph{Execution} {Reservation table \textsf{LSU} (Section~\ref{sec:reservation-LSU})}
~
{\begin{lstlisting}
stage ID:
  new variant = _ZX_2(variant);
stage RR:
  new address = GPR[registerZ];
  new target = GPR[registerY];
  new shift = (target & 31) * 8;
  new index = (target >> 5) & 15;
  new where = (registerGj << 4) + index;
  new dri = where;
stage E1:
  new argument1 = XVR[where];
  new bytemask = 15;
  new targetbytes = bits2bytes(bytemask) << shift;
  CS.XMF = 1;
  new result1 = MEM.load(address, 0xF & bytemask, variant, dri);
  result1 = ((result1 << shift) & targetbytes) | (argument1 & ~targetbytes);
stage CR:
  XVR[where] = result1;
\end{lstlisting}}
\newpage

\paragraph{Encoding} Format \textsf{LSU\_XPL16RBB.O} (Section~\ref{sec:format-LSU-XPL16RBB.O}), \verb|lsucode5 = 010|\\

\bitfields{ 1/P, 2/00, 2/-- --, 27/offset27, 1/1, 2/01, 3/011, 2/variant, 2/registerGj, 4/0111, 2/11, 3/010, 1/--, 6/registerY, 6/registerZ }


\paragraph{Syntax} \texttt{xplw}\emph{variant} \emph{registerGj}, \emph{registerY} = \emph{offset27}[\emph{registerZ}]

\paragraph{Description} The \emph{registerGj} is preloaded from the memory word at the effective address.


\paragraph{Modifier(s)}
\noindent\begin{itemize}
\item variant (cf. variant in section \ref{par:modifier-variant} on page \pageref{par:modifier-variant})
\end{itemize}

\paragraph{Execution} {Reservation table \textsf{LSU.X} (Section~\ref{sec:reservation-LSU.X})}
~
{\begin{lstlisting}
stage ID:
  new offset = _SX_27(offset27);
  new variant = _ZX_2(variant);
stage RR:
  new base = GPR[registerZ];
  new target = GPR[registerY];
  new address = base + offset;
  new shift = (target & 31) * 8;
  new index = (target >> 5) & 15;
  new where = (registerGj << 4) + index;
  new dri = where;
stage E1:
  new argument1 = XVR[where];
  new bytemask = 15;
  new targetbytes = bits2bytes(bytemask) << shift;
  CS.XMF = 1;
  new result1 = MEM.load(address, 0xF & bytemask, variant, dri);
  result1 = ((result1 << shift) & targetbytes) | (argument1 & ~targetbytes);
stage CR:
  XVR[where] = result1;
\end{lstlisting}}
\newpage

\paragraph{Encoding} Format \textsf{LSU\_XPL16RBB.D} (Section~\ref{sec:format-LSU-XPL16RBB.D}), \verb|lsucode5 = 010|\\

\bitfields{ 1/P, 2/00, 2/-- --, 27/extend27, 1/1, 2/00, 2/-- --, 27/offset27, 1/1, 2/01, 3/011, 2/variant, 2/registerGj, 4/0111, 2/11, 3/010, 1/--, 6/registerY, 6/registerZ }


\paragraph{Syntax} \texttt{xplw}\emph{variant} \emph{registerGj}, \emph{registerY} = \emph{extend27\_\discretionary{}{}{}offset27}[\emph{registerZ}]

\paragraph{Description} The \emph{registerGj} is preloaded from the memory word at the effective address.


\paragraph{Modifier(s)}
\noindent\begin{itemize}
\item variant (cf. variant in section \ref{par:modifier-variant} on page \pageref{par:modifier-variant})
\end{itemize}

\paragraph{Execution} {Reservation table \textsf{LSU.Y} (Section~\ref{sec:reservation-LSU.Y})}
~
{\begin{lstlisting}
stage ID:
  new offset = _SX_54(extend27_offset27);
  new variant = _ZX_2(variant);
stage RR:
  new base = GPR[registerZ];
  new target = GPR[registerY];
  new address = base + offset;
  new shift = (target & 31) * 8;
  new index = (target >> 5) & 15;
  new where = (registerGj << 4) + index;
  new dri = where;
stage E1:
  new argument1 = XVR[where];
  new bytemask = 15;
  new targetbytes = bits2bytes(bytemask) << shift;
  CS.XMF = 1;
  new result1 = MEM.load(address, 0xF & bytemask, variant, dri);
  result1 = ((result1 << shift) & targetbytes) | (argument1 & ~targetbytes);
stage CR:
  XVR[where] = result1;
\end{lstlisting}}
\newpage

\paragraph{Encoding} Format \textsf{LSU\_XPL32RBB} (Section~\ref{sec:format-LSU-XPL32RBB}), \verb|lsucode5 = 010|\\

\bitfields{ 1/P, 2/01, 3/011, 2/variant, 1/registerGk, 5/01111, 2/11, 3/010, 1/--, 6/registerY, 6/registerZ }


\paragraph{Syntax} \texttt{xplw}\emph{variant} \emph{registerGk}, \emph{registerY} = [\emph{registerZ}]

\paragraph{Description} The \emph{registerGk} is preloaded from the memory word at the effective address.


\paragraph{Modifier(s)}
\noindent\begin{itemize}
\item variant (cf. variant in section \ref{par:modifier-variant} on page \pageref{par:modifier-variant})
\end{itemize}

\paragraph{Execution} {Reservation table \textsf{LSU} (Section~\ref{sec:reservation-LSU})}
~
{\begin{lstlisting}
stage ID:
  new variant = _ZX_2(variant);
stage RR:
  new address = GPR[registerZ];
  new target = GPR[registerY];
  new shift = (target & 31) * 8;
  new index = (target >> 5) & 31;
  new where = (registerGk << 5) + index;
  new dri = where;
stage E1:
  new argument1 = XVR[where];
  new bytemask = 15;
  new targetbytes = bits2bytes(bytemask) << shift;
  CS.XMF = 1;
  new result1 = MEM.load(address, 0xF & bytemask, variant, dri);
  result1 = ((result1 << shift) & targetbytes) | (argument1 & ~targetbytes);
stage CR:
  XVR[where] = result1;
\end{lstlisting}}
\newpage

\paragraph{Encoding} Format \textsf{LSU\_XPL32RBB.O} (Section~\ref{sec:format-LSU-XPL32RBB.O}), \verb|lsucode5 = 010|\\

\bitfields{ 1/P, 2/00, 2/-- --, 27/offset27, 1/1, 2/01, 3/011, 2/variant, 1/registerGk, 5/01111, 2/11, 3/010, 1/--, 6/registerY, 6/registerZ }


\paragraph{Syntax} \texttt{xplw}\emph{variant} \emph{registerGk}, \emph{registerY} = \emph{offset27}[\emph{registerZ}]

\paragraph{Description} The \emph{registerGk} is preloaded from the memory word at the effective address.


\paragraph{Modifier(s)}
\noindent\begin{itemize}
\item variant (cf. variant in section \ref{par:modifier-variant} on page \pageref{par:modifier-variant})
\end{itemize}

\paragraph{Execution} {Reservation table \textsf{LSU.X} (Section~\ref{sec:reservation-LSU.X})}
~
{\begin{lstlisting}
stage ID:
  new offset = _SX_27(offset27);
  new variant = _ZX_2(variant);
stage RR:
  new base = GPR[registerZ];
  new target = GPR[registerY];
  new address = base + offset;
  new shift = (target & 31) * 8;
  new index = (target >> 5) & 31;
  new where = (registerGk << 5) + index;
  new dri = where;
stage E1:
  new argument1 = XVR[where];
  new bytemask = 15;
  new targetbytes = bits2bytes(bytemask) << shift;
  CS.XMF = 1;
  new result1 = MEM.load(address, 0xF & bytemask, variant, dri);
  result1 = ((result1 << shift) & targetbytes) | (argument1 & ~targetbytes);
stage CR:
  XVR[where] = result1;
\end{lstlisting}}
\newpage

\paragraph{Encoding} Format \textsf{LSU\_XPL32RBB.D} (Section~\ref{sec:format-LSU-XPL32RBB.D}), \verb|lsucode5 = 010|\\

\bitfields{ 1/P, 2/00, 2/-- --, 27/extend27, 1/1, 2/00, 2/-- --, 27/offset27, 1/1, 2/01, 3/011, 2/variant, 1/registerGk, 5/01111, 2/11, 3/010, 1/--, 6/registerY, 6/registerZ }


\paragraph{Syntax} \texttt{xplw}\emph{variant} \emph{registerGk}, \emph{registerY} = \emph{extend27\_\discretionary{}{}{}offset27}[\emph{registerZ}]

\paragraph{Description} The \emph{registerGk} is preloaded from the memory word at the effective address.


\paragraph{Modifier(s)}
\noindent\begin{itemize}
\item variant (cf. variant in section \ref{par:modifier-variant} on page \pageref{par:modifier-variant})
\end{itemize}

\paragraph{Execution} {Reservation table \textsf{LSU.Y} (Section~\ref{sec:reservation-LSU.Y})}
~
{\begin{lstlisting}
stage ID:
  new offset = _SX_54(extend27_offset27);
  new variant = _ZX_2(variant);
stage RR:
  new base = GPR[registerZ];
  new target = GPR[registerY];
  new address = base + offset;
  new shift = (target & 31) * 8;
  new index = (target >> 5) & 31;
  new where = (registerGk << 5) + index;
  new dri = where;
stage E1:
  new argument1 = XVR[where];
  new bytemask = 15;
  new targetbytes = bits2bytes(bytemask) << shift;
  CS.XMF = 1;
  new result1 = MEM.load(address, 0xF & bytemask, variant, dri);
  result1 = ((result1 << shift) & targetbytes) | (argument1 & ~targetbytes);
stage CR:
  XVR[where] = result1;
\end{lstlisting}}
\newpage

\paragraph{Encoding} Format \textsf{LSU\_XPL64RBB} (Section~\ref{sec:format-LSU-XPL64RBB}), \verb|lsucode5 = 010|\\

\bitfields{ 1/P, 2/01, 3/011, 2/variant, 1/registerGl, 5/11111, 2/11, 3/010, 1/--, 6/registerY, 6/registerZ }


\paragraph{Syntax} \texttt{xplw}\emph{variant} \emph{registerGl}, \emph{registerY} = [\emph{registerZ}]

\paragraph{Description} The \emph{registerGl} is preloaded from the memory word at the effective address.


\paragraph{Modifier(s)}
\noindent\begin{itemize}
\item variant (cf. variant in section \ref{par:modifier-variant} on page \pageref{par:modifier-variant})
\end{itemize}

\paragraph{Execution} {Reservation table \textsf{LSU} (Section~\ref{sec:reservation-LSU})}
~
{\begin{lstlisting}
stage ID:
  new variant = _ZX_2(variant);
stage RR:
  new address = GPR[registerZ];
  new target = GPR[registerY];
  new shift = (target & 31) * 8;
  new index = (target >> 5) & 63;
  new where = (registerGl << 6) + index;
  new dri = where;
stage E1:
  new argument1 = XVR[where];
  new bytemask = 15;
  new targetbytes = bits2bytes(bytemask) << shift;
  CS.XMF = 1;
  new result1 = MEM.load(address, 0xF & bytemask, variant, dri);
  result1 = ((result1 << shift) & targetbytes) | (argument1 & ~targetbytes);
stage CR:
  XVR[where] = result1;
\end{lstlisting}}
\newpage

\paragraph{Encoding} Format \textsf{LSU\_XPL64RBB.O} (Section~\ref{sec:format-LSU-XPL64RBB.O}), \verb|lsucode5 = 010|\\

\bitfields{ 1/P, 2/00, 2/-- --, 27/offset27, 1/1, 2/01, 3/011, 2/variant, 1/registerGl, 5/11111, 2/11, 3/010, 1/--, 6/registerY, 6/registerZ }


\paragraph{Syntax} \texttt{xplw}\emph{variant} \emph{registerGl}, \emph{registerY} = \emph{offset27}[\emph{registerZ}]

\paragraph{Description} The \emph{registerGl} is preloaded from the memory word at the effective address.


\paragraph{Modifier(s)}
\noindent\begin{itemize}
\item variant (cf. variant in section \ref{par:modifier-variant} on page \pageref{par:modifier-variant})
\end{itemize}

\paragraph{Execution} {Reservation table \textsf{LSU.X} (Section~\ref{sec:reservation-LSU.X})}
~
{\begin{lstlisting}
stage ID:
  new offset = _SX_27(offset27);
  new variant = _ZX_2(variant);
stage RR:
  new base = GPR[registerZ];
  new target = GPR[registerY];
  new address = base + offset;
  new shift = (target & 31) * 8;
  new index = (target >> 5) & 63;
  new where = (registerGl << 6) + index;
  new dri = where;
stage E1:
  new argument1 = XVR[where];
  new bytemask = 15;
  new targetbytes = bits2bytes(bytemask) << shift;
  CS.XMF = 1;
  new result1 = MEM.load(address, 0xF & bytemask, variant, dri);
  result1 = ((result1 << shift) & targetbytes) | (argument1 & ~targetbytes);
stage CR:
  XVR[where] = result1;
\end{lstlisting}}
\newpage

\paragraph{Encoding} Format \textsf{LSU\_XPL64RBB.D} (Section~\ref{sec:format-LSU-XPL64RBB.D}), \verb|lsucode5 = 010|\\

\bitfields{ 1/P, 2/00, 2/-- --, 27/extend27, 1/1, 2/00, 2/-- --, 27/offset27, 1/1, 2/01, 3/011, 2/variant, 1/registerGl, 5/11111, 2/11, 3/010, 1/--, 6/registerY, 6/registerZ }


\paragraph{Syntax} \texttt{xplw}\emph{variant} \emph{registerGl}, \emph{registerY} = \emph{extend27\_\discretionary{}{}{}offset27}[\emph{registerZ}]

\paragraph{Description} The \emph{registerGl} is preloaded from the memory word at the effective address.


\paragraph{Modifier(s)}
\noindent\begin{itemize}
\item variant (cf. variant in section \ref{par:modifier-variant} on page \pageref{par:modifier-variant})
\end{itemize}

\paragraph{Execution} {Reservation table \textsf{LSU.Y} (Section~\ref{sec:reservation-LSU.Y})}
~
{\begin{lstlisting}
stage ID:
  new offset = _SX_54(extend27_offset27);
  new variant = _ZX_2(variant);
stage RR:
  new base = GPR[registerZ];
  new target = GPR[registerY];
  new address = base + offset;
  new shift = (target & 31) * 8;
  new index = (target >> 5) & 63;
  new where = (registerGl << 6) + index;
  new dri = where;
stage E1:
  new argument1 = XVR[where];
  new bytemask = 15;
  new targetbytes = bits2bytes(bytemask) << shift;
  CS.XMF = 1;
  new result1 = MEM.load(address, 0xF & bytemask, variant, dri);
  result1 = ((result1 << shift) & targetbytes) | (argument1 & ~targetbytes);
stage CR:
  XVR[where] = result1;
\end{lstlisting}}
\newpage
\subsubsection[{\small XSO: Extension Store Octuple Word}]{XSO\hfill Extension Store Octuple Word}
\label{instr:XSO}\index{xso}
 
\paragraph{Encoding} Format \textsf{LSU\_XSOBO} (Section~\ref{sec:format-LSU-XSOBO}), \verb|steering = 01|\\

\bitfields{ 1/P, 2/01, 5/10101, 6/registerE, 2/01, 10/signed10, 6/registerZ }


\paragraph{Syntax} \texttt{xso} \emph{signed10}[\emph{registerZ}] = \emph{registerE}

\paragraph{Description} The effective address is computed by adding the \emph{signed10} to the \emph{registerZ}.
The \emph{registerE} is stored into the memory octuple word at the effective address.


\paragraph{Execution} {Reservation table \textsf{LSU\_MEMW\_ACCR} (Section~\ref{sec:reservation-LSU-MEMW-ACCR})}
~
{\begin{lstlisting}
stage ID:
  new offset = _SX_10(signed10);
stage RR:
  new bytemask = 0xFFFFFFFF;
  new dri = registerE;
  new base = GPR[registerZ];
  new address = base + offset;
stage E1:
  new argument3 = XVR[registerE];
  CS.XMF = 1;
stage E3:
  MEM.store(address, 0xFFFFFFFF & bytemask, 0, argument3, dri);
\end{lstlisting}}
\newpage

\paragraph{Encoding} Format \textsf{LSU\_XSOBO.X} (Section~\ref{sec:format-LSU-XSOBO.X}), \verb|steering = 01|\\

\bitfields{ 1/P, 2/00, 2/-- --, 27/upper27, 1/1, 2/01, 5/10101, 6/registerE, 2/01, 10/lower10, 6/registerZ }


\paragraph{Syntax} \texttt{xso} \emph{upper27\_\discretionary{}{}{}lower10}[\emph{registerZ}] = \emph{registerE}

\paragraph{Description} The effective address is computed by adding the \emph{upper27\_\discretionary{}{}{}lower10} to the \emph{registerZ}.
The \emph{registerE} is stored into the memory octuple word at the effective address.


\paragraph{Execution} {Reservation table \textsf{LSU\_MEMW\_ACCR.X} (Section~\ref{sec:reservation-LSU-MEMW-ACCR.X})}
~
{\begin{lstlisting}
stage ID:
  new offset = _SX_37(upper27_lower10);
stage RR:
  new bytemask = 0xFFFFFFFF;
  new dri = registerE;
  new base = GPR[registerZ];
  new address = base + offset;
stage E1:
  new argument3 = XVR[registerE];
  CS.XMF = 1;
stage E3:
  MEM.store(address, 0xFFFFFFFF & bytemask, 0, argument3, dri);
\end{lstlisting}}
\newpage

\paragraph{Encoding} Format \textsf{LSU\_XSOBO.Y} (Section~\ref{sec:format-LSU-XSOBO.Y}), \verb|steering = 01|\\

\bitfields{ 1/P, 2/00, 2/-- --, 27/extend27, 1/1, 2/00, 2/-- --, 27/upper27, 1/1, 2/01, 5/10101, 6/registerE, 2/01, 10/lower10, 6/registerZ }


\paragraph{Syntax} \texttt{xso} \emph{extend27\_\discretionary{}{}{}upper27\_\discretionary{}{}{}lower10}[\emph{registerZ}] = \emph{registerE}

\paragraph{Description} The effective address is computed by adding the \emph{extend27\_\discretionary{}{}{}upper27\_\discretionary{}{}{}lower10} to the \emph{registerZ}.
The \emph{registerE} is stored into the memory octuple word at the effective address.


\paragraph{Execution} {Reservation table \textsf{LSU\_MEMW\_ACCR.Y} (Section~\ref{sec:reservation-LSU-MEMW-ACCR.Y})}
~
{\begin{lstlisting}
stage ID:
  new offset = _WX_64(extend27_upper27_lower10);
stage RR:
  new bytemask = 0xFFFFFFFF;
  new dri = registerE;
  new base = GPR[registerZ];
  new address = base + offset;
stage E1:
  new argument3 = XVR[registerE];
  CS.XMF = 1;
stage E3:
  MEM.store(address, 0xFFFFFFFF & bytemask, 0, argument3, dri);
\end{lstlisting}}
\newpage

\paragraph{Encoding} Format \textsf{LSU\_XSOBI} (Section~\ref{sec:format-LSU-XSOBI}), \verb|steering = 01|\\

\bitfields{ 1/P, 2/01, 5/10101, 6/registerE, 2/11, 3/111, 1/--, 6/registerY, 6/registerZ }


\paragraph{Syntax} \texttt{xso} \emph{registerY}[\emph{registerZ}] = \emph{registerE}

\paragraph{Description} The effective address is computed by adding the \emph{registerZ} to the \emph{registerY} and scaled by the \emph{}.
The \emph{registerE} is stored into the memory octuple word at the effective address.


\paragraph{Execution} {Reservation table \textsf{LSU\_MEMW\_ACCR} (Section~\ref{sec:reservation-LSU-MEMW-ACCR})}
~
{\begin{lstlisting}
stage RR:
  new bytemask = 0xFFFFFFFF;
  new doscale = 0;
  new dri = registerE;
  new index = GPR[registerY];
  new base = GPR[registerZ];
  new scaling = doscale ? 32 : 1;
  new address = base + (index * scaling);
stage E1:
  new argument3 = XVR[registerE];
  CS.XMF = 1;
stage E3:
  MEM.store(address, 0xFFFFFFFF & bytemask, 0, argument3, dri);
\end{lstlisting}}
\newpage

\paragraph{Encoding} Format \textsf{LSU\_XSOC4BO} (Section~\ref{sec:format-LSU-XSOC4BO}), \verb|steering = 01|\\

\bitfields{ 1/P, 2/01, 5/10111, 4/registerEq, 2/qindex, 2/01, 10/signed10, 6/registerZ }


\paragraph{Syntax} \texttt{xso}\emph{qindex} \emph{signed10}[\emph{registerZ}] = \emph{registerEq}

\paragraph{Description} The effective address is computed by adding the \emph{signed10} to the \emph{registerZ}.
The \emph{registerEq} is stored into the memory octuple word at the effective address.


\paragraph{Modifier(s)}
\noindent\begin{itemize}
\item qindex (cf. qindex in section \ref{par:modifier-qindex} on page \pageref{par:modifier-qindex})
\end{itemize}

\paragraph{Execution} {Reservation table \textsf{LSU\_MEMW\_ACCR} (Section~\ref{sec:reservation-LSU-MEMW-ACCR})}
~
{\begin{lstlisting}
stage ID:
  new qindex = _ZX_2(qindex);
  new offset = _SX_10(signed10);
stage RR:
  new bytemask = 0xFFFFFFFF;
  new dri = registerEq << 2;
  new base = GPR[registerZ];
  new address = base + offset;
stage E1:
  new argument3_0 = XMR[registerEq]:0;
  new argument3_0_0 = argument3_0.64[0];
  new argument3_0_1 = argument3_0.64[1];
  new argument3_0_2 = argument3_0.64[2];
  new argument3_0_3 = argument3_0.64[3];
  new argument3_1 = XMR[registerEq]:1;
  new argument3_1_0 = argument3_1.64[0];
  new argument3_1_1 = argument3_1.64[1];
  new argument3_1_2 = argument3_1.64[2];
  new argument3_1_3 = argument3_1.64[3];
  new argument3_2 = XMR[registerEq]:2;
  new argument3_2_0 = argument3_2.64[0];
  new argument3_2_1 = argument3_2.64[1];
  new argument3_2_2 = argument3_2.64[2];
  new argument3_2_3 = argument3_2.64[3];
  new argument3_3 = XMR[registerEq]:3;
  new argument3_3_0 = argument3_3.64[0];
  new argument3_3_1 = argument3_3.64[1];
  new argument3_3_2 = argument3_3.64[2];
  new argument3_3_3 = argument3_3.64[3];
  new argument3 = 0;
  if (qindex == 0) {
    argument3 = join_64_x4(argument3_0_0, argument3_1_0, argument3_2_0, argument3_3_0);
  }
  if (qindex == 1) {
    argument3 = join_64_x4(argument3_0_1, argument3_1_1, argument3_2_1, argument3_3_1);
  }
  if (qindex == 2) {
    argument3 = join_64_x4(argument3_0_2, argument3_1_2, argument3_2_2, argument3_3_2);
  }
  if (qindex == 3) {
    argument3 = join_64_x4(argument3_0_3, argument3_1_3, argument3_2_3, argument3_3_3);
  }
  CS.XMF = 1;
stage E3:
  MEM.store(address, 0xFFFFFFFF & bytemask, 0, argument3, dri);
\end{lstlisting}}
\newpage

\paragraph{Encoding} Format \textsf{LSU\_XSOC4BO.X} (Section~\ref{sec:format-LSU-XSOC4BO.X}), \verb|steering = 01|\\

\bitfields{ 1/P, 2/00, 2/-- --, 27/upper27, 1/1, 2/01, 5/10111, 4/registerEq, 2/qindex, 2/01, 10/lower10, 6/registerZ }


\paragraph{Syntax} \texttt{xso}\emph{qindex} \emph{upper27\_\discretionary{}{}{}lower10}[\emph{registerZ}] = \emph{registerEq}

\paragraph{Description} The effective address is computed by adding the \emph{upper27\_\discretionary{}{}{}lower10} to the \emph{registerZ}.
The \emph{registerEq} is stored into the memory octuple word at the effective address.


\paragraph{Modifier(s)}
\noindent\begin{itemize}
\item qindex (cf. qindex in section \ref{par:modifier-qindex} on page \pageref{par:modifier-qindex})
\end{itemize}

\paragraph{Execution} {Reservation table \textsf{LSU\_MEMW\_ACCR.X} (Section~\ref{sec:reservation-LSU-MEMW-ACCR.X})}
~
{\begin{lstlisting}
stage ID:
  new qindex = _ZX_2(qindex);
  new offset = _SX_37(upper27_lower10);
stage RR:
  new bytemask = 0xFFFFFFFF;
  new dri = registerEq << 2;
  new base = GPR[registerZ];
  new address = base + offset;
stage E1:
  new argument3_0 = XMR[registerEq]:0;
  new argument3_0_0 = argument3_0.64[0];
  new argument3_0_1 = argument3_0.64[1];
  new argument3_0_2 = argument3_0.64[2];
  new argument3_0_3 = argument3_0.64[3];
  new argument3_1 = XMR[registerEq]:1;
  new argument3_1_0 = argument3_1.64[0];
  new argument3_1_1 = argument3_1.64[1];
  new argument3_1_2 = argument3_1.64[2];
  new argument3_1_3 = argument3_1.64[3];
  new argument3_2 = XMR[registerEq]:2;
  new argument3_2_0 = argument3_2.64[0];
  new argument3_2_1 = argument3_2.64[1];
  new argument3_2_2 = argument3_2.64[2];
  new argument3_2_3 = argument3_2.64[3];
  new argument3_3 = XMR[registerEq]:3;
  new argument3_3_0 = argument3_3.64[0];
  new argument3_3_1 = argument3_3.64[1];
  new argument3_3_2 = argument3_3.64[2];
  new argument3_3_3 = argument3_3.64[3];
  new argument3 = 0;
  if (qindex == 0) {
    argument3 = join_64_x4(argument3_0_0, argument3_1_0, argument3_2_0, argument3_3_0);
  }
  if (qindex == 1) {
    argument3 = join_64_x4(argument3_0_1, argument3_1_1, argument3_2_1, argument3_3_1);
  }
  if (qindex == 2) {
    argument3 = join_64_x4(argument3_0_2, argument3_1_2, argument3_2_2, argument3_3_2);
  }
  if (qindex == 3) {
    argument3 = join_64_x4(argument3_0_3, argument3_1_3, argument3_2_3, argument3_3_3);
  }
  CS.XMF = 1;
stage E3:
  MEM.store(address, 0xFFFFFFFF & bytemask, 0, argument3, dri);
\end{lstlisting}}
\newpage

\paragraph{Encoding} Format \textsf{LSU\_XSOC4BO.Y} (Section~\ref{sec:format-LSU-XSOC4BO.Y}), \verb|steering = 01|\\

\bitfields{ 1/P, 2/00, 2/-- --, 27/extend27, 1/1, 2/00, 2/-- --, 27/upper27, 1/1, 2/01, 5/10111, 4/registerEq, 2/qindex, 2/01, 10/lower10, 6/registerZ }


\paragraph{Syntax} \texttt{xso}\emph{qindex} \emph{extend27\_\discretionary{}{}{}upper27\_\discretionary{}{}{}lower10}[\emph{registerZ}] = \emph{registerEq}

\paragraph{Description} The effective address is computed by adding the \emph{extend27\_\discretionary{}{}{}upper27\_\discretionary{}{}{}lower10} to the \emph{registerZ}.
The \emph{registerEq} is stored into the memory octuple word at the effective address.


\paragraph{Modifier(s)}
\noindent\begin{itemize}
\item qindex (cf. qindex in section \ref{par:modifier-qindex} on page \pageref{par:modifier-qindex})
\end{itemize}

\paragraph{Execution} {Reservation table \textsf{LSU\_MEMW\_ACCR.Y} (Section~\ref{sec:reservation-LSU-MEMW-ACCR.Y})}
~
{\begin{lstlisting}
stage ID:
  new qindex = _ZX_2(qindex);
  new offset = _WX_64(extend27_upper27_lower10);
stage RR:
  new bytemask = 0xFFFFFFFF;
  new dri = registerEq << 2;
  new base = GPR[registerZ];
  new address = base + offset;
stage E1:
  new argument3_0 = XMR[registerEq]:0;
  new argument3_0_0 = argument3_0.64[0];
  new argument3_0_1 = argument3_0.64[1];
  new argument3_0_2 = argument3_0.64[2];
  new argument3_0_3 = argument3_0.64[3];
  new argument3_1 = XMR[registerEq]:1;
  new argument3_1_0 = argument3_1.64[0];
  new argument3_1_1 = argument3_1.64[1];
  new argument3_1_2 = argument3_1.64[2];
  new argument3_1_3 = argument3_1.64[3];
  new argument3_2 = XMR[registerEq]:2;
  new argument3_2_0 = argument3_2.64[0];
  new argument3_2_1 = argument3_2.64[1];
  new argument3_2_2 = argument3_2.64[2];
  new argument3_2_3 = argument3_2.64[3];
  new argument3_3 = XMR[registerEq]:3;
  new argument3_3_0 = argument3_3.64[0];
  new argument3_3_1 = argument3_3.64[1];
  new argument3_3_2 = argument3_3.64[2];
  new argument3_3_3 = argument3_3.64[3];
  new argument3 = 0;
  if (qindex == 0) {
    argument3 = join_64_x4(argument3_0_0, argument3_1_0, argument3_2_0, argument3_3_0);
  }
  if (qindex == 1) {
    argument3 = join_64_x4(argument3_0_1, argument3_1_1, argument3_2_1, argument3_3_1);
  }
  if (qindex == 2) {
    argument3 = join_64_x4(argument3_0_2, argument3_1_2, argument3_2_2, argument3_3_2);
  }
  if (qindex == 3) {
    argument3 = join_64_x4(argument3_0_3, argument3_1_3, argument3_2_3, argument3_3_3);
  }
  CS.XMF = 1;
stage E3:
  MEM.store(address, 0xFFFFFFFF & bytemask, 0, argument3, dri);
\end{lstlisting}}
\newpage

\paragraph{Encoding} Format \textsf{LSU\_XSOC4BI} (Section~\ref{sec:format-LSU-XSOC4BI}), \verb|steering = 01|\\

\bitfields{ 1/P, 2/01, 5/10111, 4/registerEq, 2/qindex, 2/11, 3/111, 1/--, 6/registerY, 6/registerZ }


\paragraph{Syntax} \texttt{xso}\emph{qindex} \emph{registerY}[\emph{registerZ}] = \emph{registerEq}

\paragraph{Description} The effective address is computed by adding the \emph{registerZ} to the \emph{registerY} and scaled by the \emph{qindex}.
The \emph{registerEq} is stored into the memory octuple word at the effective address.


\paragraph{Modifier(s)}
\noindent\begin{itemize}
\item qindex (cf. qindex in section \ref{par:modifier-qindex} on page \pageref{par:modifier-qindex})
\end{itemize}

\paragraph{Execution} {Reservation table \textsf{LSU\_MEMW\_ACCR} (Section~\ref{sec:reservation-LSU-MEMW-ACCR})}
~
{\begin{lstlisting}
stage ID:
  new qindex = _ZX_2(qindex);
stage RR:
  new bytemask = 0xFFFFFFFF;
  new doscale = 0;
  new dri = registerEq >> 2;
  new index = GPR[registerY];
  new base = GPR[registerZ];
  new scaling = doscale ? 32 : 1;
  new address = base + (index * scaling);
stage E1:
  new argument3_0 = XMR[registerEq]:0;
  new argument3_0_0 = argument3_0.64[0];
  new argument3_0_1 = argument3_0.64[1];
  new argument3_0_2 = argument3_0.64[2];
  new argument3_0_3 = argument3_0.64[3];
  new argument3_1 = XMR[registerEq]:1;
  new argument3_1_0 = argument3_1.64[0];
  new argument3_1_1 = argument3_1.64[1];
  new argument3_1_2 = argument3_1.64[2];
  new argument3_1_3 = argument3_1.64[3];
  new argument3_2 = XMR[registerEq]:2;
  new argument3_2_0 = argument3_2.64[0];
  new argument3_2_1 = argument3_2.64[1];
  new argument3_2_2 = argument3_2.64[2];
  new argument3_2_3 = argument3_2.64[3];
  new argument3_3 = XMR[registerEq]:3;
  new argument3_3_0 = argument3_3.64[0];
  new argument3_3_1 = argument3_3.64[1];
  new argument3_3_2 = argument3_3.64[2];
  new argument3_3_3 = argument3_3.64[3];
  new argument3 = 0;
  if (qindex == 0) {
    argument3 = join_64_x4(argument3_0_0, argument3_1_0, argument3_2_0, argument3_3_0);
  }
  if (qindex == 1) {
    argument3 = join_64_x4(argument3_0_1, argument3_1_1, argument3_2_1, argument3_3_1);
  }
  if (qindex == 2) {
    argument3 = join_64_x4(argument3_0_2, argument3_1_2, argument3_2_2, argument3_3_2);
  }
  if (qindex == 3) {
    argument3 = join_64_x4(argument3_0_3, argument3_1_3, argument3_2_3, argument3_3_3);
  }
  CS.XMF = 1;
stage E3:
  MEM.store(address, 0xFFFFFFFF & bytemask, 0, argument3, dri);
\end{lstlisting}}
\newpage
\subsection{ALU Instructions}
\label{sec:ALU-instructions}

\subsubsection[{\small ABDBO: Absolute Difference Byte Octuple}]{ABDBO\hfill Absolute Difference Byte Octuple}
\label{instr:ABDBO}\index{abdbo}
 
\paragraph{Encoding} Format \textsf{ALU\_BOWRRSE} (Section~\ref{sec:format-ALU-BOWRRSE}), \verb|exucode2 = 1100|\\

\bitfields{ 1/P, 2/11, 1/1, 4/1100, 5/registerWe, 1/0, 2/10, 3/010, 1/1, 5/registerYe, 1/--, 5/registerZe, 1/1 }


\paragraph{Syntax} \texttt{abdbo} \emph{registerWe} = \emph{registerZe}, \emph{registerYe}

\paragraph{Description} The \emph{registerZe} and the \emph{registerYe} are interpreted as eight bytes packed into 64 bits. The \emph{registerZe} bytes are subtracted from the \emph{registerYe} bytes. The absolute value of the eight resulting bytes are packed and stored into the \emph{registerWe}.


\paragraph{Execution} {Reservation table \textsf{ALU\_LITE} (Section~\ref{sec:reservation-ALU-LITE})}
~
{\begin{lstlisting}
stage RR:
  new argument3 = GPR[registerYe<<1];
  new argument2 = GPR[registerZe<<1];
  for (I = 0; I < 8; I += 1) {
    result1.8[I] = _ABS(_SX_8(argument3.8[I]) - _SX_8(argument2.8[I]))
  };
stage E1:
  GPR[registerWe<<1] = result1;
\end{lstlisting}}
\newpage

\paragraph{Encoding} Format \textsf{ALU\_BOWRRSO} (Section~\ref{sec:format-ALU-BOWRRSO}), \verb|exucode2 = 1100|\\

\bitfields{ 1/P, 2/11, 1/1, 4/1100, 5/registerWo, 1/1, 2/10, 3/010, 1/1, 5/registerYo, 1/--, 5/registerZo, 1/1 }


\paragraph{Syntax} \texttt{abdbo} \emph{registerWo} = \emph{registerZo}, \emph{registerYo}

\paragraph{Description} The \emph{registerZo} and the \emph{registerYo} are interpreted as eight bytes packed into 64 bits. The \emph{registerZo} bytes are subtracted from the \emph{registerYo} bytes. The absolute value of the eight resulting bytes are packed and stored into the \emph{registerWo}.


\paragraph{Execution} {Reservation table \textsf{ALU\_LITE} (Section~\ref{sec:reservation-ALU-LITE})}
~
{\begin{lstlisting}
stage RR:
  new argument3 = GPR[(registerYo<<1)+1];
  new argument2 = GPR[(registerZo<<1)+1];
  for (I = 0; I < 8; I += 1) {
    result1.8[I] = _ABS(_SX_8(argument3.8[I]) - _SX_8(argument2.8[I]))
  };
stage E1:
  GPR[(registerWo<<1)+1] = result1;
\end{lstlisting}}
\newpage

\paragraph{Encoding} Format \textsf{ALU\_BOWRRSE.M} (Section~\ref{sec:format-ALU-BOWRRSE.M}), \verb|exucode2 = 1100|\\

\bitfields{ 1/P, 2/00, 2/-- --, 27/upper27, 1/1, 2/11, 1/1, 4/1100, 5/registerWe, 1/0, 2/10, 3/010, 1/1, 1/splat32, 5/lower5, 5/registerZe, 1/1 }


\paragraph{Syntax} \texttt{abdbo} \emph{registerWe} = \emph{registerZe}, \emph{upper27\_\discretionary{}{}{}lower5}\emph{splat32}

\paragraph{Description} The \emph{registerZe} and the \emph{upper27\_\discretionary{}{}{}lower5} are interpreted as eight bytes packed into 64 bits. The \emph{registerZe} bytes are subtracted from the \emph{upper27\_\discretionary{}{}{}lower5} bytes. The absolute value of the eight resulting bytes are packed and stored into the \emph{registerWe}.


\paragraph{Modifier(s)}
\noindent\begin{itemize}
\item splat32 (cf. splat32 in section \ref{par:modifier-splat32} on page \pageref{par:modifier-splat32})
\end{itemize}

\paragraph{Execution} {Reservation table \textsf{ALU\_LITE.X} (Section~\ref{sec:reservation-ALU-LITE.X})}
~
{\begin{lstlisting}
stage ID:
  new splat32 = _ZX_1(splat32);
stage RR:
  new argument3 = splat32 ? (_WX_32(upper27_lower5) << 32) | _ZX_32(_WX_32(upper27_lower5)) : _SX_32(_WX_32(upper27_lower5));
  new argument2 = GPR[registerZe<<1];
  for (I = 0; I < 8; I += 1) {
    result1.8[I] = _ABS(_SX_8(argument3.8[I]) - _SX_8(argument2.8[I]))
  };
stage E1:
  GPR[registerWe<<1] = result1;
\end{lstlisting}}
\newpage

\paragraph{Encoding} Format \textsf{ALU\_BOWRRSO.M} (Section~\ref{sec:format-ALU-BOWRRSO.M}), \verb|exucode2 = 1100|\\

\bitfields{ 1/P, 2/00, 2/-- --, 27/upper27, 1/1, 2/11, 1/1, 4/1100, 5/registerWo, 1/1, 2/10, 3/010, 1/1, 1/splat32, 5/lower5, 5/registerZo, 1/1 }


\paragraph{Syntax} \texttt{abdbo} \emph{registerWo} = \emph{registerZo}, \emph{upper27\_\discretionary{}{}{}lower5}\emph{splat32}

\paragraph{Description} The \emph{registerZo} and the \emph{upper27\_\discretionary{}{}{}lower5} are interpreted as eight bytes packed into 64 bits. The \emph{registerZo} bytes are subtracted from the \emph{upper27\_\discretionary{}{}{}lower5} bytes. The absolute value of the eight resulting bytes are packed and stored into the \emph{registerWo}.


\paragraph{Modifier(s)}
\noindent\begin{itemize}
\item splat32 (cf. splat32 in section \ref{par:modifier-splat32} on page \pageref{par:modifier-splat32})
\end{itemize}

\paragraph{Execution} {Reservation table \textsf{ALU\_LITE.X} (Section~\ref{sec:reservation-ALU-LITE.X})}
~
{\begin{lstlisting}
stage ID:
  new splat32 = _ZX_1(splat32);
stage RR:
  new argument3 = splat32 ? (_WX_32(upper27_lower5) << 32) | _ZX_32(_WX_32(upper27_lower5)) : _SX_32(_WX_32(upper27_lower5));
  new argument2 = GPR[(registerZo<<1)+1];
  for (I = 0; I < 8; I += 1) {
    result1.8[I] = _ABS(_SX_8(argument3.8[I]) - _SX_8(argument2.8[I]))
  };
stage E1:
  GPR[(registerWo<<1)+1] = result1;
\end{lstlisting}}
\newpage
\subsubsection[{\small ABDD: Absolute Difference Double Word}]{ABDD\hfill Absolute Difference Double Word}
\label{instr:ABDD}\index{abdd}
 
\paragraph{Encoding} Format \textsf{ALU\_DWRRS} (Section~\ref{sec:format-ALU-DWRRS}), \verb|exucode2 = 1100|\\

\bitfields{ 1/P, 2/11, 1/0, 4/1100, 6/registerW, 2/10, 3/010, 1/0, 6/registerY, 6/registerZ }


\paragraph{Syntax} \texttt{abdd} \emph{registerW} = \emph{registerZ}, \emph{registerY}

\paragraph{Description} The \emph{registerZ} is subtracted from the \emph{registerY}. The absolute value of the result is stored into the \emph{registerW}.


\paragraph{Execution} {Reservation table \textsf{ALU\_TINY} (Section~\ref{sec:reservation-ALU-TINY})}
~
{\begin{lstlisting}
stage RR:
  new argument3 = GPR[registerY];
  new argument2 = GPR[registerZ];
  new result1 = _ABS(_SX_64(argument3) - _SX_64(argument2));
stage E1:
  GPR[registerW] = result1;
\end{lstlisting}}
\newpage

\paragraph{Encoding} Format \textsf{ALU\_DWRRS.M} (Section~\ref{sec:format-ALU-DWRRS.M}), \verb|exucode2 = 1100|\\

\bitfields{ 1/P, 2/00, 2/-- --, 27/upper27, 1/1, 2/11, 1/0, 4/1100, 6/registerW, 2/10, 3/010, 1/0, 1/splat32, 5/lower5, 6/registerZ }


\paragraph{Syntax} \texttt{abdd} \emph{registerW} = \emph{registerZ}, \emph{upper27\_\discretionary{}{}{}lower5}\emph{splat32}

\paragraph{Description} The \emph{registerZ} is subtracted from the \emph{upper27\_\discretionary{}{}{}lower5}. The absolute value of the result is stored into the \emph{registerW}.


\paragraph{Modifier(s)}
\noindent\begin{itemize}
\item splat32 (cf. splat32 in section \ref{par:modifier-splat32} on page \pageref{par:modifier-splat32})
\end{itemize}

\paragraph{Execution} {Reservation table \textsf{ALU\_TINY.X} (Section~\ref{sec:reservation-ALU-TINY.X})}
~
{\begin{lstlisting}
stage ID:
  new splat32 = _ZX_1(splat32);
stage RR:
  new argument3 = splat32 ? (_WX_32(upper27_lower5) << 32) | _ZX_32(_WX_32(upper27_lower5)) : _SX_32(_WX_32(upper27_lower5));
  new argument2 = GPR[registerZ];
  new result1 = _ABS(_SX_64(argument3) - _SX_64(argument2));
stage E1:
  GPR[registerW] = result1;
\end{lstlisting}}
\newpage
\subsubsection[{\small ABDHQ: Absolute Difference Half Word Quadruple}]{ABDHQ\hfill Absolute Difference Half Word Quadruple}
\label{instr:ABDHQ}\index{abdhq}
 
\paragraph{Encoding} Format \textsf{ALU\_HQWRRSE} (Section~\ref{sec:format-ALU-HQWRRSE}), \verb|exucode2 = 1100|\\

\bitfields{ 1/P, 2/11, 1/1, 4/1100, 5/registerWe, 1/0, 2/10, 3/010, 1/1, 5/registerYe, 1/--, 5/registerZe, 1/0 }


\paragraph{Syntax} \texttt{abdhq} \emph{registerWe} = \emph{registerZe}, \emph{registerYe}

\paragraph{Description} The \emph{registerZe} and the \emph{registerYe} are interpreted as four half words packed into 64 bits. The \emph{registerZe} half words are subtracted from the \emph{registerYe} half words. The absolute value of the four resulting half words are packed and stored into the \emph{registerWe}.


\paragraph{Execution} {Reservation table \textsf{ALU\_LITE} (Section~\ref{sec:reservation-ALU-LITE})}
~
{\begin{lstlisting}
stage RR:
  new argument3 = GPR[registerYe<<1];
  new argument2 = GPR[registerZe<<1];
  for (I = 0; I < 4; I += 1) {
    result1.16[I] = _ABS(_SX_16(argument3.16[I]) - _SX_16(argument2.16[I]))
  };
stage E1:
  GPR[registerWe<<1] = result1;
\end{lstlisting}}
\newpage

\paragraph{Encoding} Format \textsf{ALU\_HQWRRSO} (Section~\ref{sec:format-ALU-HQWRRSO}), \verb|exucode2 = 1100|\\

\bitfields{ 1/P, 2/11, 1/1, 4/1100, 5/registerWo, 1/1, 2/10, 3/010, 1/1, 5/registerYo, 1/--, 5/registerZo, 1/0 }


\paragraph{Syntax} \texttt{abdhq} \emph{registerWo} = \emph{registerZo}, \emph{registerYo}

\paragraph{Description} The \emph{registerZo} and the \emph{registerYo} are interpreted as four half words packed into 64 bits. The \emph{registerZo} half words are subtracted from the \emph{registerYo} half words. The absolute value of the four resulting half words are packed and stored into the \emph{registerWo}.


\paragraph{Execution} {Reservation table \textsf{ALU\_LITE} (Section~\ref{sec:reservation-ALU-LITE})}
~
{\begin{lstlisting}
stage RR:
  new argument3 = GPR[(registerYo<<1)+1];
  new argument2 = GPR[(registerZo<<1)+1];
  for (I = 0; I < 4; I += 1) {
    result1.16[I] = _ABS(_SX_16(argument3.16[I]) - _SX_16(argument2.16[I]))
  };
stage E1:
  GPR[(registerWo<<1)+1] = result1;
\end{lstlisting}}
\newpage

\paragraph{Encoding} Format \textsf{ALU\_HQWRRSE.M} (Section~\ref{sec:format-ALU-HQWRRSE.M}), \verb|exucode2 = 1100|\\

\bitfields{ 1/P, 2/00, 2/-- --, 27/upper27, 1/1, 2/11, 1/1, 4/1100, 5/registerWe, 1/0, 2/10, 3/010, 1/1, 1/splat32, 5/lower5, 5/registerZe, 1/0 }


\paragraph{Syntax} \texttt{abdhq} \emph{registerWe} = \emph{registerZe}, \emph{upper27\_\discretionary{}{}{}lower5}\emph{splat32}

\paragraph{Description} The \emph{registerZe} and the \emph{upper27\_\discretionary{}{}{}lower5} are interpreted as four half words packed into 64 bits. The \emph{registerZe} half words are subtracted from the \emph{upper27\_\discretionary{}{}{}lower5} half words. The absolute value of the four resulting half words are packed and stored into the \emph{registerWe}.


\paragraph{Modifier(s)}
\noindent\begin{itemize}
\item splat32 (cf. splat32 in section \ref{par:modifier-splat32} on page \pageref{par:modifier-splat32})
\end{itemize}

\paragraph{Execution} {Reservation table \textsf{ALU\_LITE.X} (Section~\ref{sec:reservation-ALU-LITE.X})}
~
{\begin{lstlisting}
stage ID:
  new splat32 = _ZX_1(splat32);
stage RR:
  new argument3 = splat32 ? (_WX_32(upper27_lower5) << 32) | _ZX_32(_WX_32(upper27_lower5)) : _SX_32(_WX_32(upper27_lower5));
  new argument2 = GPR[registerZe<<1];
  for (I = 0; I < 4; I += 1) {
    result1.16[I] = _ABS(_SX_16(argument3.16[I]) - _SX_16(argument2.16[I]))
  };
stage E1:
  GPR[registerWe<<1] = result1;
\end{lstlisting}}
\newpage

\paragraph{Encoding} Format \textsf{ALU\_HQWRRSO.M} (Section~\ref{sec:format-ALU-HQWRRSO.M}), \verb|exucode2 = 1100|\\

\bitfields{ 1/P, 2/00, 2/-- --, 27/upper27, 1/1, 2/11, 1/1, 4/1100, 5/registerWo, 1/1, 2/10, 3/010, 1/1, 1/splat32, 5/lower5, 5/registerZo, 1/0 }


\paragraph{Syntax} \texttt{abdhq} \emph{registerWo} = \emph{registerZo}, \emph{upper27\_\discretionary{}{}{}lower5}\emph{splat32}

\paragraph{Description} The \emph{registerZo} and the \emph{upper27\_\discretionary{}{}{}lower5} are interpreted as four half words packed into 64 bits. The \emph{registerZo} half words are subtracted from the \emph{upper27\_\discretionary{}{}{}lower5} half words. The absolute value of the four resulting half words are packed and stored into the \emph{registerWo}.


\paragraph{Modifier(s)}
\noindent\begin{itemize}
\item splat32 (cf. splat32 in section \ref{par:modifier-splat32} on page \pageref{par:modifier-splat32})
\end{itemize}

\paragraph{Execution} {Reservation table \textsf{ALU\_LITE.X} (Section~\ref{sec:reservation-ALU-LITE.X})}
~
{\begin{lstlisting}
stage ID:
  new splat32 = _ZX_1(splat32);
stage RR:
  new argument3 = splat32 ? (_WX_32(upper27_lower5) << 32) | _ZX_32(_WX_32(upper27_lower5)) : _SX_32(_WX_32(upper27_lower5));
  new argument2 = GPR[(registerZo<<1)+1];
  for (I = 0; I < 4; I += 1) {
    result1.16[I] = _ABS(_SX_16(argument3.16[I]) - _SX_16(argument2.16[I]))
  };
stage E1:
  GPR[(registerWo<<1)+1] = result1;
\end{lstlisting}}
\newpage
\subsubsection[{\small ABDSBO: Absolute Difference Saturated Byte Octuple}]{ABDSBO\hfill Absolute Difference Saturated Byte Octuple}
\label{instr:ABDSBO}\index{abdsbo}
 
\paragraph{Encoding} Format \textsf{ALU\_BOWRRSE} (Section~\ref{sec:format-ALU-BOWRRSE}), \verb|exucode2 = 1110|\\

\bitfields{ 1/P, 2/11, 1/1, 4/1110, 5/registerWe, 1/0, 2/10, 3/010, 1/1, 5/registerYe, 1/--, 5/registerZe, 1/1 }


\paragraph{Syntax} \texttt{abdsbo} \emph{registerWe} = \emph{registerZe}, \emph{registerYe}

\paragraph{Description} The \emph{registerZe} and the \emph{registerYe} are interpreted as eight bytes packed into 64 bits. The word quadruple \emph{registerZe} is subtracted from the word quadruple \emph{registerYe}. The absolute values of the eight resulting bytes are saturated and stored into the \emph{registerWe}.


\paragraph{Execution} {Reservation table \textsf{ALU\_LITE} (Section~\ref{sec:reservation-ALU-LITE})}
~
{\begin{lstlisting}
stage RR:
  new argument3 = GPR[registerYe<<1];
  new argument2 = GPR[registerZe<<1];
  for (I = 0; I < 8; I += 1) {
    result1.8[I] = _SAT_8(_ABS(_SX_8(argument3.8[I]) - _SX_8(argument2.8[I])))
  };
stage E1:
  GPR[registerWe<<1] = result1;
\end{lstlisting}}
\newpage

\paragraph{Encoding} Format \textsf{ALU\_BOWRRSO} (Section~\ref{sec:format-ALU-BOWRRSO}), \verb|exucode2 = 1110|\\

\bitfields{ 1/P, 2/11, 1/1, 4/1110, 5/registerWo, 1/1, 2/10, 3/010, 1/1, 5/registerYo, 1/--, 5/registerZo, 1/1 }


\paragraph{Syntax} \texttt{abdsbo} \emph{registerWo} = \emph{registerZo}, \emph{registerYo}

\paragraph{Description} The \emph{registerZo} and the \emph{registerYo} are interpreted as eight bytes packed into 64 bits. The word quadruple \emph{registerZo} is subtracted from the word quadruple \emph{registerYo}. The absolute values of the eight resulting bytes are saturated and stored into the \emph{registerWo}.


\paragraph{Execution} {Reservation table \textsf{ALU\_LITE} (Section~\ref{sec:reservation-ALU-LITE})}
~
{\begin{lstlisting}
stage RR:
  new argument3 = GPR[(registerYo<<1)+1];
  new argument2 = GPR[(registerZo<<1)+1];
  for (I = 0; I < 8; I += 1) {
    result1.8[I] = _SAT_8(_ABS(_SX_8(argument3.8[I]) - _SX_8(argument2.8[I])))
  };
stage E1:
  GPR[(registerWo<<1)+1] = result1;
\end{lstlisting}}
\newpage

\paragraph{Encoding} Format \textsf{ALU\_BOWRRSE.M} (Section~\ref{sec:format-ALU-BOWRRSE.M}), \verb|exucode2 = 1110|\\

\bitfields{ 1/P, 2/00, 2/-- --, 27/upper27, 1/1, 2/11, 1/1, 4/1110, 5/registerWe, 1/0, 2/10, 3/010, 1/1, 1/splat32, 5/lower5, 5/registerZe, 1/1 }


\paragraph{Syntax} \texttt{abdsbo} \emph{registerWe} = \emph{registerZe}, \emph{upper27\_\discretionary{}{}{}lower5}\emph{splat32}

\paragraph{Description} The \emph{registerZe} and the \emph{upper27\_\discretionary{}{}{}lower5} are interpreted as eight bytes packed into 64 bits. The word quadruple \emph{registerZe} is subtracted from the word quadruple \emph{upper27\_\discretionary{}{}{}lower5}. The absolute values of the eight resulting bytes are saturated and stored into the \emph{registerWe}.


\paragraph{Modifier(s)}
\noindent\begin{itemize}
\item splat32 (cf. splat32 in section \ref{par:modifier-splat32} on page \pageref{par:modifier-splat32})
\end{itemize}

\paragraph{Execution} {Reservation table \textsf{ALU\_LITE.X} (Section~\ref{sec:reservation-ALU-LITE.X})}
~
{\begin{lstlisting}
stage ID:
  new splat32 = _ZX_1(splat32);
stage RR:
  new argument3 = splat32 ? (_WX_32(upper27_lower5) << 32) | _ZX_32(_WX_32(upper27_lower5)) : _SX_32(_WX_32(upper27_lower5));
  new argument2 = GPR[registerZe<<1];
  for (I = 0; I < 8; I += 1) {
    result1.8[I] = _SAT_8(_ABS(_SX_8(argument3.8[I]) - _SX_8(argument2.8[I])))
  };
stage E1:
  GPR[registerWe<<1] = result1;
\end{lstlisting}}
\newpage

\paragraph{Encoding} Format \textsf{ALU\_BOWRRSO.M} (Section~\ref{sec:format-ALU-BOWRRSO.M}), \verb|exucode2 = 1110|\\

\bitfields{ 1/P, 2/00, 2/-- --, 27/upper27, 1/1, 2/11, 1/1, 4/1110, 5/registerWo, 1/1, 2/10, 3/010, 1/1, 1/splat32, 5/lower5, 5/registerZo, 1/1 }


\paragraph{Syntax} \texttt{abdsbo} \emph{registerWo} = \emph{registerZo}, \emph{upper27\_\discretionary{}{}{}lower5}\emph{splat32}

\paragraph{Description} The \emph{registerZo} and the \emph{upper27\_\discretionary{}{}{}lower5} are interpreted as eight bytes packed into 64 bits. The word quadruple \emph{registerZo} is subtracted from the word quadruple \emph{upper27\_\discretionary{}{}{}lower5}. The absolute values of the eight resulting bytes are saturated and stored into the \emph{registerWo}.


\paragraph{Modifier(s)}
\noindent\begin{itemize}
\item splat32 (cf. splat32 in section \ref{par:modifier-splat32} on page \pageref{par:modifier-splat32})
\end{itemize}

\paragraph{Execution} {Reservation table \textsf{ALU\_LITE.X} (Section~\ref{sec:reservation-ALU-LITE.X})}
~
{\begin{lstlisting}
stage ID:
  new splat32 = _ZX_1(splat32);
stage RR:
  new argument3 = splat32 ? (_WX_32(upper27_lower5) << 32) | _ZX_32(_WX_32(upper27_lower5)) : _SX_32(_WX_32(upper27_lower5));
  new argument2 = GPR[(registerZo<<1)+1];
  for (I = 0; I < 8; I += 1) {
    result1.8[I] = _SAT_8(_ABS(_SX_8(argument3.8[I]) - _SX_8(argument2.8[I])))
  };
stage E1:
  GPR[(registerWo<<1)+1] = result1;
\end{lstlisting}}
\newpage
\subsubsection[{\small ABDSD: Absolute Difference Saturated Double Word}]{ABDSD\hfill Absolute Difference Saturated Double Word}
\label{instr:ABDSD}\index{abdsd}
 
\paragraph{Encoding} Format \textsf{ALU\_DWRRS} (Section~\ref{sec:format-ALU-DWRRS}), \verb|exucode2 = 1110|\\

\bitfields{ 1/P, 2/11, 1/0, 4/1110, 6/registerW, 2/10, 3/010, 1/0, 6/registerY, 6/registerZ }


\paragraph{Syntax} \texttt{abdsd} \emph{registerW} = \emph{registerZ}, \emph{registerY}

\paragraph{Description} The double word \emph{registerZ} is subtracted from the double word \emph{registerY}. The absolute value of the resulting double word is saturated and stored into the \emph{registerW}.


\paragraph{Execution} {Reservation table \textsf{ALU\_TINY} (Section~\ref{sec:reservation-ALU-TINY})}
~
{\begin{lstlisting}
stage RR:
  new argument3 = GPR[registerY];
  new argument2 = GPR[registerZ];
  new result1 = _SAT_64(_ABS(_SX_64(argument3) - _SX_64(argument2)));
stage E1:
  GPR[registerW] = result1;
\end{lstlisting}}
\newpage

\paragraph{Encoding} Format \textsf{ALU\_DWRRS.M} (Section~\ref{sec:format-ALU-DWRRS.M}), \verb|exucode2 = 1110|\\

\bitfields{ 1/P, 2/00, 2/-- --, 27/upper27, 1/1, 2/11, 1/0, 4/1110, 6/registerW, 2/10, 3/010, 1/0, 1/splat32, 5/lower5, 6/registerZ }


\paragraph{Syntax} \texttt{abdsd} \emph{registerW} = \emph{registerZ}, \emph{upper27\_\discretionary{}{}{}lower5}\emph{splat32}

\paragraph{Description} The double word \emph{registerZ} is subtracted from the double word \emph{upper27\_\discretionary{}{}{}lower5}. The absolute value of the resulting double word is saturated and stored into the \emph{registerW}.


\paragraph{Modifier(s)}
\noindent\begin{itemize}
\item splat32 (cf. splat32 in section \ref{par:modifier-splat32} on page \pageref{par:modifier-splat32})
\end{itemize}

\paragraph{Execution} {Reservation table \textsf{ALU\_TINY.X} (Section~\ref{sec:reservation-ALU-TINY.X})}
~
{\begin{lstlisting}
stage ID:
  new splat32 = _ZX_1(splat32);
stage RR:
  new argument3 = splat32 ? (_WX_32(upper27_lower5) << 32) | _ZX_32(_WX_32(upper27_lower5)) : _SX_32(_WX_32(upper27_lower5));
  new argument2 = GPR[registerZ];
  new result1 = _SAT_64(_ABS(_SX_64(argument3) - _SX_64(argument2)));
stage E1:
  GPR[registerW] = result1;
\end{lstlisting}}
\newpage
\subsubsection[{\small ABDSHQ: Absolute Difference Saturated Half Word Quadruple}]{ABDSHQ\hfill Absolute Difference Saturated Half Word Quadruple}
\label{instr:ABDSHQ}\index{abdshq}
 
\paragraph{Encoding} Format \textsf{ALU\_HQWRRSE} (Section~\ref{sec:format-ALU-HQWRRSE}), \verb|exucode2 = 1110|\\

\bitfields{ 1/P, 2/11, 1/1, 4/1110, 5/registerWe, 1/0, 2/10, 3/010, 1/1, 5/registerYe, 1/--, 5/registerZe, 1/0 }


\paragraph{Syntax} \texttt{abdshq} \emph{registerWe} = \emph{registerZe}, \emph{registerYe}

\paragraph{Description} The \emph{registerZe} and the \emph{registerYe} are interpreted as four half words packed into 64 bits. The \emph{registerZe} half words are subtracted from the \emph{registerYe} half words. The absolute values of the four resulting half words are saturated, packed and stored into the \emph{registerWe}.


\paragraph{Execution} {Reservation table \textsf{ALU\_LITE} (Section~\ref{sec:reservation-ALU-LITE})}
~
{\begin{lstlisting}
stage RR:
  new argument3 = GPR[registerYe<<1];
  new argument2 = GPR[registerZe<<1];
  for (I = 0; I < 8; I += 1) {
    result1.16[I] = _SAT_16(_ABS(_SX_16(argument3.16[I]) - _SX_16(argument2.16[I])))
  };
stage E1:
  GPR[registerWe<<1] = result1;
\end{lstlisting}}
\newpage

\paragraph{Encoding} Format \textsf{ALU\_HQWRRSO} (Section~\ref{sec:format-ALU-HQWRRSO}), \verb|exucode2 = 1110|\\

\bitfields{ 1/P, 2/11, 1/1, 4/1110, 5/registerWo, 1/1, 2/10, 3/010, 1/1, 5/registerYo, 1/--, 5/registerZo, 1/0 }


\paragraph{Syntax} \texttt{abdshq} \emph{registerWo} = \emph{registerZo}, \emph{registerYo}

\paragraph{Description} The \emph{registerZo} and the \emph{registerYo} are interpreted as four half words packed into 64 bits. The \emph{registerZo} half words are subtracted from the \emph{registerYo} half words. The absolute values of the four resulting half words are saturated, packed and stored into the \emph{registerWo}.


\paragraph{Execution} {Reservation table \textsf{ALU\_LITE} (Section~\ref{sec:reservation-ALU-LITE})}
~
{\begin{lstlisting}
stage RR:
  new argument3 = GPR[(registerYo<<1)+1];
  new argument2 = GPR[(registerZo<<1)+1];
  for (I = 0; I < 8; I += 1) {
    result1.16[I] = _SAT_16(_ABS(_SX_16(argument3.16[I]) - _SX_16(argument2.16[I])))
  };
stage E1:
  GPR[(registerWo<<1)+1] = result1;
\end{lstlisting}}
\newpage

\paragraph{Encoding} Format \textsf{ALU\_HQWRRSE.M} (Section~\ref{sec:format-ALU-HQWRRSE.M}), \verb|exucode2 = 1110|\\

\bitfields{ 1/P, 2/00, 2/-- --, 27/upper27, 1/1, 2/11, 1/1, 4/1110, 5/registerWe, 1/0, 2/10, 3/010, 1/1, 1/splat32, 5/lower5, 5/registerZe, 1/0 }


\paragraph{Syntax} \texttt{abdshq} \emph{registerWe} = \emph{registerZe}, \emph{upper27\_\discretionary{}{}{}lower5}\emph{splat32}

\paragraph{Description} The \emph{registerZe} and the \emph{upper27\_\discretionary{}{}{}lower5} are interpreted as four half words packed into 64 bits. The \emph{registerZe} half words are subtracted from the \emph{upper27\_\discretionary{}{}{}lower5} half words. The absolute values of the four resulting half words are saturated, packed and stored into the \emph{registerWe}.


\paragraph{Modifier(s)}
\noindent\begin{itemize}
\item splat32 (cf. splat32 in section \ref{par:modifier-splat32} on page \pageref{par:modifier-splat32})
\end{itemize}

\paragraph{Execution} {Reservation table \textsf{ALU\_LITE.X} (Section~\ref{sec:reservation-ALU-LITE.X})}
~
{\begin{lstlisting}
stage ID:
  new splat32 = _ZX_1(splat32);
stage RR:
  new argument3 = splat32 ? (_WX_32(upper27_lower5) << 32) | _ZX_32(_WX_32(upper27_lower5)) : _SX_32(_WX_32(upper27_lower5));
  new argument2 = GPR[registerZe<<1];
  for (I = 0; I < 8; I += 1) {
    result1.16[I] = _SAT_16(_ABS(_SX_16(argument3.16[I]) - _SX_16(argument2.16[I])))
  };
stage E1:
  GPR[registerWe<<1] = result1;
\end{lstlisting}}
\newpage

\paragraph{Encoding} Format \textsf{ALU\_HQWRRSO.M} (Section~\ref{sec:format-ALU-HQWRRSO.M}), \verb|exucode2 = 1110|\\

\bitfields{ 1/P, 2/00, 2/-- --, 27/upper27, 1/1, 2/11, 1/1, 4/1110, 5/registerWo, 1/1, 2/10, 3/010, 1/1, 1/splat32, 5/lower5, 5/registerZo, 1/0 }


\paragraph{Syntax} \texttt{abdshq} \emph{registerWo} = \emph{registerZo}, \emph{upper27\_\discretionary{}{}{}lower5}\emph{splat32}

\paragraph{Description} The \emph{registerZo} and the \emph{upper27\_\discretionary{}{}{}lower5} are interpreted as four half words packed into 64 bits. The \emph{registerZo} half words are subtracted from the \emph{upper27\_\discretionary{}{}{}lower5} half words. The absolute values of the four resulting half words are saturated, packed and stored into the \emph{registerWo}.


\paragraph{Modifier(s)}
\noindent\begin{itemize}
\item splat32 (cf. splat32 in section \ref{par:modifier-splat32} on page \pageref{par:modifier-splat32})
\end{itemize}

\paragraph{Execution} {Reservation table \textsf{ALU\_LITE.X} (Section~\ref{sec:reservation-ALU-LITE.X})}
~
{\begin{lstlisting}
stage ID:
  new splat32 = _ZX_1(splat32);
stage RR:
  new argument3 = splat32 ? (_WX_32(upper27_lower5) << 32) | _ZX_32(_WX_32(upper27_lower5)) : _SX_32(_WX_32(upper27_lower5));
  new argument2 = GPR[(registerZo<<1)+1];
  for (I = 0; I < 8; I += 1) {
    result1.16[I] = _SAT_16(_ABS(_SX_16(argument3.16[I]) - _SX_16(argument2.16[I])))
  };
stage E1:
  GPR[(registerWo<<1)+1] = result1;
\end{lstlisting}}
\newpage
\subsubsection[{\small ABDSW: Absolute Difference Saturated Word}]{ABDSW\hfill Absolute Difference Saturated Word}
\label{instr:ABDSW}\index{abdsw}
 
\paragraph{Encoding} Format \textsf{ALU\_WWRRS} (Section~\ref{sec:format-ALU-WWRRS}), \verb|exucode2 = 1110|\\

\bitfields{ 1/P, 2/11, 1/0, 4/1110, 6/registerW, 2/11, 3/010, 1/signextw, 6/registerY, 6/registerZ }


\paragraph{Syntax} \texttt{abdsw}\emph{signextw} \emph{registerW} = \emph{registerZ}, \emph{registerY}

\paragraph{Description} The word \emph{registerZ} is subtracted from the word \emph{registerY}. The absolute value of the resulting word is saturated and stored into the \emph{registerW}.


\paragraph{Modifier(s)}
\noindent\begin{itemize}
\item signextw (cf. signextw in section \ref{par:modifier-signextw} on page \pageref{par:modifier-signextw})
\end{itemize}

\paragraph{Execution} {Reservation table \textsf{ALU\_TINY} (Section~\ref{sec:reservation-ALU-TINY})}
~
{\begin{lstlisting}
stage ID:
  new signextw = _ZX_1(signextw);
stage RR:
  new argument3 = GPR[registerY];
  new argument2 = GPR[registerZ];
  new result1 = _SAT_32(_ABS(_SX_32(argument3) - _SX_32(argument2)));
stage E1:
  GPR[registerW] = signextw ? _SX_32(result1) : _ZX_32(result1);
\end{lstlisting}}
\newpage

\paragraph{Encoding} Format \textsf{ALU\_WWRRS.M} (Section~\ref{sec:format-ALU-WWRRS.M}), \verb|exucode2 = 1110|\\

\bitfields{ 1/P, 2/00, 2/-- --, 27/upper27, 1/1, 2/11, 1/0, 4/1110, 6/registerW, 2/11, 3/010, 1/signextw, 1/splat32, 5/lower5, 6/registerZ }


\paragraph{Syntax} \texttt{abdsw}\emph{signextw} \emph{registerW} = \emph{registerZ}, \emph{upper27\_\discretionary{}{}{}lower5}

\paragraph{Description} The word \emph{registerZ} is subtracted from the word \emph{upper27\_\discretionary{}{}{}lower5}. The absolute value of the resulting word is saturated and stored into the \emph{registerW}.


\paragraph{Modifier(s)}
\noindent\begin{itemize}
\item signextw (cf. signextw in section \ref{par:modifier-signextw} on page \pageref{par:modifier-signextw})
\end{itemize}

\paragraph{Execution} {Reservation table \textsf{ALU\_TINY.X} (Section~\ref{sec:reservation-ALU-TINY.X})}
~
{\begin{lstlisting}
stage ID:
  new signextw = _ZX_1(signextw);
stage RR:
  new argument3 = _WX_32(upper27_lower5);
  new argument2 = GPR[registerZ];
  new result1 = _SAT_32(_ABS(_SX_32(argument3) - _SX_32(argument2)));
stage E1:
  GPR[registerW] = signextw ? _SX_32(result1) : _ZX_32(result1);
\end{lstlisting}}
\newpage
\subsubsection[{\small ABDSWP: Absolute Difference Saturated Word Pair}]{ABDSWP\hfill Absolute Difference Saturated Word Pair}
\label{instr:ABDSWP}\index{abdswp}
 
\paragraph{Encoding} Format \textsf{ALU\_WPWRRSE} (Section~\ref{sec:format-ALU-WPWRRSE}), \verb|exucode2 = 1110|\\

\bitfields{ 1/P, 2/11, 1/1, 4/1110, 5/registerWe, 1/0, 2/10, 3/010, 1/0, 5/registerYe, 1/--, 5/registerZe, 1/1 }


\paragraph{Syntax} \texttt{abdswp} \emph{registerWe} = \emph{registerZe}, \emph{registerYe}

\paragraph{Description} The \emph{registerZe} and the \emph{registerYe} are interpreted as two words packed into 64 bits. The \emph{registerZe} words are subtracted from the \emph{registerYe} words. The absolute values of the two resulting words are saturated, packed and stored into the \emph{registerWe}.


\paragraph{Execution} {Reservation table \textsf{ALU\_LITE} (Section~\ref{sec:reservation-ALU-LITE})}
~
{\begin{lstlisting}
stage RR:
  new argument3 = GPR[registerYe<<1];
  new argument2 = GPR[registerZe<<1];
  for (I = 0; I < 2; I += 1) {
    result1.32[I] = _SAT_32(_ABS(_SX_32(argument3.32[I]) - _SX_32(argument2.32[I])))
  };
stage E1:
  GPR[registerWe<<1] = result1;
\end{lstlisting}}
\newpage

\paragraph{Encoding} Format \textsf{ALU\_WPWRRSO} (Section~\ref{sec:format-ALU-WPWRRSO}), \verb|exucode2 = 1110|\\

\bitfields{ 1/P, 2/11, 1/1, 4/1110, 5/registerWo, 1/1, 2/10, 3/010, 1/0, 5/registerYo, 1/--, 5/registerZo, 1/1 }


\paragraph{Syntax} \texttt{abdswp} \emph{registerWo} = \emph{registerZo}, \emph{registerYo}

\paragraph{Description} The \emph{registerZo} and the \emph{registerYo} are interpreted as two words packed into 64 bits. The \emph{registerZo} words are subtracted from the \emph{registerYo} words. The absolute values of the two resulting words are saturated, packed and stored into the \emph{registerWo}.


\paragraph{Execution} {Reservation table \textsf{ALU\_LITE} (Section~\ref{sec:reservation-ALU-LITE})}
~
{\begin{lstlisting}
stage RR:
  new argument3 = GPR[(registerYo<<1)+1];
  new argument2 = GPR[(registerZo<<1)+1];
  for (I = 0; I < 2; I += 1) {
    result1.32[I] = _SAT_32(_ABS(_SX_32(argument3.32[I]) - _SX_32(argument2.32[I])))
  };
stage E1:
  GPR[(registerWo<<1)+1] = result1;
\end{lstlisting}}
\newpage

\paragraph{Encoding} Format \textsf{ALU\_WPWRRSE.M} (Section~\ref{sec:format-ALU-WPWRRSE.M}), \verb|exucode2 = 1110|\\

\bitfields{ 1/P, 2/00, 2/-- --, 27/upper27, 1/1, 2/11, 1/1, 4/1110, 5/registerWe, 1/0, 2/10, 3/010, 1/0, 1/splat32, 5/lower5, 5/registerZe, 1/1 }


\paragraph{Syntax} \texttt{abdswp} \emph{registerWe} = \emph{registerZe}, \emph{upper27\_\discretionary{}{}{}lower5}\emph{splat32}

\paragraph{Description} The \emph{registerZe} and the \emph{upper27\_\discretionary{}{}{}lower5} are interpreted as two words packed into 64 bits. The \emph{registerZe} words are subtracted from the \emph{upper27\_\discretionary{}{}{}lower5} words. The absolute values of the two resulting words are saturated, packed and stored into the \emph{registerWe}.


\paragraph{Modifier(s)}
\noindent\begin{itemize}
\item splat32 (cf. splat32 in section \ref{par:modifier-splat32} on page \pageref{par:modifier-splat32})
\end{itemize}

\paragraph{Execution} {Reservation table \textsf{ALU\_LITE.X} (Section~\ref{sec:reservation-ALU-LITE.X})}
~
{\begin{lstlisting}
stage ID:
  new splat32 = _ZX_1(splat32);
stage RR:
  new argument3 = splat32 ? (_WX_32(upper27_lower5) << 32) | _ZX_32(_WX_32(upper27_lower5)) : _SX_32(_WX_32(upper27_lower5));
  new argument2 = GPR[registerZe<<1];
  for (I = 0; I < 2; I += 1) {
    result1.32[I] = _SAT_32(_ABS(_SX_32(argument3.32[I]) - _SX_32(argument2.32[I])))
  };
stage E1:
  GPR[registerWe<<1] = result1;
\end{lstlisting}}
\newpage

\paragraph{Encoding} Format \textsf{ALU\_WPWRRSO.M} (Section~\ref{sec:format-ALU-WPWRRSO.M}), \verb|exucode2 = 1110|\\

\bitfields{ 1/P, 2/00, 2/-- --, 27/upper27, 1/1, 2/11, 1/1, 4/1110, 5/registerWo, 1/1, 2/10, 3/010, 1/0, 1/splat32, 5/lower5, 5/registerZo, 1/1 }


\paragraph{Syntax} \texttt{abdswp} \emph{registerWo} = \emph{registerZo}, \emph{upper27\_\discretionary{}{}{}lower5}\emph{splat32}

\paragraph{Description} The \emph{registerZo} and the \emph{upper27\_\discretionary{}{}{}lower5} are interpreted as two words packed into 64 bits. The \emph{registerZo} words are subtracted from the \emph{upper27\_\discretionary{}{}{}lower5} words. The absolute values of the two resulting words are saturated, packed and stored into the \emph{registerWo}.


\paragraph{Modifier(s)}
\noindent\begin{itemize}
\item splat32 (cf. splat32 in section \ref{par:modifier-splat32} on page \pageref{par:modifier-splat32})
\end{itemize}

\paragraph{Execution} {Reservation table \textsf{ALU\_LITE.X} (Section~\ref{sec:reservation-ALU-LITE.X})}
~
{\begin{lstlisting}
stage ID:
  new splat32 = _ZX_1(splat32);
stage RR:
  new argument3 = splat32 ? (_WX_32(upper27_lower5) << 32) | _ZX_32(_WX_32(upper27_lower5)) : _SX_32(_WX_32(upper27_lower5));
  new argument2 = GPR[(registerZo<<1)+1];
  for (I = 0; I < 2; I += 1) {
    result1.32[I] = _SAT_32(_ABS(_SX_32(argument3.32[I]) - _SX_32(argument2.32[I])))
  };
stage E1:
  GPR[(registerWo<<1)+1] = result1;
\end{lstlisting}}
\newpage
\subsubsection[{\small ABDUBO: Absolute Difference Unsigned Byte Octuple}]{ABDUBO\hfill Absolute Difference Unsigned Byte Octuple}
\label{instr:ABDUBO}\index{abdubo}
 
\paragraph{Encoding} Format \textsf{ALU\_BOWRRSE} (Section~\ref{sec:format-ALU-BOWRRSE}), \verb|exucode2 = 1101|\\

\bitfields{ 1/P, 2/11, 1/1, 4/1101, 5/registerWe, 1/0, 2/10, 3/010, 1/1, 5/registerYe, 1/--, 5/registerZe, 1/1 }


\paragraph{Syntax} \texttt{abdubo} \emph{registerWe} = \emph{registerZe}, \emph{registerYe}

\paragraph{Description} The \emph{registerZe} and the \emph{registerYe} are interpreted as eight unsigned bytes packed into 64 bits. The \emph{registerZe} bytes are subtracted from the \emph{registerYe} bytes. The absolute value of the eight resulting bytes are packed and stored into the \emph{registerWe}.


\paragraph{Execution} {Reservation table \textsf{ALU\_LITE} (Section~\ref{sec:reservation-ALU-LITE})}
~
{\begin{lstlisting}
stage RR:
  new argument3 = GPR[registerYe<<1];
  new argument2 = GPR[registerZe<<1];
  for (I = 0; I < 8; I += 1) {
    result1.8[I] = _ABS(_ZX_8(argument3.8[I]) - _ZX_8(argument2.8[I]))
  };
stage E1:
  GPR[registerWe<<1] = result1;
\end{lstlisting}}
\newpage

\paragraph{Encoding} Format \textsf{ALU\_BOWRRSO} (Section~\ref{sec:format-ALU-BOWRRSO}), \verb|exucode2 = 1101|\\

\bitfields{ 1/P, 2/11, 1/1, 4/1101, 5/registerWo, 1/1, 2/10, 3/010, 1/1, 5/registerYo, 1/--, 5/registerZo, 1/1 }


\paragraph{Syntax} \texttt{abdubo} \emph{registerWo} = \emph{registerZo}, \emph{registerYo}

\paragraph{Description} The \emph{registerZo} and the \emph{registerYo} are interpreted as eight unsigned bytes packed into 64 bits. The \emph{registerZo} bytes are subtracted from the \emph{registerYo} bytes. The absolute value of the eight resulting bytes are packed and stored into the \emph{registerWo}.


\paragraph{Execution} {Reservation table \textsf{ALU\_LITE} (Section~\ref{sec:reservation-ALU-LITE})}
~
{\begin{lstlisting}
stage RR:
  new argument3 = GPR[(registerYo<<1)+1];
  new argument2 = GPR[(registerZo<<1)+1];
  for (I = 0; I < 8; I += 1) {
    result1.8[I] = _ABS(_ZX_8(argument3.8[I]) - _ZX_8(argument2.8[I]))
  };
stage E1:
  GPR[(registerWo<<1)+1] = result1;
\end{lstlisting}}
\newpage

\paragraph{Encoding} Format \textsf{ALU\_BOWRRSE.M} (Section~\ref{sec:format-ALU-BOWRRSE.M}), \verb|exucode2 = 1101|\\

\bitfields{ 1/P, 2/00, 2/-- --, 27/upper27, 1/1, 2/11, 1/1, 4/1101, 5/registerWe, 1/0, 2/10, 3/010, 1/1, 1/splat32, 5/lower5, 5/registerZe, 1/1 }


\paragraph{Syntax} \texttt{abdubo} \emph{registerWe} = \emph{registerZe}, \emph{upper27\_\discretionary{}{}{}lower5}\emph{splat32}

\paragraph{Description} The \emph{registerZe} and the \emph{upper27\_\discretionary{}{}{}lower5} are interpreted as eight unsigned bytes packed into 64 bits. The \emph{registerZe} bytes are subtracted from the \emph{upper27\_\discretionary{}{}{}lower5} bytes. The absolute value of the eight resulting bytes are packed and stored into the \emph{registerWe}.


\paragraph{Modifier(s)}
\noindent\begin{itemize}
\item splat32 (cf. splat32 in section \ref{par:modifier-splat32} on page \pageref{par:modifier-splat32})
\end{itemize}

\paragraph{Execution} {Reservation table \textsf{ALU\_LITE.X} (Section~\ref{sec:reservation-ALU-LITE.X})}
~
{\begin{lstlisting}
stage ID:
  new splat32 = _ZX_1(splat32);
stage RR:
  new argument3 = splat32 ? (_WX_32(upper27_lower5) << 32) | _ZX_32(_WX_32(upper27_lower5)) : _SX_32(_WX_32(upper27_lower5));
  new argument2 = GPR[registerZe<<1];
  for (I = 0; I < 8; I += 1) {
    result1.8[I] = _ABS(_ZX_8(argument3.8[I]) - _ZX_8(argument2.8[I]))
  };
stage E1:
  GPR[registerWe<<1] = result1;
\end{lstlisting}}
\newpage

\paragraph{Encoding} Format \textsf{ALU\_BOWRRSO.M} (Section~\ref{sec:format-ALU-BOWRRSO.M}), \verb|exucode2 = 1101|\\

\bitfields{ 1/P, 2/00, 2/-- --, 27/upper27, 1/1, 2/11, 1/1, 4/1101, 5/registerWo, 1/1, 2/10, 3/010, 1/1, 1/splat32, 5/lower5, 5/registerZo, 1/1 }


\paragraph{Syntax} \texttt{abdubo} \emph{registerWo} = \emph{registerZo}, \emph{upper27\_\discretionary{}{}{}lower5}\emph{splat32}

\paragraph{Description} The \emph{registerZo} and the \emph{upper27\_\discretionary{}{}{}lower5} are interpreted as eight unsigned bytes packed into 64 bits. The \emph{registerZo} bytes are subtracted from the \emph{upper27\_\discretionary{}{}{}lower5} bytes. The absolute value of the eight resulting bytes are packed and stored into the \emph{registerWo}.


\paragraph{Modifier(s)}
\noindent\begin{itemize}
\item splat32 (cf. splat32 in section \ref{par:modifier-splat32} on page \pageref{par:modifier-splat32})
\end{itemize}

\paragraph{Execution} {Reservation table \textsf{ALU\_LITE.X} (Section~\ref{sec:reservation-ALU-LITE.X})}
~
{\begin{lstlisting}
stage ID:
  new splat32 = _ZX_1(splat32);
stage RR:
  new argument3 = splat32 ? (_WX_32(upper27_lower5) << 32) | _ZX_32(_WX_32(upper27_lower5)) : _SX_32(_WX_32(upper27_lower5));
  new argument2 = GPR[(registerZo<<1)+1];
  for (I = 0; I < 8; I += 1) {
    result1.8[I] = _ABS(_ZX_8(argument3.8[I]) - _ZX_8(argument2.8[I]))
  };
stage E1:
  GPR[(registerWo<<1)+1] = result1;
\end{lstlisting}}
\newpage
\subsubsection[{\small ABDUD: Absolute Difference Unsigned Double Word}]{ABDUD\hfill Absolute Difference Unsigned Double Word}
\label{instr:ABDUD}\index{abdud}
 
\paragraph{Encoding} Format \textsf{ALU\_DWRRS} (Section~\ref{sec:format-ALU-DWRRS}), \verb|exucode2 = 1101|\\

\bitfields{ 1/P, 2/11, 1/0, 4/1101, 6/registerW, 2/10, 3/010, 1/0, 6/registerY, 6/registerZ }


\paragraph{Syntax} \texttt{abdud} \emph{registerW} = \emph{registerZ}, \emph{registerY}

\paragraph{Description} The \emph{registerZ} is subtracted from the \emph{registerY}, assuming unsigned integers. The absolute value of the result is stored into the \emph{registerW}.


\paragraph{Execution} {Reservation table \textsf{ALU\_TINY} (Section~\ref{sec:reservation-ALU-TINY})}
~
{\begin{lstlisting}
stage RR:
  new argument3 = GPR[registerY];
  new argument2 = GPR[registerZ];
  new result1 = _ABS(_ZX_64(argument3) - _ZX_64(argument2));
stage E1:
  GPR[registerW] = result1;
\end{lstlisting}}
\newpage

\paragraph{Encoding} Format \textsf{ALU\_DWRRS.M} (Section~\ref{sec:format-ALU-DWRRS.M}), \verb|exucode2 = 1101|\\

\bitfields{ 1/P, 2/00, 2/-- --, 27/upper27, 1/1, 2/11, 1/0, 4/1101, 6/registerW, 2/10, 3/010, 1/0, 1/splat32, 5/lower5, 6/registerZ }


\paragraph{Syntax} \texttt{abdud} \emph{registerW} = \emph{registerZ}, \emph{upper27\_\discretionary{}{}{}lower5}\emph{splat32}

\paragraph{Description} The \emph{registerZ} is subtracted from the \emph{upper27\_\discretionary{}{}{}lower5}, assuming unsigned integers. The absolute value of the result is stored into the \emph{registerW}.


\paragraph{Modifier(s)}
\noindent\begin{itemize}
\item splat32 (cf. splat32 in section \ref{par:modifier-splat32} on page \pageref{par:modifier-splat32})
\end{itemize}

\paragraph{Execution} {Reservation table \textsf{ALU\_TINY.X} (Section~\ref{sec:reservation-ALU-TINY.X})}
~
{\begin{lstlisting}
stage ID:
  new splat32 = _ZX_1(splat32);
stage RR:
  new argument3 = splat32 ? (_WX_32(upper27_lower5) << 32) | _ZX_32(_WX_32(upper27_lower5)) : _SX_32(_WX_32(upper27_lower5));
  new argument2 = GPR[registerZ];
  new result1 = _ABS(_ZX_64(argument3) - _ZX_64(argument2));
stage E1:
  GPR[registerW] = result1;
\end{lstlisting}}
\newpage
\subsubsection[{\small ABDUHQ: Absolute Difference Unsigned Half Word Quadruple}]{ABDUHQ\hfill Absolute Difference Unsigned Half Word Quadruple}
\label{instr:ABDUHQ}\index{abduhq}
 
\paragraph{Encoding} Format \textsf{ALU\_HQWRRSE} (Section~\ref{sec:format-ALU-HQWRRSE}), \verb|exucode2 = 1101|\\

\bitfields{ 1/P, 2/11, 1/1, 4/1101, 5/registerWe, 1/0, 2/10, 3/010, 1/1, 5/registerYe, 1/--, 5/registerZe, 1/0 }


\paragraph{Syntax} \texttt{abduhq} \emph{registerWe} = \emph{registerZe}, \emph{registerYe}

\paragraph{Description} The \emph{registerZe} and the \emph{registerYe} are interpreted as four unsigned half words packed into 64 bits. The \emph{registerZe} half words are subtracted from the \emph{registerYe} half words. The absolute value of the four resulting half words are packed and stored into the \emph{registerWe}.


\paragraph{Execution} {Reservation table \textsf{ALU\_LITE} (Section~\ref{sec:reservation-ALU-LITE})}
~
{\begin{lstlisting}
stage RR:
  new argument3 = GPR[registerYe<<1];
  new argument2 = GPR[registerZe<<1];
  for (I = 0; I < 4; I += 1) {
    result1.16[I] = _ABS(_ZX_16(argument3.16[I]) - _ZX_16(argument2.16[I]))
  };
stage E1:
  GPR[registerWe<<1] = result1;
\end{lstlisting}}
\newpage

\paragraph{Encoding} Format \textsf{ALU\_HQWRRSO} (Section~\ref{sec:format-ALU-HQWRRSO}), \verb|exucode2 = 1101|\\

\bitfields{ 1/P, 2/11, 1/1, 4/1101, 5/registerWo, 1/1, 2/10, 3/010, 1/1, 5/registerYo, 1/--, 5/registerZo, 1/0 }


\paragraph{Syntax} \texttt{abduhq} \emph{registerWo} = \emph{registerZo}, \emph{registerYo}

\paragraph{Description} The \emph{registerZo} and the \emph{registerYo} are interpreted as four unsigned half words packed into 64 bits. The \emph{registerZo} half words are subtracted from the \emph{registerYo} half words. The absolute value of the four resulting half words are packed and stored into the \emph{registerWo}.


\paragraph{Execution} {Reservation table \textsf{ALU\_LITE} (Section~\ref{sec:reservation-ALU-LITE})}
~
{\begin{lstlisting}
stage RR:
  new argument3 = GPR[(registerYo<<1)+1];
  new argument2 = GPR[(registerZo<<1)+1];
  for (I = 0; I < 4; I += 1) {
    result1.16[I] = _ABS(_ZX_16(argument3.16[I]) - _ZX_16(argument2.16[I]))
  };
stage E1:
  GPR[(registerWo<<1)+1] = result1;
\end{lstlisting}}
\newpage

\paragraph{Encoding} Format \textsf{ALU\_HQWRRSE.M} (Section~\ref{sec:format-ALU-HQWRRSE.M}), \verb|exucode2 = 1101|\\

\bitfields{ 1/P, 2/00, 2/-- --, 27/upper27, 1/1, 2/11, 1/1, 4/1101, 5/registerWe, 1/0, 2/10, 3/010, 1/1, 1/splat32, 5/lower5, 5/registerZe, 1/0 }


\paragraph{Syntax} \texttt{abduhq} \emph{registerWe} = \emph{registerZe}, \emph{upper27\_\discretionary{}{}{}lower5}\emph{splat32}

\paragraph{Description} The \emph{registerZe} and the \emph{upper27\_\discretionary{}{}{}lower5} are interpreted as four unsigned half words packed into 64 bits. The \emph{registerZe} half words are subtracted from the \emph{upper27\_\discretionary{}{}{}lower5} half words. The absolute value of the four resulting half words are packed and stored into the \emph{registerWe}.


\paragraph{Modifier(s)}
\noindent\begin{itemize}
\item splat32 (cf. splat32 in section \ref{par:modifier-splat32} on page \pageref{par:modifier-splat32})
\end{itemize}

\paragraph{Execution} {Reservation table \textsf{ALU\_LITE.X} (Section~\ref{sec:reservation-ALU-LITE.X})}
~
{\begin{lstlisting}
stage ID:
  new splat32 = _ZX_1(splat32);
stage RR:
  new argument3 = splat32 ? (_WX_32(upper27_lower5) << 32) | _ZX_32(_WX_32(upper27_lower5)) : _SX_32(_WX_32(upper27_lower5));
  new argument2 = GPR[registerZe<<1];
  for (I = 0; I < 4; I += 1) {
    result1.16[I] = _ABS(_ZX_16(argument3.16[I]) - _ZX_16(argument2.16[I]))
  };
stage E1:
  GPR[registerWe<<1] = result1;
\end{lstlisting}}
\newpage

\paragraph{Encoding} Format \textsf{ALU\_HQWRRSO.M} (Section~\ref{sec:format-ALU-HQWRRSO.M}), \verb|exucode2 = 1101|\\

\bitfields{ 1/P, 2/00, 2/-- --, 27/upper27, 1/1, 2/11, 1/1, 4/1101, 5/registerWo, 1/1, 2/10, 3/010, 1/1, 1/splat32, 5/lower5, 5/registerZo, 1/0 }


\paragraph{Syntax} \texttt{abduhq} \emph{registerWo} = \emph{registerZo}, \emph{upper27\_\discretionary{}{}{}lower5}\emph{splat32}

\paragraph{Description} The \emph{registerZo} and the \emph{upper27\_\discretionary{}{}{}lower5} are interpreted as four unsigned half words packed into 64 bits. The \emph{registerZo} half words are subtracted from the \emph{upper27\_\discretionary{}{}{}lower5} half words. The absolute value of the four resulting half words are packed and stored into the \emph{registerWo}.


\paragraph{Modifier(s)}
\noindent\begin{itemize}
\item splat32 (cf. splat32 in section \ref{par:modifier-splat32} on page \pageref{par:modifier-splat32})
\end{itemize}

\paragraph{Execution} {Reservation table \textsf{ALU\_LITE.X} (Section~\ref{sec:reservation-ALU-LITE.X})}
~
{\begin{lstlisting}
stage ID:
  new splat32 = _ZX_1(splat32);
stage RR:
  new argument3 = splat32 ? (_WX_32(upper27_lower5) << 32) | _ZX_32(_WX_32(upper27_lower5)) : _SX_32(_WX_32(upper27_lower5));
  new argument2 = GPR[(registerZo<<1)+1];
  for (I = 0; I < 4; I += 1) {
    result1.16[I] = _ABS(_ZX_16(argument3.16[I]) - _ZX_16(argument2.16[I]))
  };
stage E1:
  GPR[(registerWo<<1)+1] = result1;
\end{lstlisting}}
\newpage
\subsubsection[{\small ABDUW: Absolute Difference Unsigned Word}]{ABDUW\hfill Absolute Difference Unsigned Word}
\label{instr:ABDUW}\index{abduw}
 
\paragraph{Encoding} Format \textsf{ALU\_WWRRS} (Section~\ref{sec:format-ALU-WWRRS}), \verb|exucode2 = 1101|\\

\bitfields{ 1/P, 2/11, 1/0, 4/1101, 6/registerW, 2/11, 3/010, 1/signextw, 6/registerY, 6/registerZ }


\paragraph{Syntax} \texttt{abduw}\emph{signextw} \emph{registerW} = \emph{registerZ}, \emph{registerY}

\paragraph{Description} The \emph{registerZ} is subtracted from the \emph{registerY}, assuming unsigned integers. The absolute value of the result with the 32 upper bits cleared is stored into the \emph{registerW}.


\paragraph{Modifier(s)}
\noindent\begin{itemize}
\item signextw (cf. signextw in section \ref{par:modifier-signextw} on page \pageref{par:modifier-signextw})
\end{itemize}

\paragraph{Execution} {Reservation table \textsf{ALU\_TINY} (Section~\ref{sec:reservation-ALU-TINY})}
~
{\begin{lstlisting}
stage ID:
  new signextw = _ZX_1(signextw);
stage RR:
  new argument3 = GPR[registerY];
  new argument2 = GPR[registerZ];
  new result1 = _ABS(_ZX_32(argument3) - _ZX_32(argument2));
stage E1:
  GPR[registerW] = signextw ? _SX_32(result1) : _ZX_32(result1);
\end{lstlisting}}
\newpage

\paragraph{Encoding} Format \textsf{ALU\_WWRRS.M} (Section~\ref{sec:format-ALU-WWRRS.M}), \verb|exucode2 = 1101|\\

\bitfields{ 1/P, 2/00, 2/-- --, 27/upper27, 1/1, 2/11, 1/0, 4/1101, 6/registerW, 2/11, 3/010, 1/signextw, 1/splat32, 5/lower5, 6/registerZ }


\paragraph{Syntax} \texttt{abduw}\emph{signextw} \emph{registerW} = \emph{registerZ}, \emph{upper27\_\discretionary{}{}{}lower5}

\paragraph{Description} The \emph{registerZ} is subtracted from the \emph{upper27\_\discretionary{}{}{}lower5}, assuming unsigned integers. The absolute value of the result with the 32 upper bits cleared is stored into the \emph{registerW}.


\paragraph{Modifier(s)}
\noindent\begin{itemize}
\item signextw (cf. signextw in section \ref{par:modifier-signextw} on page \pageref{par:modifier-signextw})
\end{itemize}

\paragraph{Execution} {Reservation table \textsf{ALU\_TINY.X} (Section~\ref{sec:reservation-ALU-TINY.X})}
~
{\begin{lstlisting}
stage ID:
  new signextw = _ZX_1(signextw);
stage RR:
  new argument3 = _WX_32(upper27_lower5);
  new argument2 = GPR[registerZ];
  new result1 = _ABS(_ZX_32(argument3) - _ZX_32(argument2));
stage E1:
  GPR[registerW] = signextw ? _SX_32(result1) : _ZX_32(result1);
\end{lstlisting}}
\newpage
\subsubsection[{\small ABDUWP: Absolute Difference Unsigned Word Pair}]{ABDUWP\hfill Absolute Difference Unsigned Word Pair}
\label{instr:ABDUWP}\index{abduwp}
 
\paragraph{Encoding} Format \textsf{ALU\_WPWRRSE} (Section~\ref{sec:format-ALU-WPWRRSE}), \verb|exucode2 = 1101|\\

\bitfields{ 1/P, 2/11, 1/1, 4/1101, 5/registerWe, 1/0, 2/10, 3/010, 1/0, 5/registerYe, 1/--, 5/registerZe, 1/1 }


\paragraph{Syntax} \texttt{abduwp} \emph{registerWe} = \emph{registerZe}, \emph{registerYe}

\paragraph{Description} The \emph{registerZe} and the \emph{registerYe} are interpreted as two unsigned words packed into 64 bits. The \emph{registerZe} words are subtracted from the \emph{registerYe} words. The absolute value of the two resulting words are packed and stored into the \emph{registerWe}.


\paragraph{Execution} {Reservation table \textsf{ALU\_LITE} (Section~\ref{sec:reservation-ALU-LITE})}
~
{\begin{lstlisting}
stage RR:
  new argument3 = GPR[registerYe<<1];
  new argument2 = GPR[registerZe<<1];
  for (I = 0; I < 2; I += 1) {
    result1.32[I] = _ABS(_ZX_32(argument3.32[I]) - _ZX_32(argument2.32[I]))
  };
stage E1:
  GPR[registerWe<<1] = result1;
\end{lstlisting}}
\newpage

\paragraph{Encoding} Format \textsf{ALU\_WPWRRSO} (Section~\ref{sec:format-ALU-WPWRRSO}), \verb|exucode2 = 1101|\\

\bitfields{ 1/P, 2/11, 1/1, 4/1101, 5/registerWo, 1/1, 2/10, 3/010, 1/0, 5/registerYo, 1/--, 5/registerZo, 1/1 }


\paragraph{Syntax} \texttt{abduwp} \emph{registerWo} = \emph{registerZo}, \emph{registerYo}

\paragraph{Description} The \emph{registerZo} and the \emph{registerYo} are interpreted as two unsigned words packed into 64 bits. The \emph{registerZo} words are subtracted from the \emph{registerYo} words. The absolute value of the two resulting words are packed and stored into the \emph{registerWo}.


\paragraph{Execution} {Reservation table \textsf{ALU\_LITE} (Section~\ref{sec:reservation-ALU-LITE})}
~
{\begin{lstlisting}
stage RR:
  new argument3 = GPR[(registerYo<<1)+1];
  new argument2 = GPR[(registerZo<<1)+1];
  for (I = 0; I < 2; I += 1) {
    result1.32[I] = _ABS(_ZX_32(argument3.32[I]) - _ZX_32(argument2.32[I]))
  };
stage E1:
  GPR[(registerWo<<1)+1] = result1;
\end{lstlisting}}
\newpage

\paragraph{Encoding} Format \textsf{ALU\_WPWRRSE.M} (Section~\ref{sec:format-ALU-WPWRRSE.M}), \verb|exucode2 = 1101|\\

\bitfields{ 1/P, 2/00, 2/-- --, 27/upper27, 1/1, 2/11, 1/1, 4/1101, 5/registerWe, 1/0, 2/10, 3/010, 1/0, 1/splat32, 5/lower5, 5/registerZe, 1/1 }


\paragraph{Syntax} \texttt{abduwp} \emph{registerWe} = \emph{registerZe}, \emph{upper27\_\discretionary{}{}{}lower5}\emph{splat32}

\paragraph{Description} The \emph{registerZe} and the \emph{upper27\_\discretionary{}{}{}lower5} are interpreted as two unsigned words packed into 64 bits. The \emph{registerZe} words are subtracted from the \emph{upper27\_\discretionary{}{}{}lower5} words. The absolute value of the two resulting words are packed and stored into the \emph{registerWe}.


\paragraph{Modifier(s)}
\noindent\begin{itemize}
\item splat32 (cf. splat32 in section \ref{par:modifier-splat32} on page \pageref{par:modifier-splat32})
\end{itemize}

\paragraph{Execution} {Reservation table \textsf{ALU\_LITE.X} (Section~\ref{sec:reservation-ALU-LITE.X})}
~
{\begin{lstlisting}
stage ID:
  new splat32 = _ZX_1(splat32);
stage RR:
  new argument3 = splat32 ? (_WX_32(upper27_lower5) << 32) | _ZX_32(_WX_32(upper27_lower5)) : _SX_32(_WX_32(upper27_lower5));
  new argument2 = GPR[registerZe<<1];
  for (I = 0; I < 2; I += 1) {
    result1.32[I] = _ABS(_ZX_32(argument3.32[I]) - _ZX_32(argument2.32[I]))
  };
stage E1:
  GPR[registerWe<<1] = result1;
\end{lstlisting}}
\newpage

\paragraph{Encoding} Format \textsf{ALU\_WPWRRSO.M} (Section~\ref{sec:format-ALU-WPWRRSO.M}), \verb|exucode2 = 1101|\\

\bitfields{ 1/P, 2/00, 2/-- --, 27/upper27, 1/1, 2/11, 1/1, 4/1101, 5/registerWo, 1/1, 2/10, 3/010, 1/0, 1/splat32, 5/lower5, 5/registerZo, 1/1 }


\paragraph{Syntax} \texttt{abduwp} \emph{registerWo} = \emph{registerZo}, \emph{upper27\_\discretionary{}{}{}lower5}\emph{splat32}

\paragraph{Description} The \emph{registerZo} and the \emph{upper27\_\discretionary{}{}{}lower5} are interpreted as two unsigned words packed into 64 bits. The \emph{registerZo} words are subtracted from the \emph{upper27\_\discretionary{}{}{}lower5} words. The absolute value of the two resulting words are packed and stored into the \emph{registerWo}.


\paragraph{Modifier(s)}
\noindent\begin{itemize}
\item splat32 (cf. splat32 in section \ref{par:modifier-splat32} on page \pageref{par:modifier-splat32})
\end{itemize}

\paragraph{Execution} {Reservation table \textsf{ALU\_LITE.X} (Section~\ref{sec:reservation-ALU-LITE.X})}
~
{\begin{lstlisting}
stage ID:
  new splat32 = _ZX_1(splat32);
stage RR:
  new argument3 = splat32 ? (_WX_32(upper27_lower5) << 32) | _ZX_32(_WX_32(upper27_lower5)) : _SX_32(_WX_32(upper27_lower5));
  new argument2 = GPR[(registerZo<<1)+1];
  for (I = 0; I < 2; I += 1) {
    result1.32[I] = _ABS(_ZX_32(argument3.32[I]) - _ZX_32(argument2.32[I]))
  };
stage E1:
  GPR[(registerWo<<1)+1] = result1;
\end{lstlisting}}
\newpage
\subsubsection[{\small ABDW: Absolute Difference Word}]{ABDW\hfill Absolute Difference Word}
\label{instr:ABDW}\index{abdw}
 
\paragraph{Encoding} Format \textsf{ALU\_WWRRS} (Section~\ref{sec:format-ALU-WWRRS}), \verb|exucode2 = 1100|\\

\bitfields{ 1/P, 2/11, 1/0, 4/1100, 6/registerW, 2/11, 3/010, 1/signextw, 6/registerY, 6/registerZ }


\paragraph{Syntax} \texttt{abdw}\emph{signextw} \emph{registerW} = \emph{registerZ}, \emph{registerY}

\paragraph{Description} The \emph{registerZ} is subtracted from the \emph{registerY}. The absolute value of the result with the 32 upper bits cleared is stored into the \emph{registerW}.


\paragraph{Modifier(s)}
\noindent\begin{itemize}
\item signextw (cf. signextw in section \ref{par:modifier-signextw} on page \pageref{par:modifier-signextw})
\end{itemize}

\paragraph{Execution} {Reservation table \textsf{ALU\_TINY} (Section~\ref{sec:reservation-ALU-TINY})}
~
{\begin{lstlisting}
stage ID:
  new signextw = _ZX_1(signextw);
stage RR:
  new argument3 = GPR[registerY];
  new argument2 = GPR[registerZ];
  new result1 = _ABS(_SX_32(argument3) - _SX_32(argument2));
stage E1:
  GPR[registerW] = signextw ? _SX_32(result1) : _ZX_32(result1);
\end{lstlisting}}
\newpage

\paragraph{Encoding} Format \textsf{ALU\_WWRRS.M} (Section~\ref{sec:format-ALU-WWRRS.M}), \verb|exucode2 = 1100|\\

\bitfields{ 1/P, 2/00, 2/-- --, 27/upper27, 1/1, 2/11, 1/0, 4/1100, 6/registerW, 2/11, 3/010, 1/signextw, 1/splat32, 5/lower5, 6/registerZ }


\paragraph{Syntax} \texttt{abdw}\emph{signextw} \emph{registerW} = \emph{registerZ}, \emph{upper27\_\discretionary{}{}{}lower5}

\paragraph{Description} The \emph{registerZ} is subtracted from the \emph{upper27\_\discretionary{}{}{}lower5}. The absolute value of the result with the 32 upper bits cleared is stored into the \emph{registerW}.


\paragraph{Modifier(s)}
\noindent\begin{itemize}
\item signextw (cf. signextw in section \ref{par:modifier-signextw} on page \pageref{par:modifier-signextw})
\end{itemize}

\paragraph{Execution} {Reservation table \textsf{ALU\_TINY.X} (Section~\ref{sec:reservation-ALU-TINY.X})}
~
{\begin{lstlisting}
stage ID:
  new signextw = _ZX_1(signextw);
stage RR:
  new argument3 = _WX_32(upper27_lower5);
  new argument2 = GPR[registerZ];
  new result1 = _ABS(_SX_32(argument3) - _SX_32(argument2));
stage E1:
  GPR[registerW] = signextw ? _SX_32(result1) : _ZX_32(result1);
\end{lstlisting}}
\newpage
\subsubsection[{\small ABDWP: Absolute Difference Word Pair}]{ABDWP\hfill Absolute Difference Word Pair}
\label{instr:ABDWP}\index{abdwp}
 
\paragraph{Encoding} Format \textsf{ALU\_WPWRRSE} (Section~\ref{sec:format-ALU-WPWRRSE}), \verb|exucode2 = 1100|\\

\bitfields{ 1/P, 2/11, 1/1, 4/1100, 5/registerWe, 1/0, 2/10, 3/010, 1/0, 5/registerYe, 1/--, 5/registerZe, 1/1 }


\paragraph{Syntax} \texttt{abdwp} \emph{registerWe} = \emph{registerZe}, \emph{registerYe}

\paragraph{Description} The \emph{registerZe} and the \emph{registerYe} are interpreted as two words packed into 64 bits. The \emph{registerZe} words are subtracted from the \emph{registerYe} words. The absolute value of the two resulting words are packed and stored into the \emph{registerWe}.


\paragraph{Execution} {Reservation table \textsf{ALU\_LITE} (Section~\ref{sec:reservation-ALU-LITE})}
~
{\begin{lstlisting}
stage RR:
  new argument3 = GPR[registerYe<<1];
  new argument2 = GPR[registerZe<<1];
  for (I = 0; I < 2; I += 1) {
    result1.32[I] = _ABS(_SX_32(argument3.32[I]) - _SX_32(argument2.32[I]))
  };
stage E1:
  GPR[registerWe<<1] = result1;
\end{lstlisting}}
\newpage

\paragraph{Encoding} Format \textsf{ALU\_WPWRRSO} (Section~\ref{sec:format-ALU-WPWRRSO}), \verb|exucode2 = 1100|\\

\bitfields{ 1/P, 2/11, 1/1, 4/1100, 5/registerWo, 1/1, 2/10, 3/010, 1/0, 5/registerYo, 1/--, 5/registerZo, 1/1 }


\paragraph{Syntax} \texttt{abdwp} \emph{registerWo} = \emph{registerZo}, \emph{registerYo}

\paragraph{Description} The \emph{registerZo} and the \emph{registerYo} are interpreted as two words packed into 64 bits. The \emph{registerZo} words are subtracted from the \emph{registerYo} words. The absolute value of the two resulting words are packed and stored into the \emph{registerWo}.


\paragraph{Execution} {Reservation table \textsf{ALU\_LITE} (Section~\ref{sec:reservation-ALU-LITE})}
~
{\begin{lstlisting}
stage RR:
  new argument3 = GPR[(registerYo<<1)+1];
  new argument2 = GPR[(registerZo<<1)+1];
  for (I = 0; I < 2; I += 1) {
    result1.32[I] = _ABS(_SX_32(argument3.32[I]) - _SX_32(argument2.32[I]))
  };
stage E1:
  GPR[(registerWo<<1)+1] = result1;
\end{lstlisting}}
\newpage

\paragraph{Encoding} Format \textsf{ALU\_WPWRRSE.M} (Section~\ref{sec:format-ALU-WPWRRSE.M}), \verb|exucode2 = 1100|\\

\bitfields{ 1/P, 2/00, 2/-- --, 27/upper27, 1/1, 2/11, 1/1, 4/1100, 5/registerWe, 1/0, 2/10, 3/010, 1/0, 1/splat32, 5/lower5, 5/registerZe, 1/1 }


\paragraph{Syntax} \texttt{abdwp} \emph{registerWe} = \emph{registerZe}, \emph{upper27\_\discretionary{}{}{}lower5}\emph{splat32}

\paragraph{Description} The \emph{registerZe} and the \emph{upper27\_\discretionary{}{}{}lower5} are interpreted as two words packed into 64 bits. The \emph{registerZe} words are subtracted from the \emph{upper27\_\discretionary{}{}{}lower5} words. The absolute value of the two resulting words are packed and stored into the \emph{registerWe}.


\paragraph{Modifier(s)}
\noindent\begin{itemize}
\item splat32 (cf. splat32 in section \ref{par:modifier-splat32} on page \pageref{par:modifier-splat32})
\end{itemize}

\paragraph{Execution} {Reservation table \textsf{ALU\_LITE.X} (Section~\ref{sec:reservation-ALU-LITE.X})}
~
{\begin{lstlisting}
stage ID:
  new splat32 = _ZX_1(splat32);
stage RR:
  new argument3 = splat32 ? (_WX_32(upper27_lower5) << 32) | _ZX_32(_WX_32(upper27_lower5)) : _SX_32(_WX_32(upper27_lower5));
  new argument2 = GPR[registerZe<<1];
  for (I = 0; I < 2; I += 1) {
    result1.32[I] = _ABS(_SX_32(argument3.32[I]) - _SX_32(argument2.32[I]))
  };
stage E1:
  GPR[registerWe<<1] = result1;
\end{lstlisting}}
\newpage

\paragraph{Encoding} Format \textsf{ALU\_WPWRRSO.M} (Section~\ref{sec:format-ALU-WPWRRSO.M}), \verb|exucode2 = 1100|\\

\bitfields{ 1/P, 2/00, 2/-- --, 27/upper27, 1/1, 2/11, 1/1, 4/1100, 5/registerWo, 1/1, 2/10, 3/010, 1/0, 1/splat32, 5/lower5, 5/registerZo, 1/1 }


\paragraph{Syntax} \texttt{abdwp} \emph{registerWo} = \emph{registerZo}, \emph{upper27\_\discretionary{}{}{}lower5}\emph{splat32}

\paragraph{Description} The \emph{registerZo} and the \emph{upper27\_\discretionary{}{}{}lower5} are interpreted as two words packed into 64 bits. The \emph{registerZo} words are subtracted from the \emph{upper27\_\discretionary{}{}{}lower5} words. The absolute value of the two resulting words are packed and stored into the \emph{registerWo}.


\paragraph{Modifier(s)}
\noindent\begin{itemize}
\item splat32 (cf. splat32 in section \ref{par:modifier-splat32} on page \pageref{par:modifier-splat32})
\end{itemize}

\paragraph{Execution} {Reservation table \textsf{ALU\_LITE.X} (Section~\ref{sec:reservation-ALU-LITE.X})}
~
{\begin{lstlisting}
stage ID:
  new splat32 = _ZX_1(splat32);
stage RR:
  new argument3 = splat32 ? (_WX_32(upper27_lower5) << 32) | _ZX_32(_WX_32(upper27_lower5)) : _SX_32(_WX_32(upper27_lower5));
  new argument2 = GPR[(registerZo<<1)+1];
  for (I = 0; I < 2; I += 1) {
    result1.32[I] = _ABS(_SX_32(argument3.32[I]) - _SX_32(argument2.32[I]))
  };
stage E1:
  GPR[(registerWo<<1)+1] = result1;
\end{lstlisting}}
\newpage
\subsubsection[{\small ABSBO: Absolute Value Byte Octuple}]{ABSBO\hfill Absolute Value Byte Octuple}
\label{instr:ABSBO}\index{absbo}
 
\paragraph{Encoding} Format \textsf{ALU\_BOBWRE} (Section~\ref{sec:format-ALU-BOBWRE}), \verb|alucode8 = 0001|\\

\bitfields{ 1/P, 2/11, 1/1, 4/1111, 5/registerWe, 1/0, 2/10, 3/101, 1/1, 4/0001, 1/--, 1/--, 5/registerZe, 1/1 }


\paragraph{Syntax} \texttt{absbo} \emph{registerWe} = \emph{registerZe}

\paragraph{Description} The absolute values of the \emph{registerZe} interpreted as a byte octuple are stored into the \emph{registerWe}.


\paragraph{Execution} {Reservation table \textsf{ALU\_LITE} (Section~\ref{sec:reservation-ALU-LITE})}
~
{\begin{lstlisting}
stage RR:
  new argument2 = GPR[registerZe<<1];
  for (I = 0; I < 8; I += 1) {
    result1.8[I] = _ABS(_SX_8(argument2.8[I]))
  };
stage E1:
  GPR[registerWe<<1] = result1;
\end{lstlisting}}
\newpage

\paragraph{Encoding} Format \textsf{ALU\_BOBWRO} (Section~\ref{sec:format-ALU-BOBWRO}), \verb|alucode8 = 0001|\\

\bitfields{ 1/P, 2/11, 1/1, 4/1111, 5/registerWo, 1/1, 2/10, 3/101, 1/1, 4/0001, 1/--, 1/--, 5/registerZo, 1/1 }


\paragraph{Syntax} \texttt{absbo} \emph{registerWo} = \emph{registerZo}

\paragraph{Description} The absolute values of the \emph{registerZo} interpreted as a byte octuple are stored into the \emph{registerWo}.


\paragraph{Execution} {Reservation table \textsf{ALU\_LITE} (Section~\ref{sec:reservation-ALU-LITE})}
~
{\begin{lstlisting}
stage RR:
  new argument2 = GPR[(registerZo<<1)+1];
  for (I = 0; I < 8; I += 1) {
    result1.8[I] = _ABS(_SX_8(argument2.8[I]))
  };
stage E1:
  GPR[(registerWo<<1)+1] = result1;
\end{lstlisting}}
\newpage
\subsubsection[{\small ABSD: Absolute Value Double Word}]{ABSD\hfill Absolute Value Double Word}
\label{instr:ABSD}\index{absd}
 
\paragraph{Encoding} Format \textsf{ALU\_DBWR} (Section~\ref{sec:format-ALU-DBWR}), \verb|alucode8 = 0001|\\

\bitfields{ 1/P, 2/11, 1/0, 4/1111, 6/registerW, 2/10, 3/101, 1/0, 4/0001, 1/0, 1/0, 6/registerZ }


\paragraph{Syntax} \texttt{absd} \emph{registerW} = \emph{registerZ}

\paragraph{Description} The absolute value of the \emph{registerZ} is stored into the \emph{registerW}.


\paragraph{Execution} {Reservation table \textsf{ALU\_TINY} (Section~\ref{sec:reservation-ALU-TINY})}
~
{\begin{lstlisting}
stage RR:
  new argument2 = GPR[registerZ];
  new result1 = _ABS(_SX_64(argument2));
stage E1:
  GPR[registerW] = result1;
\end{lstlisting}}
\newpage
\subsubsection[{\small ABSHQ: Absolute Value Half Word Quadruple}]{ABSHQ\hfill Absolute Value Half Word Quadruple}
\label{instr:ABSHQ}\index{abshq}
 
\paragraph{Encoding} Format \textsf{ALU\_HQBWRE} (Section~\ref{sec:format-ALU-HQBWRE}), \verb|alucode8 = 0001|\\

\bitfields{ 1/P, 2/11, 1/1, 4/1111, 5/registerWe, 1/0, 2/10, 3/101, 1/1, 4/0001, 1/--, 1/--, 5/registerZe, 1/0 }


\paragraph{Syntax} \texttt{abshq} \emph{registerWe} = \emph{registerZe}

\paragraph{Description} The absolute values of the \emph{registerZe} interpreted as a half word quadruple are stored into the \emph{registerWe}.


\paragraph{Execution} {Reservation table \textsf{ALU\_LITE} (Section~\ref{sec:reservation-ALU-LITE})}
~
{\begin{lstlisting}
stage RR:
  new argument2 = GPR[registerZe<<1];
  for (I = 0; I < 4; I += 1) {
    result1.16[I] = _ABS(_SX_16(argument2.16[I]))
  };
stage E1:
  GPR[registerWe<<1] = result1;
\end{lstlisting}}
\newpage

\paragraph{Encoding} Format \textsf{ALU\_HQBWRO} (Section~\ref{sec:format-ALU-HQBWRO}), \verb|alucode8 = 0001|\\

\bitfields{ 1/P, 2/11, 1/1, 4/1111, 5/registerWo, 1/1, 2/10, 3/101, 1/1, 4/0001, 1/--, 1/--, 5/registerZo, 1/0 }


\paragraph{Syntax} \texttt{abshq} \emph{registerWo} = \emph{registerZo}

\paragraph{Description} The absolute values of the \emph{registerZo} interpreted as a half word quadruple are stored into the \emph{registerWo}.


\paragraph{Execution} {Reservation table \textsf{ALU\_LITE} (Section~\ref{sec:reservation-ALU-LITE})}
~
{\begin{lstlisting}
stage RR:
  new argument2 = GPR[(registerZo<<1)+1];
  for (I = 0; I < 4; I += 1) {
    result1.16[I] = _ABS(_SX_16(argument2.16[I]))
  };
stage E1:
  GPR[(registerWo<<1)+1] = result1;
\end{lstlisting}}
\newpage
\subsubsection[{\small ABSSBO: Absolute Value Saturated Byte Octuple}]{ABSSBO\hfill Absolute Value Saturated Byte Octuple}
\label{instr:ABSSBO}\index{abssbo}
 
\paragraph{Encoding} Format \textsf{ALU\_BOBWRE} (Section~\ref{sec:format-ALU-BOBWRE}), \verb|alucode8 = 0011|\\

\bitfields{ 1/P, 2/11, 1/1, 4/1111, 5/registerWe, 1/0, 2/10, 3/101, 1/1, 4/0011, 1/--, 1/--, 5/registerZe, 1/1 }


\paragraph{Syntax} \texttt{abssbo} \emph{registerWe} = \emph{registerZe}

\paragraph{Description} The saturated absolute values of the \emph{registerZe} interpreted as a byte octuple are stored into the \emph{registerWe}.


\paragraph{Execution} {Reservation table \textsf{ALU\_LITE} (Section~\ref{sec:reservation-ALU-LITE})}
~
{\begin{lstlisting}
stage RR:
  new argument2 = GPR[registerZe<<1];
  for (I = 0; I < 8; I += 1) {
    result1.8[I] = _SAT_8(_ABS(_SX_8(argument2.8[I])))
  };
stage E1:
  GPR[registerWe<<1] = result1;
\end{lstlisting}}
\newpage

\paragraph{Encoding} Format \textsf{ALU\_BOBWRO} (Section~\ref{sec:format-ALU-BOBWRO}), \verb|alucode8 = 0011|\\

\bitfields{ 1/P, 2/11, 1/1, 4/1111, 5/registerWo, 1/1, 2/10, 3/101, 1/1, 4/0011, 1/--, 1/--, 5/registerZo, 1/1 }


\paragraph{Syntax} \texttt{abssbo} \emph{registerWo} = \emph{registerZo}

\paragraph{Description} The saturated absolute values of the \emph{registerZo} interpreted as a byte octuple are stored into the \emph{registerWo}.


\paragraph{Execution} {Reservation table \textsf{ALU\_LITE} (Section~\ref{sec:reservation-ALU-LITE})}
~
{\begin{lstlisting}
stage RR:
  new argument2 = GPR[(registerZo<<1)+1];
  for (I = 0; I < 8; I += 1) {
    result1.8[I] = _SAT_8(_ABS(_SX_8(argument2.8[I])))
  };
stage E1:
  GPR[(registerWo<<1)+1] = result1;
\end{lstlisting}}
\newpage
\subsubsection[{\small ABSSD: Absolute Value Saturated Double Word}]{ABSSD\hfill Absolute Value Saturated Double Word}
\label{instr:ABSSD}\index{abssd}
 
\paragraph{Encoding} Format \textsf{ALU\_DBWR} (Section~\ref{sec:format-ALU-DBWR}), \verb|alucode8 = 0011|\\

\bitfields{ 1/P, 2/11, 1/0, 4/1111, 6/registerW, 2/10, 3/101, 1/0, 4/0011, 1/0, 1/0, 6/registerZ }


\paragraph{Syntax} \texttt{abssd} \emph{registerW} = \emph{registerZ}

\paragraph{Description} The absolute value of the \emph{registerZ} is 64-bit saturated and stored into the \emph{registerW}.


\paragraph{Execution} {Reservation table \textsf{ALU\_TINY} (Section~\ref{sec:reservation-ALU-TINY})}
~
{\begin{lstlisting}
stage RR:
  new argument2 = GPR[registerZ];
  new result1 = _SAT_64(_ABS(_SX_64(argument2)));
stage E1:
  GPR[registerW] = result1;
\end{lstlisting}}
\newpage
\subsubsection[{\small ABSSHQ: Absolute Value Saturated Half Word Quadruple}]{ABSSHQ\hfill Absolute Value Saturated Half Word Quadruple}
\label{instr:ABSSHQ}\index{absshq}
 
\paragraph{Encoding} Format \textsf{ALU\_HQBWRE} (Section~\ref{sec:format-ALU-HQBWRE}), \verb|alucode8 = 0011|\\

\bitfields{ 1/P, 2/11, 1/1, 4/1111, 5/registerWe, 1/0, 2/10, 3/101, 1/1, 4/0011, 1/--, 1/--, 5/registerZe, 1/0 }


\paragraph{Syntax} \texttt{absshq} \emph{registerWe} = \emph{registerZe}

\paragraph{Description} The saturated absolute values of the \emph{registerZe} interpreted as a half word quadruple are stored into the \emph{registerWe}.


\paragraph{Execution} {Reservation table \textsf{ALU\_LITE} (Section~\ref{sec:reservation-ALU-LITE})}
~
{\begin{lstlisting}
stage RR:
  new argument2 = GPR[registerZe<<1];
  for (I = 0; I < 4; I += 1) {
    result1.16[I] = _SAT_16(_ABS(_SX_16(argument2.16[I])))
  };
stage E1:
  GPR[registerWe<<1] = result1;
\end{lstlisting}}
\newpage

\paragraph{Encoding} Format \textsf{ALU\_HQBWRO} (Section~\ref{sec:format-ALU-HQBWRO}), \verb|alucode8 = 0011|\\

\bitfields{ 1/P, 2/11, 1/1, 4/1111, 5/registerWo, 1/1, 2/10, 3/101, 1/1, 4/0011, 1/--, 1/--, 5/registerZo, 1/0 }


\paragraph{Syntax} \texttt{absshq} \emph{registerWo} = \emph{registerZo}

\paragraph{Description} The saturated absolute values of the \emph{registerZo} interpreted as a half word quadruple are stored into the \emph{registerWo}.


\paragraph{Execution} {Reservation table \textsf{ALU\_LITE} (Section~\ref{sec:reservation-ALU-LITE})}
~
{\begin{lstlisting}
stage RR:
  new argument2 = GPR[(registerZo<<1)+1];
  for (I = 0; I < 4; I += 1) {
    result1.16[I] = _SAT_16(_ABS(_SX_16(argument2.16[I])))
  };
stage E1:
  GPR[(registerWo<<1)+1] = result1;
\end{lstlisting}}
\newpage
\subsubsection[{\small ABSSW: Absolute Value Saturated Word}]{ABSSW\hfill Absolute Value Saturated Word}
\label{instr:ABSSW}\index{abssw}
 
\paragraph{Encoding} Format \textsf{ALU\_BWRW} (Section~\ref{sec:format-ALU-BWRW}), \verb|alucode8 = 0011|\\

\bitfields{ 1/P, 2/11, 1/0, 4/1111, 6/registerW, 2/11, 3/101, 1/signextw, 4/0011, 1/0, 1/0, 6/registerZ }


\paragraph{Syntax} \texttt{abssw} \emph{registerW} = \emph{registerZ}

\paragraph{Description} The absolute value of the \emph{registerZ} is 32-bit saturated and stored into the \emph{registerW}.


\paragraph{Execution} {Reservation table \textsf{ALU\_TINY} (Section~\ref{sec:reservation-ALU-TINY})}
~
{\begin{lstlisting}
stage RR:
  new argument2 = GPR[registerZ];
  new result1 = _SAT_32(_ABS(_SX_32(argument2)));
stage E1:
  GPR[registerW] = _ZX_32(result1);
\end{lstlisting}}
\newpage
\subsubsection[{\small ABSSWP: Absolute Value Saturated Word Pair}]{ABSSWP\hfill Absolute Value Saturated Word Pair}
\label{instr:ABSSWP}\index{absswp}
 
\paragraph{Encoding} Format \textsf{ALU\_WPBWRE} (Section~\ref{sec:format-ALU-WPBWRE}), \verb|alucode8 = 0011|\\

\bitfields{ 1/P, 2/11, 1/1, 4/1111, 5/registerWe, 1/0, 2/10, 3/101, 1/0, 4/0011, 1/--, 1/--, 5/registerZe, 1/1 }


\paragraph{Syntax} \texttt{absswp} \emph{registerWe} = \emph{registerZe}

\paragraph{Description} The saturated absolute values of the \emph{registerZe} interpreted as a word pair are stored into the \emph{registerWe}.


\paragraph{Execution} {Reservation table \textsf{ALU\_LITE} (Section~\ref{sec:reservation-ALU-LITE})}
~
{\begin{lstlisting}
stage RR:
  new argument2 = GPR[registerZe<<1];
  for (I = 0; I < 2; I += 1) {
    result1.32[I] = _SAT_32(_ABS(_SX_32(argument2.32[I])))
  };
stage E1:
  GPR[registerWe<<1] = result1;
\end{lstlisting}}
\newpage

\paragraph{Encoding} Format \textsf{ALU\_WPBWRO} (Section~\ref{sec:format-ALU-WPBWRO}), \verb|alucode8 = 0011|\\

\bitfields{ 1/P, 2/11, 1/1, 4/1111, 5/registerWo, 1/1, 2/10, 3/101, 1/0, 4/0011, 1/--, 1/--, 5/registerZo, 1/1 }


\paragraph{Syntax} \texttt{absswp} \emph{registerWo} = \emph{registerZo}

\paragraph{Description} The saturated absolute values of the \emph{registerZo} interpreted as a word pair are stored into the \emph{registerWo}.


\paragraph{Execution} {Reservation table \textsf{ALU\_LITE} (Section~\ref{sec:reservation-ALU-LITE})}
~
{\begin{lstlisting}
stage RR:
  new argument2 = GPR[(registerZo<<1)+1];
  for (I = 0; I < 2; I += 1) {
    result1.32[I] = _SAT_32(_ABS(_SX_32(argument2.32[I])))
  };
stage E1:
  GPR[(registerWo<<1)+1] = result1;
\end{lstlisting}}
\newpage
\subsubsection[{\small ABSW: Absolute Value Word}]{ABSW\hfill Absolute Value Word}
\label{instr:ABSW}\index{absw}
 
\paragraph{Encoding} Format \textsf{ALU\_BWRW} (Section~\ref{sec:format-ALU-BWRW}), \verb|alucode8 = 0001|\\

\bitfields{ 1/P, 2/11, 1/0, 4/1111, 6/registerW, 2/11, 3/101, 1/signextw, 4/0001, 1/0, 1/0, 6/registerZ }


\paragraph{Syntax} \texttt{absw} \emph{registerW} = \emph{registerZ}

\paragraph{Description} The absolute value of the \emph{registerZ} is stored into the \emph{registerW}.


\paragraph{Execution} {Reservation table \textsf{ALU\_TINY} (Section~\ref{sec:reservation-ALU-TINY})}
~
{\begin{lstlisting}
stage RR:
  new argument2 = GPR[registerZ];
  new result1 = _ABS(_SX_32(argument2));
stage E1:
  GPR[registerW] = _ZX_32(result1);
\end{lstlisting}}
\newpage
\subsubsection[{\small ABSWP: Absolute Value Word Pair}]{ABSWP\hfill Absolute Value Word Pair}
\label{instr:ABSWP}\index{abswp}
 
\paragraph{Encoding} Format \textsf{ALU\_WPBWRE} (Section~\ref{sec:format-ALU-WPBWRE}), \verb|alucode8 = 0001|\\

\bitfields{ 1/P, 2/11, 1/1, 4/1111, 5/registerWe, 1/0, 2/10, 3/101, 1/0, 4/0001, 1/--, 1/--, 5/registerZe, 1/1 }


\paragraph{Syntax} \texttt{abswp} \emph{registerWe} = \emph{registerZe}

\paragraph{Description} The absolute values of the \emph{registerZe} interpreted as a word pair are stored into the \emph{registerWe}.


\paragraph{Execution} {Reservation table \textsf{ALU\_LITE} (Section~\ref{sec:reservation-ALU-LITE})}
~
{\begin{lstlisting}
stage RR:
  new argument2 = GPR[registerZe<<1];
  for (I = 0; I < 2; I += 1) {
    result1.32[I] = _ABS(_SX_32(argument2.32[I]))
  };
stage E1:
  GPR[registerWe<<1] = result1;
\end{lstlisting}}
\newpage

\paragraph{Encoding} Format \textsf{ALU\_WPBWRO} (Section~\ref{sec:format-ALU-WPBWRO}), \verb|alucode8 = 0001|\\

\bitfields{ 1/P, 2/11, 1/1, 4/1111, 5/registerWo, 1/1, 2/10, 3/101, 1/0, 4/0001, 1/--, 1/--, 5/registerZo, 1/1 }


\paragraph{Syntax} \texttt{abswp} \emph{registerWo} = \emph{registerZo}

\paragraph{Description} The absolute values of the \emph{registerZo} interpreted as a word pair are stored into the \emph{registerWo}.


\paragraph{Execution} {Reservation table \textsf{ALU\_LITE} (Section~\ref{sec:reservation-ALU-LITE})}
~
{\begin{lstlisting}
stage RR:
  new argument2 = GPR[(registerZo<<1)+1];
  for (I = 0; I < 2; I += 1) {
    result1.32[I] = _ABS(_SX_32(argument2.32[I]))
  };
stage E1:
  GPR[(registerWo<<1)+1] = result1;
\end{lstlisting}}
\newpage
\subsubsection[{\small ADDBO: Add Byte Octuple}]{ADDBO\hfill Add Byte Octuple}
\label{instr:ADDBO}\index{addbo}
 
\paragraph{Encoding} Format \textsf{ALU\_BOWRR0E} (Section~\ref{sec:format-ALU-BOWRR0E}), \verb|exucode2 = 0010|\\

\bitfields{ 1/P, 2/11, 1/1, 4/0010, 5/registerWe, 1/0, 2/10, 3/000, 1/1, 5/registerYe, 1/--, 5/registerZe, 1/1 }


\paragraph{Syntax} \texttt{addbo} \emph{registerWe} = \emph{registerZe}, \emph{registerYe}

\paragraph{Description} The \emph{registerYe} and the \emph{registerZe} are interpreted as eight bytes packed into 64 bits. The \emph{registerYe} bytes and the \emph{registerZe} bytes are added. The eight resulting bytes are packed and stored into the \emph{registerWe}.


\paragraph{Execution} {Reservation table \textsf{ALU\_LITE} (Section~\ref{sec:reservation-ALU-LITE})}
~
{\begin{lstlisting}
stage RR:
  new argument3 = GPR[registerYe<<1];
  new argument2 = GPR[registerZe<<1];
  for (I = 0; I < 8; I += 1) {
    result1.8[I] = argument3.8[I] + argument2.8[I]
  };
stage E1:
  GPR[registerWe<<1] = result1;
\end{lstlisting}}
\newpage

\paragraph{Encoding} Format \textsf{ALU\_BOWRR0O} (Section~\ref{sec:format-ALU-BOWRR0O}), \verb|exucode2 = 0010|\\

\bitfields{ 1/P, 2/11, 1/1, 4/0010, 5/registerWo, 1/1, 2/10, 3/000, 1/1, 5/registerYo, 1/--, 5/registerZo, 1/1 }


\paragraph{Syntax} \texttt{addbo} \emph{registerWo} = \emph{registerZo}, \emph{registerYo}

\paragraph{Description} The \emph{registerYo} and the \emph{registerZo} are interpreted as eight bytes packed into 64 bits. The \emph{registerYo} bytes and the \emph{registerZo} bytes are added. The eight resulting bytes are packed and stored into the \emph{registerWo}.


\paragraph{Execution} {Reservation table \textsf{ALU\_LITE} (Section~\ref{sec:reservation-ALU-LITE})}
~
{\begin{lstlisting}
stage RR:
  new argument3 = GPR[(registerYo<<1)+1];
  new argument2 = GPR[(registerZo<<1)+1];
  for (I = 0; I < 8; I += 1) {
    result1.8[I] = argument3.8[I] + argument2.8[I]
  };
stage E1:
  GPR[(registerWo<<1)+1] = result1;
\end{lstlisting}}
\newpage

\paragraph{Encoding} Format \textsf{ALU\_BOWRR0E.M} (Section~\ref{sec:format-ALU-BOWRR0E.M}), \verb|exucode2 = 0010|\\

\bitfields{ 1/P, 2/00, 2/-- --, 27/upper27, 1/1, 2/11, 1/1, 4/0010, 5/registerWe, 1/0, 2/10, 3/000, 1/1, 1/splat32, 5/lower5, 5/registerZe, 1/1 }


\paragraph{Syntax} \texttt{addbo} \emph{registerWe} = \emph{registerZe}, \emph{upper27\_\discretionary{}{}{}lower5}\emph{splat32}

\paragraph{Description} The \emph{upper27\_\discretionary{}{}{}lower5} and the \emph{registerZe} are interpreted as eight bytes packed into 64 bits. The \emph{upper27\_\discretionary{}{}{}lower5} bytes and the \emph{registerZe} bytes are added. The eight resulting bytes are packed and stored into the \emph{registerWe}.


\paragraph{Modifier(s)}
\noindent\begin{itemize}
\item splat32 (cf. splat32 in section \ref{par:modifier-splat32} on page \pageref{par:modifier-splat32})
\end{itemize}

\paragraph{Execution} {Reservation table \textsf{ALU\_LITE.X} (Section~\ref{sec:reservation-ALU-LITE.X})}
~
{\begin{lstlisting}
stage ID:
  new splat32 = _ZX_1(splat32);
stage RR:
  new argument3 = splat32 ? (_WX_32(upper27_lower5) << 32) | _ZX_32(_WX_32(upper27_lower5)) : _SX_32(_WX_32(upper27_lower5));
  new argument2 = GPR[registerZe<<1];
  for (I = 0; I < 8; I += 1) {
    result1.8[I] = argument3.8[I] + argument2.8[I]
  };
stage E1:
  GPR[registerWe<<1] = result1;
\end{lstlisting}}
\newpage

\paragraph{Encoding} Format \textsf{ALU\_BOWRR0O.M} (Section~\ref{sec:format-ALU-BOWRR0O.M}), \verb|exucode2 = 0010|\\

\bitfields{ 1/P, 2/00, 2/-- --, 27/upper27, 1/1, 2/11, 1/1, 4/0010, 5/registerWo, 1/1, 2/10, 3/000, 1/1, 1/splat32, 5/lower5, 5/registerZo, 1/1 }


\paragraph{Syntax} \texttt{addbo} \emph{registerWo} = \emph{registerZo}, \emph{upper27\_\discretionary{}{}{}lower5}\emph{splat32}

\paragraph{Description} The \emph{upper27\_\discretionary{}{}{}lower5} and the \emph{registerZo} are interpreted as eight bytes packed into 64 bits. The \emph{upper27\_\discretionary{}{}{}lower5} bytes and the \emph{registerZo} bytes are added. The eight resulting bytes are packed and stored into the \emph{registerWo}.


\paragraph{Modifier(s)}
\noindent\begin{itemize}
\item splat32 (cf. splat32 in section \ref{par:modifier-splat32} on page \pageref{par:modifier-splat32})
\end{itemize}

\paragraph{Execution} {Reservation table \textsf{ALU\_LITE.X} (Section~\ref{sec:reservation-ALU-LITE.X})}
~
{\begin{lstlisting}
stage ID:
  new splat32 = _ZX_1(splat32);
stage RR:
  new argument3 = splat32 ? (_WX_32(upper27_lower5) << 32) | _ZX_32(_WX_32(upper27_lower5)) : _SX_32(_WX_32(upper27_lower5));
  new argument2 = GPR[(registerZo<<1)+1];
  for (I = 0; I < 8; I += 1) {
    result1.8[I] = argument3.8[I] + argument2.8[I]
  };
stage E1:
  GPR[(registerWo<<1)+1] = result1;
\end{lstlisting}}
\newpage
\subsubsection[{\small ADDD: Add Double Word to Double Word}]{ADDD\hfill Add Double Word to Double Word}
\label{instr:ADDD}\index{addd}
 
\paragraph{Encoding} Format \textsf{ALU\_DWRR0} (Section~\ref{sec:format-ALU-DWRR0}), \verb|exucode2 = 0010|\\

\bitfields{ 1/P, 2/11, 1/0, 4/0010, 6/registerW, 2/10, 3/000, 1/0, 6/registerY, 6/registerZ }


\paragraph{Syntax} \texttt{addd} \emph{registerW} = \emph{registerZ}, \emph{registerY}

\paragraph{Description} The \emph{registerY} and the \emph{registerZ} are added. The result is stored into the \emph{registerW}.


\paragraph{Execution} {Reservation table \textsf{ALU\_TINY} (Section~\ref{sec:reservation-ALU-TINY})}
~
{\begin{lstlisting}
stage RR:
  new argument3 = GPR[registerY];
  new argument2 = GPR[registerZ];
  new result1 = argument3 + argument2;
stage E1:
  GPR[registerW] = result1;
\end{lstlisting}}
\newpage

\paragraph{Encoding} Format \textsf{ALU\_DWRR0.M} (Section~\ref{sec:format-ALU-DWRR0.M}), \verb|exucode2 = 0010|\\

\bitfields{ 1/P, 2/00, 2/-- --, 27/upper27, 1/1, 2/11, 1/0, 4/0010, 6/registerW, 2/10, 3/000, 1/0, 1/splat32, 5/lower5, 6/registerZ }


\paragraph{Syntax} \texttt{addd} \emph{registerW} = \emph{registerZ}, \emph{upper27\_\discretionary{}{}{}lower5}\emph{splat32}

\paragraph{Description} The \emph{upper27\_\discretionary{}{}{}lower5} and the \emph{registerZ} are added. The result is stored into the \emph{registerW}.


\paragraph{Modifier(s)}
\noindent\begin{itemize}
\item splat32 (cf. splat32 in section \ref{par:modifier-splat32} on page \pageref{par:modifier-splat32})
\end{itemize}

\paragraph{Execution} {Reservation table \textsf{ALU\_TINY.X} (Section~\ref{sec:reservation-ALU-TINY.X})}
~
{\begin{lstlisting}
stage ID:
  new splat32 = _ZX_1(splat32);
stage RR:
  new argument3 = splat32 ? (_WX_32(upper27_lower5) << 32) | _ZX_32(_WX_32(upper27_lower5)) : _SX_32(_WX_32(upper27_lower5));
  new argument2 = GPR[registerZ];
  new result1 = argument3 + argument2;
stage E1:
  GPR[registerW] = result1;
\end{lstlisting}}
\newpage

\paragraph{Encoding} Format \textsf{ALU\_DWRI} (Section~\ref{sec:format-ALU-DWRI}), \verb|exucode2 = 0010|\\

\bitfields{ 1/P, 2/11, 1/0, 4/0010, 6/registerW, 2/00, 10/signed10, 6/registerZ }


\paragraph{Syntax} \texttt{addd} \emph{registerW} = \emph{registerZ}, \emph{signed10}

\paragraph{Description} The \emph{signed10} and the \emph{registerZ} are added. The result is stored into the \emph{registerW}.


\paragraph{Execution} {Reservation table \textsf{ALU\_TINY} (Section~\ref{sec:reservation-ALU-TINY})}
~
{\begin{lstlisting}
stage RR:
  new argument3 = _SX_10(signed10);
  new argument2 = GPR[registerZ];
  new result1 = argument3 + argument2;
stage E1:
  GPR[registerW] = result1;
\end{lstlisting}}
\newpage

\paragraph{Encoding} Format \textsf{ALU\_DWRI.X} (Section~\ref{sec:format-ALU-DWRI.X}), \verb|exucode2 = 0010|\\

\bitfields{ 1/P, 2/00, 2/-- --, 27/upper27, 1/1, 2/11, 1/0, 4/0010, 6/registerW, 2/00, 10/lower10, 6/registerZ }


\paragraph{Syntax} \texttt{addd} \emph{registerW} = \emph{registerZ}, \emph{upper27\_\discretionary{}{}{}lower10}

\paragraph{Description} The \emph{upper27\_\discretionary{}{}{}lower10} and the \emph{registerZ} are added. The result is stored into the \emph{registerW}.


\paragraph{Execution} {Reservation table \textsf{ALU\_TINY.X} (Section~\ref{sec:reservation-ALU-TINY.X})}
~
{\begin{lstlisting}
stage RR:
  new argument3 = _SX_37(upper27_lower10);
  new argument2 = GPR[registerZ];
  new result1 = argument3 + argument2;
stage E1:
  GPR[registerW] = result1;
\end{lstlisting}}
\newpage

\paragraph{Encoding} Format \textsf{ALU\_DWRI.Y} (Section~\ref{sec:format-ALU-DWRI.Y}), \verb|exucode2 = 0010|\\

\bitfields{ 1/P, 2/00, 2/-- --, 27/extend27, 1/1, 2/00, 2/-- --, 27/upper27, 1/1, 2/11, 1/0, 4/0010, 6/registerW, 2/00, 10/lower10, 6/registerZ }


\paragraph{Syntax} \texttt{addd} \emph{registerW} = \emph{registerZ}, \emph{extend27\_\discretionary{}{}{}upper27\_\discretionary{}{}{}lower10}

\paragraph{Description} The \emph{extend27\_\discretionary{}{}{}upper27\_\discretionary{}{}{}lower10} and the \emph{registerZ} are added. The result is stored into the \emph{registerW}.


\paragraph{Execution} {Reservation table \textsf{ALU\_TINY.Y} (Section~\ref{sec:reservation-ALU-TINY.Y})}
~
{\begin{lstlisting}
stage RR:
  new argument3 = _WX_64(extend27_upper27_lower10);
  new argument2 = GPR[registerZ];
  new result1 = argument3 + argument2;
stage E1:
  GPR[registerW] = result1;
\end{lstlisting}}
\newpage
\subsubsection[{\small ADDHQ: Add Half Word Quadruple}]{ADDHQ\hfill Add Half Word Quadruple}
\label{instr:ADDHQ}\index{addhq}
 
\paragraph{Encoding} Format \textsf{ALU\_HQWRR0E} (Section~\ref{sec:format-ALU-HQWRR0E}), \verb|exucode2 = 0010|\\

\bitfields{ 1/P, 2/11, 1/1, 4/0010, 5/registerWe, 1/0, 2/10, 3/000, 1/1, 5/registerYe, 1/--, 5/registerZe, 1/0 }


\paragraph{Syntax} \texttt{addhq} \emph{registerWe} = \emph{registerZe}, \emph{registerYe}

\paragraph{Description} The \emph{registerYe} and the \emph{registerZe} are interpreted as four half words packed into 64 bits. The \emph{registerYe} half words and the \emph{registerZe} half words are added. The four resulting half words are packed and stored into the \emph{registerWe}.


\paragraph{Execution} {Reservation table \textsf{ALU\_LITE} (Section~\ref{sec:reservation-ALU-LITE})}
~
{\begin{lstlisting}
stage RR:
  new argument3 = GPR[registerYe<<1];
  new argument2 = GPR[registerZe<<1];
  for (I = 0; I < 4; I += 1) {
    result1.16[I] = argument3.16[I] + argument2.16[I]
  };
stage E1:
  GPR[registerWe<<1] = result1;
\end{lstlisting}}
\newpage

\paragraph{Encoding} Format \textsf{ALU\_HQWRR0O} (Section~\ref{sec:format-ALU-HQWRR0O}), \verb|exucode2 = 0010|\\

\bitfields{ 1/P, 2/11, 1/1, 4/0010, 5/registerWo, 1/1, 2/10, 3/000, 1/1, 5/registerYo, 1/--, 5/registerZo, 1/0 }


\paragraph{Syntax} \texttt{addhq} \emph{registerWo} = \emph{registerZo}, \emph{registerYo}

\paragraph{Description} The \emph{registerYo} and the \emph{registerZo} are interpreted as four half words packed into 64 bits. The \emph{registerYo} half words and the \emph{registerZo} half words are added. The four resulting half words are packed and stored into the \emph{registerWo}.


\paragraph{Execution} {Reservation table \textsf{ALU\_LITE} (Section~\ref{sec:reservation-ALU-LITE})}
~
{\begin{lstlisting}
stage RR:
  new argument3 = GPR[(registerYo<<1)+1];
  new argument2 = GPR[(registerZo<<1)+1];
  for (I = 0; I < 4; I += 1) {
    result1.16[I] = argument3.16[I] + argument2.16[I]
  };
stage E1:
  GPR[(registerWo<<1)+1] = result1;
\end{lstlisting}}
\newpage

\paragraph{Encoding} Format \textsf{ALU\_HQWRR0E.M} (Section~\ref{sec:format-ALU-HQWRR0E.M}), \verb|exucode2 = 0010|\\

\bitfields{ 1/P, 2/00, 2/-- --, 27/upper27, 1/1, 2/11, 1/1, 4/0010, 5/registerWe, 1/0, 2/10, 3/000, 1/1, 1/splat32, 5/lower5, 5/registerZe, 1/0 }


\paragraph{Syntax} \texttt{addhq} \emph{registerWe} = \emph{registerZe}, \emph{upper27\_\discretionary{}{}{}lower5}\emph{splat32}

\paragraph{Description} The \emph{upper27\_\discretionary{}{}{}lower5} and the \emph{registerZe} are interpreted as four half words packed into 64 bits. The \emph{upper27\_\discretionary{}{}{}lower5} half words and the \emph{registerZe} half words are added. The four resulting half words are packed and stored into the \emph{registerWe}.


\paragraph{Modifier(s)}
\noindent\begin{itemize}
\item splat32 (cf. splat32 in section \ref{par:modifier-splat32} on page \pageref{par:modifier-splat32})
\end{itemize}

\paragraph{Execution} {Reservation table \textsf{ALU\_LITE.X} (Section~\ref{sec:reservation-ALU-LITE.X})}
~
{\begin{lstlisting}
stage ID:
  new splat32 = _ZX_1(splat32);
stage RR:
  new argument3 = splat32 ? (_WX_32(upper27_lower5) << 32) | _ZX_32(_WX_32(upper27_lower5)) : _SX_32(_WX_32(upper27_lower5));
  new argument2 = GPR[registerZe<<1];
  for (I = 0; I < 4; I += 1) {
    result1.16[I] = argument3.16[I] + argument2.16[I]
  };
stage E1:
  GPR[registerWe<<1] = result1;
\end{lstlisting}}
\newpage

\paragraph{Encoding} Format \textsf{ALU\_HQWRR0O.M} (Section~\ref{sec:format-ALU-HQWRR0O.M}), \verb|exucode2 = 0010|\\

\bitfields{ 1/P, 2/00, 2/-- --, 27/upper27, 1/1, 2/11, 1/1, 4/0010, 5/registerWo, 1/1, 2/10, 3/000, 1/1, 1/splat32, 5/lower5, 5/registerZo, 1/0 }


\paragraph{Syntax} \texttt{addhq} \emph{registerWo} = \emph{registerZo}, \emph{upper27\_\discretionary{}{}{}lower5}\emph{splat32}

\paragraph{Description} The \emph{upper27\_\discretionary{}{}{}lower5} and the \emph{registerZo} are interpreted as four half words packed into 64 bits. The \emph{upper27\_\discretionary{}{}{}lower5} half words and the \emph{registerZo} half words are added. The four resulting half words are packed and stored into the \emph{registerWo}.


\paragraph{Modifier(s)}
\noindent\begin{itemize}
\item splat32 (cf. splat32 in section \ref{par:modifier-splat32} on page \pageref{par:modifier-splat32})
\end{itemize}

\paragraph{Execution} {Reservation table \textsf{ALU\_LITE.X} (Section~\ref{sec:reservation-ALU-LITE.X})}
~
{\begin{lstlisting}
stage ID:
  new splat32 = _ZX_1(splat32);
stage RR:
  new argument3 = splat32 ? (_WX_32(upper27_lower5) << 32) | _ZX_32(_WX_32(upper27_lower5)) : _SX_32(_WX_32(upper27_lower5));
  new argument2 = GPR[(registerZo<<1)+1];
  for (I = 0; I < 4; I += 1) {
    result1.16[I] = argument3.16[I] + argument2.16[I]
  };
stage E1:
  GPR[(registerWo<<1)+1] = result1;
\end{lstlisting}}
\newpage
\subsubsection[{\small ADDQ: Add Quadruple Word to Quadruple Word}]{ADDQ\hfill Add Quadruple Word to Quadruple Word}
\label{instr:ADDQ}\index{addq}
 
\paragraph{Encoding} Format \textsf{ALU\_QWRR0} (Section~\ref{sec:format-ALU-QWRR0}), \verb|exucode2 = 0010|\\

\bitfields{ 1/P, 2/11, 1/0, 4/0010, 5/registerM, 1/--, 2/10, 3/000, 1/1, 5/registerO, 1/--, 5/registerP, 1/-- }


\paragraph{Syntax} \texttt{addq} \emph{registerM} = \emph{registerP}, \emph{registerO}

\paragraph{Description} The \emph{registerO} and the \emph{registerP} are added. The result is stored into the \emph{registerM}.


\paragraph{Execution} {Reservation table \textsf{ALU\_LITE} (Section~\ref{sec:reservation-ALU-LITE})}
~
{\begin{lstlisting}
stage RR:
  new argument3 = PGR[registerO];
  new argument2 = PGR[registerP];
  new result1 = argument3 + argument2;
stage E1:
  PGR[registerM] = result1;
\end{lstlisting}}
\newpage

\paragraph{Encoding} Format \textsf{ALU\_QWRR0.M} (Section~\ref{sec:format-ALU-QWRR0.M}), \verb|exucode2 = 0010|\\

\bitfields{ 1/P, 2/00, 2/-- --, 27/upper27, 1/1, 2/11, 1/0, 4/0010, 5/registerM, 1/--, 2/10, 3/000, 1/1, 1/splat32, 5/lower5, 5/registerP, 1/1 }


\paragraph{Syntax} \texttt{addq} \emph{registerM} = \emph{registerP}, \emph{upper27\_\discretionary{}{}{}lower5}\emph{splat32}

\paragraph{Description} The \emph{upper27\_\discretionary{}{}{}lower5} and the \emph{registerP} are added. The result is stored into the \emph{registerM}.


\paragraph{Modifier(s)}
\noindent\begin{itemize}
\item splat32 (cf. splat32 in section \ref{par:modifier-splat32} on page \pageref{par:modifier-splat32})
\end{itemize}

\paragraph{Execution} {Reservation table \textsf{ALU\_LITE.X} (Section~\ref{sec:reservation-ALU-LITE.X})}
~
{\begin{lstlisting}
stage ID:
  new splat32 = _ZX_1(splat32);
stage RR:
  new argument3 = splat32 ? (_WX_32(upper27_lower5) << 32) | _ZX_32(_WX_32(upper27_lower5)) : _SX_32(_WX_32(upper27_lower5));
  new argument2 = PGR[registerP];
  new result1 = argument3 + argument2;
stage E1:
  PGR[registerM] = result1;
\end{lstlisting}}
\newpage
\subsubsection[{\small ADDSBO: Add Saturated Byte Octuple}]{ADDSBO\hfill Add Saturated Byte Octuple}
\label{instr:ADDSBO}\index{addsbo}
 
\paragraph{Encoding} Format \textsf{ALU\_BOWRR1E} (Section~\ref{sec:format-ALU-BOWRR1E}), \verb|exucode2 = 1000|\\

\bitfields{ 1/P, 2/11, 1/1, 4/1000, 5/registerWe, 1/0, 2/10, 3/101, 1/1, 5/registerYe, 1/--, 5/registerZe, 1/1 }


\paragraph{Syntax} \texttt{addsbo} \emph{registerWe} = \emph{registerZe}, \emph{registerYe}

\paragraph{Description} The \emph{registerYe} and the \emph{registerZe} are interpreted as eight bytes packed into 64 bits. The \emph{registerYe} bytes and the \emph{registerZe} bytes are added with saturation to 8 bits. The eight resulting bytes are packed and stored into the \emph{registerWe}.


\paragraph{Execution} {Reservation table \textsf{ALU\_LITE} (Section~\ref{sec:reservation-ALU-LITE})}
~
{\begin{lstlisting}
stage RR:
  new argument3 = GPR[registerYe<<1];
  new argument2 = GPR[registerZe<<1];
  for (I = 0; I < 8; I += 1) {
    result1.8[I] = _SAT_8(_SX_8(argument3.8[I]) + _SX_8(argument2.8[I]))
  };
stage E1:
  GPR[registerWe<<1] = result1;
\end{lstlisting}}
\newpage

\paragraph{Encoding} Format \textsf{ALU\_BOWRR1O} (Section~\ref{sec:format-ALU-BOWRR1O}), \verb|exucode2 = 1000|\\

\bitfields{ 1/P, 2/11, 1/1, 4/1000, 5/registerWo, 1/1, 2/10, 3/101, 1/1, 5/registerYo, 1/--, 5/registerZo, 1/1 }


\paragraph{Syntax} \texttt{addsbo} \emph{registerWo} = \emph{registerZo}, \emph{registerYo}

\paragraph{Description} The \emph{registerYo} and the \emph{registerZo} are interpreted as eight bytes packed into 64 bits. The \emph{registerYo} bytes and the \emph{registerZo} bytes are added with saturation to 8 bits. The eight resulting bytes are packed and stored into the \emph{registerWo}.


\paragraph{Execution} {Reservation table \textsf{ALU\_LITE} (Section~\ref{sec:reservation-ALU-LITE})}
~
{\begin{lstlisting}
stage RR:
  new argument3 = GPR[(registerYo<<1)+1];
  new argument2 = GPR[(registerZo<<1)+1];
  for (I = 0; I < 8; I += 1) {
    result1.8[I] = _SAT_8(_SX_8(argument3.8[I]) + _SX_8(argument2.8[I]))
  };
stage E1:
  GPR[(registerWo<<1)+1] = result1;
\end{lstlisting}}
\newpage

\paragraph{Encoding} Format \textsf{ALU\_BOWRR1E.M} (Section~\ref{sec:format-ALU-BOWRR1E.M}), \verb|exucode2 = 1000|\\

\bitfields{ 1/P, 2/00, 2/-- --, 27/upper27, 1/1, 2/11, 1/1, 4/1000, 5/registerWe, 1/0, 2/10, 3/101, 1/1, 1/splat32, 5/lower5, 5/registerZe, 1/1 }


\paragraph{Syntax} \texttt{addsbo} \emph{registerWe} = \emph{registerZe}, \emph{upper27\_\discretionary{}{}{}lower5}\emph{splat32}

\paragraph{Description} The \emph{upper27\_\discretionary{}{}{}lower5} and the \emph{registerZe} are interpreted as eight bytes packed into 64 bits. The \emph{upper27\_\discretionary{}{}{}lower5} bytes and the \emph{registerZe} bytes are added with saturation to 8 bits. The eight resulting bytes are packed and stored into the \emph{registerWe}.


\paragraph{Modifier(s)}
\noindent\begin{itemize}
\item splat32 (cf. splat32 in section \ref{par:modifier-splat32} on page \pageref{par:modifier-splat32})
\end{itemize}

\paragraph{Execution} {Reservation table \textsf{ALU\_LITE.X} (Section~\ref{sec:reservation-ALU-LITE.X})}
~
{\begin{lstlisting}
stage ID:
  new splat32 = _ZX_1(splat32);
stage RR:
  new argument3 = splat32 ? (_WX_32(upper27_lower5) << 32) | _ZX_32(_WX_32(upper27_lower5)) : _SX_32(_WX_32(upper27_lower5));
  new argument2 = GPR[registerZe<<1];
  for (I = 0; I < 8; I += 1) {
    result1.8[I] = _SAT_8(_SX_8(argument3.8[I]) + _SX_8(argument2.8[I]))
  };
stage E1:
  GPR[registerWe<<1] = result1;
\end{lstlisting}}
\newpage

\paragraph{Encoding} Format \textsf{ALU\_BOWRR1O.M} (Section~\ref{sec:format-ALU-BOWRR1O.M}), \verb|exucode2 = 1000|\\

\bitfields{ 1/P, 2/00, 2/-- --, 27/upper27, 1/1, 2/11, 1/1, 4/1000, 5/registerWo, 1/1, 2/10, 3/101, 1/1, 1/splat32, 5/lower5, 5/registerZo, 1/1 }


\paragraph{Syntax} \texttt{addsbo} \emph{registerWo} = \emph{registerZo}, \emph{upper27\_\discretionary{}{}{}lower5}\emph{splat32}

\paragraph{Description} The \emph{upper27\_\discretionary{}{}{}lower5} and the \emph{registerZo} are interpreted as eight bytes packed into 64 bits. The \emph{upper27\_\discretionary{}{}{}lower5} bytes and the \emph{registerZo} bytes are added with saturation to 8 bits. The eight resulting bytes are packed and stored into the \emph{registerWo}.


\paragraph{Modifier(s)}
\noindent\begin{itemize}
\item splat32 (cf. splat32 in section \ref{par:modifier-splat32} on page \pageref{par:modifier-splat32})
\end{itemize}

\paragraph{Execution} {Reservation table \textsf{ALU\_LITE.X} (Section~\ref{sec:reservation-ALU-LITE.X})}
~
{\begin{lstlisting}
stage ID:
  new splat32 = _ZX_1(splat32);
stage RR:
  new argument3 = splat32 ? (_WX_32(upper27_lower5) << 32) | _ZX_32(_WX_32(upper27_lower5)) : _SX_32(_WX_32(upper27_lower5));
  new argument2 = GPR[(registerZo<<1)+1];
  for (I = 0; I < 8; I += 1) {
    result1.8[I] = _SAT_8(_SX_8(argument3.8[I]) + _SX_8(argument2.8[I]))
  };
stage E1:
  GPR[(registerWo<<1)+1] = result1;
\end{lstlisting}}
\newpage
\subsubsection[{\small ADDSD: Add Saturated Double Word}]{ADDSD\hfill Add Saturated Double Word}
\label{instr:ADDSD}\index{addsd}
 
\paragraph{Encoding} Format \textsf{ALU\_DWRR1} (Section~\ref{sec:format-ALU-DWRR1}), \verb|exucode2 = 1000|\\

\bitfields{ 1/P, 2/11, 1/0, 4/1000, 6/registerW, 2/10, 3/101, 1/0, 6/registerY, 6/registerZ }


\paragraph{Syntax} \texttt{addsd} \emph{registerW} = \emph{registerZ}, \emph{registerY}

\paragraph{Description} The \emph{registerY} and the \emph{registerZ} are added with 64-bit saturation.


\paragraph{Execution} {Reservation table \textsf{ALU\_TINY} (Section~\ref{sec:reservation-ALU-TINY})}
~
{\begin{lstlisting}
stage RR:
  new argument3 = GPR[registerY];
  new argument2 = GPR[registerZ];
  new result1 = _SAT_64(_SX_64(argument3) + _SX_64(argument2));
stage E1:
  GPR[registerW] = result1;
\end{lstlisting}}
\newpage

\paragraph{Encoding} Format \textsf{ALU\_DWRR1.M} (Section~\ref{sec:format-ALU-DWRR1.M}), \verb|exucode2 = 1000|\\

\bitfields{ 1/P, 2/00, 2/-- --, 27/upper27, 1/1, 2/11, 1/0, 4/1000, 6/registerW, 2/10, 3/101, 1/0, 1/splat32, 5/lower5, 6/registerZ }


\paragraph{Syntax} \texttt{addsd} \emph{registerW} = \emph{registerZ}, \emph{upper27\_\discretionary{}{}{}lower5}\emph{splat32}

\paragraph{Description} The \emph{upper27\_\discretionary{}{}{}lower5} and the \emph{registerZ} are added with 64-bit saturation.


\paragraph{Modifier(s)}
\noindent\begin{itemize}
\item splat32 (cf. splat32 in section \ref{par:modifier-splat32} on page \pageref{par:modifier-splat32})
\end{itemize}

\paragraph{Execution} {Reservation table \textsf{ALU\_TINY.X} (Section~\ref{sec:reservation-ALU-TINY.X})}
~
{\begin{lstlisting}
stage ID:
  new splat32 = _ZX_1(splat32);
stage RR:
  new argument3 = splat32 ? (_WX_32(upper27_lower5) << 32) | _ZX_32(_WX_32(upper27_lower5)) : _SX_32(_WX_32(upper27_lower5));
  new argument2 = GPR[registerZ];
  new result1 = _SAT_64(_SX_64(argument3) + _SX_64(argument2));
stage E1:
  GPR[registerW] = result1;
\end{lstlisting}}
\newpage
\subsubsection[{\small ADDSHQ: Add Saturated Half Word Quadruple}]{ADDSHQ\hfill Add Saturated Half Word Quadruple}
\label{instr:ADDSHQ}\index{addshq}
 
\paragraph{Encoding} Format \textsf{ALU\_HQWRR1E} (Section~\ref{sec:format-ALU-HQWRR1E}), \verb|exucode2 = 1000|\\

\bitfields{ 1/P, 2/11, 1/1, 4/1000, 5/registerWe, 1/0, 2/10, 3/101, 1/1, 5/registerYe, 1/--, 5/registerZe, 1/0 }


\paragraph{Syntax} \texttt{addshq} \emph{registerWe} = \emph{registerZe}, \emph{registerYe}

\paragraph{Description} The \emph{registerYe} and the \emph{registerZe} are interpreted as four half words packed into 64 bits. The \emph{registerYe} half words and the \emph{registerZe} half words are added with saturation to 16 bits. The four resulting half words are packed and stored into the \emph{registerWe}.


\paragraph{Execution} {Reservation table \textsf{ALU\_LITE} (Section~\ref{sec:reservation-ALU-LITE})}
~
{\begin{lstlisting}
stage RR:
  new argument3 = GPR[registerYe<<1];
  new argument2 = GPR[registerZe<<1];
  for (I = 0; I < 4; I += 1) {
    result1.16[I] = _SAT_16(_SX_16(argument3.16[I]) + _SX_16(argument2.16[I]))
  };
stage E1:
  GPR[registerWe<<1] = result1;
\end{lstlisting}}
\newpage

\paragraph{Encoding} Format \textsf{ALU\_HQWRR1O} (Section~\ref{sec:format-ALU-HQWRR1O}), \verb|exucode2 = 1000|\\

\bitfields{ 1/P, 2/11, 1/1, 4/1000, 5/registerWo, 1/1, 2/10, 3/101, 1/1, 5/registerYo, 1/--, 5/registerZo, 1/0 }


\paragraph{Syntax} \texttt{addshq} \emph{registerWo} = \emph{registerZo}, \emph{registerYo}

\paragraph{Description} The \emph{registerYo} and the \emph{registerZo} are interpreted as four half words packed into 64 bits. The \emph{registerYo} half words and the \emph{registerZo} half words are added with saturation to 16 bits. The four resulting half words are packed and stored into the \emph{registerWo}.


\paragraph{Execution} {Reservation table \textsf{ALU\_LITE} (Section~\ref{sec:reservation-ALU-LITE})}
~
{\begin{lstlisting}
stage RR:
  new argument3 = GPR[(registerYo<<1)+1];
  new argument2 = GPR[(registerZo<<1)+1];
  for (I = 0; I < 4; I += 1) {
    result1.16[I] = _SAT_16(_SX_16(argument3.16[I]) + _SX_16(argument2.16[I]))
  };
stage E1:
  GPR[(registerWo<<1)+1] = result1;
\end{lstlisting}}
\newpage

\paragraph{Encoding} Format \textsf{ALU\_HQWRR1E.M} (Section~\ref{sec:format-ALU-HQWRR1E.M}), \verb|exucode2 = 1000|\\

\bitfields{ 1/P, 2/00, 2/-- --, 27/upper27, 1/1, 2/11, 1/1, 4/1000, 5/registerWe, 1/0, 2/10, 3/101, 1/1, 1/splat32, 5/lower5, 5/registerZe, 1/0 }


\paragraph{Syntax} \texttt{addshq} \emph{registerWe} = \emph{registerZe}, \emph{upper27\_\discretionary{}{}{}lower5}\emph{splat32}

\paragraph{Description} The \emph{upper27\_\discretionary{}{}{}lower5} and the \emph{registerZe} are interpreted as four half words packed into 64 bits. The \emph{upper27\_\discretionary{}{}{}lower5} half words and the \emph{registerZe} half words are added with saturation to 16 bits. The four resulting half words are packed and stored into the \emph{registerWe}.


\paragraph{Modifier(s)}
\noindent\begin{itemize}
\item splat32 (cf. splat32 in section \ref{par:modifier-splat32} on page \pageref{par:modifier-splat32})
\end{itemize}

\paragraph{Execution} {Reservation table \textsf{ALU\_LITE.X} (Section~\ref{sec:reservation-ALU-LITE.X})}
~
{\begin{lstlisting}
stage ID:
  new splat32 = _ZX_1(splat32);
stage RR:
  new argument3 = splat32 ? (_WX_32(upper27_lower5) << 32) | _ZX_32(_WX_32(upper27_lower5)) : _SX_32(_WX_32(upper27_lower5));
  new argument2 = GPR[registerZe<<1];
  for (I = 0; I < 4; I += 1) {
    result1.16[I] = _SAT_16(_SX_16(argument3.16[I]) + _SX_16(argument2.16[I]))
  };
stage E1:
  GPR[registerWe<<1] = result1;
\end{lstlisting}}
\newpage

\paragraph{Encoding} Format \textsf{ALU\_HQWRR1O.M} (Section~\ref{sec:format-ALU-HQWRR1O.M}), \verb|exucode2 = 1000|\\

\bitfields{ 1/P, 2/00, 2/-- --, 27/upper27, 1/1, 2/11, 1/1, 4/1000, 5/registerWo, 1/1, 2/10, 3/101, 1/1, 1/splat32, 5/lower5, 5/registerZo, 1/0 }


\paragraph{Syntax} \texttt{addshq} \emph{registerWo} = \emph{registerZo}, \emph{upper27\_\discretionary{}{}{}lower5}\emph{splat32}

\paragraph{Description} The \emph{upper27\_\discretionary{}{}{}lower5} and the \emph{registerZo} are interpreted as four half words packed into 64 bits. The \emph{upper27\_\discretionary{}{}{}lower5} half words and the \emph{registerZo} half words are added with saturation to 16 bits. The four resulting half words are packed and stored into the \emph{registerWo}.


\paragraph{Modifier(s)}
\noindent\begin{itemize}
\item splat32 (cf. splat32 in section \ref{par:modifier-splat32} on page \pageref{par:modifier-splat32})
\end{itemize}

\paragraph{Execution} {Reservation table \textsf{ALU\_LITE.X} (Section~\ref{sec:reservation-ALU-LITE.X})}
~
{\begin{lstlisting}
stage ID:
  new splat32 = _ZX_1(splat32);
stage RR:
  new argument3 = splat32 ? (_WX_32(upper27_lower5) << 32) | _ZX_32(_WX_32(upper27_lower5)) : _SX_32(_WX_32(upper27_lower5));
  new argument2 = GPR[(registerZo<<1)+1];
  for (I = 0; I < 4; I += 1) {
    result1.16[I] = _SAT_16(_SX_16(argument3.16[I]) + _SX_16(argument2.16[I]))
  };
stage E1:
  GPR[(registerWo<<1)+1] = result1;
\end{lstlisting}}
\newpage
\subsubsection[{\small ADDSW: Add Saturated Word}]{ADDSW\hfill Add Saturated Word}
\label{instr:ADDSW}\index{addsw}
 
\paragraph{Encoding} Format \textsf{ALU\_WWRR1} (Section~\ref{sec:format-ALU-WWRR1}), \verb|exucode2 = 1000|\\

\bitfields{ 1/P, 2/11, 1/0, 4/1000, 6/registerW, 2/11, 3/101, 1/signextw, 6/registerY, 6/registerZ }


\paragraph{Syntax} \texttt{addsw}\emph{signextw} \emph{registerW} = \emph{registerZ}, \emph{registerY}

\paragraph{Description} The \emph{registerY} and the \emph{registerZ} are added with 32-bit saturation. The result is stored into the \emph{registerW}.


\paragraph{Modifier(s)}
\noindent\begin{itemize}
\item signextw (cf. signextw in section \ref{par:modifier-signextw} on page \pageref{par:modifier-signextw})
\end{itemize}

\paragraph{Execution} {Reservation table \textsf{ALU\_TINY} (Section~\ref{sec:reservation-ALU-TINY})}
~
{\begin{lstlisting}
stage ID:
  new signextw = _ZX_1(signextw);
stage RR:
  new argument3 = GPR[registerY];
  new argument2 = GPR[registerZ];
  new result1 = _SAT_32(_SX_32(argument3) + _SX_32(argument2));
stage E1:
  GPR[registerW] = signextw ? _SX_32(result1) : _ZX_32(result1);
\end{lstlisting}}
\newpage

\paragraph{Encoding} Format \textsf{ALU\_WWRR1.W} (Section~\ref{sec:format-ALU-WWRR1.W}), \verb|exucode2 = 1000|\\

\bitfields{ 1/P, 2/00, 2/-- --, 27/upper27, 1/1, 2/11, 1/0, 4/1000, 6/registerW, 2/11, 3/101, 1/signextw, 1/0, 5/lower5, 6/registerZ }


\paragraph{Syntax} \texttt{addsw}\emph{signextw} \emph{registerW} = \emph{registerZ}, \emph{upper27\_\discretionary{}{}{}lower5}

\paragraph{Description} The \emph{upper27\_\discretionary{}{}{}lower5} and the \emph{registerZ} are added with 32-bit saturation. The result is stored into the \emph{registerW}.


\paragraph{Modifier(s)}
\noindent\begin{itemize}
\item signextw (cf. signextw in section \ref{par:modifier-signextw} on page \pageref{par:modifier-signextw})
\end{itemize}

\paragraph{Execution} {Reservation table \textsf{ALU\_TINY.X} (Section~\ref{sec:reservation-ALU-TINY.X})}
~
{\begin{lstlisting}
stage ID:
  new signextw = _ZX_1(signextw);
stage RR:
  new argument3 = _WX_32(upper27_lower5);
  new argument2 = GPR[registerZ];
  new result1 = _SAT_32(_SX_32(argument3) + _SX_32(argument2));
stage E1:
  GPR[registerW] = signextw ? _SX_32(result1) : _ZX_32(result1);
\end{lstlisting}}
\newpage
\subsubsection[{\small ADDSWP: Add Saturated Word Pair}]{ADDSWP\hfill Add Saturated Word Pair}
\label{instr:ADDSWP}\index{addswp}
 
\paragraph{Encoding} Format \textsf{ALU\_WPWRR1E} (Section~\ref{sec:format-ALU-WPWRR1E}), \verb|exucode2 = 1000|\\

\bitfields{ 1/P, 2/11, 1/1, 4/1000, 5/registerWe, 1/0, 2/10, 3/101, 1/0, 5/registerYe, 1/--, 5/registerZe, 1/1 }


\paragraph{Syntax} \texttt{addswp} \emph{registerWe} = \emph{registerZe}, \emph{registerYe}

\paragraph{Description} The \emph{registerYe} and the \emph{registerZe} are interpreted as two words packed into 64 bits. The \emph{registerYe} words and the \emph{registerZe} words are added with saturation to 32 bits. The two resulting words are packed and stored into the \emph{registerWe}.


\paragraph{Execution} {Reservation table \textsf{ALU\_LITE} (Section~\ref{sec:reservation-ALU-LITE})}
~
{\begin{lstlisting}
stage RR:
  new argument3 = GPR[registerYe<<1];
  new argument2 = GPR[registerZe<<1];
  for (I = 0; I < 2; I += 1) {
    result1.32[I] = _SAT_32(_SX_32(argument3.32[I]) + _SX_32(argument2.32[I]))
  };
stage E1:
  GPR[registerWe<<1] = result1;
\end{lstlisting}}
\newpage

\paragraph{Encoding} Format \textsf{ALU\_WPWRR1O} (Section~\ref{sec:format-ALU-WPWRR1O}), \verb|exucode2 = 1000|\\

\bitfields{ 1/P, 2/11, 1/1, 4/1000, 5/registerWo, 1/1, 2/10, 3/101, 1/0, 5/registerYo, 1/--, 5/registerZo, 1/1 }


\paragraph{Syntax} \texttt{addswp} \emph{registerWo} = \emph{registerZo}, \emph{registerYo}

\paragraph{Description} The \emph{registerYo} and the \emph{registerZo} are interpreted as two words packed into 64 bits. The \emph{registerYo} words and the \emph{registerZo} words are added with saturation to 32 bits. The two resulting words are packed and stored into the \emph{registerWo}.


\paragraph{Execution} {Reservation table \textsf{ALU\_LITE} (Section~\ref{sec:reservation-ALU-LITE})}
~
{\begin{lstlisting}
stage RR:
  new argument3 = GPR[(registerYo<<1)+1];
  new argument2 = GPR[(registerZo<<1)+1];
  for (I = 0; I < 2; I += 1) {
    result1.32[I] = _SAT_32(_SX_32(argument3.32[I]) + _SX_32(argument2.32[I]))
  };
stage E1:
  GPR[(registerWo<<1)+1] = result1;
\end{lstlisting}}
\newpage

\paragraph{Encoding} Format \textsf{ALU\_WPWRR1E.M} (Section~\ref{sec:format-ALU-WPWRR1E.M}), \verb|exucode2 = 1000|\\

\bitfields{ 1/P, 2/00, 2/-- --, 27/upper27, 1/1, 2/11, 1/1, 4/1000, 5/registerWe, 1/0, 2/10, 3/101, 1/0, 1/splat32, 5/lower5, 5/registerZe, 1/1 }


\paragraph{Syntax} \texttt{addswp} \emph{registerWe} = \emph{registerZe}, \emph{upper27\_\discretionary{}{}{}lower5}\emph{splat32}

\paragraph{Description} The \emph{upper27\_\discretionary{}{}{}lower5} and the \emph{registerZe} are interpreted as two words packed into 64 bits. The \emph{upper27\_\discretionary{}{}{}lower5} words and the \emph{registerZe} words are added with saturation to 32 bits. The two resulting words are packed and stored into the \emph{registerWe}.


\paragraph{Modifier(s)}
\noindent\begin{itemize}
\item splat32 (cf. splat32 in section \ref{par:modifier-splat32} on page \pageref{par:modifier-splat32})
\end{itemize}

\paragraph{Execution} {Reservation table \textsf{ALU\_LITE.X} (Section~\ref{sec:reservation-ALU-LITE.X})}
~
{\begin{lstlisting}
stage ID:
  new splat32 = _ZX_1(splat32);
stage RR:
  new argument3 = splat32 ? (_WX_32(upper27_lower5) << 32) | _ZX_32(_WX_32(upper27_lower5)) : _SX_32(_WX_32(upper27_lower5));
  new argument2 = GPR[registerZe<<1];
  for (I = 0; I < 2; I += 1) {
    result1.32[I] = _SAT_32(_SX_32(argument3.32[I]) + _SX_32(argument2.32[I]))
  };
stage E1:
  GPR[registerWe<<1] = result1;
\end{lstlisting}}
\newpage

\paragraph{Encoding} Format \textsf{ALU\_WPWRR1O.M} (Section~\ref{sec:format-ALU-WPWRR1O.M}), \verb|exucode2 = 1000|\\

\bitfields{ 1/P, 2/00, 2/-- --, 27/upper27, 1/1, 2/11, 1/1, 4/1000, 5/registerWo, 1/1, 2/10, 3/101, 1/0, 1/splat32, 5/lower5, 5/registerZo, 1/1 }


\paragraph{Syntax} \texttt{addswp} \emph{registerWo} = \emph{registerZo}, \emph{upper27\_\discretionary{}{}{}lower5}\emph{splat32}

\paragraph{Description} The \emph{upper27\_\discretionary{}{}{}lower5} and the \emph{registerZo} are interpreted as two words packed into 64 bits. The \emph{upper27\_\discretionary{}{}{}lower5} words and the \emph{registerZo} words are added with saturation to 32 bits. The two resulting words are packed and stored into the \emph{registerWo}.


\paragraph{Modifier(s)}
\noindent\begin{itemize}
\item splat32 (cf. splat32 in section \ref{par:modifier-splat32} on page \pageref{par:modifier-splat32})
\end{itemize}

\paragraph{Execution} {Reservation table \textsf{ALU\_LITE.X} (Section~\ref{sec:reservation-ALU-LITE.X})}
~
{\begin{lstlisting}
stage ID:
  new splat32 = _ZX_1(splat32);
stage RR:
  new argument3 = splat32 ? (_WX_32(upper27_lower5) << 32) | _ZX_32(_WX_32(upper27_lower5)) : _SX_32(_WX_32(upper27_lower5));
  new argument2 = GPR[(registerZo<<1)+1];
  for (I = 0; I < 2; I += 1) {
    result1.32[I] = _SAT_32(_SX_32(argument3.32[I]) + _SX_32(argument2.32[I]))
  };
stage E1:
  GPR[(registerWo<<1)+1] = result1;
\end{lstlisting}}
\newpage
\subsubsection[{\small ADDUSBO: Add Unsigned Saturated Byte Octuple}]{ADDUSBO\hfill Add Unsigned Saturated Byte Octuple}
\label{instr:ADDUSBO}\index{addusbo}
 
\paragraph{Encoding} Format \textsf{ALU\_BOWRR1E} (Section~\ref{sec:format-ALU-BOWRR1E}), \verb|exucode2 = 1010|\\

\bitfields{ 1/P, 2/11, 1/1, 4/1010, 5/registerWe, 1/0, 2/10, 3/101, 1/1, 5/registerYe, 1/--, 5/registerZe, 1/1 }


\paragraph{Syntax} \texttt{addusbo} \emph{registerWe} = \emph{registerZe}, \emph{registerYe}

\paragraph{Description} The \emph{registerYe} and the \emph{registerZe} are interpreted as eight bytes packed into 64 bits. The \emph{registerYe} bytes and the \emph{registerZe} bytes are added with unsigned saturation to 8 bits. The eight resulting bytes are packed and stored into the \emph{registerWe}.


\paragraph{Execution} {Reservation table \textsf{ALU\_LITE} (Section~\ref{sec:reservation-ALU-LITE})}
~
{\begin{lstlisting}
stage RR:
  new argument3 = GPR[registerYe<<1];
  new argument2 = GPR[registerZe<<1];
  for (I = 0; I < 8; I += 1) {
    result1.8[I] = _SATU_8(_ZX_8(argument3.8[I]) + _ZX_8(argument2.8[I]))
  };
stage E1:
  GPR[registerWe<<1] = result1;
\end{lstlisting}}
\newpage

\paragraph{Encoding} Format \textsf{ALU\_BOWRR1O} (Section~\ref{sec:format-ALU-BOWRR1O}), \verb|exucode2 = 1010|\\

\bitfields{ 1/P, 2/11, 1/1, 4/1010, 5/registerWo, 1/1, 2/10, 3/101, 1/1, 5/registerYo, 1/--, 5/registerZo, 1/1 }


\paragraph{Syntax} \texttt{addusbo} \emph{registerWo} = \emph{registerZo}, \emph{registerYo}

\paragraph{Description} The \emph{registerYo} and the \emph{registerZo} are interpreted as eight bytes packed into 64 bits. The \emph{registerYo} bytes and the \emph{registerZo} bytes are added with unsigned saturation to 8 bits. The eight resulting bytes are packed and stored into the \emph{registerWo}.


\paragraph{Execution} {Reservation table \textsf{ALU\_LITE} (Section~\ref{sec:reservation-ALU-LITE})}
~
{\begin{lstlisting}
stage RR:
  new argument3 = GPR[(registerYo<<1)+1];
  new argument2 = GPR[(registerZo<<1)+1];
  for (I = 0; I < 8; I += 1) {
    result1.8[I] = _SATU_8(_ZX_8(argument3.8[I]) + _ZX_8(argument2.8[I]))
  };
stage E1:
  GPR[(registerWo<<1)+1] = result1;
\end{lstlisting}}
\newpage

\paragraph{Encoding} Format \textsf{ALU\_BOWRR1E.M} (Section~\ref{sec:format-ALU-BOWRR1E.M}), \verb|exucode2 = 1010|\\

\bitfields{ 1/P, 2/00, 2/-- --, 27/upper27, 1/1, 2/11, 1/1, 4/1010, 5/registerWe, 1/0, 2/10, 3/101, 1/1, 1/splat32, 5/lower5, 5/registerZe, 1/1 }


\paragraph{Syntax} \texttt{addusbo} \emph{registerWe} = \emph{registerZe}, \emph{upper27\_\discretionary{}{}{}lower5}\emph{splat32}

\paragraph{Description} The \emph{upper27\_\discretionary{}{}{}lower5} and the \emph{registerZe} are interpreted as eight bytes packed into 64 bits. The \emph{upper27\_\discretionary{}{}{}lower5} bytes and the \emph{registerZe} bytes are added with unsigned saturation to 8 bits. The eight resulting bytes are packed and stored into the \emph{registerWe}.


\paragraph{Modifier(s)}
\noindent\begin{itemize}
\item splat32 (cf. splat32 in section \ref{par:modifier-splat32} on page \pageref{par:modifier-splat32})
\end{itemize}

\paragraph{Execution} {Reservation table \textsf{ALU\_LITE.X} (Section~\ref{sec:reservation-ALU-LITE.X})}
~
{\begin{lstlisting}
stage ID:
  new splat32 = _ZX_1(splat32);
stage RR:
  new argument3 = splat32 ? (_WX_32(upper27_lower5) << 32) | _ZX_32(_WX_32(upper27_lower5)) : _SX_32(_WX_32(upper27_lower5));
  new argument2 = GPR[registerZe<<1];
  for (I = 0; I < 8; I += 1) {
    result1.8[I] = _SATU_8(_ZX_8(argument3.8[I]) + _ZX_8(argument2.8[I]))
  };
stage E1:
  GPR[registerWe<<1] = result1;
\end{lstlisting}}
\newpage

\paragraph{Encoding} Format \textsf{ALU\_BOWRR1O.M} (Section~\ref{sec:format-ALU-BOWRR1O.M}), \verb|exucode2 = 1010|\\

\bitfields{ 1/P, 2/00, 2/-- --, 27/upper27, 1/1, 2/11, 1/1, 4/1010, 5/registerWo, 1/1, 2/10, 3/101, 1/1, 1/splat32, 5/lower5, 5/registerZo, 1/1 }


\paragraph{Syntax} \texttt{addusbo} \emph{registerWo} = \emph{registerZo}, \emph{upper27\_\discretionary{}{}{}lower5}\emph{splat32}

\paragraph{Description} The \emph{upper27\_\discretionary{}{}{}lower5} and the \emph{registerZo} are interpreted as eight bytes packed into 64 bits. The \emph{upper27\_\discretionary{}{}{}lower5} bytes and the \emph{registerZo} bytes are added with unsigned saturation to 8 bits. The eight resulting bytes are packed and stored into the \emph{registerWo}.


\paragraph{Modifier(s)}
\noindent\begin{itemize}
\item splat32 (cf. splat32 in section \ref{par:modifier-splat32} on page \pageref{par:modifier-splat32})
\end{itemize}

\paragraph{Execution} {Reservation table \textsf{ALU\_LITE.X} (Section~\ref{sec:reservation-ALU-LITE.X})}
~
{\begin{lstlisting}
stage ID:
  new splat32 = _ZX_1(splat32);
stage RR:
  new argument3 = splat32 ? (_WX_32(upper27_lower5) << 32) | _ZX_32(_WX_32(upper27_lower5)) : _SX_32(_WX_32(upper27_lower5));
  new argument2 = GPR[(registerZo<<1)+1];
  for (I = 0; I < 8; I += 1) {
    result1.8[I] = _SATU_8(_ZX_8(argument3.8[I]) + _ZX_8(argument2.8[I]))
  };
stage E1:
  GPR[(registerWo<<1)+1] = result1;
\end{lstlisting}}
\newpage
\subsubsection[{\small ADDUSD: Add Unsigned Saturated Double Word}]{ADDUSD\hfill Add Unsigned Saturated Double Word}
\label{instr:ADDUSD}\index{addusd}
 
\paragraph{Encoding} Format \textsf{ALU\_DWRR1} (Section~\ref{sec:format-ALU-DWRR1}), \verb|exucode2 = 1010|\\

\bitfields{ 1/P, 2/11, 1/0, 4/1010, 6/registerW, 2/10, 3/101, 1/0, 6/registerY, 6/registerZ }


\paragraph{Syntax} \texttt{addusd} \emph{registerW} = \emph{registerZ}, \emph{registerY}

\paragraph{Description} The \emph{registerY} and the \emph{registerZ} are added with 64-bit unsigned saturation.


\paragraph{Execution} {Reservation table \textsf{ALU\_TINY} (Section~\ref{sec:reservation-ALU-TINY})}
~
{\begin{lstlisting}
stage RR:
  new argument3 = GPR[registerY];
  new argument2 = GPR[registerZ];
  new result1 = _SATU_64(_ZX_64(argument3) + _ZX_64(argument2));
stage E1:
  GPR[registerW] = result1;
\end{lstlisting}}
\newpage

\paragraph{Encoding} Format \textsf{ALU\_DWRR1.M} (Section~\ref{sec:format-ALU-DWRR1.M}), \verb|exucode2 = 1010|\\

\bitfields{ 1/P, 2/00, 2/-- --, 27/upper27, 1/1, 2/11, 1/0, 4/1010, 6/registerW, 2/10, 3/101, 1/0, 1/splat32, 5/lower5, 6/registerZ }


\paragraph{Syntax} \texttt{addusd} \emph{registerW} = \emph{registerZ}, \emph{upper27\_\discretionary{}{}{}lower5}\emph{splat32}

\paragraph{Description} The \emph{upper27\_\discretionary{}{}{}lower5} and the \emph{registerZ} are added with 64-bit unsigned saturation.


\paragraph{Modifier(s)}
\noindent\begin{itemize}
\item splat32 (cf. splat32 in section \ref{par:modifier-splat32} on page \pageref{par:modifier-splat32})
\end{itemize}

\paragraph{Execution} {Reservation table \textsf{ALU\_TINY.X} (Section~\ref{sec:reservation-ALU-TINY.X})}
~
{\begin{lstlisting}
stage ID:
  new splat32 = _ZX_1(splat32);
stage RR:
  new argument3 = splat32 ? (_WX_32(upper27_lower5) << 32) | _ZX_32(_WX_32(upper27_lower5)) : _SX_32(_WX_32(upper27_lower5));
  new argument2 = GPR[registerZ];
  new result1 = _SATU_64(_ZX_64(argument3) + _ZX_64(argument2));
stage E1:
  GPR[registerW] = result1;
\end{lstlisting}}
\newpage
\subsubsection[{\small ADDUSHQ: Add Unsigned Saturated Half Word Quadruple}]{ADDUSHQ\hfill Add Unsigned Saturated Half Word Quadruple}
\label{instr:ADDUSHQ}\index{addushq}
 
\paragraph{Encoding} Format \textsf{ALU\_HQWRR1E} (Section~\ref{sec:format-ALU-HQWRR1E}), \verb|exucode2 = 1010|\\

\bitfields{ 1/P, 2/11, 1/1, 4/1010, 5/registerWe, 1/0, 2/10, 3/101, 1/1, 5/registerYe, 1/--, 5/registerZe, 1/0 }


\paragraph{Syntax} \texttt{addushq} \emph{registerWe} = \emph{registerZe}, \emph{registerYe}

\paragraph{Description} The \emph{registerYe} and the \emph{registerZe} are interpreted as four half words packed into 64 bits. The \emph{registerYe} half words and the \emph{registerZe} half words are added with unsigned saturation to 16 bits. The four resulting half words are packed and stored into the \emph{registerWe}.


\paragraph{Execution} {Reservation table \textsf{ALU\_LITE} (Section~\ref{sec:reservation-ALU-LITE})}
~
{\begin{lstlisting}
stage RR:
  new argument3 = GPR[registerYe<<1];
  new argument2 = GPR[registerZe<<1];
  for (I = 0; I < 4; I += 1) {
    result1.16[I] = _SATU_16(_ZX_16(argument3.16[I]) + _ZX_16(argument2.16[I]))
  };
stage E1:
  GPR[registerWe<<1] = result1;
\end{lstlisting}}
\newpage

\paragraph{Encoding} Format \textsf{ALU\_HQWRR1O} (Section~\ref{sec:format-ALU-HQWRR1O}), \verb|exucode2 = 1010|\\

\bitfields{ 1/P, 2/11, 1/1, 4/1010, 5/registerWo, 1/1, 2/10, 3/101, 1/1, 5/registerYo, 1/--, 5/registerZo, 1/0 }


\paragraph{Syntax} \texttt{addushq} \emph{registerWo} = \emph{registerZo}, \emph{registerYo}

\paragraph{Description} The \emph{registerYo} and the \emph{registerZo} are interpreted as four half words packed into 64 bits. The \emph{registerYo} half words and the \emph{registerZo} half words are added with unsigned saturation to 16 bits. The four resulting half words are packed and stored into the \emph{registerWo}.


\paragraph{Execution} {Reservation table \textsf{ALU\_LITE} (Section~\ref{sec:reservation-ALU-LITE})}
~
{\begin{lstlisting}
stage RR:
  new argument3 = GPR[(registerYo<<1)+1];
  new argument2 = GPR[(registerZo<<1)+1];
  for (I = 0; I < 4; I += 1) {
    result1.16[I] = _SATU_16(_ZX_16(argument3.16[I]) + _ZX_16(argument2.16[I]))
  };
stage E1:
  GPR[(registerWo<<1)+1] = result1;
\end{lstlisting}}
\newpage

\paragraph{Encoding} Format \textsf{ALU\_HQWRR1E.M} (Section~\ref{sec:format-ALU-HQWRR1E.M}), \verb|exucode2 = 1010|\\

\bitfields{ 1/P, 2/00, 2/-- --, 27/upper27, 1/1, 2/11, 1/1, 4/1010, 5/registerWe, 1/0, 2/10, 3/101, 1/1, 1/splat32, 5/lower5, 5/registerZe, 1/0 }


\paragraph{Syntax} \texttt{addushq} \emph{registerWe} = \emph{registerZe}, \emph{upper27\_\discretionary{}{}{}lower5}\emph{splat32}

\paragraph{Description} The \emph{upper27\_\discretionary{}{}{}lower5} and the \emph{registerZe} are interpreted as four half words packed into 64 bits. The \emph{upper27\_\discretionary{}{}{}lower5} half words and the \emph{registerZe} half words are added with unsigned saturation to 16 bits. The four resulting half words are packed and stored into the \emph{registerWe}.


\paragraph{Modifier(s)}
\noindent\begin{itemize}
\item splat32 (cf. splat32 in section \ref{par:modifier-splat32} on page \pageref{par:modifier-splat32})
\end{itemize}

\paragraph{Execution} {Reservation table \textsf{ALU\_LITE.X} (Section~\ref{sec:reservation-ALU-LITE.X})}
~
{\begin{lstlisting}
stage ID:
  new splat32 = _ZX_1(splat32);
stage RR:
  new argument3 = splat32 ? (_WX_32(upper27_lower5) << 32) | _ZX_32(_WX_32(upper27_lower5)) : _SX_32(_WX_32(upper27_lower5));
  new argument2 = GPR[registerZe<<1];
  for (I = 0; I < 4; I += 1) {
    result1.16[I] = _SATU_16(_ZX_16(argument3.16[I]) + _ZX_16(argument2.16[I]))
  };
stage E1:
  GPR[registerWe<<1] = result1;
\end{lstlisting}}
\newpage

\paragraph{Encoding} Format \textsf{ALU\_HQWRR1O.M} (Section~\ref{sec:format-ALU-HQWRR1O.M}), \verb|exucode2 = 1010|\\

\bitfields{ 1/P, 2/00, 2/-- --, 27/upper27, 1/1, 2/11, 1/1, 4/1010, 5/registerWo, 1/1, 2/10, 3/101, 1/1, 1/splat32, 5/lower5, 5/registerZo, 1/0 }


\paragraph{Syntax} \texttt{addushq} \emph{registerWo} = \emph{registerZo}, \emph{upper27\_\discretionary{}{}{}lower5}\emph{splat32}

\paragraph{Description} The \emph{upper27\_\discretionary{}{}{}lower5} and the \emph{registerZo} are interpreted as four half words packed into 64 bits. The \emph{upper27\_\discretionary{}{}{}lower5} half words and the \emph{registerZo} half words are added with unsigned saturation to 16 bits. The four resulting half words are packed and stored into the \emph{registerWo}.


\paragraph{Modifier(s)}
\noindent\begin{itemize}
\item splat32 (cf. splat32 in section \ref{par:modifier-splat32} on page \pageref{par:modifier-splat32})
\end{itemize}

\paragraph{Execution} {Reservation table \textsf{ALU\_LITE.X} (Section~\ref{sec:reservation-ALU-LITE.X})}
~
{\begin{lstlisting}
stage ID:
  new splat32 = _ZX_1(splat32);
stage RR:
  new argument3 = splat32 ? (_WX_32(upper27_lower5) << 32) | _ZX_32(_WX_32(upper27_lower5)) : _SX_32(_WX_32(upper27_lower5));
  new argument2 = GPR[(registerZo<<1)+1];
  for (I = 0; I < 4; I += 1) {
    result1.16[I] = _SATU_16(_ZX_16(argument3.16[I]) + _ZX_16(argument2.16[I]))
  };
stage E1:
  GPR[(registerWo<<1)+1] = result1;
\end{lstlisting}}
\newpage
\subsubsection[{\small ADDUSW: Add Unsigned Saturated Word}]{ADDUSW\hfill Add Unsigned Saturated Word}
\label{instr:ADDUSW}\index{addusw}
 
\paragraph{Encoding} Format \textsf{ALU\_WWRR1} (Section~\ref{sec:format-ALU-WWRR1}), \verb|exucode2 = 1010|\\

\bitfields{ 1/P, 2/11, 1/0, 4/1010, 6/registerW, 2/11, 3/101, 1/signextw, 6/registerY, 6/registerZ }


\paragraph{Syntax} \texttt{addusw}\emph{signextw} \emph{registerW} = \emph{registerZ}, \emph{registerY}

\paragraph{Description} The \emph{registerY} and the \emph{registerZ} are added with 32-bit unsigned saturation. The result with the 32 upper bits cleared is stored into the \emph{registerW}.


\paragraph{Modifier(s)}
\noindent\begin{itemize}
\item signextw (cf. signextw in section \ref{par:modifier-signextw} on page \pageref{par:modifier-signextw})
\end{itemize}

\paragraph{Execution} {Reservation table \textsf{ALU\_TINY} (Section~\ref{sec:reservation-ALU-TINY})}
~
{\begin{lstlisting}
stage ID:
  new signextw = _ZX_1(signextw);
stage RR:
  new argument3 = GPR[registerY];
  new argument2 = GPR[registerZ];
  new result1 = _SATU_32(_ZX_32(argument3) + _ZX_32(argument2));
stage E1:
  GPR[registerW] = signextw ? _SX_32(result1) : _ZX_32(result1);
\end{lstlisting}}
\newpage

\paragraph{Encoding} Format \textsf{ALU\_WWRR1.W} (Section~\ref{sec:format-ALU-WWRR1.W}), \verb|exucode2 = 1010|\\

\bitfields{ 1/P, 2/00, 2/-- --, 27/upper27, 1/1, 2/11, 1/0, 4/1010, 6/registerW, 2/11, 3/101, 1/signextw, 1/0, 5/lower5, 6/registerZ }


\paragraph{Syntax} \texttt{addusw}\emph{signextw} \emph{registerW} = \emph{registerZ}, \emph{upper27\_\discretionary{}{}{}lower5}

\paragraph{Description} The \emph{upper27\_\discretionary{}{}{}lower5} and the \emph{registerZ} are added with 32-bit unsigned saturation. The result with the 32 upper bits cleared is stored into the \emph{registerW}.


\paragraph{Modifier(s)}
\noindent\begin{itemize}
\item signextw (cf. signextw in section \ref{par:modifier-signextw} on page \pageref{par:modifier-signextw})
\end{itemize}

\paragraph{Execution} {Reservation table \textsf{ALU\_TINY.X} (Section~\ref{sec:reservation-ALU-TINY.X})}
~
{\begin{lstlisting}
stage ID:
  new signextw = _ZX_1(signextw);
stage RR:
  new argument3 = _WX_32(upper27_lower5);
  new argument2 = GPR[registerZ];
  new result1 = _SATU_32(_ZX_32(argument3) + _ZX_32(argument2));
stage E1:
  GPR[registerW] = signextw ? _SX_32(result1) : _ZX_32(result1);
\end{lstlisting}}
\newpage
\subsubsection[{\small ADDUSWP: Add Unsigned Saturated Word Pair}]{ADDUSWP\hfill Add Unsigned Saturated Word Pair}
\label{instr:ADDUSWP}\index{adduswp}
 
\paragraph{Encoding} Format \textsf{ALU\_WPWRR1E} (Section~\ref{sec:format-ALU-WPWRR1E}), \verb|exucode2 = 1010|\\

\bitfields{ 1/P, 2/11, 1/1, 4/1010, 5/registerWe, 1/0, 2/10, 3/101, 1/0, 5/registerYe, 1/--, 5/registerZe, 1/1 }


\paragraph{Syntax} \texttt{adduswp} \emph{registerWe} = \emph{registerZe}, \emph{registerYe}

\paragraph{Description} The \emph{registerYe} and the \emph{registerZe} are interpreted as two words packed into 64 bits. The \emph{registerYe} words and the \emph{registerZe} words are added with unsigned saturation to 32 bits. The two resulting words are packed and stored into the \emph{registerWe}.


\paragraph{Execution} {Reservation table \textsf{ALU\_LITE} (Section~\ref{sec:reservation-ALU-LITE})}
~
{\begin{lstlisting}
stage RR:
  new argument3 = GPR[registerYe<<1];
  new argument2 = GPR[registerZe<<1];
  for (I = 0; I < 2; I += 1) {
    result1.32[I] = _SATU_32(_ZX_32(argument3.32[I]) + _ZX_32(argument2.32[I]))
  };
stage E1:
  GPR[registerWe<<1] = result1;
\end{lstlisting}}
\newpage

\paragraph{Encoding} Format \textsf{ALU\_WPWRR1O} (Section~\ref{sec:format-ALU-WPWRR1O}), \verb|exucode2 = 1010|\\

\bitfields{ 1/P, 2/11, 1/1, 4/1010, 5/registerWo, 1/1, 2/10, 3/101, 1/0, 5/registerYo, 1/--, 5/registerZo, 1/1 }


\paragraph{Syntax} \texttt{adduswp} \emph{registerWo} = \emph{registerZo}, \emph{registerYo}

\paragraph{Description} The \emph{registerYo} and the \emph{registerZo} are interpreted as two words packed into 64 bits. The \emph{registerYo} words and the \emph{registerZo} words are added with unsigned saturation to 32 bits. The two resulting words are packed and stored into the \emph{registerWo}.


\paragraph{Execution} {Reservation table \textsf{ALU\_LITE} (Section~\ref{sec:reservation-ALU-LITE})}
~
{\begin{lstlisting}
stage RR:
  new argument3 = GPR[(registerYo<<1)+1];
  new argument2 = GPR[(registerZo<<1)+1];
  for (I = 0; I < 2; I += 1) {
    result1.32[I] = _SATU_32(_ZX_32(argument3.32[I]) + _ZX_32(argument2.32[I]))
  };
stage E1:
  GPR[(registerWo<<1)+1] = result1;
\end{lstlisting}}
\newpage

\paragraph{Encoding} Format \textsf{ALU\_WPWRR1E.M} (Section~\ref{sec:format-ALU-WPWRR1E.M}), \verb|exucode2 = 1010|\\

\bitfields{ 1/P, 2/00, 2/-- --, 27/upper27, 1/1, 2/11, 1/1, 4/1010, 5/registerWe, 1/0, 2/10, 3/101, 1/0, 1/splat32, 5/lower5, 5/registerZe, 1/1 }


\paragraph{Syntax} \texttt{adduswp} \emph{registerWe} = \emph{registerZe}, \emph{upper27\_\discretionary{}{}{}lower5}\emph{splat32}

\paragraph{Description} The \emph{upper27\_\discretionary{}{}{}lower5} and the \emph{registerZe} are interpreted as two words packed into 64 bits. The \emph{upper27\_\discretionary{}{}{}lower5} words and the \emph{registerZe} words are added with unsigned saturation to 32 bits. The two resulting words are packed and stored into the \emph{registerWe}.


\paragraph{Modifier(s)}
\noindent\begin{itemize}
\item splat32 (cf. splat32 in section \ref{par:modifier-splat32} on page \pageref{par:modifier-splat32})
\end{itemize}

\paragraph{Execution} {Reservation table \textsf{ALU\_LITE.X} (Section~\ref{sec:reservation-ALU-LITE.X})}
~
{\begin{lstlisting}
stage ID:
  new splat32 = _ZX_1(splat32);
stage RR:
  new argument3 = splat32 ? (_WX_32(upper27_lower5) << 32) | _ZX_32(_WX_32(upper27_lower5)) : _SX_32(_WX_32(upper27_lower5));
  new argument2 = GPR[registerZe<<1];
  for (I = 0; I < 2; I += 1) {
    result1.32[I] = _SATU_32(_ZX_32(argument3.32[I]) + _ZX_32(argument2.32[I]))
  };
stage E1:
  GPR[registerWe<<1] = result1;
\end{lstlisting}}
\newpage

\paragraph{Encoding} Format \textsf{ALU\_WPWRR1O.M} (Section~\ref{sec:format-ALU-WPWRR1O.M}), \verb|exucode2 = 1010|\\

\bitfields{ 1/P, 2/00, 2/-- --, 27/upper27, 1/1, 2/11, 1/1, 4/1010, 5/registerWo, 1/1, 2/10, 3/101, 1/0, 1/splat32, 5/lower5, 5/registerZo, 1/1 }


\paragraph{Syntax} \texttt{adduswp} \emph{registerWo} = \emph{registerZo}, \emph{upper27\_\discretionary{}{}{}lower5}\emph{splat32}

\paragraph{Description} The \emph{upper27\_\discretionary{}{}{}lower5} and the \emph{registerZo} are interpreted as two words packed into 64 bits. The \emph{upper27\_\discretionary{}{}{}lower5} words and the \emph{registerZo} words are added with unsigned saturation to 32 bits. The two resulting words are packed and stored into the \emph{registerWo}.


\paragraph{Modifier(s)}
\noindent\begin{itemize}
\item splat32 (cf. splat32 in section \ref{par:modifier-splat32} on page \pageref{par:modifier-splat32})
\end{itemize}

\paragraph{Execution} {Reservation table \textsf{ALU\_LITE.X} (Section~\ref{sec:reservation-ALU-LITE.X})}
~
{\begin{lstlisting}
stage ID:
  new splat32 = _ZX_1(splat32);
stage RR:
  new argument3 = splat32 ? (_WX_32(upper27_lower5) << 32) | _ZX_32(_WX_32(upper27_lower5)) : _SX_32(_WX_32(upper27_lower5));
  new argument2 = GPR[(registerZo<<1)+1];
  for (I = 0; I < 2; I += 1) {
    result1.32[I] = _SATU_32(_ZX_32(argument3.32[I]) + _ZX_32(argument2.32[I]))
  };
stage E1:
  GPR[(registerWo<<1)+1] = result1;
\end{lstlisting}}
\newpage
\subsubsection[{\small ADDW: Add Word to Word}]{ADDW\hfill Add Word to Word}
\label{instr:ADDW}\index{addw}
 
\paragraph{Encoding} Format \textsf{ALU\_WWRR0} (Section~\ref{sec:format-ALU-WWRR0}), \verb|exucode2 = 0010|\\

\bitfields{ 1/P, 2/11, 1/0, 4/0010, 6/registerW, 2/11, 3/000, 1/signextw, 6/registerY, 6/registerZ }


\paragraph{Syntax} \texttt{addw}\emph{signextw} \emph{registerW} = \emph{registerZ}, \emph{registerY}

\paragraph{Description} The \emph{registerY} and the \emph{registerZ} are added. The result with the 32 upper bits cleared is stored into the \emph{registerW}.


\paragraph{Modifier(s)}
\noindent\begin{itemize}
\item signextw (cf. signextw in section \ref{par:modifier-signextw} on page \pageref{par:modifier-signextw})
\end{itemize}

\paragraph{Execution} {Reservation table \textsf{ALU\_TINY} (Section~\ref{sec:reservation-ALU-TINY})}
~
{\begin{lstlisting}
stage ID:
  new signextw = _ZX_1(signextw);
stage RR:
  new argument3 = GPR[registerY];
  new argument2 = GPR[registerZ];
  new result1 = argument3 + argument2;
stage E1:
  GPR[registerW] = signextw ? _SX_32(result1) : _ZX_32(result1);
\end{lstlisting}}
\newpage

\paragraph{Encoding} Format \textsf{ALU\_WWRR0.W} (Section~\ref{sec:format-ALU-WWRR0.W}), \verb|exucode2 = 0010|\\

\bitfields{ 1/P, 2/00, 2/-- --, 27/upper27, 1/1, 2/11, 1/0, 4/0010, 6/registerW, 2/11, 3/000, 1/signextw, 1/0, 5/lower5, 6/registerZ }


\paragraph{Syntax} \texttt{addw}\emph{signextw} \emph{registerW} = \emph{registerZ}, \emph{upper27\_\discretionary{}{}{}lower5}

\paragraph{Description} The \emph{upper27\_\discretionary{}{}{}lower5} and the \emph{registerZ} are added. The result with the 32 upper bits cleared is stored into the \emph{registerW}.


\paragraph{Modifier(s)}
\noindent\begin{itemize}
\item signextw (cf. signextw in section \ref{par:modifier-signextw} on page \pageref{par:modifier-signextw})
\end{itemize}

\paragraph{Execution} {Reservation table \textsf{ALU\_TINY.X} (Section~\ref{sec:reservation-ALU-TINY.X})}
~
{\begin{lstlisting}
stage ID:
  new signextw = _ZX_1(signextw);
stage RR:
  new argument3 = _WX_32(upper27_lower5);
  new argument2 = GPR[registerZ];
  new result1 = argument3 + argument2;
stage E1:
  GPR[registerW] = signextw ? _SX_32(result1) : _ZX_32(result1);
\end{lstlisting}}
\newpage
\subsubsection[{\small ADDWP: Add Word Pair}]{ADDWP\hfill Add Word Pair}
\label{instr:ADDWP}\index{addwp}
 
\paragraph{Encoding} Format \textsf{ALU\_WPWRR0E} (Section~\ref{sec:format-ALU-WPWRR0E}), \verb|exucode2 = 0010|\\

\bitfields{ 1/P, 2/11, 1/1, 4/0010, 5/registerWe, 1/0, 2/10, 3/000, 1/0, 5/registerYe, 1/--, 5/registerZe, 1/1 }


\paragraph{Syntax} \texttt{addwp} \emph{registerWe} = \emph{registerZe}, \emph{registerYe}

\paragraph{Description} The \emph{registerYe} and the \emph{registerZe} are interpreted as two words packed into 64 bits. The \emph{registerYe} words and the \emph{registerZe} words are added. The two resulting words are packed and stored into the \emph{registerWe}.


\paragraph{Execution} {Reservation table \textsf{ALU\_LITE} (Section~\ref{sec:reservation-ALU-LITE})}
~
{\begin{lstlisting}
stage RR:
  new argument3 = GPR[registerYe<<1];
  new argument2 = GPR[registerZe<<1];
  for (I = 0; I < 2; I += 1) {
    result1.32[I] = argument3.32[I] + argument2.32[I]
  };
stage E1:
  GPR[registerWe<<1] = result1;
\end{lstlisting}}
\newpage

\paragraph{Encoding} Format \textsf{ALU\_WPWRR0O} (Section~\ref{sec:format-ALU-WPWRR0O}), \verb|exucode2 = 0010|\\

\bitfields{ 1/P, 2/11, 1/1, 4/0010, 5/registerWo, 1/1, 2/10, 3/000, 1/0, 5/registerYo, 1/--, 5/registerZo, 1/1 }


\paragraph{Syntax} \texttt{addwp} \emph{registerWo} = \emph{registerZo}, \emph{registerYo}

\paragraph{Description} The \emph{registerYo} and the \emph{registerZo} are interpreted as two words packed into 64 bits. The \emph{registerYo} words and the \emph{registerZo} words are added. The two resulting words are packed and stored into the \emph{registerWo}.


\paragraph{Execution} {Reservation table \textsf{ALU\_LITE} (Section~\ref{sec:reservation-ALU-LITE})}
~
{\begin{lstlisting}
stage RR:
  new argument3 = GPR[(registerYo<<1)+1];
  new argument2 = GPR[(registerZo<<1)+1];
  for (I = 0; I < 2; I += 1) {
    result1.32[I] = argument3.32[I] + argument2.32[I]
  };
stage E1:
  GPR[(registerWo<<1)+1] = result1;
\end{lstlisting}}
\newpage

\paragraph{Encoding} Format \textsf{ALU\_WPWRR0E.M} (Section~\ref{sec:format-ALU-WPWRR0E.M}), \verb|exucode2 = 0010|\\

\bitfields{ 1/P, 2/00, 2/-- --, 27/upper27, 1/1, 2/11, 1/1, 4/0010, 5/registerWe, 1/0, 2/10, 3/000, 1/0, 1/splat32, 5/lower5, 5/registerZe, 1/1 }


\paragraph{Syntax} \texttt{addwp} \emph{registerWe} = \emph{registerZe}, \emph{upper27\_\discretionary{}{}{}lower5}\emph{splat32}

\paragraph{Description} The \emph{upper27\_\discretionary{}{}{}lower5} and the \emph{registerZe} are interpreted as two words packed into 64 bits. The \emph{upper27\_\discretionary{}{}{}lower5} words and the \emph{registerZe} words are added. The two resulting words are packed and stored into the \emph{registerWe}.


\paragraph{Modifier(s)}
\noindent\begin{itemize}
\item splat32 (cf. splat32 in section \ref{par:modifier-splat32} on page \pageref{par:modifier-splat32})
\end{itemize}

\paragraph{Execution} {Reservation table \textsf{ALU\_LITE.X} (Section~\ref{sec:reservation-ALU-LITE.X})}
~
{\begin{lstlisting}
stage ID:
  new splat32 = _ZX_1(splat32);
stage RR:
  new argument3 = splat32 ? (_WX_32(upper27_lower5) << 32) | _ZX_32(_WX_32(upper27_lower5)) : _SX_32(_WX_32(upper27_lower5));
  new argument2 = GPR[registerZe<<1];
  for (I = 0; I < 2; I += 1) {
    result1.32[I] = argument3.32[I] + argument2.32[I]
  };
stage E1:
  GPR[registerWe<<1] = result1;
\end{lstlisting}}
\newpage

\paragraph{Encoding} Format \textsf{ALU\_WPWRR0O.M} (Section~\ref{sec:format-ALU-WPWRR0O.M}), \verb|exucode2 = 0010|\\

\bitfields{ 1/P, 2/00, 2/-- --, 27/upper27, 1/1, 2/11, 1/1, 4/0010, 5/registerWo, 1/1, 2/10, 3/000, 1/0, 1/splat32, 5/lower5, 5/registerZo, 1/1 }


\paragraph{Syntax} \texttt{addwp} \emph{registerWo} = \emph{registerZo}, \emph{upper27\_\discretionary{}{}{}lower5}\emph{splat32}

\paragraph{Description} The \emph{upper27\_\discretionary{}{}{}lower5} and the \emph{registerZo} are interpreted as two words packed into 64 bits. The \emph{upper27\_\discretionary{}{}{}lower5} words and the \emph{registerZo} words are added. The two resulting words are packed and stored into the \emph{registerWo}.


\paragraph{Modifier(s)}
\noindent\begin{itemize}
\item splat32 (cf. splat32 in section \ref{par:modifier-splat32} on page \pageref{par:modifier-splat32})
\end{itemize}

\paragraph{Execution} {Reservation table \textsf{ALU\_LITE.X} (Section~\ref{sec:reservation-ALU-LITE.X})}
~
{\begin{lstlisting}
stage ID:
  new splat32 = _ZX_1(splat32);
stage RR:
  new argument3 = splat32 ? (_WX_32(upper27_lower5) << 32) | _ZX_32(_WX_32(upper27_lower5)) : _SX_32(_WX_32(upper27_lower5));
  new argument2 = GPR[(registerZo<<1)+1];
  for (I = 0; I < 2; I += 1) {
    result1.32[I] = argument3.32[I] + argument2.32[I]
  };
stage E1:
  GPR[(registerWo<<1)+1] = result1;
\end{lstlisting}}
\newpage
\subsubsection[{\small ADDX16BO: Add to Times 16 Byte Octuple}]{ADDX16BO\hfill Add to Times 16 Byte Octuple}
\label{instr:ADDX16BO}\index{addx16bo}
 
\paragraph{Encoding} Format \textsf{ALU\_BOWRR1E} (Section~\ref{sec:format-ALU-BOWRR1E}), \verb|exucode2 = 0011|\\

\bitfields{ 1/P, 2/11, 1/1, 4/0011, 5/registerWe, 1/0, 2/10, 3/101, 1/1, 5/registerYe, 1/--, 5/registerZe, 1/1 }


\paragraph{Syntax} \texttt{addx16bo} \emph{registerWe} = \emph{registerZe}, \emph{registerYe}

\paragraph{Description} The \emph{registerYe} and the \emph{registerZe} are interpreted as eight bytes packed into 64 bits. The \emph{registerZe} bytes are shifted left by four bits and added to the \emph{registerYe} bytes. The eight resulting bytes are packed and stored into the \emph{registerWe}.


\paragraph{Execution} {Reservation table \textsf{ALU\_LITE} (Section~\ref{sec:reservation-ALU-LITE})}
~
{\begin{lstlisting}
stage RR:
  new argument3 = GPR[registerYe<<1];
  new argument2 = GPR[registerZe<<1];
  for (I = 0; I < 8; I += 1) {
    result1.8[I] = _SX_8(argument3.8[I]) + (_SX_8(argument2.8[I]) << 4)
  };
stage E1:
  GPR[registerWe<<1] = result1;
\end{lstlisting}}
\newpage

\paragraph{Encoding} Format \textsf{ALU\_BOWRR1O} (Section~\ref{sec:format-ALU-BOWRR1O}), \verb|exucode2 = 0011|\\

\bitfields{ 1/P, 2/11, 1/1, 4/0011, 5/registerWo, 1/1, 2/10, 3/101, 1/1, 5/registerYo, 1/--, 5/registerZo, 1/1 }


\paragraph{Syntax} \texttt{addx16bo} \emph{registerWo} = \emph{registerZo}, \emph{registerYo}

\paragraph{Description} The \emph{registerYo} and the \emph{registerZo} are interpreted as eight bytes packed into 64 bits. The \emph{registerZo} bytes are shifted left by four bits and added to the \emph{registerYo} bytes. The eight resulting bytes are packed and stored into the \emph{registerWo}.


\paragraph{Execution} {Reservation table \textsf{ALU\_LITE} (Section~\ref{sec:reservation-ALU-LITE})}
~
{\begin{lstlisting}
stage RR:
  new argument3 = GPR[(registerYo<<1)+1];
  new argument2 = GPR[(registerZo<<1)+1];
  for (I = 0; I < 8; I += 1) {
    result1.8[I] = _SX_8(argument3.8[I]) + (_SX_8(argument2.8[I]) << 4)
  };
stage E1:
  GPR[(registerWo<<1)+1] = result1;
\end{lstlisting}}
\newpage

\paragraph{Encoding} Format \textsf{ALU\_BOWRR1E.M} (Section~\ref{sec:format-ALU-BOWRR1E.M}), \verb|exucode2 = 0011|\\

\bitfields{ 1/P, 2/00, 2/-- --, 27/upper27, 1/1, 2/11, 1/1, 4/0011, 5/registerWe, 1/0, 2/10, 3/101, 1/1, 1/splat32, 5/lower5, 5/registerZe, 1/1 }


\paragraph{Syntax} \texttt{addx16bo} \emph{registerWe} = \emph{registerZe}, \emph{upper27\_\discretionary{}{}{}lower5}\emph{splat32}

\paragraph{Description} The \emph{upper27\_\discretionary{}{}{}lower5} and the \emph{registerZe} are interpreted as eight bytes packed into 64 bits. The \emph{registerZe} bytes are shifted left by four bits and added to the \emph{upper27\_\discretionary{}{}{}lower5} bytes. The eight resulting bytes are packed and stored into the \emph{registerWe}.


\paragraph{Modifier(s)}
\noindent\begin{itemize}
\item splat32 (cf. splat32 in section \ref{par:modifier-splat32} on page \pageref{par:modifier-splat32})
\end{itemize}

\paragraph{Execution} {Reservation table \textsf{ALU\_LITE.X} (Section~\ref{sec:reservation-ALU-LITE.X})}
~
{\begin{lstlisting}
stage ID:
  new splat32 = _ZX_1(splat32);
stage RR:
  new argument3 = splat32 ? (_WX_32(upper27_lower5) << 32) | _ZX_32(_WX_32(upper27_lower5)) : _SX_32(_WX_32(upper27_lower5));
  new argument2 = GPR[registerZe<<1];
  for (I = 0; I < 8; I += 1) {
    result1.8[I] = _SX_8(argument3.8[I]) + (_SX_8(argument2.8[I]) << 4)
  };
stage E1:
  GPR[registerWe<<1] = result1;
\end{lstlisting}}
\newpage

\paragraph{Encoding} Format \textsf{ALU\_BOWRR1O.M} (Section~\ref{sec:format-ALU-BOWRR1O.M}), \verb|exucode2 = 0011|\\

\bitfields{ 1/P, 2/00, 2/-- --, 27/upper27, 1/1, 2/11, 1/1, 4/0011, 5/registerWo, 1/1, 2/10, 3/101, 1/1, 1/splat32, 5/lower5, 5/registerZo, 1/1 }


\paragraph{Syntax} \texttt{addx16bo} \emph{registerWo} = \emph{registerZo}, \emph{upper27\_\discretionary{}{}{}lower5}\emph{splat32}

\paragraph{Description} The \emph{upper27\_\discretionary{}{}{}lower5} and the \emph{registerZo} are interpreted as eight bytes packed into 64 bits. The \emph{registerZo} bytes are shifted left by four bits and added to the \emph{upper27\_\discretionary{}{}{}lower5} bytes. The eight resulting bytes are packed and stored into the \emph{registerWo}.


\paragraph{Modifier(s)}
\noindent\begin{itemize}
\item splat32 (cf. splat32 in section \ref{par:modifier-splat32} on page \pageref{par:modifier-splat32})
\end{itemize}

\paragraph{Execution} {Reservation table \textsf{ALU\_LITE.X} (Section~\ref{sec:reservation-ALU-LITE.X})}
~
{\begin{lstlisting}
stage ID:
  new splat32 = _ZX_1(splat32);
stage RR:
  new argument3 = splat32 ? (_WX_32(upper27_lower5) << 32) | _ZX_32(_WX_32(upper27_lower5)) : _SX_32(_WX_32(upper27_lower5));
  new argument2 = GPR[(registerZo<<1)+1];
  for (I = 0; I < 8; I += 1) {
    result1.8[I] = _SX_8(argument3.8[I]) + (_SX_8(argument2.8[I]) << 4)
  };
stage E1:
  GPR[(registerWo<<1)+1] = result1;
\end{lstlisting}}
\newpage
\subsubsection[{\small ADDX16D: Add Double Word to Double Word Times 16}]{ADDX16D\hfill Add Double Word to Double Word Times 16}
\label{instr:ADDX16D}\index{addx16d}
 
\paragraph{Encoding} Format \textsf{ALU\_DWRR1} (Section~\ref{sec:format-ALU-DWRR1}), \verb|exucode2 = 0011|\\

\bitfields{ 1/P, 2/11, 1/0, 4/0011, 6/registerW, 2/10, 3/101, 1/0, 6/registerY, 6/registerZ }


\paragraph{Syntax} \texttt{addx16d} \emph{registerW} = \emph{registerZ}, \emph{registerY}

\paragraph{Description} The \emph{registerZ} is shifted left by four bits and added to the \emph{registerY}. The result is stored into the \emph{registerW}.


\paragraph{Execution} {Reservation table \textsf{ALU\_TINY} (Section~\ref{sec:reservation-ALU-TINY})}
~
{\begin{lstlisting}
stage RR:
  new argument3 = GPR[registerY];
  new argument2 = GPR[registerZ];
  new result1 = argument3 + (argument2 << 4);
stage E1:
  GPR[registerW] = result1;
\end{lstlisting}}
\newpage

\paragraph{Encoding} Format \textsf{ALU\_DWRR1.M} (Section~\ref{sec:format-ALU-DWRR1.M}), \verb|exucode2 = 0011|\\

\bitfields{ 1/P, 2/00, 2/-- --, 27/upper27, 1/1, 2/11, 1/0, 4/0011, 6/registerW, 2/10, 3/101, 1/0, 1/splat32, 5/lower5, 6/registerZ }


\paragraph{Syntax} \texttt{addx16d} \emph{registerW} = \emph{registerZ}, \emph{upper27\_\discretionary{}{}{}lower5}\emph{splat32}

\paragraph{Description} The \emph{registerZ} is shifted left by four bits and added to the \emph{upper27\_\discretionary{}{}{}lower5}. The result is stored into the \emph{registerW}.


\paragraph{Modifier(s)}
\noindent\begin{itemize}
\item splat32 (cf. splat32 in section \ref{par:modifier-splat32} on page \pageref{par:modifier-splat32})
\end{itemize}

\paragraph{Execution} {Reservation table \textsf{ALU\_TINY.X} (Section~\ref{sec:reservation-ALU-TINY.X})}
~
{\begin{lstlisting}
stage ID:
  new splat32 = _ZX_1(splat32);
stage RR:
  new argument3 = splat32 ? (_WX_32(upper27_lower5) << 32) | _ZX_32(_WX_32(upper27_lower5)) : _SX_32(_WX_32(upper27_lower5));
  new argument2 = GPR[registerZ];
  new result1 = argument3 + (argument2 << 4);
stage E1:
  GPR[registerW] = result1;
\end{lstlisting}}
\newpage
\subsubsection[{\small ADDX16HQ: Add to Times 16 Half Word Quadruple}]{ADDX16HQ\hfill Add to Times 16 Half Word Quadruple}
\label{instr:ADDX16HQ}\index{addx16hq}
 
\paragraph{Encoding} Format \textsf{ALU\_HQWRR1E} (Section~\ref{sec:format-ALU-HQWRR1E}), \verb|exucode2 = 0011|\\

\bitfields{ 1/P, 2/11, 1/1, 4/0011, 5/registerWe, 1/0, 2/10, 3/101, 1/1, 5/registerYe, 1/--, 5/registerZe, 1/0 }


\paragraph{Syntax} \texttt{addx16hq} \emph{registerWe} = \emph{registerZe}, \emph{registerYe}

\paragraph{Description} The \emph{registerYe} and the \emph{registerZe} are interpreted as four half words packed into 64 bits. The \emph{registerZe} half words are shifted left by eight bits and added to the \emph{registerYe} half words. The four resulting half words are packed and stored into the \emph{registerWe}.


\paragraph{Execution} {Reservation table \textsf{ALU\_LITE} (Section~\ref{sec:reservation-ALU-LITE})}
~
{\begin{lstlisting}
stage RR:
  new argument3 = GPR[registerYe<<1];
  new argument2 = GPR[registerZe<<1];
  for (I = 0; I < 4; I += 1) {
    result1.16[I] = _SX_16(argument3.16[I]) + (_SX_16(argument2.16[I]) << 4)
  };
stage E1:
  GPR[registerWe<<1] = result1;
\end{lstlisting}}
\newpage

\paragraph{Encoding} Format \textsf{ALU\_HQWRR1O} (Section~\ref{sec:format-ALU-HQWRR1O}), \verb|exucode2 = 0011|\\

\bitfields{ 1/P, 2/11, 1/1, 4/0011, 5/registerWo, 1/1, 2/10, 3/101, 1/1, 5/registerYo, 1/--, 5/registerZo, 1/0 }


\paragraph{Syntax} \texttt{addx16hq} \emph{registerWo} = \emph{registerZo}, \emph{registerYo}

\paragraph{Description} The \emph{registerYo} and the \emph{registerZo} are interpreted as four half words packed into 64 bits. The \emph{registerZo} half words are shifted left by eight bits and added to the \emph{registerYo} half words. The four resulting half words are packed and stored into the \emph{registerWo}.


\paragraph{Execution} {Reservation table \textsf{ALU\_LITE} (Section~\ref{sec:reservation-ALU-LITE})}
~
{\begin{lstlisting}
stage RR:
  new argument3 = GPR[(registerYo<<1)+1];
  new argument2 = GPR[(registerZo<<1)+1];
  for (I = 0; I < 4; I += 1) {
    result1.16[I] = _SX_16(argument3.16[I]) + (_SX_16(argument2.16[I]) << 4)
  };
stage E1:
  GPR[(registerWo<<1)+1] = result1;
\end{lstlisting}}
\newpage

\paragraph{Encoding} Format \textsf{ALU\_HQWRR1E.M} (Section~\ref{sec:format-ALU-HQWRR1E.M}), \verb|exucode2 = 0011|\\

\bitfields{ 1/P, 2/00, 2/-- --, 27/upper27, 1/1, 2/11, 1/1, 4/0011, 5/registerWe, 1/0, 2/10, 3/101, 1/1, 1/splat32, 5/lower5, 5/registerZe, 1/0 }


\paragraph{Syntax} \texttt{addx16hq} \emph{registerWe} = \emph{registerZe}, \emph{upper27\_\discretionary{}{}{}lower5}\emph{splat32}

\paragraph{Description} The \emph{upper27\_\discretionary{}{}{}lower5} and the \emph{registerZe} are interpreted as four half words packed into 64 bits. The \emph{registerZe} half words are shifted left by eight bits and added to the \emph{upper27\_\discretionary{}{}{}lower5} half words. The four resulting half words are packed and stored into the \emph{registerWe}.


\paragraph{Modifier(s)}
\noindent\begin{itemize}
\item splat32 (cf. splat32 in section \ref{par:modifier-splat32} on page \pageref{par:modifier-splat32})
\end{itemize}

\paragraph{Execution} {Reservation table \textsf{ALU\_LITE.X} (Section~\ref{sec:reservation-ALU-LITE.X})}
~
{\begin{lstlisting}
stage ID:
  new splat32 = _ZX_1(splat32);
stage RR:
  new argument3 = splat32 ? (_WX_32(upper27_lower5) << 32) | _ZX_32(_WX_32(upper27_lower5)) : _SX_32(_WX_32(upper27_lower5));
  new argument2 = GPR[registerZe<<1];
  for (I = 0; I < 4; I += 1) {
    result1.16[I] = _SX_16(argument3.16[I]) + (_SX_16(argument2.16[I]) << 4)
  };
stage E1:
  GPR[registerWe<<1] = result1;
\end{lstlisting}}
\newpage

\paragraph{Encoding} Format \textsf{ALU\_HQWRR1O.M} (Section~\ref{sec:format-ALU-HQWRR1O.M}), \verb|exucode2 = 0011|\\

\bitfields{ 1/P, 2/00, 2/-- --, 27/upper27, 1/1, 2/11, 1/1, 4/0011, 5/registerWo, 1/1, 2/10, 3/101, 1/1, 1/splat32, 5/lower5, 5/registerZo, 1/0 }


\paragraph{Syntax} \texttt{addx16hq} \emph{registerWo} = \emph{registerZo}, \emph{upper27\_\discretionary{}{}{}lower5}\emph{splat32}

\paragraph{Description} The \emph{upper27\_\discretionary{}{}{}lower5} and the \emph{registerZo} are interpreted as four half words packed into 64 bits. The \emph{registerZo} half words are shifted left by eight bits and added to the \emph{upper27\_\discretionary{}{}{}lower5} half words. The four resulting half words are packed and stored into the \emph{registerWo}.


\paragraph{Modifier(s)}
\noindent\begin{itemize}
\item splat32 (cf. splat32 in section \ref{par:modifier-splat32} on page \pageref{par:modifier-splat32})
\end{itemize}

\paragraph{Execution} {Reservation table \textsf{ALU\_LITE.X} (Section~\ref{sec:reservation-ALU-LITE.X})}
~
{\begin{lstlisting}
stage ID:
  new splat32 = _ZX_1(splat32);
stage RR:
  new argument3 = splat32 ? (_WX_32(upper27_lower5) << 32) | _ZX_32(_WX_32(upper27_lower5)) : _SX_32(_WX_32(upper27_lower5));
  new argument2 = GPR[(registerZo<<1)+1];
  for (I = 0; I < 4; I += 1) {
    result1.16[I] = _SX_16(argument3.16[I]) + (_SX_16(argument2.16[I]) << 4)
  };
stage E1:
  GPR[(registerWo<<1)+1] = result1;
\end{lstlisting}}
\newpage
\subsubsection[{\small ADDX16W: Add Word to Word Times 16}]{ADDX16W\hfill Add Word to Word Times 16}
\label{instr:ADDX16W}\index{addx16w}
 
\paragraph{Encoding} Format \textsf{ALU\_WWRR1} (Section~\ref{sec:format-ALU-WWRR1}), \verb|exucode2 = 0011|\\

\bitfields{ 1/P, 2/11, 1/0, 4/0011, 6/registerW, 2/11, 3/101, 1/signextw, 6/registerY, 6/registerZ }


\paragraph{Syntax} \texttt{addx16w}\emph{signextw} \emph{registerW} = \emph{registerZ}, \emph{registerY}

\paragraph{Description} The \emph{registerZ} is shifted left by four bits and added to the \emph{registerY}. The result with the 32 upper bits cleared is stored into the \emph{registerW}.


\paragraph{Modifier(s)}
\noindent\begin{itemize}
\item signextw (cf. signextw in section \ref{par:modifier-signextw} on page \pageref{par:modifier-signextw})
\end{itemize}

\paragraph{Execution} {Reservation table \textsf{ALU\_TINY} (Section~\ref{sec:reservation-ALU-TINY})}
~
{\begin{lstlisting}
stage ID:
  new signextw = _ZX_1(signextw);
stage RR:
  new argument3 = GPR[registerY];
  new argument2 = GPR[registerZ];
  new result1 = argument3 + (argument2 << 4);
stage E1:
  GPR[registerW] = signextw ? _SX_32(result1) : _ZX_32(result1);
\end{lstlisting}}
\newpage

\paragraph{Encoding} Format \textsf{ALU\_WWRR1.W} (Section~\ref{sec:format-ALU-WWRR1.W}), \verb|exucode2 = 0011|\\

\bitfields{ 1/P, 2/00, 2/-- --, 27/upper27, 1/1, 2/11, 1/0, 4/0011, 6/registerW, 2/11, 3/101, 1/signextw, 1/0, 5/lower5, 6/registerZ }


\paragraph{Syntax} \texttt{addx16w}\emph{signextw} \emph{registerW} = \emph{registerZ}, \emph{upper27\_\discretionary{}{}{}lower5}

\paragraph{Description} The \emph{registerZ} is shifted left by four bits and added to the \emph{upper27\_\discretionary{}{}{}lower5}. The result with the 32 upper bits cleared is stored into the \emph{registerW}.


\paragraph{Modifier(s)}
\noindent\begin{itemize}
\item signextw (cf. signextw in section \ref{par:modifier-signextw} on page \pageref{par:modifier-signextw})
\end{itemize}

\paragraph{Execution} {Reservation table \textsf{ALU\_TINY.X} (Section~\ref{sec:reservation-ALU-TINY.X})}
~
{\begin{lstlisting}
stage ID:
  new signextw = _ZX_1(signextw);
stage RR:
  new argument3 = _WX_32(upper27_lower5);
  new argument2 = GPR[registerZ];
  new result1 = argument3 + (argument2 << 4);
stage E1:
  GPR[registerW] = signextw ? _SX_32(result1) : _ZX_32(result1);
\end{lstlisting}}
\newpage
\subsubsection[{\small ADDX16WP: Add to Times 16 Word Pair}]{ADDX16WP\hfill Add to Times 16 Word Pair}
\label{instr:ADDX16WP}\index{addx16wp}
 
\paragraph{Encoding} Format \textsf{ALU\_WPWRR1E} (Section~\ref{sec:format-ALU-WPWRR1E}), \verb|exucode2 = 0011|\\

\bitfields{ 1/P, 2/11, 1/1, 4/0011, 5/registerWe, 1/0, 2/10, 3/101, 1/0, 5/registerYe, 1/--, 5/registerZe, 1/1 }


\paragraph{Syntax} \texttt{addx16wp} \emph{registerWe} = \emph{registerZe}, \emph{registerYe}

\paragraph{Description} The \emph{registerYe} and the \emph{registerZe} are interpreted as two words packed into 64 bits. The \emph{registerZe} words are shifted left by four bits and added to the \emph{registerYe} words. The two resulting words are packed and stored into the \emph{registerWe}.


\paragraph{Execution} {Reservation table \textsf{ALU\_LITE} (Section~\ref{sec:reservation-ALU-LITE})}
~
{\begin{lstlisting}
stage RR:
  new argument3 = GPR[registerYe<<1];
  new argument2 = GPR[registerZe<<1];
  for (I = 0; I < 2; I += 1) {
    result1.32[I] = _SX_32(argument3.32[I]) + (_SX_32(argument2.32[I]) << 4)
  };
stage E1:
  GPR[registerWe<<1] = result1;
\end{lstlisting}}
\newpage

\paragraph{Encoding} Format \textsf{ALU\_WPWRR1O} (Section~\ref{sec:format-ALU-WPWRR1O}), \verb|exucode2 = 0011|\\

\bitfields{ 1/P, 2/11, 1/1, 4/0011, 5/registerWo, 1/1, 2/10, 3/101, 1/0, 5/registerYo, 1/--, 5/registerZo, 1/1 }


\paragraph{Syntax} \texttt{addx16wp} \emph{registerWo} = \emph{registerZo}, \emph{registerYo}

\paragraph{Description} The \emph{registerYo} and the \emph{registerZo} are interpreted as two words packed into 64 bits. The \emph{registerZo} words are shifted left by four bits and added to the \emph{registerYo} words. The two resulting words are packed and stored into the \emph{registerWo}.


\paragraph{Execution} {Reservation table \textsf{ALU\_LITE} (Section~\ref{sec:reservation-ALU-LITE})}
~
{\begin{lstlisting}
stage RR:
  new argument3 = GPR[(registerYo<<1)+1];
  new argument2 = GPR[(registerZo<<1)+1];
  for (I = 0; I < 2; I += 1) {
    result1.32[I] = _SX_32(argument3.32[I]) + (_SX_32(argument2.32[I]) << 4)
  };
stage E1:
  GPR[(registerWo<<1)+1] = result1;
\end{lstlisting}}
\newpage

\paragraph{Encoding} Format \textsf{ALU\_WPWRR1E.M} (Section~\ref{sec:format-ALU-WPWRR1E.M}), \verb|exucode2 = 0011|\\

\bitfields{ 1/P, 2/00, 2/-- --, 27/upper27, 1/1, 2/11, 1/1, 4/0011, 5/registerWe, 1/0, 2/10, 3/101, 1/0, 1/splat32, 5/lower5, 5/registerZe, 1/1 }


\paragraph{Syntax} \texttt{addx16wp} \emph{registerWe} = \emph{registerZe}, \emph{upper27\_\discretionary{}{}{}lower5}\emph{splat32}

\paragraph{Description} The \emph{upper27\_\discretionary{}{}{}lower5} and the \emph{registerZe} are interpreted as two words packed into 64 bits. The \emph{registerZe} words are shifted left by four bits and added to the \emph{upper27\_\discretionary{}{}{}lower5} words. The two resulting words are packed and stored into the \emph{registerWe}.


\paragraph{Modifier(s)}
\noindent\begin{itemize}
\item splat32 (cf. splat32 in section \ref{par:modifier-splat32} on page \pageref{par:modifier-splat32})
\end{itemize}

\paragraph{Execution} {Reservation table \textsf{ALU\_LITE.X} (Section~\ref{sec:reservation-ALU-LITE.X})}
~
{\begin{lstlisting}
stage ID:
  new splat32 = _ZX_1(splat32);
stage RR:
  new argument3 = splat32 ? (_WX_32(upper27_lower5) << 32) | _ZX_32(_WX_32(upper27_lower5)) : _SX_32(_WX_32(upper27_lower5));
  new argument2 = GPR[registerZe<<1];
  for (I = 0; I < 2; I += 1) {
    result1.32[I] = _SX_32(argument3.32[I]) + (_SX_32(argument2.32[I]) << 4)
  };
stage E1:
  GPR[registerWe<<1] = result1;
\end{lstlisting}}
\newpage

\paragraph{Encoding} Format \textsf{ALU\_WPWRR1O.M} (Section~\ref{sec:format-ALU-WPWRR1O.M}), \verb|exucode2 = 0011|\\

\bitfields{ 1/P, 2/00, 2/-- --, 27/upper27, 1/1, 2/11, 1/1, 4/0011, 5/registerWo, 1/1, 2/10, 3/101, 1/0, 1/splat32, 5/lower5, 5/registerZo, 1/1 }


\paragraph{Syntax} \texttt{addx16wp} \emph{registerWo} = \emph{registerZo}, \emph{upper27\_\discretionary{}{}{}lower5}\emph{splat32}

\paragraph{Description} The \emph{upper27\_\discretionary{}{}{}lower5} and the \emph{registerZo} are interpreted as two words packed into 64 bits. The \emph{registerZo} words are shifted left by four bits and added to the \emph{upper27\_\discretionary{}{}{}lower5} words. The two resulting words are packed and stored into the \emph{registerWo}.


\paragraph{Modifier(s)}
\noindent\begin{itemize}
\item splat32 (cf. splat32 in section \ref{par:modifier-splat32} on page \pageref{par:modifier-splat32})
\end{itemize}

\paragraph{Execution} {Reservation table \textsf{ALU\_LITE.X} (Section~\ref{sec:reservation-ALU-LITE.X})}
~
{\begin{lstlisting}
stage ID:
  new splat32 = _ZX_1(splat32);
stage RR:
  new argument3 = splat32 ? (_WX_32(upper27_lower5) << 32) | _ZX_32(_WX_32(upper27_lower5)) : _SX_32(_WX_32(upper27_lower5));
  new argument2 = GPR[(registerZo<<1)+1];
  for (I = 0; I < 2; I += 1) {
    result1.32[I] = _SX_32(argument3.32[I]) + (_SX_32(argument2.32[I]) << 4)
  };
stage E1:
  GPR[(registerWo<<1)+1] = result1;
\end{lstlisting}}
\newpage
\subsubsection[{\small ADDX2BO: Add to Times 2 Byte Octuple}]{ADDX2BO\hfill Add to Times 2 Byte Octuple}
\label{instr:ADDX2BO}\index{addx2bo}
 
\paragraph{Encoding} Format \textsf{ALU\_BOWRR1E} (Section~\ref{sec:format-ALU-BOWRR1E}), \verb|exucode2 = 0000|\\

\bitfields{ 1/P, 2/11, 1/1, 4/0000, 5/registerWe, 1/0, 2/10, 3/101, 1/1, 5/registerYe, 1/--, 5/registerZe, 1/1 }


\paragraph{Syntax} \texttt{addx2bo} \emph{registerWe} = \emph{registerZe}, \emph{registerYe}

\paragraph{Description} The \emph{registerYe} and the \emph{registerZe} are interpreted as eight bytes packed into 64 bits. The \emph{registerZe} bytes are shifted left by one bit and added to the \emph{registerYe} bytes. The eight resulting bytes are packed and stored into the \emph{registerWe}.


\paragraph{Execution} {Reservation table \textsf{ALU\_LITE} (Section~\ref{sec:reservation-ALU-LITE})}
~
{\begin{lstlisting}
stage RR:
  new argument3 = GPR[registerYe<<1];
  new argument2 = GPR[registerZe<<1];
  for (I = 0; I < 8; I += 1) {
    result1.8[I] = _SX_8(argument3.8[I]) + (_SX_8(argument2.8[I]) << 1)
  };
stage E1:
  GPR[registerWe<<1] = result1;
\end{lstlisting}}
\newpage

\paragraph{Encoding} Format \textsf{ALU\_BOWRR1O} (Section~\ref{sec:format-ALU-BOWRR1O}), \verb|exucode2 = 0000|\\

\bitfields{ 1/P, 2/11, 1/1, 4/0000, 5/registerWo, 1/1, 2/10, 3/101, 1/1, 5/registerYo, 1/--, 5/registerZo, 1/1 }


\paragraph{Syntax} \texttt{addx2bo} \emph{registerWo} = \emph{registerZo}, \emph{registerYo}

\paragraph{Description} The \emph{registerYo} and the \emph{registerZo} are interpreted as eight bytes packed into 64 bits. The \emph{registerZo} bytes are shifted left by one bit and added to the \emph{registerYo} bytes. The eight resulting bytes are packed and stored into the \emph{registerWo}.


\paragraph{Execution} {Reservation table \textsf{ALU\_LITE} (Section~\ref{sec:reservation-ALU-LITE})}
~
{\begin{lstlisting}
stage RR:
  new argument3 = GPR[(registerYo<<1)+1];
  new argument2 = GPR[(registerZo<<1)+1];
  for (I = 0; I < 8; I += 1) {
    result1.8[I] = _SX_8(argument3.8[I]) + (_SX_8(argument2.8[I]) << 1)
  };
stage E1:
  GPR[(registerWo<<1)+1] = result1;
\end{lstlisting}}
\newpage

\paragraph{Encoding} Format \textsf{ALU\_BOWRR1E.M} (Section~\ref{sec:format-ALU-BOWRR1E.M}), \verb|exucode2 = 0000|\\

\bitfields{ 1/P, 2/00, 2/-- --, 27/upper27, 1/1, 2/11, 1/1, 4/0000, 5/registerWe, 1/0, 2/10, 3/101, 1/1, 1/splat32, 5/lower5, 5/registerZe, 1/1 }


\paragraph{Syntax} \texttt{addx2bo} \emph{registerWe} = \emph{registerZe}, \emph{upper27\_\discretionary{}{}{}lower5}\emph{splat32}

\paragraph{Description} The \emph{upper27\_\discretionary{}{}{}lower5} and the \emph{registerZe} are interpreted as eight bytes packed into 64 bits. The \emph{registerZe} bytes are shifted left by one bit and added to the \emph{upper27\_\discretionary{}{}{}lower5} bytes. The eight resulting bytes are packed and stored into the \emph{registerWe}.


\paragraph{Modifier(s)}
\noindent\begin{itemize}
\item splat32 (cf. splat32 in section \ref{par:modifier-splat32} on page \pageref{par:modifier-splat32})
\end{itemize}

\paragraph{Execution} {Reservation table \textsf{ALU\_LITE.X} (Section~\ref{sec:reservation-ALU-LITE.X})}
~
{\begin{lstlisting}
stage ID:
  new splat32 = _ZX_1(splat32);
stage RR:
  new argument3 = splat32 ? (_WX_32(upper27_lower5) << 32) | _ZX_32(_WX_32(upper27_lower5)) : _SX_32(_WX_32(upper27_lower5));
  new argument2 = GPR[registerZe<<1];
  for (I = 0; I < 8; I += 1) {
    result1.8[I] = _SX_8(argument3.8[I]) + (_SX_8(argument2.8[I]) << 1)
  };
stage E1:
  GPR[registerWe<<1] = result1;
\end{lstlisting}}
\newpage

\paragraph{Encoding} Format \textsf{ALU\_BOWRR1O.M} (Section~\ref{sec:format-ALU-BOWRR1O.M}), \verb|exucode2 = 0000|\\

\bitfields{ 1/P, 2/00, 2/-- --, 27/upper27, 1/1, 2/11, 1/1, 4/0000, 5/registerWo, 1/1, 2/10, 3/101, 1/1, 1/splat32, 5/lower5, 5/registerZo, 1/1 }


\paragraph{Syntax} \texttt{addx2bo} \emph{registerWo} = \emph{registerZo}, \emph{upper27\_\discretionary{}{}{}lower5}\emph{splat32}

\paragraph{Description} The \emph{upper27\_\discretionary{}{}{}lower5} and the \emph{registerZo} are interpreted as eight bytes packed into 64 bits. The \emph{registerZo} bytes are shifted left by one bit and added to the \emph{upper27\_\discretionary{}{}{}lower5} bytes. The eight resulting bytes are packed and stored into the \emph{registerWo}.


\paragraph{Modifier(s)}
\noindent\begin{itemize}
\item splat32 (cf. splat32 in section \ref{par:modifier-splat32} on page \pageref{par:modifier-splat32})
\end{itemize}

\paragraph{Execution} {Reservation table \textsf{ALU\_LITE.X} (Section~\ref{sec:reservation-ALU-LITE.X})}
~
{\begin{lstlisting}
stage ID:
  new splat32 = _ZX_1(splat32);
stage RR:
  new argument3 = splat32 ? (_WX_32(upper27_lower5) << 32) | _ZX_32(_WX_32(upper27_lower5)) : _SX_32(_WX_32(upper27_lower5));
  new argument2 = GPR[(registerZo<<1)+1];
  for (I = 0; I < 8; I += 1) {
    result1.8[I] = _SX_8(argument3.8[I]) + (_SX_8(argument2.8[I]) << 1)
  };
stage E1:
  GPR[(registerWo<<1)+1] = result1;
\end{lstlisting}}
\newpage
\subsubsection[{\small ADDX2D: Add to Times 2 Double Word}]{ADDX2D\hfill Add to Times 2 Double Word}
\label{instr:ADDX2D}\index{addx2d}
 
\paragraph{Encoding} Format \textsf{ALU\_DWRR1} (Section~\ref{sec:format-ALU-DWRR1}), \verb|exucode2 = 0000|\\

\bitfields{ 1/P, 2/11, 1/0, 4/0000, 6/registerW, 2/10, 3/101, 1/0, 6/registerY, 6/registerZ }


\paragraph{Syntax} \texttt{addx2d} \emph{registerW} = \emph{registerZ}, \emph{registerY}

\paragraph{Description} The \emph{registerZ} is shifted left by one bit and added to the \emph{registerY}. The result is stored into the \emph{registerW}.


\paragraph{Execution} {Reservation table \textsf{ALU\_TINY} (Section~\ref{sec:reservation-ALU-TINY})}
~
{\begin{lstlisting}
stage RR:
  new argument3 = GPR[registerY];
  new argument2 = GPR[registerZ];
  new result1 = argument3 + (argument2 << 1);
stage E1:
  GPR[registerW] = result1;
\end{lstlisting}}
\newpage

\paragraph{Encoding} Format \textsf{ALU\_DWRR1.M} (Section~\ref{sec:format-ALU-DWRR1.M}), \verb|exucode2 = 0000|\\

\bitfields{ 1/P, 2/00, 2/-- --, 27/upper27, 1/1, 2/11, 1/0, 4/0000, 6/registerW, 2/10, 3/101, 1/0, 1/splat32, 5/lower5, 6/registerZ }


\paragraph{Syntax} \texttt{addx2d} \emph{registerW} = \emph{registerZ}, \emph{upper27\_\discretionary{}{}{}lower5}\emph{splat32}

\paragraph{Description} The \emph{registerZ} is shifted left by one bit and added to the \emph{upper27\_\discretionary{}{}{}lower5}. The result is stored into the \emph{registerW}.


\paragraph{Modifier(s)}
\noindent\begin{itemize}
\item splat32 (cf. splat32 in section \ref{par:modifier-splat32} on page \pageref{par:modifier-splat32})
\end{itemize}

\paragraph{Execution} {Reservation table \textsf{ALU\_TINY.X} (Section~\ref{sec:reservation-ALU-TINY.X})}
~
{\begin{lstlisting}
stage ID:
  new splat32 = _ZX_1(splat32);
stage RR:
  new argument3 = splat32 ? (_WX_32(upper27_lower5) << 32) | _ZX_32(_WX_32(upper27_lower5)) : _SX_32(_WX_32(upper27_lower5));
  new argument2 = GPR[registerZ];
  new result1 = argument3 + (argument2 << 1);
stage E1:
  GPR[registerW] = result1;
\end{lstlisting}}
\newpage
\subsubsection[{\small ADDX2HQ: Add to Times 2 Half Word Quadruple}]{ADDX2HQ\hfill Add to Times 2 Half Word Quadruple}
\label{instr:ADDX2HQ}\index{addx2hq}
 
\paragraph{Encoding} Format \textsf{ALU\_HQWRR1E} (Section~\ref{sec:format-ALU-HQWRR1E}), \verb|exucode2 = 0000|\\

\bitfields{ 1/P, 2/11, 1/1, 4/0000, 5/registerWe, 1/0, 2/10, 3/101, 1/1, 5/registerYe, 1/--, 5/registerZe, 1/0 }


\paragraph{Syntax} \texttt{addx2hq} \emph{registerWe} = \emph{registerZe}, \emph{registerYe}

\paragraph{Description} The \emph{registerYe} and the \emph{registerZe} are interpreted as four half words packed into 64 bits. The \emph{registerZe} half words are shifted left by one bit and added to the \emph{registerYe} half words. The four resulting half words are packed and stored into the \emph{registerWe}.


\paragraph{Execution} {Reservation table \textsf{ALU\_LITE} (Section~\ref{sec:reservation-ALU-LITE})}
~
{\begin{lstlisting}
stage RR:
  new argument3 = GPR[registerYe<<1];
  new argument2 = GPR[registerZe<<1];
  for (I = 0; I < 4; I += 1) {
    result1.16[I] = _SX_16(argument3.16[I]) + (_SX_16(argument2.16[I]) << 1)
  };
stage E1:
  GPR[registerWe<<1] = result1;
\end{lstlisting}}
\newpage

\paragraph{Encoding} Format \textsf{ALU\_HQWRR1O} (Section~\ref{sec:format-ALU-HQWRR1O}), \verb|exucode2 = 0000|\\

\bitfields{ 1/P, 2/11, 1/1, 4/0000, 5/registerWo, 1/1, 2/10, 3/101, 1/1, 5/registerYo, 1/--, 5/registerZo, 1/0 }


\paragraph{Syntax} \texttt{addx2hq} \emph{registerWo} = \emph{registerZo}, \emph{registerYo}

\paragraph{Description} The \emph{registerYo} and the \emph{registerZo} are interpreted as four half words packed into 64 bits. The \emph{registerZo} half words are shifted left by one bit and added to the \emph{registerYo} half words. The four resulting half words are packed and stored into the \emph{registerWo}.


\paragraph{Execution} {Reservation table \textsf{ALU\_LITE} (Section~\ref{sec:reservation-ALU-LITE})}
~
{\begin{lstlisting}
stage RR:
  new argument3 = GPR[(registerYo<<1)+1];
  new argument2 = GPR[(registerZo<<1)+1];
  for (I = 0; I < 4; I += 1) {
    result1.16[I] = _SX_16(argument3.16[I]) + (_SX_16(argument2.16[I]) << 1)
  };
stage E1:
  GPR[(registerWo<<1)+1] = result1;
\end{lstlisting}}
\newpage

\paragraph{Encoding} Format \textsf{ALU\_HQWRR1E.M} (Section~\ref{sec:format-ALU-HQWRR1E.M}), \verb|exucode2 = 0000|\\

\bitfields{ 1/P, 2/00, 2/-- --, 27/upper27, 1/1, 2/11, 1/1, 4/0000, 5/registerWe, 1/0, 2/10, 3/101, 1/1, 1/splat32, 5/lower5, 5/registerZe, 1/0 }


\paragraph{Syntax} \texttt{addx2hq} \emph{registerWe} = \emph{registerZe}, \emph{upper27\_\discretionary{}{}{}lower5}\emph{splat32}

\paragraph{Description} The \emph{upper27\_\discretionary{}{}{}lower5} and the \emph{registerZe} are interpreted as four half words packed into 64 bits. The \emph{registerZe} half words are shifted left by one bit and added to the \emph{upper27\_\discretionary{}{}{}lower5} half words. The four resulting half words are packed and stored into the \emph{registerWe}.


\paragraph{Modifier(s)}
\noindent\begin{itemize}
\item splat32 (cf. splat32 in section \ref{par:modifier-splat32} on page \pageref{par:modifier-splat32})
\end{itemize}

\paragraph{Execution} {Reservation table \textsf{ALU\_LITE.X} (Section~\ref{sec:reservation-ALU-LITE.X})}
~
{\begin{lstlisting}
stage ID:
  new splat32 = _ZX_1(splat32);
stage RR:
  new argument3 = splat32 ? (_WX_32(upper27_lower5) << 32) | _ZX_32(_WX_32(upper27_lower5)) : _SX_32(_WX_32(upper27_lower5));
  new argument2 = GPR[registerZe<<1];
  for (I = 0; I < 4; I += 1) {
    result1.16[I] = _SX_16(argument3.16[I]) + (_SX_16(argument2.16[I]) << 1)
  };
stage E1:
  GPR[registerWe<<1] = result1;
\end{lstlisting}}
\newpage

\paragraph{Encoding} Format \textsf{ALU\_HQWRR1O.M} (Section~\ref{sec:format-ALU-HQWRR1O.M}), \verb|exucode2 = 0000|\\

\bitfields{ 1/P, 2/00, 2/-- --, 27/upper27, 1/1, 2/11, 1/1, 4/0000, 5/registerWo, 1/1, 2/10, 3/101, 1/1, 1/splat32, 5/lower5, 5/registerZo, 1/0 }


\paragraph{Syntax} \texttt{addx2hq} \emph{registerWo} = \emph{registerZo}, \emph{upper27\_\discretionary{}{}{}lower5}\emph{splat32}

\paragraph{Description} The \emph{upper27\_\discretionary{}{}{}lower5} and the \emph{registerZo} are interpreted as four half words packed into 64 bits. The \emph{registerZo} half words are shifted left by one bit and added to the \emph{upper27\_\discretionary{}{}{}lower5} half words. The four resulting half words are packed and stored into the \emph{registerWo}.


\paragraph{Modifier(s)}
\noindent\begin{itemize}
\item splat32 (cf. splat32 in section \ref{par:modifier-splat32} on page \pageref{par:modifier-splat32})
\end{itemize}

\paragraph{Execution} {Reservation table \textsf{ALU\_LITE.X} (Section~\ref{sec:reservation-ALU-LITE.X})}
~
{\begin{lstlisting}
stage ID:
  new splat32 = _ZX_1(splat32);
stage RR:
  new argument3 = splat32 ? (_WX_32(upper27_lower5) << 32) | _ZX_32(_WX_32(upper27_lower5)) : _SX_32(_WX_32(upper27_lower5));
  new argument2 = GPR[(registerZo<<1)+1];
  for (I = 0; I < 4; I += 1) {
    result1.16[I] = _SX_16(argument3.16[I]) + (_SX_16(argument2.16[I]) << 1)
  };
stage E1:
  GPR[(registerWo<<1)+1] = result1;
\end{lstlisting}}
\newpage
\subsubsection[{\small ADDX2W: Add to Times 2 Word}]{ADDX2W\hfill Add to Times 2 Word}
\label{instr:ADDX2W}\index{addx2w}
 
\paragraph{Encoding} Format \textsf{ALU\_WWRR1} (Section~\ref{sec:format-ALU-WWRR1}), \verb|exucode2 = 0000|\\

\bitfields{ 1/P, 2/11, 1/0, 4/0000, 6/registerW, 2/11, 3/101, 1/signextw, 6/registerY, 6/registerZ }


\paragraph{Syntax} \texttt{addx2w}\emph{signextw} \emph{registerW} = \emph{registerZ}, \emph{registerY}

\paragraph{Description} The \emph{registerZ} is shifted left by one bit and added to the \emph{registerY}. The result with the 32 upper bits cleared is stored into the \emph{registerW}.


\paragraph{Modifier(s)}
\noindent\begin{itemize}
\item signextw (cf. signextw in section \ref{par:modifier-signextw} on page \pageref{par:modifier-signextw})
\end{itemize}

\paragraph{Execution} {Reservation table \textsf{ALU\_TINY} (Section~\ref{sec:reservation-ALU-TINY})}
~
{\begin{lstlisting}
stage ID:
  new signextw = _ZX_1(signextw);
stage RR:
  new argument3 = GPR[registerY];
  new argument2 = GPR[registerZ];
  new result1 = argument3 + (argument2 << 1);
stage E1:
  GPR[registerW] = signextw ? _SX_32(result1) : _ZX_32(result1);
\end{lstlisting}}
\newpage

\paragraph{Encoding} Format \textsf{ALU\_WWRR1.W} (Section~\ref{sec:format-ALU-WWRR1.W}), \verb|exucode2 = 0000|\\

\bitfields{ 1/P, 2/00, 2/-- --, 27/upper27, 1/1, 2/11, 1/0, 4/0000, 6/registerW, 2/11, 3/101, 1/signextw, 1/0, 5/lower5, 6/registerZ }


\paragraph{Syntax} \texttt{addx2w}\emph{signextw} \emph{registerW} = \emph{registerZ}, \emph{upper27\_\discretionary{}{}{}lower5}

\paragraph{Description} The \emph{registerZ} is shifted left by one bit and added to the \emph{upper27\_\discretionary{}{}{}lower5}. The result with the 32 upper bits cleared is stored into the \emph{registerW}.


\paragraph{Modifier(s)}
\noindent\begin{itemize}
\item signextw (cf. signextw in section \ref{par:modifier-signextw} on page \pageref{par:modifier-signextw})
\end{itemize}

\paragraph{Execution} {Reservation table \textsf{ALU\_TINY.X} (Section~\ref{sec:reservation-ALU-TINY.X})}
~
{\begin{lstlisting}
stage ID:
  new signextw = _ZX_1(signextw);
stage RR:
  new argument3 = _WX_32(upper27_lower5);
  new argument2 = GPR[registerZ];
  new result1 = argument3 + (argument2 << 1);
stage E1:
  GPR[registerW] = signextw ? _SX_32(result1) : _ZX_32(result1);
\end{lstlisting}}
\newpage
\subsubsection[{\small ADDX2WP: Add to Times 2 Word Pair}]{ADDX2WP\hfill Add to Times 2 Word Pair}
\label{instr:ADDX2WP}\index{addx2wp}
 
\paragraph{Encoding} Format \textsf{ALU\_WPWRR1E} (Section~\ref{sec:format-ALU-WPWRR1E}), \verb|exucode2 = 0000|\\

\bitfields{ 1/P, 2/11, 1/1, 4/0000, 5/registerWe, 1/0, 2/10, 3/101, 1/0, 5/registerYe, 1/--, 5/registerZe, 1/1 }


\paragraph{Syntax} \texttt{addx2wp} \emph{registerWe} = \emph{registerZe}, \emph{registerYe}

\paragraph{Description} The \emph{registerYe} and the \emph{registerZe} are interpreted as two words packed into 64 bits. The \emph{registerZe} words are shifted left by one bit and added to the \emph{registerYe} words. The two resulting words are packed and stored into the \emph{registerWe}.


\paragraph{Execution} {Reservation table \textsf{ALU\_LITE} (Section~\ref{sec:reservation-ALU-LITE})}
~
{\begin{lstlisting}
stage RR:
  new argument3 = GPR[registerYe<<1];
  new argument2 = GPR[registerZe<<1];
  for (I = 0; I < 2; I += 1) {
    result1.32[I] = _SX_32(argument3.32[I]) + (_SX_32(argument2.32[I]) << 1)
  };
stage E1:
  GPR[registerWe<<1] = result1;
\end{lstlisting}}
\newpage

\paragraph{Encoding} Format \textsf{ALU\_WPWRR1O} (Section~\ref{sec:format-ALU-WPWRR1O}), \verb|exucode2 = 0000|\\

\bitfields{ 1/P, 2/11, 1/1, 4/0000, 5/registerWo, 1/1, 2/10, 3/101, 1/0, 5/registerYo, 1/--, 5/registerZo, 1/1 }


\paragraph{Syntax} \texttt{addx2wp} \emph{registerWo} = \emph{registerZo}, \emph{registerYo}

\paragraph{Description} The \emph{registerYo} and the \emph{registerZo} are interpreted as two words packed into 64 bits. The \emph{registerZo} words are shifted left by one bit and added to the \emph{registerYo} words. The two resulting words are packed and stored into the \emph{registerWo}.


\paragraph{Execution} {Reservation table \textsf{ALU\_LITE} (Section~\ref{sec:reservation-ALU-LITE})}
~
{\begin{lstlisting}
stage RR:
  new argument3 = GPR[(registerYo<<1)+1];
  new argument2 = GPR[(registerZo<<1)+1];
  for (I = 0; I < 2; I += 1) {
    result1.32[I] = _SX_32(argument3.32[I]) + (_SX_32(argument2.32[I]) << 1)
  };
stage E1:
  GPR[(registerWo<<1)+1] = result1;
\end{lstlisting}}
\newpage

\paragraph{Encoding} Format \textsf{ALU\_WPWRR1E.M} (Section~\ref{sec:format-ALU-WPWRR1E.M}), \verb|exucode2 = 0000|\\

\bitfields{ 1/P, 2/00, 2/-- --, 27/upper27, 1/1, 2/11, 1/1, 4/0000, 5/registerWe, 1/0, 2/10, 3/101, 1/0, 1/splat32, 5/lower5, 5/registerZe, 1/1 }


\paragraph{Syntax} \texttt{addx2wp} \emph{registerWe} = \emph{registerZe}, \emph{upper27\_\discretionary{}{}{}lower5}\emph{splat32}

\paragraph{Description} The \emph{upper27\_\discretionary{}{}{}lower5} and the \emph{registerZe} are interpreted as two words packed into 64 bits. The \emph{registerZe} words are shifted left by one bit and added to the \emph{upper27\_\discretionary{}{}{}lower5} words. The two resulting words are packed and stored into the \emph{registerWe}.


\paragraph{Modifier(s)}
\noindent\begin{itemize}
\item splat32 (cf. splat32 in section \ref{par:modifier-splat32} on page \pageref{par:modifier-splat32})
\end{itemize}

\paragraph{Execution} {Reservation table \textsf{ALU\_LITE.X} (Section~\ref{sec:reservation-ALU-LITE.X})}
~
{\begin{lstlisting}
stage ID:
  new splat32 = _ZX_1(splat32);
stage RR:
  new argument3 = splat32 ? (_WX_32(upper27_lower5) << 32) | _ZX_32(_WX_32(upper27_lower5)) : _SX_32(_WX_32(upper27_lower5));
  new argument2 = GPR[registerZe<<1];
  for (I = 0; I < 2; I += 1) {
    result1.32[I] = _SX_32(argument3.32[I]) + (_SX_32(argument2.32[I]) << 1)
  };
stage E1:
  GPR[registerWe<<1] = result1;
\end{lstlisting}}
\newpage

\paragraph{Encoding} Format \textsf{ALU\_WPWRR1O.M} (Section~\ref{sec:format-ALU-WPWRR1O.M}), \verb|exucode2 = 0000|\\

\bitfields{ 1/P, 2/00, 2/-- --, 27/upper27, 1/1, 2/11, 1/1, 4/0000, 5/registerWo, 1/1, 2/10, 3/101, 1/0, 1/splat32, 5/lower5, 5/registerZo, 1/1 }


\paragraph{Syntax} \texttt{addx2wp} \emph{registerWo} = \emph{registerZo}, \emph{upper27\_\discretionary{}{}{}lower5}\emph{splat32}

\paragraph{Description} The \emph{upper27\_\discretionary{}{}{}lower5} and the \emph{registerZo} are interpreted as two words packed into 64 bits. The \emph{registerZo} words are shifted left by one bit and added to the \emph{upper27\_\discretionary{}{}{}lower5} words. The two resulting words are packed and stored into the \emph{registerWo}.


\paragraph{Modifier(s)}
\noindent\begin{itemize}
\item splat32 (cf. splat32 in section \ref{par:modifier-splat32} on page \pageref{par:modifier-splat32})
\end{itemize}

\paragraph{Execution} {Reservation table \textsf{ALU\_LITE.X} (Section~\ref{sec:reservation-ALU-LITE.X})}
~
{\begin{lstlisting}
stage ID:
  new splat32 = _ZX_1(splat32);
stage RR:
  new argument3 = splat32 ? (_WX_32(upper27_lower5) << 32) | _ZX_32(_WX_32(upper27_lower5)) : _SX_32(_WX_32(upper27_lower5));
  new argument2 = GPR[(registerZo<<1)+1];
  for (I = 0; I < 2; I += 1) {
    result1.32[I] = _SX_32(argument3.32[I]) + (_SX_32(argument2.32[I]) << 1)
  };
stage E1:
  GPR[(registerWo<<1)+1] = result1;
\end{lstlisting}}
\newpage
\subsubsection[{\small ADDX32D: Add to Times 32 Double Word}]{ADDX32D\hfill Add to Times 32 Double Word}
\label{instr:ADDX32D}\index{addx32d}
 
\paragraph{Encoding} Format \textsf{ALU\_DWRR1} (Section~\ref{sec:format-ALU-DWRR1}), \verb|exucode2 = 0100|\\

\bitfields{ 1/P, 2/11, 1/0, 4/0100, 6/registerW, 2/10, 3/101, 1/0, 6/registerY, 6/registerZ }


\paragraph{Syntax} \texttt{addx32d} \emph{registerW} = \emph{registerZ}, \emph{registerY}

\paragraph{Description} The \emph{registerZ} is shifted left by five bits and added to the \emph{registerY}. The result is stored into the \emph{registerW}.


\paragraph{Execution} {Reservation table \textsf{ALU\_TINY} (Section~\ref{sec:reservation-ALU-TINY})}
~
{\begin{lstlisting}
stage RR:
  new argument3 = GPR[registerY];
  new argument2 = GPR[registerZ];
  new result1 = argument3 + (argument2 << 5);
stage E1:
  GPR[registerW] = result1;
\end{lstlisting}}
\newpage

\paragraph{Encoding} Format \textsf{ALU\_DWRR1.M} (Section~\ref{sec:format-ALU-DWRR1.M}), \verb|exucode2 = 0100|\\

\bitfields{ 1/P, 2/00, 2/-- --, 27/upper27, 1/1, 2/11, 1/0, 4/0100, 6/registerW, 2/10, 3/101, 1/0, 1/splat32, 5/lower5, 6/registerZ }


\paragraph{Syntax} \texttt{addx32d} \emph{registerW} = \emph{registerZ}, \emph{upper27\_\discretionary{}{}{}lower5}\emph{splat32}

\paragraph{Description} The \emph{registerZ} is shifted left by five bits and added to the \emph{upper27\_\discretionary{}{}{}lower5}. The result is stored into the \emph{registerW}.


\paragraph{Modifier(s)}
\noindent\begin{itemize}
\item splat32 (cf. splat32 in section \ref{par:modifier-splat32} on page \pageref{par:modifier-splat32})
\end{itemize}

\paragraph{Execution} {Reservation table \textsf{ALU\_TINY.X} (Section~\ref{sec:reservation-ALU-TINY.X})}
~
{\begin{lstlisting}
stage ID:
  new splat32 = _ZX_1(splat32);
stage RR:
  new argument3 = splat32 ? (_WX_32(upper27_lower5) << 32) | _ZX_32(_WX_32(upper27_lower5)) : _SX_32(_WX_32(upper27_lower5));
  new argument2 = GPR[registerZ];
  new result1 = argument3 + (argument2 << 5);
stage E1:
  GPR[registerW] = result1;
\end{lstlisting}}
\newpage
\subsubsection[{\small ADDX32W: Add to Times 32 Word}]{ADDX32W\hfill Add to Times 32 Word}
\label{instr:ADDX32W}\index{addx32w}
 
\paragraph{Encoding} Format \textsf{ALU\_WWRR1} (Section~\ref{sec:format-ALU-WWRR1}), \verb|exucode2 = 0100|\\

\bitfields{ 1/P, 2/11, 1/0, 4/0100, 6/registerW, 2/11, 3/101, 1/signextw, 6/registerY, 6/registerZ }


\paragraph{Syntax} \texttt{addx32w}\emph{signextw} \emph{registerW} = \emph{registerZ}, \emph{registerY}

\paragraph{Description} The \emph{registerZ} is shifted left by five bits and added to the \emph{registerY}. The result with the 32 upper bits cleared is stored into the \emph{registerW}.


\paragraph{Modifier(s)}
\noindent\begin{itemize}
\item signextw (cf. signextw in section \ref{par:modifier-signextw} on page \pageref{par:modifier-signextw})
\end{itemize}

\paragraph{Execution} {Reservation table \textsf{ALU\_TINY} (Section~\ref{sec:reservation-ALU-TINY})}
~
{\begin{lstlisting}
stage ID:
  new signextw = _ZX_1(signextw);
stage RR:
  new argument3 = GPR[registerY];
  new argument2 = GPR[registerZ];
  new result1 = argument3 + (argument2 << 5);
stage E1:
  GPR[registerW] = signextw ? _SX_32(result1) : _ZX_32(result1);
\end{lstlisting}}
\newpage

\paragraph{Encoding} Format \textsf{ALU\_WWRR1.W} (Section~\ref{sec:format-ALU-WWRR1.W}), \verb|exucode2 = 0100|\\

\bitfields{ 1/P, 2/00, 2/-- --, 27/upper27, 1/1, 2/11, 1/0, 4/0100, 6/registerW, 2/11, 3/101, 1/signextw, 1/0, 5/lower5, 6/registerZ }


\paragraph{Syntax} \texttt{addx32w}\emph{signextw} \emph{registerW} = \emph{registerZ}, \emph{upper27\_\discretionary{}{}{}lower5}

\paragraph{Description} The \emph{registerZ} is shifted left by five bits and added to the \emph{upper27\_\discretionary{}{}{}lower5}. The result with the 32 upper bits cleared is stored into the \emph{registerW}.


\paragraph{Modifier(s)}
\noindent\begin{itemize}
\item signextw (cf. signextw in section \ref{par:modifier-signextw} on page \pageref{par:modifier-signextw})
\end{itemize}

\paragraph{Execution} {Reservation table \textsf{ALU\_TINY.X} (Section~\ref{sec:reservation-ALU-TINY.X})}
~
{\begin{lstlisting}
stage ID:
  new signextw = _ZX_1(signextw);
stage RR:
  new argument3 = _WX_32(upper27_lower5);
  new argument2 = GPR[registerZ];
  new result1 = argument3 + (argument2 << 5);
stage E1:
  GPR[registerW] = signextw ? _SX_32(result1) : _ZX_32(result1);
\end{lstlisting}}
\newpage
\subsubsection[{\small ADDX4BO: Add to Times 4 Byte Octuple}]{ADDX4BO\hfill Add to Times 4 Byte Octuple}
\label{instr:ADDX4BO}\index{addx4bo}
 
\paragraph{Encoding} Format \textsf{ALU\_BOWRR1E} (Section~\ref{sec:format-ALU-BOWRR1E}), \verb|exucode2 = 0001|\\

\bitfields{ 1/P, 2/11, 1/1, 4/0001, 5/registerWe, 1/0, 2/10, 3/101, 1/1, 5/registerYe, 1/--, 5/registerZe, 1/1 }


\paragraph{Syntax} \texttt{addx4bo} \emph{registerWe} = \emph{registerZe}, \emph{registerYe}

\paragraph{Description} The \emph{registerYe} and the \emph{registerZe} are interpreted as eight bytes packed into 64 bits. The \emph{registerZe} bytes are shifted left by two bits and added to the \emph{registerYe} bytes. The eight resulting bytes are packed and stored into the \emph{registerWe}.


\paragraph{Execution} {Reservation table \textsf{ALU\_LITE} (Section~\ref{sec:reservation-ALU-LITE})}
~
{\begin{lstlisting}
stage RR:
  new argument3 = GPR[registerYe<<1];
  new argument2 = GPR[registerZe<<1];
  for (I = 0; I < 8; I += 1) {
    result1.8[I] = _SX_8(argument3.8[I]) + (_SX_8(argument2.8[I]) << 2)
  };
stage E1:
  GPR[registerWe<<1] = result1;
\end{lstlisting}}
\newpage

\paragraph{Encoding} Format \textsf{ALU\_BOWRR1O} (Section~\ref{sec:format-ALU-BOWRR1O}), \verb|exucode2 = 0001|\\

\bitfields{ 1/P, 2/11, 1/1, 4/0001, 5/registerWo, 1/1, 2/10, 3/101, 1/1, 5/registerYo, 1/--, 5/registerZo, 1/1 }


\paragraph{Syntax} \texttt{addx4bo} \emph{registerWo} = \emph{registerZo}, \emph{registerYo}

\paragraph{Description} The \emph{registerYo} and the \emph{registerZo} are interpreted as eight bytes packed into 64 bits. The \emph{registerZo} bytes are shifted left by two bits and added to the \emph{registerYo} bytes. The eight resulting bytes are packed and stored into the \emph{registerWo}.


\paragraph{Execution} {Reservation table \textsf{ALU\_LITE} (Section~\ref{sec:reservation-ALU-LITE})}
~
{\begin{lstlisting}
stage RR:
  new argument3 = GPR[(registerYo<<1)+1];
  new argument2 = GPR[(registerZo<<1)+1];
  for (I = 0; I < 8; I += 1) {
    result1.8[I] = _SX_8(argument3.8[I]) + (_SX_8(argument2.8[I]) << 2)
  };
stage E1:
  GPR[(registerWo<<1)+1] = result1;
\end{lstlisting}}
\newpage

\paragraph{Encoding} Format \textsf{ALU\_BOWRR1E.M} (Section~\ref{sec:format-ALU-BOWRR1E.M}), \verb|exucode2 = 0001|\\

\bitfields{ 1/P, 2/00, 2/-- --, 27/upper27, 1/1, 2/11, 1/1, 4/0001, 5/registerWe, 1/0, 2/10, 3/101, 1/1, 1/splat32, 5/lower5, 5/registerZe, 1/1 }


\paragraph{Syntax} \texttt{addx4bo} \emph{registerWe} = \emph{registerZe}, \emph{upper27\_\discretionary{}{}{}lower5}\emph{splat32}

\paragraph{Description} The \emph{upper27\_\discretionary{}{}{}lower5} and the \emph{registerZe} are interpreted as eight bytes packed into 64 bits. The \emph{registerZe} bytes are shifted left by two bits and added to the \emph{upper27\_\discretionary{}{}{}lower5} bytes. The eight resulting bytes are packed and stored into the \emph{registerWe}.


\paragraph{Modifier(s)}
\noindent\begin{itemize}
\item splat32 (cf. splat32 in section \ref{par:modifier-splat32} on page \pageref{par:modifier-splat32})
\end{itemize}

\paragraph{Execution} {Reservation table \textsf{ALU\_LITE.X} (Section~\ref{sec:reservation-ALU-LITE.X})}
~
{\begin{lstlisting}
stage ID:
  new splat32 = _ZX_1(splat32);
stage RR:
  new argument3 = splat32 ? (_WX_32(upper27_lower5) << 32) | _ZX_32(_WX_32(upper27_lower5)) : _SX_32(_WX_32(upper27_lower5));
  new argument2 = GPR[registerZe<<1];
  for (I = 0; I < 8; I += 1) {
    result1.8[I] = _SX_8(argument3.8[I]) + (_SX_8(argument2.8[I]) << 2)
  };
stage E1:
  GPR[registerWe<<1] = result1;
\end{lstlisting}}
\newpage

\paragraph{Encoding} Format \textsf{ALU\_BOWRR1O.M} (Section~\ref{sec:format-ALU-BOWRR1O.M}), \verb|exucode2 = 0001|\\

\bitfields{ 1/P, 2/00, 2/-- --, 27/upper27, 1/1, 2/11, 1/1, 4/0001, 5/registerWo, 1/1, 2/10, 3/101, 1/1, 1/splat32, 5/lower5, 5/registerZo, 1/1 }


\paragraph{Syntax} \texttt{addx4bo} \emph{registerWo} = \emph{registerZo}, \emph{upper27\_\discretionary{}{}{}lower5}\emph{splat32}

\paragraph{Description} The \emph{upper27\_\discretionary{}{}{}lower5} and the \emph{registerZo} are interpreted as eight bytes packed into 64 bits. The \emph{registerZo} bytes are shifted left by two bits and added to the \emph{upper27\_\discretionary{}{}{}lower5} bytes. The eight resulting bytes are packed and stored into the \emph{registerWo}.


\paragraph{Modifier(s)}
\noindent\begin{itemize}
\item splat32 (cf. splat32 in section \ref{par:modifier-splat32} on page \pageref{par:modifier-splat32})
\end{itemize}

\paragraph{Execution} {Reservation table \textsf{ALU\_LITE.X} (Section~\ref{sec:reservation-ALU-LITE.X})}
~
{\begin{lstlisting}
stage ID:
  new splat32 = _ZX_1(splat32);
stage RR:
  new argument3 = splat32 ? (_WX_32(upper27_lower5) << 32) | _ZX_32(_WX_32(upper27_lower5)) : _SX_32(_WX_32(upper27_lower5));
  new argument2 = GPR[(registerZo<<1)+1];
  for (I = 0; I < 8; I += 1) {
    result1.8[I] = _SX_8(argument3.8[I]) + (_SX_8(argument2.8[I]) << 2)
  };
stage E1:
  GPR[(registerWo<<1)+1] = result1;
\end{lstlisting}}
\newpage
\subsubsection[{\small ADDX4D: Add to Times 4 Double Word}]{ADDX4D\hfill Add to Times 4 Double Word}
\label{instr:ADDX4D}\index{addx4d}
 
\paragraph{Encoding} Format \textsf{ALU\_DWRR1} (Section~\ref{sec:format-ALU-DWRR1}), \verb|exucode2 = 0001|\\

\bitfields{ 1/P, 2/11, 1/0, 4/0001, 6/registerW, 2/10, 3/101, 1/0, 6/registerY, 6/registerZ }


\paragraph{Syntax} \texttt{addx4d} \emph{registerW} = \emph{registerZ}, \emph{registerY}

\paragraph{Description} The \emph{registerZ} is shifted left by two bits and added to the \emph{registerY}. The result is stored into the \emph{registerW}.


\paragraph{Execution} {Reservation table \textsf{ALU\_TINY} (Section~\ref{sec:reservation-ALU-TINY})}
~
{\begin{lstlisting}
stage RR:
  new argument3 = GPR[registerY];
  new argument2 = GPR[registerZ];
  new result1 = argument3 + (argument2 << 2);
stage E1:
  GPR[registerW] = result1;
\end{lstlisting}}
\newpage

\paragraph{Encoding} Format \textsf{ALU\_DWRR1.M} (Section~\ref{sec:format-ALU-DWRR1.M}), \verb|exucode2 = 0001|\\

\bitfields{ 1/P, 2/00, 2/-- --, 27/upper27, 1/1, 2/11, 1/0, 4/0001, 6/registerW, 2/10, 3/101, 1/0, 1/splat32, 5/lower5, 6/registerZ }


\paragraph{Syntax} \texttt{addx4d} \emph{registerW} = \emph{registerZ}, \emph{upper27\_\discretionary{}{}{}lower5}\emph{splat32}

\paragraph{Description} The \emph{registerZ} is shifted left by two bits and added to the \emph{upper27\_\discretionary{}{}{}lower5}. The result is stored into the \emph{registerW}.


\paragraph{Modifier(s)}
\noindent\begin{itemize}
\item splat32 (cf. splat32 in section \ref{par:modifier-splat32} on page \pageref{par:modifier-splat32})
\end{itemize}

\paragraph{Execution} {Reservation table \textsf{ALU\_TINY.X} (Section~\ref{sec:reservation-ALU-TINY.X})}
~
{\begin{lstlisting}
stage ID:
  new splat32 = _ZX_1(splat32);
stage RR:
  new argument3 = splat32 ? (_WX_32(upper27_lower5) << 32) | _ZX_32(_WX_32(upper27_lower5)) : _SX_32(_WX_32(upper27_lower5));
  new argument2 = GPR[registerZ];
  new result1 = argument3 + (argument2 << 2);
stage E1:
  GPR[registerW] = result1;
\end{lstlisting}}
\newpage
\subsubsection[{\small ADDX4HQ: Add to Times 4 Half Word Quadruple}]{ADDX4HQ\hfill Add to Times 4 Half Word Quadruple}
\label{instr:ADDX4HQ}\index{addx4hq}
 
\paragraph{Encoding} Format \textsf{ALU\_HQWRR1E} (Section~\ref{sec:format-ALU-HQWRR1E}), \verb|exucode2 = 0001|\\

\bitfields{ 1/P, 2/11, 1/1, 4/0001, 5/registerWe, 1/0, 2/10, 3/101, 1/1, 5/registerYe, 1/--, 5/registerZe, 1/0 }


\paragraph{Syntax} \texttt{addx4hq} \emph{registerWe} = \emph{registerZe}, \emph{registerYe}

\paragraph{Description} The \emph{registerYe} and the \emph{registerZe} are interpreted as four half words packed into 64 bits. The \emph{registerZe} half words are shifted left by two bits and added to the \emph{registerYe} half words. The four resulting half words are packed and stored into the \emph{registerWe}.


\paragraph{Execution} {Reservation table \textsf{ALU\_LITE} (Section~\ref{sec:reservation-ALU-LITE})}
~
{\begin{lstlisting}
stage RR:
  new argument3 = GPR[registerYe<<1];
  new argument2 = GPR[registerZe<<1];
  for (I = 0; I < 4; I += 1) {
    result1.16[I] = _SX_16(argument3.16[I]) + (_SX_16(argument2.16[I]) << 2)
  };
stage E1:
  GPR[registerWe<<1] = result1;
\end{lstlisting}}
\newpage

\paragraph{Encoding} Format \textsf{ALU\_HQWRR1O} (Section~\ref{sec:format-ALU-HQWRR1O}), \verb|exucode2 = 0001|\\

\bitfields{ 1/P, 2/11, 1/1, 4/0001, 5/registerWo, 1/1, 2/10, 3/101, 1/1, 5/registerYo, 1/--, 5/registerZo, 1/0 }


\paragraph{Syntax} \texttt{addx4hq} \emph{registerWo} = \emph{registerZo}, \emph{registerYo}

\paragraph{Description} The \emph{registerYo} and the \emph{registerZo} are interpreted as four half words packed into 64 bits. The \emph{registerZo} half words are shifted left by two bits and added to the \emph{registerYo} half words. The four resulting half words are packed and stored into the \emph{registerWo}.


\paragraph{Execution} {Reservation table \textsf{ALU\_LITE} (Section~\ref{sec:reservation-ALU-LITE})}
~
{\begin{lstlisting}
stage RR:
  new argument3 = GPR[(registerYo<<1)+1];
  new argument2 = GPR[(registerZo<<1)+1];
  for (I = 0; I < 4; I += 1) {
    result1.16[I] = _SX_16(argument3.16[I]) + (_SX_16(argument2.16[I]) << 2)
  };
stage E1:
  GPR[(registerWo<<1)+1] = result1;
\end{lstlisting}}
\newpage

\paragraph{Encoding} Format \textsf{ALU\_HQWRR1E.M} (Section~\ref{sec:format-ALU-HQWRR1E.M}), \verb|exucode2 = 0001|\\

\bitfields{ 1/P, 2/00, 2/-- --, 27/upper27, 1/1, 2/11, 1/1, 4/0001, 5/registerWe, 1/0, 2/10, 3/101, 1/1, 1/splat32, 5/lower5, 5/registerZe, 1/0 }


\paragraph{Syntax} \texttt{addx4hq} \emph{registerWe} = \emph{registerZe}, \emph{upper27\_\discretionary{}{}{}lower5}\emph{splat32}

\paragraph{Description} The \emph{upper27\_\discretionary{}{}{}lower5} and the \emph{registerZe} are interpreted as four half words packed into 64 bits. The \emph{registerZe} half words are shifted left by two bits and added to the \emph{upper27\_\discretionary{}{}{}lower5} half words. The four resulting half words are packed and stored into the \emph{registerWe}.


\paragraph{Modifier(s)}
\noindent\begin{itemize}
\item splat32 (cf. splat32 in section \ref{par:modifier-splat32} on page \pageref{par:modifier-splat32})
\end{itemize}

\paragraph{Execution} {Reservation table \textsf{ALU\_LITE.X} (Section~\ref{sec:reservation-ALU-LITE.X})}
~
{\begin{lstlisting}
stage ID:
  new splat32 = _ZX_1(splat32);
stage RR:
  new argument3 = splat32 ? (_WX_32(upper27_lower5) << 32) | _ZX_32(_WX_32(upper27_lower5)) : _SX_32(_WX_32(upper27_lower5));
  new argument2 = GPR[registerZe<<1];
  for (I = 0; I < 4; I += 1) {
    result1.16[I] = _SX_16(argument3.16[I]) + (_SX_16(argument2.16[I]) << 2)
  };
stage E1:
  GPR[registerWe<<1] = result1;
\end{lstlisting}}
\newpage

\paragraph{Encoding} Format \textsf{ALU\_HQWRR1O.M} (Section~\ref{sec:format-ALU-HQWRR1O.M}), \verb|exucode2 = 0001|\\

\bitfields{ 1/P, 2/00, 2/-- --, 27/upper27, 1/1, 2/11, 1/1, 4/0001, 5/registerWo, 1/1, 2/10, 3/101, 1/1, 1/splat32, 5/lower5, 5/registerZo, 1/0 }


\paragraph{Syntax} \texttt{addx4hq} \emph{registerWo} = \emph{registerZo}, \emph{upper27\_\discretionary{}{}{}lower5}\emph{splat32}

\paragraph{Description} The \emph{upper27\_\discretionary{}{}{}lower5} and the \emph{registerZo} are interpreted as four half words packed into 64 bits. The \emph{registerZo} half words are shifted left by two bits and added to the \emph{upper27\_\discretionary{}{}{}lower5} half words. The four resulting half words are packed and stored into the \emph{registerWo}.


\paragraph{Modifier(s)}
\noindent\begin{itemize}
\item splat32 (cf. splat32 in section \ref{par:modifier-splat32} on page \pageref{par:modifier-splat32})
\end{itemize}

\paragraph{Execution} {Reservation table \textsf{ALU\_LITE.X} (Section~\ref{sec:reservation-ALU-LITE.X})}
~
{\begin{lstlisting}
stage ID:
  new splat32 = _ZX_1(splat32);
stage RR:
  new argument3 = splat32 ? (_WX_32(upper27_lower5) << 32) | _ZX_32(_WX_32(upper27_lower5)) : _SX_32(_WX_32(upper27_lower5));
  new argument2 = GPR[(registerZo<<1)+1];
  for (I = 0; I < 4; I += 1) {
    result1.16[I] = _SX_16(argument3.16[I]) + (_SX_16(argument2.16[I]) << 2)
  };
stage E1:
  GPR[(registerWo<<1)+1] = result1;
\end{lstlisting}}
\newpage
\subsubsection[{\small ADDX4W: Add to Times 4 Word}]{ADDX4W\hfill Add to Times 4 Word}
\label{instr:ADDX4W}\index{addx4w}
 
\paragraph{Encoding} Format \textsf{ALU\_WWRR1} (Section~\ref{sec:format-ALU-WWRR1}), \verb|exucode2 = 0001|\\

\bitfields{ 1/P, 2/11, 1/0, 4/0001, 6/registerW, 2/11, 3/101, 1/signextw, 6/registerY, 6/registerZ }


\paragraph{Syntax} \texttt{addx4w}\emph{signextw} \emph{registerW} = \emph{registerZ}, \emph{registerY}

\paragraph{Description} The \emph{registerZ} is shifted left by two bits and added to the \emph{registerY}. The result with the 32 upper bits cleared is stored into the \emph{registerW}.


\paragraph{Modifier(s)}
\noindent\begin{itemize}
\item signextw (cf. signextw in section \ref{par:modifier-signextw} on page \pageref{par:modifier-signextw})
\end{itemize}

\paragraph{Execution} {Reservation table \textsf{ALU\_TINY} (Section~\ref{sec:reservation-ALU-TINY})}
~
{\begin{lstlisting}
stage ID:
  new signextw = _ZX_1(signextw);
stage RR:
  new argument3 = GPR[registerY];
  new argument2 = GPR[registerZ];
  new result1 = argument3 + (argument2 << 2);
stage E1:
  GPR[registerW] = signextw ? _SX_32(result1) : _ZX_32(result1);
\end{lstlisting}}
\newpage

\paragraph{Encoding} Format \textsf{ALU\_WWRR1.W} (Section~\ref{sec:format-ALU-WWRR1.W}), \verb|exucode2 = 0001|\\

\bitfields{ 1/P, 2/00, 2/-- --, 27/upper27, 1/1, 2/11, 1/0, 4/0001, 6/registerW, 2/11, 3/101, 1/signextw, 1/0, 5/lower5, 6/registerZ }


\paragraph{Syntax} \texttt{addx4w}\emph{signextw} \emph{registerW} = \emph{registerZ}, \emph{upper27\_\discretionary{}{}{}lower5}

\paragraph{Description} The \emph{registerZ} is shifted left by two bits and added to the \emph{upper27\_\discretionary{}{}{}lower5}. The result with the 32 upper bits cleared is stored into the \emph{registerW}.


\paragraph{Modifier(s)}
\noindent\begin{itemize}
\item signextw (cf. signextw in section \ref{par:modifier-signextw} on page \pageref{par:modifier-signextw})
\end{itemize}

\paragraph{Execution} {Reservation table \textsf{ALU\_TINY.X} (Section~\ref{sec:reservation-ALU-TINY.X})}
~
{\begin{lstlisting}
stage ID:
  new signextw = _ZX_1(signextw);
stage RR:
  new argument3 = _WX_32(upper27_lower5);
  new argument2 = GPR[registerZ];
  new result1 = argument3 + (argument2 << 2);
stage E1:
  GPR[registerW] = signextw ? _SX_32(result1) : _ZX_32(result1);
\end{lstlisting}}
\newpage
\subsubsection[{\small ADDX4WP: Add to Times 4 Word Pair}]{ADDX4WP\hfill Add to Times 4 Word Pair}
\label{instr:ADDX4WP}\index{addx4wp}
 
\paragraph{Encoding} Format \textsf{ALU\_WPWRR1E} (Section~\ref{sec:format-ALU-WPWRR1E}), \verb|exucode2 = 0001|\\

\bitfields{ 1/P, 2/11, 1/1, 4/0001, 5/registerWe, 1/0, 2/10, 3/101, 1/0, 5/registerYe, 1/--, 5/registerZe, 1/1 }


\paragraph{Syntax} \texttt{addx4wp} \emph{registerWe} = \emph{registerZe}, \emph{registerYe}

\paragraph{Description} The \emph{registerYe} and the \emph{registerZe} are interpreted as two words packed into 64 bits. The \emph{registerZe} words are shifted left by two bits and added to the \emph{registerYe} words. The two resulting words are packed and stored into the \emph{registerWe}.


\paragraph{Execution} {Reservation table \textsf{ALU\_LITE} (Section~\ref{sec:reservation-ALU-LITE})}
~
{\begin{lstlisting}
stage RR:
  new argument3 = GPR[registerYe<<1];
  new argument2 = GPR[registerZe<<1];
  for (I = 0; I < 2; I += 1) {
    result1.32[I] = _SX_32(argument3.32[I]) + (_SX_32(argument2.32[I]) << 2)
  };
stage E1:
  GPR[registerWe<<1] = result1;
\end{lstlisting}}
\newpage

\paragraph{Encoding} Format \textsf{ALU\_WPWRR1O} (Section~\ref{sec:format-ALU-WPWRR1O}), \verb|exucode2 = 0001|\\

\bitfields{ 1/P, 2/11, 1/1, 4/0001, 5/registerWo, 1/1, 2/10, 3/101, 1/0, 5/registerYo, 1/--, 5/registerZo, 1/1 }


\paragraph{Syntax} \texttt{addx4wp} \emph{registerWo} = \emph{registerZo}, \emph{registerYo}

\paragraph{Description} The \emph{registerYo} and the \emph{registerZo} are interpreted as two words packed into 64 bits. The \emph{registerZo} words are shifted left by two bits and added to the \emph{registerYo} words. The two resulting words are packed and stored into the \emph{registerWo}.


\paragraph{Execution} {Reservation table \textsf{ALU\_LITE} (Section~\ref{sec:reservation-ALU-LITE})}
~
{\begin{lstlisting}
stage RR:
  new argument3 = GPR[(registerYo<<1)+1];
  new argument2 = GPR[(registerZo<<1)+1];
  for (I = 0; I < 2; I += 1) {
    result1.32[I] = _SX_32(argument3.32[I]) + (_SX_32(argument2.32[I]) << 2)
  };
stage E1:
  GPR[(registerWo<<1)+1] = result1;
\end{lstlisting}}
\newpage

\paragraph{Encoding} Format \textsf{ALU\_WPWRR1E.M} (Section~\ref{sec:format-ALU-WPWRR1E.M}), \verb|exucode2 = 0001|\\

\bitfields{ 1/P, 2/00, 2/-- --, 27/upper27, 1/1, 2/11, 1/1, 4/0001, 5/registerWe, 1/0, 2/10, 3/101, 1/0, 1/splat32, 5/lower5, 5/registerZe, 1/1 }


\paragraph{Syntax} \texttt{addx4wp} \emph{registerWe} = \emph{registerZe}, \emph{upper27\_\discretionary{}{}{}lower5}\emph{splat32}

\paragraph{Description} The \emph{upper27\_\discretionary{}{}{}lower5} and the \emph{registerZe} are interpreted as two words packed into 64 bits. The \emph{registerZe} words are shifted left by two bits and added to the \emph{upper27\_\discretionary{}{}{}lower5} words. The two resulting words are packed and stored into the \emph{registerWe}.


\paragraph{Modifier(s)}
\noindent\begin{itemize}
\item splat32 (cf. splat32 in section \ref{par:modifier-splat32} on page \pageref{par:modifier-splat32})
\end{itemize}

\paragraph{Execution} {Reservation table \textsf{ALU\_LITE.X} (Section~\ref{sec:reservation-ALU-LITE.X})}
~
{\begin{lstlisting}
stage ID:
  new splat32 = _ZX_1(splat32);
stage RR:
  new argument3 = splat32 ? (_WX_32(upper27_lower5) << 32) | _ZX_32(_WX_32(upper27_lower5)) : _SX_32(_WX_32(upper27_lower5));
  new argument2 = GPR[registerZe<<1];
  for (I = 0; I < 2; I += 1) {
    result1.32[I] = _SX_32(argument3.32[I]) + (_SX_32(argument2.32[I]) << 2)
  };
stage E1:
  GPR[registerWe<<1] = result1;
\end{lstlisting}}
\newpage

\paragraph{Encoding} Format \textsf{ALU\_WPWRR1O.M} (Section~\ref{sec:format-ALU-WPWRR1O.M}), \verb|exucode2 = 0001|\\

\bitfields{ 1/P, 2/00, 2/-- --, 27/upper27, 1/1, 2/11, 1/1, 4/0001, 5/registerWo, 1/1, 2/10, 3/101, 1/0, 1/splat32, 5/lower5, 5/registerZo, 1/1 }


\paragraph{Syntax} \texttt{addx4wp} \emph{registerWo} = \emph{registerZo}, \emph{upper27\_\discretionary{}{}{}lower5}\emph{splat32}

\paragraph{Description} The \emph{upper27\_\discretionary{}{}{}lower5} and the \emph{registerZo} are interpreted as two words packed into 64 bits. The \emph{registerZo} words are shifted left by two bits and added to the \emph{upper27\_\discretionary{}{}{}lower5} words. The two resulting words are packed and stored into the \emph{registerWo}.


\paragraph{Modifier(s)}
\noindent\begin{itemize}
\item splat32 (cf. splat32 in section \ref{par:modifier-splat32} on page \pageref{par:modifier-splat32})
\end{itemize}

\paragraph{Execution} {Reservation table \textsf{ALU\_LITE.X} (Section~\ref{sec:reservation-ALU-LITE.X})}
~
{\begin{lstlisting}
stage ID:
  new splat32 = _ZX_1(splat32);
stage RR:
  new argument3 = splat32 ? (_WX_32(upper27_lower5) << 32) | _ZX_32(_WX_32(upper27_lower5)) : _SX_32(_WX_32(upper27_lower5));
  new argument2 = GPR[(registerZo<<1)+1];
  for (I = 0; I < 2; I += 1) {
    result1.32[I] = _SX_32(argument3.32[I]) + (_SX_32(argument2.32[I]) << 2)
  };
stage E1:
  GPR[(registerWo<<1)+1] = result1;
\end{lstlisting}}
\newpage
\subsubsection[{\small ADDX64D: Add to Times 64 Double Word}]{ADDX64D\hfill Add to Times 64 Double Word}
\label{instr:ADDX64D}\index{addx64d}
 
\paragraph{Encoding} Format \textsf{ALU\_DWRR1} (Section~\ref{sec:format-ALU-DWRR1}), \verb|exucode2 = 0101|\\

\bitfields{ 1/P, 2/11, 1/0, 4/0101, 6/registerW, 2/10, 3/101, 1/0, 6/registerY, 6/registerZ }


\paragraph{Syntax} \texttt{addx64d} \emph{registerW} = \emph{registerZ}, \emph{registerY}

\paragraph{Description} The \emph{registerZ} is shifted left by six bits and added to the \emph{registerY}. The result is stored into the \emph{registerW}.


\paragraph{Execution} {Reservation table \textsf{ALU\_TINY} (Section~\ref{sec:reservation-ALU-TINY})}
~
{\begin{lstlisting}
stage RR:
  new argument3 = GPR[registerY];
  new argument2 = GPR[registerZ];
  new result1 = argument3 + (argument2 << 6);
stage E1:
  GPR[registerW] = result1;
\end{lstlisting}}
\newpage

\paragraph{Encoding} Format \textsf{ALU\_DWRR1.M} (Section~\ref{sec:format-ALU-DWRR1.M}), \verb|exucode2 = 0101|\\

\bitfields{ 1/P, 2/00, 2/-- --, 27/upper27, 1/1, 2/11, 1/0, 4/0101, 6/registerW, 2/10, 3/101, 1/0, 1/splat32, 5/lower5, 6/registerZ }


\paragraph{Syntax} \texttt{addx64d} \emph{registerW} = \emph{registerZ}, \emph{upper27\_\discretionary{}{}{}lower5}\emph{splat32}

\paragraph{Description} The \emph{registerZ} is shifted left by six bits and added to the \emph{upper27\_\discretionary{}{}{}lower5}. The result is stored into the \emph{registerW}.


\paragraph{Modifier(s)}
\noindent\begin{itemize}
\item splat32 (cf. splat32 in section \ref{par:modifier-splat32} on page \pageref{par:modifier-splat32})
\end{itemize}

\paragraph{Execution} {Reservation table \textsf{ALU\_TINY.X} (Section~\ref{sec:reservation-ALU-TINY.X})}
~
{\begin{lstlisting}
stage ID:
  new splat32 = _ZX_1(splat32);
stage RR:
  new argument3 = splat32 ? (_WX_32(upper27_lower5) << 32) | _ZX_32(_WX_32(upper27_lower5)) : _SX_32(_WX_32(upper27_lower5));
  new argument2 = GPR[registerZ];
  new result1 = argument3 + (argument2 << 6);
stage E1:
  GPR[registerW] = result1;
\end{lstlisting}}
\newpage
\subsubsection[{\small ADDX64W: Add to Times 64 Word}]{ADDX64W\hfill Add to Times 64 Word}
\label{instr:ADDX64W}\index{addx64w}
 
\paragraph{Encoding} Format \textsf{ALU\_WWRR1} (Section~\ref{sec:format-ALU-WWRR1}), \verb|exucode2 = 0101|\\

\bitfields{ 1/P, 2/11, 1/0, 4/0101, 6/registerW, 2/11, 3/101, 1/signextw, 6/registerY, 6/registerZ }


\paragraph{Syntax} \texttt{addx64w}\emph{signextw} \emph{registerW} = \emph{registerZ}, \emph{registerY}

\paragraph{Description} The \emph{registerZ} is shifted left by six bits and added to the \emph{registerY}. The result with the 32 upper bits cleared is stored into the \emph{registerW}.


\paragraph{Modifier(s)}
\noindent\begin{itemize}
\item signextw (cf. signextw in section \ref{par:modifier-signextw} on page \pageref{par:modifier-signextw})
\end{itemize}

\paragraph{Execution} {Reservation table \textsf{ALU\_TINY} (Section~\ref{sec:reservation-ALU-TINY})}
~
{\begin{lstlisting}
stage ID:
  new signextw = _ZX_1(signextw);
stage RR:
  new argument3 = GPR[registerY];
  new argument2 = GPR[registerZ];
  new result1 = argument3 + (argument2 << 6);
stage E1:
  GPR[registerW] = signextw ? _SX_32(result1) : _ZX_32(result1);
\end{lstlisting}}
\newpage

\paragraph{Encoding} Format \textsf{ALU\_WWRR1.W} (Section~\ref{sec:format-ALU-WWRR1.W}), \verb|exucode2 = 0101|\\

\bitfields{ 1/P, 2/00, 2/-- --, 27/upper27, 1/1, 2/11, 1/0, 4/0101, 6/registerW, 2/11, 3/101, 1/signextw, 1/0, 5/lower5, 6/registerZ }


\paragraph{Syntax} \texttt{addx64w}\emph{signextw} \emph{registerW} = \emph{registerZ}, \emph{upper27\_\discretionary{}{}{}lower5}

\paragraph{Description} The \emph{registerZ} is shifted left by six bits and added to the \emph{upper27\_\discretionary{}{}{}lower5}. The result with the 32 upper bits cleared is stored into the \emph{registerW}.


\paragraph{Modifier(s)}
\noindent\begin{itemize}
\item signextw (cf. signextw in section \ref{par:modifier-signextw} on page \pageref{par:modifier-signextw})
\end{itemize}

\paragraph{Execution} {Reservation table \textsf{ALU\_TINY.X} (Section~\ref{sec:reservation-ALU-TINY.X})}
~
{\begin{lstlisting}
stage ID:
  new signextw = _ZX_1(signextw);
stage RR:
  new argument3 = _WX_32(upper27_lower5);
  new argument2 = GPR[registerZ];
  new result1 = argument3 + (argument2 << 6);
stage E1:
  GPR[registerW] = signextw ? _SX_32(result1) : _ZX_32(result1);
\end{lstlisting}}
\newpage
\subsubsection[{\small ADDX8BO: Add to Times 8 Byte Octuple}]{ADDX8BO\hfill Add to Times 8 Byte Octuple}
\label{instr:ADDX8BO}\index{addx8bo}
 
\paragraph{Encoding} Format \textsf{ALU\_BOWRR1E} (Section~\ref{sec:format-ALU-BOWRR1E}), \verb|exucode2 = 0010|\\

\bitfields{ 1/P, 2/11, 1/1, 4/0010, 5/registerWe, 1/0, 2/10, 3/101, 1/1, 5/registerYe, 1/--, 5/registerZe, 1/1 }


\paragraph{Syntax} \texttt{addx8bo} \emph{registerWe} = \emph{registerZe}, \emph{registerYe}

\paragraph{Description} The \emph{registerYe} and the \emph{registerZe} are interpreted as eight bytes packed into 64 bits. The \emph{registerZe} bytes are shifted left by three bits and added to the \emph{registerYe} bytes. The eight resulting bytes are packed and stored into the \emph{registerWe}.


\paragraph{Execution} {Reservation table \textsf{ALU\_LITE} (Section~\ref{sec:reservation-ALU-LITE})}
~
{\begin{lstlisting}
stage RR:
  new argument3 = GPR[registerYe<<1];
  new argument2 = GPR[registerZe<<1];
  for (I = 0; I < 8; I += 1) {
    result1.8[I] = _SX_8(argument3.8[I]) + (_SX_8(argument2.8[I]) << 3)
  };
stage E1:
  GPR[registerWe<<1] = result1;
\end{lstlisting}}
\newpage

\paragraph{Encoding} Format \textsf{ALU\_BOWRR1O} (Section~\ref{sec:format-ALU-BOWRR1O}), \verb|exucode2 = 0010|\\

\bitfields{ 1/P, 2/11, 1/1, 4/0010, 5/registerWo, 1/1, 2/10, 3/101, 1/1, 5/registerYo, 1/--, 5/registerZo, 1/1 }


\paragraph{Syntax} \texttt{addx8bo} \emph{registerWo} = \emph{registerZo}, \emph{registerYo}

\paragraph{Description} The \emph{registerYo} and the \emph{registerZo} are interpreted as eight bytes packed into 64 bits. The \emph{registerZo} bytes are shifted left by three bits and added to the \emph{registerYo} bytes. The eight resulting bytes are packed and stored into the \emph{registerWo}.


\paragraph{Execution} {Reservation table \textsf{ALU\_LITE} (Section~\ref{sec:reservation-ALU-LITE})}
~
{\begin{lstlisting}
stage RR:
  new argument3 = GPR[(registerYo<<1)+1];
  new argument2 = GPR[(registerZo<<1)+1];
  for (I = 0; I < 8; I += 1) {
    result1.8[I] = _SX_8(argument3.8[I]) + (_SX_8(argument2.8[I]) << 3)
  };
stage E1:
  GPR[(registerWo<<1)+1] = result1;
\end{lstlisting}}
\newpage

\paragraph{Encoding} Format \textsf{ALU\_BOWRR1E.M} (Section~\ref{sec:format-ALU-BOWRR1E.M}), \verb|exucode2 = 0010|\\

\bitfields{ 1/P, 2/00, 2/-- --, 27/upper27, 1/1, 2/11, 1/1, 4/0010, 5/registerWe, 1/0, 2/10, 3/101, 1/1, 1/splat32, 5/lower5, 5/registerZe, 1/1 }


\paragraph{Syntax} \texttt{addx8bo} \emph{registerWe} = \emph{registerZe}, \emph{upper27\_\discretionary{}{}{}lower5}\emph{splat32}

\paragraph{Description} The \emph{upper27\_\discretionary{}{}{}lower5} and the \emph{registerZe} are interpreted as eight bytes packed into 64 bits. The \emph{registerZe} bytes are shifted left by three bits and added to the \emph{upper27\_\discretionary{}{}{}lower5} bytes. The eight resulting bytes are packed and stored into the \emph{registerWe}.


\paragraph{Modifier(s)}
\noindent\begin{itemize}
\item splat32 (cf. splat32 in section \ref{par:modifier-splat32} on page \pageref{par:modifier-splat32})
\end{itemize}

\paragraph{Execution} {Reservation table \textsf{ALU\_LITE.X} (Section~\ref{sec:reservation-ALU-LITE.X})}
~
{\begin{lstlisting}
stage ID:
  new splat32 = _ZX_1(splat32);
stage RR:
  new argument3 = splat32 ? (_WX_32(upper27_lower5) << 32) | _ZX_32(_WX_32(upper27_lower5)) : _SX_32(_WX_32(upper27_lower5));
  new argument2 = GPR[registerZe<<1];
  for (I = 0; I < 8; I += 1) {
    result1.8[I] = _SX_8(argument3.8[I]) + (_SX_8(argument2.8[I]) << 3)
  };
stage E1:
  GPR[registerWe<<1] = result1;
\end{lstlisting}}
\newpage

\paragraph{Encoding} Format \textsf{ALU\_BOWRR1O.M} (Section~\ref{sec:format-ALU-BOWRR1O.M}), \verb|exucode2 = 0010|\\

\bitfields{ 1/P, 2/00, 2/-- --, 27/upper27, 1/1, 2/11, 1/1, 4/0010, 5/registerWo, 1/1, 2/10, 3/101, 1/1, 1/splat32, 5/lower5, 5/registerZo, 1/1 }


\paragraph{Syntax} \texttt{addx8bo} \emph{registerWo} = \emph{registerZo}, \emph{upper27\_\discretionary{}{}{}lower5}\emph{splat32}

\paragraph{Description} The \emph{upper27\_\discretionary{}{}{}lower5} and the \emph{registerZo} are interpreted as eight bytes packed into 64 bits. The \emph{registerZo} bytes are shifted left by three bits and added to the \emph{upper27\_\discretionary{}{}{}lower5} bytes. The eight resulting bytes are packed and stored into the \emph{registerWo}.


\paragraph{Modifier(s)}
\noindent\begin{itemize}
\item splat32 (cf. splat32 in section \ref{par:modifier-splat32} on page \pageref{par:modifier-splat32})
\end{itemize}

\paragraph{Execution} {Reservation table \textsf{ALU\_LITE.X} (Section~\ref{sec:reservation-ALU-LITE.X})}
~
{\begin{lstlisting}
stage ID:
  new splat32 = _ZX_1(splat32);
stage RR:
  new argument3 = splat32 ? (_WX_32(upper27_lower5) << 32) | _ZX_32(_WX_32(upper27_lower5)) : _SX_32(_WX_32(upper27_lower5));
  new argument2 = GPR[(registerZo<<1)+1];
  for (I = 0; I < 8; I += 1) {
    result1.8[I] = _SX_8(argument3.8[I]) + (_SX_8(argument2.8[I]) << 3)
  };
stage E1:
  GPR[(registerWo<<1)+1] = result1;
\end{lstlisting}}
\newpage
\subsubsection[{\small ADDX8D: Add to Times 8 Double Word}]{ADDX8D\hfill Add to Times 8 Double Word}
\label{instr:ADDX8D}\index{addx8d}
 
\paragraph{Encoding} Format \textsf{ALU\_DWRR1} (Section~\ref{sec:format-ALU-DWRR1}), \verb|exucode2 = 0010|\\

\bitfields{ 1/P, 2/11, 1/0, 4/0010, 6/registerW, 2/10, 3/101, 1/0, 6/registerY, 6/registerZ }


\paragraph{Syntax} \texttt{addx8d} \emph{registerW} = \emph{registerZ}, \emph{registerY}

\paragraph{Description} The \emph{registerZ} is shifted left by three bits and added to the \emph{registerY}. The result is stored into the \emph{registerW}.


\paragraph{Execution} {Reservation table \textsf{ALU\_TINY} (Section~\ref{sec:reservation-ALU-TINY})}
~
{\begin{lstlisting}
stage RR:
  new argument3 = GPR[registerY];
  new argument2 = GPR[registerZ];
  new result1 = argument3 + (argument2 << 3);
stage E1:
  GPR[registerW] = result1;
\end{lstlisting}}
\newpage

\paragraph{Encoding} Format \textsf{ALU\_DWRR1.M} (Section~\ref{sec:format-ALU-DWRR1.M}), \verb|exucode2 = 0010|\\

\bitfields{ 1/P, 2/00, 2/-- --, 27/upper27, 1/1, 2/11, 1/0, 4/0010, 6/registerW, 2/10, 3/101, 1/0, 1/splat32, 5/lower5, 6/registerZ }


\paragraph{Syntax} \texttt{addx8d} \emph{registerW} = \emph{registerZ}, \emph{upper27\_\discretionary{}{}{}lower5}\emph{splat32}

\paragraph{Description} The \emph{registerZ} is shifted left by three bits and added to the \emph{upper27\_\discretionary{}{}{}lower5}. The result is stored into the \emph{registerW}.


\paragraph{Modifier(s)}
\noindent\begin{itemize}
\item splat32 (cf. splat32 in section \ref{par:modifier-splat32} on page \pageref{par:modifier-splat32})
\end{itemize}

\paragraph{Execution} {Reservation table \textsf{ALU\_TINY.X} (Section~\ref{sec:reservation-ALU-TINY.X})}
~
{\begin{lstlisting}
stage ID:
  new splat32 = _ZX_1(splat32);
stage RR:
  new argument3 = splat32 ? (_WX_32(upper27_lower5) << 32) | _ZX_32(_WX_32(upper27_lower5)) : _SX_32(_WX_32(upper27_lower5));
  new argument2 = GPR[registerZ];
  new result1 = argument3 + (argument2 << 3);
stage E1:
  GPR[registerW] = result1;
\end{lstlisting}}
\newpage
\subsubsection[{\small ADDX8HQ: Add to Times 8 Half Word Quadruple}]{ADDX8HQ\hfill Add to Times 8 Half Word Quadruple}
\label{instr:ADDX8HQ}\index{addx8hq}
 
\paragraph{Encoding} Format \textsf{ALU\_HQWRR1E} (Section~\ref{sec:format-ALU-HQWRR1E}), \verb|exucode2 = 0010|\\

\bitfields{ 1/P, 2/11, 1/1, 4/0010, 5/registerWe, 1/0, 2/10, 3/101, 1/1, 5/registerYe, 1/--, 5/registerZe, 1/0 }


\paragraph{Syntax} \texttt{addx8hq} \emph{registerWe} = \emph{registerZe}, \emph{registerYe}

\paragraph{Description} The \emph{registerYe} and the \emph{registerZe} are interpreted as four half words packed into 64 bits. The \emph{registerZe} half words are shifted left by three bits and added to the \emph{registerYe} half words. The four resulting half words are packed and stored into the \emph{registerWe}.


\paragraph{Execution} {Reservation table \textsf{ALU\_LITE} (Section~\ref{sec:reservation-ALU-LITE})}
~
{\begin{lstlisting}
stage RR:
  new argument3 = GPR[registerYe<<1];
  new argument2 = GPR[registerZe<<1];
  for (I = 0; I < 4; I += 1) {
    result1.16[I] = _SX_16(argument3.16[I]) + (_SX_16(argument2.16[I]) << 3)
  };
stage E1:
  GPR[registerWe<<1] = result1;
\end{lstlisting}}
\newpage

\paragraph{Encoding} Format \textsf{ALU\_HQWRR1O} (Section~\ref{sec:format-ALU-HQWRR1O}), \verb|exucode2 = 0010|\\

\bitfields{ 1/P, 2/11, 1/1, 4/0010, 5/registerWo, 1/1, 2/10, 3/101, 1/1, 5/registerYo, 1/--, 5/registerZo, 1/0 }


\paragraph{Syntax} \texttt{addx8hq} \emph{registerWo} = \emph{registerZo}, \emph{registerYo}

\paragraph{Description} The \emph{registerYo} and the \emph{registerZo} are interpreted as four half words packed into 64 bits. The \emph{registerZo} half words are shifted left by three bits and added to the \emph{registerYo} half words. The four resulting half words are packed and stored into the \emph{registerWo}.


\paragraph{Execution} {Reservation table \textsf{ALU\_LITE} (Section~\ref{sec:reservation-ALU-LITE})}
~
{\begin{lstlisting}
stage RR:
  new argument3 = GPR[(registerYo<<1)+1];
  new argument2 = GPR[(registerZo<<1)+1];
  for (I = 0; I < 4; I += 1) {
    result1.16[I] = _SX_16(argument3.16[I]) + (_SX_16(argument2.16[I]) << 3)
  };
stage E1:
  GPR[(registerWo<<1)+1] = result1;
\end{lstlisting}}
\newpage

\paragraph{Encoding} Format \textsf{ALU\_HQWRR1E.M} (Section~\ref{sec:format-ALU-HQWRR1E.M}), \verb|exucode2 = 0010|\\

\bitfields{ 1/P, 2/00, 2/-- --, 27/upper27, 1/1, 2/11, 1/1, 4/0010, 5/registerWe, 1/0, 2/10, 3/101, 1/1, 1/splat32, 5/lower5, 5/registerZe, 1/0 }


\paragraph{Syntax} \texttt{addx8hq} \emph{registerWe} = \emph{registerZe}, \emph{upper27\_\discretionary{}{}{}lower5}\emph{splat32}

\paragraph{Description} The \emph{upper27\_\discretionary{}{}{}lower5} and the \emph{registerZe} are interpreted as four half words packed into 64 bits. The \emph{registerZe} half words are shifted left by three bits and added to the \emph{upper27\_\discretionary{}{}{}lower5} half words. The four resulting half words are packed and stored into the \emph{registerWe}.


\paragraph{Modifier(s)}
\noindent\begin{itemize}
\item splat32 (cf. splat32 in section \ref{par:modifier-splat32} on page \pageref{par:modifier-splat32})
\end{itemize}

\paragraph{Execution} {Reservation table \textsf{ALU\_LITE.X} (Section~\ref{sec:reservation-ALU-LITE.X})}
~
{\begin{lstlisting}
stage ID:
  new splat32 = _ZX_1(splat32);
stage RR:
  new argument3 = splat32 ? (_WX_32(upper27_lower5) << 32) | _ZX_32(_WX_32(upper27_lower5)) : _SX_32(_WX_32(upper27_lower5));
  new argument2 = GPR[registerZe<<1];
  for (I = 0; I < 4; I += 1) {
    result1.16[I] = _SX_16(argument3.16[I]) + (_SX_16(argument2.16[I]) << 3)
  };
stage E1:
  GPR[registerWe<<1] = result1;
\end{lstlisting}}
\newpage

\paragraph{Encoding} Format \textsf{ALU\_HQWRR1O.M} (Section~\ref{sec:format-ALU-HQWRR1O.M}), \verb|exucode2 = 0010|\\

\bitfields{ 1/P, 2/00, 2/-- --, 27/upper27, 1/1, 2/11, 1/1, 4/0010, 5/registerWo, 1/1, 2/10, 3/101, 1/1, 1/splat32, 5/lower5, 5/registerZo, 1/0 }


\paragraph{Syntax} \texttt{addx8hq} \emph{registerWo} = \emph{registerZo}, \emph{upper27\_\discretionary{}{}{}lower5}\emph{splat32}

\paragraph{Description} The \emph{upper27\_\discretionary{}{}{}lower5} and the \emph{registerZo} are interpreted as four half words packed into 64 bits. The \emph{registerZo} half words are shifted left by three bits and added to the \emph{upper27\_\discretionary{}{}{}lower5} half words. The four resulting half words are packed and stored into the \emph{registerWo}.


\paragraph{Modifier(s)}
\noindent\begin{itemize}
\item splat32 (cf. splat32 in section \ref{par:modifier-splat32} on page \pageref{par:modifier-splat32})
\end{itemize}

\paragraph{Execution} {Reservation table \textsf{ALU\_LITE.X} (Section~\ref{sec:reservation-ALU-LITE.X})}
~
{\begin{lstlisting}
stage ID:
  new splat32 = _ZX_1(splat32);
stage RR:
  new argument3 = splat32 ? (_WX_32(upper27_lower5) << 32) | _ZX_32(_WX_32(upper27_lower5)) : _SX_32(_WX_32(upper27_lower5));
  new argument2 = GPR[(registerZo<<1)+1];
  for (I = 0; I < 4; I += 1) {
    result1.16[I] = _SX_16(argument3.16[I]) + (_SX_16(argument2.16[I]) << 3)
  };
stage E1:
  GPR[(registerWo<<1)+1] = result1;
\end{lstlisting}}
\newpage
\subsubsection[{\small ADDX8W: Add to Times 8 Word}]{ADDX8W\hfill Add to Times 8 Word}
\label{instr:ADDX8W}\index{addx8w}
 
\paragraph{Encoding} Format \textsf{ALU\_WWRR1} (Section~\ref{sec:format-ALU-WWRR1}), \verb|exucode2 = 0010|\\

\bitfields{ 1/P, 2/11, 1/0, 4/0010, 6/registerW, 2/11, 3/101, 1/signextw, 6/registerY, 6/registerZ }


\paragraph{Syntax} \texttt{addx8w}\emph{signextw} \emph{registerW} = \emph{registerZ}, \emph{registerY}

\paragraph{Description} The \emph{registerZ} is shifted left by three bits and added to the \emph{registerY}. The result with the 32 upper bits cleared is stored into the \emph{registerW}.


\paragraph{Modifier(s)}
\noindent\begin{itemize}
\item signextw (cf. signextw in section \ref{par:modifier-signextw} on page \pageref{par:modifier-signextw})
\end{itemize}

\paragraph{Execution} {Reservation table \textsf{ALU\_TINY} (Section~\ref{sec:reservation-ALU-TINY})}
~
{\begin{lstlisting}
stage ID:
  new signextw = _ZX_1(signextw);
stage RR:
  new argument3 = GPR[registerY];
  new argument2 = GPR[registerZ];
  new result1 = argument3 + (argument2 << 3);
stage E1:
  GPR[registerW] = signextw ? _SX_32(result1) : _ZX_32(result1);
\end{lstlisting}}
\newpage

\paragraph{Encoding} Format \textsf{ALU\_WWRR1.W} (Section~\ref{sec:format-ALU-WWRR1.W}), \verb|exucode2 = 0010|\\

\bitfields{ 1/P, 2/00, 2/-- --, 27/upper27, 1/1, 2/11, 1/0, 4/0010, 6/registerW, 2/11, 3/101, 1/signextw, 1/0, 5/lower5, 6/registerZ }


\paragraph{Syntax} \texttt{addx8w}\emph{signextw} \emph{registerW} = \emph{registerZ}, \emph{upper27\_\discretionary{}{}{}lower5}

\paragraph{Description} The \emph{registerZ} is shifted left by three bits and added to the \emph{upper27\_\discretionary{}{}{}lower5}. The result with the 32 upper bits cleared is stored into the \emph{registerW}.


\paragraph{Modifier(s)}
\noindent\begin{itemize}
\item signextw (cf. signextw in section \ref{par:modifier-signextw} on page \pageref{par:modifier-signextw})
\end{itemize}

\paragraph{Execution} {Reservation table \textsf{ALU\_TINY.X} (Section~\ref{sec:reservation-ALU-TINY.X})}
~
{\begin{lstlisting}
stage ID:
  new signextw = _ZX_1(signextw);
stage RR:
  new argument3 = _WX_32(upper27_lower5);
  new argument2 = GPR[registerZ];
  new result1 = argument3 + (argument2 << 3);
stage E1:
  GPR[registerW] = signextw ? _SX_32(result1) : _ZX_32(result1);
\end{lstlisting}}
\newpage
\subsubsection[{\small ADDX8WP: Add to Times 8 Word Pair}]{ADDX8WP\hfill Add to Times 8 Word Pair}
\label{instr:ADDX8WP}\index{addx8wp}
 
\paragraph{Encoding} Format \textsf{ALU\_WPWRR1E} (Section~\ref{sec:format-ALU-WPWRR1E}), \verb|exucode2 = 0010|\\

\bitfields{ 1/P, 2/11, 1/1, 4/0010, 5/registerWe, 1/0, 2/10, 3/101, 1/0, 5/registerYe, 1/--, 5/registerZe, 1/1 }


\paragraph{Syntax} \texttt{addx8wp} \emph{registerWe} = \emph{registerZe}, \emph{registerYe}

\paragraph{Description} The \emph{registerYe} and the \emph{registerZe} are interpreted as two words packed into 64 bits. The \emph{registerZe} words are shifted left by three bits and added to the \emph{registerYe} words. The two resulting words are packed and stored into the \emph{registerWe}.


\paragraph{Execution} {Reservation table \textsf{ALU\_LITE} (Section~\ref{sec:reservation-ALU-LITE})}
~
{\begin{lstlisting}
stage RR:
  new argument3 = GPR[registerYe<<1];
  new argument2 = GPR[registerZe<<1];
  for (I = 0; I < 2; I += 1) {
    result1.32[I] = _SX_32(argument3.32[I]) + (_SX_32(argument2.32[I]) << 3)
  };
stage E1:
  GPR[registerWe<<1] = result1;
\end{lstlisting}}
\newpage

\paragraph{Encoding} Format \textsf{ALU\_WPWRR1O} (Section~\ref{sec:format-ALU-WPWRR1O}), \verb|exucode2 = 0010|\\

\bitfields{ 1/P, 2/11, 1/1, 4/0010, 5/registerWo, 1/1, 2/10, 3/101, 1/0, 5/registerYo, 1/--, 5/registerZo, 1/1 }


\paragraph{Syntax} \texttt{addx8wp} \emph{registerWo} = \emph{registerZo}, \emph{registerYo}

\paragraph{Description} The \emph{registerYo} and the \emph{registerZo} are interpreted as two words packed into 64 bits. The \emph{registerZo} words are shifted left by three bits and added to the \emph{registerYo} words. The two resulting words are packed and stored into the \emph{registerWo}.


\paragraph{Execution} {Reservation table \textsf{ALU\_LITE} (Section~\ref{sec:reservation-ALU-LITE})}
~
{\begin{lstlisting}
stage RR:
  new argument3 = GPR[(registerYo<<1)+1];
  new argument2 = GPR[(registerZo<<1)+1];
  for (I = 0; I < 2; I += 1) {
    result1.32[I] = _SX_32(argument3.32[I]) + (_SX_32(argument2.32[I]) << 3)
  };
stage E1:
  GPR[(registerWo<<1)+1] = result1;
\end{lstlisting}}
\newpage

\paragraph{Encoding} Format \textsf{ALU\_WPWRR1E.M} (Section~\ref{sec:format-ALU-WPWRR1E.M}), \verb|exucode2 = 0010|\\

\bitfields{ 1/P, 2/00, 2/-- --, 27/upper27, 1/1, 2/11, 1/1, 4/0010, 5/registerWe, 1/0, 2/10, 3/101, 1/0, 1/splat32, 5/lower5, 5/registerZe, 1/1 }


\paragraph{Syntax} \texttt{addx8wp} \emph{registerWe} = \emph{registerZe}, \emph{upper27\_\discretionary{}{}{}lower5}\emph{splat32}

\paragraph{Description} The \emph{upper27\_\discretionary{}{}{}lower5} and the \emph{registerZe} are interpreted as two words packed into 64 bits. The \emph{registerZe} words are shifted left by three bits and added to the \emph{upper27\_\discretionary{}{}{}lower5} words. The two resulting words are packed and stored into the \emph{registerWe}.


\paragraph{Modifier(s)}
\noindent\begin{itemize}
\item splat32 (cf. splat32 in section \ref{par:modifier-splat32} on page \pageref{par:modifier-splat32})
\end{itemize}

\paragraph{Execution} {Reservation table \textsf{ALU\_LITE.X} (Section~\ref{sec:reservation-ALU-LITE.X})}
~
{\begin{lstlisting}
stage ID:
  new splat32 = _ZX_1(splat32);
stage RR:
  new argument3 = splat32 ? (_WX_32(upper27_lower5) << 32) | _ZX_32(_WX_32(upper27_lower5)) : _SX_32(_WX_32(upper27_lower5));
  new argument2 = GPR[registerZe<<1];
  for (I = 0; I < 2; I += 1) {
    result1.32[I] = _SX_32(argument3.32[I]) + (_SX_32(argument2.32[I]) << 3)
  };
stage E1:
  GPR[registerWe<<1] = result1;
\end{lstlisting}}
\newpage

\paragraph{Encoding} Format \textsf{ALU\_WPWRR1O.M} (Section~\ref{sec:format-ALU-WPWRR1O.M}), \verb|exucode2 = 0010|\\

\bitfields{ 1/P, 2/00, 2/-- --, 27/upper27, 1/1, 2/11, 1/1, 4/0010, 5/registerWo, 1/1, 2/10, 3/101, 1/0, 1/splat32, 5/lower5, 5/registerZo, 1/1 }


\paragraph{Syntax} \texttt{addx8wp} \emph{registerWo} = \emph{registerZo}, \emph{upper27\_\discretionary{}{}{}lower5}\emph{splat32}

\paragraph{Description} The \emph{upper27\_\discretionary{}{}{}lower5} and the \emph{registerZo} are interpreted as two words packed into 64 bits. The \emph{registerZo} words are shifted left by three bits and added to the \emph{upper27\_\discretionary{}{}{}lower5} words. The two resulting words are packed and stored into the \emph{registerWo}.


\paragraph{Modifier(s)}
\noindent\begin{itemize}
\item splat32 (cf. splat32 in section \ref{par:modifier-splat32} on page \pageref{par:modifier-splat32})
\end{itemize}

\paragraph{Execution} {Reservation table \textsf{ALU\_LITE.X} (Section~\ref{sec:reservation-ALU-LITE.X})}
~
{\begin{lstlisting}
stage ID:
  new splat32 = _ZX_1(splat32);
stage RR:
  new argument3 = splat32 ? (_WX_32(upper27_lower5) << 32) | _ZX_32(_WX_32(upper27_lower5)) : _SX_32(_WX_32(upper27_lower5));
  new argument2 = GPR[(registerZo<<1)+1];
  for (I = 0; I < 2; I += 1) {
    result1.32[I] = _SX_32(argument3.32[I]) + (_SX_32(argument2.32[I]) << 3)
  };
stage E1:
  GPR[(registerWo<<1)+1] = result1;
\end{lstlisting}}
\newpage
\subsubsection[{\small ANDD: Bitwise And Double Word}]{ANDD\hfill Bitwise And Double Word}
\label{instr:ANDD}\index{andd}
 
\paragraph{Encoding} Format \textsf{ALU\_DWRR0} (Section~\ref{sec:format-ALU-DWRR0}), \verb|exucode2 = 1000|\\

\bitfields{ 1/P, 2/11, 1/0, 4/1000, 6/registerW, 2/10, 3/000, 1/0, 6/registerY, 6/registerZ }


\paragraph{Syntax} \texttt{andd} \emph{registerW} = \emph{registerZ}, \emph{registerY}

\paragraph{Description} Bitwise and of the \emph{registerZ} with the \emph{registerY}. The result is stored into the \emph{registerW}.


\paragraph{Execution} {Reservation table \textsf{ALU\_TINY} (Section~\ref{sec:reservation-ALU-TINY})}
~
{\begin{lstlisting}
stage RR:
  new argument3 = GPR[registerY];
  new argument2 = GPR[registerZ];
  new result1 = argument2 & argument3;
stage E1:
  GPR[registerW] = result1;
\end{lstlisting}}
\newpage

\paragraph{Encoding} Format \textsf{ALU\_DWRR0.M} (Section~\ref{sec:format-ALU-DWRR0.M}), \verb|exucode2 = 1000|\\

\bitfields{ 1/P, 2/00, 2/-- --, 27/upper27, 1/1, 2/11, 1/0, 4/1000, 6/registerW, 2/10, 3/000, 1/0, 1/splat32, 5/lower5, 6/registerZ }


\paragraph{Syntax} \texttt{andd} \emph{registerW} = \emph{registerZ}, \emph{upper27\_\discretionary{}{}{}lower5}\emph{splat32}

\paragraph{Description} Bitwise and of the \emph{registerZ} with the \emph{upper27\_\discretionary{}{}{}lower5}. The result is stored into the \emph{registerW}.


\paragraph{Modifier(s)}
\noindent\begin{itemize}
\item splat32 (cf. splat32 in section \ref{par:modifier-splat32} on page \pageref{par:modifier-splat32})
\end{itemize}

\paragraph{Execution} {Reservation table \textsf{ALU\_TINY.X} (Section~\ref{sec:reservation-ALU-TINY.X})}
~
{\begin{lstlisting}
stage ID:
  new splat32 = _ZX_1(splat32);
stage RR:
  new argument3 = splat32 ? (_WX_32(upper27_lower5) << 32) | _ZX_32(_WX_32(upper27_lower5)) : _SX_32(_WX_32(upper27_lower5));
  new argument2 = GPR[registerZ];
  new result1 = argument2 & argument3;
stage E1:
  GPR[registerW] = result1;
\end{lstlisting}}
\newpage

\paragraph{Encoding} Format \textsf{ALU\_DWRI} (Section~\ref{sec:format-ALU-DWRI}), \verb|exucode2 = 1000|\\

\bitfields{ 1/P, 2/11, 1/0, 4/1000, 6/registerW, 2/00, 10/signed10, 6/registerZ }


\paragraph{Syntax} \texttt{andd} \emph{registerW} = \emph{registerZ}, \emph{signed10}

\paragraph{Description} Bitwise and of the \emph{registerZ} with the \emph{signed10}. The result is stored into the \emph{registerW}.


\paragraph{Execution} {Reservation table \textsf{ALU\_TINY} (Section~\ref{sec:reservation-ALU-TINY})}
~
{\begin{lstlisting}
stage RR:
  new argument3 = _SX_10(signed10);
  new argument2 = GPR[registerZ];
  new result1 = argument2 & argument3;
stage E1:
  GPR[registerW] = result1;
\end{lstlisting}}
\newpage

\paragraph{Encoding} Format \textsf{ALU\_DWRI.X} (Section~\ref{sec:format-ALU-DWRI.X}), \verb|exucode2 = 1000|\\

\bitfields{ 1/P, 2/00, 2/-- --, 27/upper27, 1/1, 2/11, 1/0, 4/1000, 6/registerW, 2/00, 10/lower10, 6/registerZ }


\paragraph{Syntax} \texttt{andd} \emph{registerW} = \emph{registerZ}, \emph{upper27\_\discretionary{}{}{}lower10}

\paragraph{Description} Bitwise and of the \emph{registerZ} with the \emph{upper27\_\discretionary{}{}{}lower10}. The result is stored into the \emph{registerW}.


\paragraph{Execution} {Reservation table \textsf{ALU\_TINY.X} (Section~\ref{sec:reservation-ALU-TINY.X})}
~
{\begin{lstlisting}
stage RR:
  new argument3 = _SX_37(upper27_lower10);
  new argument2 = GPR[registerZ];
  new result1 = argument2 & argument3;
stage E1:
  GPR[registerW] = result1;
\end{lstlisting}}
\newpage

\paragraph{Encoding} Format \textsf{ALU\_DWRI.Y} (Section~\ref{sec:format-ALU-DWRI.Y}), \verb|exucode2 = 1000|\\

\bitfields{ 1/P, 2/00, 2/-- --, 27/extend27, 1/1, 2/00, 2/-- --, 27/upper27, 1/1, 2/11, 1/0, 4/1000, 6/registerW, 2/00, 10/lower10, 6/registerZ }


\paragraph{Syntax} \texttt{andd} \emph{registerW} = \emph{registerZ}, \emph{extend27\_\discretionary{}{}{}upper27\_\discretionary{}{}{}lower10}

\paragraph{Description} Bitwise and of the \emph{registerZ} with the \emph{extend27\_\discretionary{}{}{}upper27\_\discretionary{}{}{}lower10}. The result is stored into the \emph{registerW}.


\paragraph{Execution} {Reservation table \textsf{ALU\_TINY.Y} (Section~\ref{sec:reservation-ALU-TINY.Y})}
~
{\begin{lstlisting}
stage RR:
  new argument3 = _WX_64(extend27_upper27_lower10);
  new argument2 = GPR[registerZ];
  new result1 = argument2 & argument3;
stage E1:
  GPR[registerW] = result1;
\end{lstlisting}}
\newpage
\subsubsection[{\small ANDND: Bitwise And with Not Operand Double Word}]{ANDND\hfill Bitwise And with Not Operand Double Word}
\label{instr:ANDND}\index{andnd}
 
\paragraph{Encoding} Format \textsf{ALU\_DWRR0} (Section~\ref{sec:format-ALU-DWRR0}), \verb|exucode2 = 1110|\\

\bitfields{ 1/P, 2/11, 1/0, 4/1110, 6/registerW, 2/10, 3/000, 1/0, 6/registerY, 6/registerZ }


\paragraph{Syntax} \texttt{andnd} \emph{registerW} = \emph{registerZ}, \emph{registerY}

\paragraph{Description} Bitwise and of not the \emph{registerZ} with the \emph{registerY}. The result is stored into the \emph{registerW}.


\paragraph{Execution} {Reservation table \textsf{ALU\_TINY} (Section~\ref{sec:reservation-ALU-TINY})}
~
{\begin{lstlisting}
stage RR:
  new argument3 = GPR[registerY];
  new argument2 = GPR[registerZ];
  new result1 = ~argument2 & argument3;
stage E1:
  GPR[registerW] = result1;
\end{lstlisting}}
\newpage

\paragraph{Encoding} Format \textsf{ALU\_DWRR0.M} (Section~\ref{sec:format-ALU-DWRR0.M}), \verb|exucode2 = 1110|\\

\bitfields{ 1/P, 2/00, 2/-- --, 27/upper27, 1/1, 2/11, 1/0, 4/1110, 6/registerW, 2/10, 3/000, 1/0, 1/splat32, 5/lower5, 6/registerZ }


\paragraph{Syntax} \texttt{andnd} \emph{registerW} = \emph{registerZ}, \emph{upper27\_\discretionary{}{}{}lower5}\emph{splat32}

\paragraph{Description} Bitwise and of not the \emph{registerZ} with the \emph{upper27\_\discretionary{}{}{}lower5}. The result is stored into the \emph{registerW}.


\paragraph{Modifier(s)}
\noindent\begin{itemize}
\item splat32 (cf. splat32 in section \ref{par:modifier-splat32} on page \pageref{par:modifier-splat32})
\end{itemize}

\paragraph{Execution} {Reservation table \textsf{ALU\_TINY.X} (Section~\ref{sec:reservation-ALU-TINY.X})}
~
{\begin{lstlisting}
stage ID:
  new splat32 = _ZX_1(splat32);
stage RR:
  new argument3 = splat32 ? (_WX_32(upper27_lower5) << 32) | _ZX_32(_WX_32(upper27_lower5)) : _SX_32(_WX_32(upper27_lower5));
  new argument2 = GPR[registerZ];
  new result1 = ~argument2 & argument3;
stage E1:
  GPR[registerW] = result1;
\end{lstlisting}}
\newpage

\paragraph{Encoding} Format \textsf{ALU\_DWRI} (Section~\ref{sec:format-ALU-DWRI}), \verb|exucode2 = 1110|\\

\bitfields{ 1/P, 2/11, 1/0, 4/1110, 6/registerW, 2/00, 10/signed10, 6/registerZ }


\paragraph{Syntax} \texttt{andnd} \emph{registerW} = \emph{registerZ}, \emph{signed10}

\paragraph{Description} Bitwise and of not the \emph{registerZ} with the \emph{signed10}. The result is stored into the \emph{registerW}.


\paragraph{Execution} {Reservation table \textsf{ALU\_TINY} (Section~\ref{sec:reservation-ALU-TINY})}
~
{\begin{lstlisting}
stage RR:
  new argument3 = _SX_10(signed10);
  new argument2 = GPR[registerZ];
  new result1 = ~argument2 & argument3;
stage E1:
  GPR[registerW] = result1;
\end{lstlisting}}
\newpage

\paragraph{Encoding} Format \textsf{ALU\_DWRI.X} (Section~\ref{sec:format-ALU-DWRI.X}), \verb|exucode2 = 1110|\\

\bitfields{ 1/P, 2/00, 2/-- --, 27/upper27, 1/1, 2/11, 1/0, 4/1110, 6/registerW, 2/00, 10/lower10, 6/registerZ }


\paragraph{Syntax} \texttt{andnd} \emph{registerW} = \emph{registerZ}, \emph{upper27\_\discretionary{}{}{}lower10}

\paragraph{Description} Bitwise and of not the \emph{registerZ} with the \emph{upper27\_\discretionary{}{}{}lower10}. The result is stored into the \emph{registerW}.


\paragraph{Execution} {Reservation table \textsf{ALU\_TINY.X} (Section~\ref{sec:reservation-ALU-TINY.X})}
~
{\begin{lstlisting}
stage RR:
  new argument3 = _SX_37(upper27_lower10);
  new argument2 = GPR[registerZ];
  new result1 = ~argument2 & argument3;
stage E1:
  GPR[registerW] = result1;
\end{lstlisting}}
\newpage

\paragraph{Encoding} Format \textsf{ALU\_DWRI.Y} (Section~\ref{sec:format-ALU-DWRI.Y}), \verb|exucode2 = 1110|\\

\bitfields{ 1/P, 2/00, 2/-- --, 27/extend27, 1/1, 2/00, 2/-- --, 27/upper27, 1/1, 2/11, 1/0, 4/1110, 6/registerW, 2/00, 10/lower10, 6/registerZ }


\paragraph{Syntax} \texttt{andnd} \emph{registerW} = \emph{registerZ}, \emph{extend27\_\discretionary{}{}{}upper27\_\discretionary{}{}{}lower10}

\paragraph{Description} Bitwise and of not the \emph{registerZ} with the \emph{extend27\_\discretionary{}{}{}upper27\_\discretionary{}{}{}lower10}. The result is stored into the \emph{registerW}.


\paragraph{Execution} {Reservation table \textsf{ALU\_TINY.Y} (Section~\ref{sec:reservation-ALU-TINY.Y})}
~
{\begin{lstlisting}
stage RR:
  new argument3 = _WX_64(extend27_upper27_lower10);
  new argument2 = GPR[registerZ];
  new result1 = ~argument2 & argument3;
stage E1:
  GPR[registerW] = result1;
\end{lstlisting}}
\newpage
\subsubsection[{\small ANDNQ: Bitwise And with Not Operand Quadruple Word}]{ANDNQ\hfill Bitwise And with Not Operand Quadruple Word}
\label{instr:ANDNQ}\index{andnq}
 
\paragraph{Encoding} Format \textsf{ALU\_QWRR0} (Section~\ref{sec:format-ALU-QWRR0}), \verb|exucode2 = 1110|\\

\bitfields{ 1/P, 2/11, 1/0, 4/1110, 5/registerM, 1/--, 2/10, 3/000, 1/1, 5/registerO, 1/--, 5/registerP, 1/-- }


\paragraph{Syntax} \texttt{andnq} \emph{registerM} = \emph{registerP}, \emph{registerO}

\paragraph{Description} Bitwise and of not the \emph{registerP} with the \emph{registerO}. The result is stored into the \emph{registerM}.


\paragraph{Execution} {Reservation table \textsf{ALU\_LITE} (Section~\ref{sec:reservation-ALU-LITE})}
~
{\begin{lstlisting}
stage RR:
  new argument3 = PGR[registerO];
  new argument2 = PGR[registerP];
  new result1 = ~argument2 & argument3;
stage E1:
  PGR[registerM] = result1;
\end{lstlisting}}
\newpage

\paragraph{Encoding} Format \textsf{ALU\_QWRR0.M} (Section~\ref{sec:format-ALU-QWRR0.M}), \verb|exucode2 = 1110|\\

\bitfields{ 1/P, 2/00, 2/-- --, 27/upper27, 1/1, 2/11, 1/0, 4/1110, 5/registerM, 1/--, 2/10, 3/000, 1/1, 1/splat32, 5/lower5, 5/registerP, 1/1 }


\paragraph{Syntax} \texttt{andnq} \emph{registerM} = \emph{registerP}, \emph{upper27\_\discretionary{}{}{}lower5}\emph{splat32}

\paragraph{Description} Bitwise and of not the \emph{registerP} with the \emph{upper27\_\discretionary{}{}{}lower5}. The result is stored into the \emph{registerM}.


\paragraph{Modifier(s)}
\noindent\begin{itemize}
\item splat32 (cf. splat32 in section \ref{par:modifier-splat32} on page \pageref{par:modifier-splat32})
\end{itemize}

\paragraph{Execution} {Reservation table \textsf{ALU\_LITE.X} (Section~\ref{sec:reservation-ALU-LITE.X})}
~
{\begin{lstlisting}
stage ID:
  new splat32 = _ZX_1(splat32);
stage RR:
  new argument3 = splat32 ? (_WX_32(upper27_lower5) << 32) | _ZX_32(_WX_32(upper27_lower5)) : _SX_32(_WX_32(upper27_lower5));
  new argument2 = PGR[registerP];
  new result1 = ~argument2 & argument3;
stage E1:
  PGR[registerM] = result1;
\end{lstlisting}}
\newpage
\subsubsection[{\small ANDNW: Bitwise And with Not Operand Word}]{ANDNW\hfill Bitwise And with Not Operand Word}
\label{instr:ANDNW}\index{andnw}
 
\paragraph{Encoding} Format \textsf{ALU\_WWRR0} (Section~\ref{sec:format-ALU-WWRR0}), \verb|exucode2 = 1110|\\

\bitfields{ 1/P, 2/11, 1/0, 4/1110, 6/registerW, 2/11, 3/000, 1/signextw, 6/registerY, 6/registerZ }


\paragraph{Syntax} \texttt{andnw}\emph{signextw} \emph{registerW} = \emph{registerZ}, \emph{registerY}

\paragraph{Description} Bitwise and of not the \emph{registerZ} with the \emph{registerY}. The result with the 32 upper bits cleared is stored into the \emph{registerW}.


\paragraph{Modifier(s)}
\noindent\begin{itemize}
\item signextw (cf. signextw in section \ref{par:modifier-signextw} on page \pageref{par:modifier-signextw})
\end{itemize}

\paragraph{Execution} {Reservation table \textsf{ALU\_TINY} (Section~\ref{sec:reservation-ALU-TINY})}
~
{\begin{lstlisting}
stage ID:
  new signextw = _ZX_1(signextw);
stage RR:
  new argument3 = GPR[registerY];
  new argument2 = GPR[registerZ];
  new result1 = ~argument2 & argument3;
stage E1:
  GPR[registerW] = signextw ? _SX_32(result1) : _ZX_32(result1);
\end{lstlisting}}
\newpage

\paragraph{Encoding} Format \textsf{ALU\_WWRR0.W} (Section~\ref{sec:format-ALU-WWRR0.W}), \verb|exucode2 = 1110|\\

\bitfields{ 1/P, 2/00, 2/-- --, 27/upper27, 1/1, 2/11, 1/0, 4/1110, 6/registerW, 2/11, 3/000, 1/signextw, 1/0, 5/lower5, 6/registerZ }


\paragraph{Syntax} \texttt{andnw}\emph{signextw} \emph{registerW} = \emph{registerZ}, \emph{upper27\_\discretionary{}{}{}lower5}

\paragraph{Description} Bitwise and of not the \emph{registerZ} with the \emph{upper27\_\discretionary{}{}{}lower5}. The result with the 32 upper bits cleared is stored into the \emph{registerW}.


\paragraph{Modifier(s)}
\noindent\begin{itemize}
\item signextw (cf. signextw in section \ref{par:modifier-signextw} on page \pageref{par:modifier-signextw})
\end{itemize}

\paragraph{Execution} {Reservation table \textsf{ALU\_TINY.X} (Section~\ref{sec:reservation-ALU-TINY.X})}
~
{\begin{lstlisting}
stage ID:
  new signextw = _ZX_1(signextw);
stage RR:
  new argument3 = _WX_32(upper27_lower5);
  new argument2 = GPR[registerZ];
  new result1 = ~argument2 & argument3;
stage E1:
  GPR[registerW] = signextw ? _SX_32(result1) : _ZX_32(result1);
\end{lstlisting}}
\newpage
\subsubsection[{\small ANDQ: Bitwise And Quadruple Word}]{ANDQ\hfill Bitwise And Quadruple Word}
\label{instr:ANDQ}\index{andq}
 
\paragraph{Encoding} Format \textsf{ALU\_QWRR0} (Section~\ref{sec:format-ALU-QWRR0}), \verb|exucode2 = 1000|\\

\bitfields{ 1/P, 2/11, 1/0, 4/1000, 5/registerM, 1/--, 2/10, 3/000, 1/1, 5/registerO, 1/--, 5/registerP, 1/-- }


\paragraph{Syntax} \texttt{andq} \emph{registerM} = \emph{registerP}, \emph{registerO}

\paragraph{Description} Bitwise and of the \emph{registerP} with the \emph{registerO}. The result is stored into the \emph{registerM}.


\paragraph{Execution} {Reservation table \textsf{ALU\_LITE} (Section~\ref{sec:reservation-ALU-LITE})}
~
{\begin{lstlisting}
stage RR:
  new argument3 = PGR[registerO];
  new argument2 = PGR[registerP];
  new result1 = argument2 & argument3;
stage E1:
  PGR[registerM] = result1;
\end{lstlisting}}
\newpage

\paragraph{Encoding} Format \textsf{ALU\_QWRR0.M} (Section~\ref{sec:format-ALU-QWRR0.M}), \verb|exucode2 = 1000|\\

\bitfields{ 1/P, 2/00, 2/-- --, 27/upper27, 1/1, 2/11, 1/0, 4/1000, 5/registerM, 1/--, 2/10, 3/000, 1/1, 1/splat32, 5/lower5, 5/registerP, 1/1 }


\paragraph{Syntax} \texttt{andq} \emph{registerM} = \emph{registerP}, \emph{upper27\_\discretionary{}{}{}lower5}\emph{splat32}

\paragraph{Description} Bitwise and of the \emph{registerP} with the \emph{upper27\_\discretionary{}{}{}lower5}. The result is stored into the \emph{registerM}.


\paragraph{Modifier(s)}
\noindent\begin{itemize}
\item splat32 (cf. splat32 in section \ref{par:modifier-splat32} on page \pageref{par:modifier-splat32})
\end{itemize}

\paragraph{Execution} {Reservation table \textsf{ALU\_LITE.X} (Section~\ref{sec:reservation-ALU-LITE.X})}
~
{\begin{lstlisting}
stage ID:
  new splat32 = _ZX_1(splat32);
stage RR:
  new argument3 = splat32 ? (_WX_32(upper27_lower5) << 32) | _ZX_32(_WX_32(upper27_lower5)) : _SX_32(_WX_32(upper27_lower5));
  new argument2 = PGR[registerP];
  new result1 = argument2 & argument3;
stage E1:
  PGR[registerM] = result1;
\end{lstlisting}}
\newpage
\subsubsection[{\small ANDW: Bitwise And Word}]{ANDW\hfill Bitwise And Word}
\label{instr:ANDW}\index{andw}
 
\paragraph{Encoding} Format \textsf{ALU\_WWRR0} (Section~\ref{sec:format-ALU-WWRR0}), \verb|exucode2 = 1000|\\

\bitfields{ 1/P, 2/11, 1/0, 4/1000, 6/registerW, 2/11, 3/000, 1/signextw, 6/registerY, 6/registerZ }


\paragraph{Syntax} \texttt{andw}\emph{signextw} \emph{registerW} = \emph{registerZ}, \emph{registerY}

\paragraph{Description} Bitwise and of the \emph{registerZ} with the \emph{registerY}. The result with the 32 upper bits cleared is stored into the \emph{registerW}.


\paragraph{Modifier(s)}
\noindent\begin{itemize}
\item signextw (cf. signextw in section \ref{par:modifier-signextw} on page \pageref{par:modifier-signextw})
\end{itemize}

\paragraph{Execution} {Reservation table \textsf{ALU\_TINY} (Section~\ref{sec:reservation-ALU-TINY})}
~
{\begin{lstlisting}
stage ID:
  new signextw = _ZX_1(signextw);
stage RR:
  new argument3 = GPR[registerY];
  new argument2 = GPR[registerZ];
  new result1 = argument2 & argument3;
stage E1:
  GPR[registerW] = signextw ? _SX_32(result1) : _ZX_32(result1);
\end{lstlisting}}
\newpage

\paragraph{Encoding} Format \textsf{ALU\_WWRR0.W} (Section~\ref{sec:format-ALU-WWRR0.W}), \verb|exucode2 = 1000|\\

\bitfields{ 1/P, 2/00, 2/-- --, 27/upper27, 1/1, 2/11, 1/0, 4/1000, 6/registerW, 2/11, 3/000, 1/signextw, 1/0, 5/lower5, 6/registerZ }


\paragraph{Syntax} \texttt{andw}\emph{signextw} \emph{registerW} = \emph{registerZ}, \emph{upper27\_\discretionary{}{}{}lower5}

\paragraph{Description} Bitwise and of the \emph{registerZ} with the \emph{upper27\_\discretionary{}{}{}lower5}. The result with the 32 upper bits cleared is stored into the \emph{registerW}.


\paragraph{Modifier(s)}
\noindent\begin{itemize}
\item signextw (cf. signextw in section \ref{par:modifier-signextw} on page \pageref{par:modifier-signextw})
\end{itemize}

\paragraph{Execution} {Reservation table \textsf{ALU\_TINY.X} (Section~\ref{sec:reservation-ALU-TINY.X})}
~
{\begin{lstlisting}
stage ID:
  new signextw = _ZX_1(signextw);
stage RR:
  new argument3 = _WX_32(upper27_lower5);
  new argument2 = GPR[registerZ];
  new result1 = argument2 & argument3;
stage E1:
  GPR[registerW] = signextw ? _SX_32(result1) : _ZX_32(result1);
\end{lstlisting}}
\newpage
\subsubsection[{\small AVGBO: Average Byte Octuple}]{AVGBO\hfill Average Byte Octuple}
\label{instr:AVGBO}\index{avgbo}
 
\paragraph{Encoding} Format \textsf{ALU\_BOAVGE} (Section~\ref{sec:format-ALU-BOAVGE}), \verb|exucode2 = 1100|\\

\bitfields{ 1/P, 2/11, 1/1, 4/1100, 5/registerWe, 1/0, 2/10, 3/001, 1/1, 5/registerYe, 1/--, 5/registerZe, 1/1 }


\paragraph{Syntax} \texttt{avgbo} \emph{registerWe} = \emph{registerZe}, \emph{registerYe}

\paragraph{Description} The \emph{registerYe} and the \emph{registerZe} are interpreted as eight bytes packed into 64 bits. The \emph{registerYe} bytes and the \emph{registerZe} bytes are added. Each sum is shifted right arithmetic one bit. The eight resulting bytes are packed and stored into the \emph{registerWe}.


\paragraph{Execution} {Reservation table \textsf{ALU\_LITE} (Section~\ref{sec:reservation-ALU-LITE})}
~
{\begin{lstlisting}
stage RR:
  new argument3 = GPR[registerYe<<1];
  new argument2 = GPR[registerZe<<1];
  for (I = 0; I < 8; I += 1) {
    result1.8[I] = (_SX_8(argument3.8[I]) + _SX_8(argument2.8[I])) >> 1
  };
stage E1:
  GPR[registerWe<<1] = result1;
\end{lstlisting}}
\newpage

\paragraph{Encoding} Format \textsf{ALU\_BOAVGO} (Section~\ref{sec:format-ALU-BOAVGO}), \verb|exucode2 = 1100|\\

\bitfields{ 1/P, 2/11, 1/1, 4/1100, 5/registerWo, 1/1, 2/10, 3/001, 1/1, 5/registerYo, 1/--, 5/registerZo, 1/1 }


\paragraph{Syntax} \texttt{avgbo} \emph{registerWo} = \emph{registerZo}, \emph{registerYo}

\paragraph{Description} The \emph{registerYo} and the \emph{registerZo} are interpreted as eight bytes packed into 64 bits. The \emph{registerYo} bytes and the \emph{registerZo} bytes are added. Each sum is shifted right arithmetic one bit. The eight resulting bytes are packed and stored into the \emph{registerWo}.


\paragraph{Execution} {Reservation table \textsf{ALU\_LITE} (Section~\ref{sec:reservation-ALU-LITE})}
~
{\begin{lstlisting}
stage RR:
  new argument3 = GPR[(registerYo<<1)+1];
  new argument2 = GPR[(registerZo<<1)+1];
  for (I = 0; I < 8; I += 1) {
    result1.8[I] = (_SX_8(argument3.8[I]) + _SX_8(argument2.8[I])) >> 1
  };
stage E1:
  GPR[(registerWo<<1)+1] = result1;
\end{lstlisting}}
\newpage

\paragraph{Encoding} Format \textsf{ALU\_BOAVGE.M} (Section~\ref{sec:format-ALU-BOAVGE.M}), \verb|exucode2 = 1100|\\

\bitfields{ 1/P, 2/00, 2/-- --, 27/upper27, 1/1, 2/11, 1/1, 4/1100, 5/registerWe, 1/0, 2/10, 3/001, 1/1, 1/splat32, 5/lower5, 5/registerZe, 1/1 }


\paragraph{Syntax} \texttt{avgbo} \emph{registerWe} = \emph{registerZe}, \emph{upper27\_\discretionary{}{}{}lower5}\emph{splat32}

\paragraph{Description} The \emph{upper27\_\discretionary{}{}{}lower5} and the \emph{registerZe} are interpreted as eight bytes packed into 64 bits. The \emph{upper27\_\discretionary{}{}{}lower5} bytes and the \emph{registerZe} bytes are added. Each sum is shifted right arithmetic one bit. The eight resulting bytes are packed and stored into the \emph{registerWe}.


\paragraph{Modifier(s)}
\noindent\begin{itemize}
\item splat32 (cf. splat32 in section \ref{par:modifier-splat32} on page \pageref{par:modifier-splat32})
\end{itemize}

\paragraph{Execution} {Reservation table \textsf{ALU\_LITE.X} (Section~\ref{sec:reservation-ALU-LITE.X})}
~
{\begin{lstlisting}
stage ID:
  new splat32 = _ZX_1(splat32);
stage RR:
  new argument3 = splat32 ? (_WX_32(upper27_lower5) << 32) | _ZX_32(_WX_32(upper27_lower5)) : _SX_32(_WX_32(upper27_lower5));
  new argument2 = GPR[registerZe<<1];
  for (I = 0; I < 8; I += 1) {
    result1.8[I] = (_SX_8(argument3.8[I]) + _SX_8(argument2.8[I])) >> 1
  };
stage E1:
  GPR[registerWe<<1] = result1;
\end{lstlisting}}
\newpage

\paragraph{Encoding} Format \textsf{ALU\_BOAVGO.M} (Section~\ref{sec:format-ALU-BOAVGO.M}), \verb|exucode2 = 1100|\\

\bitfields{ 1/P, 2/00, 2/-- --, 27/upper27, 1/1, 2/11, 1/1, 4/1100, 5/registerWo, 1/1, 2/10, 3/001, 1/1, 1/splat32, 5/lower5, 5/registerZo, 1/1 }


\paragraph{Syntax} \texttt{avgbo} \emph{registerWo} = \emph{registerZo}, \emph{upper27\_\discretionary{}{}{}lower5}\emph{splat32}

\paragraph{Description} The \emph{upper27\_\discretionary{}{}{}lower5} and the \emph{registerZo} are interpreted as eight bytes packed into 64 bits. The \emph{upper27\_\discretionary{}{}{}lower5} bytes and the \emph{registerZo} bytes are added. Each sum is shifted right arithmetic one bit. The eight resulting bytes are packed and stored into the \emph{registerWo}.


\paragraph{Modifier(s)}
\noindent\begin{itemize}
\item splat32 (cf. splat32 in section \ref{par:modifier-splat32} on page \pageref{par:modifier-splat32})
\end{itemize}

\paragraph{Execution} {Reservation table \textsf{ALU\_LITE.X} (Section~\ref{sec:reservation-ALU-LITE.X})}
~
{\begin{lstlisting}
stage ID:
  new splat32 = _ZX_1(splat32);
stage RR:
  new argument3 = splat32 ? (_WX_32(upper27_lower5) << 32) | _ZX_32(_WX_32(upper27_lower5)) : _SX_32(_WX_32(upper27_lower5));
  new argument2 = GPR[(registerZo<<1)+1];
  for (I = 0; I < 8; I += 1) {
    result1.8[I] = (_SX_8(argument3.8[I]) + _SX_8(argument2.8[I])) >> 1
  };
stage E1:
  GPR[(registerWo<<1)+1] = result1;
\end{lstlisting}}
\newpage
\subsubsection[{\small AVGHQ: Average Half Word Quadruple}]{AVGHQ\hfill Average Half Word Quadruple}
\label{instr:AVGHQ}\index{avghq}
 
\paragraph{Encoding} Format \textsf{ALU\_HQAVGE} (Section~\ref{sec:format-ALU-HQAVGE}), \verb|exucode2 = 1100|\\

\bitfields{ 1/P, 2/11, 1/1, 4/1100, 5/registerWe, 1/0, 2/10, 3/001, 1/1, 5/registerYe, 1/--, 5/registerZe, 1/0 }


\paragraph{Syntax} \texttt{avghq} \emph{registerWe} = \emph{registerZe}, \emph{registerYe}

\paragraph{Description} The \emph{registerYe} and the \emph{registerZe} are interpreted as four half words packed into 64 bits. The \emph{registerYe} half words and the \emph{registerZe} half words are added. Each sum is shifted right arithmetic one bit. The four resulting half words are packed and stored into the \emph{registerWe}.


\paragraph{Execution} {Reservation table \textsf{ALU\_LITE} (Section~\ref{sec:reservation-ALU-LITE})}
~
{\begin{lstlisting}
stage RR:
  new argument3 = GPR[registerYe<<1];
  new argument2 = GPR[registerZe<<1];
  for (I = 0; I < 4; I += 1) {
    result1.16[I] = (_SX_16(argument3.16[I]) + _SX_16(argument2.16[I])) >> 1
  };
stage E1:
  GPR[registerWe<<1] = result1;
\end{lstlisting}}
\newpage

\paragraph{Encoding} Format \textsf{ALU\_HQAVGO} (Section~\ref{sec:format-ALU-HQAVGO}), \verb|exucode2 = 1100|\\

\bitfields{ 1/P, 2/11, 1/1, 4/1100, 5/registerWo, 1/1, 2/10, 3/001, 1/1, 5/registerYo, 1/--, 5/registerZo, 1/0 }


\paragraph{Syntax} \texttt{avghq} \emph{registerWo} = \emph{registerZo}, \emph{registerYo}

\paragraph{Description} The \emph{registerYo} and the \emph{registerZo} are interpreted as four half words packed into 64 bits. The \emph{registerYo} half words and the \emph{registerZo} half words are added. Each sum is shifted right arithmetic one bit. The four resulting half words are packed and stored into the \emph{registerWo}.


\paragraph{Execution} {Reservation table \textsf{ALU\_LITE} (Section~\ref{sec:reservation-ALU-LITE})}
~
{\begin{lstlisting}
stage RR:
  new argument3 = GPR[(registerYo<<1)+1];
  new argument2 = GPR[(registerZo<<1)+1];
  for (I = 0; I < 4; I += 1) {
    result1.16[I] = (_SX_16(argument3.16[I]) + _SX_16(argument2.16[I])) >> 1
  };
stage E1:
  GPR[(registerWo<<1)+1] = result1;
\end{lstlisting}}
\newpage

\paragraph{Encoding} Format \textsf{ALU\_HQAVGE.M} (Section~\ref{sec:format-ALU-HQAVGE.M}), \verb|exucode2 = 1100|\\

\bitfields{ 1/P, 2/00, 2/-- --, 27/upper27, 1/1, 2/11, 1/1, 4/1100, 5/registerWe, 1/0, 2/10, 3/001, 1/1, 1/splat32, 5/lower5, 5/registerZe, 1/0 }


\paragraph{Syntax} \texttt{avghq} \emph{registerWe} = \emph{registerZe}, \emph{upper27\_\discretionary{}{}{}lower5}\emph{splat32}

\paragraph{Description} The \emph{upper27\_\discretionary{}{}{}lower5} and the \emph{registerZe} are interpreted as four half words packed into 64 bits. The \emph{upper27\_\discretionary{}{}{}lower5} half words and the \emph{registerZe} half words are added. Each sum is shifted right arithmetic one bit. The four resulting half words are packed and stored into the \emph{registerWe}.


\paragraph{Modifier(s)}
\noindent\begin{itemize}
\item splat32 (cf. splat32 in section \ref{par:modifier-splat32} on page \pageref{par:modifier-splat32})
\end{itemize}

\paragraph{Execution} {Reservation table \textsf{ALU\_LITE.X} (Section~\ref{sec:reservation-ALU-LITE.X})}
~
{\begin{lstlisting}
stage ID:
  new splat32 = _ZX_1(splat32);
stage RR:
  new argument3 = splat32 ? (_WX_32(upper27_lower5) << 32) | _ZX_32(_WX_32(upper27_lower5)) : _SX_32(_WX_32(upper27_lower5));
  new argument2 = GPR[registerZe<<1];
  for (I = 0; I < 4; I += 1) {
    result1.16[I] = (_SX_16(argument3.16[I]) + _SX_16(argument2.16[I])) >> 1
  };
stage E1:
  GPR[registerWe<<1] = result1;
\end{lstlisting}}
\newpage

\paragraph{Encoding} Format \textsf{ALU\_HQAVGO.M} (Section~\ref{sec:format-ALU-HQAVGO.M}), \verb|exucode2 = 1100|\\

\bitfields{ 1/P, 2/00, 2/-- --, 27/upper27, 1/1, 2/11, 1/1, 4/1100, 5/registerWo, 1/1, 2/10, 3/001, 1/1, 1/splat32, 5/lower5, 5/registerZo, 1/0 }


\paragraph{Syntax} \texttt{avghq} \emph{registerWo} = \emph{registerZo}, \emph{upper27\_\discretionary{}{}{}lower5}\emph{splat32}

\paragraph{Description} The \emph{upper27\_\discretionary{}{}{}lower5} and the \emph{registerZo} are interpreted as four half words packed into 64 bits. The \emph{upper27\_\discretionary{}{}{}lower5} half words and the \emph{registerZo} half words are added. Each sum is shifted right arithmetic one bit. The four resulting half words are packed and stored into the \emph{registerWo}.


\paragraph{Modifier(s)}
\noindent\begin{itemize}
\item splat32 (cf. splat32 in section \ref{par:modifier-splat32} on page \pageref{par:modifier-splat32})
\end{itemize}

\paragraph{Execution} {Reservation table \textsf{ALU\_LITE.X} (Section~\ref{sec:reservation-ALU-LITE.X})}
~
{\begin{lstlisting}
stage ID:
  new splat32 = _ZX_1(splat32);
stage RR:
  new argument3 = splat32 ? (_WX_32(upper27_lower5) << 32) | _ZX_32(_WX_32(upper27_lower5)) : _SX_32(_WX_32(upper27_lower5));
  new argument2 = GPR[(registerZo<<1)+1];
  for (I = 0; I < 4; I += 1) {
    result1.16[I] = (_SX_16(argument3.16[I]) + _SX_16(argument2.16[I])) >> 1
  };
stage E1:
  GPR[(registerWo<<1)+1] = result1;
\end{lstlisting}}
\newpage
\subsubsection[{\small AVGRBO: Average Rounded Byte Octuple}]{AVGRBO\hfill Average Rounded Byte Octuple}
\label{instr:AVGRBO}\index{avgrbo}
 
\paragraph{Encoding} Format \textsf{ALU\_BOAVGE} (Section~\ref{sec:format-ALU-BOAVGE}), \verb|exucode2 = 1110|\\

\bitfields{ 1/P, 2/11, 1/1, 4/1110, 5/registerWe, 1/0, 2/10, 3/001, 1/1, 5/registerYe, 1/--, 5/registerZe, 1/1 }


\paragraph{Syntax} \texttt{avgrbo} \emph{registerWe} = \emph{registerZe}, \emph{registerYe}

\paragraph{Description} The \emph{registerYe} and the \emph{registerZe} are interpreted as eight bytes packed into 64 bits. The \emph{registerYe} bytes and the \emph{registerZe} bytes are added. Each sum is incremented and shifted right arithmetic one bit. The eight resulting bytes are packed and stored into the \emph{registerWe}.


\paragraph{Execution} {Reservation table \textsf{ALU\_LITE} (Section~\ref{sec:reservation-ALU-LITE})}
~
{\begin{lstlisting}
stage RR:
  new argument3 = GPR[registerYe<<1];
  new argument2 = GPR[registerZe<<1];
  for (I = 0; I < 8; I += 1) {
    result1.8[I] = ((_SX_8(argument3.8[I]) + _SX_8(argument2.8[I])) + 1) >> 1
  };
stage E1:
  GPR[registerWe<<1] = result1;
\end{lstlisting}}
\newpage

\paragraph{Encoding} Format \textsf{ALU\_BOAVGO} (Section~\ref{sec:format-ALU-BOAVGO}), \verb|exucode2 = 1110|\\

\bitfields{ 1/P, 2/11, 1/1, 4/1110, 5/registerWo, 1/1, 2/10, 3/001, 1/1, 5/registerYo, 1/--, 5/registerZo, 1/1 }


\paragraph{Syntax} \texttt{avgrbo} \emph{registerWo} = \emph{registerZo}, \emph{registerYo}

\paragraph{Description} The \emph{registerYo} and the \emph{registerZo} are interpreted as eight bytes packed into 64 bits. The \emph{registerYo} bytes and the \emph{registerZo} bytes are added. Each sum is incremented and shifted right arithmetic one bit. The eight resulting bytes are packed and stored into the \emph{registerWo}.


\paragraph{Execution} {Reservation table \textsf{ALU\_LITE} (Section~\ref{sec:reservation-ALU-LITE})}
~
{\begin{lstlisting}
stage RR:
  new argument3 = GPR[(registerYo<<1)+1];
  new argument2 = GPR[(registerZo<<1)+1];
  for (I = 0; I < 8; I += 1) {
    result1.8[I] = ((_SX_8(argument3.8[I]) + _SX_8(argument2.8[I])) + 1) >> 1
  };
stage E1:
  GPR[(registerWo<<1)+1] = result1;
\end{lstlisting}}
\newpage

\paragraph{Encoding} Format \textsf{ALU\_BOAVGE.M} (Section~\ref{sec:format-ALU-BOAVGE.M}), \verb|exucode2 = 1110|\\

\bitfields{ 1/P, 2/00, 2/-- --, 27/upper27, 1/1, 2/11, 1/1, 4/1110, 5/registerWe, 1/0, 2/10, 3/001, 1/1, 1/splat32, 5/lower5, 5/registerZe, 1/1 }


\paragraph{Syntax} \texttt{avgrbo} \emph{registerWe} = \emph{registerZe}, \emph{upper27\_\discretionary{}{}{}lower5}\emph{splat32}

\paragraph{Description} The \emph{upper27\_\discretionary{}{}{}lower5} and the \emph{registerZe} are interpreted as eight bytes packed into 64 bits. The \emph{upper27\_\discretionary{}{}{}lower5} bytes and the \emph{registerZe} bytes are added. Each sum is incremented and shifted right arithmetic one bit. The eight resulting bytes are packed and stored into the \emph{registerWe}.


\paragraph{Modifier(s)}
\noindent\begin{itemize}
\item splat32 (cf. splat32 in section \ref{par:modifier-splat32} on page \pageref{par:modifier-splat32})
\end{itemize}

\paragraph{Execution} {Reservation table \textsf{ALU\_LITE.X} (Section~\ref{sec:reservation-ALU-LITE.X})}
~
{\begin{lstlisting}
stage ID:
  new splat32 = _ZX_1(splat32);
stage RR:
  new argument3 = splat32 ? (_WX_32(upper27_lower5) << 32) | _ZX_32(_WX_32(upper27_lower5)) : _SX_32(_WX_32(upper27_lower5));
  new argument2 = GPR[registerZe<<1];
  for (I = 0; I < 8; I += 1) {
    result1.8[I] = ((_SX_8(argument3.8[I]) + _SX_8(argument2.8[I])) + 1) >> 1
  };
stage E1:
  GPR[registerWe<<1] = result1;
\end{lstlisting}}
\newpage

\paragraph{Encoding} Format \textsf{ALU\_BOAVGO.M} (Section~\ref{sec:format-ALU-BOAVGO.M}), \verb|exucode2 = 1110|\\

\bitfields{ 1/P, 2/00, 2/-- --, 27/upper27, 1/1, 2/11, 1/1, 4/1110, 5/registerWo, 1/1, 2/10, 3/001, 1/1, 1/splat32, 5/lower5, 5/registerZo, 1/1 }


\paragraph{Syntax} \texttt{avgrbo} \emph{registerWo} = \emph{registerZo}, \emph{upper27\_\discretionary{}{}{}lower5}\emph{splat32}

\paragraph{Description} The \emph{upper27\_\discretionary{}{}{}lower5} and the \emph{registerZo} are interpreted as eight bytes packed into 64 bits. The \emph{upper27\_\discretionary{}{}{}lower5} bytes and the \emph{registerZo} bytes are added. Each sum is incremented and shifted right arithmetic one bit. The eight resulting bytes are packed and stored into the \emph{registerWo}.


\paragraph{Modifier(s)}
\noindent\begin{itemize}
\item splat32 (cf. splat32 in section \ref{par:modifier-splat32} on page \pageref{par:modifier-splat32})
\end{itemize}

\paragraph{Execution} {Reservation table \textsf{ALU\_LITE.X} (Section~\ref{sec:reservation-ALU-LITE.X})}
~
{\begin{lstlisting}
stage ID:
  new splat32 = _ZX_1(splat32);
stage RR:
  new argument3 = splat32 ? (_WX_32(upper27_lower5) << 32) | _ZX_32(_WX_32(upper27_lower5)) : _SX_32(_WX_32(upper27_lower5));
  new argument2 = GPR[(registerZo<<1)+1];
  for (I = 0; I < 8; I += 1) {
    result1.8[I] = ((_SX_8(argument3.8[I]) + _SX_8(argument2.8[I])) + 1) >> 1
  };
stage E1:
  GPR[(registerWo<<1)+1] = result1;
\end{lstlisting}}
\newpage
\subsubsection[{\small AVGRHQ: Average Rounded Half Word Quadruple}]{AVGRHQ\hfill Average Rounded Half Word Quadruple}
\label{instr:AVGRHQ}\index{avgrhq}
 
\paragraph{Encoding} Format \textsf{ALU\_HQAVGE} (Section~\ref{sec:format-ALU-HQAVGE}), \verb|exucode2 = 1110|\\

\bitfields{ 1/P, 2/11, 1/1, 4/1110, 5/registerWe, 1/0, 2/10, 3/001, 1/1, 5/registerYe, 1/--, 5/registerZe, 1/0 }


\paragraph{Syntax} \texttt{avgrhq} \emph{registerWe} = \emph{registerZe}, \emph{registerYe}

\paragraph{Description} The \emph{registerYe} and the \emph{registerZe} are interpreted as four half words packed into 64 bits. The \emph{registerYe} half words and the \emph{registerZe} half words are added. Each sum is incremented and shifted right arithmetic one bit. The four resulting half words are packed and stored into the \emph{registerWe}.


\paragraph{Execution} {Reservation table \textsf{ALU\_LITE} (Section~\ref{sec:reservation-ALU-LITE})}
~
{\begin{lstlisting}
stage RR:
  new argument3 = GPR[registerYe<<1];
  new argument2 = GPR[registerZe<<1];
  for (I = 0; I < 4; I += 1) {
    result1.16[I] = ((_SX_16(argument3.16[I]) + _SX_16(argument2.16[I])) + 1) >> 1
  };
stage E1:
  GPR[registerWe<<1] = result1;
\end{lstlisting}}
\newpage

\paragraph{Encoding} Format \textsf{ALU\_HQAVGO} (Section~\ref{sec:format-ALU-HQAVGO}), \verb|exucode2 = 1110|\\

\bitfields{ 1/P, 2/11, 1/1, 4/1110, 5/registerWo, 1/1, 2/10, 3/001, 1/1, 5/registerYo, 1/--, 5/registerZo, 1/0 }


\paragraph{Syntax} \texttt{avgrhq} \emph{registerWo} = \emph{registerZo}, \emph{registerYo}

\paragraph{Description} The \emph{registerYo} and the \emph{registerZo} are interpreted as four half words packed into 64 bits. The \emph{registerYo} half words and the \emph{registerZo} half words are added. Each sum is incremented and shifted right arithmetic one bit. The four resulting half words are packed and stored into the \emph{registerWo}.


\paragraph{Execution} {Reservation table \textsf{ALU\_LITE} (Section~\ref{sec:reservation-ALU-LITE})}
~
{\begin{lstlisting}
stage RR:
  new argument3 = GPR[(registerYo<<1)+1];
  new argument2 = GPR[(registerZo<<1)+1];
  for (I = 0; I < 4; I += 1) {
    result1.16[I] = ((_SX_16(argument3.16[I]) + _SX_16(argument2.16[I])) + 1) >> 1
  };
stage E1:
  GPR[(registerWo<<1)+1] = result1;
\end{lstlisting}}
\newpage

\paragraph{Encoding} Format \textsf{ALU\_HQAVGE.M} (Section~\ref{sec:format-ALU-HQAVGE.M}), \verb|exucode2 = 1110|\\

\bitfields{ 1/P, 2/00, 2/-- --, 27/upper27, 1/1, 2/11, 1/1, 4/1110, 5/registerWe, 1/0, 2/10, 3/001, 1/1, 1/splat32, 5/lower5, 5/registerZe, 1/0 }


\paragraph{Syntax} \texttt{avgrhq} \emph{registerWe} = \emph{registerZe}, \emph{upper27\_\discretionary{}{}{}lower5}\emph{splat32}

\paragraph{Description} The \emph{upper27\_\discretionary{}{}{}lower5} and the \emph{registerZe} are interpreted as four half words packed into 64 bits. The \emph{upper27\_\discretionary{}{}{}lower5} half words and the \emph{registerZe} half words are added. Each sum is incremented and shifted right arithmetic one bit. The four resulting half words are packed and stored into the \emph{registerWe}.


\paragraph{Modifier(s)}
\noindent\begin{itemize}
\item splat32 (cf. splat32 in section \ref{par:modifier-splat32} on page \pageref{par:modifier-splat32})
\end{itemize}

\paragraph{Execution} {Reservation table \textsf{ALU\_LITE.X} (Section~\ref{sec:reservation-ALU-LITE.X})}
~
{\begin{lstlisting}
stage ID:
  new splat32 = _ZX_1(splat32);
stage RR:
  new argument3 = splat32 ? (_WX_32(upper27_lower5) << 32) | _ZX_32(_WX_32(upper27_lower5)) : _SX_32(_WX_32(upper27_lower5));
  new argument2 = GPR[registerZe<<1];
  for (I = 0; I < 4; I += 1) {
    result1.16[I] = ((_SX_16(argument3.16[I]) + _SX_16(argument2.16[I])) + 1) >> 1
  };
stage E1:
  GPR[registerWe<<1] = result1;
\end{lstlisting}}
\newpage

\paragraph{Encoding} Format \textsf{ALU\_HQAVGO.M} (Section~\ref{sec:format-ALU-HQAVGO.M}), \verb|exucode2 = 1110|\\

\bitfields{ 1/P, 2/00, 2/-- --, 27/upper27, 1/1, 2/11, 1/1, 4/1110, 5/registerWo, 1/1, 2/10, 3/001, 1/1, 1/splat32, 5/lower5, 5/registerZo, 1/0 }


\paragraph{Syntax} \texttt{avgrhq} \emph{registerWo} = \emph{registerZo}, \emph{upper27\_\discretionary{}{}{}lower5}\emph{splat32}

\paragraph{Description} The \emph{upper27\_\discretionary{}{}{}lower5} and the \emph{registerZo} are interpreted as four half words packed into 64 bits. The \emph{upper27\_\discretionary{}{}{}lower5} half words and the \emph{registerZo} half words are added. Each sum is incremented and shifted right arithmetic one bit. The four resulting half words are packed and stored into the \emph{registerWo}.


\paragraph{Modifier(s)}
\noindent\begin{itemize}
\item splat32 (cf. splat32 in section \ref{par:modifier-splat32} on page \pageref{par:modifier-splat32})
\end{itemize}

\paragraph{Execution} {Reservation table \textsf{ALU\_LITE.X} (Section~\ref{sec:reservation-ALU-LITE.X})}
~
{\begin{lstlisting}
stage ID:
  new splat32 = _ZX_1(splat32);
stage RR:
  new argument3 = splat32 ? (_WX_32(upper27_lower5) << 32) | _ZX_32(_WX_32(upper27_lower5)) : _SX_32(_WX_32(upper27_lower5));
  new argument2 = GPR[(registerZo<<1)+1];
  for (I = 0; I < 4; I += 1) {
    result1.16[I] = ((_SX_16(argument3.16[I]) + _SX_16(argument2.16[I])) + 1) >> 1
  };
stage E1:
  GPR[(registerWo<<1)+1] = result1;
\end{lstlisting}}
\newpage
\subsubsection[{\small AVGRUBO: Average Rounded Unsigned Byte Octuple}]{AVGRUBO\hfill Average Rounded Unsigned Byte Octuple}
\label{instr:AVGRUBO}\index{avgrubo}
 
\paragraph{Encoding} Format \textsf{ALU\_BOAVGE} (Section~\ref{sec:format-ALU-BOAVGE}), \verb|exucode2 = 1111|\\

\bitfields{ 1/P, 2/11, 1/1, 4/1111, 5/registerWe, 1/0, 2/10, 3/001, 1/1, 5/registerYe, 1/--, 5/registerZe, 1/1 }


\paragraph{Syntax} \texttt{avgrubo} \emph{registerWe} = \emph{registerZe}, \emph{registerYe}

\paragraph{Description} The \emph{registerYe} and the \emph{registerZe} are interpreted as eight bytes packed into 64 bits. The \emph{registerYe} bytes and the \emph{registerZe} bytes are added. Each sum is incremented and shifted right logical one bit. The eight resulting bytes are packed and stored into the \emph{registerWe}.


\paragraph{Execution} {Reservation table \textsf{ALU\_LITE} (Section~\ref{sec:reservation-ALU-LITE})}
~
{\begin{lstlisting}
stage RR:
  new argument3 = GPR[registerYe<<1];
  new argument2 = GPR[registerZe<<1];
  for (I = 0; I < 8; I += 1) {
    result1.8[I] = ((_ZX_8(argument3.8[I]) + _ZX_8(argument2.8[I])) + 1) >> 1
  };
stage E1:
  GPR[registerWe<<1] = result1;
\end{lstlisting}}
\newpage

\paragraph{Encoding} Format \textsf{ALU\_BOAVGO} (Section~\ref{sec:format-ALU-BOAVGO}), \verb|exucode2 = 1111|\\

\bitfields{ 1/P, 2/11, 1/1, 4/1111, 5/registerWo, 1/1, 2/10, 3/001, 1/1, 5/registerYo, 1/--, 5/registerZo, 1/1 }


\paragraph{Syntax} \texttt{avgrubo} \emph{registerWo} = \emph{registerZo}, \emph{registerYo}

\paragraph{Description} The \emph{registerYo} and the \emph{registerZo} are interpreted as eight bytes packed into 64 bits. The \emph{registerYo} bytes and the \emph{registerZo} bytes are added. Each sum is incremented and shifted right logical one bit. The eight resulting bytes are packed and stored into the \emph{registerWo}.


\paragraph{Execution} {Reservation table \textsf{ALU\_LITE} (Section~\ref{sec:reservation-ALU-LITE})}
~
{\begin{lstlisting}
stage RR:
  new argument3 = GPR[(registerYo<<1)+1];
  new argument2 = GPR[(registerZo<<1)+1];
  for (I = 0; I < 8; I += 1) {
    result1.8[I] = ((_ZX_8(argument3.8[I]) + _ZX_8(argument2.8[I])) + 1) >> 1
  };
stage E1:
  GPR[(registerWo<<1)+1] = result1;
\end{lstlisting}}
\newpage

\paragraph{Encoding} Format \textsf{ALU\_BOAVGE.M} (Section~\ref{sec:format-ALU-BOAVGE.M}), \verb|exucode2 = 1111|\\

\bitfields{ 1/P, 2/00, 2/-- --, 27/upper27, 1/1, 2/11, 1/1, 4/1111, 5/registerWe, 1/0, 2/10, 3/001, 1/1, 1/splat32, 5/lower5, 5/registerZe, 1/1 }


\paragraph{Syntax} \texttt{avgrubo} \emph{registerWe} = \emph{registerZe}, \emph{upper27\_\discretionary{}{}{}lower5}\emph{splat32}

\paragraph{Description} The \emph{upper27\_\discretionary{}{}{}lower5} and the \emph{registerZe} are interpreted as eight bytes packed into 64 bits. The \emph{upper27\_\discretionary{}{}{}lower5} bytes and the \emph{registerZe} bytes are added. Each sum is incremented and shifted right logical one bit. The eight resulting bytes are packed and stored into the \emph{registerWe}.


\paragraph{Modifier(s)}
\noindent\begin{itemize}
\item splat32 (cf. splat32 in section \ref{par:modifier-splat32} on page \pageref{par:modifier-splat32})
\end{itemize}

\paragraph{Execution} {Reservation table \textsf{ALU\_LITE.X} (Section~\ref{sec:reservation-ALU-LITE.X})}
~
{\begin{lstlisting}
stage ID:
  new splat32 = _ZX_1(splat32);
stage RR:
  new argument3 = splat32 ? (_WX_32(upper27_lower5) << 32) | _ZX_32(_WX_32(upper27_lower5)) : _SX_32(_WX_32(upper27_lower5));
  new argument2 = GPR[registerZe<<1];
  for (I = 0; I < 8; I += 1) {
    result1.8[I] = ((_ZX_8(argument3.8[I]) + _ZX_8(argument2.8[I])) + 1) >> 1
  };
stage E1:
  GPR[registerWe<<1] = result1;
\end{lstlisting}}
\newpage

\paragraph{Encoding} Format \textsf{ALU\_BOAVGO.M} (Section~\ref{sec:format-ALU-BOAVGO.M}), \verb|exucode2 = 1111|\\

\bitfields{ 1/P, 2/00, 2/-- --, 27/upper27, 1/1, 2/11, 1/1, 4/1111, 5/registerWo, 1/1, 2/10, 3/001, 1/1, 1/splat32, 5/lower5, 5/registerZo, 1/1 }


\paragraph{Syntax} \texttt{avgrubo} \emph{registerWo} = \emph{registerZo}, \emph{upper27\_\discretionary{}{}{}lower5}\emph{splat32}

\paragraph{Description} The \emph{upper27\_\discretionary{}{}{}lower5} and the \emph{registerZo} are interpreted as eight bytes packed into 64 bits. The \emph{upper27\_\discretionary{}{}{}lower5} bytes and the \emph{registerZo} bytes are added. Each sum is incremented and shifted right logical one bit. The eight resulting bytes are packed and stored into the \emph{registerWo}.


\paragraph{Modifier(s)}
\noindent\begin{itemize}
\item splat32 (cf. splat32 in section \ref{par:modifier-splat32} on page \pageref{par:modifier-splat32})
\end{itemize}

\paragraph{Execution} {Reservation table \textsf{ALU\_LITE.X} (Section~\ref{sec:reservation-ALU-LITE.X})}
~
{\begin{lstlisting}
stage ID:
  new splat32 = _ZX_1(splat32);
stage RR:
  new argument3 = splat32 ? (_WX_32(upper27_lower5) << 32) | _ZX_32(_WX_32(upper27_lower5)) : _SX_32(_WX_32(upper27_lower5));
  new argument2 = GPR[(registerZo<<1)+1];
  for (I = 0; I < 8; I += 1) {
    result1.8[I] = ((_ZX_8(argument3.8[I]) + _ZX_8(argument2.8[I])) + 1) >> 1
  };
stage E1:
  GPR[(registerWo<<1)+1] = result1;
\end{lstlisting}}
\newpage
\subsubsection[{\small AVGRUHQ: Average Rounded Unsigned Half Word Quadruple}]{AVGRUHQ\hfill Average Rounded Unsigned Half Word Quadruple}
\label{instr:AVGRUHQ}\index{avgruhq}
 
\paragraph{Encoding} Format \textsf{ALU\_HQAVGE} (Section~\ref{sec:format-ALU-HQAVGE}), \verb|exucode2 = 1111|\\

\bitfields{ 1/P, 2/11, 1/1, 4/1111, 5/registerWe, 1/0, 2/10, 3/001, 1/1, 5/registerYe, 1/--, 5/registerZe, 1/0 }


\paragraph{Syntax} \texttt{avgruhq} \emph{registerWe} = \emph{registerZe}, \emph{registerYe}

\paragraph{Description} The \emph{registerYe} and the \emph{registerZe} are interpreted as four half words packed into 64 bits. The \emph{registerYe} half words and the \emph{registerZe} half words are added. Each sum is incremented and shifted right logical one bit. The four resulting half words are packed and stored into the \emph{registerWe}.


\paragraph{Execution} {Reservation table \textsf{ALU\_LITE} (Section~\ref{sec:reservation-ALU-LITE})}
~
{\begin{lstlisting}
stage RR:
  new argument3 = GPR[registerYe<<1];
  new argument2 = GPR[registerZe<<1];
  for (I = 0; I < 4; I += 1) {
    result1.16[I] = ((_ZX_16(argument3.16[I]) + _ZX_16(argument2.16[I])) + 1) >> 1
  };
stage E1:
  GPR[registerWe<<1] = result1;
\end{lstlisting}}
\newpage

\paragraph{Encoding} Format \textsf{ALU\_HQAVGO} (Section~\ref{sec:format-ALU-HQAVGO}), \verb|exucode2 = 1111|\\

\bitfields{ 1/P, 2/11, 1/1, 4/1111, 5/registerWo, 1/1, 2/10, 3/001, 1/1, 5/registerYo, 1/--, 5/registerZo, 1/0 }


\paragraph{Syntax} \texttt{avgruhq} \emph{registerWo} = \emph{registerZo}, \emph{registerYo}

\paragraph{Description} The \emph{registerYo} and the \emph{registerZo} are interpreted as four half words packed into 64 bits. The \emph{registerYo} half words and the \emph{registerZo} half words are added. Each sum is incremented and shifted right logical one bit. The four resulting half words are packed and stored into the \emph{registerWo}.


\paragraph{Execution} {Reservation table \textsf{ALU\_LITE} (Section~\ref{sec:reservation-ALU-LITE})}
~
{\begin{lstlisting}
stage RR:
  new argument3 = GPR[(registerYo<<1)+1];
  new argument2 = GPR[(registerZo<<1)+1];
  for (I = 0; I < 4; I += 1) {
    result1.16[I] = ((_ZX_16(argument3.16[I]) + _ZX_16(argument2.16[I])) + 1) >> 1
  };
stage E1:
  GPR[(registerWo<<1)+1] = result1;
\end{lstlisting}}
\newpage

\paragraph{Encoding} Format \textsf{ALU\_HQAVGE.M} (Section~\ref{sec:format-ALU-HQAVGE.M}), \verb|exucode2 = 1111|\\

\bitfields{ 1/P, 2/00, 2/-- --, 27/upper27, 1/1, 2/11, 1/1, 4/1111, 5/registerWe, 1/0, 2/10, 3/001, 1/1, 1/splat32, 5/lower5, 5/registerZe, 1/0 }


\paragraph{Syntax} \texttt{avgruhq} \emph{registerWe} = \emph{registerZe}, \emph{upper27\_\discretionary{}{}{}lower5}\emph{splat32}

\paragraph{Description} The \emph{upper27\_\discretionary{}{}{}lower5} and the \emph{registerZe} are interpreted as four half words packed into 64 bits. The \emph{upper27\_\discretionary{}{}{}lower5} half words and the \emph{registerZe} half words are added. Each sum is incremented and shifted right logical one bit. The four resulting half words are packed and stored into the \emph{registerWe}.


\paragraph{Modifier(s)}
\noindent\begin{itemize}
\item splat32 (cf. splat32 in section \ref{par:modifier-splat32} on page \pageref{par:modifier-splat32})
\end{itemize}

\paragraph{Execution} {Reservation table \textsf{ALU\_LITE.X} (Section~\ref{sec:reservation-ALU-LITE.X})}
~
{\begin{lstlisting}
stage ID:
  new splat32 = _ZX_1(splat32);
stage RR:
  new argument3 = splat32 ? (_WX_32(upper27_lower5) << 32) | _ZX_32(_WX_32(upper27_lower5)) : _SX_32(_WX_32(upper27_lower5));
  new argument2 = GPR[registerZe<<1];
  for (I = 0; I < 4; I += 1) {
    result1.16[I] = ((_ZX_16(argument3.16[I]) + _ZX_16(argument2.16[I])) + 1) >> 1
  };
stage E1:
  GPR[registerWe<<1] = result1;
\end{lstlisting}}
\newpage

\paragraph{Encoding} Format \textsf{ALU\_HQAVGO.M} (Section~\ref{sec:format-ALU-HQAVGO.M}), \verb|exucode2 = 1111|\\

\bitfields{ 1/P, 2/00, 2/-- --, 27/upper27, 1/1, 2/11, 1/1, 4/1111, 5/registerWo, 1/1, 2/10, 3/001, 1/1, 1/splat32, 5/lower5, 5/registerZo, 1/0 }


\paragraph{Syntax} \texttt{avgruhq} \emph{registerWo} = \emph{registerZo}, \emph{upper27\_\discretionary{}{}{}lower5}\emph{splat32}

\paragraph{Description} The \emph{upper27\_\discretionary{}{}{}lower5} and the \emph{registerZo} are interpreted as four half words packed into 64 bits. The \emph{upper27\_\discretionary{}{}{}lower5} half words and the \emph{registerZo} half words are added. Each sum is incremented and shifted right logical one bit. The four resulting half words are packed and stored into the \emph{registerWo}.


\paragraph{Modifier(s)}
\noindent\begin{itemize}
\item splat32 (cf. splat32 in section \ref{par:modifier-splat32} on page \pageref{par:modifier-splat32})
\end{itemize}

\paragraph{Execution} {Reservation table \textsf{ALU\_LITE.X} (Section~\ref{sec:reservation-ALU-LITE.X})}
~
{\begin{lstlisting}
stage ID:
  new splat32 = _ZX_1(splat32);
stage RR:
  new argument3 = splat32 ? (_WX_32(upper27_lower5) << 32) | _ZX_32(_WX_32(upper27_lower5)) : _SX_32(_WX_32(upper27_lower5));
  new argument2 = GPR[(registerZo<<1)+1];
  for (I = 0; I < 4; I += 1) {
    result1.16[I] = ((_ZX_16(argument3.16[I]) + _ZX_16(argument2.16[I])) + 1) >> 1
  };
stage E1:
  GPR[(registerWo<<1)+1] = result1;
\end{lstlisting}}
\newpage
\subsubsection[{\small AVGRUW: Average Rounded Unsigned Word}]{AVGRUW\hfill Average Rounded Unsigned Word}
\label{instr:AVGRUW}\index{avgruw}
 
\paragraph{Encoding} Format \textsf{ALU\_WAVG} (Section~\ref{sec:format-ALU-WAVG}), \verb|exucode2 = 1111|\\

\bitfields{ 1/P, 2/11, 1/0, 4/1111, 6/registerW, 2/11, 3/001, 1/signextw, 6/registerY, 6/registerZ }


\paragraph{Syntax} \texttt{avgruw}\emph{signextw} \emph{registerW} = \emph{registerZ}, \emph{registerY}

\paragraph{Description} The \emph{registerY} and the \emph{registerZ} are added. The result is incremented, shifted right logical one bit and stored into the \emph{registerW}.


\paragraph{Modifier(s)}
\noindent\begin{itemize}
\item signextw (cf. signextw in section \ref{par:modifier-signextw} on page \pageref{par:modifier-signextw})
\end{itemize}

\paragraph{Execution} {Reservation table \textsf{ALU\_TINY} (Section~\ref{sec:reservation-ALU-TINY})}
~
{\begin{lstlisting}
stage ID:
  new signextw = _ZX_1(signextw);
stage RR:
  new argument3 = GPR[registerY];
  new argument2 = GPR[registerZ];
  new result1 = ((argument3.32[0] + argument2.32[0]) + 1) >> 1;
stage E1:
  GPR[registerW] = signextw ? _SX_32(result1) : _ZX_32(result1);
\end{lstlisting}}
\newpage

\paragraph{Encoding} Format \textsf{ALU\_WAVG.W} (Section~\ref{sec:format-ALU-WAVG.W}), \verb|exucode2 = 1111|\\

\bitfields{ 1/P, 2/00, 2/-- --, 27/upper27, 1/1, 2/11, 1/0, 4/1111, 6/registerW, 2/11, 3/001, 1/signextw, 1/0, 5/lower5, 6/registerZ }


\paragraph{Syntax} \texttt{avgruw}\emph{signextw} \emph{registerW} = \emph{registerZ}, \emph{upper27\_\discretionary{}{}{}lower5}

\paragraph{Description} The \emph{upper27\_\discretionary{}{}{}lower5} and the \emph{registerZ} are added. The result is incremented, shifted right logical one bit and stored into the \emph{registerW}.


\paragraph{Modifier(s)}
\noindent\begin{itemize}
\item signextw (cf. signextw in section \ref{par:modifier-signextw} on page \pageref{par:modifier-signextw})
\end{itemize}

\paragraph{Execution} {Reservation table \textsf{ALU\_TINY.X} (Section~\ref{sec:reservation-ALU-TINY.X})}
~
{\begin{lstlisting}
stage ID:
  new signextw = _ZX_1(signextw);
stage RR:
  new argument3 = _WX_32(upper27_lower5);
  new argument2 = GPR[registerZ];
  new result1 = ((argument3.32[0] + argument2.32[0]) + 1) >> 1;
stage E1:
  GPR[registerW] = signextw ? _SX_32(result1) : _ZX_32(result1);
\end{lstlisting}}
\newpage
\subsubsection[{\small AVGRUWP: Average Rounded Unsigned Word Pair}]{AVGRUWP\hfill Average Rounded Unsigned Word Pair}
\label{instr:AVGRUWP}\index{avgruwp}
 
\paragraph{Encoding} Format \textsf{ALU\_WPAVGE} (Section~\ref{sec:format-ALU-WPAVGE}), \verb|exucode2 = 1111|\\

\bitfields{ 1/P, 2/11, 1/1, 4/1111, 5/registerWe, 1/0, 2/10, 3/001, 1/0, 5/registerYe, 1/--, 5/registerZe, 1/1 }


\paragraph{Syntax} \texttt{avgruwp} \emph{registerWe} = \emph{registerZe}, \emph{registerYe}

\paragraph{Description} The \emph{registerYe} and the \emph{registerZe} are interpreted as two words packed into 64 bits. The \emph{registerYe} words and the \emph{registerZe} words are added. Each sum is incremented and shifted right logical one bit. The two resulting words are packed and stored into the \emph{registerWe}.


\paragraph{Execution} {Reservation table \textsf{ALU\_LITE} (Section~\ref{sec:reservation-ALU-LITE})}
~
{\begin{lstlisting}
stage RR:
  new argument3 = GPR[registerYe<<1];
  new argument2 = GPR[registerZe<<1];
  for (I = 0; I < 2; I += 1) {
    result1.32[I] = ((_ZX_32(argument3.32[I]) + _ZX_32(argument2.32[I])) + 1) >> 1
  };
stage E1:
  GPR[registerWe<<1] = result1;
\end{lstlisting}}
\newpage

\paragraph{Encoding} Format \textsf{ALU\_WPAVGO} (Section~\ref{sec:format-ALU-WPAVGO}), \verb|exucode2 = 1111|\\

\bitfields{ 1/P, 2/11, 1/1, 4/1111, 5/registerWo, 1/1, 2/10, 3/001, 1/0, 5/registerYo, 1/--, 5/registerZo, 1/1 }


\paragraph{Syntax} \texttt{avgruwp} \emph{registerWo} = \emph{registerZo}, \emph{registerYo}

\paragraph{Description} The \emph{registerYo} and the \emph{registerZo} are interpreted as two words packed into 64 bits. The \emph{registerYo} words and the \emph{registerZo} words are added. Each sum is incremented and shifted right logical one bit. The two resulting words are packed and stored into the \emph{registerWo}.


\paragraph{Execution} {Reservation table \textsf{ALU\_LITE} (Section~\ref{sec:reservation-ALU-LITE})}
~
{\begin{lstlisting}
stage RR:
  new argument3 = GPR[(registerYo<<1)+1];
  new argument2 = GPR[(registerZo<<1)+1];
  for (I = 0; I < 2; I += 1) {
    result1.32[I] = ((_ZX_32(argument3.32[I]) + _ZX_32(argument2.32[I])) + 1) >> 1
  };
stage E1:
  GPR[(registerWo<<1)+1] = result1;
\end{lstlisting}}
\newpage

\paragraph{Encoding} Format \textsf{ALU\_WPAVGE.M} (Section~\ref{sec:format-ALU-WPAVGE.M}), \verb|exucode2 = 1111|\\

\bitfields{ 1/P, 2/00, 2/-- --, 27/upper27, 1/1, 2/11, 1/1, 4/1111, 5/registerWe, 1/0, 2/10, 3/001, 1/0, 1/splat32, 5/lower5, 5/registerZe, 1/1 }


\paragraph{Syntax} \texttt{avgruwp} \emph{registerWe} = \emph{registerZe}, \emph{upper27\_\discretionary{}{}{}lower5}\emph{splat32}

\paragraph{Description} The \emph{upper27\_\discretionary{}{}{}lower5} and the \emph{registerZe} are interpreted as two words packed into 64 bits. The \emph{upper27\_\discretionary{}{}{}lower5} words and the \emph{registerZe} words are added. Each sum is incremented and shifted right logical one bit. The two resulting words are packed and stored into the \emph{registerWe}.


\paragraph{Modifier(s)}
\noindent\begin{itemize}
\item splat32 (cf. splat32 in section \ref{par:modifier-splat32} on page \pageref{par:modifier-splat32})
\end{itemize}

\paragraph{Execution} {Reservation table \textsf{ALU\_LITE.X} (Section~\ref{sec:reservation-ALU-LITE.X})}
~
{\begin{lstlisting}
stage ID:
  new splat32 = _ZX_1(splat32);
stage RR:
  new argument3 = splat32 ? (_WX_32(upper27_lower5) << 32) | _ZX_32(_WX_32(upper27_lower5)) : _SX_32(_WX_32(upper27_lower5));
  new argument2 = GPR[registerZe<<1];
  for (I = 0; I < 2; I += 1) {
    result1.32[I] = ((_ZX_32(argument3.32[I]) + _ZX_32(argument2.32[I])) + 1) >> 1
  };
stage E1:
  GPR[registerWe<<1] = result1;
\end{lstlisting}}
\newpage

\paragraph{Encoding} Format \textsf{ALU\_WPAVGO.M} (Section~\ref{sec:format-ALU-WPAVGO.M}), \verb|exucode2 = 1111|\\

\bitfields{ 1/P, 2/00, 2/-- --, 27/upper27, 1/1, 2/11, 1/1, 4/1111, 5/registerWo, 1/1, 2/10, 3/001, 1/0, 1/splat32, 5/lower5, 5/registerZo, 1/1 }


\paragraph{Syntax} \texttt{avgruwp} \emph{registerWo} = \emph{registerZo}, \emph{upper27\_\discretionary{}{}{}lower5}\emph{splat32}

\paragraph{Description} The \emph{upper27\_\discretionary{}{}{}lower5} and the \emph{registerZo} are interpreted as two words packed into 64 bits. The \emph{upper27\_\discretionary{}{}{}lower5} words and the \emph{registerZo} words are added. Each sum is incremented and shifted right logical one bit. The two resulting words are packed and stored into the \emph{registerWo}.


\paragraph{Modifier(s)}
\noindent\begin{itemize}
\item splat32 (cf. splat32 in section \ref{par:modifier-splat32} on page \pageref{par:modifier-splat32})
\end{itemize}

\paragraph{Execution} {Reservation table \textsf{ALU\_LITE.X} (Section~\ref{sec:reservation-ALU-LITE.X})}
~
{\begin{lstlisting}
stage ID:
  new splat32 = _ZX_1(splat32);
stage RR:
  new argument3 = splat32 ? (_WX_32(upper27_lower5) << 32) | _ZX_32(_WX_32(upper27_lower5)) : _SX_32(_WX_32(upper27_lower5));
  new argument2 = GPR[(registerZo<<1)+1];
  for (I = 0; I < 2; I += 1) {
    result1.32[I] = ((_ZX_32(argument3.32[I]) + _ZX_32(argument2.32[I])) + 1) >> 1
  };
stage E1:
  GPR[(registerWo<<1)+1] = result1;
\end{lstlisting}}
\newpage
\subsubsection[{\small AVGRW: Average Rounded Word}]{AVGRW\hfill Average Rounded Word}
\label{instr:AVGRW}\index{avgrw}
 
\paragraph{Encoding} Format \textsf{ALU\_WAVG} (Section~\ref{sec:format-ALU-WAVG}), \verb|exucode2 = 1110|\\

\bitfields{ 1/P, 2/11, 1/0, 4/1110, 6/registerW, 2/11, 3/001, 1/signextw, 6/registerY, 6/registerZ }


\paragraph{Syntax} \texttt{avgrw}\emph{signextw} \emph{registerW} = \emph{registerZ}, \emph{registerY}

\paragraph{Description} The \emph{registerY} and the \emph{registerZ} are added. The result is incremented, shifted right arithmetic one bit and stored into the \emph{registerW}.


\paragraph{Modifier(s)}
\noindent\begin{itemize}
\item signextw (cf. signextw in section \ref{par:modifier-signextw} on page \pageref{par:modifier-signextw})
\end{itemize}

\paragraph{Execution} {Reservation table \textsf{ALU\_TINY} (Section~\ref{sec:reservation-ALU-TINY})}
~
{\begin{lstlisting}
stage ID:
  new signextw = _ZX_1(signextw);
stage RR:
  new argument3 = GPR[registerY];
  new argument2 = GPR[registerZ];
  new result1 = ((_SX_32(argument3) + _SX_32(argument2)) + 1) >> 1;
stage E1:
  GPR[registerW] = signextw ? _SX_32(result1) : _ZX_32(result1);
\end{lstlisting}}
\newpage

\paragraph{Encoding} Format \textsf{ALU\_WAVG.W} (Section~\ref{sec:format-ALU-WAVG.W}), \verb|exucode2 = 1110|\\

\bitfields{ 1/P, 2/00, 2/-- --, 27/upper27, 1/1, 2/11, 1/0, 4/1110, 6/registerW, 2/11, 3/001, 1/signextw, 1/0, 5/lower5, 6/registerZ }


\paragraph{Syntax} \texttt{avgrw}\emph{signextw} \emph{registerW} = \emph{registerZ}, \emph{upper27\_\discretionary{}{}{}lower5}

\paragraph{Description} The \emph{upper27\_\discretionary{}{}{}lower5} and the \emph{registerZ} are added. The result is incremented, shifted right arithmetic one bit and stored into the \emph{registerW}.


\paragraph{Modifier(s)}
\noindent\begin{itemize}
\item signextw (cf. signextw in section \ref{par:modifier-signextw} on page \pageref{par:modifier-signextw})
\end{itemize}

\paragraph{Execution} {Reservation table \textsf{ALU\_TINY.X} (Section~\ref{sec:reservation-ALU-TINY.X})}
~
{\begin{lstlisting}
stage ID:
  new signextw = _ZX_1(signextw);
stage RR:
  new argument3 = _WX_32(upper27_lower5);
  new argument2 = GPR[registerZ];
  new result1 = ((_SX_32(argument3) + _SX_32(argument2)) + 1) >> 1;
stage E1:
  GPR[registerW] = signextw ? _SX_32(result1) : _ZX_32(result1);
\end{lstlisting}}
\newpage
\subsubsection[{\small AVGRWP: Average Rounded Word Pair}]{AVGRWP\hfill Average Rounded Word Pair}
\label{instr:AVGRWP}\index{avgrwp}
 
\paragraph{Encoding} Format \textsf{ALU\_WPAVGE} (Section~\ref{sec:format-ALU-WPAVGE}), \verb|exucode2 = 1110|\\

\bitfields{ 1/P, 2/11, 1/1, 4/1110, 5/registerWe, 1/0, 2/10, 3/001, 1/0, 5/registerYe, 1/--, 5/registerZe, 1/1 }


\paragraph{Syntax} \texttt{avgrwp} \emph{registerWe} = \emph{registerZe}, \emph{registerYe}

\paragraph{Description} The \emph{registerYe} and the \emph{registerZe} are interpreted as two words packed into 64 bits. The \emph{registerYe} words and the \emph{registerZe} words are added. Each sum is incremented and shifted right arithmetic one bit. The two resulting words are packed and stored into the \emph{registerWe}.


\paragraph{Execution} {Reservation table \textsf{ALU\_LITE} (Section~\ref{sec:reservation-ALU-LITE})}
~
{\begin{lstlisting}
stage RR:
  new argument3 = GPR[registerYe<<1];
  new argument2 = GPR[registerZe<<1];
  for (I = 0; I < 2; I += 1) {
    result1.32[I] = ((_SX_32(argument3.32[I]) + _SX_32(argument2.32[I])) + 1) >> 1
  };
stage E1:
  GPR[registerWe<<1] = result1;
\end{lstlisting}}
\newpage

\paragraph{Encoding} Format \textsf{ALU\_WPAVGO} (Section~\ref{sec:format-ALU-WPAVGO}), \verb|exucode2 = 1110|\\

\bitfields{ 1/P, 2/11, 1/1, 4/1110, 5/registerWo, 1/1, 2/10, 3/001, 1/0, 5/registerYo, 1/--, 5/registerZo, 1/1 }


\paragraph{Syntax} \texttt{avgrwp} \emph{registerWo} = \emph{registerZo}, \emph{registerYo}

\paragraph{Description} The \emph{registerYo} and the \emph{registerZo} are interpreted as two words packed into 64 bits. The \emph{registerYo} words and the \emph{registerZo} words are added. Each sum is incremented and shifted right arithmetic one bit. The two resulting words are packed and stored into the \emph{registerWo}.


\paragraph{Execution} {Reservation table \textsf{ALU\_LITE} (Section~\ref{sec:reservation-ALU-LITE})}
~
{\begin{lstlisting}
stage RR:
  new argument3 = GPR[(registerYo<<1)+1];
  new argument2 = GPR[(registerZo<<1)+1];
  for (I = 0; I < 2; I += 1) {
    result1.32[I] = ((_SX_32(argument3.32[I]) + _SX_32(argument2.32[I])) + 1) >> 1
  };
stage E1:
  GPR[(registerWo<<1)+1] = result1;
\end{lstlisting}}
\newpage

\paragraph{Encoding} Format \textsf{ALU\_WPAVGE.M} (Section~\ref{sec:format-ALU-WPAVGE.M}), \verb|exucode2 = 1110|\\

\bitfields{ 1/P, 2/00, 2/-- --, 27/upper27, 1/1, 2/11, 1/1, 4/1110, 5/registerWe, 1/0, 2/10, 3/001, 1/0, 1/splat32, 5/lower5, 5/registerZe, 1/1 }


\paragraph{Syntax} \texttt{avgrwp} \emph{registerWe} = \emph{registerZe}, \emph{upper27\_\discretionary{}{}{}lower5}\emph{splat32}

\paragraph{Description} The \emph{upper27\_\discretionary{}{}{}lower5} and the \emph{registerZe} are interpreted as two words packed into 64 bits. The \emph{upper27\_\discretionary{}{}{}lower5} words and the \emph{registerZe} words are added. Each sum is incremented and shifted right arithmetic one bit. The two resulting words are packed and stored into the \emph{registerWe}.


\paragraph{Modifier(s)}
\noindent\begin{itemize}
\item splat32 (cf. splat32 in section \ref{par:modifier-splat32} on page \pageref{par:modifier-splat32})
\end{itemize}

\paragraph{Execution} {Reservation table \textsf{ALU\_LITE.X} (Section~\ref{sec:reservation-ALU-LITE.X})}
~
{\begin{lstlisting}
stage ID:
  new splat32 = _ZX_1(splat32);
stage RR:
  new argument3 = splat32 ? (_WX_32(upper27_lower5) << 32) | _ZX_32(_WX_32(upper27_lower5)) : _SX_32(_WX_32(upper27_lower5));
  new argument2 = GPR[registerZe<<1];
  for (I = 0; I < 2; I += 1) {
    result1.32[I] = ((_SX_32(argument3.32[I]) + _SX_32(argument2.32[I])) + 1) >> 1
  };
stage E1:
  GPR[registerWe<<1] = result1;
\end{lstlisting}}
\newpage

\paragraph{Encoding} Format \textsf{ALU\_WPAVGO.M} (Section~\ref{sec:format-ALU-WPAVGO.M}), \verb|exucode2 = 1110|\\

\bitfields{ 1/P, 2/00, 2/-- --, 27/upper27, 1/1, 2/11, 1/1, 4/1110, 5/registerWo, 1/1, 2/10, 3/001, 1/0, 1/splat32, 5/lower5, 5/registerZo, 1/1 }


\paragraph{Syntax} \texttt{avgrwp} \emph{registerWo} = \emph{registerZo}, \emph{upper27\_\discretionary{}{}{}lower5}\emph{splat32}

\paragraph{Description} The \emph{upper27\_\discretionary{}{}{}lower5} and the \emph{registerZo} are interpreted as two words packed into 64 bits. The \emph{upper27\_\discretionary{}{}{}lower5} words and the \emph{registerZo} words are added. Each sum is incremented and shifted right arithmetic one bit. The two resulting words are packed and stored into the \emph{registerWo}.


\paragraph{Modifier(s)}
\noindent\begin{itemize}
\item splat32 (cf. splat32 in section \ref{par:modifier-splat32} on page \pageref{par:modifier-splat32})
\end{itemize}

\paragraph{Execution} {Reservation table \textsf{ALU\_LITE.X} (Section~\ref{sec:reservation-ALU-LITE.X})}
~
{\begin{lstlisting}
stage ID:
  new splat32 = _ZX_1(splat32);
stage RR:
  new argument3 = splat32 ? (_WX_32(upper27_lower5) << 32) | _ZX_32(_WX_32(upper27_lower5)) : _SX_32(_WX_32(upper27_lower5));
  new argument2 = GPR[(registerZo<<1)+1];
  for (I = 0; I < 2; I += 1) {
    result1.32[I] = ((_SX_32(argument3.32[I]) + _SX_32(argument2.32[I])) + 1) >> 1
  };
stage E1:
  GPR[(registerWo<<1)+1] = result1;
\end{lstlisting}}
\newpage
\subsubsection[{\small AVGUBO: Average Unsigned Byte Octuple}]{AVGUBO\hfill Average Unsigned Byte Octuple}
\label{instr:AVGUBO}\index{avgubo}
 
\paragraph{Encoding} Format \textsf{ALU\_BOAVGE} (Section~\ref{sec:format-ALU-BOAVGE}), \verb|exucode2 = 1101|\\

\bitfields{ 1/P, 2/11, 1/1, 4/1101, 5/registerWe, 1/0, 2/10, 3/001, 1/1, 5/registerYe, 1/--, 5/registerZe, 1/1 }


\paragraph{Syntax} \texttt{avgubo} \emph{registerWe} = \emph{registerZe}, \emph{registerYe}

\paragraph{Description} The \emph{registerYe} and the \emph{registerZe} are interpreted as eight bytes packed into 64 bits. The \emph{registerYe} bytes and the \emph{registerZe} bytes are added. Each sum is shifted right logical one bit. The eight resulting bytes are packed and stored into the \emph{registerWe}.


\paragraph{Execution} {Reservation table \textsf{ALU\_LITE} (Section~\ref{sec:reservation-ALU-LITE})}
~
{\begin{lstlisting}
stage RR:
  new argument3 = GPR[registerYe<<1];
  new argument2 = GPR[registerZe<<1];
  for (I = 0; I < 8; I += 1) {
    result1.8[I] = (_ZX_8(argument3.8[I]) + _ZX_8(argument2.8[I])) >> 1
  };
stage E1:
  GPR[registerWe<<1] = result1;
\end{lstlisting}}
\newpage

\paragraph{Encoding} Format \textsf{ALU\_BOAVGO} (Section~\ref{sec:format-ALU-BOAVGO}), \verb|exucode2 = 1101|\\

\bitfields{ 1/P, 2/11, 1/1, 4/1101, 5/registerWo, 1/1, 2/10, 3/001, 1/1, 5/registerYo, 1/--, 5/registerZo, 1/1 }


\paragraph{Syntax} \texttt{avgubo} \emph{registerWo} = \emph{registerZo}, \emph{registerYo}

\paragraph{Description} The \emph{registerYo} and the \emph{registerZo} are interpreted as eight bytes packed into 64 bits. The \emph{registerYo} bytes and the \emph{registerZo} bytes are added. Each sum is shifted right logical one bit. The eight resulting bytes are packed and stored into the \emph{registerWo}.


\paragraph{Execution} {Reservation table \textsf{ALU\_LITE} (Section~\ref{sec:reservation-ALU-LITE})}
~
{\begin{lstlisting}
stage RR:
  new argument3 = GPR[(registerYo<<1)+1];
  new argument2 = GPR[(registerZo<<1)+1];
  for (I = 0; I < 8; I += 1) {
    result1.8[I] = (_ZX_8(argument3.8[I]) + _ZX_8(argument2.8[I])) >> 1
  };
stage E1:
  GPR[(registerWo<<1)+1] = result1;
\end{lstlisting}}
\newpage

\paragraph{Encoding} Format \textsf{ALU\_BOAVGE.M} (Section~\ref{sec:format-ALU-BOAVGE.M}), \verb|exucode2 = 1101|\\

\bitfields{ 1/P, 2/00, 2/-- --, 27/upper27, 1/1, 2/11, 1/1, 4/1101, 5/registerWe, 1/0, 2/10, 3/001, 1/1, 1/splat32, 5/lower5, 5/registerZe, 1/1 }


\paragraph{Syntax} \texttt{avgubo} \emph{registerWe} = \emph{registerZe}, \emph{upper27\_\discretionary{}{}{}lower5}\emph{splat32}

\paragraph{Description} The \emph{upper27\_\discretionary{}{}{}lower5} and the \emph{registerZe} are interpreted as eight bytes packed into 64 bits. The \emph{upper27\_\discretionary{}{}{}lower5} bytes and the \emph{registerZe} bytes are added. Each sum is shifted right logical one bit. The eight resulting bytes are packed and stored into the \emph{registerWe}.


\paragraph{Modifier(s)}
\noindent\begin{itemize}
\item splat32 (cf. splat32 in section \ref{par:modifier-splat32} on page \pageref{par:modifier-splat32})
\end{itemize}

\paragraph{Execution} {Reservation table \textsf{ALU\_LITE.X} (Section~\ref{sec:reservation-ALU-LITE.X})}
~
{\begin{lstlisting}
stage ID:
  new splat32 = _ZX_1(splat32);
stage RR:
  new argument3 = splat32 ? (_WX_32(upper27_lower5) << 32) | _ZX_32(_WX_32(upper27_lower5)) : _SX_32(_WX_32(upper27_lower5));
  new argument2 = GPR[registerZe<<1];
  for (I = 0; I < 8; I += 1) {
    result1.8[I] = (_ZX_8(argument3.8[I]) + _ZX_8(argument2.8[I])) >> 1
  };
stage E1:
  GPR[registerWe<<1] = result1;
\end{lstlisting}}
\newpage

\paragraph{Encoding} Format \textsf{ALU\_BOAVGO.M} (Section~\ref{sec:format-ALU-BOAVGO.M}), \verb|exucode2 = 1101|\\

\bitfields{ 1/P, 2/00, 2/-- --, 27/upper27, 1/1, 2/11, 1/1, 4/1101, 5/registerWo, 1/1, 2/10, 3/001, 1/1, 1/splat32, 5/lower5, 5/registerZo, 1/1 }


\paragraph{Syntax} \texttt{avgubo} \emph{registerWo} = \emph{registerZo}, \emph{upper27\_\discretionary{}{}{}lower5}\emph{splat32}

\paragraph{Description} The \emph{upper27\_\discretionary{}{}{}lower5} and the \emph{registerZo} are interpreted as eight bytes packed into 64 bits. The \emph{upper27\_\discretionary{}{}{}lower5} bytes and the \emph{registerZo} bytes are added. Each sum is shifted right logical one bit. The eight resulting bytes are packed and stored into the \emph{registerWo}.


\paragraph{Modifier(s)}
\noindent\begin{itemize}
\item splat32 (cf. splat32 in section \ref{par:modifier-splat32} on page \pageref{par:modifier-splat32})
\end{itemize}

\paragraph{Execution} {Reservation table \textsf{ALU\_LITE.X} (Section~\ref{sec:reservation-ALU-LITE.X})}
~
{\begin{lstlisting}
stage ID:
  new splat32 = _ZX_1(splat32);
stage RR:
  new argument3 = splat32 ? (_WX_32(upper27_lower5) << 32) | _ZX_32(_WX_32(upper27_lower5)) : _SX_32(_WX_32(upper27_lower5));
  new argument2 = GPR[(registerZo<<1)+1];
  for (I = 0; I < 8; I += 1) {
    result1.8[I] = (_ZX_8(argument3.8[I]) + _ZX_8(argument2.8[I])) >> 1
  };
stage E1:
  GPR[(registerWo<<1)+1] = result1;
\end{lstlisting}}
\newpage
\subsubsection[{\small AVGUHQ: Average Unsigned Half Word Quadruple}]{AVGUHQ\hfill Average Unsigned Half Word Quadruple}
\label{instr:AVGUHQ}\index{avguhq}
 
\paragraph{Encoding} Format \textsf{ALU\_HQAVGE} (Section~\ref{sec:format-ALU-HQAVGE}), \verb|exucode2 = 1101|\\

\bitfields{ 1/P, 2/11, 1/1, 4/1101, 5/registerWe, 1/0, 2/10, 3/001, 1/1, 5/registerYe, 1/--, 5/registerZe, 1/0 }


\paragraph{Syntax} \texttt{avguhq} \emph{registerWe} = \emph{registerZe}, \emph{registerYe}

\paragraph{Description} The \emph{registerYe} and the \emph{registerZe} are interpreted as four half words packed into 64 bits. The \emph{registerYe} half words and the \emph{registerZe} half words are added. Each sum is shifted right logical one bit. The four resulting half words are packed and stored into the \emph{registerWe}.


\paragraph{Execution} {Reservation table \textsf{ALU\_LITE} (Section~\ref{sec:reservation-ALU-LITE})}
~
{\begin{lstlisting}
stage RR:
  new argument3 = GPR[registerYe<<1];
  new argument2 = GPR[registerZe<<1];
  for (I = 0; I < 4; I += 1) {
    result1.16[I] = (_ZX_16(argument3.16[I]) + _ZX_16(argument2.16[I])) >> 1
  };
stage E1:
  GPR[registerWe<<1] = result1;
\end{lstlisting}}
\newpage

\paragraph{Encoding} Format \textsf{ALU\_HQAVGO} (Section~\ref{sec:format-ALU-HQAVGO}), \verb|exucode2 = 1101|\\

\bitfields{ 1/P, 2/11, 1/1, 4/1101, 5/registerWo, 1/1, 2/10, 3/001, 1/1, 5/registerYo, 1/--, 5/registerZo, 1/0 }


\paragraph{Syntax} \texttt{avguhq} \emph{registerWo} = \emph{registerZo}, \emph{registerYo}

\paragraph{Description} The \emph{registerYo} and the \emph{registerZo} are interpreted as four half words packed into 64 bits. The \emph{registerYo} half words and the \emph{registerZo} half words are added. Each sum is shifted right logical one bit. The four resulting half words are packed and stored into the \emph{registerWo}.


\paragraph{Execution} {Reservation table \textsf{ALU\_LITE} (Section~\ref{sec:reservation-ALU-LITE})}
~
{\begin{lstlisting}
stage RR:
  new argument3 = GPR[(registerYo<<1)+1];
  new argument2 = GPR[(registerZo<<1)+1];
  for (I = 0; I < 4; I += 1) {
    result1.16[I] = (_ZX_16(argument3.16[I]) + _ZX_16(argument2.16[I])) >> 1
  };
stage E1:
  GPR[(registerWo<<1)+1] = result1;
\end{lstlisting}}
\newpage

\paragraph{Encoding} Format \textsf{ALU\_HQAVGE.M} (Section~\ref{sec:format-ALU-HQAVGE.M}), \verb|exucode2 = 1101|\\

\bitfields{ 1/P, 2/00, 2/-- --, 27/upper27, 1/1, 2/11, 1/1, 4/1101, 5/registerWe, 1/0, 2/10, 3/001, 1/1, 1/splat32, 5/lower5, 5/registerZe, 1/0 }


\paragraph{Syntax} \texttt{avguhq} \emph{registerWe} = \emph{registerZe}, \emph{upper27\_\discretionary{}{}{}lower5}\emph{splat32}

\paragraph{Description} The \emph{upper27\_\discretionary{}{}{}lower5} and the \emph{registerZe} are interpreted as four half words packed into 64 bits. The \emph{upper27\_\discretionary{}{}{}lower5} half words and the \emph{registerZe} half words are added. Each sum is shifted right logical one bit. The four resulting half words are packed and stored into the \emph{registerWe}.


\paragraph{Modifier(s)}
\noindent\begin{itemize}
\item splat32 (cf. splat32 in section \ref{par:modifier-splat32} on page \pageref{par:modifier-splat32})
\end{itemize}

\paragraph{Execution} {Reservation table \textsf{ALU\_LITE.X} (Section~\ref{sec:reservation-ALU-LITE.X})}
~
{\begin{lstlisting}
stage ID:
  new splat32 = _ZX_1(splat32);
stage RR:
  new argument3 = splat32 ? (_WX_32(upper27_lower5) << 32) | _ZX_32(_WX_32(upper27_lower5)) : _SX_32(_WX_32(upper27_lower5));
  new argument2 = GPR[registerZe<<1];
  for (I = 0; I < 4; I += 1) {
    result1.16[I] = (_ZX_16(argument3.16[I]) + _ZX_16(argument2.16[I])) >> 1
  };
stage E1:
  GPR[registerWe<<1] = result1;
\end{lstlisting}}
\newpage

\paragraph{Encoding} Format \textsf{ALU\_HQAVGO.M} (Section~\ref{sec:format-ALU-HQAVGO.M}), \verb|exucode2 = 1101|\\

\bitfields{ 1/P, 2/00, 2/-- --, 27/upper27, 1/1, 2/11, 1/1, 4/1101, 5/registerWo, 1/1, 2/10, 3/001, 1/1, 1/splat32, 5/lower5, 5/registerZo, 1/0 }


\paragraph{Syntax} \texttt{avguhq} \emph{registerWo} = \emph{registerZo}, \emph{upper27\_\discretionary{}{}{}lower5}\emph{splat32}

\paragraph{Description} The \emph{upper27\_\discretionary{}{}{}lower5} and the \emph{registerZo} are interpreted as four half words packed into 64 bits. The \emph{upper27\_\discretionary{}{}{}lower5} half words and the \emph{registerZo} half words are added. Each sum is shifted right logical one bit. The four resulting half words are packed and stored into the \emph{registerWo}.


\paragraph{Modifier(s)}
\noindent\begin{itemize}
\item splat32 (cf. splat32 in section \ref{par:modifier-splat32} on page \pageref{par:modifier-splat32})
\end{itemize}

\paragraph{Execution} {Reservation table \textsf{ALU\_LITE.X} (Section~\ref{sec:reservation-ALU-LITE.X})}
~
{\begin{lstlisting}
stage ID:
  new splat32 = _ZX_1(splat32);
stage RR:
  new argument3 = splat32 ? (_WX_32(upper27_lower5) << 32) | _ZX_32(_WX_32(upper27_lower5)) : _SX_32(_WX_32(upper27_lower5));
  new argument2 = GPR[(registerZo<<1)+1];
  for (I = 0; I < 4; I += 1) {
    result1.16[I] = (_ZX_16(argument3.16[I]) + _ZX_16(argument2.16[I])) >> 1
  };
stage E1:
  GPR[(registerWo<<1)+1] = result1;
\end{lstlisting}}
\newpage
\subsubsection[{\small AVGUW: Average Unsigned Word}]{AVGUW\hfill Average Unsigned Word}
\label{instr:AVGUW}\index{avguw}
 
\paragraph{Encoding} Format \textsf{ALU\_WAVG} (Section~\ref{sec:format-ALU-WAVG}), \verb|exucode2 = 1101|\\

\bitfields{ 1/P, 2/11, 1/0, 4/1101, 6/registerW, 2/11, 3/001, 1/signextw, 6/registerY, 6/registerZ }


\paragraph{Syntax} \texttt{avguw}\emph{signextw} \emph{registerW} = \emph{registerZ}, \emph{registerY}

\paragraph{Description} The \emph{registerY} and the \emph{registerZ} are added. The result is shifted right logical one bit and stored into the \emph{registerW}.


\paragraph{Modifier(s)}
\noindent\begin{itemize}
\item signextw (cf. signextw in section \ref{par:modifier-signextw} on page \pageref{par:modifier-signextw})
\end{itemize}

\paragraph{Execution} {Reservation table \textsf{ALU\_TINY} (Section~\ref{sec:reservation-ALU-TINY})}
~
{\begin{lstlisting}
stage ID:
  new signextw = _ZX_1(signextw);
stage RR:
  new argument3 = GPR[registerY];
  new argument2 = GPR[registerZ];
  new result1 = (argument3.32[0] + argument2.32[0]) >> 1;
stage E1:
  GPR[registerW] = signextw ? _SX_32(result1) : _ZX_32(result1);
\end{lstlisting}}
\newpage

\paragraph{Encoding} Format \textsf{ALU\_WAVG.W} (Section~\ref{sec:format-ALU-WAVG.W}), \verb|exucode2 = 1101|\\

\bitfields{ 1/P, 2/00, 2/-- --, 27/upper27, 1/1, 2/11, 1/0, 4/1101, 6/registerW, 2/11, 3/001, 1/signextw, 1/0, 5/lower5, 6/registerZ }


\paragraph{Syntax} \texttt{avguw}\emph{signextw} \emph{registerW} = \emph{registerZ}, \emph{upper27\_\discretionary{}{}{}lower5}

\paragraph{Description} The \emph{upper27\_\discretionary{}{}{}lower5} and the \emph{registerZ} are added. The result is shifted right logical one bit and stored into the \emph{registerW}.


\paragraph{Modifier(s)}
\noindent\begin{itemize}
\item signextw (cf. signextw in section \ref{par:modifier-signextw} on page \pageref{par:modifier-signextw})
\end{itemize}

\paragraph{Execution} {Reservation table \textsf{ALU\_TINY.X} (Section~\ref{sec:reservation-ALU-TINY.X})}
~
{\begin{lstlisting}
stage ID:
  new signextw = _ZX_1(signextw);
stage RR:
  new argument3 = _WX_32(upper27_lower5);
  new argument2 = GPR[registerZ];
  new result1 = (argument3.32[0] + argument2.32[0]) >> 1;
stage E1:
  GPR[registerW] = signextw ? _SX_32(result1) : _ZX_32(result1);
\end{lstlisting}}
\newpage
\subsubsection[{\small AVGUWP: Average Unsigned Word Pair}]{AVGUWP\hfill Average Unsigned Word Pair}
\label{instr:AVGUWP}\index{avguwp}
 
\paragraph{Encoding} Format \textsf{ALU\_WPAVGE} (Section~\ref{sec:format-ALU-WPAVGE}), \verb|exucode2 = 1101|\\

\bitfields{ 1/P, 2/11, 1/1, 4/1101, 5/registerWe, 1/0, 2/10, 3/001, 1/0, 5/registerYe, 1/--, 5/registerZe, 1/1 }


\paragraph{Syntax} \texttt{avguwp} \emph{registerWe} = \emph{registerZe}, \emph{registerYe}

\paragraph{Description} The \emph{registerYe} and the \emph{registerZe} are interpreted as two words packed into 64 bits. The \emph{registerYe} words and the \emph{registerZe} words are added. Each sum is shifted right logical one bit. The two resulting words are packed and stored into the \emph{registerWe}.


\paragraph{Execution} {Reservation table \textsf{ALU\_LITE} (Section~\ref{sec:reservation-ALU-LITE})}
~
{\begin{lstlisting}
stage RR:
  new argument3 = GPR[registerYe<<1];
  new argument2 = GPR[registerZe<<1];
  for (I = 0; I < 2; I += 1) {
    result1.32[I] = (_ZX_32(argument3.32[I]) + _ZX_32(argument2.32[I])) >> 1
  };
stage E1:
  GPR[registerWe<<1] = result1;
\end{lstlisting}}
\newpage

\paragraph{Encoding} Format \textsf{ALU\_WPAVGO} (Section~\ref{sec:format-ALU-WPAVGO}), \verb|exucode2 = 1101|\\

\bitfields{ 1/P, 2/11, 1/1, 4/1101, 5/registerWo, 1/1, 2/10, 3/001, 1/0, 5/registerYo, 1/--, 5/registerZo, 1/1 }


\paragraph{Syntax} \texttt{avguwp} \emph{registerWo} = \emph{registerZo}, \emph{registerYo}

\paragraph{Description} The \emph{registerYo} and the \emph{registerZo} are interpreted as two words packed into 64 bits. The \emph{registerYo} words and the \emph{registerZo} words are added. Each sum is shifted right logical one bit. The two resulting words are packed and stored into the \emph{registerWo}.


\paragraph{Execution} {Reservation table \textsf{ALU\_LITE} (Section~\ref{sec:reservation-ALU-LITE})}
~
{\begin{lstlisting}
stage RR:
  new argument3 = GPR[(registerYo<<1)+1];
  new argument2 = GPR[(registerZo<<1)+1];
  for (I = 0; I < 2; I += 1) {
    result1.32[I] = (_ZX_32(argument3.32[I]) + _ZX_32(argument2.32[I])) >> 1
  };
stage E1:
  GPR[(registerWo<<1)+1] = result1;
\end{lstlisting}}
\newpage

\paragraph{Encoding} Format \textsf{ALU\_WPAVGE.M} (Section~\ref{sec:format-ALU-WPAVGE.M}), \verb|exucode2 = 1101|\\

\bitfields{ 1/P, 2/00, 2/-- --, 27/upper27, 1/1, 2/11, 1/1, 4/1101, 5/registerWe, 1/0, 2/10, 3/001, 1/0, 1/splat32, 5/lower5, 5/registerZe, 1/1 }


\paragraph{Syntax} \texttt{avguwp} \emph{registerWe} = \emph{registerZe}, \emph{upper27\_\discretionary{}{}{}lower5}\emph{splat32}

\paragraph{Description} The \emph{upper27\_\discretionary{}{}{}lower5} and the \emph{registerZe} are interpreted as two words packed into 64 bits. The \emph{upper27\_\discretionary{}{}{}lower5} words and the \emph{registerZe} words are added. Each sum is shifted right logical one bit. The two resulting words are packed and stored into the \emph{registerWe}.


\paragraph{Modifier(s)}
\noindent\begin{itemize}
\item splat32 (cf. splat32 in section \ref{par:modifier-splat32} on page \pageref{par:modifier-splat32})
\end{itemize}

\paragraph{Execution} {Reservation table \textsf{ALU\_LITE.X} (Section~\ref{sec:reservation-ALU-LITE.X})}
~
{\begin{lstlisting}
stage ID:
  new splat32 = _ZX_1(splat32);
stage RR:
  new argument3 = splat32 ? (_WX_32(upper27_lower5) << 32) | _ZX_32(_WX_32(upper27_lower5)) : _SX_32(_WX_32(upper27_lower5));
  new argument2 = GPR[registerZe<<1];
  for (I = 0; I < 2; I += 1) {
    result1.32[I] = (_ZX_32(argument3.32[I]) + _ZX_32(argument2.32[I])) >> 1
  };
stage E1:
  GPR[registerWe<<1] = result1;
\end{lstlisting}}
\newpage

\paragraph{Encoding} Format \textsf{ALU\_WPAVGO.M} (Section~\ref{sec:format-ALU-WPAVGO.M}), \verb|exucode2 = 1101|\\

\bitfields{ 1/P, 2/00, 2/-- --, 27/upper27, 1/1, 2/11, 1/1, 4/1101, 5/registerWo, 1/1, 2/10, 3/001, 1/0, 1/splat32, 5/lower5, 5/registerZo, 1/1 }


\paragraph{Syntax} \texttt{avguwp} \emph{registerWo} = \emph{registerZo}, \emph{upper27\_\discretionary{}{}{}lower5}\emph{splat32}

\paragraph{Description} The \emph{upper27\_\discretionary{}{}{}lower5} and the \emph{registerZo} are interpreted as two words packed into 64 bits. The \emph{upper27\_\discretionary{}{}{}lower5} words and the \emph{registerZo} words are added. Each sum is shifted right logical one bit. The two resulting words are packed and stored into the \emph{registerWo}.


\paragraph{Modifier(s)}
\noindent\begin{itemize}
\item splat32 (cf. splat32 in section \ref{par:modifier-splat32} on page \pageref{par:modifier-splat32})
\end{itemize}

\paragraph{Execution} {Reservation table \textsf{ALU\_LITE.X} (Section~\ref{sec:reservation-ALU-LITE.X})}
~
{\begin{lstlisting}
stage ID:
  new splat32 = _ZX_1(splat32);
stage RR:
  new argument3 = splat32 ? (_WX_32(upper27_lower5) << 32) | _ZX_32(_WX_32(upper27_lower5)) : _SX_32(_WX_32(upper27_lower5));
  new argument2 = GPR[(registerZo<<1)+1];
  for (I = 0; I < 2; I += 1) {
    result1.32[I] = (_ZX_32(argument3.32[I]) + _ZX_32(argument2.32[I])) >> 1
  };
stage E1:
  GPR[(registerWo<<1)+1] = result1;
\end{lstlisting}}
\newpage
\subsubsection[{\small AVGW: Average Word}]{AVGW\hfill Average Word}
\label{instr:AVGW}\index{avgw}
 
\paragraph{Encoding} Format \textsf{ALU\_WAVG} (Section~\ref{sec:format-ALU-WAVG}), \verb|exucode2 = 1100|\\

\bitfields{ 1/P, 2/11, 1/0, 4/1100, 6/registerW, 2/11, 3/001, 1/signextw, 6/registerY, 6/registerZ }


\paragraph{Syntax} \texttt{avgw}\emph{signextw} \emph{registerW} = \emph{registerZ}, \emph{registerY}

\paragraph{Description} The \emph{registerY} and the \emph{registerZ} are added. The result is shifted right arithmetic one bit and stored into the \emph{registerW}.


\paragraph{Modifier(s)}
\noindent\begin{itemize}
\item signextw (cf. signextw in section \ref{par:modifier-signextw} on page \pageref{par:modifier-signextw})
\end{itemize}

\paragraph{Execution} {Reservation table \textsf{ALU\_TINY} (Section~\ref{sec:reservation-ALU-TINY})}
~
{\begin{lstlisting}
stage ID:
  new signextw = _ZX_1(signextw);
stage RR:
  new argument3 = GPR[registerY];
  new argument2 = GPR[registerZ];
  new result1 = (_SX_32(argument3) + _SX_32(argument2)) >> 1;
stage E1:
  GPR[registerW] = signextw ? _SX_32(result1) : _ZX_32(result1);
\end{lstlisting}}
\newpage

\paragraph{Encoding} Format \textsf{ALU\_WAVG.W} (Section~\ref{sec:format-ALU-WAVG.W}), \verb|exucode2 = 1100|\\

\bitfields{ 1/P, 2/00, 2/-- --, 27/upper27, 1/1, 2/11, 1/0, 4/1100, 6/registerW, 2/11, 3/001, 1/signextw, 1/0, 5/lower5, 6/registerZ }


\paragraph{Syntax} \texttt{avgw}\emph{signextw} \emph{registerW} = \emph{registerZ}, \emph{upper27\_\discretionary{}{}{}lower5}

\paragraph{Description} The \emph{upper27\_\discretionary{}{}{}lower5} and the \emph{registerZ} are added. The result is shifted right arithmetic one bit and stored into the \emph{registerW}.


\paragraph{Modifier(s)}
\noindent\begin{itemize}
\item signextw (cf. signextw in section \ref{par:modifier-signextw} on page \pageref{par:modifier-signextw})
\end{itemize}

\paragraph{Execution} {Reservation table \textsf{ALU\_TINY.X} (Section~\ref{sec:reservation-ALU-TINY.X})}
~
{\begin{lstlisting}
stage ID:
  new signextw = _ZX_1(signextw);
stage RR:
  new argument3 = _WX_32(upper27_lower5);
  new argument2 = GPR[registerZ];
  new result1 = (_SX_32(argument3) + _SX_32(argument2)) >> 1;
stage E1:
  GPR[registerW] = signextw ? _SX_32(result1) : _ZX_32(result1);
\end{lstlisting}}
\newpage
\subsubsection[{\small AVGWP: Average Word Pair}]{AVGWP\hfill Average Word Pair}
\label{instr:AVGWP}\index{avgwp}
 
\paragraph{Encoding} Format \textsf{ALU\_WPAVGE} (Section~\ref{sec:format-ALU-WPAVGE}), \verb|exucode2 = 1100|\\

\bitfields{ 1/P, 2/11, 1/1, 4/1100, 5/registerWe, 1/0, 2/10, 3/001, 1/0, 5/registerYe, 1/--, 5/registerZe, 1/1 }


\paragraph{Syntax} \texttt{avgwp} \emph{registerWe} = \emph{registerZe}, \emph{registerYe}

\paragraph{Description} The \emph{registerYe} and the \emph{registerZe} are interpreted as two words packed into 64 bits. The \emph{registerYe} words and the \emph{registerZe} words are added. Each sum is shifted right arithmetic one bit. The two resulting words are packed and stored into the \emph{registerWe}.


\paragraph{Execution} {Reservation table \textsf{ALU\_LITE} (Section~\ref{sec:reservation-ALU-LITE})}
~
{\begin{lstlisting}
stage RR:
  new argument3 = GPR[registerYe<<1];
  new argument2 = GPR[registerZe<<1];
  for (I = 0; I < 2; I += 1) {
    result1.32[I] = (_SX_32(argument3.32[I]) + _SX_32(argument2.32[I])) >> 1
  };
stage E1:
  GPR[registerWe<<1] = result1;
\end{lstlisting}}
\newpage

\paragraph{Encoding} Format \textsf{ALU\_WPAVGO} (Section~\ref{sec:format-ALU-WPAVGO}), \verb|exucode2 = 1100|\\

\bitfields{ 1/P, 2/11, 1/1, 4/1100, 5/registerWo, 1/1, 2/10, 3/001, 1/0, 5/registerYo, 1/--, 5/registerZo, 1/1 }


\paragraph{Syntax} \texttt{avgwp} \emph{registerWo} = \emph{registerZo}, \emph{registerYo}

\paragraph{Description} The \emph{registerYo} and the \emph{registerZo} are interpreted as two words packed into 64 bits. The \emph{registerYo} words and the \emph{registerZo} words are added. Each sum is shifted right arithmetic one bit. The two resulting words are packed and stored into the \emph{registerWo}.


\paragraph{Execution} {Reservation table \textsf{ALU\_LITE} (Section~\ref{sec:reservation-ALU-LITE})}
~
{\begin{lstlisting}
stage RR:
  new argument3 = GPR[(registerYo<<1)+1];
  new argument2 = GPR[(registerZo<<1)+1];
  for (I = 0; I < 2; I += 1) {
    result1.32[I] = (_SX_32(argument3.32[I]) + _SX_32(argument2.32[I])) >> 1
  };
stage E1:
  GPR[(registerWo<<1)+1] = result1;
\end{lstlisting}}
\newpage

\paragraph{Encoding} Format \textsf{ALU\_WPAVGE.M} (Section~\ref{sec:format-ALU-WPAVGE.M}), \verb|exucode2 = 1100|\\

\bitfields{ 1/P, 2/00, 2/-- --, 27/upper27, 1/1, 2/11, 1/1, 4/1100, 5/registerWe, 1/0, 2/10, 3/001, 1/0, 1/splat32, 5/lower5, 5/registerZe, 1/1 }


\paragraph{Syntax} \texttt{avgwp} \emph{registerWe} = \emph{registerZe}, \emph{upper27\_\discretionary{}{}{}lower5}\emph{splat32}

\paragraph{Description} The \emph{upper27\_\discretionary{}{}{}lower5} and the \emph{registerZe} are interpreted as two words packed into 64 bits. The \emph{upper27\_\discretionary{}{}{}lower5} words and the \emph{registerZe} words are added. Each sum is shifted right arithmetic one bit. The two resulting words are packed and stored into the \emph{registerWe}.


\paragraph{Modifier(s)}
\noindent\begin{itemize}
\item splat32 (cf. splat32 in section \ref{par:modifier-splat32} on page \pageref{par:modifier-splat32})
\end{itemize}

\paragraph{Execution} {Reservation table \textsf{ALU\_LITE.X} (Section~\ref{sec:reservation-ALU-LITE.X})}
~
{\begin{lstlisting}
stage ID:
  new splat32 = _ZX_1(splat32);
stage RR:
  new argument3 = splat32 ? (_WX_32(upper27_lower5) << 32) | _ZX_32(_WX_32(upper27_lower5)) : _SX_32(_WX_32(upper27_lower5));
  new argument2 = GPR[registerZe<<1];
  for (I = 0; I < 2; I += 1) {
    result1.32[I] = (_SX_32(argument3.32[I]) + _SX_32(argument2.32[I])) >> 1
  };
stage E1:
  GPR[registerWe<<1] = result1;
\end{lstlisting}}
\newpage

\paragraph{Encoding} Format \textsf{ALU\_WPAVGO.M} (Section~\ref{sec:format-ALU-WPAVGO.M}), \verb|exucode2 = 1100|\\

\bitfields{ 1/P, 2/00, 2/-- --, 27/upper27, 1/1, 2/11, 1/1, 4/1100, 5/registerWo, 1/1, 2/10, 3/001, 1/0, 1/splat32, 5/lower5, 5/registerZo, 1/1 }


\paragraph{Syntax} \texttt{avgwp} \emph{registerWo} = \emph{registerZo}, \emph{upper27\_\discretionary{}{}{}lower5}\emph{splat32}

\paragraph{Description} The \emph{upper27\_\discretionary{}{}{}lower5} and the \emph{registerZo} are interpreted as two words packed into 64 bits. The \emph{upper27\_\discretionary{}{}{}lower5} words and the \emph{registerZo} words are added. Each sum is shifted right arithmetic one bit. The two resulting words are packed and stored into the \emph{registerWo}.


\paragraph{Modifier(s)}
\noindent\begin{itemize}
\item splat32 (cf. splat32 in section \ref{par:modifier-splat32} on page \pageref{par:modifier-splat32})
\end{itemize}

\paragraph{Execution} {Reservation table \textsf{ALU\_LITE.X} (Section~\ref{sec:reservation-ALU-LITE.X})}
~
{\begin{lstlisting}
stage ID:
  new splat32 = _ZX_1(splat32);
stage RR:
  new argument3 = splat32 ? (_WX_32(upper27_lower5) << 32) | _ZX_32(_WX_32(upper27_lower5)) : _SX_32(_WX_32(upper27_lower5));
  new argument2 = GPR[(registerZo<<1)+1];
  for (I = 0; I < 2; I += 1) {
    result1.32[I] = (_SX_32(argument3.32[I]) + _SX_32(argument2.32[I])) >> 1
  };
stage E1:
  GPR[(registerWo<<1)+1] = result1;
\end{lstlisting}}
\newpage
\subsubsection[{\small CATDQ: Catenate Double to Quadruple Word}]{CATDQ\hfill Catenate Double to Quadruple Word}
\label{instr:CATDQ}\index{catdq}
 
\paragraph{Encoding} Format \textsf{ALU\_CATDQ} (Section~\ref{sec:format-ALU-CATDQ}), \verb|exucode2 = 1111|\\

\bitfields{ 1/P, 2/11, 1/0, 4/1111, 5/registerM, 1/--, 2/10, 3/010, 1/1, 6/registerY, 6/registerZ }


\paragraph{Syntax} \texttt{catdq} \emph{registerM} = \emph{registerZ}, \emph{registerY}

\paragraph{Description} The \emph{registerM} operand receives the concatenation of the \emph{registerZ} and \emph{registerY} values.


\paragraph{Execution} {Reservation table \textsf{ALU\_LITE} (Section~\ref{sec:reservation-ALU-LITE})}
~
{\begin{lstlisting}
stage RR:
  new argument3 = GPR[registerY];
  new argument2 = GPR[registerZ];
  new result1 = _ZX_64(argument2) | (argument3 << 64);
stage E1:
  PGR[registerM] = result1;
\end{lstlisting}}
\newpage
\subsubsection[{\small CBSD: Count Bit Set Double Word}]{CBSD\hfill Count Bit Set Double Word}
\label{instr:CBSD}\index{cbsd}
 
\paragraph{Encoding} Format \textsf{ALU\_DBWR} (Section~\ref{sec:format-ALU-DBWR}), \verb|alucode8 = 0110|\\

\bitfields{ 1/P, 2/11, 1/0, 4/1111, 6/registerW, 2/10, 3/101, 1/0, 4/0110, 1/0, 1/0, 6/registerZ }


\paragraph{Syntax} \texttt{cbsd} \emph{registerW} = \emph{registerZ}

\paragraph{Description} Bit set population count value of the \emph{registerZ} is stored into the \emph{registerW}.


\paragraph{Execution} {Reservation table \textsf{ALU\_TINY} (Section~\ref{sec:reservation-ALU-TINY})}
~
{\begin{lstlisting}
stage RR:
  new argument2 = GPR[registerZ];
  new result1 = _ZX_64(_CBS_64(argument2));
stage E1:
  GPR[registerW] = result1;
\end{lstlisting}}
\newpage
\subsubsection[{\small CBSHQ: Count Bit Sets Half Word Quadruple}]{CBSHQ\hfill Count Bit Sets Half Word Quadruple}
\label{instr:CBSHQ}\index{cbshq}
 
\paragraph{Encoding} Format \textsf{ALU\_HQBWRE} (Section~\ref{sec:format-ALU-HQBWRE}), \verb|alucode8 = 0110|\\

\bitfields{ 1/P, 2/11, 1/1, 4/1111, 5/registerWe, 1/0, 2/10, 3/101, 1/1, 4/0110, 1/--, 1/--, 5/registerZe, 1/0 }


\paragraph{Syntax} \texttt{cbshq} \emph{registerWe} = \emph{registerZe}

\paragraph{Description} Bit set population count values of the \emph{registerZe} interpreted as half word quadruple are stored into the \emph{registerWe}.


\paragraph{Execution} {Reservation table \textsf{ALU\_LITE} (Section~\ref{sec:reservation-ALU-LITE})}
~
{\begin{lstlisting}
stage RR:
  new argument2 = GPR[registerZe<<1];
  for (I = 0; I < 4; I += 1) {
    result1.16[I] = _ZX_16(_CBS_16(argument2.16[I]))
  };
stage E1:
  GPR[registerWe<<1] = result1;
\end{lstlisting}}
\newpage

\paragraph{Encoding} Format \textsf{ALU\_HQBWRO} (Section~\ref{sec:format-ALU-HQBWRO}), \verb|alucode8 = 0110|\\

\bitfields{ 1/P, 2/11, 1/1, 4/1111, 5/registerWo, 1/1, 2/10, 3/101, 1/1, 4/0110, 1/--, 1/--, 5/registerZo, 1/0 }


\paragraph{Syntax} \texttt{cbshq} \emph{registerWo} = \emph{registerZo}

\paragraph{Description} Bit set population count values of the \emph{registerZo} interpreted as half word quadruple are stored into the \emph{registerWo}.


\paragraph{Execution} {Reservation table \textsf{ALU\_LITE} (Section~\ref{sec:reservation-ALU-LITE})}
~
{\begin{lstlisting}
stage RR:
  new argument2 = GPR[(registerZo<<1)+1];
  for (I = 0; I < 4; I += 1) {
    result1.16[I] = _ZX_16(_CBS_16(argument2.16[I]))
  };
stage E1:
  GPR[(registerWo<<1)+1] = result1;
\end{lstlisting}}
\newpage
\subsubsection[{\small CBSW: Count Bit Set Word}]{CBSW\hfill Count Bit Set Word}
\label{instr:CBSW}\index{cbsw}
 
\paragraph{Encoding} Format \textsf{ALU\_BWRW} (Section~\ref{sec:format-ALU-BWRW}), \verb|alucode8 = 0110|\\

\bitfields{ 1/P, 2/11, 1/0, 4/1111, 6/registerW, 2/11, 3/101, 1/signextw, 4/0110, 1/0, 1/0, 6/registerZ }


\paragraph{Syntax} \texttt{cbsw} \emph{registerW} = \emph{registerZ}

\paragraph{Description} Bit set population count value of the \emph{registerZ} is stored into the \emph{registerW}.


\paragraph{Execution} {Reservation table \textsf{ALU\_TINY} (Section~\ref{sec:reservation-ALU-TINY})}
~
{\begin{lstlisting}
stage RR:
  new argument2 = GPR[registerZ];
  new result1 = _ZX_32(_CBS_32(argument2));
stage E1:
  GPR[registerW] = _ZX_32(result1);
\end{lstlisting}}
\newpage
\subsubsection[{\small CBSWP: Count Bit Sets Word Pair}]{CBSWP\hfill Count Bit Sets Word Pair}
\label{instr:CBSWP}\index{cbswp}
 
\paragraph{Encoding} Format \textsf{ALU\_WPBWRE} (Section~\ref{sec:format-ALU-WPBWRE}), \verb|alucode8 = 0110|\\

\bitfields{ 1/P, 2/11, 1/1, 4/1111, 5/registerWe, 1/0, 2/10, 3/101, 1/0, 4/0110, 1/--, 1/--, 5/registerZe, 1/1 }


\paragraph{Syntax} \texttt{cbswp} \emph{registerWe} = \emph{registerZe}

\paragraph{Description} Bit set population count values of the \emph{registerZe} interpreted as word pair are stored into the \emph{registerWe}.


\paragraph{Execution} {Reservation table \textsf{ALU\_LITE} (Section~\ref{sec:reservation-ALU-LITE})}
~
{\begin{lstlisting}
stage RR:
  new argument2 = GPR[registerZe<<1];
  for (I = 0; I < 2; I += 1) {
    result1.32[I] = _ZX_32(_CBS_32(argument2.32[I]))
  };
stage E1:
  GPR[registerWe<<1] = result1;
\end{lstlisting}}
\newpage

\paragraph{Encoding} Format \textsf{ALU\_WPBWRO} (Section~\ref{sec:format-ALU-WPBWRO}), \verb|alucode8 = 0110|\\

\bitfields{ 1/P, 2/11, 1/1, 4/1111, 5/registerWo, 1/1, 2/10, 3/101, 1/0, 4/0110, 1/--, 1/--, 5/registerZo, 1/1 }


\paragraph{Syntax} \texttt{cbswp} \emph{registerWo} = \emph{registerZo}

\paragraph{Description} Bit set population count values of the \emph{registerZo} interpreted as word pair are stored into the \emph{registerWo}.


\paragraph{Execution} {Reservation table \textsf{ALU\_LITE} (Section~\ref{sec:reservation-ALU-LITE})}
~
{\begin{lstlisting}
stage RR:
  new argument2 = GPR[(registerZo<<1)+1];
  for (I = 0; I < 2; I += 1) {
    result1.32[I] = _ZX_32(_CBS_32(argument2.32[I]))
  };
stage E1:
  GPR[(registerWo<<1)+1] = result1;
\end{lstlisting}}
\newpage
\subsubsection[{\small CLSD: Count Leading Sign Double Word}]{CLSD\hfill Count Leading Sign Double Word}
\label{instr:CLSD}\index{clsd}
 
\paragraph{Encoding} Format \textsf{ALU\_DBWR} (Section~\ref{sec:format-ALU-DBWR}), \verb|alucode8 = 0101|\\

\bitfields{ 1/P, 2/11, 1/0, 4/1111, 6/registerW, 2/10, 3/101, 1/0, 4/0101, 1/0, 1/0, 6/registerZ }


\paragraph{Syntax} \texttt{clsd} \emph{registerW} = \emph{registerZ}

\paragraph{Description} The leading sign count value of the \emph{registerZ} is stored into the \emph{registerW}. When applied to zero or -1, the result is 63.


\paragraph{Execution} {Reservation table \textsf{ALU\_TINY} (Section~\ref{sec:reservation-ALU-TINY})}
~
{\begin{lstlisting}
stage RR:
  new argument2 = GPR[registerZ];
  new result1 = _ZX_64(_CLS_64(argument2));
stage E1:
  GPR[registerW] = result1;
\end{lstlisting}}
\newpage
\subsubsection[{\small CLSHQ: Count Leading Signs Half Word Quadruple}]{CLSHQ\hfill Count Leading Signs Half Word Quadruple}
\label{instr:CLSHQ}\index{clshq}
 
\paragraph{Encoding} Format \textsf{ALU\_HQBWRE} (Section~\ref{sec:format-ALU-HQBWRE}), \verb|alucode8 = 0101|\\

\bitfields{ 1/P, 2/11, 1/1, 4/1111, 5/registerWe, 1/0, 2/10, 3/101, 1/1, 4/0101, 1/--, 1/--, 5/registerZe, 1/0 }


\paragraph{Syntax} \texttt{clshq} \emph{registerWe} = \emph{registerZe}

\paragraph{Description} The leading sign count values of the \emph{registerZe} interpreted as half word quadruple are stored into the \emph{registerWe}. When applied to zero or -1 at the half word level, the result is 31.


\paragraph{Execution} {Reservation table \textsf{ALU\_LITE} (Section~\ref{sec:reservation-ALU-LITE})}
~
{\begin{lstlisting}
stage RR:
  new argument2 = GPR[registerZe<<1];
  for (I = 0; I < 4; I += 1) {
    result1.16[I] = _ZX_16(_CLS_16(argument2.16[I]))
  };
stage E1:
  GPR[registerWe<<1] = result1;
\end{lstlisting}}
\newpage

\paragraph{Encoding} Format \textsf{ALU\_HQBWRO} (Section~\ref{sec:format-ALU-HQBWRO}), \verb|alucode8 = 0101|\\

\bitfields{ 1/P, 2/11, 1/1, 4/1111, 5/registerWo, 1/1, 2/10, 3/101, 1/1, 4/0101, 1/--, 1/--, 5/registerZo, 1/0 }


\paragraph{Syntax} \texttt{clshq} \emph{registerWo} = \emph{registerZo}

\paragraph{Description} The leading sign count values of the \emph{registerZo} interpreted as half word quadruple are stored into the \emph{registerWo}. When applied to zero or -1 at the half word level, the result is 31.


\paragraph{Execution} {Reservation table \textsf{ALU\_LITE} (Section~\ref{sec:reservation-ALU-LITE})}
~
{\begin{lstlisting}
stage RR:
  new argument2 = GPR[(registerZo<<1)+1];
  for (I = 0; I < 4; I += 1) {
    result1.16[I] = _ZX_16(_CLS_16(argument2.16[I]))
  };
stage E1:
  GPR[(registerWo<<1)+1] = result1;
\end{lstlisting}}
\newpage
\subsubsection[{\small CLSW: Count Leading Sign Word}]{CLSW\hfill Count Leading Sign Word}
\label{instr:CLSW}\index{clsw}
 
\paragraph{Encoding} Format \textsf{ALU\_BWRW} (Section~\ref{sec:format-ALU-BWRW}), \verb|alucode8 = 0101|\\

\bitfields{ 1/P, 2/11, 1/0, 4/1111, 6/registerW, 2/11, 3/101, 1/signextw, 4/0101, 1/0, 1/0, 6/registerZ }


\paragraph{Syntax} \texttt{clsw} \emph{registerW} = \emph{registerZ}

\paragraph{Description} The leading sign count value of the \emph{registerZ} is stored into the \emph{registerW}. When applied to zero or -1, the result is 31.


\paragraph{Execution} {Reservation table \textsf{ALU\_TINY} (Section~\ref{sec:reservation-ALU-TINY})}
~
{\begin{lstlisting}
stage RR:
  new argument2 = GPR[registerZ];
  new result1 = _ZX_32(_CLS_32(argument2));
stage E1:
  GPR[registerW] = _ZX_32(result1);
\end{lstlisting}}
\newpage
\subsubsection[{\small CLSWP: Count Leading Signs Word Pair}]{CLSWP\hfill Count Leading Signs Word Pair}
\label{instr:CLSWP}\index{clswp}
 
\paragraph{Encoding} Format \textsf{ALU\_WPBWRE} (Section~\ref{sec:format-ALU-WPBWRE}), \verb|alucode8 = 0101|\\

\bitfields{ 1/P, 2/11, 1/1, 4/1111, 5/registerWe, 1/0, 2/10, 3/101, 1/0, 4/0101, 1/--, 1/--, 5/registerZe, 1/1 }


\paragraph{Syntax} \texttt{clswp} \emph{registerWe} = \emph{registerZe}

\paragraph{Description} The leading sign count values of the \emph{registerZe} interpreted as word pair are stored into the \emph{registerWe}. When applied to zero or -1 at the word level, the result is 31.


\paragraph{Execution} {Reservation table \textsf{ALU\_LITE} (Section~\ref{sec:reservation-ALU-LITE})}
~
{\begin{lstlisting}
stage RR:
  new argument2 = GPR[registerZe<<1];
  for (I = 0; I < 2; I += 1) {
    result1.32[I] = _ZX_32(_CLS_32(argument2.32[I]))
  };
stage E1:
  GPR[registerWe<<1] = result1;
\end{lstlisting}}
\newpage

\paragraph{Encoding} Format \textsf{ALU\_WPBWRO} (Section~\ref{sec:format-ALU-WPBWRO}), \verb|alucode8 = 0101|\\

\bitfields{ 1/P, 2/11, 1/1, 4/1111, 5/registerWo, 1/1, 2/10, 3/101, 1/0, 4/0101, 1/--, 1/--, 5/registerZo, 1/1 }


\paragraph{Syntax} \texttt{clswp} \emph{registerWo} = \emph{registerZo}

\paragraph{Description} The leading sign count values of the \emph{registerZo} interpreted as word pair are stored into the \emph{registerWo}. When applied to zero or -1 at the word level, the result is 31.


\paragraph{Execution} {Reservation table \textsf{ALU\_LITE} (Section~\ref{sec:reservation-ALU-LITE})}
~
{\begin{lstlisting}
stage RR:
  new argument2 = GPR[(registerZo<<1)+1];
  for (I = 0; I < 2; I += 1) {
    result1.32[I] = _ZX_32(_CLS_32(argument2.32[I]))
  };
stage E1:
  GPR[(registerWo<<1)+1] = result1;
\end{lstlisting}}
\newpage
\subsubsection[{\small CLZD: Count Leading Zero Double Word}]{CLZD\hfill Count Leading Zero Double Word}
\label{instr:CLZD}\index{clzd}
 
\paragraph{Encoding} Format \textsf{ALU\_DBWR} (Section~\ref{sec:format-ALU-DBWR}), \verb|alucode8 = 0100|\\

\bitfields{ 1/P, 2/11, 1/0, 4/1111, 6/registerW, 2/10, 3/101, 1/0, 4/0100, 1/0, 1/0, 6/registerZ }


\paragraph{Syntax} \texttt{clzd} \emph{registerW} = \emph{registerZ}

\paragraph{Description} The leading zero count value of the \emph{registerZ} is stored into the \emph{registerW}. When applied to zero, the result is 64.


\paragraph{Execution} {Reservation table \textsf{ALU\_TINY} (Section~\ref{sec:reservation-ALU-TINY})}
~
{\begin{lstlisting}
stage RR:
  new argument2 = GPR[registerZ];
  new result1 = _ZX_64(_CLZ_64(argument2));
stage E1:
  GPR[registerW] = result1;
\end{lstlisting}}
\newpage
\subsubsection[{\small CLZHQ: Count Leading Zeroes Half Word Quadruple}]{CLZHQ\hfill Count Leading Zeroes Half Word Quadruple}
\label{instr:CLZHQ}\index{clzhq}
 
\paragraph{Encoding} Format \textsf{ALU\_HQBWRE} (Section~\ref{sec:format-ALU-HQBWRE}), \verb|alucode8 = 0100|\\

\bitfields{ 1/P, 2/11, 1/1, 4/1111, 5/registerWe, 1/0, 2/10, 3/101, 1/1, 4/0100, 1/--, 1/--, 5/registerZe, 1/0 }


\paragraph{Syntax} \texttt{clzhq} \emph{registerWe} = \emph{registerZe}

\paragraph{Description} The leading zero count values of the \emph{registerZe} interpreted as a half word quadruple are stored into the \emph{registerWe}. When applied to zero at the half word level, the result is 16.


\paragraph{Execution} {Reservation table \textsf{ALU\_LITE} (Section~\ref{sec:reservation-ALU-LITE})}
~
{\begin{lstlisting}
stage RR:
  new argument2 = GPR[registerZe<<1];
  for (I = 0; I < 4; I += 1) {
    result1.16[I] = _ZX_16(_CLZ_16(argument2.16[I]))
  };
stage E1:
  GPR[registerWe<<1] = result1;
\end{lstlisting}}
\newpage

\paragraph{Encoding} Format \textsf{ALU\_HQBWRO} (Section~\ref{sec:format-ALU-HQBWRO}), \verb|alucode8 = 0100|\\

\bitfields{ 1/P, 2/11, 1/1, 4/1111, 5/registerWo, 1/1, 2/10, 3/101, 1/1, 4/0100, 1/--, 1/--, 5/registerZo, 1/0 }


\paragraph{Syntax} \texttt{clzhq} \emph{registerWo} = \emph{registerZo}

\paragraph{Description} The leading zero count values of the \emph{registerZo} interpreted as a half word quadruple are stored into the \emph{registerWo}. When applied to zero at the half word level, the result is 16.


\paragraph{Execution} {Reservation table \textsf{ALU\_LITE} (Section~\ref{sec:reservation-ALU-LITE})}
~
{\begin{lstlisting}
stage RR:
  new argument2 = GPR[(registerZo<<1)+1];
  for (I = 0; I < 4; I += 1) {
    result1.16[I] = _ZX_16(_CLZ_16(argument2.16[I]))
  };
stage E1:
  GPR[(registerWo<<1)+1] = result1;
\end{lstlisting}}
\newpage
\subsubsection[{\small CLZW: Count Leading Zero Word}]{CLZW\hfill Count Leading Zero Word}
\label{instr:CLZW}\index{clzw}
 
\paragraph{Encoding} Format \textsf{ALU\_BWRW} (Section~\ref{sec:format-ALU-BWRW}), \verb|alucode8 = 0100|\\

\bitfields{ 1/P, 2/11, 1/0, 4/1111, 6/registerW, 2/11, 3/101, 1/signextw, 4/0100, 1/0, 1/0, 6/registerZ }


\paragraph{Syntax} \texttt{clzw} \emph{registerW} = \emph{registerZ}

\paragraph{Description} The leading zero count value of the \emph{registerZ} is stored into the \emph{registerW}. When applied to zero, the result is 32.


\paragraph{Execution} {Reservation table \textsf{ALU\_TINY} (Section~\ref{sec:reservation-ALU-TINY})}
~
{\begin{lstlisting}
stage RR:
  new argument2 = GPR[registerZ];
  new result1 = _ZX_32(_CLZ_32(argument2));
stage E1:
  GPR[registerW] = _ZX_32(result1);
\end{lstlisting}}
\newpage
\subsubsection[{\small CLZWP: Count Leading Zeroes Word Pair}]{CLZWP\hfill Count Leading Zeroes Word Pair}
\label{instr:CLZWP}\index{clzwp}
 
\paragraph{Encoding} Format \textsf{ALU\_WPBWRE} (Section~\ref{sec:format-ALU-WPBWRE}), \verb|alucode8 = 0100|\\

\bitfields{ 1/P, 2/11, 1/1, 4/1111, 5/registerWe, 1/0, 2/10, 3/101, 1/0, 4/0100, 1/--, 1/--, 5/registerZe, 1/1 }


\paragraph{Syntax} \texttt{clzwp} \emph{registerWe} = \emph{registerZe}

\paragraph{Description} The leading zero count values of the \emph{registerZe} interpreted as a word pair are stored into the \emph{registerWe}. When applied to zero at the word level, the result is 32.


\paragraph{Execution} {Reservation table \textsf{ALU\_LITE} (Section~\ref{sec:reservation-ALU-LITE})}
~
{\begin{lstlisting}
stage RR:
  new argument2 = GPR[registerZe<<1];
  for (I = 0; I < 2; I += 1) {
    result1.32[I] = _ZX_32(_CLZ_32(argument2.32[I]))
  };
stage E1:
  GPR[registerWe<<1] = result1;
\end{lstlisting}}
\newpage

\paragraph{Encoding} Format \textsf{ALU\_WPBWRO} (Section~\ref{sec:format-ALU-WPBWRO}), \verb|alucode8 = 0100|\\

\bitfields{ 1/P, 2/11, 1/1, 4/1111, 5/registerWo, 1/1, 2/10, 3/101, 1/0, 4/0100, 1/--, 1/--, 5/registerZo, 1/1 }


\paragraph{Syntax} \texttt{clzwp} \emph{registerWo} = \emph{registerZo}

\paragraph{Description} The leading zero count values of the \emph{registerZo} interpreted as a word pair are stored into the \emph{registerWo}. When applied to zero at the word level, the result is 32.


\paragraph{Execution} {Reservation table \textsf{ALU\_LITE} (Section~\ref{sec:reservation-ALU-LITE})}
~
{\begin{lstlisting}
stage RR:
  new argument2 = GPR[(registerZo<<1)+1];
  for (I = 0; I < 2; I += 1) {
    result1.32[I] = _ZX_32(_CLZ_32(argument2.32[I]))
  };
stage E1:
  GPR[(registerWo<<1)+1] = result1;
\end{lstlisting}}
\newpage
\subsubsection[{\small CMOVEBO: Conditional Move Byte Octuple}]{CMOVEBO\hfill Conditional Move Byte Octuple}
\label{instr:CMOVEBO}\index{cmovebo}
 
\paragraph{Encoding} Format \textsf{ALU\_BOCMWRRE} (Section~\ref{sec:format-ALU-BOCMWRRE}), \verb|exubit24 = 0|\\

\bitfields{ 1/P, 2/11, 1/1, 3/lanecond, 1/0, 5/registerWe, 1/0, 2/10, 3/011, 1/1, 5/registerYe, 1/--, 5/registerZe, 1/1 }


\paragraph{Syntax} \texttt{cmovebo}\emph{lanecond} \emph{registerZe}? \emph{registerWe} = \emph{registerYe}

\paragraph{Description} The \emph{registerWe}, \emph{registerZe} and the \emph{registerYe} are interpreted as eight bytes packed into 64 bits. Each \emph{registerZe} word is tested using the condition \emph{lanecond}. If true, the corresponding byte of the \emph{registerYe} is stored into the corresponding word of the \emph{registerWe}.


\paragraph{Modifier(s)}
\noindent\begin{itemize}
\item lanecond (cf. lanecond in section \ref{par:modifier-lanecond} on page \pageref{par:modifier-lanecond})
\end{itemize}

\paragraph{Execution} {Reservation table \textsf{ALU\_LITE} (Section~\ref{sec:reservation-ALU-LITE})}
~
{\begin{lstlisting}
stage ID:
  new argument4 = _ZX_3(lanecond);
stage RR:
  new argument3 = GPR[registerYe<<1];
  new argument2 = GPR[registerZe<<1];
stage E1:
  if (lanecond_8(argument4, argument2.8[0])) {
    GPR[registerWe<<1] = argument3;
  }
  if (lanecond_8(argument4, argument2.8[1])) {
    GPR[registerWe<<1] = argument3;
  }
  if (lanecond_8(argument4, argument2.8[2])) {
    GPR[registerWe<<1] = argument3;
  }
  if (lanecond_8(argument4, argument2.8[3])) {
    GPR[registerWe<<1] = argument3;
  }
  if (lanecond_8(argument4, argument2.8[4])) {
    GPR[registerWe<<1] = argument3;
  }
  if (lanecond_8(argument4, argument2.8[5])) {
    GPR[registerWe<<1] = argument3;
  }
  if (lanecond_8(argument4, argument2.8[6])) {
    GPR[registerWe<<1] = argument3;
  }
  if (lanecond_8(argument4, argument2.8[7])) {
    GPR[registerWe<<1] = argument3;
  }
\end{lstlisting}}
\newpage

\paragraph{Encoding} Format \textsf{ALU\_BOCMWRRO} (Section~\ref{sec:format-ALU-BOCMWRRO}), \verb|exubit24 = 0|\\

\bitfields{ 1/P, 2/11, 1/1, 3/lanecond, 1/0, 5/registerWo, 1/1, 2/10, 3/011, 1/1, 5/registerYo, 1/--, 5/registerZo, 1/1 }


\paragraph{Syntax} \texttt{cmovebo}\emph{lanecond} \emph{registerZo}? \emph{registerWo} = \emph{registerYo}

\paragraph{Description} The \emph{registerWo}, \emph{registerZo} and the \emph{registerYo} are interpreted as eight bytes packed into 64 bits. Each \emph{registerZo} word is tested using the condition \emph{lanecond}. If true, the corresponding byte of the \emph{registerYo} is stored into the corresponding word of the \emph{registerWo}.


\paragraph{Modifier(s)}
\noindent\begin{itemize}
\item lanecond (cf. lanecond in section \ref{par:modifier-lanecond} on page \pageref{par:modifier-lanecond})
\end{itemize}

\paragraph{Execution} {Reservation table \textsf{ALU\_LITE} (Section~\ref{sec:reservation-ALU-LITE})}
~
{\begin{lstlisting}
stage ID:
  new argument4 = _ZX_3(lanecond);
stage RR:
  new argument3 = GPR[(registerYo<<1)+1];
  new argument2 = GPR[(registerZo<<1)+1];
stage E1:
  if (lanecond_8(argument4, argument2.8[0])) {
    GPR[(registerWo<<1)+1] = argument3;
  }
  if (lanecond_8(argument4, argument2.8[1])) {
    GPR[(registerWo<<1)+1] = argument3;
  }
  if (lanecond_8(argument4, argument2.8[2])) {
    GPR[(registerWo<<1)+1] = argument3;
  }
  if (lanecond_8(argument4, argument2.8[3])) {
    GPR[(registerWo<<1)+1] = argument3;
  }
  if (lanecond_8(argument4, argument2.8[4])) {
    GPR[(registerWo<<1)+1] = argument3;
  }
  if (lanecond_8(argument4, argument2.8[5])) {
    GPR[(registerWo<<1)+1] = argument3;
  }
  if (lanecond_8(argument4, argument2.8[6])) {
    GPR[(registerWo<<1)+1] = argument3;
  }
  if (lanecond_8(argument4, argument2.8[7])) {
    GPR[(registerWo<<1)+1] = argument3;
  }
\end{lstlisting}}
\newpage

\paragraph{Encoding} Format \textsf{ALU\_BOCMWRRE.M} (Section~\ref{sec:format-ALU-BOCMWRRE.M}), \verb|exubit24 = 0|\\

\bitfields{ 1/P, 2/00, 2/-- --, 27/upper27, 1/1, 2/11, 1/1, 3/lanecond, 1/0, 5/registerWe, 1/0, 2/10, 3/011, 1/1, 1/splat32, 5/lower5, 5/registerZe, 1/1 }


\paragraph{Syntax} \texttt{cmovebo}\emph{lanecond} \emph{registerZe}? \emph{registerWe} = \emph{upper27\_\discretionary{}{}{}lower5}\emph{splat32}

\paragraph{Description} The \emph{registerWe}, \emph{registerZe} and the \emph{upper27\_\discretionary{}{}{}lower5} are interpreted as eight bytes packed into 64 bits. Each \emph{registerZe} word is tested using the condition \emph{lanecond}. If true, the corresponding byte of the \emph{upper27\_\discretionary{}{}{}lower5} is stored into the corresponding word of the \emph{registerWe}.


\paragraph{Modifier(s)}
\noindent\begin{itemize}
\item lanecond (cf. lanecond in section \ref{par:modifier-lanecond} on page \pageref{par:modifier-lanecond})
\item splat32 (cf. splat32 in section \ref{par:modifier-splat32} on page \pageref{par:modifier-splat32})
\end{itemize}

\paragraph{Execution} {Reservation table \textsf{ALU\_LITE.X} (Section~\ref{sec:reservation-ALU-LITE.X})}
~
{\begin{lstlisting}
stage ID:
  new splat32 = _ZX_1(splat32);
  new argument4 = _ZX_3(lanecond);
stage RR:
  new argument3 = splat32 ? (_WX_32(upper27_lower5) << 32) | _ZX_32(_WX_32(upper27_lower5)) : _SX_32(_WX_32(upper27_lower5));
  new argument2 = GPR[registerZe<<1];
stage E1:
  if (lanecond_8(argument4, argument2.8[0])) {
    GPR[registerWe<<1] = argument3;
  }
  if (lanecond_8(argument4, argument2.8[1])) {
    GPR[registerWe<<1] = argument3;
  }
  if (lanecond_8(argument4, argument2.8[2])) {
    GPR[registerWe<<1] = argument3;
  }
  if (lanecond_8(argument4, argument2.8[3])) {
    GPR[registerWe<<1] = argument3;
  }
  if (lanecond_8(argument4, argument2.8[4])) {
    GPR[registerWe<<1] = argument3;
  }
  if (lanecond_8(argument4, argument2.8[5])) {
    GPR[registerWe<<1] = argument3;
  }
  if (lanecond_8(argument4, argument2.8[6])) {
    GPR[registerWe<<1] = argument3;
  }
  if (lanecond_8(argument4, argument2.8[7])) {
    GPR[registerWe<<1] = argument3;
  }
\end{lstlisting}}
\newpage

\paragraph{Encoding} Format \textsf{ALU\_BOCMWRRO.M} (Section~\ref{sec:format-ALU-BOCMWRRO.M}), \verb|exubit24 = 0|\\

\bitfields{ 1/P, 2/00, 2/-- --, 27/upper27, 1/1, 2/11, 1/1, 3/lanecond, 1/0, 5/registerWo, 1/1, 2/10, 3/011, 1/1, 1/splat32, 5/lower5, 5/registerZo, 1/1 }


\paragraph{Syntax} \texttt{cmovebo}\emph{lanecond} \emph{registerZo}? \emph{registerWo} = \emph{upper27\_\discretionary{}{}{}lower5}\emph{splat32}

\paragraph{Description} The \emph{registerWo}, \emph{registerZo} and the \emph{upper27\_\discretionary{}{}{}lower5} are interpreted as eight bytes packed into 64 bits. Each \emph{registerZo} word is tested using the condition \emph{lanecond}. If true, the corresponding byte of the \emph{upper27\_\discretionary{}{}{}lower5} is stored into the corresponding word of the \emph{registerWo}.


\paragraph{Modifier(s)}
\noindent\begin{itemize}
\item lanecond (cf. lanecond in section \ref{par:modifier-lanecond} on page \pageref{par:modifier-lanecond})
\item splat32 (cf. splat32 in section \ref{par:modifier-splat32} on page \pageref{par:modifier-splat32})
\end{itemize}

\paragraph{Execution} {Reservation table \textsf{ALU\_LITE.X} (Section~\ref{sec:reservation-ALU-LITE.X})}
~
{\begin{lstlisting}
stage ID:
  new splat32 = _ZX_1(splat32);
  new argument4 = _ZX_3(lanecond);
stage RR:
  new argument3 = splat32 ? (_WX_32(upper27_lower5) << 32) | _ZX_32(_WX_32(upper27_lower5)) : _SX_32(_WX_32(upper27_lower5));
  new argument2 = GPR[(registerZo<<1)+1];
stage E1:
  if (lanecond_8(argument4, argument2.8[0])) {
    GPR[(registerWo<<1)+1] = argument3;
  }
  if (lanecond_8(argument4, argument2.8[1])) {
    GPR[(registerWo<<1)+1] = argument3;
  }
  if (lanecond_8(argument4, argument2.8[2])) {
    GPR[(registerWo<<1)+1] = argument3;
  }
  if (lanecond_8(argument4, argument2.8[3])) {
    GPR[(registerWo<<1)+1] = argument3;
  }
  if (lanecond_8(argument4, argument2.8[4])) {
    GPR[(registerWo<<1)+1] = argument3;
  }
  if (lanecond_8(argument4, argument2.8[5])) {
    GPR[(registerWo<<1)+1] = argument3;
  }
  if (lanecond_8(argument4, argument2.8[6])) {
    GPR[(registerWo<<1)+1] = argument3;
  }
  if (lanecond_8(argument4, argument2.8[7])) {
    GPR[(registerWo<<1)+1] = argument3;
  }
\end{lstlisting}}
\newpage
\subsubsection[{\small CMOVED: Conditional Move Double Word}]{CMOVED\hfill Conditional Move Double Word}
\label{instr:CMOVED}\index{cmoved}
 
\paragraph{Encoding} Format \textsf{ALU\_DCMWRR} (Section~\ref{sec:format-ALU-DCMWRR}), \verb|steering = 11|\\

\bitfields{ 1/P, 2/11, 1/0, 4/cmovecond, 6/registerW, 2/10, 3/011, 1/0, 6/registerY, 6/registerZ }


\paragraph{Syntax} \texttt{cmoved}\emph{cmovecond} \emph{registerZ}? \emph{registerW} = \emph{registerY}

\paragraph{Description} The \emph{registerZ} double word is tested using the condition \emph{cmovecond}. If true, the \emph{registerY} double word is stored into the \emph{registerW} double word. Else the \emph{registerW} double word is unchanged.


\paragraph{Modifier(s)}
\noindent\begin{itemize}
\item cmovecond (cf. bcucond in section \ref{par:modifier-bcucond} on page \pageref{par:modifier-bcucond})
\end{itemize}

\paragraph{Execution} {Reservation table \textsf{ALU\_LITE} (Section~\ref{sec:reservation-ALU-LITE})}
~
{\begin{lstlisting}
stage ID:
  new argument4 = _ZX_4(cmovecond);
stage RR:
  new argument3 = GPR[registerY];
  new argument2 = GPR[registerZ];
stage E1:
  if (bcucond(argument4, argument2)) {
    GPR[registerW] = argument3;
  }
\end{lstlisting}}
\newpage

\paragraph{Encoding} Format \textsf{ALU\_DCMWRR.M} (Section~\ref{sec:format-ALU-DCMWRR.M}), \verb|steering = 11|\\

\bitfields{ 1/P, 2/00, 2/-- --, 27/upper27, 1/1, 2/11, 1/0, 4/cmovecond, 6/registerW, 2/10, 3/011, 1/0, 1/splat32, 5/lower5, 6/registerZ }


\paragraph{Syntax} \texttt{cmoved}\emph{cmovecond} \emph{registerZ}? \emph{registerW} = \emph{upper27\_\discretionary{}{}{}lower5}\emph{splat32}

\paragraph{Description} The \emph{registerZ} double word is tested using the condition \emph{cmovecond}. If true, the \emph{upper27\_\discretionary{}{}{}lower5} double word is stored into the \emph{registerW} double word. Else the \emph{registerW} double word is unchanged.


\paragraph{Modifier(s)}
\noindent\begin{itemize}
\item cmovecond (cf. bcucond in section \ref{par:modifier-bcucond} on page \pageref{par:modifier-bcucond})
\item splat32 (cf. splat32 in section \ref{par:modifier-splat32} on page \pageref{par:modifier-splat32})
\end{itemize}

\paragraph{Execution} {Reservation table \textsf{ALU\_LITE.X} (Section~\ref{sec:reservation-ALU-LITE.X})}
~
{\begin{lstlisting}
stage ID:
  new splat32 = _ZX_1(splat32);
  new argument4 = _ZX_4(cmovecond);
stage RR:
  new argument3 = splat32 ? (_WX_32(upper27_lower5) << 32) | _ZX_32(_WX_32(upper27_lower5)) : _SX_32(_WX_32(upper27_lower5));
  new argument2 = GPR[registerZ];
stage E1:
  if (bcucond(argument4, argument2)) {
    GPR[registerW] = argument3;
  }
\end{lstlisting}}
\newpage
\subsubsection[{\small CMOVEHQ: Conditional Move Half Word Quadruple}]{CMOVEHQ\hfill Conditional Move Half Word Quadruple}
\label{instr:CMOVEHQ}\index{cmovehq}
 
\paragraph{Encoding} Format \textsf{ALU\_HQCMWRRE} (Section~\ref{sec:format-ALU-HQCMWRRE}), \verb|exubit24 = 0|\\

\bitfields{ 1/P, 2/11, 1/1, 3/lanecond, 1/0, 5/registerWe, 1/0, 2/10, 3/011, 1/1, 5/registerYe, 1/--, 5/registerZe, 1/0 }


\paragraph{Syntax} \texttt{cmovehq}\emph{lanecond} \emph{registerZe}? \emph{registerWe} = \emph{registerYe}

\paragraph{Description} The \emph{registerWe}, \emph{registerZe} and the \emph{registerYe} are interpreted as four half words packed into 64 bits. Each \emph{registerZe} half word is tested using the condition \emph{lanecond}. If true, the corresponding half word of the \emph{registerYe} is stored into the corresponding half word of the \emph{registerWe}.


\paragraph{Modifier(s)}
\noindent\begin{itemize}
\item lanecond (cf. lanecond in section \ref{par:modifier-lanecond} on page \pageref{par:modifier-lanecond})
\end{itemize}

\paragraph{Execution} {Reservation table \textsf{ALU\_LITE} (Section~\ref{sec:reservation-ALU-LITE})}
~
{\begin{lstlisting}
stage ID:
  new argument4 = _ZX_3(lanecond);
stage RR:
  new argument3 = GPR[registerYe<<1];
  new argument2 = GPR[registerZe<<1];
stage E1:
  if (lanecond_16(argument4, argument2.16[0])) {
    GPR[registerWe<<1] = argument3;
  }
  if (lanecond_16(argument4, argument2.16[1])) {
    GPR[registerWe<<1] = argument3;
  }
  if (lanecond_16(argument4, argument2.16[2])) {
    GPR[registerWe<<1] = argument3;
  }
  if (lanecond_16(argument4, argument2.16[3])) {
    GPR[registerWe<<1] = argument3;
  }
\end{lstlisting}}
\newpage

\paragraph{Encoding} Format \textsf{ALU\_HQCMWRRO} (Section~\ref{sec:format-ALU-HQCMWRRO}), \verb|exubit24 = 0|\\

\bitfields{ 1/P, 2/11, 1/1, 3/lanecond, 1/0, 5/registerWo, 1/1, 2/10, 3/011, 1/1, 5/registerYo, 1/--, 5/registerZo, 1/0 }


\paragraph{Syntax} \texttt{cmovehq}\emph{lanecond} \emph{registerZo}? \emph{registerWo} = \emph{registerYo}

\paragraph{Description} The \emph{registerWo}, \emph{registerZo} and the \emph{registerYo} are interpreted as four half words packed into 64 bits. Each \emph{registerZo} half word is tested using the condition \emph{lanecond}. If true, the corresponding half word of the \emph{registerYo} is stored into the corresponding half word of the \emph{registerWo}.


\paragraph{Modifier(s)}
\noindent\begin{itemize}
\item lanecond (cf. lanecond in section \ref{par:modifier-lanecond} on page \pageref{par:modifier-lanecond})
\end{itemize}

\paragraph{Execution} {Reservation table \textsf{ALU\_LITE} (Section~\ref{sec:reservation-ALU-LITE})}
~
{\begin{lstlisting}
stage ID:
  new argument4 = _ZX_3(lanecond);
stage RR:
  new argument3 = GPR[(registerYo<<1)+1];
  new argument2 = GPR[(registerZo<<1)+1];
stage E1:
  if (lanecond_16(argument4, argument2.16[0])) {
    GPR[(registerWo<<1)+1] = argument3;
  }
  if (lanecond_16(argument4, argument2.16[1])) {
    GPR[(registerWo<<1)+1] = argument3;
  }
  if (lanecond_16(argument4, argument2.16[2])) {
    GPR[(registerWo<<1)+1] = argument3;
  }
  if (lanecond_16(argument4, argument2.16[3])) {
    GPR[(registerWo<<1)+1] = argument3;
  }
\end{lstlisting}}
\newpage

\paragraph{Encoding} Format \textsf{ALU\_HQCMWRRE.M} (Section~\ref{sec:format-ALU-HQCMWRRE.M}), \verb|exubit24 = 0|\\

\bitfields{ 1/P, 2/00, 2/-- --, 27/upper27, 1/1, 2/11, 1/1, 3/lanecond, 1/0, 5/registerWe, 1/0, 2/10, 3/011, 1/1, 1/splat32, 5/lower5, 5/registerZe, 1/0 }


\paragraph{Syntax} \texttt{cmovehq}\emph{lanecond} \emph{registerZe}? \emph{registerWe} = \emph{upper27\_\discretionary{}{}{}lower5}\emph{splat32}

\paragraph{Description} The \emph{registerWe}, \emph{registerZe} and the \emph{upper27\_\discretionary{}{}{}lower5} are interpreted as four half words packed into 64 bits. Each \emph{registerZe} half word is tested using the condition \emph{lanecond}. If true, the corresponding half word of the \emph{upper27\_\discretionary{}{}{}lower5} is stored into the corresponding half word of the \emph{registerWe}.


\paragraph{Modifier(s)}
\noindent\begin{itemize}
\item lanecond (cf. lanecond in section \ref{par:modifier-lanecond} on page \pageref{par:modifier-lanecond})
\item splat32 (cf. splat32 in section \ref{par:modifier-splat32} on page \pageref{par:modifier-splat32})
\end{itemize}

\paragraph{Execution} {Reservation table \textsf{ALU\_LITE.X} (Section~\ref{sec:reservation-ALU-LITE.X})}
~
{\begin{lstlisting}
stage ID:
  new splat32 = _ZX_1(splat32);
  new argument4 = _ZX_3(lanecond);
stage RR:
  new argument3 = splat32 ? (_WX_32(upper27_lower5) << 32) | _ZX_32(_WX_32(upper27_lower5)) : _SX_32(_WX_32(upper27_lower5));
  new argument2 = GPR[registerZe<<1];
stage E1:
  if (lanecond_16(argument4, argument2.16[0])) {
    GPR[registerWe<<1] = argument3;
  }
  if (lanecond_16(argument4, argument2.16[1])) {
    GPR[registerWe<<1] = argument3;
  }
  if (lanecond_16(argument4, argument2.16[2])) {
    GPR[registerWe<<1] = argument3;
  }
  if (lanecond_16(argument4, argument2.16[3])) {
    GPR[registerWe<<1] = argument3;
  }
\end{lstlisting}}
\newpage

\paragraph{Encoding} Format \textsf{ALU\_HQCMWRRO.M} (Section~\ref{sec:format-ALU-HQCMWRRO.M}), \verb|exubit24 = 0|\\

\bitfields{ 1/P, 2/00, 2/-- --, 27/upper27, 1/1, 2/11, 1/1, 3/lanecond, 1/0, 5/registerWo, 1/1, 2/10, 3/011, 1/1, 1/splat32, 5/lower5, 5/registerZo, 1/0 }


\paragraph{Syntax} \texttt{cmovehq}\emph{lanecond} \emph{registerZo}? \emph{registerWo} = \emph{upper27\_\discretionary{}{}{}lower5}\emph{splat32}

\paragraph{Description} The \emph{registerWo}, \emph{registerZo} and the \emph{upper27\_\discretionary{}{}{}lower5} are interpreted as four half words packed into 64 bits. Each \emph{registerZo} half word is tested using the condition \emph{lanecond}. If true, the corresponding half word of the \emph{upper27\_\discretionary{}{}{}lower5} is stored into the corresponding half word of the \emph{registerWo}.


\paragraph{Modifier(s)}
\noindent\begin{itemize}
\item lanecond (cf. lanecond in section \ref{par:modifier-lanecond} on page \pageref{par:modifier-lanecond})
\item splat32 (cf. splat32 in section \ref{par:modifier-splat32} on page \pageref{par:modifier-splat32})
\end{itemize}

\paragraph{Execution} {Reservation table \textsf{ALU\_LITE.X} (Section~\ref{sec:reservation-ALU-LITE.X})}
~
{\begin{lstlisting}
stage ID:
  new splat32 = _ZX_1(splat32);
  new argument4 = _ZX_3(lanecond);
stage RR:
  new argument3 = splat32 ? (_WX_32(upper27_lower5) << 32) | _ZX_32(_WX_32(upper27_lower5)) : _SX_32(_WX_32(upper27_lower5));
  new argument2 = GPR[(registerZo<<1)+1];
stage E1:
  if (lanecond_16(argument4, argument2.16[0])) {
    GPR[(registerWo<<1)+1] = argument3;
  }
  if (lanecond_16(argument4, argument2.16[1])) {
    GPR[(registerWo<<1)+1] = argument3;
  }
  if (lanecond_16(argument4, argument2.16[2])) {
    GPR[(registerWo<<1)+1] = argument3;
  }
  if (lanecond_16(argument4, argument2.16[3])) {
    GPR[(registerWo<<1)+1] = argument3;
  }
\end{lstlisting}}
\newpage
\subsubsection[{\small CMOVEQ: Conditional Move Quadruple Word}]{CMOVEQ\hfill Conditional Move Quadruple Word}
\label{instr:CMOVEQ}\index{cmoveq}
 
\paragraph{Encoding} Format \textsf{ALU\_QCMWRR} (Section~\ref{sec:format-ALU-QCMWRR}), \verb|steering = 11|\\

\bitfields{ 1/P, 2/11, 1/0, 4/cmovecond, 5/registerM, 1/--, 2/10, 3/011, 1/1, 5/registerO, 1/--, 6/registerZ }


\paragraph{Syntax} \texttt{cmoveq}\emph{cmovecond} \emph{registerZ}? \emph{registerM} = \emph{registerO}

\paragraph{Description} The \emph{registerZ} double word is tested using the condition \emph{cmovecond}. If true, the \emph{registerO} quadruple word is stored into the \emph{registerM} quadruple word. Else the \emph{registerM} quadruple word is unchanged.


\paragraph{Modifier(s)}
\noindent\begin{itemize}
\item cmovecond (cf. bcucond in section \ref{par:modifier-bcucond} on page \pageref{par:modifier-bcucond})
\end{itemize}

\paragraph{Execution} {Reservation table \textsf{ALU\_LITE} (Section~\ref{sec:reservation-ALU-LITE})}
~
{\begin{lstlisting}
stage ID:
  new argument4 = _ZX_4(cmovecond);
stage RR:
  new argument3 = PGR[registerO];
  new argument2 = GPR[registerZ];
stage E1:
  if (bcucond(argument4, argument2)) {
    PGR[registerM] = argument3;
  }
\end{lstlisting}}
\newpage
\subsubsection[{\small CMOVEWP: Conditional Move Word Pair}]{CMOVEWP\hfill Conditional Move Word Pair}
\label{instr:CMOVEWP}\index{cmovewp}
 
\paragraph{Encoding} Format \textsf{ALU\_WPCMWRRE} (Section~\ref{sec:format-ALU-WPCMWRRE}), \verb|exubit24 = 0|\\

\bitfields{ 1/P, 2/11, 1/1, 3/lanecond, 1/0, 5/registerWe, 1/0, 2/10, 3/011, 1/0, 5/registerYe, 1/--, 5/registerZe, 1/1 }


\paragraph{Syntax} \texttt{cmovewp}\emph{lanecond} \emph{registerZe}? \emph{registerWe} = \emph{registerYe}

\paragraph{Description} The \emph{registerWe}, \emph{registerZe} and the \emph{registerYe} are interpreted as two words packed into 64 bits. Each \emph{registerZe} word is tested using the condition \emph{lanecond}. If true, the corresponding word of the \emph{registerYe} is stored into the corresponding word of the \emph{registerWe}.


\paragraph{Modifier(s)}
\noindent\begin{itemize}
\item lanecond (cf. lanecond in section \ref{par:modifier-lanecond} on page \pageref{par:modifier-lanecond})
\end{itemize}

\paragraph{Execution} {Reservation table \textsf{ALU\_LITE} (Section~\ref{sec:reservation-ALU-LITE})}
~
{\begin{lstlisting}
stage ID:
  new argument4 = _ZX_3(lanecond);
stage RR:
  new argument3 = GPR[registerYe<<1];
  new argument2 = GPR[registerZe<<1];
stage E1:
  if (lanecond_32(argument4, argument2.32[0])) {
    GPR[registerWe<<1] = argument3;
  }
  if (lanecond_32(argument4, argument2.32[1])) {
    GPR[registerWe<<1] = argument3;
  }
\end{lstlisting}}
\newpage

\paragraph{Encoding} Format \textsf{ALU\_WPCMWRRO} (Section~\ref{sec:format-ALU-WPCMWRRO}), \verb|exubit24 = 0|\\

\bitfields{ 1/P, 2/11, 1/1, 3/lanecond, 1/0, 5/registerWo, 1/1, 2/10, 3/011, 1/0, 5/registerYo, 1/--, 5/registerZo, 1/1 }


\paragraph{Syntax} \texttt{cmovewp}\emph{lanecond} \emph{registerZo}? \emph{registerWo} = \emph{registerYo}

\paragraph{Description} The \emph{registerWo}, \emph{registerZo} and the \emph{registerYo} are interpreted as two words packed into 64 bits. Each \emph{registerZo} word is tested using the condition \emph{lanecond}. If true, the corresponding word of the \emph{registerYo} is stored into the corresponding word of the \emph{registerWo}.


\paragraph{Modifier(s)}
\noindent\begin{itemize}
\item lanecond (cf. lanecond in section \ref{par:modifier-lanecond} on page \pageref{par:modifier-lanecond})
\end{itemize}

\paragraph{Execution} {Reservation table \textsf{ALU\_LITE} (Section~\ref{sec:reservation-ALU-LITE})}
~
{\begin{lstlisting}
stage ID:
  new argument4 = _ZX_3(lanecond);
stage RR:
  new argument3 = GPR[(registerYo<<1)+1];
  new argument2 = GPR[(registerZo<<1)+1];
stage E1:
  if (lanecond_32(argument4, argument2.32[0])) {
    GPR[(registerWo<<1)+1] = argument3;
  }
  if (lanecond_32(argument4, argument2.32[1])) {
    GPR[(registerWo<<1)+1] = argument3;
  }
\end{lstlisting}}
\newpage

\paragraph{Encoding} Format \textsf{ALU\_WPCMWRRE.M} (Section~\ref{sec:format-ALU-WPCMWRRE.M}), \verb|exubit24 = 0|\\

\bitfields{ 1/P, 2/00, 2/-- --, 27/upper27, 1/1, 2/11, 1/1, 3/lanecond, 1/0, 5/registerWe, 1/0, 2/10, 3/011, 1/0, 1/splat32, 5/lower5, 5/registerZe, 1/1 }


\paragraph{Syntax} \texttt{cmovewp}\emph{lanecond} \emph{registerZe}? \emph{registerWe} = \emph{upper27\_\discretionary{}{}{}lower5}\emph{splat32}

\paragraph{Description} The \emph{registerWe}, \emph{registerZe} and the \emph{upper27\_\discretionary{}{}{}lower5} are interpreted as two words packed into 64 bits. Each \emph{registerZe} word is tested using the condition \emph{lanecond}. If true, the corresponding word of the \emph{upper27\_\discretionary{}{}{}lower5} is stored into the corresponding word of the \emph{registerWe}.


\paragraph{Modifier(s)}
\noindent\begin{itemize}
\item lanecond (cf. lanecond in section \ref{par:modifier-lanecond} on page \pageref{par:modifier-lanecond})
\item splat32 (cf. splat32 in section \ref{par:modifier-splat32} on page \pageref{par:modifier-splat32})
\end{itemize}

\paragraph{Execution} {Reservation table \textsf{ALU\_LITE.X} (Section~\ref{sec:reservation-ALU-LITE.X})}
~
{\begin{lstlisting}
stage ID:
  new splat32 = _ZX_1(splat32);
  new argument4 = _ZX_3(lanecond);
stage RR:
  new argument3 = splat32 ? (_WX_32(upper27_lower5) << 32) | _ZX_32(_WX_32(upper27_lower5)) : _SX_32(_WX_32(upper27_lower5));
  new argument2 = GPR[registerZe<<1];
stage E1:
  if (lanecond_32(argument4, argument2.32[0])) {
    GPR[registerWe<<1] = argument3;
  }
  if (lanecond_32(argument4, argument2.32[1])) {
    GPR[registerWe<<1] = argument3;
  }
\end{lstlisting}}
\newpage

\paragraph{Encoding} Format \textsf{ALU\_WPCMWRRO.M} (Section~\ref{sec:format-ALU-WPCMWRRO.M}), \verb|exubit24 = 0|\\

\bitfields{ 1/P, 2/00, 2/-- --, 27/upper27, 1/1, 2/11, 1/1, 3/lanecond, 1/0, 5/registerWo, 1/1, 2/10, 3/011, 1/0, 1/splat32, 5/lower5, 5/registerZo, 1/1 }


\paragraph{Syntax} \texttt{cmovewp}\emph{lanecond} \emph{registerZo}? \emph{registerWo} = \emph{upper27\_\discretionary{}{}{}lower5}\emph{splat32}

\paragraph{Description} The \emph{registerWo}, \emph{registerZo} and the \emph{upper27\_\discretionary{}{}{}lower5} are interpreted as two words packed into 64 bits. Each \emph{registerZo} word is tested using the condition \emph{lanecond}. If true, the corresponding word of the \emph{upper27\_\discretionary{}{}{}lower5} is stored into the corresponding word of the \emph{registerWo}.


\paragraph{Modifier(s)}
\noindent\begin{itemize}
\item lanecond (cf. lanecond in section \ref{par:modifier-lanecond} on page \pageref{par:modifier-lanecond})
\item splat32 (cf. splat32 in section \ref{par:modifier-splat32} on page \pageref{par:modifier-splat32})
\end{itemize}

\paragraph{Execution} {Reservation table \textsf{ALU\_LITE.X} (Section~\ref{sec:reservation-ALU-LITE.X})}
~
{\begin{lstlisting}
stage ID:
  new splat32 = _ZX_1(splat32);
  new argument4 = _ZX_3(lanecond);
stage RR:
  new argument3 = splat32 ? (_WX_32(upper27_lower5) << 32) | _ZX_32(_WX_32(upper27_lower5)) : _SX_32(_WX_32(upper27_lower5));
  new argument2 = GPR[(registerZo<<1)+1];
stage E1:
  if (lanecond_32(argument4, argument2.32[0])) {
    GPR[(registerWo<<1)+1] = argument3;
  }
  if (lanecond_32(argument4, argument2.32[1])) {
    GPR[(registerWo<<1)+1] = argument3;
  }
\end{lstlisting}}
\newpage
\subsubsection[{\small COMPBO: Compare Byte Octuple}]{COMPBO\hfill Compare Byte Octuple}
\label{instr:COMPBO}\index{compbo}
 
\paragraph{Encoding} Format \textsf{ALU\_BOCWRRE} (Section~\ref{sec:format-ALU-BOCWRRE}), \verb|alucode3 = 001|\\

\bitfields{ 1/P, 2/11, 1/1, 4/intcomp, 5/registerWe, 1/0, 2/10, 3/001, 1/1, 5/registerYe, 1/--, 5/registerZe, 1/1 }


\paragraph{Syntax} \texttt{compbo}\emph{intcomp} \emph{registerWe} = \emph{registerZe}, \emph{registerYe}

\paragraph{Description} The byte hexadecuple \emph{registerZe} is compared to the byte hexadecuple \emph{registerYe} using condition \emph{intcomp}. The boolean result bits are packed and stored into the \emph{registerWe}.


\paragraph{Modifier(s)}
\noindent\begin{itemize}
\item intcomp (cf. intcomp in section \ref{par:modifier-intcomp} on page \pageref{par:modifier-intcomp})
\end{itemize}

\paragraph{Execution} {Reservation table \textsf{ALU\_LITE} (Section~\ref{sec:reservation-ALU-LITE})}
~
{\begin{lstlisting}
stage ID:
  new argument4 = _ZX_4(intcomp);
stage RR:
  new argument3 = GPR[registerYe<<1];
  new argument2 = GPR[registerZe<<1];
  for (I = 0; I < 16; I += 1) {
    result1.1[I] = intcomp_8(argument4, argument2.8[I], argument3.8[I])
  };
stage E1:
  GPR[registerWe<<1] = result1;
\end{lstlisting}}
\newpage

\paragraph{Encoding} Format \textsf{ALU\_BOCWRRO} (Section~\ref{sec:format-ALU-BOCWRRO}), \verb|alucode3 = 001|\\

\bitfields{ 1/P, 2/11, 1/1, 4/intcomp, 5/registerWo, 1/1, 2/10, 3/001, 1/1, 5/registerYo, 1/--, 5/registerZo, 1/1 }


\paragraph{Syntax} \texttt{compbo}\emph{intcomp} \emph{registerWo} = \emph{registerZo}, \emph{registerYo}

\paragraph{Description} The byte hexadecuple \emph{registerZo} is compared to the byte hexadecuple \emph{registerYo} using condition \emph{intcomp}. The boolean result bits are packed and stored into the \emph{registerWo}.


\paragraph{Modifier(s)}
\noindent\begin{itemize}
\item intcomp (cf. intcomp in section \ref{par:modifier-intcomp} on page \pageref{par:modifier-intcomp})
\end{itemize}

\paragraph{Execution} {Reservation table \textsf{ALU\_LITE} (Section~\ref{sec:reservation-ALU-LITE})}
~
{\begin{lstlisting}
stage ID:
  new argument4 = _ZX_4(intcomp);
stage RR:
  new argument3 = GPR[(registerYo<<1)+1];
  new argument2 = GPR[(registerZo<<1)+1];
  for (I = 0; I < 16; I += 1) {
    result1.1[I] = intcomp_8(argument4, argument2.8[I], argument3.8[I])
  };
stage E1:
  GPR[(registerWo<<1)+1] = result1;
\end{lstlisting}}
\newpage

\paragraph{Encoding} Format \textsf{ALU\_BOCWRRE.M} (Section~\ref{sec:format-ALU-BOCWRRE.M}), \verb|alucode3 = 001|\\

\bitfields{ 1/P, 2/00, 2/-- --, 27/upper27, 1/1, 2/11, 1/1, 4/intcomp, 5/registerWe, 1/0, 2/10, 3/001, 1/1, 1/splat32, 5/lower5, 5/registerZe, 1/1 }


\paragraph{Syntax} \texttt{compbo}\emph{intcomp} \emph{registerWe} = \emph{registerZe}, \emph{upper27\_\discretionary{}{}{}lower5}\emph{splat32}

\paragraph{Description} The byte hexadecuple \emph{registerZe} is compared to the byte hexadecuple \emph{upper27\_\discretionary{}{}{}lower5} using condition \emph{intcomp}. The boolean result bits are packed and stored into the \emph{registerWe}.


\paragraph{Modifier(s)}
\noindent\begin{itemize}
\item intcomp (cf. intcomp in section \ref{par:modifier-intcomp} on page \pageref{par:modifier-intcomp})
\item splat32 (cf. splat32 in section \ref{par:modifier-splat32} on page \pageref{par:modifier-splat32})
\end{itemize}

\paragraph{Execution} {Reservation table \textsf{ALU\_LITE.X} (Section~\ref{sec:reservation-ALU-LITE.X})}
~
{\begin{lstlisting}
stage ID:
  new splat32 = _ZX_1(splat32);
  new argument4 = _ZX_4(intcomp);
stage RR:
  new argument3 = splat32 ? (_WX_32(upper27_lower5) << 32) | _ZX_32(_WX_32(upper27_lower5)) : _SX_32(_WX_32(upper27_lower5));
  new argument2 = GPR[registerZe<<1];
  for (I = 0; I < 16; I += 1) {
    result1.1[I] = intcomp_8(argument4, argument2.8[I], argument3.8[I])
  };
stage E1:
  GPR[registerWe<<1] = result1;
\end{lstlisting}}
\newpage

\paragraph{Encoding} Format \textsf{ALU\_BOCWRRO.M} (Section~\ref{sec:format-ALU-BOCWRRO.M}), \verb|alucode3 = 001|\\

\bitfields{ 1/P, 2/00, 2/-- --, 27/upper27, 1/1, 2/11, 1/1, 4/intcomp, 5/registerWo, 1/1, 2/10, 3/001, 1/1, 1/splat32, 5/lower5, 5/registerZo, 1/1 }


\paragraph{Syntax} \texttt{compbo}\emph{intcomp} \emph{registerWo} = \emph{registerZo}, \emph{upper27\_\discretionary{}{}{}lower5}\emph{splat32}

\paragraph{Description} The byte hexadecuple \emph{registerZo} is compared to the byte hexadecuple \emph{upper27\_\discretionary{}{}{}lower5} using condition \emph{intcomp}. The boolean result bits are packed and stored into the \emph{registerWo}.


\paragraph{Modifier(s)}
\noindent\begin{itemize}
\item intcomp (cf. intcomp in section \ref{par:modifier-intcomp} on page \pageref{par:modifier-intcomp})
\item splat32 (cf. splat32 in section \ref{par:modifier-splat32} on page \pageref{par:modifier-splat32})
\end{itemize}

\paragraph{Execution} {Reservation table \textsf{ALU\_LITE.X} (Section~\ref{sec:reservation-ALU-LITE.X})}
~
{\begin{lstlisting}
stage ID:
  new splat32 = _ZX_1(splat32);
  new argument4 = _ZX_4(intcomp);
stage RR:
  new argument3 = splat32 ? (_WX_32(upper27_lower5) << 32) | _ZX_32(_WX_32(upper27_lower5)) : _SX_32(_WX_32(upper27_lower5));
  new argument2 = GPR[(registerZo<<1)+1];
  for (I = 0; I < 16; I += 1) {
    result1.1[I] = intcomp_8(argument4, argument2.8[I], argument3.8[I])
  };
stage E1:
  GPR[(registerWo<<1)+1] = result1;
\end{lstlisting}}
\newpage
\subsubsection[{\small COMPD: Compare Double Word}]{COMPD\hfill Compare Double Word}
\label{instr:COMPD}\index{compd}
 
\paragraph{Encoding} Format \textsf{ALU\_DCWRR} (Section~\ref{sec:format-ALU-DCWRR}), \verb|steering = 11|\\

\bitfields{ 1/P, 2/11, 1/0, 4/intcomp, 6/registerW, 2/10, 3/001, 1/0, 6/registerY, 6/registerZ }


\paragraph{Syntax} \texttt{compd}\emph{intcomp} \emph{registerW} = \emph{registerZ}, \emph{registerY}

\paragraph{Description} The double word \emph{registerZ} is compared to the double word \emph{registerY} using condition \emph{intcomp}. The boolean result extended to double word is stored into the \emph{registerW}.


\paragraph{Modifier(s)}
\noindent\begin{itemize}
\item intcomp (cf. intcomp in section \ref{par:modifier-intcomp} on page \pageref{par:modifier-intcomp})
\end{itemize}

\paragraph{Execution} {Reservation table \textsf{ALU\_TINY} (Section~\ref{sec:reservation-ALU-TINY})}
~
{\begin{lstlisting}
stage ID:
  new argument4 = _ZX_4(intcomp);
stage RR:
  new argument3 = GPR[registerY];
  new argument2 = GPR[registerZ];
  new result1 = intcomp_64(argument4, argument2, argument3);
stage E1:
  GPR[registerW] = result1;
\end{lstlisting}}
\newpage

\paragraph{Encoding} Format \textsf{ALU\_DCWRR.W} (Section~\ref{sec:format-ALU-DCWRR.W}), \verb|steering = 11|\\

\bitfields{ 1/P, 2/00, 2/-- --, 27/upper27, 1/1, 2/11, 1/0, 4/intcomp, 6/registerW, 2/10, 3/001, 1/0, 1/0, 5/lower5, 6/registerZ }


\paragraph{Syntax} \texttt{compd}\emph{intcomp} \emph{registerW} = \emph{registerZ}, \emph{upper27\_\discretionary{}{}{}lower5}

\paragraph{Description} The double word \emph{registerZ} is compared to the double word \emph{upper27\_\discretionary{}{}{}lower5} using condition \emph{intcomp}. The boolean result extended to double word is stored into the \emph{registerW}.


\paragraph{Modifier(s)}
\noindent\begin{itemize}
\item intcomp (cf. intcomp in section \ref{par:modifier-intcomp} on page \pageref{par:modifier-intcomp})
\end{itemize}

\paragraph{Execution} {Reservation table \textsf{ALU\_TINY.X} (Section~\ref{sec:reservation-ALU-TINY.X})}
~
{\begin{lstlisting}
stage ID:
  new argument4 = _ZX_4(intcomp);
stage RR:
  new argument3 = _WX_32(upper27_lower5);
  new argument2 = GPR[registerZ];
  new result1 = intcomp_64(argument4, argument2, argument3);
stage E1:
  GPR[registerW] = result1;
\end{lstlisting}}
\newpage
\subsubsection[{\small COMPHQ: Compare Half Word Quadruple}]{COMPHQ\hfill Compare Half Word Quadruple}
\label{instr:COMPHQ}\index{comphq}
 
\paragraph{Encoding} Format \textsf{ALU\_HQCWRRE} (Section~\ref{sec:format-ALU-HQCWRRE}), \verb|alucode3 = 001|\\

\bitfields{ 1/P, 2/11, 1/1, 4/intcomp, 5/registerWe, 1/0, 2/10, 3/001, 1/1, 5/registerYe, 1/--, 5/registerZe, 1/0 }


\paragraph{Syntax} \texttt{comphq}\emph{intcomp} \emph{registerWe} = \emph{registerZe}, \emph{registerYe}

\paragraph{Description} The half word quadruple \emph{registerZe} is compared to the half word quadruple \emph{registerYe} using condition \emph{intcomp}. The four boolean result bits are packed and stored into the \emph{registerWe}.


\paragraph{Modifier(s)}
\noindent\begin{itemize}
\item intcomp (cf. intcomp in section \ref{par:modifier-intcomp} on page \pageref{par:modifier-intcomp})
\end{itemize}

\paragraph{Execution} {Reservation table \textsf{ALU\_LITE} (Section~\ref{sec:reservation-ALU-LITE})}
~
{\begin{lstlisting}
stage ID:
  new argument4 = _ZX_4(intcomp);
stage RR:
  new argument3 = GPR[registerYe<<1];
  new argument2 = GPR[registerZe<<1];
  for (I = 0; I < 4; I += 1) {
    result1.1[I] = intcomp_16(argument4, argument2.16[I], argument3.16[I])
  };
stage E1:
  GPR[registerWe<<1] = result1;
\end{lstlisting}}
\newpage

\paragraph{Encoding} Format \textsf{ALU\_HQCWRRO} (Section~\ref{sec:format-ALU-HQCWRRO}), \verb|alucode3 = 001|\\

\bitfields{ 1/P, 2/11, 1/1, 4/intcomp, 5/registerWo, 1/1, 2/10, 3/001, 1/1, 5/registerYo, 1/--, 5/registerZo, 1/0 }


\paragraph{Syntax} \texttt{comphq}\emph{intcomp} \emph{registerWo} = \emph{registerZo}, \emph{registerYo}

\paragraph{Description} The half word quadruple \emph{registerZo} is compared to the half word quadruple \emph{registerYo} using condition \emph{intcomp}. The four boolean result bits are packed and stored into the \emph{registerWo}.


\paragraph{Modifier(s)}
\noindent\begin{itemize}
\item intcomp (cf. intcomp in section \ref{par:modifier-intcomp} on page \pageref{par:modifier-intcomp})
\end{itemize}

\paragraph{Execution} {Reservation table \textsf{ALU\_LITE} (Section~\ref{sec:reservation-ALU-LITE})}
~
{\begin{lstlisting}
stage ID:
  new argument4 = _ZX_4(intcomp);
stage RR:
  new argument3 = GPR[(registerYo<<1)+1];
  new argument2 = GPR[(registerZo<<1)+1];
  for (I = 0; I < 4; I += 1) {
    result1.1[I] = intcomp_16(argument4, argument2.16[I], argument3.16[I])
  };
stage E1:
  GPR[(registerWo<<1)+1] = result1;
\end{lstlisting}}
\newpage

\paragraph{Encoding} Format \textsf{ALU\_HQCWRRE.M} (Section~\ref{sec:format-ALU-HQCWRRE.M}), \verb|alucode3 = 001|\\

\bitfields{ 1/P, 2/00, 2/-- --, 27/upper27, 1/1, 2/11, 1/1, 4/intcomp, 5/registerWe, 1/0, 2/10, 3/001, 1/1, 1/splat32, 5/lower5, 5/registerZe, 1/0 }


\paragraph{Syntax} \texttt{comphq}\emph{intcomp} \emph{registerWe} = \emph{registerZe}, \emph{upper27\_\discretionary{}{}{}lower5}\emph{splat32}

\paragraph{Description} The half word quadruple \emph{registerZe} is compared to the half word quadruple \emph{upper27\_\discretionary{}{}{}lower5} using condition \emph{intcomp}. The four boolean result bits are packed and stored into the \emph{registerWe}.


\paragraph{Modifier(s)}
\noindent\begin{itemize}
\item intcomp (cf. intcomp in section \ref{par:modifier-intcomp} on page \pageref{par:modifier-intcomp})
\item splat32 (cf. splat32 in section \ref{par:modifier-splat32} on page \pageref{par:modifier-splat32})
\end{itemize}

\paragraph{Execution} {Reservation table \textsf{ALU\_LITE.X} (Section~\ref{sec:reservation-ALU-LITE.X})}
~
{\begin{lstlisting}
stage ID:
  new splat32 = _ZX_1(splat32);
  new argument4 = _ZX_4(intcomp);
stage RR:
  new argument3 = splat32 ? (_WX_32(upper27_lower5) << 32) | _ZX_32(_WX_32(upper27_lower5)) : _SX_32(_WX_32(upper27_lower5));
  new argument2 = GPR[registerZe<<1];
  for (I = 0; I < 4; I += 1) {
    result1.1[I] = intcomp_16(argument4, argument2.16[I], argument3.16[I])
  };
stage E1:
  GPR[registerWe<<1] = result1;
\end{lstlisting}}
\newpage

\paragraph{Encoding} Format \textsf{ALU\_HQCWRRO.M} (Section~\ref{sec:format-ALU-HQCWRRO.M}), \verb|alucode3 = 001|\\

\bitfields{ 1/P, 2/00, 2/-- --, 27/upper27, 1/1, 2/11, 1/1, 4/intcomp, 5/registerWo, 1/1, 2/10, 3/001, 1/1, 1/splat32, 5/lower5, 5/registerZo, 1/0 }


\paragraph{Syntax} \texttt{comphq}\emph{intcomp} \emph{registerWo} = \emph{registerZo}, \emph{upper27\_\discretionary{}{}{}lower5}\emph{splat32}

\paragraph{Description} The half word quadruple \emph{registerZo} is compared to the half word quadruple \emph{upper27\_\discretionary{}{}{}lower5} using condition \emph{intcomp}. The four boolean result bits are packed and stored into the \emph{registerWo}.


\paragraph{Modifier(s)}
\noindent\begin{itemize}
\item intcomp (cf. intcomp in section \ref{par:modifier-intcomp} on page \pageref{par:modifier-intcomp})
\item splat32 (cf. splat32 in section \ref{par:modifier-splat32} on page \pageref{par:modifier-splat32})
\end{itemize}

\paragraph{Execution} {Reservation table \textsf{ALU\_LITE.X} (Section~\ref{sec:reservation-ALU-LITE.X})}
~
{\begin{lstlisting}
stage ID:
  new splat32 = _ZX_1(splat32);
  new argument4 = _ZX_4(intcomp);
stage RR:
  new argument3 = splat32 ? (_WX_32(upper27_lower5) << 32) | _ZX_32(_WX_32(upper27_lower5)) : _SX_32(_WX_32(upper27_lower5));
  new argument2 = GPR[(registerZo<<1)+1];
  for (I = 0; I < 4; I += 1) {
    result1.1[I] = intcomp_16(argument4, argument2.16[I], argument3.16[I])
  };
stage E1:
  GPR[(registerWo<<1)+1] = result1;
\end{lstlisting}}
\newpage
\subsubsection[{\small COMPNBO: Compare Negate Byte Octuple}]{COMPNBO\hfill Compare Negate Byte Octuple}
\label{instr:COMPNBO}\index{compnbo}
 
\paragraph{Encoding} Format \textsf{ALU\_BOCNWRRE} (Section~\ref{sec:format-ALU-BOCNWRRE}), \verb|alucode3 = 010|\\

\bitfields{ 1/P, 2/11, 1/1, 4/intcomp, 5/registerWe, 1/0, 2/10, 3/010, 1/1, 5/registerYe, 1/--, 5/registerZe, 1/1 }


\paragraph{Syntax} \texttt{compnbo}\emph{intcomp} \emph{registerWe} = \emph{registerZe}, \emph{registerYe}

\paragraph{Description} The \emph{registerZe} and the \emph{registerYe} are interpreted as eight bytes packed into 64 bits. The \emph{registerZe} bytes are compared to the \emph{registerYe} bytes using condition \emph{intcomp}. The eight resulting boolean bytes are negated, packed and stored into the \emph{registerWe}.


\paragraph{Modifier(s)}
\noindent\begin{itemize}
\item intcomp (cf. intcomp in section \ref{par:modifier-intcomp} on page \pageref{par:modifier-intcomp})
\end{itemize}

\paragraph{Execution} {Reservation table \textsf{ALU\_LITE} (Section~\ref{sec:reservation-ALU-LITE})}
~
{\begin{lstlisting}
stage ID:
  new argument4 = _ZX_4(intcomp);
stage RR:
  new argument3 = GPR[registerYe<<1];
  new argument2 = GPR[registerZe<<1];
  for (I = 0; I < 8; I += 1) {
    result1.8[I] = -intcomp_8(argument4, argument2.8[I], argument3.8[I])
  };
stage E1:
  GPR[registerWe<<1] = result1;
\end{lstlisting}}
\newpage

\paragraph{Encoding} Format \textsf{ALU\_BOCNWRRO} (Section~\ref{sec:format-ALU-BOCNWRRO}), \verb|alucode3 = 010|\\

\bitfields{ 1/P, 2/11, 1/1, 4/intcomp, 5/registerWo, 1/1, 2/10, 3/010, 1/1, 5/registerYo, 1/--, 5/registerZo, 1/1 }


\paragraph{Syntax} \texttt{compnbo}\emph{intcomp} \emph{registerWo} = \emph{registerZo}, \emph{registerYo}

\paragraph{Description} The \emph{registerZo} and the \emph{registerYo} are interpreted as eight bytes packed into 64 bits. The \emph{registerZo} bytes are compared to the \emph{registerYo} bytes using condition \emph{intcomp}. The eight resulting boolean bytes are negated, packed and stored into the \emph{registerWo}.


\paragraph{Modifier(s)}
\noindent\begin{itemize}
\item intcomp (cf. intcomp in section \ref{par:modifier-intcomp} on page \pageref{par:modifier-intcomp})
\end{itemize}

\paragraph{Execution} {Reservation table \textsf{ALU\_LITE} (Section~\ref{sec:reservation-ALU-LITE})}
~
{\begin{lstlisting}
stage ID:
  new argument4 = _ZX_4(intcomp);
stage RR:
  new argument3 = GPR[(registerYo<<1)+1];
  new argument2 = GPR[(registerZo<<1)+1];
  for (I = 0; I < 8; I += 1) {
    result1.8[I] = -intcomp_8(argument4, argument2.8[I], argument3.8[I])
  };
stage E1:
  GPR[(registerWo<<1)+1] = result1;
\end{lstlisting}}
\newpage

\paragraph{Encoding} Format \textsf{ALU\_BOCNWRRE.M} (Section~\ref{sec:format-ALU-BOCNWRRE.M}), \verb|alucode3 = 010|\\

\bitfields{ 1/P, 2/00, 2/-- --, 27/upper27, 1/1, 2/11, 1/1, 4/intcomp, 5/registerWe, 1/0, 2/10, 3/010, 1/1, 1/splat32, 5/lower5, 5/registerZe, 1/1 }


\paragraph{Syntax} \texttt{compnbo}\emph{intcomp} \emph{registerWe} = \emph{registerZe}, \emph{upper27\_\discretionary{}{}{}lower5}\emph{splat32}

\paragraph{Description} The \emph{registerZe} and the \emph{upper27\_\discretionary{}{}{}lower5} are interpreted as eight bytes packed into 64 bits. The \emph{registerZe} bytes are compared to the \emph{upper27\_\discretionary{}{}{}lower5} bytes using condition \emph{intcomp}. The eight resulting boolean bytes are negated, packed and stored into the \emph{registerWe}.


\paragraph{Modifier(s)}
\noindent\begin{itemize}
\item intcomp (cf. intcomp in section \ref{par:modifier-intcomp} on page \pageref{par:modifier-intcomp})
\item splat32 (cf. splat32 in section \ref{par:modifier-splat32} on page \pageref{par:modifier-splat32})
\end{itemize}

\paragraph{Execution} {Reservation table \textsf{ALU\_LITE.X} (Section~\ref{sec:reservation-ALU-LITE.X})}
~
{\begin{lstlisting}
stage ID:
  new splat32 = _ZX_1(splat32);
  new argument4 = _ZX_4(intcomp);
stage RR:
  new argument3 = splat32 ? (_WX_32(upper27_lower5) << 32) | _ZX_32(_WX_32(upper27_lower5)) : _SX_32(_WX_32(upper27_lower5));
  new argument2 = GPR[registerZe<<1];
  for (I = 0; I < 8; I += 1) {
    result1.8[I] = -intcomp_8(argument4, argument2.8[I], argument3.8[I])
  };
stage E1:
  GPR[registerWe<<1] = result1;
\end{lstlisting}}
\newpage

\paragraph{Encoding} Format \textsf{ALU\_BOCNWRRO.M} (Section~\ref{sec:format-ALU-BOCNWRRO.M}), \verb|alucode3 = 010|\\

\bitfields{ 1/P, 2/00, 2/-- --, 27/upper27, 1/1, 2/11, 1/1, 4/intcomp, 5/registerWo, 1/1, 2/10, 3/010, 1/1, 1/splat32, 5/lower5, 5/registerZo, 1/1 }


\paragraph{Syntax} \texttt{compnbo}\emph{intcomp} \emph{registerWo} = \emph{registerZo}, \emph{upper27\_\discretionary{}{}{}lower5}\emph{splat32}

\paragraph{Description} The \emph{registerZo} and the \emph{upper27\_\discretionary{}{}{}lower5} are interpreted as eight bytes packed into 64 bits. The \emph{registerZo} bytes are compared to the \emph{upper27\_\discretionary{}{}{}lower5} bytes using condition \emph{intcomp}. The eight resulting boolean bytes are negated, packed and stored into the \emph{registerWo}.


\paragraph{Modifier(s)}
\noindent\begin{itemize}
\item intcomp (cf. intcomp in section \ref{par:modifier-intcomp} on page \pageref{par:modifier-intcomp})
\item splat32 (cf. splat32 in section \ref{par:modifier-splat32} on page \pageref{par:modifier-splat32})
\end{itemize}

\paragraph{Execution} {Reservation table \textsf{ALU\_LITE.X} (Section~\ref{sec:reservation-ALU-LITE.X})}
~
{\begin{lstlisting}
stage ID:
  new splat32 = _ZX_1(splat32);
  new argument4 = _ZX_4(intcomp);
stage RR:
  new argument3 = splat32 ? (_WX_32(upper27_lower5) << 32) | _ZX_32(_WX_32(upper27_lower5)) : _SX_32(_WX_32(upper27_lower5));
  new argument2 = GPR[(registerZo<<1)+1];
  for (I = 0; I < 8; I += 1) {
    result1.8[I] = -intcomp_8(argument4, argument2.8[I], argument3.8[I])
  };
stage E1:
  GPR[(registerWo<<1)+1] = result1;
\end{lstlisting}}
\newpage
\subsubsection[{\small COMPND: Compare Negate Double Word}]{COMPND\hfill Compare Negate Double Word}
\label{instr:COMPND}\index{compnd}
 
\paragraph{Encoding} Format \textsf{ALU\_DCNWRRE} (Section~\ref{sec:format-ALU-DCNWRRE}), \verb|alucode3 = 010|\\

\bitfields{ 1/P, 2/11, 1/1, 4/intcomp, 5/registerWe, 1/0, 2/10, 3/010, 1/0, 5/registerYe, 1/--, 5/registerZe, 1/0 }


\paragraph{Syntax} \texttt{compnd}\emph{intcomp} \emph{registerWe} = \emph{registerZe}, \emph{registerYe}

\paragraph{Description} The double word \emph{registerZe} is compared to the double word \emph{registerYe} using condition \emph{intcomp}. The negated boolean result extended to double word is stored into the \emph{registerWe}.


\paragraph{Modifier(s)}
\noindent\begin{itemize}
\item intcomp (cf. intcomp in section \ref{par:modifier-intcomp} on page \pageref{par:modifier-intcomp})
\end{itemize}

\paragraph{Execution} {Reservation table \textsf{ALU\_LITE} (Section~\ref{sec:reservation-ALU-LITE})}
~
{\begin{lstlisting}
stage ID:
  new argument4 = _ZX_4(intcomp);
stage RR:
  new argument3 = GPR[registerYe<<1];
  new argument2 = GPR[registerZe<<1];
  new result1 = -intcomp_64(argument4, argument2, argument3);
stage E1:
  GPR[registerWe<<1] = result1;
\end{lstlisting}}
\newpage

\paragraph{Encoding} Format \textsf{ALU\_DCNWRRO} (Section~\ref{sec:format-ALU-DCNWRRO}), \verb|alucode3 = 010|\\

\bitfields{ 1/P, 2/11, 1/1, 4/intcomp, 5/registerWo, 1/1, 2/10, 3/010, 1/0, 5/registerYo, 1/--, 5/registerZo, 1/0 }


\paragraph{Syntax} \texttt{compnd}\emph{intcomp} \emph{registerWo} = \emph{registerZo}, \emph{registerYo}

\paragraph{Description} The double word \emph{registerZo} is compared to the double word \emph{registerYo} using condition \emph{intcomp}. The negated boolean result extended to double word is stored into the \emph{registerWo}.


\paragraph{Modifier(s)}
\noindent\begin{itemize}
\item intcomp (cf. intcomp in section \ref{par:modifier-intcomp} on page \pageref{par:modifier-intcomp})
\end{itemize}

\paragraph{Execution} {Reservation table \textsf{ALU\_LITE} (Section~\ref{sec:reservation-ALU-LITE})}
~
{\begin{lstlisting}
stage ID:
  new argument4 = _ZX_4(intcomp);
stage RR:
  new argument3 = GPR[(registerYo<<1)+1];
  new argument2 = GPR[(registerZo<<1)+1];
  new result1 = -intcomp_64(argument4, argument2, argument3);
stage E1:
  GPR[(registerWo<<1)+1] = result1;
\end{lstlisting}}
\newpage

\paragraph{Encoding} Format \textsf{ALU\_DCNWRRE.M} (Section~\ref{sec:format-ALU-DCNWRRE.M}), \verb|alucode3 = 010|\\

\bitfields{ 1/P, 2/00, 2/-- --, 27/upper27, 1/1, 2/11, 1/1, 4/intcomp, 5/registerWe, 1/0, 2/10, 3/010, 1/0, 1/splat32, 5/lower5, 5/registerZe, 1/0 }


\paragraph{Syntax} \texttt{compnd}\emph{intcomp} \emph{registerWe} = \emph{registerZe}, \emph{upper27\_\discretionary{}{}{}lower5}\emph{splat32}

\paragraph{Description} The double word \emph{registerZe} is compared to the double word \emph{upper27\_\discretionary{}{}{}lower5} using condition \emph{intcomp}. The negated boolean result extended to double word is stored into the \emph{registerWe}.


\paragraph{Modifier(s)}
\noindent\begin{itemize}
\item intcomp (cf. intcomp in section \ref{par:modifier-intcomp} on page \pageref{par:modifier-intcomp})
\item splat32 (cf. splat32 in section \ref{par:modifier-splat32} on page \pageref{par:modifier-splat32})
\end{itemize}

\paragraph{Execution} {Reservation table \textsf{ALU\_LITE.X} (Section~\ref{sec:reservation-ALU-LITE.X})}
~
{\begin{lstlisting}
stage ID:
  new splat32 = _ZX_1(splat32);
  new argument4 = _ZX_4(intcomp);
stage RR:
  new argument3 = splat32 ? (_WX_32(upper27_lower5) << 32) | _ZX_32(_WX_32(upper27_lower5)) : _SX_32(_WX_32(upper27_lower5));
  new argument2 = GPR[registerZe<<1];
  new result1 = -intcomp_64(argument4, argument2, argument3);
stage E1:
  GPR[registerWe<<1] = result1;
\end{lstlisting}}
\newpage

\paragraph{Encoding} Format \textsf{ALU\_DCNWRRO.M} (Section~\ref{sec:format-ALU-DCNWRRO.M}), \verb|alucode3 = 010|\\

\bitfields{ 1/P, 2/00, 2/-- --, 27/upper27, 1/1, 2/11, 1/1, 4/intcomp, 5/registerWo, 1/1, 2/10, 3/010, 1/0, 1/splat32, 5/lower5, 5/registerZo, 1/0 }


\paragraph{Syntax} \texttt{compnd}\emph{intcomp} \emph{registerWo} = \emph{registerZo}, \emph{upper27\_\discretionary{}{}{}lower5}\emph{splat32}

\paragraph{Description} The double word \emph{registerZo} is compared to the double word \emph{upper27\_\discretionary{}{}{}lower5} using condition \emph{intcomp}. The negated boolean result extended to double word is stored into the \emph{registerWo}.


\paragraph{Modifier(s)}
\noindent\begin{itemize}
\item intcomp (cf. intcomp in section \ref{par:modifier-intcomp} on page \pageref{par:modifier-intcomp})
\item splat32 (cf. splat32 in section \ref{par:modifier-splat32} on page \pageref{par:modifier-splat32})
\end{itemize}

\paragraph{Execution} {Reservation table \textsf{ALU\_LITE.X} (Section~\ref{sec:reservation-ALU-LITE.X})}
~
{\begin{lstlisting}
stage ID:
  new splat32 = _ZX_1(splat32);
  new argument4 = _ZX_4(intcomp);
stage RR:
  new argument3 = splat32 ? (_WX_32(upper27_lower5) << 32) | _ZX_32(_WX_32(upper27_lower5)) : _SX_32(_WX_32(upper27_lower5));
  new argument2 = GPR[(registerZo<<1)+1];
  new result1 = -intcomp_64(argument4, argument2, argument3);
stage E1:
  GPR[(registerWo<<1)+1] = result1;
\end{lstlisting}}
\newpage
\subsubsection[{\small COMPNHQ: Compare Negate Half Word Quadruple}]{COMPNHQ\hfill Compare Negate Half Word Quadruple}
\label{instr:COMPNHQ}\index{compnhq}
 
\paragraph{Encoding} Format \textsf{ALU\_HQCNWRRE} (Section~\ref{sec:format-ALU-HQCNWRRE}), \verb|alucode3 = 010|\\

\bitfields{ 1/P, 2/11, 1/1, 4/intcomp, 5/registerWe, 1/0, 2/10, 3/010, 1/1, 5/registerYe, 1/--, 5/registerZe, 1/0 }


\paragraph{Syntax} \texttt{compnhq}\emph{intcomp} \emph{registerWe} = \emph{registerZe}, \emph{registerYe}

\paragraph{Description} The \emph{registerZe} and the \emph{registerYe} are interpreted as four half words packed into 64 bits. The \emph{registerZe} half words are compared to the \emph{registerYe} half words using condition \emph{intcomp}. The four resulting boolean half word are negated, packed and stored into the \emph{registerWe}.


\paragraph{Modifier(s)}
\noindent\begin{itemize}
\item intcomp (cf. intcomp in section \ref{par:modifier-intcomp} on page \pageref{par:modifier-intcomp})
\end{itemize}

\paragraph{Execution} {Reservation table \textsf{ALU\_LITE} (Section~\ref{sec:reservation-ALU-LITE})}
~
{\begin{lstlisting}
stage ID:
  new argument4 = _ZX_4(intcomp);
stage RR:
  new argument3 = GPR[registerYe<<1];
  new argument2 = GPR[registerZe<<1];
  for (I = 0; I < 4; I += 1) {
    result1.16[I] = -intcomp_16(argument4, argument2.16[I], argument3.16[I])
  };
stage E1:
  GPR[registerWe<<1] = result1;
\end{lstlisting}}
\newpage

\paragraph{Encoding} Format \textsf{ALU\_HQCNWRRO} (Section~\ref{sec:format-ALU-HQCNWRRO}), \verb|alucode3 = 010|\\

\bitfields{ 1/P, 2/11, 1/1, 4/intcomp, 5/registerWo, 1/1, 2/10, 3/010, 1/1, 5/registerYo, 1/--, 5/registerZo, 1/0 }


\paragraph{Syntax} \texttt{compnhq}\emph{intcomp} \emph{registerWo} = \emph{registerZo}, \emph{registerYo}

\paragraph{Description} The \emph{registerZo} and the \emph{registerYo} are interpreted as four half words packed into 64 bits. The \emph{registerZo} half words are compared to the \emph{registerYo} half words using condition \emph{intcomp}. The four resulting boolean half word are negated, packed and stored into the \emph{registerWo}.


\paragraph{Modifier(s)}
\noindent\begin{itemize}
\item intcomp (cf. intcomp in section \ref{par:modifier-intcomp} on page \pageref{par:modifier-intcomp})
\end{itemize}

\paragraph{Execution} {Reservation table \textsf{ALU\_LITE} (Section~\ref{sec:reservation-ALU-LITE})}
~
{\begin{lstlisting}
stage ID:
  new argument4 = _ZX_4(intcomp);
stage RR:
  new argument3 = GPR[(registerYo<<1)+1];
  new argument2 = GPR[(registerZo<<1)+1];
  for (I = 0; I < 4; I += 1) {
    result1.16[I] = -intcomp_16(argument4, argument2.16[I], argument3.16[I])
  };
stage E1:
  GPR[(registerWo<<1)+1] = result1;
\end{lstlisting}}
\newpage

\paragraph{Encoding} Format \textsf{ALU\_HQCNWRRE.M} (Section~\ref{sec:format-ALU-HQCNWRRE.M}), \verb|alucode3 = 010|\\

\bitfields{ 1/P, 2/00, 2/-- --, 27/upper27, 1/1, 2/11, 1/1, 4/intcomp, 5/registerWe, 1/0, 2/10, 3/010, 1/1, 1/splat32, 5/lower5, 5/registerZe, 1/0 }


\paragraph{Syntax} \texttt{compnhq}\emph{intcomp} \emph{registerWe} = \emph{registerZe}, \emph{upper27\_\discretionary{}{}{}lower5}\emph{splat32}

\paragraph{Description} The \emph{registerZe} and the \emph{upper27\_\discretionary{}{}{}lower5} are interpreted as four half words packed into 64 bits. The \emph{registerZe} half words are compared to the \emph{upper27\_\discretionary{}{}{}lower5} half words using condition \emph{intcomp}. The four resulting boolean half word are negated, packed and stored into the \emph{registerWe}.


\paragraph{Modifier(s)}
\noindent\begin{itemize}
\item intcomp (cf. intcomp in section \ref{par:modifier-intcomp} on page \pageref{par:modifier-intcomp})
\item splat32 (cf. splat32 in section \ref{par:modifier-splat32} on page \pageref{par:modifier-splat32})
\end{itemize}

\paragraph{Execution} {Reservation table \textsf{ALU\_LITE.X} (Section~\ref{sec:reservation-ALU-LITE.X})}
~
{\begin{lstlisting}
stage ID:
  new splat32 = _ZX_1(splat32);
  new argument4 = _ZX_4(intcomp);
stage RR:
  new argument3 = splat32 ? (_WX_32(upper27_lower5) << 32) | _ZX_32(_WX_32(upper27_lower5)) : _SX_32(_WX_32(upper27_lower5));
  new argument2 = GPR[registerZe<<1];
  for (I = 0; I < 4; I += 1) {
    result1.16[I] = -intcomp_16(argument4, argument2.16[I], argument3.16[I])
  };
stage E1:
  GPR[registerWe<<1] = result1;
\end{lstlisting}}
\newpage

\paragraph{Encoding} Format \textsf{ALU\_HQCNWRRO.M} (Section~\ref{sec:format-ALU-HQCNWRRO.M}), \verb|alucode3 = 010|\\

\bitfields{ 1/P, 2/00, 2/-- --, 27/upper27, 1/1, 2/11, 1/1, 4/intcomp, 5/registerWo, 1/1, 2/10, 3/010, 1/1, 1/splat32, 5/lower5, 5/registerZo, 1/0 }


\paragraph{Syntax} \texttt{compnhq}\emph{intcomp} \emph{registerWo} = \emph{registerZo}, \emph{upper27\_\discretionary{}{}{}lower5}\emph{splat32}

\paragraph{Description} The \emph{registerZo} and the \emph{upper27\_\discretionary{}{}{}lower5} are interpreted as four half words packed into 64 bits. The \emph{registerZo} half words are compared to the \emph{upper27\_\discretionary{}{}{}lower5} half words using condition \emph{intcomp}. The four resulting boolean half word are negated, packed and stored into the \emph{registerWo}.


\paragraph{Modifier(s)}
\noindent\begin{itemize}
\item intcomp (cf. intcomp in section \ref{par:modifier-intcomp} on page \pageref{par:modifier-intcomp})
\item splat32 (cf. splat32 in section \ref{par:modifier-splat32} on page \pageref{par:modifier-splat32})
\end{itemize}

\paragraph{Execution} {Reservation table \textsf{ALU\_LITE.X} (Section~\ref{sec:reservation-ALU-LITE.X})}
~
{\begin{lstlisting}
stage ID:
  new splat32 = _ZX_1(splat32);
  new argument4 = _ZX_4(intcomp);
stage RR:
  new argument3 = splat32 ? (_WX_32(upper27_lower5) << 32) | _ZX_32(_WX_32(upper27_lower5)) : _SX_32(_WX_32(upper27_lower5));
  new argument2 = GPR[(registerZo<<1)+1];
  for (I = 0; I < 4; I += 1) {
    result1.16[I] = -intcomp_16(argument4, argument2.16[I], argument3.16[I])
  };
stage E1:
  GPR[(registerWo<<1)+1] = result1;
\end{lstlisting}}
\newpage
\subsubsection[{\small COMPNWP: Compare Negate Word Pair}]{COMPNWP\hfill Compare Negate Word Pair}
\label{instr:COMPNWP}\index{compnwp}
 
\paragraph{Encoding} Format \textsf{ALU\_WPCNWRRE} (Section~\ref{sec:format-ALU-WPCNWRRE}), \verb|alucode3 = 010|\\

\bitfields{ 1/P, 2/11, 1/1, 4/intcomp, 5/registerWe, 1/0, 2/10, 3/010, 1/0, 5/registerYe, 1/--, 5/registerZe, 1/1 }


\paragraph{Syntax} \texttt{compnwp}\emph{intcomp} \emph{registerWe} = \emph{registerZe}, \emph{registerYe}

\paragraph{Description} The \emph{registerZe} and the \emph{registerYe} are interpreted as two words packed into 64 bits. The \emph{registerZe} words are compared to the \emph{registerYe} words using condition \emph{intcomp}. The two resulting boolean words are negated, packed and stored into the \emph{registerWe}.


\paragraph{Modifier(s)}
\noindent\begin{itemize}
\item intcomp (cf. intcomp in section \ref{par:modifier-intcomp} on page \pageref{par:modifier-intcomp})
\end{itemize}

\paragraph{Execution} {Reservation table \textsf{ALU\_LITE} (Section~\ref{sec:reservation-ALU-LITE})}
~
{\begin{lstlisting}
stage ID:
  new argument4 = _ZX_4(intcomp);
stage RR:
  new argument3 = GPR[registerYe<<1];
  new argument2 = GPR[registerZe<<1];
  for (I = 0; I < 2; I += 1) {
    result1.32[I] = -intcomp_32(argument4, argument2.32[I], argument3.32[I])
  };
stage E1:
  GPR[registerWe<<1] = result1;
\end{lstlisting}}
\newpage

\paragraph{Encoding} Format \textsf{ALU\_WPCNWRRO} (Section~\ref{sec:format-ALU-WPCNWRRO}), \verb|alucode3 = 010|\\

\bitfields{ 1/P, 2/11, 1/1, 4/intcomp, 5/registerWo, 1/1, 2/10, 3/010, 1/0, 5/registerYo, 1/--, 5/registerZo, 1/1 }


\paragraph{Syntax} \texttt{compnwp}\emph{intcomp} \emph{registerWo} = \emph{registerZo}, \emph{registerYo}

\paragraph{Description} The \emph{registerZo} and the \emph{registerYo} are interpreted as two words packed into 64 bits. The \emph{registerZo} words are compared to the \emph{registerYo} words using condition \emph{intcomp}. The two resulting boolean words are negated, packed and stored into the \emph{registerWo}.


\paragraph{Modifier(s)}
\noindent\begin{itemize}
\item intcomp (cf. intcomp in section \ref{par:modifier-intcomp} on page \pageref{par:modifier-intcomp})
\end{itemize}

\paragraph{Execution} {Reservation table \textsf{ALU\_LITE} (Section~\ref{sec:reservation-ALU-LITE})}
~
{\begin{lstlisting}
stage ID:
  new argument4 = _ZX_4(intcomp);
stage RR:
  new argument3 = GPR[(registerYo<<1)+1];
  new argument2 = GPR[(registerZo<<1)+1];
  for (I = 0; I < 2; I += 1) {
    result1.32[I] = -intcomp_32(argument4, argument2.32[I], argument3.32[I])
  };
stage E1:
  GPR[(registerWo<<1)+1] = result1;
\end{lstlisting}}
\newpage

\paragraph{Encoding} Format \textsf{ALU\_WPCNWRRE.M} (Section~\ref{sec:format-ALU-WPCNWRRE.M}), \verb|alucode3 = 010|\\

\bitfields{ 1/P, 2/00, 2/-- --, 27/upper27, 1/1, 2/11, 1/1, 4/intcomp, 5/registerWe, 1/0, 2/10, 3/010, 1/0, 1/splat32, 5/lower5, 5/registerZe, 1/1 }


\paragraph{Syntax} \texttt{compnwp}\emph{intcomp} \emph{registerWe} = \emph{registerZe}, \emph{upper27\_\discretionary{}{}{}lower5}\emph{splat32}

\paragraph{Description} The \emph{registerZe} and the \emph{upper27\_\discretionary{}{}{}lower5} are interpreted as two words packed into 64 bits. The \emph{registerZe} words are compared to the \emph{upper27\_\discretionary{}{}{}lower5} words using condition \emph{intcomp}. The two resulting boolean words are negated, packed and stored into the \emph{registerWe}.


\paragraph{Modifier(s)}
\noindent\begin{itemize}
\item intcomp (cf. intcomp in section \ref{par:modifier-intcomp} on page \pageref{par:modifier-intcomp})
\item splat32 (cf. splat32 in section \ref{par:modifier-splat32} on page \pageref{par:modifier-splat32})
\end{itemize}

\paragraph{Execution} {Reservation table \textsf{ALU\_LITE.X} (Section~\ref{sec:reservation-ALU-LITE.X})}
~
{\begin{lstlisting}
stage ID:
  new splat32 = _ZX_1(splat32);
  new argument4 = _ZX_4(intcomp);
stage RR:
  new argument3 = splat32 ? (_WX_32(upper27_lower5) << 32) | _ZX_32(_WX_32(upper27_lower5)) : _SX_32(_WX_32(upper27_lower5));
  new argument2 = GPR[registerZe<<1];
  for (I = 0; I < 2; I += 1) {
    result1.32[I] = -intcomp_32(argument4, argument2.32[I], argument3.32[I])
  };
stage E1:
  GPR[registerWe<<1] = result1;
\end{lstlisting}}
\newpage

\paragraph{Encoding} Format \textsf{ALU\_WPCNWRRO.M} (Section~\ref{sec:format-ALU-WPCNWRRO.M}), \verb|alucode3 = 010|\\

\bitfields{ 1/P, 2/00, 2/-- --, 27/upper27, 1/1, 2/11, 1/1, 4/intcomp, 5/registerWo, 1/1, 2/10, 3/010, 1/0, 1/splat32, 5/lower5, 5/registerZo, 1/1 }


\paragraph{Syntax} \texttt{compnwp}\emph{intcomp} \emph{registerWo} = \emph{registerZo}, \emph{upper27\_\discretionary{}{}{}lower5}\emph{splat32}

\paragraph{Description} The \emph{registerZo} and the \emph{upper27\_\discretionary{}{}{}lower5} are interpreted as two words packed into 64 bits. The \emph{registerZo} words are compared to the \emph{upper27\_\discretionary{}{}{}lower5} words using condition \emph{intcomp}. The two resulting boolean words are negated, packed and stored into the \emph{registerWo}.


\paragraph{Modifier(s)}
\noindent\begin{itemize}
\item intcomp (cf. intcomp in section \ref{par:modifier-intcomp} on page \pageref{par:modifier-intcomp})
\item splat32 (cf. splat32 in section \ref{par:modifier-splat32} on page \pageref{par:modifier-splat32})
\end{itemize}

\paragraph{Execution} {Reservation table \textsf{ALU\_LITE.X} (Section~\ref{sec:reservation-ALU-LITE.X})}
~
{\begin{lstlisting}
stage ID:
  new splat32 = _ZX_1(splat32);
  new argument4 = _ZX_4(intcomp);
stage RR:
  new argument3 = splat32 ? (_WX_32(upper27_lower5) << 32) | _ZX_32(_WX_32(upper27_lower5)) : _SX_32(_WX_32(upper27_lower5));
  new argument2 = GPR[(registerZo<<1)+1];
  for (I = 0; I < 2; I += 1) {
    result1.32[I] = -intcomp_32(argument4, argument2.32[I], argument3.32[I])
  };
stage E1:
  GPR[(registerWo<<1)+1] = result1;
\end{lstlisting}}
\newpage
\subsubsection[{\small COMPQ: Compare Quadruple Word}]{COMPQ\hfill Compare Quadruple Word}
\label{instr:COMPQ}\index{compq}
 
\paragraph{Encoding} Format \textsf{ALU\_QCWRR} (Section~\ref{sec:format-ALU-QCWRR}), \verb|steering = 11|\\

\bitfields{ 1/P, 2/11, 1/0, 4/intcomp, 5/registerM, 1/--, 2/10, 3/001, 1/1, 5/registerO, 1/--, 5/registerP, 1/-- }


\paragraph{Syntax} \texttt{compq}\emph{intcomp} \emph{registerM} = \emph{registerP}, \emph{registerO}

\paragraph{Description} The quadruple word \emph{registerP} is compared to the quadruple word \emph{registerO} using condition \emph{intcomp}. The boolean result extended to double word is stored into the \emph{registerM}.


\paragraph{Modifier(s)}
\noindent\begin{itemize}
\item intcomp (cf. intcomp in section \ref{par:modifier-intcomp} on page \pageref{par:modifier-intcomp})
\end{itemize}

\paragraph{Execution} {Reservation table \textsf{ALU\_LITE} (Section~\ref{sec:reservation-ALU-LITE})}
~
{\begin{lstlisting}
stage ID:
  new argument4 = _ZX_4(intcomp);
stage RR:
  new argument3 = PGR[registerO];
  new argument2 = PGR[registerP];
  new result1 = intcomp_128(argument4, argument2, argument3);
stage E1:
  PGR[registerM] = result1;
\end{lstlisting}}
\newpage
\subsubsection[{\small COMPW: Compare Word}]{COMPW\hfill Compare Word}
\label{instr:COMPW}\index{compw}
 
\paragraph{Encoding} Format \textsf{ALU\_WCWRR} (Section~\ref{sec:format-ALU-WCWRR}), \verb|steering = 11|\\

\bitfields{ 1/P, 2/11, 1/0, 4/intcomp, 6/registerW, 2/11, 3/001, 1/signextw, 6/registerY, 6/registerZ }


\paragraph{Syntax} \texttt{compw}\emph{intcomp}\emph{signextw} \emph{registerW} = \emph{registerZ}, \emph{registerY}

\paragraph{Description} The word \emph{registerZ} is compared to the word \emph{registerY} using condition \emph{intcomp}. The boolean result extended to word is stored into the \emph{registerW}.


\paragraph{Modifier(s)}
\noindent\begin{itemize}
\item intcomp (cf. intcomp in section \ref{par:modifier-intcomp} on page \pageref{par:modifier-intcomp})
\item signextw (cf. signextw in section \ref{par:modifier-signextw} on page \pageref{par:modifier-signextw})
\end{itemize}

\paragraph{Execution} {Reservation table \textsf{ALU\_TINY} (Section~\ref{sec:reservation-ALU-TINY})}
~
{\begin{lstlisting}
stage ID:
  new signextw = _ZX_1(signextw);
  new argument4 = _ZX_4(intcomp);
stage RR:
  new argument3 = GPR[registerY];
  new argument2 = GPR[registerZ];
  new result1 = intcomp_32(argument4, argument2, argument3);
stage E1:
  GPR[registerW] = signextw ? _SX_32(result1) : _ZX_32(result1);
\end{lstlisting}}
\newpage

\paragraph{Encoding} Format \textsf{ALU\_WCWRR.W} (Section~\ref{sec:format-ALU-WCWRR.W}), \verb|steering = 11|\\

\bitfields{ 1/P, 2/00, 2/-- --, 27/upper27, 1/1, 2/11, 1/0, 4/intcomp, 6/registerW, 2/11, 3/001, 1/signextw, 1/0, 5/lower5, 6/registerZ }


\paragraph{Syntax} \texttt{compw}\emph{intcomp}\emph{signextw} \emph{registerW} = \emph{registerZ}, \emph{upper27\_\discretionary{}{}{}lower5}

\paragraph{Description} The word \emph{registerZ} is compared to the word \emph{upper27\_\discretionary{}{}{}lower5} using condition \emph{intcomp}. The boolean result extended to word is stored into the \emph{registerW}.


\paragraph{Modifier(s)}
\noindent\begin{itemize}
\item intcomp (cf. intcomp in section \ref{par:modifier-intcomp} on page \pageref{par:modifier-intcomp})
\item signextw (cf. signextw in section \ref{par:modifier-signextw} on page \pageref{par:modifier-signextw})
\end{itemize}

\paragraph{Execution} {Reservation table \textsf{ALU\_TINY.X} (Section~\ref{sec:reservation-ALU-TINY.X})}
~
{\begin{lstlisting}
stage ID:
  new signextw = _ZX_1(signextw);
  new argument4 = _ZX_4(intcomp);
stage RR:
  new argument3 = _WX_32(upper27_lower5);
  new argument2 = GPR[registerZ];
  new result1 = intcomp_32(argument4, argument2, argument3);
stage E1:
  GPR[registerW] = signextw ? _SX_32(result1) : _ZX_32(result1);
\end{lstlisting}}
\newpage
\subsubsection[{\small COMPWP: Compare Word Pair}]{COMPWP\hfill Compare Word Pair}
\label{instr:COMPWP}\index{compwp}
 
\paragraph{Encoding} Format \textsf{ALU\_WPCWRRE} (Section~\ref{sec:format-ALU-WPCWRRE}), \verb|alucode3 = 001|\\

\bitfields{ 1/P, 2/11, 1/1, 4/intcomp, 5/registerWe, 1/0, 2/10, 3/001, 1/0, 5/registerYe, 1/--, 5/registerZe, 1/1 }


\paragraph{Syntax} \texttt{compwp}\emph{intcomp} \emph{registerWe} = \emph{registerZe}, \emph{registerYe}

\paragraph{Description} The word pair \emph{registerZe} is compared to the word pair \emph{registerYe} using condition \emph{intcomp}. The two boolean result bits are packed and stored into the \emph{registerWe}.


\paragraph{Modifier(s)}
\noindent\begin{itemize}
\item intcomp (cf. intcomp in section \ref{par:modifier-intcomp} on page \pageref{par:modifier-intcomp})
\end{itemize}

\paragraph{Execution} {Reservation table \textsf{ALU\_TINY} (Section~\ref{sec:reservation-ALU-TINY})}
~
{\begin{lstlisting}
stage ID:
  new argument4 = _ZX_4(intcomp);
stage RR:
  new argument3 = GPR[registerYe<<1];
  new argument2 = GPR[registerZe<<1];
  for (I = 0; I < 2; I += 1) {
    result1.1[I] = intcomp_32(argument4, argument2.32[I], argument3.32[I])
  };
stage E1:
  GPR[registerWe<<1] = result1;
\end{lstlisting}}
\newpage

\paragraph{Encoding} Format \textsf{ALU\_WPCWRRO} (Section~\ref{sec:format-ALU-WPCWRRO}), \verb|alucode3 = 001|\\

\bitfields{ 1/P, 2/11, 1/1, 4/intcomp, 5/registerWo, 1/1, 2/10, 3/001, 1/0, 5/registerYo, 1/--, 5/registerZo, 1/1 }


\paragraph{Syntax} \texttt{compwp}\emph{intcomp} \emph{registerWo} = \emph{registerZo}, \emph{registerYo}

\paragraph{Description} The word pair \emph{registerZo} is compared to the word pair \emph{registerYo} using condition \emph{intcomp}. The two boolean result bits are packed and stored into the \emph{registerWo}.


\paragraph{Modifier(s)}
\noindent\begin{itemize}
\item intcomp (cf. intcomp in section \ref{par:modifier-intcomp} on page \pageref{par:modifier-intcomp})
\end{itemize}

\paragraph{Execution} {Reservation table \textsf{ALU\_TINY} (Section~\ref{sec:reservation-ALU-TINY})}
~
{\begin{lstlisting}
stage ID:
  new argument4 = _ZX_4(intcomp);
stage RR:
  new argument3 = GPR[(registerYo<<1)+1];
  new argument2 = GPR[(registerZo<<1)+1];
  for (I = 0; I < 2; I += 1) {
    result1.1[I] = intcomp_32(argument4, argument2.32[I], argument3.32[I])
  };
stage E1:
  GPR[(registerWo<<1)+1] = result1;
\end{lstlisting}}
\newpage

\paragraph{Encoding} Format \textsf{ALU\_WPCWRRE.M} (Section~\ref{sec:format-ALU-WPCWRRE.M}), \verb|alucode3 = 001|\\

\bitfields{ 1/P, 2/00, 2/-- --, 27/upper27, 1/1, 2/11, 1/1, 4/intcomp, 5/registerWe, 1/0, 2/10, 3/001, 1/0, 1/splat32, 5/lower5, 5/registerZe, 1/1 }


\paragraph{Syntax} \texttt{compwp}\emph{intcomp} \emph{registerWe} = \emph{registerZe}, \emph{upper27\_\discretionary{}{}{}lower5}\emph{splat32}

\paragraph{Description} The word pair \emph{registerZe} is compared to the word pair \emph{upper27\_\discretionary{}{}{}lower5} using condition \emph{intcomp}. The two boolean result bits are packed and stored into the \emph{registerWe}.


\paragraph{Modifier(s)}
\noindent\begin{itemize}
\item intcomp (cf. intcomp in section \ref{par:modifier-intcomp} on page \pageref{par:modifier-intcomp})
\item splat32 (cf. splat32 in section \ref{par:modifier-splat32} on page \pageref{par:modifier-splat32})
\end{itemize}

\paragraph{Execution} {Reservation table \textsf{ALU\_TINY.X} (Section~\ref{sec:reservation-ALU-TINY.X})}
~
{\begin{lstlisting}
stage ID:
  new splat32 = _ZX_1(splat32);
  new argument4 = _ZX_4(intcomp);
stage RR:
  new argument3 = splat32 ? (_WX_32(upper27_lower5) << 32) | _ZX_32(_WX_32(upper27_lower5)) : _SX_32(_WX_32(upper27_lower5));
  new argument2 = GPR[registerZe<<1];
  for (I = 0; I < 2; I += 1) {
    result1.1[I] = intcomp_32(argument4, argument2.32[I], argument3.32[I])
  };
stage E1:
  GPR[registerWe<<1] = result1;
\end{lstlisting}}
\newpage

\paragraph{Encoding} Format \textsf{ALU\_WPCWRRO.M} (Section~\ref{sec:format-ALU-WPCWRRO.M}), \verb|alucode3 = 001|\\

\bitfields{ 1/P, 2/00, 2/-- --, 27/upper27, 1/1, 2/11, 1/1, 4/intcomp, 5/registerWo, 1/1, 2/10, 3/001, 1/0, 1/splat32, 5/lower5, 5/registerZo, 1/1 }


\paragraph{Syntax} \texttt{compwp}\emph{intcomp} \emph{registerWo} = \emph{registerZo}, \emph{upper27\_\discretionary{}{}{}lower5}\emph{splat32}

\paragraph{Description} The word pair \emph{registerZo} is compared to the word pair \emph{upper27\_\discretionary{}{}{}lower5} using condition \emph{intcomp}. The two boolean result bits are packed and stored into the \emph{registerWo}.


\paragraph{Modifier(s)}
\noindent\begin{itemize}
\item intcomp (cf. intcomp in section \ref{par:modifier-intcomp} on page \pageref{par:modifier-intcomp})
\item splat32 (cf. splat32 in section \ref{par:modifier-splat32} on page \pageref{par:modifier-splat32})
\end{itemize}

\paragraph{Execution} {Reservation table \textsf{ALU\_TINY.X} (Section~\ref{sec:reservation-ALU-TINY.X})}
~
{\begin{lstlisting}
stage ID:
  new splat32 = _ZX_1(splat32);
  new argument4 = _ZX_4(intcomp);
stage RR:
  new argument3 = splat32 ? (_WX_32(upper27_lower5) << 32) | _ZX_32(_WX_32(upper27_lower5)) : _SX_32(_WX_32(upper27_lower5));
  new argument2 = GPR[(registerZo<<1)+1];
  for (I = 0; I < 2; I += 1) {
    result1.1[I] = intcomp_32(argument4, argument2.32[I], argument3.32[I])
  };
stage E1:
  GPR[(registerWo<<1)+1] = result1;
\end{lstlisting}}
\newpage
\subsubsection[{\small CRCBELLW: Cyclic Redundancy Check Big-Endian Least Significant with Least Significant Word}]{CRCBELLW\hfill Cyclic Redundancy Check Big-Endian Least Significant with Least Significant Word}
\label{instr:CRCBELLW}\index{crcbellw}
 
\paragraph{Encoding} Format \textsf{ALU\_WCWRRR} (Section~\ref{sec:format-ALU-WCWRRR}), \verb|exucode2 = 0101|\\

\bitfields{ 1/P, 2/11, 1/0, 4/0101, 6/registerW, 2/11, 3/010, 1/1, 6/registerY, 6/registerZ }


\paragraph{Syntax} \texttt{crcbellw} \emph{registerW} = \emph{registerZ}, \emph{registerY}

\paragraph{Description} The 4 least significant bytes of the \emph{registerW} are XOR-ed with the 4 least significant bytes of the \emph{registerZ} considered in reverse order. The result is reduced modulo the 4-byte polynomial stored in the 4 least significant bytes of the \emph{registerY}. The result is stored back into the \emph{registerW}.


\paragraph{Execution} {Reservation table \textsf{ALU\_LITE} (Section~\ref{sec:reservation-ALU-LITE})}
~
{\begin{lstlisting}
stage RR:
  new argument3 = GPR[registerY];
  new argument2 = GPR[registerZ];
stage E1:
  new argument1 = GPR[registerW];
  new argument1_data = argument1.32[0] ^ _SWAP_16(_SWAP_8(argument2));
  new result1 = crc32_be_u32(argument1_data.32[0], argument3.32[0]);
stage E2:
  GPR[registerW] = result1;
\end{lstlisting}}
\newpage

\paragraph{Encoding} Format \textsf{ALU\_WCWRRR.W} (Section~\ref{sec:format-ALU-WCWRRR.W}), \verb|exucode2 = 0101|\\

\bitfields{ 1/P, 2/00, 2/-- --, 27/upper27, 1/1, 2/11, 1/0, 4/0101, 6/registerW, 2/11, 3/010, 1/1, 1/0, 5/lower5, 6/registerZ }


\paragraph{Syntax} \texttt{crcbellw} \emph{registerW} = \emph{registerZ}, \emph{upper27\_\discretionary{}{}{}lower5}

\paragraph{Description} The 4 least significant bytes of the \emph{registerW} are XOR-ed with the 4 least significant bytes of the \emph{registerZ} considered in reverse order. The result is reduced modulo the 4-byte polynomial stored in the 4 least significant bytes of the \emph{upper27\_\discretionary{}{}{}lower5}. The result is stored back into the \emph{registerW}.


\paragraph{Execution} {Reservation table \textsf{ALU\_LITE.X} (Section~\ref{sec:reservation-ALU-LITE.X})}
~
{\begin{lstlisting}
stage RR:
  new argument3 = _WX_32(upper27_lower5);
  new argument2 = GPR[registerZ];
stage E1:
  new argument1 = GPR[registerW];
  new argument1_data = argument1.32[0] ^ _SWAP_16(_SWAP_8(argument2));
  new result1 = crc32_be_u32(argument1_data.32[0], argument3.32[0]);
stage E2:
  GPR[registerW] = result1;
\end{lstlisting}}
\newpage
\subsubsection[{\small CRCBELMW: Cyclic Redundancy Check Big-Endian Least Significant with Most Significant Word}]{CRCBELMW\hfill Cyclic Redundancy Check Big-Endian Least Significant with Most Significant Word}
\label{instr:CRCBELMW}\index{crcbelmw}
 
\paragraph{Encoding} Format \textsf{ALU\_WCWRRR} (Section~\ref{sec:format-ALU-WCWRRR}), \verb|exucode2 = 0100|\\

\bitfields{ 1/P, 2/11, 1/0, 4/0100, 6/registerW, 2/11, 3/010, 1/1, 6/registerY, 6/registerZ }


\paragraph{Syntax} \texttt{crcbelmw} \emph{registerW} = \emph{registerZ}, \emph{registerY}

\paragraph{Description} The 4 least significant bytes of the \emph{registerW} are XOR-ed with the 4 most significant bytes of the \emph{registerZ} considered in reverse order. The result is reduced modulo the 4-byte polynomial stored in the 4 least significant bytes of the \emph{registerY}. The result is stored back into the \emph{registerW}.


\paragraph{Execution} {Reservation table \textsf{ALU\_LITE} (Section~\ref{sec:reservation-ALU-LITE})}
~
{\begin{lstlisting}
stage RR:
  new argument3 = GPR[registerY];
  new argument2 = GPR[registerZ];
stage E1:
  new argument1 = GPR[registerW];
  new argument1_data = argument1.32[0] ^ _SWAP_16(_SWAP_8(argument2 >> 32));
  new result1 = crc32_be_u32(argument1_data.32[0], argument3.32[0]);
stage E2:
  GPR[registerW] = result1;
\end{lstlisting}}
\newpage

\paragraph{Encoding} Format \textsf{ALU\_WCWRRR.W} (Section~\ref{sec:format-ALU-WCWRRR.W}), \verb|exucode2 = 0100|\\

\bitfields{ 1/P, 2/00, 2/-- --, 27/upper27, 1/1, 2/11, 1/0, 4/0100, 6/registerW, 2/11, 3/010, 1/1, 1/0, 5/lower5, 6/registerZ }


\paragraph{Syntax} \texttt{crcbelmw} \emph{registerW} = \emph{registerZ}, \emph{upper27\_\discretionary{}{}{}lower5}

\paragraph{Description} The 4 least significant bytes of the \emph{registerW} are XOR-ed with the 4 most significant bytes of the \emph{registerZ} considered in reverse order. The result is reduced modulo the 4-byte polynomial stored in the 4 least significant bytes of the \emph{upper27\_\discretionary{}{}{}lower5}. The result is stored back into the \emph{registerW}.


\paragraph{Execution} {Reservation table \textsf{ALU\_LITE.X} (Section~\ref{sec:reservation-ALU-LITE.X})}
~
{\begin{lstlisting}
stage RR:
  new argument3 = _WX_32(upper27_lower5);
  new argument2 = GPR[registerZ];
stage E1:
  new argument1 = GPR[registerW];
  new argument1_data = argument1.32[0] ^ _SWAP_16(_SWAP_8(argument2 >> 32));
  new result1 = crc32_be_u32(argument1_data.32[0], argument3.32[0]);
stage E2:
  GPR[registerW] = result1;
\end{lstlisting}}
\newpage
\subsubsection[{\small CRCLELLW: Cyclic Redundancy Check Little-Endian Least Significant with Least Significant Word}]{CRCLELLW\hfill Cyclic Redundancy Check Little-Endian Least Significant with Least Significant Word}
\label{instr:CRCLELLW}\index{crclellw}
 
\paragraph{Encoding} Format \textsf{ALU\_WCWRRR} (Section~\ref{sec:format-ALU-WCWRRR}), \verb|exucode2 = 0111|\\

\bitfields{ 1/P, 2/11, 1/0, 4/0111, 6/registerW, 2/11, 3/010, 1/1, 6/registerY, 6/registerZ }


\paragraph{Syntax} \texttt{crclellw} \emph{registerW} = \emph{registerZ}, \emph{registerY}

\paragraph{Description} The 4 least significant bytes of the \emph{registerW} are XOR-ed with the 4 least significant bytes of the \emph{registerZ}. The result is bit reversed and then reduced modulo the 4-byte polynomial stored in the 4 least significant bytes of the \emph{registerY} bit reversed beforehand. The result is bit reversed and stored back into the \emph{registerW}.


\paragraph{Execution} {Reservation table \textsf{ALU\_LITE} (Section~\ref{sec:reservation-ALU-LITE})}
~
{\begin{lstlisting}
stage RR:
  new argument3 = GPR[registerY];
  new argument2 = GPR[registerZ];
stage E1:
  new argument1 = GPR[registerW];
  new argument1_data = argument1.32[0] ^ argument2.32[0];
  new result1 = reflect_32(crc32_be_u32(reflect_32(argument1_data.32[0]), reflect_32(argument3.32[0])));
stage E2:
  GPR[registerW] = result1;
\end{lstlisting}}
\newpage

\paragraph{Encoding} Format \textsf{ALU\_WCWRRR.W} (Section~\ref{sec:format-ALU-WCWRRR.W}), \verb|exucode2 = 0111|\\

\bitfields{ 1/P, 2/00, 2/-- --, 27/upper27, 1/1, 2/11, 1/0, 4/0111, 6/registerW, 2/11, 3/010, 1/1, 1/0, 5/lower5, 6/registerZ }


\paragraph{Syntax} \texttt{crclellw} \emph{registerW} = \emph{registerZ}, \emph{upper27\_\discretionary{}{}{}lower5}

\paragraph{Description} The 4 least significant bytes of the \emph{registerW} are XOR-ed with the 4 least significant bytes of the \emph{registerZ}. The result is bit reversed and then reduced modulo the 4-byte polynomial stored in the 4 least significant bytes of the \emph{upper27\_\discretionary{}{}{}lower5} bit reversed beforehand. The result is bit reversed and stored back into the \emph{registerW}.


\paragraph{Execution} {Reservation table \textsf{ALU\_LITE.X} (Section~\ref{sec:reservation-ALU-LITE.X})}
~
{\begin{lstlisting}
stage RR:
  new argument3 = _WX_32(upper27_lower5);
  new argument2 = GPR[registerZ];
stage E1:
  new argument1 = GPR[registerW];
  new argument1_data = argument1.32[0] ^ argument2.32[0];
  new result1 = reflect_32(crc32_be_u32(reflect_32(argument1_data.32[0]), reflect_32(argument3.32[0])));
stage E2:
  GPR[registerW] = result1;
\end{lstlisting}}
\newpage
\subsubsection[{\small CRCLELMW: Cyclic Redundancy Check Little-Endian Least Significant with Most Significant Word}]{CRCLELMW\hfill Cyclic Redundancy Check Little-Endian Least Significant with Most Significant Word}
\label{instr:CRCLELMW}\index{crclelmw}
 
\paragraph{Encoding} Format \textsf{ALU\_WCWRRR} (Section~\ref{sec:format-ALU-WCWRRR}), \verb|exucode2 = 0110|\\

\bitfields{ 1/P, 2/11, 1/0, 4/0110, 6/registerW, 2/11, 3/010, 1/1, 6/registerY, 6/registerZ }


\paragraph{Syntax} \texttt{crclelmw} \emph{registerW} = \emph{registerZ}, \emph{registerY}

\paragraph{Description} The 4 least significant bytes of the \emph{registerW} are XOR-ed with the 4 most significant bytes of the \emph{registerZ}. The result is bit reversed and then reduced modulo the 4-byte polynomial stored in the 4 least significant bytes of the \emph{registerY} bit reversed beforehand. The result is bit reversed and stored back into the \emph{registerW}.


\paragraph{Execution} {Reservation table \textsf{ALU\_LITE} (Section~\ref{sec:reservation-ALU-LITE})}
~
{\begin{lstlisting}
stage RR:
  new argument3 = GPR[registerY];
  new argument2 = GPR[registerZ];
stage E1:
  new argument1 = GPR[registerW];
  new argument1_data = argument1.32[0] ^ argument2.32[1];
  new result1 = reflect_32(crc32_be_u32(reflect_32(argument1_data.32[0]), reflect_32(argument3.32[0])));
stage E2:
  GPR[registerW] = result1;
\end{lstlisting}}
\newpage

\paragraph{Encoding} Format \textsf{ALU\_WCWRRR.W} (Section~\ref{sec:format-ALU-WCWRRR.W}), \verb|exucode2 = 0110|\\

\bitfields{ 1/P, 2/00, 2/-- --, 27/upper27, 1/1, 2/11, 1/0, 4/0110, 6/registerW, 2/11, 3/010, 1/1, 1/0, 5/lower5, 6/registerZ }


\paragraph{Syntax} \texttt{crclelmw} \emph{registerW} = \emph{registerZ}, \emph{upper27\_\discretionary{}{}{}lower5}

\paragraph{Description} The 4 least significant bytes of the \emph{registerW} are XOR-ed with the 4 most significant bytes of the \emph{registerZ}. The result is bit reversed and then reduced modulo the 4-byte polynomial stored in the 4 least significant bytes of the \emph{upper27\_\discretionary{}{}{}lower5} bit reversed beforehand. The result is bit reversed and stored back into the \emph{registerW}.


\paragraph{Execution} {Reservation table \textsf{ALU\_LITE.X} (Section~\ref{sec:reservation-ALU-LITE.X})}
~
{\begin{lstlisting}
stage RR:
  new argument3 = _WX_32(upper27_lower5);
  new argument2 = GPR[registerZ];
stage E1:
  new argument1 = GPR[registerW];
  new argument1_data = argument1.32[0] ^ argument2.32[1];
  new result1 = reflect_32(crc32_be_u32(reflect_32(argument1_data.32[0]), reflect_32(argument3.32[0])));
stage E2:
  GPR[registerW] = result1;
\end{lstlisting}}
\newpage
\subsubsection[{\small CTZD: Count Trailing Zero Double Word}]{CTZD\hfill Count Trailing Zero Double Word}
\label{instr:CTZD}\index{ctzd}
 
\paragraph{Encoding} Format \textsf{ALU\_DBWR} (Section~\ref{sec:format-ALU-DBWR}), \verb|alucode8 = 0111|\\

\bitfields{ 1/P, 2/11, 1/0, 4/1111, 6/registerW, 2/10, 3/101, 1/0, 4/0111, 1/0, 1/0, 6/registerZ }


\paragraph{Syntax} \texttt{ctzd} \emph{registerW} = \emph{registerZ}

\paragraph{Description} The trailing zero count value of the \emph{registerZ} is stored into the \emph{registerW}. When applied to zero, the result is 64.


\paragraph{Execution} {Reservation table \textsf{ALU\_TINY} (Section~\ref{sec:reservation-ALU-TINY})}
~
{\begin{lstlisting}
stage RR:
  new argument2 = GPR[registerZ];
  new result1 = _ZX_64(_CTZ_64(argument2));
stage E1:
  GPR[registerW] = result1;
\end{lstlisting}}
\newpage
\subsubsection[{\small CTZHQ: Count Trailing Zeroes Half Word Quadruple}]{CTZHQ\hfill Count Trailing Zeroes Half Word Quadruple}
\label{instr:CTZHQ}\index{ctzhq}
 
\paragraph{Encoding} Format \textsf{ALU\_HQBWRE} (Section~\ref{sec:format-ALU-HQBWRE}), \verb|alucode8 = 0111|\\

\bitfields{ 1/P, 2/11, 1/1, 4/1111, 5/registerWe, 1/0, 2/10, 3/101, 1/1, 4/0111, 1/--, 1/--, 5/registerZe, 1/0 }


\paragraph{Syntax} \texttt{ctzhq} \emph{registerWe} = \emph{registerZe}

\paragraph{Description} The trailing zero count values of the \emph{registerZe} interpreted as half word quadruple are stored into the \emph{registerWe}. When applied to zero at the half word level, the result is 16.


\paragraph{Execution} {Reservation table \textsf{ALU\_LITE} (Section~\ref{sec:reservation-ALU-LITE})}
~
{\begin{lstlisting}
stage RR:
  new argument2 = GPR[registerZe<<1];
  for (I = 0; I < 4; I += 1) {
    result1.16[I] = _ZX_16(_CTZ_16(argument2.16[I]))
  };
stage E1:
  GPR[registerWe<<1] = result1;
\end{lstlisting}}
\newpage

\paragraph{Encoding} Format \textsf{ALU\_HQBWRO} (Section~\ref{sec:format-ALU-HQBWRO}), \verb|alucode8 = 0111|\\

\bitfields{ 1/P, 2/11, 1/1, 4/1111, 5/registerWo, 1/1, 2/10, 3/101, 1/1, 4/0111, 1/--, 1/--, 5/registerZo, 1/0 }


\paragraph{Syntax} \texttt{ctzhq} \emph{registerWo} = \emph{registerZo}

\paragraph{Description} The trailing zero count values of the \emph{registerZo} interpreted as half word quadruple are stored into the \emph{registerWo}. When applied to zero at the half word level, the result is 16.


\paragraph{Execution} {Reservation table \textsf{ALU\_LITE} (Section~\ref{sec:reservation-ALU-LITE})}
~
{\begin{lstlisting}
stage RR:
  new argument2 = GPR[(registerZo<<1)+1];
  for (I = 0; I < 4; I += 1) {
    result1.16[I] = _ZX_16(_CTZ_16(argument2.16[I]))
  };
stage E1:
  GPR[(registerWo<<1)+1] = result1;
\end{lstlisting}}
\newpage
\subsubsection[{\small CTZW: Count Trailing Zero Word}]{CTZW\hfill Count Trailing Zero Word}
\label{instr:CTZW}\index{ctzw}
 
\paragraph{Encoding} Format \textsf{ALU\_BWRW} (Section~\ref{sec:format-ALU-BWRW}), \verb|alucode8 = 0111|\\

\bitfields{ 1/P, 2/11, 1/0, 4/1111, 6/registerW, 2/11, 3/101, 1/signextw, 4/0111, 1/0, 1/0, 6/registerZ }


\paragraph{Syntax} \texttt{ctzw} \emph{registerW} = \emph{registerZ}

\paragraph{Description} The trailing zero count value of the \emph{registerZ} is stored into the \emph{registerW}. When applied to zero, the result is 32.


\paragraph{Execution} {Reservation table \textsf{ALU\_TINY} (Section~\ref{sec:reservation-ALU-TINY})}
~
{\begin{lstlisting}
stage RR:
  new argument2 = GPR[registerZ];
  new result1 = _ZX_32(_CTZ_32(argument2));
stage E1:
  GPR[registerW] = _ZX_32(result1);
\end{lstlisting}}
\newpage
\subsubsection[{\small CTZWP: Count Trailing Zeroes Word Pair}]{CTZWP\hfill Count Trailing Zeroes Word Pair}
\label{instr:CTZWP}\index{ctzwp}
 
\paragraph{Encoding} Format \textsf{ALU\_WPBWRE} (Section~\ref{sec:format-ALU-WPBWRE}), \verb|alucode8 = 0111|\\

\bitfields{ 1/P, 2/11, 1/1, 4/1111, 5/registerWe, 1/0, 2/10, 3/101, 1/0, 4/0111, 1/--, 1/--, 5/registerZe, 1/1 }


\paragraph{Syntax} \texttt{ctzwp} \emph{registerWe} = \emph{registerZe}

\paragraph{Description} The trailing zero count values of the \emph{registerZe} interpreted as word pair are stored into the \emph{registerWe}. When applied to zero at the word level, the result is 32.


\paragraph{Execution} {Reservation table \textsf{ALU\_LITE} (Section~\ref{sec:reservation-ALU-LITE})}
~
{\begin{lstlisting}
stage RR:
  new argument2 = GPR[registerZe<<1];
  for (I = 0; I < 2; I += 1) {
    result1.32[I] = _ZX_32(_CTZ_32(argument2.32[I]))
  };
stage E1:
  GPR[registerWe<<1] = result1;
\end{lstlisting}}
\newpage

\paragraph{Encoding} Format \textsf{ALU\_WPBWRO} (Section~\ref{sec:format-ALU-WPBWRO}), \verb|alucode8 = 0111|\\

\bitfields{ 1/P, 2/11, 1/1, 4/1111, 5/registerWo, 1/1, 2/10, 3/101, 1/0, 4/0111, 1/--, 1/--, 5/registerZo, 1/1 }


\paragraph{Syntax} \texttt{ctzwp} \emph{registerWo} = \emph{registerZo}

\paragraph{Description} The trailing zero count values of the \emph{registerZo} interpreted as word pair are stored into the \emph{registerWo}. When applied to zero at the word level, the result is 32.


\paragraph{Execution} {Reservation table \textsf{ALU\_LITE} (Section~\ref{sec:reservation-ALU-LITE})}
~
{\begin{lstlisting}
stage RR:
  new argument2 = GPR[(registerZo<<1)+1];
  for (I = 0; I < 2; I += 1) {
    result1.32[I] = _ZX_32(_CTZ_32(argument2.32[I]))
  };
stage E1:
  GPR[(registerWo<<1)+1] = result1;
\end{lstlisting}}
\newpage
\subsubsection[{\small DIVMODD: Integer Division Modulus Double Word}]{DIVMODD\hfill Integer Division Modulus Double Word}
\label{instr:DIVMODD}\index{divmodd}
 
\paragraph{Encoding} Format \textsf{ALU\_DDMWRR} (Section~\ref{sec:format-ALU-DDMWRR}), \verb|exucode2 = 1100|\\

\bitfields{ 1/P, 2/11, 1/0, 4/1100, 5/registerM, 1/--, 2/10, 3/101, 1/0, 6/registerY, 6/registerZ }


\paragraph{Syntax} \texttt{divmodd} \emph{registerM} = \emph{registerZ}, \emph{registerY}

\paragraph{Description} The dividend and the divisor assuming 64-bit signed integers in the \emph{registerZ} and the \emph{registerY}. The quotient and the remainder are computed and the corresponding 64 bits integers and packed into the 128 bits of the \emph{registerM}. Division by zero returns zero. Modulus by zero returns zero. Division of LONG\_MIN by -1 returns LONG\_MIN. Modulus of LONG\_MIN by -1 returns zero.


\paragraph{Execution} {Reservation table \textsf{ALU\_FULL} (Section~\ref{sec:reservation-ALU-FULL})}
~
{\begin{lstlisting}
stage RR:
  new argument3 = GPR[registerY];
  new argument2 = GPR[registerZ];
  if (argument3.64[0] == 0) {
    new result1 = 0;
  } else if ((argument2.64[0] == 0x8000000000000000) && (argument3.64[0] == 0xFFFFFFFFFFFFFFFF)) {
    result1.64[0] = argument2.64[0];
    result1.64[1] = 0;
  } else {
    new dividend = _SX_64(argument2.64[0]);
    new divisor = _SX_64(argument3.64[0]);
    result1.64[0] = dividend / divisor;
    result1.64[1] = dividend % divisor;
  }
stage SF:
  PGR[registerM] = result1;
\end{lstlisting}}
\newpage
\subsubsection[{\small DIVMODUD: Integer Division Modulus Unsigned Double Word}]{DIVMODUD\hfill Integer Division Modulus Unsigned Double Word}
\label{instr:DIVMODUD}\index{divmodud}
 
\paragraph{Encoding} Format \textsf{ALU\_DDMWRR} (Section~\ref{sec:format-ALU-DDMWRR}), \verb|exucode2 = 1101|\\

\bitfields{ 1/P, 2/11, 1/0, 4/1101, 5/registerM, 1/--, 2/10, 3/101, 1/0, 6/registerY, 6/registerZ }


\paragraph{Syntax} \texttt{divmodud} \emph{registerM} = \emph{registerZ}, \emph{registerY}

\paragraph{Description} The dividend and the divisor assuming 64-bit unsigned integers in the \emph{registerZ} and the \emph{registerY}. The quotient and the remainder are computed and the corresponding 64-bit integers and packed into the 128 bits of the \emph{registerM}. Quotient of division by zero returns zero. Modulus by zero returns zero.


\paragraph{Execution} {Reservation table \textsf{ALU\_FULL} (Section~\ref{sec:reservation-ALU-FULL})}
~
{\begin{lstlisting}
stage RR:
  new argument3 = GPR[registerY];
  new argument2 = GPR[registerZ];
  if (argument3.64[0] == 0) {
    new result1 = 0;
  } else {
    new dividend = _ZX_64(argument2.64[0]);
    new divisor = _ZX_64(argument3.64[0]);
    result1.64[0] = dividend / divisor;
    result1.64[1] = dividend % divisor;
  }
stage SF:
  PGR[registerM] = result1;
\end{lstlisting}}
\newpage
\subsubsection[{\small DIVMODUW: Integer Division Modulus Unsigned Word}]{DIVMODUW\hfill Integer Division Modulus Unsigned Word}
\label{instr:DIVMODUW}\index{divmoduw}
 
\paragraph{Encoding} Format \textsf{ALU\_WDMWRR} (Section~\ref{sec:format-ALU-WDMWRR}), \verb|exucode2 = 1101|\\

\bitfields{ 1/P, 2/11, 1/0, 4/1101, 5/registerM, 1/--, 2/11, 3/101, 1/signextw, 6/registerY, 6/registerZ }


\paragraph{Syntax} \texttt{divmoduw}\emph{signextw} \emph{registerM} = \emph{registerZ}, \emph{registerY}

\paragraph{Description} The dividend and the divisor assuming 32-bit unsigned integers in the \emph{registerZ} and the \emph{registerY}. The quotient and the remainder are computed and the corresponding 32-bit integers are zero-extended and packed into the 128 bits of the \emph{registerM}. Quotient of division by zero returns zero. Modulus by zero returns zero.


\paragraph{Modifier(s)}
\noindent\begin{itemize}
\item signextw (cf. signextw in section \ref{par:modifier-signextw} on page \pageref{par:modifier-signextw})
\end{itemize}

\paragraph{Execution} {Reservation table \textsf{ALU\_FULL} (Section~\ref{sec:reservation-ALU-FULL})}
~
{\begin{lstlisting}
stage ID:
  new signextw = _ZX_1(signextw);
stage RR:
  new argument3 = GPR[registerY];
  new argument2 = GPR[registerZ];
  if (argument3.32[0] == 0) {
    new result1 = 0;
  } else {
    new dividend = _ZX_32(argument2.32[0]);
    new divisor = _ZX_32(argument3.32[0]);
    result1.64[0] = _ZX_32(dividend / divisor);
    result1.64[1] = _ZX_32(dividend % divisor);
  }
  if (signextw) {
    result1.64[0] = _SX_32(result1.64[0]);
    result1.64[1] = _SX_32(result1.64[1]);
  }
stage SF:
  PGR[registerM] = result1;
\end{lstlisting}}
\newpage
\subsubsection[{\small DIVMODW: Integer Division Modulus Word}]{DIVMODW\hfill Integer Division Modulus Word}
\label{instr:DIVMODW}\index{divmodw}
 
\paragraph{Encoding} Format \textsf{ALU\_WDMWRR} (Section~\ref{sec:format-ALU-WDMWRR}), \verb|exucode2 = 1100|\\

\bitfields{ 1/P, 2/11, 1/0, 4/1100, 5/registerM, 1/--, 2/11, 3/101, 1/signextw, 6/registerY, 6/registerZ }


\paragraph{Syntax} \texttt{divmodw}\emph{signextw} \emph{registerM} = \emph{registerZ}, \emph{registerY}

\paragraph{Description} The dividend and the divisor assuming 32-bit signed integers in the \emph{registerZ} and the \emph{registerY}. The quotient and the remainder are computed and the corresponding 32-bit integers are zero-extended and packed into the 128 bits of the \emph{registerM}. Division by zero returns zero. Modulus by zero returns zero. Division of INT\_MIN by -1 returns INT\_MIN. Modulus of INT\_MIN by -1 returns zero.


\paragraph{Modifier(s)}
\noindent\begin{itemize}
\item signextw (cf. signextw in section \ref{par:modifier-signextw} on page \pageref{par:modifier-signextw})
\end{itemize}

\paragraph{Execution} {Reservation table \textsf{ALU\_FULL} (Section~\ref{sec:reservation-ALU-FULL})}
~
{\begin{lstlisting}
stage ID:
  new signextw = _ZX_1(signextw);
stage RR:
  new argument3 = GPR[registerY];
  new argument2 = GPR[registerZ];
  if (argument3.32[0] == 0) {
    new result1 = 0;
  } else if ((argument2.32[0] == 0x80000000) && (argument3.32[0] == 0xFFFFFFFF)) {
    result1.64[0] = _ZX_32(argument2.32[0]);
    result1.64[1] = 0;
  } else {
    new dividend = _SX_32(argument2.32[0]);
    new divisor = _SX_32(argument3.32[0]);
    result1.64[0] = _ZX_32(dividend / divisor);
    result1.64[1] = _ZX_32(dividend % divisor);
  }
  if (signextw) {
    result1.64[0] = _SX_32(result1.64[0]);
    result1.64[1] = _SX_32(result1.64[1]);
  }
stage SF:
  PGR[registerM] = result1;
\end{lstlisting}}
\newpage
\subsubsection[{\small EORD: Bitwise Exclusive Or Double Word}]{EORD\hfill Bitwise Exclusive Or Double Word}
\label{instr:EORD}\index{eord}
 
\paragraph{Encoding} Format \textsf{ALU\_DWRR0} (Section~\ref{sec:format-ALU-DWRR0}), \verb|exucode2 = 1100|\\

\bitfields{ 1/P, 2/11, 1/0, 4/1100, 6/registerW, 2/10, 3/000, 1/0, 6/registerY, 6/registerZ }


\paragraph{Syntax} \texttt{eord} \emph{registerW} = \emph{registerZ}, \emph{registerY}

\paragraph{Description} Bitwise exclusive or of the \emph{registerZ} with the \emph{registerY}. The result is stored into the \emph{registerW}.


\paragraph{Execution} {Reservation table \textsf{ALU\_TINY} (Section~\ref{sec:reservation-ALU-TINY})}
~
{\begin{lstlisting}
stage RR:
  new argument3 = GPR[registerY];
  new argument2 = GPR[registerZ];
  new result1 = argument2 ^ argument3;
stage E1:
  GPR[registerW] = result1;
\end{lstlisting}}
\newpage

\paragraph{Encoding} Format \textsf{ALU\_DWRR0.M} (Section~\ref{sec:format-ALU-DWRR0.M}), \verb|exucode2 = 1100|\\

\bitfields{ 1/P, 2/00, 2/-- --, 27/upper27, 1/1, 2/11, 1/0, 4/1100, 6/registerW, 2/10, 3/000, 1/0, 1/splat32, 5/lower5, 6/registerZ }


\paragraph{Syntax} \texttt{eord} \emph{registerW} = \emph{registerZ}, \emph{upper27\_\discretionary{}{}{}lower5}\emph{splat32}

\paragraph{Description} Bitwise exclusive or of the \emph{registerZ} with the \emph{upper27\_\discretionary{}{}{}lower5}. The result is stored into the \emph{registerW}.


\paragraph{Modifier(s)}
\noindent\begin{itemize}
\item splat32 (cf. splat32 in section \ref{par:modifier-splat32} on page \pageref{par:modifier-splat32})
\end{itemize}

\paragraph{Execution} {Reservation table \textsf{ALU\_TINY.X} (Section~\ref{sec:reservation-ALU-TINY.X})}
~
{\begin{lstlisting}
stage ID:
  new splat32 = _ZX_1(splat32);
stage RR:
  new argument3 = splat32 ? (_WX_32(upper27_lower5) << 32) | _ZX_32(_WX_32(upper27_lower5)) : _SX_32(_WX_32(upper27_lower5));
  new argument2 = GPR[registerZ];
  new result1 = argument2 ^ argument3;
stage E1:
  GPR[registerW] = result1;
\end{lstlisting}}
\newpage

\paragraph{Encoding} Format \textsf{ALU\_DWRI} (Section~\ref{sec:format-ALU-DWRI}), \verb|exucode2 = 1100|\\

\bitfields{ 1/P, 2/11, 1/0, 4/1100, 6/registerW, 2/00, 10/signed10, 6/registerZ }


\paragraph{Syntax} \texttt{eord} \emph{registerW} = \emph{registerZ}, \emph{signed10}

\paragraph{Description} Bitwise exclusive or of the \emph{registerZ} with the \emph{signed10}. The result is stored into the \emph{registerW}.


\paragraph{Execution} {Reservation table \textsf{ALU\_TINY} (Section~\ref{sec:reservation-ALU-TINY})}
~
{\begin{lstlisting}
stage RR:
  new argument3 = _SX_10(signed10);
  new argument2 = GPR[registerZ];
  new result1 = argument2 ^ argument3;
stage E1:
  GPR[registerW] = result1;
\end{lstlisting}}
\newpage

\paragraph{Encoding} Format \textsf{ALU\_DWRI.X} (Section~\ref{sec:format-ALU-DWRI.X}), \verb|exucode2 = 1100|\\

\bitfields{ 1/P, 2/00, 2/-- --, 27/upper27, 1/1, 2/11, 1/0, 4/1100, 6/registerW, 2/00, 10/lower10, 6/registerZ }


\paragraph{Syntax} \texttt{eord} \emph{registerW} = \emph{registerZ}, \emph{upper27\_\discretionary{}{}{}lower10}

\paragraph{Description} Bitwise exclusive or of the \emph{registerZ} with the \emph{upper27\_\discretionary{}{}{}lower10}. The result is stored into the \emph{registerW}.


\paragraph{Execution} {Reservation table \textsf{ALU\_TINY.X} (Section~\ref{sec:reservation-ALU-TINY.X})}
~
{\begin{lstlisting}
stage RR:
  new argument3 = _SX_37(upper27_lower10);
  new argument2 = GPR[registerZ];
  new result1 = argument2 ^ argument3;
stage E1:
  GPR[registerW] = result1;
\end{lstlisting}}
\newpage

\paragraph{Encoding} Format \textsf{ALU\_DWRI.Y} (Section~\ref{sec:format-ALU-DWRI.Y}), \verb|exucode2 = 1100|\\

\bitfields{ 1/P, 2/00, 2/-- --, 27/extend27, 1/1, 2/00, 2/-- --, 27/upper27, 1/1, 2/11, 1/0, 4/1100, 6/registerW, 2/00, 10/lower10, 6/registerZ }


\paragraph{Syntax} \texttt{eord} \emph{registerW} = \emph{registerZ}, \emph{extend27\_\discretionary{}{}{}upper27\_\discretionary{}{}{}lower10}

\paragraph{Description} Bitwise exclusive or of the \emph{registerZ} with the \emph{extend27\_\discretionary{}{}{}upper27\_\discretionary{}{}{}lower10}. The result is stored into the \emph{registerW}.


\paragraph{Execution} {Reservation table \textsf{ALU\_TINY.Y} (Section~\ref{sec:reservation-ALU-TINY.Y})}
~
{\begin{lstlisting}
stage RR:
  new argument3 = _WX_64(extend27_upper27_lower10);
  new argument2 = GPR[registerZ];
  new result1 = argument2 ^ argument3;
stage E1:
  GPR[registerW] = result1;
\end{lstlisting}}
\newpage
\subsubsection[{\small EORQ: Bitwise Exclusive Or Quadruple Word}]{EORQ\hfill Bitwise Exclusive Or Quadruple Word}
\label{instr:EORQ}\index{eorq}
 
\paragraph{Encoding} Format \textsf{ALU\_QWRR0} (Section~\ref{sec:format-ALU-QWRR0}), \verb|exucode2 = 1100|\\

\bitfields{ 1/P, 2/11, 1/0, 4/1100, 5/registerM, 1/--, 2/10, 3/000, 1/1, 5/registerO, 1/--, 5/registerP, 1/-- }


\paragraph{Syntax} \texttt{eorq} \emph{registerM} = \emph{registerP}, \emph{registerO}

\paragraph{Description} Bitwise exclusive or of the \emph{registerP} with the \emph{registerO}. The result is stored into the \emph{registerM}.


\paragraph{Execution} {Reservation table \textsf{ALU\_LITE} (Section~\ref{sec:reservation-ALU-LITE})}
~
{\begin{lstlisting}
stage RR:
  new argument3 = PGR[registerO];
  new argument2 = PGR[registerP];
  new result1 = argument2 ^ argument3;
stage E1:
  PGR[registerM] = result1;
\end{lstlisting}}
\newpage

\paragraph{Encoding} Format \textsf{ALU\_QWRR0.M} (Section~\ref{sec:format-ALU-QWRR0.M}), \verb|exucode2 = 1100|\\

\bitfields{ 1/P, 2/00, 2/-- --, 27/upper27, 1/1, 2/11, 1/0, 4/1100, 5/registerM, 1/--, 2/10, 3/000, 1/1, 1/splat32, 5/lower5, 5/registerP, 1/1 }


\paragraph{Syntax} \texttt{eorq} \emph{registerM} = \emph{registerP}, \emph{upper27\_\discretionary{}{}{}lower5}\emph{splat32}

\paragraph{Description} Bitwise exclusive or of the \emph{registerP} with the \emph{upper27\_\discretionary{}{}{}lower5}. The result is stored into the \emph{registerM}.


\paragraph{Modifier(s)}
\noindent\begin{itemize}
\item splat32 (cf. splat32 in section \ref{par:modifier-splat32} on page \pageref{par:modifier-splat32})
\end{itemize}

\paragraph{Execution} {Reservation table \textsf{ALU\_LITE.X} (Section~\ref{sec:reservation-ALU-LITE.X})}
~
{\begin{lstlisting}
stage ID:
  new splat32 = _ZX_1(splat32);
stage RR:
  new argument3 = splat32 ? (_WX_32(upper27_lower5) << 32) | _ZX_32(_WX_32(upper27_lower5)) : _SX_32(_WX_32(upper27_lower5));
  new argument2 = PGR[registerP];
  new result1 = argument2 ^ argument3;
stage E1:
  PGR[registerM] = result1;
\end{lstlisting}}
\newpage
\subsubsection[{\small EORW: Bitwise Exclusive Or Word}]{EORW\hfill Bitwise Exclusive Or Word}
\label{instr:EORW}\index{eorw}
 
\paragraph{Encoding} Format \textsf{ALU\_WWRR0} (Section~\ref{sec:format-ALU-WWRR0}), \verb|exucode2 = 1100|\\

\bitfields{ 1/P, 2/11, 1/0, 4/1100, 6/registerW, 2/11, 3/000, 1/signextw, 6/registerY, 6/registerZ }


\paragraph{Syntax} \texttt{eorw}\emph{signextw} \emph{registerW} = \emph{registerZ}, \emph{registerY}

\paragraph{Description} Bitwise exclusive or of the \emph{registerZ} with the \emph{registerY}. The result with the 32 upper bits cleared is stored into the \emph{registerW}.


\paragraph{Modifier(s)}
\noindent\begin{itemize}
\item signextw (cf. signextw in section \ref{par:modifier-signextw} on page \pageref{par:modifier-signextw})
\end{itemize}

\paragraph{Execution} {Reservation table \textsf{ALU\_TINY} (Section~\ref{sec:reservation-ALU-TINY})}
~
{\begin{lstlisting}
stage ID:
  new signextw = _ZX_1(signextw);
stage RR:
  new argument3 = GPR[registerY];
  new argument2 = GPR[registerZ];
  new result1 = argument2 ^ argument3;
stage E1:
  GPR[registerW] = signextw ? _SX_32(result1) : _ZX_32(result1);
\end{lstlisting}}
\newpage

\paragraph{Encoding} Format \textsf{ALU\_WWRR0.W} (Section~\ref{sec:format-ALU-WWRR0.W}), \verb|exucode2 = 1100|\\

\bitfields{ 1/P, 2/00, 2/-- --, 27/upper27, 1/1, 2/11, 1/0, 4/1100, 6/registerW, 2/11, 3/000, 1/signextw, 1/0, 5/lower5, 6/registerZ }


\paragraph{Syntax} \texttt{eorw}\emph{signextw} \emph{registerW} = \emph{registerZ}, \emph{upper27\_\discretionary{}{}{}lower5}

\paragraph{Description} Bitwise exclusive or of the \emph{registerZ} with the \emph{upper27\_\discretionary{}{}{}lower5}. The result with the 32 upper bits cleared is stored into the \emph{registerW}.


\paragraph{Modifier(s)}
\noindent\begin{itemize}
\item signextw (cf. signextw in section \ref{par:modifier-signextw} on page \pageref{par:modifier-signextw})
\end{itemize}

\paragraph{Execution} {Reservation table \textsf{ALU\_TINY.X} (Section~\ref{sec:reservation-ALU-TINY.X})}
~
{\begin{lstlisting}
stage ID:
  new signextw = _ZX_1(signextw);
stage RR:
  new argument3 = _WX_32(upper27_lower5);
  new argument2 = GPR[registerZ];
  new result1 = argument2 ^ argument3;
stage E1:
  GPR[registerW] = signextw ? _SX_32(result1) : _ZX_32(result1);
\end{lstlisting}}
\newpage
\subsubsection[{\small EXTFS: Extract Bit Field and Sign Extend to Double Word}]{EXTFS\hfill Extract Bit Field and Sign Extend to Double Word}
\label{instr:EXTFS}\index{extfs}
 
\paragraph{Encoding} Format \textsf{ALU\_BFWR} (Section~\ref{sec:format-ALU-BFWR}), \verb|wricode1 = 10|\\

\bitfields{ 1/P, 2/11, 1/0, 2/10, 2/stopbit2, 6/registerW, 2/01, 4/stopbit4, 6/startbit, 6/registerZ }


\paragraph{Syntax} \texttt{extfs} \emph{registerW} = \emph{registerZ}, \emph{stopbit2\_\discretionary{}{}{}stopbit4}, \emph{startbit}

\paragraph{Description} The start bit position and the stop bit position are extracted from the \emph{startbit} and the \emph{stopbit2\_\discretionary{}{}{}stopbit4} respectively. The bit-field of the \emph{registerZ} extending from the start bit position to the stop bit position is extracted, sign-extended from the stop bit position, shifted right by the start bit position and stored into the \emph{registerW}. Behavior is undefined if the start bit position is greater than the stop bit position.


\paragraph{Execution} {Reservation table \textsf{ALU\_TINY} (Section~\ref{sec:reservation-ALU-TINY})}
~
{\begin{lstlisting}
stage ID:
  new startbit = _ZX_6(startbit);
  new stopbit = _ZX_6(stopbit2_stopbit4);
stage RR:
  new bias = startbit <= stopbit;
  new mask = _ZX_64(((2 << stopbit) - ((2 - bias) << startbit)) + (bias - 1));
  new argument2 = GPR[registerZ];
  new masked1 = _ZX_64(argument2 & mask);
  new masked2 = _SX_64(argument2 | ~mask);
  new negative = argument2 & (mask & ~(mask >> 1));
  new masked = negative == 0 ? masked1 : masked2;
  new result1 = masked >> startbit;
stage E1:
  GPR[registerW] = result1;
\end{lstlisting}}
\newpage
\subsubsection[{\small EXTFZ: Extract Bit Field and Zero Extend to Double Word}]{EXTFZ\hfill Extract Bit Field and Zero Extend to Double Word}
\label{instr:EXTFZ}\index{extfz}
 
\paragraph{Encoding} Format \textsf{ALU\_BFWR} (Section~\ref{sec:format-ALU-BFWR}), \verb|wricode1 = 01|\\

\bitfields{ 1/P, 2/11, 1/0, 2/01, 2/stopbit2, 6/registerW, 2/01, 4/stopbit4, 6/startbit, 6/registerZ }


\paragraph{Syntax} \texttt{extfz} \emph{registerW} = \emph{registerZ}, \emph{stopbit2\_\discretionary{}{}{}stopbit4}, \emph{startbit}

\paragraph{Description} The start bit position and the stop bit position are extracted from the \emph{startbit} and the \emph{stopbit2\_\discretionary{}{}{}stopbit4} respectively. The bit-field of the \emph{registerZ} extending from the start bit position to the stop bit position is extracted, zero-extended after the stop bit position, shifted right by the start bit position and stored into the \emph{registerW}. Behavior is undefined if the start bit position is greater than the stop bit position.


\paragraph{Execution} {Reservation table \textsf{ALU\_TINY} (Section~\ref{sec:reservation-ALU-TINY})}
~
{\begin{lstlisting}
stage ID:
  new startbit = _ZX_6(startbit);
  new stopbit = _ZX_6(stopbit2_stopbit4);
stage RR:
  new bias = startbit <= stopbit;
  new mask = _ZX_64(((2 << stopbit) - ((2 - bias) << startbit)) + (bias - 1));
  new argument2 = GPR[registerZ];
  new masked = argument2 & mask;
  new result1 = masked >> startbit;
stage E1:
  GPR[registerW] = result1;
\end{lstlisting}}
\newpage
\subsubsection[{\small EXTLQBHO: Extend Lane Q7 Byte to Half Word Octuple}]{EXTLQBHO\hfill Extend Lane Q7 Byte to Half Word Octuple}
\label{instr:EXTLQBHO}\index{extlqbho}
 
\paragraph{Encoding} Format \textsf{ALU\_HOEXTLWR} (Section~\ref{sec:format-ALU-HOEXTLWR}), \verb|alucode8 = 1110|\\

\bitfields{ 1/P, 2/11, 1/1, 4/1111, 5/registerM, 1/oddlanes, 2/10, 3/101, 1/1, 4/1110, 1/--, 1/--, 5/registerP, 1/0 }


\paragraph{Syntax} \texttt{extlqbho}\emph{oddlanes} \emph{registerM} = \emph{registerP}

\paragraph{Description} The \emph{registerP} is interpreted as sixteen bytes packed into 128 bits. The even or odd bytes of the \emph{registerP} are selected, depending on \emph{oddlanes}. These eight bytes interpreted as Q7 numbers are mapped to Q15 numbers, packed and stored into the \emph{registerM}.


\paragraph{Modifier(s)}
\noindent\begin{itemize}
\item oddlanes (cf. oddlanes in section \ref{par:modifier-oddlanes} on page \pageref{par:modifier-oddlanes})
\end{itemize}

\paragraph{Execution} {Reservation table \textsf{ALU\_LITE} (Section~\ref{sec:reservation-ALU-LITE})}
~
{\begin{lstlisting}
stage ID:
  new oddlanes = _ZX_1(oddlanes);
stage RR:
  new argument2 = PGR[registerP];
  if (oddlanes) {
    argument2 = argument2 >> 8;
  }
  for (I = 0; I < 8; I += 1) {
    result1.16[I] = argument2.16[I] << 8
  };
stage E1:
  PGR[registerM] = result1;
\end{lstlisting}}
\newpage
\subsubsection[{\small EXTLQHWQ: Extend Lane Q15 Half Word to Word Quadruple}]{EXTLQHWQ\hfill Extend Lane Q15 Half Word to Word Quadruple}
\label{instr:EXTLQHWQ}\index{extlqhwq}
 
\paragraph{Encoding} Format \textsf{ALU\_WQEXTLWR} (Section~\ref{sec:format-ALU-WQEXTLWR}), \verb|alucode8 = 1110|\\

\bitfields{ 1/P, 2/11, 1/1, 4/1111, 5/registerM, 1/oddlanes, 2/10, 3/101, 1/0, 4/1110, 1/--, 1/--, 5/registerP, 1/1 }


\paragraph{Syntax} \texttt{extlqhwq}\emph{oddlanes} \emph{registerM} = \emph{registerP}

\paragraph{Description} The \emph{registerP} is interpreted as eight half words packed into 128 bits. The even or odd half words of the \emph{registerP} are selected, depending on \emph{oddlanes}. These four half words interpreted as Q15 numbers are mapped to Q31 numbers, packed and stored into the \emph{registerM}.


\paragraph{Modifier(s)}
\noindent\begin{itemize}
\item oddlanes (cf. oddlanes in section \ref{par:modifier-oddlanes} on page \pageref{par:modifier-oddlanes})
\end{itemize}

\paragraph{Execution} {Reservation table \textsf{ALU\_LITE} (Section~\ref{sec:reservation-ALU-LITE})}
~
{\begin{lstlisting}
stage ID:
  new oddlanes = _ZX_1(oddlanes);
stage RR:
  new argument2 = PGR[registerP];
  if (oddlanes) {
    argument2 = argument2 >> 16;
  }
  for (I = 0; I < 4; I += 1) {
    result1.32[I] = argument2.32[I] << 16
  };
stage E1:
  PGR[registerM] = result1;
\end{lstlisting}}
\newpage
\subsubsection[{\small EXTLQNBX: Extend Lane Q3 Nibble to Byte Hexadecuple}]{EXTLQNBX\hfill Extend Lane Q3 Nibble to Byte Hexadecuple}
\label{instr:EXTLQNBX}\index{extlqnbx}
 
\paragraph{Encoding} Format \textsf{ALU\_BXEXTLWR} (Section~\ref{sec:format-ALU-BXEXTLWR}), \verb|alucode8 = 1110|\\

\bitfields{ 1/P, 2/11, 1/1, 4/1111, 5/registerM, 1/oddlanes, 2/10, 3/101, 1/1, 4/1110, 1/--, 1/--, 5/registerP, 1/1 }


\paragraph{Syntax} \texttt{extlqnbx}\emph{oddlanes} \emph{registerM} = \emph{registerP}

\paragraph{Description} The \emph{registerP} is interpreted as thirty-two nibbles packed into 128 bits. The even or odd nibbles of the \emph{registerP} are selected, depending on \emph{oddlanes}. These sixteen nibbles interpreted as Q4 numbers are mapped to Q7 numbers, packed and stored into the \emph{registerM}.


\paragraph{Modifier(s)}
\noindent\begin{itemize}
\item oddlanes (cf. oddlanes in section \ref{par:modifier-oddlanes} on page \pageref{par:modifier-oddlanes})
\end{itemize}

\paragraph{Execution} {Reservation table \textsf{ALU\_LITE} (Section~\ref{sec:reservation-ALU-LITE})}
~
{\begin{lstlisting}
stage ID:
  new oddlanes = _ZX_1(oddlanes);
stage RR:
  new argument2 = PGR[registerP];
  if (oddlanes) {
    argument2 = argument2 >> 4;
  }
  for (I = 0; I < 16; I += 1) {
    result1.8[I] = argument2.8[I] << 4
  };
stage E1:
  PGR[registerM] = result1;
\end{lstlisting}}
\newpage
\subsubsection[{\small EXTLQWDP: Extend Lane Q31 Word to Double Word Pair}]{EXTLQWDP\hfill Extend Lane Q31 Word to Double Word Pair}
\label{instr:EXTLQWDP}\index{extlqwdp}
 
\paragraph{Encoding} Format \textsf{ALU\_DPEXTLWR} (Section~\ref{sec:format-ALU-DPEXTLWR}), \verb|alucode8 = 1110|\\

\bitfields{ 1/P, 2/11, 1/1, 4/1111, 5/registerM, 1/oddlanes, 2/10, 3/101, 1/0, 4/1110, 1/--, 1/--, 5/registerP, 1/0 }


\paragraph{Syntax} \texttt{extlqwdp}\emph{oddlanes} \emph{registerM} = \emph{registerP}

\paragraph{Description} The \emph{registerP} is interpreted as four words packed into 128 bits. The even or odd words of the \emph{registerP} are selected, depending on \emph{oddlanes}. These two words interpreted as Q31 numbers are mapped to Q63 numbers, packed and stored into the \emph{registerM}.


\paragraph{Modifier(s)}
\noindent\begin{itemize}
\item oddlanes (cf. oddlanes in section \ref{par:modifier-oddlanes} on page \pageref{par:modifier-oddlanes})
\end{itemize}

\paragraph{Execution} {Reservation table \textsf{ALU\_LITE} (Section~\ref{sec:reservation-ALU-LITE})}
~
{\begin{lstlisting}
stage ID:
  new oddlanes = _ZX_1(oddlanes);
stage RR:
  new argument2 = PGR[registerP];
  if (oddlanes) {
    argument2 = argument2 >> 32;
  }
  for (I = 0; I < 2; I += 1) {
    result1.64[I] = argument2.64[I] << 32
  };
stage E1:
  PGR[registerM] = result1;
\end{lstlisting}}
\newpage
\subsubsection[{\small EXTLSBHO: Extend Lane with Sign Byte to Half Word Octuple}]{EXTLSBHO\hfill Extend Lane with Sign Byte to Half Word Octuple}
\label{instr:EXTLSBHO}\index{extlsbho}
 
\paragraph{Encoding} Format \textsf{ALU\_HOEXTLWR} (Section~\ref{sec:format-ALU-HOEXTLWR}), \verb|alucode8 = 1101|\\

\bitfields{ 1/P, 2/11, 1/1, 4/1111, 5/registerM, 1/oddlanes, 2/10, 3/101, 1/1, 4/1101, 1/--, 1/--, 5/registerP, 1/0 }


\paragraph{Syntax} \texttt{extlsbho}\emph{oddlanes} \emph{registerM} = \emph{registerP}

\paragraph{Description} The \emph{registerP} is interpreted as sixteen bytes packed into 128 bits. The even or odd bytes of the \emph{registerP} are selected, depending on \emph{oddlanes}. These eight bytes are sign-extended to 16 bits, packed and stored into the \emph{registerM}.


\paragraph{Modifier(s)}
\noindent\begin{itemize}
\item oddlanes (cf. oddlanes in section \ref{par:modifier-oddlanes} on page \pageref{par:modifier-oddlanes})
\end{itemize}

\paragraph{Execution} {Reservation table \textsf{ALU\_LITE} (Section~\ref{sec:reservation-ALU-LITE})}
~
{\begin{lstlisting}
stage ID:
  new oddlanes = _ZX_1(oddlanes);
stage RR:
  new argument2 = PGR[registerP];
  if (oddlanes) {
    argument2 = argument2 >> 8;
  }
  for (I = 0; I < 8; I += 1) {
    result1.16[I] = _SX_8(argument2.16[I])
  };
stage E1:
  PGR[registerM] = result1;
\end{lstlisting}}
\newpage
\subsubsection[{\small EXTLSHWQ: Extend Lane with Sign Half Word to Word Quadruple}]{EXTLSHWQ\hfill Extend Lane with Sign Half Word to Word Quadruple}
\label{instr:EXTLSHWQ}\index{extlshwq}
 
\paragraph{Encoding} Format \textsf{ALU\_WQEXTLWR} (Section~\ref{sec:format-ALU-WQEXTLWR}), \verb|alucode8 = 1101|\\

\bitfields{ 1/P, 2/11, 1/1, 4/1111, 5/registerM, 1/oddlanes, 2/10, 3/101, 1/0, 4/1101, 1/--, 1/--, 5/registerP, 1/1 }


\paragraph{Syntax} \texttt{extlshwq}\emph{oddlanes} \emph{registerM} = \emph{registerP}

\paragraph{Description} The \emph{registerP} is interpreted as eight half words packed into 128 bits. The even or odd half words of the \emph{registerP} are selected, depending on \emph{oddlanes}. These four half words are sign-extended to 32 bits, packed and stored into the \emph{registerM}.


\paragraph{Modifier(s)}
\noindent\begin{itemize}
\item oddlanes (cf. oddlanes in section \ref{par:modifier-oddlanes} on page \pageref{par:modifier-oddlanes})
\end{itemize}

\paragraph{Execution} {Reservation table \textsf{ALU\_LITE} (Section~\ref{sec:reservation-ALU-LITE})}
~
{\begin{lstlisting}
stage ID:
  new oddlanes = _ZX_1(oddlanes);
stage RR:
  new argument2 = PGR[registerP];
  if (oddlanes) {
    argument2 = argument2 >> 16;
  }
  for (I = 0; I < 4; I += 1) {
    result1.32[I] = _SX_16(argument2.32[I])
  };
stage E1:
  PGR[registerM] = result1;
\end{lstlisting}}
\newpage
\subsubsection[{\small EXTLSNBX: Extend Lane with Sign Nibble to Byte Hexadecuple}]{EXTLSNBX\hfill Extend Lane with Sign Nibble to Byte Hexadecuple}
\label{instr:EXTLSNBX}\index{extlsnbx}
 
\paragraph{Encoding} Format \textsf{ALU\_BXEXTLWR} (Section~\ref{sec:format-ALU-BXEXTLWR}), \verb|alucode8 = 1101|\\

\bitfields{ 1/P, 2/11, 1/1, 4/1111, 5/registerM, 1/oddlanes, 2/10, 3/101, 1/1, 4/1101, 1/--, 1/--, 5/registerP, 1/1 }


\paragraph{Syntax} \texttt{extlsnbx}\emph{oddlanes} \emph{registerM} = \emph{registerP}

\paragraph{Description} The \emph{registerP} is interpreted as thirty-two nibbles packed into 128 bits. The even or odd nibbles of the \emph{registerP} are selected, depending on \emph{oddlanes}. These sixteen nibbles are sign-extended to 8 bits, packed and stored into the \emph{registerM}.


\paragraph{Modifier(s)}
\noindent\begin{itemize}
\item oddlanes (cf. oddlanes in section \ref{par:modifier-oddlanes} on page \pageref{par:modifier-oddlanes})
\end{itemize}

\paragraph{Execution} {Reservation table \textsf{ALU\_LITE} (Section~\ref{sec:reservation-ALU-LITE})}
~
{\begin{lstlisting}
stage ID:
  new oddlanes = _ZX_1(oddlanes);
stage RR:
  new argument2 = PGR[registerP];
  if (oddlanes) {
    argument2 = argument2 >> 4;
  }
  for (I = 0; I < 16; I += 1) {
    result1.8[I] = _SX_4(argument2.8[I])
  };
stage E1:
  PGR[registerM] = result1;
\end{lstlisting}}
\newpage
\subsubsection[{\small EXTLSWDP: Extend Lane with Sign Word to Double Word Pair}]{EXTLSWDP\hfill Extend Lane with Sign Word to Double Word Pair}
\label{instr:EXTLSWDP}\index{extlswdp}
 
\paragraph{Encoding} Format \textsf{ALU\_DPEXTLWR} (Section~\ref{sec:format-ALU-DPEXTLWR}), \verb|alucode8 = 1101|\\

\bitfields{ 1/P, 2/11, 1/1, 4/1111, 5/registerM, 1/oddlanes, 2/10, 3/101, 1/0, 4/1101, 1/--, 1/--, 5/registerP, 1/0 }


\paragraph{Syntax} \texttt{extlswdp}\emph{oddlanes} \emph{registerM} = \emph{registerP}

\paragraph{Description} The \emph{registerP} is interpreted as four words packed into 128 bits. The even or odd words of the \emph{registerP} are selected, depending on \emph{oddlanes}. These two words are sign-extended to 64 bits, packed and stored into the \emph{registerM}.


\paragraph{Modifier(s)}
\noindent\begin{itemize}
\item oddlanes (cf. oddlanes in section \ref{par:modifier-oddlanes} on page \pageref{par:modifier-oddlanes})
\end{itemize}

\paragraph{Execution} {Reservation table \textsf{ALU\_LITE} (Section~\ref{sec:reservation-ALU-LITE})}
~
{\begin{lstlisting}
stage ID:
  new oddlanes = _ZX_1(oddlanes);
stage RR:
  new argument2 = PGR[registerP];
  if (oddlanes) {
    argument2 = argument2 >> 32;
  }
  for (I = 0; I < 2; I += 1) {
    result1.64[I] = _SX_32(argument2.64[I])
  };
stage E1:
  PGR[registerM] = result1;
\end{lstlisting}}
\newpage
\subsubsection[{\small EXTLZBHO: Extend Lane with Zero Byte to Half Word Octuple}]{EXTLZBHO\hfill Extend Lane with Zero Byte to Half Word Octuple}
\label{instr:EXTLZBHO}\index{extlzbho}
 
\paragraph{Encoding} Format \textsf{ALU\_HOEXTLWR} (Section~\ref{sec:format-ALU-HOEXTLWR}), \verb|alucode8 = 1100|\\

\bitfields{ 1/P, 2/11, 1/1, 4/1111, 5/registerM, 1/oddlanes, 2/10, 3/101, 1/1, 4/1100, 1/--, 1/--, 5/registerP, 1/0 }


\paragraph{Syntax} \texttt{extlzbho}\emph{oddlanes} \emph{registerM} = \emph{registerP}

\paragraph{Description} The \emph{registerP} is interpreted as sixteen bytes packed into 128 bits. The even or odd bytes of the \emph{registerP} are selected, depending on \emph{oddlanes}. These eight bytes are zero-extended to 16 bits, packed and stored into the \emph{registerM}.


\paragraph{Modifier(s)}
\noindent\begin{itemize}
\item oddlanes (cf. oddlanes in section \ref{par:modifier-oddlanes} on page \pageref{par:modifier-oddlanes})
\end{itemize}

\paragraph{Execution} {Reservation table \textsf{ALU\_LITE} (Section~\ref{sec:reservation-ALU-LITE})}
~
{\begin{lstlisting}
stage ID:
  new oddlanes = _ZX_1(oddlanes);
stage RR:
  new argument2 = PGR[registerP];
  if (oddlanes) {
    argument2 = argument2 >> 8;
  }
  for (I = 0; I < 8; I += 1) {
    result1.16[I] = _ZX_8(argument2.16[I])
  };
stage E1:
  PGR[registerM] = result1;
\end{lstlisting}}
\newpage
\subsubsection[{\small EXTLZHWQ: Extend Lane with Zero Half Word to Word Quadruple}]{EXTLZHWQ\hfill Extend Lane with Zero Half Word to Word Quadruple}
\label{instr:EXTLZHWQ}\index{extlzhwq}
 
\paragraph{Encoding} Format \textsf{ALU\_WQEXTLWR} (Section~\ref{sec:format-ALU-WQEXTLWR}), \verb|alucode8 = 1100|\\

\bitfields{ 1/P, 2/11, 1/1, 4/1111, 5/registerM, 1/oddlanes, 2/10, 3/101, 1/0, 4/1100, 1/--, 1/--, 5/registerP, 1/1 }


\paragraph{Syntax} \texttt{extlzhwq}\emph{oddlanes} \emph{registerM} = \emph{registerP}

\paragraph{Description} The \emph{registerP} is interpreted as eight half words packed into 128 bits. The even or odd half words of the \emph{registerP} are selected, depending on \emph{oddlanes}. These four half words are zero-extended to 32 bits, packed and stored into the \emph{registerM}.


\paragraph{Modifier(s)}
\noindent\begin{itemize}
\item oddlanes (cf. oddlanes in section \ref{par:modifier-oddlanes} on page \pageref{par:modifier-oddlanes})
\end{itemize}

\paragraph{Execution} {Reservation table \textsf{ALU\_LITE} (Section~\ref{sec:reservation-ALU-LITE})}
~
{\begin{lstlisting}
stage ID:
  new oddlanes = _ZX_1(oddlanes);
stage RR:
  new argument2 = PGR[registerP];
  if (oddlanes) {
    argument2 = argument2 >> 16;
  }
  for (I = 0; I < 4; I += 1) {
    result1.32[I] = _ZX_16(argument2.32[I])
  };
stage E1:
  PGR[registerM] = result1;
\end{lstlisting}}
\newpage
\subsubsection[{\small EXTLZNBX: Extend Lane with Zero Nibble to Byte Hexadecuple}]{EXTLZNBX\hfill Extend Lane with Zero Nibble to Byte Hexadecuple}
\label{instr:EXTLZNBX}\index{extlznbx}
 
\paragraph{Encoding} Format \textsf{ALU\_BXEXTLWR} (Section~\ref{sec:format-ALU-BXEXTLWR}), \verb|alucode8 = 1100|\\

\bitfields{ 1/P, 2/11, 1/1, 4/1111, 5/registerM, 1/oddlanes, 2/10, 3/101, 1/1, 4/1100, 1/--, 1/--, 5/registerP, 1/1 }


\paragraph{Syntax} \texttt{extlznbx}\emph{oddlanes} \emph{registerM} = \emph{registerP}

\paragraph{Description} The \emph{registerP} is interpreted as thirty-two nibbles packed into 128 bits. The even or odd nibbles of the \emph{registerP} are selected, depending on \emph{oddlanes}. These sixteen nibbles are zero-extended to 8 bits, packed and stored into the \emph{registerM}.


\paragraph{Modifier(s)}
\noindent\begin{itemize}
\item oddlanes (cf. oddlanes in section \ref{par:modifier-oddlanes} on page \pageref{par:modifier-oddlanes})
\end{itemize}

\paragraph{Execution} {Reservation table \textsf{ALU\_LITE} (Section~\ref{sec:reservation-ALU-LITE})}
~
{\begin{lstlisting}
stage ID:
  new oddlanes = _ZX_1(oddlanes);
stage RR:
  new argument2 = PGR[registerP];
  if (oddlanes) {
    argument2 = argument2 >> 4;
  }
  for (I = 0; I < 16; I += 1) {
    result1.8[I] = _ZX_4(argument2.8[I])
  };
stage E1:
  PGR[registerM] = result1;
\end{lstlisting}}
\newpage
\subsubsection[{\small EXTLZWDP: Extend Lane with Zero Word to Double Word Pair}]{EXTLZWDP\hfill Extend Lane with Zero Word to Double Word Pair}
\label{instr:EXTLZWDP}\index{extlzwdp}
 
\paragraph{Encoding} Format \textsf{ALU\_DPEXTLWR} (Section~\ref{sec:format-ALU-DPEXTLWR}), \verb|alucode8 = 1100|\\

\bitfields{ 1/P, 2/11, 1/1, 4/1111, 5/registerM, 1/oddlanes, 2/10, 3/101, 1/0, 4/1100, 1/--, 1/--, 5/registerP, 1/0 }


\paragraph{Syntax} \texttt{extlzwdp}\emph{oddlanes} \emph{registerM} = \emph{registerP}

\paragraph{Description} The \emph{registerP} is interpreted as four words packed into 128 bits. The even or odd words of the \emph{registerP} are selected, depending on \emph{oddlanes}. These two words are zero-extended to 64 bits, packed and stored into the \emph{registerM}.


\paragraph{Modifier(s)}
\noindent\begin{itemize}
\item oddlanes (cf. oddlanes in section \ref{par:modifier-oddlanes} on page \pageref{par:modifier-oddlanes})
\end{itemize}

\paragraph{Execution} {Reservation table \textsf{ALU\_LITE} (Section~\ref{sec:reservation-ALU-LITE})}
~
{\begin{lstlisting}
stage ID:
  new oddlanes = _ZX_1(oddlanes);
stage RR:
  new argument2 = PGR[registerP];
  if (oddlanes) {
    argument2 = argument2 >> 32;
  }
  for (I = 0; I < 2; I += 1) {
    result1.64[I] = _ZX_32(argument2.64[I])
  };
stage E1:
  PGR[registerM] = result1;
\end{lstlisting}}
\newpage
\subsubsection[{\small FCOMPHQ: Floating-Point Compare Half Word Quadruple}]{FCOMPHQ\hfill Floating-Point Compare Half Word Quadruple}
\label{instr:FCOMPHQ}\index{fcomphq}
 
\paragraph{Encoding} Format \textsf{ALU\_FHQCWRRE} (Section~\ref{sec:format-ALU-FHQCWRRE}), \verb|alucode3 = 010|\\

\bitfields{ 1/P, 2/11, 1/1, 1/0, 3/floatcomp, 5/registerWe, 1/0, 2/11, 3/010, 1/1, 5/registerYe, 1/--, 5/registerZe, 1/1 }


\paragraph{Syntax} \texttt{fcomphq}\emph{floatcomp} \emph{registerWe} = \emph{registerZe}, \emph{registerYe}

\paragraph{Description} The \emph{registerZe} and the \emph{registerYe} are interpreted as four binary 16 floating-point numbers packed into 64 bits. The \emph{registerZe} half words are compared to the \emph{registerYe} half words using condition \emph{floatcomp}. The four boolean result bits are packed and stored into the \emph{registerWe}.


\paragraph{Modifier(s)}
\noindent\begin{itemize}
\item floatcomp (cf. floatcomp in section \ref{par:modifier-floatcomp} on page \pageref{par:modifier-floatcomp})
\end{itemize}

\paragraph{Execution} {Reservation table \textsf{ALU\_LITE} (Section~\ref{sec:reservation-ALU-LITE})}
~
{\begin{lstlisting}
stage ID:
  new argument4 = _ZX_3(floatcomp);
stage RR:
  new argument3 = GPR[registerYe<<1];
  new argument2 = GPR[registerZe<<1];
  for (I = 0; I < 4; I += 1) {
    result1.1[I] = floatcomp_16(argument4, argument2.16[I], argument3.16[I])
  };
stage E1:
  GPR[registerWe<<1] = result1;
\end{lstlisting}}
\newpage

\paragraph{Encoding} Format \textsf{ALU\_FHQCWRRO} (Section~\ref{sec:format-ALU-FHQCWRRO}), \verb|alucode3 = 010|\\

\bitfields{ 1/P, 2/11, 1/1, 1/0, 3/floatcomp, 5/registerWo, 1/1, 2/11, 3/010, 1/1, 5/registerYo, 1/--, 5/registerZo, 1/1 }


\paragraph{Syntax} \texttt{fcomphq}\emph{floatcomp} \emph{registerWo} = \emph{registerZo}, \emph{registerYo}

\paragraph{Description} The \emph{registerZo} and the \emph{registerYo} are interpreted as four binary 16 floating-point numbers packed into 64 bits. The \emph{registerZo} half words are compared to the \emph{registerYo} half words using condition \emph{floatcomp}. The four boolean result bits are packed and stored into the \emph{registerWo}.


\paragraph{Modifier(s)}
\noindent\begin{itemize}
\item floatcomp (cf. floatcomp in section \ref{par:modifier-floatcomp} on page \pageref{par:modifier-floatcomp})
\end{itemize}

\paragraph{Execution} {Reservation table \textsf{ALU\_LITE} (Section~\ref{sec:reservation-ALU-LITE})}
~
{\begin{lstlisting}
stage ID:
  new argument4 = _ZX_3(floatcomp);
stage RR:
  new argument3 = GPR[(registerYo<<1)+1];
  new argument2 = GPR[(registerZo<<1)+1];
  for (I = 0; I < 4; I += 1) {
    result1.1[I] = floatcomp_16(argument4, argument2.16[I], argument3.16[I])
  };
stage E1:
  GPR[(registerWo<<1)+1] = result1;
\end{lstlisting}}
\newpage

\paragraph{Encoding} Format \textsf{ALU\_FHQCWRRE.M} (Section~\ref{sec:format-ALU-FHQCWRRE.M}), \verb|alucode3 = 010|\\

\bitfields{ 1/P, 2/00, 2/-- --, 27/upper27, 1/1, 2/11, 1/1, 1/0, 3/floatcomp, 5/registerWe, 1/0, 2/11, 3/010, 1/1, 1/splat32, 5/lower5, 5/registerZe, 1/1 }


\paragraph{Syntax} \texttt{fcomphq}\emph{floatcomp} \emph{registerWe} = \emph{registerZe}, \emph{upper27\_\discretionary{}{}{}lower5}\emph{splat32}

\paragraph{Description} The \emph{registerZe} and the \emph{upper27\_\discretionary{}{}{}lower5} are interpreted as four binary 16 floating-point numbers packed into 64 bits. The \emph{registerZe} half words are compared to the \emph{upper27\_\discretionary{}{}{}lower5} half words using condition \emph{floatcomp}. The four boolean result bits are packed and stored into the \emph{registerWe}.


\paragraph{Modifier(s)}
\noindent\begin{itemize}
\item floatcomp (cf. floatcomp in section \ref{par:modifier-floatcomp} on page \pageref{par:modifier-floatcomp})
\item splat32 (cf. splat32 in section \ref{par:modifier-splat32} on page \pageref{par:modifier-splat32})
\end{itemize}

\paragraph{Execution} {Reservation table \textsf{ALU\_LITE.X} (Section~\ref{sec:reservation-ALU-LITE.X})}
~
{\begin{lstlisting}
stage ID:
  new splat32 = _ZX_1(splat32);
  new argument4 = _ZX_3(floatcomp);
stage RR:
  new argument3 = splat32 ? (_WX_32(upper27_lower5) << 32) | _ZX_32(_WX_32(upper27_lower5)) : _SX_32(_WX_32(upper27_lower5));
  new argument2 = GPR[registerZe<<1];
  for (I = 0; I < 4; I += 1) {
    result1.1[I] = floatcomp_16(argument4, argument2.16[I], argument3.16[I])
  };
stage E1:
  GPR[registerWe<<1] = result1;
\end{lstlisting}}
\newpage

\paragraph{Encoding} Format \textsf{ALU\_FHQCWRRO.M} (Section~\ref{sec:format-ALU-FHQCWRRO.M}), \verb|alucode3 = 010|\\

\bitfields{ 1/P, 2/00, 2/-- --, 27/upper27, 1/1, 2/11, 1/1, 1/0, 3/floatcomp, 5/registerWo, 1/1, 2/11, 3/010, 1/1, 1/splat32, 5/lower5, 5/registerZo, 1/1 }


\paragraph{Syntax} \texttt{fcomphq}\emph{floatcomp} \emph{registerWo} = \emph{registerZo}, \emph{upper27\_\discretionary{}{}{}lower5}\emph{splat32}

\paragraph{Description} The \emph{registerZo} and the \emph{upper27\_\discretionary{}{}{}lower5} are interpreted as four binary 16 floating-point numbers packed into 64 bits. The \emph{registerZo} half words are compared to the \emph{upper27\_\discretionary{}{}{}lower5} half words using condition \emph{floatcomp}. The four boolean result bits are packed and stored into the \emph{registerWo}.


\paragraph{Modifier(s)}
\noindent\begin{itemize}
\item floatcomp (cf. floatcomp in section \ref{par:modifier-floatcomp} on page \pageref{par:modifier-floatcomp})
\item splat32 (cf. splat32 in section \ref{par:modifier-splat32} on page \pageref{par:modifier-splat32})
\end{itemize}

\paragraph{Execution} {Reservation table \textsf{ALU\_LITE.X} (Section~\ref{sec:reservation-ALU-LITE.X})}
~
{\begin{lstlisting}
stage ID:
  new splat32 = _ZX_1(splat32);
  new argument4 = _ZX_3(floatcomp);
stage RR:
  new argument3 = splat32 ? (_WX_32(upper27_lower5) << 32) | _ZX_32(_WX_32(upper27_lower5)) : _SX_32(_WX_32(upper27_lower5));
  new argument2 = GPR[(registerZo<<1)+1];
  for (I = 0; I < 4; I += 1) {
    result1.1[I] = floatcomp_16(argument4, argument2.16[I], argument3.16[I])
  };
stage E1:
  GPR[(registerWo<<1)+1] = result1;
\end{lstlisting}}
\newpage
\subsubsection[{\small FCOMPWP: Floating-Point Compare Word Pair}]{FCOMPWP\hfill Floating-Point Compare Word Pair}
\label{instr:FCOMPWP}\index{fcompwp}
 
\paragraph{Encoding} Format \textsf{ALU\_FWPCWRRE} (Section~\ref{sec:format-ALU-FWPCWRRE}), \verb|alucode3 = 010|\\

\bitfields{ 1/P, 2/11, 1/1, 1/0, 3/floatcomp, 5/registerWe, 1/0, 2/11, 3/010, 1/1, 5/registerYe, 1/--, 5/registerZe, 1/0 }


\paragraph{Syntax} \texttt{fcompwp}\emph{floatcomp} \emph{registerWe} = \emph{registerZe}, \emph{registerYe}

\paragraph{Description} The \emph{registerZe} and the \emph{registerYe} are interpreted as two binary 32 floating-point numbers packed into 64 bits. The \emph{registerZe} words are compared to the \emph{registerYe} words using condition \emph{floatcomp}. The two boolean result bits are packed and stored into the \emph{registerWe}.


\paragraph{Modifier(s)}
\noindent\begin{itemize}
\item floatcomp (cf. floatcomp in section \ref{par:modifier-floatcomp} on page \pageref{par:modifier-floatcomp})
\end{itemize}

\paragraph{Execution} {Reservation table \textsf{ALU\_LITE} (Section~\ref{sec:reservation-ALU-LITE})}
~
{\begin{lstlisting}
stage ID:
  new argument4 = _ZX_3(floatcomp);
stage RR:
  new argument3 = GPR[registerYe<<1];
  new argument2 = GPR[registerZe<<1];
  for (I = 0; I < 4; I += 1) {
    result1.1[I] = floatcomp_32(argument4, argument2.32[I], argument3.32[I])
  };
stage E1:
  GPR[registerWe<<1] = result1;
\end{lstlisting}}
\newpage

\paragraph{Encoding} Format \textsf{ALU\_FWPCWRRO} (Section~\ref{sec:format-ALU-FWPCWRRO}), \verb|alucode3 = 010|\\

\bitfields{ 1/P, 2/11, 1/1, 1/0, 3/floatcomp, 5/registerWo, 1/1, 2/11, 3/010, 1/1, 5/registerYo, 1/--, 5/registerZo, 1/0 }


\paragraph{Syntax} \texttt{fcompwp}\emph{floatcomp} \emph{registerWo} = \emph{registerZo}, \emph{registerYo}

\paragraph{Description} The \emph{registerZo} and the \emph{registerYo} are interpreted as two binary 32 floating-point numbers packed into 64 bits. The \emph{registerZo} words are compared to the \emph{registerYo} words using condition \emph{floatcomp}. The two boolean result bits are packed and stored into the \emph{registerWo}.


\paragraph{Modifier(s)}
\noindent\begin{itemize}
\item floatcomp (cf. floatcomp in section \ref{par:modifier-floatcomp} on page \pageref{par:modifier-floatcomp})
\end{itemize}

\paragraph{Execution} {Reservation table \textsf{ALU\_LITE} (Section~\ref{sec:reservation-ALU-LITE})}
~
{\begin{lstlisting}
stage ID:
  new argument4 = _ZX_3(floatcomp);
stage RR:
  new argument3 = GPR[(registerYo<<1)+1];
  new argument2 = GPR[(registerZo<<1)+1];
  for (I = 0; I < 4; I += 1) {
    result1.1[I] = floatcomp_32(argument4, argument2.32[I], argument3.32[I])
  };
stage E1:
  GPR[(registerWo<<1)+1] = result1;
\end{lstlisting}}
\newpage

\paragraph{Encoding} Format \textsf{ALU\_FWPCWRRE.M} (Section~\ref{sec:format-ALU-FWPCWRRE.M}), \verb|alucode3 = 010|\\

\bitfields{ 1/P, 2/00, 2/-- --, 27/upper27, 1/1, 2/11, 1/1, 1/0, 3/floatcomp, 5/registerWe, 1/0, 2/11, 3/010, 1/1, 1/splat32, 5/lower5, 5/registerZe, 1/0 }


\paragraph{Syntax} \texttt{fcompwp}\emph{floatcomp} \emph{registerWe} = \emph{registerZe}, \emph{upper27\_\discretionary{}{}{}lower5}\emph{splat32}

\paragraph{Description} The \emph{registerZe} and the \emph{upper27\_\discretionary{}{}{}lower5} are interpreted as two binary 32 floating-point numbers packed into 64 bits. The \emph{registerZe} words are compared to the \emph{upper27\_\discretionary{}{}{}lower5} words using condition \emph{floatcomp}. The two boolean result bits are packed and stored into the \emph{registerWe}.


\paragraph{Modifier(s)}
\noindent\begin{itemize}
\item floatcomp (cf. floatcomp in section \ref{par:modifier-floatcomp} on page \pageref{par:modifier-floatcomp})
\item splat32 (cf. splat32 in section \ref{par:modifier-splat32} on page \pageref{par:modifier-splat32})
\end{itemize}

\paragraph{Execution} {Reservation table \textsf{ALU\_LITE.X} (Section~\ref{sec:reservation-ALU-LITE.X})}
~
{\begin{lstlisting}
stage ID:
  new splat32 = _ZX_1(splat32);
  new argument4 = _ZX_3(floatcomp);
stage RR:
  new argument3 = splat32 ? (_WX_32(upper27_lower5) << 32) | _ZX_32(_WX_32(upper27_lower5)) : _SX_32(_WX_32(upper27_lower5));
  new argument2 = GPR[registerZe<<1];
  for (I = 0; I < 4; I += 1) {
    result1.1[I] = floatcomp_32(argument4, argument2.32[I], argument3.32[I])
  };
stage E1:
  GPR[registerWe<<1] = result1;
\end{lstlisting}}
\newpage

\paragraph{Encoding} Format \textsf{ALU\_FWPCWRRO.M} (Section~\ref{sec:format-ALU-FWPCWRRO.M}), \verb|alucode3 = 010|\\

\bitfields{ 1/P, 2/00, 2/-- --, 27/upper27, 1/1, 2/11, 1/1, 1/0, 3/floatcomp, 5/registerWo, 1/1, 2/11, 3/010, 1/1, 1/splat32, 5/lower5, 5/registerZo, 1/0 }


\paragraph{Syntax} \texttt{fcompwp}\emph{floatcomp} \emph{registerWo} = \emph{registerZo}, \emph{upper27\_\discretionary{}{}{}lower5}\emph{splat32}

\paragraph{Description} The \emph{registerZo} and the \emph{upper27\_\discretionary{}{}{}lower5} are interpreted as two binary 32 floating-point numbers packed into 64 bits. The \emph{registerZo} words are compared to the \emph{upper27\_\discretionary{}{}{}lower5} words using condition \emph{floatcomp}. The two boolean result bits are packed and stored into the \emph{registerWo}.


\paragraph{Modifier(s)}
\noindent\begin{itemize}
\item floatcomp (cf. floatcomp in section \ref{par:modifier-floatcomp} on page \pageref{par:modifier-floatcomp})
\item splat32 (cf. splat32 in section \ref{par:modifier-splat32} on page \pageref{par:modifier-splat32})
\end{itemize}

\paragraph{Execution} {Reservation table \textsf{ALU\_LITE.X} (Section~\ref{sec:reservation-ALU-LITE.X})}
~
{\begin{lstlisting}
stage ID:
  new splat32 = _ZX_1(splat32);
  new argument4 = _ZX_3(floatcomp);
stage RR:
  new argument3 = splat32 ? (_WX_32(upper27_lower5) << 32) | _ZX_32(_WX_32(upper27_lower5)) : _SX_32(_WX_32(upper27_lower5));
  new argument2 = GPR[(registerZo<<1)+1];
  for (I = 0; I < 4; I += 1) {
    result1.1[I] = floatcomp_32(argument4, argument2.32[I], argument3.32[I])
  };
stage E1:
  GPR[(registerWo<<1)+1] = result1;
\end{lstlisting}}
\newpage
\subsubsection[{\small FNARROWWHQ: Floating Point Narrow Word to Half Word Quadruple}]{FNARROWWHQ\hfill Floating Point Narrow Word to Half Word Quadruple}
\label{instr:FNARROWWHQ}\index{fnarrowwhq}
 
\paragraph{Encoding} Format \textsf{ALU\_FHQWR} (Section~\ref{sec:format-ALU-FHQWR}), \verb|fpucode6 = 1100|\\

\bitfields{ 1/P, 2/11, 1/1, 1/1, 3/floatmode, 6/registerW, 2/11, 3/001, 1/1, 4/1100, 1/ziplanes, 1/0, 5/registerP, 1/1 }


\paragraph{Syntax} \texttt{fnarrowwhq}\emph{ziplanes}\emph{floatmode} \emph{registerW} = \emph{registerP}

\paragraph{Description} The \emph{registerP} is interpreted as four binary 32 floating-point numbers packed into 128 bits. The \emph{registerP} words are converted to binary 16 floating-point numbers according to the rounding mode. If \emph{ziplanes}, the resulting half words are interleaved by perfect shuffling. The four resulting half words are packed and stored back into the \emph{registerW}. This instruction may raise exception bits in the CS register.


\paragraph{Modifier(s)}
\noindent\begin{itemize}
\item floatmode (cf. floatmode in section \ref{par:modifier-floatmode} on page \pageref{par:modifier-floatmode})
\item ziplanes (cf. ziplanes in section \ref{par:modifier-ziplanes} on page \pageref{par:modifier-ziplanes})
\end{itemize}

\paragraph{Execution} {Reservation table \textsf{ALU\_LITE} (Section~\ref{sec:reservation-ALU-LITE})}
~
{\begin{lstlisting}
stage ID:
  new ziplanes = _ZX_1(ziplanes);
  new floatmode = _ZX_3(floatmode);
stage RR:
  new argument2 = PGR[registerP];
  new RM = decode_riscv_float_rounding_mode(floatmode);
  if (ziplanes) {
    result1.16[0] = f32_to_f16(RM, argument2.32[0]);
    result1.16[1] = f32_to_f16(RM, argument2.32[2]);
    result1.16[2] = f32_to_f16(RM, argument2.32[1]);
  } else {
    result1.16[0] = f32_to_f16(RM, argument2.32[0]);
    result1.16[1] = f32_to_f16(RM, argument2.32[1]);
    result1.16[2] = f32_to_f16(RM, argument2.32[2]);
  }
  result1.16[3] = f32_to_f16(RM, argument2.32[3]);
stage E1:
  GPR[registerW] = result1;
  commit_float_exception_flags();
\end{lstlisting}}
\newpage
\subsubsection[{\small FRACTDWP: Fraction Double Word to Word Pair}]{FRACTDWP\hfill Fraction Double Word to Word Pair}
\label{instr:FRACTDWP}\index{fractdwp}
 
\paragraph{Encoding} Format \textsf{ALU\_WPNWRR} (Section~\ref{sec:format-ALU-WPNWRR}), \verb|alucode8 = 1001|\\

\bitfields{ 1/P, 2/11, 1/1, 4/1111, 6/registerW, 2/10, 3/101, 1/0, 4/1001, 1/--, 1/0, 5/registerP, 1/1 }


\paragraph{Syntax} \texttt{fractdwp} \emph{registerW} = \emph{registerP}

\paragraph{Description} The \emph{registerP} is interpreted as two double words packed into 128 bits. The \emph{registerP} double words are converted to 32-bit words by removing the low bits. The two resulting words are packed and stored back into the \emph{registerW}.


\paragraph{Execution} {Reservation table \textsf{ALU\_LITE} (Section~\ref{sec:reservation-ALU-LITE})}
~
{\begin{lstlisting}
stage RR:
  new argument2 = PGR[registerP];
  result1.32[0] = argument2.32[1];
  result1.32[1] = argument2.32[3];
stage E1:
  GPR[registerW] = result1;
\end{lstlisting}}
\newpage
\subsubsection[{\small FRACTHBO: Fraction Half Word to Byte Octuple}]{FRACTHBO\hfill Fraction Half Word to Byte Octuple}
\label{instr:FRACTHBO}\index{fracthbo}
 
\paragraph{Encoding} Format \textsf{ALU\_BONWRR} (Section~\ref{sec:format-ALU-BONWRR}), \verb|alucode8 = 1001|\\

\bitfields{ 1/P, 2/11, 1/1, 4/1111, 6/registerW, 2/10, 3/101, 1/1, 4/1001, 1/ziplanes, 1/0, 5/registerP, 1/1 }


\paragraph{Syntax} \texttt{fracthbo}\emph{ziplanes} \emph{registerW} = \emph{registerP}

\paragraph{Description} The \emph{registerP} is interpreted as eight half words packed into 128 bits. The \emph{registerP} half words are converted to bytes by removing the low bits. If \emph{ziplanes}, the resulting bytes are interleaved perfect shuffling. The eight resulting bytes are packed and stored back into the \emph{registerW}.


\paragraph{Modifier(s)}
\noindent\begin{itemize}
\item ziplanes (cf. ziplanes in section \ref{par:modifier-ziplanes} on page \pageref{par:modifier-ziplanes})
\end{itemize}

\paragraph{Execution} {Reservation table \textsf{ALU\_LITE} (Section~\ref{sec:reservation-ALU-LITE})}
~
{\begin{lstlisting}
stage ID:
  new ziplanes = _ZX_1(ziplanes);
stage RR:
  new argument2 = PGR[registerP];
  if (ziplanes) {
    result1.8[0] = argument2.8[1];
    result1.8[1] = argument2.8[9];
    result1.8[2] = argument2.8[3];
    result1.8[3] = argument2.8[11];
    result1.8[4] = argument2.8[5];
    result1.8[5] = argument2.8[13];
    result1.8[6] = argument2.8[7];
  } else {
    result1.8[0] = argument2.8[1];
    result1.8[1] = argument2.8[3];
    result1.8[2] = argument2.8[5];
    result1.8[3] = argument2.8[7];
    result1.8[4] = argument2.8[9];
    result1.8[5] = argument2.8[11];
    result1.8[6] = argument2.8[13];
  }
  result1.8[7] = argument2.8[15];
stage E1:
  GPR[registerW] = result1;
\end{lstlisting}}
\newpage
\subsubsection[{\small FRACTWHQ: Fraction Word to Half Word Quadruple}]{FRACTWHQ\hfill Fraction Word to Half Word Quadruple}
\label{instr:FRACTWHQ}\index{fractwhq}
 
\paragraph{Encoding} Format \textsf{ALU\_HQNWRR} (Section~\ref{sec:format-ALU-HQNWRR}), \verb|alucode8 = 1001|\\

\bitfields{ 1/P, 2/11, 1/1, 4/1111, 6/registerW, 2/10, 3/101, 1/1, 4/1001, 1/ziplanes, 1/0, 5/registerP, 1/0 }


\paragraph{Syntax} \texttt{fractwhq}\emph{ziplanes} \emph{registerW} = \emph{registerP}

\paragraph{Description} The \emph{ziplanes} and the \emph{registerP} are interpreted as four words packed into 128 bits. The \emph{ziplanes} and the \emph{registerP} words are converted to 16-bit half words by removing the low bits. If \emph{ziplanes}, the resulting half words are interleaved with the \emph{registerP} half words in even lanes and the \emph{ziplanes} half words in odd lanes. The four resulting half words are packed and stored back into the \emph{registerW}.


\paragraph{Modifier(s)}
\noindent\begin{itemize}
\item ziplanes (cf. ziplanes in section \ref{par:modifier-ziplanes} on page \pageref{par:modifier-ziplanes})
\end{itemize}

\paragraph{Execution} {Reservation table \textsf{ALU\_LITE} (Section~\ref{sec:reservation-ALU-LITE})}
~
{\begin{lstlisting}
stage ID:
  new ziplanes = _ZX_1(ziplanes);
stage RR:
  new argument2 = PGR[registerP];
  if (ziplanes) {
    result1.16[0] = argument2.16[1];
    result1.16[1] = argument2.16[5];
    result1.16[2] = argument2.16[3];
  } else {
    result1.16[0] = argument2.16[1];
    result1.16[1] = argument2.16[3];
    result1.16[2] = argument2.16[5];
  }
  result1.16[3] = argument2.16[7];
stage E1:
  GPR[registerW] = result1;
\end{lstlisting}}
\newpage
\subsubsection[{\small INSF: Insert Bit Field into Word}]{INSF\hfill Insert Bit Field into Word}
\label{instr:INSF}\index{insf}
 
\paragraph{Encoding} Format \textsf{ALU\_BFWR} (Section~\ref{sec:format-ALU-BFWR}), \verb|wricode1 = 00|\\

\bitfields{ 1/P, 2/11, 1/0, 2/00, 2/stopbit2, 6/registerW, 2/01, 4/stopbit4, 6/startbit, 6/registerZ }


\paragraph{Syntax} \texttt{insf} \emph{registerW} = \emph{registerZ}, \emph{stopbit2\_\discretionary{}{}{}stopbit4}, \emph{startbit}

\paragraph{Description} The start bit position and the stop bit position are extracted from the \emph{startbit} and the \emph{stopbit2\_\discretionary{}{}{}stopbit4} respectively. A bit-mask is computed with ones set between the start bit position and the stop bit position inclusive. The \emph{registerZ} is shifted left by the start bit position and inserted into the \emph{registerW} according to the bit-mask. Behavior is undefined if the start bit position is greater than the stop bit position.


\paragraph{Execution} {Reservation table \textsf{ALU\_TINY} (Section~\ref{sec:reservation-ALU-TINY})}
~
{\begin{lstlisting}
stage ID:
  new startbit = _ZX_6(startbit);
  new stopbit = _ZX_6(stopbit2_stopbit4);
stage RR:
  new bias = startbit <= stopbit;
  new mask = _ZX_64(((2 << stopbit) - ((2 - bias) << startbit)) + (bias - 1));
  new argument2 = GPR[registerZ];
  new argument1 = GPR[registerW];
  new result1 = ((argument2 << startbit) & mask) | (argument1 & ~mask);
stage E1:
  GPR[registerW] = result1;
\end{lstlisting}}
\newpage
\subsubsection[{\small IORD: Bitwise Inclusive Or Double Word}]{IORD\hfill Bitwise Inclusive Or Double Word}
\label{instr:IORD}\index{iord}
 
\paragraph{Encoding} Format \textsf{ALU\_DWRR0} (Section~\ref{sec:format-ALU-DWRR0}), \verb|exucode2 = 1010|\\

\bitfields{ 1/P, 2/11, 1/0, 4/1010, 6/registerW, 2/10, 3/000, 1/0, 6/registerY, 6/registerZ }


\paragraph{Syntax} \texttt{iord} \emph{registerW} = \emph{registerZ}, \emph{registerY}

\paragraph{Description} Bitwise inclusive or of the \emph{registerZ} with the \emph{registerY}. The result is stored into the \emph{registerW}.


\paragraph{Execution} {Reservation table \textsf{ALU\_TINY} (Section~\ref{sec:reservation-ALU-TINY})}
~
{\begin{lstlisting}
stage RR:
  new argument3 = GPR[registerY];
  new argument2 = GPR[registerZ];
  new result1 = argument2 | argument3;
stage E1:
  GPR[registerW] = result1;
\end{lstlisting}}
\newpage

\paragraph{Encoding} Format \textsf{ALU\_DWRR0.M} (Section~\ref{sec:format-ALU-DWRR0.M}), \verb|exucode2 = 1010|\\

\bitfields{ 1/P, 2/00, 2/-- --, 27/upper27, 1/1, 2/11, 1/0, 4/1010, 6/registerW, 2/10, 3/000, 1/0, 1/splat32, 5/lower5, 6/registerZ }


\paragraph{Syntax} \texttt{iord} \emph{registerW} = \emph{registerZ}, \emph{upper27\_\discretionary{}{}{}lower5}\emph{splat32}

\paragraph{Description} Bitwise inclusive or of the \emph{registerZ} with the \emph{upper27\_\discretionary{}{}{}lower5}. The result is stored into the \emph{registerW}.


\paragraph{Modifier(s)}
\noindent\begin{itemize}
\item splat32 (cf. splat32 in section \ref{par:modifier-splat32} on page \pageref{par:modifier-splat32})
\end{itemize}

\paragraph{Execution} {Reservation table \textsf{ALU\_TINY.X} (Section~\ref{sec:reservation-ALU-TINY.X})}
~
{\begin{lstlisting}
stage ID:
  new splat32 = _ZX_1(splat32);
stage RR:
  new argument3 = splat32 ? (_WX_32(upper27_lower5) << 32) | _ZX_32(_WX_32(upper27_lower5)) : _SX_32(_WX_32(upper27_lower5));
  new argument2 = GPR[registerZ];
  new result1 = argument2 | argument3;
stage E1:
  GPR[registerW] = result1;
\end{lstlisting}}
\newpage

\paragraph{Encoding} Format \textsf{ALU\_DWRI} (Section~\ref{sec:format-ALU-DWRI}), \verb|exucode2 = 1010|\\

\bitfields{ 1/P, 2/11, 1/0, 4/1010, 6/registerW, 2/00, 10/signed10, 6/registerZ }


\paragraph{Syntax} \texttt{iord} \emph{registerW} = \emph{registerZ}, \emph{signed10}

\paragraph{Description} Bitwise inclusive or of the \emph{registerZ} with the \emph{signed10}. The result is stored into the \emph{registerW}.


\paragraph{Execution} {Reservation table \textsf{ALU\_TINY} (Section~\ref{sec:reservation-ALU-TINY})}
~
{\begin{lstlisting}
stage RR:
  new argument3 = _SX_10(signed10);
  new argument2 = GPR[registerZ];
  new result1 = argument2 | argument3;
stage E1:
  GPR[registerW] = result1;
\end{lstlisting}}
\newpage

\paragraph{Encoding} Format \textsf{ALU\_DWRI.X} (Section~\ref{sec:format-ALU-DWRI.X}), \verb|exucode2 = 1010|\\

\bitfields{ 1/P, 2/00, 2/-- --, 27/upper27, 1/1, 2/11, 1/0, 4/1010, 6/registerW, 2/00, 10/lower10, 6/registerZ }


\paragraph{Syntax} \texttt{iord} \emph{registerW} = \emph{registerZ}, \emph{upper27\_\discretionary{}{}{}lower10}

\paragraph{Description} Bitwise inclusive or of the \emph{registerZ} with the \emph{upper27\_\discretionary{}{}{}lower10}. The result is stored into the \emph{registerW}.


\paragraph{Execution} {Reservation table \textsf{ALU\_TINY.X} (Section~\ref{sec:reservation-ALU-TINY.X})}
~
{\begin{lstlisting}
stage RR:
  new argument3 = _SX_37(upper27_lower10);
  new argument2 = GPR[registerZ];
  new result1 = argument2 | argument3;
stage E1:
  GPR[registerW] = result1;
\end{lstlisting}}
\newpage

\paragraph{Encoding} Format \textsf{ALU\_DWRI.Y} (Section~\ref{sec:format-ALU-DWRI.Y}), \verb|exucode2 = 1010|\\

\bitfields{ 1/P, 2/00, 2/-- --, 27/extend27, 1/1, 2/00, 2/-- --, 27/upper27, 1/1, 2/11, 1/0, 4/1010, 6/registerW, 2/00, 10/lower10, 6/registerZ }


\paragraph{Syntax} \texttt{iord} \emph{registerW} = \emph{registerZ}, \emph{extend27\_\discretionary{}{}{}upper27\_\discretionary{}{}{}lower10}

\paragraph{Description} Bitwise inclusive or of the \emph{registerZ} with the \emph{extend27\_\discretionary{}{}{}upper27\_\discretionary{}{}{}lower10}. The result is stored into the \emph{registerW}.


\paragraph{Execution} {Reservation table \textsf{ALU\_TINY.Y} (Section~\ref{sec:reservation-ALU-TINY.Y})}
~
{\begin{lstlisting}
stage RR:
  new argument3 = _WX_64(extend27_upper27_lower10);
  new argument2 = GPR[registerZ];
  new result1 = argument2 | argument3;
stage E1:
  GPR[registerW] = result1;
\end{lstlisting}}
\newpage
\subsubsection[{\small IORND: Bitwise Inclusive Or with Not Operand Double Word}]{IORND\hfill Bitwise Inclusive Or with Not Operand Double Word}
\label{instr:IORND}\index{iornd}
 
\paragraph{Encoding} Format \textsf{ALU\_DWRR0} (Section~\ref{sec:format-ALU-DWRR0}), \verb|exucode2 = 1111|\\

\bitfields{ 1/P, 2/11, 1/0, 4/1111, 6/registerW, 2/10, 3/000, 1/0, 6/registerY, 6/registerZ }


\paragraph{Syntax} \texttt{iornd} \emph{registerW} = \emph{registerZ}, \emph{registerY}

\paragraph{Description} Bitwise inclusive or of not the \emph{registerZ} with the \emph{registerY}. The result is stored into the \emph{registerW}.


\paragraph{Execution} {Reservation table \textsf{ALU\_TINY} (Section~\ref{sec:reservation-ALU-TINY})}
~
{\begin{lstlisting}
stage RR:
  new argument3 = GPR[registerY];
  new argument2 = GPR[registerZ];
  new result1 = ~argument2 | argument3;
stage E1:
  GPR[registerW] = result1;
\end{lstlisting}}
\newpage

\paragraph{Encoding} Format \textsf{ALU\_DWRR0.M} (Section~\ref{sec:format-ALU-DWRR0.M}), \verb|exucode2 = 1111|\\

\bitfields{ 1/P, 2/00, 2/-- --, 27/upper27, 1/1, 2/11, 1/0, 4/1111, 6/registerW, 2/10, 3/000, 1/0, 1/splat32, 5/lower5, 6/registerZ }


\paragraph{Syntax} \texttt{iornd} \emph{registerW} = \emph{registerZ}, \emph{upper27\_\discretionary{}{}{}lower5}\emph{splat32}

\paragraph{Description} Bitwise inclusive or of not the \emph{registerZ} with the \emph{upper27\_\discretionary{}{}{}lower5}. The result is stored into the \emph{registerW}.


\paragraph{Modifier(s)}
\noindent\begin{itemize}
\item splat32 (cf. splat32 in section \ref{par:modifier-splat32} on page \pageref{par:modifier-splat32})
\end{itemize}

\paragraph{Execution} {Reservation table \textsf{ALU\_TINY.X} (Section~\ref{sec:reservation-ALU-TINY.X})}
~
{\begin{lstlisting}
stage ID:
  new splat32 = _ZX_1(splat32);
stage RR:
  new argument3 = splat32 ? (_WX_32(upper27_lower5) << 32) | _ZX_32(_WX_32(upper27_lower5)) : _SX_32(_WX_32(upper27_lower5));
  new argument2 = GPR[registerZ];
  new result1 = ~argument2 | argument3;
stage E1:
  GPR[registerW] = result1;
\end{lstlisting}}
\newpage

\paragraph{Encoding} Format \textsf{ALU\_DWRI} (Section~\ref{sec:format-ALU-DWRI}), \verb|exucode2 = 1111|\\

\bitfields{ 1/P, 2/11, 1/0, 4/1111, 6/registerW, 2/00, 10/signed10, 6/registerZ }


\paragraph{Syntax} \texttt{iornd} \emph{registerW} = \emph{registerZ}, \emph{signed10}

\paragraph{Description} Bitwise inclusive or of not the \emph{registerZ} with the \emph{signed10}. The result is stored into the \emph{registerW}.


\paragraph{Execution} {Reservation table \textsf{ALU\_TINY} (Section~\ref{sec:reservation-ALU-TINY})}
~
{\begin{lstlisting}
stage RR:
  new argument3 = _SX_10(signed10);
  new argument2 = GPR[registerZ];
  new result1 = ~argument2 | argument3;
stage E1:
  GPR[registerW] = result1;
\end{lstlisting}}
\newpage

\paragraph{Encoding} Format \textsf{ALU\_DWRI.X} (Section~\ref{sec:format-ALU-DWRI.X}), \verb|exucode2 = 1111|\\

\bitfields{ 1/P, 2/00, 2/-- --, 27/upper27, 1/1, 2/11, 1/0, 4/1111, 6/registerW, 2/00, 10/lower10, 6/registerZ }


\paragraph{Syntax} \texttt{iornd} \emph{registerW} = \emph{registerZ}, \emph{upper27\_\discretionary{}{}{}lower10}

\paragraph{Description} Bitwise inclusive or of not the \emph{registerZ} with the \emph{upper27\_\discretionary{}{}{}lower10}. The result is stored into the \emph{registerW}.


\paragraph{Execution} {Reservation table \textsf{ALU\_TINY.X} (Section~\ref{sec:reservation-ALU-TINY.X})}
~
{\begin{lstlisting}
stage RR:
  new argument3 = _SX_37(upper27_lower10);
  new argument2 = GPR[registerZ];
  new result1 = ~argument2 | argument3;
stage E1:
  GPR[registerW] = result1;
\end{lstlisting}}
\newpage

\paragraph{Encoding} Format \textsf{ALU\_DWRI.Y} (Section~\ref{sec:format-ALU-DWRI.Y}), \verb|exucode2 = 1111|\\

\bitfields{ 1/P, 2/00, 2/-- --, 27/extend27, 1/1, 2/00, 2/-- --, 27/upper27, 1/1, 2/11, 1/0, 4/1111, 6/registerW, 2/00, 10/lower10, 6/registerZ }


\paragraph{Syntax} \texttt{iornd} \emph{registerW} = \emph{registerZ}, \emph{extend27\_\discretionary{}{}{}upper27\_\discretionary{}{}{}lower10}

\paragraph{Description} Bitwise inclusive or of not the \emph{registerZ} with the \emph{extend27\_\discretionary{}{}{}upper27\_\discretionary{}{}{}lower10}. The result is stored into the \emph{registerW}.


\paragraph{Execution} {Reservation table \textsf{ALU\_TINY.Y} (Section~\ref{sec:reservation-ALU-TINY.Y})}
~
{\begin{lstlisting}
stage RR:
  new argument3 = _WX_64(extend27_upper27_lower10);
  new argument2 = GPR[registerZ];
  new result1 = ~argument2 | argument3;
stage E1:
  GPR[registerW] = result1;
\end{lstlisting}}
\newpage
\subsubsection[{\small IORNQ: Bitwise Inclusive Or with Not Operand Quadruple Word}]{IORNQ\hfill Bitwise Inclusive Or with Not Operand Quadruple Word}
\label{instr:IORNQ}\index{iornq}
 
\paragraph{Encoding} Format \textsf{ALU\_QWRR0} (Section~\ref{sec:format-ALU-QWRR0}), \verb|exucode2 = 1111|\\

\bitfields{ 1/P, 2/11, 1/0, 4/1111, 5/registerM, 1/--, 2/10, 3/000, 1/1, 5/registerO, 1/--, 5/registerP, 1/-- }


\paragraph{Syntax} \texttt{iornq} \emph{registerM} = \emph{registerP}, \emph{registerO}

\paragraph{Description} Bitwise inclusive or of not the \emph{registerP} with the \emph{registerO}. The result is stored into the \emph{registerM}.


\paragraph{Execution} {Reservation table \textsf{ALU\_LITE} (Section~\ref{sec:reservation-ALU-LITE})}
~
{\begin{lstlisting}
stage RR:
  new argument3 = PGR[registerO];
  new argument2 = PGR[registerP];
  new result1 = ~argument2 | argument3;
stage E1:
  PGR[registerM] = result1;
\end{lstlisting}}
\newpage

\paragraph{Encoding} Format \textsf{ALU\_QWRR0.M} (Section~\ref{sec:format-ALU-QWRR0.M}), \verb|exucode2 = 1111|\\

\bitfields{ 1/P, 2/00, 2/-- --, 27/upper27, 1/1, 2/11, 1/0, 4/1111, 5/registerM, 1/--, 2/10, 3/000, 1/1, 1/splat32, 5/lower5, 5/registerP, 1/1 }


\paragraph{Syntax} \texttt{iornq} \emph{registerM} = \emph{registerP}, \emph{upper27\_\discretionary{}{}{}lower5}\emph{splat32}

\paragraph{Description} Bitwise inclusive or of not the \emph{registerP} with the \emph{upper27\_\discretionary{}{}{}lower5}. The result is stored into the \emph{registerM}.


\paragraph{Modifier(s)}
\noindent\begin{itemize}
\item splat32 (cf. splat32 in section \ref{par:modifier-splat32} on page \pageref{par:modifier-splat32})
\end{itemize}

\paragraph{Execution} {Reservation table \textsf{ALU\_LITE.X} (Section~\ref{sec:reservation-ALU-LITE.X})}
~
{\begin{lstlisting}
stage ID:
  new splat32 = _ZX_1(splat32);
stage RR:
  new argument3 = splat32 ? (_WX_32(upper27_lower5) << 32) | _ZX_32(_WX_32(upper27_lower5)) : _SX_32(_WX_32(upper27_lower5));
  new argument2 = PGR[registerP];
  new result1 = ~argument2 | argument3;
stage E1:
  PGR[registerM] = result1;
\end{lstlisting}}
\newpage
\subsubsection[{\small IORNW: Bitwise Inclusive Or with Not Operand Double Word}]{IORNW\hfill Bitwise Inclusive Or with Not Operand Double Word}
\label{instr:IORNW}\index{iornw}
 
\paragraph{Encoding} Format \textsf{ALU\_WWRR0} (Section~\ref{sec:format-ALU-WWRR0}), \verb|exucode2 = 1111|\\

\bitfields{ 1/P, 2/11, 1/0, 4/1111, 6/registerW, 2/11, 3/000, 1/signextw, 6/registerY, 6/registerZ }


\paragraph{Syntax} \texttt{iornw}\emph{signextw} \emph{registerW} = \emph{registerZ}, \emph{registerY}

\paragraph{Description} Bitwise inclusive or of not the \emph{registerZ} with the \emph{registerY}. The result with the 32 upper bits cleared is stored into the \emph{registerW}.


\paragraph{Modifier(s)}
\noindent\begin{itemize}
\item signextw (cf. signextw in section \ref{par:modifier-signextw} on page \pageref{par:modifier-signextw})
\end{itemize}

\paragraph{Execution} {Reservation table \textsf{ALU\_TINY} (Section~\ref{sec:reservation-ALU-TINY})}
~
{\begin{lstlisting}
stage ID:
  new signextw = _ZX_1(signextw);
stage RR:
  new argument3 = GPR[registerY];
  new argument2 = GPR[registerZ];
  new result1 = ~argument2 | argument3;
stage E1:
  GPR[registerW] = signextw ? _SX_32(result1) : _ZX_32(result1);
\end{lstlisting}}
\newpage

\paragraph{Encoding} Format \textsf{ALU\_WWRR0.W} (Section~\ref{sec:format-ALU-WWRR0.W}), \verb|exucode2 = 1111|\\

\bitfields{ 1/P, 2/00, 2/-- --, 27/upper27, 1/1, 2/11, 1/0, 4/1111, 6/registerW, 2/11, 3/000, 1/signextw, 1/0, 5/lower5, 6/registerZ }


\paragraph{Syntax} \texttt{iornw}\emph{signextw} \emph{registerW} = \emph{registerZ}, \emph{upper27\_\discretionary{}{}{}lower5}

\paragraph{Description} Bitwise inclusive or of not the \emph{registerZ} with the \emph{upper27\_\discretionary{}{}{}lower5}. The result with the 32 upper bits cleared is stored into the \emph{registerW}.


\paragraph{Modifier(s)}
\noindent\begin{itemize}
\item signextw (cf. signextw in section \ref{par:modifier-signextw} on page \pageref{par:modifier-signextw})
\end{itemize}

\paragraph{Execution} {Reservation table \textsf{ALU\_TINY.X} (Section~\ref{sec:reservation-ALU-TINY.X})}
~
{\begin{lstlisting}
stage ID:
  new signextw = _ZX_1(signextw);
stage RR:
  new argument3 = _WX_32(upper27_lower5);
  new argument2 = GPR[registerZ];
  new result1 = ~argument2 | argument3;
stage E1:
  GPR[registerW] = signextw ? _SX_32(result1) : _ZX_32(result1);
\end{lstlisting}}
\newpage
\subsubsection[{\small IORQ: Bitwise Inclusive Or Quadruple Word}]{IORQ\hfill Bitwise Inclusive Or Quadruple Word}
\label{instr:IORQ}\index{iorq}
 
\paragraph{Encoding} Format \textsf{ALU\_QWRR0} (Section~\ref{sec:format-ALU-QWRR0}), \verb|exucode2 = 1010|\\

\bitfields{ 1/P, 2/11, 1/0, 4/1010, 5/registerM, 1/--, 2/10, 3/000, 1/1, 5/registerO, 1/--, 5/registerP, 1/-- }


\paragraph{Syntax} \texttt{iorq} \emph{registerM} = \emph{registerP}, \emph{registerO}

\paragraph{Description} Bitwise inclusive or of the \emph{registerP} with the \emph{registerO}. The result is stored into the \emph{registerM}.


\paragraph{Execution} {Reservation table \textsf{ALU\_LITE} (Section~\ref{sec:reservation-ALU-LITE})}
~
{\begin{lstlisting}
stage RR:
  new argument3 = PGR[registerO];
  new argument2 = PGR[registerP];
  new result1 = argument2 | argument3;
stage E1:
  PGR[registerM] = result1;
\end{lstlisting}}
\newpage

\paragraph{Encoding} Format \textsf{ALU\_QWRR0.M} (Section~\ref{sec:format-ALU-QWRR0.M}), \verb|exucode2 = 1010|\\

\bitfields{ 1/P, 2/00, 2/-- --, 27/upper27, 1/1, 2/11, 1/0, 4/1010, 5/registerM, 1/--, 2/10, 3/000, 1/1, 1/splat32, 5/lower5, 5/registerP, 1/1 }


\paragraph{Syntax} \texttt{iorq} \emph{registerM} = \emph{registerP}, \emph{upper27\_\discretionary{}{}{}lower5}\emph{splat32}

\paragraph{Description} Bitwise inclusive or of the \emph{registerP} with the \emph{upper27\_\discretionary{}{}{}lower5}. The result is stored into the \emph{registerM}.


\paragraph{Modifier(s)}
\noindent\begin{itemize}
\item splat32 (cf. splat32 in section \ref{par:modifier-splat32} on page \pageref{par:modifier-splat32})
\end{itemize}

\paragraph{Execution} {Reservation table \textsf{ALU\_LITE.X} (Section~\ref{sec:reservation-ALU-LITE.X})}
~
{\begin{lstlisting}
stage ID:
  new splat32 = _ZX_1(splat32);
stage RR:
  new argument3 = splat32 ? (_WX_32(upper27_lower5) << 32) | _ZX_32(_WX_32(upper27_lower5)) : _SX_32(_WX_32(upper27_lower5));
  new argument2 = PGR[registerP];
  new result1 = argument2 | argument3;
stage E1:
  PGR[registerM] = result1;
\end{lstlisting}}
\newpage
\subsubsection[{\small IORW: Bitwise Inclusive Or Word}]{IORW\hfill Bitwise Inclusive Or Word}
\label{instr:IORW}\index{iorw}
 
\paragraph{Encoding} Format \textsf{ALU\_WWRR0} (Section~\ref{sec:format-ALU-WWRR0}), \verb|exucode2 = 1010|\\

\bitfields{ 1/P, 2/11, 1/0, 4/1010, 6/registerW, 2/11, 3/000, 1/signextw, 6/registerY, 6/registerZ }


\paragraph{Syntax} \texttt{iorw}\emph{signextw} \emph{registerW} = \emph{registerZ}, \emph{registerY}

\paragraph{Description} Bitwise inclusive or of the \emph{registerZ} with the \emph{registerY}. The result with the 32 upper bits cleared is stored into the \emph{registerW}.


\paragraph{Modifier(s)}
\noindent\begin{itemize}
\item signextw (cf. signextw in section \ref{par:modifier-signextw} on page \pageref{par:modifier-signextw})
\end{itemize}

\paragraph{Execution} {Reservation table \textsf{ALU\_TINY} (Section~\ref{sec:reservation-ALU-TINY})}
~
{\begin{lstlisting}
stage ID:
  new signextw = _ZX_1(signextw);
stage RR:
  new argument3 = GPR[registerY];
  new argument2 = GPR[registerZ];
  new result1 = argument2 | argument3;
stage E1:
  GPR[registerW] = signextw ? _SX_32(result1) : _ZX_32(result1);
\end{lstlisting}}
\newpage

\paragraph{Encoding} Format \textsf{ALU\_WWRR0.W} (Section~\ref{sec:format-ALU-WWRR0.W}), \verb|exucode2 = 1010|\\

\bitfields{ 1/P, 2/00, 2/-- --, 27/upper27, 1/1, 2/11, 1/0, 4/1010, 6/registerW, 2/11, 3/000, 1/signextw, 1/0, 5/lower5, 6/registerZ }


\paragraph{Syntax} \texttt{iorw}\emph{signextw} \emph{registerW} = \emph{registerZ}, \emph{upper27\_\discretionary{}{}{}lower5}

\paragraph{Description} Bitwise inclusive or of the \emph{registerZ} with the \emph{upper27\_\discretionary{}{}{}lower5}. The result with the 32 upper bits cleared is stored into the \emph{registerW}.


\paragraph{Modifier(s)}
\noindent\begin{itemize}
\item signextw (cf. signextw in section \ref{par:modifier-signextw} on page \pageref{par:modifier-signextw})
\end{itemize}

\paragraph{Execution} {Reservation table \textsf{ALU\_TINY.X} (Section~\ref{sec:reservation-ALU-TINY.X})}
~
{\begin{lstlisting}
stage ID:
  new signextw = _ZX_1(signextw);
stage RR:
  new argument3 = _WX_32(upper27_lower5);
  new argument2 = GPR[registerZ];
  new result1 = argument2 | argument3;
stage E1:
  GPR[registerW] = signextw ? _SX_32(result1) : _ZX_32(result1);
\end{lstlisting}}
\newpage
\subsubsection[{\small LANDD: Logical And Double Word}]{LANDD\hfill Logical And Double Word}
\label{instr:LANDD}\index{landd}
 
\paragraph{Encoding} Format \textsf{ALU\_DWRR2} (Section~\ref{sec:format-ALU-DWRR2}), \verb|exucode2 = 0000|\\

\bitfields{ 1/P, 2/11, 1/0, 4/0000, 6/registerW, 2/10, 3/010, 1/0, 6/registerY, 6/registerZ }


\paragraph{Syntax} \texttt{landd} \emph{registerW} = \emph{registerZ}, \emph{registerY}

\paragraph{Description} The \emph{registerY} and the \emph{registerZ} are converted to booleans, which are AND-ed. The result is stored into the \emph{registerW}.


\paragraph{Execution} {Reservation table \textsf{ALU\_TINY} (Section~\ref{sec:reservation-ALU-TINY})}
~
{\begin{lstlisting}
stage RR:
  new argument3 = GPR[registerY];
  new argument2 = GPR[registerZ];
  new result1 = _ZX_64(argument3) && _ZX_64(argument2);
stage E1:
  GPR[registerW] = result1;
\end{lstlisting}}
\newpage
\subsubsection[{\small LANDW: Logical And Word}]{LANDW\hfill Logical And Word}
\label{instr:LANDW}\index{landw}
 
\paragraph{Encoding} Format \textsf{ALU\_WWRR2} (Section~\ref{sec:format-ALU-WWRR2}), \verb|exucode2 = 0000|\\

\bitfields{ 1/P, 2/11, 1/0, 4/0000, 6/registerW, 2/11, 3/010, 1/0, 6/registerY, 6/registerZ }


\paragraph{Syntax} \texttt{landw} \emph{registerW} = \emph{registerZ}, \emph{registerY}

\paragraph{Description} The \emph{registerY} and the \emph{registerZ} are converted to booleans, which are AND-ed. The result is stored into the \emph{registerW}.


\paragraph{Execution} {Reservation table \textsf{ALU\_TINY} (Section~\ref{sec:reservation-ALU-TINY})}
~
{\begin{lstlisting}
stage RR:
  new argument3 = GPR[registerY];
  new argument2 = GPR[registerZ];
  new result1 = argument3.32[0] && argument2.32[0];
stage E1:
  GPR[registerW] = result1;
\end{lstlisting}}
\newpage
\subsubsection[{\small LIORD: Logical Or Double Word}]{LIORD\hfill Logical Or Double Word}
\label{instr:LIORD}\index{liord}
 
\paragraph{Encoding} Format \textsf{ALU\_DWRR2} (Section~\ref{sec:format-ALU-DWRR2}), \verb|exucode2 = 0010|\\

\bitfields{ 1/P, 2/11, 1/0, 4/0010, 6/registerW, 2/10, 3/010, 1/0, 6/registerY, 6/registerZ }


\paragraph{Syntax} \texttt{liord} \emph{registerW} = \emph{registerZ}, \emph{registerY}

\paragraph{Description} The \emph{registerY} and the \emph{registerZ} are converted to booleans, which are OR-ed. The result is stored into the \emph{registerW}.


\paragraph{Execution} {Reservation table \textsf{ALU\_TINY} (Section~\ref{sec:reservation-ALU-TINY})}
~
{\begin{lstlisting}
stage RR:
  new argument3 = GPR[registerY];
  new argument2 = GPR[registerZ];
  new result1 = _ZX_64(argument3) || _ZX_64(argument2);
stage E1:
  GPR[registerW] = result1;
\end{lstlisting}}
\newpage
\subsubsection[{\small LIORW: Logical Or Word}]{LIORW\hfill Logical Or Word}
\label{instr:LIORW}\index{liorw}
 
\paragraph{Encoding} Format \textsf{ALU\_WWRR2} (Section~\ref{sec:format-ALU-WWRR2}), \verb|exucode2 = 0010|\\

\bitfields{ 1/P, 2/11, 1/0, 4/0010, 6/registerW, 2/11, 3/010, 1/0, 6/registerY, 6/registerZ }


\paragraph{Syntax} \texttt{liorw} \emph{registerW} = \emph{registerZ}, \emph{registerY}

\paragraph{Description} The \emph{registerY} and the \emph{registerZ} are converted to booleans, which are OR-ed. The result is stored into the \emph{registerW}.


\paragraph{Execution} {Reservation table \textsf{ALU\_TINY} (Section~\ref{sec:reservation-ALU-TINY})}
~
{\begin{lstlisting}
stage RR:
  new argument3 = GPR[registerY];
  new argument2 = GPR[registerZ];
  new result1 = argument3.32[0] || argument2.32[0];
stage E1:
  GPR[registerW] = result1;
\end{lstlisting}}
\newpage
\subsubsection[{\small LNANDD: Logical Not And Double Word}]{LNANDD\hfill Logical Not And Double Word}
\label{instr:LNANDD}\index{lnandd}
 
\paragraph{Encoding} Format \textsf{ALU\_DWRR2} (Section~\ref{sec:format-ALU-DWRR2}), \verb|exucode2 = 0001|\\

\bitfields{ 1/P, 2/11, 1/0, 4/0001, 6/registerW, 2/10, 3/010, 1/0, 6/registerY, 6/registerZ }


\paragraph{Syntax} \texttt{lnandd} \emph{registerW} = \emph{registerZ}, \emph{registerY}

\paragraph{Description} The \emph{registerY} and the \emph{registerZ} are converted to booleans, which are NOT AND-ed. The result is stored into the \emph{registerW}.


\paragraph{Execution} {Reservation table \textsf{ALU\_TINY} (Section~\ref{sec:reservation-ALU-TINY})}
~
{\begin{lstlisting}
stage RR:
  new argument3 = GPR[registerY];
  new argument2 = GPR[registerZ];
  new result1 = !(_ZX_64(argument3) && _ZX_64(argument2));
stage E1:
  GPR[registerW] = result1;
\end{lstlisting}}
\newpage
\subsubsection[{\small LNANDW: Logical Not And Word}]{LNANDW\hfill Logical Not And Word}
\label{instr:LNANDW}\index{lnandw}
 
\paragraph{Encoding} Format \textsf{ALU\_WWRR2} (Section~\ref{sec:format-ALU-WWRR2}), \verb|exucode2 = 0001|\\

\bitfields{ 1/P, 2/11, 1/0, 4/0001, 6/registerW, 2/11, 3/010, 1/0, 6/registerY, 6/registerZ }


\paragraph{Syntax} \texttt{lnandw} \emph{registerW} = \emph{registerZ}, \emph{registerY}

\paragraph{Description} The \emph{registerY} and the \emph{registerZ} are converted to booleans, which are NOT AND-ed. The result is stored into the \emph{registerW}.


\paragraph{Execution} {Reservation table \textsf{ALU\_TINY} (Section~\ref{sec:reservation-ALU-TINY})}
~
{\begin{lstlisting}
stage RR:
  new argument3 = GPR[registerY];
  new argument2 = GPR[registerZ];
  new result1 = !(argument3.32[0] && argument2.32[0]);
stage E1:
  GPR[registerW] = result1;
\end{lstlisting}}
\newpage
\subsubsection[{\small LNIORD: Logical Not Or Double Word}]{LNIORD\hfill Logical Not Or Double Word}
\label{instr:LNIORD}\index{lniord}
 
\paragraph{Encoding} Format \textsf{ALU\_DWRR2} (Section~\ref{sec:format-ALU-DWRR2}), \verb|exucode2 = 0011|\\

\bitfields{ 1/P, 2/11, 1/0, 4/0011, 6/registerW, 2/10, 3/010, 1/0, 6/registerY, 6/registerZ }


\paragraph{Syntax} \texttt{lniord} \emph{registerW} = \emph{registerZ}, \emph{registerY}

\paragraph{Description} The \emph{registerY} and the \emph{registerZ} are converted to booleans, which are not OR-ed. The result is stored into the \emph{registerW}.


\paragraph{Execution} {Reservation table \textsf{ALU\_TINY} (Section~\ref{sec:reservation-ALU-TINY})}
~
{\begin{lstlisting}
stage RR:
  new argument3 = GPR[registerY];
  new argument2 = GPR[registerZ];
  new result1 = !(_ZX_64(argument3) || _ZX_64(argument2));
stage E1:
  GPR[registerW] = result1;
\end{lstlisting}}
\newpage
\subsubsection[{\small LNIORW: Logical Not Or Word}]{LNIORW\hfill Logical Not Or Word}
\label{instr:LNIORW}\index{lniorw}
 
\paragraph{Encoding} Format \textsf{ALU\_WWRR2} (Section~\ref{sec:format-ALU-WWRR2}), \verb|exucode2 = 0011|\\

\bitfields{ 1/P, 2/11, 1/0, 4/0011, 6/registerW, 2/11, 3/010, 1/0, 6/registerY, 6/registerZ }


\paragraph{Syntax} \texttt{lniorw} \emph{registerW} = \emph{registerZ}, \emph{registerY}

\paragraph{Description} The \emph{registerY} and the \emph{registerZ} are converted to booleans, which are not OR-ed. The result is stored into the \emph{registerW}.


\paragraph{Execution} {Reservation table \textsf{ALU\_TINY} (Section~\ref{sec:reservation-ALU-TINY})}
~
{\begin{lstlisting}
stage RR:
  new argument3 = GPR[registerY];
  new argument2 = GPR[registerZ];
  new result1 = !(argument3.32[0] || argument2.32[0]);
stage E1:
  GPR[registerW] = result1;
\end{lstlisting}}
\newpage
\subsubsection[{\small MADDBHO: Multiply Add Extend Byte to Half Word Octuple}]{MADDBHO\hfill Multiply Add Extend Byte to Half Word Octuple}
\label{instr:MADDBHO}\index{maddbho}
 
\paragraph{Encoding} Format \textsf{ALU\_HOMEWRRRE} (Section~\ref{sec:format-ALU-HOMEWRRRE}), \verb|alucode2 = 10|\\

\bitfields{ 1/P, 2/11, 1/1, 2/10, 2/widemult, 5/registerM, 1/0, 2/10, 3/110, 1/1, 5/registerYe, 1/1, 5/registerZe, 1/0 }


\paragraph{Syntax} \texttt{maddbho}\emph{widemult} \emph{registerM} = \emph{registerZe}, \emph{registerYe}

\paragraph{Description} The \emph{registerYe} and the \emph{registerZe} are interpreted as eight bytes packed into 64 bits. These \emph{registerYe} and \emph{registerZe} bytes are sign-extended or zero-extended from 8 bits, depending on the \emph{widemult}. The \emph{registerYe} bytes are multiplied by the \emph{registerZe} bytes, producing eight 16-bit results. The eight resulting half words are packed, added to and stored into the \emph{registerM}.


\paragraph{Modifier(s)}
\noindent\begin{itemize}
\item widemult (cf. widemult in section \ref{par:modifier-widemult} on page \pageref{par:modifier-widemult})
\end{itemize}

\paragraph{Execution} {Reservation table \textsf{ALU\_LITE} (Section~\ref{sec:reservation-ALU-LITE})}
~
{\begin{lstlisting}
stage ID:
  new widemult = _ZX_2(widemult);
stage RR:
  new argument3 = GPR[registerYe<<1];
  new argument2 = GPR[registerZe<<1];
stage E1:
  new argument1 = PGR[registerM];
  for (I = 0; I < 8; I += 1) {
    if (widemult == 0) {
      new product = _SX_8(argument2.8[I]) * _SX_8(argument3.8[I]);
    }
    if (widemult == 1) {
      product = _ZX_8(argument2.8[I]) * _ZX_8(argument3.8[I]);
    }
    if (widemult == 2) {
      product = _SX_8(argument2.8[I]) * _ZX_8(argument3.8[I]);
    }
    result1.16[I] = argument1.16[I] + product
  };
stage E2:
  PGR[registerM] = result1;
\end{lstlisting}}
\newpage

\paragraph{Encoding} Format \textsf{ALU\_HOMEWRRRO} (Section~\ref{sec:format-ALU-HOMEWRRRO}), \verb|alucode2 = 10|\\

\bitfields{ 1/P, 2/11, 1/1, 2/10, 2/widemult, 5/registerM, 1/1, 2/10, 3/110, 1/1, 5/registerYo, 1/1, 5/registerZo, 1/0 }


\paragraph{Syntax} \texttt{maddbho}\emph{widemult} \emph{registerM} = \emph{registerZo}, \emph{registerYo}

\paragraph{Description} The \emph{registerYo} and the \emph{registerZo} are interpreted as eight bytes packed into 64 bits. These \emph{registerYo} and \emph{registerZo} bytes are sign-extended or zero-extended from 8 bits, depending on the \emph{widemult}. The \emph{registerYo} bytes are multiplied by the \emph{registerZo} bytes, producing eight 16-bit results. The eight resulting half words are packed, added to and stored into the \emph{registerM}.


\paragraph{Modifier(s)}
\noindent\begin{itemize}
\item widemult (cf. widemult in section \ref{par:modifier-widemult} on page \pageref{par:modifier-widemult})
\end{itemize}

\paragraph{Execution} {Reservation table \textsf{ALU\_LITE} (Section~\ref{sec:reservation-ALU-LITE})}
~
{\begin{lstlisting}
stage ID:
  new widemult = _ZX_2(widemult);
stage RR:
  new argument3 = GPR[(registerYo<<1)+1];
  new argument2 = GPR[(registerZo<<1)+1];
stage E1:
  new argument1 = PGR[registerM];
  for (I = 0; I < 8; I += 1) {
    if (widemult == 0) {
      new product = _SX_8(argument2.8[I]) * _SX_8(argument3.8[I]);
    }
    if (widemult == 1) {
      product = _ZX_8(argument2.8[I]) * _ZX_8(argument3.8[I]);
    }
    if (widemult == 2) {
      product = _SX_8(argument2.8[I]) * _ZX_8(argument3.8[I]);
    }
    result1.16[I] = argument1.16[I] + product
  };
stage E2:
  PGR[registerM] = result1;
\end{lstlisting}}
\newpage
\subsubsection[{\small MADDD: Multiply Add Double Word}]{MADDD\hfill Multiply Add Double Word}
\label{instr:MADDD}\index{maddd}
 
\paragraph{Encoding} Format \textsf{ALU\_DMWRRR} (Section~\ref{sec:format-ALU-DMWRRR}), \verb|alucode2 = 10|\\

\bitfields{ 1/P, 2/11, 1/0, 2/10, 2/highmult, 6/registerW, 2/10, 3/110, 1/0, 6/registerY, 6/registerZ }


\paragraph{Syntax} \texttt{maddd}\emph{highmult} \emph{registerW} = \emph{registerZ}, \emph{registerY}

\paragraph{Description} The \emph{registerY} and the \emph{registerZ} double words are sign-extended or zero-extended from 64 bits, depending on the \emph{highmult}. The \emph{registerY} double word is multiplied by the \emph{registerZ} double word, and the product is optionally shifted right by 64 bits, depending on the \emph{highmult}. The resulting double word is added to and stored into the \emph{registerW}.


\paragraph{Modifier(s)}
\noindent\begin{itemize}
\item highmult (cf. highmult in section \ref{par:modifier-highmult} on page \pageref{par:modifier-highmult})
\end{itemize}

\paragraph{Execution} {Reservation table \textsf{ALU\_LITE} (Section~\ref{sec:reservation-ALU-LITE})}
~
{\begin{lstlisting}
stage ID:
  new highmult = _ZX_2(highmult);
stage RR:
  new argument3 = GPR[registerY];
  new argument2 = GPR[registerZ];
stage E1:
  new argument1 = GPR[registerW];
  if ((highmult == 0) || (highmult == 3)) {
    new product = _SX_64(argument2) * _SX_64(argument3);
  }
  if (highmult == 1) {
    product = _ZX_64(argument2) * _ZX_64(argument3);
  }
  if (highmult == 2) {
    product = _SX_64(argument2) * _ZX_64(argument3);
  }
  if (highmult) {
    new result1 = argument1 + product.64[1];
  } else {
    result1 = argument1 + product.64[0];
  }
stage E2:
  GPR[registerW] = result1;
\end{lstlisting}}
\newpage

\paragraph{Encoding} Format \textsf{ALU\_DMWRRR.M} (Section~\ref{sec:format-ALU-DMWRRR.M}), \verb|alucode2 = 10|\\

\bitfields{ 1/P, 2/00, 2/-- --, 27/upper27, 1/1, 2/11, 1/0, 2/10, 2/highmult, 6/registerW, 2/10, 3/110, 1/0, 1/splat32, 5/lower5, 6/registerZ }


\paragraph{Syntax} \texttt{maddd}\emph{highmult} \emph{registerW} = \emph{registerZ}, \emph{upper27\_\discretionary{}{}{}lower5}\emph{splat32}

\paragraph{Description} The \emph{upper27\_\discretionary{}{}{}lower5} and the \emph{registerZ} double words are sign-extended or zero-extended from 64 bits, depending on the \emph{highmult}. The \emph{upper27\_\discretionary{}{}{}lower5} double word is multiplied by the \emph{registerZ} double word, and the product is optionally shifted right by 64 bits, depending on the \emph{highmult}. The resulting double word is added to and stored into the \emph{registerW}.


\paragraph{Modifier(s)}
\noindent\begin{itemize}
\item highmult (cf. highmult in section \ref{par:modifier-highmult} on page \pageref{par:modifier-highmult})
\item splat32 (cf. splat32 in section \ref{par:modifier-splat32} on page \pageref{par:modifier-splat32})
\end{itemize}

\paragraph{Execution} {Reservation table \textsf{ALU\_LITE.X} (Section~\ref{sec:reservation-ALU-LITE.X})}
~
{\begin{lstlisting}
stage ID:
  new splat32 = _ZX_1(splat32);
  new highmult = _ZX_2(highmult);
stage RR:
  new argument3 = splat32 ? (_WX_32(upper27_lower5) << 32) | _ZX_32(_WX_32(upper27_lower5)) : _SX_32(_WX_32(upper27_lower5));
  new argument2 = GPR[registerZ];
stage E1:
  new argument1 = GPR[registerW];
  if ((highmult == 0) || (highmult == 3)) {
    new product = _SX_64(argument2) * _SX_64(argument3);
  }
  if (highmult == 1) {
    product = _ZX_64(argument2) * _ZX_64(argument3);
  }
  if (highmult == 2) {
    product = _SX_64(argument2) * _ZX_64(argument3);
  }
  if (highmult) {
    new result1 = argument1 + product.64[1];
  } else {
    result1 = argument1 + product.64[0];
  }
stage E2:
  GPR[registerW] = result1;
\end{lstlisting}}
\newpage
\subsubsection[{\small MADDDQ: Multiply Add Extend Double Word to Quadruple Word}]{MADDDQ\hfill Multiply Add Extend Double Word to Quadruple Word}
\label{instr:MADDDQ}\index{madddq}
 
\paragraph{Encoding} Format \textsf{ALU\_QMEWRRR} (Section~\ref{sec:format-ALU-QMEWRRR}), \verb|alucode2 = 10|\\

\bitfields{ 1/P, 2/11, 1/0, 2/10, 2/widemult, 5/registerM, 1/--, 2/10, 3/111, 1/1, 6/registerY, 6/registerZ }


\paragraph{Syntax} \texttt{madddq}\emph{widemult} \emph{registerM} = \emph{registerZ}, \emph{registerY}

\paragraph{Description} The \emph{registerY} and the \emph{registerZ} double words are sign-extended or zero-extended from 64 bits, depending on the \emph{widemult}. The \emph{registerY} is multiplied by the \emph{registerZ}, producing a 128-bit result. The resulting quadruple word is added to and stored into the \emph{registerM}.


\paragraph{Modifier(s)}
\noindent\begin{itemize}
\item widemult (cf. widemult in section \ref{par:modifier-widemult} on page \pageref{par:modifier-widemult})
\end{itemize}

\paragraph{Execution} {Reservation table \textsf{ALU\_LITE} (Section~\ref{sec:reservation-ALU-LITE})}
~
{\begin{lstlisting}
stage ID:
  new widemult = _ZX_2(widemult);
stage RR:
  new argument3 = GPR[registerY];
  new argument2 = GPR[registerZ];
stage E1:
  new argument1 = PGR[registerM];
  if (widemult == 0) {
    new product = _SX_64(argument2) * _SX_64(argument3);
  }
  if (widemult == 1) {
    product = _ZX_64(argument2) * _ZX_64(argument3);
  }
  if (widemult == 2) {
    product = _SX_64(argument2) * _ZX_64(argument3);
  }
  new result1 = argument1 + product;
stage E2:
  PGR[registerM] = result1;
\end{lstlisting}}
\newpage
\subsubsection[{\small MADDHQ: Multiply Add Half Word Quadruple}]{MADDHQ\hfill Multiply Add Half Word Quadruple}
\label{instr:MADDHQ}\index{maddhq}
 
\paragraph{Encoding} Format \textsf{ALU\_HQMWRRRE} (Section~\ref{sec:format-ALU-HQMWRRRE}), \verb|alucode2 = 10|\\

\bitfields{ 1/P, 2/11, 1/1, 2/10, 2/highmult, 5/registerWe, 1/0, 2/10, 3/110, 1/1, 5/registerYe, 1/0, 5/registerZe, 1/0 }


\paragraph{Syntax} \texttt{maddhq}\emph{highmult} \emph{registerWe} = \emph{registerZe}, \emph{registerYe}

\paragraph{Description} The \emph{registerYe} and the \emph{registerZe} are interpreted as four half words packed into 64 bits. The \emph{registerYe} and the \emph{registerZe} half words are sign-extended or zero-extended from 16 bits, depending on the \emph{highmult}. The \emph{registerYe} half words are multiplied by the \emph{registerZe} half words, and the product is optionally shifted right by 16 bits, depending on the \emph{highmult}. The four resulting half words are added to the \emph{registerWe} half words, packed and stored back into the \emph{registerWe}.


\paragraph{Modifier(s)}
\noindent\begin{itemize}
\item highmult (cf. highmult in section \ref{par:modifier-highmult} on page \pageref{par:modifier-highmult})
\end{itemize}

\paragraph{Execution} {Reservation table \textsf{ALU\_LITE} (Section~\ref{sec:reservation-ALU-LITE})}
~
{\begin{lstlisting}
stage ID:
  new highmult = _ZX_2(highmult);
stage RR:
  new argument3 = GPR[registerYe<<1];
  new argument2 = GPR[registerZe<<1];
stage E1:
  new argument1 = GPR[registerWe<<1];
  for (I = 0; I < 4; I += 1) {
    if ((highmult == 0) || (highmult == 3)) {
      new product = _SX_16(argument2.16[I]) * _SX_16(argument3.16[I]);
    }
    if (highmult == 1) {
      product = _ZX_16(argument2.16[I]) * _ZX_16(argument3.16[I]);
    }
    if (highmult == 2) {
      product = _SX_16(argument2.16[I]) * _ZX_16(argument3.16[I]);
    }
    if (highmult) {
      result1.16[I] = argument1.16[I] + product.16[1];
    } else {
      result1.16[I] = argument1.16[I] + product.16[0];
    }
  };
stage E2:
  GPR[registerWe<<1] = result1;
\end{lstlisting}}
\newpage

\paragraph{Encoding} Format \textsf{ALU\_HQMWRRRO} (Section~\ref{sec:format-ALU-HQMWRRRO}), \verb|alucode2 = 10|\\

\bitfields{ 1/P, 2/11, 1/1, 2/10, 2/highmult, 5/registerWo, 1/1, 2/10, 3/110, 1/1, 5/registerYo, 1/0, 5/registerZo, 1/0 }


\paragraph{Syntax} \texttt{maddhq}\emph{highmult} \emph{registerWo} = \emph{registerZo}, \emph{registerYo}

\paragraph{Description} The \emph{registerYo} and the \emph{registerZo} are interpreted as four half words packed into 64 bits. The \emph{registerYo} and the \emph{registerZo} half words are sign-extended or zero-extended from 16 bits, depending on the \emph{highmult}. The \emph{registerYo} half words are multiplied by the \emph{registerZo} half words, and the product is optionally shifted right by 16 bits, depending on the \emph{highmult}. The four resulting half words are added to the \emph{registerWo} half words, packed and stored back into the \emph{registerWo}.


\paragraph{Modifier(s)}
\noindent\begin{itemize}
\item highmult (cf. highmult in section \ref{par:modifier-highmult} on page \pageref{par:modifier-highmult})
\end{itemize}

\paragraph{Execution} {Reservation table \textsf{ALU\_LITE} (Section~\ref{sec:reservation-ALU-LITE})}
~
{\begin{lstlisting}
stage ID:
  new highmult = _ZX_2(highmult);
stage RR:
  new argument3 = GPR[(registerYo<<1)+1];
  new argument2 = GPR[(registerZo<<1)+1];
stage E1:
  new argument1 = GPR[(registerWo<<1)+1];
  for (I = 0; I < 4; I += 1) {
    if ((highmult == 0) || (highmult == 3)) {
      new product = _SX_16(argument2.16[I]) * _SX_16(argument3.16[I]);
    }
    if (highmult == 1) {
      product = _ZX_16(argument2.16[I]) * _ZX_16(argument3.16[I]);
    }
    if (highmult == 2) {
      product = _SX_16(argument2.16[I]) * _ZX_16(argument3.16[I]);
    }
    if (highmult) {
      result1.16[I] = argument1.16[I] + product.16[1];
    } else {
      result1.16[I] = argument1.16[I] + product.16[0];
    }
  };
stage E2:
  GPR[(registerWo<<1)+1] = result1;
\end{lstlisting}}
\newpage
\subsubsection[{\small MADDHWQ: Multiply Add Extend Half Word to Word Quadruple}]{MADDHWQ\hfill Multiply Add Extend Half Word to Word Quadruple}
\label{instr:MADDHWQ}\index{maddhwq}
 
\paragraph{Encoding} Format \textsf{ALU\_WQMEWRRRE} (Section~\ref{sec:format-ALU-WQMEWRRRE}), \verb|alucode2 = 10|\\

\bitfields{ 1/P, 2/11, 1/1, 2/10, 2/widemult, 5/registerM, 1/0, 2/10, 3/110, 1/0, 5/registerYe, 1/1, 5/registerZe, 1/1 }


\paragraph{Syntax} \texttt{maddhwq}\emph{widemult} \emph{registerM} = \emph{registerZe}, \emph{registerYe}

\paragraph{Description} The \emph{registerYe} and the \emph{registerZe} are interpreted as four half words packed into 64 bits. These \emph{registerYe} and \emph{registerZe} half words are sign-extended or zero-extended from 16 bits, depending on the \emph{widemult}. The \emph{registerYe} half words are multiplied by the \emph{registerZe} half words, producing four 32-bit results. The four resulting words are packed, added to and stored into the \emph{registerM}.


\paragraph{Modifier(s)}
\noindent\begin{itemize}
\item widemult (cf. widemult in section \ref{par:modifier-widemult} on page \pageref{par:modifier-widemult})
\end{itemize}

\paragraph{Execution} {Reservation table \textsf{ALU\_LITE} (Section~\ref{sec:reservation-ALU-LITE})}
~
{\begin{lstlisting}
stage ID:
  new widemult = _ZX_2(widemult);
stage RR:
  new argument3 = GPR[registerYe<<1];
  new argument2 = GPR[registerZe<<1];
stage E1:
  new argument1 = PGR[registerM];
  for (I = 0; I < 4; I += 1) {
    if (widemult == 0) {
      new product = _SX_16(argument2.16[I]) * _SX_16(argument3.16[I]);
    }
    if (widemult == 1) {
      product = _ZX_16(argument2.16[I]) * _ZX_16(argument3.16[I]);
    }
    if (widemult == 2) {
      product = _SX_16(argument2.16[I]) * _ZX_16(argument3.16[I]);
    }
    result1.32[I] = argument1.32[I] + product
  };
stage E2:
  PGR[registerM] = result1;
\end{lstlisting}}
\newpage

\paragraph{Encoding} Format \textsf{ALU\_WQMEWRRRO} (Section~\ref{sec:format-ALU-WQMEWRRRO}), \verb|alucode2 = 10|\\

\bitfields{ 1/P, 2/11, 1/1, 2/10, 2/widemult, 5/registerM, 1/1, 2/10, 3/110, 1/0, 5/registerYo, 1/1, 5/registerZo, 1/1 }


\paragraph{Syntax} \texttt{maddhwq}\emph{widemult} \emph{registerM} = \emph{registerZo}, \emph{registerYo}

\paragraph{Description} The \emph{registerYo} and the \emph{registerZo} are interpreted as four half words packed into 64 bits. These \emph{registerYo} and \emph{registerZo} half words are sign-extended or zero-extended from 16 bits, depending on the \emph{widemult}. The \emph{registerYo} half words are multiplied by the \emph{registerZo} half words, producing four 32-bit results. The four resulting words are packed, added to and stored into the \emph{registerM}.


\paragraph{Modifier(s)}
\noindent\begin{itemize}
\item widemult (cf. widemult in section \ref{par:modifier-widemult} on page \pageref{par:modifier-widemult})
\end{itemize}

\paragraph{Execution} {Reservation table \textsf{ALU\_LITE} (Section~\ref{sec:reservation-ALU-LITE})}
~
{\begin{lstlisting}
stage ID:
  new widemult = _ZX_2(widemult);
stage RR:
  new argument3 = GPR[(registerYo<<1)+1];
  new argument2 = GPR[(registerZo<<1)+1];
stage E1:
  new argument1 = PGR[registerM];
  for (I = 0; I < 4; I += 1) {
    if (widemult == 0) {
      new product = _SX_16(argument2.16[I]) * _SX_16(argument3.16[I]);
    }
    if (widemult == 1) {
      product = _ZX_16(argument2.16[I]) * _ZX_16(argument3.16[I]);
    }
    if (widemult == 2) {
      product = _SX_16(argument2.16[I]) * _ZX_16(argument3.16[I]);
    }
    result1.32[I] = argument1.32[I] + product
  };
stage E2:
  PGR[registerM] = result1;
\end{lstlisting}}
\newpage
\subsubsection[{\small MADDW: Multiply Add Word}]{MADDW\hfill Multiply Add Word}
\label{instr:MADDW}\index{maddw}
 
\paragraph{Encoding} Format \textsf{ALU\_WMWRRR} (Section~\ref{sec:format-ALU-WMWRRR}), \verb|alucode2 = 10|\\

\bitfields{ 1/P, 2/11, 1/0, 2/10, 2/highmult, 6/registerW, 2/11, 3/110, 1/signextw, 6/registerY, 6/registerZ }


\paragraph{Syntax} \texttt{maddw}\emph{highmult}\emph{signextw} \emph{registerW} = \emph{registerZ}, \emph{registerY}

\paragraph{Description} The \emph{registerY} and the \emph{registerZ} words are sign-extended or zero-extended from 32 bits, depending on the \emph{highmult}. The \emph{registerY} word is multiplied by the \emph{registerZ} word, and the product is optionally shifted right by 32 bits, depending on the \emph{highmult}. The resulting word is added to and stored into the \emph{registerW}.


\paragraph{Modifier(s)}
\noindent\begin{itemize}
\item highmult (cf. highmult in section \ref{par:modifier-highmult} on page \pageref{par:modifier-highmult})
\item signextw (cf. signextw in section \ref{par:modifier-signextw} on page \pageref{par:modifier-signextw})
\end{itemize}

\paragraph{Execution} {Reservation table \textsf{ALU\_LITE} (Section~\ref{sec:reservation-ALU-LITE})}
~
{\begin{lstlisting}
stage ID:
  new signextw = _ZX_1(signextw);
  new highmult = _ZX_2(highmult);
stage RR:
  new argument3 = GPR[registerY];
  new argument2 = GPR[registerZ];
stage E1:
  new argument1 = GPR[registerW];
  if ((highmult == 0) || (highmult == 3)) {
    new product = _SX_32(argument2) * _SX_32(argument3);
  }
  if (highmult == 1) {
    product = _ZX_32(argument2) * _ZX_32(argument3);
  }
  if (highmult == 2) {
    product = _SX_32(argument2) * _ZX_32(argument3);
  }
  if (highmult) {
    new result1 = argument1 + product.32[1];
  } else {
    result1 = argument1 + product.32[0];
  }
stage E2:
  GPR[registerW] = signextw ? _SX_32(result1) : _ZX_32(result1);
\end{lstlisting}}
\newpage

\paragraph{Encoding} Format \textsf{ALU\_WMWRRR.W} (Section~\ref{sec:format-ALU-WMWRRR.W}), \verb|alucode2 = 10|\\

\bitfields{ 1/P, 2/00, 2/-- --, 27/upper27, 1/1, 2/11, 1/0, 2/10, 2/highmult, 6/registerW, 2/11, 3/110, 1/signextw, 1/0, 5/lower5, 6/registerZ }


\paragraph{Syntax} \texttt{maddw}\emph{highmult}\emph{signextw} \emph{registerW} = \emph{registerZ}, \emph{upper27\_\discretionary{}{}{}lower5}

\paragraph{Description} The \emph{upper27\_\discretionary{}{}{}lower5} and the \emph{registerZ} words are sign-extended or zero-extended from 32 bits, depending on the \emph{highmult}. The \emph{upper27\_\discretionary{}{}{}lower5} word is multiplied by the \emph{registerZ} word, and the product is optionally shifted right by 32 bits, depending on the \emph{highmult}. The resulting word is added to and stored into the \emph{registerW}.


\paragraph{Modifier(s)}
\noindent\begin{itemize}
\item highmult (cf. highmult in section \ref{par:modifier-highmult} on page \pageref{par:modifier-highmult})
\item signextw (cf. signextw in section \ref{par:modifier-signextw} on page \pageref{par:modifier-signextw})
\end{itemize}

\paragraph{Execution} {Reservation table \textsf{ALU\_LITE.X} (Section~\ref{sec:reservation-ALU-LITE.X})}
~
{\begin{lstlisting}
stage ID:
  new signextw = _ZX_1(signextw);
  new highmult = _ZX_2(highmult);
stage RR:
  new argument3 = _WX_32(upper27_lower5);
  new argument2 = GPR[registerZ];
stage E1:
  new argument1 = GPR[registerW];
  if ((highmult == 0) || (highmult == 3)) {
    new product = _SX_32(argument2) * _SX_32(argument3);
  }
  if (highmult == 1) {
    product = _ZX_32(argument2) * _ZX_32(argument3);
  }
  if (highmult == 2) {
    product = _SX_32(argument2) * _ZX_32(argument3);
  }
  if (highmult) {
    new result1 = argument1 + product.32[1];
  } else {
    result1 = argument1 + product.32[0];
  }
stage E2:
  GPR[registerW] = signextw ? _SX_32(result1) : _ZX_32(result1);
\end{lstlisting}}
\newpage
\subsubsection[{\small MADDWD: Multiply Add Extend Word to Double Word}]{MADDWD\hfill Multiply Add Extend Word to Double Word}
\label{instr:MADDWD}\index{maddwd}
 
\paragraph{Encoding} Format \textsf{ALU\_DMEWRRR} (Section~\ref{sec:format-ALU-DMEWRRR}), \verb|alucode2 = 10|\\

\bitfields{ 1/P, 2/11, 1/0, 2/10, 2/widemult, 6/registerW, 2/10, 3/111, 1/0, 6/registerY, 6/registerZ }


\paragraph{Syntax} \texttt{maddwd}\emph{widemult} \emph{registerW} = \emph{registerZ}, \emph{registerY}

\paragraph{Description} The \emph{registerY} and the \emph{registerZ} words are sign-extended or zero-extended from 32 bits, depending on the \emph{widemult}. The \emph{registerY} is multiplied by the \emph{registerZ}, producing a 64-bit result. The resulting double word is added to and stored into the \emph{registerW}.


\paragraph{Modifier(s)}
\noindent\begin{itemize}
\item widemult (cf. widemult in section \ref{par:modifier-widemult} on page \pageref{par:modifier-widemult})
\end{itemize}

\paragraph{Execution} {Reservation table \textsf{ALU\_LITE} (Section~\ref{sec:reservation-ALU-LITE})}
~
{\begin{lstlisting}
stage ID:
  new widemult = _ZX_2(widemult);
stage RR:
  new argument3 = GPR[registerY];
  new argument2 = GPR[registerZ];
stage E1:
  new argument1 = GPR[registerW];
  if (widemult == 0) {
    new product = _SX_32(argument2) * _SX_32(argument3);
  }
  if (widemult == 1) {
    product = _ZX_32(argument2) * _ZX_32(argument3);
  }
  if (widemult == 2) {
    product = _SX_32(argument2) * _ZX_32(argument3);
  }
  new result1 = argument1 + product;
stage E2:
  GPR[registerW] = result1;
\end{lstlisting}}
\newpage

\paragraph{Encoding} Format \textsf{ALU\_DMEWRRR.M} (Section~\ref{sec:format-ALU-DMEWRRR.M}), \verb|alucode2 = 10|\\

\bitfields{ 1/P, 2/00, 2/-- --, 27/upper27, 1/1, 2/11, 1/0, 2/10, 2/widemult, 6/registerW, 2/10, 3/111, 1/0, 1/splat32, 5/lower5, 6/registerZ }


\paragraph{Syntax} \texttt{maddwd}\emph{widemult} \emph{registerW} = \emph{registerZ}, \emph{upper27\_\discretionary{}{}{}lower5}\emph{splat32}

\paragraph{Description} The \emph{upper27\_\discretionary{}{}{}lower5} and the \emph{registerZ} words are sign-extended or zero-extended from 32 bits, depending on the \emph{widemult}. The \emph{upper27\_\discretionary{}{}{}lower5} is multiplied by the \emph{registerZ}, producing a 64-bit result. The resulting double word is added to and stored into the \emph{registerW}.


\paragraph{Modifier(s)}
\noindent\begin{itemize}
\item widemult (cf. widemult in section \ref{par:modifier-widemult} on page \pageref{par:modifier-widemult})
\item splat32 (cf. splat32 in section \ref{par:modifier-splat32} on page \pageref{par:modifier-splat32})
\end{itemize}

\paragraph{Execution} {Reservation table \textsf{ALU\_LITE.X} (Section~\ref{sec:reservation-ALU-LITE.X})}
~
{\begin{lstlisting}
stage ID:
  new splat32 = _ZX_1(splat32);
  new widemult = _ZX_2(widemult);
stage RR:
  new argument3 = splat32 ? (_WX_32(upper27_lower5) << 32) | _ZX_32(_WX_32(upper27_lower5)) : _SX_32(_WX_32(upper27_lower5));
  new argument2 = GPR[registerZ];
stage E1:
  new argument1 = GPR[registerW];
  if (widemult == 0) {
    new product = _SX_32(argument2) * _SX_32(argument3);
  }
  if (widemult == 1) {
    product = _ZX_32(argument2) * _ZX_32(argument3);
  }
  if (widemult == 2) {
    product = _SX_32(argument2) * _ZX_32(argument3);
  }
  new result1 = argument1 + product;
stage E2:
  GPR[registerW] = result1;
\end{lstlisting}}
\newpage
\subsubsection[{\small MADDWDP: Multiply Add Extend Word to Double Word Pair}]{MADDWDP\hfill Multiply Add Extend Word to Double Word Pair}
\label{instr:MADDWDP}\index{maddwdp}
 
\paragraph{Encoding} Format \textsf{ALU\_DPMEWRRRE} (Section~\ref{sec:format-ALU-DPMEWRRRE}), \verb|alucode2 = 10|\\

\bitfields{ 1/P, 2/11, 1/1, 2/10, 2/widemult, 5/registerM, 1/0, 2/10, 3/110, 1/0, 5/registerYe, 1/1, 5/registerZe, 1/0 }


\paragraph{Syntax} \texttt{maddwdp}\emph{widemult} \emph{registerM} = \emph{registerZe}, \emph{registerYe}

\paragraph{Description} The \emph{registerYe} and the \emph{registerZe} are interpreted as two words packed into 64 bits. These \emph{registerYe} and \emph{registerZe} words are sign-extended or zero-extended from 32 bits, depending on the \emph{widemult}. The \emph{registerYe} words are multiplied by the \emph{registerZe} words, producing two 64-bit results. The two resulting double words are packed, added to and stored into the \emph{registerM}.


\paragraph{Modifier(s)}
\noindent\begin{itemize}
\item widemult (cf. widemult in section \ref{par:modifier-widemult} on page \pageref{par:modifier-widemult})
\end{itemize}

\paragraph{Execution} {Reservation table \textsf{ALU\_LITE} (Section~\ref{sec:reservation-ALU-LITE})}
~
{\begin{lstlisting}
stage ID:
  new widemult = _ZX_2(widemult);
stage RR:
  new argument3 = GPR[registerYe<<1];
  new argument2 = GPR[registerZe<<1];
stage E1:
  new argument1 = PGR[registerM];
  for (I = 0; I < 2; I += 1) {
    if (widemult == 0) {
      new product = _SX_32(argument2.32[I]) * _SX_32(argument3.32[I]);
    }
    if (widemult == 1) {
      product = _ZX_32(argument2.32[I]) * _ZX_32(argument3.32[I]);
    }
    if (widemult == 2) {
      product = _SX_32(argument2.32[I]) * _ZX_32(argument3.32[I]);
    }
    result1.64[I] = argument1.64[I] + product
  };
stage E2:
  PGR[registerM] = result1;
\end{lstlisting}}
\newpage

\paragraph{Encoding} Format \textsf{ALU\_DPMEWRRRO} (Section~\ref{sec:format-ALU-DPMEWRRRO}), \verb|alucode2 = 10|\\

\bitfields{ 1/P, 2/11, 1/1, 2/10, 2/widemult, 5/registerM, 1/1, 2/10, 3/110, 1/0, 5/registerYo, 1/1, 5/registerZo, 1/0 }


\paragraph{Syntax} \texttt{maddwdp}\emph{widemult} \emph{registerM} = \emph{registerZo}, \emph{registerYo}

\paragraph{Description} The \emph{registerYo} and the \emph{registerZo} are interpreted as two words packed into 64 bits. These \emph{registerYo} and \emph{registerZo} words are sign-extended or zero-extended from 32 bits, depending on the \emph{widemult}. The \emph{registerYo} words are multiplied by the \emph{registerZo} words, producing two 64-bit results. The two resulting double words are packed, added to and stored into the \emph{registerM}.


\paragraph{Modifier(s)}
\noindent\begin{itemize}
\item widemult (cf. widemult in section \ref{par:modifier-widemult} on page \pageref{par:modifier-widemult})
\end{itemize}

\paragraph{Execution} {Reservation table \textsf{ALU\_LITE} (Section~\ref{sec:reservation-ALU-LITE})}
~
{\begin{lstlisting}
stage ID:
  new widemult = _ZX_2(widemult);
stage RR:
  new argument3 = GPR[(registerYo<<1)+1];
  new argument2 = GPR[(registerZo<<1)+1];
stage E1:
  new argument1 = PGR[registerM];
  for (I = 0; I < 2; I += 1) {
    if (widemult == 0) {
      new product = _SX_32(argument2.32[I]) * _SX_32(argument3.32[I]);
    }
    if (widemult == 1) {
      product = _ZX_32(argument2.32[I]) * _ZX_32(argument3.32[I]);
    }
    if (widemult == 2) {
      product = _SX_32(argument2.32[I]) * _ZX_32(argument3.32[I]);
    }
    result1.64[I] = argument1.64[I] + product
  };
stage E2:
  PGR[registerM] = result1;
\end{lstlisting}}
\newpage
\subsubsection[{\small MADDWP: Multiply Add Word Pair}]{MADDWP\hfill Multiply Add Word Pair}
\label{instr:MADDWP}\index{maddwp}
 
\paragraph{Encoding} Format \textsf{ALU\_WPMWRRRE} (Section~\ref{sec:format-ALU-WPMWRRRE}), \verb|alucode2 = 10|\\

\bitfields{ 1/P, 2/11, 1/1, 2/10, 2/highmult, 5/registerWe, 1/0, 2/10, 3/110, 1/0, 5/registerYe, 1/0, 5/registerZe, 1/1 }


\paragraph{Syntax} \texttt{maddwp}\emph{highmult} \emph{registerWe} = \emph{registerZe}, \emph{registerYe}

\paragraph{Description} The \emph{registerYe} and the \emph{registerZe} are interpreted as two words packed into 64 bits. The \emph{registerYe} and the \emph{registerZe} words are sign-extended or zero-extended from 32 bits, depending on the \emph{highmult}. The \emph{registerYe} words are multiplied by the \emph{registerZe} words, and the product is optionally shifted right by 32 bits, depending on the \emph{highmult}. The two resulting words are added to the \emph{registerWe} words, packed and stored back into the \emph{registerWe}.


\paragraph{Modifier(s)}
\noindent\begin{itemize}
\item highmult (cf. highmult in section \ref{par:modifier-highmult} on page \pageref{par:modifier-highmult})
\end{itemize}

\paragraph{Execution} {Reservation table \textsf{ALU\_LITE} (Section~\ref{sec:reservation-ALU-LITE})}
~
{\begin{lstlisting}
stage ID:
  new highmult = _ZX_2(highmult);
stage RR:
  new argument3 = GPR[registerYe<<1];
  new argument2 = GPR[registerZe<<1];
stage E1:
  new argument1 = GPR[registerWe<<1];
  for (I = 0; I < 2; I += 1) {
    if ((highmult == 0) || (highmult == 3)) {
      new product = _SX_32(argument2.32[I]) * _SX_32(argument3.32[I]);
    }
    if (highmult == 1) {
      product = _ZX_32(argument2.32[I]) * _ZX_32(argument3.32[I]);
    }
    if (highmult == 2) {
      product = _SX_32(argument2.32[I]) * _ZX_32(argument3.32[I]);
    }
    if (highmult) {
      result1.32[I] = argument1.32[I] + product.32[1];
    } else {
      result1.32[I] = argument1.32[I] + product.32[0];
    }
  };
stage E2:
  GPR[registerWe<<1] = result1;
\end{lstlisting}}
\newpage

\paragraph{Encoding} Format \textsf{ALU\_WPMWRRRO} (Section~\ref{sec:format-ALU-WPMWRRRO}), \verb|alucode2 = 10|\\

\bitfields{ 1/P, 2/11, 1/1, 2/10, 2/highmult, 5/registerWo, 1/1, 2/10, 3/110, 1/0, 5/registerYo, 1/0, 5/registerZo, 1/1 }


\paragraph{Syntax} \texttt{maddwp}\emph{highmult} \emph{registerWo} = \emph{registerZo}, \emph{registerYo}

\paragraph{Description} The \emph{registerYo} and the \emph{registerZo} are interpreted as two words packed into 64 bits. The \emph{registerYo} and the \emph{registerZo} words are sign-extended or zero-extended from 32 bits, depending on the \emph{highmult}. The \emph{registerYo} words are multiplied by the \emph{registerZo} words, and the product is optionally shifted right by 32 bits, depending on the \emph{highmult}. The two resulting words are added to the \emph{registerWo} words, packed and stored back into the \emph{registerWo}.


\paragraph{Modifier(s)}
\noindent\begin{itemize}
\item highmult (cf. highmult in section \ref{par:modifier-highmult} on page \pageref{par:modifier-highmult})
\end{itemize}

\paragraph{Execution} {Reservation table \textsf{ALU\_LITE} (Section~\ref{sec:reservation-ALU-LITE})}
~
{\begin{lstlisting}
stage ID:
  new highmult = _ZX_2(highmult);
stage RR:
  new argument3 = GPR[(registerYo<<1)+1];
  new argument2 = GPR[(registerZo<<1)+1];
stage E1:
  new argument1 = GPR[(registerWo<<1)+1];
  for (I = 0; I < 2; I += 1) {
    if ((highmult == 0) || (highmult == 3)) {
      new product = _SX_32(argument2.32[I]) * _SX_32(argument3.32[I]);
    }
    if (highmult == 1) {
      product = _ZX_32(argument2.32[I]) * _ZX_32(argument3.32[I]);
    }
    if (highmult == 2) {
      product = _SX_32(argument2.32[I]) * _ZX_32(argument3.32[I]);
    }
    if (highmult) {
      result1.32[I] = argument1.32[I] + product.32[1];
    } else {
      result1.32[I] = argument1.32[I] + product.32[0];
    }
  };
stage E2:
  GPR[(registerWo<<1)+1] = result1;
\end{lstlisting}}
\newpage
\subsubsection[{\small MADDXBHO: Multiply Add Extend Byte to Half Word Octuple}]{MADDXBHO\hfill Multiply Add Extend Byte to Half Word Octuple}
\label{instr:MADDXBHO}\index{maddxbho}
 
\paragraph{Encoding} Format \textsf{ALU\_HOMXWRRR} (Section~\ref{sec:format-ALU-HOMXWRRR}), \verb|alucode2 = 10|\\

\bitfields{ 1/P, 2/11, 1/1, 2/10, 2/widemult, 5/registerM, 1/oddlanes, 2/10, 3/111, 1/1, 5/registerO, 1/--, 5/registerP, 1/0 }


\paragraph{Syntax} \texttt{maddxbho}\emph{oddlanes}\emph{widemult} \emph{registerM} = \emph{registerP}, \emph{registerO}

\paragraph{Description} The \emph{registerO} and the \emph{registerP} are interpreted as sixteen bytes packed into 128 bits. The even or odd bytes of the \emph{registerO} and the \emph{registerP} are selected, depending on \emph{oddlanes}. These \emph{registerO} and \emph{registerP} bytes are sign-extended or zero-extended from 8 bits, depending on the \emph{widemult}. The \emph{registerO} bytes are multiplied by the \emph{registerP} bytes, producing eight 16-bit results. The eight resulting half words are packed, added to and stored into the \emph{registerM}.


\paragraph{Modifier(s)}
\noindent\begin{itemize}
\item widemult (cf. widemult in section \ref{par:modifier-widemult} on page \pageref{par:modifier-widemult})
\item oddlanes (cf. oddlanes in section \ref{par:modifier-oddlanes} on page \pageref{par:modifier-oddlanes})
\end{itemize}

\paragraph{Execution} {Reservation table \textsf{ALU\_LITE} (Section~\ref{sec:reservation-ALU-LITE})}
~
{\begin{lstlisting}
stage ID:
  new oddlanes = _ZX_1(oddlanes);
  new widemult = _ZX_2(widemult);
stage RR:
  new argument3 = PGR[registerO];
  new argument2 = PGR[registerP];
stage E1:
  new argument1 = PGR[registerM];
  if (oddlanes) {
    new argument2 = argument2 >> 8;
    new argument3 = argument3 >> 8;
  }
  for (I = 0; I < 8; I += 1) {
    if (widemult == 0) {
      new product = _SX_8(argument2.16[I]) * _SX_8(argument3.16[I]);
    }
    if (widemult == 1) {
      product = _ZX_8(argument2.16[I]) * _ZX_8(argument3.16[I]);
    }
    if (widemult == 2) {
      product = _SX_8(argument2.16[I]) * _ZX_8(argument3.16[I]);
    }
    result1.16[I] = argument1.16[I] + product
  };
stage E2:
  PGR[registerM] = result1;
\end{lstlisting}}
\newpage
\subsubsection[{\small MADDXHWQ: Multiply Add Extend Half Word to Word Quadruple}]{MADDXHWQ\hfill Multiply Add Extend Half Word to Word Quadruple}
\label{instr:MADDXHWQ}\index{maddxhwq}
 
\paragraph{Encoding} Format \textsf{ALU\_WQMXWRRR} (Section~\ref{sec:format-ALU-WQMXWRRR}), \verb|alucode2 = 10|\\

\bitfields{ 1/P, 2/11, 1/1, 2/10, 2/widemult, 5/registerM, 1/oddlanes, 2/10, 3/111, 1/0, 5/registerO, 1/--, 5/registerP, 1/1 }


\paragraph{Syntax} \texttt{maddxhwq}\emph{oddlanes}\emph{widemult} \emph{registerM} = \emph{registerP}, \emph{registerO}

\paragraph{Description} The \emph{registerO} and the \emph{registerP} are interpreted as eight half words packed into 128 bits. The even or odd half words of the \emph{registerO} and the \emph{registerP} are selected, depending on \emph{oddlanes}. These \emph{registerO} and \emph{registerP} half words are sign-extended or zero-extended from 16 bits, depending on the \emph{widemult}. The \emph{registerO} half words are multiplied by the \emph{registerP} half words, producing four 32-bit results. The four resulting words are packed, added to and stored into the \emph{registerM}.


\paragraph{Modifier(s)}
\noindent\begin{itemize}
\item widemult (cf. widemult in section \ref{par:modifier-widemult} on page \pageref{par:modifier-widemult})
\item oddlanes (cf. oddlanes in section \ref{par:modifier-oddlanes} on page \pageref{par:modifier-oddlanes})
\end{itemize}

\paragraph{Execution} {Reservation table \textsf{ALU\_LITE} (Section~\ref{sec:reservation-ALU-LITE})}
~
{\begin{lstlisting}
stage ID:
  new oddlanes = _ZX_1(oddlanes);
  new widemult = _ZX_2(widemult);
stage RR:
  new argument3 = PGR[registerO];
  new argument2 = PGR[registerP];
stage E1:
  new argument1 = PGR[registerM];
  if (oddlanes) {
    new argument2 = argument2 >> 16;
    new argument3 = argument3 >> 16;
  }
  for (I = 0; I < 4; I += 1) {
    if (widemult == 0) {
      new product = _SX_16(argument2.32[I]) * _SX_16(argument3.32[I]);
    }
    if (widemult == 1) {
      product = _ZX_16(argument2.32[I]) * _ZX_16(argument3.32[I]);
    }
    if (widemult == 2) {
      product = _SX_16(argument2.32[I]) * _ZX_16(argument3.32[I]);
    }
    result1.32[I] = argument1.32[I] + product
  };
stage E2:
  PGR[registerM] = result1;
\end{lstlisting}}
\newpage
\subsubsection[{\small MADDXWDP: Multiply Add Extend Word to Double Word Pair}]{MADDXWDP\hfill Multiply Add Extend Word to Double Word Pair}
\label{instr:MADDXWDP}\index{maddxwdp}
 
\paragraph{Encoding} Format \textsf{ALU\_DPMXWRRR} (Section~\ref{sec:format-ALU-DPMXWRRR}), \verb|alucode2 = 10|\\

\bitfields{ 1/P, 2/11, 1/1, 2/10, 2/widemult, 5/registerM, 1/oddlanes, 2/10, 3/111, 1/0, 5/registerO, 1/--, 5/registerP, 1/0 }


\paragraph{Syntax} \texttt{maddxwdp}\emph{oddlanes}\emph{widemult} \emph{registerM} = \emph{registerP}, \emph{registerO}

\paragraph{Description} The \emph{registerO} and the \emph{registerP} are interpreted as four words packed into 128 bits. The even or odd words of the \emph{registerO} and the \emph{registerP} are selected, depending on \emph{oddlanes}. These \emph{registerO} and \emph{registerP} words are sign-extended or zero-extended from 32 bits, depending on the \emph{widemult}. The \emph{registerO} words are multiplied by the \emph{registerP} words, producing two 64-bit results. The two resulting double words are packed, added to and stored into the \emph{registerM}.


\paragraph{Modifier(s)}
\noindent\begin{itemize}
\item widemult (cf. widemult in section \ref{par:modifier-widemult} on page \pageref{par:modifier-widemult})
\item oddlanes (cf. oddlanes in section \ref{par:modifier-oddlanes} on page \pageref{par:modifier-oddlanes})
\end{itemize}

\paragraph{Execution} {Reservation table \textsf{ALU\_LITE} (Section~\ref{sec:reservation-ALU-LITE})}
~
{\begin{lstlisting}
stage ID:
  new oddlanes = _ZX_1(oddlanes);
  new widemult = _ZX_2(widemult);
stage RR:
  new argument3 = PGR[registerO];
  new argument2 = PGR[registerP];
stage E1:
  new argument1 = PGR[registerM];
  if (oddlanes) {
    new argument2 = argument2 >> 32;
    new argument3 = argument3 >> 32;
  }
  for (I = 0; I < 2; I += 1) {
    if (widemult == 0) {
      new product = _SX_32(argument2.64[I]) * _SX_32(argument3.64[I]);
    }
    if (widemult == 1) {
      product = _ZX_32(argument2.64[I]) * _ZX_32(argument3.64[I]);
    }
    if (widemult == 2) {
      product = _SX_32(argument2.64[I]) * _ZX_32(argument3.64[I]);
    }
    result1.64[I] = argument1.64[I] + product
  };
stage E2:
  PGR[registerM] = result1;
\end{lstlisting}}
\newpage
\subsubsection[{\small MAKE: Make Word from Immediate}]{MAKE\hfill Make Word from Immediate}
\label{instr:MAKE}\index{make}
 
\paragraph{Encoding} Format \textsf{ALU\_DWI} (Section~\ref{sec:format-ALU-DWI}), \verb|exucode2 = 0000|\\

\bitfields{ 1/P, 2/11, 1/0, 4/0000, 6/registerW, 2/00, 16/signed16 }


\paragraph{Syntax} \texttt{make} \emph{registerW} = \emph{signed16}

\paragraph{Description} The \emph{signed16} is stored into the \emph{registerW}.


\paragraph{Execution} {Reservation table \textsf{ALU\_TINY} (Section~\ref{sec:reservation-ALU-TINY})}
~
{\begin{lstlisting}
stage RR:
  new argument2 = _SX_16(signed16);
  new result1 = _SX_64(argument2);
stage E1:
  GPR[registerW] = result1;
\end{lstlisting}}
\newpage

\paragraph{Encoding} Format \textsf{ALU\_DWI.X} (Section~\ref{sec:format-ALU-DWI.X}), \verb|exucode2 = 0000|\\

\bitfields{ 1/P, 2/00, 2/-- --, 27/upper27, 1/1, 2/11, 1/0, 4/0000, 6/registerW, 2/00, 10/lower10, 6/extend6 }


\paragraph{Syntax} \texttt{make} \emph{registerW} = \emph{extend6\_\discretionary{}{}{}upper27\_\discretionary{}{}{}lower10}

\paragraph{Description} The \emph{extend6\_\discretionary{}{}{}upper27\_\discretionary{}{}{}lower10} is stored into the \emph{registerW}.


\paragraph{Execution} {Reservation table \textsf{ALU\_TINY.X} (Section~\ref{sec:reservation-ALU-TINY.X})}
~
{\begin{lstlisting}
stage RR:
  new argument2 = _SX_43(extend6_upper27_lower10);
  new result1 = _SX_64(argument2);
stage E1:
  GPR[registerW] = result1;
\end{lstlisting}}
\newpage

\paragraph{Encoding} Format \textsf{ALU\_DWI.Y} (Section~\ref{sec:format-ALU-DWI.Y}), \verb|exucode2 = 0000|\\

\bitfields{ 1/P, 2/00, 2/-- --, 27/extend27, 1/1, 2/00, 2/-- --, 27/upper27, 1/1, 2/11, 1/0, 4/0000, 6/registerW, 2/00, 10/lower10, 6/-- -- -- -- -- -- }


\paragraph{Syntax} \texttt{make} \emph{registerW} = \emph{extend27\_\discretionary{}{}{}upper27\_\discretionary{}{}{}lower10}

\paragraph{Description} The \emph{extend27\_\discretionary{}{}{}upper27\_\discretionary{}{}{}lower10} is stored into the \emph{registerW}.


\paragraph{Execution} {Reservation table \textsf{ALU\_TINY.Y} (Section~\ref{sec:reservation-ALU-TINY.Y})}
~
{\begin{lstlisting}
stage RR:
  new argument2 = _WX_64(extend27_upper27_lower10);
  new result1 = _SX_64(argument2);
stage E1:
  GPR[registerW] = result1;
\end{lstlisting}}
\newpage
\subsubsection[{\small MAXBO: Maximum Byte Octuple}]{MAXBO\hfill Maximum Byte Octuple}
\label{instr:MAXBO}\index{maxbo}
 
\paragraph{Encoding} Format \textsf{ALU\_BOWRR0E} (Section~\ref{sec:format-ALU-BOWRR0E}), \verb|exucode2 = 0101|\\

\bitfields{ 1/P, 2/11, 1/1, 4/0101, 5/registerWe, 1/0, 2/10, 3/000, 1/1, 5/registerYe, 1/--, 5/registerZe, 1/1 }


\paragraph{Syntax} \texttt{maxbo} \emph{registerWe} = \emph{registerZe}, \emph{registerYe}

\paragraph{Description} The \emph{registerYe} and the \emph{registerZe} are interpreted as eight bytes packed into 64 bits. The maximum of the \emph{registerYe} bytes and the \emph{registerZe} bytes are computed. The eight resulting bytes are packed and stored into the \emph{registerWe}.


\paragraph{Execution} {Reservation table \textsf{ALU\_LITE} (Section~\ref{sec:reservation-ALU-LITE})}
~
{\begin{lstlisting}
stage RR:
  new argument3 = GPR[registerYe<<1];
  new argument2 = GPR[registerZe<<1];
  for (I = 0; I < 8; I += 1) {
    result1.8[I] = _MAX(_SX_8(argument3.8[I]), _SX_8(argument2.8[I]))
  };
stage E1:
  GPR[registerWe<<1] = result1;
\end{lstlisting}}
\newpage

\paragraph{Encoding} Format \textsf{ALU\_BOWRR0O} (Section~\ref{sec:format-ALU-BOWRR0O}), \verb|exucode2 = 0101|\\

\bitfields{ 1/P, 2/11, 1/1, 4/0101, 5/registerWo, 1/1, 2/10, 3/000, 1/1, 5/registerYo, 1/--, 5/registerZo, 1/1 }


\paragraph{Syntax} \texttt{maxbo} \emph{registerWo} = \emph{registerZo}, \emph{registerYo}

\paragraph{Description} The \emph{registerYo} and the \emph{registerZo} are interpreted as eight bytes packed into 64 bits. The maximum of the \emph{registerYo} bytes and the \emph{registerZo} bytes are computed. The eight resulting bytes are packed and stored into the \emph{registerWo}.


\paragraph{Execution} {Reservation table \textsf{ALU\_LITE} (Section~\ref{sec:reservation-ALU-LITE})}
~
{\begin{lstlisting}
stage RR:
  new argument3 = GPR[(registerYo<<1)+1];
  new argument2 = GPR[(registerZo<<1)+1];
  for (I = 0; I < 8; I += 1) {
    result1.8[I] = _MAX(_SX_8(argument3.8[I]), _SX_8(argument2.8[I]))
  };
stage E1:
  GPR[(registerWo<<1)+1] = result1;
\end{lstlisting}}
\newpage

\paragraph{Encoding} Format \textsf{ALU\_BOWRR0E.M} (Section~\ref{sec:format-ALU-BOWRR0E.M}), \verb|exucode2 = 0101|\\

\bitfields{ 1/P, 2/00, 2/-- --, 27/upper27, 1/1, 2/11, 1/1, 4/0101, 5/registerWe, 1/0, 2/10, 3/000, 1/1, 1/splat32, 5/lower5, 5/registerZe, 1/1 }


\paragraph{Syntax} \texttt{maxbo} \emph{registerWe} = \emph{registerZe}, \emph{upper27\_\discretionary{}{}{}lower5}\emph{splat32}

\paragraph{Description} The \emph{upper27\_\discretionary{}{}{}lower5} and the \emph{registerZe} are interpreted as eight bytes packed into 64 bits. The maximum of the \emph{upper27\_\discretionary{}{}{}lower5} bytes and the \emph{registerZe} bytes are computed. The eight resulting bytes are packed and stored into the \emph{registerWe}.


\paragraph{Modifier(s)}
\noindent\begin{itemize}
\item splat32 (cf. splat32 in section \ref{par:modifier-splat32} on page \pageref{par:modifier-splat32})
\end{itemize}

\paragraph{Execution} {Reservation table \textsf{ALU\_LITE.X} (Section~\ref{sec:reservation-ALU-LITE.X})}
~
{\begin{lstlisting}
stage ID:
  new splat32 = _ZX_1(splat32);
stage RR:
  new argument3 = splat32 ? (_WX_32(upper27_lower5) << 32) | _ZX_32(_WX_32(upper27_lower5)) : _SX_32(_WX_32(upper27_lower5));
  new argument2 = GPR[registerZe<<1];
  for (I = 0; I < 8; I += 1) {
    result1.8[I] = _MAX(_SX_8(argument3.8[I]), _SX_8(argument2.8[I]))
  };
stage E1:
  GPR[registerWe<<1] = result1;
\end{lstlisting}}
\newpage

\paragraph{Encoding} Format \textsf{ALU\_BOWRR0O.M} (Section~\ref{sec:format-ALU-BOWRR0O.M}), \verb|exucode2 = 0101|\\

\bitfields{ 1/P, 2/00, 2/-- --, 27/upper27, 1/1, 2/11, 1/1, 4/0101, 5/registerWo, 1/1, 2/10, 3/000, 1/1, 1/splat32, 5/lower5, 5/registerZo, 1/1 }


\paragraph{Syntax} \texttt{maxbo} \emph{registerWo} = \emph{registerZo}, \emph{upper27\_\discretionary{}{}{}lower5}\emph{splat32}

\paragraph{Description} The \emph{upper27\_\discretionary{}{}{}lower5} and the \emph{registerZo} are interpreted as eight bytes packed into 64 bits. The maximum of the \emph{upper27\_\discretionary{}{}{}lower5} bytes and the \emph{registerZo} bytes are computed. The eight resulting bytes are packed and stored into the \emph{registerWo}.


\paragraph{Modifier(s)}
\noindent\begin{itemize}
\item splat32 (cf. splat32 in section \ref{par:modifier-splat32} on page \pageref{par:modifier-splat32})
\end{itemize}

\paragraph{Execution} {Reservation table \textsf{ALU\_LITE.X} (Section~\ref{sec:reservation-ALU-LITE.X})}
~
{\begin{lstlisting}
stage ID:
  new splat32 = _ZX_1(splat32);
stage RR:
  new argument3 = splat32 ? (_WX_32(upper27_lower5) << 32) | _ZX_32(_WX_32(upper27_lower5)) : _SX_32(_WX_32(upper27_lower5));
  new argument2 = GPR[(registerZo<<1)+1];
  for (I = 0; I < 8; I += 1) {
    result1.8[I] = _MAX(_SX_8(argument3.8[I]), _SX_8(argument2.8[I]))
  };
stage E1:
  GPR[(registerWo<<1)+1] = result1;
\end{lstlisting}}
\newpage
\subsubsection[{\small MAXD: Maximum Double Word}]{MAXD\hfill Maximum Double Word}
\label{instr:MAXD}\index{maxd}
 
\paragraph{Encoding} Format \textsf{ALU\_DWRR0} (Section~\ref{sec:format-ALU-DWRR0}), \verb|exucode2 = 0101|\\

\bitfields{ 1/P, 2/11, 1/0, 4/0101, 6/registerW, 2/10, 3/000, 1/0, 6/registerY, 6/registerZ }


\paragraph{Syntax} \texttt{maxd} \emph{registerW} = \emph{registerZ}, \emph{registerY}

\paragraph{Description} The maximum of the \emph{registerY} and the \emph{registerZ} is computed. The result is stored into the \emph{registerW}.


\paragraph{Execution} {Reservation table \textsf{ALU\_TINY} (Section~\ref{sec:reservation-ALU-TINY})}
~
{\begin{lstlisting}
stage RR:
  new argument3 = GPR[registerY];
  new argument2 = GPR[registerZ];
  new result1 = _MAX(_SX_64(argument3), _SX_64(argument2));
stage E1:
  GPR[registerW] = result1;
\end{lstlisting}}
\newpage

\paragraph{Encoding} Format \textsf{ALU\_DWRR0.M} (Section~\ref{sec:format-ALU-DWRR0.M}), \verb|exucode2 = 0101|\\

\bitfields{ 1/P, 2/00, 2/-- --, 27/upper27, 1/1, 2/11, 1/0, 4/0101, 6/registerW, 2/10, 3/000, 1/0, 1/splat32, 5/lower5, 6/registerZ }


\paragraph{Syntax} \texttt{maxd} \emph{registerW} = \emph{registerZ}, \emph{upper27\_\discretionary{}{}{}lower5}\emph{splat32}

\paragraph{Description} The maximum of the \emph{upper27\_\discretionary{}{}{}lower5} and the \emph{registerZ} is computed. The result is stored into the \emph{registerW}.


\paragraph{Modifier(s)}
\noindent\begin{itemize}
\item splat32 (cf. splat32 in section \ref{par:modifier-splat32} on page \pageref{par:modifier-splat32})
\end{itemize}

\paragraph{Execution} {Reservation table \textsf{ALU\_TINY.X} (Section~\ref{sec:reservation-ALU-TINY.X})}
~
{\begin{lstlisting}
stage ID:
  new splat32 = _ZX_1(splat32);
stage RR:
  new argument3 = splat32 ? (_WX_32(upper27_lower5) << 32) | _ZX_32(_WX_32(upper27_lower5)) : _SX_32(_WX_32(upper27_lower5));
  new argument2 = GPR[registerZ];
  new result1 = _MAX(_SX_64(argument3), _SX_64(argument2));
stage E1:
  GPR[registerW] = result1;
\end{lstlisting}}
\newpage

\paragraph{Encoding} Format \textsf{ALU\_DWRI} (Section~\ref{sec:format-ALU-DWRI}), \verb|exucode2 = 0101|\\

\bitfields{ 1/P, 2/11, 1/0, 4/0101, 6/registerW, 2/00, 10/signed10, 6/registerZ }


\paragraph{Syntax} \texttt{maxd} \emph{registerW} = \emph{registerZ}, \emph{signed10}

\paragraph{Description} The maximum of the \emph{signed10} and the \emph{registerZ} is computed. The result is stored into the \emph{registerW}.


\paragraph{Execution} {Reservation table \textsf{ALU\_TINY} (Section~\ref{sec:reservation-ALU-TINY})}
~
{\begin{lstlisting}
stage RR:
  new argument3 = _SX_10(signed10);
  new argument2 = GPR[registerZ];
  new result1 = _MAX(_SX_64(argument3), _SX_64(argument2));
stage E1:
  GPR[registerW] = result1;
\end{lstlisting}}
\newpage

\paragraph{Encoding} Format \textsf{ALU\_DWRI.X} (Section~\ref{sec:format-ALU-DWRI.X}), \verb|exucode2 = 0101|\\

\bitfields{ 1/P, 2/00, 2/-- --, 27/upper27, 1/1, 2/11, 1/0, 4/0101, 6/registerW, 2/00, 10/lower10, 6/registerZ }


\paragraph{Syntax} \texttt{maxd} \emph{registerW} = \emph{registerZ}, \emph{upper27\_\discretionary{}{}{}lower10}

\paragraph{Description} The maximum of the \emph{upper27\_\discretionary{}{}{}lower10} and the \emph{registerZ} is computed. The result is stored into the \emph{registerW}.


\paragraph{Execution} {Reservation table \textsf{ALU\_TINY.X} (Section~\ref{sec:reservation-ALU-TINY.X})}
~
{\begin{lstlisting}
stage RR:
  new argument3 = _SX_37(upper27_lower10);
  new argument2 = GPR[registerZ];
  new result1 = _MAX(_SX_64(argument3), _SX_64(argument2));
stage E1:
  GPR[registerW] = result1;
\end{lstlisting}}
\newpage

\paragraph{Encoding} Format \textsf{ALU\_DWRI.Y} (Section~\ref{sec:format-ALU-DWRI.Y}), \verb|exucode2 = 0101|\\

\bitfields{ 1/P, 2/00, 2/-- --, 27/extend27, 1/1, 2/00, 2/-- --, 27/upper27, 1/1, 2/11, 1/0, 4/0101, 6/registerW, 2/00, 10/lower10, 6/registerZ }


\paragraph{Syntax} \texttt{maxd} \emph{registerW} = \emph{registerZ}, \emph{extend27\_\discretionary{}{}{}upper27\_\discretionary{}{}{}lower10}

\paragraph{Description} The maximum of the \emph{extend27\_\discretionary{}{}{}upper27\_\discretionary{}{}{}lower10} and the \emph{registerZ} is computed. The result is stored into the \emph{registerW}.


\paragraph{Execution} {Reservation table \textsf{ALU\_TINY.Y} (Section~\ref{sec:reservation-ALU-TINY.Y})}
~
{\begin{lstlisting}
stage RR:
  new argument3 = _WX_64(extend27_upper27_lower10);
  new argument2 = GPR[registerZ];
  new result1 = _MAX(_SX_64(argument3), _SX_64(argument2));
stage E1:
  GPR[registerW] = result1;
\end{lstlisting}}
\newpage
\subsubsection[{\small MAXHQ: Maximum Half Word Quadruple}]{MAXHQ\hfill Maximum Half Word Quadruple}
\label{instr:MAXHQ}\index{maxhq}
 
\paragraph{Encoding} Format \textsf{ALU\_HQWRR0E} (Section~\ref{sec:format-ALU-HQWRR0E}), \verb|exucode2 = 0101|\\

\bitfields{ 1/P, 2/11, 1/1, 4/0101, 5/registerWe, 1/0, 2/10, 3/000, 1/1, 5/registerYe, 1/--, 5/registerZe, 1/0 }


\paragraph{Syntax} \texttt{maxhq} \emph{registerWe} = \emph{registerZe}, \emph{registerYe}

\paragraph{Description} The \emph{registerYe} and the \emph{registerZe} are interpreted as four half words packed into 64 bits. The maximum of the \emph{registerYe} half words and the \emph{registerZe} half words are computed. The four resulting half words are packed and stored into the \emph{registerWe}.


\paragraph{Execution} {Reservation table \textsf{ALU\_LITE} (Section~\ref{sec:reservation-ALU-LITE})}
~
{\begin{lstlisting}
stage RR:
  new argument3 = GPR[registerYe<<1];
  new argument2 = GPR[registerZe<<1];
  for (I = 0; I < 4; I += 1) {
    result1.16[I] = _MAX(_SX_16(argument3.16[I]), _SX_16(argument2.16[I]))
  };
stage E1:
  GPR[registerWe<<1] = result1;
\end{lstlisting}}
\newpage

\paragraph{Encoding} Format \textsf{ALU\_HQWRR0O} (Section~\ref{sec:format-ALU-HQWRR0O}), \verb|exucode2 = 0101|\\

\bitfields{ 1/P, 2/11, 1/1, 4/0101, 5/registerWo, 1/1, 2/10, 3/000, 1/1, 5/registerYo, 1/--, 5/registerZo, 1/0 }


\paragraph{Syntax} \texttt{maxhq} \emph{registerWo} = \emph{registerZo}, \emph{registerYo}

\paragraph{Description} The \emph{registerYo} and the \emph{registerZo} are interpreted as four half words packed into 64 bits. The maximum of the \emph{registerYo} half words and the \emph{registerZo} half words are computed. The four resulting half words are packed and stored into the \emph{registerWo}.


\paragraph{Execution} {Reservation table \textsf{ALU\_LITE} (Section~\ref{sec:reservation-ALU-LITE})}
~
{\begin{lstlisting}
stage RR:
  new argument3 = GPR[(registerYo<<1)+1];
  new argument2 = GPR[(registerZo<<1)+1];
  for (I = 0; I < 4; I += 1) {
    result1.16[I] = _MAX(_SX_16(argument3.16[I]), _SX_16(argument2.16[I]))
  };
stage E1:
  GPR[(registerWo<<1)+1] = result1;
\end{lstlisting}}
\newpage

\paragraph{Encoding} Format \textsf{ALU\_HQWRR0E.M} (Section~\ref{sec:format-ALU-HQWRR0E.M}), \verb|exucode2 = 0101|\\

\bitfields{ 1/P, 2/00, 2/-- --, 27/upper27, 1/1, 2/11, 1/1, 4/0101, 5/registerWe, 1/0, 2/10, 3/000, 1/1, 1/splat32, 5/lower5, 5/registerZe, 1/0 }


\paragraph{Syntax} \texttt{maxhq} \emph{registerWe} = \emph{registerZe}, \emph{upper27\_\discretionary{}{}{}lower5}\emph{splat32}

\paragraph{Description} The \emph{upper27\_\discretionary{}{}{}lower5} and the \emph{registerZe} are interpreted as four half words packed into 64 bits. The maximum of the \emph{upper27\_\discretionary{}{}{}lower5} half words and the \emph{registerZe} half words are computed. The four resulting half words are packed and stored into the \emph{registerWe}.


\paragraph{Modifier(s)}
\noindent\begin{itemize}
\item splat32 (cf. splat32 in section \ref{par:modifier-splat32} on page \pageref{par:modifier-splat32})
\end{itemize}

\paragraph{Execution} {Reservation table \textsf{ALU\_LITE.X} (Section~\ref{sec:reservation-ALU-LITE.X})}
~
{\begin{lstlisting}
stage ID:
  new splat32 = _ZX_1(splat32);
stage RR:
  new argument3 = splat32 ? (_WX_32(upper27_lower5) << 32) | _ZX_32(_WX_32(upper27_lower5)) : _SX_32(_WX_32(upper27_lower5));
  new argument2 = GPR[registerZe<<1];
  for (I = 0; I < 4; I += 1) {
    result1.16[I] = _MAX(_SX_16(argument3.16[I]), _SX_16(argument2.16[I]))
  };
stage E1:
  GPR[registerWe<<1] = result1;
\end{lstlisting}}
\newpage

\paragraph{Encoding} Format \textsf{ALU\_HQWRR0O.M} (Section~\ref{sec:format-ALU-HQWRR0O.M}), \verb|exucode2 = 0101|\\

\bitfields{ 1/P, 2/00, 2/-- --, 27/upper27, 1/1, 2/11, 1/1, 4/0101, 5/registerWo, 1/1, 2/10, 3/000, 1/1, 1/splat32, 5/lower5, 5/registerZo, 1/0 }


\paragraph{Syntax} \texttt{maxhq} \emph{registerWo} = \emph{registerZo}, \emph{upper27\_\discretionary{}{}{}lower5}\emph{splat32}

\paragraph{Description} The \emph{upper27\_\discretionary{}{}{}lower5} and the \emph{registerZo} are interpreted as four half words packed into 64 bits. The maximum of the \emph{upper27\_\discretionary{}{}{}lower5} half words and the \emph{registerZo} half words are computed. The four resulting half words are packed and stored into the \emph{registerWo}.


\paragraph{Modifier(s)}
\noindent\begin{itemize}
\item splat32 (cf. splat32 in section \ref{par:modifier-splat32} on page \pageref{par:modifier-splat32})
\end{itemize}

\paragraph{Execution} {Reservation table \textsf{ALU\_LITE.X} (Section~\ref{sec:reservation-ALU-LITE.X})}
~
{\begin{lstlisting}
stage ID:
  new splat32 = _ZX_1(splat32);
stage RR:
  new argument3 = splat32 ? (_WX_32(upper27_lower5) << 32) | _ZX_32(_WX_32(upper27_lower5)) : _SX_32(_WX_32(upper27_lower5));
  new argument2 = GPR[(registerZo<<1)+1];
  for (I = 0; I < 4; I += 1) {
    result1.16[I] = _MAX(_SX_16(argument3.16[I]), _SX_16(argument2.16[I]))
  };
stage E1:
  GPR[(registerWo<<1)+1] = result1;
\end{lstlisting}}
\newpage
\subsubsection[{\small MAXUBO: Maximum Unsigned Byte Octuple}]{MAXUBO\hfill Maximum Unsigned Byte Octuple}
\label{instr:MAXUBO}\index{maxubo}
 
\paragraph{Encoding} Format \textsf{ALU\_BOWRR0E} (Section~\ref{sec:format-ALU-BOWRR0E}), \verb|exucode2 = 0111|\\

\bitfields{ 1/P, 2/11, 1/1, 4/0111, 5/registerWe, 1/0, 2/10, 3/000, 1/1, 5/registerYe, 1/--, 5/registerZe, 1/1 }


\paragraph{Syntax} \texttt{maxubo} \emph{registerWe} = \emph{registerZe}, \emph{registerYe}

\paragraph{Description} The \emph{registerYe} and the \emph{registerZe} are interpreted as eight bytes packed into 64 bits. The unsigned maximum of the \emph{registerYe} bytes and the \emph{registerZe} bytes are computed. The eight resulting bytes are packed and stored into the \emph{registerWe}.


\paragraph{Execution} {Reservation table \textsf{ALU\_LITE} (Section~\ref{sec:reservation-ALU-LITE})}
~
{\begin{lstlisting}
stage RR:
  new argument3 = GPR[registerYe<<1];
  new argument2 = GPR[registerZe<<1];
  for (I = 0; I < 8; I += 1) {
    result1.8[I] = _MAX(_ZX_8(argument3.8[I]), _ZX_8(argument2.8[I]))
  };
stage E1:
  GPR[registerWe<<1] = result1;
\end{lstlisting}}
\newpage

\paragraph{Encoding} Format \textsf{ALU\_BOWRR0O} (Section~\ref{sec:format-ALU-BOWRR0O}), \verb|exucode2 = 0111|\\

\bitfields{ 1/P, 2/11, 1/1, 4/0111, 5/registerWo, 1/1, 2/10, 3/000, 1/1, 5/registerYo, 1/--, 5/registerZo, 1/1 }


\paragraph{Syntax} \texttt{maxubo} \emph{registerWo} = \emph{registerZo}, \emph{registerYo}

\paragraph{Description} The \emph{registerYo} and the \emph{registerZo} are interpreted as eight bytes packed into 64 bits. The unsigned maximum of the \emph{registerYo} bytes and the \emph{registerZo} bytes are computed. The eight resulting bytes are packed and stored into the \emph{registerWo}.


\paragraph{Execution} {Reservation table \textsf{ALU\_LITE} (Section~\ref{sec:reservation-ALU-LITE})}
~
{\begin{lstlisting}
stage RR:
  new argument3 = GPR[(registerYo<<1)+1];
  new argument2 = GPR[(registerZo<<1)+1];
  for (I = 0; I < 8; I += 1) {
    result1.8[I] = _MAX(_ZX_8(argument3.8[I]), _ZX_8(argument2.8[I]))
  };
stage E1:
  GPR[(registerWo<<1)+1] = result1;
\end{lstlisting}}
\newpage

\paragraph{Encoding} Format \textsf{ALU\_BOWRR0E.M} (Section~\ref{sec:format-ALU-BOWRR0E.M}), \verb|exucode2 = 0111|\\

\bitfields{ 1/P, 2/00, 2/-- --, 27/upper27, 1/1, 2/11, 1/1, 4/0111, 5/registerWe, 1/0, 2/10, 3/000, 1/1, 1/splat32, 5/lower5, 5/registerZe, 1/1 }


\paragraph{Syntax} \texttt{maxubo} \emph{registerWe} = \emph{registerZe}, \emph{upper27\_\discretionary{}{}{}lower5}\emph{splat32}

\paragraph{Description} The \emph{upper27\_\discretionary{}{}{}lower5} and the \emph{registerZe} are interpreted as eight bytes packed into 64 bits. The unsigned maximum of the \emph{upper27\_\discretionary{}{}{}lower5} bytes and the \emph{registerZe} bytes are computed. The eight resulting bytes are packed and stored into the \emph{registerWe}.


\paragraph{Modifier(s)}
\noindent\begin{itemize}
\item splat32 (cf. splat32 in section \ref{par:modifier-splat32} on page \pageref{par:modifier-splat32})
\end{itemize}

\paragraph{Execution} {Reservation table \textsf{ALU\_LITE.X} (Section~\ref{sec:reservation-ALU-LITE.X})}
~
{\begin{lstlisting}
stage ID:
  new splat32 = _ZX_1(splat32);
stage RR:
  new argument3 = splat32 ? (_WX_32(upper27_lower5) << 32) | _ZX_32(_WX_32(upper27_lower5)) : _SX_32(_WX_32(upper27_lower5));
  new argument2 = GPR[registerZe<<1];
  for (I = 0; I < 8; I += 1) {
    result1.8[I] = _MAX(_ZX_8(argument3.8[I]), _ZX_8(argument2.8[I]))
  };
stage E1:
  GPR[registerWe<<1] = result1;
\end{lstlisting}}
\newpage

\paragraph{Encoding} Format \textsf{ALU\_BOWRR0O.M} (Section~\ref{sec:format-ALU-BOWRR0O.M}), \verb|exucode2 = 0111|\\

\bitfields{ 1/P, 2/00, 2/-- --, 27/upper27, 1/1, 2/11, 1/1, 4/0111, 5/registerWo, 1/1, 2/10, 3/000, 1/1, 1/splat32, 5/lower5, 5/registerZo, 1/1 }


\paragraph{Syntax} \texttt{maxubo} \emph{registerWo} = \emph{registerZo}, \emph{upper27\_\discretionary{}{}{}lower5}\emph{splat32}

\paragraph{Description} The \emph{upper27\_\discretionary{}{}{}lower5} and the \emph{registerZo} are interpreted as eight bytes packed into 64 bits. The unsigned maximum of the \emph{upper27\_\discretionary{}{}{}lower5} bytes and the \emph{registerZo} bytes are computed. The eight resulting bytes are packed and stored into the \emph{registerWo}.


\paragraph{Modifier(s)}
\noindent\begin{itemize}
\item splat32 (cf. splat32 in section \ref{par:modifier-splat32} on page \pageref{par:modifier-splat32})
\end{itemize}

\paragraph{Execution} {Reservation table \textsf{ALU\_LITE.X} (Section~\ref{sec:reservation-ALU-LITE.X})}
~
{\begin{lstlisting}
stage ID:
  new splat32 = _ZX_1(splat32);
stage RR:
  new argument3 = splat32 ? (_WX_32(upper27_lower5) << 32) | _ZX_32(_WX_32(upper27_lower5)) : _SX_32(_WX_32(upper27_lower5));
  new argument2 = GPR[(registerZo<<1)+1];
  for (I = 0; I < 8; I += 1) {
    result1.8[I] = _MAX(_ZX_8(argument3.8[I]), _ZX_8(argument2.8[I]))
  };
stage E1:
  GPR[(registerWo<<1)+1] = result1;
\end{lstlisting}}
\newpage
\subsubsection[{\small MAXUD: Maximum Unsigned Double Word}]{MAXUD\hfill Maximum Unsigned Double Word}
\label{instr:MAXUD}\index{maxud}
 
\paragraph{Encoding} Format \textsf{ALU\_DWRR0} (Section~\ref{sec:format-ALU-DWRR0}), \verb|exucode2 = 0111|\\

\bitfields{ 1/P, 2/11, 1/0, 4/0111, 6/registerW, 2/10, 3/000, 1/0, 6/registerY, 6/registerZ }


\paragraph{Syntax} \texttt{maxud} \emph{registerW} = \emph{registerZ}, \emph{registerY}

\paragraph{Description} The unsigned maximum of the \emph{registerY} and the \emph{registerZ} is computed. The result is stored into the \emph{registerW}.


\paragraph{Execution} {Reservation table \textsf{ALU\_TINY} (Section~\ref{sec:reservation-ALU-TINY})}
~
{\begin{lstlisting}
stage RR:
  new argument3 = GPR[registerY];
  new argument2 = GPR[registerZ];
  new result1 = _MAX(_ZX_64(argument3), _ZX_64(argument2));
stage E1:
  GPR[registerW] = result1;
\end{lstlisting}}
\newpage

\paragraph{Encoding} Format \textsf{ALU\_DWRR0.M} (Section~\ref{sec:format-ALU-DWRR0.M}), \verb|exucode2 = 0111|\\

\bitfields{ 1/P, 2/00, 2/-- --, 27/upper27, 1/1, 2/11, 1/0, 4/0111, 6/registerW, 2/10, 3/000, 1/0, 1/splat32, 5/lower5, 6/registerZ }


\paragraph{Syntax} \texttt{maxud} \emph{registerW} = \emph{registerZ}, \emph{upper27\_\discretionary{}{}{}lower5}\emph{splat32}

\paragraph{Description} The unsigned maximum of the \emph{upper27\_\discretionary{}{}{}lower5} and the \emph{registerZ} is computed. The result is stored into the \emph{registerW}.


\paragraph{Modifier(s)}
\noindent\begin{itemize}
\item splat32 (cf. splat32 in section \ref{par:modifier-splat32} on page \pageref{par:modifier-splat32})
\end{itemize}

\paragraph{Execution} {Reservation table \textsf{ALU\_TINY.X} (Section~\ref{sec:reservation-ALU-TINY.X})}
~
{\begin{lstlisting}
stage ID:
  new splat32 = _ZX_1(splat32);
stage RR:
  new argument3 = splat32 ? (_WX_32(upper27_lower5) << 32) | _ZX_32(_WX_32(upper27_lower5)) : _SX_32(_WX_32(upper27_lower5));
  new argument2 = GPR[registerZ];
  new result1 = _MAX(_ZX_64(argument3), _ZX_64(argument2));
stage E1:
  GPR[registerW] = result1;
\end{lstlisting}}
\newpage

\paragraph{Encoding} Format \textsf{ALU\_DWRI} (Section~\ref{sec:format-ALU-DWRI}), \verb|exucode2 = 0111|\\

\bitfields{ 1/P, 2/11, 1/0, 4/0111, 6/registerW, 2/00, 10/signed10, 6/registerZ }


\paragraph{Syntax} \texttt{maxud} \emph{registerW} = \emph{registerZ}, \emph{signed10}

\paragraph{Description} The unsigned maximum of the \emph{signed10} and the \emph{registerZ} is computed. The result is stored into the \emph{registerW}.


\paragraph{Execution} {Reservation table \textsf{ALU\_TINY} (Section~\ref{sec:reservation-ALU-TINY})}
~
{\begin{lstlisting}
stage RR:
  new argument3 = _SX_10(signed10);
  new argument2 = GPR[registerZ];
  new result1 = _MAX(_ZX_64(argument3), _ZX_64(argument2));
stage E1:
  GPR[registerW] = result1;
\end{lstlisting}}
\newpage

\paragraph{Encoding} Format \textsf{ALU\_DWRI.X} (Section~\ref{sec:format-ALU-DWRI.X}), \verb|exucode2 = 0111|\\

\bitfields{ 1/P, 2/00, 2/-- --, 27/upper27, 1/1, 2/11, 1/0, 4/0111, 6/registerW, 2/00, 10/lower10, 6/registerZ }


\paragraph{Syntax} \texttt{maxud} \emph{registerW} = \emph{registerZ}, \emph{upper27\_\discretionary{}{}{}lower10}

\paragraph{Description} The unsigned maximum of the \emph{upper27\_\discretionary{}{}{}lower10} and the \emph{registerZ} is computed. The result is stored into the \emph{registerW}.


\paragraph{Execution} {Reservation table \textsf{ALU\_TINY.X} (Section~\ref{sec:reservation-ALU-TINY.X})}
~
{\begin{lstlisting}
stage RR:
  new argument3 = _SX_37(upper27_lower10);
  new argument2 = GPR[registerZ];
  new result1 = _MAX(_ZX_64(argument3), _ZX_64(argument2));
stage E1:
  GPR[registerW] = result1;
\end{lstlisting}}
\newpage

\paragraph{Encoding} Format \textsf{ALU\_DWRI.Y} (Section~\ref{sec:format-ALU-DWRI.Y}), \verb|exucode2 = 0111|\\

\bitfields{ 1/P, 2/00, 2/-- --, 27/extend27, 1/1, 2/00, 2/-- --, 27/upper27, 1/1, 2/11, 1/0, 4/0111, 6/registerW, 2/00, 10/lower10, 6/registerZ }


\paragraph{Syntax} \texttt{maxud} \emph{registerW} = \emph{registerZ}, \emph{extend27\_\discretionary{}{}{}upper27\_\discretionary{}{}{}lower10}

\paragraph{Description} The unsigned maximum of the \emph{extend27\_\discretionary{}{}{}upper27\_\discretionary{}{}{}lower10} and the \emph{registerZ} is computed. The result is stored into the \emph{registerW}.


\paragraph{Execution} {Reservation table \textsf{ALU\_TINY.Y} (Section~\ref{sec:reservation-ALU-TINY.Y})}
~
{\begin{lstlisting}
stage RR:
  new argument3 = _WX_64(extend27_upper27_lower10);
  new argument2 = GPR[registerZ];
  new result1 = _MAX(_ZX_64(argument3), _ZX_64(argument2));
stage E1:
  GPR[registerW] = result1;
\end{lstlisting}}
\newpage
\subsubsection[{\small MAXUHQ: Maximum Unsigned Half Word Quadruple}]{MAXUHQ\hfill Maximum Unsigned Half Word Quadruple}
\label{instr:MAXUHQ}\index{maxuhq}
 
\paragraph{Encoding} Format \textsf{ALU\_HQWRR0E} (Section~\ref{sec:format-ALU-HQWRR0E}), \verb|exucode2 = 0111|\\

\bitfields{ 1/P, 2/11, 1/1, 4/0111, 5/registerWe, 1/0, 2/10, 3/000, 1/1, 5/registerYe, 1/--, 5/registerZe, 1/0 }


\paragraph{Syntax} \texttt{maxuhq} \emph{registerWe} = \emph{registerZe}, \emph{registerYe}

\paragraph{Description} The \emph{registerYe} and the \emph{registerZe} are interpreted as four half words packed into 64 bits. The unsigned maximum of the \emph{registerYe} half words and the \emph{registerZe} half words are computed. The four resulting half words are packed and stored into the \emph{registerWe}.


\paragraph{Execution} {Reservation table \textsf{ALU\_LITE} (Section~\ref{sec:reservation-ALU-LITE})}
~
{\begin{lstlisting}
stage RR:
  new argument3 = GPR[registerYe<<1];
  new argument2 = GPR[registerZe<<1];
  for (I = 0; I < 4; I += 1) {
    result1.16[I] = _MAX(_ZX_16(argument3.16[I]), _ZX_16(argument2.16[I]))
  };
stage E1:
  GPR[registerWe<<1] = result1;
\end{lstlisting}}
\newpage

\paragraph{Encoding} Format \textsf{ALU\_HQWRR0O} (Section~\ref{sec:format-ALU-HQWRR0O}), \verb|exucode2 = 0111|\\

\bitfields{ 1/P, 2/11, 1/1, 4/0111, 5/registerWo, 1/1, 2/10, 3/000, 1/1, 5/registerYo, 1/--, 5/registerZo, 1/0 }


\paragraph{Syntax} \texttt{maxuhq} \emph{registerWo} = \emph{registerZo}, \emph{registerYo}

\paragraph{Description} The \emph{registerYo} and the \emph{registerZo} are interpreted as four half words packed into 64 bits. The unsigned maximum of the \emph{registerYo} half words and the \emph{registerZo} half words are computed. The four resulting half words are packed and stored into the \emph{registerWo}.


\paragraph{Execution} {Reservation table \textsf{ALU\_LITE} (Section~\ref{sec:reservation-ALU-LITE})}
~
{\begin{lstlisting}
stage RR:
  new argument3 = GPR[(registerYo<<1)+1];
  new argument2 = GPR[(registerZo<<1)+1];
  for (I = 0; I < 4; I += 1) {
    result1.16[I] = _MAX(_ZX_16(argument3.16[I]), _ZX_16(argument2.16[I]))
  };
stage E1:
  GPR[(registerWo<<1)+1] = result1;
\end{lstlisting}}
\newpage

\paragraph{Encoding} Format \textsf{ALU\_HQWRR0E.M} (Section~\ref{sec:format-ALU-HQWRR0E.M}), \verb|exucode2 = 0111|\\

\bitfields{ 1/P, 2/00, 2/-- --, 27/upper27, 1/1, 2/11, 1/1, 4/0111, 5/registerWe, 1/0, 2/10, 3/000, 1/1, 1/splat32, 5/lower5, 5/registerZe, 1/0 }


\paragraph{Syntax} \texttt{maxuhq} \emph{registerWe} = \emph{registerZe}, \emph{upper27\_\discretionary{}{}{}lower5}\emph{splat32}

\paragraph{Description} The \emph{upper27\_\discretionary{}{}{}lower5} and the \emph{registerZe} are interpreted as four half words packed into 64 bits. The unsigned maximum of the \emph{upper27\_\discretionary{}{}{}lower5} half words and the \emph{registerZe} half words are computed. The four resulting half words are packed and stored into the \emph{registerWe}.


\paragraph{Modifier(s)}
\noindent\begin{itemize}
\item splat32 (cf. splat32 in section \ref{par:modifier-splat32} on page \pageref{par:modifier-splat32})
\end{itemize}

\paragraph{Execution} {Reservation table \textsf{ALU\_LITE.X} (Section~\ref{sec:reservation-ALU-LITE.X})}
~
{\begin{lstlisting}
stage ID:
  new splat32 = _ZX_1(splat32);
stage RR:
  new argument3 = splat32 ? (_WX_32(upper27_lower5) << 32) | _ZX_32(_WX_32(upper27_lower5)) : _SX_32(_WX_32(upper27_lower5));
  new argument2 = GPR[registerZe<<1];
  for (I = 0; I < 4; I += 1) {
    result1.16[I] = _MAX(_ZX_16(argument3.16[I]), _ZX_16(argument2.16[I]))
  };
stage E1:
  GPR[registerWe<<1] = result1;
\end{lstlisting}}
\newpage

\paragraph{Encoding} Format \textsf{ALU\_HQWRR0O.M} (Section~\ref{sec:format-ALU-HQWRR0O.M}), \verb|exucode2 = 0111|\\

\bitfields{ 1/P, 2/00, 2/-- --, 27/upper27, 1/1, 2/11, 1/1, 4/0111, 5/registerWo, 1/1, 2/10, 3/000, 1/1, 1/splat32, 5/lower5, 5/registerZo, 1/0 }


\paragraph{Syntax} \texttt{maxuhq} \emph{registerWo} = \emph{registerZo}, \emph{upper27\_\discretionary{}{}{}lower5}\emph{splat32}

\paragraph{Description} The \emph{upper27\_\discretionary{}{}{}lower5} and the \emph{registerZo} are interpreted as four half words packed into 64 bits. The unsigned maximum of the \emph{upper27\_\discretionary{}{}{}lower5} half words and the \emph{registerZo} half words are computed. The four resulting half words are packed and stored into the \emph{registerWo}.


\paragraph{Modifier(s)}
\noindent\begin{itemize}
\item splat32 (cf. splat32 in section \ref{par:modifier-splat32} on page \pageref{par:modifier-splat32})
\end{itemize}

\paragraph{Execution} {Reservation table \textsf{ALU\_LITE.X} (Section~\ref{sec:reservation-ALU-LITE.X})}
~
{\begin{lstlisting}
stage ID:
  new splat32 = _ZX_1(splat32);
stage RR:
  new argument3 = splat32 ? (_WX_32(upper27_lower5) << 32) | _ZX_32(_WX_32(upper27_lower5)) : _SX_32(_WX_32(upper27_lower5));
  new argument2 = GPR[(registerZo<<1)+1];
  for (I = 0; I < 4; I += 1) {
    result1.16[I] = _MAX(_ZX_16(argument3.16[I]), _ZX_16(argument2.16[I]))
  };
stage E1:
  GPR[(registerWo<<1)+1] = result1;
\end{lstlisting}}
\newpage
\subsubsection[{\small MAXUW: Maximum Unsigned Word}]{MAXUW\hfill Maximum Unsigned Word}
\label{instr:MAXUW}\index{maxuw}
 
\paragraph{Encoding} Format \textsf{ALU\_WWRR0} (Section~\ref{sec:format-ALU-WWRR0}), \verb|exucode2 = 0111|\\

\bitfields{ 1/P, 2/11, 1/0, 4/0111, 6/registerW, 2/11, 3/000, 1/signextw, 6/registerY, 6/registerZ }


\paragraph{Syntax} \texttt{maxuw}\emph{signextw} \emph{registerW} = \emph{registerZ}, \emph{registerY}

\paragraph{Description} The unsigned maximum of the \emph{registerY} and the \emph{registerZ} is computed. The result with the 32 upper bits cleared is stored into the \emph{registerW}.


\paragraph{Modifier(s)}
\noindent\begin{itemize}
\item signextw (cf. signextw in section \ref{par:modifier-signextw} on page \pageref{par:modifier-signextw})
\end{itemize}

\paragraph{Execution} {Reservation table \textsf{ALU\_TINY} (Section~\ref{sec:reservation-ALU-TINY})}
~
{\begin{lstlisting}
stage ID:
  new signextw = _ZX_1(signextw);
stage RR:
  new argument3 = GPR[registerY];
  new argument2 = GPR[registerZ];
  new result1 = _MAX(argument3.32[0], argument2.32[0]);
stage E1:
  GPR[registerW] = signextw ? _SX_32(result1) : _ZX_32(result1);
\end{lstlisting}}
\newpage

\paragraph{Encoding} Format \textsf{ALU\_WWRR0.W} (Section~\ref{sec:format-ALU-WWRR0.W}), \verb|exucode2 = 0111|\\

\bitfields{ 1/P, 2/00, 2/-- --, 27/upper27, 1/1, 2/11, 1/0, 4/0111, 6/registerW, 2/11, 3/000, 1/signextw, 1/0, 5/lower5, 6/registerZ }


\paragraph{Syntax} \texttt{maxuw}\emph{signextw} \emph{registerW} = \emph{registerZ}, \emph{upper27\_\discretionary{}{}{}lower5}

\paragraph{Description} The unsigned maximum of the \emph{upper27\_\discretionary{}{}{}lower5} and the \emph{registerZ} is computed. The result with the 32 upper bits cleared is stored into the \emph{registerW}.


\paragraph{Modifier(s)}
\noindent\begin{itemize}
\item signextw (cf. signextw in section \ref{par:modifier-signextw} on page \pageref{par:modifier-signextw})
\end{itemize}

\paragraph{Execution} {Reservation table \textsf{ALU\_TINY.X} (Section~\ref{sec:reservation-ALU-TINY.X})}
~
{\begin{lstlisting}
stage ID:
  new signextw = _ZX_1(signextw);
stage RR:
  new argument3 = _WX_32(upper27_lower5);
  new argument2 = GPR[registerZ];
  new result1 = _MAX(argument3.32[0], argument2.32[0]);
stage E1:
  GPR[registerW] = signextw ? _SX_32(result1) : _ZX_32(result1);
\end{lstlisting}}
\newpage
\subsubsection[{\small MAXUWP: Maximum Unsigned Word Pair}]{MAXUWP\hfill Maximum Unsigned Word Pair}
\label{instr:MAXUWP}\index{maxuwp}
 
\paragraph{Encoding} Format \textsf{ALU\_WPWRR0E} (Section~\ref{sec:format-ALU-WPWRR0E}), \verb|exucode2 = 0111|\\

\bitfields{ 1/P, 2/11, 1/1, 4/0111, 5/registerWe, 1/0, 2/10, 3/000, 1/0, 5/registerYe, 1/--, 5/registerZe, 1/1 }


\paragraph{Syntax} \texttt{maxuwp} \emph{registerWe} = \emph{registerZe}, \emph{registerYe}

\paragraph{Description} The \emph{registerYe} and the \emph{registerZe} are interpreted as two words packed into 64 bits. The unsigned maximum of the \emph{registerYe} words and the \emph{registerZe} words are computed. The two resulting words are packed and stored into the \emph{registerWe}.


\paragraph{Execution} {Reservation table \textsf{ALU\_LITE} (Section~\ref{sec:reservation-ALU-LITE})}
~
{\begin{lstlisting}
stage RR:
  new argument3 = GPR[registerYe<<1];
  new argument2 = GPR[registerZe<<1];
  for (I = 0; I < 2; I += 1) {
    result1.32[I] = _MAX(_ZX_32(argument3.32[I]), _ZX_32(argument2.32[I]))
  };
stage E1:
  GPR[registerWe<<1] = result1;
\end{lstlisting}}
\newpage

\paragraph{Encoding} Format \textsf{ALU\_WPWRR0O} (Section~\ref{sec:format-ALU-WPWRR0O}), \verb|exucode2 = 0111|\\

\bitfields{ 1/P, 2/11, 1/1, 4/0111, 5/registerWo, 1/1, 2/10, 3/000, 1/0, 5/registerYo, 1/--, 5/registerZo, 1/1 }


\paragraph{Syntax} \texttt{maxuwp} \emph{registerWo} = \emph{registerZo}, \emph{registerYo}

\paragraph{Description} The \emph{registerYo} and the \emph{registerZo} are interpreted as two words packed into 64 bits. The unsigned maximum of the \emph{registerYo} words and the \emph{registerZo} words are computed. The two resulting words are packed and stored into the \emph{registerWo}.


\paragraph{Execution} {Reservation table \textsf{ALU\_LITE} (Section~\ref{sec:reservation-ALU-LITE})}
~
{\begin{lstlisting}
stage RR:
  new argument3 = GPR[(registerYo<<1)+1];
  new argument2 = GPR[(registerZo<<1)+1];
  for (I = 0; I < 2; I += 1) {
    result1.32[I] = _MAX(_ZX_32(argument3.32[I]), _ZX_32(argument2.32[I]))
  };
stage E1:
  GPR[(registerWo<<1)+1] = result1;
\end{lstlisting}}
\newpage

\paragraph{Encoding} Format \textsf{ALU\_WPWRR0E.M} (Section~\ref{sec:format-ALU-WPWRR0E.M}), \verb|exucode2 = 0111|\\

\bitfields{ 1/P, 2/00, 2/-- --, 27/upper27, 1/1, 2/11, 1/1, 4/0111, 5/registerWe, 1/0, 2/10, 3/000, 1/0, 1/splat32, 5/lower5, 5/registerZe, 1/1 }


\paragraph{Syntax} \texttt{maxuwp} \emph{registerWe} = \emph{registerZe}, \emph{upper27\_\discretionary{}{}{}lower5}\emph{splat32}

\paragraph{Description} The \emph{upper27\_\discretionary{}{}{}lower5} and the \emph{registerZe} are interpreted as two words packed into 64 bits. The unsigned maximum of the \emph{upper27\_\discretionary{}{}{}lower5} words and the \emph{registerZe} words are computed. The two resulting words are packed and stored into the \emph{registerWe}.


\paragraph{Modifier(s)}
\noindent\begin{itemize}
\item splat32 (cf. splat32 in section \ref{par:modifier-splat32} on page \pageref{par:modifier-splat32})
\end{itemize}

\paragraph{Execution} {Reservation table \textsf{ALU\_LITE.X} (Section~\ref{sec:reservation-ALU-LITE.X})}
~
{\begin{lstlisting}
stage ID:
  new splat32 = _ZX_1(splat32);
stage RR:
  new argument3 = splat32 ? (_WX_32(upper27_lower5) << 32) | _ZX_32(_WX_32(upper27_lower5)) : _SX_32(_WX_32(upper27_lower5));
  new argument2 = GPR[registerZe<<1];
  for (I = 0; I < 2; I += 1) {
    result1.32[I] = _MAX(_ZX_32(argument3.32[I]), _ZX_32(argument2.32[I]))
  };
stage E1:
  GPR[registerWe<<1] = result1;
\end{lstlisting}}
\newpage

\paragraph{Encoding} Format \textsf{ALU\_WPWRR0O.M} (Section~\ref{sec:format-ALU-WPWRR0O.M}), \verb|exucode2 = 0111|\\

\bitfields{ 1/P, 2/00, 2/-- --, 27/upper27, 1/1, 2/11, 1/1, 4/0111, 5/registerWo, 1/1, 2/10, 3/000, 1/0, 1/splat32, 5/lower5, 5/registerZo, 1/1 }


\paragraph{Syntax} \texttt{maxuwp} \emph{registerWo} = \emph{registerZo}, \emph{upper27\_\discretionary{}{}{}lower5}\emph{splat32}

\paragraph{Description} The \emph{upper27\_\discretionary{}{}{}lower5} and the \emph{registerZo} are interpreted as two words packed into 64 bits. The unsigned maximum of the \emph{upper27\_\discretionary{}{}{}lower5} words and the \emph{registerZo} words are computed. The two resulting words are packed and stored into the \emph{registerWo}.


\paragraph{Modifier(s)}
\noindent\begin{itemize}
\item splat32 (cf. splat32 in section \ref{par:modifier-splat32} on page \pageref{par:modifier-splat32})
\end{itemize}

\paragraph{Execution} {Reservation table \textsf{ALU\_LITE.X} (Section~\ref{sec:reservation-ALU-LITE.X})}
~
{\begin{lstlisting}
stage ID:
  new splat32 = _ZX_1(splat32);
stage RR:
  new argument3 = splat32 ? (_WX_32(upper27_lower5) << 32) | _ZX_32(_WX_32(upper27_lower5)) : _SX_32(_WX_32(upper27_lower5));
  new argument2 = GPR[(registerZo<<1)+1];
  for (I = 0; I < 2; I += 1) {
    result1.32[I] = _MAX(_ZX_32(argument3.32[I]), _ZX_32(argument2.32[I]))
  };
stage E1:
  GPR[(registerWo<<1)+1] = result1;
\end{lstlisting}}
\newpage
\subsubsection[{\small MAXW: Maximum Word}]{MAXW\hfill Maximum Word}
\label{instr:MAXW}\index{maxw}
 
\paragraph{Encoding} Format \textsf{ALU\_WWRR0} (Section~\ref{sec:format-ALU-WWRR0}), \verb|exucode2 = 0101|\\

\bitfields{ 1/P, 2/11, 1/0, 4/0101, 6/registerW, 2/11, 3/000, 1/signextw, 6/registerY, 6/registerZ }


\paragraph{Syntax} \texttt{maxw}\emph{signextw} \emph{registerW} = \emph{registerZ}, \emph{registerY}

\paragraph{Description} The maximum of the \emph{registerY} and the \emph{registerZ} is computed. The result with the 32 upper bits cleared is stored into the \emph{registerW}.


\paragraph{Modifier(s)}
\noindent\begin{itemize}
\item signextw (cf. signextw in section \ref{par:modifier-signextw} on page \pageref{par:modifier-signextw})
\end{itemize}

\paragraph{Execution} {Reservation table \textsf{ALU\_TINY} (Section~\ref{sec:reservation-ALU-TINY})}
~
{\begin{lstlisting}
stage ID:
  new signextw = _ZX_1(signextw);
stage RR:
  new argument3 = GPR[registerY];
  new argument2 = GPR[registerZ];
  new result1 = _MAX(_SX_32(argument3), _SX_32(argument2));
stage E1:
  GPR[registerW] = signextw ? _SX_32(result1) : _ZX_32(result1);
\end{lstlisting}}
\newpage

\paragraph{Encoding} Format \textsf{ALU\_WWRR0.W} (Section~\ref{sec:format-ALU-WWRR0.W}), \verb|exucode2 = 0101|\\

\bitfields{ 1/P, 2/00, 2/-- --, 27/upper27, 1/1, 2/11, 1/0, 4/0101, 6/registerW, 2/11, 3/000, 1/signextw, 1/0, 5/lower5, 6/registerZ }


\paragraph{Syntax} \texttt{maxw}\emph{signextw} \emph{registerW} = \emph{registerZ}, \emph{upper27\_\discretionary{}{}{}lower5}

\paragraph{Description} The maximum of the \emph{upper27\_\discretionary{}{}{}lower5} and the \emph{registerZ} is computed. The result with the 32 upper bits cleared is stored into the \emph{registerW}.


\paragraph{Modifier(s)}
\noindent\begin{itemize}
\item signextw (cf. signextw in section \ref{par:modifier-signextw} on page \pageref{par:modifier-signextw})
\end{itemize}

\paragraph{Execution} {Reservation table \textsf{ALU\_TINY.X} (Section~\ref{sec:reservation-ALU-TINY.X})}
~
{\begin{lstlisting}
stage ID:
  new signextw = _ZX_1(signextw);
stage RR:
  new argument3 = _WX_32(upper27_lower5);
  new argument2 = GPR[registerZ];
  new result1 = _MAX(_SX_32(argument3), _SX_32(argument2));
stage E1:
  GPR[registerW] = signextw ? _SX_32(result1) : _ZX_32(result1);
\end{lstlisting}}
\newpage
\subsubsection[{\small MAXWP: Maximum Word Pair}]{MAXWP\hfill Maximum Word Pair}
\label{instr:MAXWP}\index{maxwp}
 
\paragraph{Encoding} Format \textsf{ALU\_WPWRR0E} (Section~\ref{sec:format-ALU-WPWRR0E}), \verb|exucode2 = 0101|\\

\bitfields{ 1/P, 2/11, 1/1, 4/0101, 5/registerWe, 1/0, 2/10, 3/000, 1/0, 5/registerYe, 1/--, 5/registerZe, 1/1 }


\paragraph{Syntax} \texttt{maxwp} \emph{registerWe} = \emph{registerZe}, \emph{registerYe}

\paragraph{Description} The \emph{registerYe} and the \emph{registerZe} are interpreted as two words packed into 64 bits. The maximum of the \emph{registerYe} words and the \emph{registerZe} words are computed. The two resulting words are packed and stored into the \emph{registerWe}.


\paragraph{Execution} {Reservation table \textsf{ALU\_LITE} (Section~\ref{sec:reservation-ALU-LITE})}
~
{\begin{lstlisting}
stage RR:
  new argument3 = GPR[registerYe<<1];
  new argument2 = GPR[registerZe<<1];
  for (I = 0; I < 2; I += 1) {
    result1.32[I] = _MAX(_SX_32(argument3.32[I]), _SX_32(argument2.32[I]))
  };
stage E1:
  GPR[registerWe<<1] = result1;
\end{lstlisting}}
\newpage

\paragraph{Encoding} Format \textsf{ALU\_WPWRR0O} (Section~\ref{sec:format-ALU-WPWRR0O}), \verb|exucode2 = 0101|\\

\bitfields{ 1/P, 2/11, 1/1, 4/0101, 5/registerWo, 1/1, 2/10, 3/000, 1/0, 5/registerYo, 1/--, 5/registerZo, 1/1 }


\paragraph{Syntax} \texttt{maxwp} \emph{registerWo} = \emph{registerZo}, \emph{registerYo}

\paragraph{Description} The \emph{registerYo} and the \emph{registerZo} are interpreted as two words packed into 64 bits. The maximum of the \emph{registerYo} words and the \emph{registerZo} words are computed. The two resulting words are packed and stored into the \emph{registerWo}.


\paragraph{Execution} {Reservation table \textsf{ALU\_LITE} (Section~\ref{sec:reservation-ALU-LITE})}
~
{\begin{lstlisting}
stage RR:
  new argument3 = GPR[(registerYo<<1)+1];
  new argument2 = GPR[(registerZo<<1)+1];
  for (I = 0; I < 2; I += 1) {
    result1.32[I] = _MAX(_SX_32(argument3.32[I]), _SX_32(argument2.32[I]))
  };
stage E1:
  GPR[(registerWo<<1)+1] = result1;
\end{lstlisting}}
\newpage

\paragraph{Encoding} Format \textsf{ALU\_WPWRR0E.M} (Section~\ref{sec:format-ALU-WPWRR0E.M}), \verb|exucode2 = 0101|\\

\bitfields{ 1/P, 2/00, 2/-- --, 27/upper27, 1/1, 2/11, 1/1, 4/0101, 5/registerWe, 1/0, 2/10, 3/000, 1/0, 1/splat32, 5/lower5, 5/registerZe, 1/1 }


\paragraph{Syntax} \texttt{maxwp} \emph{registerWe} = \emph{registerZe}, \emph{upper27\_\discretionary{}{}{}lower5}\emph{splat32}

\paragraph{Description} The \emph{upper27\_\discretionary{}{}{}lower5} and the \emph{registerZe} are interpreted as two words packed into 64 bits. The maximum of the \emph{upper27\_\discretionary{}{}{}lower5} words and the \emph{registerZe} words are computed. The two resulting words are packed and stored into the \emph{registerWe}.


\paragraph{Modifier(s)}
\noindent\begin{itemize}
\item splat32 (cf. splat32 in section \ref{par:modifier-splat32} on page \pageref{par:modifier-splat32})
\end{itemize}

\paragraph{Execution} {Reservation table \textsf{ALU\_LITE.X} (Section~\ref{sec:reservation-ALU-LITE.X})}
~
{\begin{lstlisting}
stage ID:
  new splat32 = _ZX_1(splat32);
stage RR:
  new argument3 = splat32 ? (_WX_32(upper27_lower5) << 32) | _ZX_32(_WX_32(upper27_lower5)) : _SX_32(_WX_32(upper27_lower5));
  new argument2 = GPR[registerZe<<1];
  for (I = 0; I < 2; I += 1) {
    result1.32[I] = _MAX(_SX_32(argument3.32[I]), _SX_32(argument2.32[I]))
  };
stage E1:
  GPR[registerWe<<1] = result1;
\end{lstlisting}}
\newpage

\paragraph{Encoding} Format \textsf{ALU\_WPWRR0O.M} (Section~\ref{sec:format-ALU-WPWRR0O.M}), \verb|exucode2 = 0101|\\

\bitfields{ 1/P, 2/00, 2/-- --, 27/upper27, 1/1, 2/11, 1/1, 4/0101, 5/registerWo, 1/1, 2/10, 3/000, 1/0, 1/splat32, 5/lower5, 5/registerZo, 1/1 }


\paragraph{Syntax} \texttt{maxwp} \emph{registerWo} = \emph{registerZo}, \emph{upper27\_\discretionary{}{}{}lower5}\emph{splat32}

\paragraph{Description} The \emph{upper27\_\discretionary{}{}{}lower5} and the \emph{registerZo} are interpreted as two words packed into 64 bits. The maximum of the \emph{upper27\_\discretionary{}{}{}lower5} words and the \emph{registerZo} words are computed. The two resulting words are packed and stored into the \emph{registerWo}.


\paragraph{Modifier(s)}
\noindent\begin{itemize}
\item splat32 (cf. splat32 in section \ref{par:modifier-splat32} on page \pageref{par:modifier-splat32})
\end{itemize}

\paragraph{Execution} {Reservation table \textsf{ALU\_LITE.X} (Section~\ref{sec:reservation-ALU-LITE.X})}
~
{\begin{lstlisting}
stage ID:
  new splat32 = _ZX_1(splat32);
stage RR:
  new argument3 = splat32 ? (_WX_32(upper27_lower5) << 32) | _ZX_32(_WX_32(upper27_lower5)) : _SX_32(_WX_32(upper27_lower5));
  new argument2 = GPR[(registerZo<<1)+1];
  for (I = 0; I < 2; I += 1) {
    result1.32[I] = _MAX(_SX_32(argument3.32[I]), _SX_32(argument2.32[I]))
  };
stage E1:
  GPR[(registerWo<<1)+1] = result1;
\end{lstlisting}}
\newpage
\subsubsection[{\small MINBO: Minimum Byte Octuple}]{MINBO\hfill Minimum Byte Octuple}
\label{instr:MINBO}\index{minbo}
 
\paragraph{Encoding} Format \textsf{ALU\_BOWRR0E} (Section~\ref{sec:format-ALU-BOWRR0E}), \verb|exucode2 = 0100|\\

\bitfields{ 1/P, 2/11, 1/1, 4/0100, 5/registerWe, 1/0, 2/10, 3/000, 1/1, 5/registerYe, 1/--, 5/registerZe, 1/1 }


\paragraph{Syntax} \texttt{minbo} \emph{registerWe} = \emph{registerZe}, \emph{registerYe}

\paragraph{Description} The \emph{registerYe} and the \emph{registerZe} are interpreted as eight bytes packed into 64 bits. The minimum of the \emph{registerYe} bytes and the \emph{registerZe} bytes are computed. The eight resulting bytes are packed and stored into the \emph{registerWe}.


\paragraph{Execution} {Reservation table \textsf{ALU\_LITE} (Section~\ref{sec:reservation-ALU-LITE})}
~
{\begin{lstlisting}
stage RR:
  new argument3 = GPR[registerYe<<1];
  new argument2 = GPR[registerZe<<1];
  for (I = 0; I < 8; I += 1) {
    result1.8[I] = _MIN(_SX_8(argument3.8[I]), _SX_8(argument2.8[I]))
  };
stage E1:
  GPR[registerWe<<1] = result1;
\end{lstlisting}}
\newpage

\paragraph{Encoding} Format \textsf{ALU\_BOWRR0O} (Section~\ref{sec:format-ALU-BOWRR0O}), \verb|exucode2 = 0100|\\

\bitfields{ 1/P, 2/11, 1/1, 4/0100, 5/registerWo, 1/1, 2/10, 3/000, 1/1, 5/registerYo, 1/--, 5/registerZo, 1/1 }


\paragraph{Syntax} \texttt{minbo} \emph{registerWo} = \emph{registerZo}, \emph{registerYo}

\paragraph{Description} The \emph{registerYo} and the \emph{registerZo} are interpreted as eight bytes packed into 64 bits. The minimum of the \emph{registerYo} bytes and the \emph{registerZo} bytes are computed. The eight resulting bytes are packed and stored into the \emph{registerWo}.


\paragraph{Execution} {Reservation table \textsf{ALU\_LITE} (Section~\ref{sec:reservation-ALU-LITE})}
~
{\begin{lstlisting}
stage RR:
  new argument3 = GPR[(registerYo<<1)+1];
  new argument2 = GPR[(registerZo<<1)+1];
  for (I = 0; I < 8; I += 1) {
    result1.8[I] = _MIN(_SX_8(argument3.8[I]), _SX_8(argument2.8[I]))
  };
stage E1:
  GPR[(registerWo<<1)+1] = result1;
\end{lstlisting}}
\newpage

\paragraph{Encoding} Format \textsf{ALU\_BOWRR0E.M} (Section~\ref{sec:format-ALU-BOWRR0E.M}), \verb|exucode2 = 0100|\\

\bitfields{ 1/P, 2/00, 2/-- --, 27/upper27, 1/1, 2/11, 1/1, 4/0100, 5/registerWe, 1/0, 2/10, 3/000, 1/1, 1/splat32, 5/lower5, 5/registerZe, 1/1 }


\paragraph{Syntax} \texttt{minbo} \emph{registerWe} = \emph{registerZe}, \emph{upper27\_\discretionary{}{}{}lower5}\emph{splat32}

\paragraph{Description} The \emph{upper27\_\discretionary{}{}{}lower5} and the \emph{registerZe} are interpreted as eight bytes packed into 64 bits. The minimum of the \emph{upper27\_\discretionary{}{}{}lower5} bytes and the \emph{registerZe} bytes are computed. The eight resulting bytes are packed and stored into the \emph{registerWe}.


\paragraph{Modifier(s)}
\noindent\begin{itemize}
\item splat32 (cf. splat32 in section \ref{par:modifier-splat32} on page \pageref{par:modifier-splat32})
\end{itemize}

\paragraph{Execution} {Reservation table \textsf{ALU\_LITE.X} (Section~\ref{sec:reservation-ALU-LITE.X})}
~
{\begin{lstlisting}
stage ID:
  new splat32 = _ZX_1(splat32);
stage RR:
  new argument3 = splat32 ? (_WX_32(upper27_lower5) << 32) | _ZX_32(_WX_32(upper27_lower5)) : _SX_32(_WX_32(upper27_lower5));
  new argument2 = GPR[registerZe<<1];
  for (I = 0; I < 8; I += 1) {
    result1.8[I] = _MIN(_SX_8(argument3.8[I]), _SX_8(argument2.8[I]))
  };
stage E1:
  GPR[registerWe<<1] = result1;
\end{lstlisting}}
\newpage

\paragraph{Encoding} Format \textsf{ALU\_BOWRR0O.M} (Section~\ref{sec:format-ALU-BOWRR0O.M}), \verb|exucode2 = 0100|\\

\bitfields{ 1/P, 2/00, 2/-- --, 27/upper27, 1/1, 2/11, 1/1, 4/0100, 5/registerWo, 1/1, 2/10, 3/000, 1/1, 1/splat32, 5/lower5, 5/registerZo, 1/1 }


\paragraph{Syntax} \texttt{minbo} \emph{registerWo} = \emph{registerZo}, \emph{upper27\_\discretionary{}{}{}lower5}\emph{splat32}

\paragraph{Description} The \emph{upper27\_\discretionary{}{}{}lower5} and the \emph{registerZo} are interpreted as eight bytes packed into 64 bits. The minimum of the \emph{upper27\_\discretionary{}{}{}lower5} bytes and the \emph{registerZo} bytes are computed. The eight resulting bytes are packed and stored into the \emph{registerWo}.


\paragraph{Modifier(s)}
\noindent\begin{itemize}
\item splat32 (cf. splat32 in section \ref{par:modifier-splat32} on page \pageref{par:modifier-splat32})
\end{itemize}

\paragraph{Execution} {Reservation table \textsf{ALU\_LITE.X} (Section~\ref{sec:reservation-ALU-LITE.X})}
~
{\begin{lstlisting}
stage ID:
  new splat32 = _ZX_1(splat32);
stage RR:
  new argument3 = splat32 ? (_WX_32(upper27_lower5) << 32) | _ZX_32(_WX_32(upper27_lower5)) : _SX_32(_WX_32(upper27_lower5));
  new argument2 = GPR[(registerZo<<1)+1];
  for (I = 0; I < 8; I += 1) {
    result1.8[I] = _MIN(_SX_8(argument3.8[I]), _SX_8(argument2.8[I]))
  };
stage E1:
  GPR[(registerWo<<1)+1] = result1;
\end{lstlisting}}
\newpage
\subsubsection[{\small MIND: Minimum Double Word}]{MIND\hfill Minimum Double Word}
\label{instr:MIND}\index{mind}
 
\paragraph{Encoding} Format \textsf{ALU\_DWRR0} (Section~\ref{sec:format-ALU-DWRR0}), \verb|exucode2 = 0100|\\

\bitfields{ 1/P, 2/11, 1/0, 4/0100, 6/registerW, 2/10, 3/000, 1/0, 6/registerY, 6/registerZ }


\paragraph{Syntax} \texttt{mind} \emph{registerW} = \emph{registerZ}, \emph{registerY}

\paragraph{Description} The minimum of the \emph{registerY} and the \emph{registerZ} is computed. The result is stored into the \emph{registerW}.


\paragraph{Execution} {Reservation table \textsf{ALU\_TINY} (Section~\ref{sec:reservation-ALU-TINY})}
~
{\begin{lstlisting}
stage RR:
  new argument3 = GPR[registerY];
  new argument2 = GPR[registerZ];
  new result1 = _MIN(_SX_64(argument3), _SX_64(argument2));
stage E1:
  GPR[registerW] = result1;
\end{lstlisting}}
\newpage

\paragraph{Encoding} Format \textsf{ALU\_DWRR0.M} (Section~\ref{sec:format-ALU-DWRR0.M}), \verb|exucode2 = 0100|\\

\bitfields{ 1/P, 2/00, 2/-- --, 27/upper27, 1/1, 2/11, 1/0, 4/0100, 6/registerW, 2/10, 3/000, 1/0, 1/splat32, 5/lower5, 6/registerZ }


\paragraph{Syntax} \texttt{mind} \emph{registerW} = \emph{registerZ}, \emph{upper27\_\discretionary{}{}{}lower5}\emph{splat32}

\paragraph{Description} The minimum of the \emph{upper27\_\discretionary{}{}{}lower5} and the \emph{registerZ} is computed. The result is stored into the \emph{registerW}.


\paragraph{Modifier(s)}
\noindent\begin{itemize}
\item splat32 (cf. splat32 in section \ref{par:modifier-splat32} on page \pageref{par:modifier-splat32})
\end{itemize}

\paragraph{Execution} {Reservation table \textsf{ALU\_TINY.X} (Section~\ref{sec:reservation-ALU-TINY.X})}
~
{\begin{lstlisting}
stage ID:
  new splat32 = _ZX_1(splat32);
stage RR:
  new argument3 = splat32 ? (_WX_32(upper27_lower5) << 32) | _ZX_32(_WX_32(upper27_lower5)) : _SX_32(_WX_32(upper27_lower5));
  new argument2 = GPR[registerZ];
  new result1 = _MIN(_SX_64(argument3), _SX_64(argument2));
stage E1:
  GPR[registerW] = result1;
\end{lstlisting}}
\newpage

\paragraph{Encoding} Format \textsf{ALU\_DWRI} (Section~\ref{sec:format-ALU-DWRI}), \verb|exucode2 = 0100|\\

\bitfields{ 1/P, 2/11, 1/0, 4/0100, 6/registerW, 2/00, 10/signed10, 6/registerZ }


\paragraph{Syntax} \texttt{mind} \emph{registerW} = \emph{registerZ}, \emph{signed10}

\paragraph{Description} The minimum of the \emph{signed10} and the \emph{registerZ} is computed. The result is stored into the \emph{registerW}.


\paragraph{Execution} {Reservation table \textsf{ALU\_TINY} (Section~\ref{sec:reservation-ALU-TINY})}
~
{\begin{lstlisting}
stage RR:
  new argument3 = _SX_10(signed10);
  new argument2 = GPR[registerZ];
  new result1 = _MIN(_SX_64(argument3), _SX_64(argument2));
stage E1:
  GPR[registerW] = result1;
\end{lstlisting}}
\newpage

\paragraph{Encoding} Format \textsf{ALU\_DWRI.X} (Section~\ref{sec:format-ALU-DWRI.X}), \verb|exucode2 = 0100|\\

\bitfields{ 1/P, 2/00, 2/-- --, 27/upper27, 1/1, 2/11, 1/0, 4/0100, 6/registerW, 2/00, 10/lower10, 6/registerZ }


\paragraph{Syntax} \texttt{mind} \emph{registerW} = \emph{registerZ}, \emph{upper27\_\discretionary{}{}{}lower10}

\paragraph{Description} The minimum of the \emph{upper27\_\discretionary{}{}{}lower10} and the \emph{registerZ} is computed. The result is stored into the \emph{registerW}.


\paragraph{Execution} {Reservation table \textsf{ALU\_TINY.X} (Section~\ref{sec:reservation-ALU-TINY.X})}
~
{\begin{lstlisting}
stage RR:
  new argument3 = _SX_37(upper27_lower10);
  new argument2 = GPR[registerZ];
  new result1 = _MIN(_SX_64(argument3), _SX_64(argument2));
stage E1:
  GPR[registerW] = result1;
\end{lstlisting}}
\newpage

\paragraph{Encoding} Format \textsf{ALU\_DWRI.Y} (Section~\ref{sec:format-ALU-DWRI.Y}), \verb|exucode2 = 0100|\\

\bitfields{ 1/P, 2/00, 2/-- --, 27/extend27, 1/1, 2/00, 2/-- --, 27/upper27, 1/1, 2/11, 1/0, 4/0100, 6/registerW, 2/00, 10/lower10, 6/registerZ }


\paragraph{Syntax} \texttt{mind} \emph{registerW} = \emph{registerZ}, \emph{extend27\_\discretionary{}{}{}upper27\_\discretionary{}{}{}lower10}

\paragraph{Description} The minimum of the \emph{extend27\_\discretionary{}{}{}upper27\_\discretionary{}{}{}lower10} and the \emph{registerZ} is computed. The result is stored into the \emph{registerW}.


\paragraph{Execution} {Reservation table \textsf{ALU\_TINY.Y} (Section~\ref{sec:reservation-ALU-TINY.Y})}
~
{\begin{lstlisting}
stage RR:
  new argument3 = _WX_64(extend27_upper27_lower10);
  new argument2 = GPR[registerZ];
  new result1 = _MIN(_SX_64(argument3), _SX_64(argument2));
stage E1:
  GPR[registerW] = result1;
\end{lstlisting}}
\newpage
\subsubsection[{\small MINHQ: Minimum Half Word Quadruple}]{MINHQ\hfill Minimum Half Word Quadruple}
\label{instr:MINHQ}\index{minhq}
 
\paragraph{Encoding} Format \textsf{ALU\_HQWRR0E} (Section~\ref{sec:format-ALU-HQWRR0E}), \verb|exucode2 = 0100|\\

\bitfields{ 1/P, 2/11, 1/1, 4/0100, 5/registerWe, 1/0, 2/10, 3/000, 1/1, 5/registerYe, 1/--, 5/registerZe, 1/0 }


\paragraph{Syntax} \texttt{minhq} \emph{registerWe} = \emph{registerZe}, \emph{registerYe}

\paragraph{Description} The \emph{registerYe} and the \emph{registerZe} are interpreted as four half words packed into 64 bits. The minimum of the \emph{registerYe} half words and the \emph{registerZe} half words are computed. The four resulting half words are packed and stored into the \emph{registerWe}.


\paragraph{Execution} {Reservation table \textsf{ALU\_LITE} (Section~\ref{sec:reservation-ALU-LITE})}
~
{\begin{lstlisting}
stage RR:
  new argument3 = GPR[registerYe<<1];
  new argument2 = GPR[registerZe<<1];
  for (I = 0; I < 4; I += 1) {
    result1.16[I] = _MIN(_SX_16(argument3.16[I]), _SX_16(argument2.16[I]))
  };
stage E1:
  GPR[registerWe<<1] = result1;
\end{lstlisting}}
\newpage

\paragraph{Encoding} Format \textsf{ALU\_HQWRR0O} (Section~\ref{sec:format-ALU-HQWRR0O}), \verb|exucode2 = 0100|\\

\bitfields{ 1/P, 2/11, 1/1, 4/0100, 5/registerWo, 1/1, 2/10, 3/000, 1/1, 5/registerYo, 1/--, 5/registerZo, 1/0 }


\paragraph{Syntax} \texttt{minhq} \emph{registerWo} = \emph{registerZo}, \emph{registerYo}

\paragraph{Description} The \emph{registerYo} and the \emph{registerZo} are interpreted as four half words packed into 64 bits. The minimum of the \emph{registerYo} half words and the \emph{registerZo} half words are computed. The four resulting half words are packed and stored into the \emph{registerWo}.


\paragraph{Execution} {Reservation table \textsf{ALU\_LITE} (Section~\ref{sec:reservation-ALU-LITE})}
~
{\begin{lstlisting}
stage RR:
  new argument3 = GPR[(registerYo<<1)+1];
  new argument2 = GPR[(registerZo<<1)+1];
  for (I = 0; I < 4; I += 1) {
    result1.16[I] = _MIN(_SX_16(argument3.16[I]), _SX_16(argument2.16[I]))
  };
stage E1:
  GPR[(registerWo<<1)+1] = result1;
\end{lstlisting}}
\newpage

\paragraph{Encoding} Format \textsf{ALU\_HQWRR0E.M} (Section~\ref{sec:format-ALU-HQWRR0E.M}), \verb|exucode2 = 0100|\\

\bitfields{ 1/P, 2/00, 2/-- --, 27/upper27, 1/1, 2/11, 1/1, 4/0100, 5/registerWe, 1/0, 2/10, 3/000, 1/1, 1/splat32, 5/lower5, 5/registerZe, 1/0 }


\paragraph{Syntax} \texttt{minhq} \emph{registerWe} = \emph{registerZe}, \emph{upper27\_\discretionary{}{}{}lower5}\emph{splat32}

\paragraph{Description} The \emph{upper27\_\discretionary{}{}{}lower5} and the \emph{registerZe} are interpreted as four half words packed into 64 bits. The minimum of the \emph{upper27\_\discretionary{}{}{}lower5} half words and the \emph{registerZe} half words are computed. The four resulting half words are packed and stored into the \emph{registerWe}.


\paragraph{Modifier(s)}
\noindent\begin{itemize}
\item splat32 (cf. splat32 in section \ref{par:modifier-splat32} on page \pageref{par:modifier-splat32})
\end{itemize}

\paragraph{Execution} {Reservation table \textsf{ALU\_LITE.X} (Section~\ref{sec:reservation-ALU-LITE.X})}
~
{\begin{lstlisting}
stage ID:
  new splat32 = _ZX_1(splat32);
stage RR:
  new argument3 = splat32 ? (_WX_32(upper27_lower5) << 32) | _ZX_32(_WX_32(upper27_lower5)) : _SX_32(_WX_32(upper27_lower5));
  new argument2 = GPR[registerZe<<1];
  for (I = 0; I < 4; I += 1) {
    result1.16[I] = _MIN(_SX_16(argument3.16[I]), _SX_16(argument2.16[I]))
  };
stage E1:
  GPR[registerWe<<1] = result1;
\end{lstlisting}}
\newpage

\paragraph{Encoding} Format \textsf{ALU\_HQWRR0O.M} (Section~\ref{sec:format-ALU-HQWRR0O.M}), \verb|exucode2 = 0100|\\

\bitfields{ 1/P, 2/00, 2/-- --, 27/upper27, 1/1, 2/11, 1/1, 4/0100, 5/registerWo, 1/1, 2/10, 3/000, 1/1, 1/splat32, 5/lower5, 5/registerZo, 1/0 }


\paragraph{Syntax} \texttt{minhq} \emph{registerWo} = \emph{registerZo}, \emph{upper27\_\discretionary{}{}{}lower5}\emph{splat32}

\paragraph{Description} The \emph{upper27\_\discretionary{}{}{}lower5} and the \emph{registerZo} are interpreted as four half words packed into 64 bits. The minimum of the \emph{upper27\_\discretionary{}{}{}lower5} half words and the \emph{registerZo} half words are computed. The four resulting half words are packed and stored into the \emph{registerWo}.


\paragraph{Modifier(s)}
\noindent\begin{itemize}
\item splat32 (cf. splat32 in section \ref{par:modifier-splat32} on page \pageref{par:modifier-splat32})
\end{itemize}

\paragraph{Execution} {Reservation table \textsf{ALU\_LITE.X} (Section~\ref{sec:reservation-ALU-LITE.X})}
~
{\begin{lstlisting}
stage ID:
  new splat32 = _ZX_1(splat32);
stage RR:
  new argument3 = splat32 ? (_WX_32(upper27_lower5) << 32) | _ZX_32(_WX_32(upper27_lower5)) : _SX_32(_WX_32(upper27_lower5));
  new argument2 = GPR[(registerZo<<1)+1];
  for (I = 0; I < 4; I += 1) {
    result1.16[I] = _MIN(_SX_16(argument3.16[I]), _SX_16(argument2.16[I]))
  };
stage E1:
  GPR[(registerWo<<1)+1] = result1;
\end{lstlisting}}
\newpage
\subsubsection[{\small MINUBO: Minimum Unsigned Byte Octuple}]{MINUBO\hfill Minimum Unsigned Byte Octuple}
\label{instr:MINUBO}\index{minubo}
 
\paragraph{Encoding} Format \textsf{ALU\_BOWRR0E} (Section~\ref{sec:format-ALU-BOWRR0E}), \verb|exucode2 = 0110|\\

\bitfields{ 1/P, 2/11, 1/1, 4/0110, 5/registerWe, 1/0, 2/10, 3/000, 1/1, 5/registerYe, 1/--, 5/registerZe, 1/1 }


\paragraph{Syntax} \texttt{minubo} \emph{registerWe} = \emph{registerZe}, \emph{registerYe}

\paragraph{Description} The \emph{registerYe} and the \emph{registerZe} are interpreted as eight bytes packed into 64 bits. The unsigned minimum of the \emph{registerYe} bytes and the \emph{registerZe} bytes are computed. The eight resulting bytes are packed and stored into the \emph{registerWe}.


\paragraph{Execution} {Reservation table \textsf{ALU\_LITE} (Section~\ref{sec:reservation-ALU-LITE})}
~
{\begin{lstlisting}
stage RR:
  new argument3 = GPR[registerYe<<1];
  new argument2 = GPR[registerZe<<1];
  for (I = 0; I < 8; I += 1) {
    result1.8[I] = _MIN(_ZX_8(argument3.8[I]), _ZX_8(argument2.8[I]))
  };
stage E1:
  GPR[registerWe<<1] = result1;
\end{lstlisting}}
\newpage

\paragraph{Encoding} Format \textsf{ALU\_BOWRR0O} (Section~\ref{sec:format-ALU-BOWRR0O}), \verb|exucode2 = 0110|\\

\bitfields{ 1/P, 2/11, 1/1, 4/0110, 5/registerWo, 1/1, 2/10, 3/000, 1/1, 5/registerYo, 1/--, 5/registerZo, 1/1 }


\paragraph{Syntax} \texttt{minubo} \emph{registerWo} = \emph{registerZo}, \emph{registerYo}

\paragraph{Description} The \emph{registerYo} and the \emph{registerZo} are interpreted as eight bytes packed into 64 bits. The unsigned minimum of the \emph{registerYo} bytes and the \emph{registerZo} bytes are computed. The eight resulting bytes are packed and stored into the \emph{registerWo}.


\paragraph{Execution} {Reservation table \textsf{ALU\_LITE} (Section~\ref{sec:reservation-ALU-LITE})}
~
{\begin{lstlisting}
stage RR:
  new argument3 = GPR[(registerYo<<1)+1];
  new argument2 = GPR[(registerZo<<1)+1];
  for (I = 0; I < 8; I += 1) {
    result1.8[I] = _MIN(_ZX_8(argument3.8[I]), _ZX_8(argument2.8[I]))
  };
stage E1:
  GPR[(registerWo<<1)+1] = result1;
\end{lstlisting}}
\newpage

\paragraph{Encoding} Format \textsf{ALU\_BOWRR0E.M} (Section~\ref{sec:format-ALU-BOWRR0E.M}), \verb|exucode2 = 0110|\\

\bitfields{ 1/P, 2/00, 2/-- --, 27/upper27, 1/1, 2/11, 1/1, 4/0110, 5/registerWe, 1/0, 2/10, 3/000, 1/1, 1/splat32, 5/lower5, 5/registerZe, 1/1 }


\paragraph{Syntax} \texttt{minubo} \emph{registerWe} = \emph{registerZe}, \emph{upper27\_\discretionary{}{}{}lower5}\emph{splat32}

\paragraph{Description} The \emph{upper27\_\discretionary{}{}{}lower5} and the \emph{registerZe} are interpreted as eight bytes packed into 64 bits. The unsigned minimum of the \emph{upper27\_\discretionary{}{}{}lower5} bytes and the \emph{registerZe} bytes are computed. The eight resulting bytes are packed and stored into the \emph{registerWe}.


\paragraph{Modifier(s)}
\noindent\begin{itemize}
\item splat32 (cf. splat32 in section \ref{par:modifier-splat32} on page \pageref{par:modifier-splat32})
\end{itemize}

\paragraph{Execution} {Reservation table \textsf{ALU\_LITE.X} (Section~\ref{sec:reservation-ALU-LITE.X})}
~
{\begin{lstlisting}
stage ID:
  new splat32 = _ZX_1(splat32);
stage RR:
  new argument3 = splat32 ? (_WX_32(upper27_lower5) << 32) | _ZX_32(_WX_32(upper27_lower5)) : _SX_32(_WX_32(upper27_lower5));
  new argument2 = GPR[registerZe<<1];
  for (I = 0; I < 8; I += 1) {
    result1.8[I] = _MIN(_ZX_8(argument3.8[I]), _ZX_8(argument2.8[I]))
  };
stage E1:
  GPR[registerWe<<1] = result1;
\end{lstlisting}}
\newpage

\paragraph{Encoding} Format \textsf{ALU\_BOWRR0O.M} (Section~\ref{sec:format-ALU-BOWRR0O.M}), \verb|exucode2 = 0110|\\

\bitfields{ 1/P, 2/00, 2/-- --, 27/upper27, 1/1, 2/11, 1/1, 4/0110, 5/registerWo, 1/1, 2/10, 3/000, 1/1, 1/splat32, 5/lower5, 5/registerZo, 1/1 }


\paragraph{Syntax} \texttt{minubo} \emph{registerWo} = \emph{registerZo}, \emph{upper27\_\discretionary{}{}{}lower5}\emph{splat32}

\paragraph{Description} The \emph{upper27\_\discretionary{}{}{}lower5} and the \emph{registerZo} are interpreted as eight bytes packed into 64 bits. The unsigned minimum of the \emph{upper27\_\discretionary{}{}{}lower5} bytes and the \emph{registerZo} bytes are computed. The eight resulting bytes are packed and stored into the \emph{registerWo}.


\paragraph{Modifier(s)}
\noindent\begin{itemize}
\item splat32 (cf. splat32 in section \ref{par:modifier-splat32} on page \pageref{par:modifier-splat32})
\end{itemize}

\paragraph{Execution} {Reservation table \textsf{ALU\_LITE.X} (Section~\ref{sec:reservation-ALU-LITE.X})}
~
{\begin{lstlisting}
stage ID:
  new splat32 = _ZX_1(splat32);
stage RR:
  new argument3 = splat32 ? (_WX_32(upper27_lower5) << 32) | _ZX_32(_WX_32(upper27_lower5)) : _SX_32(_WX_32(upper27_lower5));
  new argument2 = GPR[(registerZo<<1)+1];
  for (I = 0; I < 8; I += 1) {
    result1.8[I] = _MIN(_ZX_8(argument3.8[I]), _ZX_8(argument2.8[I]))
  };
stage E1:
  GPR[(registerWo<<1)+1] = result1;
\end{lstlisting}}
\newpage
\subsubsection[{\small MINUD: Minimum Unsigned Double Word}]{MINUD\hfill Minimum Unsigned Double Word}
\label{instr:MINUD}\index{minud}
 
\paragraph{Encoding} Format \textsf{ALU\_DWRR0} (Section~\ref{sec:format-ALU-DWRR0}), \verb|exucode2 = 0110|\\

\bitfields{ 1/P, 2/11, 1/0, 4/0110, 6/registerW, 2/10, 3/000, 1/0, 6/registerY, 6/registerZ }


\paragraph{Syntax} \texttt{minud} \emph{registerW} = \emph{registerZ}, \emph{registerY}

\paragraph{Description} The unsigned minimum of the \emph{registerY} and the \emph{registerZ} is computed. The result is stored into the \emph{registerW}.


\paragraph{Execution} {Reservation table \textsf{ALU\_TINY} (Section~\ref{sec:reservation-ALU-TINY})}
~
{\begin{lstlisting}
stage RR:
  new argument3 = GPR[registerY];
  new argument2 = GPR[registerZ];
  new result1 = _MIN(_ZX_64(argument3), _ZX_64(argument2));
stage E1:
  GPR[registerW] = result1;
\end{lstlisting}}
\newpage

\paragraph{Encoding} Format \textsf{ALU\_DWRR0.M} (Section~\ref{sec:format-ALU-DWRR0.M}), \verb|exucode2 = 0110|\\

\bitfields{ 1/P, 2/00, 2/-- --, 27/upper27, 1/1, 2/11, 1/0, 4/0110, 6/registerW, 2/10, 3/000, 1/0, 1/splat32, 5/lower5, 6/registerZ }


\paragraph{Syntax} \texttt{minud} \emph{registerW} = \emph{registerZ}, \emph{upper27\_\discretionary{}{}{}lower5}\emph{splat32}

\paragraph{Description} The unsigned minimum of the \emph{upper27\_\discretionary{}{}{}lower5} and the \emph{registerZ} is computed. The result is stored into the \emph{registerW}.


\paragraph{Modifier(s)}
\noindent\begin{itemize}
\item splat32 (cf. splat32 in section \ref{par:modifier-splat32} on page \pageref{par:modifier-splat32})
\end{itemize}

\paragraph{Execution} {Reservation table \textsf{ALU\_TINY.X} (Section~\ref{sec:reservation-ALU-TINY.X})}
~
{\begin{lstlisting}
stage ID:
  new splat32 = _ZX_1(splat32);
stage RR:
  new argument3 = splat32 ? (_WX_32(upper27_lower5) << 32) | _ZX_32(_WX_32(upper27_lower5)) : _SX_32(_WX_32(upper27_lower5));
  new argument2 = GPR[registerZ];
  new result1 = _MIN(_ZX_64(argument3), _ZX_64(argument2));
stage E1:
  GPR[registerW] = result1;
\end{lstlisting}}
\newpage

\paragraph{Encoding} Format \textsf{ALU\_DWRI} (Section~\ref{sec:format-ALU-DWRI}), \verb|exucode2 = 0110|\\

\bitfields{ 1/P, 2/11, 1/0, 4/0110, 6/registerW, 2/00, 10/signed10, 6/registerZ }


\paragraph{Syntax} \texttt{minud} \emph{registerW} = \emph{registerZ}, \emph{signed10}

\paragraph{Description} The unsigned minimum of the \emph{signed10} and the \emph{registerZ} is computed. The result is stored into the \emph{registerW}.


\paragraph{Execution} {Reservation table \textsf{ALU\_TINY} (Section~\ref{sec:reservation-ALU-TINY})}
~
{\begin{lstlisting}
stage RR:
  new argument3 = _SX_10(signed10);
  new argument2 = GPR[registerZ];
  new result1 = _MIN(_ZX_64(argument3), _ZX_64(argument2));
stage E1:
  GPR[registerW] = result1;
\end{lstlisting}}
\newpage

\paragraph{Encoding} Format \textsf{ALU\_DWRI.X} (Section~\ref{sec:format-ALU-DWRI.X}), \verb|exucode2 = 0110|\\

\bitfields{ 1/P, 2/00, 2/-- --, 27/upper27, 1/1, 2/11, 1/0, 4/0110, 6/registerW, 2/00, 10/lower10, 6/registerZ }


\paragraph{Syntax} \texttt{minud} \emph{registerW} = \emph{registerZ}, \emph{upper27\_\discretionary{}{}{}lower10}

\paragraph{Description} The unsigned minimum of the \emph{upper27\_\discretionary{}{}{}lower10} and the \emph{registerZ} is computed. The result is stored into the \emph{registerW}.


\paragraph{Execution} {Reservation table \textsf{ALU\_TINY.X} (Section~\ref{sec:reservation-ALU-TINY.X})}
~
{\begin{lstlisting}
stage RR:
  new argument3 = _SX_37(upper27_lower10);
  new argument2 = GPR[registerZ];
  new result1 = _MIN(_ZX_64(argument3), _ZX_64(argument2));
stage E1:
  GPR[registerW] = result1;
\end{lstlisting}}
\newpage

\paragraph{Encoding} Format \textsf{ALU\_DWRI.Y} (Section~\ref{sec:format-ALU-DWRI.Y}), \verb|exucode2 = 0110|\\

\bitfields{ 1/P, 2/00, 2/-- --, 27/extend27, 1/1, 2/00, 2/-- --, 27/upper27, 1/1, 2/11, 1/0, 4/0110, 6/registerW, 2/00, 10/lower10, 6/registerZ }


\paragraph{Syntax} \texttt{minud} \emph{registerW} = \emph{registerZ}, \emph{extend27\_\discretionary{}{}{}upper27\_\discretionary{}{}{}lower10}

\paragraph{Description} The unsigned minimum of the \emph{extend27\_\discretionary{}{}{}upper27\_\discretionary{}{}{}lower10} and the \emph{registerZ} is computed. The result is stored into the \emph{registerW}.


\paragraph{Execution} {Reservation table \textsf{ALU\_TINY.Y} (Section~\ref{sec:reservation-ALU-TINY.Y})}
~
{\begin{lstlisting}
stage RR:
  new argument3 = _WX_64(extend27_upper27_lower10);
  new argument2 = GPR[registerZ];
  new result1 = _MIN(_ZX_64(argument3), _ZX_64(argument2));
stage E1:
  GPR[registerW] = result1;
\end{lstlisting}}
\newpage
\subsubsection[{\small MINUHQ: Minimum Unsigned Half Word Octuple}]{MINUHQ\hfill Minimum Unsigned Half Word Octuple}
\label{instr:MINUHQ}\index{minuhq}
 
\paragraph{Encoding} Format \textsf{ALU\_HQWRR0E} (Section~\ref{sec:format-ALU-HQWRR0E}), \verb|exucode2 = 0110|\\

\bitfields{ 1/P, 2/11, 1/1, 4/0110, 5/registerWe, 1/0, 2/10, 3/000, 1/1, 5/registerYe, 1/--, 5/registerZe, 1/0 }


\paragraph{Syntax} \texttt{minuhq} \emph{registerWe} = \emph{registerZe}, \emph{registerYe}

\paragraph{Description} The \emph{registerYe} and the \emph{registerZe} are interpreted as four half words packed into 64 bits. The unsigned minimum of the \emph{registerYe} half words and the \emph{registerZe} half words are computed. The four resulting half words are packed and stored into the \emph{registerWe}.


\paragraph{Execution} {Reservation table \textsf{ALU\_LITE} (Section~\ref{sec:reservation-ALU-LITE})}
~
{\begin{lstlisting}
stage RR:
  new argument3 = GPR[registerYe<<1];
  new argument2 = GPR[registerZe<<1];
  for (I = 0; I < 4; I += 1) {
    result1.16[I] = _MIN(_ZX_16(argument3.16[I]), _ZX_16(argument2.16[I]))
  };
stage E1:
  GPR[registerWe<<1] = result1;
\end{lstlisting}}
\newpage

\paragraph{Encoding} Format \textsf{ALU\_HQWRR0O} (Section~\ref{sec:format-ALU-HQWRR0O}), \verb|exucode2 = 0110|\\

\bitfields{ 1/P, 2/11, 1/1, 4/0110, 5/registerWo, 1/1, 2/10, 3/000, 1/1, 5/registerYo, 1/--, 5/registerZo, 1/0 }


\paragraph{Syntax} \texttt{minuhq} \emph{registerWo} = \emph{registerZo}, \emph{registerYo}

\paragraph{Description} The \emph{registerYo} and the \emph{registerZo} are interpreted as four half words packed into 64 bits. The unsigned minimum of the \emph{registerYo} half words and the \emph{registerZo} half words are computed. The four resulting half words are packed and stored into the \emph{registerWo}.


\paragraph{Execution} {Reservation table \textsf{ALU\_LITE} (Section~\ref{sec:reservation-ALU-LITE})}
~
{\begin{lstlisting}
stage RR:
  new argument3 = GPR[(registerYo<<1)+1];
  new argument2 = GPR[(registerZo<<1)+1];
  for (I = 0; I < 4; I += 1) {
    result1.16[I] = _MIN(_ZX_16(argument3.16[I]), _ZX_16(argument2.16[I]))
  };
stage E1:
  GPR[(registerWo<<1)+1] = result1;
\end{lstlisting}}
\newpage

\paragraph{Encoding} Format \textsf{ALU\_HQWRR0E.M} (Section~\ref{sec:format-ALU-HQWRR0E.M}), \verb|exucode2 = 0110|\\

\bitfields{ 1/P, 2/00, 2/-- --, 27/upper27, 1/1, 2/11, 1/1, 4/0110, 5/registerWe, 1/0, 2/10, 3/000, 1/1, 1/splat32, 5/lower5, 5/registerZe, 1/0 }


\paragraph{Syntax} \texttt{minuhq} \emph{registerWe} = \emph{registerZe}, \emph{upper27\_\discretionary{}{}{}lower5}\emph{splat32}

\paragraph{Description} The \emph{upper27\_\discretionary{}{}{}lower5} and the \emph{registerZe} are interpreted as four half words packed into 64 bits. The unsigned minimum of the \emph{upper27\_\discretionary{}{}{}lower5} half words and the \emph{registerZe} half words are computed. The four resulting half words are packed and stored into the \emph{registerWe}.


\paragraph{Modifier(s)}
\noindent\begin{itemize}
\item splat32 (cf. splat32 in section \ref{par:modifier-splat32} on page \pageref{par:modifier-splat32})
\end{itemize}

\paragraph{Execution} {Reservation table \textsf{ALU\_LITE.X} (Section~\ref{sec:reservation-ALU-LITE.X})}
~
{\begin{lstlisting}
stage ID:
  new splat32 = _ZX_1(splat32);
stage RR:
  new argument3 = splat32 ? (_WX_32(upper27_lower5) << 32) | _ZX_32(_WX_32(upper27_lower5)) : _SX_32(_WX_32(upper27_lower5));
  new argument2 = GPR[registerZe<<1];
  for (I = 0; I < 4; I += 1) {
    result1.16[I] = _MIN(_ZX_16(argument3.16[I]), _ZX_16(argument2.16[I]))
  };
stage E1:
  GPR[registerWe<<1] = result1;
\end{lstlisting}}
\newpage

\paragraph{Encoding} Format \textsf{ALU\_HQWRR0O.M} (Section~\ref{sec:format-ALU-HQWRR0O.M}), \verb|exucode2 = 0110|\\

\bitfields{ 1/P, 2/00, 2/-- --, 27/upper27, 1/1, 2/11, 1/1, 4/0110, 5/registerWo, 1/1, 2/10, 3/000, 1/1, 1/splat32, 5/lower5, 5/registerZo, 1/0 }


\paragraph{Syntax} \texttt{minuhq} \emph{registerWo} = \emph{registerZo}, \emph{upper27\_\discretionary{}{}{}lower5}\emph{splat32}

\paragraph{Description} The \emph{upper27\_\discretionary{}{}{}lower5} and the \emph{registerZo} are interpreted as four half words packed into 64 bits. The unsigned minimum of the \emph{upper27\_\discretionary{}{}{}lower5} half words and the \emph{registerZo} half words are computed. The four resulting half words are packed and stored into the \emph{registerWo}.


\paragraph{Modifier(s)}
\noindent\begin{itemize}
\item splat32 (cf. splat32 in section \ref{par:modifier-splat32} on page \pageref{par:modifier-splat32})
\end{itemize}

\paragraph{Execution} {Reservation table \textsf{ALU\_LITE.X} (Section~\ref{sec:reservation-ALU-LITE.X})}
~
{\begin{lstlisting}
stage ID:
  new splat32 = _ZX_1(splat32);
stage RR:
  new argument3 = splat32 ? (_WX_32(upper27_lower5) << 32) | _ZX_32(_WX_32(upper27_lower5)) : _SX_32(_WX_32(upper27_lower5));
  new argument2 = GPR[(registerZo<<1)+1];
  for (I = 0; I < 4; I += 1) {
    result1.16[I] = _MIN(_ZX_16(argument3.16[I]), _ZX_16(argument2.16[I]))
  };
stage E1:
  GPR[(registerWo<<1)+1] = result1;
\end{lstlisting}}
\newpage
\subsubsection[{\small MINUW: Minimum Unsigned Word}]{MINUW\hfill Minimum Unsigned Word}
\label{instr:MINUW}\index{minuw}
 
\paragraph{Encoding} Format \textsf{ALU\_WWRR0} (Section~\ref{sec:format-ALU-WWRR0}), \verb|exucode2 = 0110|\\

\bitfields{ 1/P, 2/11, 1/0, 4/0110, 6/registerW, 2/11, 3/000, 1/signextw, 6/registerY, 6/registerZ }


\paragraph{Syntax} \texttt{minuw}\emph{signextw} \emph{registerW} = \emph{registerZ}, \emph{registerY}

\paragraph{Description} The unsigned minimum of the \emph{registerY} and the \emph{registerZ} is computed. The result with the 32 upper bits cleared is stored into the \emph{registerW}.


\paragraph{Modifier(s)}
\noindent\begin{itemize}
\item signextw (cf. signextw in section \ref{par:modifier-signextw} on page \pageref{par:modifier-signextw})
\end{itemize}

\paragraph{Execution} {Reservation table \textsf{ALU\_TINY} (Section~\ref{sec:reservation-ALU-TINY})}
~
{\begin{lstlisting}
stage ID:
  new signextw = _ZX_1(signextw);
stage RR:
  new argument3 = GPR[registerY];
  new argument2 = GPR[registerZ];
  new result1 = _MIN(argument3.32[0], argument2.32[0]);
stage E1:
  GPR[registerW] = signextw ? _SX_32(result1) : _ZX_32(result1);
\end{lstlisting}}
\newpage

\paragraph{Encoding} Format \textsf{ALU\_WWRR0.W} (Section~\ref{sec:format-ALU-WWRR0.W}), \verb|exucode2 = 0110|\\

\bitfields{ 1/P, 2/00, 2/-- --, 27/upper27, 1/1, 2/11, 1/0, 4/0110, 6/registerW, 2/11, 3/000, 1/signextw, 1/0, 5/lower5, 6/registerZ }


\paragraph{Syntax} \texttt{minuw}\emph{signextw} \emph{registerW} = \emph{registerZ}, \emph{upper27\_\discretionary{}{}{}lower5}

\paragraph{Description} The unsigned minimum of the \emph{upper27\_\discretionary{}{}{}lower5} and the \emph{registerZ} is computed. The result with the 32 upper bits cleared is stored into the \emph{registerW}.


\paragraph{Modifier(s)}
\noindent\begin{itemize}
\item signextw (cf. signextw in section \ref{par:modifier-signextw} on page \pageref{par:modifier-signextw})
\end{itemize}

\paragraph{Execution} {Reservation table \textsf{ALU\_TINY.X} (Section~\ref{sec:reservation-ALU-TINY.X})}
~
{\begin{lstlisting}
stage ID:
  new signextw = _ZX_1(signextw);
stage RR:
  new argument3 = _WX_32(upper27_lower5);
  new argument2 = GPR[registerZ];
  new result1 = _MIN(argument3.32[0], argument2.32[0]);
stage E1:
  GPR[registerW] = signextw ? _SX_32(result1) : _ZX_32(result1);
\end{lstlisting}}
\newpage
\subsubsection[{\small MINUWP: Minimum Unsigned Word Pair}]{MINUWP\hfill Minimum Unsigned Word Pair}
\label{instr:MINUWP}\index{minuwp}
 
\paragraph{Encoding} Format \textsf{ALU\_WPWRR0E} (Section~\ref{sec:format-ALU-WPWRR0E}), \verb|exucode2 = 0110|\\

\bitfields{ 1/P, 2/11, 1/1, 4/0110, 5/registerWe, 1/0, 2/10, 3/000, 1/0, 5/registerYe, 1/--, 5/registerZe, 1/1 }


\paragraph{Syntax} \texttt{minuwp} \emph{registerWe} = \emph{registerZe}, \emph{registerYe}

\paragraph{Description} The \emph{registerYe} and the \emph{registerZe} are interpreted as two words packed into 64 bits. The unsigned minimum of the \emph{registerYe} words and the \emph{registerZe} words are computed. The two resulting words are packed and stored into the \emph{registerWe}.


\paragraph{Execution} {Reservation table \textsf{ALU\_LITE} (Section~\ref{sec:reservation-ALU-LITE})}
~
{\begin{lstlisting}
stage RR:
  new argument3 = GPR[registerYe<<1];
  new argument2 = GPR[registerZe<<1];
  for (I = 0; I < 2; I += 1) {
    result1.32[I] = _MIN(_ZX_32(argument3.32[I]), _ZX_32(argument2.32[I]))
  };
stage E1:
  GPR[registerWe<<1] = result1;
\end{lstlisting}}
\newpage

\paragraph{Encoding} Format \textsf{ALU\_WPWRR0O} (Section~\ref{sec:format-ALU-WPWRR0O}), \verb|exucode2 = 0110|\\

\bitfields{ 1/P, 2/11, 1/1, 4/0110, 5/registerWo, 1/1, 2/10, 3/000, 1/0, 5/registerYo, 1/--, 5/registerZo, 1/1 }


\paragraph{Syntax} \texttt{minuwp} \emph{registerWo} = \emph{registerZo}, \emph{registerYo}

\paragraph{Description} The \emph{registerYo} and the \emph{registerZo} are interpreted as two words packed into 64 bits. The unsigned minimum of the \emph{registerYo} words and the \emph{registerZo} words are computed. The two resulting words are packed and stored into the \emph{registerWo}.


\paragraph{Execution} {Reservation table \textsf{ALU\_LITE} (Section~\ref{sec:reservation-ALU-LITE})}
~
{\begin{lstlisting}
stage RR:
  new argument3 = GPR[(registerYo<<1)+1];
  new argument2 = GPR[(registerZo<<1)+1];
  for (I = 0; I < 2; I += 1) {
    result1.32[I] = _MIN(_ZX_32(argument3.32[I]), _ZX_32(argument2.32[I]))
  };
stage E1:
  GPR[(registerWo<<1)+1] = result1;
\end{lstlisting}}
\newpage

\paragraph{Encoding} Format \textsf{ALU\_WPWRR0E.M} (Section~\ref{sec:format-ALU-WPWRR0E.M}), \verb|exucode2 = 0110|\\

\bitfields{ 1/P, 2/00, 2/-- --, 27/upper27, 1/1, 2/11, 1/1, 4/0110, 5/registerWe, 1/0, 2/10, 3/000, 1/0, 1/splat32, 5/lower5, 5/registerZe, 1/1 }


\paragraph{Syntax} \texttt{minuwp} \emph{registerWe} = \emph{registerZe}, \emph{upper27\_\discretionary{}{}{}lower5}\emph{splat32}

\paragraph{Description} The \emph{upper27\_\discretionary{}{}{}lower5} and the \emph{registerZe} are interpreted as two words packed into 64 bits. The unsigned minimum of the \emph{upper27\_\discretionary{}{}{}lower5} words and the \emph{registerZe} words are computed. The two resulting words are packed and stored into the \emph{registerWe}.


\paragraph{Modifier(s)}
\noindent\begin{itemize}
\item splat32 (cf. splat32 in section \ref{par:modifier-splat32} on page \pageref{par:modifier-splat32})
\end{itemize}

\paragraph{Execution} {Reservation table \textsf{ALU\_LITE.X} (Section~\ref{sec:reservation-ALU-LITE.X})}
~
{\begin{lstlisting}
stage ID:
  new splat32 = _ZX_1(splat32);
stage RR:
  new argument3 = splat32 ? (_WX_32(upper27_lower5) << 32) | _ZX_32(_WX_32(upper27_lower5)) : _SX_32(_WX_32(upper27_lower5));
  new argument2 = GPR[registerZe<<1];
  for (I = 0; I < 2; I += 1) {
    result1.32[I] = _MIN(_ZX_32(argument3.32[I]), _ZX_32(argument2.32[I]))
  };
stage E1:
  GPR[registerWe<<1] = result1;
\end{lstlisting}}
\newpage

\paragraph{Encoding} Format \textsf{ALU\_WPWRR0O.M} (Section~\ref{sec:format-ALU-WPWRR0O.M}), \verb|exucode2 = 0110|\\

\bitfields{ 1/P, 2/00, 2/-- --, 27/upper27, 1/1, 2/11, 1/1, 4/0110, 5/registerWo, 1/1, 2/10, 3/000, 1/0, 1/splat32, 5/lower5, 5/registerZo, 1/1 }


\paragraph{Syntax} \texttt{minuwp} \emph{registerWo} = \emph{registerZo}, \emph{upper27\_\discretionary{}{}{}lower5}\emph{splat32}

\paragraph{Description} The \emph{upper27\_\discretionary{}{}{}lower5} and the \emph{registerZo} are interpreted as two words packed into 64 bits. The unsigned minimum of the \emph{upper27\_\discretionary{}{}{}lower5} words and the \emph{registerZo} words are computed. The two resulting words are packed and stored into the \emph{registerWo}.


\paragraph{Modifier(s)}
\noindent\begin{itemize}
\item splat32 (cf. splat32 in section \ref{par:modifier-splat32} on page \pageref{par:modifier-splat32})
\end{itemize}

\paragraph{Execution} {Reservation table \textsf{ALU\_LITE.X} (Section~\ref{sec:reservation-ALU-LITE.X})}
~
{\begin{lstlisting}
stage ID:
  new splat32 = _ZX_1(splat32);
stage RR:
  new argument3 = splat32 ? (_WX_32(upper27_lower5) << 32) | _ZX_32(_WX_32(upper27_lower5)) : _SX_32(_WX_32(upper27_lower5));
  new argument2 = GPR[(registerZo<<1)+1];
  for (I = 0; I < 2; I += 1) {
    result1.32[I] = _MIN(_ZX_32(argument3.32[I]), _ZX_32(argument2.32[I]))
  };
stage E1:
  GPR[(registerWo<<1)+1] = result1;
\end{lstlisting}}
\newpage
\subsubsection[{\small MINW: Minimum Word}]{MINW\hfill Minimum Word}
\label{instr:MINW}\index{minw}
 
\paragraph{Encoding} Format \textsf{ALU\_WWRR0} (Section~\ref{sec:format-ALU-WWRR0}), \verb|exucode2 = 0100|\\

\bitfields{ 1/P, 2/11, 1/0, 4/0100, 6/registerW, 2/11, 3/000, 1/signextw, 6/registerY, 6/registerZ }


\paragraph{Syntax} \texttt{minw}\emph{signextw} \emph{registerW} = \emph{registerZ}, \emph{registerY}

\paragraph{Description} The minimum of the \emph{registerY} and the \emph{registerZ} is computed. The result with the 32 upper bits cleared is stored into the \emph{registerW}.


\paragraph{Modifier(s)}
\noindent\begin{itemize}
\item signextw (cf. signextw in section \ref{par:modifier-signextw} on page \pageref{par:modifier-signextw})
\end{itemize}

\paragraph{Execution} {Reservation table \textsf{ALU\_TINY} (Section~\ref{sec:reservation-ALU-TINY})}
~
{\begin{lstlisting}
stage ID:
  new signextw = _ZX_1(signextw);
stage RR:
  new argument3 = GPR[registerY];
  new argument2 = GPR[registerZ];
  new result1 = _MIN(_SX_32(argument3), _SX_32(argument2));
stage E1:
  GPR[registerW] = signextw ? _SX_32(result1) : _ZX_32(result1);
\end{lstlisting}}
\newpage

\paragraph{Encoding} Format \textsf{ALU\_WWRR0.W} (Section~\ref{sec:format-ALU-WWRR0.W}), \verb|exucode2 = 0100|\\

\bitfields{ 1/P, 2/00, 2/-- --, 27/upper27, 1/1, 2/11, 1/0, 4/0100, 6/registerW, 2/11, 3/000, 1/signextw, 1/0, 5/lower5, 6/registerZ }


\paragraph{Syntax} \texttt{minw}\emph{signextw} \emph{registerW} = \emph{registerZ}, \emph{upper27\_\discretionary{}{}{}lower5}

\paragraph{Description} The minimum of the \emph{upper27\_\discretionary{}{}{}lower5} and the \emph{registerZ} is computed. The result with the 32 upper bits cleared is stored into the \emph{registerW}.


\paragraph{Modifier(s)}
\noindent\begin{itemize}
\item signextw (cf. signextw in section \ref{par:modifier-signextw} on page \pageref{par:modifier-signextw})
\end{itemize}

\paragraph{Execution} {Reservation table \textsf{ALU\_TINY.X} (Section~\ref{sec:reservation-ALU-TINY.X})}
~
{\begin{lstlisting}
stage ID:
  new signextw = _ZX_1(signextw);
stage RR:
  new argument3 = _WX_32(upper27_lower5);
  new argument2 = GPR[registerZ];
  new result1 = _MIN(_SX_32(argument3), _SX_32(argument2));
stage E1:
  GPR[registerW] = signextw ? _SX_32(result1) : _ZX_32(result1);
\end{lstlisting}}
\newpage
\subsubsection[{\small MINWP: Minimum Word Pair}]{MINWP\hfill Minimum Word Pair}
\label{instr:MINWP}\index{minwp}
 
\paragraph{Encoding} Format \textsf{ALU\_WPWRR0E} (Section~\ref{sec:format-ALU-WPWRR0E}), \verb|exucode2 = 0100|\\

\bitfields{ 1/P, 2/11, 1/1, 4/0100, 5/registerWe, 1/0, 2/10, 3/000, 1/0, 5/registerYe, 1/--, 5/registerZe, 1/1 }


\paragraph{Syntax} \texttt{minwp} \emph{registerWe} = \emph{registerZe}, \emph{registerYe}

\paragraph{Description} The \emph{registerYe} and the \emph{registerZe} are interpreted as two words packed into 64 bits. The minimum of the \emph{registerYe} words and the \emph{registerZe} words are computed. The two resulting words are packed and stored into the \emph{registerWe}.


\paragraph{Execution} {Reservation table \textsf{ALU\_LITE} (Section~\ref{sec:reservation-ALU-LITE})}
~
{\begin{lstlisting}
stage RR:
  new argument3 = GPR[registerYe<<1];
  new argument2 = GPR[registerZe<<1];
  for (I = 0; I < 2; I += 1) {
    result1.32[I] = _MIN(_SX_32(argument3.32[I]), _SX_32(argument2.32[I]))
  };
stage E1:
  GPR[registerWe<<1] = result1;
\end{lstlisting}}
\newpage

\paragraph{Encoding} Format \textsf{ALU\_WPWRR0O} (Section~\ref{sec:format-ALU-WPWRR0O}), \verb|exucode2 = 0100|\\

\bitfields{ 1/P, 2/11, 1/1, 4/0100, 5/registerWo, 1/1, 2/10, 3/000, 1/0, 5/registerYo, 1/--, 5/registerZo, 1/1 }


\paragraph{Syntax} \texttt{minwp} \emph{registerWo} = \emph{registerZo}, \emph{registerYo}

\paragraph{Description} The \emph{registerYo} and the \emph{registerZo} are interpreted as two words packed into 64 bits. The minimum of the \emph{registerYo} words and the \emph{registerZo} words are computed. The two resulting words are packed and stored into the \emph{registerWo}.


\paragraph{Execution} {Reservation table \textsf{ALU\_LITE} (Section~\ref{sec:reservation-ALU-LITE})}
~
{\begin{lstlisting}
stage RR:
  new argument3 = GPR[(registerYo<<1)+1];
  new argument2 = GPR[(registerZo<<1)+1];
  for (I = 0; I < 2; I += 1) {
    result1.32[I] = _MIN(_SX_32(argument3.32[I]), _SX_32(argument2.32[I]))
  };
stage E1:
  GPR[(registerWo<<1)+1] = result1;
\end{lstlisting}}
\newpage

\paragraph{Encoding} Format \textsf{ALU\_WPWRR0E.M} (Section~\ref{sec:format-ALU-WPWRR0E.M}), \verb|exucode2 = 0100|\\

\bitfields{ 1/P, 2/00, 2/-- --, 27/upper27, 1/1, 2/11, 1/1, 4/0100, 5/registerWe, 1/0, 2/10, 3/000, 1/0, 1/splat32, 5/lower5, 5/registerZe, 1/1 }


\paragraph{Syntax} \texttt{minwp} \emph{registerWe} = \emph{registerZe}, \emph{upper27\_\discretionary{}{}{}lower5}\emph{splat32}

\paragraph{Description} The \emph{upper27\_\discretionary{}{}{}lower5} and the \emph{registerZe} are interpreted as two words packed into 64 bits. The minimum of the \emph{upper27\_\discretionary{}{}{}lower5} words and the \emph{registerZe} words are computed. The two resulting words are packed and stored into the \emph{registerWe}.


\paragraph{Modifier(s)}
\noindent\begin{itemize}
\item splat32 (cf. splat32 in section \ref{par:modifier-splat32} on page \pageref{par:modifier-splat32})
\end{itemize}

\paragraph{Execution} {Reservation table \textsf{ALU\_LITE.X} (Section~\ref{sec:reservation-ALU-LITE.X})}
~
{\begin{lstlisting}
stage ID:
  new splat32 = _ZX_1(splat32);
stage RR:
  new argument3 = splat32 ? (_WX_32(upper27_lower5) << 32) | _ZX_32(_WX_32(upper27_lower5)) : _SX_32(_WX_32(upper27_lower5));
  new argument2 = GPR[registerZe<<1];
  for (I = 0; I < 2; I += 1) {
    result1.32[I] = _MIN(_SX_32(argument3.32[I]), _SX_32(argument2.32[I]))
  };
stage E1:
  GPR[registerWe<<1] = result1;
\end{lstlisting}}
\newpage

\paragraph{Encoding} Format \textsf{ALU\_WPWRR0O.M} (Section~\ref{sec:format-ALU-WPWRR0O.M}), \verb|exucode2 = 0100|\\

\bitfields{ 1/P, 2/00, 2/-- --, 27/upper27, 1/1, 2/11, 1/1, 4/0100, 5/registerWo, 1/1, 2/10, 3/000, 1/0, 1/splat32, 5/lower5, 5/registerZo, 1/1 }


\paragraph{Syntax} \texttt{minwp} \emph{registerWo} = \emph{registerZo}, \emph{upper27\_\discretionary{}{}{}lower5}\emph{splat32}

\paragraph{Description} The \emph{upper27\_\discretionary{}{}{}lower5} and the \emph{registerZo} are interpreted as two words packed into 64 bits. The minimum of the \emph{upper27\_\discretionary{}{}{}lower5} words and the \emph{registerZo} words are computed. The two resulting words are packed and stored into the \emph{registerWo}.


\paragraph{Modifier(s)}
\noindent\begin{itemize}
\item splat32 (cf. splat32 in section \ref{par:modifier-splat32} on page \pageref{par:modifier-splat32})
\end{itemize}

\paragraph{Execution} {Reservation table \textsf{ALU\_LITE.X} (Section~\ref{sec:reservation-ALU-LITE.X})}
~
{\begin{lstlisting}
stage ID:
  new splat32 = _ZX_1(splat32);
stage RR:
  new argument3 = splat32 ? (_WX_32(upper27_lower5) << 32) | _ZX_32(_WX_32(upper27_lower5)) : _SX_32(_WX_32(upper27_lower5));
  new argument2 = GPR[(registerZo<<1)+1];
  for (I = 0; I < 2; I += 1) {
    result1.32[I] = _MIN(_SX_32(argument3.32[I]), _SX_32(argument2.32[I]))
  };
stage E1:
  GPR[(registerWo<<1)+1] = result1;
\end{lstlisting}}
\newpage
\subsubsection[{\small MSBFBHO: Multiply Subtract From Extend Byte to Half Word Octuple}]{MSBFBHO\hfill Multiply Subtract From Extend Byte to Half Word Octuple}
\label{instr:MSBFBHO}\index{msbfbho}
 
\paragraph{Encoding} Format \textsf{ALU\_HOMEWRRRE} (Section~\ref{sec:format-ALU-HOMEWRRRE}), \verb|alucode2 = 11|\\

\bitfields{ 1/P, 2/11, 1/1, 2/11, 2/widemult, 5/registerM, 1/0, 2/10, 3/110, 1/1, 5/registerYe, 1/1, 5/registerZe, 1/0 }


\paragraph{Syntax} \texttt{msbfbho}\emph{widemult} \emph{registerM} = \emph{registerZe}, \emph{registerYe}

\paragraph{Description} The \emph{registerYe} and the \emph{registerZe} are interpreted as eight bytes packed into 64 bits. These \emph{registerYe} and \emph{registerZe} bytes are sign-extended or zero-extended from 8 bits, depending on the \emph{widemult}. The \emph{registerYe} bytes are multiplied by the \emph{registerZe} bytes, producing eight 16-bit results. The eight resulting half words are packed, subtracted from and stored into the \emph{registerM}.


\paragraph{Modifier(s)}
\noindent\begin{itemize}
\item widemult (cf. widemult in section \ref{par:modifier-widemult} on page \pageref{par:modifier-widemult})
\end{itemize}

\paragraph{Execution} {Reservation table \textsf{ALU\_LITE} (Section~\ref{sec:reservation-ALU-LITE})}
~
{\begin{lstlisting}
stage ID:
  new widemult = _ZX_2(widemult);
stage RR:
  new argument3 = GPR[registerYe<<1];
  new argument2 = GPR[registerZe<<1];
stage E1:
  new argument1 = PGR[registerM];
  for (I = 0; I < 8; I += 1) {
    if (widemult == 0) {
      new product = _SX_8(argument2.8[I]) * _SX_8(argument3.8[I]);
    }
    if (widemult == 1) {
      product = _ZX_8(argument2.8[I]) * _ZX_8(argument3.8[I]);
    }
    if (widemult == 2) {
      product = _SX_8(argument2.8[I]) * _ZX_8(argument3.8[I]);
    }
    result1.16[I] = argument1.16[I] - product
  };
stage E2:
  PGR[registerM] = result1;
\end{lstlisting}}
\newpage

\paragraph{Encoding} Format \textsf{ALU\_HOMEWRRRO} (Section~\ref{sec:format-ALU-HOMEWRRRO}), \verb|alucode2 = 11|\\

\bitfields{ 1/P, 2/11, 1/1, 2/11, 2/widemult, 5/registerM, 1/1, 2/10, 3/110, 1/1, 5/registerYo, 1/1, 5/registerZo, 1/0 }


\paragraph{Syntax} \texttt{msbfbho}\emph{widemult} \emph{registerM} = \emph{registerZo}, \emph{registerYo}

\paragraph{Description} The \emph{registerYo} and the \emph{registerZo} are interpreted as eight bytes packed into 64 bits. These \emph{registerYo} and \emph{registerZo} bytes are sign-extended or zero-extended from 8 bits, depending on the \emph{widemult}. The \emph{registerYo} bytes are multiplied by the \emph{registerZo} bytes, producing eight 16-bit results. The eight resulting half words are packed, subtracted from and stored into the \emph{registerM}.


\paragraph{Modifier(s)}
\noindent\begin{itemize}
\item widemult (cf. widemult in section \ref{par:modifier-widemult} on page \pageref{par:modifier-widemult})
\end{itemize}

\paragraph{Execution} {Reservation table \textsf{ALU\_LITE} (Section~\ref{sec:reservation-ALU-LITE})}
~
{\begin{lstlisting}
stage ID:
  new widemult = _ZX_2(widemult);
stage RR:
  new argument3 = GPR[(registerYo<<1)+1];
  new argument2 = GPR[(registerZo<<1)+1];
stage E1:
  new argument1 = PGR[registerM];
  for (I = 0; I < 8; I += 1) {
    if (widemult == 0) {
      new product = _SX_8(argument2.8[I]) * _SX_8(argument3.8[I]);
    }
    if (widemult == 1) {
      product = _ZX_8(argument2.8[I]) * _ZX_8(argument3.8[I]);
    }
    if (widemult == 2) {
      product = _SX_8(argument2.8[I]) * _ZX_8(argument3.8[I]);
    }
    result1.16[I] = argument1.16[I] - product
  };
stage E2:
  PGR[registerM] = result1;
\end{lstlisting}}
\newpage
\subsubsection[{\small MSBFD: Multiply Subtract From Double Word}]{MSBFD\hfill Multiply Subtract From Double Word}
\label{instr:MSBFD}\index{msbfd}
 
\paragraph{Encoding} Format \textsf{ALU\_DMWRRR} (Section~\ref{sec:format-ALU-DMWRRR}), \verb|alucode2 = 11|\\

\bitfields{ 1/P, 2/11, 1/0, 2/11, 2/highmult, 6/registerW, 2/10, 3/110, 1/0, 6/registerY, 6/registerZ }


\paragraph{Syntax} \texttt{msbfd}\emph{highmult} \emph{registerW} = \emph{registerZ}, \emph{registerY}

\paragraph{Description} The \emph{registerY} and the \emph{registerZ} double words are sign-extended or zero-extended from 64 bits, depending on the \emph{highmult}. The \emph{registerY} double word is multiplied by the \emph{registerZ} double word, and the product is optionally shifted right by 64 bits, depending on the \emph{highmult}. The resulting double word is subtracted from and stored into the \emph{registerW}.


\paragraph{Modifier(s)}
\noindent\begin{itemize}
\item highmult (cf. highmult in section \ref{par:modifier-highmult} on page \pageref{par:modifier-highmult})
\end{itemize}

\paragraph{Execution} {Reservation table \textsf{ALU\_LITE} (Section~\ref{sec:reservation-ALU-LITE})}
~
{\begin{lstlisting}
stage ID:
  new highmult = _ZX_2(highmult);
stage RR:
  new argument3 = GPR[registerY];
  new argument2 = GPR[registerZ];
stage E1:
  new argument1 = GPR[registerW];
  if ((highmult == 0) || (highmult == 3)) {
    new product = _SX_64(argument2) * _SX_64(argument3);
  }
  if (highmult == 1) {
    product = _ZX_64(argument2) * _ZX_64(argument3);
  }
  if (highmult == 2) {
    product = _SX_64(argument2) * _ZX_64(argument3);
  }
  if (highmult) {
    new result1 = argument1 - product.64[1];
  } else {
    result1 = argument1 - product.64[0];
  }
stage E2:
  GPR[registerW] = result1;
\end{lstlisting}}
\newpage

\paragraph{Encoding} Format \textsf{ALU\_DMWRRR.M} (Section~\ref{sec:format-ALU-DMWRRR.M}), \verb|alucode2 = 11|\\

\bitfields{ 1/P, 2/00, 2/-- --, 27/upper27, 1/1, 2/11, 1/0, 2/11, 2/highmult, 6/registerW, 2/10, 3/110, 1/0, 1/splat32, 5/lower5, 6/registerZ }


\paragraph{Syntax} \texttt{msbfd}\emph{highmult} \emph{registerW} = \emph{registerZ}, \emph{upper27\_\discretionary{}{}{}lower5}\emph{splat32}

\paragraph{Description} The \emph{upper27\_\discretionary{}{}{}lower5} and the \emph{registerZ} double words are sign-extended or zero-extended from 64 bits, depending on the \emph{highmult}. The \emph{upper27\_\discretionary{}{}{}lower5} double word is multiplied by the \emph{registerZ} double word, and the product is optionally shifted right by 64 bits, depending on the \emph{highmult}. The resulting double word is subtracted from and stored into the \emph{registerW}.


\paragraph{Modifier(s)}
\noindent\begin{itemize}
\item highmult (cf. highmult in section \ref{par:modifier-highmult} on page \pageref{par:modifier-highmult})
\item splat32 (cf. splat32 in section \ref{par:modifier-splat32} on page \pageref{par:modifier-splat32})
\end{itemize}

\paragraph{Execution} {Reservation table \textsf{ALU\_LITE.X} (Section~\ref{sec:reservation-ALU-LITE.X})}
~
{\begin{lstlisting}
stage ID:
  new splat32 = _ZX_1(splat32);
  new highmult = _ZX_2(highmult);
stage RR:
  new argument3 = splat32 ? (_WX_32(upper27_lower5) << 32) | _ZX_32(_WX_32(upper27_lower5)) : _SX_32(_WX_32(upper27_lower5));
  new argument2 = GPR[registerZ];
stage E1:
  new argument1 = GPR[registerW];
  if ((highmult == 0) || (highmult == 3)) {
    new product = _SX_64(argument2) * _SX_64(argument3);
  }
  if (highmult == 1) {
    product = _ZX_64(argument2) * _ZX_64(argument3);
  }
  if (highmult == 2) {
    product = _SX_64(argument2) * _ZX_64(argument3);
  }
  if (highmult) {
    new result1 = argument1 - product.64[1];
  } else {
    result1 = argument1 - product.64[0];
  }
stage E2:
  GPR[registerW] = result1;
\end{lstlisting}}
\newpage
\subsubsection[{\small MSBFDQ: Multiply Subtract From Extend Double Word to Quadruple Word}]{MSBFDQ\hfill Multiply Subtract From Extend Double Word to Quadruple Word}
\label{instr:MSBFDQ}\index{msbfdq}
 
\paragraph{Encoding} Format \textsf{ALU\_QMEWRRR} (Section~\ref{sec:format-ALU-QMEWRRR}), \verb|alucode2 = 11|\\

\bitfields{ 1/P, 2/11, 1/0, 2/11, 2/widemult, 5/registerM, 1/--, 2/10, 3/111, 1/1, 6/registerY, 6/registerZ }


\paragraph{Syntax} \texttt{msbfdq}\emph{widemult} \emph{registerM} = \emph{registerZ}, \emph{registerY}

\paragraph{Description} The \emph{registerY} and the \emph{registerZ} double words are sign-extended or zero-extended from 64 bits, depending on the \emph{widemult}. The \emph{registerY} is multiplied by the \emph{registerZ}, producing a 128-bit result. The resulting quadruple word subtracted from and stored into the \emph{registerM}.


\paragraph{Modifier(s)}
\noindent\begin{itemize}
\item widemult (cf. widemult in section \ref{par:modifier-widemult} on page \pageref{par:modifier-widemult})
\end{itemize}

\paragraph{Execution} {Reservation table \textsf{ALU\_LITE} (Section~\ref{sec:reservation-ALU-LITE})}
~
{\begin{lstlisting}
stage ID:
  new widemult = _ZX_2(widemult);
stage RR:
  new argument3 = GPR[registerY];
  new argument2 = GPR[registerZ];
stage E1:
  new argument1 = PGR[registerM];
  if (widemult == 0) {
    new product = _SX_64(argument2) * _SX_64(argument3);
  }
  if (widemult == 1) {
    product = _ZX_64(argument2) * _ZX_64(argument3);
  }
  if (widemult == 2) {
    product = _SX_64(argument2) * _ZX_64(argument3);
  }
  new result1 = argument1 - product;
stage E2:
  PGR[registerM] = result1;
\end{lstlisting}}
\newpage
\subsubsection[{\small MSBFHQ: Multiply Subtract From Half Word Quadruple}]{MSBFHQ\hfill Multiply Subtract From Half Word Quadruple}
\label{instr:MSBFHQ}\index{msbfhq}
 
\paragraph{Encoding} Format \textsf{ALU\_HQMWRRRE} (Section~\ref{sec:format-ALU-HQMWRRRE}), \verb|alucode2 = 11|\\

\bitfields{ 1/P, 2/11, 1/1, 2/11, 2/highmult, 5/registerWe, 1/0, 2/10, 3/110, 1/1, 5/registerYe, 1/0, 5/registerZe, 1/0 }


\paragraph{Syntax} \texttt{msbfhq}\emph{highmult} \emph{registerWe} = \emph{registerZe}, \emph{registerYe}

\paragraph{Description} The \emph{registerYe} and the \emph{registerZe} are interpreted as four half words packed into 64 bits. The \emph{registerYe} and the \emph{registerZe} half words are sign-extended or zero-extended from 16 bits, depending on the \emph{highmult}. The \emph{registerYe} half words are multiplied by the \emph{registerZe} half words, and the product is optionally shifted right by 16 bits, depending on the \emph{highmult}. The four resulting half words are subtracted from the \emph{registerWe} half words, packed and stored back into the \emph{registerWe}.


\paragraph{Modifier(s)}
\noindent\begin{itemize}
\item highmult (cf. highmult in section \ref{par:modifier-highmult} on page \pageref{par:modifier-highmult})
\end{itemize}

\paragraph{Execution} {Reservation table \textsf{ALU\_LITE} (Section~\ref{sec:reservation-ALU-LITE})}
~
{\begin{lstlisting}
stage ID:
  new highmult = _ZX_2(highmult);
stage RR:
  new argument3 = GPR[registerYe<<1];
  new argument2 = GPR[registerZe<<1];
stage E1:
  new argument1 = GPR[registerWe<<1];
  for (I = 0; I < 4; I += 1) {
    if ((highmult == 0) || (highmult == 3)) {
      new product = _SX_16(argument2.16[I]) * _SX_16(argument3.16[I]);
    }
    if (highmult == 1) {
      product = _ZX_16(argument2.16[I]) * _ZX_16(argument3.16[I]);
    }
    if (highmult == 2) {
      product = _SX_16(argument2.16[I]) * _ZX_16(argument3.16[I]);
    }
    if (highmult) {
      result1.16[I] = argument1.16[I] - product.16[1];
    } else {
      result1.16[I] = argument1.16[I] - product.16[0];
    }
  };
stage E2:
  GPR[registerWe<<1] = result1;
\end{lstlisting}}
\newpage

\paragraph{Encoding} Format \textsf{ALU\_HQMWRRRO} (Section~\ref{sec:format-ALU-HQMWRRRO}), \verb|alucode2 = 11|\\

\bitfields{ 1/P, 2/11, 1/1, 2/11, 2/highmult, 5/registerWo, 1/1, 2/10, 3/110, 1/1, 5/registerYo, 1/0, 5/registerZo, 1/0 }


\paragraph{Syntax} \texttt{msbfhq}\emph{highmult} \emph{registerWo} = \emph{registerZo}, \emph{registerYo}

\paragraph{Description} The \emph{registerYo} and the \emph{registerZo} are interpreted as four half words packed into 64 bits. The \emph{registerYo} and the \emph{registerZo} half words are sign-extended or zero-extended from 16 bits, depending on the \emph{highmult}. The \emph{registerYo} half words are multiplied by the \emph{registerZo} half words, and the product is optionally shifted right by 16 bits, depending on the \emph{highmult}. The four resulting half words are subtracted from the \emph{registerWo} half words, packed and stored back into the \emph{registerWo}.


\paragraph{Modifier(s)}
\noindent\begin{itemize}
\item highmult (cf. highmult in section \ref{par:modifier-highmult} on page \pageref{par:modifier-highmult})
\end{itemize}

\paragraph{Execution} {Reservation table \textsf{ALU\_LITE} (Section~\ref{sec:reservation-ALU-LITE})}
~
{\begin{lstlisting}
stage ID:
  new highmult = _ZX_2(highmult);
stage RR:
  new argument3 = GPR[(registerYo<<1)+1];
  new argument2 = GPR[(registerZo<<1)+1];
stage E1:
  new argument1 = GPR[(registerWo<<1)+1];
  for (I = 0; I < 4; I += 1) {
    if ((highmult == 0) || (highmult == 3)) {
      new product = _SX_16(argument2.16[I]) * _SX_16(argument3.16[I]);
    }
    if (highmult == 1) {
      product = _ZX_16(argument2.16[I]) * _ZX_16(argument3.16[I]);
    }
    if (highmult == 2) {
      product = _SX_16(argument2.16[I]) * _ZX_16(argument3.16[I]);
    }
    if (highmult) {
      result1.16[I] = argument1.16[I] - product.16[1];
    } else {
      result1.16[I] = argument1.16[I] - product.16[0];
    }
  };
stage E2:
  GPR[(registerWo<<1)+1] = result1;
\end{lstlisting}}
\newpage
\subsubsection[{\small MSBFHWQ: Multiply Subtract From Extend Half Word to Word Quadruple}]{MSBFHWQ\hfill Multiply Subtract From Extend Half Word to Word Quadruple}
\label{instr:MSBFHWQ}\index{msbfhwq}
 
\paragraph{Encoding} Format \textsf{ALU\_WQMEWRRRE} (Section~\ref{sec:format-ALU-WQMEWRRRE}), \verb|alucode2 = 11|\\

\bitfields{ 1/P, 2/11, 1/1, 2/11, 2/widemult, 5/registerM, 1/0, 2/10, 3/110, 1/0, 5/registerYe, 1/1, 5/registerZe, 1/1 }


\paragraph{Syntax} \texttt{msbfhwq}\emph{widemult} \emph{registerM} = \emph{registerZe}, \emph{registerYe}

\paragraph{Description} The \emph{registerYe} and the \emph{registerZe} are interpreted as eight half words packed into 128 bits. These \emph{registerYe} and \emph{registerZe} half words are sign-extended or zero-extended from 16 bits, depending on the \emph{widemult}. The \emph{registerYe} half words are multiplied by the \emph{registerZe} half words, producing four 32-bit results. The four resulting words are packed, subtracted from and stored into the \emph{registerM}.


\paragraph{Modifier(s)}
\noindent\begin{itemize}
\item widemult (cf. widemult in section \ref{par:modifier-widemult} on page \pageref{par:modifier-widemult})
\end{itemize}

\paragraph{Execution} {Reservation table \textsf{ALU\_LITE} (Section~\ref{sec:reservation-ALU-LITE})}
~
{\begin{lstlisting}
stage ID:
  new widemult = _ZX_2(widemult);
stage RR:
  new argument3 = GPR[registerYe<<1];
  new argument2 = GPR[registerZe<<1];
stage E1:
  new argument1 = PGR[registerM];
  for (I = 0; I < 4; I += 1) {
    if (widemult == 0) {
      new product = _SX_16(argument2.16[I]) * _SX_16(argument3.16[I]);
    }
    if (widemult == 1) {
      product = _ZX_16(argument2.16[I]) * _ZX_16(argument3.16[I]);
    }
    if (widemult == 2) {
      product = _SX_16(argument2.16[I]) * _ZX_16(argument3.16[I]);
    }
    result1.32[I] = argument1.32[I] - product
  };
stage E2:
  PGR[registerM] = result1;
\end{lstlisting}}
\newpage

\paragraph{Encoding} Format \textsf{ALU\_WQMEWRRRO} (Section~\ref{sec:format-ALU-WQMEWRRRO}), \verb|alucode2 = 11|\\

\bitfields{ 1/P, 2/11, 1/1, 2/11, 2/widemult, 5/registerM, 1/1, 2/10, 3/110, 1/0, 5/registerYo, 1/1, 5/registerZo, 1/1 }


\paragraph{Syntax} \texttt{msbfhwq}\emph{widemult} \emph{registerM} = \emph{registerZo}, \emph{registerYo}

\paragraph{Description} The \emph{registerYo} and the \emph{registerZo} are interpreted as eight half words packed into 128 bits. These \emph{registerYo} and \emph{registerZo} half words are sign-extended or zero-extended from 16 bits, depending on the \emph{widemult}. The \emph{registerYo} half words are multiplied by the \emph{registerZo} half words, producing four 32-bit results. The four resulting words are packed, subtracted from and stored into the \emph{registerM}.


\paragraph{Modifier(s)}
\noindent\begin{itemize}
\item widemult (cf. widemult in section \ref{par:modifier-widemult} on page \pageref{par:modifier-widemult})
\end{itemize}

\paragraph{Execution} {Reservation table \textsf{ALU\_LITE} (Section~\ref{sec:reservation-ALU-LITE})}
~
{\begin{lstlisting}
stage ID:
  new widemult = _ZX_2(widemult);
stage RR:
  new argument3 = GPR[(registerYo<<1)+1];
  new argument2 = GPR[(registerZo<<1)+1];
stage E1:
  new argument1 = PGR[registerM];
  for (I = 0; I < 4; I += 1) {
    if (widemult == 0) {
      new product = _SX_16(argument2.16[I]) * _SX_16(argument3.16[I]);
    }
    if (widemult == 1) {
      product = _ZX_16(argument2.16[I]) * _ZX_16(argument3.16[I]);
    }
    if (widemult == 2) {
      product = _SX_16(argument2.16[I]) * _ZX_16(argument3.16[I]);
    }
    result1.32[I] = argument1.32[I] - product
  };
stage E2:
  PGR[registerM] = result1;
\end{lstlisting}}
\newpage
\subsubsection[{\small MSBFW: Multiply Subtract From Word}]{MSBFW\hfill Multiply Subtract From Word}
\label{instr:MSBFW}\index{msbfw}
 
\paragraph{Encoding} Format \textsf{ALU\_WMWRRR} (Section~\ref{sec:format-ALU-WMWRRR}), \verb|alucode2 = 11|\\

\bitfields{ 1/P, 2/11, 1/0, 2/11, 2/highmult, 6/registerW, 2/11, 3/110, 1/signextw, 6/registerY, 6/registerZ }


\paragraph{Syntax} \texttt{msbfw}\emph{highmult}\emph{signextw} \emph{registerW} = \emph{registerZ}, \emph{registerY}

\paragraph{Description} The \emph{registerY} and the \emph{registerZ} words are sign-extended or zero-extended from 32 bits, depending on the \emph{highmult}. The \emph{registerY} word is multiplied by the \emph{registerZ} word, and the product is optionally shifted right by 32 bits, depending on the \emph{highmult}. The resulting word is subtracted from and stored into the \emph{registerW}.


\paragraph{Modifier(s)}
\noindent\begin{itemize}
\item highmult (cf. highmult in section \ref{par:modifier-highmult} on page \pageref{par:modifier-highmult})
\item signextw (cf. signextw in section \ref{par:modifier-signextw} on page \pageref{par:modifier-signextw})
\end{itemize}

\paragraph{Execution} {Reservation table \textsf{ALU\_LITE} (Section~\ref{sec:reservation-ALU-LITE})}
~
{\begin{lstlisting}
stage ID:
  new signextw = _ZX_1(signextw);
  new highmult = _ZX_2(highmult);
stage RR:
  new argument3 = GPR[registerY];
  new argument2 = GPR[registerZ];
stage E1:
  new argument1 = GPR[registerW];
  if ((highmult == 0) || (highmult == 3)) {
    new product = _SX_32(argument2) * _SX_32(argument3);
  }
  if (highmult == 1) {
    product = _ZX_32(argument2) * _ZX_32(argument3);
  }
  if (highmult == 2) {
    product = _SX_32(argument2) * _ZX_32(argument3);
  }
  if (highmult) {
    new result1 = argument1 - product.32[1];
  } else {
    result1 = argument1 - product.32[0];
  }
stage E2:
  GPR[registerW] = signextw ? _SX_32(result1) : _ZX_32(result1);
\end{lstlisting}}
\newpage

\paragraph{Encoding} Format \textsf{ALU\_WMWRRR.W} (Section~\ref{sec:format-ALU-WMWRRR.W}), \verb|alucode2 = 11|\\

\bitfields{ 1/P, 2/00, 2/-- --, 27/upper27, 1/1, 2/11, 1/0, 2/11, 2/highmult, 6/registerW, 2/11, 3/110, 1/signextw, 1/0, 5/lower5, 6/registerZ }


\paragraph{Syntax} \texttt{msbfw}\emph{highmult}\emph{signextw} \emph{registerW} = \emph{registerZ}, \emph{upper27\_\discretionary{}{}{}lower5}

\paragraph{Description} The \emph{upper27\_\discretionary{}{}{}lower5} and the \emph{registerZ} words are sign-extended or zero-extended from 32 bits, depending on the \emph{highmult}. The \emph{upper27\_\discretionary{}{}{}lower5} word is multiplied by the \emph{registerZ} word, and the product is optionally shifted right by 32 bits, depending on the \emph{highmult}. The resulting word is subtracted from and stored into the \emph{registerW}.


\paragraph{Modifier(s)}
\noindent\begin{itemize}
\item highmult (cf. highmult in section \ref{par:modifier-highmult} on page \pageref{par:modifier-highmult})
\item signextw (cf. signextw in section \ref{par:modifier-signextw} on page \pageref{par:modifier-signextw})
\end{itemize}

\paragraph{Execution} {Reservation table \textsf{ALU\_LITE.X} (Section~\ref{sec:reservation-ALU-LITE.X})}
~
{\begin{lstlisting}
stage ID:
  new signextw = _ZX_1(signextw);
  new highmult = _ZX_2(highmult);
stage RR:
  new argument3 = _WX_32(upper27_lower5);
  new argument2 = GPR[registerZ];
stage E1:
  new argument1 = GPR[registerW];
  if ((highmult == 0) || (highmult == 3)) {
    new product = _SX_32(argument2) * _SX_32(argument3);
  }
  if (highmult == 1) {
    product = _ZX_32(argument2) * _ZX_32(argument3);
  }
  if (highmult == 2) {
    product = _SX_32(argument2) * _ZX_32(argument3);
  }
  if (highmult) {
    new result1 = argument1 - product.32[1];
  } else {
    result1 = argument1 - product.32[0];
  }
stage E2:
  GPR[registerW] = signextw ? _SX_32(result1) : _ZX_32(result1);
\end{lstlisting}}
\newpage
\subsubsection[{\small MSBFWD: Multiply Subtract From Extend Word to Double Word}]{MSBFWD\hfill Multiply Subtract From Extend Word to Double Word}
\label{instr:MSBFWD}\index{msbfwd}
 
\paragraph{Encoding} Format \textsf{ALU\_DMEWRRR} (Section~\ref{sec:format-ALU-DMEWRRR}), \verb|alucode2 = 11|\\

\bitfields{ 1/P, 2/11, 1/0, 2/11, 2/widemult, 6/registerW, 2/10, 3/111, 1/0, 6/registerY, 6/registerZ }


\paragraph{Syntax} \texttt{msbfwd}\emph{widemult} \emph{registerW} = \emph{registerZ}, \emph{registerY}

\paragraph{Description} The \emph{registerY} and the \emph{registerZ} words are sign-extended or zero-extended from 32 bits, depending on the \emph{widemult}. The \emph{registerY} is multiplied by the \emph{registerZ}, producing a 64-bit result. The resulting double word is subtracted from and stored into the \emph{registerW}.


\paragraph{Modifier(s)}
\noindent\begin{itemize}
\item widemult (cf. widemult in section \ref{par:modifier-widemult} on page \pageref{par:modifier-widemult})
\end{itemize}

\paragraph{Execution} {Reservation table \textsf{ALU\_LITE} (Section~\ref{sec:reservation-ALU-LITE})}
~
{\begin{lstlisting}
stage ID:
  new widemult = _ZX_2(widemult);
stage RR:
  new argument3 = GPR[registerY];
  new argument2 = GPR[registerZ];
stage E1:
  new argument1 = GPR[registerW];
  if (widemult == 0) {
    new product = _SX_32(argument2) * _SX_32(argument3);
  }
  if (widemult == 1) {
    product = _ZX_32(argument2) * _ZX_32(argument3);
  }
  if (widemult == 2) {
    product = _SX_32(argument2) * _ZX_32(argument3);
  }
  new result1 = argument1 - product;
stage E2:
  GPR[registerW] = result1;
\end{lstlisting}}
\newpage

\paragraph{Encoding} Format \textsf{ALU\_DMEWRRR.M} (Section~\ref{sec:format-ALU-DMEWRRR.M}), \verb|alucode2 = 11|\\

\bitfields{ 1/P, 2/00, 2/-- --, 27/upper27, 1/1, 2/11, 1/0, 2/11, 2/widemult, 6/registerW, 2/10, 3/111, 1/0, 1/splat32, 5/lower5, 6/registerZ }


\paragraph{Syntax} \texttt{msbfwd}\emph{widemult} \emph{registerW} = \emph{registerZ}, \emph{upper27\_\discretionary{}{}{}lower5}\emph{splat32}

\paragraph{Description} The \emph{upper27\_\discretionary{}{}{}lower5} and the \emph{registerZ} words are sign-extended or zero-extended from 32 bits, depending on the \emph{widemult}. The \emph{upper27\_\discretionary{}{}{}lower5} is multiplied by the \emph{registerZ}, producing a 64-bit result. The resulting double word is subtracted from and stored into the \emph{registerW}.


\paragraph{Modifier(s)}
\noindent\begin{itemize}
\item widemult (cf. widemult in section \ref{par:modifier-widemult} on page \pageref{par:modifier-widemult})
\item splat32 (cf. splat32 in section \ref{par:modifier-splat32} on page \pageref{par:modifier-splat32})
\end{itemize}

\paragraph{Execution} {Reservation table \textsf{ALU\_LITE.X} (Section~\ref{sec:reservation-ALU-LITE.X})}
~
{\begin{lstlisting}
stage ID:
  new splat32 = _ZX_1(splat32);
  new widemult = _ZX_2(widemult);
stage RR:
  new argument3 = splat32 ? (_WX_32(upper27_lower5) << 32) | _ZX_32(_WX_32(upper27_lower5)) : _SX_32(_WX_32(upper27_lower5));
  new argument2 = GPR[registerZ];
stage E1:
  new argument1 = GPR[registerW];
  if (widemult == 0) {
    new product = _SX_32(argument2) * _SX_32(argument3);
  }
  if (widemult == 1) {
    product = _ZX_32(argument2) * _ZX_32(argument3);
  }
  if (widemult == 2) {
    product = _SX_32(argument2) * _ZX_32(argument3);
  }
  new result1 = argument1 - product;
stage E2:
  GPR[registerW] = result1;
\end{lstlisting}}
\newpage
\subsubsection[{\small MSBFWDP: Multiply Subtract From Extend Word to Double Word Pair}]{MSBFWDP\hfill Multiply Subtract From Extend Word to Double Word Pair}
\label{instr:MSBFWDP}\index{msbfwdp}
 
\paragraph{Encoding} Format \textsf{ALU\_DPMEWRRRE} (Section~\ref{sec:format-ALU-DPMEWRRRE}), \verb|alucode2 = 11|\\

\bitfields{ 1/P, 2/11, 1/1, 2/11, 2/widemult, 5/registerM, 1/0, 2/10, 3/110, 1/0, 5/registerYe, 1/1, 5/registerZe, 1/0 }


\paragraph{Syntax} \texttt{msbfwdp}\emph{widemult} \emph{registerM} = \emph{registerZe}, \emph{registerYe}

\paragraph{Description} The \emph{registerYe} and the \emph{registerZe} are interpreted as two words packed into 65 bits. These \emph{registerYe} and \emph{registerZe} words are sign-extended or zero-extended from 32 bits, depending on the \emph{widemult}. The \emph{registerYe} words are multiplied by the \emph{registerZe} words, producing two 64-bit results. The two resulting double words are packed, subtracted from and stored into the \emph{registerM}.


\paragraph{Modifier(s)}
\noindent\begin{itemize}
\item widemult (cf. widemult in section \ref{par:modifier-widemult} on page \pageref{par:modifier-widemult})
\end{itemize}

\paragraph{Execution} {Reservation table \textsf{ALU\_LITE} (Section~\ref{sec:reservation-ALU-LITE})}
~
{\begin{lstlisting}
stage ID:
  new widemult = _ZX_2(widemult);
stage RR:
  new argument3 = GPR[registerYe<<1];
  new argument2 = GPR[registerZe<<1];
stage E1:
  new argument1 = PGR[registerM];
  for (I = 0; I < 2; I += 1) {
    if (widemult == 0) {
      new product = _SX_32(argument2.64[I]) * _SX_32(argument3.64[I]);
    }
    if (widemult == 1) {
      product = _ZX_32(argument2.64[I]) * _ZX_32(argument3.64[I]);
    }
    if (widemult == 2) {
      product = _SX_32(argument2.64[I]) * _ZX_32(argument3.64[I]);
    }
    result1.64[I] = argument1.64[I] - product
  };
stage E2:
  PGR[registerM] = result1;
\end{lstlisting}}
\newpage

\paragraph{Encoding} Format \textsf{ALU\_DPMEWRRRO} (Section~\ref{sec:format-ALU-DPMEWRRRO}), \verb|alucode2 = 11|\\

\bitfields{ 1/P, 2/11, 1/1, 2/11, 2/widemult, 5/registerM, 1/1, 2/10, 3/110, 1/0, 5/registerYo, 1/1, 5/registerZo, 1/0 }


\paragraph{Syntax} \texttt{msbfwdp}\emph{widemult} \emph{registerM} = \emph{registerZo}, \emph{registerYo}

\paragraph{Description} The \emph{registerYo} and the \emph{registerZo} are interpreted as two words packed into 65 bits. These \emph{registerYo} and \emph{registerZo} words are sign-extended or zero-extended from 32 bits, depending on the \emph{widemult}. The \emph{registerYo} words are multiplied by the \emph{registerZo} words, producing two 64-bit results. The two resulting double words are packed, subtracted from and stored into the \emph{registerM}.


\paragraph{Modifier(s)}
\noindent\begin{itemize}
\item widemult (cf. widemult in section \ref{par:modifier-widemult} on page \pageref{par:modifier-widemult})
\end{itemize}

\paragraph{Execution} {Reservation table \textsf{ALU\_LITE} (Section~\ref{sec:reservation-ALU-LITE})}
~
{\begin{lstlisting}
stage ID:
  new widemult = _ZX_2(widemult);
stage RR:
  new argument3 = GPR[(registerYo<<1)+1];
  new argument2 = GPR[(registerZo<<1)+1];
stage E1:
  new argument1 = PGR[registerM];
  for (I = 0; I < 2; I += 1) {
    if (widemult == 0) {
      new product = _SX_32(argument2.64[I]) * _SX_32(argument3.64[I]);
    }
    if (widemult == 1) {
      product = _ZX_32(argument2.64[I]) * _ZX_32(argument3.64[I]);
    }
    if (widemult == 2) {
      product = _SX_32(argument2.64[I]) * _ZX_32(argument3.64[I]);
    }
    result1.64[I] = argument1.64[I] - product
  };
stage E2:
  PGR[registerM] = result1;
\end{lstlisting}}
\newpage
\subsubsection[{\small MSBFWP: Multiply Subtract From Word Pair}]{MSBFWP\hfill Multiply Subtract From Word Pair}
\label{instr:MSBFWP}\index{msbfwp}
 
\paragraph{Encoding} Format \textsf{ALU\_WPMWRRRE} (Section~\ref{sec:format-ALU-WPMWRRRE}), \verb|alucode2 = 11|\\

\bitfields{ 1/P, 2/11, 1/1, 2/11, 2/highmult, 5/registerWe, 1/0, 2/10, 3/110, 1/0, 5/registerYe, 1/0, 5/registerZe, 1/1 }


\paragraph{Syntax} \texttt{msbfwp}\emph{highmult} \emph{registerWe} = \emph{registerZe}, \emph{registerYe}

\paragraph{Description} The \emph{registerYe} and the \emph{registerZe} are interpreted as two words packed into 64 bits. The \emph{registerYe} and the \emph{registerZe} words are sign-extended or zero-extended from 32 bits, depending on the \emph{highmult}. The \emph{registerYe} words are multiplied by the \emph{registerZe} words, and the product is optionally shifted right by 32 bits, depending on the \emph{highmult}. The two resulting words are subtracted from the \emph{registerWe} words, packed and stored back into the \emph{registerWe}.


\paragraph{Modifier(s)}
\noindent\begin{itemize}
\item highmult (cf. highmult in section \ref{par:modifier-highmult} on page \pageref{par:modifier-highmult})
\end{itemize}

\paragraph{Execution} {Reservation table \textsf{ALU\_LITE} (Section~\ref{sec:reservation-ALU-LITE})}
~
{\begin{lstlisting}
stage ID:
  new highmult = _ZX_2(highmult);
stage RR:
  new argument3 = GPR[registerYe<<1];
  new argument2 = GPR[registerZe<<1];
stage E1:
  new argument1 = GPR[registerWe<<1];
  for (I = 0; I < 2; I += 1) {
    if ((highmult == 0) || (highmult == 3)) {
      new product = _SX_32(argument2.32[I]) * _SX_32(argument3.32[I]);
    }
    if (highmult == 1) {
      product = _ZX_32(argument2.32[I]) * _ZX_32(argument3.32[I]);
    }
    if (highmult == 2) {
      product = _SX_32(argument2.32[I]) * _ZX_32(argument3.32[I]);
    }
    if (highmult) {
      result1.32[I] = argument1.32[I] - product.32[1];
    } else {
      result1.32[I] = argument1.32[I] - product.32[0];
    }
  };
stage E2:
  GPR[registerWe<<1] = result1;
\end{lstlisting}}
\newpage

\paragraph{Encoding} Format \textsf{ALU\_WPMWRRRO} (Section~\ref{sec:format-ALU-WPMWRRRO}), \verb|alucode2 = 11|\\

\bitfields{ 1/P, 2/11, 1/1, 2/11, 2/highmult, 5/registerWo, 1/1, 2/10, 3/110, 1/0, 5/registerYo, 1/0, 5/registerZo, 1/1 }


\paragraph{Syntax} \texttt{msbfwp}\emph{highmult} \emph{registerWo} = \emph{registerZo}, \emph{registerYo}

\paragraph{Description} The \emph{registerYo} and the \emph{registerZo} are interpreted as two words packed into 64 bits. The \emph{registerYo} and the \emph{registerZo} words are sign-extended or zero-extended from 32 bits, depending on the \emph{highmult}. The \emph{registerYo} words are multiplied by the \emph{registerZo} words, and the product is optionally shifted right by 32 bits, depending on the \emph{highmult}. The two resulting words are subtracted from the \emph{registerWo} words, packed and stored back into the \emph{registerWo}.


\paragraph{Modifier(s)}
\noindent\begin{itemize}
\item highmult (cf. highmult in section \ref{par:modifier-highmult} on page \pageref{par:modifier-highmult})
\end{itemize}

\paragraph{Execution} {Reservation table \textsf{ALU\_LITE} (Section~\ref{sec:reservation-ALU-LITE})}
~
{\begin{lstlisting}
stage ID:
  new highmult = _ZX_2(highmult);
stage RR:
  new argument3 = GPR[(registerYo<<1)+1];
  new argument2 = GPR[(registerZo<<1)+1];
stage E1:
  new argument1 = GPR[(registerWo<<1)+1];
  for (I = 0; I < 2; I += 1) {
    if ((highmult == 0) || (highmult == 3)) {
      new product = _SX_32(argument2.32[I]) * _SX_32(argument3.32[I]);
    }
    if (highmult == 1) {
      product = _ZX_32(argument2.32[I]) * _ZX_32(argument3.32[I]);
    }
    if (highmult == 2) {
      product = _SX_32(argument2.32[I]) * _ZX_32(argument3.32[I]);
    }
    if (highmult) {
      result1.32[I] = argument1.32[I] - product.32[1];
    } else {
      result1.32[I] = argument1.32[I] - product.32[0];
    }
  };
stage E2:
  GPR[(registerWo<<1)+1] = result1;
\end{lstlisting}}
\newpage
\subsubsection[{\small MSBFXBHO: Multiply Subtract From Extend Byte to Half Word Octuple}]{MSBFXBHO\hfill Multiply Subtract From Extend Byte to Half Word Octuple}
\label{instr:MSBFXBHO}\index{msbfxbho}
 
\paragraph{Encoding} Format \textsf{ALU\_HOMXWRRR} (Section~\ref{sec:format-ALU-HOMXWRRR}), \verb|alucode2 = 11|\\

\bitfields{ 1/P, 2/11, 1/1, 2/11, 2/widemult, 5/registerM, 1/oddlanes, 2/10, 3/111, 1/1, 5/registerO, 1/--, 5/registerP, 1/0 }


\paragraph{Syntax} \texttt{msbfxbho}\emph{oddlanes}\emph{widemult} \emph{registerM} = \emph{registerP}, \emph{registerO}

\paragraph{Description} The \emph{registerO} and the \emph{registerP} are interpreted as sixteen bytes packed into 128 bits. The even or odd bytes of the \emph{registerO} and the \emph{registerP} are selected, depending on \emph{oddlanes}. These \emph{registerO} and \emph{registerP} bytes are sign-extended or zero-extended from 8 bits, depending on the \emph{widemult}. The \emph{registerO} bytes are multiplied by the \emph{registerP} bytes, producing eight 16-bit results. The eight resulting half words are packed, subtracted from and stored into the \emph{registerM}.


\paragraph{Modifier(s)}
\noindent\begin{itemize}
\item widemult (cf. widemult in section \ref{par:modifier-widemult} on page \pageref{par:modifier-widemult})
\item oddlanes (cf. oddlanes in section \ref{par:modifier-oddlanes} on page \pageref{par:modifier-oddlanes})
\end{itemize}

\paragraph{Execution} {Reservation table \textsf{ALU\_LITE} (Section~\ref{sec:reservation-ALU-LITE})}
~
{\begin{lstlisting}
stage ID:
  new oddlanes = _ZX_1(oddlanes);
  new widemult = _ZX_2(widemult);
stage RR:
  new argument3 = PGR[registerO];
  new argument2 = PGR[registerP];
stage E1:
  new argument1 = PGR[registerM];
  if (oddlanes) {
    new argument2 = argument2 >> 8;
    new argument3 = argument3 >> 8;
  }
  for (I = 0; I < 8; I += 1) {
    if (widemult == 0) {
      new product = _SX_8(argument2.16[I]) * _SX_8(argument3.16[I]);
    }
    if (widemult == 1) {
      product = _ZX_8(argument2.16[I]) * _ZX_8(argument3.16[I]);
    }
    if (widemult == 2) {
      product = _SX_8(argument2.16[I]) * _ZX_8(argument3.16[I]);
    }
    result1.16[I] = argument1.16[I] - product
  };
stage E2:
  PGR[registerM] = result1;
\end{lstlisting}}
\newpage
\subsubsection[{\small MSBFXHWQ: Multiply Subtract From Extend Half Word to Word Quadruple}]{MSBFXHWQ\hfill Multiply Subtract From Extend Half Word to Word Quadruple}
\label{instr:MSBFXHWQ}\index{msbfxhwq}
 
\paragraph{Encoding} Format \textsf{ALU\_WQMXWRRR} (Section~\ref{sec:format-ALU-WQMXWRRR}), \verb|alucode2 = 11|\\

\bitfields{ 1/P, 2/11, 1/1, 2/11, 2/widemult, 5/registerM, 1/oddlanes, 2/10, 3/111, 1/0, 5/registerO, 1/--, 5/registerP, 1/1 }


\paragraph{Syntax} \texttt{msbfxhwq}\emph{oddlanes}\emph{widemult} \emph{registerM} = \emph{registerP}, \emph{registerO}

\paragraph{Description} The \emph{registerO} and the \emph{registerP} are interpreted as eight half words packed into 128 bits. The even or odd half words of the \emph{registerO} and the \emph{registerP} are selected, depending on \emph{oddlanes}. These \emph{registerO} and \emph{registerP} half words are sign-extended or zero-extended from 16 bits, depending on the \emph{widemult}. The \emph{registerO} half words are multiplied by the \emph{registerP} half words, producing four 32-bit results. The four resulting words are packed, subtracted from and stored into the \emph{registerM}.


\paragraph{Modifier(s)}
\noindent\begin{itemize}
\item widemult (cf. widemult in section \ref{par:modifier-widemult} on page \pageref{par:modifier-widemult})
\item oddlanes (cf. oddlanes in section \ref{par:modifier-oddlanes} on page \pageref{par:modifier-oddlanes})
\end{itemize}

\paragraph{Execution} {Reservation table \textsf{ALU\_LITE} (Section~\ref{sec:reservation-ALU-LITE})}
~
{\begin{lstlisting}
stage ID:
  new oddlanes = _ZX_1(oddlanes);
  new widemult = _ZX_2(widemult);
stage RR:
  new argument3 = PGR[registerO];
  new argument2 = PGR[registerP];
stage E1:
  new argument1 = PGR[registerM];
  if (oddlanes) {
    new argument2 = argument2 >> 16;
    new argument3 = argument3 >> 16;
  }
  for (I = 0; I < 4; I += 1) {
    if (widemult == 0) {
      new product = _SX_16(argument2.32[I]) * _SX_16(argument3.32[I]);
    }
    if (widemult == 1) {
      product = _ZX_16(argument2.32[I]) * _ZX_16(argument3.32[I]);
    }
    if (widemult == 2) {
      product = _SX_16(argument2.32[I]) * _ZX_16(argument3.32[I]);
    }
    result1.32[I] = argument1.32[I] - product
  };
stage E2:
  PGR[registerM] = result1;
\end{lstlisting}}
\newpage
\subsubsection[{\small MSBFXWDP: Multiply Subtract From Extend Word to Double Word Pair}]{MSBFXWDP\hfill Multiply Subtract From Extend Word to Double Word Pair}
\label{instr:MSBFXWDP}\index{msbfxwdp}
 
\paragraph{Encoding} Format \textsf{ALU\_DPMXWRRR} (Section~\ref{sec:format-ALU-DPMXWRRR}), \verb|alucode2 = 11|\\

\bitfields{ 1/P, 2/11, 1/1, 2/11, 2/widemult, 5/registerM, 1/oddlanes, 2/10, 3/111, 1/0, 5/registerO, 1/--, 5/registerP, 1/0 }


\paragraph{Syntax} \texttt{msbfxwdp}\emph{oddlanes}\emph{widemult} \emph{registerM} = \emph{registerP}, \emph{registerO}

\paragraph{Description} The \emph{registerO} and the \emph{registerP} are interpreted as four words packed into 128 bits. The even or odd words of the \emph{registerO} and the \emph{registerP} are selected, depending on \emph{oddlanes}. These \emph{registerO} and \emph{registerP} words are sign-extended or zero-extended from 32 bits, depending on the \emph{widemult}. The \emph{registerO} words are multiplied by the \emph{registerP} words, producing two 64-bit results. The two resulting double words are packed, subtracted from and stored into the \emph{registerM}.


\paragraph{Modifier(s)}
\noindent\begin{itemize}
\item widemult (cf. widemult in section \ref{par:modifier-widemult} on page \pageref{par:modifier-widemult})
\item oddlanes (cf. oddlanes in section \ref{par:modifier-oddlanes} on page \pageref{par:modifier-oddlanes})
\end{itemize}

\paragraph{Execution} {Reservation table \textsf{ALU\_LITE} (Section~\ref{sec:reservation-ALU-LITE})}
~
{\begin{lstlisting}
stage ID:
  new oddlanes = _ZX_1(oddlanes);
  new widemult = _ZX_2(widemult);
stage RR:
  new argument3 = PGR[registerO];
  new argument2 = PGR[registerP];
stage E1:
  new argument1 = PGR[registerM];
  if (oddlanes) {
    new argument2 = argument2 >> 32;
    new argument3 = argument3 >> 32;
  }
  for (I = 0; I < 2; I += 1) {
    if (widemult == 0) {
      new product = _SX_32(argument2.64[I]) * _SX_32(argument3.64[I]);
    }
    if (widemult == 1) {
      product = _ZX_32(argument2.64[I]) * _ZX_32(argument3.64[I]);
    }
    if (widemult == 2) {
      product = _SX_32(argument2.64[I]) * _ZX_32(argument3.64[I]);
    }
    result1.64[I] = argument1.64[I] - product
  };
stage E2:
  PGR[registerM] = result1;
\end{lstlisting}}
\newpage
\subsubsection[{\small MULBHO: Multiply Extend Byte to Half Word Octuple}]{MULBHO\hfill Multiply Extend Byte to Half Word Octuple}
\label{instr:MULBHO}\index{mulbho}
 
\paragraph{Encoding} Format \textsf{ALU\_HOMEWRRE} (Section~\ref{sec:format-ALU-HOMEWRRE}), \verb|alucode2 = 00|\\

\bitfields{ 1/P, 2/11, 1/1, 2/00, 2/widemult, 5/registerM, 1/0, 2/10, 3/110, 1/1, 5/registerYe, 1/1, 5/registerZe, 1/0 }


\paragraph{Syntax} \texttt{mulbho}\emph{widemult} \emph{registerM} = \emph{registerZe}, \emph{registerYe}

\paragraph{Description} The \emph{registerYe} and the \emph{registerZe} are interpreted as eight bytes packed into 64 bits. These \emph{registerYe} and \emph{registerZe} bytes are sign-extended or zero-extended from 8 bits, depending on the \emph{widemult}. The \emph{registerYe} bytes are multiplied by the \emph{registerZe} bytes, producing eight 16-bit results. The eight resulting half words are packed and stored into the \emph{registerM}.


\paragraph{Modifier(s)}
\noindent\begin{itemize}
\item widemult (cf. widemult in section \ref{par:modifier-widemult} on page \pageref{par:modifier-widemult})
\end{itemize}

\paragraph{Execution} {Reservation table \textsf{ALU\_LITE} (Section~\ref{sec:reservation-ALU-LITE})}
~
{\begin{lstlisting}
stage ID:
  new widemult = _ZX_2(widemult);
stage RR:
  new argument3 = GPR[registerYe<<1];
  new argument2 = GPR[registerZe<<1];
  for (I = 0; I < 8; I += 1) {
    if (widemult == 0) {
      new product = _SX_8(argument2.8[I]) * _SX_8(argument3.8[I]);
    }
    if (widemult == 1) {
      product = _ZX_8(argument2.8[I]) * _ZX_8(argument3.8[I]);
    }
    if (widemult == 2) {
      product = _SX_8(argument2.8[I]) * _ZX_8(argument3.8[I]);
    }
    result1.16[I] = product
  };
stage E2:
  PGR[registerM] = result1;
\end{lstlisting}}
\newpage

\paragraph{Encoding} Format \textsf{ALU\_HOMEWRRO} (Section~\ref{sec:format-ALU-HOMEWRRO}), \verb|alucode2 = 00|\\

\bitfields{ 1/P, 2/11, 1/1, 2/00, 2/widemult, 5/registerM, 1/1, 2/10, 3/110, 1/1, 5/registerYo, 1/1, 5/registerZo, 1/0 }


\paragraph{Syntax} \texttt{mulbho}\emph{widemult} \emph{registerM} = \emph{registerZo}, \emph{registerYo}

\paragraph{Description} The \emph{registerYo} and the \emph{registerZo} are interpreted as eight bytes packed into 64 bits. These \emph{registerYo} and \emph{registerZo} bytes are sign-extended or zero-extended from 8 bits, depending on the \emph{widemult}. The \emph{registerYo} bytes are multiplied by the \emph{registerZo} bytes, producing eight 16-bit results. The eight resulting half words are packed and stored into the \emph{registerM}.


\paragraph{Modifier(s)}
\noindent\begin{itemize}
\item widemult (cf. widemult in section \ref{par:modifier-widemult} on page \pageref{par:modifier-widemult})
\end{itemize}

\paragraph{Execution} {Reservation table \textsf{ALU\_LITE} (Section~\ref{sec:reservation-ALU-LITE})}
~
{\begin{lstlisting}
stage ID:
  new widemult = _ZX_2(widemult);
stage RR:
  new argument3 = GPR[(registerYo<<1)+1];
  new argument2 = GPR[(registerZo<<1)+1];
  for (I = 0; I < 8; I += 1) {
    if (widemult == 0) {
      new product = _SX_8(argument2.8[I]) * _SX_8(argument3.8[I]);
    }
    if (widemult == 1) {
      product = _ZX_8(argument2.8[I]) * _ZX_8(argument3.8[I]);
    }
    if (widemult == 2) {
      product = _SX_8(argument2.8[I]) * _ZX_8(argument3.8[I]);
    }
    result1.16[I] = product
  };
stage E2:
  PGR[registerM] = result1;
\end{lstlisting}}
\newpage
\subsubsection[{\small MULD: Multiply Double Word}]{MULD\hfill Multiply Double Word}
\label{instr:MULD}\index{muld}
 
\paragraph{Encoding} Format \textsf{ALU\_DMWRR} (Section~\ref{sec:format-ALU-DMWRR}), \verb|alucode2 = 00|\\

\bitfields{ 1/P, 2/11, 1/0, 2/00, 2/highmult, 6/registerW, 2/10, 3/110, 1/0, 6/registerY, 6/registerZ }


\paragraph{Syntax} \texttt{muld}\emph{highmult} \emph{registerW} = \emph{registerZ}, \emph{registerY}

\paragraph{Description} The \emph{registerY} and the \emph{registerZ} double words are sign-extended or zero-extended from 64 bits, depending on the \emph{highmult}. The \emph{registerY} double word is multiplied by the \emph{registerZ} double word, and the product is optionally shifted right by 64 bits, depending on the \emph{highmult}. The resulting double word is stored into the \emph{registerW}.


\paragraph{Modifier(s)}
\noindent\begin{itemize}
\item highmult (cf. highmult in section \ref{par:modifier-highmult} on page \pageref{par:modifier-highmult})
\end{itemize}

\paragraph{Execution} {Reservation table \textsf{ALU\_LITE} (Section~\ref{sec:reservation-ALU-LITE})}
~
{\begin{lstlisting}
stage ID:
  new highmult = _ZX_2(highmult);
stage RR:
  new argument3 = GPR[registerY];
  new argument2 = GPR[registerZ];
  if ((highmult == 0) || (highmult == 3)) {
    new product = _SX_64(argument2) * _SX_64(argument3);
  }
  if (highmult == 1) {
    product = _ZX_64(argument2) * _ZX_64(argument3);
  }
  if (highmult == 2) {
    product = _SX_64(argument2) * _ZX_64(argument3);
  }
  if (highmult) {
    new result1 = product.64[1];
  } else {
    result1 = product.64[0];
  }
stage E2:
  GPR[registerW] = result1;
\end{lstlisting}}
\newpage

\paragraph{Encoding} Format \textsf{ALU\_DMWRR.M} (Section~\ref{sec:format-ALU-DMWRR.M}), \verb|alucode2 = 00|\\

\bitfields{ 1/P, 2/00, 2/-- --, 27/upper27, 1/1, 2/11, 1/0, 2/00, 2/highmult, 6/registerW, 2/10, 3/110, 1/0, 1/splat32, 5/lower5, 6/registerZ }


\paragraph{Syntax} \texttt{muld}\emph{highmult} \emph{registerW} = \emph{registerZ}, \emph{upper27\_\discretionary{}{}{}lower5}\emph{splat32}

\paragraph{Description} The \emph{upper27\_\discretionary{}{}{}lower5} and the \emph{registerZ} double words are sign-extended or zero-extended from 64 bits, depending on the \emph{highmult}. The \emph{upper27\_\discretionary{}{}{}lower5} double word is multiplied by the \emph{registerZ} double word, and the product is optionally shifted right by 64 bits, depending on the \emph{highmult}. The resulting double word is stored into the \emph{registerW}.


\paragraph{Modifier(s)}
\noindent\begin{itemize}
\item highmult (cf. highmult in section \ref{par:modifier-highmult} on page \pageref{par:modifier-highmult})
\item splat32 (cf. splat32 in section \ref{par:modifier-splat32} on page \pageref{par:modifier-splat32})
\end{itemize}

\paragraph{Execution} {Reservation table \textsf{ALU\_LITE.X} (Section~\ref{sec:reservation-ALU-LITE.X})}
~
{\begin{lstlisting}
stage ID:
  new splat32 = _ZX_1(splat32);
  new highmult = _ZX_2(highmult);
stage RR:
  new argument3 = splat32 ? (_WX_32(upper27_lower5) << 32) | _ZX_32(_WX_32(upper27_lower5)) : _SX_32(_WX_32(upper27_lower5));
  new argument2 = GPR[registerZ];
  if ((highmult == 0) || (highmult == 3)) {
    new product = _SX_64(argument2) * _SX_64(argument3);
  }
  if (highmult == 1) {
    product = _ZX_64(argument2) * _ZX_64(argument3);
  }
  if (highmult == 2) {
    product = _SX_64(argument2) * _ZX_64(argument3);
  }
  if (highmult) {
    new result1 = product.64[1];
  } else {
    result1 = product.64[0];
  }
stage E2:
  GPR[registerW] = result1;
\end{lstlisting}}
\newpage
\subsubsection[{\small MULDQ: Multiply Extend Double Word to Quadruple Word}]{MULDQ\hfill Multiply Extend Double Word to Quadruple Word}
\label{instr:MULDQ}\index{muldq}
 
\paragraph{Encoding} Format \textsf{ALU\_QMEWRR} (Section~\ref{sec:format-ALU-QMEWRR}), \verb|alucode2 = 00|\\

\bitfields{ 1/P, 2/11, 1/0, 2/00, 2/widemult, 5/registerM, 1/--, 2/10, 3/111, 1/1, 6/registerY, 6/registerZ }


\paragraph{Syntax} \texttt{muldq}\emph{widemult} \emph{registerM} = \emph{registerZ}, \emph{registerY}

\paragraph{Description} The \emph{registerY} and the \emph{registerZ} double words are sign-extended or zero-extended from 64 bits, depending on the \emph{widemult}. The \emph{registerY} is multiplied by the \emph{registerZ}, producing a 128-bit result. The resulting quadruple word is stored into the \emph{registerM}.


\paragraph{Modifier(s)}
\noindent\begin{itemize}
\item widemult (cf. widemult in section \ref{par:modifier-widemult} on page \pageref{par:modifier-widemult})
\end{itemize}

\paragraph{Execution} {Reservation table \textsf{ALU\_LITE} (Section~\ref{sec:reservation-ALU-LITE})}
~
{\begin{lstlisting}
stage ID:
  new widemult = _ZX_2(widemult);
stage RR:
  new argument3 = GPR[registerY];
  new argument2 = GPR[registerZ];
  if (widemult == 0) {
    new product = _SX_64(argument2) * _SX_64(argument3);
  }
  if (widemult == 1) {
    product = _ZX_64(argument2) * _ZX_64(argument3);
  }
  if (widemult == 2) {
    product = _SX_64(argument2) * _ZX_64(argument3);
  }
  new result1 = product;
stage E2:
  PGR[registerM] = result1;
\end{lstlisting}}
\newpage
\subsubsection[{\small MULHQ: Multiply Half Word Quadruple}]{MULHQ\hfill Multiply Half Word Quadruple}
\label{instr:MULHQ}\index{mulhq}
 
\paragraph{Encoding} Format \textsf{ALU\_HQMWRRE} (Section~\ref{sec:format-ALU-HQMWRRE}), \verb|alucode2 = 00|\\

\bitfields{ 1/P, 2/11, 1/1, 2/00, 2/highmult, 5/registerWe, 1/0, 2/10, 3/110, 1/1, 5/registerYe, 1/0, 5/registerZe, 1/0 }


\paragraph{Syntax} \texttt{mulhq}\emph{highmult} \emph{registerWe} = \emph{registerZe}, \emph{registerYe}

\paragraph{Description} The \emph{registerYe} and the \emph{registerZe} are interpreted as four half words packed into 64 bits. The \emph{registerYe} and the \emph{registerZe} half words are sign-extended or zero-extended from 16 bits, depending on the \emph{highmult}. The \emph{registerYe} half words are multiplied by the \emph{registerZe} half words, and the product is optionally shifted right by 16 bits, depending on the \emph{highmult}. The four resulting half words are packed and stored back into the \emph{registerWe}.


\paragraph{Modifier(s)}
\noindent\begin{itemize}
\item highmult (cf. highmult in section \ref{par:modifier-highmult} on page \pageref{par:modifier-highmult})
\end{itemize}

\paragraph{Execution} {Reservation table \textsf{ALU\_LITE} (Section~\ref{sec:reservation-ALU-LITE})}
~
{\begin{lstlisting}
stage ID:
  new highmult = _ZX_2(highmult);
stage RR:
  new argument3 = GPR[registerYe<<1];
  new argument2 = GPR[registerZe<<1];
  for (I = 0; I < 4; I += 1) {
    if ((highmult == 0) || (highmult == 3)) {
      new product = _SX_16(argument2.16[I]) * _SX_16(argument3.16[I]);
    }
    if (highmult == 1) {
      product = _ZX_16(argument2.16[I]) * _ZX_16(argument3.16[I]);
    }
    if (highmult == 2) {
      product = _SX_16(argument2.16[I]) * _ZX_16(argument3.16[I]);
    }
    if (highmult) {
      result1.16[I] = product.16[1];
    } else {
      result1.16[I] = product.16[0];
    }
  };
stage E2:
  GPR[registerWe<<1] = result1;
\end{lstlisting}}
\newpage

\paragraph{Encoding} Format \textsf{ALU\_HQMWRRO} (Section~\ref{sec:format-ALU-HQMWRRO}), \verb|alucode2 = 00|\\

\bitfields{ 1/P, 2/11, 1/1, 2/00, 2/highmult, 5/registerWo, 1/1, 2/10, 3/110, 1/1, 5/registerYo, 1/0, 5/registerZo, 1/0 }


\paragraph{Syntax} \texttt{mulhq}\emph{highmult} \emph{registerWo} = \emph{registerZo}, \emph{registerYo}

\paragraph{Description} The \emph{registerYo} and the \emph{registerZo} are interpreted as four half words packed into 64 bits. The \emph{registerYo} and the \emph{registerZo} half words are sign-extended or zero-extended from 16 bits, depending on the \emph{highmult}. The \emph{registerYo} half words are multiplied by the \emph{registerZo} half words, and the product is optionally shifted right by 16 bits, depending on the \emph{highmult}. The four resulting half words are packed and stored back into the \emph{registerWo}.


\paragraph{Modifier(s)}
\noindent\begin{itemize}
\item highmult (cf. highmult in section \ref{par:modifier-highmult} on page \pageref{par:modifier-highmult})
\end{itemize}

\paragraph{Execution} {Reservation table \textsf{ALU\_LITE} (Section~\ref{sec:reservation-ALU-LITE})}
~
{\begin{lstlisting}
stage ID:
  new highmult = _ZX_2(highmult);
stage RR:
  new argument3 = GPR[(registerYo<<1)+1];
  new argument2 = GPR[(registerZo<<1)+1];
  for (I = 0; I < 4; I += 1) {
    if ((highmult == 0) || (highmult == 3)) {
      new product = _SX_16(argument2.16[I]) * _SX_16(argument3.16[I]);
    }
    if (highmult == 1) {
      product = _ZX_16(argument2.16[I]) * _ZX_16(argument3.16[I]);
    }
    if (highmult == 2) {
      product = _SX_16(argument2.16[I]) * _ZX_16(argument3.16[I]);
    }
    if (highmult) {
      result1.16[I] = product.16[1];
    } else {
      result1.16[I] = product.16[0];
    }
  };
stage E2:
  GPR[(registerWo<<1)+1] = result1;
\end{lstlisting}}
\newpage
\subsubsection[{\small MULHWQ: Multiply Extend Half Word to Word Quadruple}]{MULHWQ\hfill Multiply Extend Half Word to Word Quadruple}
\label{instr:MULHWQ}\index{mulhwq}
 
\paragraph{Encoding} Format \textsf{ALU\_WQMEWRRE} (Section~\ref{sec:format-ALU-WQMEWRRE}), \verb|alucode2 = 00|\\

\bitfields{ 1/P, 2/11, 1/1, 2/00, 2/widemult, 5/registerM, 1/0, 2/10, 3/110, 1/0, 5/registerYe, 1/1, 5/registerZe, 1/1 }


\paragraph{Syntax} \texttt{mulhwq}\emph{widemult} \emph{registerM} = \emph{registerZe}, \emph{registerYe}

\paragraph{Description} The \emph{registerYe} and the \emph{registerZe} are interpreted as four half words packed into 64 bits. These \emph{registerYe} and \emph{registerZe} half words are sign-extended or zero-extended from 16 bits, depending on the \emph{widemult}. The \emph{registerYe} half words are multiplied by the \emph{registerZe} half words, producing four 32-bit results. The four resulting words are packed and stored into the \emph{registerM}.


\paragraph{Modifier(s)}
\noindent\begin{itemize}
\item widemult (cf. widemult in section \ref{par:modifier-widemult} on page \pageref{par:modifier-widemult})
\end{itemize}

\paragraph{Execution} {Reservation table \textsf{ALU\_LITE} (Section~\ref{sec:reservation-ALU-LITE})}
~
{\begin{lstlisting}
stage ID:
  new widemult = _ZX_2(widemult);
stage RR:
  new argument3 = GPR[registerYe<<1];
  new argument2 = GPR[registerZe<<1];
  for (I = 0; I < 4; I += 1) {
    if (widemult == 0) {
      new product = _SX_16(argument2.16[I]) * _SX_16(argument3.16[I]);
    }
    if (widemult == 1) {
      product = _ZX_16(argument2.16[I]) * _ZX_16(argument3.16[I]);
    }
    if (widemult == 2) {
      product = _SX_16(argument2.16[I]) * _ZX_16(argument3.16[I]);
    }
    result1.32[I] = product
  };
stage E2:
  PGR[registerM] = result1;
\end{lstlisting}}
\newpage

\paragraph{Encoding} Format \textsf{ALU\_WQMEWRRO} (Section~\ref{sec:format-ALU-WQMEWRRO}), \verb|alucode2 = 00|\\

\bitfields{ 1/P, 2/11, 1/1, 2/00, 2/widemult, 5/registerM, 1/1, 2/10, 3/110, 1/0, 5/registerYo, 1/1, 5/registerZo, 1/1 }


\paragraph{Syntax} \texttt{mulhwq}\emph{widemult} \emph{registerM} = \emph{registerZo}, \emph{registerYo}

\paragraph{Description} The \emph{registerYo} and the \emph{registerZo} are interpreted as four half words packed into 64 bits. These \emph{registerYo} and \emph{registerZo} half words are sign-extended or zero-extended from 16 bits, depending on the \emph{widemult}. The \emph{registerYo} half words are multiplied by the \emph{registerZo} half words, producing four 32-bit results. The four resulting words are packed and stored into the \emph{registerM}.


\paragraph{Modifier(s)}
\noindent\begin{itemize}
\item widemult (cf. widemult in section \ref{par:modifier-widemult} on page \pageref{par:modifier-widemult})
\end{itemize}

\paragraph{Execution} {Reservation table \textsf{ALU\_LITE} (Section~\ref{sec:reservation-ALU-LITE})}
~
{\begin{lstlisting}
stage ID:
  new widemult = _ZX_2(widemult);
stage RR:
  new argument3 = GPR[(registerYo<<1)+1];
  new argument2 = GPR[(registerZo<<1)+1];
  for (I = 0; I < 4; I += 1) {
    if (widemult == 0) {
      new product = _SX_16(argument2.16[I]) * _SX_16(argument3.16[I]);
    }
    if (widemult == 1) {
      product = _ZX_16(argument2.16[I]) * _ZX_16(argument3.16[I]);
    }
    if (widemult == 2) {
      product = _SX_16(argument2.16[I]) * _ZX_16(argument3.16[I]);
    }
    result1.32[I] = product
  };
stage E2:
  PGR[registerM] = result1;
\end{lstlisting}}
\newpage
\subsubsection[{\small MULNBHO: Multiply Negate Extend Byte to Half Word Octuple}]{MULNBHO\hfill Multiply Negate Extend Byte to Half Word Octuple}
\label{instr:MULNBHO}\index{mulnbho}
 
\paragraph{Encoding} Format \textsf{ALU\_HOMEWRRE} (Section~\ref{sec:format-ALU-HOMEWRRE}), \verb|alucode2 = 01|\\

\bitfields{ 1/P, 2/11, 1/1, 2/01, 2/widemult, 5/registerM, 1/0, 2/10, 3/110, 1/1, 5/registerYe, 1/1, 5/registerZe, 1/0 }


\paragraph{Syntax} \texttt{mulnbho}\emph{widemult} \emph{registerM} = \emph{registerZe}, \emph{registerYe}

\paragraph{Description} The \emph{registerYe} and the \emph{registerZe} are interpreted as eight bytes packed into 64 bits. These \emph{registerYe} and \emph{registerZe} bytes are sign-extended or zero-extended from 8 bits, depending on the \emph{widemult}. The \emph{registerYe} bytes are multiplied by the \emph{registerZe} bytes, producing eight 16-bit results. The eight resulting half words are negated, packed and stored into the \emph{registerM}.


\paragraph{Modifier(s)}
\noindent\begin{itemize}
\item widemult (cf. widemult in section \ref{par:modifier-widemult} on page \pageref{par:modifier-widemult})
\end{itemize}

\paragraph{Execution} {Reservation table \textsf{ALU\_LITE} (Section~\ref{sec:reservation-ALU-LITE})}
~
{\begin{lstlisting}
stage ID:
  new widemult = _ZX_2(widemult);
stage RR:
  new argument3 = GPR[registerYe<<1];
  new argument2 = GPR[registerZe<<1];
  for (I = 0; I < 8; I += 1) {
    if (widemult == 0) {
      new product = _SX_8(argument2.8[I]) * _SX_8(argument3.8[I]);
    }
    if (widemult == 1) {
      product = _ZX_8(argument2.8[I]) * _ZX_8(argument3.8[I]);
    }
    if (widemult == 2) {
      product = _SX_8(argument2.8[I]) * _ZX_8(argument3.8[I]);
    }
    result1.16[I] = -product
  };
stage E2:
  PGR[registerM] = result1;
\end{lstlisting}}
\newpage

\paragraph{Encoding} Format \textsf{ALU\_HOMEWRRO} (Section~\ref{sec:format-ALU-HOMEWRRO}), \verb|alucode2 = 01|\\

\bitfields{ 1/P, 2/11, 1/1, 2/01, 2/widemult, 5/registerM, 1/1, 2/10, 3/110, 1/1, 5/registerYo, 1/1, 5/registerZo, 1/0 }


\paragraph{Syntax} \texttt{mulnbho}\emph{widemult} \emph{registerM} = \emph{registerZo}, \emph{registerYo}

\paragraph{Description} The \emph{registerYo} and the \emph{registerZo} are interpreted as eight bytes packed into 64 bits. These \emph{registerYo} and \emph{registerZo} bytes are sign-extended or zero-extended from 8 bits, depending on the \emph{widemult}. The \emph{registerYo} bytes are multiplied by the \emph{registerZo} bytes, producing eight 16-bit results. The eight resulting half words are negated, packed and stored into the \emph{registerM}.


\paragraph{Modifier(s)}
\noindent\begin{itemize}
\item widemult (cf. widemult in section \ref{par:modifier-widemult} on page \pageref{par:modifier-widemult})
\end{itemize}

\paragraph{Execution} {Reservation table \textsf{ALU\_LITE} (Section~\ref{sec:reservation-ALU-LITE})}
~
{\begin{lstlisting}
stage ID:
  new widemult = _ZX_2(widemult);
stage RR:
  new argument3 = GPR[(registerYo<<1)+1];
  new argument2 = GPR[(registerZo<<1)+1];
  for (I = 0; I < 8; I += 1) {
    if (widemult == 0) {
      new product = _SX_8(argument2.8[I]) * _SX_8(argument3.8[I]);
    }
    if (widemult == 1) {
      product = _ZX_8(argument2.8[I]) * _ZX_8(argument3.8[I]);
    }
    if (widemult == 2) {
      product = _SX_8(argument2.8[I]) * _ZX_8(argument3.8[I]);
    }
    result1.16[I] = -product
  };
stage E2:
  PGR[registerM] = result1;
\end{lstlisting}}
\newpage
\subsubsection[{\small MULND: Multiply Negate Double Word}]{MULND\hfill Multiply Negate Double Word}
\label{instr:MULND}\index{mulnd}
 
\paragraph{Encoding} Format \textsf{ALU\_DMWRR} (Section~\ref{sec:format-ALU-DMWRR}), \verb|alucode2 = 01|\\

\bitfields{ 1/P, 2/11, 1/0, 2/01, 2/highmult, 6/registerW, 2/10, 3/110, 1/0, 6/registerY, 6/registerZ }


\paragraph{Syntax} \texttt{mulnd}\emph{highmult} \emph{registerW} = \emph{registerZ}, \emph{registerY}

\paragraph{Description} The \emph{registerY} and the \emph{registerZ} double words are sign-extended or zero-extended from 64 bits, depending on the \emph{highmult}. The \emph{registerY} double word is multiplied by the \emph{registerZ} double word, and the product is optionally shifted right by 64 bits, depending on the \emph{highmult}. The resulting double word is negated and stored into the \emph{registerW}.


\paragraph{Modifier(s)}
\noindent\begin{itemize}
\item highmult (cf. highmult in section \ref{par:modifier-highmult} on page \pageref{par:modifier-highmult})
\end{itemize}

\paragraph{Execution} {Reservation table \textsf{ALU\_LITE} (Section~\ref{sec:reservation-ALU-LITE})}
~
{\begin{lstlisting}
stage ID:
  new highmult = _ZX_2(highmult);
stage RR:
  new argument3 = GPR[registerY];
  new argument2 = GPR[registerZ];
  if ((highmult == 0) || (highmult == 3)) {
    new product = _SX_64(argument2) * _SX_64(argument3);
  }
  if (highmult == 1) {
    product = _ZX_64(argument2) * _ZX_64(argument3);
  }
  if (highmult == 2) {
    product = _SX_64(argument2) * _ZX_64(argument3);
  }
  if (highmult) {
    new result1 = -product.64[1];
  } else {
    result1 = -product.64[0];
  }
stage E2:
  GPR[registerW] = result1;
\end{lstlisting}}
\newpage

\paragraph{Encoding} Format \textsf{ALU\_DMWRR.M} (Section~\ref{sec:format-ALU-DMWRR.M}), \verb|alucode2 = 01|\\

\bitfields{ 1/P, 2/00, 2/-- --, 27/upper27, 1/1, 2/11, 1/0, 2/01, 2/highmult, 6/registerW, 2/10, 3/110, 1/0, 1/splat32, 5/lower5, 6/registerZ }


\paragraph{Syntax} \texttt{mulnd}\emph{highmult} \emph{registerW} = \emph{registerZ}, \emph{upper27\_\discretionary{}{}{}lower5}\emph{splat32}

\paragraph{Description} The \emph{upper27\_\discretionary{}{}{}lower5} and the \emph{registerZ} double words are sign-extended or zero-extended from 64 bits, depending on the \emph{highmult}. The \emph{upper27\_\discretionary{}{}{}lower5} double word is multiplied by the \emph{registerZ} double word, and the product is optionally shifted right by 64 bits, depending on the \emph{highmult}. The resulting double word is negated and stored into the \emph{registerW}.


\paragraph{Modifier(s)}
\noindent\begin{itemize}
\item highmult (cf. highmult in section \ref{par:modifier-highmult} on page \pageref{par:modifier-highmult})
\item splat32 (cf. splat32 in section \ref{par:modifier-splat32} on page \pageref{par:modifier-splat32})
\end{itemize}

\paragraph{Execution} {Reservation table \textsf{ALU\_LITE.X} (Section~\ref{sec:reservation-ALU-LITE.X})}
~
{\begin{lstlisting}
stage ID:
  new splat32 = _ZX_1(splat32);
  new highmult = _ZX_2(highmult);
stage RR:
  new argument3 = splat32 ? (_WX_32(upper27_lower5) << 32) | _ZX_32(_WX_32(upper27_lower5)) : _SX_32(_WX_32(upper27_lower5));
  new argument2 = GPR[registerZ];
  if ((highmult == 0) || (highmult == 3)) {
    new product = _SX_64(argument2) * _SX_64(argument3);
  }
  if (highmult == 1) {
    product = _ZX_64(argument2) * _ZX_64(argument3);
  }
  if (highmult == 2) {
    product = _SX_64(argument2) * _ZX_64(argument3);
  }
  if (highmult) {
    new result1 = -product.64[1];
  } else {
    result1 = -product.64[0];
  }
stage E2:
  GPR[registerW] = result1;
\end{lstlisting}}
\newpage
\subsubsection[{\small MULNDQ: Multiply Negate Extend Double Word to Quadruple Word}]{MULNDQ\hfill Multiply Negate Extend Double Word to Quadruple Word}
\label{instr:MULNDQ}\index{mulndq}
 
\paragraph{Encoding} Format \textsf{ALU\_QMEWRR} (Section~\ref{sec:format-ALU-QMEWRR}), \verb|alucode2 = 01|\\

\bitfields{ 1/P, 2/11, 1/0, 2/01, 2/widemult, 5/registerM, 1/--, 2/10, 3/111, 1/1, 6/registerY, 6/registerZ }


\paragraph{Syntax} \texttt{mulndq}\emph{widemult} \emph{registerM} = \emph{registerZ}, \emph{registerY}

\paragraph{Description} The \emph{registerY} and the \emph{registerZ} double words are sign-extended or zero-extended from 64 bits, depending on the \emph{widemult}. The \emph{registerY} is multiplied by the \emph{registerZ}, producing a 128-bit result. The resulting quadruple word is negated and stored into the \emph{registerM}.


\paragraph{Modifier(s)}
\noindent\begin{itemize}
\item widemult (cf. widemult in section \ref{par:modifier-widemult} on page \pageref{par:modifier-widemult})
\end{itemize}

\paragraph{Execution} {Reservation table \textsf{ALU\_LITE} (Section~\ref{sec:reservation-ALU-LITE})}
~
{\begin{lstlisting}
stage ID:
  new widemult = _ZX_2(widemult);
stage RR:
  new argument3 = GPR[registerY];
  new argument2 = GPR[registerZ];
  if (widemult == 0) {
    new product = _SX_64(argument2) * _SX_64(argument3);
  }
  if (widemult == 1) {
    product = _ZX_64(argument2) * _ZX_64(argument3);
  }
  if (widemult == 2) {
    product = _SX_64(argument2) * _ZX_64(argument3);
  }
  new result1 = -product;
stage E2:
  PGR[registerM] = result1;
\end{lstlisting}}
\newpage
\subsubsection[{\small MULNHQ: Multiply Negate Half Word Quadruple}]{MULNHQ\hfill Multiply Negate Half Word Quadruple}
\label{instr:MULNHQ}\index{mulnhq}
 
\paragraph{Encoding} Format \textsf{ALU\_HQMWRRE} (Section~\ref{sec:format-ALU-HQMWRRE}), \verb|alucode2 = 01|\\

\bitfields{ 1/P, 2/11, 1/1, 2/01, 2/highmult, 5/registerWe, 1/0, 2/10, 3/110, 1/1, 5/registerYe, 1/0, 5/registerZe, 1/0 }


\paragraph{Syntax} \texttt{mulnhq}\emph{highmult} \emph{registerWe} = \emph{registerZe}, \emph{registerYe}

\paragraph{Description} The \emph{registerYe} and the \emph{registerZe} are interpreted as four half words packed into 64 bits. The \emph{registerYe} and the \emph{registerZe} half words are sign-extended or zero-extended from 16 bits, depending on the \emph{highmult}. The \emph{registerYe} half words are multiplied by the \emph{registerZe} half words, and the product is optionally shifted right by 16 bits, depending on the \emph{highmult}. The four resulting half words are negated, packed and stored back into the \emph{registerWe}.


\paragraph{Modifier(s)}
\noindent\begin{itemize}
\item highmult (cf. highmult in section \ref{par:modifier-highmult} on page \pageref{par:modifier-highmult})
\end{itemize}

\paragraph{Execution} {Reservation table \textsf{ALU\_LITE} (Section~\ref{sec:reservation-ALU-LITE})}
~
{\begin{lstlisting}
stage ID:
  new highmult = _ZX_2(highmult);
stage RR:
  new argument3 = GPR[registerYe<<1];
  new argument2 = GPR[registerZe<<1];
  for (I = 0; I < 4; I += 1) {
    if ((highmult == 0) || (highmult == 3)) {
      new product = _SX_16(argument2.16[I]) * _SX_16(argument3.16[I]);
    }
    if (highmult == 1) {
      product = _ZX_16(argument2.16[I]) * _ZX_16(argument3.16[I]);
    }
    if (highmult == 2) {
      product = _SX_16(argument2.16[I]) * _ZX_16(argument3.16[I]);
    }
    if (highmult) {
      result1.16[I] = -product.16[1];
    } else {
      result1.16[I] = -product.16[0];
    }
  };
stage E2:
  GPR[registerWe<<1] = result1;
\end{lstlisting}}
\newpage

\paragraph{Encoding} Format \textsf{ALU\_HQMWRRO} (Section~\ref{sec:format-ALU-HQMWRRO}), \verb|alucode2 = 01|\\

\bitfields{ 1/P, 2/11, 1/1, 2/01, 2/highmult, 5/registerWo, 1/1, 2/10, 3/110, 1/1, 5/registerYo, 1/0, 5/registerZo, 1/0 }


\paragraph{Syntax} \texttt{mulnhq}\emph{highmult} \emph{registerWo} = \emph{registerZo}, \emph{registerYo}

\paragraph{Description} The \emph{registerYo} and the \emph{registerZo} are interpreted as four half words packed into 64 bits. The \emph{registerYo} and the \emph{registerZo} half words are sign-extended or zero-extended from 16 bits, depending on the \emph{highmult}. The \emph{registerYo} half words are multiplied by the \emph{registerZo} half words, and the product is optionally shifted right by 16 bits, depending on the \emph{highmult}. The four resulting half words are negated, packed and stored back into the \emph{registerWo}.


\paragraph{Modifier(s)}
\noindent\begin{itemize}
\item highmult (cf. highmult in section \ref{par:modifier-highmult} on page \pageref{par:modifier-highmult})
\end{itemize}

\paragraph{Execution} {Reservation table \textsf{ALU\_LITE} (Section~\ref{sec:reservation-ALU-LITE})}
~
{\begin{lstlisting}
stage ID:
  new highmult = _ZX_2(highmult);
stage RR:
  new argument3 = GPR[(registerYo<<1)+1];
  new argument2 = GPR[(registerZo<<1)+1];
  for (I = 0; I < 4; I += 1) {
    if ((highmult == 0) || (highmult == 3)) {
      new product = _SX_16(argument2.16[I]) * _SX_16(argument3.16[I]);
    }
    if (highmult == 1) {
      product = _ZX_16(argument2.16[I]) * _ZX_16(argument3.16[I]);
    }
    if (highmult == 2) {
      product = _SX_16(argument2.16[I]) * _ZX_16(argument3.16[I]);
    }
    if (highmult) {
      result1.16[I] = -product.16[1];
    } else {
      result1.16[I] = -product.16[0];
    }
  };
stage E2:
  GPR[(registerWo<<1)+1] = result1;
\end{lstlisting}}
\newpage
\subsubsection[{\small MULNHWQ: Multiply Negate Extend Half Word to Word Quadruple}]{MULNHWQ\hfill Multiply Negate Extend Half Word to Word Quadruple}
\label{instr:MULNHWQ}\index{mulnhwq}
 
\paragraph{Encoding} Format \textsf{ALU\_WQMEWRRE} (Section~\ref{sec:format-ALU-WQMEWRRE}), \verb|alucode2 = 01|\\

\bitfields{ 1/P, 2/11, 1/1, 2/01, 2/widemult, 5/registerM, 1/0, 2/10, 3/110, 1/0, 5/registerYe, 1/1, 5/registerZe, 1/1 }


\paragraph{Syntax} \texttt{mulnhwq}\emph{widemult} \emph{registerM} = \emph{registerZe}, \emph{registerYe}

\paragraph{Description} The \emph{registerYe} and the \emph{registerZe} are interpreted as four half words packed into 64 bits. These \emph{registerYe} and \emph{registerZe} half words are sign-extended or zero-extended from 16 bits, depending on the \emph{widemult}. The \emph{registerYe} half words are multiplied by the \emph{registerZe} half words, producing four 32-bit results. The four resulting words are negated, packed and stored into the \emph{registerM}.


\paragraph{Modifier(s)}
\noindent\begin{itemize}
\item widemult (cf. widemult in section \ref{par:modifier-widemult} on page \pageref{par:modifier-widemult})
\end{itemize}

\paragraph{Execution} {Reservation table \textsf{ALU\_LITE} (Section~\ref{sec:reservation-ALU-LITE})}
~
{\begin{lstlisting}
stage ID:
  new widemult = _ZX_2(widemult);
stage RR:
  new argument3 = GPR[registerYe<<1];
  new argument2 = GPR[registerZe<<1];
  for (I = 0; I < 4; I += 1) {
    if (widemult == 0) {
      new product = _SX_16(argument2.16[I]) * _SX_16(argument3.16[I]);
    }
    if (widemult == 1) {
      product = _ZX_16(argument2.16[I]) * _ZX_16(argument3.16[I]);
    }
    if (widemult == 2) {
      product = _SX_16(argument2.16[I]) * _ZX_16(argument3.16[I]);
    }
    result1.32[I] = -product
  };
stage E2:
  PGR[registerM] = result1;
\end{lstlisting}}
\newpage

\paragraph{Encoding} Format \textsf{ALU\_WQMEWRRO} (Section~\ref{sec:format-ALU-WQMEWRRO}), \verb|alucode2 = 01|\\

\bitfields{ 1/P, 2/11, 1/1, 2/01, 2/widemult, 5/registerM, 1/1, 2/10, 3/110, 1/0, 5/registerYo, 1/1, 5/registerZo, 1/1 }


\paragraph{Syntax} \texttt{mulnhwq}\emph{widemult} \emph{registerM} = \emph{registerZo}, \emph{registerYo}

\paragraph{Description} The \emph{registerYo} and the \emph{registerZo} are interpreted as four half words packed into 64 bits. These \emph{registerYo} and \emph{registerZo} half words are sign-extended or zero-extended from 16 bits, depending on the \emph{widemult}. The \emph{registerYo} half words are multiplied by the \emph{registerZo} half words, producing four 32-bit results. The four resulting words are negated, packed and stored into the \emph{registerM}.


\paragraph{Modifier(s)}
\noindent\begin{itemize}
\item widemult (cf. widemult in section \ref{par:modifier-widemult} on page \pageref{par:modifier-widemult})
\end{itemize}

\paragraph{Execution} {Reservation table \textsf{ALU\_LITE} (Section~\ref{sec:reservation-ALU-LITE})}
~
{\begin{lstlisting}
stage ID:
  new widemult = _ZX_2(widemult);
stage RR:
  new argument3 = GPR[(registerYo<<1)+1];
  new argument2 = GPR[(registerZo<<1)+1];
  for (I = 0; I < 4; I += 1) {
    if (widemult == 0) {
      new product = _SX_16(argument2.16[I]) * _SX_16(argument3.16[I]);
    }
    if (widemult == 1) {
      product = _ZX_16(argument2.16[I]) * _ZX_16(argument3.16[I]);
    }
    if (widemult == 2) {
      product = _SX_16(argument2.16[I]) * _ZX_16(argument3.16[I]);
    }
    result1.32[I] = -product
  };
stage E2:
  PGR[registerM] = result1;
\end{lstlisting}}
\newpage
\subsubsection[{\small MULNW: Multiply Negate Word}]{MULNW\hfill Multiply Negate Word}
\label{instr:MULNW}\index{mulnw}
 
\paragraph{Encoding} Format \textsf{ALU\_WMWRR} (Section~\ref{sec:format-ALU-WMWRR}), \verb|alucode2 = 01|\\

\bitfields{ 1/P, 2/11, 1/0, 2/01, 2/highmult, 6/registerW, 2/11, 3/110, 1/signextw, 6/registerY, 6/registerZ }


\paragraph{Syntax} \texttt{mulnw}\emph{highmult}\emph{signextw} \emph{registerW} = \emph{registerZ}, \emph{registerY}

\paragraph{Description} The \emph{registerY} and the \emph{registerZ} words are sign-extended or zero-extended from 32 bits, depending on the \emph{highmult}. The \emph{registerY} word is multiplied by the \emph{registerZ} word, and the product is optionally shifted right by 32 bits, depending on the \emph{highmult}. The resulting word is negated and stored into the \emph{registerW}.


\paragraph{Modifier(s)}
\noindent\begin{itemize}
\item highmult (cf. highmult in section \ref{par:modifier-highmult} on page \pageref{par:modifier-highmult})
\item signextw (cf. signextw in section \ref{par:modifier-signextw} on page \pageref{par:modifier-signextw})
\end{itemize}

\paragraph{Execution} {Reservation table \textsf{ALU\_LITE} (Section~\ref{sec:reservation-ALU-LITE})}
~
{\begin{lstlisting}
stage ID:
  new signextw = _ZX_1(signextw);
  new highmult = _ZX_2(highmult);
stage RR:
  new argument3 = GPR[registerY];
  new argument2 = GPR[registerZ];
  if ((highmult == 0) || (highmult == 3)) {
    new product = _SX_32(argument2) * _SX_32(argument3);
  }
  if (highmult == 1) {
    product = _ZX_32(argument2) * _ZX_32(argument3);
  }
  if (highmult == 2) {
    product = _SX_32(argument2) * _ZX_32(argument3);
  }
  if (highmult) {
    new result1 = -product.32[1];
  } else {
    result1 = -product.32[0];
  }
stage E2:
  GPR[registerW] = signextw ? _SX_32(result1) : _ZX_32(result1);
\end{lstlisting}}
\newpage

\paragraph{Encoding} Format \textsf{ALU\_WMWRR.W} (Section~\ref{sec:format-ALU-WMWRR.W}), \verb|alucode2 = 01|\\

\bitfields{ 1/P, 2/00, 2/-- --, 27/upper27, 1/1, 2/11, 1/0, 2/01, 2/highmult, 6/registerW, 2/11, 3/110, 1/signextw, 1/0, 5/lower5, 6/registerZ }


\paragraph{Syntax} \texttt{mulnw}\emph{highmult}\emph{signextw} \emph{registerW} = \emph{registerZ}, \emph{upper27\_\discretionary{}{}{}lower5}

\paragraph{Description} The \emph{upper27\_\discretionary{}{}{}lower5} and the \emph{registerZ} words are sign-extended or zero-extended from 32 bits, depending on the \emph{highmult}. The \emph{upper27\_\discretionary{}{}{}lower5} word is multiplied by the \emph{registerZ} word, and the product is optionally shifted right by 32 bits, depending on the \emph{highmult}. The resulting word is negated and stored into the \emph{registerW}.


\paragraph{Modifier(s)}
\noindent\begin{itemize}
\item highmult (cf. highmult in section \ref{par:modifier-highmult} on page \pageref{par:modifier-highmult})
\item signextw (cf. signextw in section \ref{par:modifier-signextw} on page \pageref{par:modifier-signextw})
\end{itemize}

\paragraph{Execution} {Reservation table \textsf{ALU\_LITE.X} (Section~\ref{sec:reservation-ALU-LITE.X})}
~
{\begin{lstlisting}
stage ID:
  new signextw = _ZX_1(signextw);
  new highmult = _ZX_2(highmult);
stage RR:
  new argument3 = _WX_32(upper27_lower5);
  new argument2 = GPR[registerZ];
  if ((highmult == 0) || (highmult == 3)) {
    new product = _SX_32(argument2) * _SX_32(argument3);
  }
  if (highmult == 1) {
    product = _ZX_32(argument2) * _ZX_32(argument3);
  }
  if (highmult == 2) {
    product = _SX_32(argument2) * _ZX_32(argument3);
  }
  if (highmult) {
    new result1 = -product.32[1];
  } else {
    result1 = -product.32[0];
  }
stage E2:
  GPR[registerW] = signextw ? _SX_32(result1) : _ZX_32(result1);
\end{lstlisting}}
\newpage
\subsubsection[{\small MULNWD: Multiply Negate Extend Word to Double Word}]{MULNWD\hfill Multiply Negate Extend Word to Double Word}
\label{instr:MULNWD}\index{mulnwd}
 
\paragraph{Encoding} Format \textsf{ALU\_DMEWRR} (Section~\ref{sec:format-ALU-DMEWRR}), \verb|alucode2 = 01|\\

\bitfields{ 1/P, 2/11, 1/0, 2/01, 2/widemult, 6/registerW, 2/10, 3/111, 1/0, 6/registerY, 6/registerZ }


\paragraph{Syntax} \texttt{mulnwd}\emph{widemult} \emph{registerW} = \emph{registerZ}, \emph{registerY}

\paragraph{Description} The \emph{registerY} and the \emph{registerZ} words are sign-extended or zero-extended from 32 bits, depending on the \emph{widemult}. The \emph{registerY} is multiplied by the \emph{registerZ}, producing a 64-bit result. The resulting double word is negated and stored into the \emph{registerW}.


\paragraph{Modifier(s)}
\noindent\begin{itemize}
\item widemult (cf. widemult in section \ref{par:modifier-widemult} on page \pageref{par:modifier-widemult})
\end{itemize}

\paragraph{Execution} {Reservation table \textsf{ALU\_LITE} (Section~\ref{sec:reservation-ALU-LITE})}
~
{\begin{lstlisting}
stage ID:
  new widemult = _ZX_2(widemult);
stage RR:
  new argument3 = GPR[registerY];
  new argument2 = GPR[registerZ];
  if (widemult == 0) {
    new product = _SX_32(argument2) * _SX_32(argument3);
  }
  if (widemult == 1) {
    product = _ZX_32(argument2) * _ZX_32(argument3);
  }
  if (widemult == 2) {
    product = _SX_32(argument2) * _ZX_32(argument3);
  }
  new result1 = -product;
stage E2:
  GPR[registerW] = result1;
\end{lstlisting}}
\newpage

\paragraph{Encoding} Format \textsf{ALU\_DMEWRR.M} (Section~\ref{sec:format-ALU-DMEWRR.M}), \verb|alucode2 = 01|\\

\bitfields{ 1/P, 2/00, 2/-- --, 27/upper27, 1/1, 2/11, 1/0, 2/01, 2/widemult, 6/registerW, 2/10, 3/111, 1/0, 1/splat32, 5/lower5, 6/registerZ }


\paragraph{Syntax} \texttt{mulnwd}\emph{widemult} \emph{registerW} = \emph{registerZ}, \emph{upper27\_\discretionary{}{}{}lower5}\emph{splat32}

\paragraph{Description} The \emph{upper27\_\discretionary{}{}{}lower5} and the \emph{registerZ} words are sign-extended or zero-extended from 32 bits, depending on the \emph{widemult}. The \emph{upper27\_\discretionary{}{}{}lower5} is multiplied by the \emph{registerZ}, producing a 64-bit result. The resulting double word is negated and stored into the \emph{registerW}.


\paragraph{Modifier(s)}
\noindent\begin{itemize}
\item widemult (cf. widemult in section \ref{par:modifier-widemult} on page \pageref{par:modifier-widemult})
\item splat32 (cf. splat32 in section \ref{par:modifier-splat32} on page \pageref{par:modifier-splat32})
\end{itemize}

\paragraph{Execution} {Reservation table \textsf{ALU\_LITE.X} (Section~\ref{sec:reservation-ALU-LITE.X})}
~
{\begin{lstlisting}
stage ID:
  new splat32 = _ZX_1(splat32);
  new widemult = _ZX_2(widemult);
stage RR:
  new argument3 = splat32 ? (_WX_32(upper27_lower5) << 32) | _ZX_32(_WX_32(upper27_lower5)) : _SX_32(_WX_32(upper27_lower5));
  new argument2 = GPR[registerZ];
  if (widemult == 0) {
    new product = _SX_32(argument2) * _SX_32(argument3);
  }
  if (widemult == 1) {
    product = _ZX_32(argument2) * _ZX_32(argument3);
  }
  if (widemult == 2) {
    product = _SX_32(argument2) * _ZX_32(argument3);
  }
  new result1 = -product;
stage E2:
  GPR[registerW] = result1;
\end{lstlisting}}
\newpage
\subsubsection[{\small MULNWDP: Multiply Negate Extend Word to Double Word Pair}]{MULNWDP\hfill Multiply Negate Extend Word to Double Word Pair}
\label{instr:MULNWDP}\index{mulnwdp}
 
\paragraph{Encoding} Format \textsf{ALU\_DPMEWRRE} (Section~\ref{sec:format-ALU-DPMEWRRE}), \verb|alucode2 = 01|\\

\bitfields{ 1/P, 2/11, 1/1, 2/01, 2/widemult, 5/registerM, 1/0, 2/10, 3/110, 1/0, 5/registerYe, 1/1, 5/registerZe, 1/0 }


\paragraph{Syntax} \texttt{mulnwdp}\emph{widemult} \emph{registerM} = \emph{registerZe}, \emph{registerYe}

\paragraph{Description} The \emph{registerYe} and the \emph{registerZe} are interpreted as two words packed into 64 bits. These \emph{registerYe} and \emph{registerZe} words are sign-extended or zero-extended from 32 bits, depending on the \emph{widemult}. The \emph{registerYe} words are multiplied by the \emph{registerZe} words, producing two 64-bit results. The two resulting double words are negated, packed and stored into the \emph{registerM}.


\paragraph{Modifier(s)}
\noindent\begin{itemize}
\item widemult (cf. widemult in section \ref{par:modifier-widemult} on page \pageref{par:modifier-widemult})
\end{itemize}

\paragraph{Execution} {Reservation table \textsf{ALU\_LITE} (Section~\ref{sec:reservation-ALU-LITE})}
~
{\begin{lstlisting}
stage ID:
  new widemult = _ZX_2(widemult);
stage RR:
  new argument3 = GPR[registerYe<<1];
  new argument2 = GPR[registerZe<<1];
  for (I = 0; I < 2; I += 1) {
    if (widemult == 0) {
      new product = _SX_32(argument2.32[I]) * _SX_32(argument3.32[I]);
    }
    if (widemult == 1) {
      product = _ZX_32(argument2.32[I]) * _ZX_32(argument3.32[I]);
    }
    if (widemult == 2) {
      product = _SX_32(argument2.32[I]) * _ZX_32(argument3.32[I]);
    }
    result1.64[I] = -product
  };
stage E2:
  PGR[registerM] = result1;
\end{lstlisting}}
\newpage

\paragraph{Encoding} Format \textsf{ALU\_DPMEWRRO} (Section~\ref{sec:format-ALU-DPMEWRRO}), \verb|alucode2 = 01|\\

\bitfields{ 1/P, 2/11, 1/1, 2/01, 2/widemult, 5/registerM, 1/1, 2/10, 3/110, 1/0, 5/registerYo, 1/1, 5/registerZo, 1/0 }


\paragraph{Syntax} \texttt{mulnwdp}\emph{widemult} \emph{registerM} = \emph{registerZo}, \emph{registerYo}

\paragraph{Description} The \emph{registerYo} and the \emph{registerZo} are interpreted as two words packed into 64 bits. These \emph{registerYo} and \emph{registerZo} words are sign-extended or zero-extended from 32 bits, depending on the \emph{widemult}. The \emph{registerYo} words are multiplied by the \emph{registerZo} words, producing two 64-bit results. The two resulting double words are negated, packed and stored into the \emph{registerM}.


\paragraph{Modifier(s)}
\noindent\begin{itemize}
\item widemult (cf. widemult in section \ref{par:modifier-widemult} on page \pageref{par:modifier-widemult})
\end{itemize}

\paragraph{Execution} {Reservation table \textsf{ALU\_LITE} (Section~\ref{sec:reservation-ALU-LITE})}
~
{\begin{lstlisting}
stage ID:
  new widemult = _ZX_2(widemult);
stage RR:
  new argument3 = GPR[(registerYo<<1)+1];
  new argument2 = GPR[(registerZo<<1)+1];
  for (I = 0; I < 2; I += 1) {
    if (widemult == 0) {
      new product = _SX_32(argument2.32[I]) * _SX_32(argument3.32[I]);
    }
    if (widemult == 1) {
      product = _ZX_32(argument2.32[I]) * _ZX_32(argument3.32[I]);
    }
    if (widemult == 2) {
      product = _SX_32(argument2.32[I]) * _ZX_32(argument3.32[I]);
    }
    result1.64[I] = -product
  };
stage E2:
  PGR[registerM] = result1;
\end{lstlisting}}
\newpage
\subsubsection[{\small MULNWP: Multiply Negate Word Pair}]{MULNWP\hfill Multiply Negate Word Pair}
\label{instr:MULNWP}\index{mulnwp}
 
\paragraph{Encoding} Format \textsf{ALU\_WPMWRRE} (Section~\ref{sec:format-ALU-WPMWRRE}), \verb|alucode2 = 01|\\

\bitfields{ 1/P, 2/11, 1/1, 2/01, 2/highmult, 5/registerWe, 1/0, 2/10, 3/110, 1/0, 5/registerYe, 1/0, 5/registerZe, 1/1 }


\paragraph{Syntax} \texttt{mulnwp}\emph{highmult} \emph{registerWe} = \emph{registerZe}, \emph{registerYe}

\paragraph{Description} The \emph{registerYe} and the \emph{registerZe} are interpreted as two words packed into 64 bits. The \emph{registerYe} and the \emph{registerZe} words are sign-extended or zero-extended from 32 bits, depending on the \emph{highmult}. The \emph{registerYe} words are multiplied by the \emph{registerZe} words, and the product is optionally shifted right by 32 bits, depending on the \emph{highmult}. The two resulting words are negated, packed and stored back into the \emph{registerWe}.


\paragraph{Modifier(s)}
\noindent\begin{itemize}
\item highmult (cf. highmult in section \ref{par:modifier-highmult} on page \pageref{par:modifier-highmult})
\end{itemize}

\paragraph{Execution} {Reservation table \textsf{ALU\_LITE} (Section~\ref{sec:reservation-ALU-LITE})}
~
{\begin{lstlisting}
stage ID:
  new highmult = _ZX_2(highmult);
stage RR:
  new argument3 = GPR[registerYe<<1];
  new argument2 = GPR[registerZe<<1];
  for (I = 0; I < 2; I += 1) {
    if ((highmult == 0) || (highmult == 3)) {
      new product = _SX_32(argument2.32[I]) * _SX_32(argument3.32[I]);
    }
    if (highmult == 1) {
      product = _ZX_32(argument2.32[I]) * _ZX_32(argument3.32[I]);
    }
    if (highmult == 2) {
      product = _SX_32(argument2.32[I]) * _ZX_32(argument3.32[I]);
    }
    if (highmult) {
      result1.32[I] = -product.32[1];
    } else {
      result1.32[I] = -product.32[0];
    }
  };
stage E2:
  GPR[registerWe<<1] = result1;
\end{lstlisting}}
\newpage

\paragraph{Encoding} Format \textsf{ALU\_WPMWRRO} (Section~\ref{sec:format-ALU-WPMWRRO}), \verb|alucode2 = 01|\\

\bitfields{ 1/P, 2/11, 1/1, 2/01, 2/highmult, 5/registerWo, 1/1, 2/10, 3/110, 1/0, 5/registerYo, 1/0, 5/registerZo, 1/1 }


\paragraph{Syntax} \texttt{mulnwp}\emph{highmult} \emph{registerWo} = \emph{registerZo}, \emph{registerYo}

\paragraph{Description} The \emph{registerYo} and the \emph{registerZo} are interpreted as two words packed into 64 bits. The \emph{registerYo} and the \emph{registerZo} words are sign-extended or zero-extended from 32 bits, depending on the \emph{highmult}. The \emph{registerYo} words are multiplied by the \emph{registerZo} words, and the product is optionally shifted right by 32 bits, depending on the \emph{highmult}. The two resulting words are negated, packed and stored back into the \emph{registerWo}.


\paragraph{Modifier(s)}
\noindent\begin{itemize}
\item highmult (cf. highmult in section \ref{par:modifier-highmult} on page \pageref{par:modifier-highmult})
\end{itemize}

\paragraph{Execution} {Reservation table \textsf{ALU\_LITE} (Section~\ref{sec:reservation-ALU-LITE})}
~
{\begin{lstlisting}
stage ID:
  new highmult = _ZX_2(highmult);
stage RR:
  new argument3 = GPR[(registerYo<<1)+1];
  new argument2 = GPR[(registerZo<<1)+1];
  for (I = 0; I < 2; I += 1) {
    if ((highmult == 0) || (highmult == 3)) {
      new product = _SX_32(argument2.32[I]) * _SX_32(argument3.32[I]);
    }
    if (highmult == 1) {
      product = _ZX_32(argument2.32[I]) * _ZX_32(argument3.32[I]);
    }
    if (highmult == 2) {
      product = _SX_32(argument2.32[I]) * _ZX_32(argument3.32[I]);
    }
    if (highmult) {
      result1.32[I] = -product.32[1];
    } else {
      result1.32[I] = -product.32[0];
    }
  };
stage E2:
  GPR[(registerWo<<1)+1] = result1;
\end{lstlisting}}
\newpage
\subsubsection[{\small MULNXBHO: Multiply Negate Extend Byte to Half Word Octuple}]{MULNXBHO\hfill Multiply Negate Extend Byte to Half Word Octuple}
\label{instr:MULNXBHO}\index{mulnxbho}
 
\paragraph{Encoding} Format \textsf{ALU\_HOMXWRR} (Section~\ref{sec:format-ALU-HOMXWRR}), \verb|alucode2 = 01|\\

\bitfields{ 1/P, 2/11, 1/1, 2/01, 2/widemult, 5/registerM, 1/oddlanes, 2/10, 3/111, 1/1, 5/registerO, 1/--, 5/registerP, 1/0 }


\paragraph{Syntax} \texttt{mulnxbho}\emph{oddlanes}\emph{widemult} \emph{registerM} = \emph{registerP}, \emph{registerO}

\paragraph{Description} The \emph{registerO} and the \emph{registerP} are interpreted as sixteen bytes packed into 128 bits. The even or odd bytes of the \emph{registerO} and the \emph{registerP} are selected, depending on \emph{oddlanes}. These \emph{registerO} and \emph{registerP} bytes are sign-extended or zero-extended from 8 bits, depending on the \emph{widemult}. The \emph{registerO} bytes are multiplied by the \emph{registerP} bytes, producing eight 16-bit results. The eight resulting half words are negated, packed and stored into the \emph{registerM}.


\paragraph{Modifier(s)}
\noindent\begin{itemize}
\item widemult (cf. widemult in section \ref{par:modifier-widemult} on page \pageref{par:modifier-widemult})
\item oddlanes (cf. oddlanes in section \ref{par:modifier-oddlanes} on page \pageref{par:modifier-oddlanes})
\end{itemize}

\paragraph{Execution} {Reservation table \textsf{ALU\_LITE} (Section~\ref{sec:reservation-ALU-LITE})}
~
{\begin{lstlisting}
stage ID:
  new oddlanes = _ZX_1(oddlanes);
  new widemult = _ZX_2(widemult);
stage RR:
  new argument3 = PGR[registerO];
  new argument2 = PGR[registerP];
  if (oddlanes) {
    argument2 = argument2 >> 8;
    argument3 = argument3 >> 8;
  }
  for (I = 0; I < 8; I += 1) {
    if (widemult == 0) {
      new product = _SX_8(argument2.16[I]) * _SX_8(argument3.16[I]);
    }
    if (widemult == 1) {
      product = _ZX_8(argument2.16[I]) * _ZX_8(argument3.16[I]);
    }
    if (widemult == 2) {
      product = _SX_8(argument2.16[I]) * _ZX_8(argument3.16[I]);
    }
    result1.16[I] = -product
  };
stage E2:
  PGR[registerM] = result1;
\end{lstlisting}}
\newpage
\subsubsection[{\small MULNXHWQ: Multiply Negate Extend Half Word to Word Quadruple}]{MULNXHWQ\hfill Multiply Negate Extend Half Word to Word Quadruple}
\label{instr:MULNXHWQ}\index{mulnxhwq}
 
\paragraph{Encoding} Format \textsf{ALU\_WQMXWRR} (Section~\ref{sec:format-ALU-WQMXWRR}), \verb|alucode2 = 01|\\

\bitfields{ 1/P, 2/11, 1/1, 2/01, 2/widemult, 5/registerM, 1/oddlanes, 2/10, 3/111, 1/0, 5/registerO, 1/--, 5/registerP, 1/1 }


\paragraph{Syntax} \texttt{mulnxhwq}\emph{oddlanes}\emph{widemult} \emph{registerM} = \emph{registerP}, \emph{registerO}

\paragraph{Description} The \emph{registerO} and the \emph{registerP} are interpreted as eight half words packed into 128 bits. The even or odd half words of the \emph{registerO} and the \emph{registerP} are selected, depending on \emph{oddlanes}. These \emph{registerO} and \emph{registerP} half words are sign-extended or zero-extended from 16 bits, depending on the \emph{widemult}. The \emph{registerO} half words are multiplied by the \emph{registerP} half words, producing four 32-bit results. The four resulting words are negated, packed and stored into the \emph{registerM}.


\paragraph{Modifier(s)}
\noindent\begin{itemize}
\item widemult (cf. widemult in section \ref{par:modifier-widemult} on page \pageref{par:modifier-widemult})
\item oddlanes (cf. oddlanes in section \ref{par:modifier-oddlanes} on page \pageref{par:modifier-oddlanes})
\end{itemize}

\paragraph{Execution} {Reservation table \textsf{ALU\_LITE} (Section~\ref{sec:reservation-ALU-LITE})}
~
{\begin{lstlisting}
stage ID:
  new oddlanes = _ZX_1(oddlanes);
  new widemult = _ZX_2(widemult);
stage RR:
  new argument3 = PGR[registerO];
  new argument2 = PGR[registerP];
  if (oddlanes) {
    argument2 = argument2 >> 16;
    argument3 = argument3 >> 16;
  }
  for (I = 0; I < 4; I += 1) {
    if (widemult == 0) {
      new product = _SX_16(argument2.32[I]) * _SX_16(argument3.32[I]);
    }
    if (widemult == 1) {
      product = _ZX_16(argument2.32[I]) * _ZX_16(argument3.32[I]);
    }
    if (widemult == 2) {
      product = _SX_16(argument2.32[I]) * _ZX_16(argument3.32[I]);
    }
    result1.32[I] = -product
  };
stage E2:
  PGR[registerM] = result1;
\end{lstlisting}}
\newpage
\subsubsection[{\small MULNXWDP: Multiply Negate Extend Word to Double Word Pair}]{MULNXWDP\hfill Multiply Negate Extend Word to Double Word Pair}
\label{instr:MULNXWDP}\index{mulnxwdp}
 
\paragraph{Encoding} Format \textsf{ALU\_DPMXWRR} (Section~\ref{sec:format-ALU-DPMXWRR}), \verb|alucode2 = 01|\\

\bitfields{ 1/P, 2/11, 1/1, 2/01, 2/widemult, 5/registerM, 1/oddlanes, 2/10, 3/111, 1/0, 5/registerO, 1/--, 5/registerP, 1/0 }


\paragraph{Syntax} \texttt{mulnxwdp}\emph{oddlanes}\emph{widemult} \emph{registerM} = \emph{registerP}, \emph{registerO}

\paragraph{Description} The \emph{registerO} and the \emph{registerP} are interpreted as four words packed into 128 bits. The even or odd words of the \emph{registerO} and the \emph{registerP} are selected, depending on \emph{oddlanes}. These \emph{registerO} and \emph{registerP} words are sign-extended or zero-extended from 32 bits, depending on the \emph{widemult}. The \emph{registerO} words are multiplied by the \emph{registerP} words, producing two 64-bit results. The two resulting double words are negated, packed and stored into the \emph{registerM}.


\paragraph{Modifier(s)}
\noindent\begin{itemize}
\item widemult (cf. widemult in section \ref{par:modifier-widemult} on page \pageref{par:modifier-widemult})
\item oddlanes (cf. oddlanes in section \ref{par:modifier-oddlanes} on page \pageref{par:modifier-oddlanes})
\end{itemize}

\paragraph{Execution} {Reservation table \textsf{ALU\_LITE} (Section~\ref{sec:reservation-ALU-LITE})}
~
{\begin{lstlisting}
stage ID:
  new oddlanes = _ZX_1(oddlanes);
  new widemult = _ZX_2(widemult);
stage RR:
  new argument3 = PGR[registerO];
  new argument2 = PGR[registerP];
  if (oddlanes) {
    argument2 = argument2 >> 32;
    argument3 = argument3 >> 32;
  }
  for (I = 0; I < 2; I += 1) {
    if (widemult == 0) {
      new product = _SX_32(argument2.64[I]) * _SX_32(argument3.64[I]);
    }
    if (widemult == 1) {
      product = _ZX_32(argument2.64[I]) * _ZX_32(argument3.64[I]);
    }
    if (widemult == 2) {
      product = _SX_32(argument2.64[I]) * _ZX_32(argument3.64[I]);
    }
    result1.64[I] = -product
  };
stage E2:
  PGR[registerM] = result1;
\end{lstlisting}}
\newpage
\subsubsection[{\small MULW: Multiply Word}]{MULW\hfill Multiply Word}
\label{instr:MULW}\index{mulw}
 
\paragraph{Encoding} Format \textsf{ALU\_WMWRR} (Section~\ref{sec:format-ALU-WMWRR}), \verb|alucode2 = 00|\\

\bitfields{ 1/P, 2/11, 1/0, 2/00, 2/highmult, 6/registerW, 2/11, 3/110, 1/signextw, 6/registerY, 6/registerZ }


\paragraph{Syntax} \texttt{mulw}\emph{highmult}\emph{signextw} \emph{registerW} = \emph{registerZ}, \emph{registerY}

\paragraph{Description} The \emph{registerY} and the \emph{registerZ} words are sign-extended or zero-extended from 32 bits, depending on the \emph{highmult}. The \emph{registerY} word is multiplied by the \emph{registerZ} word, and the product is optionally shifted right by 32 bits, depending on the \emph{highmult}. The resulting word is stored into the \emph{registerW}.


\paragraph{Modifier(s)}
\noindent\begin{itemize}
\item highmult (cf. highmult in section \ref{par:modifier-highmult} on page \pageref{par:modifier-highmult})
\item signextw (cf. signextw in section \ref{par:modifier-signextw} on page \pageref{par:modifier-signextw})
\end{itemize}

\paragraph{Execution} {Reservation table \textsf{ALU\_LITE} (Section~\ref{sec:reservation-ALU-LITE})}
~
{\begin{lstlisting}
stage ID:
  new signextw = _ZX_1(signextw);
  new highmult = _ZX_2(highmult);
stage RR:
  new argument3 = GPR[registerY];
  new argument2 = GPR[registerZ];
  if ((highmult == 0) || (highmult == 3)) {
    new product = _SX_32(argument2) * _SX_32(argument3);
  }
  if (highmult == 1) {
    product = _ZX_32(argument2) * _ZX_32(argument3);
  }
  if (highmult == 2) {
    product = _SX_32(argument2) * _ZX_32(argument3);
  }
  if (highmult) {
    new result1 = product.32[1];
  } else {
    result1 = product.32[0];
  }
stage E2:
  GPR[registerW] = signextw ? _SX_32(result1) : _ZX_32(result1);
\end{lstlisting}}
\newpage

\paragraph{Encoding} Format \textsf{ALU\_WMWRR.W} (Section~\ref{sec:format-ALU-WMWRR.W}), \verb|alucode2 = 00|\\

\bitfields{ 1/P, 2/00, 2/-- --, 27/upper27, 1/1, 2/11, 1/0, 2/00, 2/highmult, 6/registerW, 2/11, 3/110, 1/signextw, 1/0, 5/lower5, 6/registerZ }


\paragraph{Syntax} \texttt{mulw}\emph{highmult}\emph{signextw} \emph{registerW} = \emph{registerZ}, \emph{upper27\_\discretionary{}{}{}lower5}

\paragraph{Description} The \emph{upper27\_\discretionary{}{}{}lower5} and the \emph{registerZ} words are sign-extended or zero-extended from 32 bits, depending on the \emph{highmult}. The \emph{upper27\_\discretionary{}{}{}lower5} word is multiplied by the \emph{registerZ} word, and the product is optionally shifted right by 32 bits, depending on the \emph{highmult}. The resulting word is stored into the \emph{registerW}.


\paragraph{Modifier(s)}
\noindent\begin{itemize}
\item highmult (cf. highmult in section \ref{par:modifier-highmult} on page \pageref{par:modifier-highmult})
\item signextw (cf. signextw in section \ref{par:modifier-signextw} on page \pageref{par:modifier-signextw})
\end{itemize}

\paragraph{Execution} {Reservation table \textsf{ALU\_LITE.X} (Section~\ref{sec:reservation-ALU-LITE.X})}
~
{\begin{lstlisting}
stage ID:
  new signextw = _ZX_1(signextw);
  new highmult = _ZX_2(highmult);
stage RR:
  new argument3 = _WX_32(upper27_lower5);
  new argument2 = GPR[registerZ];
  if ((highmult == 0) || (highmult == 3)) {
    new product = _SX_32(argument2) * _SX_32(argument3);
  }
  if (highmult == 1) {
    product = _ZX_32(argument2) * _ZX_32(argument3);
  }
  if (highmult == 2) {
    product = _SX_32(argument2) * _ZX_32(argument3);
  }
  if (highmult) {
    new result1 = product.32[1];
  } else {
    result1 = product.32[0];
  }
stage E2:
  GPR[registerW] = signextw ? _SX_32(result1) : _ZX_32(result1);
\end{lstlisting}}
\newpage
\subsubsection[{\small MULWD: Multiply Extend Word to Double Word}]{MULWD\hfill Multiply Extend Word to Double Word}
\label{instr:MULWD}\index{mulwd}
 
\paragraph{Encoding} Format \textsf{ALU\_DMEWRR} (Section~\ref{sec:format-ALU-DMEWRR}), \verb|alucode2 = 00|\\

\bitfields{ 1/P, 2/11, 1/0, 2/00, 2/widemult, 6/registerW, 2/10, 3/111, 1/0, 6/registerY, 6/registerZ }


\paragraph{Syntax} \texttt{mulwd}\emph{widemult} \emph{registerW} = \emph{registerZ}, \emph{registerY}

\paragraph{Description} The \emph{registerY} and the \emph{registerZ} words are sign-extended or zero-extended from 32 bits, depending on the \emph{widemult}. The \emph{registerY} is multiplied by the \emph{registerZ}, producing a 64-bit result. The resulting double word is stored into the \emph{registerW}.


\paragraph{Modifier(s)}
\noindent\begin{itemize}
\item widemult (cf. widemult in section \ref{par:modifier-widemult} on page \pageref{par:modifier-widemult})
\end{itemize}

\paragraph{Execution} {Reservation table \textsf{ALU\_LITE} (Section~\ref{sec:reservation-ALU-LITE})}
~
{\begin{lstlisting}
stage ID:
  new widemult = _ZX_2(widemult);
stage RR:
  new argument3 = GPR[registerY];
  new argument2 = GPR[registerZ];
  if (widemult == 0) {
    new product = _SX_32(argument2) * _SX_32(argument3);
  }
  if (widemult == 1) {
    product = _ZX_32(argument2) * _ZX_32(argument3);
  }
  if (widemult == 2) {
    product = _SX_32(argument2) * _ZX_32(argument3);
  }
  new result1 = product;
stage E2:
  GPR[registerW] = result1;
\end{lstlisting}}
\newpage

\paragraph{Encoding} Format \textsf{ALU\_DMEWRR.M} (Section~\ref{sec:format-ALU-DMEWRR.M}), \verb|alucode2 = 00|\\

\bitfields{ 1/P, 2/00, 2/-- --, 27/upper27, 1/1, 2/11, 1/0, 2/00, 2/widemult, 6/registerW, 2/10, 3/111, 1/0, 1/splat32, 5/lower5, 6/registerZ }


\paragraph{Syntax} \texttt{mulwd}\emph{widemult} \emph{registerW} = \emph{registerZ}, \emph{upper27\_\discretionary{}{}{}lower5}\emph{splat32}

\paragraph{Description} The \emph{upper27\_\discretionary{}{}{}lower5} and the \emph{registerZ} words are sign-extended or zero-extended from 32 bits, depending on the \emph{widemult}. The \emph{upper27\_\discretionary{}{}{}lower5} is multiplied by the \emph{registerZ}, producing a 64-bit result. The resulting double word is stored into the \emph{registerW}.


\paragraph{Modifier(s)}
\noindent\begin{itemize}
\item widemult (cf. widemult in section \ref{par:modifier-widemult} on page \pageref{par:modifier-widemult})
\item splat32 (cf. splat32 in section \ref{par:modifier-splat32} on page \pageref{par:modifier-splat32})
\end{itemize}

\paragraph{Execution} {Reservation table \textsf{ALU\_LITE.X} (Section~\ref{sec:reservation-ALU-LITE.X})}
~
{\begin{lstlisting}
stage ID:
  new splat32 = _ZX_1(splat32);
  new widemult = _ZX_2(widemult);
stage RR:
  new argument3 = splat32 ? (_WX_32(upper27_lower5) << 32) | _ZX_32(_WX_32(upper27_lower5)) : _SX_32(_WX_32(upper27_lower5));
  new argument2 = GPR[registerZ];
  if (widemult == 0) {
    new product = _SX_32(argument2) * _SX_32(argument3);
  }
  if (widemult == 1) {
    product = _ZX_32(argument2) * _ZX_32(argument3);
  }
  if (widemult == 2) {
    product = _SX_32(argument2) * _ZX_32(argument3);
  }
  new result1 = product;
stage E2:
  GPR[registerW] = result1;
\end{lstlisting}}
\newpage
\subsubsection[{\small MULWDP: Multiply Extend Word to Double Word Pair}]{MULWDP\hfill Multiply Extend Word to Double Word Pair}
\label{instr:MULWDP}\index{mulwdp}
 
\paragraph{Encoding} Format \textsf{ALU\_DPMEWRRE} (Section~\ref{sec:format-ALU-DPMEWRRE}), \verb|alucode2 = 00|\\

\bitfields{ 1/P, 2/11, 1/1, 2/00, 2/widemult, 5/registerM, 1/0, 2/10, 3/110, 1/0, 5/registerYe, 1/1, 5/registerZe, 1/0 }


\paragraph{Syntax} \texttt{mulwdp}\emph{widemult} \emph{registerM} = \emph{registerZe}, \emph{registerYe}

\paragraph{Description} The \emph{registerYe} and the \emph{registerZe} are interpreted as two words packed into 64 bits. These \emph{registerYe} and \emph{registerZe} words are sign-extended or zero-extended from 32 bits, depending on the \emph{widemult}. The \emph{registerYe} words are multiplied by the \emph{registerZe} words, producing two 64-bit results. The two resulting double words are packed and stored into the \emph{registerM}.


\paragraph{Modifier(s)}
\noindent\begin{itemize}
\item widemult (cf. widemult in section \ref{par:modifier-widemult} on page \pageref{par:modifier-widemult})
\end{itemize}

\paragraph{Execution} {Reservation table \textsf{ALU\_LITE} (Section~\ref{sec:reservation-ALU-LITE})}
~
{\begin{lstlisting}
stage ID:
  new widemult = _ZX_2(widemult);
stage RR:
  new argument3 = GPR[registerYe<<1];
  new argument2 = GPR[registerZe<<1];
  for (I = 0; I < 2; I += 1) {
    if (widemult == 0) {
      new product = _SX_32(argument2.32[I]) * _SX_32(argument3.32[I]);
    }
    if (widemult == 1) {
      product = _ZX_32(argument2.32[I]) * _ZX_32(argument3.32[I]);
    }
    if (widemult == 2) {
      product = _SX_32(argument2.32[I]) * _ZX_32(argument3.32[I]);
    }
    result1.64[I] = product
  };
stage E2:
  PGR[registerM] = result1;
\end{lstlisting}}
\newpage

\paragraph{Encoding} Format \textsf{ALU\_DPMEWRRO} (Section~\ref{sec:format-ALU-DPMEWRRO}), \verb|alucode2 = 00|\\

\bitfields{ 1/P, 2/11, 1/1, 2/00, 2/widemult, 5/registerM, 1/1, 2/10, 3/110, 1/0, 5/registerYo, 1/1, 5/registerZo, 1/0 }


\paragraph{Syntax} \texttt{mulwdp}\emph{widemult} \emph{registerM} = \emph{registerZo}, \emph{registerYo}

\paragraph{Description} The \emph{registerYo} and the \emph{registerZo} are interpreted as two words packed into 64 bits. These \emph{registerYo} and \emph{registerZo} words are sign-extended or zero-extended from 32 bits, depending on the \emph{widemult}. The \emph{registerYo} words are multiplied by the \emph{registerZo} words, producing two 64-bit results. The two resulting double words are packed and stored into the \emph{registerM}.


\paragraph{Modifier(s)}
\noindent\begin{itemize}
\item widemult (cf. widemult in section \ref{par:modifier-widemult} on page \pageref{par:modifier-widemult})
\end{itemize}

\paragraph{Execution} {Reservation table \textsf{ALU\_LITE} (Section~\ref{sec:reservation-ALU-LITE})}
~
{\begin{lstlisting}
stage ID:
  new widemult = _ZX_2(widemult);
stage RR:
  new argument3 = GPR[(registerYo<<1)+1];
  new argument2 = GPR[(registerZo<<1)+1];
  for (I = 0; I < 2; I += 1) {
    if (widemult == 0) {
      new product = _SX_32(argument2.32[I]) * _SX_32(argument3.32[I]);
    }
    if (widemult == 1) {
      product = _ZX_32(argument2.32[I]) * _ZX_32(argument3.32[I]);
    }
    if (widemult == 2) {
      product = _SX_32(argument2.32[I]) * _ZX_32(argument3.32[I]);
    }
    result1.64[I] = product
  };
stage E2:
  PGR[registerM] = result1;
\end{lstlisting}}
\newpage
\subsubsection[{\small MULWP: Multiply Word Pair}]{MULWP\hfill Multiply Word Pair}
\label{instr:MULWP}\index{mulwp}
 
\paragraph{Encoding} Format \textsf{ALU\_WPMWRRE} (Section~\ref{sec:format-ALU-WPMWRRE}), \verb|alucode2 = 00|\\

\bitfields{ 1/P, 2/11, 1/1, 2/00, 2/highmult, 5/registerWe, 1/0, 2/10, 3/110, 1/0, 5/registerYe, 1/0, 5/registerZe, 1/1 }


\paragraph{Syntax} \texttt{mulwp}\emph{highmult} \emph{registerWe} = \emph{registerZe}, \emph{registerYe}

\paragraph{Description} The \emph{registerYe} and the \emph{registerZe} are interpreted as two words packed into 64 bits. The \emph{registerYe} and the \emph{registerZe} words are sign-extended or zero-extended from 32 bits, depending on the \emph{highmult}. The \emph{registerYe} words are multiplied by the \emph{registerZe} words, and the product is optionally shifted right by 32 bits, depending on the \emph{highmult}. The two resulting words are packed and stored back into the \emph{registerWe}.


\paragraph{Modifier(s)}
\noindent\begin{itemize}
\item highmult (cf. highmult in section \ref{par:modifier-highmult} on page \pageref{par:modifier-highmult})
\end{itemize}

\paragraph{Execution} {Reservation table \textsf{ALU\_LITE} (Section~\ref{sec:reservation-ALU-LITE})}
~
{\begin{lstlisting}
stage ID:
  new highmult = _ZX_2(highmult);
stage RR:
  new argument3 = GPR[registerYe<<1];
  new argument2 = GPR[registerZe<<1];
  for (I = 0; I < 2; I += 1) {
    if ((highmult == 0) || (highmult == 3)) {
      new product = _SX_32(argument2.32[I]) * _SX_32(argument3.32[I]);
    }
    if (highmult == 1) {
      product = _ZX_32(argument2.32[I]) * _ZX_32(argument3.32[I]);
    }
    if (highmult == 2) {
      product = _SX_32(argument2.32[I]) * _ZX_32(argument3.32[I]);
    }
    if (highmult) {
      result1.32[I] = product.32[1];
    } else {
      result1.32[I] = product.32[0];
    }
  };
stage E2:
  GPR[registerWe<<1] = result1;
\end{lstlisting}}
\newpage

\paragraph{Encoding} Format \textsf{ALU\_WPMWRRO} (Section~\ref{sec:format-ALU-WPMWRRO}), \verb|alucode2 = 00|\\

\bitfields{ 1/P, 2/11, 1/1, 2/00, 2/highmult, 5/registerWo, 1/1, 2/10, 3/110, 1/0, 5/registerYo, 1/0, 5/registerZo, 1/1 }


\paragraph{Syntax} \texttt{mulwp}\emph{highmult} \emph{registerWo} = \emph{registerZo}, \emph{registerYo}

\paragraph{Description} The \emph{registerYo} and the \emph{registerZo} are interpreted as two words packed into 64 bits. The \emph{registerYo} and the \emph{registerZo} words are sign-extended or zero-extended from 32 bits, depending on the \emph{highmult}. The \emph{registerYo} words are multiplied by the \emph{registerZo} words, and the product is optionally shifted right by 32 bits, depending on the \emph{highmult}. The two resulting words are packed and stored back into the \emph{registerWo}.


\paragraph{Modifier(s)}
\noindent\begin{itemize}
\item highmult (cf. highmult in section \ref{par:modifier-highmult} on page \pageref{par:modifier-highmult})
\end{itemize}

\paragraph{Execution} {Reservation table \textsf{ALU\_LITE} (Section~\ref{sec:reservation-ALU-LITE})}
~
{\begin{lstlisting}
stage ID:
  new highmult = _ZX_2(highmult);
stage RR:
  new argument3 = GPR[(registerYo<<1)+1];
  new argument2 = GPR[(registerZo<<1)+1];
  for (I = 0; I < 2; I += 1) {
    if ((highmult == 0) || (highmult == 3)) {
      new product = _SX_32(argument2.32[I]) * _SX_32(argument3.32[I]);
    }
    if (highmult == 1) {
      product = _ZX_32(argument2.32[I]) * _ZX_32(argument3.32[I]);
    }
    if (highmult == 2) {
      product = _SX_32(argument2.32[I]) * _ZX_32(argument3.32[I]);
    }
    if (highmult) {
      result1.32[I] = product.32[1];
    } else {
      result1.32[I] = product.32[0];
    }
  };
stage E2:
  GPR[(registerWo<<1)+1] = result1;
\end{lstlisting}}
\newpage
\subsubsection[{\small MULXBHO: Multiply Extend Byte to Half Word Octuple}]{MULXBHO\hfill Multiply Extend Byte to Half Word Octuple}
\label{instr:MULXBHO}\index{mulxbho}
 
\paragraph{Encoding} Format \textsf{ALU\_HOMXWRR} (Section~\ref{sec:format-ALU-HOMXWRR}), \verb|alucode2 = 00|\\

\bitfields{ 1/P, 2/11, 1/1, 2/00, 2/widemult, 5/registerM, 1/oddlanes, 2/10, 3/111, 1/1, 5/registerO, 1/--, 5/registerP, 1/0 }


\paragraph{Syntax} \texttt{mulxbho}\emph{oddlanes}\emph{widemult} \emph{registerM} = \emph{registerP}, \emph{registerO}

\paragraph{Description} The \emph{registerO} and the \emph{registerP} are interpreted as sixteen bytes packed into 128 bits. The even or odd bytes of the \emph{registerO} and the \emph{registerP} are selected, depending on \emph{oddlanes}. These \emph{registerO} and \emph{registerP} bytes are sign-extended or zero-extended from 8 bits, depending on the \emph{widemult}. The \emph{registerO} bytes are multiplied by the \emph{registerP} bytes, producing eight 16-bit results. The eight resulting half words are packed and stored into the \emph{registerM}.


\paragraph{Modifier(s)}
\noindent\begin{itemize}
\item widemult (cf. widemult in section \ref{par:modifier-widemult} on page \pageref{par:modifier-widemult})
\item oddlanes (cf. oddlanes in section \ref{par:modifier-oddlanes} on page \pageref{par:modifier-oddlanes})
\end{itemize}

\paragraph{Execution} {Reservation table \textsf{ALU\_LITE} (Section~\ref{sec:reservation-ALU-LITE})}
~
{\begin{lstlisting}
stage ID:
  new oddlanes = _ZX_1(oddlanes);
  new widemult = _ZX_2(widemult);
stage RR:
  new argument3 = PGR[registerO];
  new argument2 = PGR[registerP];
  if (oddlanes) {
    argument2 = argument2 >> 8;
    argument3 = argument3 >> 8;
  }
  for (I = 0; I < 8; I += 1) {
    if (widemult == 0) {
      new product = _SX_8(argument2.16[I]) * _SX_8(argument3.16[I]);
    }
    if (widemult == 1) {
      product = _ZX_8(argument2.16[I]) * _ZX_8(argument3.16[I]);
    }
    if (widemult == 2) {
      product = _SX_8(argument2.16[I]) * _ZX_8(argument3.16[I]);
    }
    result1.16[I] = product
  };
stage E2:
  PGR[registerM] = result1;
\end{lstlisting}}
\newpage
\subsubsection[{\small MULXHWQ: Multiply Extend Half Word to Word Quadruple}]{MULXHWQ\hfill Multiply Extend Half Word to Word Quadruple}
\label{instr:MULXHWQ}\index{mulxhwq}
 
\paragraph{Encoding} Format \textsf{ALU\_WQMXWRR} (Section~\ref{sec:format-ALU-WQMXWRR}), \verb|alucode2 = 00|\\

\bitfields{ 1/P, 2/11, 1/1, 2/00, 2/widemult, 5/registerM, 1/oddlanes, 2/10, 3/111, 1/0, 5/registerO, 1/--, 5/registerP, 1/1 }


\paragraph{Syntax} \texttt{mulxhwq}\emph{oddlanes}\emph{widemult} \emph{registerM} = \emph{registerP}, \emph{registerO}

\paragraph{Description} The \emph{registerO} and the \emph{registerP} are interpreted as eight half words packed into 128 bits. The even or odd half words of the \emph{registerO} and the \emph{registerP} are selected, depending on \emph{oddlanes}. These \emph{registerO} and \emph{registerP} half words are sign-extended or zero-extended from 16 bits, depending on the \emph{widemult}. The \emph{registerO} half words are multiplied by the \emph{registerP} half words, producing four 32-bit results. The four resulting words are packed and stored into the \emph{registerM}.


\paragraph{Modifier(s)}
\noindent\begin{itemize}
\item widemult (cf. widemult in section \ref{par:modifier-widemult} on page \pageref{par:modifier-widemult})
\item oddlanes (cf. oddlanes in section \ref{par:modifier-oddlanes} on page \pageref{par:modifier-oddlanes})
\end{itemize}

\paragraph{Execution} {Reservation table \textsf{ALU\_LITE} (Section~\ref{sec:reservation-ALU-LITE})}
~
{\begin{lstlisting}
stage ID:
  new oddlanes = _ZX_1(oddlanes);
  new widemult = _ZX_2(widemult);
stage RR:
  new argument3 = PGR[registerO];
  new argument2 = PGR[registerP];
  if (oddlanes) {
    argument2 = argument2 >> 16;
    argument3 = argument3 >> 16;
  }
  for (I = 0; I < 4; I += 1) {
    if (widemult == 0) {
      new product = _SX_16(argument2.32[I]) * _SX_16(argument3.32[I]);
    }
    if (widemult == 1) {
      product = _ZX_16(argument2.32[I]) * _ZX_16(argument3.32[I]);
    }
    if (widemult == 2) {
      product = _SX_16(argument2.32[I]) * _ZX_16(argument3.32[I]);
    }
    result1.32[I] = product
  };
stage E2:
  PGR[registerM] = result1;
\end{lstlisting}}
\newpage
\subsubsection[{\small MULXWDP: Multiply Extend Word to Double Word Pair}]{MULXWDP\hfill Multiply Extend Word to Double Word Pair}
\label{instr:MULXWDP}\index{mulxwdp}
 
\paragraph{Encoding} Format \textsf{ALU\_DPMXWRR} (Section~\ref{sec:format-ALU-DPMXWRR}), \verb|alucode2 = 00|\\

\bitfields{ 1/P, 2/11, 1/1, 2/00, 2/widemult, 5/registerM, 1/oddlanes, 2/10, 3/111, 1/0, 5/registerO, 1/--, 5/registerP, 1/0 }


\paragraph{Syntax} \texttt{mulxwdp}\emph{oddlanes}\emph{widemult} \emph{registerM} = \emph{registerP}, \emph{registerO}

\paragraph{Description} The \emph{registerO} and the \emph{registerP} are interpreted as four words packed into 128 bits. The even or odd words of the \emph{registerO} and the \emph{registerP} are selected, depending on \emph{oddlanes}. These \emph{registerO} and \emph{registerP} words are sign-extended or zero-extended from 32 bits, depending on the \emph{widemult}. The \emph{registerO} words are multiplied by the \emph{registerP} words, producing two 64-bit results. The two resulting double words are packed and stored into the \emph{registerM}.


\paragraph{Modifier(s)}
\noindent\begin{itemize}
\item widemult (cf. widemult in section \ref{par:modifier-widemult} on page \pageref{par:modifier-widemult})
\item oddlanes (cf. oddlanes in section \ref{par:modifier-oddlanes} on page \pageref{par:modifier-oddlanes})
\end{itemize}

\paragraph{Execution} {Reservation table \textsf{ALU\_LITE} (Section~\ref{sec:reservation-ALU-LITE})}
~
{\begin{lstlisting}
stage ID:
  new oddlanes = _ZX_1(oddlanes);
  new widemult = _ZX_2(widemult);
stage RR:
  new argument3 = PGR[registerO];
  new argument2 = PGR[registerP];
  if (oddlanes) {
    argument2 = argument2 >> 32;
    argument3 = argument3 >> 32;
  }
  for (I = 0; I < 2; I += 1) {
    if (widemult == 0) {
      new product = _SX_32(argument2.64[I]) * _SX_32(argument3.64[I]);
    }
    if (widemult == 1) {
      product = _ZX_32(argument2.64[I]) * _ZX_32(argument3.64[I]);
    }
    if (widemult == 2) {
      product = _SX_32(argument2.64[I]) * _ZX_32(argument3.64[I]);
    }
    result1.64[I] = product
  };
stage E2:
  PGR[registerM] = result1;
\end{lstlisting}}
\newpage
\subsubsection[{\small NANDD: Bitwise Not And Double Word}]{NANDD\hfill Bitwise Not And Double Word}
\label{instr:NANDD}\index{nandd}
 
\paragraph{Encoding} Format \textsf{ALU\_DWRR0} (Section~\ref{sec:format-ALU-DWRR0}), \verb|exucode2 = 1001|\\

\bitfields{ 1/P, 2/11, 1/0, 4/1001, 6/registerW, 2/10, 3/000, 1/0, 6/registerY, 6/registerZ }


\paragraph{Syntax} \texttt{nandd} \emph{registerW} = \emph{registerZ}, \emph{registerY}

\paragraph{Description} Bitwise not of bitwise and of the \emph{registerZ} with the \emph{registerY}. The result is stored into the \emph{registerW}.


\paragraph{Execution} {Reservation table \textsf{ALU\_TINY} (Section~\ref{sec:reservation-ALU-TINY})}
~
{\begin{lstlisting}
stage RR:
  new argument3 = GPR[registerY];
  new argument2 = GPR[registerZ];
  new result1 = ~(argument2 & argument3);
stage E1:
  GPR[registerW] = result1;
\end{lstlisting}}
\newpage

\paragraph{Encoding} Format \textsf{ALU\_DWRR0.M} (Section~\ref{sec:format-ALU-DWRR0.M}), \verb|exucode2 = 1001|\\

\bitfields{ 1/P, 2/00, 2/-- --, 27/upper27, 1/1, 2/11, 1/0, 4/1001, 6/registerW, 2/10, 3/000, 1/0, 1/splat32, 5/lower5, 6/registerZ }


\paragraph{Syntax} \texttt{nandd} \emph{registerW} = \emph{registerZ}, \emph{upper27\_\discretionary{}{}{}lower5}\emph{splat32}

\paragraph{Description} Bitwise not of bitwise and of the \emph{registerZ} with the \emph{upper27\_\discretionary{}{}{}lower5}. The result is stored into the \emph{registerW}.


\paragraph{Modifier(s)}
\noindent\begin{itemize}
\item splat32 (cf. splat32 in section \ref{par:modifier-splat32} on page \pageref{par:modifier-splat32})
\end{itemize}

\paragraph{Execution} {Reservation table \textsf{ALU\_TINY.X} (Section~\ref{sec:reservation-ALU-TINY.X})}
~
{\begin{lstlisting}
stage ID:
  new splat32 = _ZX_1(splat32);
stage RR:
  new argument3 = splat32 ? (_WX_32(upper27_lower5) << 32) | _ZX_32(_WX_32(upper27_lower5)) : _SX_32(_WX_32(upper27_lower5));
  new argument2 = GPR[registerZ];
  new result1 = ~(argument2 & argument3);
stage E1:
  GPR[registerW] = result1;
\end{lstlisting}}
\newpage

\paragraph{Encoding} Format \textsf{ALU\_DWRI} (Section~\ref{sec:format-ALU-DWRI}), \verb|exucode2 = 1001|\\

\bitfields{ 1/P, 2/11, 1/0, 4/1001, 6/registerW, 2/00, 10/signed10, 6/registerZ }


\paragraph{Syntax} \texttt{nandd} \emph{registerW} = \emph{registerZ}, \emph{signed10}

\paragraph{Description} Bitwise not of bitwise and of the \emph{registerZ} with the \emph{signed10}. The result is stored into the \emph{registerW}.


\paragraph{Execution} {Reservation table \textsf{ALU\_TINY} (Section~\ref{sec:reservation-ALU-TINY})}
~
{\begin{lstlisting}
stage RR:
  new argument3 = _SX_10(signed10);
  new argument2 = GPR[registerZ];
  new result1 = ~(argument2 & argument3);
stage E1:
  GPR[registerW] = result1;
\end{lstlisting}}
\newpage

\paragraph{Encoding} Format \textsf{ALU\_DWRI.X} (Section~\ref{sec:format-ALU-DWRI.X}), \verb|exucode2 = 1001|\\

\bitfields{ 1/P, 2/00, 2/-- --, 27/upper27, 1/1, 2/11, 1/0, 4/1001, 6/registerW, 2/00, 10/lower10, 6/registerZ }


\paragraph{Syntax} \texttt{nandd} \emph{registerW} = \emph{registerZ}, \emph{upper27\_\discretionary{}{}{}lower10}

\paragraph{Description} Bitwise not of bitwise and of the \emph{registerZ} with the \emph{upper27\_\discretionary{}{}{}lower10}. The result is stored into the \emph{registerW}.


\paragraph{Execution} {Reservation table \textsf{ALU\_TINY.X} (Section~\ref{sec:reservation-ALU-TINY.X})}
~
{\begin{lstlisting}
stage RR:
  new argument3 = _SX_37(upper27_lower10);
  new argument2 = GPR[registerZ];
  new result1 = ~(argument2 & argument3);
stage E1:
  GPR[registerW] = result1;
\end{lstlisting}}
\newpage

\paragraph{Encoding} Format \textsf{ALU\_DWRI.Y} (Section~\ref{sec:format-ALU-DWRI.Y}), \verb|exucode2 = 1001|\\

\bitfields{ 1/P, 2/00, 2/-- --, 27/extend27, 1/1, 2/00, 2/-- --, 27/upper27, 1/1, 2/11, 1/0, 4/1001, 6/registerW, 2/00, 10/lower10, 6/registerZ }


\paragraph{Syntax} \texttt{nandd} \emph{registerW} = \emph{registerZ}, \emph{extend27\_\discretionary{}{}{}upper27\_\discretionary{}{}{}lower10}

\paragraph{Description} Bitwise not of bitwise and of the \emph{registerZ} with the \emph{extend27\_\discretionary{}{}{}upper27\_\discretionary{}{}{}lower10}. The result is stored into the \emph{registerW}.


\paragraph{Execution} {Reservation table \textsf{ALU\_TINY.Y} (Section~\ref{sec:reservation-ALU-TINY.Y})}
~
{\begin{lstlisting}
stage RR:
  new argument3 = _WX_64(extend27_upper27_lower10);
  new argument2 = GPR[registerZ];
  new result1 = ~(argument2 & argument3);
stage E1:
  GPR[registerW] = result1;
\end{lstlisting}}
\newpage
\subsubsection[{\small NANDQ: Bitwise Not And Quadruple Word}]{NANDQ\hfill Bitwise Not And Quadruple Word}
\label{instr:NANDQ}\index{nandq}
 
\paragraph{Encoding} Format \textsf{ALU\_QWRR0} (Section~\ref{sec:format-ALU-QWRR0}), \verb|exucode2 = 1001|\\

\bitfields{ 1/P, 2/11, 1/0, 4/1001, 5/registerM, 1/--, 2/10, 3/000, 1/1, 5/registerO, 1/--, 5/registerP, 1/-- }


\paragraph{Syntax} \texttt{nandq} \emph{registerM} = \emph{registerP}, \emph{registerO}

\paragraph{Description} Bitwise not of bitwise and of the \emph{registerP} with the \emph{registerO}. The result is stored into the \emph{registerM}.


\paragraph{Execution} {Reservation table \textsf{ALU\_LITE} (Section~\ref{sec:reservation-ALU-LITE})}
~
{\begin{lstlisting}
stage RR:
  new argument3 = PGR[registerO];
  new argument2 = PGR[registerP];
  new result1 = ~(argument2 & argument3);
stage E1:
  PGR[registerM] = result1;
\end{lstlisting}}
\newpage

\paragraph{Encoding} Format \textsf{ALU\_QWRR0.M} (Section~\ref{sec:format-ALU-QWRR0.M}), \verb|exucode2 = 1001|\\

\bitfields{ 1/P, 2/00, 2/-- --, 27/upper27, 1/1, 2/11, 1/0, 4/1001, 5/registerM, 1/--, 2/10, 3/000, 1/1, 1/splat32, 5/lower5, 5/registerP, 1/1 }


\paragraph{Syntax} \texttt{nandq} \emph{registerM} = \emph{registerP}, \emph{upper27\_\discretionary{}{}{}lower5}\emph{splat32}

\paragraph{Description} Bitwise not of bitwise and of the \emph{registerP} with the \emph{upper27\_\discretionary{}{}{}lower5}. The result is stored into the \emph{registerM}.


\paragraph{Modifier(s)}
\noindent\begin{itemize}
\item splat32 (cf. splat32 in section \ref{par:modifier-splat32} on page \pageref{par:modifier-splat32})
\end{itemize}

\paragraph{Execution} {Reservation table \textsf{ALU\_LITE.X} (Section~\ref{sec:reservation-ALU-LITE.X})}
~
{\begin{lstlisting}
stage ID:
  new splat32 = _ZX_1(splat32);
stage RR:
  new argument3 = splat32 ? (_WX_32(upper27_lower5) << 32) | _ZX_32(_WX_32(upper27_lower5)) : _SX_32(_WX_32(upper27_lower5));
  new argument2 = PGR[registerP];
  new result1 = ~(argument2 & argument3);
stage E1:
  PGR[registerM] = result1;
\end{lstlisting}}
\newpage
\subsubsection[{\small NANDW: Bitwise Not And Word}]{NANDW\hfill Bitwise Not And Word}
\label{instr:NANDW}\index{nandw}
 
\paragraph{Encoding} Format \textsf{ALU\_WWRR0} (Section~\ref{sec:format-ALU-WWRR0}), \verb|exucode2 = 1001|\\

\bitfields{ 1/P, 2/11, 1/0, 4/1001, 6/registerW, 2/11, 3/000, 1/signextw, 6/registerY, 6/registerZ }


\paragraph{Syntax} \texttt{nandw}\emph{signextw} \emph{registerW} = \emph{registerZ}, \emph{registerY}

\paragraph{Description} Bitwise not of bitwise and of the \emph{registerZ} with the \emph{registerY}. The result with the 32 upper bits cleared is stored into the \emph{registerW}.


\paragraph{Modifier(s)}
\noindent\begin{itemize}
\item signextw (cf. signextw in section \ref{par:modifier-signextw} on page \pageref{par:modifier-signextw})
\end{itemize}

\paragraph{Execution} {Reservation table \textsf{ALU\_TINY} (Section~\ref{sec:reservation-ALU-TINY})}
~
{\begin{lstlisting}
stage ID:
  new signextw = _ZX_1(signextw);
stage RR:
  new argument3 = GPR[registerY];
  new argument2 = GPR[registerZ];
  new result1 = ~(argument2 & argument3);
stage E1:
  GPR[registerW] = signextw ? _SX_32(result1) : _ZX_32(result1);
\end{lstlisting}}
\newpage

\paragraph{Encoding} Format \textsf{ALU\_WWRR0.W} (Section~\ref{sec:format-ALU-WWRR0.W}), \verb|exucode2 = 1001|\\

\bitfields{ 1/P, 2/00, 2/-- --, 27/upper27, 1/1, 2/11, 1/0, 4/1001, 6/registerW, 2/11, 3/000, 1/signextw, 1/0, 5/lower5, 6/registerZ }


\paragraph{Syntax} \texttt{nandw}\emph{signextw} \emph{registerW} = \emph{registerZ}, \emph{upper27\_\discretionary{}{}{}lower5}

\paragraph{Description} Bitwise not of bitwise and of the \emph{registerZ} with the \emph{upper27\_\discretionary{}{}{}lower5}. The result with the 32 upper bits cleared is stored into the \emph{registerW}.


\paragraph{Modifier(s)}
\noindent\begin{itemize}
\item signextw (cf. signextw in section \ref{par:modifier-signextw} on page \pageref{par:modifier-signextw})
\end{itemize}

\paragraph{Execution} {Reservation table \textsf{ALU\_TINY.X} (Section~\ref{sec:reservation-ALU-TINY.X})}
~
{\begin{lstlisting}
stage ID:
  new signextw = _ZX_1(signextw);
stage RR:
  new argument3 = _WX_32(upper27_lower5);
  new argument2 = GPR[registerZ];
  new result1 = ~(argument2 & argument3);
stage E1:
  GPR[registerW] = signextw ? _SX_32(result1) : _ZX_32(result1);
\end{lstlisting}}
\newpage
\subsubsection[{\small NEGBO: Negate Byte Octuple}]{NEGBO\hfill Negate Byte Octuple}
\label{instr:NEGBO}\index{negbo}
 
\paragraph{Encoding} Format \textsf{ALU\_BOBWRE} (Section~\ref{sec:format-ALU-BOBWRE}), \verb|alucode8 = 0000|\\

\bitfields{ 1/P, 2/11, 1/1, 4/1111, 5/registerWe, 1/0, 2/10, 3/101, 1/1, 4/0000, 1/--, 1/--, 5/registerZe, 1/1 }


\paragraph{Syntax} \texttt{negbo} \emph{registerWe} = \emph{registerZe}

\paragraph{Description} The negated values of the \emph{registerZe} interpreted as a byte octuple are stored into the \emph{registerWe}.


\paragraph{Execution} {Reservation table \textsf{ALU\_LITE} (Section~\ref{sec:reservation-ALU-LITE})}
~
{\begin{lstlisting}
stage RR:
  new argument2 = GPR[registerZe<<1];
  for (I = 0; I < 8; I += 1) {
    result1.8[I] = -_SX_8(argument2.8[I])
  };
stage E1:
  GPR[registerWe<<1] = result1;
\end{lstlisting}}
\newpage

\paragraph{Encoding} Format \textsf{ALU\_BOBWRO} (Section~\ref{sec:format-ALU-BOBWRO}), \verb|alucode8 = 0000|\\

\bitfields{ 1/P, 2/11, 1/1, 4/1111, 5/registerWo, 1/1, 2/10, 3/101, 1/1, 4/0000, 1/--, 1/--, 5/registerZo, 1/1 }


\paragraph{Syntax} \texttt{negbo} \emph{registerWo} = \emph{registerZo}

\paragraph{Description} The negated values of the \emph{registerZo} interpreted as a byte octuple are stored into the \emph{registerWo}.


\paragraph{Execution} {Reservation table \textsf{ALU\_LITE} (Section~\ref{sec:reservation-ALU-LITE})}
~
{\begin{lstlisting}
stage RR:
  new argument2 = GPR[(registerZo<<1)+1];
  for (I = 0; I < 8; I += 1) {
    result1.8[I] = -_SX_8(argument2.8[I])
  };
stage E1:
  GPR[(registerWo<<1)+1] = result1;
\end{lstlisting}}
\newpage
\subsubsection[{\small NEGD: Negate Double Word}]{NEGD\hfill Negate Double Word}
\label{instr:NEGD}\index{negd}
 
\paragraph{Encoding} Format \textsf{ALU\_DBWR} (Section~\ref{sec:format-ALU-DBWR}), \verb|alucode8 = 0000|\\

\bitfields{ 1/P, 2/11, 1/0, 4/1111, 6/registerW, 2/10, 3/101, 1/0, 4/0000, 1/0, 1/0, 6/registerZ }


\paragraph{Syntax} \texttt{negd} \emph{registerW} = \emph{registerZ}

\paragraph{Description} The 64-bit negated value of the \emph{registerZ} is stored into the \emph{registerW}.


\paragraph{Execution} {Reservation table \textsf{ALU\_TINY} (Section~\ref{sec:reservation-ALU-TINY})}
~
{\begin{lstlisting}
stage RR:
  new argument2 = GPR[registerZ];
  new result1 = -_SX_64(argument2);
stage E1:
  GPR[registerW] = result1;
\end{lstlisting}}
\newpage
\subsubsection[{\small NEGHQ: Negate Half Word Quadruple}]{NEGHQ\hfill Negate Half Word Quadruple}
\label{instr:NEGHQ}\index{neghq}
 
\paragraph{Encoding} Format \textsf{ALU\_HQBWRE} (Section~\ref{sec:format-ALU-HQBWRE}), \verb|alucode8 = 0000|\\

\bitfields{ 1/P, 2/11, 1/1, 4/1111, 5/registerWe, 1/0, 2/10, 3/101, 1/1, 4/0000, 1/--, 1/--, 5/registerZe, 1/0 }


\paragraph{Syntax} \texttt{neghq} \emph{registerWe} = \emph{registerZe}

\paragraph{Description} The negated values of the \emph{registerZe} interpreted as a half word quadruple are stored into the \emph{registerWe}.


\paragraph{Execution} {Reservation table \textsf{ALU\_LITE} (Section~\ref{sec:reservation-ALU-LITE})}
~
{\begin{lstlisting}
stage RR:
  new argument2 = GPR[registerZe<<1];
  for (I = 0; I < 4; I += 1) {
    result1.16[I] = -_SX_16(argument2.16[I])
  };
stage E1:
  GPR[registerWe<<1] = result1;
\end{lstlisting}}
\newpage

\paragraph{Encoding} Format \textsf{ALU\_HQBWRO} (Section~\ref{sec:format-ALU-HQBWRO}), \verb|alucode8 = 0000|\\

\bitfields{ 1/P, 2/11, 1/1, 4/1111, 5/registerWo, 1/1, 2/10, 3/101, 1/1, 4/0000, 1/--, 1/--, 5/registerZo, 1/0 }


\paragraph{Syntax} \texttt{neghq} \emph{registerWo} = \emph{registerZo}

\paragraph{Description} The negated values of the \emph{registerZo} interpreted as a half word quadruple are stored into the \emph{registerWo}.


\paragraph{Execution} {Reservation table \textsf{ALU\_LITE} (Section~\ref{sec:reservation-ALU-LITE})}
~
{\begin{lstlisting}
stage RR:
  new argument2 = GPR[(registerZo<<1)+1];
  for (I = 0; I < 4; I += 1) {
    result1.16[I] = -_SX_16(argument2.16[I])
  };
stage E1:
  GPR[(registerWo<<1)+1] = result1;
\end{lstlisting}}
\newpage
\subsubsection[{\small NEGQ: Negate Quadruple Word}]{NEGQ\hfill Negate Quadruple Word}
\label{instr:NEGQ}\index{negq}
 
\paragraph{Encoding} Format \textsf{ALU\_QBWR} (Section~\ref{sec:format-ALU-QBWR}), \verb|alucode8 = 0000|\\

\bitfields{ 1/P, 2/11, 1/0, 4/1111, 5/registerM, 1/--, 2/10, 3/101, 1/1, 4/0000, 1/0, 1/0, 5/registerP, 1/-- }


\paragraph{Syntax} \texttt{negq} \emph{registerM} = \emph{registerP}

\paragraph{Description} The 128-bit negated value of the \emph{registerP} is stored into the \emph{registerM}.


\paragraph{Execution} {Reservation table \textsf{ALU\_LITE} (Section~\ref{sec:reservation-ALU-LITE})}
~
{\begin{lstlisting}
stage RR:
  new argument2 = PGR[registerP];
  new result1 = -_SX_128(argument2);
stage E1:
  PGR[registerM] = result1;
\end{lstlisting}}
\newpage
\subsubsection[{\small NEGSBO: Negate Saturated Byte Octuple}]{NEGSBO\hfill Negate Saturated Byte Octuple}
\label{instr:NEGSBO}\index{negsbo}
 
\paragraph{Encoding} Format \textsf{ALU\_BOBWRE} (Section~\ref{sec:format-ALU-BOBWRE}), \verb|alucode8 = 0010|\\

\bitfields{ 1/P, 2/11, 1/1, 4/1111, 5/registerWe, 1/0, 2/10, 3/101, 1/1, 4/0010, 1/--, 1/--, 5/registerZe, 1/1 }


\paragraph{Syntax} \texttt{negsbo} \emph{registerWe} = \emph{registerZe}

\paragraph{Description} The negated saturated values of the \emph{registerZe} interpreted as a byte octuple are stored into the \emph{registerWe}.


\paragraph{Execution} {Reservation table \textsf{ALU\_LITE} (Section~\ref{sec:reservation-ALU-LITE})}
~
{\begin{lstlisting}
stage RR:
  new argument2 = GPR[registerZe<<1];
  for (I = 0; I < 8; I += 1) {
    result1.8[I] = _SAT_8(-_SX_8(argument2.8[I]))
  };
stage E1:
  GPR[registerWe<<1] = result1;
\end{lstlisting}}
\newpage

\paragraph{Encoding} Format \textsf{ALU\_BOBWRO} (Section~\ref{sec:format-ALU-BOBWRO}), \verb|alucode8 = 0010|\\

\bitfields{ 1/P, 2/11, 1/1, 4/1111, 5/registerWo, 1/1, 2/10, 3/101, 1/1, 4/0010, 1/--, 1/--, 5/registerZo, 1/1 }


\paragraph{Syntax} \texttt{negsbo} \emph{registerWo} = \emph{registerZo}

\paragraph{Description} The negated saturated values of the \emph{registerZo} interpreted as a byte octuple are stored into the \emph{registerWo}.


\paragraph{Execution} {Reservation table \textsf{ALU\_LITE} (Section~\ref{sec:reservation-ALU-LITE})}
~
{\begin{lstlisting}
stage RR:
  new argument2 = GPR[(registerZo<<1)+1];
  for (I = 0; I < 8; I += 1) {
    result1.8[I] = _SAT_8(-_SX_8(argument2.8[I]))
  };
stage E1:
  GPR[(registerWo<<1)+1] = result1;
\end{lstlisting}}
\newpage
\subsubsection[{\small NEGSD: Negate Saturated Double Word}]{NEGSD\hfill Negate Saturated Double Word}
\label{instr:NEGSD}\index{negsd}
 
\paragraph{Encoding} Format \textsf{ALU\_DBWR} (Section~\ref{sec:format-ALU-DBWR}), \verb|alucode8 = 0010|\\

\bitfields{ 1/P, 2/11, 1/0, 4/1111, 6/registerW, 2/10, 3/101, 1/0, 4/0010, 1/0, 1/0, 6/registerZ }


\paragraph{Syntax} \texttt{negsd} \emph{registerW} = \emph{registerZ}

\paragraph{Description} The 64-bit negated value of the \emph{registerZ} is saturated and stored into the \emph{registerW}.


\paragraph{Execution} {Reservation table \textsf{ALU\_TINY} (Section~\ref{sec:reservation-ALU-TINY})}
~
{\begin{lstlisting}
stage RR:
  new argument2 = GPR[registerZ];
  new result1 = _SAT_64(-_SX_64(argument2));
stage E1:
  GPR[registerW] = result1;
\end{lstlisting}}
\newpage
\subsubsection[{\small NEGSHQ: Negate Saturated Half Word Quadruple}]{NEGSHQ\hfill Negate Saturated Half Word Quadruple}
\label{instr:NEGSHQ}\index{negshq}
 
\paragraph{Encoding} Format \textsf{ALU\_HQBWRE} (Section~\ref{sec:format-ALU-HQBWRE}), \verb|alucode8 = 0010|\\

\bitfields{ 1/P, 2/11, 1/1, 4/1111, 5/registerWe, 1/0, 2/10, 3/101, 1/1, 4/0010, 1/--, 1/--, 5/registerZe, 1/0 }


\paragraph{Syntax} \texttt{negshq} \emph{registerWe} = \emph{registerZe}

\paragraph{Description} The negated saturated values of the \emph{registerZe} interpreted as a half word quadruple are stored into the \emph{registerWe}.


\paragraph{Execution} {Reservation table \textsf{ALU\_LITE} (Section~\ref{sec:reservation-ALU-LITE})}
~
{\begin{lstlisting}
stage RR:
  new argument2 = GPR[registerZe<<1];
  for (I = 0; I < 4; I += 1) {
    result1.16[I] = _SAT_16(-_SX_16(argument2.16[I]))
  };
stage E1:
  GPR[registerWe<<1] = result1;
\end{lstlisting}}
\newpage

\paragraph{Encoding} Format \textsf{ALU\_HQBWRO} (Section~\ref{sec:format-ALU-HQBWRO}), \verb|alucode8 = 0010|\\

\bitfields{ 1/P, 2/11, 1/1, 4/1111, 5/registerWo, 1/1, 2/10, 3/101, 1/1, 4/0010, 1/--, 1/--, 5/registerZo, 1/0 }


\paragraph{Syntax} \texttt{negshq} \emph{registerWo} = \emph{registerZo}

\paragraph{Description} The negated saturated values of the \emph{registerZo} interpreted as a half word quadruple are stored into the \emph{registerWo}.


\paragraph{Execution} {Reservation table \textsf{ALU\_LITE} (Section~\ref{sec:reservation-ALU-LITE})}
~
{\begin{lstlisting}
stage RR:
  new argument2 = GPR[(registerZo<<1)+1];
  for (I = 0; I < 4; I += 1) {
    result1.16[I] = _SAT_16(-_SX_16(argument2.16[I]))
  };
stage E1:
  GPR[(registerWo<<1)+1] = result1;
\end{lstlisting}}
\newpage
\subsubsection[{\small NEGSW: Negate Saturated Word}]{NEGSW\hfill Negate Saturated Word}
\label{instr:NEGSW}\index{negsw}
 
\paragraph{Encoding} Format \textsf{ALU\_BWRW} (Section~\ref{sec:format-ALU-BWRW}), \verb|alucode8 = 0010|\\

\bitfields{ 1/P, 2/11, 1/0, 4/1111, 6/registerW, 2/11, 3/101, 1/signextw, 4/0010, 1/0, 1/0, 6/registerZ }


\paragraph{Syntax} \texttt{negsw} \emph{registerW} = \emph{registerZ}

\paragraph{Description} The 32-bit negated value of the \emph{registerZ} saturated and is stored into the \emph{registerW}.


\paragraph{Execution} {Reservation table \textsf{ALU\_TINY} (Section~\ref{sec:reservation-ALU-TINY})}
~
{\begin{lstlisting}
stage RR:
  new argument2 = GPR[registerZ];
  new result1 = _SAT_32(-_SX_32(argument2));
stage E1:
  GPR[registerW] = _ZX_32(result1);
\end{lstlisting}}
\newpage
\subsubsection[{\small NEGSWP: Negate Saturated Word Pair}]{NEGSWP\hfill Negate Saturated Word Pair}
\label{instr:NEGSWP}\index{negswp}
 
\paragraph{Encoding} Format \textsf{ALU\_WPBWRE} (Section~\ref{sec:format-ALU-WPBWRE}), \verb|alucode8 = 0010|\\

\bitfields{ 1/P, 2/11, 1/1, 4/1111, 5/registerWe, 1/0, 2/10, 3/101, 1/0, 4/0010, 1/--, 1/--, 5/registerZe, 1/1 }


\paragraph{Syntax} \texttt{negswp} \emph{registerWe} = \emph{registerZe}

\paragraph{Description} The negated saturated values of the \emph{registerZe} interpreted as a word pair are stored into the \emph{registerWe}.


\paragraph{Execution} {Reservation table \textsf{ALU\_LITE} (Section~\ref{sec:reservation-ALU-LITE})}
~
{\begin{lstlisting}
stage RR:
  new argument2 = GPR[registerZe<<1];
  for (I = 0; I < 2; I += 1) {
    result1.32[I] = _SAT_32(-_SX_32(argument2.32[I]))
  };
stage E1:
  GPR[registerWe<<1] = result1;
\end{lstlisting}}
\newpage

\paragraph{Encoding} Format \textsf{ALU\_WPBWRO} (Section~\ref{sec:format-ALU-WPBWRO}), \verb|alucode8 = 0010|\\

\bitfields{ 1/P, 2/11, 1/1, 4/1111, 5/registerWo, 1/1, 2/10, 3/101, 1/0, 4/0010, 1/--, 1/--, 5/registerZo, 1/1 }


\paragraph{Syntax} \texttt{negswp} \emph{registerWo} = \emph{registerZo}

\paragraph{Description} The negated saturated values of the \emph{registerZo} interpreted as a word pair are stored into the \emph{registerWo}.


\paragraph{Execution} {Reservation table \textsf{ALU\_LITE} (Section~\ref{sec:reservation-ALU-LITE})}
~
{\begin{lstlisting}
stage RR:
  new argument2 = GPR[(registerZo<<1)+1];
  for (I = 0; I < 2; I += 1) {
    result1.32[I] = _SAT_32(-_SX_32(argument2.32[I]))
  };
stage E1:
  GPR[(registerWo<<1)+1] = result1;
\end{lstlisting}}
\newpage
\subsubsection[{\small NEGW: Negate Word}]{NEGW\hfill Negate Word}
\label{instr:NEGW}\index{negw}
 
\paragraph{Encoding} Format \textsf{ALU\_BWRW} (Section~\ref{sec:format-ALU-BWRW}), \verb|alucode8 = 0000|\\

\bitfields{ 1/P, 2/11, 1/0, 4/1111, 6/registerW, 2/11, 3/101, 1/signextw, 4/0000, 1/0, 1/0, 6/registerZ }


\paragraph{Syntax} \texttt{negw} \emph{registerW} = \emph{registerZ}

\paragraph{Description} The 32-bit negated value of the \emph{registerZ} is stored into the \emph{registerW}.


\paragraph{Execution} {Reservation table \textsf{ALU\_TINY} (Section~\ref{sec:reservation-ALU-TINY})}
~
{\begin{lstlisting}
stage RR:
  new argument2 = GPR[registerZ];
  new result1 = -_SX_32(argument2);
stage E1:
  GPR[registerW] = _ZX_32(result1);
\end{lstlisting}}
\newpage
\subsubsection[{\small NEGWP: Negate Word Pair}]{NEGWP\hfill Negate Word Pair}
\label{instr:NEGWP}\index{negwp}
 
\paragraph{Encoding} Format \textsf{ALU\_WPBWRE} (Section~\ref{sec:format-ALU-WPBWRE}), \verb|alucode8 = 0000|\\

\bitfields{ 1/P, 2/11, 1/1, 4/1111, 5/registerWe, 1/0, 2/10, 3/101, 1/0, 4/0000, 1/--, 1/--, 5/registerZe, 1/1 }


\paragraph{Syntax} \texttt{negwp} \emph{registerWe} = \emph{registerZe}

\paragraph{Description} The negated values of the \emph{registerZe} interpreted as a word pair are stored into the \emph{registerWe}.


\paragraph{Execution} {Reservation table \textsf{ALU\_LITE} (Section~\ref{sec:reservation-ALU-LITE})}
~
{\begin{lstlisting}
stage RR:
  new argument2 = GPR[registerZe<<1];
  for (I = 0; I < 2; I += 1) {
    result1.32[I] = -_SX_32(argument2.32[I])
  };
stage E1:
  GPR[registerWe<<1] = result1;
\end{lstlisting}}
\newpage

\paragraph{Encoding} Format \textsf{ALU\_WPBWRO} (Section~\ref{sec:format-ALU-WPBWRO}), \verb|alucode8 = 0000|\\

\bitfields{ 1/P, 2/11, 1/1, 4/1111, 5/registerWo, 1/1, 2/10, 3/101, 1/0, 4/0000, 1/--, 1/--, 5/registerZo, 1/1 }


\paragraph{Syntax} \texttt{negwp} \emph{registerWo} = \emph{registerZo}

\paragraph{Description} The negated values of the \emph{registerZo} interpreted as a word pair are stored into the \emph{registerWo}.


\paragraph{Execution} {Reservation table \textsf{ALU\_LITE} (Section~\ref{sec:reservation-ALU-LITE})}
~
{\begin{lstlisting}
stage RR:
  new argument2 = GPR[(registerZo<<1)+1];
  for (I = 0; I < 2; I += 1) {
    result1.32[I] = -_SX_32(argument2.32[I])
  };
stage E1:
  GPR[(registerWo<<1)+1] = result1;
\end{lstlisting}}
\newpage
\subsubsection[{\small NEORD: Bitwise Not Exclusive Or Double Word}]{NEORD\hfill Bitwise Not Exclusive Or Double Word}
\label{instr:NEORD}\index{neord}
 
\paragraph{Encoding} Format \textsf{ALU\_DWRR0} (Section~\ref{sec:format-ALU-DWRR0}), \verb|exucode2 = 1101|\\

\bitfields{ 1/P, 2/11, 1/0, 4/1101, 6/registerW, 2/10, 3/000, 1/0, 6/registerY, 6/registerZ }


\paragraph{Syntax} \texttt{neord} \emph{registerW} = \emph{registerZ}, \emph{registerY}

\paragraph{Description} Bitwise not of bitwise exclusive or of the \emph{registerZ} with the \emph{registerY}. The result is stored into the \emph{registerW}.


\paragraph{Execution} {Reservation table \textsf{ALU\_TINY} (Section~\ref{sec:reservation-ALU-TINY})}
~
{\begin{lstlisting}
stage RR:
  new argument3 = GPR[registerY];
  new argument2 = GPR[registerZ];
  new result1 = ~(argument2 ^ argument3);
stage E1:
  GPR[registerW] = result1;
\end{lstlisting}}
\newpage

\paragraph{Encoding} Format \textsf{ALU\_DWRR0.M} (Section~\ref{sec:format-ALU-DWRR0.M}), \verb|exucode2 = 1101|\\

\bitfields{ 1/P, 2/00, 2/-- --, 27/upper27, 1/1, 2/11, 1/0, 4/1101, 6/registerW, 2/10, 3/000, 1/0, 1/splat32, 5/lower5, 6/registerZ }


\paragraph{Syntax} \texttt{neord} \emph{registerW} = \emph{registerZ}, \emph{upper27\_\discretionary{}{}{}lower5}\emph{splat32}

\paragraph{Description} Bitwise not of bitwise exclusive or of the \emph{registerZ} with the \emph{upper27\_\discretionary{}{}{}lower5}. The result is stored into the \emph{registerW}.


\paragraph{Modifier(s)}
\noindent\begin{itemize}
\item splat32 (cf. splat32 in section \ref{par:modifier-splat32} on page \pageref{par:modifier-splat32})
\end{itemize}

\paragraph{Execution} {Reservation table \textsf{ALU\_TINY.X} (Section~\ref{sec:reservation-ALU-TINY.X})}
~
{\begin{lstlisting}
stage ID:
  new splat32 = _ZX_1(splat32);
stage RR:
  new argument3 = splat32 ? (_WX_32(upper27_lower5) << 32) | _ZX_32(_WX_32(upper27_lower5)) : _SX_32(_WX_32(upper27_lower5));
  new argument2 = GPR[registerZ];
  new result1 = ~(argument2 ^ argument3);
stage E1:
  GPR[registerW] = result1;
\end{lstlisting}}
\newpage

\paragraph{Encoding} Format \textsf{ALU\_DWRI} (Section~\ref{sec:format-ALU-DWRI}), \verb|exucode2 = 1101|\\

\bitfields{ 1/P, 2/11, 1/0, 4/1101, 6/registerW, 2/00, 10/signed10, 6/registerZ }


\paragraph{Syntax} \texttt{neord} \emph{registerW} = \emph{registerZ}, \emph{signed10}

\paragraph{Description} Bitwise not of bitwise exclusive or of the \emph{registerZ} with the \emph{signed10}. The result is stored into the \emph{registerW}.


\paragraph{Execution} {Reservation table \textsf{ALU\_TINY} (Section~\ref{sec:reservation-ALU-TINY})}
~
{\begin{lstlisting}
stage RR:
  new argument3 = _SX_10(signed10);
  new argument2 = GPR[registerZ];
  new result1 = ~(argument2 ^ argument3);
stage E1:
  GPR[registerW] = result1;
\end{lstlisting}}
\newpage

\paragraph{Encoding} Format \textsf{ALU\_DWRI.X} (Section~\ref{sec:format-ALU-DWRI.X}), \verb|exucode2 = 1101|\\

\bitfields{ 1/P, 2/00, 2/-- --, 27/upper27, 1/1, 2/11, 1/0, 4/1101, 6/registerW, 2/00, 10/lower10, 6/registerZ }


\paragraph{Syntax} \texttt{neord} \emph{registerW} = \emph{registerZ}, \emph{upper27\_\discretionary{}{}{}lower10}

\paragraph{Description} Bitwise not of bitwise exclusive or of the \emph{registerZ} with the \emph{upper27\_\discretionary{}{}{}lower10}. The result is stored into the \emph{registerW}.


\paragraph{Execution} {Reservation table \textsf{ALU\_TINY.X} (Section~\ref{sec:reservation-ALU-TINY.X})}
~
{\begin{lstlisting}
stage RR:
  new argument3 = _SX_37(upper27_lower10);
  new argument2 = GPR[registerZ];
  new result1 = ~(argument2 ^ argument3);
stage E1:
  GPR[registerW] = result1;
\end{lstlisting}}
\newpage

\paragraph{Encoding} Format \textsf{ALU\_DWRI.Y} (Section~\ref{sec:format-ALU-DWRI.Y}), \verb|exucode2 = 1101|\\

\bitfields{ 1/P, 2/00, 2/-- --, 27/extend27, 1/1, 2/00, 2/-- --, 27/upper27, 1/1, 2/11, 1/0, 4/1101, 6/registerW, 2/00, 10/lower10, 6/registerZ }


\paragraph{Syntax} \texttt{neord} \emph{registerW} = \emph{registerZ}, \emph{extend27\_\discretionary{}{}{}upper27\_\discretionary{}{}{}lower10}

\paragraph{Description} Bitwise not of bitwise exclusive or of the \emph{registerZ} with the \emph{extend27\_\discretionary{}{}{}upper27\_\discretionary{}{}{}lower10}. The result is stored into the \emph{registerW}.


\paragraph{Execution} {Reservation table \textsf{ALU\_TINY.Y} (Section~\ref{sec:reservation-ALU-TINY.Y})}
~
{\begin{lstlisting}
stage RR:
  new argument3 = _WX_64(extend27_upper27_lower10);
  new argument2 = GPR[registerZ];
  new result1 = ~(argument2 ^ argument3);
stage E1:
  GPR[registerW] = result1;
\end{lstlisting}}
\newpage
\subsubsection[{\small NEORQ: Bitwise Not Exclusive Or Quadruple Word}]{NEORQ\hfill Bitwise Not Exclusive Or Quadruple Word}
\label{instr:NEORQ}\index{neorq}
 
\paragraph{Encoding} Format \textsf{ALU\_QWRR0} (Section~\ref{sec:format-ALU-QWRR0}), \verb|exucode2 = 1101|\\

\bitfields{ 1/P, 2/11, 1/0, 4/1101, 5/registerM, 1/--, 2/10, 3/000, 1/1, 5/registerO, 1/--, 5/registerP, 1/-- }


\paragraph{Syntax} \texttt{neorq} \emph{registerM} = \emph{registerP}, \emph{registerO}

\paragraph{Description} Bitwise not of bitwise exclusive or of the \emph{registerP} with the \emph{registerO}. The result is stored into the \emph{registerM}.


\paragraph{Execution} {Reservation table \textsf{ALU\_LITE} (Section~\ref{sec:reservation-ALU-LITE})}
~
{\begin{lstlisting}
stage RR:
  new argument3 = PGR[registerO];
  new argument2 = PGR[registerP];
  new result1 = ~(argument2 ^ argument3);
stage E1:
  PGR[registerM] = result1;
\end{lstlisting}}
\newpage

\paragraph{Encoding} Format \textsf{ALU\_QWRR0.M} (Section~\ref{sec:format-ALU-QWRR0.M}), \verb|exucode2 = 1101|\\

\bitfields{ 1/P, 2/00, 2/-- --, 27/upper27, 1/1, 2/11, 1/0, 4/1101, 5/registerM, 1/--, 2/10, 3/000, 1/1, 1/splat32, 5/lower5, 5/registerP, 1/1 }


\paragraph{Syntax} \texttt{neorq} \emph{registerM} = \emph{registerP}, \emph{upper27\_\discretionary{}{}{}lower5}\emph{splat32}

\paragraph{Description} Bitwise not of bitwise exclusive or of the \emph{registerP} with the \emph{upper27\_\discretionary{}{}{}lower5}. The result is stored into the \emph{registerM}.


\paragraph{Modifier(s)}
\noindent\begin{itemize}
\item splat32 (cf. splat32 in section \ref{par:modifier-splat32} on page \pageref{par:modifier-splat32})
\end{itemize}

\paragraph{Execution} {Reservation table \textsf{ALU\_LITE.X} (Section~\ref{sec:reservation-ALU-LITE.X})}
~
{\begin{lstlisting}
stage ID:
  new splat32 = _ZX_1(splat32);
stage RR:
  new argument3 = splat32 ? (_WX_32(upper27_lower5) << 32) | _ZX_32(_WX_32(upper27_lower5)) : _SX_32(_WX_32(upper27_lower5));
  new argument2 = PGR[registerP];
  new result1 = ~(argument2 ^ argument3);
stage E1:
  PGR[registerM] = result1;
\end{lstlisting}}
\newpage
\subsubsection[{\small NEORW: Bitwise Not Exclusive Or Word}]{NEORW\hfill Bitwise Not Exclusive Or Word}
\label{instr:NEORW}\index{neorw}
 
\paragraph{Encoding} Format \textsf{ALU\_WWRR0} (Section~\ref{sec:format-ALU-WWRR0}), \verb|exucode2 = 1101|\\

\bitfields{ 1/P, 2/11, 1/0, 4/1101, 6/registerW, 2/11, 3/000, 1/signextw, 6/registerY, 6/registerZ }


\paragraph{Syntax} \texttt{neorw}\emph{signextw} \emph{registerW} = \emph{registerZ}, \emph{registerY}

\paragraph{Description} Bitwise not of bitwise exclusive or of the \emph{registerZ} with the \emph{registerY}. The result with the 32 upper bits cleared is stored into the \emph{registerW}.


\paragraph{Modifier(s)}
\noindent\begin{itemize}
\item signextw (cf. signextw in section \ref{par:modifier-signextw} on page \pageref{par:modifier-signextw})
\end{itemize}

\paragraph{Execution} {Reservation table \textsf{ALU\_TINY} (Section~\ref{sec:reservation-ALU-TINY})}
~
{\begin{lstlisting}
stage ID:
  new signextw = _ZX_1(signextw);
stage RR:
  new argument3 = GPR[registerY];
  new argument2 = GPR[registerZ];
  new result1 = ~(argument2 ^ argument3);
stage E1:
  GPR[registerW] = signextw ? _SX_32(result1) : _ZX_32(result1);
\end{lstlisting}}
\newpage

\paragraph{Encoding} Format \textsf{ALU\_WWRR0.W} (Section~\ref{sec:format-ALU-WWRR0.W}), \verb|exucode2 = 1101|\\

\bitfields{ 1/P, 2/00, 2/-- --, 27/upper27, 1/1, 2/11, 1/0, 4/1101, 6/registerW, 2/11, 3/000, 1/signextw, 1/0, 5/lower5, 6/registerZ }


\paragraph{Syntax} \texttt{neorw}\emph{signextw} \emph{registerW} = \emph{registerZ}, \emph{upper27\_\discretionary{}{}{}lower5}

\paragraph{Description} Bitwise not of bitwise exclusive or of the \emph{registerZ} with the \emph{upper27\_\discretionary{}{}{}lower5}. The result with the 32 upper bits cleared is stored into the \emph{registerW}.


\paragraph{Modifier(s)}
\noindent\begin{itemize}
\item signextw (cf. signextw in section \ref{par:modifier-signextw} on page \pageref{par:modifier-signextw})
\end{itemize}

\paragraph{Execution} {Reservation table \textsf{ALU\_TINY.X} (Section~\ref{sec:reservation-ALU-TINY.X})}
~
{\begin{lstlisting}
stage ID:
  new signextw = _ZX_1(signextw);
stage RR:
  new argument3 = _WX_32(upper27_lower5);
  new argument2 = GPR[registerZ];
  new result1 = ~(argument2 ^ argument3);
stage E1:
  GPR[registerW] = signextw ? _SX_32(result1) : _ZX_32(result1);
\end{lstlisting}}
\newpage
\subsubsection[{\small NIORD: Bitwise Not Inclusive Or Double Word}]{NIORD\hfill Bitwise Not Inclusive Or Double Word}
\label{instr:NIORD}\index{niord}
 
\paragraph{Encoding} Format \textsf{ALU\_DWRR0} (Section~\ref{sec:format-ALU-DWRR0}), \verb|exucode2 = 1011|\\

\bitfields{ 1/P, 2/11, 1/0, 4/1011, 6/registerW, 2/10, 3/000, 1/0, 6/registerY, 6/registerZ }


\paragraph{Syntax} \texttt{niord} \emph{registerW} = \emph{registerZ}, \emph{registerY}

\paragraph{Description} Bitwise not of bitwise inclusive or of the \emph{registerZ} with the \emph{registerY}. The result is stored into the \emph{registerW}.


\paragraph{Execution} {Reservation table \textsf{ALU\_TINY} (Section~\ref{sec:reservation-ALU-TINY})}
~
{\begin{lstlisting}
stage RR:
  new argument3 = GPR[registerY];
  new argument2 = GPR[registerZ];
  new result1 = ~(argument2 | argument3);
stage E1:
  GPR[registerW] = result1;
\end{lstlisting}}
\newpage

\paragraph{Encoding} Format \textsf{ALU\_DWRR0.M} (Section~\ref{sec:format-ALU-DWRR0.M}), \verb|exucode2 = 1011|\\

\bitfields{ 1/P, 2/00, 2/-- --, 27/upper27, 1/1, 2/11, 1/0, 4/1011, 6/registerW, 2/10, 3/000, 1/0, 1/splat32, 5/lower5, 6/registerZ }


\paragraph{Syntax} \texttt{niord} \emph{registerW} = \emph{registerZ}, \emph{upper27\_\discretionary{}{}{}lower5}\emph{splat32}

\paragraph{Description} Bitwise not of bitwise inclusive or of the \emph{registerZ} with the \emph{upper27\_\discretionary{}{}{}lower5}. The result is stored into the \emph{registerW}.


\paragraph{Modifier(s)}
\noindent\begin{itemize}
\item splat32 (cf. splat32 in section \ref{par:modifier-splat32} on page \pageref{par:modifier-splat32})
\end{itemize}

\paragraph{Execution} {Reservation table \textsf{ALU\_TINY.X} (Section~\ref{sec:reservation-ALU-TINY.X})}
~
{\begin{lstlisting}
stage ID:
  new splat32 = _ZX_1(splat32);
stage RR:
  new argument3 = splat32 ? (_WX_32(upper27_lower5) << 32) | _ZX_32(_WX_32(upper27_lower5)) : _SX_32(_WX_32(upper27_lower5));
  new argument2 = GPR[registerZ];
  new result1 = ~(argument2 | argument3);
stage E1:
  GPR[registerW] = result1;
\end{lstlisting}}
\newpage

\paragraph{Encoding} Format \textsf{ALU\_DWRI} (Section~\ref{sec:format-ALU-DWRI}), \verb|exucode2 = 1011|\\

\bitfields{ 1/P, 2/11, 1/0, 4/1011, 6/registerW, 2/00, 10/signed10, 6/registerZ }


\paragraph{Syntax} \texttt{niord} \emph{registerW} = \emph{registerZ}, \emph{signed10}

\paragraph{Description} Bitwise not of bitwise inclusive or of the \emph{registerZ} with the \emph{signed10}. The result is stored into the \emph{registerW}.


\paragraph{Execution} {Reservation table \textsf{ALU\_TINY} (Section~\ref{sec:reservation-ALU-TINY})}
~
{\begin{lstlisting}
stage RR:
  new argument3 = _SX_10(signed10);
  new argument2 = GPR[registerZ];
  new result1 = ~(argument2 | argument3);
stage E1:
  GPR[registerW] = result1;
\end{lstlisting}}
\newpage

\paragraph{Encoding} Format \textsf{ALU\_DWRI.X} (Section~\ref{sec:format-ALU-DWRI.X}), \verb|exucode2 = 1011|\\

\bitfields{ 1/P, 2/00, 2/-- --, 27/upper27, 1/1, 2/11, 1/0, 4/1011, 6/registerW, 2/00, 10/lower10, 6/registerZ }


\paragraph{Syntax} \texttt{niord} \emph{registerW} = \emph{registerZ}, \emph{upper27\_\discretionary{}{}{}lower10}

\paragraph{Description} Bitwise not of bitwise inclusive or of the \emph{registerZ} with the \emph{upper27\_\discretionary{}{}{}lower10}. The result is stored into the \emph{registerW}.


\paragraph{Execution} {Reservation table \textsf{ALU\_TINY.X} (Section~\ref{sec:reservation-ALU-TINY.X})}
~
{\begin{lstlisting}
stage RR:
  new argument3 = _SX_37(upper27_lower10);
  new argument2 = GPR[registerZ];
  new result1 = ~(argument2 | argument3);
stage E1:
  GPR[registerW] = result1;
\end{lstlisting}}
\newpage

\paragraph{Encoding} Format \textsf{ALU\_DWRI.Y} (Section~\ref{sec:format-ALU-DWRI.Y}), \verb|exucode2 = 1011|\\

\bitfields{ 1/P, 2/00, 2/-- --, 27/extend27, 1/1, 2/00, 2/-- --, 27/upper27, 1/1, 2/11, 1/0, 4/1011, 6/registerW, 2/00, 10/lower10, 6/registerZ }


\paragraph{Syntax} \texttt{niord} \emph{registerW} = \emph{registerZ}, \emph{extend27\_\discretionary{}{}{}upper27\_\discretionary{}{}{}lower10}

\paragraph{Description} Bitwise not of bitwise inclusive or of the \emph{registerZ} with the \emph{extend27\_\discretionary{}{}{}upper27\_\discretionary{}{}{}lower10}. The result is stored into the \emph{registerW}.


\paragraph{Execution} {Reservation table \textsf{ALU\_TINY.Y} (Section~\ref{sec:reservation-ALU-TINY.Y})}
~
{\begin{lstlisting}
stage RR:
  new argument3 = _WX_64(extend27_upper27_lower10);
  new argument2 = GPR[registerZ];
  new result1 = ~(argument2 | argument3);
stage E1:
  GPR[registerW] = result1;
\end{lstlisting}}
\newpage
\subsubsection[{\small NIORQ: Bitwise Not Inclusive Or Quadruple Word}]{NIORQ\hfill Bitwise Not Inclusive Or Quadruple Word}
\label{instr:NIORQ}\index{niorq}
 
\paragraph{Encoding} Format \textsf{ALU\_QWRR0} (Section~\ref{sec:format-ALU-QWRR0}), \verb|exucode2 = 1011|\\

\bitfields{ 1/P, 2/11, 1/0, 4/1011, 5/registerM, 1/--, 2/10, 3/000, 1/1, 5/registerO, 1/--, 5/registerP, 1/-- }


\paragraph{Syntax} \texttt{niorq} \emph{registerM} = \emph{registerP}, \emph{registerO}

\paragraph{Description} Bitwise not of bitwise inclusive or of the \emph{registerP} with the \emph{registerO}. The result is stored into the \emph{registerM}.


\paragraph{Execution} {Reservation table \textsf{ALU\_LITE} (Section~\ref{sec:reservation-ALU-LITE})}
~
{\begin{lstlisting}
stage RR:
  new argument3 = PGR[registerO];
  new argument2 = PGR[registerP];
  new result1 = ~(argument2 | argument3);
stage E1:
  PGR[registerM] = result1;
\end{lstlisting}}
\newpage

\paragraph{Encoding} Format \textsf{ALU\_QWRR0.M} (Section~\ref{sec:format-ALU-QWRR0.M}), \verb|exucode2 = 1011|\\

\bitfields{ 1/P, 2/00, 2/-- --, 27/upper27, 1/1, 2/11, 1/0, 4/1011, 5/registerM, 1/--, 2/10, 3/000, 1/1, 1/splat32, 5/lower5, 5/registerP, 1/1 }


\paragraph{Syntax} \texttt{niorq} \emph{registerM} = \emph{registerP}, \emph{upper27\_\discretionary{}{}{}lower5}\emph{splat32}

\paragraph{Description} Bitwise not of bitwise inclusive or of the \emph{registerP} with the \emph{upper27\_\discretionary{}{}{}lower5}. The result is stored into the \emph{registerM}.


\paragraph{Modifier(s)}
\noindent\begin{itemize}
\item splat32 (cf. splat32 in section \ref{par:modifier-splat32} on page \pageref{par:modifier-splat32})
\end{itemize}

\paragraph{Execution} {Reservation table \textsf{ALU\_LITE.X} (Section~\ref{sec:reservation-ALU-LITE.X})}
~
{\begin{lstlisting}
stage ID:
  new splat32 = _ZX_1(splat32);
stage RR:
  new argument3 = splat32 ? (_WX_32(upper27_lower5) << 32) | _ZX_32(_WX_32(upper27_lower5)) : _SX_32(_WX_32(upper27_lower5));
  new argument2 = PGR[registerP];
  new result1 = ~(argument2 | argument3);
stage E1:
  PGR[registerM] = result1;
\end{lstlisting}}
\newpage
\subsubsection[{\small NIORW: Bitwise Not Inclusive Or Word}]{NIORW\hfill Bitwise Not Inclusive Or Word}
\label{instr:NIORW}\index{niorw}
 
\paragraph{Encoding} Format \textsf{ALU\_WWRR0} (Section~\ref{sec:format-ALU-WWRR0}), \verb|exucode2 = 1011|\\

\bitfields{ 1/P, 2/11, 1/0, 4/1011, 6/registerW, 2/11, 3/000, 1/signextw, 6/registerY, 6/registerZ }


\paragraph{Syntax} \texttt{niorw}\emph{signextw} \emph{registerW} = \emph{registerZ}, \emph{registerY}

\paragraph{Description} Bitwise not of bitwise inclusive or of the \emph{registerZ} with the \emph{registerY}. The result with the 32 upper bits cleared is stored into the \emph{registerW}.


\paragraph{Modifier(s)}
\noindent\begin{itemize}
\item signextw (cf. signextw in section \ref{par:modifier-signextw} on page \pageref{par:modifier-signextw})
\end{itemize}

\paragraph{Execution} {Reservation table \textsf{ALU\_TINY} (Section~\ref{sec:reservation-ALU-TINY})}
~
{\begin{lstlisting}
stage ID:
  new signextw = _ZX_1(signextw);
stage RR:
  new argument3 = GPR[registerY];
  new argument2 = GPR[registerZ];
  new result1 = ~(argument2 | argument3);
stage E1:
  GPR[registerW] = signextw ? _SX_32(result1) : _ZX_32(result1);
\end{lstlisting}}
\newpage

\paragraph{Encoding} Format \textsf{ALU\_WWRR0.W} (Section~\ref{sec:format-ALU-WWRR0.W}), \verb|exucode2 = 1011|\\

\bitfields{ 1/P, 2/00, 2/-- --, 27/upper27, 1/1, 2/11, 1/0, 4/1011, 6/registerW, 2/11, 3/000, 1/signextw, 1/0, 5/lower5, 6/registerZ }


\paragraph{Syntax} \texttt{niorw}\emph{signextw} \emph{registerW} = \emph{registerZ}, \emph{upper27\_\discretionary{}{}{}lower5}

\paragraph{Description} Bitwise not of bitwise inclusive or of the \emph{registerZ} with the \emph{upper27\_\discretionary{}{}{}lower5}. The result with the 32 upper bits cleared is stored into the \emph{registerW}.


\paragraph{Modifier(s)}
\noindent\begin{itemize}
\item signextw (cf. signextw in section \ref{par:modifier-signextw} on page \pageref{par:modifier-signextw})
\end{itemize}

\paragraph{Execution} {Reservation table \textsf{ALU\_TINY.X} (Section~\ref{sec:reservation-ALU-TINY.X})}
~
{\begin{lstlisting}
stage ID:
  new signextw = _ZX_1(signextw);
stage RR:
  new argument3 = _WX_32(upper27_lower5);
  new argument2 = GPR[registerZ];
  new result1 = ~(argument2 | argument3);
stage E1:
  GPR[registerW] = signextw ? _SX_32(result1) : _ZX_32(result1);
\end{lstlisting}}
\newpage
\subsubsection[{\small NOP: No Operation}]{NOP\hfill No Operation}
\label{instr:NOP}\index{nop}
 
\paragraph{Encoding} Format \textsf{ALU\_NOP} (Section~\ref{sec:format-ALU-NOP}), \verb|exucode2 = 1111|\\

\bitfields{ 1/P, 2/11, 1/1, 4/1111, 5/-- -- -- -- --, 1/1, 2/11, 4/1111, 5/11111, 1/1, 5/-- -- -- -- --, 1/1 }


\paragraph{Syntax} \texttt{nop}

\paragraph{Description} No effects except extending the current bundle by one syllable.


\paragraph{Execution} {Reservation table \textsf{ALU\_TINY} (Section~\ref{sec:reservation-ALU-TINY})}
~
{\begin{lstlisting}
stage RR:
  ;
\end{lstlisting}}
\newpage
\subsubsection[{\small NOTD: Negate Double Word}]{NOTD\hfill Negate Double Word}
\label{instr:NOTD}\index{notd}
 
\paragraph{Encoding} Format \textsf{ALU\_DBWR} (Section~\ref{sec:format-ALU-DBWR}), \verb|alucode8 = 1000|\\

\bitfields{ 1/P, 2/11, 1/0, 4/1111, 6/registerW, 2/10, 3/101, 1/0, 4/1000, 1/0, 1/0, 6/registerZ }


\paragraph{Syntax} \texttt{notd} \emph{registerW} = \emph{registerZ}

\paragraph{Description} The bitwise not value of the \emph{registerZ} is stored into the \emph{registerW}.


\paragraph{Execution} {Reservation table \textsf{ALU\_TINY} (Section~\ref{sec:reservation-ALU-TINY})}
~
{\begin{lstlisting}
stage RR:
  new argument2 = GPR[registerZ];
  new result1 = ~argument2;
stage E1:
  GPR[registerW] = result1;
\end{lstlisting}}
\newpage
\subsubsection[{\small NOTQ: Negate Quadruple Word}]{NOTQ\hfill Negate Quadruple Word}
\label{instr:NOTQ}\index{notq}
 
\paragraph{Encoding} Format \textsf{ALU\_QBWR} (Section~\ref{sec:format-ALU-QBWR}), \verb|alucode8 = 1000|\\

\bitfields{ 1/P, 2/11, 1/0, 4/1111, 5/registerM, 1/--, 2/10, 3/101, 1/1, 4/1000, 1/0, 1/0, 5/registerP, 1/-- }


\paragraph{Syntax} \texttt{notq} \emph{registerM} = \emph{registerP}

\paragraph{Description} The bitwise not value of the \emph{registerP} is stored into the \emph{registerM}.


\paragraph{Execution} {Reservation table \textsf{ALU\_LITE} (Section~\ref{sec:reservation-ALU-LITE})}
~
{\begin{lstlisting}
stage RR:
  new argument2 = PGR[registerP];
  new result1 = ~argument2;
stage E1:
  PGR[registerM] = result1;
\end{lstlisting}}
\newpage
\subsubsection[{\small NOTW: Negate Word}]{NOTW\hfill Negate Word}
\label{instr:NOTW}\index{notw}
 
\paragraph{Encoding} Format \textsf{ALU\_BWRW} (Section~\ref{sec:format-ALU-BWRW}), \verb|alucode8 = 1000|\\

\bitfields{ 1/P, 2/11, 1/0, 4/1111, 6/registerW, 2/11, 3/101, 1/signextw, 4/1000, 1/0, 1/0, 6/registerZ }


\paragraph{Syntax} \texttt{notw} \emph{registerW} = \emph{registerZ}

\paragraph{Description} The bitwise not value of the \emph{registerZ} is stored into the \emph{registerW}.


\paragraph{Execution} {Reservation table \textsf{ALU\_TINY} (Section~\ref{sec:reservation-ALU-TINY})}
~
{\begin{lstlisting}
stage RR:
  new argument2 = GPR[registerZ];
  new result1 = ~argument2;
stage E1:
  GPR[registerW] = _ZX_32(result1);
\end{lstlisting}}
\newpage
\subsubsection[{\small PCREL: PC-Relative}]{PCREL\hfill PC-Relative}
\label{instr:PCREL}\index{pcrel}
 
\paragraph{Encoding} Format \textsf{ALU\_DWI} (Section~\ref{sec:format-ALU-DWI}), \verb|exucode2 = 0001|\\

\bitfields{ 1/P, 2/11, 1/0, 4/0001, 6/registerW, 2/00, 16/signed16 }


\paragraph{Syntax} \texttt{pcrel} \emph{registerW} = \emph{signed16}

\paragraph{Description} The \emph{signed16} is is added to the bundle \textbf{PC} and the result is stored into the \emph{registerW}.


\paragraph{Execution} {Reservation table \textsf{ALU\_TINY} (Section~\ref{sec:reservation-ALU-TINY})}
~
{\begin{lstlisting}
stage RR:
  new argument2 = _SX_16(signed16);
  new result1 = PC + _SX_64(argument2);
stage E1:
  GPR[registerW] = result1;
\end{lstlisting}}
\newpage

\paragraph{Encoding} Format \textsf{ALU\_DWI.X} (Section~\ref{sec:format-ALU-DWI.X}), \verb|exucode2 = 0001|\\

\bitfields{ 1/P, 2/00, 2/-- --, 27/upper27, 1/1, 2/11, 1/0, 4/0001, 6/registerW, 2/00, 10/lower10, 6/extend6 }


\paragraph{Syntax} \texttt{pcrel} \emph{registerW} = \emph{extend6\_\discretionary{}{}{}upper27\_\discretionary{}{}{}lower10}

\paragraph{Description} The \emph{extend6\_\discretionary{}{}{}upper27\_\discretionary{}{}{}lower10} is is added to the bundle \textbf{PC} and the result is stored into the \emph{registerW}.


\paragraph{Execution} {Reservation table \textsf{ALU\_TINY.X} (Section~\ref{sec:reservation-ALU-TINY.X})}
~
{\begin{lstlisting}
stage RR:
  new argument2 = _SX_43(extend6_upper27_lower10);
  new result1 = PC + _SX_64(argument2);
stage E1:
  GPR[registerW] = result1;
\end{lstlisting}}
\newpage

\paragraph{Encoding} Format \textsf{ALU\_DWI.Y} (Section~\ref{sec:format-ALU-DWI.Y}), \verb|exucode2 = 0001|\\

\bitfields{ 1/P, 2/00, 2/-- --, 27/extend27, 1/1, 2/00, 2/-- --, 27/upper27, 1/1, 2/11, 1/0, 4/0001, 6/registerW, 2/00, 10/lower10, 6/-- -- -- -- -- -- }


\paragraph{Syntax} \texttt{pcrel} \emph{registerW} = \emph{extend27\_\discretionary{}{}{}upper27\_\discretionary{}{}{}lower10}

\paragraph{Description} The \emph{extend27\_\discretionary{}{}{}upper27\_\discretionary{}{}{}lower10} is is added to the bundle \textbf{PC} and the result is stored into the \emph{registerW}.


\paragraph{Execution} {Reservation table \textsf{ALU\_TINY.Y} (Section~\ref{sec:reservation-ALU-TINY.Y})}
~
{\begin{lstlisting}
stage RR:
  new argument2 = _WX_64(extend27_upper27_lower10);
  new result1 = PC + _SX_64(argument2);
stage E1:
  GPR[registerW] = result1;
\end{lstlisting}}
\newpage
\subsubsection[{\small ROLD: Rotate Left Double Word}]{ROLD\hfill Rotate Left Double Word}
\label{instr:ROLD}\index{rold}
 
\paragraph{Encoding} Format \textsf{ALU\_DSWRR} (Section~\ref{sec:format-ALU-DSWRR}), \verb|exucode1 = 110|\\

\bitfields{ 1/P, 2/11, 1/0, 1/0, 3/110, 6/registerW, 2/10, 3/100, 1/0, 6/registerY, 6/registerZ }


\paragraph{Syntax} \texttt{rold} \emph{registerW} = \emph{registerZ}, \emph{registerY}

\paragraph{Description} The \emph{registerZ} is rotated left by the \emph{registerY} modulo 64. The result is stored into the \emph{registerW}.


\paragraph{Execution} {Reservation table \textsf{ALU\_TINY} (Section~\ref{sec:reservation-ALU-TINY})}
~
{\begin{lstlisting}
stage RR:
  new argument3 = GPR[registerY];
  new argument2 = GPR[registerZ];
  new result1 = _ROL_64(argument2.64[0], _ZX_6(argument3));
stage E1:
  GPR[registerW] = result1;
\end{lstlisting}}
\newpage

\paragraph{Encoding} Format \textsf{ALU\_DSWRI} (Section~\ref{sec:format-ALU-DSWRI}), \verb|exucode1 = 110|\\

\bitfields{ 1/P, 2/11, 1/0, 1/1, 3/110, 6/registerW, 2/10, 3/100, 1/0, 6/unsigned6, 6/registerZ }


\paragraph{Syntax} \texttt{rold} \emph{registerW} = \emph{registerZ}, \emph{unsigned6}

\paragraph{Description} The \emph{registerZ} is rotated left by the \emph{unsigned6} modulo 64. The result is stored into the \emph{registerW}.


\paragraph{Execution} {Reservation table \textsf{ALU\_TINY} (Section~\ref{sec:reservation-ALU-TINY})}
~
{\begin{lstlisting}
stage RR:
  new argument3 = _ZX_6(unsigned6);
  new argument2 = GPR[registerZ];
  new result1 = _ROL_64(argument2.64[0], _ZX_6(argument3));
stage E1:
  GPR[registerW] = result1;
\end{lstlisting}}
\newpage
\subsubsection[{\small ROLW: Rotate Left Word}]{ROLW\hfill Rotate Left Word}
\label{instr:ROLW}\index{rolw}
 
\paragraph{Encoding} Format \textsf{ALU\_WSWRR} (Section~\ref{sec:format-ALU-WSWRR}), \verb|exucode1 = 110|\\

\bitfields{ 1/P, 2/11, 1/0, 1/0, 3/110, 6/registerW, 2/11, 3/100, 1/signextw, 6/registerY, 6/registerZ }


\paragraph{Syntax} \texttt{rolw}\emph{signextw} \emph{registerW} = \emph{registerZ}, \emph{registerY}

\paragraph{Description} The \emph{registerZ} is rotated left by the \emph{registerY} modulo 32. The result is stored into the \emph{registerW}.


\paragraph{Modifier(s)}
\noindent\begin{itemize}
\item signextw (cf. signextw in section \ref{par:modifier-signextw} on page \pageref{par:modifier-signextw})
\end{itemize}

\paragraph{Execution} {Reservation table \textsf{ALU\_TINY} (Section~\ref{sec:reservation-ALU-TINY})}
~
{\begin{lstlisting}
stage ID:
  new signextw = _ZX_1(signextw);
stage RR:
  new argument3 = GPR[registerY];
  new argument2 = GPR[registerZ];
  new result1 = _ROL_32(argument2.32[0], _ZX_5(argument3));
stage E1:
  GPR[registerW] = signextw ? _SX_32(result1) : _ZX_32(result1);
\end{lstlisting}}
\newpage

\paragraph{Encoding} Format \textsf{ALU\_WSWRI} (Section~\ref{sec:format-ALU-WSWRI}), \verb|exucode1 = 110|\\

\bitfields{ 1/P, 2/11, 1/0, 1/1, 3/110, 6/registerW, 2/11, 3/100, 1/signextw, 6/unsigned6, 6/registerZ }


\paragraph{Syntax} \texttt{rolw}\emph{signextw} \emph{registerW} = \emph{registerZ}, \emph{unsigned6}

\paragraph{Description} The \emph{registerZ} is rotated left by the \emph{unsigned6} modulo 32. The result is stored into the \emph{registerW}.


\paragraph{Modifier(s)}
\noindent\begin{itemize}
\item signextw (cf. signextw in section \ref{par:modifier-signextw} on page \pageref{par:modifier-signextw})
\end{itemize}

\paragraph{Execution} {Reservation table \textsf{ALU\_TINY} (Section~\ref{sec:reservation-ALU-TINY})}
~
{\begin{lstlisting}
stage ID:
  new signextw = _ZX_1(signextw);
stage RR:
  new argument3 = _ZX_6(unsigned6);
  new argument2 = GPR[registerZ];
  new result1 = _ROL_32(argument2.32[0], _ZX_5(argument3));
stage E1:
  GPR[registerW] = signextw ? _SX_32(result1) : _ZX_32(result1);
\end{lstlisting}}
\newpage
\subsubsection[{\small ROLWP: Rotate Left Word Pair}]{ROLWP\hfill Rotate Left Word Pair}
\label{instr:ROLWP}\index{rolwp}
 
\paragraph{Encoding} Format \textsf{ALU\_WPSWRIE} (Section~\ref{sec:format-ALU-WPSWRIE}), \verb|exucode1 = 110|\\

\bitfields{ 1/P, 2/11, 1/1, 1/1, 3/110, 5/registerWe, 1/0, 2/10, 3/100, 1/0, 6/unsigned6, 5/registerZe, 1/1 }


\paragraph{Syntax} \texttt{rolwp} \emph{registerWe} = \emph{registerZe}, \emph{unsigned6}

\paragraph{Description} The \emph{registerZe} is interpreted as two words packed into 64 bits. The \emph{registerZe} words are rotated left by the \emph{unsigned6} modulo 32. The two resulting words are packed and stored into the \emph{registerWe}.


\paragraph{Execution} {Reservation table \textsf{ALU\_LITE} (Section~\ref{sec:reservation-ALU-LITE})}
~
{\begin{lstlisting}
stage RR:
  new argument3 = _ZX_6(unsigned6);
  new argument2 = GPR[registerZe<<1];
  new shift = _ZX_5(argument3);
  for (I = 0; I < 2; I += 1) {
    result1.32[I] = _ROL_32(argument2.32[I], shift)
  };
stage E1:
  GPR[registerWe<<1] = result1;
\end{lstlisting}}
\newpage

\paragraph{Encoding} Format \textsf{ALU\_WPSWRIO} (Section~\ref{sec:format-ALU-WPSWRIO}), \verb|exucode1 = 110|\\

\bitfields{ 1/P, 2/11, 1/1, 1/1, 3/110, 5/registerWo, 1/1, 2/10, 3/100, 1/0, 6/unsigned6, 5/registerZo, 1/1 }


\paragraph{Syntax} \texttt{rolwp} \emph{registerWo} = \emph{registerZo}, \emph{unsigned6}

\paragraph{Description} The \emph{registerZo} is interpreted as two words packed into 64 bits. The \emph{registerZo} words are rotated left by the \emph{unsigned6} modulo 32. The two resulting words are packed and stored into the \emph{registerWo}.


\paragraph{Execution} {Reservation table \textsf{ALU\_LITE} (Section~\ref{sec:reservation-ALU-LITE})}
~
{\begin{lstlisting}
stage RR:
  new argument3 = _ZX_6(unsigned6);
  new argument2 = GPR[(registerZo<<1)+1];
  new shift = _ZX_5(argument3);
  for (I = 0; I < 2; I += 1) {
    result1.32[I] = _ROL_32(argument2.32[I], shift)
  };
stage E1:
  GPR[(registerWo<<1)+1] = result1;
\end{lstlisting}}
\newpage

\paragraph{Encoding} Format \textsf{ALU\_WPSWRRE} (Section~\ref{sec:format-ALU-WPSWRRE}), \verb|exucode1 = 110|\\

\bitfields{ 1/P, 2/11, 1/1, 1/0, 3/110, 5/registerWe, 1/0, 2/10, 3/100, 1/0, 6/registerY, 5/registerZe, 1/1 }


\paragraph{Syntax} \texttt{rolwp} \emph{registerWe} = \emph{registerZe}, \emph{registerY}

\paragraph{Description} The \emph{registerZe} is interpreted as two words packed into 64 bits. The \emph{registerZe} words are rotated left by the \emph{registerY} modulo 32. The two resulting words are packed and stored into the \emph{registerWe}.


\paragraph{Execution} {Reservation table \textsf{ALU\_LITE} (Section~\ref{sec:reservation-ALU-LITE})}
~
{\begin{lstlisting}
stage RR:
  new argument3 = GPR[registerY];
  new argument2 = GPR[registerZe<<1];
  new shift = _ZX_5(argument3);
  for (I = 0; I < 2; I += 1) {
    result1.32[I] = _ROL_32(argument2.32[I], shift)
  };
stage E1:
  GPR[registerWe<<1] = result1;
\end{lstlisting}}
\newpage

\paragraph{Encoding} Format \textsf{ALU\_WPSWRRO} (Section~\ref{sec:format-ALU-WPSWRRO}), \verb|exucode1 = 110|\\

\bitfields{ 1/P, 2/11, 1/1, 1/0, 3/110, 5/registerWo, 1/1, 2/10, 3/100, 1/0, 6/registerY, 5/registerZo, 1/1 }


\paragraph{Syntax} \texttt{rolwp} \emph{registerWo} = \emph{registerZo}, \emph{registerY}

\paragraph{Description} The \emph{registerZo} is interpreted as two words packed into 64 bits. The \emph{registerZo} words are rotated left by the \emph{registerY} modulo 32. The two resulting words are packed and stored into the \emph{registerWo}.


\paragraph{Execution} {Reservation table \textsf{ALU\_LITE} (Section~\ref{sec:reservation-ALU-LITE})}
~
{\begin{lstlisting}
stage RR:
  new argument3 = GPR[registerY];
  new argument2 = GPR[(registerZo<<1)+1];
  new shift = _ZX_5(argument3);
  for (I = 0; I < 2; I += 1) {
    result1.32[I] = _ROL_32(argument2.32[I], shift)
  };
stage E1:
  GPR[(registerWo<<1)+1] = result1;
\end{lstlisting}}
\newpage
\subsubsection[{\small RORD: Rotate Right Double Word}]{RORD\hfill Rotate Right Double Word}
\label{instr:RORD}\index{rord}
 
\paragraph{Encoding} Format \textsf{ALU\_DSWRR} (Section~\ref{sec:format-ALU-DSWRR}), \verb|exucode1 = 111|\\

\bitfields{ 1/P, 2/11, 1/0, 1/0, 3/111, 6/registerW, 2/10, 3/100, 1/0, 6/registerY, 6/registerZ }


\paragraph{Syntax} \texttt{rord} \emph{registerW} = \emph{registerZ}, \emph{registerY}

\paragraph{Description} The \emph{registerZ} is rotated right by the \emph{registerY} modulo 64. The result is stored into the \emph{registerW}.


\paragraph{Execution} {Reservation table \textsf{ALU\_TINY} (Section~\ref{sec:reservation-ALU-TINY})}
~
{\begin{lstlisting}
stage RR:
  new argument3 = GPR[registerY];
  new argument2 = GPR[registerZ];
  new result1 = _ROR_64(argument2.64[0], _ZX_6(argument3));
stage E1:
  GPR[registerW] = result1;
\end{lstlisting}}
\newpage

\paragraph{Encoding} Format \textsf{ALU\_DSWRI} (Section~\ref{sec:format-ALU-DSWRI}), \verb|exucode1 = 111|\\

\bitfields{ 1/P, 2/11, 1/0, 1/1, 3/111, 6/registerW, 2/10, 3/100, 1/0, 6/unsigned6, 6/registerZ }


\paragraph{Syntax} \texttt{rord} \emph{registerW} = \emph{registerZ}, \emph{unsigned6}

\paragraph{Description} The \emph{registerZ} is rotated right by the \emph{unsigned6} modulo 64. The result is stored into the \emph{registerW}.


\paragraph{Execution} {Reservation table \textsf{ALU\_TINY} (Section~\ref{sec:reservation-ALU-TINY})}
~
{\begin{lstlisting}
stage RR:
  new argument3 = _ZX_6(unsigned6);
  new argument2 = GPR[registerZ];
  new result1 = _ROR_64(argument2.64[0], _ZX_6(argument3));
stage E1:
  GPR[registerW] = result1;
\end{lstlisting}}
\newpage
\subsubsection[{\small RORW: Rotate Right Word}]{RORW\hfill Rotate Right Word}
\label{instr:RORW}\index{rorw}
 
\paragraph{Encoding} Format \textsf{ALU\_WSWRR} (Section~\ref{sec:format-ALU-WSWRR}), \verb|exucode1 = 111|\\

\bitfields{ 1/P, 2/11, 1/0, 1/0, 3/111, 6/registerW, 2/11, 3/100, 1/signextw, 6/registerY, 6/registerZ }


\paragraph{Syntax} \texttt{rorw}\emph{signextw} \emph{registerW} = \emph{registerZ}, \emph{registerY}

\paragraph{Description} The \emph{registerZ} is rotated right by the \emph{registerY} modulo 32. The result is stored into the \emph{registerW}.


\paragraph{Modifier(s)}
\noindent\begin{itemize}
\item signextw (cf. signextw in section \ref{par:modifier-signextw} on page \pageref{par:modifier-signextw})
\end{itemize}

\paragraph{Execution} {Reservation table \textsf{ALU\_TINY} (Section~\ref{sec:reservation-ALU-TINY})}
~
{\begin{lstlisting}
stage ID:
  new signextw = _ZX_1(signextw);
stage RR:
  new argument3 = GPR[registerY];
  new argument2 = GPR[registerZ];
  new result1 = _ROR_32(argument2.32[0], _ZX_5(argument3));
stage E1:
  GPR[registerW] = signextw ? _SX_32(result1) : _ZX_32(result1);
\end{lstlisting}}
\newpage

\paragraph{Encoding} Format \textsf{ALU\_WSWRI} (Section~\ref{sec:format-ALU-WSWRI}), \verb|exucode1 = 111|\\

\bitfields{ 1/P, 2/11, 1/0, 1/1, 3/111, 6/registerW, 2/11, 3/100, 1/signextw, 6/unsigned6, 6/registerZ }


\paragraph{Syntax} \texttt{rorw}\emph{signextw} \emph{registerW} = \emph{registerZ}, \emph{unsigned6}

\paragraph{Description} The \emph{registerZ} is rotated right by the \emph{unsigned6} modulo 32. The result is stored into the \emph{registerW}.


\paragraph{Modifier(s)}
\noindent\begin{itemize}
\item signextw (cf. signextw in section \ref{par:modifier-signextw} on page \pageref{par:modifier-signextw})
\end{itemize}

\paragraph{Execution} {Reservation table \textsf{ALU\_TINY} (Section~\ref{sec:reservation-ALU-TINY})}
~
{\begin{lstlisting}
stage ID:
  new signextw = _ZX_1(signextw);
stage RR:
  new argument3 = _ZX_6(unsigned6);
  new argument2 = GPR[registerZ];
  new result1 = _ROR_32(argument2.32[0], _ZX_5(argument3));
stage E1:
  GPR[registerW] = signextw ? _SX_32(result1) : _ZX_32(result1);
\end{lstlisting}}
\newpage
\subsubsection[{\small RORWP: Rotate Right Word Pair}]{RORWP\hfill Rotate Right Word Pair}
\label{instr:RORWP}\index{rorwp}
 
\paragraph{Encoding} Format \textsf{ALU\_WPSWRIE} (Section~\ref{sec:format-ALU-WPSWRIE}), \verb|exucode1 = 111|\\

\bitfields{ 1/P, 2/11, 1/1, 1/1, 3/111, 5/registerWe, 1/0, 2/10, 3/100, 1/0, 6/unsigned6, 5/registerZe, 1/1 }


\paragraph{Syntax} \texttt{rorwp} \emph{registerWe} = \emph{registerZe}, \emph{unsigned6}

\paragraph{Description} The \emph{registerZe} is interpreted as two words packed into 64 bits. The \emph{registerZe} words are rotated right by the \emph{unsigned6} modulo 32. The two resulting words are packed and stored into the \emph{registerWe}.


\paragraph{Execution} {Reservation table \textsf{ALU\_LITE} (Section~\ref{sec:reservation-ALU-LITE})}
~
{\begin{lstlisting}
stage RR:
  new argument3 = _ZX_6(unsigned6);
  new argument2 = GPR[registerZe<<1];
  new shift = _ZX_5(argument3);
  for (I = 0; I < 2; I += 1) {
    result1.32[I] = _ROR_32(argument2.32[I], shift)
  };
stage E1:
  GPR[registerWe<<1] = result1;
\end{lstlisting}}
\newpage

\paragraph{Encoding} Format \textsf{ALU\_WPSWRIO} (Section~\ref{sec:format-ALU-WPSWRIO}), \verb|exucode1 = 111|\\

\bitfields{ 1/P, 2/11, 1/1, 1/1, 3/111, 5/registerWo, 1/1, 2/10, 3/100, 1/0, 6/unsigned6, 5/registerZo, 1/1 }


\paragraph{Syntax} \texttt{rorwp} \emph{registerWo} = \emph{registerZo}, \emph{unsigned6}

\paragraph{Description} The \emph{registerZo} is interpreted as two words packed into 64 bits. The \emph{registerZo} words are rotated right by the \emph{unsigned6} modulo 32. The two resulting words are packed and stored into the \emph{registerWo}.


\paragraph{Execution} {Reservation table \textsf{ALU\_LITE} (Section~\ref{sec:reservation-ALU-LITE})}
~
{\begin{lstlisting}
stage RR:
  new argument3 = _ZX_6(unsigned6);
  new argument2 = GPR[(registerZo<<1)+1];
  new shift = _ZX_5(argument3);
  for (I = 0; I < 2; I += 1) {
    result1.32[I] = _ROR_32(argument2.32[I], shift)
  };
stage E1:
  GPR[(registerWo<<1)+1] = result1;
\end{lstlisting}}
\newpage

\paragraph{Encoding} Format \textsf{ALU\_WPSWRRE} (Section~\ref{sec:format-ALU-WPSWRRE}), \verb|exucode1 = 111|\\

\bitfields{ 1/P, 2/11, 1/1, 1/0, 3/111, 5/registerWe, 1/0, 2/10, 3/100, 1/0, 6/registerY, 5/registerZe, 1/1 }


\paragraph{Syntax} \texttt{rorwp} \emph{registerWe} = \emph{registerZe}, \emph{registerY}

\paragraph{Description} The \emph{registerZe} is interpreted as two words packed into 64 bits. The \emph{registerZe} words are rotated right by the \emph{registerY} modulo 32. The two resulting words are packed and stored into the \emph{registerWe}.


\paragraph{Execution} {Reservation table \textsf{ALU\_LITE} (Section~\ref{sec:reservation-ALU-LITE})}
~
{\begin{lstlisting}
stage RR:
  new argument3 = GPR[registerY];
  new argument2 = GPR[registerZe<<1];
  new shift = _ZX_5(argument3);
  for (I = 0; I < 2; I += 1) {
    result1.32[I] = _ROR_32(argument2.32[I], shift)
  };
stage E1:
  GPR[registerWe<<1] = result1;
\end{lstlisting}}
\newpage

\paragraph{Encoding} Format \textsf{ALU\_WPSWRRO} (Section~\ref{sec:format-ALU-WPSWRRO}), \verb|exucode1 = 111|\\

\bitfields{ 1/P, 2/11, 1/1, 1/0, 3/111, 5/registerWo, 1/1, 2/10, 3/100, 1/0, 6/registerY, 5/registerZo, 1/1 }


\paragraph{Syntax} \texttt{rorwp} \emph{registerWo} = \emph{registerZo}, \emph{registerY}

\paragraph{Description} The \emph{registerZo} is interpreted as two words packed into 64 bits. The \emph{registerZo} words are rotated right by the \emph{registerY} modulo 32. The two resulting words are packed and stored into the \emph{registerWo}.


\paragraph{Execution} {Reservation table \textsf{ALU\_LITE} (Section~\ref{sec:reservation-ALU-LITE})}
~
{\begin{lstlisting}
stage RR:
  new argument3 = GPR[registerY];
  new argument2 = GPR[(registerZo<<1)+1];
  new shift = _ZX_5(argument3);
  for (I = 0; I < 2; I += 1) {
    result1.32[I] = _ROR_32(argument2.32[I], shift)
  };
stage E1:
  GPR[(registerWo<<1)+1] = result1;
\end{lstlisting}}
\newpage
\subsubsection[{\small SBFBO: Subtract From Byte Octuple}]{SBFBO\hfill Subtract From Byte Octuple}
\label{instr:SBFBO}\index{sbfbo}
 
\paragraph{Encoding} Format \textsf{ALU\_BOWRR0E} (Section~\ref{sec:format-ALU-BOWRR0E}), \verb|exucode2 = 0011|\\

\bitfields{ 1/P, 2/11, 1/1, 4/0011, 5/registerWe, 1/0, 2/10, 3/000, 1/1, 5/registerYe, 1/--, 5/registerZe, 1/1 }


\paragraph{Syntax} \texttt{sbfbo} \emph{registerWe} = \emph{registerZe}, \emph{registerYe}

\paragraph{Description} The \emph{registerYe} and the \emph{registerZe} are interpreted as eight bytes packed into 64 bits. The \emph{registerZe} bytes are subtracted from the \emph{registerYe} bytes. The eight resulting bytes are packed and stored into the \emph{registerWe}.


\paragraph{Execution} {Reservation table \textsf{ALU\_LITE} (Section~\ref{sec:reservation-ALU-LITE})}
~
{\begin{lstlisting}
stage RR:
  new argument3 = GPR[registerYe<<1];
  new argument2 = GPR[registerZe<<1];
  for (I = 0; I < 8; I += 1) {
    result1.8[I] = argument3.8[I] - argument2.8[I]
  };
stage E1:
  GPR[registerWe<<1] = result1;
\end{lstlisting}}
\newpage

\paragraph{Encoding} Format \textsf{ALU\_BOWRR0O} (Section~\ref{sec:format-ALU-BOWRR0O}), \verb|exucode2 = 0011|\\

\bitfields{ 1/P, 2/11, 1/1, 4/0011, 5/registerWo, 1/1, 2/10, 3/000, 1/1, 5/registerYo, 1/--, 5/registerZo, 1/1 }


\paragraph{Syntax} \texttt{sbfbo} \emph{registerWo} = \emph{registerZo}, \emph{registerYo}

\paragraph{Description} The \emph{registerYo} and the \emph{registerZo} are interpreted as eight bytes packed into 64 bits. The \emph{registerZo} bytes are subtracted from the \emph{registerYo} bytes. The eight resulting bytes are packed and stored into the \emph{registerWo}.


\paragraph{Execution} {Reservation table \textsf{ALU\_LITE} (Section~\ref{sec:reservation-ALU-LITE})}
~
{\begin{lstlisting}
stage RR:
  new argument3 = GPR[(registerYo<<1)+1];
  new argument2 = GPR[(registerZo<<1)+1];
  for (I = 0; I < 8; I += 1) {
    result1.8[I] = argument3.8[I] - argument2.8[I]
  };
stage E1:
  GPR[(registerWo<<1)+1] = result1;
\end{lstlisting}}
\newpage

\paragraph{Encoding} Format \textsf{ALU\_BOWRR0E.M} (Section~\ref{sec:format-ALU-BOWRR0E.M}), \verb|exucode2 = 0011|\\

\bitfields{ 1/P, 2/00, 2/-- --, 27/upper27, 1/1, 2/11, 1/1, 4/0011, 5/registerWe, 1/0, 2/10, 3/000, 1/1, 1/splat32, 5/lower5, 5/registerZe, 1/1 }


\paragraph{Syntax} \texttt{sbfbo} \emph{registerWe} = \emph{registerZe}, \emph{upper27\_\discretionary{}{}{}lower5}\emph{splat32}

\paragraph{Description} The \emph{upper27\_\discretionary{}{}{}lower5} and the \emph{registerZe} are interpreted as eight bytes packed into 64 bits. The \emph{registerZe} bytes are subtracted from the \emph{upper27\_\discretionary{}{}{}lower5} bytes. The eight resulting bytes are packed and stored into the \emph{registerWe}.


\paragraph{Modifier(s)}
\noindent\begin{itemize}
\item splat32 (cf. splat32 in section \ref{par:modifier-splat32} on page \pageref{par:modifier-splat32})
\end{itemize}

\paragraph{Execution} {Reservation table \textsf{ALU\_LITE.X} (Section~\ref{sec:reservation-ALU-LITE.X})}
~
{\begin{lstlisting}
stage ID:
  new splat32 = _ZX_1(splat32);
stage RR:
  new argument3 = splat32 ? (_WX_32(upper27_lower5) << 32) | _ZX_32(_WX_32(upper27_lower5)) : _SX_32(_WX_32(upper27_lower5));
  new argument2 = GPR[registerZe<<1];
  for (I = 0; I < 8; I += 1) {
    result1.8[I] = argument3.8[I] - argument2.8[I]
  };
stage E1:
  GPR[registerWe<<1] = result1;
\end{lstlisting}}
\newpage

\paragraph{Encoding} Format \textsf{ALU\_BOWRR0O.M} (Section~\ref{sec:format-ALU-BOWRR0O.M}), \verb|exucode2 = 0011|\\

\bitfields{ 1/P, 2/00, 2/-- --, 27/upper27, 1/1, 2/11, 1/1, 4/0011, 5/registerWo, 1/1, 2/10, 3/000, 1/1, 1/splat32, 5/lower5, 5/registerZo, 1/1 }


\paragraph{Syntax} \texttt{sbfbo} \emph{registerWo} = \emph{registerZo}, \emph{upper27\_\discretionary{}{}{}lower5}\emph{splat32}

\paragraph{Description} The \emph{upper27\_\discretionary{}{}{}lower5} and the \emph{registerZo} are interpreted as eight bytes packed into 64 bits. The \emph{registerZo} bytes are subtracted from the \emph{upper27\_\discretionary{}{}{}lower5} bytes. The eight resulting bytes are packed and stored into the \emph{registerWo}.


\paragraph{Modifier(s)}
\noindent\begin{itemize}
\item splat32 (cf. splat32 in section \ref{par:modifier-splat32} on page \pageref{par:modifier-splat32})
\end{itemize}

\paragraph{Execution} {Reservation table \textsf{ALU\_LITE.X} (Section~\ref{sec:reservation-ALU-LITE.X})}
~
{\begin{lstlisting}
stage ID:
  new splat32 = _ZX_1(splat32);
stage RR:
  new argument3 = splat32 ? (_WX_32(upper27_lower5) << 32) | _ZX_32(_WX_32(upper27_lower5)) : _SX_32(_WX_32(upper27_lower5));
  new argument2 = GPR[(registerZo<<1)+1];
  for (I = 0; I < 8; I += 1) {
    result1.8[I] = argument3.8[I] - argument2.8[I]
  };
stage E1:
  GPR[(registerWo<<1)+1] = result1;
\end{lstlisting}}
\newpage
\subsubsection[{\small SBFD: Subtract From Double Word}]{SBFD\hfill Subtract From Double Word}
\label{instr:SBFD}\index{sbfd}
 
\paragraph{Encoding} Format \textsf{ALU\_DWRR0} (Section~\ref{sec:format-ALU-DWRR0}), \verb|exucode2 = 0011|\\

\bitfields{ 1/P, 2/11, 1/0, 4/0011, 6/registerW, 2/10, 3/000, 1/0, 6/registerY, 6/registerZ }


\paragraph{Syntax} \texttt{sbfd} \emph{registerW} = \emph{registerZ}, \emph{registerY}

\paragraph{Description} The \emph{registerZ} is subtracted from the \emph{registerY}. The result is stored into the \emph{registerW}.


\paragraph{Execution} {Reservation table \textsf{ALU\_TINY} (Section~\ref{sec:reservation-ALU-TINY})}
~
{\begin{lstlisting}
stage RR:
  new argument3 = GPR[registerY];
  new argument2 = GPR[registerZ];
  new result1 = argument3 - argument2;
stage E1:
  GPR[registerW] = result1;
\end{lstlisting}}
\newpage

\paragraph{Encoding} Format \textsf{ALU\_DWRR0.M} (Section~\ref{sec:format-ALU-DWRR0.M}), \verb|exucode2 = 0011|\\

\bitfields{ 1/P, 2/00, 2/-- --, 27/upper27, 1/1, 2/11, 1/0, 4/0011, 6/registerW, 2/10, 3/000, 1/0, 1/splat32, 5/lower5, 6/registerZ }


\paragraph{Syntax} \texttt{sbfd} \emph{registerW} = \emph{registerZ}, \emph{upper27\_\discretionary{}{}{}lower5}\emph{splat32}

\paragraph{Description} The \emph{registerZ} is subtracted from the \emph{upper27\_\discretionary{}{}{}lower5}. The result is stored into the \emph{registerW}.


\paragraph{Modifier(s)}
\noindent\begin{itemize}
\item splat32 (cf. splat32 in section \ref{par:modifier-splat32} on page \pageref{par:modifier-splat32})
\end{itemize}

\paragraph{Execution} {Reservation table \textsf{ALU\_TINY.X} (Section~\ref{sec:reservation-ALU-TINY.X})}
~
{\begin{lstlisting}
stage ID:
  new splat32 = _ZX_1(splat32);
stage RR:
  new argument3 = splat32 ? (_WX_32(upper27_lower5) << 32) | _ZX_32(_WX_32(upper27_lower5)) : _SX_32(_WX_32(upper27_lower5));
  new argument2 = GPR[registerZ];
  new result1 = argument3 - argument2;
stage E1:
  GPR[registerW] = result1;
\end{lstlisting}}
\newpage

\paragraph{Encoding} Format \textsf{ALU\_DWRI} (Section~\ref{sec:format-ALU-DWRI}), \verb|exucode2 = 0011|\\

\bitfields{ 1/P, 2/11, 1/0, 4/0011, 6/registerW, 2/00, 10/signed10, 6/registerZ }


\paragraph{Syntax} \texttt{sbfd} \emph{registerW} = \emph{registerZ}, \emph{signed10}

\paragraph{Description} The \emph{registerZ} is subtracted from the \emph{signed10}. The result is stored into the \emph{registerW}.


\paragraph{Execution} {Reservation table \textsf{ALU\_TINY} (Section~\ref{sec:reservation-ALU-TINY})}
~
{\begin{lstlisting}
stage RR:
  new argument3 = _SX_10(signed10);
  new argument2 = GPR[registerZ];
  new result1 = argument3 - argument2;
stage E1:
  GPR[registerW] = result1;
\end{lstlisting}}
\newpage

\paragraph{Encoding} Format \textsf{ALU\_DWRI.X} (Section~\ref{sec:format-ALU-DWRI.X}), \verb|exucode2 = 0011|\\

\bitfields{ 1/P, 2/00, 2/-- --, 27/upper27, 1/1, 2/11, 1/0, 4/0011, 6/registerW, 2/00, 10/lower10, 6/registerZ }


\paragraph{Syntax} \texttt{sbfd} \emph{registerW} = \emph{registerZ}, \emph{upper27\_\discretionary{}{}{}lower10}

\paragraph{Description} The \emph{registerZ} is subtracted from the \emph{upper27\_\discretionary{}{}{}lower10}. The result is stored into the \emph{registerW}.


\paragraph{Execution} {Reservation table \textsf{ALU\_TINY.X} (Section~\ref{sec:reservation-ALU-TINY.X})}
~
{\begin{lstlisting}
stage RR:
  new argument3 = _SX_37(upper27_lower10);
  new argument2 = GPR[registerZ];
  new result1 = argument3 - argument2;
stage E1:
  GPR[registerW] = result1;
\end{lstlisting}}
\newpage

\paragraph{Encoding} Format \textsf{ALU\_DWRI.Y} (Section~\ref{sec:format-ALU-DWRI.Y}), \verb|exucode2 = 0011|\\

\bitfields{ 1/P, 2/00, 2/-- --, 27/extend27, 1/1, 2/00, 2/-- --, 27/upper27, 1/1, 2/11, 1/0, 4/0011, 6/registerW, 2/00, 10/lower10, 6/registerZ }


\paragraph{Syntax} \texttt{sbfd} \emph{registerW} = \emph{registerZ}, \emph{extend27\_\discretionary{}{}{}upper27\_\discretionary{}{}{}lower10}

\paragraph{Description} The \emph{registerZ} is subtracted from the \emph{extend27\_\discretionary{}{}{}upper27\_\discretionary{}{}{}lower10}. The result is stored into the \emph{registerW}.


\paragraph{Execution} {Reservation table \textsf{ALU\_TINY.Y} (Section~\ref{sec:reservation-ALU-TINY.Y})}
~
{\begin{lstlisting}
stage RR:
  new argument3 = _WX_64(extend27_upper27_lower10);
  new argument2 = GPR[registerZ];
  new result1 = argument3 - argument2;
stage E1:
  GPR[registerW] = result1;
\end{lstlisting}}
\newpage
\subsubsection[{\small SBFHQ: Subtract From Half Word Quadruple}]{SBFHQ\hfill Subtract From Half Word Quadruple}
\label{instr:SBFHQ}\index{sbfhq}
 
\paragraph{Encoding} Format \textsf{ALU\_HQWRR0E} (Section~\ref{sec:format-ALU-HQWRR0E}), \verb|exucode2 = 0011|\\

\bitfields{ 1/P, 2/11, 1/1, 4/0011, 5/registerWe, 1/0, 2/10, 3/000, 1/1, 5/registerYe, 1/--, 5/registerZe, 1/0 }


\paragraph{Syntax} \texttt{sbfhq} \emph{registerWe} = \emph{registerZe}, \emph{registerYe}

\paragraph{Description} The \emph{registerYe} and the \emph{registerZe} are interpreted as four half words packed into 64 bits. The \emph{registerZe} half words are subtracted from the \emph{registerYe} half words. The four resulting half words are packed and stored into the \emph{registerWe}.


\paragraph{Execution} {Reservation table \textsf{ALU\_LITE} (Section~\ref{sec:reservation-ALU-LITE})}
~
{\begin{lstlisting}
stage RR:
  new argument3 = GPR[registerYe<<1];
  new argument2 = GPR[registerZe<<1];
  for (I = 0; I < 4; I += 1) {
    result1.16[I] = argument3.16[I] - argument2.16[I]
  };
stage E1:
  GPR[registerWe<<1] = result1;
\end{lstlisting}}
\newpage

\paragraph{Encoding} Format \textsf{ALU\_HQWRR0O} (Section~\ref{sec:format-ALU-HQWRR0O}), \verb|exucode2 = 0011|\\

\bitfields{ 1/P, 2/11, 1/1, 4/0011, 5/registerWo, 1/1, 2/10, 3/000, 1/1, 5/registerYo, 1/--, 5/registerZo, 1/0 }


\paragraph{Syntax} \texttt{sbfhq} \emph{registerWo} = \emph{registerZo}, \emph{registerYo}

\paragraph{Description} The \emph{registerYo} and the \emph{registerZo} are interpreted as four half words packed into 64 bits. The \emph{registerZo} half words are subtracted from the \emph{registerYo} half words. The four resulting half words are packed and stored into the \emph{registerWo}.


\paragraph{Execution} {Reservation table \textsf{ALU\_LITE} (Section~\ref{sec:reservation-ALU-LITE})}
~
{\begin{lstlisting}
stage RR:
  new argument3 = GPR[(registerYo<<1)+1];
  new argument2 = GPR[(registerZo<<1)+1];
  for (I = 0; I < 4; I += 1) {
    result1.16[I] = argument3.16[I] - argument2.16[I]
  };
stage E1:
  GPR[(registerWo<<1)+1] = result1;
\end{lstlisting}}
\newpage

\paragraph{Encoding} Format \textsf{ALU\_HQWRR0E.M} (Section~\ref{sec:format-ALU-HQWRR0E.M}), \verb|exucode2 = 0011|\\

\bitfields{ 1/P, 2/00, 2/-- --, 27/upper27, 1/1, 2/11, 1/1, 4/0011, 5/registerWe, 1/0, 2/10, 3/000, 1/1, 1/splat32, 5/lower5, 5/registerZe, 1/0 }


\paragraph{Syntax} \texttt{sbfhq} \emph{registerWe} = \emph{registerZe}, \emph{upper27\_\discretionary{}{}{}lower5}\emph{splat32}

\paragraph{Description} The \emph{upper27\_\discretionary{}{}{}lower5} and the \emph{registerZe} are interpreted as four half words packed into 64 bits. The \emph{registerZe} half words are subtracted from the \emph{upper27\_\discretionary{}{}{}lower5} half words. The four resulting half words are packed and stored into the \emph{registerWe}.


\paragraph{Modifier(s)}
\noindent\begin{itemize}
\item splat32 (cf. splat32 in section \ref{par:modifier-splat32} on page \pageref{par:modifier-splat32})
\end{itemize}

\paragraph{Execution} {Reservation table \textsf{ALU\_LITE.X} (Section~\ref{sec:reservation-ALU-LITE.X})}
~
{\begin{lstlisting}
stage ID:
  new splat32 = _ZX_1(splat32);
stage RR:
  new argument3 = splat32 ? (_WX_32(upper27_lower5) << 32) | _ZX_32(_WX_32(upper27_lower5)) : _SX_32(_WX_32(upper27_lower5));
  new argument2 = GPR[registerZe<<1];
  for (I = 0; I < 4; I += 1) {
    result1.16[I] = argument3.16[I] - argument2.16[I]
  };
stage E1:
  GPR[registerWe<<1] = result1;
\end{lstlisting}}
\newpage

\paragraph{Encoding} Format \textsf{ALU\_HQWRR0O.M} (Section~\ref{sec:format-ALU-HQWRR0O.M}), \verb|exucode2 = 0011|\\

\bitfields{ 1/P, 2/00, 2/-- --, 27/upper27, 1/1, 2/11, 1/1, 4/0011, 5/registerWo, 1/1, 2/10, 3/000, 1/1, 1/splat32, 5/lower5, 5/registerZo, 1/0 }


\paragraph{Syntax} \texttt{sbfhq} \emph{registerWo} = \emph{registerZo}, \emph{upper27\_\discretionary{}{}{}lower5}\emph{splat32}

\paragraph{Description} The \emph{upper27\_\discretionary{}{}{}lower5} and the \emph{registerZo} are interpreted as four half words packed into 64 bits. The \emph{registerZo} half words are subtracted from the \emph{upper27\_\discretionary{}{}{}lower5} half words. The four resulting half words are packed and stored into the \emph{registerWo}.


\paragraph{Modifier(s)}
\noindent\begin{itemize}
\item splat32 (cf. splat32 in section \ref{par:modifier-splat32} on page \pageref{par:modifier-splat32})
\end{itemize}

\paragraph{Execution} {Reservation table \textsf{ALU\_LITE.X} (Section~\ref{sec:reservation-ALU-LITE.X})}
~
{\begin{lstlisting}
stage ID:
  new splat32 = _ZX_1(splat32);
stage RR:
  new argument3 = splat32 ? (_WX_32(upper27_lower5) << 32) | _ZX_32(_WX_32(upper27_lower5)) : _SX_32(_WX_32(upper27_lower5));
  new argument2 = GPR[(registerZo<<1)+1];
  for (I = 0; I < 4; I += 1) {
    result1.16[I] = argument3.16[I] - argument2.16[I]
  };
stage E1:
  GPR[(registerWo<<1)+1] = result1;
\end{lstlisting}}
\newpage
\subsubsection[{\small SBFQ: Subtract From Quadruple Word}]{SBFQ\hfill Subtract From Quadruple Word}
\label{instr:SBFQ}\index{sbfq}
 
\paragraph{Encoding} Format \textsf{ALU\_QWRR0} (Section~\ref{sec:format-ALU-QWRR0}), \verb|exucode2 = 0011|\\

\bitfields{ 1/P, 2/11, 1/0, 4/0011, 5/registerM, 1/--, 2/10, 3/000, 1/1, 5/registerO, 1/--, 5/registerP, 1/-- }


\paragraph{Syntax} \texttt{sbfq} \emph{registerM} = \emph{registerP}, \emph{registerO}

\paragraph{Description} The \emph{registerP} is subtracted from the \emph{registerO}. The result is stored into the \emph{registerM}.


\paragraph{Execution} {Reservation table \textsf{ALU\_LITE} (Section~\ref{sec:reservation-ALU-LITE})}
~
{\begin{lstlisting}
stage RR:
  new argument3 = PGR[registerO];
  new argument2 = PGR[registerP];
  new result1 = argument3 - argument2;
stage E1:
  PGR[registerM] = result1;
\end{lstlisting}}
\newpage

\paragraph{Encoding} Format \textsf{ALU\_QWRR0.M} (Section~\ref{sec:format-ALU-QWRR0.M}), \verb|exucode2 = 0011|\\

\bitfields{ 1/P, 2/00, 2/-- --, 27/upper27, 1/1, 2/11, 1/0, 4/0011, 5/registerM, 1/--, 2/10, 3/000, 1/1, 1/splat32, 5/lower5, 5/registerP, 1/1 }


\paragraph{Syntax} \texttt{sbfq} \emph{registerM} = \emph{registerP}, \emph{upper27\_\discretionary{}{}{}lower5}\emph{splat32}

\paragraph{Description} The \emph{registerP} is subtracted from the \emph{upper27\_\discretionary{}{}{}lower5}. The result is stored into the \emph{registerM}.


\paragraph{Modifier(s)}
\noindent\begin{itemize}
\item splat32 (cf. splat32 in section \ref{par:modifier-splat32} on page \pageref{par:modifier-splat32})
\end{itemize}

\paragraph{Execution} {Reservation table \textsf{ALU\_LITE.X} (Section~\ref{sec:reservation-ALU-LITE.X})}
~
{\begin{lstlisting}
stage ID:
  new splat32 = _ZX_1(splat32);
stage RR:
  new argument3 = splat32 ? (_WX_32(upper27_lower5) << 32) | _ZX_32(_WX_32(upper27_lower5)) : _SX_32(_WX_32(upper27_lower5));
  new argument2 = PGR[registerP];
  new result1 = argument3 - argument2;
stage E1:
  PGR[registerM] = result1;
\end{lstlisting}}
\newpage
\subsubsection[{\small SBFSBO: Subtract Saturated Byte Octuple}]{SBFSBO\hfill Subtract Saturated Byte Octuple}
\label{instr:SBFSBO}\index{sbfsbo}
 
\paragraph{Encoding} Format \textsf{ALU\_BOWRR1E} (Section~\ref{sec:format-ALU-BOWRR1E}), \verb|exucode2 = 1001|\\

\bitfields{ 1/P, 2/11, 1/1, 4/1001, 5/registerWe, 1/0, 2/10, 3/101, 1/1, 5/registerYe, 1/--, 5/registerZe, 1/1 }


\paragraph{Syntax} \texttt{sbfsbo} \emph{registerWe} = \emph{registerZe}, \emph{registerYe}

\paragraph{Description} The \emph{registerYe} and the \emph{registerZe} are interpreted as eight bytes packed into 64 bits. The \emph{registerZe} bytes are subtracted with saturation to 8 bits from the \emph{registerYe} bytes. The eight resulting bytes are packed and stored into the \emph{registerWe}.


\paragraph{Execution} {Reservation table \textsf{ALU\_LITE} (Section~\ref{sec:reservation-ALU-LITE})}
~
{\begin{lstlisting}
stage RR:
  new argument3 = GPR[registerYe<<1];
  new argument2 = GPR[registerZe<<1];
  for (I = 0; I < 8; I += 1) {
    result1.8[I] = _SAT_8(_SX_8(argument3.8[I]) - _SX_8(argument2.8[I]))
  };
stage E1:
  GPR[registerWe<<1] = result1;
\end{lstlisting}}
\newpage

\paragraph{Encoding} Format \textsf{ALU\_BOWRR1O} (Section~\ref{sec:format-ALU-BOWRR1O}), \verb|exucode2 = 1001|\\

\bitfields{ 1/P, 2/11, 1/1, 4/1001, 5/registerWo, 1/1, 2/10, 3/101, 1/1, 5/registerYo, 1/--, 5/registerZo, 1/1 }


\paragraph{Syntax} \texttt{sbfsbo} \emph{registerWo} = \emph{registerZo}, \emph{registerYo}

\paragraph{Description} The \emph{registerYo} and the \emph{registerZo} are interpreted as eight bytes packed into 64 bits. The \emph{registerZo} bytes are subtracted with saturation to 8 bits from the \emph{registerYo} bytes. The eight resulting bytes are packed and stored into the \emph{registerWo}.


\paragraph{Execution} {Reservation table \textsf{ALU\_LITE} (Section~\ref{sec:reservation-ALU-LITE})}
~
{\begin{lstlisting}
stage RR:
  new argument3 = GPR[(registerYo<<1)+1];
  new argument2 = GPR[(registerZo<<1)+1];
  for (I = 0; I < 8; I += 1) {
    result1.8[I] = _SAT_8(_SX_8(argument3.8[I]) - _SX_8(argument2.8[I]))
  };
stage E1:
  GPR[(registerWo<<1)+1] = result1;
\end{lstlisting}}
\newpage

\paragraph{Encoding} Format \textsf{ALU\_BOWRR1E.M} (Section~\ref{sec:format-ALU-BOWRR1E.M}), \verb|exucode2 = 1001|\\

\bitfields{ 1/P, 2/00, 2/-- --, 27/upper27, 1/1, 2/11, 1/1, 4/1001, 5/registerWe, 1/0, 2/10, 3/101, 1/1, 1/splat32, 5/lower5, 5/registerZe, 1/1 }


\paragraph{Syntax} \texttt{sbfsbo} \emph{registerWe} = \emph{registerZe}, \emph{upper27\_\discretionary{}{}{}lower5}\emph{splat32}

\paragraph{Description} The \emph{upper27\_\discretionary{}{}{}lower5} and the \emph{registerZe} are interpreted as eight bytes packed into 64 bits. The \emph{registerZe} bytes are subtracted with saturation to 8 bits from the \emph{upper27\_\discretionary{}{}{}lower5} bytes. The eight resulting bytes are packed and stored into the \emph{registerWe}.


\paragraph{Modifier(s)}
\noindent\begin{itemize}
\item splat32 (cf. splat32 in section \ref{par:modifier-splat32} on page \pageref{par:modifier-splat32})
\end{itemize}

\paragraph{Execution} {Reservation table \textsf{ALU\_LITE.X} (Section~\ref{sec:reservation-ALU-LITE.X})}
~
{\begin{lstlisting}
stage ID:
  new splat32 = _ZX_1(splat32);
stage RR:
  new argument3 = splat32 ? (_WX_32(upper27_lower5) << 32) | _ZX_32(_WX_32(upper27_lower5)) : _SX_32(_WX_32(upper27_lower5));
  new argument2 = GPR[registerZe<<1];
  for (I = 0; I < 8; I += 1) {
    result1.8[I] = _SAT_8(_SX_8(argument3.8[I]) - _SX_8(argument2.8[I]))
  };
stage E1:
  GPR[registerWe<<1] = result1;
\end{lstlisting}}
\newpage

\paragraph{Encoding} Format \textsf{ALU\_BOWRR1O.M} (Section~\ref{sec:format-ALU-BOWRR1O.M}), \verb|exucode2 = 1001|\\

\bitfields{ 1/P, 2/00, 2/-- --, 27/upper27, 1/1, 2/11, 1/1, 4/1001, 5/registerWo, 1/1, 2/10, 3/101, 1/1, 1/splat32, 5/lower5, 5/registerZo, 1/1 }


\paragraph{Syntax} \texttt{sbfsbo} \emph{registerWo} = \emph{registerZo}, \emph{upper27\_\discretionary{}{}{}lower5}\emph{splat32}

\paragraph{Description} The \emph{upper27\_\discretionary{}{}{}lower5} and the \emph{registerZo} are interpreted as eight bytes packed into 64 bits. The \emph{registerZo} bytes are subtracted with saturation to 8 bits from the \emph{upper27\_\discretionary{}{}{}lower5} bytes. The eight resulting bytes are packed and stored into the \emph{registerWo}.


\paragraph{Modifier(s)}
\noindent\begin{itemize}
\item splat32 (cf. splat32 in section \ref{par:modifier-splat32} on page \pageref{par:modifier-splat32})
\end{itemize}

\paragraph{Execution} {Reservation table \textsf{ALU\_LITE.X} (Section~\ref{sec:reservation-ALU-LITE.X})}
~
{\begin{lstlisting}
stage ID:
  new splat32 = _ZX_1(splat32);
stage RR:
  new argument3 = splat32 ? (_WX_32(upper27_lower5) << 32) | _ZX_32(_WX_32(upper27_lower5)) : _SX_32(_WX_32(upper27_lower5));
  new argument2 = GPR[(registerZo<<1)+1];
  for (I = 0; I < 8; I += 1) {
    result1.8[I] = _SAT_8(_SX_8(argument3.8[I]) - _SX_8(argument2.8[I]))
  };
stage E1:
  GPR[(registerWo<<1)+1] = result1;
\end{lstlisting}}
\newpage
\subsubsection[{\small SBFSD: Subtract Saturated Double Word}]{SBFSD\hfill Subtract Saturated Double Word}
\label{instr:SBFSD}\index{sbfsd}
 
\paragraph{Encoding} Format \textsf{ALU\_DWRR1} (Section~\ref{sec:format-ALU-DWRR1}), \verb|exucode2 = 1001|\\

\bitfields{ 1/P, 2/11, 1/0, 4/1001, 6/registerW, 2/10, 3/101, 1/0, 6/registerY, 6/registerZ }


\paragraph{Syntax} \texttt{sbfsd} \emph{registerW} = \emph{registerZ}, \emph{registerY}

\paragraph{Description} The \emph{registerZ} is subtracted from the \emph{registerY} with 64-bit saturation.


\paragraph{Execution} {Reservation table \textsf{ALU\_TINY} (Section~\ref{sec:reservation-ALU-TINY})}
~
{\begin{lstlisting}
stage RR:
  new argument3 = GPR[registerY];
  new argument2 = GPR[registerZ];
  new result1 = _SAT_64(_SX_64(argument3) - _SX_64(argument2));
stage E1:
  GPR[registerW] = result1;
\end{lstlisting}}
\newpage

\paragraph{Encoding} Format \textsf{ALU\_DWRR1.M} (Section~\ref{sec:format-ALU-DWRR1.M}), \verb|exucode2 = 1001|\\

\bitfields{ 1/P, 2/00, 2/-- --, 27/upper27, 1/1, 2/11, 1/0, 4/1001, 6/registerW, 2/10, 3/101, 1/0, 1/splat32, 5/lower5, 6/registerZ }


\paragraph{Syntax} \texttt{sbfsd} \emph{registerW} = \emph{registerZ}, \emph{upper27\_\discretionary{}{}{}lower5}\emph{splat32}

\paragraph{Description} The \emph{registerZ} is subtracted from the \emph{upper27\_\discretionary{}{}{}lower5} with 64-bit saturation.


\paragraph{Modifier(s)}
\noindent\begin{itemize}
\item splat32 (cf. splat32 in section \ref{par:modifier-splat32} on page \pageref{par:modifier-splat32})
\end{itemize}

\paragraph{Execution} {Reservation table \textsf{ALU\_TINY.X} (Section~\ref{sec:reservation-ALU-TINY.X})}
~
{\begin{lstlisting}
stage ID:
  new splat32 = _ZX_1(splat32);
stage RR:
  new argument3 = splat32 ? (_WX_32(upper27_lower5) << 32) | _ZX_32(_WX_32(upper27_lower5)) : _SX_32(_WX_32(upper27_lower5));
  new argument2 = GPR[registerZ];
  new result1 = _SAT_64(_SX_64(argument3) - _SX_64(argument2));
stage E1:
  GPR[registerW] = result1;
\end{lstlisting}}
\newpage
\subsubsection[{\small SBFSHQ: Subtract Saturated Half Word Quadruple}]{SBFSHQ\hfill Subtract Saturated Half Word Quadruple}
\label{instr:SBFSHQ}\index{sbfshq}
 
\paragraph{Encoding} Format \textsf{ALU\_HQWRR1E} (Section~\ref{sec:format-ALU-HQWRR1E}), \verb|exucode2 = 1001|\\

\bitfields{ 1/P, 2/11, 1/1, 4/1001, 5/registerWe, 1/0, 2/10, 3/101, 1/1, 5/registerYe, 1/--, 5/registerZe, 1/0 }


\paragraph{Syntax} \texttt{sbfshq} \emph{registerWe} = \emph{registerZe}, \emph{registerYe}

\paragraph{Description} The \emph{registerYe} and the \emph{registerZe} are interpreted as four half words packed into 64 bits. The \emph{registerZe} half words are subtracted with saturation to 16 bits from the \emph{registerYe} half words. The four resulting half words are packed and stored into the \emph{registerWe}.


\paragraph{Execution} {Reservation table \textsf{ALU\_LITE} (Section~\ref{sec:reservation-ALU-LITE})}
~
{\begin{lstlisting}
stage RR:
  new argument3 = GPR[registerYe<<1];
  new argument2 = GPR[registerZe<<1];
  for (I = 0; I < 4; I += 1) {
    result1.16[I] = _SAT_16(_SX_16(argument3.16[I]) - _SX_16(argument2.16[I]))
  };
stage E1:
  GPR[registerWe<<1] = result1;
\end{lstlisting}}
\newpage

\paragraph{Encoding} Format \textsf{ALU\_HQWRR1O} (Section~\ref{sec:format-ALU-HQWRR1O}), \verb|exucode2 = 1001|\\

\bitfields{ 1/P, 2/11, 1/1, 4/1001, 5/registerWo, 1/1, 2/10, 3/101, 1/1, 5/registerYo, 1/--, 5/registerZo, 1/0 }


\paragraph{Syntax} \texttt{sbfshq} \emph{registerWo} = \emph{registerZo}, \emph{registerYo}

\paragraph{Description} The \emph{registerYo} and the \emph{registerZo} are interpreted as four half words packed into 64 bits. The \emph{registerZo} half words are subtracted with saturation to 16 bits from the \emph{registerYo} half words. The four resulting half words are packed and stored into the \emph{registerWo}.


\paragraph{Execution} {Reservation table \textsf{ALU\_LITE} (Section~\ref{sec:reservation-ALU-LITE})}
~
{\begin{lstlisting}
stage RR:
  new argument3 = GPR[(registerYo<<1)+1];
  new argument2 = GPR[(registerZo<<1)+1];
  for (I = 0; I < 4; I += 1) {
    result1.16[I] = _SAT_16(_SX_16(argument3.16[I]) - _SX_16(argument2.16[I]))
  };
stage E1:
  GPR[(registerWo<<1)+1] = result1;
\end{lstlisting}}
\newpage

\paragraph{Encoding} Format \textsf{ALU\_HQWRR1E.M} (Section~\ref{sec:format-ALU-HQWRR1E.M}), \verb|exucode2 = 1001|\\

\bitfields{ 1/P, 2/00, 2/-- --, 27/upper27, 1/1, 2/11, 1/1, 4/1001, 5/registerWe, 1/0, 2/10, 3/101, 1/1, 1/splat32, 5/lower5, 5/registerZe, 1/0 }


\paragraph{Syntax} \texttt{sbfshq} \emph{registerWe} = \emph{registerZe}, \emph{upper27\_\discretionary{}{}{}lower5}\emph{splat32}

\paragraph{Description} The \emph{upper27\_\discretionary{}{}{}lower5} and the \emph{registerZe} are interpreted as four half words packed into 64 bits. The \emph{registerZe} half words are subtracted with saturation to 16 bits from the \emph{upper27\_\discretionary{}{}{}lower5} half words. The four resulting half words are packed and stored into the \emph{registerWe}.


\paragraph{Modifier(s)}
\noindent\begin{itemize}
\item splat32 (cf. splat32 in section \ref{par:modifier-splat32} on page \pageref{par:modifier-splat32})
\end{itemize}

\paragraph{Execution} {Reservation table \textsf{ALU\_LITE.X} (Section~\ref{sec:reservation-ALU-LITE.X})}
~
{\begin{lstlisting}
stage ID:
  new splat32 = _ZX_1(splat32);
stage RR:
  new argument3 = splat32 ? (_WX_32(upper27_lower5) << 32) | _ZX_32(_WX_32(upper27_lower5)) : _SX_32(_WX_32(upper27_lower5));
  new argument2 = GPR[registerZe<<1];
  for (I = 0; I < 4; I += 1) {
    result1.16[I] = _SAT_16(_SX_16(argument3.16[I]) - _SX_16(argument2.16[I]))
  };
stage E1:
  GPR[registerWe<<1] = result1;
\end{lstlisting}}
\newpage

\paragraph{Encoding} Format \textsf{ALU\_HQWRR1O.M} (Section~\ref{sec:format-ALU-HQWRR1O.M}), \verb|exucode2 = 1001|\\

\bitfields{ 1/P, 2/00, 2/-- --, 27/upper27, 1/1, 2/11, 1/1, 4/1001, 5/registerWo, 1/1, 2/10, 3/101, 1/1, 1/splat32, 5/lower5, 5/registerZo, 1/0 }


\paragraph{Syntax} \texttt{sbfshq} \emph{registerWo} = \emph{registerZo}, \emph{upper27\_\discretionary{}{}{}lower5}\emph{splat32}

\paragraph{Description} The \emph{upper27\_\discretionary{}{}{}lower5} and the \emph{registerZo} are interpreted as four half words packed into 64 bits. The \emph{registerZo} half words are subtracted with saturation to 16 bits from the \emph{upper27\_\discretionary{}{}{}lower5} half words. The four resulting half words are packed and stored into the \emph{registerWo}.


\paragraph{Modifier(s)}
\noindent\begin{itemize}
\item splat32 (cf. splat32 in section \ref{par:modifier-splat32} on page \pageref{par:modifier-splat32})
\end{itemize}

\paragraph{Execution} {Reservation table \textsf{ALU\_LITE.X} (Section~\ref{sec:reservation-ALU-LITE.X})}
~
{\begin{lstlisting}
stage ID:
  new splat32 = _ZX_1(splat32);
stage RR:
  new argument3 = splat32 ? (_WX_32(upper27_lower5) << 32) | _ZX_32(_WX_32(upper27_lower5)) : _SX_32(_WX_32(upper27_lower5));
  new argument2 = GPR[(registerZo<<1)+1];
  for (I = 0; I < 4; I += 1) {
    result1.16[I] = _SAT_16(_SX_16(argument3.16[I]) - _SX_16(argument2.16[I]))
  };
stage E1:
  GPR[(registerWo<<1)+1] = result1;
\end{lstlisting}}
\newpage
\subsubsection[{\small SBFSW: Subtract Saturated Word}]{SBFSW\hfill Subtract Saturated Word}
\label{instr:SBFSW}\index{sbfsw}
 
\paragraph{Encoding} Format \textsf{ALU\_WWRR1} (Section~\ref{sec:format-ALU-WWRR1}), \verb|exucode2 = 1001|\\

\bitfields{ 1/P, 2/11, 1/0, 4/1001, 6/registerW, 2/11, 3/101, 1/signextw, 6/registerY, 6/registerZ }


\paragraph{Syntax} \texttt{sbfsw}\emph{signextw} \emph{registerW} = \emph{registerZ}, \emph{registerY}

\paragraph{Description} The \emph{registerZ} is subtracted from the \emph{registerY} with 32-bit saturation. The result with the 32 upper bits cleared is stored into the \emph{registerW}.


\paragraph{Modifier(s)}
\noindent\begin{itemize}
\item signextw (cf. signextw in section \ref{par:modifier-signextw} on page \pageref{par:modifier-signextw})
\end{itemize}

\paragraph{Execution} {Reservation table \textsf{ALU\_TINY} (Section~\ref{sec:reservation-ALU-TINY})}
~
{\begin{lstlisting}
stage ID:
  new signextw = _ZX_1(signextw);
stage RR:
  new argument3 = GPR[registerY];
  new argument2 = GPR[registerZ];
  new result1 = _SAT_32(_SX_32(argument3) - _SX_32(argument2));
stage E1:
  GPR[registerW] = signextw ? _SX_32(result1) : _ZX_32(result1);
\end{lstlisting}}
\newpage

\paragraph{Encoding} Format \textsf{ALU\_WWRR1.W} (Section~\ref{sec:format-ALU-WWRR1.W}), \verb|exucode2 = 1001|\\

\bitfields{ 1/P, 2/00, 2/-- --, 27/upper27, 1/1, 2/11, 1/0, 4/1001, 6/registerW, 2/11, 3/101, 1/signextw, 1/0, 5/lower5, 6/registerZ }


\paragraph{Syntax} \texttt{sbfsw}\emph{signextw} \emph{registerW} = \emph{registerZ}, \emph{upper27\_\discretionary{}{}{}lower5}

\paragraph{Description} The \emph{registerZ} is subtracted from the \emph{upper27\_\discretionary{}{}{}lower5} with 32-bit saturation. The result with the 32 upper bits cleared is stored into the \emph{registerW}.


\paragraph{Modifier(s)}
\noindent\begin{itemize}
\item signextw (cf. signextw in section \ref{par:modifier-signextw} on page \pageref{par:modifier-signextw})
\end{itemize}

\paragraph{Execution} {Reservation table \textsf{ALU\_TINY.X} (Section~\ref{sec:reservation-ALU-TINY.X})}
~
{\begin{lstlisting}
stage ID:
  new signextw = _ZX_1(signextw);
stage RR:
  new argument3 = _WX_32(upper27_lower5);
  new argument2 = GPR[registerZ];
  new result1 = _SAT_32(_SX_32(argument3) - _SX_32(argument2));
stage E1:
  GPR[registerW] = signextw ? _SX_32(result1) : _ZX_32(result1);
\end{lstlisting}}
\newpage
\subsubsection[{\small SBFSWP: Subtract Saturated Word Pair}]{SBFSWP\hfill Subtract Saturated Word Pair}
\label{instr:SBFSWP}\index{sbfswp}
 
\paragraph{Encoding} Format \textsf{ALU\_WPWRR1E} (Section~\ref{sec:format-ALU-WPWRR1E}), \verb|exucode2 = 1001|\\

\bitfields{ 1/P, 2/11, 1/1, 4/1001, 5/registerWe, 1/0, 2/10, 3/101, 1/0, 5/registerYe, 1/--, 5/registerZe, 1/1 }


\paragraph{Syntax} \texttt{sbfswp} \emph{registerWe} = \emph{registerZe}, \emph{registerYe}

\paragraph{Description} The \emph{registerYe} and the \emph{registerZe} are interpreted as two words packed into 64 bits. The \emph{registerZe} words are subtracted with saturation to 32 bits from the \emph{registerYe} words. The two resulting words are packed and stored into the \emph{registerWe}.


\paragraph{Execution} {Reservation table \textsf{ALU\_LITE} (Section~\ref{sec:reservation-ALU-LITE})}
~
{\begin{lstlisting}
stage RR:
  new argument3 = GPR[registerYe<<1];
  new argument2 = GPR[registerZe<<1];
  for (I = 0; I < 2; I += 1) {
    result1.32[I] = _SAT_32(_SX_32(argument3.32[I]) - _SX_32(argument2.32[I]))
  };
stage E1:
  GPR[registerWe<<1] = result1;
\end{lstlisting}}
\newpage

\paragraph{Encoding} Format \textsf{ALU\_WPWRR1O} (Section~\ref{sec:format-ALU-WPWRR1O}), \verb|exucode2 = 1001|\\

\bitfields{ 1/P, 2/11, 1/1, 4/1001, 5/registerWo, 1/1, 2/10, 3/101, 1/0, 5/registerYo, 1/--, 5/registerZo, 1/1 }


\paragraph{Syntax} \texttt{sbfswp} \emph{registerWo} = \emph{registerZo}, \emph{registerYo}

\paragraph{Description} The \emph{registerYo} and the \emph{registerZo} are interpreted as two words packed into 64 bits. The \emph{registerZo} words are subtracted with saturation to 32 bits from the \emph{registerYo} words. The two resulting words are packed and stored into the \emph{registerWo}.


\paragraph{Execution} {Reservation table \textsf{ALU\_LITE} (Section~\ref{sec:reservation-ALU-LITE})}
~
{\begin{lstlisting}
stage RR:
  new argument3 = GPR[(registerYo<<1)+1];
  new argument2 = GPR[(registerZo<<1)+1];
  for (I = 0; I < 2; I += 1) {
    result1.32[I] = _SAT_32(_SX_32(argument3.32[I]) - _SX_32(argument2.32[I]))
  };
stage E1:
  GPR[(registerWo<<1)+1] = result1;
\end{lstlisting}}
\newpage

\paragraph{Encoding} Format \textsf{ALU\_WPWRR1E.M} (Section~\ref{sec:format-ALU-WPWRR1E.M}), \verb|exucode2 = 1001|\\

\bitfields{ 1/P, 2/00, 2/-- --, 27/upper27, 1/1, 2/11, 1/1, 4/1001, 5/registerWe, 1/0, 2/10, 3/101, 1/0, 1/splat32, 5/lower5, 5/registerZe, 1/1 }


\paragraph{Syntax} \texttt{sbfswp} \emph{registerWe} = \emph{registerZe}, \emph{upper27\_\discretionary{}{}{}lower5}\emph{splat32}

\paragraph{Description} The \emph{upper27\_\discretionary{}{}{}lower5} and the \emph{registerZe} are interpreted as two words packed into 64 bits. The \emph{registerZe} words are subtracted with saturation to 32 bits from the \emph{upper27\_\discretionary{}{}{}lower5} words. The two resulting words are packed and stored into the \emph{registerWe}.


\paragraph{Modifier(s)}
\noindent\begin{itemize}
\item splat32 (cf. splat32 in section \ref{par:modifier-splat32} on page \pageref{par:modifier-splat32})
\end{itemize}

\paragraph{Execution} {Reservation table \textsf{ALU\_LITE.X} (Section~\ref{sec:reservation-ALU-LITE.X})}
~
{\begin{lstlisting}
stage ID:
  new splat32 = _ZX_1(splat32);
stage RR:
  new argument3 = splat32 ? (_WX_32(upper27_lower5) << 32) | _ZX_32(_WX_32(upper27_lower5)) : _SX_32(_WX_32(upper27_lower5));
  new argument2 = GPR[registerZe<<1];
  for (I = 0; I < 2; I += 1) {
    result1.32[I] = _SAT_32(_SX_32(argument3.32[I]) - _SX_32(argument2.32[I]))
  };
stage E1:
  GPR[registerWe<<1] = result1;
\end{lstlisting}}
\newpage

\paragraph{Encoding} Format \textsf{ALU\_WPWRR1O.M} (Section~\ref{sec:format-ALU-WPWRR1O.M}), \verb|exucode2 = 1001|\\

\bitfields{ 1/P, 2/00, 2/-- --, 27/upper27, 1/1, 2/11, 1/1, 4/1001, 5/registerWo, 1/1, 2/10, 3/101, 1/0, 1/splat32, 5/lower5, 5/registerZo, 1/1 }


\paragraph{Syntax} \texttt{sbfswp} \emph{registerWo} = \emph{registerZo}, \emph{upper27\_\discretionary{}{}{}lower5}\emph{splat32}

\paragraph{Description} The \emph{upper27\_\discretionary{}{}{}lower5} and the \emph{registerZo} are interpreted as two words packed into 64 bits. The \emph{registerZo} words are subtracted with saturation to 32 bits from the \emph{upper27\_\discretionary{}{}{}lower5} words. The two resulting words are packed and stored into the \emph{registerWo}.


\paragraph{Modifier(s)}
\noindent\begin{itemize}
\item splat32 (cf. splat32 in section \ref{par:modifier-splat32} on page \pageref{par:modifier-splat32})
\end{itemize}

\paragraph{Execution} {Reservation table \textsf{ALU\_LITE.X} (Section~\ref{sec:reservation-ALU-LITE.X})}
~
{\begin{lstlisting}
stage ID:
  new splat32 = _ZX_1(splat32);
stage RR:
  new argument3 = splat32 ? (_WX_32(upper27_lower5) << 32) | _ZX_32(_WX_32(upper27_lower5)) : _SX_32(_WX_32(upper27_lower5));
  new argument2 = GPR[(registerZo<<1)+1];
  for (I = 0; I < 2; I += 1) {
    result1.32[I] = _SAT_32(_SX_32(argument3.32[I]) - _SX_32(argument2.32[I]))
  };
stage E1:
  GPR[(registerWo<<1)+1] = result1;
\end{lstlisting}}
\newpage
\subsubsection[{\small SBFUSBO: Subtract Unsigned Saturated Byte Octuple}]{SBFUSBO\hfill Subtract Unsigned Saturated Byte Octuple}
\label{instr:SBFUSBO}\index{sbfusbo}
 
\paragraph{Encoding} Format \textsf{ALU\_BOWRR1E} (Section~\ref{sec:format-ALU-BOWRR1E}), \verb|exucode2 = 1011|\\

\bitfields{ 1/P, 2/11, 1/1, 4/1011, 5/registerWe, 1/0, 2/10, 3/101, 1/1, 5/registerYe, 1/--, 5/registerZe, 1/1 }


\paragraph{Syntax} \texttt{sbfusbo} \emph{registerWe} = \emph{registerZe}, \emph{registerYe}

\paragraph{Description} The \emph{registerYe} and the \emph{registerZe} are interpreted as eight bytes packed into 64 bits. The \emph{registerZe} bytes are subtracted with unsigned saturation to 8 bits from the \emph{registerYe} bytes. The eight resulting bytes are packed and stored into the \emph{registerWe}.


\paragraph{Execution} {Reservation table \textsf{ALU\_LITE} (Section~\ref{sec:reservation-ALU-LITE})}
~
{\begin{lstlisting}
stage RR:
  new argument3 = GPR[registerYe<<1];
  new argument2 = GPR[registerZe<<1];
  for (I = 0; I < 8; I += 1) {
    result1.8[I] = _SATU_8(_ZX_8(argument3.8[I]) - _ZX_8(argument2.8[I]))
  };
stage E1:
  GPR[registerWe<<1] = result1;
\end{lstlisting}}
\newpage

\paragraph{Encoding} Format \textsf{ALU\_BOWRR1O} (Section~\ref{sec:format-ALU-BOWRR1O}), \verb|exucode2 = 1011|\\

\bitfields{ 1/P, 2/11, 1/1, 4/1011, 5/registerWo, 1/1, 2/10, 3/101, 1/1, 5/registerYo, 1/--, 5/registerZo, 1/1 }


\paragraph{Syntax} \texttt{sbfusbo} \emph{registerWo} = \emph{registerZo}, \emph{registerYo}

\paragraph{Description} The \emph{registerYo} and the \emph{registerZo} are interpreted as eight bytes packed into 64 bits. The \emph{registerZo} bytes are subtracted with unsigned saturation to 8 bits from the \emph{registerYo} bytes. The eight resulting bytes are packed and stored into the \emph{registerWo}.


\paragraph{Execution} {Reservation table \textsf{ALU\_LITE} (Section~\ref{sec:reservation-ALU-LITE})}
~
{\begin{lstlisting}
stage RR:
  new argument3 = GPR[(registerYo<<1)+1];
  new argument2 = GPR[(registerZo<<1)+1];
  for (I = 0; I < 8; I += 1) {
    result1.8[I] = _SATU_8(_ZX_8(argument3.8[I]) - _ZX_8(argument2.8[I]))
  };
stage E1:
  GPR[(registerWo<<1)+1] = result1;
\end{lstlisting}}
\newpage

\paragraph{Encoding} Format \textsf{ALU\_BOWRR1E.M} (Section~\ref{sec:format-ALU-BOWRR1E.M}), \verb|exucode2 = 1011|\\

\bitfields{ 1/P, 2/00, 2/-- --, 27/upper27, 1/1, 2/11, 1/1, 4/1011, 5/registerWe, 1/0, 2/10, 3/101, 1/1, 1/splat32, 5/lower5, 5/registerZe, 1/1 }


\paragraph{Syntax} \texttt{sbfusbo} \emph{registerWe} = \emph{registerZe}, \emph{upper27\_\discretionary{}{}{}lower5}\emph{splat32}

\paragraph{Description} The \emph{upper27\_\discretionary{}{}{}lower5} and the \emph{registerZe} are interpreted as eight bytes packed into 64 bits. The \emph{registerZe} bytes are subtracted with unsigned saturation to 8 bits from the \emph{upper27\_\discretionary{}{}{}lower5} bytes. The eight resulting bytes are packed and stored into the \emph{registerWe}.


\paragraph{Modifier(s)}
\noindent\begin{itemize}
\item splat32 (cf. splat32 in section \ref{par:modifier-splat32} on page \pageref{par:modifier-splat32})
\end{itemize}

\paragraph{Execution} {Reservation table \textsf{ALU\_LITE.X} (Section~\ref{sec:reservation-ALU-LITE.X})}
~
{\begin{lstlisting}
stage ID:
  new splat32 = _ZX_1(splat32);
stage RR:
  new argument3 = splat32 ? (_WX_32(upper27_lower5) << 32) | _ZX_32(_WX_32(upper27_lower5)) : _SX_32(_WX_32(upper27_lower5));
  new argument2 = GPR[registerZe<<1];
  for (I = 0; I < 8; I += 1) {
    result1.8[I] = _SATU_8(_ZX_8(argument3.8[I]) - _ZX_8(argument2.8[I]))
  };
stage E1:
  GPR[registerWe<<1] = result1;
\end{lstlisting}}
\newpage

\paragraph{Encoding} Format \textsf{ALU\_BOWRR1O.M} (Section~\ref{sec:format-ALU-BOWRR1O.M}), \verb|exucode2 = 1011|\\

\bitfields{ 1/P, 2/00, 2/-- --, 27/upper27, 1/1, 2/11, 1/1, 4/1011, 5/registerWo, 1/1, 2/10, 3/101, 1/1, 1/splat32, 5/lower5, 5/registerZo, 1/1 }


\paragraph{Syntax} \texttt{sbfusbo} \emph{registerWo} = \emph{registerZo}, \emph{upper27\_\discretionary{}{}{}lower5}\emph{splat32}

\paragraph{Description} The \emph{upper27\_\discretionary{}{}{}lower5} and the \emph{registerZo} are interpreted as eight bytes packed into 64 bits. The \emph{registerZo} bytes are subtracted with unsigned saturation to 8 bits from the \emph{upper27\_\discretionary{}{}{}lower5} bytes. The eight resulting bytes are packed and stored into the \emph{registerWo}.


\paragraph{Modifier(s)}
\noindent\begin{itemize}
\item splat32 (cf. splat32 in section \ref{par:modifier-splat32} on page \pageref{par:modifier-splat32})
\end{itemize}

\paragraph{Execution} {Reservation table \textsf{ALU\_LITE.X} (Section~\ref{sec:reservation-ALU-LITE.X})}
~
{\begin{lstlisting}
stage ID:
  new splat32 = _ZX_1(splat32);
stage RR:
  new argument3 = splat32 ? (_WX_32(upper27_lower5) << 32) | _ZX_32(_WX_32(upper27_lower5)) : _SX_32(_WX_32(upper27_lower5));
  new argument2 = GPR[(registerZo<<1)+1];
  for (I = 0; I < 8; I += 1) {
    result1.8[I] = _SATU_8(_ZX_8(argument3.8[I]) - _ZX_8(argument2.8[I]))
  };
stage E1:
  GPR[(registerWo<<1)+1] = result1;
\end{lstlisting}}
\newpage
\subsubsection[{\small SBFUSD: Subtract Unsigned Saturated Double Word}]{SBFUSD\hfill Subtract Unsigned Saturated Double Word}
\label{instr:SBFUSD}\index{sbfusd}
 
\paragraph{Encoding} Format \textsf{ALU\_DWRR1} (Section~\ref{sec:format-ALU-DWRR1}), \verb|exucode2 = 1011|\\

\bitfields{ 1/P, 2/11, 1/0, 4/1011, 6/registerW, 2/10, 3/101, 1/0, 6/registerY, 6/registerZ }


\paragraph{Syntax} \texttt{sbfusd} \emph{registerW} = \emph{registerZ}, \emph{registerY}

\paragraph{Description} The \emph{registerZ} is subtracted from the \emph{registerY} with 64-bit unsigned saturation.


\paragraph{Execution} {Reservation table \textsf{ALU\_TINY} (Section~\ref{sec:reservation-ALU-TINY})}
~
{\begin{lstlisting}
stage RR:
  new argument3 = GPR[registerY];
  new argument2 = GPR[registerZ];
  new result1 = _SATU_64(_ZX_64(argument3) - _ZX_64(argument2));
stage E1:
  GPR[registerW] = result1;
\end{lstlisting}}
\newpage

\paragraph{Encoding} Format \textsf{ALU\_DWRR1.M} (Section~\ref{sec:format-ALU-DWRR1.M}), \verb|exucode2 = 1011|\\

\bitfields{ 1/P, 2/00, 2/-- --, 27/upper27, 1/1, 2/11, 1/0, 4/1011, 6/registerW, 2/10, 3/101, 1/0, 1/splat32, 5/lower5, 6/registerZ }


\paragraph{Syntax} \texttt{sbfusd} \emph{registerW} = \emph{registerZ}, \emph{upper27\_\discretionary{}{}{}lower5}\emph{splat32}

\paragraph{Description} The \emph{registerZ} is subtracted from the \emph{upper27\_\discretionary{}{}{}lower5} with 64-bit unsigned saturation.


\paragraph{Modifier(s)}
\noindent\begin{itemize}
\item splat32 (cf. splat32 in section \ref{par:modifier-splat32} on page \pageref{par:modifier-splat32})
\end{itemize}

\paragraph{Execution} {Reservation table \textsf{ALU\_TINY.X} (Section~\ref{sec:reservation-ALU-TINY.X})}
~
{\begin{lstlisting}
stage ID:
  new splat32 = _ZX_1(splat32);
stage RR:
  new argument3 = splat32 ? (_WX_32(upper27_lower5) << 32) | _ZX_32(_WX_32(upper27_lower5)) : _SX_32(_WX_32(upper27_lower5));
  new argument2 = GPR[registerZ];
  new result1 = _SATU_64(_ZX_64(argument3) - _ZX_64(argument2));
stage E1:
  GPR[registerW] = result1;
\end{lstlisting}}
\newpage
\subsubsection[{\small SBFUSHQ: Subtract Unsigned Saturated Half Word Quadruple}]{SBFUSHQ\hfill Subtract Unsigned Saturated Half Word Quadruple}
\label{instr:SBFUSHQ}\index{sbfushq}
 
\paragraph{Encoding} Format \textsf{ALU\_HQWRR1E} (Section~\ref{sec:format-ALU-HQWRR1E}), \verb|exucode2 = 1011|\\

\bitfields{ 1/P, 2/11, 1/1, 4/1011, 5/registerWe, 1/0, 2/10, 3/101, 1/1, 5/registerYe, 1/--, 5/registerZe, 1/0 }


\paragraph{Syntax} \texttt{sbfushq} \emph{registerWe} = \emph{registerZe}, \emph{registerYe}

\paragraph{Description} The \emph{registerYe} and the \emph{registerZe} are interpreted as four half words packed into 64 bits. The \emph{registerZe} half words are subtracted with unsigned saturation to 16 bits from the \emph{registerYe} half words. The four resulting half words are packed and stored into the \emph{registerWe}.


\paragraph{Execution} {Reservation table \textsf{ALU\_LITE} (Section~\ref{sec:reservation-ALU-LITE})}
~
{\begin{lstlisting}
stage RR:
  new argument3 = GPR[registerYe<<1];
  new argument2 = GPR[registerZe<<1];
  for (I = 0; I < 4; I += 1) {
    result1.16[I] = _SATU_16(_ZX_16(argument3.16[I]) - _ZX_16(argument2.16[I]))
  };
stage E1:
  GPR[registerWe<<1] = result1;
\end{lstlisting}}
\newpage

\paragraph{Encoding} Format \textsf{ALU\_HQWRR1O} (Section~\ref{sec:format-ALU-HQWRR1O}), \verb|exucode2 = 1011|\\

\bitfields{ 1/P, 2/11, 1/1, 4/1011, 5/registerWo, 1/1, 2/10, 3/101, 1/1, 5/registerYo, 1/--, 5/registerZo, 1/0 }


\paragraph{Syntax} \texttt{sbfushq} \emph{registerWo} = \emph{registerZo}, \emph{registerYo}

\paragraph{Description} The \emph{registerYo} and the \emph{registerZo} are interpreted as four half words packed into 64 bits. The \emph{registerZo} half words are subtracted with unsigned saturation to 16 bits from the \emph{registerYo} half words. The four resulting half words are packed and stored into the \emph{registerWo}.


\paragraph{Execution} {Reservation table \textsf{ALU\_LITE} (Section~\ref{sec:reservation-ALU-LITE})}
~
{\begin{lstlisting}
stage RR:
  new argument3 = GPR[(registerYo<<1)+1];
  new argument2 = GPR[(registerZo<<1)+1];
  for (I = 0; I < 4; I += 1) {
    result1.16[I] = _SATU_16(_ZX_16(argument3.16[I]) - _ZX_16(argument2.16[I]))
  };
stage E1:
  GPR[(registerWo<<1)+1] = result1;
\end{lstlisting}}
\newpage

\paragraph{Encoding} Format \textsf{ALU\_HQWRR1E.M} (Section~\ref{sec:format-ALU-HQWRR1E.M}), \verb|exucode2 = 1011|\\

\bitfields{ 1/P, 2/00, 2/-- --, 27/upper27, 1/1, 2/11, 1/1, 4/1011, 5/registerWe, 1/0, 2/10, 3/101, 1/1, 1/splat32, 5/lower5, 5/registerZe, 1/0 }


\paragraph{Syntax} \texttt{sbfushq} \emph{registerWe} = \emph{registerZe}, \emph{upper27\_\discretionary{}{}{}lower5}\emph{splat32}

\paragraph{Description} The \emph{upper27\_\discretionary{}{}{}lower5} and the \emph{registerZe} are interpreted as four half words packed into 64 bits. The \emph{registerZe} half words are subtracted with unsigned saturation to 16 bits from the \emph{upper27\_\discretionary{}{}{}lower5} half words. The four resulting half words are packed and stored into the \emph{registerWe}.


\paragraph{Modifier(s)}
\noindent\begin{itemize}
\item splat32 (cf. splat32 in section \ref{par:modifier-splat32} on page \pageref{par:modifier-splat32})
\end{itemize}

\paragraph{Execution} {Reservation table \textsf{ALU\_LITE.X} (Section~\ref{sec:reservation-ALU-LITE.X})}
~
{\begin{lstlisting}
stage ID:
  new splat32 = _ZX_1(splat32);
stage RR:
  new argument3 = splat32 ? (_WX_32(upper27_lower5) << 32) | _ZX_32(_WX_32(upper27_lower5)) : _SX_32(_WX_32(upper27_lower5));
  new argument2 = GPR[registerZe<<1];
  for (I = 0; I < 4; I += 1) {
    result1.16[I] = _SATU_16(_ZX_16(argument3.16[I]) - _ZX_16(argument2.16[I]))
  };
stage E1:
  GPR[registerWe<<1] = result1;
\end{lstlisting}}
\newpage

\paragraph{Encoding} Format \textsf{ALU\_HQWRR1O.M} (Section~\ref{sec:format-ALU-HQWRR1O.M}), \verb|exucode2 = 1011|\\

\bitfields{ 1/P, 2/00, 2/-- --, 27/upper27, 1/1, 2/11, 1/1, 4/1011, 5/registerWo, 1/1, 2/10, 3/101, 1/1, 1/splat32, 5/lower5, 5/registerZo, 1/0 }


\paragraph{Syntax} \texttt{sbfushq} \emph{registerWo} = \emph{registerZo}, \emph{upper27\_\discretionary{}{}{}lower5}\emph{splat32}

\paragraph{Description} The \emph{upper27\_\discretionary{}{}{}lower5} and the \emph{registerZo} are interpreted as four half words packed into 64 bits. The \emph{registerZo} half words are subtracted with unsigned saturation to 16 bits from the \emph{upper27\_\discretionary{}{}{}lower5} half words. The four resulting half words are packed and stored into the \emph{registerWo}.


\paragraph{Modifier(s)}
\noindent\begin{itemize}
\item splat32 (cf. splat32 in section \ref{par:modifier-splat32} on page \pageref{par:modifier-splat32})
\end{itemize}

\paragraph{Execution} {Reservation table \textsf{ALU\_LITE.X} (Section~\ref{sec:reservation-ALU-LITE.X})}
~
{\begin{lstlisting}
stage ID:
  new splat32 = _ZX_1(splat32);
stage RR:
  new argument3 = splat32 ? (_WX_32(upper27_lower5) << 32) | _ZX_32(_WX_32(upper27_lower5)) : _SX_32(_WX_32(upper27_lower5));
  new argument2 = GPR[(registerZo<<1)+1];
  for (I = 0; I < 4; I += 1) {
    result1.16[I] = _SATU_16(_ZX_16(argument3.16[I]) - _ZX_16(argument2.16[I]))
  };
stage E1:
  GPR[(registerWo<<1)+1] = result1;
\end{lstlisting}}
\newpage
\subsubsection[{\small SBFUSW: Subtract Unsigned Saturated Word}]{SBFUSW\hfill Subtract Unsigned Saturated Word}
\label{instr:SBFUSW}\index{sbfusw}
 
\paragraph{Encoding} Format \textsf{ALU\_WWRR1} (Section~\ref{sec:format-ALU-WWRR1}), \verb|exucode2 = 1011|\\

\bitfields{ 1/P, 2/11, 1/0, 4/1011, 6/registerW, 2/11, 3/101, 1/signextw, 6/registerY, 6/registerZ }


\paragraph{Syntax} \texttt{sbfusw}\emph{signextw} \emph{registerW} = \emph{registerZ}, \emph{registerY}

\paragraph{Description} The \emph{registerZ} is subtracted from the \emph{registerY} with 32-bit unsigned saturation. The result with the 32 upper bits cleared is stored into the \emph{registerW}.


\paragraph{Modifier(s)}
\noindent\begin{itemize}
\item signextw (cf. signextw in section \ref{par:modifier-signextw} on page \pageref{par:modifier-signextw})
\end{itemize}

\paragraph{Execution} {Reservation table \textsf{ALU\_TINY} (Section~\ref{sec:reservation-ALU-TINY})}
~
{\begin{lstlisting}
stage ID:
  new signextw = _ZX_1(signextw);
stage RR:
  new argument3 = GPR[registerY];
  new argument2 = GPR[registerZ];
  new result1 = _SATU_32(_ZX_32(argument3) - _ZX_32(argument2));
stage E1:
  GPR[registerW] = signextw ? _SX_32(result1) : _ZX_32(result1);
\end{lstlisting}}
\newpage

\paragraph{Encoding} Format \textsf{ALU\_WWRR1.W} (Section~\ref{sec:format-ALU-WWRR1.W}), \verb|exucode2 = 1011|\\

\bitfields{ 1/P, 2/00, 2/-- --, 27/upper27, 1/1, 2/11, 1/0, 4/1011, 6/registerW, 2/11, 3/101, 1/signextw, 1/0, 5/lower5, 6/registerZ }


\paragraph{Syntax} \texttt{sbfusw}\emph{signextw} \emph{registerW} = \emph{registerZ}, \emph{upper27\_\discretionary{}{}{}lower5}

\paragraph{Description} The \emph{registerZ} is subtracted from the \emph{upper27\_\discretionary{}{}{}lower5} with 32-bit unsigned saturation. The result with the 32 upper bits cleared is stored into the \emph{registerW}.


\paragraph{Modifier(s)}
\noindent\begin{itemize}
\item signextw (cf. signextw in section \ref{par:modifier-signextw} on page \pageref{par:modifier-signextw})
\end{itemize}

\paragraph{Execution} {Reservation table \textsf{ALU\_TINY.X} (Section~\ref{sec:reservation-ALU-TINY.X})}
~
{\begin{lstlisting}
stage ID:
  new signextw = _ZX_1(signextw);
stage RR:
  new argument3 = _WX_32(upper27_lower5);
  new argument2 = GPR[registerZ];
  new result1 = _SATU_32(_ZX_32(argument3) - _ZX_32(argument2));
stage E1:
  GPR[registerW] = signextw ? _SX_32(result1) : _ZX_32(result1);
\end{lstlisting}}
\newpage
\subsubsection[{\small SBFUSWP: Subtract Unsigned Saturated Word Pair}]{SBFUSWP\hfill Subtract Unsigned Saturated Word Pair}
\label{instr:SBFUSWP}\index{sbfuswp}
 
\paragraph{Encoding} Format \textsf{ALU\_WPWRR1E} (Section~\ref{sec:format-ALU-WPWRR1E}), \verb|exucode2 = 1011|\\

\bitfields{ 1/P, 2/11, 1/1, 4/1011, 5/registerWe, 1/0, 2/10, 3/101, 1/0, 5/registerYe, 1/--, 5/registerZe, 1/1 }


\paragraph{Syntax} \texttt{sbfuswp} \emph{registerWe} = \emph{registerZe}, \emph{registerYe}

\paragraph{Description} The \emph{registerYe} and the \emph{registerZe} are interpreted as two words packed into 64 bits. The \emph{registerZe} words are subtracted with unsigned saturation to 32 bits from the \emph{registerYe} words. The two resulting words are packed and stored into the \emph{registerWe}.


\paragraph{Execution} {Reservation table \textsf{ALU\_LITE} (Section~\ref{sec:reservation-ALU-LITE})}
~
{\begin{lstlisting}
stage RR:
  new argument3 = GPR[registerYe<<1];
  new argument2 = GPR[registerZe<<1];
  for (I = 0; I < 2; I += 1) {
    result1.32[I] = _SATU_32(_ZX_32(argument3.32[I]) - _ZX_32(argument2.32[I]))
  };
stage E1:
  GPR[registerWe<<1] = result1;
\end{lstlisting}}
\newpage

\paragraph{Encoding} Format \textsf{ALU\_WPWRR1O} (Section~\ref{sec:format-ALU-WPWRR1O}), \verb|exucode2 = 1011|\\

\bitfields{ 1/P, 2/11, 1/1, 4/1011, 5/registerWo, 1/1, 2/10, 3/101, 1/0, 5/registerYo, 1/--, 5/registerZo, 1/1 }


\paragraph{Syntax} \texttt{sbfuswp} \emph{registerWo} = \emph{registerZo}, \emph{registerYo}

\paragraph{Description} The \emph{registerYo} and the \emph{registerZo} are interpreted as two words packed into 64 bits. The \emph{registerZo} words are subtracted with unsigned saturation to 32 bits from the \emph{registerYo} words. The two resulting words are packed and stored into the \emph{registerWo}.


\paragraph{Execution} {Reservation table \textsf{ALU\_LITE} (Section~\ref{sec:reservation-ALU-LITE})}
~
{\begin{lstlisting}
stage RR:
  new argument3 = GPR[(registerYo<<1)+1];
  new argument2 = GPR[(registerZo<<1)+1];
  for (I = 0; I < 2; I += 1) {
    result1.32[I] = _SATU_32(_ZX_32(argument3.32[I]) - _ZX_32(argument2.32[I]))
  };
stage E1:
  GPR[(registerWo<<1)+1] = result1;
\end{lstlisting}}
\newpage

\paragraph{Encoding} Format \textsf{ALU\_WPWRR1E.M} (Section~\ref{sec:format-ALU-WPWRR1E.M}), \verb|exucode2 = 1011|\\

\bitfields{ 1/P, 2/00, 2/-- --, 27/upper27, 1/1, 2/11, 1/1, 4/1011, 5/registerWe, 1/0, 2/10, 3/101, 1/0, 1/splat32, 5/lower5, 5/registerZe, 1/1 }


\paragraph{Syntax} \texttt{sbfuswp} \emph{registerWe} = \emph{registerZe}, \emph{upper27\_\discretionary{}{}{}lower5}\emph{splat32}

\paragraph{Description} The \emph{upper27\_\discretionary{}{}{}lower5} and the \emph{registerZe} are interpreted as two words packed into 64 bits. The \emph{registerZe} words are subtracted with unsigned saturation to 32 bits from the \emph{upper27\_\discretionary{}{}{}lower5} words. The two resulting words are packed and stored into the \emph{registerWe}.


\paragraph{Modifier(s)}
\noindent\begin{itemize}
\item splat32 (cf. splat32 in section \ref{par:modifier-splat32} on page \pageref{par:modifier-splat32})
\end{itemize}

\paragraph{Execution} {Reservation table \textsf{ALU\_LITE.X} (Section~\ref{sec:reservation-ALU-LITE.X})}
~
{\begin{lstlisting}
stage ID:
  new splat32 = _ZX_1(splat32);
stage RR:
  new argument3 = splat32 ? (_WX_32(upper27_lower5) << 32) | _ZX_32(_WX_32(upper27_lower5)) : _SX_32(_WX_32(upper27_lower5));
  new argument2 = GPR[registerZe<<1];
  for (I = 0; I < 2; I += 1) {
    result1.32[I] = _SATU_32(_ZX_32(argument3.32[I]) - _ZX_32(argument2.32[I]))
  };
stage E1:
  GPR[registerWe<<1] = result1;
\end{lstlisting}}
\newpage

\paragraph{Encoding} Format \textsf{ALU\_WPWRR1O.M} (Section~\ref{sec:format-ALU-WPWRR1O.M}), \verb|exucode2 = 1011|\\

\bitfields{ 1/P, 2/00, 2/-- --, 27/upper27, 1/1, 2/11, 1/1, 4/1011, 5/registerWo, 1/1, 2/10, 3/101, 1/0, 1/splat32, 5/lower5, 5/registerZo, 1/1 }


\paragraph{Syntax} \texttt{sbfuswp} \emph{registerWo} = \emph{registerZo}, \emph{upper27\_\discretionary{}{}{}lower5}\emph{splat32}

\paragraph{Description} The \emph{upper27\_\discretionary{}{}{}lower5} and the \emph{registerZo} are interpreted as two words packed into 64 bits. The \emph{registerZo} words are subtracted with unsigned saturation to 32 bits from the \emph{upper27\_\discretionary{}{}{}lower5} words. The two resulting words are packed and stored into the \emph{registerWo}.


\paragraph{Modifier(s)}
\noindent\begin{itemize}
\item splat32 (cf. splat32 in section \ref{par:modifier-splat32} on page \pageref{par:modifier-splat32})
\end{itemize}

\paragraph{Execution} {Reservation table \textsf{ALU\_LITE.X} (Section~\ref{sec:reservation-ALU-LITE.X})}
~
{\begin{lstlisting}
stage ID:
  new splat32 = _ZX_1(splat32);
stage RR:
  new argument3 = splat32 ? (_WX_32(upper27_lower5) << 32) | _ZX_32(_WX_32(upper27_lower5)) : _SX_32(_WX_32(upper27_lower5));
  new argument2 = GPR[(registerZo<<1)+1];
  for (I = 0; I < 2; I += 1) {
    result1.32[I] = _SATU_32(_ZX_32(argument3.32[I]) - _ZX_32(argument2.32[I]))
  };
stage E1:
  GPR[(registerWo<<1)+1] = result1;
\end{lstlisting}}
\newpage
\subsubsection[{\small SBFW: Subtract From Word}]{SBFW\hfill Subtract From Word}
\label{instr:SBFW}\index{sbfw}
 
\paragraph{Encoding} Format \textsf{ALU\_WWRR0} (Section~\ref{sec:format-ALU-WWRR0}), \verb|exucode2 = 0011|\\

\bitfields{ 1/P, 2/11, 1/0, 4/0011, 6/registerW, 2/11, 3/000, 1/signextw, 6/registerY, 6/registerZ }


\paragraph{Syntax} \texttt{sbfw}\emph{signextw} \emph{registerW} = \emph{registerZ}, \emph{registerY}

\paragraph{Description} The \emph{registerZ} is subtracted from the \emph{registerY}. The result with the 32 upper bits cleared is stored into the \emph{registerW}.


\paragraph{Modifier(s)}
\noindent\begin{itemize}
\item signextw (cf. signextw in section \ref{par:modifier-signextw} on page \pageref{par:modifier-signextw})
\end{itemize}

\paragraph{Execution} {Reservation table \textsf{ALU\_TINY} (Section~\ref{sec:reservation-ALU-TINY})}
~
{\begin{lstlisting}
stage ID:
  new signextw = _ZX_1(signextw);
stage RR:
  new argument3 = GPR[registerY];
  new argument2 = GPR[registerZ];
  new result1 = argument3 - argument2;
stage E1:
  GPR[registerW] = signextw ? _SX_32(result1) : _ZX_32(result1);
\end{lstlisting}}
\newpage

\paragraph{Encoding} Format \textsf{ALU\_WWRR0.W} (Section~\ref{sec:format-ALU-WWRR0.W}), \verb|exucode2 = 0011|\\

\bitfields{ 1/P, 2/00, 2/-- --, 27/upper27, 1/1, 2/11, 1/0, 4/0011, 6/registerW, 2/11, 3/000, 1/signextw, 1/0, 5/lower5, 6/registerZ }


\paragraph{Syntax} \texttt{sbfw}\emph{signextw} \emph{registerW} = \emph{registerZ}, \emph{upper27\_\discretionary{}{}{}lower5}

\paragraph{Description} The \emph{registerZ} is subtracted from the \emph{upper27\_\discretionary{}{}{}lower5}. The result with the 32 upper bits cleared is stored into the \emph{registerW}.


\paragraph{Modifier(s)}
\noindent\begin{itemize}
\item signextw (cf. signextw in section \ref{par:modifier-signextw} on page \pageref{par:modifier-signextw})
\end{itemize}

\paragraph{Execution} {Reservation table \textsf{ALU\_TINY.X} (Section~\ref{sec:reservation-ALU-TINY.X})}
~
{\begin{lstlisting}
stage ID:
  new signextw = _ZX_1(signextw);
stage RR:
  new argument3 = _WX_32(upper27_lower5);
  new argument2 = GPR[registerZ];
  new result1 = argument3 - argument2;
stage E1:
  GPR[registerW] = signextw ? _SX_32(result1) : _ZX_32(result1);
\end{lstlisting}}
\newpage
\subsubsection[{\small SBFWP: Subtract From Word Pair}]{SBFWP\hfill Subtract From Word Pair}
\label{instr:SBFWP}\index{sbfwp}
 
\paragraph{Encoding} Format \textsf{ALU\_WPWRR0E} (Section~\ref{sec:format-ALU-WPWRR0E}), \verb|exucode2 = 0011|\\

\bitfields{ 1/P, 2/11, 1/1, 4/0011, 5/registerWe, 1/0, 2/10, 3/000, 1/0, 5/registerYe, 1/--, 5/registerZe, 1/1 }


\paragraph{Syntax} \texttt{sbfwp} \emph{registerWe} = \emph{registerZe}, \emph{registerYe}

\paragraph{Description} The \emph{registerYe} and the \emph{registerZe} are interpreted as two words packed into 64 bits. The \emph{registerZe} words are subtracted from the \emph{registerYe} words. The two resulting words are packed and stored into the \emph{registerWe}.


\paragraph{Execution} {Reservation table \textsf{ALU\_LITE} (Section~\ref{sec:reservation-ALU-LITE})}
~
{\begin{lstlisting}
stage RR:
  new argument3 = GPR[registerYe<<1];
  new argument2 = GPR[registerZe<<1];
  for (I = 0; I < 2; I += 1) {
    result1.32[I] = argument3.32[I] - argument2.32[I]
  };
stage E1:
  GPR[registerWe<<1] = result1;
\end{lstlisting}}
\newpage

\paragraph{Encoding} Format \textsf{ALU\_WPWRR0O} (Section~\ref{sec:format-ALU-WPWRR0O}), \verb|exucode2 = 0011|\\

\bitfields{ 1/P, 2/11, 1/1, 4/0011, 5/registerWo, 1/1, 2/10, 3/000, 1/0, 5/registerYo, 1/--, 5/registerZo, 1/1 }


\paragraph{Syntax} \texttt{sbfwp} \emph{registerWo} = \emph{registerZo}, \emph{registerYo}

\paragraph{Description} The \emph{registerYo} and the \emph{registerZo} are interpreted as two words packed into 64 bits. The \emph{registerZo} words are subtracted from the \emph{registerYo} words. The two resulting words are packed and stored into the \emph{registerWo}.


\paragraph{Execution} {Reservation table \textsf{ALU\_LITE} (Section~\ref{sec:reservation-ALU-LITE})}
~
{\begin{lstlisting}
stage RR:
  new argument3 = GPR[(registerYo<<1)+1];
  new argument2 = GPR[(registerZo<<1)+1];
  for (I = 0; I < 2; I += 1) {
    result1.32[I] = argument3.32[I] - argument2.32[I]
  };
stage E1:
  GPR[(registerWo<<1)+1] = result1;
\end{lstlisting}}
\newpage

\paragraph{Encoding} Format \textsf{ALU\_WPWRR0E.M} (Section~\ref{sec:format-ALU-WPWRR0E.M}), \verb|exucode2 = 0011|\\

\bitfields{ 1/P, 2/00, 2/-- --, 27/upper27, 1/1, 2/11, 1/1, 4/0011, 5/registerWe, 1/0, 2/10, 3/000, 1/0, 1/splat32, 5/lower5, 5/registerZe, 1/1 }


\paragraph{Syntax} \texttt{sbfwp} \emph{registerWe} = \emph{registerZe}, \emph{upper27\_\discretionary{}{}{}lower5}\emph{splat32}

\paragraph{Description} The \emph{upper27\_\discretionary{}{}{}lower5} and the \emph{registerZe} are interpreted as two words packed into 64 bits. The \emph{registerZe} words are subtracted from the \emph{upper27\_\discretionary{}{}{}lower5} words. The two resulting words are packed and stored into the \emph{registerWe}.


\paragraph{Modifier(s)}
\noindent\begin{itemize}
\item splat32 (cf. splat32 in section \ref{par:modifier-splat32} on page \pageref{par:modifier-splat32})
\end{itemize}

\paragraph{Execution} {Reservation table \textsf{ALU\_LITE.X} (Section~\ref{sec:reservation-ALU-LITE.X})}
~
{\begin{lstlisting}
stage ID:
  new splat32 = _ZX_1(splat32);
stage RR:
  new argument3 = splat32 ? (_WX_32(upper27_lower5) << 32) | _ZX_32(_WX_32(upper27_lower5)) : _SX_32(_WX_32(upper27_lower5));
  new argument2 = GPR[registerZe<<1];
  for (I = 0; I < 2; I += 1) {
    result1.32[I] = argument3.32[I] - argument2.32[I]
  };
stage E1:
  GPR[registerWe<<1] = result1;
\end{lstlisting}}
\newpage

\paragraph{Encoding} Format \textsf{ALU\_WPWRR0O.M} (Section~\ref{sec:format-ALU-WPWRR0O.M}), \verb|exucode2 = 0011|\\

\bitfields{ 1/P, 2/00, 2/-- --, 27/upper27, 1/1, 2/11, 1/1, 4/0011, 5/registerWo, 1/1, 2/10, 3/000, 1/0, 1/splat32, 5/lower5, 5/registerZo, 1/1 }


\paragraph{Syntax} \texttt{sbfwp} \emph{registerWo} = \emph{registerZo}, \emph{upper27\_\discretionary{}{}{}lower5}\emph{splat32}

\paragraph{Description} The \emph{upper27\_\discretionary{}{}{}lower5} and the \emph{registerZo} are interpreted as two words packed into 64 bits. The \emph{registerZo} words are subtracted from the \emph{upper27\_\discretionary{}{}{}lower5} words. The two resulting words are packed and stored into the \emph{registerWo}.


\paragraph{Modifier(s)}
\noindent\begin{itemize}
\item splat32 (cf. splat32 in section \ref{par:modifier-splat32} on page \pageref{par:modifier-splat32})
\end{itemize}

\paragraph{Execution} {Reservation table \textsf{ALU\_LITE.X} (Section~\ref{sec:reservation-ALU-LITE.X})}
~
{\begin{lstlisting}
stage ID:
  new splat32 = _ZX_1(splat32);
stage RR:
  new argument3 = splat32 ? (_WX_32(upper27_lower5) << 32) | _ZX_32(_WX_32(upper27_lower5)) : _SX_32(_WX_32(upper27_lower5));
  new argument2 = GPR[(registerZo<<1)+1];
  for (I = 0; I < 2; I += 1) {
    result1.32[I] = argument3.32[I] - argument2.32[I]
  };
stage E1:
  GPR[(registerWo<<1)+1] = result1;
\end{lstlisting}}
\newpage
\subsubsection[{\small SBMM8D: Swapped Bit Matrix Multiplication 8\texorpdfstring{$\times$}{x}8 Double Word}]{SBMM8D\hfill Swapped Bit Matrix Multiplication 8$\times$8 Double Word}
\label{instr:SBMM8D}\index{sbmm8d}
 
\paragraph{Encoding} Format \textsf{ALU\_DBMWRR} (Section~\ref{sec:format-ALU-DBMWRR}), \verb|exucode2 = 1110|\\

\bitfields{ 1/P, 2/11, 1/0, 4/1110, 6/registerW, 2/10, 3/001, 1/0, 6/registerY, 6/registerZ }


\paragraph{Syntax} \texttt{sbmm8d} \emph{registerW} = \emph{registerZ}, \emph{registerY}

\paragraph{Description} The \emph{registerY} interpreted as a 8$\times$8 bit matrix is multiplied by the \emph{registerZ} interpreted as a 8$\times$8 bit matrix. The resulting 8$\times$8 bit matrix is stored into the \emph{registerW}.


\paragraph{Execution} {Reservation table \textsf{ALU\_TINY} (Section~\ref{sec:reservation-ALU-TINY})}
~
{\begin{lstlisting}
stage RR:
  new argument3 = GPR[registerY];
  new argument2 = GPR[registerZ];
  new result1 = _BMM_8(_ZX_64(argument3), _ZX_64(argument2));
stage E1:
  GPR[registerW] = result1;
\end{lstlisting}}
\newpage

\paragraph{Encoding} Format \textsf{ALU\_DBMWRR.M} (Section~\ref{sec:format-ALU-DBMWRR.M}), \verb|exucode2 = 1110|\\

\bitfields{ 1/P, 2/00, 2/-- --, 27/upper27, 1/1, 2/11, 1/0, 4/1110, 6/registerW, 2/10, 3/001, 1/0, 1/splat32, 5/lower5, 6/registerZ }


\paragraph{Syntax} \texttt{sbmm8d} \emph{registerW} = \emph{registerZ}, \emph{upper27\_\discretionary{}{}{}lower5}\emph{splat32}

\paragraph{Description} The \emph{upper27\_\discretionary{}{}{}lower5} interpreted as a 8$\times$8 bit matrix is multiplied by the \emph{registerZ} interpreted as a 8$\times$8 bit matrix. The resulting 8$\times$8 bit matrix is stored into the \emph{registerW}.


\paragraph{Modifier(s)}
\noindent\begin{itemize}
\item splat32 (cf. splat32 in section \ref{par:modifier-splat32} on page \pageref{par:modifier-splat32})
\end{itemize}

\paragraph{Execution} {Reservation table \textsf{ALU\_TINY.X} (Section~\ref{sec:reservation-ALU-TINY.X})}
~
{\begin{lstlisting}
stage ID:
  new splat32 = _ZX_1(splat32);
stage RR:
  new argument3 = splat32 ? (_WX_32(upper27_lower5) << 32) | _ZX_32(_WX_32(upper27_lower5)) : _SX_32(_WX_32(upper27_lower5));
  new argument2 = GPR[registerZ];
  new result1 = _BMM_8(_ZX_64(argument3), _ZX_64(argument2));
stage E1:
  GPR[registerW] = result1;
\end{lstlisting}}
\newpage

\paragraph{Encoding} Format \textsf{ALU\_DBMWRI} (Section~\ref{sec:format-ALU-DBMWRI}), \verb|exucode2 = 1110|\\

\bitfields{ 1/P, 2/11, 1/0, 4/1110, 6/registerW, 2/01, 10/signed10, 6/registerZ }


\paragraph{Syntax} \texttt{sbmm8d} \emph{registerW} = \emph{registerZ}, \emph{signed10}

\paragraph{Description} The \emph{signed10} interpreted as a 8$\times$8 bit matrix is multiplied by the \emph{registerZ} interpreted as a 8$\times$8 bit matrix. The resulting 8$\times$8 bit matrix is stored into the \emph{registerW}.


\paragraph{Execution} {Reservation table \textsf{ALU\_TINY} (Section~\ref{sec:reservation-ALU-TINY})}
~
{\begin{lstlisting}
stage RR:
  new argument3 = _SX_10(signed10);
  new argument2 = GPR[registerZ];
  new result1 = _BMM_8(_ZX_64(argument3), _ZX_64(argument2));
stage E1:
  GPR[registerW] = result1;
\end{lstlisting}}
\newpage

\paragraph{Encoding} Format \textsf{ALU\_DBMWRI.X} (Section~\ref{sec:format-ALU-DBMWRI.X}), \verb|exucode2 = 1110|\\

\bitfields{ 1/P, 2/00, 2/-- --, 27/upper27, 1/1, 2/11, 1/0, 4/1110, 6/registerW, 2/01, 10/lower10, 6/registerZ }


\paragraph{Syntax} \texttt{sbmm8d} \emph{registerW} = \emph{registerZ}, \emph{upper27\_\discretionary{}{}{}lower10}

\paragraph{Description} The \emph{upper27\_\discretionary{}{}{}lower10} interpreted as a 8$\times$8 bit matrix is multiplied by the \emph{registerZ} interpreted as a 8$\times$8 bit matrix. The resulting 8$\times$8 bit matrix is stored into the \emph{registerW}.


\paragraph{Execution} {Reservation table \textsf{ALU\_TINY.X} (Section~\ref{sec:reservation-ALU-TINY.X})}
~
{\begin{lstlisting}
stage RR:
  new argument3 = _SX_37(upper27_lower10);
  new argument2 = GPR[registerZ];
  new result1 = _BMM_8(_ZX_64(argument3), _ZX_64(argument2));
stage E1:
  GPR[registerW] = result1;
\end{lstlisting}}
\newpage

\paragraph{Encoding} Format \textsf{ALU\_DBMWRI.Y} (Section~\ref{sec:format-ALU-DBMWRI.Y}), \verb|exucode2 = 1110|\\

\bitfields{ 1/P, 2/00, 2/-- --, 27/extend27, 1/1, 2/00, 2/-- --, 27/upper27, 1/1, 2/11, 1/0, 4/1110, 6/registerW, 2/01, 10/lower10, 6/registerZ }


\paragraph{Syntax} \texttt{sbmm8d} \emph{registerW} = \emph{registerZ}, \emph{extend27\_\discretionary{}{}{}upper27\_\discretionary{}{}{}lower10}

\paragraph{Description} The \emph{extend27\_\discretionary{}{}{}upper27\_\discretionary{}{}{}lower10} interpreted as a 8$\times$8 bit matrix is multiplied by the \emph{registerZ} interpreted as a 8$\times$8 bit matrix. The resulting 8$\times$8 bit matrix is stored into the \emph{registerW}.


\paragraph{Execution} {Reservation table \textsf{ALU\_TINY.Y} (Section~\ref{sec:reservation-ALU-TINY.Y})}
~
{\begin{lstlisting}
stage RR:
  new argument3 = _WX_64(extend27_upper27_lower10);
  new argument2 = GPR[registerZ];
  new result1 = _BMM_8(_ZX_64(argument3), _ZX_64(argument2));
stage E1:
  GPR[registerW] = result1;
\end{lstlisting}}
\newpage
\subsubsection[{\small SBMM8EORD: Swapped Bit Matrix Multiplication 8\texorpdfstring{$\times$}{x}8 Exclusive OR Double Word}]{SBMM8EORD\hfill Swapped Bit Matrix Multiplication 8$\times$8 Exclusive OR Double Word}
\label{instr:SBMM8EORD}\index{sbmm8eord}
 
\paragraph{Encoding} Format \textsf{ALU\_DBMWRRR} (Section~\ref{sec:format-ALU-DBMWRRR}), \verb|exucode2 = 1100|\\

\bitfields{ 1/P, 2/11, 1/0, 4/1100, 6/registerW, 2/10, 3/001, 1/0, 6/registerY, 6/registerZ }


\paragraph{Syntax} \texttt{sbmm8eord} \emph{registerW} = \emph{registerZ}, \emph{registerY}

\paragraph{Description} The \emph{registerY} interpreted as a 8$\times$8 bit matrix is multiplied by the \emph{registerZ} interpreted as a 8$\times$8 bit matrix. The resulting $\times$8 bit matrix is XORed with the value of the \emph{registerW} and stored into the \emph{registerW}.


\paragraph{Execution} {Reservation table \textsf{ALU\_TINY} (Section~\ref{sec:reservation-ALU-TINY})}
~
{\begin{lstlisting}
stage RR:
  new argument3 = GPR[registerY];
  new argument2 = GPR[registerZ];
  new argument1 = GPR[registerW];
  new result1 = _BMM_8(_ZX_64(argument3), _ZX_64(argument2)) ^ _ZX_64(argument1);
stage E1:
  GPR[registerW] = result1;
\end{lstlisting}}
\newpage

\paragraph{Encoding} Format \textsf{ALU\_DBMWRRR.M} (Section~\ref{sec:format-ALU-DBMWRRR.M}), \verb|exucode2 = 1100|\\

\bitfields{ 1/P, 2/00, 2/-- --, 27/upper27, 1/1, 2/11, 1/0, 4/1100, 6/registerW, 2/10, 3/001, 1/0, 1/splat32, 5/lower5, 6/registerZ }


\paragraph{Syntax} \texttt{sbmm8eord} \emph{registerW} = \emph{registerZ}, \emph{upper27\_\discretionary{}{}{}lower5}\emph{splat32}

\paragraph{Description} The \emph{upper27\_\discretionary{}{}{}lower5} interpreted as a 8$\times$8 bit matrix is multiplied by the \emph{registerZ} interpreted as a 8$\times$8 bit matrix. The resulting $\times$8 bit matrix is XORed with the value of the \emph{registerW} and stored into the \emph{registerW}.


\paragraph{Modifier(s)}
\noindent\begin{itemize}
\item splat32 (cf. splat32 in section \ref{par:modifier-splat32} on page \pageref{par:modifier-splat32})
\end{itemize}

\paragraph{Execution} {Reservation table \textsf{ALU\_TINY.X} (Section~\ref{sec:reservation-ALU-TINY.X})}
~
{\begin{lstlisting}
stage ID:
  new splat32 = _ZX_1(splat32);
stage RR:
  new argument3 = splat32 ? (_WX_32(upper27_lower5) << 32) | _ZX_32(_WX_32(upper27_lower5)) : _SX_32(_WX_32(upper27_lower5));
  new argument2 = GPR[registerZ];
  new argument1 = GPR[registerW];
  new result1 = _BMM_8(_ZX_64(argument3), _ZX_64(argument2)) ^ _ZX_64(argument1);
stage E1:
  GPR[registerW] = result1;
\end{lstlisting}}
\newpage
\subsubsection[{\small SBMMT8D: Swapped Bit Matrix Multiplication Transposed 8\texorpdfstring{$\times$}{x}8 Double Word}]{SBMMT8D\hfill Swapped Bit Matrix Multiplication Transposed 8$\times$8 Double Word}
\label{instr:SBMMT8D}\index{sbmmt8d}
 
\paragraph{Encoding} Format \textsf{ALU\_DBMWRR} (Section~\ref{sec:format-ALU-DBMWRR}), \verb|exucode2 = 1111|\\

\bitfields{ 1/P, 2/11, 1/0, 4/1111, 6/registerW, 2/10, 3/001, 1/0, 6/registerY, 6/registerZ }


\paragraph{Syntax} \texttt{sbmmt8d} \emph{registerW} = \emph{registerZ}, \emph{registerY}

\paragraph{Description} The \emph{registerY} interpreted as a 8$\times$8 bit matrix is multiplied by the \emph{registerZ} interpreted as a 8$\times$8 bit matrix. The resulting 8$\times$8 bit matrix is transposed, and stored into the \emph{registerW}.


\paragraph{Execution} {Reservation table \textsf{ALU\_TINY} (Section~\ref{sec:reservation-ALU-TINY})}
~
{\begin{lstlisting}
stage RR:
  new argument3 = GPR[registerY];
  new argument2 = GPR[registerZ];
  new result1 = _BMT_8(_BMM_8(_ZX_64(argument3), _ZX_64(argument2)));
stage E1:
  GPR[registerW] = result1;
\end{lstlisting}}
\newpage

\paragraph{Encoding} Format \textsf{ALU\_DBMWRR.M} (Section~\ref{sec:format-ALU-DBMWRR.M}), \verb|exucode2 = 1111|\\

\bitfields{ 1/P, 2/00, 2/-- --, 27/upper27, 1/1, 2/11, 1/0, 4/1111, 6/registerW, 2/10, 3/001, 1/0, 1/splat32, 5/lower5, 6/registerZ }


\paragraph{Syntax} \texttt{sbmmt8d} \emph{registerW} = \emph{registerZ}, \emph{upper27\_\discretionary{}{}{}lower5}\emph{splat32}

\paragraph{Description} The \emph{upper27\_\discretionary{}{}{}lower5} interpreted as a 8$\times$8 bit matrix is multiplied by the \emph{registerZ} interpreted as a 8$\times$8 bit matrix. The resulting 8$\times$8 bit matrix is transposed, and stored into the \emph{registerW}.


\paragraph{Modifier(s)}
\noindent\begin{itemize}
\item splat32 (cf. splat32 in section \ref{par:modifier-splat32} on page \pageref{par:modifier-splat32})
\end{itemize}

\paragraph{Execution} {Reservation table \textsf{ALU\_TINY.X} (Section~\ref{sec:reservation-ALU-TINY.X})}
~
{\begin{lstlisting}
stage ID:
  new splat32 = _ZX_1(splat32);
stage RR:
  new argument3 = splat32 ? (_WX_32(upper27_lower5) << 32) | _ZX_32(_WX_32(upper27_lower5)) : _SX_32(_WX_32(upper27_lower5));
  new argument2 = GPR[registerZ];
  new result1 = _BMT_8(_BMM_8(_ZX_64(argument3), _ZX_64(argument2)));
stage E1:
  GPR[registerW] = result1;
\end{lstlisting}}
\newpage

\paragraph{Encoding} Format \textsf{ALU\_DBMWRI} (Section~\ref{sec:format-ALU-DBMWRI}), \verb|exucode2 = 1111|\\

\bitfields{ 1/P, 2/11, 1/0, 4/1111, 6/registerW, 2/01, 10/signed10, 6/registerZ }


\paragraph{Syntax} \texttt{sbmmt8d} \emph{registerW} = \emph{registerZ}, \emph{signed10}

\paragraph{Description} The \emph{signed10} interpreted as a 8$\times$8 bit matrix is multiplied by the \emph{registerZ} interpreted as a 8$\times$8 bit matrix. The resulting 8$\times$8 bit matrix is transposed, and stored into the \emph{registerW}.


\paragraph{Execution} {Reservation table \textsf{ALU\_TINY} (Section~\ref{sec:reservation-ALU-TINY})}
~
{\begin{lstlisting}
stage RR:
  new argument3 = _SX_10(signed10);
  new argument2 = GPR[registerZ];
  new result1 = _BMT_8(_BMM_8(_ZX_64(argument3), _ZX_64(argument2)));
stage E1:
  GPR[registerW] = result1;
\end{lstlisting}}
\newpage

\paragraph{Encoding} Format \textsf{ALU\_DBMWRI.X} (Section~\ref{sec:format-ALU-DBMWRI.X}), \verb|exucode2 = 1111|\\

\bitfields{ 1/P, 2/00, 2/-- --, 27/upper27, 1/1, 2/11, 1/0, 4/1111, 6/registerW, 2/01, 10/lower10, 6/registerZ }


\paragraph{Syntax} \texttt{sbmmt8d} \emph{registerW} = \emph{registerZ}, \emph{upper27\_\discretionary{}{}{}lower10}

\paragraph{Description} The \emph{upper27\_\discretionary{}{}{}lower10} interpreted as a 8$\times$8 bit matrix is multiplied by the \emph{registerZ} interpreted as a 8$\times$8 bit matrix. The resulting 8$\times$8 bit matrix is transposed, and stored into the \emph{registerW}.


\paragraph{Execution} {Reservation table \textsf{ALU\_TINY.X} (Section~\ref{sec:reservation-ALU-TINY.X})}
~
{\begin{lstlisting}
stage RR:
  new argument3 = _SX_37(upper27_lower10);
  new argument2 = GPR[registerZ];
  new result1 = _BMT_8(_BMM_8(_ZX_64(argument3), _ZX_64(argument2)));
stage E1:
  GPR[registerW] = result1;
\end{lstlisting}}
\newpage

\paragraph{Encoding} Format \textsf{ALU\_DBMWRI.Y} (Section~\ref{sec:format-ALU-DBMWRI.Y}), \verb|exucode2 = 1111|\\

\bitfields{ 1/P, 2/00, 2/-- --, 27/extend27, 1/1, 2/00, 2/-- --, 27/upper27, 1/1, 2/11, 1/0, 4/1111, 6/registerW, 2/01, 10/lower10, 6/registerZ }


\paragraph{Syntax} \texttt{sbmmt8d} \emph{registerW} = \emph{registerZ}, \emph{extend27\_\discretionary{}{}{}upper27\_\discretionary{}{}{}lower10}

\paragraph{Description} The \emph{extend27\_\discretionary{}{}{}upper27\_\discretionary{}{}{}lower10} interpreted as a 8$\times$8 bit matrix is multiplied by the \emph{registerZ} interpreted as a 8$\times$8 bit matrix. The resulting 8$\times$8 bit matrix is transposed, and stored into the \emph{registerW}.


\paragraph{Execution} {Reservation table \textsf{ALU\_TINY.Y} (Section~\ref{sec:reservation-ALU-TINY.Y})}
~
{\begin{lstlisting}
stage RR:
  new argument3 = _WX_64(extend27_upper27_lower10);
  new argument2 = GPR[registerZ];
  new result1 = _BMT_8(_BMM_8(_ZX_64(argument3), _ZX_64(argument2)));
stage E1:
  GPR[registerW] = result1;
\end{lstlisting}}
\newpage
\subsubsection[{\small SBMMT8EORD: Swapped Bit Matrix Multiplication Transposed 8\texorpdfstring{$\times$}{x}8 Exclusive OR Double Word}]{SBMMT8EORD\hfill Swapped Bit Matrix Multiplication Transposed 8$\times$8 Exclusive OR Double Word}
\label{instr:SBMMT8EORD}\index{sbmmt8eord}
 
\paragraph{Encoding} Format \textsf{ALU\_DBMWRRR} (Section~\ref{sec:format-ALU-DBMWRRR}), \verb|exucode2 = 1101|\\

\bitfields{ 1/P, 2/11, 1/0, 4/1101, 6/registerW, 2/10, 3/001, 1/0, 6/registerY, 6/registerZ }


\paragraph{Syntax} \texttt{sbmmt8eord} \emph{registerW} = \emph{registerZ}, \emph{registerY}

\paragraph{Description} The \emph{registerY} interpreted as a 8$\times$8 bit matrix is multiplied by the \emph{registerZ} interpreted as a 8$\times$8 bit matrix. The resulting 8$\times$8 bit matrix is transposed, XORed with the value of the \emph{registerW} and stored into the \emph{registerW}.


\paragraph{Execution} {Reservation table \textsf{ALU\_TINY} (Section~\ref{sec:reservation-ALU-TINY})}
~
{\begin{lstlisting}
stage RR:
  new argument3 = GPR[registerY];
  new argument2 = GPR[registerZ];
  new argument1 = GPR[registerW];
  new result1 = _BMT_8(_BMM_8(_ZX_64(argument3), _ZX_64(argument2))) ^ _ZX_64(argument1);
stage E1:
  GPR[registerW] = result1;
\end{lstlisting}}
\newpage

\paragraph{Encoding} Format \textsf{ALU\_DBMWRRR.M} (Section~\ref{sec:format-ALU-DBMWRRR.M}), \verb|exucode2 = 1101|\\

\bitfields{ 1/P, 2/00, 2/-- --, 27/upper27, 1/1, 2/11, 1/0, 4/1101, 6/registerW, 2/10, 3/001, 1/0, 1/splat32, 5/lower5, 6/registerZ }


\paragraph{Syntax} \texttt{sbmmt8eord} \emph{registerW} = \emph{registerZ}, \emph{upper27\_\discretionary{}{}{}lower5}\emph{splat32}

\paragraph{Description} The \emph{upper27\_\discretionary{}{}{}lower5} interpreted as a 8$\times$8 bit matrix is multiplied by the \emph{registerZ} interpreted as a 8$\times$8 bit matrix. The resulting 8$\times$8 bit matrix is transposed, XORed with the value of the \emph{registerW} and stored into the \emph{registerW}.


\paragraph{Modifier(s)}
\noindent\begin{itemize}
\item splat32 (cf. splat32 in section \ref{par:modifier-splat32} on page \pageref{par:modifier-splat32})
\end{itemize}

\paragraph{Execution} {Reservation table \textsf{ALU\_TINY.X} (Section~\ref{sec:reservation-ALU-TINY.X})}
~
{\begin{lstlisting}
stage ID:
  new splat32 = _ZX_1(splat32);
stage RR:
  new argument3 = splat32 ? (_WX_32(upper27_lower5) << 32) | _ZX_32(_WX_32(upper27_lower5)) : _SX_32(_WX_32(upper27_lower5));
  new argument2 = GPR[registerZ];
  new argument1 = GPR[registerW];
  new result1 = _BMT_8(_BMM_8(_ZX_64(argument3), _ZX_64(argument2))) ^ _ZX_64(argument1);
stage E1:
  GPR[registerW] = result1;
\end{lstlisting}}
\newpage
\subsubsection[{\small SIGNBO: Apply Sign Byte Octuple}]{SIGNBO\hfill Apply Sign Byte Octuple}
\label{instr:SIGNBO}\index{signbo}
 
\paragraph{Encoding} Format \textsf{ALU\_BOWRR0E} (Section~\ref{sec:format-ALU-BOWRR0E}), \verb|exucode2 = 0000|\\

\bitfields{ 1/P, 2/11, 1/1, 4/0000, 5/registerWe, 1/0, 2/10, 3/000, 1/1, 5/registerYe, 1/--, 5/registerZe, 1/1 }


\paragraph{Syntax} \texttt{signbo} \emph{registerWe} = \emph{registerZe}, \emph{registerYe}

\paragraph{Description} The \emph{registerYe} and the \emph{registerZe} are interpreted as eight signed bytes packed into 64 bits. The \emph{registerZe} bytes are multiplied by the signs of the \emph{registerYe} bytes. The eight resulting bytes are packed and stored into the \emph{registerWe}.


\paragraph{Execution} {Reservation table \textsf{ALU\_LITE} (Section~\ref{sec:reservation-ALU-LITE})}
~
{\begin{lstlisting}
stage RR:
  new argument3 = GPR[registerYe<<1];
  new argument2 = GPR[registerZe<<1];
  for (I = 0; I < 8; I += 1) {
    result1.8[I] = argument2.8[I]
    if (_SX_8(argument3.8[I]) < 0) {
      result1.8[I] = -argument2.8[I];
    }
    if (_ZX_8(argument3.8[I]) == 0) {
      result1.8[I] = 0;
    }
  };
stage E1:
  GPR[registerWe<<1] = result1;
\end{lstlisting}}
\newpage

\paragraph{Encoding} Format \textsf{ALU\_BOWRR0O} (Section~\ref{sec:format-ALU-BOWRR0O}), \verb|exucode2 = 0000|\\

\bitfields{ 1/P, 2/11, 1/1, 4/0000, 5/registerWo, 1/1, 2/10, 3/000, 1/1, 5/registerYo, 1/--, 5/registerZo, 1/1 }


\paragraph{Syntax} \texttt{signbo} \emph{registerWo} = \emph{registerZo}, \emph{registerYo}

\paragraph{Description} The \emph{registerYo} and the \emph{registerZo} are interpreted as eight signed bytes packed into 64 bits. The \emph{registerZo} bytes are multiplied by the signs of the \emph{registerYo} bytes. The eight resulting bytes are packed and stored into the \emph{registerWo}.


\paragraph{Execution} {Reservation table \textsf{ALU\_LITE} (Section~\ref{sec:reservation-ALU-LITE})}
~
{\begin{lstlisting}
stage RR:
  new argument3 = GPR[(registerYo<<1)+1];
  new argument2 = GPR[(registerZo<<1)+1];
  for (I = 0; I < 8; I += 1) {
    result1.8[I] = argument2.8[I]
    if (_SX_8(argument3.8[I]) < 0) {
      result1.8[I] = -argument2.8[I];
    }
    if (_ZX_8(argument3.8[I]) == 0) {
      result1.8[I] = 0;
    }
  };
stage E1:
  GPR[(registerWo<<1)+1] = result1;
\end{lstlisting}}
\newpage

\paragraph{Encoding} Format \textsf{ALU\_BOWRR0E.M} (Section~\ref{sec:format-ALU-BOWRR0E.M}), \verb|exucode2 = 0000|\\

\bitfields{ 1/P, 2/00, 2/-- --, 27/upper27, 1/1, 2/11, 1/1, 4/0000, 5/registerWe, 1/0, 2/10, 3/000, 1/1, 1/splat32, 5/lower5, 5/registerZe, 1/1 }


\paragraph{Syntax} \texttt{signbo} \emph{registerWe} = \emph{registerZe}, \emph{upper27\_\discretionary{}{}{}lower5}\emph{splat32}

\paragraph{Description} The \emph{upper27\_\discretionary{}{}{}lower5} and the \emph{registerZe} are interpreted as eight signed bytes packed into 64 bits. The \emph{registerZe} bytes are multiplied by the signs of the \emph{upper27\_\discretionary{}{}{}lower5} bytes. The eight resulting bytes are packed and stored into the \emph{registerWe}.


\paragraph{Modifier(s)}
\noindent\begin{itemize}
\item splat32 (cf. splat32 in section \ref{par:modifier-splat32} on page \pageref{par:modifier-splat32})
\end{itemize}

\paragraph{Execution} {Reservation table \textsf{ALU\_LITE.X} (Section~\ref{sec:reservation-ALU-LITE.X})}
~
{\begin{lstlisting}
stage ID:
  new splat32 = _ZX_1(splat32);
stage RR:
  new argument3 = splat32 ? (_WX_32(upper27_lower5) << 32) | _ZX_32(_WX_32(upper27_lower5)) : _SX_32(_WX_32(upper27_lower5));
  new argument2 = GPR[registerZe<<1];
  for (I = 0; I < 8; I += 1) {
    result1.8[I] = argument2.8[I]
    if (_SX_8(argument3.8[I]) < 0) {
      result1.8[I] = -argument2.8[I];
    }
    if (_ZX_8(argument3.8[I]) == 0) {
      result1.8[I] = 0;
    }
  };
stage E1:
  GPR[registerWe<<1] = result1;
\end{lstlisting}}
\newpage

\paragraph{Encoding} Format \textsf{ALU\_BOWRR0O.M} (Section~\ref{sec:format-ALU-BOWRR0O.M}), \verb|exucode2 = 0000|\\

\bitfields{ 1/P, 2/00, 2/-- --, 27/upper27, 1/1, 2/11, 1/1, 4/0000, 5/registerWo, 1/1, 2/10, 3/000, 1/1, 1/splat32, 5/lower5, 5/registerZo, 1/1 }


\paragraph{Syntax} \texttt{signbo} \emph{registerWo} = \emph{registerZo}, \emph{upper27\_\discretionary{}{}{}lower5}\emph{splat32}

\paragraph{Description} The \emph{upper27\_\discretionary{}{}{}lower5} and the \emph{registerZo} are interpreted as eight signed bytes packed into 64 bits. The \emph{registerZo} bytes are multiplied by the signs of the \emph{upper27\_\discretionary{}{}{}lower5} bytes. The eight resulting bytes are packed and stored into the \emph{registerWo}.


\paragraph{Modifier(s)}
\noindent\begin{itemize}
\item splat32 (cf. splat32 in section \ref{par:modifier-splat32} on page \pageref{par:modifier-splat32})
\end{itemize}

\paragraph{Execution} {Reservation table \textsf{ALU\_LITE.X} (Section~\ref{sec:reservation-ALU-LITE.X})}
~
{\begin{lstlisting}
stage ID:
  new splat32 = _ZX_1(splat32);
stage RR:
  new argument3 = splat32 ? (_WX_32(upper27_lower5) << 32) | _ZX_32(_WX_32(upper27_lower5)) : _SX_32(_WX_32(upper27_lower5));
  new argument2 = GPR[(registerZo<<1)+1];
  for (I = 0; I < 8; I += 1) {
    result1.8[I] = argument2.8[I]
    if (_SX_8(argument3.8[I]) < 0) {
      result1.8[I] = -argument2.8[I];
    }
    if (_ZX_8(argument3.8[I]) == 0) {
      result1.8[I] = 0;
    }
  };
stage E1:
  GPR[(registerWo<<1)+1] = result1;
\end{lstlisting}}
\newpage
\subsubsection[{\small SIGND: Apply Sign Double Word}]{SIGND\hfill Apply Sign Double Word}
\label{instr:SIGND}\index{signd}
 
\paragraph{Encoding} Format \textsf{ALU\_DWRR0} (Section~\ref{sec:format-ALU-DWRR0}), \verb|exucode2 = 0000|\\

\bitfields{ 1/P, 2/11, 1/0, 4/0000, 6/registerW, 2/10, 3/000, 1/0, 6/registerY, 6/registerZ }


\paragraph{Syntax} \texttt{signd} \emph{registerW} = \emph{registerZ}, \emph{registerY}

\paragraph{Description} The \emph{registerY} and the \emph{registerZ} are interpreted as signed integer double words. The \emph{registerZ} double word is multiplied by the sign of the \emph{registerY} double word. The resulting double word is stored into the \emph{registerW}.


\paragraph{Execution} {Reservation table \textsf{ALU\_TINY} (Section~\ref{sec:reservation-ALU-TINY})}
~
{\begin{lstlisting}
stage RR:
  new argument3 = GPR[registerY];
  new argument2 = GPR[registerZ];
  new result1 = argument2;
  if (_SX_64(argument3) < 0) {
    result1 = -argument2;
  }
  if (_ZX_64(argument3) == 0) {
    result1 = 0;
  }
stage E1:
  GPR[registerW] = result1;
\end{lstlisting}}
\newpage

\paragraph{Encoding} Format \textsf{ALU\_DWRR0.M} (Section~\ref{sec:format-ALU-DWRR0.M}), \verb|exucode2 = 0000|\\

\bitfields{ 1/P, 2/00, 2/-- --, 27/upper27, 1/1, 2/11, 1/0, 4/0000, 6/registerW, 2/10, 3/000, 1/0, 1/splat32, 5/lower5, 6/registerZ }


\paragraph{Syntax} \texttt{signd} \emph{registerW} = \emph{registerZ}, \emph{upper27\_\discretionary{}{}{}lower5}\emph{splat32}

\paragraph{Description} The \emph{upper27\_\discretionary{}{}{}lower5} and the \emph{registerZ} are interpreted as signed integer double words. The \emph{registerZ} double word is multiplied by the sign of the \emph{upper27\_\discretionary{}{}{}lower5} double word. The resulting double word is stored into the \emph{registerW}.


\paragraph{Modifier(s)}
\noindent\begin{itemize}
\item splat32 (cf. splat32 in section \ref{par:modifier-splat32} on page \pageref{par:modifier-splat32})
\end{itemize}

\paragraph{Execution} {Reservation table \textsf{ALU\_TINY.X} (Section~\ref{sec:reservation-ALU-TINY.X})}
~
{\begin{lstlisting}
stage ID:
  new splat32 = _ZX_1(splat32);
stage RR:
  new argument3 = splat32 ? (_WX_32(upper27_lower5) << 32) | _ZX_32(_WX_32(upper27_lower5)) : _SX_32(_WX_32(upper27_lower5));
  new argument2 = GPR[registerZ];
  new result1 = argument2;
  if (_SX_64(argument3) < 0) {
    result1 = -argument2;
  }
  if (_ZX_64(argument3) == 0) {
    result1 = 0;
  }
stage E1:
  GPR[registerW] = result1;
\end{lstlisting}}
\newpage
\subsubsection[{\small SIGNHQ: Apply Sign Half Word Quadruple}]{SIGNHQ\hfill Apply Sign Half Word Quadruple}
\label{instr:SIGNHQ}\index{signhq}
 
\paragraph{Encoding} Format \textsf{ALU\_HQWRR0E} (Section~\ref{sec:format-ALU-HQWRR0E}), \verb|exucode2 = 0000|\\

\bitfields{ 1/P, 2/11, 1/1, 4/0000, 5/registerWe, 1/0, 2/10, 3/000, 1/1, 5/registerYe, 1/--, 5/registerZe, 1/0 }


\paragraph{Syntax} \texttt{signhq} \emph{registerWe} = \emph{registerZe}, \emph{registerYe}

\paragraph{Description} The \emph{registerYe} and the \emph{registerZe} are interpreted as four signed integer half words packed into 64 bits. The \emph{registerZe} half words are multiplied by the signs of the \emph{registerYe} half words. The four resulting half words are packed and stored into the \emph{registerWe}.


\paragraph{Execution} {Reservation table \textsf{ALU\_LITE} (Section~\ref{sec:reservation-ALU-LITE})}
~
{\begin{lstlisting}
stage RR:
  new argument3 = GPR[registerYe<<1];
  new argument2 = GPR[registerZe<<1];
  for (I = 0; I < 4; I += 1) {
    result1.16[I] = argument2.16[I]
    if (_SX_16(argument3.16[I]) < 0) {
      result1.16[I] = -argument2.16[I];
    }
    if (_ZX_16(argument3.16[I]) == 0) {
      result1.16[I] = 0;
    }
  };
stage E1:
  GPR[registerWe<<1] = result1;
\end{lstlisting}}
\newpage

\paragraph{Encoding} Format \textsf{ALU\_HQWRR0O} (Section~\ref{sec:format-ALU-HQWRR0O}), \verb|exucode2 = 0000|\\

\bitfields{ 1/P, 2/11, 1/1, 4/0000, 5/registerWo, 1/1, 2/10, 3/000, 1/1, 5/registerYo, 1/--, 5/registerZo, 1/0 }


\paragraph{Syntax} \texttt{signhq} \emph{registerWo} = \emph{registerZo}, \emph{registerYo}

\paragraph{Description} The \emph{registerYo} and the \emph{registerZo} are interpreted as four signed integer half words packed into 64 bits. The \emph{registerZo} half words are multiplied by the signs of the \emph{registerYo} half words. The four resulting half words are packed and stored into the \emph{registerWo}.


\paragraph{Execution} {Reservation table \textsf{ALU\_LITE} (Section~\ref{sec:reservation-ALU-LITE})}
~
{\begin{lstlisting}
stage RR:
  new argument3 = GPR[(registerYo<<1)+1];
  new argument2 = GPR[(registerZo<<1)+1];
  for (I = 0; I < 4; I += 1) {
    result1.16[I] = argument2.16[I]
    if (_SX_16(argument3.16[I]) < 0) {
      result1.16[I] = -argument2.16[I];
    }
    if (_ZX_16(argument3.16[I]) == 0) {
      result1.16[I] = 0;
    }
  };
stage E1:
  GPR[(registerWo<<1)+1] = result1;
\end{lstlisting}}
\newpage

\paragraph{Encoding} Format \textsf{ALU\_HQWRR0E.M} (Section~\ref{sec:format-ALU-HQWRR0E.M}), \verb|exucode2 = 0000|\\

\bitfields{ 1/P, 2/00, 2/-- --, 27/upper27, 1/1, 2/11, 1/1, 4/0000, 5/registerWe, 1/0, 2/10, 3/000, 1/1, 1/splat32, 5/lower5, 5/registerZe, 1/0 }


\paragraph{Syntax} \texttt{signhq} \emph{registerWe} = \emph{registerZe}, \emph{upper27\_\discretionary{}{}{}lower5}\emph{splat32}

\paragraph{Description} The \emph{upper27\_\discretionary{}{}{}lower5} and the \emph{registerZe} are interpreted as four signed integer half words packed into 64 bits. The \emph{registerZe} half words are multiplied by the signs of the \emph{upper27\_\discretionary{}{}{}lower5} half words. The four resulting half words are packed and stored into the \emph{registerWe}.


\paragraph{Modifier(s)}
\noindent\begin{itemize}
\item splat32 (cf. splat32 in section \ref{par:modifier-splat32} on page \pageref{par:modifier-splat32})
\end{itemize}

\paragraph{Execution} {Reservation table \textsf{ALU\_LITE.X} (Section~\ref{sec:reservation-ALU-LITE.X})}
~
{\begin{lstlisting}
stage ID:
  new splat32 = _ZX_1(splat32);
stage RR:
  new argument3 = splat32 ? (_WX_32(upper27_lower5) << 32) | _ZX_32(_WX_32(upper27_lower5)) : _SX_32(_WX_32(upper27_lower5));
  new argument2 = GPR[registerZe<<1];
  for (I = 0; I < 4; I += 1) {
    result1.16[I] = argument2.16[I]
    if (_SX_16(argument3.16[I]) < 0) {
      result1.16[I] = -argument2.16[I];
    }
    if (_ZX_16(argument3.16[I]) == 0) {
      result1.16[I] = 0;
    }
  };
stage E1:
  GPR[registerWe<<1] = result1;
\end{lstlisting}}
\newpage

\paragraph{Encoding} Format \textsf{ALU\_HQWRR0O.M} (Section~\ref{sec:format-ALU-HQWRR0O.M}), \verb|exucode2 = 0000|\\

\bitfields{ 1/P, 2/00, 2/-- --, 27/upper27, 1/1, 2/11, 1/1, 4/0000, 5/registerWo, 1/1, 2/10, 3/000, 1/1, 1/splat32, 5/lower5, 5/registerZo, 1/0 }


\paragraph{Syntax} \texttt{signhq} \emph{registerWo} = \emph{registerZo}, \emph{upper27\_\discretionary{}{}{}lower5}\emph{splat32}

\paragraph{Description} The \emph{upper27\_\discretionary{}{}{}lower5} and the \emph{registerZo} are interpreted as four signed integer half words packed into 64 bits. The \emph{registerZo} half words are multiplied by the signs of the \emph{upper27\_\discretionary{}{}{}lower5} half words. The four resulting half words are packed and stored into the \emph{registerWo}.


\paragraph{Modifier(s)}
\noindent\begin{itemize}
\item splat32 (cf. splat32 in section \ref{par:modifier-splat32} on page \pageref{par:modifier-splat32})
\end{itemize}

\paragraph{Execution} {Reservation table \textsf{ALU\_LITE.X} (Section~\ref{sec:reservation-ALU-LITE.X})}
~
{\begin{lstlisting}
stage ID:
  new splat32 = _ZX_1(splat32);
stage RR:
  new argument3 = splat32 ? (_WX_32(upper27_lower5) << 32) | _ZX_32(_WX_32(upper27_lower5)) : _SX_32(_WX_32(upper27_lower5));
  new argument2 = GPR[(registerZo<<1)+1];
  for (I = 0; I < 4; I += 1) {
    result1.16[I] = argument2.16[I]
    if (_SX_16(argument3.16[I]) < 0) {
      result1.16[I] = -argument2.16[I];
    }
    if (_ZX_16(argument3.16[I]) == 0) {
      result1.16[I] = 0;
    }
  };
stage E1:
  GPR[(registerWo<<1)+1] = result1;
\end{lstlisting}}
\newpage
\subsubsection[{\small SIGNSBO: Apply Sign Saturated Byte Octuple}]{SIGNSBO\hfill Apply Sign Saturated Byte Octuple}
\label{instr:SIGNSBO}\index{signsbo}
 
\paragraph{Encoding} Format \textsf{ALU\_BOWRR0E} (Section~\ref{sec:format-ALU-BOWRR0E}), \verb|exucode2 = 0001|\\

\bitfields{ 1/P, 2/11, 1/1, 4/0001, 5/registerWe, 1/0, 2/10, 3/000, 1/1, 5/registerYe, 1/--, 5/registerZe, 1/1 }


\paragraph{Syntax} \texttt{signsbo} \emph{registerWe} = \emph{registerZe}, \emph{registerYe}

\paragraph{Description} The \emph{registerYe} and the \emph{registerZe} are interpreted as eight signed bytes packed into 64 bits. The \emph{registerZe} bytes are multiplied by the signs of the \emph{registerYe} bytes and 8-bit saturated. The 8 resulting bytes are packed and stored into the \emph{registerWe}.


\paragraph{Execution} {Reservation table \textsf{ALU\_LITE} (Section~\ref{sec:reservation-ALU-LITE})}
~
{\begin{lstlisting}
stage RR:
  new argument3 = GPR[registerYe<<1];
  new argument2 = GPR[registerZe<<1];
  for (I = 0; I < 8; I += 1) {
    result1.8[I] = argument2.8[I]
    if (_SX_8(argument3.8[I]) < 0) {
      result1.8[I] = _SAT_8(-argument2.8[I]);
    }
    if (_ZX_8(argument3.8[I]) == 0) {
      result1.8[I] = 0;
    }
  };
stage E1:
  GPR[registerWe<<1] = result1;
\end{lstlisting}}
\newpage

\paragraph{Encoding} Format \textsf{ALU\_BOWRR0O} (Section~\ref{sec:format-ALU-BOWRR0O}), \verb|exucode2 = 0001|\\

\bitfields{ 1/P, 2/11, 1/1, 4/0001, 5/registerWo, 1/1, 2/10, 3/000, 1/1, 5/registerYo, 1/--, 5/registerZo, 1/1 }


\paragraph{Syntax} \texttt{signsbo} \emph{registerWo} = \emph{registerZo}, \emph{registerYo}

\paragraph{Description} The \emph{registerYo} and the \emph{registerZo} are interpreted as eight signed bytes packed into 64 bits. The \emph{registerZo} bytes are multiplied by the signs of the \emph{registerYo} bytes and 8-bit saturated. The 8 resulting bytes are packed and stored into the \emph{registerWo}.


\paragraph{Execution} {Reservation table \textsf{ALU\_LITE} (Section~\ref{sec:reservation-ALU-LITE})}
~
{\begin{lstlisting}
stage RR:
  new argument3 = GPR[(registerYo<<1)+1];
  new argument2 = GPR[(registerZo<<1)+1];
  for (I = 0; I < 8; I += 1) {
    result1.8[I] = argument2.8[I]
    if (_SX_8(argument3.8[I]) < 0) {
      result1.8[I] = _SAT_8(-argument2.8[I]);
    }
    if (_ZX_8(argument3.8[I]) == 0) {
      result1.8[I] = 0;
    }
  };
stage E1:
  GPR[(registerWo<<1)+1] = result1;
\end{lstlisting}}
\newpage

\paragraph{Encoding} Format \textsf{ALU\_BOWRR0E.M} (Section~\ref{sec:format-ALU-BOWRR0E.M}), \verb|exucode2 = 0001|\\

\bitfields{ 1/P, 2/00, 2/-- --, 27/upper27, 1/1, 2/11, 1/1, 4/0001, 5/registerWe, 1/0, 2/10, 3/000, 1/1, 1/splat32, 5/lower5, 5/registerZe, 1/1 }


\paragraph{Syntax} \texttt{signsbo} \emph{registerWe} = \emph{registerZe}, \emph{upper27\_\discretionary{}{}{}lower5}\emph{splat32}

\paragraph{Description} The \emph{upper27\_\discretionary{}{}{}lower5} and the \emph{registerZe} are interpreted as eight signed bytes packed into 64 bits. The \emph{registerZe} bytes are multiplied by the signs of the \emph{upper27\_\discretionary{}{}{}lower5} bytes and 8-bit saturated. The 8 resulting bytes are packed and stored into the \emph{registerWe}.


\paragraph{Modifier(s)}
\noindent\begin{itemize}
\item splat32 (cf. splat32 in section \ref{par:modifier-splat32} on page \pageref{par:modifier-splat32})
\end{itemize}

\paragraph{Execution} {Reservation table \textsf{ALU\_LITE.X} (Section~\ref{sec:reservation-ALU-LITE.X})}
~
{\begin{lstlisting}
stage ID:
  new splat32 = _ZX_1(splat32);
stage RR:
  new argument3 = splat32 ? (_WX_32(upper27_lower5) << 32) | _ZX_32(_WX_32(upper27_lower5)) : _SX_32(_WX_32(upper27_lower5));
  new argument2 = GPR[registerZe<<1];
  for (I = 0; I < 8; I += 1) {
    result1.8[I] = argument2.8[I]
    if (_SX_8(argument3.8[I]) < 0) {
      result1.8[I] = _SAT_8(-argument2.8[I]);
    }
    if (_ZX_8(argument3.8[I]) == 0) {
      result1.8[I] = 0;
    }
  };
stage E1:
  GPR[registerWe<<1] = result1;
\end{lstlisting}}
\newpage

\paragraph{Encoding} Format \textsf{ALU\_BOWRR0O.M} (Section~\ref{sec:format-ALU-BOWRR0O.M}), \verb|exucode2 = 0001|\\

\bitfields{ 1/P, 2/00, 2/-- --, 27/upper27, 1/1, 2/11, 1/1, 4/0001, 5/registerWo, 1/1, 2/10, 3/000, 1/1, 1/splat32, 5/lower5, 5/registerZo, 1/1 }


\paragraph{Syntax} \texttt{signsbo} \emph{registerWo} = \emph{registerZo}, \emph{upper27\_\discretionary{}{}{}lower5}\emph{splat32}

\paragraph{Description} The \emph{upper27\_\discretionary{}{}{}lower5} and the \emph{registerZo} are interpreted as eight signed bytes packed into 64 bits. The \emph{registerZo} bytes are multiplied by the signs of the \emph{upper27\_\discretionary{}{}{}lower5} bytes and 8-bit saturated. The 8 resulting bytes are packed and stored into the \emph{registerWo}.


\paragraph{Modifier(s)}
\noindent\begin{itemize}
\item splat32 (cf. splat32 in section \ref{par:modifier-splat32} on page \pageref{par:modifier-splat32})
\end{itemize}

\paragraph{Execution} {Reservation table \textsf{ALU\_LITE.X} (Section~\ref{sec:reservation-ALU-LITE.X})}
~
{\begin{lstlisting}
stage ID:
  new splat32 = _ZX_1(splat32);
stage RR:
  new argument3 = splat32 ? (_WX_32(upper27_lower5) << 32) | _ZX_32(_WX_32(upper27_lower5)) : _SX_32(_WX_32(upper27_lower5));
  new argument2 = GPR[(registerZo<<1)+1];
  for (I = 0; I < 8; I += 1) {
    result1.8[I] = argument2.8[I]
    if (_SX_8(argument3.8[I]) < 0) {
      result1.8[I] = _SAT_8(-argument2.8[I]);
    }
    if (_ZX_8(argument3.8[I]) == 0) {
      result1.8[I] = 0;
    }
  };
stage E1:
  GPR[(registerWo<<1)+1] = result1;
\end{lstlisting}}
\newpage
\subsubsection[{\small SIGNSD: Apply Sign Saturated Double Word}]{SIGNSD\hfill Apply Sign Saturated Double Word}
\label{instr:SIGNSD}\index{signsd}
 
\paragraph{Encoding} Format \textsf{ALU\_DWRR0} (Section~\ref{sec:format-ALU-DWRR0}), \verb|exucode2 = 0001|\\

\bitfields{ 1/P, 2/11, 1/0, 4/0001, 6/registerW, 2/10, 3/000, 1/0, 6/registerY, 6/registerZ }


\paragraph{Syntax} \texttt{signsd} \emph{registerW} = \emph{registerZ}, \emph{registerY}

\paragraph{Description} The \emph{registerY} and the \emph{registerZ} are interpreted as signed integer double words. The \emph{registerZ} double words is multiplied by the sign of the \emph{registerY} double words and 64-bit saturated. The resulting double word is stored into the \emph{registerW}.


\paragraph{Execution} {Reservation table \textsf{ALU\_TINY} (Section~\ref{sec:reservation-ALU-TINY})}
~
{\begin{lstlisting}
stage RR:
  new argument3 = GPR[registerY];
  new argument2 = GPR[registerZ];
  new result1 = argument2;
  if (_SX_64(argument3) < 0) {
    result1 = _SAT_64(-argument2);
  }
  if (_ZX_64(argument3) == 0) {
    result1 = 0;
  }
stage E1:
  GPR[registerW] = result1;
\end{lstlisting}}
\newpage

\paragraph{Encoding} Format \textsf{ALU\_DWRR0.M} (Section~\ref{sec:format-ALU-DWRR0.M}), \verb|exucode2 = 0001|\\

\bitfields{ 1/P, 2/00, 2/-- --, 27/upper27, 1/1, 2/11, 1/0, 4/0001, 6/registerW, 2/10, 3/000, 1/0, 1/splat32, 5/lower5, 6/registerZ }


\paragraph{Syntax} \texttt{signsd} \emph{registerW} = \emph{registerZ}, \emph{upper27\_\discretionary{}{}{}lower5}\emph{splat32}

\paragraph{Description} The \emph{upper27\_\discretionary{}{}{}lower5} and the \emph{registerZ} are interpreted as signed integer double words. The \emph{registerZ} double words is multiplied by the sign of the \emph{upper27\_\discretionary{}{}{}lower5} double words and 64-bit saturated. The resulting double word is stored into the \emph{registerW}.


\paragraph{Modifier(s)}
\noindent\begin{itemize}
\item splat32 (cf. splat32 in section \ref{par:modifier-splat32} on page \pageref{par:modifier-splat32})
\end{itemize}

\paragraph{Execution} {Reservation table \textsf{ALU\_TINY.X} (Section~\ref{sec:reservation-ALU-TINY.X})}
~
{\begin{lstlisting}
stage ID:
  new splat32 = _ZX_1(splat32);
stage RR:
  new argument3 = splat32 ? (_WX_32(upper27_lower5) << 32) | _ZX_32(_WX_32(upper27_lower5)) : _SX_32(_WX_32(upper27_lower5));
  new argument2 = GPR[registerZ];
  new result1 = argument2;
  if (_SX_64(argument3) < 0) {
    result1 = _SAT_64(-argument2);
  }
  if (_ZX_64(argument3) == 0) {
    result1 = 0;
  }
stage E1:
  GPR[registerW] = result1;
\end{lstlisting}}
\newpage
\subsubsection[{\small SIGNSHQ: Apply Sign Saturated Half Word Quadruple}]{SIGNSHQ\hfill Apply Sign Saturated Half Word Quadruple}
\label{instr:SIGNSHQ}\index{signshq}
 
\paragraph{Encoding} Format \textsf{ALU\_HQWRR0E} (Section~\ref{sec:format-ALU-HQWRR0E}), \verb|exucode2 = 0001|\\

\bitfields{ 1/P, 2/11, 1/1, 4/0001, 5/registerWe, 1/0, 2/10, 3/000, 1/1, 5/registerYe, 1/--, 5/registerZe, 1/0 }


\paragraph{Syntax} \texttt{signshq} \emph{registerWe} = \emph{registerZe}, \emph{registerYe}

\paragraph{Description} The \emph{registerYe} and the \emph{registerZe} are interpreted as four signed integer half words packed into 64 bits. The \emph{registerZe} half words are multiplied by the signs of the \emph{registerYe} half words and 16-bit saturated. The four resulting half words are packed and stored into the \emph{registerWe}.


\paragraph{Execution} {Reservation table \textsf{ALU\_LITE} (Section~\ref{sec:reservation-ALU-LITE})}
~
{\begin{lstlisting}
stage RR:
  new argument3 = GPR[registerYe<<1];
  new argument2 = GPR[registerZe<<1];
  for (I = 0; I < 4; I += 1) {
    result1.16[I] = argument2.16[I]
    if (_SX_16(argument3.16[I]) < 0) {
      result1.16[I] = _SAT_16(-argument2.16[I]);
    }
    if (_ZX_16(argument3.16[I]) == 0) {
      result1.16[I] = 0;
    }
  };
stage E1:
  GPR[registerWe<<1] = result1;
\end{lstlisting}}
\newpage

\paragraph{Encoding} Format \textsf{ALU\_HQWRR0O} (Section~\ref{sec:format-ALU-HQWRR0O}), \verb|exucode2 = 0001|\\

\bitfields{ 1/P, 2/11, 1/1, 4/0001, 5/registerWo, 1/1, 2/10, 3/000, 1/1, 5/registerYo, 1/--, 5/registerZo, 1/0 }


\paragraph{Syntax} \texttt{signshq} \emph{registerWo} = \emph{registerZo}, \emph{registerYo}

\paragraph{Description} The \emph{registerYo} and the \emph{registerZo} are interpreted as four signed integer half words packed into 64 bits. The \emph{registerZo} half words are multiplied by the signs of the \emph{registerYo} half words and 16-bit saturated. The four resulting half words are packed and stored into the \emph{registerWo}.


\paragraph{Execution} {Reservation table \textsf{ALU\_LITE} (Section~\ref{sec:reservation-ALU-LITE})}
~
{\begin{lstlisting}
stage RR:
  new argument3 = GPR[(registerYo<<1)+1];
  new argument2 = GPR[(registerZo<<1)+1];
  for (I = 0; I < 4; I += 1) {
    result1.16[I] = argument2.16[I]
    if (_SX_16(argument3.16[I]) < 0) {
      result1.16[I] = _SAT_16(-argument2.16[I]);
    }
    if (_ZX_16(argument3.16[I]) == 0) {
      result1.16[I] = 0;
    }
  };
stage E1:
  GPR[(registerWo<<1)+1] = result1;
\end{lstlisting}}
\newpage

\paragraph{Encoding} Format \textsf{ALU\_HQWRR0E.M} (Section~\ref{sec:format-ALU-HQWRR0E.M}), \verb|exucode2 = 0001|\\

\bitfields{ 1/P, 2/00, 2/-- --, 27/upper27, 1/1, 2/11, 1/1, 4/0001, 5/registerWe, 1/0, 2/10, 3/000, 1/1, 1/splat32, 5/lower5, 5/registerZe, 1/0 }


\paragraph{Syntax} \texttt{signshq} \emph{registerWe} = \emph{registerZe}, \emph{upper27\_\discretionary{}{}{}lower5}\emph{splat32}

\paragraph{Description} The \emph{upper27\_\discretionary{}{}{}lower5} and the \emph{registerZe} are interpreted as four signed integer half words packed into 64 bits. The \emph{registerZe} half words are multiplied by the signs of the \emph{upper27\_\discretionary{}{}{}lower5} half words and 16-bit saturated. The four resulting half words are packed and stored into the \emph{registerWe}.


\paragraph{Modifier(s)}
\noindent\begin{itemize}
\item splat32 (cf. splat32 in section \ref{par:modifier-splat32} on page \pageref{par:modifier-splat32})
\end{itemize}

\paragraph{Execution} {Reservation table \textsf{ALU\_LITE.X} (Section~\ref{sec:reservation-ALU-LITE.X})}
~
{\begin{lstlisting}
stage ID:
  new splat32 = _ZX_1(splat32);
stage RR:
  new argument3 = splat32 ? (_WX_32(upper27_lower5) << 32) | _ZX_32(_WX_32(upper27_lower5)) : _SX_32(_WX_32(upper27_lower5));
  new argument2 = GPR[registerZe<<1];
  for (I = 0; I < 4; I += 1) {
    result1.16[I] = argument2.16[I]
    if (_SX_16(argument3.16[I]) < 0) {
      result1.16[I] = _SAT_16(-argument2.16[I]);
    }
    if (_ZX_16(argument3.16[I]) == 0) {
      result1.16[I] = 0;
    }
  };
stage E1:
  GPR[registerWe<<1] = result1;
\end{lstlisting}}
\newpage

\paragraph{Encoding} Format \textsf{ALU\_HQWRR0O.M} (Section~\ref{sec:format-ALU-HQWRR0O.M}), \verb|exucode2 = 0001|\\

\bitfields{ 1/P, 2/00, 2/-- --, 27/upper27, 1/1, 2/11, 1/1, 4/0001, 5/registerWo, 1/1, 2/10, 3/000, 1/1, 1/splat32, 5/lower5, 5/registerZo, 1/0 }


\paragraph{Syntax} \texttt{signshq} \emph{registerWo} = \emph{registerZo}, \emph{upper27\_\discretionary{}{}{}lower5}\emph{splat32}

\paragraph{Description} The \emph{upper27\_\discretionary{}{}{}lower5} and the \emph{registerZo} are interpreted as four signed integer half words packed into 64 bits. The \emph{registerZo} half words are multiplied by the signs of the \emph{upper27\_\discretionary{}{}{}lower5} half words and 16-bit saturated. The four resulting half words are packed and stored into the \emph{registerWo}.


\paragraph{Modifier(s)}
\noindent\begin{itemize}
\item splat32 (cf. splat32 in section \ref{par:modifier-splat32} on page \pageref{par:modifier-splat32})
\end{itemize}

\paragraph{Execution} {Reservation table \textsf{ALU\_LITE.X} (Section~\ref{sec:reservation-ALU-LITE.X})}
~
{\begin{lstlisting}
stage ID:
  new splat32 = _ZX_1(splat32);
stage RR:
  new argument3 = splat32 ? (_WX_32(upper27_lower5) << 32) | _ZX_32(_WX_32(upper27_lower5)) : _SX_32(_WX_32(upper27_lower5));
  new argument2 = GPR[(registerZo<<1)+1];
  for (I = 0; I < 4; I += 1) {
    result1.16[I] = argument2.16[I]
    if (_SX_16(argument3.16[I]) < 0) {
      result1.16[I] = _SAT_16(-argument2.16[I]);
    }
    if (_ZX_16(argument3.16[I]) == 0) {
      result1.16[I] = 0;
    }
  };
stage E1:
  GPR[(registerWo<<1)+1] = result1;
\end{lstlisting}}
\newpage
\subsubsection[{\small SIGNSW: Apply Sign Saturated Word}]{SIGNSW\hfill Apply Sign Saturated Word}
\label{instr:SIGNSW}\index{signsw}
 
\paragraph{Encoding} Format \textsf{ALU\_WWRR0} (Section~\ref{sec:format-ALU-WWRR0}), \verb|exucode2 = 0001|\\

\bitfields{ 1/P, 2/11, 1/0, 4/0001, 6/registerW, 2/11, 3/000, 1/signextw, 6/registerY, 6/registerZ }


\paragraph{Syntax} \texttt{signsw}\emph{signextw} \emph{registerW} = \emph{registerZ}, \emph{registerY}

\paragraph{Description} The \emph{registerY} and the \emph{registerZ} are interpreted as integer signed words. The \emph{registerZ} word is multiplied by the sign of the \emph{registerY} word and 32-bit saturated. The resulting word is stored into the \emph{registerW}.


\paragraph{Modifier(s)}
\noindent\begin{itemize}
\item signextw (cf. signextw in section \ref{par:modifier-signextw} on page \pageref{par:modifier-signextw})
\end{itemize}

\paragraph{Execution} {Reservation table \textsf{ALU\_TINY} (Section~\ref{sec:reservation-ALU-TINY})}
~
{\begin{lstlisting}
stage ID:
  new signextw = _ZX_1(signextw);
stage RR:
  new argument3 = GPR[registerY];
  new argument2 = GPR[registerZ];
  new result1 = argument2;
  if (_SX_32(argument3) < 0) {
    result1 = _SAT_32(-argument2);
  }
  if (_ZX_32(argument3) == 0) {
    result1 = 0;
  }
stage E1:
  GPR[registerW] = signextw ? _SX_32(result1) : _ZX_32(result1);
\end{lstlisting}}
\newpage

\paragraph{Encoding} Format \textsf{ALU\_WWRR0.W} (Section~\ref{sec:format-ALU-WWRR0.W}), \verb|exucode2 = 0001|\\

\bitfields{ 1/P, 2/00, 2/-- --, 27/upper27, 1/1, 2/11, 1/0, 4/0001, 6/registerW, 2/11, 3/000, 1/signextw, 1/0, 5/lower5, 6/registerZ }


\paragraph{Syntax} \texttt{signsw}\emph{signextw} \emph{registerW} = \emph{registerZ}, \emph{upper27\_\discretionary{}{}{}lower5}

\paragraph{Description} The \emph{upper27\_\discretionary{}{}{}lower5} and the \emph{registerZ} are interpreted as integer signed words. The \emph{registerZ} word is multiplied by the sign of the \emph{upper27\_\discretionary{}{}{}lower5} word and 32-bit saturated. The resulting word is stored into the \emph{registerW}.


\paragraph{Modifier(s)}
\noindent\begin{itemize}
\item signextw (cf. signextw in section \ref{par:modifier-signextw} on page \pageref{par:modifier-signextw})
\end{itemize}

\paragraph{Execution} {Reservation table \textsf{ALU\_TINY.X} (Section~\ref{sec:reservation-ALU-TINY.X})}
~
{\begin{lstlisting}
stage ID:
  new signextw = _ZX_1(signextw);
stage RR:
  new argument3 = _WX_32(upper27_lower5);
  new argument2 = GPR[registerZ];
  new result1 = argument2;
  if (_SX_32(argument3) < 0) {
    result1 = _SAT_32(-argument2);
  }
  if (_ZX_32(argument3) == 0) {
    result1 = 0;
  }
stage E1:
  GPR[registerW] = signextw ? _SX_32(result1) : _ZX_32(result1);
\end{lstlisting}}
\newpage
\subsubsection[{\small SIGNSWP: Apply Sign Saturated Word Pair}]{SIGNSWP\hfill Apply Sign Saturated Word Pair}
\label{instr:SIGNSWP}\index{signswp}
 
\paragraph{Encoding} Format \textsf{ALU\_WPWRR0E} (Section~\ref{sec:format-ALU-WPWRR0E}), \verb|exucode2 = 0001|\\

\bitfields{ 1/P, 2/11, 1/1, 4/0001, 5/registerWe, 1/0, 2/10, 3/000, 1/0, 5/registerYe, 1/--, 5/registerZe, 1/1 }


\paragraph{Syntax} \texttt{signswp} \emph{registerWe} = \emph{registerZe}, \emph{registerYe}

\paragraph{Description} The \emph{registerYe} and the \emph{registerZe} are interpreted as two signed integer words packed into 64 bits. The \emph{registerZe} words are multiplied by the signs of the \emph{registerYe} words and 32-bit saturated. The two resulting words are packed and stored into the \emph{registerWe}.


\paragraph{Execution} {Reservation table \textsf{ALU\_LITE} (Section~\ref{sec:reservation-ALU-LITE})}
~
{\begin{lstlisting}
stage RR:
  new argument3 = GPR[registerYe<<1];
  new argument2 = GPR[registerZe<<1];
  for (I = 0; I < 2; I += 1) {
    result1.32[I] = argument2.32[I]
    if (_SX_32(argument3.32[I]) < 0) {
      result1.32[I] = _SAT_32(-argument2.32[I]);
    }
    if (_ZX_32(argument3.32[I]) == 0) {
      result1.32[I] = 0;
    }
  };
stage E1:
  GPR[registerWe<<1] = result1;
\end{lstlisting}}
\newpage

\paragraph{Encoding} Format \textsf{ALU\_WPWRR0O} (Section~\ref{sec:format-ALU-WPWRR0O}), \verb|exucode2 = 0001|\\

\bitfields{ 1/P, 2/11, 1/1, 4/0001, 5/registerWo, 1/1, 2/10, 3/000, 1/0, 5/registerYo, 1/--, 5/registerZo, 1/1 }


\paragraph{Syntax} \texttt{signswp} \emph{registerWo} = \emph{registerZo}, \emph{registerYo}

\paragraph{Description} The \emph{registerYo} and the \emph{registerZo} are interpreted as two signed integer words packed into 64 bits. The \emph{registerZo} words are multiplied by the signs of the \emph{registerYo} words and 32-bit saturated. The two resulting words are packed and stored into the \emph{registerWo}.


\paragraph{Execution} {Reservation table \textsf{ALU\_LITE} (Section~\ref{sec:reservation-ALU-LITE})}
~
{\begin{lstlisting}
stage RR:
  new argument3 = GPR[(registerYo<<1)+1];
  new argument2 = GPR[(registerZo<<1)+1];
  for (I = 0; I < 2; I += 1) {
    result1.32[I] = argument2.32[I]
    if (_SX_32(argument3.32[I]) < 0) {
      result1.32[I] = _SAT_32(-argument2.32[I]);
    }
    if (_ZX_32(argument3.32[I]) == 0) {
      result1.32[I] = 0;
    }
  };
stage E1:
  GPR[(registerWo<<1)+1] = result1;
\end{lstlisting}}
\newpage

\paragraph{Encoding} Format \textsf{ALU\_WPWRR0E.M} (Section~\ref{sec:format-ALU-WPWRR0E.M}), \verb|exucode2 = 0001|\\

\bitfields{ 1/P, 2/00, 2/-- --, 27/upper27, 1/1, 2/11, 1/1, 4/0001, 5/registerWe, 1/0, 2/10, 3/000, 1/0, 1/splat32, 5/lower5, 5/registerZe, 1/1 }


\paragraph{Syntax} \texttt{signswp} \emph{registerWe} = \emph{registerZe}, \emph{upper27\_\discretionary{}{}{}lower5}\emph{splat32}

\paragraph{Description} The \emph{upper27\_\discretionary{}{}{}lower5} and the \emph{registerZe} are interpreted as two signed integer words packed into 64 bits. The \emph{registerZe} words are multiplied by the signs of the \emph{upper27\_\discretionary{}{}{}lower5} words and 32-bit saturated. The two resulting words are packed and stored into the \emph{registerWe}.


\paragraph{Modifier(s)}
\noindent\begin{itemize}
\item splat32 (cf. splat32 in section \ref{par:modifier-splat32} on page \pageref{par:modifier-splat32})
\end{itemize}

\paragraph{Execution} {Reservation table \textsf{ALU\_LITE.X} (Section~\ref{sec:reservation-ALU-LITE.X})}
~
{\begin{lstlisting}
stage ID:
  new splat32 = _ZX_1(splat32);
stage RR:
  new argument3 = splat32 ? (_WX_32(upper27_lower5) << 32) | _ZX_32(_WX_32(upper27_lower5)) : _SX_32(_WX_32(upper27_lower5));
  new argument2 = GPR[registerZe<<1];
  for (I = 0; I < 2; I += 1) {
    result1.32[I] = argument2.32[I]
    if (_SX_32(argument3.32[I]) < 0) {
      result1.32[I] = _SAT_32(-argument2.32[I]);
    }
    if (_ZX_32(argument3.32[I]) == 0) {
      result1.32[I] = 0;
    }
  };
stage E1:
  GPR[registerWe<<1] = result1;
\end{lstlisting}}
\newpage

\paragraph{Encoding} Format \textsf{ALU\_WPWRR0O.M} (Section~\ref{sec:format-ALU-WPWRR0O.M}), \verb|exucode2 = 0001|\\

\bitfields{ 1/P, 2/00, 2/-- --, 27/upper27, 1/1, 2/11, 1/1, 4/0001, 5/registerWo, 1/1, 2/10, 3/000, 1/0, 1/splat32, 5/lower5, 5/registerZo, 1/1 }


\paragraph{Syntax} \texttt{signswp} \emph{registerWo} = \emph{registerZo}, \emph{upper27\_\discretionary{}{}{}lower5}\emph{splat32}

\paragraph{Description} The \emph{upper27\_\discretionary{}{}{}lower5} and the \emph{registerZo} are interpreted as two signed integer words packed into 64 bits. The \emph{registerZo} words are multiplied by the signs of the \emph{upper27\_\discretionary{}{}{}lower5} words and 32-bit saturated. The two resulting words are packed and stored into the \emph{registerWo}.


\paragraph{Modifier(s)}
\noindent\begin{itemize}
\item splat32 (cf. splat32 in section \ref{par:modifier-splat32} on page \pageref{par:modifier-splat32})
\end{itemize}

\paragraph{Execution} {Reservation table \textsf{ALU\_LITE.X} (Section~\ref{sec:reservation-ALU-LITE.X})}
~
{\begin{lstlisting}
stage ID:
  new splat32 = _ZX_1(splat32);
stage RR:
  new argument3 = splat32 ? (_WX_32(upper27_lower5) << 32) | _ZX_32(_WX_32(upper27_lower5)) : _SX_32(_WX_32(upper27_lower5));
  new argument2 = GPR[(registerZo<<1)+1];
  for (I = 0; I < 2; I += 1) {
    result1.32[I] = argument2.32[I]
    if (_SX_32(argument3.32[I]) < 0) {
      result1.32[I] = _SAT_32(-argument2.32[I]);
    }
    if (_ZX_32(argument3.32[I]) == 0) {
      result1.32[I] = 0;
    }
  };
stage E1:
  GPR[(registerWo<<1)+1] = result1;
\end{lstlisting}}
\newpage
\subsubsection[{\small SIGNW: Apply Sign Word}]{SIGNW\hfill Apply Sign Word}
\label{instr:SIGNW}\index{signw}
 
\paragraph{Encoding} Format \textsf{ALU\_WWRR0} (Section~\ref{sec:format-ALU-WWRR0}), \verb|exucode2 = 0000|\\

\bitfields{ 1/P, 2/11, 1/0, 4/0000, 6/registerW, 2/11, 3/000, 1/signextw, 6/registerY, 6/registerZ }


\paragraph{Syntax} \texttt{signw}\emph{signextw} \emph{registerW} = \emph{registerZ}, \emph{registerY}

\paragraph{Description} The \emph{registerY} and the \emph{registerZ} are interpreted as integer signed words. The \emph{registerZ} word is multiplied by the sign of the \emph{registerY} word. The resulting word is stored into the \emph{registerW}.


\paragraph{Modifier(s)}
\noindent\begin{itemize}
\item signextw (cf. signextw in section \ref{par:modifier-signextw} on page \pageref{par:modifier-signextw})
\end{itemize}

\paragraph{Execution} {Reservation table \textsf{ALU\_TINY} (Section~\ref{sec:reservation-ALU-TINY})}
~
{\begin{lstlisting}
stage ID:
  new signextw = _ZX_1(signextw);
stage RR:
  new argument3 = GPR[registerY];
  new argument2 = GPR[registerZ];
  new result1 = argument2;
  if (_SX_32(argument3) < 0) {
    result1 = -argument2;
  }
  if (_ZX_32(argument3) == 0) {
    result1 = 0;
  }
stage E1:
  GPR[registerW] = signextw ? _SX_32(result1) : _ZX_32(result1);
\end{lstlisting}}
\newpage

\paragraph{Encoding} Format \textsf{ALU\_WWRR0.W} (Section~\ref{sec:format-ALU-WWRR0.W}), \verb|exucode2 = 0000|\\

\bitfields{ 1/P, 2/00, 2/-- --, 27/upper27, 1/1, 2/11, 1/0, 4/0000, 6/registerW, 2/11, 3/000, 1/signextw, 1/0, 5/lower5, 6/registerZ }


\paragraph{Syntax} \texttt{signw}\emph{signextw} \emph{registerW} = \emph{registerZ}, \emph{upper27\_\discretionary{}{}{}lower5}

\paragraph{Description} The \emph{upper27\_\discretionary{}{}{}lower5} and the \emph{registerZ} are interpreted as integer signed words. The \emph{registerZ} word is multiplied by the sign of the \emph{upper27\_\discretionary{}{}{}lower5} word. The resulting word is stored into the \emph{registerW}.


\paragraph{Modifier(s)}
\noindent\begin{itemize}
\item signextw (cf. signextw in section \ref{par:modifier-signextw} on page \pageref{par:modifier-signextw})
\end{itemize}

\paragraph{Execution} {Reservation table \textsf{ALU\_TINY.X} (Section~\ref{sec:reservation-ALU-TINY.X})}
~
{\begin{lstlisting}
stage ID:
  new signextw = _ZX_1(signextw);
stage RR:
  new argument3 = _WX_32(upper27_lower5);
  new argument2 = GPR[registerZ];
  new result1 = argument2;
  if (_SX_32(argument3) < 0) {
    result1 = -argument2;
  }
  if (_ZX_32(argument3) == 0) {
    result1 = 0;
  }
stage E1:
  GPR[registerW] = signextw ? _SX_32(result1) : _ZX_32(result1);
\end{lstlisting}}
\newpage
\subsubsection[{\small SIGNWP: Apply Sign Word Pair}]{SIGNWP\hfill Apply Sign Word Pair}
\label{instr:SIGNWP}\index{signwp}
 
\paragraph{Encoding} Format \textsf{ALU\_WPWRR0E} (Section~\ref{sec:format-ALU-WPWRR0E}), \verb|exucode2 = 0000|\\

\bitfields{ 1/P, 2/11, 1/1, 4/0000, 5/registerWe, 1/0, 2/10, 3/000, 1/0, 5/registerYe, 1/--, 5/registerZe, 1/1 }


\paragraph{Syntax} \texttt{signwp} \emph{registerWe} = \emph{registerZe}, \emph{registerYe}

\paragraph{Description} The \emph{registerYe} and the \emph{registerZe} are interpreted as two signed integer words packed into 64 bits. The \emph{registerZe} words are multiplied by the signs of the \emph{registerYe} words. The two resulting words are packed and stored into the \emph{registerWe}.


\paragraph{Execution} {Reservation table \textsf{ALU\_LITE} (Section~\ref{sec:reservation-ALU-LITE})}
~
{\begin{lstlisting}
stage RR:
  new argument3 = GPR[registerYe<<1];
  new argument2 = GPR[registerZe<<1];
  for (I = 0; I < 2; I += 1) {
    result1.32[I] = argument2.32[I]
    if (_SX_32(argument3.32[I]) < 0) {
      result1.32[I] = -argument2.32[I];
    }
    if (_ZX_32(argument3.32[I]) == 0) {
      result1.32[I] = 0;
    }
  };
stage E1:
  GPR[registerWe<<1] = result1;
\end{lstlisting}}
\newpage

\paragraph{Encoding} Format \textsf{ALU\_WPWRR0O} (Section~\ref{sec:format-ALU-WPWRR0O}), \verb|exucode2 = 0000|\\

\bitfields{ 1/P, 2/11, 1/1, 4/0000, 5/registerWo, 1/1, 2/10, 3/000, 1/0, 5/registerYo, 1/--, 5/registerZo, 1/1 }


\paragraph{Syntax} \texttt{signwp} \emph{registerWo} = \emph{registerZo}, \emph{registerYo}

\paragraph{Description} The \emph{registerYo} and the \emph{registerZo} are interpreted as two signed integer words packed into 64 bits. The \emph{registerZo} words are multiplied by the signs of the \emph{registerYo} words. The two resulting words are packed and stored into the \emph{registerWo}.


\paragraph{Execution} {Reservation table \textsf{ALU\_LITE} (Section~\ref{sec:reservation-ALU-LITE})}
~
{\begin{lstlisting}
stage RR:
  new argument3 = GPR[(registerYo<<1)+1];
  new argument2 = GPR[(registerZo<<1)+1];
  for (I = 0; I < 2; I += 1) {
    result1.32[I] = argument2.32[I]
    if (_SX_32(argument3.32[I]) < 0) {
      result1.32[I] = -argument2.32[I];
    }
    if (_ZX_32(argument3.32[I]) == 0) {
      result1.32[I] = 0;
    }
  };
stage E1:
  GPR[(registerWo<<1)+1] = result1;
\end{lstlisting}}
\newpage

\paragraph{Encoding} Format \textsf{ALU\_WPWRR0E.M} (Section~\ref{sec:format-ALU-WPWRR0E.M}), \verb|exucode2 = 0000|\\

\bitfields{ 1/P, 2/00, 2/-- --, 27/upper27, 1/1, 2/11, 1/1, 4/0000, 5/registerWe, 1/0, 2/10, 3/000, 1/0, 1/splat32, 5/lower5, 5/registerZe, 1/1 }


\paragraph{Syntax} \texttt{signwp} \emph{registerWe} = \emph{registerZe}, \emph{upper27\_\discretionary{}{}{}lower5}\emph{splat32}

\paragraph{Description} The \emph{upper27\_\discretionary{}{}{}lower5} and the \emph{registerZe} are interpreted as two signed integer words packed into 64 bits. The \emph{registerZe} words are multiplied by the signs of the \emph{upper27\_\discretionary{}{}{}lower5} words. The two resulting words are packed and stored into the \emph{registerWe}.


\paragraph{Modifier(s)}
\noindent\begin{itemize}
\item splat32 (cf. splat32 in section \ref{par:modifier-splat32} on page \pageref{par:modifier-splat32})
\end{itemize}

\paragraph{Execution} {Reservation table \textsf{ALU\_LITE.X} (Section~\ref{sec:reservation-ALU-LITE.X})}
~
{\begin{lstlisting}
stage ID:
  new splat32 = _ZX_1(splat32);
stage RR:
  new argument3 = splat32 ? (_WX_32(upper27_lower5) << 32) | _ZX_32(_WX_32(upper27_lower5)) : _SX_32(_WX_32(upper27_lower5));
  new argument2 = GPR[registerZe<<1];
  for (I = 0; I < 2; I += 1) {
    result1.32[I] = argument2.32[I]
    if (_SX_32(argument3.32[I]) < 0) {
      result1.32[I] = -argument2.32[I];
    }
    if (_ZX_32(argument3.32[I]) == 0) {
      result1.32[I] = 0;
    }
  };
stage E1:
  GPR[registerWe<<1] = result1;
\end{lstlisting}}
\newpage

\paragraph{Encoding} Format \textsf{ALU\_WPWRR0O.M} (Section~\ref{sec:format-ALU-WPWRR0O.M}), \verb|exucode2 = 0000|\\

\bitfields{ 1/P, 2/00, 2/-- --, 27/upper27, 1/1, 2/11, 1/1, 4/0000, 5/registerWo, 1/1, 2/10, 3/000, 1/0, 1/splat32, 5/lower5, 5/registerZo, 1/1 }


\paragraph{Syntax} \texttt{signwp} \emph{registerWo} = \emph{registerZo}, \emph{upper27\_\discretionary{}{}{}lower5}\emph{splat32}

\paragraph{Description} The \emph{upper27\_\discretionary{}{}{}lower5} and the \emph{registerZo} are interpreted as two signed integer words packed into 64 bits. The \emph{registerZo} words are multiplied by the signs of the \emph{upper27\_\discretionary{}{}{}lower5} words. The two resulting words are packed and stored into the \emph{registerWo}.


\paragraph{Modifier(s)}
\noindent\begin{itemize}
\item splat32 (cf. splat32 in section \ref{par:modifier-splat32} on page \pageref{par:modifier-splat32})
\end{itemize}

\paragraph{Execution} {Reservation table \textsf{ALU\_LITE.X} (Section~\ref{sec:reservation-ALU-LITE.X})}
~
{\begin{lstlisting}
stage ID:
  new splat32 = _ZX_1(splat32);
stage RR:
  new argument3 = splat32 ? (_WX_32(upper27_lower5) << 32) | _ZX_32(_WX_32(upper27_lower5)) : _SX_32(_WX_32(upper27_lower5));
  new argument2 = GPR[(registerZo<<1)+1];
  for (I = 0; I < 2; I += 1) {
    result1.32[I] = argument2.32[I]
    if (_SX_32(argument3.32[I]) < 0) {
      result1.32[I] = -argument2.32[I];
    }
    if (_ZX_32(argument3.32[I]) == 0) {
      result1.32[I] = 0;
    }
  };
stage E1:
  GPR[(registerWo<<1)+1] = result1;
\end{lstlisting}}
\newpage
\subsubsection[{\small SLLBO: Shift Left Logical Byte Octuple}]{SLLBO\hfill Shift Left Logical Byte Octuple}
\label{instr:SLLBO}\index{sllbo}
 
\paragraph{Encoding} Format \textsf{ALU\_BOSWRIE} (Section~\ref{sec:format-ALU-BOSWRIE}), \verb|exucode1 = 001|\\

\bitfields{ 1/P, 2/11, 1/1, 1/1, 3/001, 5/registerWe, 1/0, 2/10, 3/100, 1/1, 6/unsigned6, 5/registerZe, 1/1 }


\paragraph{Syntax} \texttt{sllbo} \emph{registerWe} = \emph{registerZe}, \emph{unsigned6}

\paragraph{Description} The \emph{registerZe} is interpreted as eight bytes packed into 64 bits. The \emph{registerZe} bytes are shifted left by the \emph{unsigned6} modulo 8. The vacant positions are filled in with zeros. The eight resulting bytes are packed and stored into the \emph{registerWe}.


\paragraph{Execution} {Reservation table \textsf{ALU\_LITE} (Section~\ref{sec:reservation-ALU-LITE})}
~
{\begin{lstlisting}
stage RR:
  new argument3 = _ZX_6(unsigned6);
  new argument2 = GPR[registerZe<<1];
  new shift = _ZX_3(argument3);
  for (I = 0; I < 8; I += 1) {
    result1.8[I] = argument2.8[I] << shift
  };
stage E1:
  GPR[registerWe<<1] = result1;
\end{lstlisting}}
\newpage

\paragraph{Encoding} Format \textsf{ALU\_BOSWRIO} (Section~\ref{sec:format-ALU-BOSWRIO}), \verb|exucode1 = 001|\\

\bitfields{ 1/P, 2/11, 1/1, 1/1, 3/001, 5/registerWo, 1/1, 2/10, 3/100, 1/1, 6/unsigned6, 5/registerZo, 1/1 }


\paragraph{Syntax} \texttt{sllbo} \emph{registerWo} = \emph{registerZo}, \emph{unsigned6}

\paragraph{Description} The \emph{registerZo} is interpreted as eight bytes packed into 64 bits. The \emph{registerZo} bytes are shifted left by the \emph{unsigned6} modulo 8. The vacant positions are filled in with zeros. The eight resulting bytes are packed and stored into the \emph{registerWo}.


\paragraph{Execution} {Reservation table \textsf{ALU\_LITE} (Section~\ref{sec:reservation-ALU-LITE})}
~
{\begin{lstlisting}
stage RR:
  new argument3 = _ZX_6(unsigned6);
  new argument2 = GPR[(registerZo<<1)+1];
  new shift = _ZX_3(argument3);
  for (I = 0; I < 8; I += 1) {
    result1.8[I] = argument2.8[I] << shift
  };
stage E1:
  GPR[(registerWo<<1)+1] = result1;
\end{lstlisting}}
\newpage

\paragraph{Encoding} Format \textsf{ALU\_BOSWRRE} (Section~\ref{sec:format-ALU-BOSWRRE}), \verb|exucode1 = 001|\\

\bitfields{ 1/P, 2/11, 1/1, 1/0, 3/001, 5/registerWe, 1/0, 2/10, 3/100, 1/1, 6/registerY, 5/registerZe, 1/1 }


\paragraph{Syntax} \texttt{sllbo} \emph{registerWe} = \emph{registerZe}, \emph{registerY}

\paragraph{Description} The \emph{registerZe} is interpreted as eight bytes packed into 64 bits. The \emph{registerZe} bytes are shifted left by the \emph{registerY} modulo 8. The vacant positions are filled in with zeros. The eight resulting bytes are packed and stored into the \emph{registerWe}.


\paragraph{Execution} {Reservation table \textsf{ALU\_LITE} (Section~\ref{sec:reservation-ALU-LITE})}
~
{\begin{lstlisting}
stage RR:
  new argument3 = GPR[registerY];
  new argument2 = GPR[registerZe<<1];
  new shift = _ZX_3(argument3);
  for (I = 0; I < 8; I += 1) {
    result1.8[I] = argument2.8[I] << shift
  };
stage E1:
  GPR[registerWe<<1] = result1;
\end{lstlisting}}
\newpage

\paragraph{Encoding} Format \textsf{ALU\_BOSWRRO} (Section~\ref{sec:format-ALU-BOSWRRO}), \verb|exucode1 = 001|\\

\bitfields{ 1/P, 2/11, 1/1, 1/0, 3/001, 5/registerWo, 1/1, 2/10, 3/100, 1/1, 6/registerY, 5/registerZo, 1/1 }


\paragraph{Syntax} \texttt{sllbo} \emph{registerWo} = \emph{registerZo}, \emph{registerY}

\paragraph{Description} The \emph{registerZo} is interpreted as eight bytes packed into 64 bits. The \emph{registerZo} bytes are shifted left by the \emph{registerY} modulo 8. The vacant positions are filled in with zeros. The eight resulting bytes are packed and stored into the \emph{registerWo}.


\paragraph{Execution} {Reservation table \textsf{ALU\_LITE} (Section~\ref{sec:reservation-ALU-LITE})}
~
{\begin{lstlisting}
stage RR:
  new argument3 = GPR[registerY];
  new argument2 = GPR[(registerZo<<1)+1];
  new shift = _ZX_3(argument3);
  for (I = 0; I < 8; I += 1) {
    result1.8[I] = argument2.8[I] << shift
  };
stage E1:
  GPR[(registerWo<<1)+1] = result1;
\end{lstlisting}}
\newpage
\subsubsection[{\small SLLD: Shift Left Logical Double Word}]{SLLD\hfill Shift Left Logical Double Word}
\label{instr:SLLD}\index{slld}
 
\paragraph{Encoding} Format \textsf{ALU\_DSWRI} (Section~\ref{sec:format-ALU-DSWRI}), \verb|exucode1 = 001|\\

\bitfields{ 1/P, 2/11, 1/0, 1/1, 3/001, 6/registerW, 2/10, 3/100, 1/0, 6/unsigned6, 6/registerZ }


\paragraph{Syntax} \texttt{slld} \emph{registerW} = \emph{registerZ}, \emph{unsigned6}

\paragraph{Description} The \emph{registerZ} is shifted left by the \emph{unsigned6} modulo 64. The vacant positions are filled in with zeros. The result is stored into the \emph{registerW}.


\paragraph{Execution} {Reservation table \textsf{ALU\_TINY} (Section~\ref{sec:reservation-ALU-TINY})}
~
{\begin{lstlisting}
stage RR:
  new argument3 = _ZX_6(unsigned6);
  new argument2 = GPR[registerZ];
  new result1 = _ZX_64(argument2) << _ZX_6(argument3);
stage E1:
  GPR[registerW] = result1;
\end{lstlisting}}
\newpage

\paragraph{Encoding} Format \textsf{ALU\_DSWRR} (Section~\ref{sec:format-ALU-DSWRR}), \verb|exucode1 = 001|\\

\bitfields{ 1/P, 2/11, 1/0, 1/0, 3/001, 6/registerW, 2/10, 3/100, 1/0, 6/registerY, 6/registerZ }


\paragraph{Syntax} \texttt{slld} \emph{registerW} = \emph{registerZ}, \emph{registerY}

\paragraph{Description} The \emph{registerZ} is shifted left by the \emph{registerY} modulo 64. The vacant positions are filled in with zeros. The result is stored into the \emph{registerW}.


\paragraph{Execution} {Reservation table \textsf{ALU\_TINY} (Section~\ref{sec:reservation-ALU-TINY})}
~
{\begin{lstlisting}
stage RR:
  new argument3 = GPR[registerY];
  new argument2 = GPR[registerZ];
  new result1 = _ZX_64(argument2) << _ZX_6(argument3);
stage E1:
  GPR[registerW] = result1;
\end{lstlisting}}
\newpage
\subsubsection[{\small SLLHQ: Shift Left Logical Half Word Quadruple}]{SLLHQ\hfill Shift Left Logical Half Word Quadruple}
\label{instr:SLLHQ}\index{sllhq}
 
\paragraph{Encoding} Format \textsf{ALU\_HQSWRIE} (Section~\ref{sec:format-ALU-HQSWRIE}), \verb|exucode1 = 001|\\

\bitfields{ 1/P, 2/11, 1/1, 1/1, 3/001, 5/registerWe, 1/0, 2/10, 3/100, 1/1, 6/unsigned6, 5/registerZe, 1/0 }


\paragraph{Syntax} \texttt{sllhq} \emph{registerWe} = \emph{registerZe}, \emph{unsigned6}

\paragraph{Description} The \emph{registerZe} is interpreted as four half words packed into 64 bits. The \emph{registerZe} half words are shifted left by the \emph{unsigned6} modulo 16. The vacant positions are filled in with zeros. The four resulting half words are packed and stored into the \emph{registerWe}.


\paragraph{Execution} {Reservation table \textsf{ALU\_LITE} (Section~\ref{sec:reservation-ALU-LITE})}
~
{\begin{lstlisting}
stage RR:
  new argument3 = _ZX_6(unsigned6);
  new argument2 = GPR[registerZe<<1];
  new shift = _ZX_4(argument3);
  for (I = 0; I < 4; I += 1) {
    result1.16[I] = argument2.16[I] << shift
  };
stage E1:
  GPR[registerWe<<1] = result1;
\end{lstlisting}}
\newpage

\paragraph{Encoding} Format \textsf{ALU\_HQSWRIO} (Section~\ref{sec:format-ALU-HQSWRIO}), \verb|exucode1 = 001|\\

\bitfields{ 1/P, 2/11, 1/1, 1/1, 3/001, 5/registerWo, 1/1, 2/10, 3/100, 1/1, 6/unsigned6, 5/registerZo, 1/0 }


\paragraph{Syntax} \texttt{sllhq} \emph{registerWo} = \emph{registerZo}, \emph{unsigned6}

\paragraph{Description} The \emph{registerZo} is interpreted as four half words packed into 64 bits. The \emph{registerZo} half words are shifted left by the \emph{unsigned6} modulo 16. The vacant positions are filled in with zeros. The four resulting half words are packed and stored into the \emph{registerWo}.


\paragraph{Execution} {Reservation table \textsf{ALU\_LITE} (Section~\ref{sec:reservation-ALU-LITE})}
~
{\begin{lstlisting}
stage RR:
  new argument3 = _ZX_6(unsigned6);
  new argument2 = GPR[(registerZo<<1)+1];
  new shift = _ZX_4(argument3);
  for (I = 0; I < 4; I += 1) {
    result1.16[I] = argument2.16[I] << shift
  };
stage E1:
  GPR[(registerWo<<1)+1] = result1;
\end{lstlisting}}
\newpage

\paragraph{Encoding} Format \textsf{ALU\_HQSWRRE} (Section~\ref{sec:format-ALU-HQSWRRE}), \verb|exucode1 = 001|\\

\bitfields{ 1/P, 2/11, 1/1, 1/0, 3/001, 5/registerWe, 1/0, 2/10, 3/100, 1/1, 6/registerY, 5/registerZe, 1/0 }


\paragraph{Syntax} \texttt{sllhq} \emph{registerWe} = \emph{registerZe}, \emph{registerY}

\paragraph{Description} The \emph{registerZe} is interpreted as four half words packed into 64 bits. The \emph{registerZe} half words are shifted left by the \emph{registerY} modulo 16. The vacant positions are filled in with zeros. The four resulting half words are packed and stored into the \emph{registerWe}.


\paragraph{Execution} {Reservation table \textsf{ALU\_LITE} (Section~\ref{sec:reservation-ALU-LITE})}
~
{\begin{lstlisting}
stage RR:
  new argument3 = GPR[registerY];
  new argument2 = GPR[registerZe<<1];
  new shift = _ZX_4(argument3);
  for (I = 0; I < 4; I += 1) {
    result1.16[I] = argument2.16[I] << shift
  };
stage E1:
  GPR[registerWe<<1] = result1;
\end{lstlisting}}
\newpage

\paragraph{Encoding} Format \textsf{ALU\_HQSWRRO} (Section~\ref{sec:format-ALU-HQSWRRO}), \verb|exucode1 = 001|\\

\bitfields{ 1/P, 2/11, 1/1, 1/0, 3/001, 5/registerWo, 1/1, 2/10, 3/100, 1/1, 6/registerY, 5/registerZo, 1/0 }


\paragraph{Syntax} \texttt{sllhq} \emph{registerWo} = \emph{registerZo}, \emph{registerY}

\paragraph{Description} The \emph{registerZo} is interpreted as four half words packed into 64 bits. The \emph{registerZo} half words are shifted left by the \emph{registerY} modulo 16. The vacant positions are filled in with zeros. The four resulting half words are packed and stored into the \emph{registerWo}.


\paragraph{Execution} {Reservation table \textsf{ALU\_LITE} (Section~\ref{sec:reservation-ALU-LITE})}
~
{\begin{lstlisting}
stage RR:
  new argument3 = GPR[registerY];
  new argument2 = GPR[(registerZo<<1)+1];
  new shift = _ZX_4(argument3);
  for (I = 0; I < 4; I += 1) {
    result1.16[I] = argument2.16[I] << shift
  };
stage E1:
  GPR[(registerWo<<1)+1] = result1;
\end{lstlisting}}
\newpage
\subsubsection[{\small SLLW: Shift Left Logical Word}]{SLLW\hfill Shift Left Logical Word}
\label{instr:SLLW}\index{sllw}
 
\paragraph{Encoding} Format \textsf{ALU\_WSWRR} (Section~\ref{sec:format-ALU-WSWRR}), \verb|exucode1 = 001|\\

\bitfields{ 1/P, 2/11, 1/0, 1/0, 3/001, 6/registerW, 2/11, 3/100, 1/signextw, 6/registerY, 6/registerZ }


\paragraph{Syntax} \texttt{sllw}\emph{signextw} \emph{registerW} = \emph{registerZ}, \emph{registerY}

\paragraph{Description} The \emph{registerZ} is shifted left by the \emph{registerY} modulo 32. The vacant positions are filled in with zeros. The result with the 32 upper bits cleared is stored into the \emph{registerW}.


\paragraph{Modifier(s)}
\noindent\begin{itemize}
\item signextw (cf. signextw in section \ref{par:modifier-signextw} on page \pageref{par:modifier-signextw})
\end{itemize}

\paragraph{Execution} {Reservation table \textsf{ALU\_TINY} (Section~\ref{sec:reservation-ALU-TINY})}
~
{\begin{lstlisting}
stage ID:
  new signextw = _ZX_1(signextw);
stage RR:
  new argument3 = GPR[registerY];
  new argument2 = GPR[registerZ];
  new result1 = argument2.32[0] << _ZX_5(argument3);
stage E1:
  GPR[registerW] = signextw ? _SX_32(result1) : _ZX_32(result1);
\end{lstlisting}}
\newpage

\paragraph{Encoding} Format \textsf{ALU\_WSWRI} (Section~\ref{sec:format-ALU-WSWRI}), \verb|exucode1 = 001|\\

\bitfields{ 1/P, 2/11, 1/0, 1/1, 3/001, 6/registerW, 2/11, 3/100, 1/signextw, 6/unsigned6, 6/registerZ }


\paragraph{Syntax} \texttt{sllw}\emph{signextw} \emph{registerW} = \emph{registerZ}, \emph{unsigned6}

\paragraph{Description} The \emph{registerZ} is shifted left by the \emph{unsigned6} modulo 32. The vacant positions are filled in with zeros. The result with the 32 upper bits cleared is stored into the \emph{registerW}.


\paragraph{Modifier(s)}
\noindent\begin{itemize}
\item signextw (cf. signextw in section \ref{par:modifier-signextw} on page \pageref{par:modifier-signextw})
\end{itemize}

\paragraph{Execution} {Reservation table \textsf{ALU\_TINY} (Section~\ref{sec:reservation-ALU-TINY})}
~
{\begin{lstlisting}
stage ID:
  new signextw = _ZX_1(signextw);
stage RR:
  new argument3 = _ZX_6(unsigned6);
  new argument2 = GPR[registerZ];
  new result1 = argument2.32[0] << _ZX_5(argument3);
stage E1:
  GPR[registerW] = signextw ? _SX_32(result1) : _ZX_32(result1);
\end{lstlisting}}
\newpage
\subsubsection[{\small SLLWP: Shift Left Logical Word Pair}]{SLLWP\hfill Shift Left Logical Word Pair}
\label{instr:SLLWP}\index{sllwp}
 
\paragraph{Encoding} Format \textsf{ALU\_WPSWRIE} (Section~\ref{sec:format-ALU-WPSWRIE}), \verb|exucode1 = 001|\\

\bitfields{ 1/P, 2/11, 1/1, 1/1, 3/001, 5/registerWe, 1/0, 2/10, 3/100, 1/0, 6/unsigned6, 5/registerZe, 1/1 }


\paragraph{Syntax} \texttt{sllwp} \emph{registerWe} = \emph{registerZe}, \emph{unsigned6}

\paragraph{Description} The \emph{registerZe} is interpreted as two words packed into 64 bits. The \emph{registerZe} words are shifted left by the \emph{unsigned6} modulo 32. The vacant positions are filled in with zeros. The two resulting words are packed and stored into the \emph{registerWe}.


\paragraph{Execution} {Reservation table \textsf{ALU\_LITE} (Section~\ref{sec:reservation-ALU-LITE})}
~
{\begin{lstlisting}
stage RR:
  new argument3 = _ZX_6(unsigned6);
  new argument2 = GPR[registerZe<<1];
  new shift = _ZX_5(argument3);
  for (I = 0; I < 2; I += 1) {
    result1.32[I] = argument2.32[I] << shift
  };
stage E1:
  GPR[registerWe<<1] = result1;
\end{lstlisting}}
\newpage

\paragraph{Encoding} Format \textsf{ALU\_WPSWRIO} (Section~\ref{sec:format-ALU-WPSWRIO}), \verb|exucode1 = 001|\\

\bitfields{ 1/P, 2/11, 1/1, 1/1, 3/001, 5/registerWo, 1/1, 2/10, 3/100, 1/0, 6/unsigned6, 5/registerZo, 1/1 }


\paragraph{Syntax} \texttt{sllwp} \emph{registerWo} = \emph{registerZo}, \emph{unsigned6}

\paragraph{Description} The \emph{registerZo} is interpreted as two words packed into 64 bits. The \emph{registerZo} words are shifted left by the \emph{unsigned6} modulo 32. The vacant positions are filled in with zeros. The two resulting words are packed and stored into the \emph{registerWo}.


\paragraph{Execution} {Reservation table \textsf{ALU\_LITE} (Section~\ref{sec:reservation-ALU-LITE})}
~
{\begin{lstlisting}
stage RR:
  new argument3 = _ZX_6(unsigned6);
  new argument2 = GPR[(registerZo<<1)+1];
  new shift = _ZX_5(argument3);
  for (I = 0; I < 2; I += 1) {
    result1.32[I] = argument2.32[I] << shift
  };
stage E1:
  GPR[(registerWo<<1)+1] = result1;
\end{lstlisting}}
\newpage

\paragraph{Encoding} Format \textsf{ALU\_WPSWRRE} (Section~\ref{sec:format-ALU-WPSWRRE}), \verb|exucode1 = 001|\\

\bitfields{ 1/P, 2/11, 1/1, 1/0, 3/001, 5/registerWe, 1/0, 2/10, 3/100, 1/0, 6/registerY, 5/registerZe, 1/1 }


\paragraph{Syntax} \texttt{sllwp} \emph{registerWe} = \emph{registerZe}, \emph{registerY}

\paragraph{Description} The \emph{registerZe} is interpreted as two words packed into 64 bits. The \emph{registerZe} words are shifted left by the \emph{registerY} modulo 32. The vacant positions are filled in with zeros. The two resulting words are packed and stored into the \emph{registerWe}.


\paragraph{Execution} {Reservation table \textsf{ALU\_LITE} (Section~\ref{sec:reservation-ALU-LITE})}
~
{\begin{lstlisting}
stage RR:
  new argument3 = GPR[registerY];
  new argument2 = GPR[registerZe<<1];
  new shift = _ZX_5(argument3);
  for (I = 0; I < 2; I += 1) {
    result1.32[I] = argument2.32[I] << shift
  };
stage E1:
  GPR[registerWe<<1] = result1;
\end{lstlisting}}
\newpage

\paragraph{Encoding} Format \textsf{ALU\_WPSWRRO} (Section~\ref{sec:format-ALU-WPSWRRO}), \verb|exucode1 = 001|\\

\bitfields{ 1/P, 2/11, 1/1, 1/0, 3/001, 5/registerWo, 1/1, 2/10, 3/100, 1/0, 6/registerY, 5/registerZo, 1/1 }


\paragraph{Syntax} \texttt{sllwp} \emph{registerWo} = \emph{registerZo}, \emph{registerY}

\paragraph{Description} The \emph{registerZo} is interpreted as two words packed into 64 bits. The \emph{registerZo} words are shifted left by the \emph{registerY} modulo 32. The vacant positions are filled in with zeros. The two resulting words are packed and stored into the \emph{registerWo}.


\paragraph{Execution} {Reservation table \textsf{ALU\_LITE} (Section~\ref{sec:reservation-ALU-LITE})}
~
{\begin{lstlisting}
stage RR:
  new argument3 = GPR[registerY];
  new argument2 = GPR[(registerZo<<1)+1];
  new shift = _ZX_5(argument3);
  for (I = 0; I < 2; I += 1) {
    result1.32[I] = argument2.32[I] << shift
  };
stage E1:
  GPR[(registerWo<<1)+1] = result1;
\end{lstlisting}}
\newpage
\subsubsection[{\small SLSBO: Shift Left Saturated Byte Octuple}]{SLSBO\hfill Shift Left Saturated Byte Octuple}
\label{instr:SLSBO}\index{slsbo}
 
\paragraph{Encoding} Format \textsf{ALU\_BOSWRIE} (Section~\ref{sec:format-ALU-BOSWRIE}), \verb|exucode1 = 100|\\

\bitfields{ 1/P, 2/11, 1/1, 1/1, 3/100, 5/registerWe, 1/0, 2/10, 3/100, 1/1, 6/unsigned6, 5/registerZe, 1/1 }


\paragraph{Syntax} \texttt{slsbo} \emph{registerWe} = \emph{registerZe}, \emph{unsigned6}

\paragraph{Description} The \emph{registerZe} is interpreted as eight bytes packed into 64 bits. The \emph{registerZe} bytes are shifted left by the \emph{unsigned6} modulo 8. The eight resulting bytes are saturated to 8 bits, packed and stored into the \emph{registerWe}.


\paragraph{Execution} {Reservation table \textsf{ALU\_LITE} (Section~\ref{sec:reservation-ALU-LITE})}
~
{\begin{lstlisting}
stage RR:
  new argument3 = _ZX_6(unsigned6);
  new argument2 = GPR[registerZe<<1];
  new shift = _ZX_3(argument3);
  for (I = 0; I < 8; I += 1) {
    result1.8[I] = _SAT_8(_SX_8(argument2.8[I]) << shift)
  };
stage E1:
  GPR[registerWe<<1] = result1;
\end{lstlisting}}
\newpage

\paragraph{Encoding} Format \textsf{ALU\_BOSWRIO} (Section~\ref{sec:format-ALU-BOSWRIO}), \verb|exucode1 = 100|\\

\bitfields{ 1/P, 2/11, 1/1, 1/1, 3/100, 5/registerWo, 1/1, 2/10, 3/100, 1/1, 6/unsigned6, 5/registerZo, 1/1 }


\paragraph{Syntax} \texttt{slsbo} \emph{registerWo} = \emph{registerZo}, \emph{unsigned6}

\paragraph{Description} The \emph{registerZo} is interpreted as eight bytes packed into 64 bits. The \emph{registerZo} bytes are shifted left by the \emph{unsigned6} modulo 8. The eight resulting bytes are saturated to 8 bits, packed and stored into the \emph{registerWo}.


\paragraph{Execution} {Reservation table \textsf{ALU\_LITE} (Section~\ref{sec:reservation-ALU-LITE})}
~
{\begin{lstlisting}
stage RR:
  new argument3 = _ZX_6(unsigned6);
  new argument2 = GPR[(registerZo<<1)+1];
  new shift = _ZX_3(argument3);
  for (I = 0; I < 8; I += 1) {
    result1.8[I] = _SAT_8(_SX_8(argument2.8[I]) << shift)
  };
stage E1:
  GPR[(registerWo<<1)+1] = result1;
\end{lstlisting}}
\newpage

\paragraph{Encoding} Format \textsf{ALU\_BOSWRRE} (Section~\ref{sec:format-ALU-BOSWRRE}), \verb|exucode1 = 100|\\

\bitfields{ 1/P, 2/11, 1/1, 1/0, 3/100, 5/registerWe, 1/0, 2/10, 3/100, 1/1, 6/registerY, 5/registerZe, 1/1 }


\paragraph{Syntax} \texttt{slsbo} \emph{registerWe} = \emph{registerZe}, \emph{registerY}

\paragraph{Description} The \emph{registerZe} is interpreted as eight bytes packed into 64 bits. The \emph{registerZe} bytes are shifted left by the \emph{registerY} modulo 8. The eight resulting bytes are saturated to 8 bits, packed and stored into the \emph{registerWe}.


\paragraph{Execution} {Reservation table \textsf{ALU\_LITE} (Section~\ref{sec:reservation-ALU-LITE})}
~
{\begin{lstlisting}
stage RR:
  new argument3 = GPR[registerY];
  new argument2 = GPR[registerZe<<1];
  new shift = _ZX_3(argument3);
  for (I = 0; I < 8; I += 1) {
    result1.8[I] = _SAT_8(_SX_8(argument2.8[I]) << shift)
  };
stage E1:
  GPR[registerWe<<1] = result1;
\end{lstlisting}}
\newpage

\paragraph{Encoding} Format \textsf{ALU\_BOSWRRO} (Section~\ref{sec:format-ALU-BOSWRRO}), \verb|exucode1 = 100|\\

\bitfields{ 1/P, 2/11, 1/1, 1/0, 3/100, 5/registerWo, 1/1, 2/10, 3/100, 1/1, 6/registerY, 5/registerZo, 1/1 }


\paragraph{Syntax} \texttt{slsbo} \emph{registerWo} = \emph{registerZo}, \emph{registerY}

\paragraph{Description} The \emph{registerZo} is interpreted as eight bytes packed into 64 bits. The \emph{registerZo} bytes are shifted left by the \emph{registerY} modulo 8. The eight resulting bytes are saturated to 8 bits, packed and stored into the \emph{registerWo}.


\paragraph{Execution} {Reservation table \textsf{ALU\_LITE} (Section~\ref{sec:reservation-ALU-LITE})}
~
{\begin{lstlisting}
stage RR:
  new argument3 = GPR[registerY];
  new argument2 = GPR[(registerZo<<1)+1];
  new shift = _ZX_3(argument3);
  for (I = 0; I < 8; I += 1) {
    result1.8[I] = _SAT_8(_SX_8(argument2.8[I]) << shift)
  };
stage E1:
  GPR[(registerWo<<1)+1] = result1;
\end{lstlisting}}
\newpage
\subsubsection[{\small SLSD: Shift Left Saturated Double Word}]{SLSD\hfill Shift Left Saturated Double Word}
\label{instr:SLSD}\index{slsd}
 
\paragraph{Encoding} Format \textsf{ALU\_DSWRI} (Section~\ref{sec:format-ALU-DSWRI}), \verb|exucode1 = 100|\\

\bitfields{ 1/P, 2/11, 1/0, 1/1, 3/100, 6/registerW, 2/10, 3/100, 1/0, 6/unsigned6, 6/registerZ }


\paragraph{Syntax} \texttt{slsd} \emph{registerW} = \emph{registerZ}, \emph{unsigned6}

\paragraph{Description} The \emph{registerZ} is shifted left by the \emph{unsigned6} modulo 64. The result is saturated to 64 bits and stored into the \emph{registerW}.


\paragraph{Execution} {Reservation table \textsf{ALU\_LITE} (Section~\ref{sec:reservation-ALU-LITE})}
~
{\begin{lstlisting}
stage RR:
  new argument3 = _ZX_6(unsigned6);
  new argument2 = GPR[registerZ];
  new shift = _ZX_6(argument3);
  new result1 = _SAT_64(_SX_64(argument2) << shift);
stage E1:
  GPR[registerW] = result1;
\end{lstlisting}}
\newpage

\paragraph{Encoding} Format \textsf{ALU\_DSWRR} (Section~\ref{sec:format-ALU-DSWRR}), \verb|exucode1 = 100|\\

\bitfields{ 1/P, 2/11, 1/0, 1/0, 3/100, 6/registerW, 2/10, 3/100, 1/0, 6/registerY, 6/registerZ }


\paragraph{Syntax} \texttt{slsd} \emph{registerW} = \emph{registerZ}, \emph{registerY}

\paragraph{Description} The \emph{registerZ} is shifted left by the \emph{registerY} modulo 64. The result is saturated to 64 bits and stored into the \emph{registerW}.


\paragraph{Execution} {Reservation table \textsf{ALU\_LITE} (Section~\ref{sec:reservation-ALU-LITE})}
~
{\begin{lstlisting}
stage RR:
  new argument3 = GPR[registerY];
  new argument2 = GPR[registerZ];
  new shift = _ZX_6(argument3);
  new result1 = _SAT_64(_SX_64(argument2) << shift);
stage E1:
  GPR[registerW] = result1;
\end{lstlisting}}
\newpage
\subsubsection[{\small SLSHQ: Shift Left Saturated Half Word Quadruple}]{SLSHQ\hfill Shift Left Saturated Half Word Quadruple}
\label{instr:SLSHQ}\index{slshq}
 
\paragraph{Encoding} Format \textsf{ALU\_HQSWRIE} (Section~\ref{sec:format-ALU-HQSWRIE}), \verb|exucode1 = 100|\\

\bitfields{ 1/P, 2/11, 1/1, 1/1, 3/100, 5/registerWe, 1/0, 2/10, 3/100, 1/1, 6/unsigned6, 5/registerZe, 1/0 }


\paragraph{Syntax} \texttt{slshq} \emph{registerWe} = \emph{registerZe}, \emph{unsigned6}

\paragraph{Description} The \emph{registerZe} is interpreted as four half words packed into 128 bits. The \emph{registerZe} half words are shifted left by the \emph{unsigned6} modulo 16. The four resulting half words are saturated to 16 bits, packed and stored into the \emph{registerWe}.


\paragraph{Execution} {Reservation table \textsf{ALU\_LITE} (Section~\ref{sec:reservation-ALU-LITE})}
~
{\begin{lstlisting}
stage RR:
  new argument3 = _ZX_6(unsigned6);
  new argument2 = GPR[registerZe<<1];
  new shift = _ZX_4(argument3);
  for (I = 0; I < 4; I += 1) {
    result1.16[I] = _SAT_16(_SX_16(argument2.16[I]) << shift)
  };
stage E1:
  GPR[registerWe<<1] = result1;
\end{lstlisting}}
\newpage

\paragraph{Encoding} Format \textsf{ALU\_HQSWRIO} (Section~\ref{sec:format-ALU-HQSWRIO}), \verb|exucode1 = 100|\\

\bitfields{ 1/P, 2/11, 1/1, 1/1, 3/100, 5/registerWo, 1/1, 2/10, 3/100, 1/1, 6/unsigned6, 5/registerZo, 1/0 }


\paragraph{Syntax} \texttt{slshq} \emph{registerWo} = \emph{registerZo}, \emph{unsigned6}

\paragraph{Description} The \emph{registerZo} is interpreted as four half words packed into 128 bits. The \emph{registerZo} half words are shifted left by the \emph{unsigned6} modulo 16. The four resulting half words are saturated to 16 bits, packed and stored into the \emph{registerWo}.


\paragraph{Execution} {Reservation table \textsf{ALU\_LITE} (Section~\ref{sec:reservation-ALU-LITE})}
~
{\begin{lstlisting}
stage RR:
  new argument3 = _ZX_6(unsigned6);
  new argument2 = GPR[(registerZo<<1)+1];
  new shift = _ZX_4(argument3);
  for (I = 0; I < 4; I += 1) {
    result1.16[I] = _SAT_16(_SX_16(argument2.16[I]) << shift)
  };
stage E1:
  GPR[(registerWo<<1)+1] = result1;
\end{lstlisting}}
\newpage

\paragraph{Encoding} Format \textsf{ALU\_HQSWRRE} (Section~\ref{sec:format-ALU-HQSWRRE}), \verb|exucode1 = 100|\\

\bitfields{ 1/P, 2/11, 1/1, 1/0, 3/100, 5/registerWe, 1/0, 2/10, 3/100, 1/1, 6/registerY, 5/registerZe, 1/0 }


\paragraph{Syntax} \texttt{slshq} \emph{registerWe} = \emph{registerZe}, \emph{registerY}

\paragraph{Description} The \emph{registerZe} is interpreted as four half words packed into 128 bits. The \emph{registerZe} half words are shifted left by the \emph{registerY} modulo 16. The four resulting half words are saturated to 16 bits, packed and stored into the \emph{registerWe}.


\paragraph{Execution} {Reservation table \textsf{ALU\_LITE} (Section~\ref{sec:reservation-ALU-LITE})}
~
{\begin{lstlisting}
stage RR:
  new argument3 = GPR[registerY];
  new argument2 = GPR[registerZe<<1];
  new shift = _ZX_4(argument3);
  for (I = 0; I < 4; I += 1) {
    result1.16[I] = _SAT_16(_SX_16(argument2.16[I]) << shift)
  };
stage E1:
  GPR[registerWe<<1] = result1;
\end{lstlisting}}
\newpage

\paragraph{Encoding} Format \textsf{ALU\_HQSWRRO} (Section~\ref{sec:format-ALU-HQSWRRO}), \verb|exucode1 = 100|\\

\bitfields{ 1/P, 2/11, 1/1, 1/0, 3/100, 5/registerWo, 1/1, 2/10, 3/100, 1/1, 6/registerY, 5/registerZo, 1/0 }


\paragraph{Syntax} \texttt{slshq} \emph{registerWo} = \emph{registerZo}, \emph{registerY}

\paragraph{Description} The \emph{registerZo} is interpreted as four half words packed into 128 bits. The \emph{registerZo} half words are shifted left by the \emph{registerY} modulo 16. The four resulting half words are saturated to 16 bits, packed and stored into the \emph{registerWo}.


\paragraph{Execution} {Reservation table \textsf{ALU\_LITE} (Section~\ref{sec:reservation-ALU-LITE})}
~
{\begin{lstlisting}
stage RR:
  new argument3 = GPR[registerY];
  new argument2 = GPR[(registerZo<<1)+1];
  new shift = _ZX_4(argument3);
  for (I = 0; I < 4; I += 1) {
    result1.16[I] = _SAT_16(_SX_16(argument2.16[I]) << shift)
  };
stage E1:
  GPR[(registerWo<<1)+1] = result1;
\end{lstlisting}}
\newpage
\subsubsection[{\small SLSW: Shift Left Saturated Word}]{SLSW\hfill Shift Left Saturated Word}
\label{instr:SLSW}\index{slsw}
 
\paragraph{Encoding} Format \textsf{ALU\_WSWRR} (Section~\ref{sec:format-ALU-WSWRR}), \verb|exucode1 = 100|\\

\bitfields{ 1/P, 2/11, 1/0, 1/0, 3/100, 6/registerW, 2/11, 3/100, 1/signextw, 6/registerY, 6/registerZ }


\paragraph{Syntax} \texttt{slsw}\emph{signextw} \emph{registerW} = \emph{registerZ}, \emph{registerY}

\paragraph{Description} The \emph{registerZ} is shifted left by the \emph{registerY} modulo 32 then is saturated to 32 bits. The result with the 32 upper bits cleared is stored into the \emph{registerW}.


\paragraph{Modifier(s)}
\noindent\begin{itemize}
\item signextw (cf. signextw in section \ref{par:modifier-signextw} on page \pageref{par:modifier-signextw})
\end{itemize}

\paragraph{Execution} {Reservation table \textsf{ALU\_LITE} (Section~\ref{sec:reservation-ALU-LITE})}
~
{\begin{lstlisting}
stage ID:
  new signextw = _ZX_1(signextw);
stage RR:
  new argument3 = GPR[registerY];
  new argument2 = GPR[registerZ];
  new shift = _ZX_5(argument3);
  new result1 = _SAT_32(_SX_32(argument2) << shift);
stage E1:
  GPR[registerW] = signextw ? _SX_32(result1) : _ZX_32(result1);
\end{lstlisting}}
\newpage

\paragraph{Encoding} Format \textsf{ALU\_WSWRI} (Section~\ref{sec:format-ALU-WSWRI}), \verb|exucode1 = 100|\\

\bitfields{ 1/P, 2/11, 1/0, 1/1, 3/100, 6/registerW, 2/11, 3/100, 1/signextw, 6/unsigned6, 6/registerZ }


\paragraph{Syntax} \texttt{slsw}\emph{signextw} \emph{registerW} = \emph{registerZ}, \emph{unsigned6}

\paragraph{Description} The \emph{registerZ} is shifted left by the \emph{unsigned6} modulo 32 then is saturated to 32 bits. The result with the 32 upper bits cleared is stored into the \emph{registerW}.


\paragraph{Modifier(s)}
\noindent\begin{itemize}
\item signextw (cf. signextw in section \ref{par:modifier-signextw} on page \pageref{par:modifier-signextw})
\end{itemize}

\paragraph{Execution} {Reservation table \textsf{ALU\_LITE} (Section~\ref{sec:reservation-ALU-LITE})}
~
{\begin{lstlisting}
stage ID:
  new signextw = _ZX_1(signextw);
stage RR:
  new argument3 = _ZX_6(unsigned6);
  new argument2 = GPR[registerZ];
  new shift = _ZX_5(argument3);
  new result1 = _SAT_32(_SX_32(argument2) << shift);
stage E1:
  GPR[registerW] = signextw ? _SX_32(result1) : _ZX_32(result1);
\end{lstlisting}}
\newpage
\subsubsection[{\small SLSWP: Shift Left Saturated Word Pair}]{SLSWP\hfill Shift Left Saturated Word Pair}
\label{instr:SLSWP}\index{slswp}
 
\paragraph{Encoding} Format \textsf{ALU\_WPSWRIE} (Section~\ref{sec:format-ALU-WPSWRIE}), \verb|exucode1 = 100|\\

\bitfields{ 1/P, 2/11, 1/1, 1/1, 3/100, 5/registerWe, 1/0, 2/10, 3/100, 1/0, 6/unsigned6, 5/registerZe, 1/1 }


\paragraph{Syntax} \texttt{slswp} \emph{registerWe} = \emph{registerZe}, \emph{unsigned6}

\paragraph{Description} The \emph{registerZe} is interpreted as two words packed into 64 bits. The \emph{registerZe} words are shifted left by the \emph{unsigned6} modulo 32. The two resulting words are saturated to 32 bits, packed and stored into the \emph{registerWe}.


\paragraph{Execution} {Reservation table \textsf{ALU\_LITE} (Section~\ref{sec:reservation-ALU-LITE})}
~
{\begin{lstlisting}
stage RR:
  new argument3 = _ZX_6(unsigned6);
  new argument2 = GPR[registerZe<<1];
  new shift = _ZX_5(argument3);
  for (I = 0; I < 2; I += 1) {
    result1.32[I] = _SAT_32(_SX_32(argument2.32[I]) << shift)
  };
stage E1:
  GPR[registerWe<<1] = result1;
\end{lstlisting}}
\newpage

\paragraph{Encoding} Format \textsf{ALU\_WPSWRIO} (Section~\ref{sec:format-ALU-WPSWRIO}), \verb|exucode1 = 100|\\

\bitfields{ 1/P, 2/11, 1/1, 1/1, 3/100, 5/registerWo, 1/1, 2/10, 3/100, 1/0, 6/unsigned6, 5/registerZo, 1/1 }


\paragraph{Syntax} \texttt{slswp} \emph{registerWo} = \emph{registerZo}, \emph{unsigned6}

\paragraph{Description} The \emph{registerZo} is interpreted as two words packed into 64 bits. The \emph{registerZo} words are shifted left by the \emph{unsigned6} modulo 32. The two resulting words are saturated to 32 bits, packed and stored into the \emph{registerWo}.


\paragraph{Execution} {Reservation table \textsf{ALU\_LITE} (Section~\ref{sec:reservation-ALU-LITE})}
~
{\begin{lstlisting}
stage RR:
  new argument3 = _ZX_6(unsigned6);
  new argument2 = GPR[(registerZo<<1)+1];
  new shift = _ZX_5(argument3);
  for (I = 0; I < 2; I += 1) {
    result1.32[I] = _SAT_32(_SX_32(argument2.32[I]) << shift)
  };
stage E1:
  GPR[(registerWo<<1)+1] = result1;
\end{lstlisting}}
\newpage

\paragraph{Encoding} Format \textsf{ALU\_WPSWRRE} (Section~\ref{sec:format-ALU-WPSWRRE}), \verb|exucode1 = 100|\\

\bitfields{ 1/P, 2/11, 1/1, 1/0, 3/100, 5/registerWe, 1/0, 2/10, 3/100, 1/0, 6/registerY, 5/registerZe, 1/1 }


\paragraph{Syntax} \texttt{slswp} \emph{registerWe} = \emph{registerZe}, \emph{registerY}

\paragraph{Description} The \emph{registerZe} is interpreted as two words packed into 64 bits. The \emph{registerZe} words are shifted left by the \emph{registerY} modulo 32. The two resulting words are saturated to 32 bits, packed and stored into the \emph{registerWe}.


\paragraph{Execution} {Reservation table \textsf{ALU\_LITE} (Section~\ref{sec:reservation-ALU-LITE})}
~
{\begin{lstlisting}
stage RR:
  new argument3 = GPR[registerY];
  new argument2 = GPR[registerZe<<1];
  new shift = _ZX_5(argument3);
  for (I = 0; I < 2; I += 1) {
    result1.32[I] = _SAT_32(_SX_32(argument2.32[I]) << shift)
  };
stage E1:
  GPR[registerWe<<1] = result1;
\end{lstlisting}}
\newpage

\paragraph{Encoding} Format \textsf{ALU\_WPSWRRO} (Section~\ref{sec:format-ALU-WPSWRRO}), \verb|exucode1 = 100|\\

\bitfields{ 1/P, 2/11, 1/1, 1/0, 3/100, 5/registerWo, 1/1, 2/10, 3/100, 1/0, 6/registerY, 5/registerZo, 1/1 }


\paragraph{Syntax} \texttt{slswp} \emph{registerWo} = \emph{registerZo}, \emph{registerY}

\paragraph{Description} The \emph{registerZo} is interpreted as two words packed into 64 bits. The \emph{registerZo} words are shifted left by the \emph{registerY} modulo 32. The two resulting words are saturated to 32 bits, packed and stored into the \emph{registerWo}.


\paragraph{Execution} {Reservation table \textsf{ALU\_LITE} (Section~\ref{sec:reservation-ALU-LITE})}
~
{\begin{lstlisting}
stage RR:
  new argument3 = GPR[registerY];
  new argument2 = GPR[(registerZo<<1)+1];
  new shift = _ZX_5(argument3);
  for (I = 0; I < 2; I += 1) {
    result1.32[I] = _SAT_32(_SX_32(argument2.32[I]) << shift)
  };
stage E1:
  GPR[(registerWo<<1)+1] = result1;
\end{lstlisting}}
\newpage
\subsubsection[{\small SLUSBO: Shift Left Unsigned Saturated Byte Octuple}]{SLUSBO\hfill Shift Left Unsigned Saturated Byte Octuple}
\label{instr:SLUSBO}\index{slusbo}
 
\paragraph{Encoding} Format \textsf{ALU\_BOSWRIE} (Section~\ref{sec:format-ALU-BOSWRIE}), \verb|exucode1 = 101|\\

\bitfields{ 1/P, 2/11, 1/1, 1/1, 3/101, 5/registerWe, 1/0, 2/10, 3/100, 1/1, 6/unsigned6, 5/registerZe, 1/1 }


\paragraph{Syntax} \texttt{slusbo} \emph{registerWe} = \emph{registerZe}, \emph{unsigned6}

\paragraph{Description} The \emph{registerZe} is interpreted as eight bytes packed into 64 bits. The \emph{registerZe} bytes are shifted left by the \emph{unsigned6} modulo 8. The eight resulting bytes are unsigned saturated to 8 bits, packed and stored into the \emph{registerWe}.


\paragraph{Execution} {Reservation table \textsf{ALU\_LITE} (Section~\ref{sec:reservation-ALU-LITE})}
~
{\begin{lstlisting}
stage RR:
  new argument3 = _ZX_6(unsigned6);
  new argument2 = GPR[registerZe<<1];
  new shift = _ZX_3(argument3);
  for (I = 0; I < 8; I += 1) {
    result1.8[I] = _SATU_8(_ZX_8(argument2.8[I]) << shift)
  };
stage E1:
  GPR[registerWe<<1] = result1;
\end{lstlisting}}
\newpage

\paragraph{Encoding} Format \textsf{ALU\_BOSWRIO} (Section~\ref{sec:format-ALU-BOSWRIO}), \verb|exucode1 = 101|\\

\bitfields{ 1/P, 2/11, 1/1, 1/1, 3/101, 5/registerWo, 1/1, 2/10, 3/100, 1/1, 6/unsigned6, 5/registerZo, 1/1 }


\paragraph{Syntax} \texttt{slusbo} \emph{registerWo} = \emph{registerZo}, \emph{unsigned6}

\paragraph{Description} The \emph{registerZo} is interpreted as eight bytes packed into 64 bits. The \emph{registerZo} bytes are shifted left by the \emph{unsigned6} modulo 8. The eight resulting bytes are unsigned saturated to 8 bits, packed and stored into the \emph{registerWo}.


\paragraph{Execution} {Reservation table \textsf{ALU\_LITE} (Section~\ref{sec:reservation-ALU-LITE})}
~
{\begin{lstlisting}
stage RR:
  new argument3 = _ZX_6(unsigned6);
  new argument2 = GPR[(registerZo<<1)+1];
  new shift = _ZX_3(argument3);
  for (I = 0; I < 8; I += 1) {
    result1.8[I] = _SATU_8(_ZX_8(argument2.8[I]) << shift)
  };
stage E1:
  GPR[(registerWo<<1)+1] = result1;
\end{lstlisting}}
\newpage

\paragraph{Encoding} Format \textsf{ALU\_BOSWRRE} (Section~\ref{sec:format-ALU-BOSWRRE}), \verb|exucode1 = 101|\\

\bitfields{ 1/P, 2/11, 1/1, 1/0, 3/101, 5/registerWe, 1/0, 2/10, 3/100, 1/1, 6/registerY, 5/registerZe, 1/1 }


\paragraph{Syntax} \texttt{slusbo} \emph{registerWe} = \emph{registerZe}, \emph{registerY}

\paragraph{Description} The \emph{registerZe} is interpreted as eight bytes packed into 64 bits. The \emph{registerZe} bytes are shifted left by the \emph{registerY} modulo 8. The eight resulting bytes are unsigned saturated to 8 bits, packed and stored into the \emph{registerWe}.


\paragraph{Execution} {Reservation table \textsf{ALU\_LITE} (Section~\ref{sec:reservation-ALU-LITE})}
~
{\begin{lstlisting}
stage RR:
  new argument3 = GPR[registerY];
  new argument2 = GPR[registerZe<<1];
  new shift = _ZX_3(argument3);
  for (I = 0; I < 8; I += 1) {
    result1.8[I] = _SATU_8(_ZX_8(argument2.8[I]) << shift)
  };
stage E1:
  GPR[registerWe<<1] = result1;
\end{lstlisting}}
\newpage

\paragraph{Encoding} Format \textsf{ALU\_BOSWRRO} (Section~\ref{sec:format-ALU-BOSWRRO}), \verb|exucode1 = 101|\\

\bitfields{ 1/P, 2/11, 1/1, 1/0, 3/101, 5/registerWo, 1/1, 2/10, 3/100, 1/1, 6/registerY, 5/registerZo, 1/1 }


\paragraph{Syntax} \texttt{slusbo} \emph{registerWo} = \emph{registerZo}, \emph{registerY}

\paragraph{Description} The \emph{registerZo} is interpreted as eight bytes packed into 64 bits. The \emph{registerZo} bytes are shifted left by the \emph{registerY} modulo 8. The eight resulting bytes are unsigned saturated to 8 bits, packed and stored into the \emph{registerWo}.


\paragraph{Execution} {Reservation table \textsf{ALU\_LITE} (Section~\ref{sec:reservation-ALU-LITE})}
~
{\begin{lstlisting}
stage RR:
  new argument3 = GPR[registerY];
  new argument2 = GPR[(registerZo<<1)+1];
  new shift = _ZX_3(argument3);
  for (I = 0; I < 8; I += 1) {
    result1.8[I] = _SATU_8(_ZX_8(argument2.8[I]) << shift)
  };
stage E1:
  GPR[(registerWo<<1)+1] = result1;
\end{lstlisting}}
\newpage
\subsubsection[{\small SLUSD: Shift Left Unsigned Saturated Double Word}]{SLUSD\hfill Shift Left Unsigned Saturated Double Word}
\label{instr:SLUSD}\index{slusd}
 
\paragraph{Encoding} Format \textsf{ALU\_DSWRI} (Section~\ref{sec:format-ALU-DSWRI}), \verb|exucode1 = 101|\\

\bitfields{ 1/P, 2/11, 1/0, 1/1, 3/101, 6/registerW, 2/10, 3/100, 1/0, 6/unsigned6, 6/registerZ }


\paragraph{Syntax} \texttt{slusd} \emph{registerW} = \emph{registerZ}, \emph{unsigned6}

\paragraph{Description} The \emph{registerZ} is shifted left by the \emph{unsigned6} modulo 64. The result is unsigned saturated to 64 bits and stored into the \emph{registerW}.


\paragraph{Execution} {Reservation table \textsf{ALU\_LITE} (Section~\ref{sec:reservation-ALU-LITE})}
~
{\begin{lstlisting}
stage RR:
  new argument3 = _ZX_6(unsigned6);
  new argument2 = GPR[registerZ];
  new shift = _ZX_6(argument3);
  new result1 = _SATU_64(_ZX_64(argument2) << shift);
stage E1:
  GPR[registerW] = result1;
\end{lstlisting}}
\newpage

\paragraph{Encoding} Format \textsf{ALU\_DSWRR} (Section~\ref{sec:format-ALU-DSWRR}), \verb|exucode1 = 101|\\

\bitfields{ 1/P, 2/11, 1/0, 1/0, 3/101, 6/registerW, 2/10, 3/100, 1/0, 6/registerY, 6/registerZ }


\paragraph{Syntax} \texttt{slusd} \emph{registerW} = \emph{registerZ}, \emph{registerY}

\paragraph{Description} The \emph{registerZ} is shifted left by the \emph{registerY} modulo 64. The result is unsigned saturated to 64 bits and stored into the \emph{registerW}.


\paragraph{Execution} {Reservation table \textsf{ALU\_LITE} (Section~\ref{sec:reservation-ALU-LITE})}
~
{\begin{lstlisting}
stage RR:
  new argument3 = GPR[registerY];
  new argument2 = GPR[registerZ];
  new shift = _ZX_6(argument3);
  new result1 = _SATU_64(_ZX_64(argument2) << shift);
stage E1:
  GPR[registerW] = result1;
\end{lstlisting}}
\newpage
\subsubsection[{\small SLUSHQ: Shift Left Unsigned Saturated Half Word Quadruple}]{SLUSHQ\hfill Shift Left Unsigned Saturated Half Word Quadruple}
\label{instr:SLUSHQ}\index{slushq}
 
\paragraph{Encoding} Format \textsf{ALU\_HQSWRIE} (Section~\ref{sec:format-ALU-HQSWRIE}), \verb|exucode1 = 101|\\

\bitfields{ 1/P, 2/11, 1/1, 1/1, 3/101, 5/registerWe, 1/0, 2/10, 3/100, 1/1, 6/unsigned6, 5/registerZe, 1/0 }


\paragraph{Syntax} \texttt{slushq} \emph{registerWe} = \emph{registerZe}, \emph{unsigned6}

\paragraph{Description} The \emph{registerZe} is interpreted as four half words packed into 64 bits. The \emph{registerZe} half words are shifted left by the \emph{unsigned6} modulo 16. The four resulting half words are unsigned saturated to 16 bits, packed and stored into the \emph{registerWe}.


\paragraph{Execution} {Reservation table \textsf{ALU\_LITE} (Section~\ref{sec:reservation-ALU-LITE})}
~
{\begin{lstlisting}
stage RR:
  new argument3 = _ZX_6(unsigned6);
  new argument2 = GPR[registerZe<<1];
  new shift = _ZX_4(argument3);
  for (I = 0; I < 4; I += 1) {
    result1.16[I] = _SATU_16(_ZX_16(argument2.16[I]) << shift)
  };
stage E1:
  GPR[registerWe<<1] = result1;
\end{lstlisting}}
\newpage

\paragraph{Encoding} Format \textsf{ALU\_HQSWRIO} (Section~\ref{sec:format-ALU-HQSWRIO}), \verb|exucode1 = 101|\\

\bitfields{ 1/P, 2/11, 1/1, 1/1, 3/101, 5/registerWo, 1/1, 2/10, 3/100, 1/1, 6/unsigned6, 5/registerZo, 1/0 }


\paragraph{Syntax} \texttt{slushq} \emph{registerWo} = \emph{registerZo}, \emph{unsigned6}

\paragraph{Description} The \emph{registerZo} is interpreted as four half words packed into 64 bits. The \emph{registerZo} half words are shifted left by the \emph{unsigned6} modulo 16. The four resulting half words are unsigned saturated to 16 bits, packed and stored into the \emph{registerWo}.


\paragraph{Execution} {Reservation table \textsf{ALU\_LITE} (Section~\ref{sec:reservation-ALU-LITE})}
~
{\begin{lstlisting}
stage RR:
  new argument3 = _ZX_6(unsigned6);
  new argument2 = GPR[(registerZo<<1)+1];
  new shift = _ZX_4(argument3);
  for (I = 0; I < 4; I += 1) {
    result1.16[I] = _SATU_16(_ZX_16(argument2.16[I]) << shift)
  };
stage E1:
  GPR[(registerWo<<1)+1] = result1;
\end{lstlisting}}
\newpage

\paragraph{Encoding} Format \textsf{ALU\_HQSWRRE} (Section~\ref{sec:format-ALU-HQSWRRE}), \verb|exucode1 = 101|\\

\bitfields{ 1/P, 2/11, 1/1, 1/0, 3/101, 5/registerWe, 1/0, 2/10, 3/100, 1/1, 6/registerY, 5/registerZe, 1/0 }


\paragraph{Syntax} \texttt{slushq} \emph{registerWe} = \emph{registerZe}, \emph{registerY}

\paragraph{Description} The \emph{registerZe} is interpreted as four half words packed into 64 bits. The \emph{registerZe} half words are shifted left by the \emph{registerY} modulo 16. The four resulting half words are unsigned saturated to 16 bits, packed and stored into the \emph{registerWe}.


\paragraph{Execution} {Reservation table \textsf{ALU\_LITE} (Section~\ref{sec:reservation-ALU-LITE})}
~
{\begin{lstlisting}
stage RR:
  new argument3 = GPR[registerY];
  new argument2 = GPR[registerZe<<1];
  new shift = _ZX_4(argument3);
  for (I = 0; I < 4; I += 1) {
    result1.16[I] = _SATU_16(_ZX_16(argument2.16[I]) << shift)
  };
stage E1:
  GPR[registerWe<<1] = result1;
\end{lstlisting}}
\newpage

\paragraph{Encoding} Format \textsf{ALU\_HQSWRRO} (Section~\ref{sec:format-ALU-HQSWRRO}), \verb|exucode1 = 101|\\

\bitfields{ 1/P, 2/11, 1/1, 1/0, 3/101, 5/registerWo, 1/1, 2/10, 3/100, 1/1, 6/registerY, 5/registerZo, 1/0 }


\paragraph{Syntax} \texttt{slushq} \emph{registerWo} = \emph{registerZo}, \emph{registerY}

\paragraph{Description} The \emph{registerZo} is interpreted as four half words packed into 64 bits. The \emph{registerZo} half words are shifted left by the \emph{registerY} modulo 16. The four resulting half words are unsigned saturated to 16 bits, packed and stored into the \emph{registerWo}.


\paragraph{Execution} {Reservation table \textsf{ALU\_LITE} (Section~\ref{sec:reservation-ALU-LITE})}
~
{\begin{lstlisting}
stage RR:
  new argument3 = GPR[registerY];
  new argument2 = GPR[(registerZo<<1)+1];
  new shift = _ZX_4(argument3);
  for (I = 0; I < 4; I += 1) {
    result1.16[I] = _SATU_16(_ZX_16(argument2.16[I]) << shift)
  };
stage E1:
  GPR[(registerWo<<1)+1] = result1;
\end{lstlisting}}
\newpage
\subsubsection[{\small SLUSW: Shift Left Unsigned Saturated Word}]{SLUSW\hfill Shift Left Unsigned Saturated Word}
\label{instr:SLUSW}\index{slusw}
 
\paragraph{Encoding} Format \textsf{ALU\_WSWRR} (Section~\ref{sec:format-ALU-WSWRR}), \verb|exucode1 = 101|\\

\bitfields{ 1/P, 2/11, 1/0, 1/0, 3/101, 6/registerW, 2/11, 3/100, 1/signextw, 6/registerY, 6/registerZ }


\paragraph{Syntax} \texttt{slusw}\emph{signextw} \emph{registerW} = \emph{registerZ}, \emph{registerY}

\paragraph{Description} The \emph{registerZ} is shifted left by the \emph{registerY} modulo 32 then is unsigned saturated to 32 bits. The result with the 32 upper bits cleared is stored into the \emph{registerW}.


\paragraph{Modifier(s)}
\noindent\begin{itemize}
\item signextw (cf. signextw in section \ref{par:modifier-signextw} on page \pageref{par:modifier-signextw})
\end{itemize}

\paragraph{Execution} {Reservation table \textsf{ALU\_LITE} (Section~\ref{sec:reservation-ALU-LITE})}
~
{\begin{lstlisting}
stage ID:
  new signextw = _ZX_1(signextw);
stage RR:
  new argument3 = GPR[registerY];
  new argument2 = GPR[registerZ];
  new shift = _ZX_5(argument3);
  new result1 = _SATU_32(_ZX_32(argument2) << shift);
stage E1:
  GPR[registerW] = signextw ? _SX_32(result1) : _ZX_32(result1);
\end{lstlisting}}
\newpage

\paragraph{Encoding} Format \textsf{ALU\_WSWRI} (Section~\ref{sec:format-ALU-WSWRI}), \verb|exucode1 = 101|\\

\bitfields{ 1/P, 2/11, 1/0, 1/1, 3/101, 6/registerW, 2/11, 3/100, 1/signextw, 6/unsigned6, 6/registerZ }


\paragraph{Syntax} \texttt{slusw}\emph{signextw} \emph{registerW} = \emph{registerZ}, \emph{unsigned6}

\paragraph{Description} The \emph{registerZ} is shifted left by the \emph{unsigned6} modulo 32 then is unsigned saturated to 32 bits. The result with the 32 upper bits cleared is stored into the \emph{registerW}.


\paragraph{Modifier(s)}
\noindent\begin{itemize}
\item signextw (cf. signextw in section \ref{par:modifier-signextw} on page \pageref{par:modifier-signextw})
\end{itemize}

\paragraph{Execution} {Reservation table \textsf{ALU\_LITE} (Section~\ref{sec:reservation-ALU-LITE})}
~
{\begin{lstlisting}
stage ID:
  new signextw = _ZX_1(signextw);
stage RR:
  new argument3 = _ZX_6(unsigned6);
  new argument2 = GPR[registerZ];
  new shift = _ZX_5(argument3);
  new result1 = _SATU_32(_ZX_32(argument2) << shift);
stage E1:
  GPR[registerW] = signextw ? _SX_32(result1) : _ZX_32(result1);
\end{lstlisting}}
\newpage
\subsubsection[{\small SLUSWP: Shift Left Unsigned Saturated Word Pair}]{SLUSWP\hfill Shift Left Unsigned Saturated Word Pair}
\label{instr:SLUSWP}\index{sluswp}
 
\paragraph{Encoding} Format \textsf{ALU\_WPSWRIE} (Section~\ref{sec:format-ALU-WPSWRIE}), \verb|exucode1 = 101|\\

\bitfields{ 1/P, 2/11, 1/1, 1/1, 3/101, 5/registerWe, 1/0, 2/10, 3/100, 1/0, 6/unsigned6, 5/registerZe, 1/1 }


\paragraph{Syntax} \texttt{sluswp} \emph{registerWe} = \emph{registerZe}, \emph{unsigned6}

\paragraph{Description} The \emph{registerZe} is interpreted as two words packed into 64 bits. The \emph{registerZe} words are shifted left by the \emph{unsigned6} modulo 32. The two resulting words are unsigned saturated to 32 bits, packed and stored into the \emph{registerWe}.


\paragraph{Execution} {Reservation table \textsf{ALU\_LITE} (Section~\ref{sec:reservation-ALU-LITE})}
~
{\begin{lstlisting}
stage RR:
  new argument3 = _ZX_6(unsigned6);
  new argument2 = GPR[registerZe<<1];
  new shift = _ZX_5(argument3);
  for (I = 0; I < 2; I += 1) {
    result1.32[I] = _SATU_32(_ZX_32(argument2.32[I]) << shift)
  };
stage E1:
  GPR[registerWe<<1] = result1;
\end{lstlisting}}
\newpage

\paragraph{Encoding} Format \textsf{ALU\_WPSWRIO} (Section~\ref{sec:format-ALU-WPSWRIO}), \verb|exucode1 = 101|\\

\bitfields{ 1/P, 2/11, 1/1, 1/1, 3/101, 5/registerWo, 1/1, 2/10, 3/100, 1/0, 6/unsigned6, 5/registerZo, 1/1 }


\paragraph{Syntax} \texttt{sluswp} \emph{registerWo} = \emph{registerZo}, \emph{unsigned6}

\paragraph{Description} The \emph{registerZo} is interpreted as two words packed into 64 bits. The \emph{registerZo} words are shifted left by the \emph{unsigned6} modulo 32. The two resulting words are unsigned saturated to 32 bits, packed and stored into the \emph{registerWo}.


\paragraph{Execution} {Reservation table \textsf{ALU\_LITE} (Section~\ref{sec:reservation-ALU-LITE})}
~
{\begin{lstlisting}
stage RR:
  new argument3 = _ZX_6(unsigned6);
  new argument2 = GPR[(registerZo<<1)+1];
  new shift = _ZX_5(argument3);
  for (I = 0; I < 2; I += 1) {
    result1.32[I] = _SATU_32(_ZX_32(argument2.32[I]) << shift)
  };
stage E1:
  GPR[(registerWo<<1)+1] = result1;
\end{lstlisting}}
\newpage

\paragraph{Encoding} Format \textsf{ALU\_WPSWRRE} (Section~\ref{sec:format-ALU-WPSWRRE}), \verb|exucode1 = 101|\\

\bitfields{ 1/P, 2/11, 1/1, 1/0, 3/101, 5/registerWe, 1/0, 2/10, 3/100, 1/0, 6/registerY, 5/registerZe, 1/1 }


\paragraph{Syntax} \texttt{sluswp} \emph{registerWe} = \emph{registerZe}, \emph{registerY}

\paragraph{Description} The \emph{registerZe} is interpreted as two words packed into 64 bits. The \emph{registerZe} words are shifted left by the \emph{registerY} modulo 32. The two resulting words are unsigned saturated to 32 bits, packed and stored into the \emph{registerWe}.


\paragraph{Execution} {Reservation table \textsf{ALU\_LITE} (Section~\ref{sec:reservation-ALU-LITE})}
~
{\begin{lstlisting}
stage RR:
  new argument3 = GPR[registerY];
  new argument2 = GPR[registerZe<<1];
  new shift = _ZX_5(argument3);
  for (I = 0; I < 2; I += 1) {
    result1.32[I] = _SATU_32(_ZX_32(argument2.32[I]) << shift)
  };
stage E1:
  GPR[registerWe<<1] = result1;
\end{lstlisting}}
\newpage

\paragraph{Encoding} Format \textsf{ALU\_WPSWRRO} (Section~\ref{sec:format-ALU-WPSWRRO}), \verb|exucode1 = 101|\\

\bitfields{ 1/P, 2/11, 1/1, 1/0, 3/101, 5/registerWo, 1/1, 2/10, 3/100, 1/0, 6/registerY, 5/registerZo, 1/1 }


\paragraph{Syntax} \texttt{sluswp} \emph{registerWo} = \emph{registerZo}, \emph{registerY}

\paragraph{Description} The \emph{registerZo} is interpreted as two words packed into 64 bits. The \emph{registerZo} words are shifted left by the \emph{registerY} modulo 32. The two resulting words are unsigned saturated to 32 bits, packed and stored into the \emph{registerWo}.


\paragraph{Execution} {Reservation table \textsf{ALU\_LITE} (Section~\ref{sec:reservation-ALU-LITE})}
~
{\begin{lstlisting}
stage RR:
  new argument3 = GPR[registerY];
  new argument2 = GPR[(registerZo<<1)+1];
  new shift = _ZX_5(argument3);
  for (I = 0; I < 2; I += 1) {
    result1.32[I] = _SATU_32(_ZX_32(argument2.32[I]) << shift)
  };
stage E1:
  GPR[(registerWo<<1)+1] = result1;
\end{lstlisting}}
\newpage
\subsubsection[{\small SRABO: Shift Right Arithmetic Byte Octuple}]{SRABO\hfill Shift Right Arithmetic Byte Octuple}
\label{instr:SRABO}\index{srabo}
 
\paragraph{Encoding} Format \textsf{ALU\_BOSWRIE} (Section~\ref{sec:format-ALU-BOSWRIE}), \verb|exucode1 = 010|\\

\bitfields{ 1/P, 2/11, 1/1, 1/1, 3/010, 5/registerWe, 1/0, 2/10, 3/100, 1/1, 6/unsigned6, 5/registerZe, 1/1 }


\paragraph{Syntax} \texttt{srabo} \emph{registerWe} = \emph{registerZe}, \emph{unsigned6}

\paragraph{Description} The \emph{registerZe} is interpreted as eight bytes packed into 64 bits. The \emph{registerZe} bytes are shifted right by the \emph{unsigned6} modulo 8. The leftmost bit are replicated to fill in the vacant positions, thereby extending the signs of the shifted operands. The eight resulting bytes are packed and stored into the \emph{registerWe}.


\paragraph{Execution} {Reservation table \textsf{ALU\_LITE} (Section~\ref{sec:reservation-ALU-LITE})}
~
{\begin{lstlisting}
stage RR:
  new argument3 = _ZX_6(unsigned6);
  new argument2 = GPR[registerZe<<1];
  new shift = _ZX_3(argument3);
  for (I = 0; I < 8; I += 1) {
    result1.8[I] = _SX_8(argument2.8[I]) >> shift
  };
stage E1:
  GPR[registerWe<<1] = result1;
\end{lstlisting}}
\newpage

\paragraph{Encoding} Format \textsf{ALU\_BOSWRIO} (Section~\ref{sec:format-ALU-BOSWRIO}), \verb|exucode1 = 010|\\

\bitfields{ 1/P, 2/11, 1/1, 1/1, 3/010, 5/registerWo, 1/1, 2/10, 3/100, 1/1, 6/unsigned6, 5/registerZo, 1/1 }


\paragraph{Syntax} \texttt{srabo} \emph{registerWo} = \emph{registerZo}, \emph{unsigned6}

\paragraph{Description} The \emph{registerZo} is interpreted as eight bytes packed into 64 bits. The \emph{registerZo} bytes are shifted right by the \emph{unsigned6} modulo 8. The leftmost bit are replicated to fill in the vacant positions, thereby extending the signs of the shifted operands. The eight resulting bytes are packed and stored into the \emph{registerWo}.


\paragraph{Execution} {Reservation table \textsf{ALU\_LITE} (Section~\ref{sec:reservation-ALU-LITE})}
~
{\begin{lstlisting}
stage RR:
  new argument3 = _ZX_6(unsigned6);
  new argument2 = GPR[(registerZo<<1)+1];
  new shift = _ZX_3(argument3);
  for (I = 0; I < 8; I += 1) {
    result1.8[I] = _SX_8(argument2.8[I]) >> shift
  };
stage E1:
  GPR[(registerWo<<1)+1] = result1;
\end{lstlisting}}
\newpage

\paragraph{Encoding} Format \textsf{ALU\_BOSWRRE} (Section~\ref{sec:format-ALU-BOSWRRE}), \verb|exucode1 = 010|\\

\bitfields{ 1/P, 2/11, 1/1, 1/0, 3/010, 5/registerWe, 1/0, 2/10, 3/100, 1/1, 6/registerY, 5/registerZe, 1/1 }


\paragraph{Syntax} \texttt{srabo} \emph{registerWe} = \emph{registerZe}, \emph{registerY}

\paragraph{Description} The \emph{registerZe} is interpreted as eight bytes packed into 64 bits. The \emph{registerZe} bytes are shifted right by the \emph{registerY} modulo 8. The leftmost bit are replicated to fill in the vacant positions, thereby extending the signs of the shifted operands. The eight resulting bytes are packed and stored into the \emph{registerWe}.


\paragraph{Execution} {Reservation table \textsf{ALU\_LITE} (Section~\ref{sec:reservation-ALU-LITE})}
~
{\begin{lstlisting}
stage RR:
  new argument3 = GPR[registerY];
  new argument2 = GPR[registerZe<<1];
  new shift = _ZX_3(argument3);
  for (I = 0; I < 8; I += 1) {
    result1.8[I] = _SX_8(argument2.8[I]) >> shift
  };
stage E1:
  GPR[registerWe<<1] = result1;
\end{lstlisting}}
\newpage

\paragraph{Encoding} Format \textsf{ALU\_BOSWRRO} (Section~\ref{sec:format-ALU-BOSWRRO}), \verb|exucode1 = 010|\\

\bitfields{ 1/P, 2/11, 1/1, 1/0, 3/010, 5/registerWo, 1/1, 2/10, 3/100, 1/1, 6/registerY, 5/registerZo, 1/1 }


\paragraph{Syntax} \texttt{srabo} \emph{registerWo} = \emph{registerZo}, \emph{registerY}

\paragraph{Description} The \emph{registerZo} is interpreted as eight bytes packed into 64 bits. The \emph{registerZo} bytes are shifted right by the \emph{registerY} modulo 8. The leftmost bit are replicated to fill in the vacant positions, thereby extending the signs of the shifted operands. The eight resulting bytes are packed and stored into the \emph{registerWo}.


\paragraph{Execution} {Reservation table \textsf{ALU\_LITE} (Section~\ref{sec:reservation-ALU-LITE})}
~
{\begin{lstlisting}
stage RR:
  new argument3 = GPR[registerY];
  new argument2 = GPR[(registerZo<<1)+1];
  new shift = _ZX_3(argument3);
  for (I = 0; I < 8; I += 1) {
    result1.8[I] = _SX_8(argument2.8[I]) >> shift
  };
stage E1:
  GPR[(registerWo<<1)+1] = result1;
\end{lstlisting}}
\newpage
\subsubsection[{\small SRAD: Shift Right Arithmetic Double Word}]{SRAD\hfill Shift Right Arithmetic Double Word}
\label{instr:SRAD}\index{srad}
 
\paragraph{Encoding} Format \textsf{ALU\_DSWRI} (Section~\ref{sec:format-ALU-DSWRI}), \verb|exucode1 = 010|\\

\bitfields{ 1/P, 2/11, 1/0, 1/1, 3/010, 6/registerW, 2/10, 3/100, 1/0, 6/unsigned6, 6/registerZ }


\paragraph{Syntax} \texttt{srad} \emph{registerW} = \emph{registerZ}, \emph{unsigned6}

\paragraph{Description} The \emph{registerZ} is shifted right by the \emph{unsigned6} modulo 64. The leftmost bit is replicated to fill in the vacant positions, thereby extending the sign of the shifted operand. The result is stored into the \emph{registerW}.


\paragraph{Execution} {Reservation table \textsf{ALU\_TINY} (Section~\ref{sec:reservation-ALU-TINY})}
~
{\begin{lstlisting}
stage RR:
  new argument3 = _ZX_6(unsigned6);
  new argument2 = GPR[registerZ];
  new result1 = _SX_64(argument2) >> _ZX_6(argument3);
stage E1:
  GPR[registerW] = result1;
\end{lstlisting}}
\newpage

\paragraph{Encoding} Format \textsf{ALU\_DSWRR} (Section~\ref{sec:format-ALU-DSWRR}), \verb|exucode1 = 010|\\

\bitfields{ 1/P, 2/11, 1/0, 1/0, 3/010, 6/registerW, 2/10, 3/100, 1/0, 6/registerY, 6/registerZ }


\paragraph{Syntax} \texttt{srad} \emph{registerW} = \emph{registerZ}, \emph{registerY}

\paragraph{Description} The \emph{registerZ} is shifted right by the \emph{registerY} modulo 64. The leftmost bit is replicated to fill in the vacant positions, thereby extending the sign of the shifted operand. The result is stored into the \emph{registerW}.


\paragraph{Execution} {Reservation table \textsf{ALU\_TINY} (Section~\ref{sec:reservation-ALU-TINY})}
~
{\begin{lstlisting}
stage RR:
  new argument3 = GPR[registerY];
  new argument2 = GPR[registerZ];
  new result1 = _SX_64(argument2) >> _ZX_6(argument3);
stage E1:
  GPR[registerW] = result1;
\end{lstlisting}}
\newpage
\subsubsection[{\small SRAHQ: Shift Right Arithmetic Half Word Quadruple}]{SRAHQ\hfill Shift Right Arithmetic Half Word Quadruple}
\label{instr:SRAHQ}\index{srahq}
 
\paragraph{Encoding} Format \textsf{ALU\_HQSWRIE} (Section~\ref{sec:format-ALU-HQSWRIE}), \verb|exucode1 = 010|\\

\bitfields{ 1/P, 2/11, 1/1, 1/1, 3/010, 5/registerWe, 1/0, 2/10, 3/100, 1/1, 6/unsigned6, 5/registerZe, 1/0 }


\paragraph{Syntax} \texttt{srahq} \emph{registerWe} = \emph{registerZe}, \emph{unsigned6}

\paragraph{Description} The \emph{registerZe} is interpreted as four half words packed into 128 bits. The \emph{registerZe} half words are shifted right by the \emph{unsigned6} modulo 16. The leftmost bit are replicated to fill in the vacant positions, thereby extending the signs of the shifted operands. The four resulting half words are packed and stored into the \emph{registerWe}.


\paragraph{Execution} {Reservation table \textsf{ALU\_LITE} (Section~\ref{sec:reservation-ALU-LITE})}
~
{\begin{lstlisting}
stage RR:
  new argument3 = _ZX_6(unsigned6);
  new argument2 = GPR[registerZe<<1];
  new shift = _ZX_4(argument3);
  for (I = 0; I < 4; I += 1) {
    result1.16[I] = _SX_16(argument2.16[I]) >> shift
  };
stage E1:
  GPR[registerWe<<1] = result1;
\end{lstlisting}}
\newpage

\paragraph{Encoding} Format \textsf{ALU\_HQSWRIO} (Section~\ref{sec:format-ALU-HQSWRIO}), \verb|exucode1 = 010|\\

\bitfields{ 1/P, 2/11, 1/1, 1/1, 3/010, 5/registerWo, 1/1, 2/10, 3/100, 1/1, 6/unsigned6, 5/registerZo, 1/0 }


\paragraph{Syntax} \texttt{srahq} \emph{registerWo} = \emph{registerZo}, \emph{unsigned6}

\paragraph{Description} The \emph{registerZo} is interpreted as four half words packed into 128 bits. The \emph{registerZo} half words are shifted right by the \emph{unsigned6} modulo 16. The leftmost bit are replicated to fill in the vacant positions, thereby extending the signs of the shifted operands. The four resulting half words are packed and stored into the \emph{registerWo}.


\paragraph{Execution} {Reservation table \textsf{ALU\_LITE} (Section~\ref{sec:reservation-ALU-LITE})}
~
{\begin{lstlisting}
stage RR:
  new argument3 = _ZX_6(unsigned6);
  new argument2 = GPR[(registerZo<<1)+1];
  new shift = _ZX_4(argument3);
  for (I = 0; I < 4; I += 1) {
    result1.16[I] = _SX_16(argument2.16[I]) >> shift
  };
stage E1:
  GPR[(registerWo<<1)+1] = result1;
\end{lstlisting}}
\newpage

\paragraph{Encoding} Format \textsf{ALU\_HQSWRRE} (Section~\ref{sec:format-ALU-HQSWRRE}), \verb|exucode1 = 010|\\

\bitfields{ 1/P, 2/11, 1/1, 1/0, 3/010, 5/registerWe, 1/0, 2/10, 3/100, 1/1, 6/registerY, 5/registerZe, 1/0 }


\paragraph{Syntax} \texttt{srahq} \emph{registerWe} = \emph{registerZe}, \emph{registerY}

\paragraph{Description} The \emph{registerZe} is interpreted as four half words packed into 128 bits. The \emph{registerZe} half words are shifted right by the \emph{registerY} modulo 16. The leftmost bit are replicated to fill in the vacant positions, thereby extending the signs of the shifted operands. The four resulting half words are packed and stored into the \emph{registerWe}.


\paragraph{Execution} {Reservation table \textsf{ALU\_LITE} (Section~\ref{sec:reservation-ALU-LITE})}
~
{\begin{lstlisting}
stage RR:
  new argument3 = GPR[registerY];
  new argument2 = GPR[registerZe<<1];
  new shift = _ZX_4(argument3);
  for (I = 0; I < 4; I += 1) {
    result1.16[I] = _SX_16(argument2.16[I]) >> shift
  };
stage E1:
  GPR[registerWe<<1] = result1;
\end{lstlisting}}
\newpage

\paragraph{Encoding} Format \textsf{ALU\_HQSWRRO} (Section~\ref{sec:format-ALU-HQSWRRO}), \verb|exucode1 = 010|\\

\bitfields{ 1/P, 2/11, 1/1, 1/0, 3/010, 5/registerWo, 1/1, 2/10, 3/100, 1/1, 6/registerY, 5/registerZo, 1/0 }


\paragraph{Syntax} \texttt{srahq} \emph{registerWo} = \emph{registerZo}, \emph{registerY}

\paragraph{Description} The \emph{registerZo} is interpreted as four half words packed into 128 bits. The \emph{registerZo} half words are shifted right by the \emph{registerY} modulo 16. The leftmost bit are replicated to fill in the vacant positions, thereby extending the signs of the shifted operands. The four resulting half words are packed and stored into the \emph{registerWo}.


\paragraph{Execution} {Reservation table \textsf{ALU\_LITE} (Section~\ref{sec:reservation-ALU-LITE})}
~
{\begin{lstlisting}
stage RR:
  new argument3 = GPR[registerY];
  new argument2 = GPR[(registerZo<<1)+1];
  new shift = _ZX_4(argument3);
  for (I = 0; I < 4; I += 1) {
    result1.16[I] = _SX_16(argument2.16[I]) >> shift
  };
stage E1:
  GPR[(registerWo<<1)+1] = result1;
\end{lstlisting}}
\newpage
\subsubsection[{\small SRAW: Shift Right Arithmetic Word}]{SRAW\hfill Shift Right Arithmetic Word}
\label{instr:SRAW}\index{sraw}
 
\paragraph{Encoding} Format \textsf{ALU\_WSWRR} (Section~\ref{sec:format-ALU-WSWRR}), \verb|exucode1 = 010|\\

\bitfields{ 1/P, 2/11, 1/0, 1/0, 3/010, 6/registerW, 2/11, 3/100, 1/signextw, 6/registerY, 6/registerZ }


\paragraph{Syntax} \texttt{sraw}\emph{signextw} \emph{registerW} = \emph{registerZ}, \emph{registerY}

\paragraph{Description} The \emph{registerZ} is shifted right by the \emph{registerY} modulo 32. The leftmost bit is replicated to fill in the vacant positions, thereby extending the sign of the shifted operand. The result with the 32 upper bits cleared is stored into the \emph{registerW}.


\paragraph{Modifier(s)}
\noindent\begin{itemize}
\item signextw (cf. signextw in section \ref{par:modifier-signextw} on page \pageref{par:modifier-signextw})
\end{itemize}

\paragraph{Execution} {Reservation table \textsf{ALU\_TINY} (Section~\ref{sec:reservation-ALU-TINY})}
~
{\begin{lstlisting}
stage ID:
  new signextw = _ZX_1(signextw);
stage RR:
  new argument3 = GPR[registerY];
  new argument2 = GPR[registerZ];
  new result1 = _SX_32(argument2) >> _ZX_5(argument3);
stage E1:
  GPR[registerW] = signextw ? _SX_32(result1) : _ZX_32(result1);
\end{lstlisting}}
\newpage

\paragraph{Encoding} Format \textsf{ALU\_WSWRI} (Section~\ref{sec:format-ALU-WSWRI}), \verb|exucode1 = 010|\\

\bitfields{ 1/P, 2/11, 1/0, 1/1, 3/010, 6/registerW, 2/11, 3/100, 1/signextw, 6/unsigned6, 6/registerZ }


\paragraph{Syntax} \texttt{sraw}\emph{signextw} \emph{registerW} = \emph{registerZ}, \emph{unsigned6}

\paragraph{Description} The \emph{registerZ} is shifted right by the \emph{unsigned6} modulo 32. The leftmost bit is replicated to fill in the vacant positions, thereby extending the sign of the shifted operand. The result with the 32 upper bits cleared is stored into the \emph{registerW}.


\paragraph{Modifier(s)}
\noindent\begin{itemize}
\item signextw (cf. signextw in section \ref{par:modifier-signextw} on page \pageref{par:modifier-signextw})
\end{itemize}

\paragraph{Execution} {Reservation table \textsf{ALU\_TINY} (Section~\ref{sec:reservation-ALU-TINY})}
~
{\begin{lstlisting}
stage ID:
  new signextw = _ZX_1(signextw);
stage RR:
  new argument3 = _ZX_6(unsigned6);
  new argument2 = GPR[registerZ];
  new result1 = _SX_32(argument2) >> _ZX_5(argument3);
stage E1:
  GPR[registerW] = signextw ? _SX_32(result1) : _ZX_32(result1);
\end{lstlisting}}
\newpage
\subsubsection[{\small SRAWP: Shift Right Arithmetic Word Pair}]{SRAWP\hfill Shift Right Arithmetic Word Pair}
\label{instr:SRAWP}\index{srawp}
 
\paragraph{Encoding} Format \textsf{ALU\_WPSWRIE} (Section~\ref{sec:format-ALU-WPSWRIE}), \verb|exucode1 = 010|\\

\bitfields{ 1/P, 2/11, 1/1, 1/1, 3/010, 5/registerWe, 1/0, 2/10, 3/100, 1/0, 6/unsigned6, 5/registerZe, 1/1 }


\paragraph{Syntax} \texttt{srawp} \emph{registerWe} = \emph{registerZe}, \emph{unsigned6}

\paragraph{Description} The \emph{registerZe} is interpreted as two words packed into 64 bits. The \emph{registerZe} words are shifted right by the \emph{unsigned6} modulo 32. The leftmost bit are replicated to fill in the vacant positions, thereby extending the signs of the shifted operands. The two resulting words are packed and stored into the \emph{registerWe}.


\paragraph{Execution} {Reservation table \textsf{ALU\_LITE} (Section~\ref{sec:reservation-ALU-LITE})}
~
{\begin{lstlisting}
stage RR:
  new argument3 = _ZX_6(unsigned6);
  new argument2 = GPR[registerZe<<1];
  new shift = _ZX_5(argument3);
  for (I = 0; I < 2; I += 1) {
    result1.32[I] = _SX_32(argument2.32[I]) >> shift
  };
stage E1:
  GPR[registerWe<<1] = result1;
\end{lstlisting}}
\newpage

\paragraph{Encoding} Format \textsf{ALU\_WPSWRIO} (Section~\ref{sec:format-ALU-WPSWRIO}), \verb|exucode1 = 010|\\

\bitfields{ 1/P, 2/11, 1/1, 1/1, 3/010, 5/registerWo, 1/1, 2/10, 3/100, 1/0, 6/unsigned6, 5/registerZo, 1/1 }


\paragraph{Syntax} \texttt{srawp} \emph{registerWo} = \emph{registerZo}, \emph{unsigned6}

\paragraph{Description} The \emph{registerZo} is interpreted as two words packed into 64 bits. The \emph{registerZo} words are shifted right by the \emph{unsigned6} modulo 32. The leftmost bit are replicated to fill in the vacant positions, thereby extending the signs of the shifted operands. The two resulting words are packed and stored into the \emph{registerWo}.


\paragraph{Execution} {Reservation table \textsf{ALU\_LITE} (Section~\ref{sec:reservation-ALU-LITE})}
~
{\begin{lstlisting}
stage RR:
  new argument3 = _ZX_6(unsigned6);
  new argument2 = GPR[(registerZo<<1)+1];
  new shift = _ZX_5(argument3);
  for (I = 0; I < 2; I += 1) {
    result1.32[I] = _SX_32(argument2.32[I]) >> shift
  };
stage E1:
  GPR[(registerWo<<1)+1] = result1;
\end{lstlisting}}
\newpage

\paragraph{Encoding} Format \textsf{ALU\_WPSWRRE} (Section~\ref{sec:format-ALU-WPSWRRE}), \verb|exucode1 = 010|\\

\bitfields{ 1/P, 2/11, 1/1, 1/0, 3/010, 5/registerWe, 1/0, 2/10, 3/100, 1/0, 6/registerY, 5/registerZe, 1/1 }


\paragraph{Syntax} \texttt{srawp} \emph{registerWe} = \emph{registerZe}, \emph{registerY}

\paragraph{Description} The \emph{registerZe} is interpreted as two words packed into 64 bits. The \emph{registerZe} words are shifted right by the \emph{registerY} modulo 32. The leftmost bit are replicated to fill in the vacant positions, thereby extending the signs of the shifted operands. The two resulting words are packed and stored into the \emph{registerWe}.


\paragraph{Execution} {Reservation table \textsf{ALU\_LITE} (Section~\ref{sec:reservation-ALU-LITE})}
~
{\begin{lstlisting}
stage RR:
  new argument3 = GPR[registerY];
  new argument2 = GPR[registerZe<<1];
  new shift = _ZX_5(argument3);
  for (I = 0; I < 2; I += 1) {
    result1.32[I] = _SX_32(argument2.32[I]) >> shift
  };
stage E1:
  GPR[registerWe<<1] = result1;
\end{lstlisting}}
\newpage

\paragraph{Encoding} Format \textsf{ALU\_WPSWRRO} (Section~\ref{sec:format-ALU-WPSWRRO}), \verb|exucode1 = 010|\\

\bitfields{ 1/P, 2/11, 1/1, 1/0, 3/010, 5/registerWo, 1/1, 2/10, 3/100, 1/0, 6/registerY, 5/registerZo, 1/1 }


\paragraph{Syntax} \texttt{srawp} \emph{registerWo} = \emph{registerZo}, \emph{registerY}

\paragraph{Description} The \emph{registerZo} is interpreted as two words packed into 64 bits. The \emph{registerZo} words are shifted right by the \emph{registerY} modulo 32. The leftmost bit are replicated to fill in the vacant positions, thereby extending the signs of the shifted operands. The two resulting words are packed and stored into the \emph{registerWo}.


\paragraph{Execution} {Reservation table \textsf{ALU\_LITE} (Section~\ref{sec:reservation-ALU-LITE})}
~
{\begin{lstlisting}
stage RR:
  new argument3 = GPR[registerY];
  new argument2 = GPR[(registerZo<<1)+1];
  new shift = _ZX_5(argument3);
  for (I = 0; I < 2; I += 1) {
    result1.32[I] = _SX_32(argument2.32[I]) >> shift
  };
stage E1:
  GPR[(registerWo<<1)+1] = result1;
\end{lstlisting}}
\newpage
\subsubsection[{\small SRLBO: Shift Right Logical Byte Octuple}]{SRLBO\hfill Shift Right Logical Byte Octuple}
\label{instr:SRLBO}\index{srlbo}
 
\paragraph{Encoding} Format \textsf{ALU\_BOSWRIE} (Section~\ref{sec:format-ALU-BOSWRIE}), \verb|exucode1 = 011|\\

\bitfields{ 1/P, 2/11, 1/1, 1/1, 3/011, 5/registerWe, 1/0, 2/10, 3/100, 1/1, 6/unsigned6, 5/registerZe, 1/1 }


\paragraph{Syntax} \texttt{srlbo} \emph{registerWe} = \emph{registerZe}, \emph{unsigned6}

\paragraph{Description} The \emph{registerZe} is interpreted as eight bytes packed into 64 bits. The \emph{registerZe} bytes are shifted right by the \emph{unsigned6} modulo 8. The vacant positions are filled in with zeros. The eight resulting bytes are packed and stored into the \emph{registerWe}.


\paragraph{Execution} {Reservation table \textsf{ALU\_LITE} (Section~\ref{sec:reservation-ALU-LITE})}
~
{\begin{lstlisting}
stage RR:
  new argument3 = _ZX_6(unsigned6);
  new argument2 = GPR[registerZe<<1];
  new shift = _ZX_3(argument3);
  for (I = 0; I < 8; I += 1) {
    result1.8[I] = _ZX_8(argument2.8[I]) >> shift
  };
stage E1:
  GPR[registerWe<<1] = result1;
\end{lstlisting}}
\newpage

\paragraph{Encoding} Format \textsf{ALU\_BOSWRIO} (Section~\ref{sec:format-ALU-BOSWRIO}), \verb|exucode1 = 011|\\

\bitfields{ 1/P, 2/11, 1/1, 1/1, 3/011, 5/registerWo, 1/1, 2/10, 3/100, 1/1, 6/unsigned6, 5/registerZo, 1/1 }


\paragraph{Syntax} \texttt{srlbo} \emph{registerWo} = \emph{registerZo}, \emph{unsigned6}

\paragraph{Description} The \emph{registerZo} is interpreted as eight bytes packed into 64 bits. The \emph{registerZo} bytes are shifted right by the \emph{unsigned6} modulo 8. The vacant positions are filled in with zeros. The eight resulting bytes are packed and stored into the \emph{registerWo}.


\paragraph{Execution} {Reservation table \textsf{ALU\_LITE} (Section~\ref{sec:reservation-ALU-LITE})}
~
{\begin{lstlisting}
stage RR:
  new argument3 = _ZX_6(unsigned6);
  new argument2 = GPR[(registerZo<<1)+1];
  new shift = _ZX_3(argument3);
  for (I = 0; I < 8; I += 1) {
    result1.8[I] = _ZX_8(argument2.8[I]) >> shift
  };
stage E1:
  GPR[(registerWo<<1)+1] = result1;
\end{lstlisting}}
\newpage

\paragraph{Encoding} Format \textsf{ALU\_BOSWRRE} (Section~\ref{sec:format-ALU-BOSWRRE}), \verb|exucode1 = 011|\\

\bitfields{ 1/P, 2/11, 1/1, 1/0, 3/011, 5/registerWe, 1/0, 2/10, 3/100, 1/1, 6/registerY, 5/registerZe, 1/1 }


\paragraph{Syntax} \texttt{srlbo} \emph{registerWe} = \emph{registerZe}, \emph{registerY}

\paragraph{Description} The \emph{registerZe} is interpreted as eight bytes packed into 64 bits. The \emph{registerZe} bytes are shifted right by the \emph{registerY} modulo 8. The vacant positions are filled in with zeros. The eight resulting bytes are packed and stored into the \emph{registerWe}.


\paragraph{Execution} {Reservation table \textsf{ALU\_LITE} (Section~\ref{sec:reservation-ALU-LITE})}
~
{\begin{lstlisting}
stage RR:
  new argument3 = GPR[registerY];
  new argument2 = GPR[registerZe<<1];
  new shift = _ZX_3(argument3);
  for (I = 0; I < 8; I += 1) {
    result1.8[I] = _ZX_8(argument2.8[I]) >> shift
  };
stage E1:
  GPR[registerWe<<1] = result1;
\end{lstlisting}}
\newpage

\paragraph{Encoding} Format \textsf{ALU\_BOSWRRO} (Section~\ref{sec:format-ALU-BOSWRRO}), \verb|exucode1 = 011|\\

\bitfields{ 1/P, 2/11, 1/1, 1/0, 3/011, 5/registerWo, 1/1, 2/10, 3/100, 1/1, 6/registerY, 5/registerZo, 1/1 }


\paragraph{Syntax} \texttt{srlbo} \emph{registerWo} = \emph{registerZo}, \emph{registerY}

\paragraph{Description} The \emph{registerZo} is interpreted as eight bytes packed into 64 bits. The \emph{registerZo} bytes are shifted right by the \emph{registerY} modulo 8. The vacant positions are filled in with zeros. The eight resulting bytes are packed and stored into the \emph{registerWo}.


\paragraph{Execution} {Reservation table \textsf{ALU\_LITE} (Section~\ref{sec:reservation-ALU-LITE})}
~
{\begin{lstlisting}
stage RR:
  new argument3 = GPR[registerY];
  new argument2 = GPR[(registerZo<<1)+1];
  new shift = _ZX_3(argument3);
  for (I = 0; I < 8; I += 1) {
    result1.8[I] = _ZX_8(argument2.8[I]) >> shift
  };
stage E1:
  GPR[(registerWo<<1)+1] = result1;
\end{lstlisting}}
\newpage
\subsubsection[{\small SRLD: Shift Right Logical Double Word}]{SRLD\hfill Shift Right Logical Double Word}
\label{instr:SRLD}\index{srld}
 
\paragraph{Encoding} Format \textsf{ALU\_DSWRI} (Section~\ref{sec:format-ALU-DSWRI}), \verb|exucode1 = 011|\\

\bitfields{ 1/P, 2/11, 1/0, 1/1, 3/011, 6/registerW, 2/10, 3/100, 1/0, 6/unsigned6, 6/registerZ }


\paragraph{Syntax} \texttt{srld} \emph{registerW} = \emph{registerZ}, \emph{unsigned6}

\paragraph{Description} The \emph{registerZ} is shifted right by the \emph{unsigned6} modulo 64. The vacant positions are filled in with zeros. The result is stored into the \emph{registerW}.


\paragraph{Execution} {Reservation table \textsf{ALU\_TINY} (Section~\ref{sec:reservation-ALU-TINY})}
~
{\begin{lstlisting}
stage RR:
  new argument3 = _ZX_6(unsigned6);
  new argument2 = GPR[registerZ];
  new result1 = _ZX_64(argument2) >> _ZX_6(argument3);
stage E1:
  GPR[registerW] = result1;
\end{lstlisting}}
\newpage

\paragraph{Encoding} Format \textsf{ALU\_DSWRR} (Section~\ref{sec:format-ALU-DSWRR}), \verb|exucode1 = 011|\\

\bitfields{ 1/P, 2/11, 1/0, 1/0, 3/011, 6/registerW, 2/10, 3/100, 1/0, 6/registerY, 6/registerZ }


\paragraph{Syntax} \texttt{srld} \emph{registerW} = \emph{registerZ}, \emph{registerY}

\paragraph{Description} The \emph{registerZ} is shifted right by the \emph{registerY} modulo 64. The vacant positions are filled in with zeros. The result is stored into the \emph{registerW}.


\paragraph{Execution} {Reservation table \textsf{ALU\_TINY} (Section~\ref{sec:reservation-ALU-TINY})}
~
{\begin{lstlisting}
stage RR:
  new argument3 = GPR[registerY];
  new argument2 = GPR[registerZ];
  new result1 = _ZX_64(argument2) >> _ZX_6(argument3);
stage E1:
  GPR[registerW] = result1;
\end{lstlisting}}
\newpage
\subsubsection[{\small SRLHQ: Shift Right Logical Half Word Quadruple}]{SRLHQ\hfill Shift Right Logical Half Word Quadruple}
\label{instr:SRLHQ}\index{srlhq}
 
\paragraph{Encoding} Format \textsf{ALU\_HQSWRIE} (Section~\ref{sec:format-ALU-HQSWRIE}), \verb|exucode1 = 011|\\

\bitfields{ 1/P, 2/11, 1/1, 1/1, 3/011, 5/registerWe, 1/0, 2/10, 3/100, 1/1, 6/unsigned6, 5/registerZe, 1/0 }


\paragraph{Syntax} \texttt{srlhq} \emph{registerWe} = \emph{registerZe}, \emph{unsigned6}

\paragraph{Description} The \emph{registerZe} is interpreted as four half words packed into 64 bits. The \emph{registerZe} half words are shifted right by the \emph{unsigned6} modulo 16. The vacant positions are filled in with zeros. The four resulting half words are packed and stored into the \emph{registerWe}.


\paragraph{Execution} {Reservation table \textsf{ALU\_LITE} (Section~\ref{sec:reservation-ALU-LITE})}
~
{\begin{lstlisting}
stage RR:
  new argument3 = _ZX_6(unsigned6);
  new argument2 = GPR[registerZe<<1];
  new shift = _ZX_4(argument3);
  for (I = 0; I < 4; I += 1) {
    result1.16[I] = _ZX_16(argument2.16[I]) >> shift
  };
stage E1:
  GPR[registerWe<<1] = result1;
\end{lstlisting}}
\newpage

\paragraph{Encoding} Format \textsf{ALU\_HQSWRIO} (Section~\ref{sec:format-ALU-HQSWRIO}), \verb|exucode1 = 011|\\

\bitfields{ 1/P, 2/11, 1/1, 1/1, 3/011, 5/registerWo, 1/1, 2/10, 3/100, 1/1, 6/unsigned6, 5/registerZo, 1/0 }


\paragraph{Syntax} \texttt{srlhq} \emph{registerWo} = \emph{registerZo}, \emph{unsigned6}

\paragraph{Description} The \emph{registerZo} is interpreted as four half words packed into 64 bits. The \emph{registerZo} half words are shifted right by the \emph{unsigned6} modulo 16. The vacant positions are filled in with zeros. The four resulting half words are packed and stored into the \emph{registerWo}.


\paragraph{Execution} {Reservation table \textsf{ALU\_LITE} (Section~\ref{sec:reservation-ALU-LITE})}
~
{\begin{lstlisting}
stage RR:
  new argument3 = _ZX_6(unsigned6);
  new argument2 = GPR[(registerZo<<1)+1];
  new shift = _ZX_4(argument3);
  for (I = 0; I < 4; I += 1) {
    result1.16[I] = _ZX_16(argument2.16[I]) >> shift
  };
stage E1:
  GPR[(registerWo<<1)+1] = result1;
\end{lstlisting}}
\newpage

\paragraph{Encoding} Format \textsf{ALU\_HQSWRRE} (Section~\ref{sec:format-ALU-HQSWRRE}), \verb|exucode1 = 011|\\

\bitfields{ 1/P, 2/11, 1/1, 1/0, 3/011, 5/registerWe, 1/0, 2/10, 3/100, 1/1, 6/registerY, 5/registerZe, 1/0 }


\paragraph{Syntax} \texttt{srlhq} \emph{registerWe} = \emph{registerZe}, \emph{registerY}

\paragraph{Description} The \emph{registerZe} is interpreted as four half words packed into 64 bits. The \emph{registerZe} half words are shifted right by the \emph{registerY} modulo 16. The vacant positions are filled in with zeros. The four resulting half words are packed and stored into the \emph{registerWe}.


\paragraph{Execution} {Reservation table \textsf{ALU\_LITE} (Section~\ref{sec:reservation-ALU-LITE})}
~
{\begin{lstlisting}
stage RR:
  new argument3 = GPR[registerY];
  new argument2 = GPR[registerZe<<1];
  new shift = _ZX_4(argument3);
  for (I = 0; I < 4; I += 1) {
    result1.16[I] = _ZX_16(argument2.16[I]) >> shift
  };
stage E1:
  GPR[registerWe<<1] = result1;
\end{lstlisting}}
\newpage

\paragraph{Encoding} Format \textsf{ALU\_HQSWRRO} (Section~\ref{sec:format-ALU-HQSWRRO}), \verb|exucode1 = 011|\\

\bitfields{ 1/P, 2/11, 1/1, 1/0, 3/011, 5/registerWo, 1/1, 2/10, 3/100, 1/1, 6/registerY, 5/registerZo, 1/0 }


\paragraph{Syntax} \texttt{srlhq} \emph{registerWo} = \emph{registerZo}, \emph{registerY}

\paragraph{Description} The \emph{registerZo} is interpreted as four half words packed into 64 bits. The \emph{registerZo} half words are shifted right by the \emph{registerY} modulo 16. The vacant positions are filled in with zeros. The four resulting half words are packed and stored into the \emph{registerWo}.


\paragraph{Execution} {Reservation table \textsf{ALU\_LITE} (Section~\ref{sec:reservation-ALU-LITE})}
~
{\begin{lstlisting}
stage RR:
  new argument3 = GPR[registerY];
  new argument2 = GPR[(registerZo<<1)+1];
  new shift = _ZX_4(argument3);
  for (I = 0; I < 4; I += 1) {
    result1.16[I] = _ZX_16(argument2.16[I]) >> shift
  };
stage E1:
  GPR[(registerWo<<1)+1] = result1;
\end{lstlisting}}
\newpage
\subsubsection[{\small SRLW: Shift Right Logical Word}]{SRLW\hfill Shift Right Logical Word}
\label{instr:SRLW}\index{srlw}
 
\paragraph{Encoding} Format \textsf{ALU\_WSWRR} (Section~\ref{sec:format-ALU-WSWRR}), \verb|exucode1 = 011|\\

\bitfields{ 1/P, 2/11, 1/0, 1/0, 3/011, 6/registerW, 2/11, 3/100, 1/signextw, 6/registerY, 6/registerZ }


\paragraph{Syntax} \texttt{srlw}\emph{signextw} \emph{registerW} = \emph{registerZ}, \emph{registerY}

\paragraph{Description} The \emph{registerZ} is shifted right by the \emph{registerY} modulo 32. The vacant positions are filled in with zeros. The result with the 32 upper bits cleared is stored into the \emph{registerW}.


\paragraph{Modifier(s)}
\noindent\begin{itemize}
\item signextw (cf. signextw in section \ref{par:modifier-signextw} on page \pageref{par:modifier-signextw})
\end{itemize}

\paragraph{Execution} {Reservation table \textsf{ALU\_TINY} (Section~\ref{sec:reservation-ALU-TINY})}
~
{\begin{lstlisting}
stage ID:
  new signextw = _ZX_1(signextw);
stage RR:
  new argument3 = GPR[registerY];
  new argument2 = GPR[registerZ];
  new result1 = argument2.32[0] >> _ZX_5(argument3);
stage E1:
  GPR[registerW] = signextw ? _SX_32(result1) : _ZX_32(result1);
\end{lstlisting}}
\newpage

\paragraph{Encoding} Format \textsf{ALU\_WSWRI} (Section~\ref{sec:format-ALU-WSWRI}), \verb|exucode1 = 011|\\

\bitfields{ 1/P, 2/11, 1/0, 1/1, 3/011, 6/registerW, 2/11, 3/100, 1/signextw, 6/unsigned6, 6/registerZ }


\paragraph{Syntax} \texttt{srlw}\emph{signextw} \emph{registerW} = \emph{registerZ}, \emph{unsigned6}

\paragraph{Description} The \emph{registerZ} is shifted right by the \emph{unsigned6} modulo 32. The vacant positions are filled in with zeros. The result with the 32 upper bits cleared is stored into the \emph{registerW}.


\paragraph{Modifier(s)}
\noindent\begin{itemize}
\item signextw (cf. signextw in section \ref{par:modifier-signextw} on page \pageref{par:modifier-signextw})
\end{itemize}

\paragraph{Execution} {Reservation table \textsf{ALU\_TINY} (Section~\ref{sec:reservation-ALU-TINY})}
~
{\begin{lstlisting}
stage ID:
  new signextw = _ZX_1(signextw);
stage RR:
  new argument3 = _ZX_6(unsigned6);
  new argument2 = GPR[registerZ];
  new result1 = argument2.32[0] >> _ZX_5(argument3);
stage E1:
  GPR[registerW] = signextw ? _SX_32(result1) : _ZX_32(result1);
\end{lstlisting}}
\newpage
\subsubsection[{\small SRLWP: Shift Right Logical Word Pair}]{SRLWP\hfill Shift Right Logical Word Pair}
\label{instr:SRLWP}\index{srlwp}
 
\paragraph{Encoding} Format \textsf{ALU\_WPSWRIE} (Section~\ref{sec:format-ALU-WPSWRIE}), \verb|exucode1 = 011|\\

\bitfields{ 1/P, 2/11, 1/1, 1/1, 3/011, 5/registerWe, 1/0, 2/10, 3/100, 1/0, 6/unsigned6, 5/registerZe, 1/1 }


\paragraph{Syntax} \texttt{srlwp} \emph{registerWe} = \emph{registerZe}, \emph{unsigned6}

\paragraph{Description} The \emph{registerZe} is interpreted as two words packed into 64 bits. The \emph{registerZe} words are shifted right by the \emph{unsigned6} modulo 32. The vacant positions are filled in with zeros. The two resulting words are packed and stored into the \emph{registerWe}.


\paragraph{Execution} {Reservation table \textsf{ALU\_LITE} (Section~\ref{sec:reservation-ALU-LITE})}
~
{\begin{lstlisting}
stage RR:
  new argument3 = _ZX_6(unsigned6);
  new argument2 = GPR[registerZe<<1];
  new shift = _ZX_5(argument3);
  for (I = 0; I < 4; I += 1) {
    result1.32[I] = _ZX_32(argument2.32[I]) >> shift
  };
stage E1:
  GPR[registerWe<<1] = result1;
\end{lstlisting}}
\newpage

\paragraph{Encoding} Format \textsf{ALU\_WPSWRIO} (Section~\ref{sec:format-ALU-WPSWRIO}), \verb|exucode1 = 011|\\

\bitfields{ 1/P, 2/11, 1/1, 1/1, 3/011, 5/registerWo, 1/1, 2/10, 3/100, 1/0, 6/unsigned6, 5/registerZo, 1/1 }


\paragraph{Syntax} \texttt{srlwp} \emph{registerWo} = \emph{registerZo}, \emph{unsigned6}

\paragraph{Description} The \emph{registerZo} is interpreted as two words packed into 64 bits. The \emph{registerZo} words are shifted right by the \emph{unsigned6} modulo 32. The vacant positions are filled in with zeros. The two resulting words are packed and stored into the \emph{registerWo}.


\paragraph{Execution} {Reservation table \textsf{ALU\_LITE} (Section~\ref{sec:reservation-ALU-LITE})}
~
{\begin{lstlisting}
stage RR:
  new argument3 = _ZX_6(unsigned6);
  new argument2 = GPR[(registerZo<<1)+1];
  new shift = _ZX_5(argument3);
  for (I = 0; I < 4; I += 1) {
    result1.32[I] = _ZX_32(argument2.32[I]) >> shift
  };
stage E1:
  GPR[(registerWo<<1)+1] = result1;
\end{lstlisting}}
\newpage

\paragraph{Encoding} Format \textsf{ALU\_WPSWRRE} (Section~\ref{sec:format-ALU-WPSWRRE}), \verb|exucode1 = 011|\\

\bitfields{ 1/P, 2/11, 1/1, 1/0, 3/011, 5/registerWe, 1/0, 2/10, 3/100, 1/0, 6/registerY, 5/registerZe, 1/1 }


\paragraph{Syntax} \texttt{srlwp} \emph{registerWe} = \emph{registerZe}, \emph{registerY}

\paragraph{Description} The \emph{registerZe} is interpreted as two words packed into 64 bits. The \emph{registerZe} words are shifted right by the \emph{registerY} modulo 32. The vacant positions are filled in with zeros. The two resulting words are packed and stored into the \emph{registerWe}.


\paragraph{Execution} {Reservation table \textsf{ALU\_LITE} (Section~\ref{sec:reservation-ALU-LITE})}
~
{\begin{lstlisting}
stage RR:
  new argument3 = GPR[registerY];
  new argument2 = GPR[registerZe<<1];
  new shift = _ZX_5(argument3);
  for (I = 0; I < 4; I += 1) {
    result1.32[I] = _ZX_32(argument2.32[I]) >> shift
  };
stage E1:
  GPR[registerWe<<1] = result1;
\end{lstlisting}}
\newpage

\paragraph{Encoding} Format \textsf{ALU\_WPSWRRO} (Section~\ref{sec:format-ALU-WPSWRRO}), \verb|exucode1 = 011|\\

\bitfields{ 1/P, 2/11, 1/1, 1/0, 3/011, 5/registerWo, 1/1, 2/10, 3/100, 1/0, 6/registerY, 5/registerZo, 1/1 }


\paragraph{Syntax} \texttt{srlwp} \emph{registerWo} = \emph{registerZo}, \emph{registerY}

\paragraph{Description} The \emph{registerZo} is interpreted as two words packed into 64 bits. The \emph{registerZo} words are shifted right by the \emph{registerY} modulo 32. The vacant positions are filled in with zeros. The two resulting words are packed and stored into the \emph{registerWo}.


\paragraph{Execution} {Reservation table \textsf{ALU\_LITE} (Section~\ref{sec:reservation-ALU-LITE})}
~
{\begin{lstlisting}
stage RR:
  new argument3 = GPR[registerY];
  new argument2 = GPR[(registerZo<<1)+1];
  new shift = _ZX_5(argument3);
  for (I = 0; I < 4; I += 1) {
    result1.32[I] = _ZX_32(argument2.32[I]) >> shift
  };
stage E1:
  GPR[(registerWo<<1)+1] = result1;
\end{lstlisting}}
\newpage
\subsubsection[{\small SRSBO: Shift Right Symmetric Byte Octuple}]{SRSBO\hfill Shift Right Symmetric Byte Octuple}
\label{instr:SRSBO}\index{srsbo}
 
\paragraph{Encoding} Format \textsf{ALU\_BOSWRIE} (Section~\ref{sec:format-ALU-BOSWRIE}), \verb|exucode1 = 000|\\

\bitfields{ 1/P, 2/11, 1/1, 1/1, 3/000, 5/registerWe, 1/0, 2/10, 3/100, 1/1, 6/unsigned6, 5/registerZe, 1/1 }


\paragraph{Syntax} \texttt{srsbo} \emph{registerWe} = \emph{registerZe}, \emph{unsigned6}

\paragraph{Description} The \emph{registerZe} is interpreted as eight bytes packed into 64 bits. If non-negative, each \emph{registerZe} word is shifted right by the \emph{unsigned6} modulo 8. Else, each \emph{registerZe} word is biased by (2**(\emph{unsigned6} modulo 8)) - 1 and shifted right by the \emph{unsigned6} modulo 8. Each result is stored into the corresponding \emph{registerWe} byte. This instruction implements a signed division of the \emph{registerZe} bytes by the power of two given by the \emph{unsigned6} modulo 8.


\paragraph{Execution} {Reservation table \textsf{ALU\_LITE} (Section~\ref{sec:reservation-ALU-LITE})}
~
{\begin{lstlisting}
stage RR:
  new argument3 = _ZX_6(unsigned6);
  new argument2 = GPR[registerZe<<1];
  new shift = _ZX_3(argument3);
  new bias = (1 << shift) - 1;
  for (I = 0; I < 8; I += 1) {
    result1.8[I] = (_SX_8(argument2.8[I]) + _SX_8(argument2.8[I]) < 0 ? bias : 0) >> shift
  };
stage E1:
  GPR[registerWe<<1] = result1;
\end{lstlisting}}
\newpage

\paragraph{Encoding} Format \textsf{ALU\_BOSWRIO} (Section~\ref{sec:format-ALU-BOSWRIO}), \verb|exucode1 = 000|\\

\bitfields{ 1/P, 2/11, 1/1, 1/1, 3/000, 5/registerWo, 1/1, 2/10, 3/100, 1/1, 6/unsigned6, 5/registerZo, 1/1 }


\paragraph{Syntax} \texttt{srsbo} \emph{registerWo} = \emph{registerZo}, \emph{unsigned6}

\paragraph{Description} The \emph{registerZo} is interpreted as eight bytes packed into 64 bits. If non-negative, each \emph{registerZo} word is shifted right by the \emph{unsigned6} modulo 8. Else, each \emph{registerZo} word is biased by (2**(\emph{unsigned6} modulo 8)) - 1 and shifted right by the \emph{unsigned6} modulo 8. Each result is stored into the corresponding \emph{registerWo} byte. This instruction implements a signed division of the \emph{registerZo} bytes by the power of two given by the \emph{unsigned6} modulo 8.


\paragraph{Execution} {Reservation table \textsf{ALU\_LITE} (Section~\ref{sec:reservation-ALU-LITE})}
~
{\begin{lstlisting}
stage RR:
  new argument3 = _ZX_6(unsigned6);
  new argument2 = GPR[(registerZo<<1)+1];
  new shift = _ZX_3(argument3);
  new bias = (1 << shift) - 1;
  for (I = 0; I < 8; I += 1) {
    result1.8[I] = (_SX_8(argument2.8[I]) + _SX_8(argument2.8[I]) < 0 ? bias : 0) >> shift
  };
stage E1:
  GPR[(registerWo<<1)+1] = result1;
\end{lstlisting}}
\newpage

\paragraph{Encoding} Format \textsf{ALU\_BOSWRRE} (Section~\ref{sec:format-ALU-BOSWRRE}), \verb|exucode1 = 000|\\

\bitfields{ 1/P, 2/11, 1/1, 1/0, 3/000, 5/registerWe, 1/0, 2/10, 3/100, 1/1, 6/registerY, 5/registerZe, 1/1 }


\paragraph{Syntax} \texttt{srsbo} \emph{registerWe} = \emph{registerZe}, \emph{registerY}

\paragraph{Description} The \emph{registerZe} is interpreted as eight bytes packed into 64 bits. If non-negative, each \emph{registerZe} word is shifted right by the \emph{registerY} modulo 8. Else, each \emph{registerZe} word is biased by (2**(\emph{registerY} modulo 8)) - 1 and shifted right by the \emph{registerY} modulo 8. Each result is stored into the corresponding \emph{registerWe} byte. This instruction implements a signed division of the \emph{registerZe} bytes by the power of two given by the \emph{registerY} modulo 8.


\paragraph{Execution} {Reservation table \textsf{ALU\_LITE} (Section~\ref{sec:reservation-ALU-LITE})}
~
{\begin{lstlisting}
stage RR:
  new argument3 = GPR[registerY];
  new argument2 = GPR[registerZe<<1];
  new shift = _ZX_3(argument3);
  new bias = (1 << shift) - 1;
  for (I = 0; I < 8; I += 1) {
    result1.8[I] = (_SX_8(argument2.8[I]) + _SX_8(argument2.8[I]) < 0 ? bias : 0) >> shift
  };
stage E1:
  GPR[registerWe<<1] = result1;
\end{lstlisting}}
\newpage

\paragraph{Encoding} Format \textsf{ALU\_BOSWRRO} (Section~\ref{sec:format-ALU-BOSWRRO}), \verb|exucode1 = 000|\\

\bitfields{ 1/P, 2/11, 1/1, 1/0, 3/000, 5/registerWo, 1/1, 2/10, 3/100, 1/1, 6/registerY, 5/registerZo, 1/1 }


\paragraph{Syntax} \texttt{srsbo} \emph{registerWo} = \emph{registerZo}, \emph{registerY}

\paragraph{Description} The \emph{registerZo} is interpreted as eight bytes packed into 64 bits. If non-negative, each \emph{registerZo} word is shifted right by the \emph{registerY} modulo 8. Else, each \emph{registerZo} word is biased by (2**(\emph{registerY} modulo 8)) - 1 and shifted right by the \emph{registerY} modulo 8. Each result is stored into the corresponding \emph{registerWo} byte. This instruction implements a signed division of the \emph{registerZo} bytes by the power of two given by the \emph{registerY} modulo 8.


\paragraph{Execution} {Reservation table \textsf{ALU\_LITE} (Section~\ref{sec:reservation-ALU-LITE})}
~
{\begin{lstlisting}
stage RR:
  new argument3 = GPR[registerY];
  new argument2 = GPR[(registerZo<<1)+1];
  new shift = _ZX_3(argument3);
  new bias = (1 << shift) - 1;
  for (I = 0; I < 8; I += 1) {
    result1.8[I] = (_SX_8(argument2.8[I]) + _SX_8(argument2.8[I]) < 0 ? bias : 0) >> shift
  };
stage E1:
  GPR[(registerWo<<1)+1] = result1;
\end{lstlisting}}
\newpage
\subsubsection[{\small SRSD: Shift Right Symmetric Double Word}]{SRSD\hfill Shift Right Symmetric Double Word}
\label{instr:SRSD}\index{srsd}
 
\paragraph{Encoding} Format \textsf{ALU\_DSWRI} (Section~\ref{sec:format-ALU-DSWRI}), \verb|exucode1 = 000|\\

\bitfields{ 1/P, 2/11, 1/0, 1/1, 3/000, 6/registerW, 2/10, 3/100, 1/0, 6/unsigned6, 6/registerZ }


\paragraph{Syntax} \texttt{srsd} \emph{registerW} = \emph{registerZ}, \emph{unsigned6}

\paragraph{Description} If non-negative, the \emph{registerZ} is shifted right by the \emph{unsigned6} modulo 64. Else, the \emph{registerZ} biased by (2**(\emph{unsigned6} modulo 64)) - 1 and shifted right by the \emph{unsigned6} modulo 64. The result is stored into the \emph{registerW}. This instruction implements a signed division of the \emph{registerZ} double word by a power of two.


\paragraph{Execution} {Reservation table \textsf{ALU\_LITE} (Section~\ref{sec:reservation-ALU-LITE})}
~
{\begin{lstlisting}
stage RR:
  new argument3 = _ZX_6(unsigned6);
  new argument2 = GPR[registerZ];
  new shift = _ZX_6(argument3);
  new bias = _SX_64(argument2) < 0 ? (1 << shift) - 1 : 0;
  new result1 = (_SX_64(argument2) + bias) >> shift;
stage E1:
  GPR[registerW] = result1;
\end{lstlisting}}
\newpage

\paragraph{Encoding} Format \textsf{ALU\_DSWRR} (Section~\ref{sec:format-ALU-DSWRR}), \verb|exucode1 = 000|\\

\bitfields{ 1/P, 2/11, 1/0, 1/0, 3/000, 6/registerW, 2/10, 3/100, 1/0, 6/registerY, 6/registerZ }


\paragraph{Syntax} \texttt{srsd} \emph{registerW} = \emph{registerZ}, \emph{registerY}

\paragraph{Description} If non-negative, the \emph{registerZ} is shifted right by the \emph{registerY} modulo 64. Else, the \emph{registerZ} biased by (2**(\emph{registerY} modulo 64)) - 1 and shifted right by the \emph{registerY} modulo 64. The result is stored into the \emph{registerW}. This instruction implements a signed division of the \emph{registerZ} double word by a power of two.


\paragraph{Execution} {Reservation table \textsf{ALU\_LITE} (Section~\ref{sec:reservation-ALU-LITE})}
~
{\begin{lstlisting}
stage RR:
  new argument3 = GPR[registerY];
  new argument2 = GPR[registerZ];
  new shift = _ZX_6(argument3);
  new bias = _SX_64(argument2) < 0 ? (1 << shift) - 1 : 0;
  new result1 = (_SX_64(argument2) + bias) >> shift;
stage E1:
  GPR[registerW] = result1;
\end{lstlisting}}
\newpage
\subsubsection[{\small SRSHQ: Shift Right Symmetric Half Word Quadruple}]{SRSHQ\hfill Shift Right Symmetric Half Word Quadruple}
\label{instr:SRSHQ}\index{srshq}
 
\paragraph{Encoding} Format \textsf{ALU\_HQSWRIE} (Section~\ref{sec:format-ALU-HQSWRIE}), \verb|exucode1 = 000|\\

\bitfields{ 1/P, 2/11, 1/1, 1/1, 3/000, 5/registerWe, 1/0, 2/10, 3/100, 1/1, 6/unsigned6, 5/registerZe, 1/0 }


\paragraph{Syntax} \texttt{srshq} \emph{registerWe} = \emph{registerZe}, \emph{unsigned6}

\paragraph{Description} The \emph{registerZe} is interpreted as four half words packed into 64 bits. If non-negative, each \emph{registerZe} half word is shifted right by the \emph{unsigned6} modulo 16. Else, each \emph{registerZe} half word is biased by (2**(\emph{unsigned6} modulo 16)) - 1 and shifted right by the \emph{unsigned6} modulo 16. Each result is stored into the corresponding \emph{registerWe} half word. This instruction implements a signed division of the \emph{registerZe} half words by the power of two given by the \emph{unsigned6} modulo 16.


\paragraph{Execution} {Reservation table \textsf{ALU\_LITE} (Section~\ref{sec:reservation-ALU-LITE})}
~
{\begin{lstlisting}
stage RR:
  new argument3 = _ZX_6(unsigned6);
  new argument2 = GPR[registerZe<<1];
  new shift = _ZX_4(argument3);
  new bias = (1 << shift) - 1;
  for (I = 0; I < 4; I += 1) {
    result1.16[I] = (_SX_16(argument2.16[I]) + _SX_16(argument2.16[I]) < 0 ? bias : 0) >> shift
  };
stage E1:
  GPR[registerWe<<1] = result1;
\end{lstlisting}}
\newpage

\paragraph{Encoding} Format \textsf{ALU\_HQSWRIO} (Section~\ref{sec:format-ALU-HQSWRIO}), \verb|exucode1 = 000|\\

\bitfields{ 1/P, 2/11, 1/1, 1/1, 3/000, 5/registerWo, 1/1, 2/10, 3/100, 1/1, 6/unsigned6, 5/registerZo, 1/0 }


\paragraph{Syntax} \texttt{srshq} \emph{registerWo} = \emph{registerZo}, \emph{unsigned6}

\paragraph{Description} The \emph{registerZo} is interpreted as four half words packed into 64 bits. If non-negative, each \emph{registerZo} half word is shifted right by the \emph{unsigned6} modulo 16. Else, each \emph{registerZo} half word is biased by (2**(\emph{unsigned6} modulo 16)) - 1 and shifted right by the \emph{unsigned6} modulo 16. Each result is stored into the corresponding \emph{registerWo} half word. This instruction implements a signed division of the \emph{registerZo} half words by the power of two given by the \emph{unsigned6} modulo 16.


\paragraph{Execution} {Reservation table \textsf{ALU\_LITE} (Section~\ref{sec:reservation-ALU-LITE})}
~
{\begin{lstlisting}
stage RR:
  new argument3 = _ZX_6(unsigned6);
  new argument2 = GPR[(registerZo<<1)+1];
  new shift = _ZX_4(argument3);
  new bias = (1 << shift) - 1;
  for (I = 0; I < 4; I += 1) {
    result1.16[I] = (_SX_16(argument2.16[I]) + _SX_16(argument2.16[I]) < 0 ? bias : 0) >> shift
  };
stage E1:
  GPR[(registerWo<<1)+1] = result1;
\end{lstlisting}}
\newpage

\paragraph{Encoding} Format \textsf{ALU\_HQSWRRE} (Section~\ref{sec:format-ALU-HQSWRRE}), \verb|exucode1 = 000|\\

\bitfields{ 1/P, 2/11, 1/1, 1/0, 3/000, 5/registerWe, 1/0, 2/10, 3/100, 1/1, 6/registerY, 5/registerZe, 1/0 }


\paragraph{Syntax} \texttt{srshq} \emph{registerWe} = \emph{registerZe}, \emph{registerY}

\paragraph{Description} The \emph{registerZe} is interpreted as four half words packed into 64 bits. If non-negative, each \emph{registerZe} half word is shifted right by the \emph{registerY} modulo 16. Else, each \emph{registerZe} half word is biased by (2**(\emph{registerY} modulo 16)) - 1 and shifted right by the \emph{registerY} modulo 16. Each result is stored into the corresponding \emph{registerWe} half word. This instruction implements a signed division of the \emph{registerZe} half words by the power of two given by the \emph{registerY} modulo 16.


\paragraph{Execution} {Reservation table \textsf{ALU\_LITE} (Section~\ref{sec:reservation-ALU-LITE})}
~
{\begin{lstlisting}
stage RR:
  new argument3 = GPR[registerY];
  new argument2 = GPR[registerZe<<1];
  new shift = _ZX_4(argument3);
  new bias = (1 << shift) - 1;
  for (I = 0; I < 4; I += 1) {
    result1.16[I] = (_SX_16(argument2.16[I]) + _SX_16(argument2.16[I]) < 0 ? bias : 0) >> shift
  };
stage E1:
  GPR[registerWe<<1] = result1;
\end{lstlisting}}
\newpage

\paragraph{Encoding} Format \textsf{ALU\_HQSWRRO} (Section~\ref{sec:format-ALU-HQSWRRO}), \verb|exucode1 = 000|\\

\bitfields{ 1/P, 2/11, 1/1, 1/0, 3/000, 5/registerWo, 1/1, 2/10, 3/100, 1/1, 6/registerY, 5/registerZo, 1/0 }


\paragraph{Syntax} \texttt{srshq} \emph{registerWo} = \emph{registerZo}, \emph{registerY}

\paragraph{Description} The \emph{registerZo} is interpreted as four half words packed into 64 bits. If non-negative, each \emph{registerZo} half word is shifted right by the \emph{registerY} modulo 16. Else, each \emph{registerZo} half word is biased by (2**(\emph{registerY} modulo 16)) - 1 and shifted right by the \emph{registerY} modulo 16. Each result is stored into the corresponding \emph{registerWo} half word. This instruction implements a signed division of the \emph{registerZo} half words by the power of two given by the \emph{registerY} modulo 16.


\paragraph{Execution} {Reservation table \textsf{ALU\_LITE} (Section~\ref{sec:reservation-ALU-LITE})}
~
{\begin{lstlisting}
stage RR:
  new argument3 = GPR[registerY];
  new argument2 = GPR[(registerZo<<1)+1];
  new shift = _ZX_4(argument3);
  new bias = (1 << shift) - 1;
  for (I = 0; I < 4; I += 1) {
    result1.16[I] = (_SX_16(argument2.16[I]) + _SX_16(argument2.16[I]) < 0 ? bias : 0) >> shift
  };
stage E1:
  GPR[(registerWo<<1)+1] = result1;
\end{lstlisting}}
\newpage
\subsubsection[{\small SRSW: Shift Right Symmetric Word}]{SRSW\hfill Shift Right Symmetric Word}
\label{instr:SRSW}\index{srsw}
 
\paragraph{Encoding} Format \textsf{ALU\_WSWRR} (Section~\ref{sec:format-ALU-WSWRR}), \verb|exucode1 = 000|\\

\bitfields{ 1/P, 2/11, 1/0, 1/0, 3/000, 6/registerW, 2/11, 3/100, 1/signextw, 6/registerY, 6/registerZ }


\paragraph{Syntax} \texttt{srsw}\emph{signextw} \emph{registerW} = \emph{registerZ}, \emph{registerY}

\paragraph{Description} If non-negative, the \emph{registerZ} is shifted right by the \emph{registerY} modulo 32. Else, the \emph{registerZ} biased by (2**(\emph{registerY} modulo 32)) - 1 and shifted right by the \emph{registerY} modulo 32. The result with the 32 upper bits cleared is stored into the \emph{registerW}. This instruction implements a signed division of the \emph{registerZ} least significant word by a power of two.


\paragraph{Modifier(s)}
\noindent\begin{itemize}
\item signextw (cf. signextw in section \ref{par:modifier-signextw} on page \pageref{par:modifier-signextw})
\end{itemize}

\paragraph{Execution} {Reservation table \textsf{ALU\_LITE} (Section~\ref{sec:reservation-ALU-LITE})}
~
{\begin{lstlisting}
stage ID:
  new signextw = _ZX_1(signextw);
stage RR:
  new argument3 = GPR[registerY];
  new argument2 = GPR[registerZ];
  new shift = _ZX_5(argument3);
  new bias = _SX_32(argument2) < 0 ? (1 << shift) - 1 : 0;
  new result1 = (_SX_32(argument2) + bias) >> shift;
stage E1:
  GPR[registerW] = signextw ? _SX_32(result1) : _ZX_32(result1);
\end{lstlisting}}
\newpage

\paragraph{Encoding} Format \textsf{ALU\_WSWRI} (Section~\ref{sec:format-ALU-WSWRI}), \verb|exucode1 = 000|\\

\bitfields{ 1/P, 2/11, 1/0, 1/1, 3/000, 6/registerW, 2/11, 3/100, 1/signextw, 6/unsigned6, 6/registerZ }


\paragraph{Syntax} \texttt{srsw}\emph{signextw} \emph{registerW} = \emph{registerZ}, \emph{unsigned6}

\paragraph{Description} If non-negative, the \emph{registerZ} is shifted right by the \emph{unsigned6} modulo 32. Else, the \emph{registerZ} biased by (2**(\emph{unsigned6} modulo 32)) - 1 and shifted right by the \emph{unsigned6} modulo 32. The result with the 32 upper bits cleared is stored into the \emph{registerW}. This instruction implements a signed division of the \emph{registerZ} least significant word by a power of two.


\paragraph{Modifier(s)}
\noindent\begin{itemize}
\item signextw (cf. signextw in section \ref{par:modifier-signextw} on page \pageref{par:modifier-signextw})
\end{itemize}

\paragraph{Execution} {Reservation table \textsf{ALU\_LITE} (Section~\ref{sec:reservation-ALU-LITE})}
~
{\begin{lstlisting}
stage ID:
  new signextw = _ZX_1(signextw);
stage RR:
  new argument3 = _ZX_6(unsigned6);
  new argument2 = GPR[registerZ];
  new shift = _ZX_5(argument3);
  new bias = _SX_32(argument2) < 0 ? (1 << shift) - 1 : 0;
  new result1 = (_SX_32(argument2) + bias) >> shift;
stage E1:
  GPR[registerW] = signextw ? _SX_32(result1) : _ZX_32(result1);
\end{lstlisting}}
\newpage
\subsubsection[{\small SRSWP: Shift Right Symmetric Word Pair}]{SRSWP\hfill Shift Right Symmetric Word Pair}
\label{instr:SRSWP}\index{srswp}
 
\paragraph{Encoding} Format \textsf{ALU\_WPSWRIE} (Section~\ref{sec:format-ALU-WPSWRIE}), \verb|exucode1 = 000|\\

\bitfields{ 1/P, 2/11, 1/1, 1/1, 3/000, 5/registerWe, 1/0, 2/10, 3/100, 1/0, 6/unsigned6, 5/registerZe, 1/1 }


\paragraph{Syntax} \texttt{srswp} \emph{registerWe} = \emph{registerZe}, \emph{unsigned6}

\paragraph{Description} The \emph{registerZe} is interpreted as two words packed into 64 bits. If non-negative, each \emph{registerZe} word is shifted right by the \emph{unsigned6} modulo 32. Else, each \emph{registerZe} word is biased by (2**(\emph{unsigned6} modulo 32)) - 1 and shifted right by the \emph{unsigned6} modulo 32. The two resulting words are packed and stored into the \emph{registerWe}. This instruction implements a signed division of the \emph{registerZe} words by the power of two given by the \emph{unsigned6} modulo 32.


\paragraph{Execution} {Reservation table \textsf{ALU\_LITE} (Section~\ref{sec:reservation-ALU-LITE})}
~
{\begin{lstlisting}
stage RR:
  new argument3 = _ZX_6(unsigned6);
  new argument2 = GPR[registerZe<<1];
  new shift = _ZX_5(argument3);
  new bias = (1 << shift) - 1;
  for (I = 0; I < 2; I += 1) {
    result1.32[I] = (_SX_32(argument2.32[I]) + _SX_32(argument2.32[I]) < 0 ? bias : 0) >> shift
  };
stage E1:
  GPR[registerWe<<1] = result1;
\end{lstlisting}}
\newpage

\paragraph{Encoding} Format \textsf{ALU\_WPSWRIO} (Section~\ref{sec:format-ALU-WPSWRIO}), \verb|exucode1 = 000|\\

\bitfields{ 1/P, 2/11, 1/1, 1/1, 3/000, 5/registerWo, 1/1, 2/10, 3/100, 1/0, 6/unsigned6, 5/registerZo, 1/1 }


\paragraph{Syntax} \texttt{srswp} \emph{registerWo} = \emph{registerZo}, \emph{unsigned6}

\paragraph{Description} The \emph{registerZo} is interpreted as two words packed into 64 bits. If non-negative, each \emph{registerZo} word is shifted right by the \emph{unsigned6} modulo 32. Else, each \emph{registerZo} word is biased by (2**(\emph{unsigned6} modulo 32)) - 1 and shifted right by the \emph{unsigned6} modulo 32. The two resulting words are packed and stored into the \emph{registerWo}. This instruction implements a signed division of the \emph{registerZo} words by the power of two given by the \emph{unsigned6} modulo 32.


\paragraph{Execution} {Reservation table \textsf{ALU\_LITE} (Section~\ref{sec:reservation-ALU-LITE})}
~
{\begin{lstlisting}
stage RR:
  new argument3 = _ZX_6(unsigned6);
  new argument2 = GPR[(registerZo<<1)+1];
  new shift = _ZX_5(argument3);
  new bias = (1 << shift) - 1;
  for (I = 0; I < 2; I += 1) {
    result1.32[I] = (_SX_32(argument2.32[I]) + _SX_32(argument2.32[I]) < 0 ? bias : 0) >> shift
  };
stage E1:
  GPR[(registerWo<<1)+1] = result1;
\end{lstlisting}}
\newpage

\paragraph{Encoding} Format \textsf{ALU\_WPSWRRE} (Section~\ref{sec:format-ALU-WPSWRRE}), \verb|exucode1 = 000|\\

\bitfields{ 1/P, 2/11, 1/1, 1/0, 3/000, 5/registerWe, 1/0, 2/10, 3/100, 1/0, 6/registerY, 5/registerZe, 1/1 }


\paragraph{Syntax} \texttt{srswp} \emph{registerWe} = \emph{registerZe}, \emph{registerY}

\paragraph{Description} The \emph{registerZe} is interpreted as two words packed into 64 bits. If non-negative, each \emph{registerZe} word is shifted right by the \emph{registerY} modulo 32. Else, each \emph{registerZe} word is biased by (2**(\emph{registerY} modulo 32)) - 1 and shifted right by the \emph{registerY} modulo 32. The two resulting words are packed and stored into the \emph{registerWe}. This instruction implements a signed division of the \emph{registerZe} words by the power of two given by the \emph{registerY} modulo 32.


\paragraph{Execution} {Reservation table \textsf{ALU\_LITE} (Section~\ref{sec:reservation-ALU-LITE})}
~
{\begin{lstlisting}
stage RR:
  new argument3 = GPR[registerY];
  new argument2 = GPR[registerZe<<1];
  new shift = _ZX_5(argument3);
  new bias = (1 << shift) - 1;
  for (I = 0; I < 2; I += 1) {
    result1.32[I] = (_SX_32(argument2.32[I]) + _SX_32(argument2.32[I]) < 0 ? bias : 0) >> shift
  };
stage E1:
  GPR[registerWe<<1] = result1;
\end{lstlisting}}
\newpage

\paragraph{Encoding} Format \textsf{ALU\_WPSWRRO} (Section~\ref{sec:format-ALU-WPSWRRO}), \verb|exucode1 = 000|\\

\bitfields{ 1/P, 2/11, 1/1, 1/0, 3/000, 5/registerWo, 1/1, 2/10, 3/100, 1/0, 6/registerY, 5/registerZo, 1/1 }


\paragraph{Syntax} \texttt{srswp} \emph{registerWo} = \emph{registerZo}, \emph{registerY}

\paragraph{Description} The \emph{registerZo} is interpreted as two words packed into 64 bits. If non-negative, each \emph{registerZo} word is shifted right by the \emph{registerY} modulo 32. Else, each \emph{registerZo} word is biased by (2**(\emph{registerY} modulo 32)) - 1 and shifted right by the \emph{registerY} modulo 32. The two resulting words are packed and stored into the \emph{registerWo}. This instruction implements a signed division of the \emph{registerZo} words by the power of two given by the \emph{registerY} modulo 32.


\paragraph{Execution} {Reservation table \textsf{ALU\_LITE} (Section~\ref{sec:reservation-ALU-LITE})}
~
{\begin{lstlisting}
stage RR:
  new argument3 = GPR[registerY];
  new argument2 = GPR[(registerZo<<1)+1];
  new shift = _ZX_5(argument3);
  new bias = (1 << shift) - 1;
  for (I = 0; I < 2; I += 1) {
    result1.32[I] = (_SX_32(argument2.32[I]) + _SX_32(argument2.32[I]) < 0 ? bias : 0) >> shift
  };
stage E1:
  GPR[(registerWo<<1)+1] = result1;
\end{lstlisting}}
\newpage
\subsubsection[{\small STSUD: Subtract, Test and Shift Unsigned Double Word}]{STSUD\hfill Subtract, Test and Shift Unsigned Double Word}
\label{instr:STSUD}\index{stsud}
 
\paragraph{Encoding} Format \textsf{ALU\_DWRRS} (Section~\ref{sec:format-ALU-DWRRS}), \verb|exucode2 = 1111|\\

\bitfields{ 1/P, 2/11, 1/0, 4/1111, 6/registerW, 2/10, 3/010, 1/0, 6/registerY, 6/registerZ }


\paragraph{Syntax} \texttt{stsud} \emph{registerW} = \emph{registerZ}, \emph{registerY}

\paragraph{Description} The \emph{registerZ} is subtracted from the \emph{registerY} If positive, the result is shifted left one bit and OR-ed with 0x1. Else, the \emph{registerY} is shifted left one bit. The result is stored into the \emph{registerW}. This instruction is provided to support double word unsigned integer division.


\paragraph{Execution} {Reservation table \textsf{ALU\_TINY} (Section~\ref{sec:reservation-ALU-TINY})}
~
{\begin{lstlisting}
stage RR:
  new argument3 = GPR[registerY];
  new argument2 = GPR[registerZ];
  new result1 = _ZX_64(argument3) >= _ZX_64(argument2) ? ((_ZX_64(argument3) - _ZX_64(argument2)) << 1) | 1 : _ZX_64(argument3) << 1;
stage E1:
  GPR[registerW] = result1;
\end{lstlisting}}
\newpage
\subsubsection[{\small STSUHQ: Subtract, Test and Shift Unsigned Half Word Quadruple}]{STSUHQ\hfill Subtract, Test and Shift Unsigned Half Word Quadruple}
\label{instr:STSUHQ}\index{stsuhq}
 
\paragraph{Encoding} Format \textsf{ALU\_HQWRRSE} (Section~\ref{sec:format-ALU-HQWRRSE}), \verb|exucode2 = 1111|\\

\bitfields{ 1/P, 2/11, 1/1, 4/1111, 5/registerWe, 1/0, 2/10, 3/010, 1/1, 5/registerYe, 1/--, 5/registerZe, 1/0 }


\paragraph{Syntax} \texttt{stsuhq} \emph{registerWe} = \emph{registerZe}, \emph{registerYe}

\paragraph{Description} The \emph{registerZe} interpreted as a half word quadruple is subtracted from the \emph{registerYe} interpreted as a half word quadruple. For each result, if the difference is positive it is shifted left one bit and OR-ed with 0x1. Else, the result is the corresponding half word of the \emph{registerYe} is shifted left one bit. The four resulting half words are packed and stored into the \emph{registerWe}.


\paragraph{Execution} {Reservation table \textsf{ALU\_LITE} (Section~\ref{sec:reservation-ALU-LITE})}
~
{\begin{lstlisting}
stage RR:
  new argument3 = GPR[registerYe<<1];
  new argument2 = GPR[registerZe<<1];
  for (I = 0; I < 4; I += 1) {
    result1.16[I] = _ZX_16(argument3.16[I]) >= _ZX_16(argument2.16[I]) ? ((_ZX_16(argument3.16[I]) - _ZX_16(argument2.16[I])) << 1) | 1 : _ZX_16(argument3.16[I]) << 1
  };
stage E1:
  GPR[registerWe<<1] = result1;
\end{lstlisting}}
\newpage

\paragraph{Encoding} Format \textsf{ALU\_HQWRRSO} (Section~\ref{sec:format-ALU-HQWRRSO}), \verb|exucode2 = 1111|\\

\bitfields{ 1/P, 2/11, 1/1, 4/1111, 5/registerWo, 1/1, 2/10, 3/010, 1/1, 5/registerYo, 1/--, 5/registerZo, 1/0 }


\paragraph{Syntax} \texttt{stsuhq} \emph{registerWo} = \emph{registerZo}, \emph{registerYo}

\paragraph{Description} The \emph{registerZo} interpreted as a half word quadruple is subtracted from the \emph{registerYo} interpreted as a half word quadruple. For each result, if the difference is positive it is shifted left one bit and OR-ed with 0x1. Else, the result is the corresponding half word of the \emph{registerYo} is shifted left one bit. The four resulting half words are packed and stored into the \emph{registerWo}.


\paragraph{Execution} {Reservation table \textsf{ALU\_LITE} (Section~\ref{sec:reservation-ALU-LITE})}
~
{\begin{lstlisting}
stage RR:
  new argument3 = GPR[(registerYo<<1)+1];
  new argument2 = GPR[(registerZo<<1)+1];
  for (I = 0; I < 4; I += 1) {
    result1.16[I] = _ZX_16(argument3.16[I]) >= _ZX_16(argument2.16[I]) ? ((_ZX_16(argument3.16[I]) - _ZX_16(argument2.16[I])) << 1) | 1 : _ZX_16(argument3.16[I]) << 1
  };
stage E1:
  GPR[(registerWo<<1)+1] = result1;
\end{lstlisting}}
\newpage
\subsubsection[{\small STSUW: Subtract, Test and Shift Unsigned Word}]{STSUW\hfill Subtract, Test and Shift Unsigned Word}
\label{instr:STSUW}\index{stsuw}
 
\paragraph{Encoding} Format \textsf{ALU\_WWRRS} (Section~\ref{sec:format-ALU-WWRRS}), \verb|exucode2 = 1111|\\

\bitfields{ 1/P, 2/11, 1/0, 4/1111, 6/registerW, 2/11, 3/010, 1/signextw, 6/registerY, 6/registerZ }


\paragraph{Syntax} \texttt{stsuw}\emph{signextw} \emph{registerW} = \emph{registerZ}, \emph{registerY}

\paragraph{Description} The \emph{registerZ} is subtracted from the \emph{registerY}. If positive, the result is shifted left one bit and OR-ed with 0x1. Else, the \emph{registerY} is shifted left one bit. The result with the 32 upper bits cleared is stored into the \emph{registerW}. This instruction is provided to support unsigned integer division.


\paragraph{Modifier(s)}
\noindent\begin{itemize}
\item signextw (cf. signextw in section \ref{par:modifier-signextw} on page \pageref{par:modifier-signextw})
\end{itemize}

\paragraph{Execution} {Reservation table \textsf{ALU\_TINY} (Section~\ref{sec:reservation-ALU-TINY})}
~
{\begin{lstlisting}
stage ID:
  new signextw = _ZX_1(signextw);
stage RR:
  new argument3 = GPR[registerY];
  new argument2 = GPR[registerZ];
  new result1 = argument3.32[0] >= argument2.32[0] ? ((argument3.32[0] - argument2.32[0]) << 1) | 1 : argument3.32[0] << 1;
stage E1:
  GPR[registerW] = signextw ? _SX_32(result1) : _ZX_32(result1);
\end{lstlisting}}
\newpage
\subsubsection[{\small STSUWP: Subtract, Test and Shift Unsigned Word Pair}]{STSUWP\hfill Subtract, Test and Shift Unsigned Word Pair}
\label{instr:STSUWP}\index{stsuwp}
 
\paragraph{Encoding} Format \textsf{ALU\_WPWRRSE} (Section~\ref{sec:format-ALU-WPWRRSE}), \verb|exucode2 = 1111|\\

\bitfields{ 1/P, 2/11, 1/1, 4/1111, 5/registerWe, 1/0, 2/10, 3/010, 1/0, 5/registerYe, 1/--, 5/registerZe, 1/1 }


\paragraph{Syntax} \texttt{stsuwp} \emph{registerWe} = \emph{registerZe}, \emph{registerYe}

\paragraph{Description} The \emph{registerZe} interpreted as a word pair is subtracted from the \emph{registerYe} interpreted as a word pair. For each result, if the difference is positive it is shifted left one bit and OR-ed with 0x1. Else, the result is the corresponding word of the \emph{registerYe} is shifted left one bit. The two resulting word are packed and stored into the \emph{registerWe}.


\paragraph{Execution} {Reservation table \textsf{ALU\_LITE} (Section~\ref{sec:reservation-ALU-LITE})}
~
{\begin{lstlisting}
stage RR:
  new argument3 = GPR[registerYe<<1];
  new argument2 = GPR[registerZe<<1];
  for (I = 0; I < 2; I += 1) {
    result1.32[I] = _ZX_32(argument3.32[I]) >= _ZX_32(argument2.32[I]) ? ((_ZX_32(argument3.32[I]) - _ZX_32(argument2.32[I])) << 1) | 1 : _ZX_32(argument3.32[I]) << 1
  };
stage E1:
  GPR[registerWe<<1] = result1;
\end{lstlisting}}
\newpage

\paragraph{Encoding} Format \textsf{ALU\_WPWRRSO} (Section~\ref{sec:format-ALU-WPWRRSO}), \verb|exucode2 = 1111|\\

\bitfields{ 1/P, 2/11, 1/1, 4/1111, 5/registerWo, 1/1, 2/10, 3/010, 1/0, 5/registerYo, 1/--, 5/registerZo, 1/1 }


\paragraph{Syntax} \texttt{stsuwp} \emph{registerWo} = \emph{registerZo}, \emph{registerYo}

\paragraph{Description} The \emph{registerZo} interpreted as a word pair is subtracted from the \emph{registerYo} interpreted as a word pair. For each result, if the difference is positive it is shifted left one bit and OR-ed with 0x1. Else, the result is the corresponding word of the \emph{registerYo} is shifted left one bit. The two resulting word are packed and stored into the \emph{registerWo}.


\paragraph{Execution} {Reservation table \textsf{ALU\_LITE} (Section~\ref{sec:reservation-ALU-LITE})}
~
{\begin{lstlisting}
stage RR:
  new argument3 = GPR[(registerYo<<1)+1];
  new argument2 = GPR[(registerZo<<1)+1];
  for (I = 0; I < 2; I += 1) {
    result1.32[I] = _ZX_32(argument3.32[I]) >= _ZX_32(argument2.32[I]) ? ((_ZX_32(argument3.32[I]) - _ZX_32(argument2.32[I])) << 1) | 1 : _ZX_32(argument3.32[I]) << 1
  };
stage E1:
  GPR[(registerWo<<1)+1] = result1;
\end{lstlisting}}
\newpage
\subsubsection[{\small TRUNCDWP: Truncate Double Word to Word Pair}]{TRUNCDWP\hfill Truncate Double Word to Word Pair}
\label{instr:TRUNCDWP}\index{truncdwp}
 
\paragraph{Encoding} Format \textsf{ALU\_WPNWRR} (Section~\ref{sec:format-ALU-WPNWRR}), \verb|alucode8 = 1000|\\

\bitfields{ 1/P, 2/11, 1/1, 4/1111, 6/registerW, 2/10, 3/101, 1/0, 4/1000, 1/--, 1/0, 5/registerP, 1/1 }


\paragraph{Syntax} \texttt{truncdwp} \emph{registerW} = \emph{registerP}

\paragraph{Description} The \emph{registerP} is interpreted as two double words packed into 128 bits. The \emph{registerP} double words are converted to 32-bit words by removing the high bits. The two resulting words are packed and stored back into the \emph{registerW}.


\paragraph{Execution} {Reservation table \textsf{ALU\_LITE} (Section~\ref{sec:reservation-ALU-LITE})}
~
{\begin{lstlisting}
stage RR:
  new argument2 = PGR[registerP];
  result1.32[0] = argument2.64[0];
  result1.32[1] = argument2.64[1];
stage E1:
  GPR[registerW] = result1;
\end{lstlisting}}
\newpage
\subsubsection[{\small TRUNCHBO: Truncate Half Word to Byte Octuple}]{TRUNCHBO\hfill Truncate Half Word to Byte Octuple}
\label{instr:TRUNCHBO}\index{trunchbo}
 
\paragraph{Encoding} Format \textsf{ALU\_BONWRR} (Section~\ref{sec:format-ALU-BONWRR}), \verb|alucode8 = 1000|\\

\bitfields{ 1/P, 2/11, 1/1, 4/1111, 6/registerW, 2/10, 3/101, 1/1, 4/1000, 1/ziplanes, 1/0, 5/registerP, 1/1 }


\paragraph{Syntax} \texttt{trunchbo}\emph{ziplanes} \emph{registerW} = \emph{registerP}

\paragraph{Description} The \emph{registerP} is interpreted as eight half words packed into 128 bits. The \emph{registerP} half words are converted to bytes by removing the high bits. If \emph{ziplanes}, the resulting bytes are interleaved perfect shuffling. The eight resulting bytes are packed and stored back into the \emph{registerW}.


\paragraph{Modifier(s)}
\noindent\begin{itemize}
\item ziplanes (cf. ziplanes in section \ref{par:modifier-ziplanes} on page \pageref{par:modifier-ziplanes})
\end{itemize}

\paragraph{Execution} {Reservation table \textsf{ALU\_LITE} (Section~\ref{sec:reservation-ALU-LITE})}
~
{\begin{lstlisting}
stage ID:
  new ziplanes = _ZX_1(ziplanes);
stage RR:
  new argument2 = PGR[registerP];
  if (ziplanes) {
    result1.8[0] = argument2.16[0];
    result1.8[1] = argument2.16[4];
    result1.8[2] = argument2.16[1];
    result1.8[3] = argument2.16[5];
    result1.8[4] = argument2.16[2];
    result1.8[5] = argument2.16[6];
    result1.8[6] = argument2.16[3];
  } else {
    result1.8[0] = argument2.16[0];
    result1.8[1] = argument2.16[1];
    result1.8[2] = argument2.16[2];
    result1.8[3] = argument2.16[3];
    result1.8[4] = argument2.16[4];
    result1.8[5] = argument2.16[5];
    result1.8[6] = argument2.16[6];
  }
  result1.8[7] = argument2.16[7];
stage E1:
  GPR[registerW] = result1;
\end{lstlisting}}
\newpage
\subsubsection[{\small TRUNCWHQ: Truncate Word to Half Word Quadruple}]{TRUNCWHQ\hfill Truncate Word to Half Word Quadruple}
\label{instr:TRUNCWHQ}\index{truncwhq}
 
\paragraph{Encoding} Format \textsf{ALU\_HQNWRR} (Section~\ref{sec:format-ALU-HQNWRR}), \verb|alucode8 = 1000|\\

\bitfields{ 1/P, 2/11, 1/1, 4/1111, 6/registerW, 2/10, 3/101, 1/1, 4/1000, 1/ziplanes, 1/0, 5/registerP, 1/0 }


\paragraph{Syntax} \texttt{truncwhq}\emph{ziplanes} \emph{registerW} = \emph{registerP}

\paragraph{Description} The \emph{registerP} is interpreted as four words packed into 128 bits. The \emph{registerP} words are converted to 16-bit half words by removing the high bits. If \emph{ziplanes}, the resulting half words are interleaved by perfect shuffling. The four resulting half words are packed and stored back into the \emph{registerW}.


\paragraph{Modifier(s)}
\noindent\begin{itemize}
\item ziplanes (cf. ziplanes in section \ref{par:modifier-ziplanes} on page \pageref{par:modifier-ziplanes})
\end{itemize}

\paragraph{Execution} {Reservation table \textsf{ALU\_LITE} (Section~\ref{sec:reservation-ALU-LITE})}
~
{\begin{lstlisting}
stage ID:
  new ziplanes = _ZX_1(ziplanes);
stage RR:
  new argument2 = PGR[registerP];
  if (ziplanes) {
    result1.16[0] = argument2.32[0];
    result1.16[1] = argument2.32[2];
    result1.16[2] = argument2.32[1];
  } else {
    result1.16[0] = argument2.32[0];
    result1.16[1] = argument2.32[1];
    result1.16[2] = argument2.32[2];
  }
  result1.16[3] = argument2.32[3];
stage E1:
  GPR[registerW] = result1;
\end{lstlisting}}
\newpage
\subsubsection[{\small WIDENQBHO: Widen Q7 Byte to Half Word Octuple}]{WIDENQBHO\hfill Widen Q7 Byte to Half Word Octuple}
\label{instr:WIDENQBHO}\index{widenqbho}
 
\paragraph{Encoding} Format \textsf{ALU\_HOWIDEWR} (Section~\ref{sec:format-ALU-HOWIDEWR}), \verb|alucode8 = 1010|\\

\bitfields{ 1/P, 2/11, 1/1, 4/1111, 5/registerM, 1/--, 2/10, 3/101, 1/1, 4/1010, 1/mostsig, 1/1, 5/registerP, 1/0 }


\paragraph{Syntax} \texttt{widenqbho}\emph{mostsig} \emph{registerM} = \emph{registerP}

\paragraph{Description} The \emph{registerP} is interpreted as sixteen bytes packed into 128 bits. The least or most significant bytes of the \emph{registerP} are selected, depending on \emph{mostsig}. These eight bytes interpreted as Q7 numbers are mapped to Q15 numbers, packed and stored into the \emph{registerM}.


\paragraph{Modifier(s)}
\noindent\begin{itemize}
\item mostsig (cf. mostsig in section \ref{par:modifier-mostsig} on page \pageref{par:modifier-mostsig})
\end{itemize}

\paragraph{Execution} {Reservation table \textsf{ALU\_LITE} (Section~\ref{sec:reservation-ALU-LITE})}
~
{\begin{lstlisting}
stage ID:
  new mostsig = _ZX_1(mostsig);
stage RR:
  new argument2 = PGR[registerP];
  if (mostsig) {
    argument2 = argument2 >> 64;
  }
  for (I = 0; I < 8; I += 1) {
    result1.16[I] = argument2.8[I] << 8
  };
stage E1:
  PGR[registerM] = result1;
\end{lstlisting}}
\newpage
\subsubsection[{\small WIDENQHWQ: Widen Q15 Half Word to Word Quadruple}]{WIDENQHWQ\hfill Widen Q15 Half Word to Word Quadruple}
\label{instr:WIDENQHWQ}\index{widenqhwq}
 
\paragraph{Encoding} Format \textsf{ALU\_WQWIDEWR} (Section~\ref{sec:format-ALU-WQWIDEWR}), \verb|alucode8 = 1010|\\

\bitfields{ 1/P, 2/11, 1/1, 4/1111, 5/registerM, 1/--, 2/10, 3/101, 1/0, 4/1010, 1/mostsig, 1/1, 5/registerP, 1/1 }


\paragraph{Syntax} \texttt{widenqhwq}\emph{mostsig} \emph{registerM} = \emph{registerP}

\paragraph{Description} The \emph{registerP} is interpreted as eight half words packed into 128 bits. The least or most significant half words of the \emph{registerP} are selected, depending on \emph{mostsig}. These four half words interpreted as Q15 numbers are mapped to Q31 numbers, packed and stored into the \emph{registerM}.


\paragraph{Modifier(s)}
\noindent\begin{itemize}
\item mostsig (cf. mostsig in section \ref{par:modifier-mostsig} on page \pageref{par:modifier-mostsig})
\end{itemize}

\paragraph{Execution} {Reservation table \textsf{ALU\_LITE} (Section~\ref{sec:reservation-ALU-LITE})}
~
{\begin{lstlisting}
stage ID:
  new mostsig = _ZX_1(mostsig);
stage RR:
  new argument2 = PGR[registerP];
  if (mostsig) {
    argument2 = argument2 >> 64;
  }
  for (I = 0; I < 4; I += 1) {
    result1.32[I] = argument2.16[I] << 16
  };
stage E1:
  PGR[registerM] = result1;
\end{lstlisting}}
\newpage
\subsubsection[{\small WIDENQWDP: Widen Q31 Word to Double Word Pair}]{WIDENQWDP\hfill Widen Q31 Word to Double Word Pair}
\label{instr:WIDENQWDP}\index{widenqwdp}
 
\paragraph{Encoding} Format \textsf{ALU\_DPWIDEWR} (Section~\ref{sec:format-ALU-DPWIDEWR}), \verb|alucode8 = 1010|\\

\bitfields{ 1/P, 2/11, 1/1, 4/1111, 5/registerM, 1/--, 2/10, 3/101, 1/0, 4/1010, 1/mostsig, 1/1, 5/registerP, 1/0 }


\paragraph{Syntax} \texttt{widenqwdp}\emph{mostsig} \emph{registerM} = \emph{registerP}

\paragraph{Description} The \emph{registerP} is interpreted as four words packed into 128 bits. The least or most significant words of the \emph{registerP} are selected, depending on \emph{mostsig}. These two words interpreted as Q31 numbers are mapped to Q63 numbers, packed and stored into the \emph{registerM}.


\paragraph{Modifier(s)}
\noindent\begin{itemize}
\item mostsig (cf. mostsig in section \ref{par:modifier-mostsig} on page \pageref{par:modifier-mostsig})
\end{itemize}

\paragraph{Execution} {Reservation table \textsf{ALU\_LITE} (Section~\ref{sec:reservation-ALU-LITE})}
~
{\begin{lstlisting}
stage ID:
  new mostsig = _ZX_1(mostsig);
stage RR:
  new argument2 = PGR[registerP];
  if (mostsig) {
    argument2 = argument2 >> 64;
  }
  for (I = 0; I < 2; I += 1) {
    result1.64[I] = argument2.32[I] << 32
  };
stage E1:
  PGR[registerM] = result1;
\end{lstlisting}}
\newpage
\subsubsection[{\small WIDENSBHO: Sign Widen Byte to Half Word Octuple}]{WIDENSBHO\hfill Sign Widen Byte to Half Word Octuple}
\label{instr:WIDENSBHO}\index{widensbho}
 
\paragraph{Encoding} Format \textsf{ALU\_HOWIDEWR} (Section~\ref{sec:format-ALU-HOWIDEWR}), \verb|alucode8 = 1001|\\

\bitfields{ 1/P, 2/11, 1/1, 4/1111, 5/registerM, 1/--, 2/10, 3/101, 1/1, 4/1001, 1/mostsig, 1/1, 5/registerP, 1/0 }


\paragraph{Syntax} \texttt{widensbho}\emph{mostsig} \emph{registerM} = \emph{registerP}

\paragraph{Description} The \emph{registerP} is interpreted as sixteen bytes packed into 128 bits. The least or most significant bytes of the \emph{registerP} are selected, depending on \emph{mostsig}. These eight bytes are sign-extended to 16 bits, packed and stored into the \emph{registerM}.


\paragraph{Modifier(s)}
\noindent\begin{itemize}
\item mostsig (cf. mostsig in section \ref{par:modifier-mostsig} on page \pageref{par:modifier-mostsig})
\end{itemize}

\paragraph{Execution} {Reservation table \textsf{ALU\_LITE} (Section~\ref{sec:reservation-ALU-LITE})}
~
{\begin{lstlisting}
stage ID:
  new mostsig = _ZX_1(mostsig);
stage RR:
  new argument2 = PGR[registerP];
  if (mostsig) {
    argument2 = argument2 >> 64;
  }
  for (I = 0; I < 8; I += 1) {
    result1.16[I] = _SX_8(argument2.8[I])
  };
stage E1:
  PGR[registerM] = result1;
\end{lstlisting}}
\newpage
\subsubsection[{\small WIDENSHWQ: Sign Widen Half Word to Word Quadruple}]{WIDENSHWQ\hfill Sign Widen Half Word to Word Quadruple}
\label{instr:WIDENSHWQ}\index{widenshwq}
 
\paragraph{Encoding} Format \textsf{ALU\_WQWIDEWR} (Section~\ref{sec:format-ALU-WQWIDEWR}), \verb|alucode8 = 1001|\\

\bitfields{ 1/P, 2/11, 1/1, 4/1111, 5/registerM, 1/--, 2/10, 3/101, 1/0, 4/1001, 1/mostsig, 1/1, 5/registerP, 1/1 }


\paragraph{Syntax} \texttt{widenshwq}\emph{mostsig} \emph{registerM} = \emph{registerP}

\paragraph{Description} The \emph{registerP} is interpreted as eight half words packed into 128 bits. The least or most significant half words of the \emph{registerP} are selected, depending on \emph{mostsig}. These four half words are sign-extended to 32 bits, packed and stored into the \emph{registerM}.


\paragraph{Modifier(s)}
\noindent\begin{itemize}
\item mostsig (cf. mostsig in section \ref{par:modifier-mostsig} on page \pageref{par:modifier-mostsig})
\end{itemize}

\paragraph{Execution} {Reservation table \textsf{ALU\_LITE} (Section~\ref{sec:reservation-ALU-LITE})}
~
{\begin{lstlisting}
stage ID:
  new mostsig = _ZX_1(mostsig);
stage RR:
  new argument2 = PGR[registerP];
  if (mostsig) {
    argument2 = argument2 >> 64;
  }
  for (I = 0; I < 4; I += 1) {
    result1.32[I] = _SX_16(argument2.16[I])
  };
stage E1:
  PGR[registerM] = result1;
\end{lstlisting}}
\newpage
\subsubsection[{\small WIDENSWDP: Sign Widen Word to Double Word Pair}]{WIDENSWDP\hfill Sign Widen Word to Double Word Pair}
\label{instr:WIDENSWDP}\index{widenswdp}
 
\paragraph{Encoding} Format \textsf{ALU\_DPWIDEWR} (Section~\ref{sec:format-ALU-DPWIDEWR}), \verb|alucode8 = 1001|\\

\bitfields{ 1/P, 2/11, 1/1, 4/1111, 5/registerM, 1/--, 2/10, 3/101, 1/0, 4/1001, 1/mostsig, 1/1, 5/registerP, 1/0 }


\paragraph{Syntax} \texttt{widenswdp}\emph{mostsig} \emph{registerM} = \emph{registerP}

\paragraph{Description} The \emph{registerP} is interpreted as four words packed into 128 bits. The least or most significant words of the \emph{registerP} are selected, depending on \emph{mostsig}. These two words are sign-extended to 64 bits, packed and stored into the \emph{registerM}.


\paragraph{Modifier(s)}
\noindent\begin{itemize}
\item mostsig (cf. mostsig in section \ref{par:modifier-mostsig} on page \pageref{par:modifier-mostsig})
\end{itemize}

\paragraph{Execution} {Reservation table \textsf{ALU\_LITE} (Section~\ref{sec:reservation-ALU-LITE})}
~
{\begin{lstlisting}
stage ID:
  new mostsig = _ZX_1(mostsig);
stage RR:
  new argument2 = PGR[registerP];
  if (mostsig) {
    argument2 = argument2 >> 64;
  }
  for (I = 0; I < 2; I += 1) {
    result1.64[I] = _SX_32(argument2.32[I])
  };
stage E1:
  PGR[registerM] = result1;
\end{lstlisting}}
\newpage
\subsubsection[{\small WIDENZBHO: Zero Widen Byte to Half Word Octuple}]{WIDENZBHO\hfill Zero Widen Byte to Half Word Octuple}
\label{instr:WIDENZBHO}\index{widenzbho}
 
\paragraph{Encoding} Format \textsf{ALU\_HOWIDEWR} (Section~\ref{sec:format-ALU-HOWIDEWR}), \verb|alucode8 = 1000|\\

\bitfields{ 1/P, 2/11, 1/1, 4/1111, 5/registerM, 1/--, 2/10, 3/101, 1/1, 4/1000, 1/mostsig, 1/1, 5/registerP, 1/0 }


\paragraph{Syntax} \texttt{widenzbho}\emph{mostsig} \emph{registerM} = \emph{registerP}

\paragraph{Description} The \emph{registerP} is interpreted as sixteen bytes packed into 128 bits. The least or most significant bytes of the \emph{registerP} are selected, depending on \emph{mostsig}. These eight bytes are zero-extended to 16 bits, packed and stored into the \emph{registerM}.


\paragraph{Modifier(s)}
\noindent\begin{itemize}
\item mostsig (cf. mostsig in section \ref{par:modifier-mostsig} on page \pageref{par:modifier-mostsig})
\end{itemize}

\paragraph{Execution} {Reservation table \textsf{ALU\_LITE} (Section~\ref{sec:reservation-ALU-LITE})}
~
{\begin{lstlisting}
stage ID:
  new mostsig = _ZX_1(mostsig);
stage RR:
  new argument2 = PGR[registerP];
  if (mostsig) {
    argument2 = argument2 >> 64;
  }
  for (I = 0; I < 8; I += 1) {
    result1.16[I] = _ZX_8(argument2.8[I])
  };
stage E1:
  PGR[registerM] = result1;
\end{lstlisting}}
\newpage
\subsubsection[{\small WIDENZHWQ: Zero Widen Half Word to Word Quadruple}]{WIDENZHWQ\hfill Zero Widen Half Word to Word Quadruple}
\label{instr:WIDENZHWQ}\index{widenzhwq}
 
\paragraph{Encoding} Format \textsf{ALU\_WQWIDEWR} (Section~\ref{sec:format-ALU-WQWIDEWR}), \verb|alucode8 = 1000|\\

\bitfields{ 1/P, 2/11, 1/1, 4/1111, 5/registerM, 1/--, 2/10, 3/101, 1/0, 4/1000, 1/mostsig, 1/1, 5/registerP, 1/1 }


\paragraph{Syntax} \texttt{widenzhwq}\emph{mostsig} \emph{registerM} = \emph{registerP}

\paragraph{Description} The \emph{registerP} is interpreted as eight half words packed into 128 bits. The least or most significant half words of the \emph{registerP} are selected, depending on \emph{mostsig}. These four half words are zero-extended to 32 bits, packed and stored into the \emph{registerM}.


\paragraph{Modifier(s)}
\noindent\begin{itemize}
\item mostsig (cf. mostsig in section \ref{par:modifier-mostsig} on page \pageref{par:modifier-mostsig})
\end{itemize}

\paragraph{Execution} {Reservation table \textsf{ALU\_LITE} (Section~\ref{sec:reservation-ALU-LITE})}
~
{\begin{lstlisting}
stage ID:
  new mostsig = _ZX_1(mostsig);
stage RR:
  new argument2 = PGR[registerP];
  if (mostsig) {
    argument2 = argument2 >> 64;
  }
  for (I = 0; I < 4; I += 1) {
    result1.32[I] = _ZX_16(argument2.16[I])
  };
stage E1:
  PGR[registerM] = result1;
\end{lstlisting}}
\newpage
\subsubsection[{\small WIDENZWDP: Zero Widen Word to Double Word Pair}]{WIDENZWDP\hfill Zero Widen Word to Double Word Pair}
\label{instr:WIDENZWDP}\index{widenzwdp}
 
\paragraph{Encoding} Format \textsf{ALU\_DPWIDEWR} (Section~\ref{sec:format-ALU-DPWIDEWR}), \verb|alucode8 = 1000|\\

\bitfields{ 1/P, 2/11, 1/1, 4/1111, 5/registerM, 1/--, 2/10, 3/101, 1/0, 4/1000, 1/mostsig, 1/1, 5/registerP, 1/0 }


\paragraph{Syntax} \texttt{widenzwdp}\emph{mostsig} \emph{registerM} = \emph{registerP}

\paragraph{Description} The \emph{registerP} is interpreted as four words packed into 128 bits. The least or most significant words of the \emph{registerP} are selected, depending on \emph{mostsig}. These two words are zero-extended to 64 bits, packed and stored into the \emph{registerM}.


\paragraph{Modifier(s)}
\noindent\begin{itemize}
\item mostsig (cf. mostsig in section \ref{par:modifier-mostsig} on page \pageref{par:modifier-mostsig})
\end{itemize}

\paragraph{Execution} {Reservation table \textsf{ALU\_LITE} (Section~\ref{sec:reservation-ALU-LITE})}
~
{\begin{lstlisting}
stage ID:
  new mostsig = _ZX_1(mostsig);
stage RR:
  new argument2 = PGR[registerP];
  if (mostsig) {
    argument2 = argument2 >> 64;
  }
  for (I = 0; I < 2; I += 1) {
    result1.64[I] = _ZX_32(argument2.32[I])
  };
stage E1:
  PGR[registerM] = result1;
\end{lstlisting}}
\newpage
\subsubsection[{\small XMOVETD: Move To Extension Double Word}]{XMOVETD\hfill Move To Extension Double Word}
\label{instr:XMOVETD}\index{xmovetd}
 
\paragraph{Encoding} Format \textsf{ALU\_XMOVETD0} (Section~\ref{sec:format-ALU-XMOVETD0}), \verb|alucode3 = 010|\\

\bitfields{ 1/P, 2/11, 1/0, 2/10, 2/00, 6/registerAx, 2/10, 3/010, 1/1, 6/-- -- -- -- -- --, 6/registerZ }


\paragraph{Syntax} \texttt{xmovetd} \emph{registerAx} = \emph{registerZ}

\paragraph{Description} The \emph{registerAx} operand receives the \emph{registerZ} value.


\paragraph{Execution} {Reservation table \textsf{ALU\_LITE\_MISC} (Section~\ref{sec:reservation-ALU-LITE-MISC})}
~
{\begin{lstlisting}
stage RR:
  new argument2 = GPR[registerZ];
  new result1 = _ZX_64(argument2);
stage E1:
  CS.XMF = 1;
  XCR[registerAx<<2] = result1;
\end{lstlisting}}
\newpage

\paragraph{Encoding} Format \textsf{ALU\_XMOVETD1} (Section~\ref{sec:format-ALU-XMOVETD1}), \verb|alucode3 = 010|\\

\bitfields{ 1/P, 2/11, 1/0, 2/10, 2/01, 6/registerAy, 2/10, 3/010, 1/1, 6/-- -- -- -- -- --, 6/registerZ }


\paragraph{Syntax} \texttt{xmovetd} \emph{registerAy} = \emph{registerZ}

\paragraph{Description} The \emph{registerAy} operand receives the \emph{registerZ} value.


\paragraph{Execution} {Reservation table \textsf{ALU\_LITE\_MISC} (Section~\ref{sec:reservation-ALU-LITE-MISC})}
~
{\begin{lstlisting}
stage RR:
  new argument2 = GPR[registerZ];
  new result1 = _ZX_64(argument2);
stage E1:
  CS.XMF = 1;
  XCR[(registerAy<<2)+1] = result1;
\end{lstlisting}}
\newpage

\paragraph{Encoding} Format \textsf{ALU\_XMOVETD2} (Section~\ref{sec:format-ALU-XMOVETD2}), \verb|alucode3 = 010|\\

\bitfields{ 1/P, 2/11, 1/0, 2/10, 2/10, 6/registerAz, 2/10, 3/010, 1/1, 6/-- -- -- -- -- --, 6/registerZ }


\paragraph{Syntax} \texttt{xmovetd} \emph{registerAz} = \emph{registerZ}

\paragraph{Description} The \emph{registerAz} operand receives the \emph{registerZ} value.


\paragraph{Execution} {Reservation table \textsf{ALU\_LITE\_MISC} (Section~\ref{sec:reservation-ALU-LITE-MISC})}
~
{\begin{lstlisting}
stage RR:
  new argument2 = GPR[registerZ];
  new result1 = _ZX_64(argument2);
stage E1:
  CS.XMF = 1;
  XCR[(registerAz<<2)+2] = result1;
\end{lstlisting}}
\newpage

\paragraph{Encoding} Format \textsf{ALU\_XMOVETD3} (Section~\ref{sec:format-ALU-XMOVETD3}), \verb|alucode3 = 010|\\

\bitfields{ 1/P, 2/11, 1/0, 2/10, 2/11, 6/registerAt, 2/10, 3/010, 1/1, 6/-- -- -- -- -- --, 6/registerZ }


\paragraph{Syntax} \texttt{xmovetd} \emph{registerAt} = \emph{registerZ}

\paragraph{Description} The \emph{registerAt} operand receives the \emph{registerZ} value.


\paragraph{Execution} {Reservation table \textsf{ALU\_LITE\_MISC} (Section~\ref{sec:reservation-ALU-LITE-MISC})}
~
{\begin{lstlisting}
stage RR:
  new argument2 = GPR[registerZ];
  new result1 = _ZX_64(argument2);
stage E1:
  CS.XMF = 1;
  XCR[(registerAt<<2)+3] = result1;
\end{lstlisting}}
\newpage
\subsubsection[{\small XMOVETO: Move To Extension Octuple Word}]{XMOVETO\hfill Move To Extension Octuple Word}
\label{instr:XMOVETO}\index{xmoveto}
 
\paragraph{Encoding} Format \textsf{ALU\_XMOVETO} (Section~\ref{sec:format-ALU-XMOVETO}), \verb|alucode3 = 010|\\

\bitfields{ 1/P, 2/11, 1/0, 4/1110, 6/registerA, 2/10, 3/010, 1/1, 5/registerO, 1/--, 5/registerP, 1/-- }


\paragraph{Syntax} \texttt{xmoveto} \emph{registerA} = \emph{registerP}, \emph{registerO}

\paragraph{Description} The \emph{registerA} operand receives the concatenation of the \emph{registerP} value and the \emph{registerO} value.


\paragraph{Execution} {Reservation table \textsf{ALU\_LITE\_MISC} (Section~\ref{sec:reservation-ALU-LITE-MISC})}
~
{\begin{lstlisting}
stage RR:
  new argument3 = PGR[registerO];
  new argument2 = PGR[registerP];
  new result1 = 0;
  result1.128[0] = argument2;
  result1.128[1] = argument3;
stage E1:
  CS.XMF = 1;
  XVR[registerA] = result1;
\end{lstlisting}}
\newpage
\subsubsection[{\small XMOVETQ: Move To Extension Quadruple Word}]{XMOVETQ\hfill Move To Extension Quadruple Word}
\label{instr:XMOVETQ}\index{xmovetq}
 
\paragraph{Encoding} Format \textsf{ALU\_XMOVETQ0} (Section~\ref{sec:format-ALU-XMOVETQ0}), \verb|alucode3 = 010|\\

\bitfields{ 1/P, 2/11, 1/0, 3/110, 1/0, 6/registerAE, 2/10, 3/010, 1/1, 6/registerY, 6/registerZ }


\paragraph{Syntax} \texttt{xmovetq} \emph{registerAE} = \emph{registerZ}, \emph{registerY}

\paragraph{Description} The \emph{registerAE} operand receives the concatenation of the \emph{registerZ} value and the \emph{registerY} value.


\paragraph{Execution} {Reservation table \textsf{ALU\_LITE\_MISC} (Section~\ref{sec:reservation-ALU-LITE-MISC})}
~
{\begin{lstlisting}
stage RR:
  new argument3 = GPR[registerY];
  new argument2 = GPR[registerZ];
  new result1 = 0;
  result1.64[0] = argument2;
  result1.64[1] = argument3;
stage E1:
  CS.XMF = 1;
  XBR[registerAE<<1] = result1;
\end{lstlisting}}
\newpage

\paragraph{Encoding} Format \textsf{ALU\_XMOVETQ1} (Section~\ref{sec:format-ALU-XMOVETQ1}), \verb|alucode3 = 010|\\

\bitfields{ 1/P, 2/11, 1/0, 3/110, 1/1, 6/registerAO, 2/10, 3/010, 1/1, 6/registerY, 6/registerZ }


\paragraph{Syntax} \texttt{xmovetq} \emph{registerAO} = \emph{registerZ}, \emph{registerY}

\paragraph{Description} The \emph{registerAO} operand receives the concatenation of the \emph{registerZ} value and the \emph{registerY} value.


\paragraph{Execution} {Reservation table \textsf{ALU\_LITE\_MISC} (Section~\ref{sec:reservation-ALU-LITE-MISC})}
~
{\begin{lstlisting}
stage RR:
  new argument3 = GPR[registerY];
  new argument2 = GPR[registerZ];
  new result1 = 0;
  result1.64[0] = argument2;
  result1.64[1] = argument3;
stage E1:
  CS.XMF = 1;
  XBR[(registerAO<<1)+1] = result1;
\end{lstlisting}}
\newpage
\subsection{FPU Instructions}
\label{sec:FPU-instructions}

\subsubsection[{\small FABSD: Floating-Point Absolute Value Double Word}]{FABSD\hfill Floating-Point Absolute Value Double Word}
\label{instr:FABSD}\index{fabsd}
 
\paragraph{Encoding} Format \textsf{ALU\_FDMO1} (Section~\ref{sec:format-ALU-FDMO1}), \verb|fpucode6 = 0011|\\

\bitfields{ 1/P, 2/11, 1/1, 1/1, 3/111, 6/registerW, 2/00, 3/000, 1/0, 4/0011, 1/--, 1/0, 6/registerZ }


\paragraph{Syntax} \texttt{fabsd} \emph{registerW} = \emph{registerZ}

\paragraph{Description} The absolute value the \emph{registerZ}, assuming binary 64 floating-point number, is stored into the \emph{registerW}. Implemented as a bitwise operation, does not raise FP exceptions.


\paragraph{Execution} {Reservation table \textsf{ALU\_LITE} (Section~\ref{sec:reservation-ALU-LITE})}
~
{\begin{lstlisting}
stage RR:
  new argument2 = GPR[registerZ];
  new result1 = argument2 & 0x7FFFFFFFFFFFFFFF;
stage E1:
  GPR[registerW] = result1;
\end{lstlisting}}
\newpage
\subsubsection[{\small FABSH: Floating-Point Absolute Value Half Word}]{FABSH\hfill Floating-Point Absolute Value Half Word}
\label{instr:FABSH}\index{fabsh}
 
\paragraph{Encoding} Format \textsf{ALU\_FHMO1} (Section~\ref{sec:format-ALU-FHMO1}), \verb|fpucode6 = 0011|\\

\bitfields{ 1/P, 2/11, 1/1, 1/1, 3/111, 6/registerW, 2/01, 3/000, 1/1, 4/0011, 1/--, 1/0, 6/registerZ }


\paragraph{Syntax} \texttt{fabsh} \emph{registerW} = \emph{registerZ}

\paragraph{Description} The absolute value the \emph{registerZ}, assuming binary 16 floating-point number, is stored into the \emph{registerW}. Implemented as a bitwise operation, does not raise FP exceptions.


\paragraph{Execution} {Reservation table \textsf{ALU\_LITE} (Section~\ref{sec:reservation-ALU-LITE})}
~
{\begin{lstlisting}
stage RR:
  new argument2 = GPR[registerZ];
  new result1 = _ZX_16(argument2 & 0x7FFF);
stage E1:
  GPR[registerW] = result1;
\end{lstlisting}}
\newpage
\subsubsection[{\small FABSHQ: Floating-Point Absolute Value Half Word Quadruple}]{FABSHQ\hfill Floating-Point Absolute Value Half Word Quadruple}
\label{instr:FABSHQ}\index{fabshq}
 
\paragraph{Encoding} Format \textsf{ALU\_FHQMO1E} (Section~\ref{sec:format-ALU-FHQMO1E}), \verb|fpucode6 = 0011|\\

\bitfields{ 1/P, 2/11, 1/1, 1/1, 3/111, 5/registerWe, 1/0, 2/11, 3/000, 1/1, 4/0011, 1/--, 1/0, 5/registerZe, 1/1 }


\paragraph{Syntax} \texttt{fabshq} \emph{registerWe} = \emph{registerZe}

\paragraph{Description} The \emph{registerZe} is interpreted as four binary 16 floating-point numbers packed into 64 bits. The four absolute values the \emph{registerZe} half words are packed and stored into the \emph{registerWe}.


\paragraph{Execution} {Reservation table \textsf{ALU\_LITE} (Section~\ref{sec:reservation-ALU-LITE})}
~
{\begin{lstlisting}
stage RR:
  new argument2 = GPR[registerZe<<1];
  for (I = 0; I < 4; I += 1) {
    result1.16[I] = argument2.16[I] & 0x7FFF
  };
stage E1:
  GPR[registerWe<<1] = result1;
\end{lstlisting}}
\newpage

\paragraph{Encoding} Format \textsf{ALU\_FHQMO1O} (Section~\ref{sec:format-ALU-FHQMO1O}), \verb|fpucode6 = 0011|\\

\bitfields{ 1/P, 2/11, 1/1, 1/1, 3/111, 5/registerWo, 1/1, 2/11, 3/000, 1/1, 4/0011, 1/--, 1/0, 5/registerZo, 1/1 }


\paragraph{Syntax} \texttt{fabshq} \emph{registerWo} = \emph{registerZo}

\paragraph{Description} The \emph{registerZo} is interpreted as four binary 16 floating-point numbers packed into 64 bits. The four absolute values the \emph{registerZo} half words are packed and stored into the \emph{registerWo}.


\paragraph{Execution} {Reservation table \textsf{ALU\_LITE} (Section~\ref{sec:reservation-ALU-LITE})}
~
{\begin{lstlisting}
stage RR:
  new argument2 = GPR[(registerZo<<1)+1];
  for (I = 0; I < 4; I += 1) {
    result1.16[I] = argument2.16[I] & 0x7FFF
  };
stage E1:
  GPR[(registerWo<<1)+1] = result1;
\end{lstlisting}}
\newpage
\subsubsection[{\small FABSW: Floating-Point Absolute Value Word}]{FABSW\hfill Floating-Point Absolute Value Word}
\label{instr:FABSW}\index{fabsw}
 
\paragraph{Encoding} Format \textsf{ALU\_FWMO1} (Section~\ref{sec:format-ALU-FWMO1}), \verb|fpucode6 = 0011|\\

\bitfields{ 1/P, 2/11, 1/1, 1/1, 3/111, 6/registerW, 2/01, 3/000, 1/0, 4/0011, 1/--, 1/0, 6/registerZ }


\paragraph{Syntax} \texttt{fabsw} \emph{registerW} = \emph{registerZ}

\paragraph{Description} The absolute value the \emph{registerZ}, assuming binary 32 floating-point number, is stored into the \emph{registerW}. Implemented as a bitwise operation, does not raise FP exceptions.


\paragraph{Execution} {Reservation table \textsf{ALU\_LITE} (Section~\ref{sec:reservation-ALU-LITE})}
~
{\begin{lstlisting}
stage RR:
  new argument2 = GPR[registerZ];
  new result1 = _ZX_32(argument2 & 0x7FFFFFFF);
stage E1:
  GPR[registerW] = result1;
\end{lstlisting}}
\newpage
\subsubsection[{\small FABSWP: Floating-Point Absolute Value Word Pair}]{FABSWP\hfill Floating-Point Absolute Value Word Pair}
\label{instr:FABSWP}\index{fabswp}
 
\paragraph{Encoding} Format \textsf{ALU\_FWPMO1} (Section~\ref{sec:format-ALU-FWPMO1}), \verb|fpucode6 = 0011|\\

\bitfields{ 1/P, 2/11, 1/1, 1/1, 3/111, 6/registerW, 2/00, 3/000, 1/0, 4/0011, 1/--, 1/1, 6/registerZ }


\paragraph{Syntax} \texttt{fabswp} \emph{registerW} = \emph{registerZ}

\paragraph{Description} The \emph{registerZ} is interpreted as two binary 32 floating-point numbers packed into 64 bits. The two absolute values the \emph{registerZ} words are packed and stored into the \emph{registerW}.


\paragraph{Execution} {Reservation table \textsf{ALU\_LITE} (Section~\ref{sec:reservation-ALU-LITE})}
~
{\begin{lstlisting}
stage RR:
  new argument2 = GPR[registerZ];
  for (I = 0; I < 2; I += 1) {
    result1.32[I] = argument2.32[I] & 0x7FFFFFFF
  };
stage E1:
  GPR[registerW] = result1;
\end{lstlisting}}
\newpage
\subsubsection[{\small FADDD: Floating-Point Add Double Word}]{FADDD\hfill Floating-Point Add Double Word}
\label{instr:FADDD}\index{faddd}
 
\paragraph{Encoding} Format \textsf{ALU\_FDADD} (Section~\ref{sec:format-ALU-FDADD}), \verb|alucode3 = 000|\\

\bitfields{ 1/P, 2/11, 1/1, 1/0, 3/floatmode, 6/registerW, 2/00, 3/000, 1/0, 6/registerY, 6/registerZ }


\paragraph{Syntax} \texttt{faddd}\emph{floatmode} \emph{registerW} = \emph{registerZ}, \emph{registerY}

\paragraph{Description} The \emph{registerZ} is added to the \emph{registerY}, assuming binary 64 floating-point numbers. The result is rounded to the binary 64 floating-point format according to the rounding mode and stored into the \emph{registerW}. This instruction may raise exception bits in the CS register.


\paragraph{Modifier(s)}
\noindent\begin{itemize}
\item floatmode (cf. floatmode in section \ref{par:modifier-floatmode} on page \pageref{par:modifier-floatmode})
\end{itemize}

\paragraph{Execution} {Reservation table \textsf{ALU\_LITE} (Section~\ref{sec:reservation-ALU-LITE})}
~
{\begin{lstlisting}
stage ID:
  new floatmode = _ZX_3(floatmode);
stage RR:
  new argument3 = GPR[registerY];
  new argument2 = GPR[registerZ];
  new RM = decode_riscv_float_rounding_mode(floatmode);
  new result1 = f64_add(RM, argument2, argument3);
stage E4:
  GPR[registerW] = result1;
  commit_float_exception_flags();
\end{lstlisting}}
\newpage
\subsubsection[{\small FADDH: Floating-Point Add Half Word}]{FADDH\hfill Floating-Point Add Half Word}
\label{instr:FADDH}\index{faddh}
 
\paragraph{Encoding} Format \textsf{ALU\_FHADD} (Section~\ref{sec:format-ALU-FHADD}), \verb|alucode3 = 000|\\

\bitfields{ 1/P, 2/11, 1/1, 1/0, 3/floatmode, 6/registerW, 2/01, 3/000, 1/1, 6/registerY, 6/registerZ }


\paragraph{Syntax} \texttt{faddh}\emph{floatmode} \emph{registerW} = \emph{registerZ}, \emph{registerY}

\paragraph{Description} The \emph{registerZ} is added to the \emph{registerY}, assuming binary 16 floating-point numbers. The result is rounded to the binary 16 floating-point format according to the rounding mode and stored into the \emph{registerW}. This instruction may raise exception bits in the CS register.


\paragraph{Modifier(s)}
\noindent\begin{itemize}
\item floatmode (cf. floatmode in section \ref{par:modifier-floatmode} on page \pageref{par:modifier-floatmode})
\end{itemize}

\paragraph{Execution} {Reservation table \textsf{ALU\_LITE} (Section~\ref{sec:reservation-ALU-LITE})}
~
{\begin{lstlisting}
stage ID:
  new floatmode = _ZX_3(floatmode);
stage RR:
  new argument3 = GPR[registerY];
  new argument2 = GPR[registerZ];
  new RM = decode_riscv_float_rounding_mode(floatmode);
  new result1 = f16_add(RM, argument2, argument3);
stage E3:
  GPR[registerW] = result1;
  commit_float_exception_flags();
\end{lstlisting}}
\newpage
\subsubsection[{\small FADDHQ: Floating-Point Add Half Word Quadruple}]{FADDHQ\hfill Floating-Point Add Half Word Quadruple}
\label{instr:FADDHQ}\index{faddhq}
 
\paragraph{Encoding} Format \textsf{ALU\_FHQADDE} (Section~\ref{sec:format-ALU-FHQADDE}), \verb|alucode3 = 000|\\

\bitfields{ 1/P, 2/11, 1/1, 1/0, 3/floatmode, 5/registerWe, 1/0, 2/11, 3/000, 1/1, 5/registerYe, 1/--, 5/registerZe, 1/1 }


\paragraph{Syntax} \texttt{faddhq}\emph{floatmode} \emph{registerWe} = \emph{registerZe}, \emph{registerYe}

\paragraph{Description} The \emph{registerZe} and the \emph{registerYe} are interpreted as four binary 16 floating-point numbers packed into 64 bits. The \emph{registerZe} half words are added to the \emph{registerYe} half words. The four resulting half words are packed and stored into the \emph{registerWe}.


\paragraph{Modifier(s)}
\noindent\begin{itemize}
\item floatmode (cf. floatmode in section \ref{par:modifier-floatmode} on page \pageref{par:modifier-floatmode})
\end{itemize}

\paragraph{Execution} {Reservation table \textsf{ALU\_LITE} (Section~\ref{sec:reservation-ALU-LITE})}
~
{\begin{lstlisting}
stage ID:
  new floatmode = _ZX_3(floatmode);
stage RR:
  new argument3 = GPR[registerYe<<1];
  new argument2 = GPR[registerZe<<1];
  new RM = decode_riscv_float_rounding_mode(floatmode);
  for (I = 0; I < 4; I += 1) {
    result1.16[I] = f16_add(RM, argument2.16[I], argument3.16[I])
  };
stage E4:
  GPR[registerWe<<1] = result1;
  commit_float_exception_flags();
\end{lstlisting}}
\newpage

\paragraph{Encoding} Format \textsf{ALU\_FHQADDO} (Section~\ref{sec:format-ALU-FHQADDO}), \verb|alucode3 = 000|\\

\bitfields{ 1/P, 2/11, 1/1, 1/0, 3/floatmode, 5/registerWo, 1/1, 2/11, 3/000, 1/1, 5/registerYo, 1/--, 5/registerZo, 1/1 }


\paragraph{Syntax} \texttt{faddhq}\emph{floatmode} \emph{registerWo} = \emph{registerZo}, \emph{registerYo}

\paragraph{Description} The \emph{registerZo} and the \emph{registerYo} are interpreted as four binary 16 floating-point numbers packed into 64 bits. The \emph{registerZo} half words are added to the \emph{registerYo} half words. The four resulting half words are packed and stored into the \emph{registerWo}.


\paragraph{Modifier(s)}
\noindent\begin{itemize}
\item floatmode (cf. floatmode in section \ref{par:modifier-floatmode} on page \pageref{par:modifier-floatmode})
\end{itemize}

\paragraph{Execution} {Reservation table \textsf{ALU\_LITE} (Section~\ref{sec:reservation-ALU-LITE})}
~
{\begin{lstlisting}
stage ID:
  new floatmode = _ZX_3(floatmode);
stage RR:
  new argument3 = GPR[(registerYo<<1)+1];
  new argument2 = GPR[(registerZo<<1)+1];
  new RM = decode_riscv_float_rounding_mode(floatmode);
  for (I = 0; I < 4; I += 1) {
    result1.16[I] = f16_add(RM, argument2.16[I], argument3.16[I])
  };
stage E4:
  GPR[(registerWo<<1)+1] = result1;
  commit_float_exception_flags();
\end{lstlisting}}
\newpage
\subsubsection[{\small FADDW: Floating-Point Add Word}]{FADDW\hfill Floating-Point Add Word}
\label{instr:FADDW}\index{faddw}
 
\paragraph{Encoding} Format \textsf{ALU\_FWADD} (Section~\ref{sec:format-ALU-FWADD}), \verb|alucode3 = 000|\\

\bitfields{ 1/P, 2/11, 1/1, 1/0, 3/floatmode, 6/registerW, 2/01, 3/000, 1/0, 6/registerY, 6/registerZ }


\paragraph{Syntax} \texttt{faddw}\emph{floatmode} \emph{registerW} = \emph{registerZ}, \emph{registerY}

\paragraph{Description} The \emph{registerZ} is added to the \emph{registerY}, assuming binary 32 floating-point numbers. The result is rounded to the binary 32 floating-point format according to the rounding mode and stored into the \emph{registerW}. This instruction may raise exception bits in the CS register.


\paragraph{Modifier(s)}
\noindent\begin{itemize}
\item floatmode (cf. floatmode in section \ref{par:modifier-floatmode} on page \pageref{par:modifier-floatmode})
\end{itemize}

\paragraph{Execution} {Reservation table \textsf{ALU\_LITE} (Section~\ref{sec:reservation-ALU-LITE})}
~
{\begin{lstlisting}
stage ID:
  new floatmode = _ZX_3(floatmode);
stage RR:
  new argument3 = GPR[registerY];
  new argument2 = GPR[registerZ];
  new RM = decode_riscv_float_rounding_mode(floatmode);
  new result1 = f32_add(RM, argument2, argument3);
stage E4:
  GPR[registerW] = result1;
  commit_float_exception_flags();
\end{lstlisting}}
\newpage
\subsubsection[{\small FADDWC: Floating-Point Add to Word Complex}]{FADDWC\hfill Floating-Point Add to Word Complex}
\label{instr:FADDWC}\index{faddwc}
 
\paragraph{Encoding} Format \textsf{ALU\_FWCOP} (Section~\ref{sec:format-ALU-FWCOP}), \verb|alucode6 = 00|\\

\bitfields{ 1/P, 2/11, 1/1, 1/imultiply, 3/floatmode, 6/registerW, 2/00, 2/00, 1/conjugate, 1/1, 6/registerY, 6/registerZ }


\paragraph{Syntax} \texttt{faddwc}\emph{conjugate}\emph{imultiply}\emph{floatmode} \emph{registerW} = \emph{registerZ}, \emph{registerY}

\paragraph{Description} The \emph{registerY} and the \emph{registerZ} are interpreted as binary 32 floating-point complex numbers. If \emph{conjugate} is set, the complex conjugate of the \emph{registerZ} is used in the operation. The real and imaginary parts are respectively added to and rounded according to the rounding mode. If \emph{imultiply} is set, the result is multiplied by the imaginary unit. The resulting binary 32 floating-point complex number is stored back into the \emph{registerW}. This instruction may raise exception bits in the CS register.


\paragraph{Modifier(s)}
\noindent\begin{itemize}
\item floatmode (cf. floatmode in section \ref{par:modifier-floatmode} on page \pageref{par:modifier-floatmode})
\item imultiply (cf. imultiply in section \ref{par:modifier-imultiply} on page \pageref{par:modifier-imultiply})
\item conjugate (cf. conjugate in section \ref{par:modifier-conjugate} on page \pageref{par:modifier-conjugate})
\end{itemize}

\paragraph{Execution} {Reservation table \textsf{ALU\_LITE} (Section~\ref{sec:reservation-ALU-LITE})}
~
{\begin{lstlisting}
stage ID:
  new conjugate = _ZX_1(conjugate);
  new imultiply = _ZX_1(imultiply);
  new floatmode = _ZX_3(floatmode);
stage RR:
  new argument3 = GPR[registerY];
  new argument2 = GPR[registerZ];
  new argument1 = GPR[registerW];
  new RM = decode_riscv_float_rounding_mode(floatmode);
  new fconj = conjugate ? 0x80000000 : 0;
  result1.32[0] = f32_add(RM, argument2.32[0], argument3.32[0]);
  result1.32[1] = f32_add(RM, argument2.32[1] ^ fconj, argument3.32[1]);
  if (imultiply) {
    new result1 = _SWAP_32(result1);
    result1.32[0] = result1.32[0] ^ 0x80000000;
  }
stage E4:
  GPR[registerW] = result1;
  commit_float_exception_flags();
\end{lstlisting}}
\newpage
\subsubsection[{\small FADDWP: Floating-Point Add Word Pair}]{FADDWP\hfill Floating-Point Add Word Pair}
\label{instr:FADDWP}\index{faddwp}
 
\paragraph{Encoding} Format \textsf{ALU\_FWPADDE} (Section~\ref{sec:format-ALU-FWPADDE}), \verb|alucode3 = 000|\\

\bitfields{ 1/P, 2/11, 1/1, 1/0, 3/floatmode, 5/registerWe, 1/0, 2/11, 3/000, 1/1, 5/registerYe, 1/--, 5/registerZe, 1/0 }


\paragraph{Syntax} \texttt{faddwp}\emph{floatmode} \emph{registerWe} = \emph{registerZe}, \emph{registerYe}

\paragraph{Description} The \emph{registerZe} and the \emph{registerYe} are interpreted as two binary 32 floating-point numbers packed into 64 bits. The \emph{registerZe} words are added to the \emph{registerYe} words. The two resulting words are packed and stored into the \emph{registerWe}.


\paragraph{Modifier(s)}
\noindent\begin{itemize}
\item floatmode (cf. floatmode in section \ref{par:modifier-floatmode} on page \pageref{par:modifier-floatmode})
\end{itemize}

\paragraph{Execution} {Reservation table \textsf{ALU\_LITE} (Section~\ref{sec:reservation-ALU-LITE})}
~
{\begin{lstlisting}
stage ID:
  new floatmode = _ZX_3(floatmode);
stage RR:
  new argument3 = GPR[registerYe<<1];
  new argument2 = GPR[registerZe<<1];
  new RM = decode_riscv_float_rounding_mode(floatmode);
  for (I = 0; I < 2; I += 1) {
    result1.32[I] = f32_add(RM, argument2.32[I], argument3.32[I])
  };
stage E4:
  GPR[registerWe<<1] = result1;
  commit_float_exception_flags();
\end{lstlisting}}
\newpage

\paragraph{Encoding} Format \textsf{ALU\_FWPADDO} (Section~\ref{sec:format-ALU-FWPADDO}), \verb|alucode3 = 000|\\

\bitfields{ 1/P, 2/11, 1/1, 1/0, 3/floatmode, 5/registerWo, 1/1, 2/11, 3/000, 1/1, 5/registerYo, 1/--, 5/registerZo, 1/0 }


\paragraph{Syntax} \texttt{faddwp}\emph{floatmode} \emph{registerWo} = \emph{registerZo}, \emph{registerYo}

\paragraph{Description} The \emph{registerZo} and the \emph{registerYo} are interpreted as two binary 32 floating-point numbers packed into 64 bits. The \emph{registerZo} words are added to the \emph{registerYo} words. The two resulting words are packed and stored into the \emph{registerWo}.


\paragraph{Modifier(s)}
\noindent\begin{itemize}
\item floatmode (cf. floatmode in section \ref{par:modifier-floatmode} on page \pageref{par:modifier-floatmode})
\end{itemize}

\paragraph{Execution} {Reservation table \textsf{ALU\_LITE} (Section~\ref{sec:reservation-ALU-LITE})}
~
{\begin{lstlisting}
stage ID:
  new floatmode = _ZX_3(floatmode);
stage RR:
  new argument3 = GPR[(registerYo<<1)+1];
  new argument2 = GPR[(registerZo<<1)+1];
  new RM = decode_riscv_float_rounding_mode(floatmode);
  for (I = 0; I < 2; I += 1) {
    result1.32[I] = f32_add(RM, argument2.32[I], argument3.32[I])
  };
stage E4:
  GPR[(registerWo<<1)+1] = result1;
  commit_float_exception_flags();
\end{lstlisting}}
\newpage
\subsubsection[{\small FCLASSD: Floating-Point Classify Double Word}]{FCLASSD\hfill Floating-Point Classify Double Word}
\label{instr:FCLASSD}\index{fclassd}
 
\paragraph{Encoding} Format \textsf{ALU\_FDMO1} (Section~\ref{sec:format-ALU-FDMO1}), \verb|fpucode6 = 0110|\\

\bitfields{ 1/P, 2/11, 1/1, 1/1, 3/111, 6/registerW, 2/00, 3/000, 1/0, 4/0110, 1/--, 1/0, 6/registerZ }


\paragraph{Syntax} \texttt{fclassd} \emph{registerW} = \emph{registerZ}

\paragraph{Description} Classify the \emph{registerZ}, assuming a binary 64 floating-point number. is stored into the \emph{registerW}.


\paragraph{Execution} {Reservation table \textsf{ALU\_LITE} (Section~\ref{sec:reservation-ALU-LITE})}
~
{\begin{lstlisting}
stage RR:
  new argument2 = GPR[registerZ];
  new result1 = f64_classify(argument2);
stage E1:
  GPR[registerW] = result1;
\end{lstlisting}}
\newpage
\subsubsection[{\small FCLASSH: Floating-Point Classify Half Word}]{FCLASSH\hfill Floating-Point Classify Half Word}
\label{instr:FCLASSH}\index{fclassh}
 
\paragraph{Encoding} Format \textsf{ALU\_FHMO1} (Section~\ref{sec:format-ALU-FHMO1}), \verb|fpucode6 = 0110|\\

\bitfields{ 1/P, 2/11, 1/1, 1/1, 3/111, 6/registerW, 2/01, 3/000, 1/1, 4/0110, 1/--, 1/0, 6/registerZ }


\paragraph{Syntax} \texttt{fclassh} \emph{registerW} = \emph{registerZ}

\paragraph{Description} Classify the \emph{registerZ}, assuming a binary 16 floating-point number. is stored into the \emph{registerW}.


\paragraph{Execution} {Reservation table \textsf{ALU\_LITE} (Section~\ref{sec:reservation-ALU-LITE})}
~
{\begin{lstlisting}
stage RR:
  new argument2 = GPR[registerZ];
  new result1 = f16_classify(argument2);
stage E1:
  GPR[registerW] = result1;
\end{lstlisting}}
\newpage
\subsubsection[{\small FCLASSHQ: Floating-Point Classify Half Word Quadruple}]{FCLASSHQ\hfill Floating-Point Classify Half Word Quadruple}
\label{instr:FCLASSHQ}\index{fclasshq}
 
\paragraph{Encoding} Format \textsf{ALU\_FHQMO1E} (Section~\ref{sec:format-ALU-FHQMO1E}), \verb|fpucode6 = 0110|\\

\bitfields{ 1/P, 2/11, 1/1, 1/1, 3/111, 5/registerWe, 1/0, 2/11, 3/000, 1/1, 4/0110, 1/--, 1/0, 5/registerZe, 1/1 }


\paragraph{Syntax} \texttt{fclasshq} \emph{registerWe} = \emph{registerZe}

\paragraph{Description} The \emph{registerZe} is interpreted as four binary 16 floating-point numbers packed into 64 bits. The four binary 16 floating-point number are classified. The resulting four half words are packed and stored into the \emph{registerWe}.


\paragraph{Execution} {Reservation table \textsf{ALU\_LITE} (Section~\ref{sec:reservation-ALU-LITE})}
~
{\begin{lstlisting}
stage RR:
  new argument2 = GPR[registerZe<<1];
  for (I = 0; I < 4; I += 1) {
    result1.16[I] = f16_classify(argument2.16[I])
  };
stage E1:
  GPR[registerWe<<1] = result1;
\end{lstlisting}}
\newpage

\paragraph{Encoding} Format \textsf{ALU\_FHQMO1O} (Section~\ref{sec:format-ALU-FHQMO1O}), \verb|fpucode6 = 0110|\\

\bitfields{ 1/P, 2/11, 1/1, 1/1, 3/111, 5/registerWo, 1/1, 2/11, 3/000, 1/1, 4/0110, 1/--, 1/0, 5/registerZo, 1/1 }


\paragraph{Syntax} \texttt{fclasshq} \emph{registerWo} = \emph{registerZo}

\paragraph{Description} The \emph{registerZo} is interpreted as four binary 16 floating-point numbers packed into 64 bits. The four binary 16 floating-point number are classified. The resulting four half words are packed and stored into the \emph{registerWo}.


\paragraph{Execution} {Reservation table \textsf{ALU\_LITE} (Section~\ref{sec:reservation-ALU-LITE})}
~
{\begin{lstlisting}
stage RR:
  new argument2 = GPR[(registerZo<<1)+1];
  for (I = 0; I < 4; I += 1) {
    result1.16[I] = f16_classify(argument2.16[I])
  };
stage E1:
  GPR[(registerWo<<1)+1] = result1;
\end{lstlisting}}
\newpage
\subsubsection[{\small FCLASSW: Floating-Point Classify Word}]{FCLASSW\hfill Floating-Point Classify Word}
\label{instr:FCLASSW}\index{fclassw}
 
\paragraph{Encoding} Format \textsf{ALU\_FWMO1} (Section~\ref{sec:format-ALU-FWMO1}), \verb|fpucode6 = 0110|\\

\bitfields{ 1/P, 2/11, 1/1, 1/1, 3/111, 6/registerW, 2/01, 3/000, 1/0, 4/0110, 1/--, 1/0, 6/registerZ }


\paragraph{Syntax} \texttt{fclassw} \emph{registerW} = \emph{registerZ}

\paragraph{Description} Classify the \emph{registerZ}, assuming a binary 32 floating-point number. is stored into the \emph{registerW}.


\paragraph{Execution} {Reservation table \textsf{ALU\_LITE} (Section~\ref{sec:reservation-ALU-LITE})}
~
{\begin{lstlisting}
stage RR:
  new argument2 = GPR[registerZ];
  new result1 = f32_classify(argument2);
stage E1:
  GPR[registerW] = result1;
\end{lstlisting}}
\newpage
\subsubsection[{\small FCLASSWP: Floating-Point Classify Word Pair}]{FCLASSWP\hfill Floating-Point Classify Word Pair}
\label{instr:FCLASSWP}\index{fclasswp}
 
\paragraph{Encoding} Format \textsf{ALU\_FWPMO1} (Section~\ref{sec:format-ALU-FWPMO1}), \verb|fpucode6 = 0110|\\

\bitfields{ 1/P, 2/11, 1/1, 1/1, 3/111, 6/registerW, 2/00, 3/000, 1/0, 4/0110, 1/--, 1/1, 6/registerZ }


\paragraph{Syntax} \texttt{fclasswp} \emph{registerW} = \emph{registerZ}

\paragraph{Description} The \emph{registerZ} is interpreted as two binary 32 floating-point numbers packed into 64 bits. The two binary 32 floating-point number are classified. The resulting two words are packed and stored into the \emph{registerW}.


\paragraph{Execution} {Reservation table \textsf{ALU\_LITE} (Section~\ref{sec:reservation-ALU-LITE})}
~
{\begin{lstlisting}
stage RR:
  new argument2 = GPR[registerZ];
  for (I = 0; I < 2; I += 1) {
    result1.32[I] = f32_classify(argument2.32[I])
  };
stage E1:
  GPR[registerW] = result1;
\end{lstlisting}}
\newpage
\subsubsection[{\small FCOMPD: Floating-Point Compare Double Word}]{FCOMPD\hfill Floating-Point Compare Double Word}
\label{instr:FCOMPD}\index{fcompd}
 
\paragraph{Encoding} Format \textsf{ALU\_FDCWRR} (Section~\ref{sec:format-ALU-FDCWRR}), \verb|alucode3 = 010|\\

\bitfields{ 1/P, 2/11, 1/1, 1/0, 3/floatcomp, 6/registerW, 2/00, 3/010, 1/0, 6/registerY, 6/registerZ }


\paragraph{Syntax} \texttt{fcompd}\emph{floatcomp} \emph{registerW} = \emph{registerZ}, \emph{registerY}

\paragraph{Description} The \emph{registerZ} is compared to the \emph{registerY} using condition \emph{floatcomp}, assuming binary 64 floating-point numbers. The boolean result is stored into the \emph{registerW}. This instruction does not raise an invalid exception when handling a signaling NaN.


\paragraph{Modifier(s)}
\noindent\begin{itemize}
\item floatcomp (cf. floatcomp in section \ref{par:modifier-floatcomp} on page \pageref{par:modifier-floatcomp})
\end{itemize}

\paragraph{Execution} {Reservation table \textsf{ALU\_LITE} (Section~\ref{sec:reservation-ALU-LITE})}
~
{\begin{lstlisting}
stage ID:
  new argument4 = _ZX_3(floatcomp);
stage RR:
  new argument3 = GPR[registerY];
  new argument2 = GPR[registerZ];
  new result1 = floatcomp_64(argument4, argument2, argument3);
stage E1:
  GPR[registerW] = result1;
\end{lstlisting}}
\newpage

\paragraph{Encoding} Format \textsf{ALU\_FDCWRR.W} (Section~\ref{sec:format-ALU-FDCWRR.W}), \verb|alucode3 = 010|\\

\bitfields{ 1/P, 2/00, 2/-- --, 27/upper27, 1/1, 2/11, 1/1, 1/0, 3/floatcomp, 6/registerW, 2/00, 3/010, 1/0, 1/0, 5/lower5, 6/registerZ }


\paragraph{Syntax} \texttt{fcompd}\emph{floatcomp} \emph{registerW} = \emph{registerZ}, \emph{upper27\_\discretionary{}{}{}lower5}

\paragraph{Description} The \emph{registerZ} is compared to the \emph{upper27\_\discretionary{}{}{}lower5} using condition \emph{floatcomp}, assuming binary 64 floating-point numbers. The boolean result is stored into the \emph{registerW}. This instruction does not raise an invalid exception when handling a signaling NaN.


\paragraph{Modifier(s)}
\noindent\begin{itemize}
\item floatcomp (cf. floatcomp in section \ref{par:modifier-floatcomp} on page \pageref{par:modifier-floatcomp})
\end{itemize}

\paragraph{Execution} {Reservation table \textsf{ALU\_LITE.X} (Section~\ref{sec:reservation-ALU-LITE.X})}
~
{\begin{lstlisting}
stage ID:
  new argument4 = _ZX_3(floatcomp);
stage RR:
  new argument3 = _WX_32(upper27_lower5);
  new argument2 = GPR[registerZ];
  new result1 = floatcomp_64(argument4, argument2, argument3);
stage E1:
  GPR[registerW] = result1;
\end{lstlisting}}
\newpage
\subsubsection[{\small FCOMPH: Floating-Point Compare Half Word}]{FCOMPH\hfill Floating-Point Compare Half Word}
\label{instr:FCOMPH}\index{fcomph}
 
\paragraph{Encoding} Format \textsf{ALU\_FHCWRR} (Section~\ref{sec:format-ALU-FHCWRR}), \verb|alucode3 = 010|\\

\bitfields{ 1/P, 2/11, 1/1, 1/0, 3/floatcomp, 6/registerW, 2/01, 3/010, 1/1, 6/registerY, 6/registerZ }


\paragraph{Syntax} \texttt{fcomph}\emph{floatcomp} \emph{registerW} = \emph{registerZ}, \emph{registerY}

\paragraph{Description} The \emph{registerZ} is compared to the \emph{registerY} using condition \emph{floatcomp}, assuming binary 16 floating-point numbers. The boolean result is stored into the \emph{registerW}. This instruction does not raise an invalid exception when handling a signaling NaN.


\paragraph{Modifier(s)}
\noindent\begin{itemize}
\item floatcomp (cf. floatcomp in section \ref{par:modifier-floatcomp} on page \pageref{par:modifier-floatcomp})
\end{itemize}

\paragraph{Execution} {Reservation table \textsf{ALU\_LITE} (Section~\ref{sec:reservation-ALU-LITE})}
~
{\begin{lstlisting}
stage ID:
  new argument4 = _ZX_3(floatcomp);
stage RR:
  new argument3 = GPR[registerY];
  new argument2 = GPR[registerZ];
  new result1 = floatcomp_16(argument4, argument2, argument3);
stage E1:
  GPR[registerW] = _ZX_32(result1);
\end{lstlisting}}
\newpage

\paragraph{Encoding} Format \textsf{ALU\_FHCWRR.W} (Section~\ref{sec:format-ALU-FHCWRR.W}), \verb|alucode3 = 010|\\

\bitfields{ 1/P, 2/00, 2/-- --, 27/upper27, 1/1, 2/11, 1/1, 1/0, 3/floatcomp, 6/registerW, 2/01, 3/010, 1/1, 1/0, 5/lower5, 6/registerZ }


\paragraph{Syntax} \texttt{fcomph}\emph{floatcomp} \emph{registerW} = \emph{registerZ}, \emph{upper27\_\discretionary{}{}{}lower5}

\paragraph{Description} The \emph{registerZ} is compared to the \emph{upper27\_\discretionary{}{}{}lower5} using condition \emph{floatcomp}, assuming binary 16 floating-point numbers. The boolean result is stored into the \emph{registerW}. This instruction does not raise an invalid exception when handling a signaling NaN.


\paragraph{Modifier(s)}
\noindent\begin{itemize}
\item floatcomp (cf. floatcomp in section \ref{par:modifier-floatcomp} on page \pageref{par:modifier-floatcomp})
\end{itemize}

\paragraph{Execution} {Reservation table \textsf{ALU\_LITE.X} (Section~\ref{sec:reservation-ALU-LITE.X})}
~
{\begin{lstlisting}
stage ID:
  new argument4 = _ZX_3(floatcomp);
stage RR:
  new argument3 = _WX_32(upper27_lower5);
  new argument2 = GPR[registerZ];
  new result1 = floatcomp_16(argument4, argument2, argument3);
stage E1:
  GPR[registerW] = _ZX_32(result1);
\end{lstlisting}}
\newpage
\subsubsection[{\small FCOMPND: Floating-Point Compare Negate Double Word}]{FCOMPND\hfill Floating-Point Compare Negate Double Word}
\label{instr:FCOMPND}\index{fcompnd}
 
\paragraph{Encoding} Format \textsf{ALU\_FDCNWRRE} (Section~\ref{sec:format-ALU-FDCNWRRE}), \verb|alucode3 = 010|\\

\bitfields{ 1/P, 2/11, 1/1, 1/1, 3/floatcomp, 5/registerWe, 1/0, 2/11, 3/010, 1/0, 5/registerYe, 1/0, 5/registerZe, 1/1 }


\paragraph{Syntax} \texttt{fcompnd}\emph{floatcomp} \emph{registerWe} = \emph{registerZe}, \emph{registerYe}

\paragraph{Description} The \emph{registerZe} and the \emph{registerYe} are interpreted as two binary 64 floating-point numbers. The \emph{registerZe} double word is compared to the \emph{registerYe} double word using condition \emph{floatcomp}. The negated boolean result extended to double words is stored into the \emph{registerWe}. This instruction does not raise an invalid exception when handling a signaling NaN.


\paragraph{Modifier(s)}
\noindent\begin{itemize}
\item floatcomp (cf. floatcomp in section \ref{par:modifier-floatcomp} on page \pageref{par:modifier-floatcomp})
\end{itemize}

\paragraph{Execution} {Reservation table \textsf{ALU\_LITE} (Section~\ref{sec:reservation-ALU-LITE})}
~
{\begin{lstlisting}
stage ID:
  new argument4 = _ZX_3(floatcomp);
stage RR:
  new argument3 = GPR[registerYe<<1];
  new argument2 = GPR[registerZe<<1];
  new result1 = -floatcomp_64(argument4, argument2, argument3);
stage E1:
  GPR[registerWe<<1] = result1;
\end{lstlisting}}
\newpage

\paragraph{Encoding} Format \textsf{ALU\_FDCNWRRO} (Section~\ref{sec:format-ALU-FDCNWRRO}), \verb|alucode3 = 010|\\

\bitfields{ 1/P, 2/11, 1/1, 1/1, 3/floatcomp, 5/registerWo, 1/1, 2/11, 3/010, 1/0, 5/registerYo, 1/0, 5/registerZo, 1/1 }


\paragraph{Syntax} \texttt{fcompnd}\emph{floatcomp} \emph{registerWo} = \emph{registerZo}, \emph{registerYo}

\paragraph{Description} The \emph{registerZo} and the \emph{registerYo} are interpreted as two binary 64 floating-point numbers. The \emph{registerZo} double word is compared to the \emph{registerYo} double word using condition \emph{floatcomp}. The negated boolean result extended to double words is stored into the \emph{registerWo}. This instruction does not raise an invalid exception when handling a signaling NaN.


\paragraph{Modifier(s)}
\noindent\begin{itemize}
\item floatcomp (cf. floatcomp in section \ref{par:modifier-floatcomp} on page \pageref{par:modifier-floatcomp})
\end{itemize}

\paragraph{Execution} {Reservation table \textsf{ALU\_LITE} (Section~\ref{sec:reservation-ALU-LITE})}
~
{\begin{lstlisting}
stage ID:
  new argument4 = _ZX_3(floatcomp);
stage RR:
  new argument3 = GPR[(registerYo<<1)+1];
  new argument2 = GPR[(registerZo<<1)+1];
  new result1 = -floatcomp_64(argument4, argument2, argument3);
stage E1:
  GPR[(registerWo<<1)+1] = result1;
\end{lstlisting}}
\newpage

\paragraph{Encoding} Format \textsf{ALU\_FDCNWRRE.M} (Section~\ref{sec:format-ALU-FDCNWRRE.M}), \verb|alucode3 = 011|\\

\bitfields{ 1/P, 2/00, 2/-- --, 27/upper27, 1/1, 2/11, 1/1, 1/1, 3/floatcomp, 5/registerWe, 1/0, 2/11, 3/011, 1/0, 1/splat32, 5/lower5, 5/registerZe, 1/1 }


\paragraph{Syntax} \texttt{fcompnd}\emph{floatcomp} \emph{registerWe} = \emph{registerZe}, \emph{upper27\_\discretionary{}{}{}lower5}\emph{splat32}

\paragraph{Description} The \emph{registerZe} and the \emph{upper27\_\discretionary{}{}{}lower5} are interpreted as two binary 64 floating-point numbers. The \emph{registerZe} double word is compared to the \emph{upper27\_\discretionary{}{}{}lower5} double word using condition \emph{floatcomp}. The negated boolean result extended to double words is stored into the \emph{registerWe}. This instruction does not raise an invalid exception when handling a signaling NaN.


\paragraph{Modifier(s)}
\noindent\begin{itemize}
\item floatcomp (cf. floatcomp in section \ref{par:modifier-floatcomp} on page \pageref{par:modifier-floatcomp})
\item splat32 (cf. splat32 in section \ref{par:modifier-splat32} on page \pageref{par:modifier-splat32})
\end{itemize}

\paragraph{Execution} {Reservation table \textsf{ALU\_LITE.X} (Section~\ref{sec:reservation-ALU-LITE.X})}
~
{\begin{lstlisting}
stage ID:
  new splat32 = _ZX_1(splat32);
  new argument4 = _ZX_3(floatcomp);
stage RR:
  new argument3 = splat32 ? (_WX_32(upper27_lower5) << 32) | _ZX_32(_WX_32(upper27_lower5)) : _SX_32(_WX_32(upper27_lower5));
  new argument2 = GPR[registerZe<<1];
  new result1 = -floatcomp_64(argument4, argument2, argument3);
stage E1:
  GPR[registerWe<<1] = result1;
\end{lstlisting}}
\newpage

\paragraph{Encoding} Format \textsf{ALU\_FDCNWRRO.M} (Section~\ref{sec:format-ALU-FDCNWRRO.M}), \verb|alucode3 = 011|\\

\bitfields{ 1/P, 2/00, 2/-- --, 27/upper27, 1/1, 2/11, 1/1, 1/1, 3/floatcomp, 5/registerWo, 1/1, 2/11, 3/011, 1/0, 1/splat32, 5/lower5, 5/registerZo, 1/1 }


\paragraph{Syntax} \texttt{fcompnd}\emph{floatcomp} \emph{registerWo} = \emph{registerZo}, \emph{upper27\_\discretionary{}{}{}lower5}\emph{splat32}

\paragraph{Description} The \emph{registerZo} and the \emph{upper27\_\discretionary{}{}{}lower5} are interpreted as two binary 64 floating-point numbers. The \emph{registerZo} double word is compared to the \emph{upper27\_\discretionary{}{}{}lower5} double word using condition \emph{floatcomp}. The negated boolean result extended to double words is stored into the \emph{registerWo}. This instruction does not raise an invalid exception when handling a signaling NaN.


\paragraph{Modifier(s)}
\noindent\begin{itemize}
\item floatcomp (cf. floatcomp in section \ref{par:modifier-floatcomp} on page \pageref{par:modifier-floatcomp})
\item splat32 (cf. splat32 in section \ref{par:modifier-splat32} on page \pageref{par:modifier-splat32})
\end{itemize}

\paragraph{Execution} {Reservation table \textsf{ALU\_LITE.X} (Section~\ref{sec:reservation-ALU-LITE.X})}
~
{\begin{lstlisting}
stage ID:
  new splat32 = _ZX_1(splat32);
  new argument4 = _ZX_3(floatcomp);
stage RR:
  new argument3 = splat32 ? (_WX_32(upper27_lower5) << 32) | _ZX_32(_WX_32(upper27_lower5)) : _SX_32(_WX_32(upper27_lower5));
  new argument2 = GPR[(registerZo<<1)+1];
  new result1 = -floatcomp_64(argument4, argument2, argument3);
stage E1:
  GPR[(registerWo<<1)+1] = result1;
\end{lstlisting}}
\newpage
\subsubsection[{\small FCOMPNHQ: Floating-Point Compare Negate Half Word Quadruple}]{FCOMPNHQ\hfill Floating-Point Compare Negate Half Word Quadruple}
\label{instr:FCOMPNHQ}\index{fcompnhq}
 
\paragraph{Encoding} Format \textsf{ALU\_FHQCNWRRE} (Section~\ref{sec:format-ALU-FHQCNWRRE}), \verb|alucode3 = 010|\\

\bitfields{ 1/P, 2/11, 1/1, 1/1, 3/floatcomp, 5/registerWe, 1/0, 2/11, 3/010, 1/1, 5/registerYe, 1/0, 5/registerZe, 1/1 }


\paragraph{Syntax} \texttt{fcompnhq}\emph{floatcomp} \emph{registerWe} = \emph{registerZe}, \emph{registerYe}

\paragraph{Description} The \emph{registerZe} and the \emph{registerYe} are interpreted as four binary 16 floating-point numbers packed into 64 bits. The \emph{registerZe} half words are compared to the \emph{registerYe} half words using condition \emph{floatcomp}. The four negated boolean results extended to half words are packed and stored into the \emph{registerWe}. This instruction does not raise an invalid exception when handling a signaling NaN.


\paragraph{Modifier(s)}
\noindent\begin{itemize}
\item floatcomp (cf. floatcomp in section \ref{par:modifier-floatcomp} on page \pageref{par:modifier-floatcomp})
\end{itemize}

\paragraph{Execution} {Reservation table \textsf{ALU\_LITE} (Section~\ref{sec:reservation-ALU-LITE})}
~
{\begin{lstlisting}
stage ID:
  new argument4 = _ZX_3(floatcomp);
stage RR:
  new argument3 = GPR[registerYe<<1];
  new argument2 = GPR[registerZe<<1];
  for (I = 0; I < 4; I += 1) {
    result1.16[I] = -floatcomp_16(argument4, argument2.16[I], argument3.16[I])
  };
stage E1:
  GPR[registerWe<<1] = result1;
\end{lstlisting}}
\newpage

\paragraph{Encoding} Format \textsf{ALU\_FHQCNWRRO} (Section~\ref{sec:format-ALU-FHQCNWRRO}), \verb|alucode3 = 010|\\

\bitfields{ 1/P, 2/11, 1/1, 1/1, 3/floatcomp, 5/registerWo, 1/1, 2/11, 3/010, 1/1, 5/registerYo, 1/0, 5/registerZo, 1/1 }


\paragraph{Syntax} \texttt{fcompnhq}\emph{floatcomp} \emph{registerWo} = \emph{registerZo}, \emph{registerYo}

\paragraph{Description} The \emph{registerZo} and the \emph{registerYo} are interpreted as four binary 16 floating-point numbers packed into 64 bits. The \emph{registerZo} half words are compared to the \emph{registerYo} half words using condition \emph{floatcomp}. The four negated boolean results extended to half words are packed and stored into the \emph{registerWo}. This instruction does not raise an invalid exception when handling a signaling NaN.


\paragraph{Modifier(s)}
\noindent\begin{itemize}
\item floatcomp (cf. floatcomp in section \ref{par:modifier-floatcomp} on page \pageref{par:modifier-floatcomp})
\end{itemize}

\paragraph{Execution} {Reservation table \textsf{ALU\_LITE} (Section~\ref{sec:reservation-ALU-LITE})}
~
{\begin{lstlisting}
stage ID:
  new argument4 = _ZX_3(floatcomp);
stage RR:
  new argument3 = GPR[(registerYo<<1)+1];
  new argument2 = GPR[(registerZo<<1)+1];
  for (I = 0; I < 4; I += 1) {
    result1.16[I] = -floatcomp_16(argument4, argument2.16[I], argument3.16[I])
  };
stage E1:
  GPR[(registerWo<<1)+1] = result1;
\end{lstlisting}}
\newpage

\paragraph{Encoding} Format \textsf{ALU\_FHQCNWRRE.M} (Section~\ref{sec:format-ALU-FHQCNWRRE.M}), \verb|alucode3 = 010|\\

\bitfields{ 1/P, 2/00, 2/-- --, 27/upper27, 1/1, 2/11, 1/1, 1/1, 3/floatcomp, 5/registerWe, 1/0, 2/11, 3/010, 1/1, 1/splat32, 5/lower5, 5/registerZe, 1/1 }


\paragraph{Syntax} \texttt{fcompnhq}\emph{floatcomp} \emph{registerWe} = \emph{registerZe}, \emph{upper27\_\discretionary{}{}{}lower5}\emph{splat32}

\paragraph{Description} The \emph{registerZe} and the \emph{upper27\_\discretionary{}{}{}lower5} are interpreted as four binary 16 floating-point numbers packed into 64 bits. The \emph{registerZe} half words are compared to the \emph{upper27\_\discretionary{}{}{}lower5} half words using condition \emph{floatcomp}. The four negated boolean results extended to half words are packed and stored into the \emph{registerWe}. This instruction does not raise an invalid exception when handling a signaling NaN.


\paragraph{Modifier(s)}
\noindent\begin{itemize}
\item floatcomp (cf. floatcomp in section \ref{par:modifier-floatcomp} on page \pageref{par:modifier-floatcomp})
\item splat32 (cf. splat32 in section \ref{par:modifier-splat32} on page \pageref{par:modifier-splat32})
\end{itemize}

\paragraph{Execution} {Reservation table \textsf{ALU\_LITE.X} (Section~\ref{sec:reservation-ALU-LITE.X})}
~
{\begin{lstlisting}
stage ID:
  new splat32 = _ZX_1(splat32);
  new argument4 = _ZX_3(floatcomp);
stage RR:
  new argument3 = splat32 ? (_WX_32(upper27_lower5) << 32) | _ZX_32(_WX_32(upper27_lower5)) : _SX_32(_WX_32(upper27_lower5));
  new argument2 = GPR[registerZe<<1];
  for (I = 0; I < 4; I += 1) {
    result1.16[I] = -floatcomp_16(argument4, argument2.16[I], argument3.16[I])
  };
stage E1:
  GPR[registerWe<<1] = result1;
\end{lstlisting}}
\newpage

\paragraph{Encoding} Format \textsf{ALU\_FHQCNWRRO.M} (Section~\ref{sec:format-ALU-FHQCNWRRO.M}), \verb|alucode3 = 010|\\

\bitfields{ 1/P, 2/00, 2/-- --, 27/upper27, 1/1, 2/11, 1/1, 1/1, 3/floatcomp, 5/registerWo, 1/1, 2/11, 3/010, 1/1, 1/splat32, 5/lower5, 5/registerZo, 1/1 }


\paragraph{Syntax} \texttt{fcompnhq}\emph{floatcomp} \emph{registerWo} = \emph{registerZo}, \emph{upper27\_\discretionary{}{}{}lower5}\emph{splat32}

\paragraph{Description} The \emph{registerZo} and the \emph{upper27\_\discretionary{}{}{}lower5} are interpreted as four binary 16 floating-point numbers packed into 64 bits. The \emph{registerZo} half words are compared to the \emph{upper27\_\discretionary{}{}{}lower5} half words using condition \emph{floatcomp}. The four negated boolean results extended to half words are packed and stored into the \emph{registerWo}. This instruction does not raise an invalid exception when handling a signaling NaN.


\paragraph{Modifier(s)}
\noindent\begin{itemize}
\item floatcomp (cf. floatcomp in section \ref{par:modifier-floatcomp} on page \pageref{par:modifier-floatcomp})
\item splat32 (cf. splat32 in section \ref{par:modifier-splat32} on page \pageref{par:modifier-splat32})
\end{itemize}

\paragraph{Execution} {Reservation table \textsf{ALU\_LITE.X} (Section~\ref{sec:reservation-ALU-LITE.X})}
~
{\begin{lstlisting}
stage ID:
  new splat32 = _ZX_1(splat32);
  new argument4 = _ZX_3(floatcomp);
stage RR:
  new argument3 = splat32 ? (_WX_32(upper27_lower5) << 32) | _ZX_32(_WX_32(upper27_lower5)) : _SX_32(_WX_32(upper27_lower5));
  new argument2 = GPR[(registerZo<<1)+1];
  for (I = 0; I < 4; I += 1) {
    result1.16[I] = -floatcomp_16(argument4, argument2.16[I], argument3.16[I])
  };
stage E1:
  GPR[(registerWo<<1)+1] = result1;
\end{lstlisting}}
\newpage
\subsubsection[{\small FCOMPNWP: Floating-Point Compare Negate Word Pair}]{FCOMPNWP\hfill Floating-Point Compare Negate Word Pair}
\label{instr:FCOMPNWP}\index{fcompnwp}
 
\paragraph{Encoding} Format \textsf{ALU\_FWPCNWRRE} (Section~\ref{sec:format-ALU-FWPCNWRRE}), \verb|alucode3 = 010|\\

\bitfields{ 1/P, 2/11, 1/1, 1/1, 3/floatcomp, 5/registerWe, 1/0, 2/11, 3/010, 1/1, 5/registerYe, 1/0, 5/registerZe, 1/0 }


\paragraph{Syntax} \texttt{fcompnwp}\emph{floatcomp} \emph{registerWe} = \emph{registerZe}, \emph{registerYe}

\paragraph{Description} The \emph{registerZe} and the \emph{registerYe} are interpreted as two binary 32 floating-point numbers packed into 64 bits. The \emph{registerZe} words are compared to the \emph{registerYe} words using condition \emph{floatcomp}. The two negated boolean results extended to words are packed and stored into the \emph{registerWe}. This instruction does not raise an invalid exception when handling a signaling NaN.


\paragraph{Modifier(s)}
\noindent\begin{itemize}
\item floatcomp (cf. floatcomp in section \ref{par:modifier-floatcomp} on page \pageref{par:modifier-floatcomp})
\end{itemize}

\paragraph{Execution} {Reservation table \textsf{ALU\_LITE} (Section~\ref{sec:reservation-ALU-LITE})}
~
{\begin{lstlisting}
stage ID:
  new argument4 = _ZX_3(floatcomp);
stage RR:
  new argument3 = GPR[registerYe<<1];
  new argument2 = GPR[registerZe<<1];
  for (I = 0; I < 2; I += 1) {
    result1.32[I] = -floatcomp_32(argument4, argument2.32[I], argument3.32[I])
  };
stage E1:
  GPR[registerWe<<1] = result1;
\end{lstlisting}}
\newpage

\paragraph{Encoding} Format \textsf{ALU\_FWPCNWRRO} (Section~\ref{sec:format-ALU-FWPCNWRRO}), \verb|alucode3 = 010|\\

\bitfields{ 1/P, 2/11, 1/1, 1/1, 3/floatcomp, 5/registerWo, 1/1, 2/11, 3/010, 1/1, 5/registerYo, 1/0, 5/registerZo, 1/0 }


\paragraph{Syntax} \texttt{fcompnwp}\emph{floatcomp} \emph{registerWo} = \emph{registerZo}, \emph{registerYo}

\paragraph{Description} The \emph{registerZo} and the \emph{registerYo} are interpreted as two binary 32 floating-point numbers packed into 64 bits. The \emph{registerZo} words are compared to the \emph{registerYo} words using condition \emph{floatcomp}. The two negated boolean results extended to words are packed and stored into the \emph{registerWo}. This instruction does not raise an invalid exception when handling a signaling NaN.


\paragraph{Modifier(s)}
\noindent\begin{itemize}
\item floatcomp (cf. floatcomp in section \ref{par:modifier-floatcomp} on page \pageref{par:modifier-floatcomp})
\end{itemize}

\paragraph{Execution} {Reservation table \textsf{ALU\_LITE} (Section~\ref{sec:reservation-ALU-LITE})}
~
{\begin{lstlisting}
stage ID:
  new argument4 = _ZX_3(floatcomp);
stage RR:
  new argument3 = GPR[(registerYo<<1)+1];
  new argument2 = GPR[(registerZo<<1)+1];
  for (I = 0; I < 2; I += 1) {
    result1.32[I] = -floatcomp_32(argument4, argument2.32[I], argument3.32[I])
  };
stage E1:
  GPR[(registerWo<<1)+1] = result1;
\end{lstlisting}}
\newpage

\paragraph{Encoding} Format \textsf{ALU\_FWPCNWRRE.M} (Section~\ref{sec:format-ALU-FWPCNWRRE.M}), \verb|alucode3 = 001|\\

\bitfields{ 1/P, 2/00, 2/-- --, 27/upper27, 1/1, 2/11, 1/1, 1/1, 3/floatcomp, 5/registerWe, 1/0, 2/11, 3/001, 1/1, 1/splat32, 5/lower5, 5/registerZe, 1/0 }


\paragraph{Syntax} \texttt{fcompnwp}\emph{floatcomp} \emph{registerWe} = \emph{registerZe}, \emph{upper27\_\discretionary{}{}{}lower5}\emph{splat32}

\paragraph{Description} The \emph{registerZe} and the \emph{upper27\_\discretionary{}{}{}lower5} are interpreted as two binary 32 floating-point numbers packed into 64 bits. The \emph{registerZe} words are compared to the \emph{upper27\_\discretionary{}{}{}lower5} words using condition \emph{floatcomp}. The two negated boolean results extended to words are packed and stored into the \emph{registerWe}. This instruction does not raise an invalid exception when handling a signaling NaN.


\paragraph{Modifier(s)}
\noindent\begin{itemize}
\item floatcomp (cf. floatcomp in section \ref{par:modifier-floatcomp} on page \pageref{par:modifier-floatcomp})
\item splat32 (cf. splat32 in section \ref{par:modifier-splat32} on page \pageref{par:modifier-splat32})
\end{itemize}

\paragraph{Execution} {Reservation table \textsf{ALU\_LITE.X} (Section~\ref{sec:reservation-ALU-LITE.X})}
~
{\begin{lstlisting}
stage ID:
  new splat32 = _ZX_1(splat32);
  new argument4 = _ZX_3(floatcomp);
stage RR:
  new argument3 = splat32 ? (_WX_32(upper27_lower5) << 32) | _ZX_32(_WX_32(upper27_lower5)) : _SX_32(_WX_32(upper27_lower5));
  new argument2 = GPR[registerZe<<1];
  for (I = 0; I < 2; I += 1) {
    result1.32[I] = -floatcomp_32(argument4, argument2.32[I], argument3.32[I])
  };
stage E1:
  GPR[registerWe<<1] = result1;
\end{lstlisting}}
\newpage

\paragraph{Encoding} Format \textsf{ALU\_FWPCNWRRO.M} (Section~\ref{sec:format-ALU-FWPCNWRRO.M}), \verb|alucode3 = 001|\\

\bitfields{ 1/P, 2/00, 2/-- --, 27/upper27, 1/1, 2/11, 1/1, 1/1, 3/floatcomp, 5/registerWo, 1/1, 2/11, 3/001, 1/1, 1/splat32, 5/lower5, 5/registerZo, 1/0 }


\paragraph{Syntax} \texttt{fcompnwp}\emph{floatcomp} \emph{registerWo} = \emph{registerZo}, \emph{upper27\_\discretionary{}{}{}lower5}\emph{splat32}

\paragraph{Description} The \emph{registerZo} and the \emph{upper27\_\discretionary{}{}{}lower5} are interpreted as two binary 32 floating-point numbers packed into 64 bits. The \emph{registerZo} words are compared to the \emph{upper27\_\discretionary{}{}{}lower5} words using condition \emph{floatcomp}. The two negated boolean results extended to words are packed and stored into the \emph{registerWo}. This instruction does not raise an invalid exception when handling a signaling NaN.


\paragraph{Modifier(s)}
\noindent\begin{itemize}
\item floatcomp (cf. floatcomp in section \ref{par:modifier-floatcomp} on page \pageref{par:modifier-floatcomp})
\item splat32 (cf. splat32 in section \ref{par:modifier-splat32} on page \pageref{par:modifier-splat32})
\end{itemize}

\paragraph{Execution} {Reservation table \textsf{ALU\_LITE.X} (Section~\ref{sec:reservation-ALU-LITE.X})}
~
{\begin{lstlisting}
stage ID:
  new splat32 = _ZX_1(splat32);
  new argument4 = _ZX_3(floatcomp);
stage RR:
  new argument3 = splat32 ? (_WX_32(upper27_lower5) << 32) | _ZX_32(_WX_32(upper27_lower5)) : _SX_32(_WX_32(upper27_lower5));
  new argument2 = GPR[(registerZo<<1)+1];
  for (I = 0; I < 2; I += 1) {
    result1.32[I] = -floatcomp_32(argument4, argument2.32[I], argument3.32[I])
  };
stage E1:
  GPR[(registerWo<<1)+1] = result1;
\end{lstlisting}}
\newpage
\subsubsection[{\small FCOMPW: Floating-Point Compare Word}]{FCOMPW\hfill Floating-Point Compare Word}
\label{instr:FCOMPW}\index{fcompw}
 
\paragraph{Encoding} Format \textsf{ALU\_FWCWRR} (Section~\ref{sec:format-ALU-FWCWRR}), \verb|alucode3 = 010|\\

\bitfields{ 1/P, 2/11, 1/1, 1/0, 3/floatcomp, 6/registerW, 2/01, 3/010, 1/0, 6/registerY, 6/registerZ }


\paragraph{Syntax} \texttt{fcompw}\emph{floatcomp} \emph{registerW} = \emph{registerZ}, \emph{registerY}

\paragraph{Description} The \emph{registerZ} is compared to the \emph{registerY} using condition \emph{floatcomp}, assuming binary 32 floating-point numbers. The boolean result is stored into the \emph{registerW}. This instruction does not raise an invalid exception when handling a signaling NaN.


\paragraph{Modifier(s)}
\noindent\begin{itemize}
\item floatcomp (cf. floatcomp in section \ref{par:modifier-floatcomp} on page \pageref{par:modifier-floatcomp})
\end{itemize}

\paragraph{Execution} {Reservation table \textsf{ALU\_LITE} (Section~\ref{sec:reservation-ALU-LITE})}
~
{\begin{lstlisting}
stage ID:
  new argument4 = _ZX_3(floatcomp);
stage RR:
  new argument3 = GPR[registerY];
  new argument2 = GPR[registerZ];
  new result1 = floatcomp_32(argument4, argument2, argument3);
stage E1:
  GPR[registerW] = _ZX_32(result1);
\end{lstlisting}}
\newpage

\paragraph{Encoding} Format \textsf{ALU\_FWCWRR.W} (Section~\ref{sec:format-ALU-FWCWRR.W}), \verb|alucode3 = 010|\\

\bitfields{ 1/P, 2/00, 2/-- --, 27/upper27, 1/1, 2/11, 1/1, 1/0, 3/floatcomp, 6/registerW, 2/01, 3/010, 1/0, 1/0, 5/lower5, 6/registerZ }


\paragraph{Syntax} \texttt{fcompw}\emph{floatcomp} \emph{registerW} = \emph{registerZ}, \emph{upper27\_\discretionary{}{}{}lower5}

\paragraph{Description} The \emph{registerZ} is compared to the \emph{upper27\_\discretionary{}{}{}lower5} using condition \emph{floatcomp}, assuming binary 32 floating-point numbers. The boolean result is stored into the \emph{registerW}. This instruction does not raise an invalid exception when handling a signaling NaN.


\paragraph{Modifier(s)}
\noindent\begin{itemize}
\item floatcomp (cf. floatcomp in section \ref{par:modifier-floatcomp} on page \pageref{par:modifier-floatcomp})
\end{itemize}

\paragraph{Execution} {Reservation table \textsf{ALU\_LITE.X} (Section~\ref{sec:reservation-ALU-LITE.X})}
~
{\begin{lstlisting}
stage ID:
  new argument4 = _ZX_3(floatcomp);
stage RR:
  new argument3 = _WX_32(upper27_lower5);
  new argument2 = GPR[registerZ];
  new result1 = floatcomp_32(argument4, argument2, argument3);
stage E1:
  GPR[registerW] = _ZX_32(result1);
\end{lstlisting}}
\newpage
\subsubsection[{\small FDIVD: Floating-Point Divide Double Word}]{FDIVD\hfill Floating-Point Divide Double Word}
\label{instr:FDIVD}\index{fdivd}
 
\paragraph{Encoding} Format \textsf{ALU\_FDMUL} (Section~\ref{sec:format-ALU-FDMUL}), \verb|alucode3 = 111|\\

\bitfields{ 1/P, 2/11, 1/1, 1/fnegate, 3/floatmode, 6/registerW, 2/00, 3/111, 1/0, 6/registerY, 6/registerZ }


\paragraph{Syntax} \texttt{fdivd}\emph{fnegate}\emph{floatmode} \emph{registerW} = \emph{registerZ}, \emph{registerY}

\paragraph{Description} The \emph{registerZ} is divided by the \emph{registerY}, assuming binary 64 floating-point numbers. If \emph{fnegate} is set, the \emph{registerZ} double word is negated prior to the operation. The result is rounded to the binary 64 floating-point format according to the rounding mode and stored into the \emph{registerW}. This instruction may raise exception bits in the CS register.


\paragraph{Modifier(s)}
\noindent\begin{itemize}
\item floatmode (cf. floatmode in section \ref{par:modifier-floatmode} on page \pageref{par:modifier-floatmode})
\item fnegate (cf. fnegate in section \ref{par:modifier-fnegate} on page \pageref{par:modifier-fnegate})
\end{itemize}

\paragraph{Execution} {Reservation table \textsf{ALU\_FULL} (Section~\ref{sec:reservation-ALU-FULL})}
~
{\begin{lstlisting}
stage ID:
  new fnegate = _ZX_1(fnegate);
  new floatmode = _ZX_3(floatmode);
stage RR:
  new argument3 = GPR[registerY];
  new argument2 = GPR[registerZ];
  new RM = decode_riscv_float_rounding_mode(floatmode);
  new fsign = fnegate ? 0x8000000000000000 : 0;
  new result1 = f64_div(RM, argument2 ^ fsign, argument3);
stage E4:
  GPR[registerW] = result1;
  commit_float_exception_flags();
\end{lstlisting}}
\newpage
\subsubsection[{\small FDIVH: Floating-Point Divide Half Word}]{FDIVH\hfill Floating-Point Divide Half Word}
\label{instr:FDIVH}\index{fdivh}
 
\paragraph{Encoding} Format \textsf{ALU\_FHMUL} (Section~\ref{sec:format-ALU-FHMUL}), \verb|alucode3 = 111|\\

\bitfields{ 1/P, 2/11, 1/1, 1/fnegate, 3/floatmode, 6/registerW, 2/01, 3/111, 1/1, 6/registerY, 6/registerZ }


\paragraph{Syntax} \texttt{fdivh}\emph{fnegate}\emph{floatmode} \emph{registerW} = \emph{registerZ}, \emph{registerY}

\paragraph{Description} The \emph{registerZ} is divided by the \emph{registerY}, assuming binary 16 floating-point numbers. If \emph{fnegate} is set, the \emph{registerZ} half word is negated prior to the operation. The result is rounded to the binary 16 floating-point format according to the rounding mode and stored into the \emph{registerW}. This instruction may raise exception bits in the CS register.


\paragraph{Modifier(s)}
\noindent\begin{itemize}
\item floatmode (cf. floatmode in section \ref{par:modifier-floatmode} on page \pageref{par:modifier-floatmode})
\item fnegate (cf. fnegate in section \ref{par:modifier-fnegate} on page \pageref{par:modifier-fnegate})
\end{itemize}

\paragraph{Execution} {Reservation table \textsf{ALU\_FULL} (Section~\ref{sec:reservation-ALU-FULL})}
~
{\begin{lstlisting}
stage ID:
  new fnegate = _ZX_1(fnegate);
  new floatmode = _ZX_3(floatmode);
stage RR:
  new argument3 = GPR[registerY];
  new argument2 = GPR[registerZ];
  new RM = decode_riscv_float_rounding_mode(floatmode);
  new fsign = fnegate ? 0x8000 : 0;
  new result1 = f16_div(RM, argument2 ^ fsign, argument3);
stage E3:
  GPR[registerW] = result1;
  commit_float_exception_flags();
\end{lstlisting}}
\newpage
\subsubsection[{\small FDIVW: Floating-Point Divide Word}]{FDIVW\hfill Floating-Point Divide Word}
\label{instr:FDIVW}\index{fdivw}
 
\paragraph{Encoding} Format \textsf{ALU\_FWMUL} (Section~\ref{sec:format-ALU-FWMUL}), \verb|alucode3 = 111|\\

\bitfields{ 1/P, 2/11, 1/1, 1/fnegate, 3/floatmode, 6/registerW, 2/01, 3/111, 1/0, 6/registerY, 6/registerZ }


\paragraph{Syntax} \texttt{fdivw}\emph{fnegate}\emph{floatmode} \emph{registerW} = \emph{registerZ}, \emph{registerY}

\paragraph{Description} The \emph{registerZ} is divided by the \emph{registerY}, assuming binary 32 floating-point numbers. If \emph{fnegate} is set, the \emph{registerZ} word is negated prior to the operation. The result is rounded to the binary 32 floating-point format according to the rounding mode and stored into the \emph{registerW}. This instruction may raise exception bits in the CS register.


\paragraph{Modifier(s)}
\noindent\begin{itemize}
\item floatmode (cf. floatmode in section \ref{par:modifier-floatmode} on page \pageref{par:modifier-floatmode})
\item fnegate (cf. fnegate in section \ref{par:modifier-fnegate} on page \pageref{par:modifier-fnegate})
\end{itemize}

\paragraph{Execution} {Reservation table \textsf{ALU\_FULL} (Section~\ref{sec:reservation-ALU-FULL})}
~
{\begin{lstlisting}
stage ID:
  new fnegate = _ZX_1(fnegate);
  new floatmode = _ZX_3(floatmode);
stage RR:
  new argument3 = GPR[registerY];
  new argument2 = GPR[registerZ];
  new RM = decode_riscv_float_rounding_mode(floatmode);
  new fsign = fnegate ? 0x80000000 : 0;
  new result1 = f32_div(RM, argument2 ^ fsign, argument3);
stage E4:
  GPR[registerW] = result1;
  commit_float_exception_flags();
\end{lstlisting}}
\newpage
\subsubsection[{\small FEXTLHWQ: Floating Point Extend Lane Half Word to Word Quadruple}]{FEXTLHWQ\hfill Floating Point Extend Lane Half Word to Word Quadruple}
\label{instr:FEXTLHWQ}\index{fextlhwq}
 
\paragraph{Encoding} Format \textsf{ALU\_FWQEXTLWR} (Section~\ref{sec:format-ALU-FWQEXTLWR}), \verb|fpucode6 = 1001|\\

\bitfields{ 1/P, 2/11, 1/1, 1/1, 3/111, 5/registerM, 1/oddlanes, 2/11, 3/000, 1/1, 4/1001, 1/--, 1/0, 5/registerP, 1/0 }


\paragraph{Syntax} \texttt{fextlhwq}\emph{oddlanes} \emph{registerM} = \emph{registerP}

\paragraph{Description} The \emph{registerP} interpreted as eight binary 16 floating-point numbers packed into 128 bits. The even or odd half words of the \emph{registerP} are selected, depending on \emph{oddlanes}. The four selected half words are converted to binary 32 floating-point numbers, packed and stored into the \emph{registerM}. This instruction may raise invalid exception bit in the CS register.


\paragraph{Modifier(s)}
\noindent\begin{itemize}
\item oddlanes (cf. oddlanes in section \ref{par:modifier-oddlanes} on page \pageref{par:modifier-oddlanes})
\end{itemize}

\paragraph{Execution} {Reservation table \textsf{ALU\_LITE} (Section~\ref{sec:reservation-ALU-LITE})}
~
{\begin{lstlisting}
stage ID:
  new oddlanes = _ZX_1(oddlanes);
stage RR:
  new argument2 = PGR[registerP];
  if (oddlanes) {
    argument2 = argument2 >> 16;
  }
  for (I = 0; I < 4; I += 1) {
    result1.32[I] = f16_to_f32(_ZX_16(argument2.32[I]))
  };
stage E1:
  PGR[registerM] = result1;
\end{lstlisting}}
\newpage
\subsubsection[{\small FFMAD: Floating-Point Fused Multiply Add Double Word}]{FFMAD\hfill Floating-Point Fused Multiply Add Double Word}
\label{instr:FFMAD}\index{ffmad}
 
\paragraph{Encoding} Format \textsf{ALU\_FDFMA} (Section~\ref{sec:format-ALU-FDFMA}), \verb|alucode3 = 100|\\

\bitfields{ 1/P, 2/11, 1/1, 1/fnegate, 3/floatmode, 6/registerW, 2/00, 3/100, 1/0, 6/registerY, 6/registerZ }


\paragraph{Syntax} \texttt{ffmad}\emph{fnegate}\emph{floatmode} \emph{registerW} = \emph{registerZ}, \emph{registerY}

\paragraph{Description} The \emph{registerY} is multiplied by the \emph{registerZ}, assuming binary 64 floating-point numbers, and the product is added to the \emph{registerW}. If \emph{fnegate} is set, the \emph{registerW} and the \emph{registerZ} double words are negated prior to the operation. The result is rounded to the binary 64 floating-point format according to the rounding mode and stored into the \emph{registerW}. This instruction may raise exception bits in the CS register.


\paragraph{Modifier(s)}
\noindent\begin{itemize}
\item floatmode (cf. floatmode in section \ref{par:modifier-floatmode} on page \pageref{par:modifier-floatmode})
\item fnegate (cf. fnegate in section \ref{par:modifier-fnegate} on page \pageref{par:modifier-fnegate})
\end{itemize}

\paragraph{Execution} {Reservation table \textsf{ALU\_LITE} (Section~\ref{sec:reservation-ALU-LITE})}
~
{\begin{lstlisting}
stage ID:
  new fnegate = _ZX_1(fnegate);
  new floatmode = _ZX_3(floatmode);
stage RR:
  new argument3 = GPR[registerY];
  new argument2 = GPR[registerZ];
  new argument1 = GPR[registerW];
  new RM = decode_riscv_float_rounding_mode(floatmode);
  new fsign = fnegate ? 0x8000000000000000 : 0;
  new result1 = f64_mulAdd(RM, argument3, argument2 ^ fsign, argument1 ^ fsign);
stage E4:
  GPR[registerW] = result1;
  commit_float_exception_flags();
\end{lstlisting}}
\newpage
\subsubsection[{\small FFMAH: Floating-Point Fused Multiply Add Half Word}]{FFMAH\hfill Floating-Point Fused Multiply Add Half Word}
\label{instr:FFMAH}\index{ffmah}
 
\paragraph{Encoding} Format \textsf{ALU\_FHFMA} (Section~\ref{sec:format-ALU-FHFMA}), \verb|alucode3 = 100|\\

\bitfields{ 1/P, 2/11, 1/1, 1/fnegate, 3/floatmode, 6/registerW, 2/01, 3/100, 1/1, 6/registerY, 6/registerZ }


\paragraph{Syntax} \texttt{ffmah}\emph{fnegate}\emph{floatmode} \emph{registerW} = \emph{registerZ}, \emph{registerY}

\paragraph{Description} The \emph{registerY} is multiplied by the \emph{registerZ}, assuming binary 16 floating-point numbers, and the product is added to the \emph{registerW}. If \emph{fnegate} is set, the \emph{registerW} and the \emph{registerZ} half words are negated prior to the operation. The result is rounded to the binary 16 floating-point format according to the rounding mode and stored into the \emph{registerW}. This instruction may raise exception bits in the CS register.


\paragraph{Modifier(s)}
\noindent\begin{itemize}
\item floatmode (cf. floatmode in section \ref{par:modifier-floatmode} on page \pageref{par:modifier-floatmode})
\item fnegate (cf. fnegate in section \ref{par:modifier-fnegate} on page \pageref{par:modifier-fnegate})
\end{itemize}

\paragraph{Execution} {Reservation table \textsf{ALU\_LITE} (Section~\ref{sec:reservation-ALU-LITE})}
~
{\begin{lstlisting}
stage ID:
  new fnegate = _ZX_1(fnegate);
  new floatmode = _ZX_3(floatmode);
stage RR:
  new argument3 = GPR[registerY];
  new argument2 = GPR[registerZ];
  new argument1 = GPR[registerW];
  new RM = decode_riscv_float_rounding_mode(floatmode);
  new fsign = fnegate ? 0x8000 : 0;
  new result1 = f16_mulAdd(RM, argument3, argument2 ^ fsign, argument1 ^ fsign);
stage E3:
  GPR[registerW] = result1;
  commit_float_exception_flags();
\end{lstlisting}}
\newpage
\subsubsection[{\small FFMAHQ: Floating-Point Fused Multiply Add Half Word Quadruple}]{FFMAHQ\hfill Floating-Point Fused Multiply Add Half Word Quadruple}
\label{instr:FFMAHQ}\index{ffmahq}
 
\paragraph{Encoding} Format \textsf{ALU\_FHQFMAE} (Section~\ref{sec:format-ALU-FHQFMAE}), \verb|alucode3 = 100|\\

\bitfields{ 1/P, 2/11, 1/1, 1/fnegate, 3/floatmode, 5/registerWe, 1/0, 2/11, 3/100, 1/1, 5/registerYe, 1/0, 5/registerZe, 1/1 }


\paragraph{Syntax} \texttt{ffmahq}\emph{fnegate}\emph{floatmode} \emph{registerWe} = \emph{registerZe}, \emph{registerYe}

\paragraph{Description} The \emph{registerWe}, the \emph{registerZe} and the \emph{registerYe} are interpreted as four binary 16 floating-point numbers packed into 64 bits. If \emph{fnegate} is set, the \emph{registerWe} and the \emph{registerZe} half words are negated prior to the operation. The \emph{registerYe} half words are multiplied by the \emph{registerZe} half words, assuming binary 16 floating-point numbers, and the product half words are added to the \emph{registerWe} half words. The four results are rounded to the binary 16 floating-point format according to the rounding mode and stored into the \emph{registerWe}. This instruction may raise exception bits in the CS register.


\paragraph{Modifier(s)}
\noindent\begin{itemize}
\item floatmode (cf. floatmode in section \ref{par:modifier-floatmode} on page \pageref{par:modifier-floatmode})
\item fnegate (cf. fnegate in section \ref{par:modifier-fnegate} on page \pageref{par:modifier-fnegate})
\end{itemize}

\paragraph{Execution} {Reservation table \textsf{ALU\_LITE} (Section~\ref{sec:reservation-ALU-LITE})}
~
{\begin{lstlisting}
stage ID:
  new fnegate = _ZX_1(fnegate);
  new floatmode = _ZX_3(floatmode);
stage RR:
  new argument3 = GPR[registerYe<<1];
  new argument2 = GPR[registerZe<<1];
  new argument1 = GPR[registerWe<<1];
  new RM = decode_riscv_float_rounding_mode(floatmode);
  new fsign = fnegate ? 0x8000 : 0;
  for (I = 0; I < 4; I += 1) {
    result1.16[I] = f16_mulAdd(RM, argument3.16[I], argument2.16[I] ^ fsign, argument1.16[I] ^ fsign)
  };
stage E4:
  GPR[registerWe<<1] = result1;
  commit_float_exception_flags();
\end{lstlisting}}
\newpage

\paragraph{Encoding} Format \textsf{ALU\_FHQFMAO} (Section~\ref{sec:format-ALU-FHQFMAO}), \verb|alucode3 = 100|\\

\bitfields{ 1/P, 2/11, 1/1, 1/fnegate, 3/floatmode, 5/registerWo, 1/1, 2/11, 3/100, 1/1, 5/registerYo, 1/0, 5/registerZo, 1/1 }


\paragraph{Syntax} \texttt{ffmahq}\emph{fnegate}\emph{floatmode} \emph{registerWo} = \emph{registerZo}, \emph{registerYo}

\paragraph{Description} The \emph{registerWo}, the \emph{registerZo} and the \emph{registerYo} are interpreted as four binary 16 floating-point numbers packed into 64 bits. If \emph{fnegate} is set, the \emph{registerWo} and the \emph{registerZo} half words are negated prior to the operation. The \emph{registerYo} half words are multiplied by the \emph{registerZo} half words, assuming binary 16 floating-point numbers, and the product half words are added to the \emph{registerWo} half words. The four results are rounded to the binary 16 floating-point format according to the rounding mode and stored into the \emph{registerWo}. This instruction may raise exception bits in the CS register.


\paragraph{Modifier(s)}
\noindent\begin{itemize}
\item floatmode (cf. floatmode in section \ref{par:modifier-floatmode} on page \pageref{par:modifier-floatmode})
\item fnegate (cf. fnegate in section \ref{par:modifier-fnegate} on page \pageref{par:modifier-fnegate})
\end{itemize}

\paragraph{Execution} {Reservation table \textsf{ALU\_LITE} (Section~\ref{sec:reservation-ALU-LITE})}
~
{\begin{lstlisting}
stage ID:
  new fnegate = _ZX_1(fnegate);
  new floatmode = _ZX_3(floatmode);
stage RR:
  new argument3 = GPR[(registerYo<<1)+1];
  new argument2 = GPR[(registerZo<<1)+1];
  new argument1 = GPR[(registerWo<<1)+1];
  new RM = decode_riscv_float_rounding_mode(floatmode);
  new fsign = fnegate ? 0x8000 : 0;
  for (I = 0; I < 4; I += 1) {
    result1.16[I] = f16_mulAdd(RM, argument3.16[I], argument2.16[I] ^ fsign, argument1.16[I] ^ fsign)
  };
stage E4:
  GPR[(registerWo<<1)+1] = result1;
  commit_float_exception_flags();
\end{lstlisting}}
\newpage
\subsubsection[{\small FFMAW: Floating-Point Fused Multiply Add Word}]{FFMAW\hfill Floating-Point Fused Multiply Add Word}
\label{instr:FFMAW}\index{ffmaw}
 
\paragraph{Encoding} Format \textsf{ALU\_FWFMA} (Section~\ref{sec:format-ALU-FWFMA}), \verb|alucode3 = 100|\\

\bitfields{ 1/P, 2/11, 1/1, 1/fnegate, 3/floatmode, 6/registerW, 2/01, 3/100, 1/0, 6/registerY, 6/registerZ }


\paragraph{Syntax} \texttt{ffmaw}\emph{fnegate}\emph{floatmode} \emph{registerW} = \emph{registerZ}, \emph{registerY}

\paragraph{Description} The \emph{registerY} is multiplied by the \emph{registerZ}, assuming binary 32 floating-point numbers, and the product is added to the \emph{registerW}. If \emph{fnegate} is set, the \emph{registerW} and the \emph{registerZ} words are negated prior to the operation. The result is rounded to the binary 32 floating-point format according to the rounding mode and stored into the \emph{registerW}. This instruction may raise exception bits in the CS register.


\paragraph{Modifier(s)}
\noindent\begin{itemize}
\item floatmode (cf. floatmode in section \ref{par:modifier-floatmode} on page \pageref{par:modifier-floatmode})
\item fnegate (cf. fnegate in section \ref{par:modifier-fnegate} on page \pageref{par:modifier-fnegate})
\end{itemize}

\paragraph{Execution} {Reservation table \textsf{ALU\_LITE} (Section~\ref{sec:reservation-ALU-LITE})}
~
{\begin{lstlisting}
stage ID:
  new fnegate = _ZX_1(fnegate);
  new floatmode = _ZX_3(floatmode);
stage RR:
  new argument3 = GPR[registerY];
  new argument2 = GPR[registerZ];
  new argument1 = GPR[registerW];
  new RM = decode_riscv_float_rounding_mode(floatmode);
  new fsign = fnegate ? 0x80000000 : 0;
  new result1 = f32_mulAdd(RM, argument3, argument2 ^ fsign, argument1 ^ fsign);
stage E4:
  GPR[registerW] = result1;
  commit_float_exception_flags();
\end{lstlisting}}
\newpage
\subsubsection[{\small FFMAWC: Floating-Point Multiply Add Word Complex}]{FFMAWC\hfill Floating-Point Multiply Add Word Complex}
\label{instr:FFMAWC}\index{ffmawc}
 
\paragraph{Encoding} Format \textsf{ALU\_FWCOP} (Section~\ref{sec:format-ALU-FWCOP}), \verb|alucode6 = 11|\\

\bitfields{ 1/P, 2/11, 1/1, 1/imultiply, 3/floatmode, 6/registerW, 2/00, 2/11, 1/conjugate, 1/1, 6/registerY, 6/registerZ }


\paragraph{Syntax} \texttt{ffmawc}\emph{conjugate}\emph{imultiply}\emph{floatmode} \emph{registerW} = \emph{registerZ}, \emph{registerY}

\paragraph{Description} The \emph{registerW}, the \emph{registerY} and the \emph{registerZ} are interpreted as binary 32 floating-point complex numbers. If \emph{conjugate} is set, the complex conjugate of the \emph{registerZ} is used in the multiplication. The complex product is added to the \emph{registerW} and rounded according to the rounding mode. If \emph{imultiply} is set, the result is multiplied by the imaginary unit. The resulting binary 32 floating-point complex number is stored back into the \emph{registerW}. This instruction may raise exception bits in the CS register.


\paragraph{Modifier(s)}
\noindent\begin{itemize}
\item floatmode (cf. floatmode in section \ref{par:modifier-floatmode} on page \pageref{par:modifier-floatmode})
\item imultiply (cf. imultiply in section \ref{par:modifier-imultiply} on page \pageref{par:modifier-imultiply})
\item conjugate (cf. conjugate in section \ref{par:modifier-conjugate} on page \pageref{par:modifier-conjugate})
\end{itemize}

\paragraph{Execution} {Reservation table \textsf{ALU\_LITE} (Section~\ref{sec:reservation-ALU-LITE})}
~
{\begin{lstlisting}
stage ID:
  new conjugate = _ZX_1(conjugate);
  new imultiply = _ZX_1(imultiply);
  new floatmode = _ZX_3(floatmode);
stage RR:
  new argument3 = GPR[registerY];
  new argument2 = GPR[registerZ];
  new argument1 = GPR[registerW];
  new RM = decode_riscv_float_rounding_mode(floatmode);
  if (conjugate) {
    argument2.64[0] = fconj_32_32(argument2.64[0]);
  }
  result1.64[0] = ffmac_32_32(RM, argument2.64[0], argument3.64[0], argument1.64[0]);
  if (imultiply) {
    new result1 = _SWAP_32(result1);
    result1.32[0] = result1.32[0] ^ 0x80000000;
  }
stage E4:
  GPR[registerW] = result1;
  commit_float_exception_flags();
\end{lstlisting}}
\newpage
\subsubsection[{\small FFMAWP: Floating-Point Fused Multiply Add Word Pair}]{FFMAWP\hfill Floating-Point Fused Multiply Add Word Pair}
\label{instr:FFMAWP}\index{ffmawp}
 
\paragraph{Encoding} Format \textsf{ALU\_FWPFMAE} (Section~\ref{sec:format-ALU-FWPFMAE}), \verb|alucode3 = 100|\\

\bitfields{ 1/P, 2/11, 1/1, 1/fnegate, 3/floatmode, 5/registerWe, 1/0, 2/11, 3/100, 1/1, 5/registerYe, 1/0, 5/registerZe, 1/0 }


\paragraph{Syntax} \texttt{ffmawp}\emph{fnegate}\emph{floatmode} \emph{registerWe} = \emph{registerZe}, \emph{registerYe}

\paragraph{Description} The \emph{registerWe}, the \emph{registerZe} and the \emph{registerYe} are interpreted as two binary 32 floating-point numbers packed into 64 bits. If \emph{fnegate} is set, the \emph{registerWe} and the \emph{registerZe} words are negated prior to the operation. The \emph{registerYe} words are multiplied by the \emph{registerZe} words, assuming binary 32 floating-point numbers, and the product words are added to the \emph{registerWe} words. The two results are rounded to the binary 32 floating-point format according to the rounding mode and stored into the \emph{registerWe}. This instruction may raise exception bits in the CS register.


\paragraph{Modifier(s)}
\noindent\begin{itemize}
\item floatmode (cf. floatmode in section \ref{par:modifier-floatmode} on page \pageref{par:modifier-floatmode})
\item fnegate (cf. fnegate in section \ref{par:modifier-fnegate} on page \pageref{par:modifier-fnegate})
\end{itemize}

\paragraph{Execution} {Reservation table \textsf{ALU\_LITE} (Section~\ref{sec:reservation-ALU-LITE})}
~
{\begin{lstlisting}
stage ID:
  new fnegate = _ZX_1(fnegate);
  new floatmode = _ZX_3(floatmode);
stage RR:
  new argument3 = GPR[registerYe<<1];
  new argument2 = GPR[registerZe<<1];
  new argument1 = GPR[registerWe<<1];
  new RM = decode_riscv_float_rounding_mode(floatmode);
  new fsign = fnegate ? 0x80000000 : 0;
  for (I = 0; I < 2; I += 1) {
    result1.32[I] = f32_mulAdd(RM, argument3.32[I], argument2.32[I] ^ fsign, argument1.32[I] ^ fsign)
  };
stage E4:
  GPR[registerWe<<1] = result1;
  commit_float_exception_flags();
\end{lstlisting}}
\newpage

\paragraph{Encoding} Format \textsf{ALU\_FWPFMAO} (Section~\ref{sec:format-ALU-FWPFMAO}), \verb|alucode3 = 100|\\

\bitfields{ 1/P, 2/11, 1/1, 1/fnegate, 3/floatmode, 5/registerWo, 1/1, 2/11, 3/100, 1/1, 5/registerYo, 1/0, 5/registerZo, 1/0 }


\paragraph{Syntax} \texttt{ffmawp}\emph{fnegate}\emph{floatmode} \emph{registerWo} = \emph{registerZo}, \emph{registerYo}

\paragraph{Description} The \emph{registerWo}, the \emph{registerZo} and the \emph{registerYo} are interpreted as two binary 32 floating-point numbers packed into 64 bits. If \emph{fnegate} is set, the \emph{registerWo} and the \emph{registerZo} words are negated prior to the operation. The \emph{registerYo} words are multiplied by the \emph{registerZo} words, assuming binary 32 floating-point numbers, and the product words are added to the \emph{registerWo} words. The two results are rounded to the binary 32 floating-point format according to the rounding mode and stored into the \emph{registerWo}. This instruction may raise exception bits in the CS register.


\paragraph{Modifier(s)}
\noindent\begin{itemize}
\item floatmode (cf. floatmode in section \ref{par:modifier-floatmode} on page \pageref{par:modifier-floatmode})
\item fnegate (cf. fnegate in section \ref{par:modifier-fnegate} on page \pageref{par:modifier-fnegate})
\end{itemize}

\paragraph{Execution} {Reservation table \textsf{ALU\_LITE} (Section~\ref{sec:reservation-ALU-LITE})}
~
{\begin{lstlisting}
stage ID:
  new fnegate = _ZX_1(fnegate);
  new floatmode = _ZX_3(floatmode);
stage RR:
  new argument3 = GPR[(registerYo<<1)+1];
  new argument2 = GPR[(registerZo<<1)+1];
  new argument1 = GPR[(registerWo<<1)+1];
  new RM = decode_riscv_float_rounding_mode(floatmode);
  new fsign = fnegate ? 0x80000000 : 0;
  for (I = 0; I < 2; I += 1) {
    result1.32[I] = f32_mulAdd(RM, argument3.32[I], argument2.32[I] ^ fsign, argument1.32[I] ^ fsign)
  };
stage E4:
  GPR[(registerWo<<1)+1] = result1;
  commit_float_exception_flags();
\end{lstlisting}}
\newpage
\subsubsection[{\small FFMSD: Floating-Point Fused Multiply Subtract From Double Word}]{FFMSD\hfill Floating-Point Fused Multiply Subtract From Double Word}
\label{instr:FFMSD}\index{ffmsd}
 
\paragraph{Encoding} Format \textsf{ALU\_FDFMA} (Section~\ref{sec:format-ALU-FDFMA}), \verb|alucode3 = 101|\\

\bitfields{ 1/P, 2/11, 1/1, 1/fnegate, 3/floatmode, 6/registerW, 2/00, 3/101, 1/0, 6/registerY, 6/registerZ }


\paragraph{Syntax} \texttt{ffmsd}\emph{fnegate}\emph{floatmode} \emph{registerW} = \emph{registerZ}, \emph{registerY}

\paragraph{Description} The \emph{registerY} is multiplied by the \emph{registerZ}, assuming binary 64 floating-point numbers, and the product is subtracted from the \emph{registerW}. If \emph{fnegate} is set, the \emph{registerW} and the \emph{registerZ} double words are negated prior to the operation. The result is rounded to the binary 64 floating-point format according to the rounding mode and stored into the \emph{registerW}. This instruction may raise exception bits in the CS register.


\paragraph{Modifier(s)}
\noindent\begin{itemize}
\item floatmode (cf. floatmode in section \ref{par:modifier-floatmode} on page \pageref{par:modifier-floatmode})
\item fnegate (cf. fnegate in section \ref{par:modifier-fnegate} on page \pageref{par:modifier-fnegate})
\end{itemize}

\paragraph{Execution} {Reservation table \textsf{ALU\_LITE} (Section~\ref{sec:reservation-ALU-LITE})}
~
{\begin{lstlisting}
stage ID:
  new fnegate = _ZX_1(fnegate);
  new floatmode = _ZX_3(floatmode);
stage RR:
  new argument3 = GPR[registerY];
  new argument2 = GPR[registerZ];
  new argument1 = GPR[registerW];
  new RM = decode_riscv_float_rounding_mode(floatmode);
  new fsign = fnegate ? 0x8000000000000000 : 0;
  new result1 = f64_mulnAdd(RM, argument3, argument2 ^ fsign, argument1 ^ fsign);
stage E4:
  GPR[registerW] = result1;
  commit_float_exception_flags();
\end{lstlisting}}
\newpage
\subsubsection[{\small FFMSH: Floating-Point Fused Multiply Subtract From Half Word}]{FFMSH\hfill Floating-Point Fused Multiply Subtract From Half Word}
\label{instr:FFMSH}\index{ffmsh}
 
\paragraph{Encoding} Format \textsf{ALU\_FHFMA} (Section~\ref{sec:format-ALU-FHFMA}), \verb|alucode3 = 101|\\

\bitfields{ 1/P, 2/11, 1/1, 1/fnegate, 3/floatmode, 6/registerW, 2/01, 3/101, 1/1, 6/registerY, 6/registerZ }


\paragraph{Syntax} \texttt{ffmsh}\emph{fnegate}\emph{floatmode} \emph{registerW} = \emph{registerZ}, \emph{registerY}

\paragraph{Description} The \emph{registerY} is multiplied by the \emph{registerZ}, assuming binary 16 floating-point numbers, and the product is subtracted from the \emph{registerW}. If \emph{fnegate} is set, the \emph{registerW} and the \emph{registerZ} half words are negated prior to the operation. The result is rounded to the binary 16 floating-point format according to the rounding mode and stored into the \emph{registerW}. This instruction may raise exception bits in the CS register.


\paragraph{Modifier(s)}
\noindent\begin{itemize}
\item floatmode (cf. floatmode in section \ref{par:modifier-floatmode} on page \pageref{par:modifier-floatmode})
\item fnegate (cf. fnegate in section \ref{par:modifier-fnegate} on page \pageref{par:modifier-fnegate})
\end{itemize}

\paragraph{Execution} {Reservation table \textsf{ALU\_LITE} (Section~\ref{sec:reservation-ALU-LITE})}
~
{\begin{lstlisting}
stage ID:
  new fnegate = _ZX_1(fnegate);
  new floatmode = _ZX_3(floatmode);
stage RR:
  new argument3 = GPR[registerY];
  new argument2 = GPR[registerZ];
  new argument1 = GPR[registerW];
  new RM = decode_riscv_float_rounding_mode(floatmode);
  new fsign = fnegate ? 0x8000 : 0;
  new result1 = f16_mulnAdd(RM, argument3, argument2 ^ fsign, argument1 ^ fsign);
stage E3:
  GPR[registerW] = result1;
  commit_float_exception_flags();
\end{lstlisting}}
\newpage
\subsubsection[{\small FFMSHQ: Floating-Point Fused Multiply Subtract From Half Word Quadruple}]{FFMSHQ\hfill Floating-Point Fused Multiply Subtract From Half Word Quadruple}
\label{instr:FFMSHQ}\index{ffmshq}
 
\paragraph{Encoding} Format \textsf{ALU\_FHQFMAE} (Section~\ref{sec:format-ALU-FHQFMAE}), \verb|alucode3 = 101|\\

\bitfields{ 1/P, 2/11, 1/1, 1/fnegate, 3/floatmode, 5/registerWe, 1/0, 2/11, 3/101, 1/1, 5/registerYe, 1/0, 5/registerZe, 1/1 }


\paragraph{Syntax} \texttt{ffmshq}\emph{fnegate}\emph{floatmode} \emph{registerWe} = \emph{registerZe}, \emph{registerYe}

\paragraph{Description} The \emph{registerWe}, the \emph{registerZe} and the \emph{registerYe} are interpreted as four binary 16 floating-point numbers packed into 64 bits. If \emph{fnegate} is set, the \emph{registerWe} and the \emph{registerZe} words are negated prior to the operation. The \emph{registerYe} half words are multiplied by the \emph{registerZe} half words, assuming binary 16 floating-point numbers, and the product half words are subtracted from the \emph{registerWe} half words. The four results are rounded to the binary 16 floating-point format according to the rounding mode and stored into the \emph{registerWe}. This instruction may raise exception bits in the CS register.


\paragraph{Modifier(s)}
\noindent\begin{itemize}
\item floatmode (cf. floatmode in section \ref{par:modifier-floatmode} on page \pageref{par:modifier-floatmode})
\item fnegate (cf. fnegate in section \ref{par:modifier-fnegate} on page \pageref{par:modifier-fnegate})
\end{itemize}

\paragraph{Execution} {Reservation table \textsf{ALU\_LITE} (Section~\ref{sec:reservation-ALU-LITE})}
~
{\begin{lstlisting}
stage ID:
  new fnegate = _ZX_1(fnegate);
  new floatmode = _ZX_3(floatmode);
stage RR:
  new argument3 = GPR[registerYe<<1];
  new argument2 = GPR[registerZe<<1];
  new argument1 = GPR[registerWe<<1];
  new RM = decode_riscv_float_rounding_mode(floatmode);
  new fsign = fnegate ? 0x8000 : 0;
  for (I = 0; I < 4; I += 1) {
    result1.16[I] = f16_mulnAdd(RM, argument3.16[I], argument2.16[I] ^ fsign, argument1.16[I] ^ fsign)
  };
stage E4:
  GPR[registerWe<<1] = result1;
  commit_float_exception_flags();
\end{lstlisting}}
\newpage

\paragraph{Encoding} Format \textsf{ALU\_FHQFMAO} (Section~\ref{sec:format-ALU-FHQFMAO}), \verb|alucode3 = 101|\\

\bitfields{ 1/P, 2/11, 1/1, 1/fnegate, 3/floatmode, 5/registerWo, 1/1, 2/11, 3/101, 1/1, 5/registerYo, 1/0, 5/registerZo, 1/1 }


\paragraph{Syntax} \texttt{ffmshq}\emph{fnegate}\emph{floatmode} \emph{registerWo} = \emph{registerZo}, \emph{registerYo}

\paragraph{Description} The \emph{registerWo}, the \emph{registerZo} and the \emph{registerYo} are interpreted as four binary 16 floating-point numbers packed into 64 bits. If \emph{fnegate} is set, the \emph{registerWo} and the \emph{registerZo} words are negated prior to the operation. The \emph{registerYo} half words are multiplied by the \emph{registerZo} half words, assuming binary 16 floating-point numbers, and the product half words are subtracted from the \emph{registerWo} half words. The four results are rounded to the binary 16 floating-point format according to the rounding mode and stored into the \emph{registerWo}. This instruction may raise exception bits in the CS register.


\paragraph{Modifier(s)}
\noindent\begin{itemize}
\item floatmode (cf. floatmode in section \ref{par:modifier-floatmode} on page \pageref{par:modifier-floatmode})
\item fnegate (cf. fnegate in section \ref{par:modifier-fnegate} on page \pageref{par:modifier-fnegate})
\end{itemize}

\paragraph{Execution} {Reservation table \textsf{ALU\_LITE} (Section~\ref{sec:reservation-ALU-LITE})}
~
{\begin{lstlisting}
stage ID:
  new fnegate = _ZX_1(fnegate);
  new floatmode = _ZX_3(floatmode);
stage RR:
  new argument3 = GPR[(registerYo<<1)+1];
  new argument2 = GPR[(registerZo<<1)+1];
  new argument1 = GPR[(registerWo<<1)+1];
  new RM = decode_riscv_float_rounding_mode(floatmode);
  new fsign = fnegate ? 0x8000 : 0;
  for (I = 0; I < 4; I += 1) {
    result1.16[I] = f16_mulnAdd(RM, argument3.16[I], argument2.16[I] ^ fsign, argument1.16[I] ^ fsign)
  };
stage E4:
  GPR[(registerWo<<1)+1] = result1;
  commit_float_exception_flags();
\end{lstlisting}}
\newpage
\subsubsection[{\small FFMSW: Floating-Point Fused Multiply Subtract From Word}]{FFMSW\hfill Floating-Point Fused Multiply Subtract From Word}
\label{instr:FFMSW}\index{ffmsw}
 
\paragraph{Encoding} Format \textsf{ALU\_FWFMA} (Section~\ref{sec:format-ALU-FWFMA}), \verb|alucode3 = 101|\\

\bitfields{ 1/P, 2/11, 1/1, 1/fnegate, 3/floatmode, 6/registerW, 2/01, 3/101, 1/0, 6/registerY, 6/registerZ }


\paragraph{Syntax} \texttt{ffmsw}\emph{fnegate}\emph{floatmode} \emph{registerW} = \emph{registerZ}, \emph{registerY}

\paragraph{Description} The \emph{registerY} is multiplied by the \emph{registerZ}, assuming binary 32 floating-point numbers, and the product is subtracted from the \emph{registerW}. If \emph{fnegate} is set, the \emph{registerW} and the \emph{registerZ} words are negated prior to the operation. The result is rounded to the binary 32 floating-point format according to the rounding mode and stored into the \emph{registerW}. This instruction may raise exception bits in the CS register.


\paragraph{Modifier(s)}
\noindent\begin{itemize}
\item floatmode (cf. floatmode in section \ref{par:modifier-floatmode} on page \pageref{par:modifier-floatmode})
\item fnegate (cf. fnegate in section \ref{par:modifier-fnegate} on page \pageref{par:modifier-fnegate})
\end{itemize}

\paragraph{Execution} {Reservation table \textsf{ALU\_LITE} (Section~\ref{sec:reservation-ALU-LITE})}
~
{\begin{lstlisting}
stage ID:
  new fnegate = _ZX_1(fnegate);
  new floatmode = _ZX_3(floatmode);
stage RR:
  new argument3 = GPR[registerY];
  new argument2 = GPR[registerZ];
  new argument1 = GPR[registerW];
  new RM = decode_riscv_float_rounding_mode(floatmode);
  new fsign = fnegate ? 0x80000000 : 0;
  new result1 = f32_mulnAdd(RM, argument3, argument2 ^ fsign, argument1 ^ fsign);
stage E4:
  GPR[registerW] = result1;
  commit_float_exception_flags();
\end{lstlisting}}
\newpage
\subsubsection[{\small FFMSWP: Floating-Point Fused Multiply Subtract From Word Pair}]{FFMSWP\hfill Floating-Point Fused Multiply Subtract From Word Pair}
\label{instr:FFMSWP}\index{ffmswp}
 
\paragraph{Encoding} Format \textsf{ALU\_FWPFMAE} (Section~\ref{sec:format-ALU-FWPFMAE}), \verb|alucode3 = 101|\\

\bitfields{ 1/P, 2/11, 1/1, 1/fnegate, 3/floatmode, 5/registerWe, 1/0, 2/11, 3/101, 1/1, 5/registerYe, 1/0, 5/registerZe, 1/0 }


\paragraph{Syntax} \texttt{ffmswp}\emph{fnegate}\emph{floatmode} \emph{registerWe} = \emph{registerZe}, \emph{registerYe}

\paragraph{Description} The \emph{registerWe}, the \emph{registerZe} and the \emph{registerYe} are interpreted as two binary 32 floating-point numbers packed into 128 bits. If \emph{fnegate} is set, the \emph{registerWe} and the \emph{registerZe} words are negated prior to the operation. The \emph{registerYe} words are multiplied by the \emph{registerZe} words, assuming binary 32 floating-point numbers, and the product words are subtracted from the \emph{registerWe} words. The two results are rounded to the binary 32 floating-point format according to the rounding mode and stored into the \emph{registerWe}. This instruction may raise exception bits in the CS register.


\paragraph{Modifier(s)}
\noindent\begin{itemize}
\item floatmode (cf. floatmode in section \ref{par:modifier-floatmode} on page \pageref{par:modifier-floatmode})
\item fnegate (cf. fnegate in section \ref{par:modifier-fnegate} on page \pageref{par:modifier-fnegate})
\end{itemize}

\paragraph{Execution} {Reservation table \textsf{ALU\_LITE} (Section~\ref{sec:reservation-ALU-LITE})}
~
{\begin{lstlisting}
stage ID:
  new fnegate = _ZX_1(fnegate);
  new floatmode = _ZX_3(floatmode);
stage RR:
  new argument3 = GPR[registerYe<<1];
  new argument2 = GPR[registerZe<<1];
  new argument1 = GPR[registerWe<<1];
  new RM = decode_riscv_float_rounding_mode(floatmode);
  new fsign = fnegate ? 0x80000000 : 0;
  for (I = 0; I < 2; I += 1) {
    result1.32[I] = f32_mulnAdd(RM, argument3.32[I], argument2.32[I] ^ fsign, argument1.32[I] ^ fsign)
  };
stage E4:
  GPR[registerWe<<1] = result1;
  commit_float_exception_flags();
\end{lstlisting}}
\newpage

\paragraph{Encoding} Format \textsf{ALU\_FWPFMAO} (Section~\ref{sec:format-ALU-FWPFMAO}), \verb|alucode3 = 101|\\

\bitfields{ 1/P, 2/11, 1/1, 1/fnegate, 3/floatmode, 5/registerWo, 1/1, 2/11, 3/101, 1/1, 5/registerYo, 1/0, 5/registerZo, 1/0 }


\paragraph{Syntax} \texttt{ffmswp}\emph{fnegate}\emph{floatmode} \emph{registerWo} = \emph{registerZo}, \emph{registerYo}

\paragraph{Description} The \emph{registerWo}, the \emph{registerZo} and the \emph{registerYo} are interpreted as two binary 32 floating-point numbers packed into 128 bits. If \emph{fnegate} is set, the \emph{registerWo} and the \emph{registerZo} words are negated prior to the operation. The \emph{registerYo} words are multiplied by the \emph{registerZo} words, assuming binary 32 floating-point numbers, and the product words are subtracted from the \emph{registerWo} words. The two results are rounded to the binary 32 floating-point format according to the rounding mode and stored into the \emph{registerWo}. This instruction may raise exception bits in the CS register.


\paragraph{Modifier(s)}
\noindent\begin{itemize}
\item floatmode (cf. floatmode in section \ref{par:modifier-floatmode} on page \pageref{par:modifier-floatmode})
\item fnegate (cf. fnegate in section \ref{par:modifier-fnegate} on page \pageref{par:modifier-fnegate})
\end{itemize}

\paragraph{Execution} {Reservation table \textsf{ALU\_LITE} (Section~\ref{sec:reservation-ALU-LITE})}
~
{\begin{lstlisting}
stage ID:
  new fnegate = _ZX_1(fnegate);
  new floatmode = _ZX_3(floatmode);
stage RR:
  new argument3 = GPR[(registerYo<<1)+1];
  new argument2 = GPR[(registerZo<<1)+1];
  new argument1 = GPR[(registerWo<<1)+1];
  new RM = decode_riscv_float_rounding_mode(floatmode);
  new fsign = fnegate ? 0x80000000 : 0;
  for (I = 0; I < 2; I += 1) {
    result1.32[I] = f32_mulnAdd(RM, argument3.32[I], argument2.32[I] ^ fsign, argument1.32[I] ^ fsign)
  };
stage E4:
  GPR[(registerWo<<1)+1] = result1;
  commit_float_exception_flags();
\end{lstlisting}}
\newpage
\subsubsection[{\small FIXEDD: Floating Point Conversion to Fixed-Point Double Word}]{FIXEDD\hfill Floating Point Conversion to Fixed-Point Double Word}
\label{instr:FIXEDD}\index{fixedd}
 
\paragraph{Encoding} Format \textsf{ALU\_FDFIXED} (Section~\ref{sec:format-ALU-FDFIXED}), \verb|fpucode6 = 0100|\\

\bitfields{ 1/P, 2/11, 1/1, 1/1, 3/floatmode, 6/registerW, 2/00, 3/001, 1/0, 4/0100, 1/--, 1/0, 6/registerZ }


\paragraph{Syntax} \texttt{fixedd}\emph{floatmode} \emph{registerW} = \emph{registerZ}

\paragraph{Description} The \emph{registerZ} interpreted as a 64-bit floating-point number is converted to a 64-bit signed integer according to the rounding mode. This instruction may raise exception bits in the CS register.


\paragraph{Modifier(s)}
\noindent\begin{itemize}
\item floatmode (cf. floatmode in section \ref{par:modifier-floatmode} on page \pageref{par:modifier-floatmode})
\end{itemize}

\paragraph{Execution} {Reservation table \textsf{ALU\_LITE} (Section~\ref{sec:reservation-ALU-LITE})}
~
{\begin{lstlisting}
stage ID:
  new floatmode = _ZX_3(floatmode);
stage RR:
  new argument2 = GPR[registerZ];
  new RM = decode_riscv_float_rounding_mode(floatmode);
  new result1 = f64_to_i64(RM, argument2);
stage E4:
  GPR[registerW] = result1;
  commit_float_exception_flags();
\end{lstlisting}}
\newpage
\subsubsection[{\small FIXEDDW: Floating Point Conversion from Double Word to Fixed-Point Word}]{FIXEDDW\hfill Floating Point Conversion from Double Word to Fixed-Point Word}
\label{instr:FIXEDDW}\index{fixeddw}
 
\paragraph{Encoding} Format \textsf{ALU\_FWFIXED} (Section~\ref{sec:format-ALU-FWFIXED}), \verb|fpucode6 = 0110|\\

\bitfields{ 1/P, 2/11, 1/1, 1/1, 3/floatmode, 6/registerW, 2/01, 3/001, 1/0, 4/0110, 1/--, 1/0, 6/registerZ }


\paragraph{Syntax} \texttt{fixeddw}\emph{floatmode} \emph{registerW} = \emph{registerZ}

\paragraph{Description} The \emph{registerZ} interpreted as a 64-bit floating-point number is converted to a 32-bit signed integer according to the rounding mode. This instruction may raise exception bits in the CS register.


\paragraph{Modifier(s)}
\noindent\begin{itemize}
\item floatmode (cf. floatmode in section \ref{par:modifier-floatmode} on page \pageref{par:modifier-floatmode})
\end{itemize}

\paragraph{Execution} {Reservation table \textsf{ALU\_LITE} (Section~\ref{sec:reservation-ALU-LITE})}
~
{\begin{lstlisting}
stage ID:
  new floatmode = _ZX_3(floatmode);
stage RR:
  new argument2 = GPR[registerZ];
  new RM = decode_riscv_float_rounding_mode(floatmode);
  new result1 = f64_to_i32(RM, argument2);
stage E4:
  GPR[registerW] = result1;
  commit_float_exception_flags();
\end{lstlisting}}
\newpage
\subsubsection[{\small FIXEDHQ: Floating Point Conversion to Fixed-Point Half Word Quadruple}]{FIXEDHQ\hfill Floating Point Conversion to Fixed-Point Half Word Quadruple}
\label{instr:FIXEDHQ}\index{fixedhq}
 
\paragraph{Encoding} Format \textsf{ALU\_FHQFIXEDE} (Section~\ref{sec:format-ALU-FHQFIXEDE}), \verb|fpucode5 = 00100|\\

\bitfields{ 1/P, 2/11, 1/1, 1/1, 3/floatmode, 5/registerWe, 1/0, 2/11, 3/001, 1/1, 5/00100, 1/0, 5/registerZe, 1/1 }


\paragraph{Syntax} \texttt{fixedhq}\emph{floatmode} \emph{registerWe} = \emph{registerZe}

\paragraph{Description} The \emph{registerZe} is interpreted as four binary 16 floating-point numbers packed into 64 bits. The \emph{registerZe} half words are converted to 16-bit signed integers according to the rounding mode. The four resulting half words are packed and stored into the \emph{registerWe}. This instruction may raise exception bits in the CS register.


\paragraph{Modifier(s)}
\noindent\begin{itemize}
\item floatmode (cf. floatmode in section \ref{par:modifier-floatmode} on page \pageref{par:modifier-floatmode})
\end{itemize}

\paragraph{Execution} {Reservation table \textsf{ALU\_LITE} (Section~\ref{sec:reservation-ALU-LITE})}
~
{\begin{lstlisting}
stage ID:
  new floatmode = _ZX_3(floatmode);
stage RR:
  new argument2 = GPR[registerZe<<1];
  new RM = decode_riscv_float_rounding_mode(floatmode);
  for (I = 0; I < 4; I += 1) {
    result1.16[I] = f16_to_i16(RM, argument2.16[I])
  };
stage E4:
  GPR[registerWe<<1] = result1;
  commit_float_exception_flags();
\end{lstlisting}}
\newpage

\paragraph{Encoding} Format \textsf{ALU\_FHQFIXEDO} (Section~\ref{sec:format-ALU-FHQFIXEDO}), \verb|fpucode5 = 00100|\\

\bitfields{ 1/P, 2/11, 1/1, 1/1, 3/floatmode, 5/registerWo, 1/1, 2/11, 3/001, 1/1, 5/00100, 1/0, 5/registerZo, 1/1 }


\paragraph{Syntax} \texttt{fixedhq}\emph{floatmode} \emph{registerWo} = \emph{registerZo}

\paragraph{Description} The \emph{registerZo} is interpreted as four binary 16 floating-point numbers packed into 64 bits. The \emph{registerZo} half words are converted to 16-bit signed integers according to the rounding mode. The four resulting half words are packed and stored into the \emph{registerWo}. This instruction may raise exception bits in the CS register.


\paragraph{Modifier(s)}
\noindent\begin{itemize}
\item floatmode (cf. floatmode in section \ref{par:modifier-floatmode} on page \pageref{par:modifier-floatmode})
\end{itemize}

\paragraph{Execution} {Reservation table \textsf{ALU\_LITE} (Section~\ref{sec:reservation-ALU-LITE})}
~
{\begin{lstlisting}
stage ID:
  new floatmode = _ZX_3(floatmode);
stage RR:
  new argument2 = GPR[(registerZo<<1)+1];
  new RM = decode_riscv_float_rounding_mode(floatmode);
  for (I = 0; I < 4; I += 1) {
    result1.16[I] = f16_to_i16(RM, argument2.16[I])
  };
stage E4:
  GPR[(registerWo<<1)+1] = result1;
  commit_float_exception_flags();
\end{lstlisting}}
\newpage
\subsubsection[{\small FIXEDUD: Floating Point Conversion to Unsigned Fixed-Point Double Word}]{FIXEDUD\hfill Floating Point Conversion to Unsigned Fixed-Point Double Word}
\label{instr:FIXEDUD}\index{fixedud}
 
\paragraph{Encoding} Format \textsf{ALU\_FDFIXED} (Section~\ref{sec:format-ALU-FDFIXED}), \verb|fpucode6 = 0101|\\

\bitfields{ 1/P, 2/11, 1/1, 1/1, 3/floatmode, 6/registerW, 2/00, 3/001, 1/0, 4/0101, 1/--, 1/0, 6/registerZ }


\paragraph{Syntax} \texttt{fixedud}\emph{floatmode} \emph{registerW} = \emph{registerZ}

\paragraph{Description} The \emph{registerZ} interpreted as a 64-bit floating-point number is converted to a 64-bit unsigned integer according to the rounding mode. This instruction may raise exception bits in the CS register.


\paragraph{Modifier(s)}
\noindent\begin{itemize}
\item floatmode (cf. floatmode in section \ref{par:modifier-floatmode} on page \pageref{par:modifier-floatmode})
\end{itemize}

\paragraph{Execution} {Reservation table \textsf{ALU\_LITE} (Section~\ref{sec:reservation-ALU-LITE})}
~
{\begin{lstlisting}
stage ID:
  new floatmode = _ZX_3(floatmode);
stage RR:
  new argument2 = GPR[registerZ];
  new RM = decode_riscv_float_rounding_mode(floatmode);
  new result1 = f64_to_ui64(RM, argument2);
stage E4:
  GPR[registerW] = result1;
  commit_float_exception_flags();
\end{lstlisting}}
\newpage
\subsubsection[{\small FIXEDUDW: Floating Point Conversion from Double Word to Unsigned Fixed-Point Word}]{FIXEDUDW\hfill Floating Point Conversion from Double Word to Unsigned Fixed-Point Word}
\label{instr:FIXEDUDW}\index{fixedudw}
 
\paragraph{Encoding} Format \textsf{ALU\_FWFIXED} (Section~\ref{sec:format-ALU-FWFIXED}), \verb|fpucode6 = 0111|\\

\bitfields{ 1/P, 2/11, 1/1, 1/1, 3/floatmode, 6/registerW, 2/01, 3/001, 1/0, 4/0111, 1/--, 1/0, 6/registerZ }


\paragraph{Syntax} \texttt{fixedudw}\emph{floatmode} \emph{registerW} = \emph{registerZ}

\paragraph{Description} The \emph{registerZ} interpreted as a 64-bit floating-point number is converted to a 32-bit unsigned integer according to the rounding mode. This instruction may raise exception bits in the CS register.


\paragraph{Modifier(s)}
\noindent\begin{itemize}
\item floatmode (cf. floatmode in section \ref{par:modifier-floatmode} on page \pageref{par:modifier-floatmode})
\end{itemize}

\paragraph{Execution} {Reservation table \textsf{ALU\_LITE} (Section~\ref{sec:reservation-ALU-LITE})}
~
{\begin{lstlisting}
stage ID:
  new floatmode = _ZX_3(floatmode);
stage RR:
  new argument2 = GPR[registerZ];
  new RM = decode_riscv_float_rounding_mode(floatmode);
  new result1 = f64_to_ui32(RM, argument2);
stage E4:
  GPR[registerW] = result1;
  commit_float_exception_flags();
\end{lstlisting}}
\newpage
\subsubsection[{\small FIXEDUHQ: Floating Point Conversion to Unsigned Fixed-Point Half Word Quadruple}]{FIXEDUHQ\hfill Floating Point Conversion to Unsigned Fixed-Point Half Word Quadruple}
\label{instr:FIXEDUHQ}\index{fixeduhq}
 
\paragraph{Encoding} Format \textsf{ALU\_FHQFIXEDE} (Section~\ref{sec:format-ALU-FHQFIXEDE}), \verb|fpucode5 = 00101|\\

\bitfields{ 1/P, 2/11, 1/1, 1/1, 3/floatmode, 5/registerWe, 1/0, 2/11, 3/001, 1/1, 5/00101, 1/0, 5/registerZe, 1/1 }


\paragraph{Syntax} \texttt{fixeduhq}\emph{floatmode} \emph{registerWe} = \emph{registerZe}

\paragraph{Description} The \emph{registerZe} is interpreted as four binary 16 floating-point numbers packed into 64 bits. The \emph{registerZe} half words are converted to 16-bit unsigned integers according to the rounding mode. The four resulting half words are packed and stored into the \emph{registerWe}. This instruction may raise exception bits in the CS register.


\paragraph{Modifier(s)}
\noindent\begin{itemize}
\item floatmode (cf. floatmode in section \ref{par:modifier-floatmode} on page \pageref{par:modifier-floatmode})
\end{itemize}

\paragraph{Execution} {Reservation table \textsf{ALU\_LITE} (Section~\ref{sec:reservation-ALU-LITE})}
~
{\begin{lstlisting}
stage ID:
  new floatmode = _ZX_3(floatmode);
stage RR:
  new argument2 = GPR[registerZe<<1];
  new RM = decode_riscv_float_rounding_mode(floatmode);
  for (I = 0; I < 4; I += 1) {
    result1.16[I] = f16_to_ui16(RM, argument2.16[I])
  };
stage E4:
  GPR[registerWe<<1] = result1;
  commit_float_exception_flags();
\end{lstlisting}}
\newpage

\paragraph{Encoding} Format \textsf{ALU\_FHQFIXEDO} (Section~\ref{sec:format-ALU-FHQFIXEDO}), \verb|fpucode5 = 00101|\\

\bitfields{ 1/P, 2/11, 1/1, 1/1, 3/floatmode, 5/registerWo, 1/1, 2/11, 3/001, 1/1, 5/00101, 1/0, 5/registerZo, 1/1 }


\paragraph{Syntax} \texttt{fixeduhq}\emph{floatmode} \emph{registerWo} = \emph{registerZo}

\paragraph{Description} The \emph{registerZo} is interpreted as four binary 16 floating-point numbers packed into 64 bits. The \emph{registerZo} half words are converted to 16-bit unsigned integers according to the rounding mode. The four resulting half words are packed and stored into the \emph{registerWo}. This instruction may raise exception bits in the CS register.


\paragraph{Modifier(s)}
\noindent\begin{itemize}
\item floatmode (cf. floatmode in section \ref{par:modifier-floatmode} on page \pageref{par:modifier-floatmode})
\end{itemize}

\paragraph{Execution} {Reservation table \textsf{ALU\_LITE} (Section~\ref{sec:reservation-ALU-LITE})}
~
{\begin{lstlisting}
stage ID:
  new floatmode = _ZX_3(floatmode);
stage RR:
  new argument2 = GPR[(registerZo<<1)+1];
  new RM = decode_riscv_float_rounding_mode(floatmode);
  for (I = 0; I < 4; I += 1) {
    result1.16[I] = f16_to_ui16(RM, argument2.16[I])
  };
stage E4:
  GPR[(registerWo<<1)+1] = result1;
  commit_float_exception_flags();
\end{lstlisting}}
\newpage
\subsubsection[{\small FIXEDUW: Floating Point Conversion to Unsigned Fixed-Point}]{FIXEDUW\hfill Floating Point Conversion to Unsigned Fixed-Point}
\label{instr:FIXEDUW}\index{fixeduw}
 
\paragraph{Encoding} Format \textsf{ALU\_FWFIXED} (Section~\ref{sec:format-ALU-FWFIXED}), \verb|fpucode6 = 0101|\\

\bitfields{ 1/P, 2/11, 1/1, 1/1, 3/floatmode, 6/registerW, 2/01, 3/001, 1/0, 4/0101, 1/--, 1/0, 6/registerZ }


\paragraph{Syntax} \texttt{fixeduw}\emph{floatmode} \emph{registerW} = \emph{registerZ}

\paragraph{Description} The \emph{registerZ} interpreted as a 32-bit floating-point number is converted to a 32-bit unsigned integer according to the rounding mode. This instruction may raise exception bits in the CS register.


\paragraph{Modifier(s)}
\noindent\begin{itemize}
\item floatmode (cf. floatmode in section \ref{par:modifier-floatmode} on page \pageref{par:modifier-floatmode})
\end{itemize}

\paragraph{Execution} {Reservation table \textsf{ALU\_LITE} (Section~\ref{sec:reservation-ALU-LITE})}
~
{\begin{lstlisting}
stage ID:
  new floatmode = _ZX_3(floatmode);
stage RR:
  new argument2 = GPR[registerZ];
  new RM = decode_riscv_float_rounding_mode(floatmode);
  new result1 = f32_to_ui32(RM, argument2);
stage E4:
  GPR[registerW] = result1;
  commit_float_exception_flags();
\end{lstlisting}}
\newpage
\subsubsection[{\small FIXEDUWD: Floating Point Conversion from Word to Unsigned Fixed-Point Double Word}]{FIXEDUWD\hfill Floating Point Conversion from Word to Unsigned Fixed-Point Double Word}
\label{instr:FIXEDUWD}\index{fixeduwd}
 
\paragraph{Encoding} Format \textsf{ALU\_FDFIXED} (Section~\ref{sec:format-ALU-FDFIXED}), \verb|fpucode6 = 0111|\\

\bitfields{ 1/P, 2/11, 1/1, 1/1, 3/floatmode, 6/registerW, 2/00, 3/001, 1/0, 4/0111, 1/--, 1/0, 6/registerZ }


\paragraph{Syntax} \texttt{fixeduwd}\emph{floatmode} \emph{registerW} = \emph{registerZ}

\paragraph{Description} The \emph{registerZ} interpreted as a 32-bit floating-point number is converted to a 64-bit unsigned integer according to the rounding mode. This instruction may raise exception bits in the CS register.


\paragraph{Modifier(s)}
\noindent\begin{itemize}
\item floatmode (cf. floatmode in section \ref{par:modifier-floatmode} on page \pageref{par:modifier-floatmode})
\end{itemize}

\paragraph{Execution} {Reservation table \textsf{ALU\_LITE} (Section~\ref{sec:reservation-ALU-LITE})}
~
{\begin{lstlisting}
stage ID:
  new floatmode = _ZX_3(floatmode);
stage RR:
  new argument2 = GPR[registerZ];
  new RM = decode_riscv_float_rounding_mode(floatmode);
  new result1 = f32_to_ui64(RM, argument2);
stage E4:
  GPR[registerW] = result1;
  commit_float_exception_flags();
\end{lstlisting}}
\newpage
\subsubsection[{\small FIXEDUWP: Floating Point Conversion to Unsigned Fixed-Point Word Pair}]{FIXEDUWP\hfill Floating Point Conversion to Unsigned Fixed-Point Word Pair}
\label{instr:FIXEDUWP}\index{fixeduwp}
 
\paragraph{Encoding} Format \textsf{ALU\_FWPFIXED} (Section~\ref{sec:format-ALU-FWPFIXED}), \verb|fpucode6 = 0101|\\

\bitfields{ 1/P, 2/11, 1/1, 1/1, 3/floatmode, 6/registerW, 2/00, 3/001, 1/0, 4/0101, 1/--, 1/1, 6/registerZ }


\paragraph{Syntax} \texttt{fixeduwp}\emph{floatmode} \emph{registerW} = \emph{registerZ}

\paragraph{Description} The \emph{registerZ} is interpreted as two binary 32 floating-point numbers packed into 64 bits. The \emph{registerZ} words are converted to 32-bit unsigned integers according to the rounding mode. The two resulting words are packed and stored into the \emph{registerW}. This instruction may raise exception bits in the CS register.


\paragraph{Modifier(s)}
\noindent\begin{itemize}
\item floatmode (cf. floatmode in section \ref{par:modifier-floatmode} on page \pageref{par:modifier-floatmode})
\end{itemize}

\paragraph{Execution} {Reservation table \textsf{ALU\_LITE} (Section~\ref{sec:reservation-ALU-LITE})}
~
{\begin{lstlisting}
stage ID:
  new floatmode = _ZX_3(floatmode);
stage RR:
  new argument2 = GPR[registerZ];
  new RM = decode_riscv_float_rounding_mode(floatmode);
  for (I = 0; I < 2; I += 1) {
    result1.32[I] = f32_to_ui32(RM, argument2.32[I])
  };
stage E4:
  GPR[registerW] = result1;
  commit_float_exception_flags();
\end{lstlisting}}
\newpage
\subsubsection[{\small FIXEDW: Floating Point Conversion to Fixed-Point}]{FIXEDW\hfill Floating Point Conversion to Fixed-Point}
\label{instr:FIXEDW}\index{fixedw}
 
\paragraph{Encoding} Format \textsf{ALU\_FWFIXED} (Section~\ref{sec:format-ALU-FWFIXED}), \verb|fpucode6 = 0100|\\

\bitfields{ 1/P, 2/11, 1/1, 1/1, 3/floatmode, 6/registerW, 2/01, 3/001, 1/0, 4/0100, 1/--, 1/0, 6/registerZ }


\paragraph{Syntax} \texttt{fixedw}\emph{floatmode} \emph{registerW} = \emph{registerZ}

\paragraph{Description} The \emph{registerZ} interpreted as a 32-bit floating-point number is converted to a 32-bit signed integer according to the rounding mode. This instruction may raise exception bits in the CS register.


\paragraph{Modifier(s)}
\noindent\begin{itemize}
\item floatmode (cf. floatmode in section \ref{par:modifier-floatmode} on page \pageref{par:modifier-floatmode})
\end{itemize}

\paragraph{Execution} {Reservation table \textsf{ALU\_LITE} (Section~\ref{sec:reservation-ALU-LITE})}
~
{\begin{lstlisting}
stage ID:
  new floatmode = _ZX_3(floatmode);
stage RR:
  new argument2 = GPR[registerZ];
  new RM = decode_riscv_float_rounding_mode(floatmode);
  new result1 = f32_to_i32(RM, argument2);
stage E4:
  GPR[registerW] = result1;
  commit_float_exception_flags();
\end{lstlisting}}
\newpage
\subsubsection[{\small FIXEDWD: Floating Point Conversion from Word to Fixed-Point Double Word}]{FIXEDWD\hfill Floating Point Conversion from Word to Fixed-Point Double Word}
\label{instr:FIXEDWD}\index{fixedwd}
 
\paragraph{Encoding} Format \textsf{ALU\_FDFIXED} (Section~\ref{sec:format-ALU-FDFIXED}), \verb|fpucode6 = 0110|\\

\bitfields{ 1/P, 2/11, 1/1, 1/1, 3/floatmode, 6/registerW, 2/00, 3/001, 1/0, 4/0110, 1/--, 1/0, 6/registerZ }


\paragraph{Syntax} \texttt{fixedwd}\emph{floatmode} \emph{registerW} = \emph{registerZ}

\paragraph{Description} The \emph{registerZ} interpreted as a 32-bit floating-point number is converted to a 64-bit signed integer according to the rounding mode. This instruction may raise exception bits in the CS register.


\paragraph{Modifier(s)}
\noindent\begin{itemize}
\item floatmode (cf. floatmode in section \ref{par:modifier-floatmode} on page \pageref{par:modifier-floatmode})
\end{itemize}

\paragraph{Execution} {Reservation table \textsf{ALU\_LITE} (Section~\ref{sec:reservation-ALU-LITE})}
~
{\begin{lstlisting}
stage ID:
  new floatmode = _ZX_3(floatmode);
stage RR:
  new argument2 = GPR[registerZ];
  new RM = decode_riscv_float_rounding_mode(floatmode);
  new result1 = f32_to_i64(RM, argument2);
stage E4:
  GPR[registerW] = result1;
  commit_float_exception_flags();
\end{lstlisting}}
\newpage
\subsubsection[{\small FIXEDWP: Floating Point Conversion to Fixed-Point Word Pair}]{FIXEDWP\hfill Floating Point Conversion to Fixed-Point Word Pair}
\label{instr:FIXEDWP}\index{fixedwp}
 
\paragraph{Encoding} Format \textsf{ALU\_FWPFIXED} (Section~\ref{sec:format-ALU-FWPFIXED}), \verb|fpucode6 = 0100|\\

\bitfields{ 1/P, 2/11, 1/1, 1/1, 3/floatmode, 6/registerW, 2/00, 3/001, 1/0, 4/0100, 1/--, 1/1, 6/registerZ }


\paragraph{Syntax} \texttt{fixedwp}\emph{floatmode} \emph{registerW} = \emph{registerZ}

\paragraph{Description} The \emph{registerZ} is interpreted as two binary 32 floating-point numbers packed into 64 bits. The \emph{registerZ} words are converted to 32-bit signed integers according to the rounding mode. The two resulting words are packed and stored into the \emph{registerW}. This instruction may raise exception bits in the CS register.


\paragraph{Modifier(s)}
\noindent\begin{itemize}
\item floatmode (cf. floatmode in section \ref{par:modifier-floatmode} on page \pageref{par:modifier-floatmode})
\end{itemize}

\paragraph{Execution} {Reservation table \textsf{ALU\_LITE} (Section~\ref{sec:reservation-ALU-LITE})}
~
{\begin{lstlisting}
stage ID:
  new floatmode = _ZX_3(floatmode);
stage RR:
  new argument2 = GPR[registerZ];
  new RM = decode_riscv_float_rounding_mode(floatmode);
  for (I = 0; I < 2; I += 1) {
    result1.32[I] = f32_to_i32(RM, argument2.32[I])
  };
stage E4:
  GPR[registerW] = result1;
  commit_float_exception_flags();
\end{lstlisting}}
\newpage
\subsubsection[{\small FLOATD: Floating Point Conversion from Fixed-Point Double Word}]{FLOATD\hfill Floating Point Conversion from Fixed-Point Double Word}
\label{instr:FLOATD}\index{floatd}
 
\paragraph{Encoding} Format \textsf{ALU\_FDFLOAT} (Section~\ref{sec:format-ALU-FDFLOAT}), \verb|fpucode6 = 0000|\\

\bitfields{ 1/P, 2/11, 1/1, 1/1, 3/floatmode, 6/registerW, 2/00, 3/001, 1/0, 4/0000, 1/--, 1/0, 6/registerZ }


\paragraph{Syntax} \texttt{floatd}\emph{floatmode} \emph{registerW} = \emph{registerZ}

\paragraph{Description} The \emph{registerZ} interpreted as a 64-bit signed integer is converted to a 64-bit floating-point number according to the rounding mode. This instruction may raise exception bits in the CS register.


\paragraph{Modifier(s)}
\noindent\begin{itemize}
\item floatmode (cf. floatmode in section \ref{par:modifier-floatmode} on page \pageref{par:modifier-floatmode})
\end{itemize}

\paragraph{Execution} {Reservation table \textsf{ALU\_LITE} (Section~\ref{sec:reservation-ALU-LITE})}
~
{\begin{lstlisting}
stage ID:
  new floatmode = _ZX_3(floatmode);
stage RR:
  new argument2 = GPR[registerZ];
  new RM = decode_riscv_float_rounding_mode(floatmode);
  new result1 = i64_to_f64(RM, argument2);
stage E4:
  GPR[registerW] = result1;
  commit_float_exception_flags();
\end{lstlisting}}
\newpage
\subsubsection[{\small FLOATDW: Floating Point Conversion from Fixed-Point Double Word to Word}]{FLOATDW\hfill Floating Point Conversion from Fixed-Point Double Word to Word}
\label{instr:FLOATDW}\index{floatdw}
 
\paragraph{Encoding} Format \textsf{ALU\_FWFLOAT} (Section~\ref{sec:format-ALU-FWFLOAT}), \verb|fpucode6 = 0010|\\

\bitfields{ 1/P, 2/11, 1/1, 1/1, 3/floatmode, 6/registerW, 2/01, 3/001, 1/0, 4/0010, 1/--, 1/0, 6/registerZ }


\paragraph{Syntax} \texttt{floatdw}\emph{floatmode} \emph{registerW} = \emph{registerZ}

\paragraph{Description} The \emph{registerZ} interpreted as a 64-bit signed integer is converted to a 32-bit floating-point number according to the rounding mode. This instruction may raise exception bits in the CS register.


\paragraph{Modifier(s)}
\noindent\begin{itemize}
\item floatmode (cf. floatmode in section \ref{par:modifier-floatmode} on page \pageref{par:modifier-floatmode})
\end{itemize}

\paragraph{Execution} {Reservation table \textsf{ALU\_LITE} (Section~\ref{sec:reservation-ALU-LITE})}
~
{\begin{lstlisting}
stage ID:
  new floatmode = _ZX_3(floatmode);
stage RR:
  new argument2 = GPR[registerZ];
  new RM = decode_riscv_float_rounding_mode(floatmode);
  new result1 = i64_to_f32(RM, argument2);
stage E4:
  GPR[registerW] = result1;
  commit_float_exception_flags();
\end{lstlisting}}
\newpage
\subsubsection[{\small FLOATHQ: Floating Point Conversion from Fixed-Point Half Word Quadruple}]{FLOATHQ\hfill Floating Point Conversion from Fixed-Point Half Word Quadruple}
\label{instr:FLOATHQ}\index{floathq}
 
\paragraph{Encoding} Format \textsf{ALU\_FHQFLOATE} (Section~\ref{sec:format-ALU-FHQFLOATE}), \verb|fpucode6 = 0000|\\

\bitfields{ 1/P, 2/11, 1/1, 1/1, 3/floatmode, 5/registerWe, 1/0, 2/11, 3/001, 1/1, 4/0000, 1/--, 1/0, 5/registerZe, 1/1 }


\paragraph{Syntax} \texttt{floathq}\emph{floatmode} \emph{registerWe} = \emph{registerZe}

\paragraph{Description} The \emph{registerZe} is interpreted as four 16-bit signed integers packed into 64 bits. The \emph{registerZe} half words are converted to binary 16 floating-point numbers according to the rounding mode. The four resulting half words are packed and stored into the \emph{registerWe}. This instruction may raise exception bits in the CS register.


\paragraph{Modifier(s)}
\noindent\begin{itemize}
\item floatmode (cf. floatmode in section \ref{par:modifier-floatmode} on page \pageref{par:modifier-floatmode})
\end{itemize}

\paragraph{Execution} {Reservation table \textsf{ALU\_LITE} (Section~\ref{sec:reservation-ALU-LITE})}
~
{\begin{lstlisting}
stage ID:
  new floatmode = _ZX_3(floatmode);
stage RR:
  new argument2 = GPR[registerZe<<1];
  new RM = decode_riscv_float_rounding_mode(floatmode);
  for (I = 0; I < 4; I += 1) {
    result1.16[I] = i16_to_f16(RM, argument2.16[I])
  };
stage E4:
  GPR[registerWe<<1] = result1;
  commit_float_exception_flags();
\end{lstlisting}}
\newpage

\paragraph{Encoding} Format \textsf{ALU\_FHQFLOATO} (Section~\ref{sec:format-ALU-FHQFLOATO}), \verb|fpucode6 = 0000|\\

\bitfields{ 1/P, 2/11, 1/1, 1/1, 3/floatmode, 5/registerWo, 1/1, 2/11, 3/001, 1/1, 4/0000, 1/--, 1/0, 5/registerZo, 1/1 }


\paragraph{Syntax} \texttt{floathq}\emph{floatmode} \emph{registerWo} = \emph{registerZo}

\paragraph{Description} The \emph{registerZo} is interpreted as four 16-bit signed integers packed into 64 bits. The \emph{registerZo} half words are converted to binary 16 floating-point numbers according to the rounding mode. The four resulting half words are packed and stored into the \emph{registerWo}. This instruction may raise exception bits in the CS register.


\paragraph{Modifier(s)}
\noindent\begin{itemize}
\item floatmode (cf. floatmode in section \ref{par:modifier-floatmode} on page \pageref{par:modifier-floatmode})
\end{itemize}

\paragraph{Execution} {Reservation table \textsf{ALU\_LITE} (Section~\ref{sec:reservation-ALU-LITE})}
~
{\begin{lstlisting}
stage ID:
  new floatmode = _ZX_3(floatmode);
stage RR:
  new argument2 = GPR[(registerZo<<1)+1];
  new RM = decode_riscv_float_rounding_mode(floatmode);
  for (I = 0; I < 4; I += 1) {
    result1.16[I] = i16_to_f16(RM, argument2.16[I])
  };
stage E4:
  GPR[(registerWo<<1)+1] = result1;
  commit_float_exception_flags();
\end{lstlisting}}
\newpage
\subsubsection[{\small FLOATUD: Floating Point Conversion from Unsigned Fixed-Point Double Word}]{FLOATUD\hfill Floating Point Conversion from Unsigned Fixed-Point Double Word}
\label{instr:FLOATUD}\index{floatud}
 
\paragraph{Encoding} Format \textsf{ALU\_FDFLOAT} (Section~\ref{sec:format-ALU-FDFLOAT}), \verb|fpucode6 = 0001|\\

\bitfields{ 1/P, 2/11, 1/1, 1/1, 3/floatmode, 6/registerW, 2/00, 3/001, 1/0, 4/0001, 1/--, 1/0, 6/registerZ }


\paragraph{Syntax} \texttt{floatud}\emph{floatmode} \emph{registerW} = \emph{registerZ}

\paragraph{Description} The \emph{registerZ} interpreted as a 64-bit unsigned integer is converted to a 64-bit floating-point number according to the rounding mode. This instruction may raise exception bits in the CS register.


\paragraph{Modifier(s)}
\noindent\begin{itemize}
\item floatmode (cf. floatmode in section \ref{par:modifier-floatmode} on page \pageref{par:modifier-floatmode})
\end{itemize}

\paragraph{Execution} {Reservation table \textsf{ALU\_LITE} (Section~\ref{sec:reservation-ALU-LITE})}
~
{\begin{lstlisting}
stage ID:
  new floatmode = _ZX_3(floatmode);
stage RR:
  new argument2 = GPR[registerZ];
  new RM = decode_riscv_float_rounding_mode(floatmode);
  new result1 = ui64_to_f64(RM, argument2);
stage E4:
  GPR[registerW] = result1;
  commit_float_exception_flags();
\end{lstlisting}}
\newpage
\subsubsection[{\small FLOATUDW: Floating Point Conversion from Unsigned Fixed-Point Double Word to Word}]{FLOATUDW\hfill Floating Point Conversion from Unsigned Fixed-Point Double Word to Word}
\label{instr:FLOATUDW}\index{floatudw}
 
\paragraph{Encoding} Format \textsf{ALU\_FWFLOAT} (Section~\ref{sec:format-ALU-FWFLOAT}), \verb|fpucode6 = 0011|\\

\bitfields{ 1/P, 2/11, 1/1, 1/1, 3/floatmode, 6/registerW, 2/01, 3/001, 1/0, 4/0011, 1/--, 1/0, 6/registerZ }


\paragraph{Syntax} \texttt{floatudw}\emph{floatmode} \emph{registerW} = \emph{registerZ}

\paragraph{Description} The \emph{registerZ} interpreted as a 64-bit unsigned integer is converted to a 32-bit floating-point number according to the rounding mode. This instruction may raise exception bits in the CS register.


\paragraph{Modifier(s)}
\noindent\begin{itemize}
\item floatmode (cf. floatmode in section \ref{par:modifier-floatmode} on page \pageref{par:modifier-floatmode})
\end{itemize}

\paragraph{Execution} {Reservation table \textsf{ALU\_LITE} (Section~\ref{sec:reservation-ALU-LITE})}
~
{\begin{lstlisting}
stage ID:
  new floatmode = _ZX_3(floatmode);
stage RR:
  new argument2 = GPR[registerZ];
  new RM = decode_riscv_float_rounding_mode(floatmode);
  new result1 = ui64_to_f32(RM, argument2);
stage E4:
  GPR[registerW] = result1;
  commit_float_exception_flags();
\end{lstlisting}}
\newpage
\subsubsection[{\small FLOATUHQ: Floating Point Conversion from Unsigned Fixed-Point Half Word Quadruple}]{FLOATUHQ\hfill Floating Point Conversion from Unsigned Fixed-Point Half Word Quadruple}
\label{instr:FLOATUHQ}\index{floatuhq}
 
\paragraph{Encoding} Format \textsf{ALU\_FHQFLOATE} (Section~\ref{sec:format-ALU-FHQFLOATE}), \verb|fpucode6 = 0001|\\

\bitfields{ 1/P, 2/11, 1/1, 1/1, 3/floatmode, 5/registerWe, 1/0, 2/11, 3/001, 1/1, 4/0001, 1/--, 1/0, 5/registerZe, 1/1 }


\paragraph{Syntax} \texttt{floatuhq}\emph{floatmode} \emph{registerWe} = \emph{registerZe}

\paragraph{Description} The \emph{registerZe} is interpreted as four 16-bit unsigned integers packed into 64 bits. The \emph{registerZe} half words are converted to binary 16 floating-point numbers according to the rounding mode. The four resulting half words are packed and stored into the \emph{registerWe}. This instruction may raise exception bits in the CS register.


\paragraph{Modifier(s)}
\noindent\begin{itemize}
\item floatmode (cf. floatmode in section \ref{par:modifier-floatmode} on page \pageref{par:modifier-floatmode})
\end{itemize}

\paragraph{Execution} {Reservation table \textsf{ALU\_LITE} (Section~\ref{sec:reservation-ALU-LITE})}
~
{\begin{lstlisting}
stage ID:
  new floatmode = _ZX_3(floatmode);
stage RR:
  new argument2 = GPR[registerZe<<1];
  new RM = decode_riscv_float_rounding_mode(floatmode);
  for (I = 0; I < 8; I += 1) {
    result1.16[I] = ui16_to_f16(RM, argument2.16[I])
  };
stage E4:
  GPR[registerWe<<1] = result1;
  commit_float_exception_flags();
\end{lstlisting}}
\newpage

\paragraph{Encoding} Format \textsf{ALU\_FHQFLOATO} (Section~\ref{sec:format-ALU-FHQFLOATO}), \verb|fpucode6 = 0001|\\

\bitfields{ 1/P, 2/11, 1/1, 1/1, 3/floatmode, 5/registerWo, 1/1, 2/11, 3/001, 1/1, 4/0001, 1/--, 1/0, 5/registerZo, 1/1 }


\paragraph{Syntax} \texttt{floatuhq}\emph{floatmode} \emph{registerWo} = \emph{registerZo}

\paragraph{Description} The \emph{registerZo} is interpreted as four 16-bit unsigned integers packed into 64 bits. The \emph{registerZo} half words are converted to binary 16 floating-point numbers according to the rounding mode. The four resulting half words are packed and stored into the \emph{registerWo}. This instruction may raise exception bits in the CS register.


\paragraph{Modifier(s)}
\noindent\begin{itemize}
\item floatmode (cf. floatmode in section \ref{par:modifier-floatmode} on page \pageref{par:modifier-floatmode})
\end{itemize}

\paragraph{Execution} {Reservation table \textsf{ALU\_LITE} (Section~\ref{sec:reservation-ALU-LITE})}
~
{\begin{lstlisting}
stage ID:
  new floatmode = _ZX_3(floatmode);
stage RR:
  new argument2 = GPR[(registerZo<<1)+1];
  new RM = decode_riscv_float_rounding_mode(floatmode);
  for (I = 0; I < 8; I += 1) {
    result1.16[I] = ui16_to_f16(RM, argument2.16[I])
  };
stage E4:
  GPR[(registerWo<<1)+1] = result1;
  commit_float_exception_flags();
\end{lstlisting}}
\newpage
\subsubsection[{\small FLOATUW: Floating Point Conversion from Unsigned Fixed-Point}]{FLOATUW\hfill Floating Point Conversion from Unsigned Fixed-Point}
\label{instr:FLOATUW}\index{floatuw}
 
\paragraph{Encoding} Format \textsf{ALU\_FWFLOAT} (Section~\ref{sec:format-ALU-FWFLOAT}), \verb|fpucode6 = 0001|\\

\bitfields{ 1/P, 2/11, 1/1, 1/1, 3/floatmode, 6/registerW, 2/01, 3/001, 1/0, 4/0001, 1/--, 1/0, 6/registerZ }


\paragraph{Syntax} \texttt{floatuw}\emph{floatmode} \emph{registerW} = \emph{registerZ}

\paragraph{Description} The \emph{registerZ} interpreted as a 32-bit unsigned integer is converted to a 32-bit floating-point number according to the rounding mode. This instruction may raise exception bits in the CS register.


\paragraph{Modifier(s)}
\noindent\begin{itemize}
\item floatmode (cf. floatmode in section \ref{par:modifier-floatmode} on page \pageref{par:modifier-floatmode})
\end{itemize}

\paragraph{Execution} {Reservation table \textsf{ALU\_LITE} (Section~\ref{sec:reservation-ALU-LITE})}
~
{\begin{lstlisting}
stage ID:
  new floatmode = _ZX_3(floatmode);
stage RR:
  new argument2 = GPR[registerZ];
  new RM = decode_riscv_float_rounding_mode(floatmode);
  new result1 = ui32_to_f32(RM, argument2);
stage E4:
  GPR[registerW] = result1;
  commit_float_exception_flags();
\end{lstlisting}}
\newpage
\subsubsection[{\small FLOATUWD: Floating Point Conversion from Unsigned Fixed-Point Word to Double Word}]{FLOATUWD\hfill Floating Point Conversion from Unsigned Fixed-Point Word to Double Word}
\label{instr:FLOATUWD}\index{floatuwd}
 
\paragraph{Encoding} Format \textsf{ALU\_FDFLOAT} (Section~\ref{sec:format-ALU-FDFLOAT}), \verb|fpucode6 = 0011|\\

\bitfields{ 1/P, 2/11, 1/1, 1/1, 3/floatmode, 6/registerW, 2/00, 3/001, 1/0, 4/0011, 1/--, 1/0, 6/registerZ }


\paragraph{Syntax} \texttt{floatuwd}\emph{floatmode} \emph{registerW} = \emph{registerZ}

\paragraph{Description} The \emph{registerZ} interpreted as a 32-bit unsigned integer is converted to a 64-bit floating-point number according to the rounding mode. This instruction may raise exception bits in the CS register.


\paragraph{Modifier(s)}
\noindent\begin{itemize}
\item floatmode (cf. floatmode in section \ref{par:modifier-floatmode} on page \pageref{par:modifier-floatmode})
\end{itemize}

\paragraph{Execution} {Reservation table \textsf{ALU\_LITE} (Section~\ref{sec:reservation-ALU-LITE})}
~
{\begin{lstlisting}
stage ID:
  new floatmode = _ZX_3(floatmode);
stage RR:
  new argument2 = GPR[registerZ];
  new RM = decode_riscv_float_rounding_mode(floatmode);
  new result1 = ui32_to_f64(RM, argument2);
stage E4:
  GPR[registerW] = result1;
  commit_float_exception_flags();
\end{lstlisting}}
\newpage
\subsubsection[{\small FLOATUWP: Floating Point Conversion from Unsigned Fixed-Point Word Pair}]{FLOATUWP\hfill Floating Point Conversion from Unsigned Fixed-Point Word Pair}
\label{instr:FLOATUWP}\index{floatuwp}
 
\paragraph{Encoding} Format \textsf{ALU\_FWPFLOAT} (Section~\ref{sec:format-ALU-FWPFLOAT}), \verb|fpucode6 = 0001|\\

\bitfields{ 1/P, 2/11, 1/1, 1/1, 3/floatmode, 6/registerW, 2/00, 3/001, 1/0, 4/0001, 1/--, 1/1, 6/registerZ }


\paragraph{Syntax} \texttt{floatuwp}\emph{floatmode} \emph{registerW} = \emph{registerZ}

\paragraph{Description} The \emph{registerZ} is interpreted as two 32-bit unsigned integers packed into 64 bits. The \emph{registerZ} words are converted to binary 32 floating-point numbers according to the rounding mode. The two resulting words are packed and stored into the \emph{registerW}. This instruction may raise exception bits in the CS register.


\paragraph{Modifier(s)}
\noindent\begin{itemize}
\item floatmode (cf. floatmode in section \ref{par:modifier-floatmode} on page \pageref{par:modifier-floatmode})
\end{itemize}

\paragraph{Execution} {Reservation table \textsf{ALU\_LITE} (Section~\ref{sec:reservation-ALU-LITE})}
~
{\begin{lstlisting}
stage ID:
  new floatmode = _ZX_3(floatmode);
stage RR:
  new argument2 = GPR[registerZ];
  new RM = decode_riscv_float_rounding_mode(floatmode);
  for (I = 0; I < 2; I += 1) {
    result1.32[I] = ui32_to_f32(RM, argument2.32[I])
  };
stage E4:
  GPR[registerW] = result1;
  commit_float_exception_flags();
\end{lstlisting}}
\newpage
\subsubsection[{\small FLOATW: Floating Point Conversion from Fixed-Point}]{FLOATW\hfill Floating Point Conversion from Fixed-Point}
\label{instr:FLOATW}\index{floatw}
 
\paragraph{Encoding} Format \textsf{ALU\_FWFLOAT} (Section~\ref{sec:format-ALU-FWFLOAT}), \verb|fpucode6 = 0000|\\

\bitfields{ 1/P, 2/11, 1/1, 1/1, 3/floatmode, 6/registerW, 2/01, 3/001, 1/0, 4/0000, 1/--, 1/0, 6/registerZ }


\paragraph{Syntax} \texttt{floatw}\emph{floatmode} \emph{registerW} = \emph{registerZ}

\paragraph{Description} The \emph{registerZ} interpreted as a 32-bit signed integer is converted to a 32-bit floating-point number according to the rounding mode. This instruction may raise exception bits in the CS register.


\paragraph{Modifier(s)}
\noindent\begin{itemize}
\item floatmode (cf. floatmode in section \ref{par:modifier-floatmode} on page \pageref{par:modifier-floatmode})
\end{itemize}

\paragraph{Execution} {Reservation table \textsf{ALU\_LITE} (Section~\ref{sec:reservation-ALU-LITE})}
~
{\begin{lstlisting}
stage ID:
  new floatmode = _ZX_3(floatmode);
stage RR:
  new argument2 = GPR[registerZ];
  new RM = decode_riscv_float_rounding_mode(floatmode);
  new result1 = i32_to_f32(RM, argument2);
stage E4:
  GPR[registerW] = result1;
  commit_float_exception_flags();
\end{lstlisting}}
\newpage
\subsubsection[{\small FLOATWD: Floating Point Conversion from Fixed-Point Word to Double Word}]{FLOATWD\hfill Floating Point Conversion from Fixed-Point Word to Double Word}
\label{instr:FLOATWD}\index{floatwd}
 
\paragraph{Encoding} Format \textsf{ALU\_FDFLOAT} (Section~\ref{sec:format-ALU-FDFLOAT}), \verb|fpucode6 = 0010|\\

\bitfields{ 1/P, 2/11, 1/1, 1/1, 3/floatmode, 6/registerW, 2/00, 3/001, 1/0, 4/0010, 1/--, 1/0, 6/registerZ }


\paragraph{Syntax} \texttt{floatwd}\emph{floatmode} \emph{registerW} = \emph{registerZ}

\paragraph{Description} The \emph{registerZ} interpreted as a 32-bit signed integer is converted to a 64-bit floating-point number according to the rounding mode. This instruction may raise exception bits in the CS register.


\paragraph{Modifier(s)}
\noindent\begin{itemize}
\item floatmode (cf. floatmode in section \ref{par:modifier-floatmode} on page \pageref{par:modifier-floatmode})
\end{itemize}

\paragraph{Execution} {Reservation table \textsf{ALU\_LITE} (Section~\ref{sec:reservation-ALU-LITE})}
~
{\begin{lstlisting}
stage ID:
  new floatmode = _ZX_3(floatmode);
stage RR:
  new argument2 = GPR[registerZ];
  new RM = decode_riscv_float_rounding_mode(floatmode);
  new result1 = i32_to_f64(RM, argument2);
stage E4:
  GPR[registerW] = result1;
  commit_float_exception_flags();
\end{lstlisting}}
\newpage
\subsubsection[{\small FLOATWP: Floating Point Conversion from Fixed-Point Word Pair}]{FLOATWP\hfill Floating Point Conversion from Fixed-Point Word Pair}
\label{instr:FLOATWP}\index{floatwp}
 
\paragraph{Encoding} Format \textsf{ALU\_FWPFLOAT} (Section~\ref{sec:format-ALU-FWPFLOAT}), \verb|fpucode6 = 0000|\\

\bitfields{ 1/P, 2/11, 1/1, 1/1, 3/floatmode, 6/registerW, 2/00, 3/001, 1/0, 4/0000, 1/--, 1/1, 6/registerZ }


\paragraph{Syntax} \texttt{floatwp}\emph{floatmode} \emph{registerW} = \emph{registerZ}

\paragraph{Description} The \emph{registerZ} is interpreted as two 32-bit signed integers packed into 64 bits. The \emph{registerZ} words are converted to binary 32 floating-point numbers according to the rounding mode. The two resulting words are packed and stored into the \emph{registerW}. This instruction may raise exception bits in the CS register.


\paragraph{Modifier(s)}
\noindent\begin{itemize}
\item floatmode (cf. floatmode in section \ref{par:modifier-floatmode} on page \pageref{par:modifier-floatmode})
\end{itemize}

\paragraph{Execution} {Reservation table \textsf{ALU\_LITE} (Section~\ref{sec:reservation-ALU-LITE})}
~
{\begin{lstlisting}
stage ID:
  new floatmode = _ZX_3(floatmode);
stage RR:
  new argument2 = GPR[registerZ];
  new RM = decode_riscv_float_rounding_mode(floatmode);
  for (I = 0; I < 2; I += 1) {
    result1.32[I] = i32_to_f32(RM, argument2.32[I])
  };
stage E4:
  GPR[registerW] = result1;
  commit_float_exception_flags();
\end{lstlisting}}
\newpage
\subsubsection[{\small FMAXD: Floating-Point Maximum Double Word}]{FMAXD\hfill Floating-Point Maximum Double Word}
\label{instr:FMAXD}\index{fmaxd}
 
\paragraph{Encoding} Format \textsf{ALU\_FDMMWR} (Section~\ref{sec:format-ALU-FDMMWR}), \verb|floatmode = 001|\\

\bitfields{ 1/P, 2/11, 1/1, 1/1, 3/001, 6/registerW, 2/00, 3/000, 1/0, 6/registerY, 6/registerZ }


\paragraph{Syntax} \texttt{fmaxd} \emph{registerW} = \emph{registerZ}, \emph{registerY}

\paragraph{Description} The \emph{registerZ} is compared to the \emph{registerY}, assuming binary 64 floating-point numbers. The maximum value is stored into the \emph{registerW}.


\paragraph{Execution} {Reservation table \textsf{ALU\_LITE} (Section~\ref{sec:reservation-ALU-LITE})}
~
{\begin{lstlisting}
stage RR:
  new argument3 = GPR[registerY];
  new argument2 = GPR[registerZ];
  new result1 = f64_max(argument2, argument3);
stage E1:
  GPR[registerW] = result1;
\end{lstlisting}}
\newpage
\subsubsection[{\small FMAXH: Floating-Point Maximum Half Word}]{FMAXH\hfill Floating-Point Maximum Half Word}
\label{instr:FMAXH}\index{fmaxh}
 
\paragraph{Encoding} Format \textsf{ALU\_FHMMWR} (Section~\ref{sec:format-ALU-FHMMWR}), \verb|floatmode = 001|\\

\bitfields{ 1/P, 2/11, 1/1, 1/1, 3/001, 6/registerW, 2/01, 3/011, 1/1, 6/registerY, 6/registerZ }


\paragraph{Syntax} \texttt{fmaxh} \emph{registerW} = \emph{registerZ}, \emph{registerY}

\paragraph{Description} The \emph{registerZ} is compared to the \emph{registerY}, assuming binary 16 floating-point numbers. The maximum value is stored into the \emph{registerW}.


\paragraph{Execution} {Reservation table \textsf{ALU\_LITE} (Section~\ref{sec:reservation-ALU-LITE})}
~
{\begin{lstlisting}
stage RR:
  new argument3 = GPR[registerY];
  new argument2 = GPR[registerZ];
  new result1 = f16_max(argument2, argument3);
stage E1:
  GPR[registerW] = result1;
\end{lstlisting}}
\newpage
\subsubsection[{\small FMAXHQ: Floating-Point Maximum Half Word Quadruple}]{FMAXHQ\hfill Floating-Point Maximum Half Word Quadruple}
\label{instr:FMAXHQ}\index{fmaxhq}
 
\paragraph{Encoding} Format \textsf{ALU\_FHQMMWRE} (Section~\ref{sec:format-ALU-FHQMMWRE}), \verb|floatmode = 001|\\

\bitfields{ 1/P, 2/11, 1/1, 1/1, 3/001, 5/registerWe, 1/0, 2/11, 3/000, 1/1, 5/registerYe, 1/0, 5/registerZe, 1/1 }


\paragraph{Syntax} \texttt{fmaxhq} \emph{registerWe} = \emph{registerZe}, \emph{registerYe}

\paragraph{Description} The \emph{registerZe} and the \emph{registerYe} are interpreted as four binary 16 floating-point numbers packed into 64 bits. The \emph{registerZe} half words are compared to the \emph{registerYe} half words. The four maximum half words are packed and stored into the \emph{registerWe}.


\paragraph{Execution} {Reservation table \textsf{ALU\_LITE} (Section~\ref{sec:reservation-ALU-LITE})}
~
{\begin{lstlisting}
stage RR:
  new argument3 = GPR[registerYe<<1];
  new argument2 = GPR[registerZe<<1];
  for (I = 0; I < 4; I += 1) {
    result1.16[I] = f16_max(argument2.16[I], argument3.16[I])
  };
stage E1:
  GPR[registerWe<<1] = result1;
\end{lstlisting}}
\newpage

\paragraph{Encoding} Format \textsf{ALU\_FHQMMWRO} (Section~\ref{sec:format-ALU-FHQMMWRO}), \verb|floatmode = 001|\\

\bitfields{ 1/P, 2/11, 1/1, 1/1, 3/001, 5/registerWo, 1/1, 2/11, 3/000, 1/1, 5/registerYo, 1/0, 5/registerZo, 1/1 }


\paragraph{Syntax} \texttt{fmaxhq} \emph{registerWo} = \emph{registerZo}, \emph{registerYo}

\paragraph{Description} The \emph{registerZo} and the \emph{registerYo} are interpreted as four binary 16 floating-point numbers packed into 64 bits. The \emph{registerZo} half words are compared to the \emph{registerYo} half words. The four maximum half words are packed and stored into the \emph{registerWo}.


\paragraph{Execution} {Reservation table \textsf{ALU\_LITE} (Section~\ref{sec:reservation-ALU-LITE})}
~
{\begin{lstlisting}
stage RR:
  new argument3 = GPR[(registerYo<<1)+1];
  new argument2 = GPR[(registerZo<<1)+1];
  for (I = 0; I < 4; I += 1) {
    result1.16[I] = f16_max(argument2.16[I], argument3.16[I])
  };
stage E1:
  GPR[(registerWo<<1)+1] = result1;
\end{lstlisting}}
\newpage
\subsubsection[{\small FMAXND: Floating-Point Maximum Number Double Word}]{FMAXND\hfill Floating-Point Maximum Number Double Word}
\label{instr:FMAXND}\index{fmaxnd}
 
\paragraph{Encoding} Format \textsf{ALU\_FDMMWR} (Section~\ref{sec:format-ALU-FDMMWR}), \verb|floatmode = 011|\\

\bitfields{ 1/P, 2/11, 1/1, 1/1, 3/011, 6/registerW, 2/00, 3/000, 1/0, 6/registerY, 6/registerZ }


\paragraph{Syntax} \texttt{fmaxnd} \emph{registerW} = \emph{registerZ}, \emph{registerY}

\paragraph{Description} The \emph{registerZ} is compared to the \emph{registerY}, assuming binary 64 floating-point numbers. The maximum number value is stored into the \emph{registerW}.


\paragraph{Execution} {Reservation table \textsf{ALU\_LITE} (Section~\ref{sec:reservation-ALU-LITE})}
~
{\begin{lstlisting}
stage RR:
  new argument3 = GPR[registerY];
  new argument2 = GPR[registerZ];
  new result1 = f64_maxNum(argument2, argument3);
stage E1:
  GPR[registerW] = result1;
\end{lstlisting}}
\newpage
\subsubsection[{\small FMAXNH: Floating-Point Maximum Number Half Word}]{FMAXNH\hfill Floating-Point Maximum Number Half Word}
\label{instr:FMAXNH}\index{fmaxnh}
 
\paragraph{Encoding} Format \textsf{ALU\_FHMMWR} (Section~\ref{sec:format-ALU-FHMMWR}), \verb|floatmode = 011|\\

\bitfields{ 1/P, 2/11, 1/1, 1/1, 3/011, 6/registerW, 2/01, 3/011, 1/1, 6/registerY, 6/registerZ }


\paragraph{Syntax} \texttt{fmaxnh} \emph{registerW} = \emph{registerZ}, \emph{registerY}

\paragraph{Description} The \emph{registerZ} is compared to the \emph{registerY}, assuming binary 16 floating-point numbers. The maximum number value is stored into the \emph{registerW}.


\paragraph{Execution} {Reservation table \textsf{ALU\_LITE} (Section~\ref{sec:reservation-ALU-LITE})}
~
{\begin{lstlisting}
stage RR:
  new argument3 = GPR[registerY];
  new argument2 = GPR[registerZ];
  new result1 = f16_maxNum(argument2, argument3);
stage E1:
  GPR[registerW] = result1;
\end{lstlisting}}
\newpage
\subsubsection[{\small FMAXNHQ: Floating-Point Maximum Number Half Word Quadruple}]{FMAXNHQ\hfill Floating-Point Maximum Number Half Word Quadruple}
\label{instr:FMAXNHQ}\index{fmaxnhq}
 
\paragraph{Encoding} Format \textsf{ALU\_FHQMMWRE} (Section~\ref{sec:format-ALU-FHQMMWRE}), \verb|floatmode = 011|\\

\bitfields{ 1/P, 2/11, 1/1, 1/1, 3/011, 5/registerWe, 1/0, 2/11, 3/000, 1/1, 5/registerYe, 1/0, 5/registerZe, 1/1 }


\paragraph{Syntax} \texttt{fmaxnhq} \emph{registerWe} = \emph{registerZe}, \emph{registerYe}

\paragraph{Description} The \emph{registerZe} and the \emph{registerYe} are interpreted as four binary 16 floating-point numbers packed into 64 bits. The \emph{registerZe} half words are compared to the \emph{registerYe} half words. The four maximum number half words are packed and stored into the \emph{registerWe}.


\paragraph{Execution} {Reservation table \textsf{ALU\_LITE} (Section~\ref{sec:reservation-ALU-LITE})}
~
{\begin{lstlisting}
stage RR:
  new argument3 = GPR[registerYe<<1];
  new argument2 = GPR[registerZe<<1];
  for (I = 0; I < 4; I += 1) {
    result1.16[I] = f16_maxNum(argument2.16[I], argument3.16[I])
  };
stage E1:
  GPR[registerWe<<1] = result1;
\end{lstlisting}}
\newpage

\paragraph{Encoding} Format \textsf{ALU\_FHQMMWRO} (Section~\ref{sec:format-ALU-FHQMMWRO}), \verb|floatmode = 011|\\

\bitfields{ 1/P, 2/11, 1/1, 1/1, 3/011, 5/registerWo, 1/1, 2/11, 3/000, 1/1, 5/registerYo, 1/0, 5/registerZo, 1/1 }


\paragraph{Syntax} \texttt{fmaxnhq} \emph{registerWo} = \emph{registerZo}, \emph{registerYo}

\paragraph{Description} The \emph{registerZo} and the \emph{registerYo} are interpreted as four binary 16 floating-point numbers packed into 64 bits. The \emph{registerZo} half words are compared to the \emph{registerYo} half words. The four maximum number half words are packed and stored into the \emph{registerWo}.


\paragraph{Execution} {Reservation table \textsf{ALU\_LITE} (Section~\ref{sec:reservation-ALU-LITE})}
~
{\begin{lstlisting}
stage RR:
  new argument3 = GPR[(registerYo<<1)+1];
  new argument2 = GPR[(registerZo<<1)+1];
  for (I = 0; I < 4; I += 1) {
    result1.16[I] = f16_maxNum(argument2.16[I], argument3.16[I])
  };
stage E1:
  GPR[(registerWo<<1)+1] = result1;
\end{lstlisting}}
\newpage
\subsubsection[{\small FMAXNW: Floating-Point Maximum Number Word}]{FMAXNW\hfill Floating-Point Maximum Number Word}
\label{instr:FMAXNW}\index{fmaxnw}
 
\paragraph{Encoding} Format \textsf{ALU\_FWMMWR} (Section~\ref{sec:format-ALU-FWMMWR}), \verb|floatmode = 011|\\

\bitfields{ 1/P, 2/11, 1/1, 1/1, 3/011, 6/registerW, 2/01, 3/000, 1/0, 6/registerY, 6/registerZ }


\paragraph{Syntax} \texttt{fmaxnw} \emph{registerW} = \emph{registerZ}, \emph{registerY}

\paragraph{Description} The \emph{registerZ} is compared to the \emph{registerY}, assuming binary 32 floating-point numbers. The maximum number value is stored into the \emph{registerW}.


\paragraph{Execution} {Reservation table \textsf{ALU\_LITE} (Section~\ref{sec:reservation-ALU-LITE})}
~
{\begin{lstlisting}
stage RR:
  new argument3 = GPR[registerY];
  new argument2 = GPR[registerZ];
  new result1 = f32_maxNum(argument2, argument3);
stage E1:
  GPR[registerW] = result1;
\end{lstlisting}}
\newpage
\subsubsection[{\small FMAXNWP: Floating-Point Maximum Number Word Pair}]{FMAXNWP\hfill Floating-Point Maximum Number Word Pair}
\label{instr:FMAXNWP}\index{fmaxnwp}
 
\paragraph{Encoding} Format \textsf{ALU\_FWPMMWRE} (Section~\ref{sec:format-ALU-FWPMMWRE}), \verb|floatmode = 011|\\

\bitfields{ 1/P, 2/11, 1/1, 1/1, 3/011, 5/registerWe, 1/0, 2/11, 3/000, 1/1, 5/registerYe, 1/0, 5/registerZe, 1/0 }


\paragraph{Syntax} \texttt{fmaxnwp} \emph{registerWe} = \emph{registerZe}, \emph{registerYe}

\paragraph{Description} The \emph{registerZe} and the \emph{registerYe} are interpreted as two binary 32 floating-point numbers packed into 64 bits. The \emph{registerZe} words are compared to the \emph{registerYe} words. The two maximum number words are packed and stored into the \emph{registerWe}.


\paragraph{Execution} {Reservation table \textsf{ALU\_LITE} (Section~\ref{sec:reservation-ALU-LITE})}
~
{\begin{lstlisting}
stage RR:
  new argument3 = GPR[registerYe<<1];
  new argument2 = GPR[registerZe<<1];
  for (I = 0; I < 2; I += 1) {
    result1.32[I] = f32_maxNum(argument2.32[I], argument3.32[I])
  };
stage E1:
  GPR[registerWe<<1] = result1;
\end{lstlisting}}
\newpage

\paragraph{Encoding} Format \textsf{ALU\_FWPMMWRO} (Section~\ref{sec:format-ALU-FWPMMWRO}), \verb|floatmode = 011|\\

\bitfields{ 1/P, 2/11, 1/1, 1/1, 3/011, 5/registerWo, 1/1, 2/11, 3/000, 1/1, 5/registerYo, 1/0, 5/registerZo, 1/0 }


\paragraph{Syntax} \texttt{fmaxnwp} \emph{registerWo} = \emph{registerZo}, \emph{registerYo}

\paragraph{Description} The \emph{registerZo} and the \emph{registerYo} are interpreted as two binary 32 floating-point numbers packed into 64 bits. The \emph{registerZo} words are compared to the \emph{registerYo} words. The two maximum number words are packed and stored into the \emph{registerWo}.


\paragraph{Execution} {Reservation table \textsf{ALU\_LITE} (Section~\ref{sec:reservation-ALU-LITE})}
~
{\begin{lstlisting}
stage RR:
  new argument3 = GPR[(registerYo<<1)+1];
  new argument2 = GPR[(registerZo<<1)+1];
  for (I = 0; I < 2; I += 1) {
    result1.32[I] = f32_maxNum(argument2.32[I], argument3.32[I])
  };
stage E1:
  GPR[(registerWo<<1)+1] = result1;
\end{lstlisting}}
\newpage
\subsubsection[{\small FMAXW: Floating-Point Maximum Word}]{FMAXW\hfill Floating-Point Maximum Word}
\label{instr:FMAXW}\index{fmaxw}
 
\paragraph{Encoding} Format \textsf{ALU\_FWMMWR} (Section~\ref{sec:format-ALU-FWMMWR}), \verb|floatmode = 001|\\

\bitfields{ 1/P, 2/11, 1/1, 1/1, 3/001, 6/registerW, 2/01, 3/000, 1/0, 6/registerY, 6/registerZ }


\paragraph{Syntax} \texttt{fmaxw} \emph{registerW} = \emph{registerZ}, \emph{registerY}

\paragraph{Description} The \emph{registerZ} is compared to the \emph{registerY}, assuming binary 32 floating-point numbers. The maximum value is stored into the \emph{registerW}.


\paragraph{Execution} {Reservation table \textsf{ALU\_LITE} (Section~\ref{sec:reservation-ALU-LITE})}
~
{\begin{lstlisting}
stage RR:
  new argument3 = GPR[registerY];
  new argument2 = GPR[registerZ];
  new result1 = f32_max(argument2, argument3);
stage E1:
  GPR[registerW] = result1;
\end{lstlisting}}
\newpage
\subsubsection[{\small FMAXWP: Floating-Point Maximum Word Pair}]{FMAXWP\hfill Floating-Point Maximum Word Pair}
\label{instr:FMAXWP}\index{fmaxwp}
 
\paragraph{Encoding} Format \textsf{ALU\_FWPMMWRE} (Section~\ref{sec:format-ALU-FWPMMWRE}), \verb|floatmode = 001|\\

\bitfields{ 1/P, 2/11, 1/1, 1/1, 3/001, 5/registerWe, 1/0, 2/11, 3/000, 1/1, 5/registerYe, 1/0, 5/registerZe, 1/0 }


\paragraph{Syntax} \texttt{fmaxwp} \emph{registerWe} = \emph{registerZe}, \emph{registerYe}

\paragraph{Description} The \emph{registerZe} and the \emph{registerYe} are interpreted as two binary 32 floating-point numbers packed into 64 bits. The \emph{registerZe} words are compared to the \emph{registerYe} words. The two maximum words are packed and stored into the \emph{registerWe}.


\paragraph{Execution} {Reservation table \textsf{ALU\_LITE} (Section~\ref{sec:reservation-ALU-LITE})}
~
{\begin{lstlisting}
stage RR:
  new argument3 = GPR[registerYe<<1];
  new argument2 = GPR[registerZe<<1];
  for (I = 0; I < 2; I += 1) {
    result1.32[I] = f32_max(argument2.32[I], argument3.32[I])
  };
stage E1:
  GPR[registerWe<<1] = result1;
\end{lstlisting}}
\newpage

\paragraph{Encoding} Format \textsf{ALU\_FWPMMWRO} (Section~\ref{sec:format-ALU-FWPMMWRO}), \verb|floatmode = 001|\\

\bitfields{ 1/P, 2/11, 1/1, 1/1, 3/001, 5/registerWo, 1/1, 2/11, 3/000, 1/1, 5/registerYo, 1/0, 5/registerZo, 1/0 }


\paragraph{Syntax} \texttt{fmaxwp} \emph{registerWo} = \emph{registerZo}, \emph{registerYo}

\paragraph{Description} The \emph{registerZo} and the \emph{registerYo} are interpreted as two binary 32 floating-point numbers packed into 64 bits. The \emph{registerZo} words are compared to the \emph{registerYo} words. The two maximum words are packed and stored into the \emph{registerWo}.


\paragraph{Execution} {Reservation table \textsf{ALU\_LITE} (Section~\ref{sec:reservation-ALU-LITE})}
~
{\begin{lstlisting}
stage RR:
  new argument3 = GPR[(registerYo<<1)+1];
  new argument2 = GPR[(registerZo<<1)+1];
  for (I = 0; I < 2; I += 1) {
    result1.32[I] = f32_max(argument2.32[I], argument3.32[I])
  };
stage E1:
  GPR[(registerWo<<1)+1] = result1;
\end{lstlisting}}
\newpage
\subsubsection[{\small FMIND: Floating-Point Minimum Double Word}]{FMIND\hfill Floating-Point Minimum Double Word}
\label{instr:FMIND}\index{fmind}
 
\paragraph{Encoding} Format \textsf{ALU\_FDMMWR} (Section~\ref{sec:format-ALU-FDMMWR}), \verb|floatmode = 000|\\

\bitfields{ 1/P, 2/11, 1/1, 1/1, 3/000, 6/registerW, 2/00, 3/000, 1/0, 6/registerY, 6/registerZ }


\paragraph{Syntax} \texttt{fmind} \emph{registerW} = \emph{registerZ}, \emph{registerY}

\paragraph{Description} The \emph{registerZ} is compared to the \emph{registerY}, assuming binary 64 floating-point numbers. The minimum is stored into the \emph{registerW}.


\paragraph{Execution} {Reservation table \textsf{ALU\_LITE} (Section~\ref{sec:reservation-ALU-LITE})}
~
{\begin{lstlisting}
stage RR:
  new argument3 = GPR[registerY];
  new argument2 = GPR[registerZ];
  new result1 = f64_min(argument2, argument3);
stage E1:
  GPR[registerW] = result1;
\end{lstlisting}}
\newpage
\subsubsection[{\small FMINH: Floating-Point Minimum Half Word}]{FMINH\hfill Floating-Point Minimum Half Word}
\label{instr:FMINH}\index{fminh}
 
\paragraph{Encoding} Format \textsf{ALU\_FHMMWR} (Section~\ref{sec:format-ALU-FHMMWR}), \verb|floatmode = 000|\\

\bitfields{ 1/P, 2/11, 1/1, 1/1, 3/000, 6/registerW, 2/01, 3/011, 1/1, 6/registerY, 6/registerZ }


\paragraph{Syntax} \texttt{fminh} \emph{registerW} = \emph{registerZ}, \emph{registerY}

\paragraph{Description} The \emph{registerZ} is compared to the \emph{registerY}, assuming binary 16 floating-point numbers. The minimum value is stored into the \emph{registerW}.


\paragraph{Execution} {Reservation table \textsf{ALU\_LITE} (Section~\ref{sec:reservation-ALU-LITE})}
~
{\begin{lstlisting}
stage RR:
  new argument3 = GPR[registerY];
  new argument2 = GPR[registerZ];
  new result1 = f16_min(argument2, argument3);
stage E1:
  GPR[registerW] = result1;
\end{lstlisting}}
\newpage
\subsubsection[{\small FMINHQ: Floating-Point Minimum Half Word Quadruple}]{FMINHQ\hfill Floating-Point Minimum Half Word Quadruple}
\label{instr:FMINHQ}\index{fminhq}
 
\paragraph{Encoding} Format \textsf{ALU\_FHQMMWRE} (Section~\ref{sec:format-ALU-FHQMMWRE}), \verb|floatmode = 000|\\

\bitfields{ 1/P, 2/11, 1/1, 1/1, 3/000, 5/registerWe, 1/0, 2/11, 3/000, 1/1, 5/registerYe, 1/0, 5/registerZe, 1/1 }


\paragraph{Syntax} \texttt{fminhq} \emph{registerWe} = \emph{registerZe}, \emph{registerYe}

\paragraph{Description} The \emph{registerZe} and the \emph{registerYe} are interpreted as four binary 16 floating-point numbers packed into 64 bits. The \emph{registerZe} half words are compared to the \emph{registerYe} half words. The four minimum half words are packed and stored into the \emph{registerWe}.


\paragraph{Execution} {Reservation table \textsf{ALU\_LITE} (Section~\ref{sec:reservation-ALU-LITE})}
~
{\begin{lstlisting}
stage RR:
  new argument3 = GPR[registerYe<<1];
  new argument2 = GPR[registerZe<<1];
  for (I = 0; I < 4; I += 1) {
    result1.16[I] = f16_min(argument2.16[I], argument3.16[I])
  };
stage E1:
  GPR[registerWe<<1] = result1;
\end{lstlisting}}
\newpage

\paragraph{Encoding} Format \textsf{ALU\_FHQMMWRO} (Section~\ref{sec:format-ALU-FHQMMWRO}), \verb|floatmode = 000|\\

\bitfields{ 1/P, 2/11, 1/1, 1/1, 3/000, 5/registerWo, 1/1, 2/11, 3/000, 1/1, 5/registerYo, 1/0, 5/registerZo, 1/1 }


\paragraph{Syntax} \texttt{fminhq} \emph{registerWo} = \emph{registerZo}, \emph{registerYo}

\paragraph{Description} The \emph{registerZo} and the \emph{registerYo} are interpreted as four binary 16 floating-point numbers packed into 64 bits. The \emph{registerZo} half words are compared to the \emph{registerYo} half words. The four minimum half words are packed and stored into the \emph{registerWo}.


\paragraph{Execution} {Reservation table \textsf{ALU\_LITE} (Section~\ref{sec:reservation-ALU-LITE})}
~
{\begin{lstlisting}
stage RR:
  new argument3 = GPR[(registerYo<<1)+1];
  new argument2 = GPR[(registerZo<<1)+1];
  for (I = 0; I < 4; I += 1) {
    result1.16[I] = f16_min(argument2.16[I], argument3.16[I])
  };
stage E1:
  GPR[(registerWo<<1)+1] = result1;
\end{lstlisting}}
\newpage
\subsubsection[{\small FMINND: Floating-Point Minimum Number Double Word}]{FMINND\hfill Floating-Point Minimum Number Double Word}
\label{instr:FMINND}\index{fminnd}
 
\paragraph{Encoding} Format \textsf{ALU\_FDMMWR} (Section~\ref{sec:format-ALU-FDMMWR}), \verb|floatmode = 010|\\

\bitfields{ 1/P, 2/11, 1/1, 1/1, 3/010, 6/registerW, 2/00, 3/000, 1/0, 6/registerY, 6/registerZ }


\paragraph{Syntax} \texttt{fminnd} \emph{registerW} = \emph{registerZ}, \emph{registerY}

\paragraph{Description} The \emph{registerZ} is compared to the \emph{registerY}, assuming binary 64 floating-point numbers. The minimum value is stored into the \emph{registerW}.


\paragraph{Execution} {Reservation table \textsf{ALU\_LITE} (Section~\ref{sec:reservation-ALU-LITE})}
~
{\begin{lstlisting}
stage RR:
  new argument3 = GPR[registerY];
  new argument2 = GPR[registerZ];
  new result1 = f64_minNum(argument2, argument3);
stage E1:
  GPR[registerW] = result1;
\end{lstlisting}}
\newpage
\subsubsection[{\small FMINNH: Floating-Point Minimum Number Half Word}]{FMINNH\hfill Floating-Point Minimum Number Half Word}
\label{instr:FMINNH}\index{fminnh}
 
\paragraph{Encoding} Format \textsf{ALU\_FHMMWR} (Section~\ref{sec:format-ALU-FHMMWR}), \verb|floatmode = 010|\\

\bitfields{ 1/P, 2/11, 1/1, 1/1, 3/010, 6/registerW, 2/01, 3/011, 1/1, 6/registerY, 6/registerZ }


\paragraph{Syntax} \texttt{fminnh} \emph{registerW} = \emph{registerZ}, \emph{registerY}

\paragraph{Description} The \emph{registerZ} is compared to the \emph{registerY}, assuming binary 16 floating-point numbers. The minimum number value is stored into the \emph{registerW}.


\paragraph{Execution} {Reservation table \textsf{ALU\_LITE} (Section~\ref{sec:reservation-ALU-LITE})}
~
{\begin{lstlisting}
stage RR:
  new argument3 = GPR[registerY];
  new argument2 = GPR[registerZ];
  new result1 = f16_minNum(argument2, argument3);
stage E1:
  GPR[registerW] = result1;
\end{lstlisting}}
\newpage
\subsubsection[{\small FMINNHQ: Floating-Point Minimum Number Half Word Quadruple}]{FMINNHQ\hfill Floating-Point Minimum Number Half Word Quadruple}
\label{instr:FMINNHQ}\index{fminnhq}
 
\paragraph{Encoding} Format \textsf{ALU\_FHQMMWRE} (Section~\ref{sec:format-ALU-FHQMMWRE}), \verb|floatmode = 010|\\

\bitfields{ 1/P, 2/11, 1/1, 1/1, 3/010, 5/registerWe, 1/0, 2/11, 3/000, 1/1, 5/registerYe, 1/0, 5/registerZe, 1/1 }


\paragraph{Syntax} \texttt{fminnhq} \emph{registerWe} = \emph{registerZe}, \emph{registerYe}

\paragraph{Description} The \emph{registerZe} and the \emph{registerYe} are interpreted as four binary 16 floating-point numbers packed into 64 bits. The \emph{registerZe} half words are compared to the \emph{registerYe} half words. The four minimum number half words are packed and stored into the \emph{registerWe}.


\paragraph{Execution} {Reservation table \textsf{ALU\_LITE} (Section~\ref{sec:reservation-ALU-LITE})}
~
{\begin{lstlisting}
stage RR:
  new argument3 = GPR[registerYe<<1];
  new argument2 = GPR[registerZe<<1];
  for (I = 0; I < 4; I += 1) {
    result1.16[I] = f16_minNum(argument2.16[I], argument3.16[I])
  };
stage E1:
  GPR[registerWe<<1] = result1;
\end{lstlisting}}
\newpage

\paragraph{Encoding} Format \textsf{ALU\_FHQMMWRO} (Section~\ref{sec:format-ALU-FHQMMWRO}), \verb|floatmode = 010|\\

\bitfields{ 1/P, 2/11, 1/1, 1/1, 3/010, 5/registerWo, 1/1, 2/11, 3/000, 1/1, 5/registerYo, 1/0, 5/registerZo, 1/1 }


\paragraph{Syntax} \texttt{fminnhq} \emph{registerWo} = \emph{registerZo}, \emph{registerYo}

\paragraph{Description} The \emph{registerZo} and the \emph{registerYo} are interpreted as four binary 16 floating-point numbers packed into 64 bits. The \emph{registerZo} half words are compared to the \emph{registerYo} half words. The four minimum number half words are packed and stored into the \emph{registerWo}.


\paragraph{Execution} {Reservation table \textsf{ALU\_LITE} (Section~\ref{sec:reservation-ALU-LITE})}
~
{\begin{lstlisting}
stage RR:
  new argument3 = GPR[(registerYo<<1)+1];
  new argument2 = GPR[(registerZo<<1)+1];
  for (I = 0; I < 4; I += 1) {
    result1.16[I] = f16_minNum(argument2.16[I], argument3.16[I])
  };
stage E1:
  GPR[(registerWo<<1)+1] = result1;
\end{lstlisting}}
\newpage
\subsubsection[{\small FMINNW: Floating-Point Minimum Number Word}]{FMINNW\hfill Floating-Point Minimum Number Word}
\label{instr:FMINNW}\index{fminnw}
 
\paragraph{Encoding} Format \textsf{ALU\_FWMMWR} (Section~\ref{sec:format-ALU-FWMMWR}), \verb|floatmode = 010|\\

\bitfields{ 1/P, 2/11, 1/1, 1/1, 3/010, 6/registerW, 2/01, 3/000, 1/0, 6/registerY, 6/registerZ }


\paragraph{Syntax} \texttt{fminnw} \emph{registerW} = \emph{registerZ}, \emph{registerY}

\paragraph{Description} The \emph{registerZ} is compared to the \emph{registerY}, assuming binary 32 floating-point numbers. The minimum number value is stored into the \emph{registerW}.


\paragraph{Execution} {Reservation table \textsf{ALU\_LITE} (Section~\ref{sec:reservation-ALU-LITE})}
~
{\begin{lstlisting}
stage RR:
  new argument3 = GPR[registerY];
  new argument2 = GPR[registerZ];
  new result1 = f32_minNum(argument2, argument3);
stage E1:
  GPR[registerW] = result1;
\end{lstlisting}}
\newpage
\subsubsection[{\small FMINNWP: Floating-Point Minimum Number Word Pair}]{FMINNWP\hfill Floating-Point Minimum Number Word Pair}
\label{instr:FMINNWP}\index{fminnwp}
 
\paragraph{Encoding} Format \textsf{ALU\_FWPMMWRE} (Section~\ref{sec:format-ALU-FWPMMWRE}), \verb|floatmode = 010|\\

\bitfields{ 1/P, 2/11, 1/1, 1/1, 3/010, 5/registerWe, 1/0, 2/11, 3/000, 1/1, 5/registerYe, 1/0, 5/registerZe, 1/0 }


\paragraph{Syntax} \texttt{fminnwp} \emph{registerWe} = \emph{registerZe}, \emph{registerYe}

\paragraph{Description} The \emph{registerZe} and the \emph{registerYe} are interpreted as two binary 32 floating-point numbers packed into 64 bits. The \emph{registerZe} words are compared to the \emph{registerYe} words. The two minimum number words are packed and stored into the \emph{registerWe}.


\paragraph{Execution} {Reservation table \textsf{ALU\_LITE} (Section~\ref{sec:reservation-ALU-LITE})}
~
{\begin{lstlisting}
stage RR:
  new argument3 = GPR[registerYe<<1];
  new argument2 = GPR[registerZe<<1];
  for (I = 0; I < 2; I += 1) {
    result1.32[I] = f32_minNum(argument2.32[I], argument3.32[I])
  };
stage E1:
  GPR[registerWe<<1] = result1;
\end{lstlisting}}
\newpage

\paragraph{Encoding} Format \textsf{ALU\_FWPMMWRO} (Section~\ref{sec:format-ALU-FWPMMWRO}), \verb|floatmode = 010|\\

\bitfields{ 1/P, 2/11, 1/1, 1/1, 3/010, 5/registerWo, 1/1, 2/11, 3/000, 1/1, 5/registerYo, 1/0, 5/registerZo, 1/0 }


\paragraph{Syntax} \texttt{fminnwp} \emph{registerWo} = \emph{registerZo}, \emph{registerYo}

\paragraph{Description} The \emph{registerZo} and the \emph{registerYo} are interpreted as two binary 32 floating-point numbers packed into 64 bits. The \emph{registerZo} words are compared to the \emph{registerYo} words. The two minimum number words are packed and stored into the \emph{registerWo}.


\paragraph{Execution} {Reservation table \textsf{ALU\_LITE} (Section~\ref{sec:reservation-ALU-LITE})}
~
{\begin{lstlisting}
stage RR:
  new argument3 = GPR[(registerYo<<1)+1];
  new argument2 = GPR[(registerZo<<1)+1];
  for (I = 0; I < 2; I += 1) {
    result1.32[I] = f32_minNum(argument2.32[I], argument3.32[I])
  };
stage E1:
  GPR[(registerWo<<1)+1] = result1;
\end{lstlisting}}
\newpage
\subsubsection[{\small FMINW: Floating-Point Minimum Word}]{FMINW\hfill Floating-Point Minimum Word}
\label{instr:FMINW}\index{fminw}
 
\paragraph{Encoding} Format \textsf{ALU\_FWMMWR} (Section~\ref{sec:format-ALU-FWMMWR}), \verb|floatmode = 000|\\

\bitfields{ 1/P, 2/11, 1/1, 1/1, 3/000, 6/registerW, 2/01, 3/000, 1/0, 6/registerY, 6/registerZ }


\paragraph{Syntax} \texttt{fminw} \emph{registerW} = \emph{registerZ}, \emph{registerY}

\paragraph{Description} The \emph{registerZ} is compared to the \emph{registerY}, assuming binary 32 floating-point numbers. The minimum value is stored into the \emph{registerW}.


\paragraph{Execution} {Reservation table \textsf{ALU\_LITE} (Section~\ref{sec:reservation-ALU-LITE})}
~
{\begin{lstlisting}
stage RR:
  new argument3 = GPR[registerY];
  new argument2 = GPR[registerZ];
  new result1 = f32_min(argument2, argument3);
stage E1:
  GPR[registerW] = result1;
\end{lstlisting}}
\newpage
\subsubsection[{\small FMINWP: Floating-Point Minimum Word Pair}]{FMINWP\hfill Floating-Point Minimum Word Pair}
\label{instr:FMINWP}\index{fminwp}
 
\paragraph{Encoding} Format \textsf{ALU\_FWPMMWRE} (Section~\ref{sec:format-ALU-FWPMMWRE}), \verb|floatmode = 000|\\

\bitfields{ 1/P, 2/11, 1/1, 1/1, 3/000, 5/registerWe, 1/0, 2/11, 3/000, 1/1, 5/registerYe, 1/0, 5/registerZe, 1/0 }


\paragraph{Syntax} \texttt{fminwp} \emph{registerWe} = \emph{registerZe}, \emph{registerYe}

\paragraph{Description} The \emph{registerZe} and the \emph{registerYe} are interpreted as two binary 32 floating-point numbers packed into 64 bits. The \emph{registerZe} words are compared to the \emph{registerYe} words. The two minimum words are packed and stored into the \emph{registerWe}.


\paragraph{Execution} {Reservation table \textsf{ALU\_LITE} (Section~\ref{sec:reservation-ALU-LITE})}
~
{\begin{lstlisting}
stage RR:
  new argument3 = GPR[registerYe<<1];
  new argument2 = GPR[registerZe<<1];
  for (I = 0; I < 2; I += 1) {
    result1.32[I] = f32_min(argument2.32[I], argument3.32[I])
  };
stage E1:
  GPR[registerWe<<1] = result1;
\end{lstlisting}}
\newpage

\paragraph{Encoding} Format \textsf{ALU\_FWPMMWRO} (Section~\ref{sec:format-ALU-FWPMMWRO}), \verb|floatmode = 000|\\

\bitfields{ 1/P, 2/11, 1/1, 1/1, 3/000, 5/registerWo, 1/1, 2/11, 3/000, 1/1, 5/registerYo, 1/0, 5/registerZo, 1/0 }


\paragraph{Syntax} \texttt{fminwp} \emph{registerWo} = \emph{registerZo}, \emph{registerYo}

\paragraph{Description} The \emph{registerZo} and the \emph{registerYo} are interpreted as two binary 32 floating-point numbers packed into 64 bits. The \emph{registerZo} words are compared to the \emph{registerYo} words. The two minimum words are packed and stored into the \emph{registerWo}.


\paragraph{Execution} {Reservation table \textsf{ALU\_LITE} (Section~\ref{sec:reservation-ALU-LITE})}
~
{\begin{lstlisting}
stage RR:
  new argument3 = GPR[(registerYo<<1)+1];
  new argument2 = GPR[(registerZo<<1)+1];
  for (I = 0; I < 2; I += 1) {
    result1.32[I] = f32_min(argument2.32[I], argument3.32[I])
  };
stage E1:
  GPR[(registerWo<<1)+1] = result1;
\end{lstlisting}}
\newpage
\subsubsection[{\small FMULD: Floating-Point Multiply Double Word}]{FMULD\hfill Floating-Point Multiply Double Word}
\label{instr:FMULD}\index{fmuld}
 
\paragraph{Encoding} Format \textsf{ALU\_FDMUL} (Section~\ref{sec:format-ALU-FDMUL}), \verb|alucode3 = 110|\\

\bitfields{ 1/P, 2/11, 1/1, 1/fnegate, 3/floatmode, 6/registerW, 2/00, 3/110, 1/0, 6/registerY, 6/registerZ }


\paragraph{Syntax} \texttt{fmuld}\emph{fnegate}\emph{floatmode} \emph{registerW} = \emph{registerZ}, \emph{registerY}

\paragraph{Description} The \emph{registerZ} is multiplied by the \emph{registerY}, assuming binary 64 floating-point numbers. If \emph{fnegate} is set, the \emph{registerZ} double word is negated prior to the operation. The result is rounded to the binary 64 floating-point format according to the rounding mode and stored into the \emph{registerW}. This instruction may raise exception bits in the CS register.


\paragraph{Modifier(s)}
\noindent\begin{itemize}
\item floatmode (cf. floatmode in section \ref{par:modifier-floatmode} on page \pageref{par:modifier-floatmode})
\item fnegate (cf. fnegate in section \ref{par:modifier-fnegate} on page \pageref{par:modifier-fnegate})
\end{itemize}

\paragraph{Execution} {Reservation table \textsf{ALU\_LITE} (Section~\ref{sec:reservation-ALU-LITE})}
~
{\begin{lstlisting}
stage ID:
  new fnegate = _ZX_1(fnegate);
  new floatmode = _ZX_3(floatmode);
stage RR:
  new argument3 = GPR[registerY];
  new argument2 = GPR[registerZ];
  new RM = decode_riscv_float_rounding_mode(floatmode);
  new fsign = fnegate ? 0x8000000000000000 : 0;
  new result1 = f64_mul(RM, argument2 ^ fsign, argument3);
stage E4:
  GPR[registerW] = result1;
  commit_float_exception_flags();
\end{lstlisting}}
\newpage
\subsubsection[{\small FMULH: Floating-Point Multiply Half Word}]{FMULH\hfill Floating-Point Multiply Half Word}
\label{instr:FMULH}\index{fmulh}
 
\paragraph{Encoding} Format \textsf{ALU\_FHMUL} (Section~\ref{sec:format-ALU-FHMUL}), \verb|alucode3 = 110|\\

\bitfields{ 1/P, 2/11, 1/1, 1/fnegate, 3/floatmode, 6/registerW, 2/01, 3/110, 1/1, 6/registerY, 6/registerZ }


\paragraph{Syntax} \texttt{fmulh}\emph{fnegate}\emph{floatmode} \emph{registerW} = \emph{registerZ}, \emph{registerY}

\paragraph{Description} The \emph{registerZ} is multiplied by the \emph{registerY}, assuming binary 16 floating-point numbers. If \emph{fnegate} is set, the \emph{registerZ} half word is negated prior to the operation. The result is rounded to the binary 16 floating-point format according to the rounding mode and stored into the \emph{registerW}. This instruction may raise exception bits in the CS register.


\paragraph{Modifier(s)}
\noindent\begin{itemize}
\item floatmode (cf. floatmode in section \ref{par:modifier-floatmode} on page \pageref{par:modifier-floatmode})
\item fnegate (cf. fnegate in section \ref{par:modifier-fnegate} on page \pageref{par:modifier-fnegate})
\end{itemize}

\paragraph{Execution} {Reservation table \textsf{ALU\_LITE} (Section~\ref{sec:reservation-ALU-LITE})}
~
{\begin{lstlisting}
stage ID:
  new fnegate = _ZX_1(fnegate);
  new floatmode = _ZX_3(floatmode);
stage RR:
  new argument3 = GPR[registerY];
  new argument2 = GPR[registerZ];
  new RM = decode_riscv_float_rounding_mode(floatmode);
  new fsign = fnegate ? 0x8000 : 0;
  new result1 = f16_mul(RM, argument2 ^ fsign, argument3);
stage E3:
  GPR[registerW] = result1;
  commit_float_exception_flags();
\end{lstlisting}}
\newpage
\subsubsection[{\small FMULHQ: Floating-Point Multiply Half Word Quadruple}]{FMULHQ\hfill Floating-Point Multiply Half Word Quadruple}
\label{instr:FMULHQ}\index{fmulhq}
 
\paragraph{Encoding} Format \textsf{ALU\_FHQMULE} (Section~\ref{sec:format-ALU-FHQMULE}), \verb|alucode3 = 110|\\

\bitfields{ 1/P, 2/11, 1/1, 1/fnegate, 3/floatmode, 5/registerWe, 1/0, 2/11, 3/110, 1/1, 5/registerYe, 1/0, 5/registerZe, 1/1 }


\paragraph{Syntax} \texttt{fmulhq}\emph{fnegate}\emph{floatmode} \emph{registerWe} = \emph{registerZe}, \emph{registerYe}

\paragraph{Description} The \emph{registerZe} and the \emph{registerYe} are interpreted as four binary 16 floating-point numbers packed into 64 bits. If \emph{fnegate} is set, the \emph{registerZe} half words are negated prior to the operation. The \emph{registerZe} half words are multiplied by the \emph{registerYe} half words. The four resulting half words are packed and stored into the \emph{registerWe}.


\paragraph{Modifier(s)}
\noindent\begin{itemize}
\item floatmode (cf. floatmode in section \ref{par:modifier-floatmode} on page \pageref{par:modifier-floatmode})
\item fnegate (cf. fnegate in section \ref{par:modifier-fnegate} on page \pageref{par:modifier-fnegate})
\end{itemize}

\paragraph{Execution} {Reservation table \textsf{ALU\_LITE} (Section~\ref{sec:reservation-ALU-LITE})}
~
{\begin{lstlisting}
stage ID:
  new fnegate = _ZX_1(fnegate);
  new floatmode = _ZX_3(floatmode);
stage RR:
  new argument3 = GPR[registerYe<<1];
  new argument2 = GPR[registerZe<<1];
  new RM = decode_riscv_float_rounding_mode(floatmode);
  new fsign = fnegate ? 0x8000 : 0;
  for (I = 0; I < 4; I += 1) {
    result1.16[I] = f16_mul(RM, argument2.16[I] ^ fsign, argument3.16[I])
  };
stage E4:
  GPR[registerWe<<1] = result1;
  commit_float_exception_flags();
\end{lstlisting}}
\newpage

\paragraph{Encoding} Format \textsf{ALU\_FHQMULO} (Section~\ref{sec:format-ALU-FHQMULO}), \verb|alucode3 = 110|\\

\bitfields{ 1/P, 2/11, 1/1, 1/fnegate, 3/floatmode, 5/registerWo, 1/1, 2/11, 3/110, 1/1, 5/registerYo, 1/0, 5/registerZo, 1/1 }


\paragraph{Syntax} \texttt{fmulhq}\emph{fnegate}\emph{floatmode} \emph{registerWo} = \emph{registerZo}, \emph{registerYo}

\paragraph{Description} The \emph{registerZo} and the \emph{registerYo} are interpreted as four binary 16 floating-point numbers packed into 64 bits. If \emph{fnegate} is set, the \emph{registerZo} half words are negated prior to the operation. The \emph{registerZo} half words are multiplied by the \emph{registerYo} half words. The four resulting half words are packed and stored into the \emph{registerWo}.


\paragraph{Modifier(s)}
\noindent\begin{itemize}
\item floatmode (cf. floatmode in section \ref{par:modifier-floatmode} on page \pageref{par:modifier-floatmode})
\item fnegate (cf. fnegate in section \ref{par:modifier-fnegate} on page \pageref{par:modifier-fnegate})
\end{itemize}

\paragraph{Execution} {Reservation table \textsf{ALU\_LITE} (Section~\ref{sec:reservation-ALU-LITE})}
~
{\begin{lstlisting}
stage ID:
  new fnegate = _ZX_1(fnegate);
  new floatmode = _ZX_3(floatmode);
stage RR:
  new argument3 = GPR[(registerYo<<1)+1];
  new argument2 = GPR[(registerZo<<1)+1];
  new RM = decode_riscv_float_rounding_mode(floatmode);
  new fsign = fnegate ? 0x8000 : 0;
  for (I = 0; I < 4; I += 1) {
    result1.16[I] = f16_mul(RM, argument2.16[I] ^ fsign, argument3.16[I])
  };
stage E4:
  GPR[(registerWo<<1)+1] = result1;
  commit_float_exception_flags();
\end{lstlisting}}
\newpage
\subsubsection[{\small FMULW: Floating-Point Multiply Word}]{FMULW\hfill Floating-Point Multiply Word}
\label{instr:FMULW}\index{fmulw}
 
\paragraph{Encoding} Format \textsf{ALU\_FWMUL} (Section~\ref{sec:format-ALU-FWMUL}), \verb|alucode3 = 110|\\

\bitfields{ 1/P, 2/11, 1/1, 1/fnegate, 3/floatmode, 6/registerW, 2/01, 3/110, 1/0, 6/registerY, 6/registerZ }


\paragraph{Syntax} \texttt{fmulw}\emph{fnegate}\emph{floatmode} \emph{registerW} = \emph{registerZ}, \emph{registerY}

\paragraph{Description} The \emph{registerZ} is multiplied by the \emph{registerY}, assuming binary 32 floating-point numbers. If \emph{fnegate} is set, the \emph{registerZ} word is negated prior to the operation. The result is rounded to the binary 32 floating-point format according to the rounding mode and stored into the \emph{registerW}. This instruction may raise exception bits in the CS register.


\paragraph{Modifier(s)}
\noindent\begin{itemize}
\item floatmode (cf. floatmode in section \ref{par:modifier-floatmode} on page \pageref{par:modifier-floatmode})
\item fnegate (cf. fnegate in section \ref{par:modifier-fnegate} on page \pageref{par:modifier-fnegate})
\end{itemize}

\paragraph{Execution} {Reservation table \textsf{ALU\_LITE} (Section~\ref{sec:reservation-ALU-LITE})}
~
{\begin{lstlisting}
stage ID:
  new fnegate = _ZX_1(fnegate);
  new floatmode = _ZX_3(floatmode);
stage RR:
  new argument3 = GPR[registerY];
  new argument2 = GPR[registerZ];
  new RM = decode_riscv_float_rounding_mode(floatmode);
  new fsign = fnegate ? 0x80000000 : 0;
  new result1 = f32_mul(RM, argument2 ^ fsign, argument3);
stage E4:
  GPR[registerW] = result1;
  commit_float_exception_flags();
\end{lstlisting}}
\newpage
\subsubsection[{\small FMULWC: Floating-Point Multiply Word Complex}]{FMULWC\hfill Floating-Point Multiply Word Complex}
\label{instr:FMULWC}\index{fmulwc}
 
\paragraph{Encoding} Format \textsf{ALU\_FWCOP} (Section~\ref{sec:format-ALU-FWCOP}), \verb|alucode6 = 10|\\

\bitfields{ 1/P, 2/11, 1/1, 1/imultiply, 3/floatmode, 6/registerW, 2/00, 2/10, 1/conjugate, 1/1, 6/registerY, 6/registerZ }


\paragraph{Syntax} \texttt{fmulwc}\emph{conjugate}\emph{imultiply}\emph{floatmode} \emph{registerW} = \emph{registerZ}, \emph{registerY}

\paragraph{Description} The \emph{registerY} and the \emph{registerZ} are interpreted as binary 32 floating-point complex numbers. If \emph{conjugate} is set, the complex conjugate of the \emph{registerZ} is used in the multiplication. The complex product is and rounded according to the rounding mode. If \emph{imultiply} is set, the result is multiplied by the imaginary unit. The resulting binary 32 floating-point complex number is stored back into the \emph{registerW}. This instruction may raise exception bits in the CS register.


\paragraph{Modifier(s)}
\noindent\begin{itemize}
\item floatmode (cf. floatmode in section \ref{par:modifier-floatmode} on page \pageref{par:modifier-floatmode})
\item imultiply (cf. imultiply in section \ref{par:modifier-imultiply} on page \pageref{par:modifier-imultiply})
\item conjugate (cf. conjugate in section \ref{par:modifier-conjugate} on page \pageref{par:modifier-conjugate})
\end{itemize}

\paragraph{Execution} {Reservation table \textsf{ALU\_LITE} (Section~\ref{sec:reservation-ALU-LITE})}
~
{\begin{lstlisting}
stage ID:
  new conjugate = _ZX_1(conjugate);
  new imultiply = _ZX_1(imultiply);
  new floatmode = _ZX_3(floatmode);
stage RR:
  new argument3 = GPR[registerY];
  new argument2 = GPR[registerZ];
  new argument1 = GPR[registerW];
  new RM = decode_riscv_float_rounding_mode(floatmode);
  if (conjugate) {
    argument2.64[0] = fconj_32_32(argument2.64[0]);
  }
  result1.64[0] = fmulc_32_32(RM, argument2.64[0], argument3.64[0]);
  if (imultiply) {
    new result1 = _SWAP_32(result1);
    result1.32[0] = result1.32[0] ^ 0x80000000;
  }
stage E4:
  GPR[registerW] = result1;
  commit_float_exception_flags();
\end{lstlisting}}
\newpage
\subsubsection[{\small FMULWP: Floating-Point Multiply Word Pair}]{FMULWP\hfill Floating-Point Multiply Word Pair}
\label{instr:FMULWP}\index{fmulwp}
 
\paragraph{Encoding} Format \textsf{ALU\_FWPMULE} (Section~\ref{sec:format-ALU-FWPMULE}), \verb|alucode3 = 110|\\

\bitfields{ 1/P, 2/11, 1/1, 1/fnegate, 3/floatmode, 5/registerWe, 1/0, 2/11, 3/110, 1/1, 5/registerYe, 1/0, 5/registerZe, 1/0 }


\paragraph{Syntax} \texttt{fmulwp}\emph{fnegate}\emph{floatmode} \emph{registerWe} = \emph{registerZe}, \emph{registerYe}

\paragraph{Description} The \emph{registerZe} and the \emph{registerYe} are interpreted as two binary 32 floating-point numbers packed into 64 bits. If \emph{fnegate} is set, the \emph{registerZe} words are negated prior to the operation. The \emph{registerZe} words are multiplied by the \emph{registerYe} words. The two resulting words are packed and stored into the \emph{registerWe}.


\paragraph{Modifier(s)}
\noindent\begin{itemize}
\item floatmode (cf. floatmode in section \ref{par:modifier-floatmode} on page \pageref{par:modifier-floatmode})
\item fnegate (cf. fnegate in section \ref{par:modifier-fnegate} on page \pageref{par:modifier-fnegate})
\end{itemize}

\paragraph{Execution} {Reservation table \textsf{ALU\_LITE} (Section~\ref{sec:reservation-ALU-LITE})}
~
{\begin{lstlisting}
stage ID:
  new fnegate = _ZX_1(fnegate);
  new floatmode = _ZX_3(floatmode);
stage RR:
  new argument3 = GPR[registerYe<<1];
  new argument2 = GPR[registerZe<<1];
  new RM = decode_riscv_float_rounding_mode(floatmode);
  new fsign = fnegate ? 0x80000000 : 0;
  for (I = 0; I < 2; I += 1) {
    result1.32[I] = f32_mul(RM, argument2.32[I] ^ fsign, argument3.32[I])
  };
stage E4:
  GPR[registerWe<<1] = result1;
  commit_float_exception_flags();
\end{lstlisting}}
\newpage

\paragraph{Encoding} Format \textsf{ALU\_FWPMULO} (Section~\ref{sec:format-ALU-FWPMULO}), \verb|alucode3 = 110|\\

\bitfields{ 1/P, 2/11, 1/1, 1/fnegate, 3/floatmode, 5/registerWo, 1/1, 2/11, 3/110, 1/1, 5/registerYo, 1/0, 5/registerZo, 1/0 }


\paragraph{Syntax} \texttt{fmulwp}\emph{fnegate}\emph{floatmode} \emph{registerWo} = \emph{registerZo}, \emph{registerYo}

\paragraph{Description} The \emph{registerZo} and the \emph{registerYo} are interpreted as two binary 32 floating-point numbers packed into 64 bits. If \emph{fnegate} is set, the \emph{registerZo} words are negated prior to the operation. The \emph{registerZo} words are multiplied by the \emph{registerYo} words. The two resulting words are packed and stored into the \emph{registerWo}.


\paragraph{Modifier(s)}
\noindent\begin{itemize}
\item floatmode (cf. floatmode in section \ref{par:modifier-floatmode} on page \pageref{par:modifier-floatmode})
\item fnegate (cf. fnegate in section \ref{par:modifier-fnegate} on page \pageref{par:modifier-fnegate})
\end{itemize}

\paragraph{Execution} {Reservation table \textsf{ALU\_LITE} (Section~\ref{sec:reservation-ALU-LITE})}
~
{\begin{lstlisting}
stage ID:
  new fnegate = _ZX_1(fnegate);
  new floatmode = _ZX_3(floatmode);
stage RR:
  new argument3 = GPR[(registerYo<<1)+1];
  new argument2 = GPR[(registerZo<<1)+1];
  new RM = decode_riscv_float_rounding_mode(floatmode);
  new fsign = fnegate ? 0x80000000 : 0;
  for (I = 0; I < 2; I += 1) {
    result1.32[I] = f32_mul(RM, argument2.32[I] ^ fsign, argument3.32[I])
  };
stage E4:
  GPR[(registerWo<<1)+1] = result1;
  commit_float_exception_flags();
\end{lstlisting}}
\newpage
\subsubsection[{\small FNARROWDW: Floating Point Narrow Double Word to Word}]{FNARROWDW\hfill Floating Point Narrow Double Word to Word}
\label{instr:FNARROWDW}\index{fnarrowdw}
 
\paragraph{Encoding} Format \textsf{ALU\_FWWR} (Section~\ref{sec:format-ALU-FWWR}), \verb|fpucode6 = 1100|\\

\bitfields{ 1/P, 2/11, 1/1, 1/1, 3/floatmode, 6/registerW, 2/01, 3/001, 1/0, 4/1100, 1/--, 1/0, 6/registerZ }


\paragraph{Syntax} \texttt{fnarrowdw}\emph{floatmode} \emph{registerW} = \emph{registerZ}

\paragraph{Description} The \emph{registerZ} interpreted as a binary 64 floating-point number is converted to a binary 32 floating-point number according to the rounding mode and stored into the \emph{registerW}. This instruction may raise exception bits in the CS register.


\paragraph{Modifier(s)}
\noindent\begin{itemize}
\item floatmode (cf. floatmode in section \ref{par:modifier-floatmode} on page \pageref{par:modifier-floatmode})
\end{itemize}

\paragraph{Execution} {Reservation table \textsf{ALU\_LITE} (Section~\ref{sec:reservation-ALU-LITE})}
~
{\begin{lstlisting}
stage ID:
  new floatmode = _ZX_3(floatmode);
stage RR:
  new argument2 = GPR[registerZ];
  new RM = decode_riscv_float_rounding_mode(floatmode);
  new result1 = f64_to_f32(RM, argument2.64[0]);
stage E1:
  GPR[registerW] = result1;
  commit_float_exception_flags();
\end{lstlisting}}
\newpage
\subsubsection[{\small FNARROWDWP: Floating Point Narrow Double Word to Word Pair}]{FNARROWDWP\hfill Floating Point Narrow Double Word to Word Pair}
\label{instr:FNARROWDWP}\index{fnarrowdwp}
 
\paragraph{Encoding} Format \textsf{ALU\_FWPWR} (Section~\ref{sec:format-ALU-FWPWR}), \verb|fpucode6 = 1100|\\

\bitfields{ 1/P, 2/11, 1/1, 1/1, 3/floatmode, 6/registerW, 2/11, 3/001, 1/1, 4/1100, 1/ziplanes, 1/0, 5/registerP, 1/0 }


\paragraph{Syntax} \texttt{fnarrowdwp}\emph{floatmode} \emph{registerW} = \emph{registerP}

\paragraph{Description} The \emph{registerP} is interpreted as a pair of binary 64 floating-point numbers packed into 128 bits. The \emph{registerP} double words are converted to binary 32 floating-point numbers according to the rounding mode. The resulting words are packed and stored into the \emph{registerW}. This instruction may raise exception bits in the CS register.


\paragraph{Modifier(s)}
\noindent\begin{itemize}
\item floatmode (cf. floatmode in section \ref{par:modifier-floatmode} on page \pageref{par:modifier-floatmode})
\end{itemize}

\paragraph{Execution} {Reservation table \textsf{ALU\_LITE} (Section~\ref{sec:reservation-ALU-LITE})}
~
{\begin{lstlisting}
stage ID:
  new floatmode = _ZX_3(floatmode);
stage RR:
  new argument2 = PGR[registerP];
  new RM = decode_riscv_float_rounding_mode(floatmode);
  result1.32[0] = f64_to_f32(RM, argument2.64[0]);
  result1.32[1] = f64_to_f32(RM, argument2.64[1]);
stage E1:
  GPR[registerW] = result1;
  commit_float_exception_flags();
\end{lstlisting}}
\newpage
\subsubsection[{\small FNARROWWH: Floating Point Narrow Word to Half Word}]{FNARROWWH\hfill Floating Point Narrow Word to Half Word}
\label{instr:FNARROWWH}\index{fnarrowwh}
 
\paragraph{Encoding} Format \textsf{ALU\_FHWR} (Section~\ref{sec:format-ALU-FHWR}), \verb|fpucode6 = 1100|\\

\bitfields{ 1/P, 2/11, 1/1, 1/1, 3/floatmode, 6/registerW, 2/01, 3/001, 1/1, 4/1100, 1/--, 1/0, 6/registerZ }


\paragraph{Syntax} \texttt{fnarrowwh}\emph{floatmode} \emph{registerW} = \emph{registerZ}

\paragraph{Description} The \emph{registerZ} interpreted as a binary 32 floating-point number is converted to a binary 16 floating-point number according to the rounding mode and stored into the \emph{registerW}. This instruction may raise exception bits in the CS register.


\paragraph{Modifier(s)}
\noindent\begin{itemize}
\item floatmode (cf. floatmode in section \ref{par:modifier-floatmode} on page \pageref{par:modifier-floatmode})
\end{itemize}

\paragraph{Execution} {Reservation table \textsf{ALU\_LITE} (Section~\ref{sec:reservation-ALU-LITE})}
~
{\begin{lstlisting}
stage ID:
  new floatmode = _ZX_3(floatmode);
stage RR:
  new argument2 = GPR[registerZ];
  new RM = decode_riscv_float_rounding_mode(floatmode);
  new result1 = f32_to_f16(RM, argument2.32[0]);
stage E1:
  GPR[registerW] = result1;
  commit_float_exception_flags();
\end{lstlisting}}
\newpage
\subsubsection[{\small FNEGD: Floating-Point Negate Double Word}]{FNEGD\hfill Floating-Point Negate Double Word}
\label{instr:FNEGD}\index{fnegd}
 
\paragraph{Encoding} Format \textsf{ALU\_FDMO1} (Section~\ref{sec:format-ALU-FDMO1}), \verb|fpucode6 = 0010|\\

\bitfields{ 1/P, 2/11, 1/1, 1/1, 3/111, 6/registerW, 2/00, 3/000, 1/0, 4/0010, 1/--, 1/0, 6/registerZ }


\paragraph{Syntax} \texttt{fnegd} \emph{registerW} = \emph{registerZ}

\paragraph{Description} The negated value the \emph{registerZ}, assuming binary 64 floating-point number, is stored into the \emph{registerW}. Implemented as a bitwise operation, does not raise FP exceptions.


\paragraph{Execution} {Reservation table \textsf{ALU\_LITE} (Section~\ref{sec:reservation-ALU-LITE})}
~
{\begin{lstlisting}
stage RR:
  new argument2 = GPR[registerZ];
  new result1 = argument2 ^ 0x8000000000000000;
stage E1:
  GPR[registerW] = result1;
\end{lstlisting}}
\newpage
\subsubsection[{\small FNEGH: Floating-Point Negate Half Word}]{FNEGH\hfill Floating-Point Negate Half Word}
\label{instr:FNEGH}\index{fnegh}
 
\paragraph{Encoding} Format \textsf{ALU\_FHMO1} (Section~\ref{sec:format-ALU-FHMO1}), \verb|fpucode6 = 0010|\\

\bitfields{ 1/P, 2/11, 1/1, 1/1, 3/111, 6/registerW, 2/01, 3/000, 1/1, 4/0010, 1/--, 1/0, 6/registerZ }


\paragraph{Syntax} \texttt{fnegh} \emph{registerW} = \emph{registerZ}

\paragraph{Description} The negated value the \emph{registerZ}, assuming binary 16 floating-point number, is stored into the \emph{registerW}. Implemented as a bitwise operation, does not raise FP exceptions.


\paragraph{Execution} {Reservation table \textsf{ALU\_LITE} (Section~\ref{sec:reservation-ALU-LITE})}
~
{\begin{lstlisting}
stage RR:
  new argument2 = GPR[registerZ];
  new result1 = _ZX_32(argument2 ^ 0x8000);
stage E1:
  GPR[registerW] = result1;
\end{lstlisting}}
\newpage
\subsubsection[{\small FNEGHQ: Floating-Point Negate Half Word Quadruple}]{FNEGHQ\hfill Floating-Point Negate Half Word Quadruple}
\label{instr:FNEGHQ}\index{fneghq}
 
\paragraph{Encoding} Format \textsf{ALU\_FHQMO1E} (Section~\ref{sec:format-ALU-FHQMO1E}), \verb|fpucode6 = 0010|\\

\bitfields{ 1/P, 2/11, 1/1, 1/1, 3/111, 5/registerWe, 1/0, 2/11, 3/000, 1/1, 4/0010, 1/--, 1/0, 5/registerZe, 1/1 }


\paragraph{Syntax} \texttt{fneghq} \emph{registerWe} = \emph{registerZe}

\paragraph{Description} The \emph{registerZe} is interpreted as four binary 16 floating-point numbers packed into 64 bits. The four negated values the \emph{registerZe} half words are packed and stored into the \emph{registerWe}.


\paragraph{Execution} {Reservation table \textsf{ALU\_LITE} (Section~\ref{sec:reservation-ALU-LITE})}
~
{\begin{lstlisting}
stage RR:
  new argument2 = GPR[registerZe<<1];
  for (I = 0; I < 4; I += 1) {
    result1.16[I] = argument2.16[I] ^ 0x8000
  };
stage E1:
  GPR[registerWe<<1] = result1;
\end{lstlisting}}
\newpage

\paragraph{Encoding} Format \textsf{ALU\_FHQMO1O} (Section~\ref{sec:format-ALU-FHQMO1O}), \verb|fpucode6 = 0010|\\

\bitfields{ 1/P, 2/11, 1/1, 1/1, 3/111, 5/registerWo, 1/1, 2/11, 3/000, 1/1, 4/0010, 1/--, 1/0, 5/registerZo, 1/1 }


\paragraph{Syntax} \texttt{fneghq} \emph{registerWo} = \emph{registerZo}

\paragraph{Description} The \emph{registerZo} is interpreted as four binary 16 floating-point numbers packed into 64 bits. The four negated values the \emph{registerZo} half words are packed and stored into the \emph{registerWo}.


\paragraph{Execution} {Reservation table \textsf{ALU\_LITE} (Section~\ref{sec:reservation-ALU-LITE})}
~
{\begin{lstlisting}
stage RR:
  new argument2 = GPR[(registerZo<<1)+1];
  for (I = 0; I < 4; I += 1) {
    result1.16[I] = argument2.16[I] ^ 0x8000
  };
stage E1:
  GPR[(registerWo<<1)+1] = result1;
\end{lstlisting}}
\newpage
\subsubsection[{\small FNEGW: Floating-Point Negate Word}]{FNEGW\hfill Floating-Point Negate Word}
\label{instr:FNEGW}\index{fnegw}
 
\paragraph{Encoding} Format \textsf{ALU\_FWMO1} (Section~\ref{sec:format-ALU-FWMO1}), \verb|fpucode6 = 0010|\\

\bitfields{ 1/P, 2/11, 1/1, 1/1, 3/111, 6/registerW, 2/01, 3/000, 1/0, 4/0010, 1/--, 1/0, 6/registerZ }


\paragraph{Syntax} \texttt{fnegw} \emph{registerW} = \emph{registerZ}

\paragraph{Description} The negated value the \emph{registerZ}, assuming binary 32 floating-point number, is stored into the \emph{registerW}. Implemented as a bitwise operation, does not raise FP exceptions.


\paragraph{Execution} {Reservation table \textsf{ALU\_LITE} (Section~\ref{sec:reservation-ALU-LITE})}
~
{\begin{lstlisting}
stage RR:
  new argument2 = GPR[registerZ];
  new result1 = _ZX_32(argument2 ^ 0x80000000);
stage E1:
  GPR[registerW] = result1;
\end{lstlisting}}
\newpage
\subsubsection[{\small FNEGWP: Floating-Point Negate Word Pair}]{FNEGWP\hfill Floating-Point Negate Word Pair}
\label{instr:FNEGWP}\index{fnegwp}
 
\paragraph{Encoding} Format \textsf{ALU\_FWPMO1} (Section~\ref{sec:format-ALU-FWPMO1}), \verb|fpucode6 = 0010|\\

\bitfields{ 1/P, 2/11, 1/1, 1/1, 3/111, 6/registerW, 2/00, 3/000, 1/0, 4/0010, 1/--, 1/1, 6/registerZ }


\paragraph{Syntax} \texttt{fnegwp} \emph{registerW} = \emph{registerZ}

\paragraph{Description} The \emph{registerZ} is interpreted as two binary 32 floating-point numbers packed into 64 bits. The two negated values the \emph{registerZ} words are packed and stored into the \emph{registerW}.


\paragraph{Execution} {Reservation table \textsf{ALU\_LITE} (Section~\ref{sec:reservation-ALU-LITE})}
~
{\begin{lstlisting}
stage RR:
  new argument2 = GPR[registerZ];
  for (I = 0; I < 2; I += 1) {
    result1.32[I] = argument2.32[I] ^ 0x80000000
  };
stage E1:
  GPR[registerW] = result1;
\end{lstlisting}}
\newpage
\subsubsection[{\small FRINTD: Floating Point Round to Int Double Word}]{FRINTD\hfill Floating Point Round to Int Double Word}
\label{instr:FRINTD}\index{frintd}
 
\paragraph{Encoding} Format \textsf{ALU\_FDMO0} (Section~\ref{sec:format-ALU-FDMO0}), \verb|fpucode6 = 1001|\\

\bitfields{ 1/P, 2/11, 1/1, 1/1, 3/floatmode, 6/registerW, 2/00, 3/001, 1/0, 4/1001, 1/--, 1/0, 6/registerZ }


\paragraph{Syntax} \texttt{frintd}\emph{floatmode} \emph{registerW} = \emph{registerZ}

\paragraph{Description} The the \emph{registerZ} interpreted as a 64-bit floating-point number is rounded to integer in the 64-bit floating-point reprsentation according to the rounding mode and stored into the \emph{registerW}. This instruction may raise exception bits in the CS register.


\paragraph{Modifier(s)}
\noindent\begin{itemize}
\item floatmode (cf. floatmode in section \ref{par:modifier-floatmode} on page \pageref{par:modifier-floatmode})
\end{itemize}

\paragraph{Execution} {Reservation table \textsf{ALU\_LITE} (Section~\ref{sec:reservation-ALU-LITE})}
~
{\begin{lstlisting}
stage ID:
  new floatmode = _ZX_3(floatmode);
stage RR:
  new argument2 = GPR[registerZ];
  new RM = decode_riscv_float_rounding_mode(floatmode);
  new result1 = f64_rint(RM, argument2);
stage E4:
  GPR[registerW] = result1;
  commit_float_exception_flags();
\end{lstlisting}}
\newpage
\subsubsection[{\small FRINTH: Floating Point Round to Int Half Word}]{FRINTH\hfill Floating Point Round to Int Half Word}
\label{instr:FRINTH}\index{frinth}
 
\paragraph{Encoding} Format \textsf{ALU\_FHMO0} (Section~\ref{sec:format-ALU-FHMO0}), \verb|fpucode6 = 1001|\\

\bitfields{ 1/P, 2/11, 1/1, 1/1, 3/floatmode, 6/registerW, 2/01, 3/001, 1/1, 4/1001, 1/--, 1/0, 6/registerZ }


\paragraph{Syntax} \texttt{frinth}\emph{floatmode} \emph{registerW} = \emph{registerZ}

\paragraph{Description} The the \emph{registerZ} interpreted as a 16-bit floating-point number is rounded to integer in the 16-bit floating-point reprsentation according to the rounding mode and stored into the \emph{registerW}. This instruction may raise exception bits in the CS register.


\paragraph{Modifier(s)}
\noindent\begin{itemize}
\item floatmode (cf. floatmode in section \ref{par:modifier-floatmode} on page \pageref{par:modifier-floatmode})
\end{itemize}

\paragraph{Execution} {Reservation table \textsf{ALU\_LITE} (Section~\ref{sec:reservation-ALU-LITE})}
~
{\begin{lstlisting}
stage ID:
  new floatmode = _ZX_3(floatmode);
stage RR:
  new argument2 = GPR[registerZ];
  new RM = decode_riscv_float_rounding_mode(floatmode);
  new result1 = f16_rint(RM, argument2);
stage SF:
  GPR[registerW] = result1;
  commit_float_exception_flags();
\end{lstlisting}}
\newpage
\subsubsection[{\small FRINTW: Floating Point Round to Int Word}]{FRINTW\hfill Floating Point Round to Int Word}
\label{instr:FRINTW}\index{frintw}
 
\paragraph{Encoding} Format \textsf{ALU\_FWMO0} (Section~\ref{sec:format-ALU-FWMO0}), \verb|fpucode6 = 1001|\\

\bitfields{ 1/P, 2/11, 1/1, 1/1, 3/floatmode, 6/registerW, 2/01, 3/001, 1/0, 4/1001, 1/--, 1/0, 6/registerZ }


\paragraph{Syntax} \texttt{frintw}\emph{floatmode} \emph{registerW} = \emph{registerZ}

\paragraph{Description} The the \emph{registerZ} interpreted as a 32-bit floating-point number is rounded to integer in the 32-bit floating-point reprsentation according to the rounding mode and stored into the \emph{registerW}. This instruction may raise exception bits in the CS register.


\paragraph{Modifier(s)}
\noindent\begin{itemize}
\item floatmode (cf. floatmode in section \ref{par:modifier-floatmode} on page \pageref{par:modifier-floatmode})
\end{itemize}

\paragraph{Execution} {Reservation table \textsf{ALU\_LITE} (Section~\ref{sec:reservation-ALU-LITE})}
~
{\begin{lstlisting}
stage ID:
  new floatmode = _ZX_3(floatmode);
stage RR:
  new argument2 = GPR[registerZ];
  new RM = decode_riscv_float_rounding_mode(floatmode);
  new result1 = f32_rint(RM, argument2);
stage SF:
  GPR[registerW] = result1;
  commit_float_exception_flags();
\end{lstlisting}}
\newpage
\subsubsection[{\small FSBFD: Floating-Point Subtract From Double Word}]{FSBFD\hfill Floating-Point Subtract From Double Word}
\label{instr:FSBFD}\index{fsbfd}
 
\paragraph{Encoding} Format \textsf{ALU\_FDADD} (Section~\ref{sec:format-ALU-FDADD}), \verb|alucode3 = 001|\\

\bitfields{ 1/P, 2/11, 1/1, 1/0, 3/floatmode, 6/registerW, 2/00, 3/001, 1/0, 6/registerY, 6/registerZ }


\paragraph{Syntax} \texttt{fsbfd}\emph{floatmode} \emph{registerW} = \emph{registerZ}, \emph{registerY}

\paragraph{Description} The \emph{registerZ} is subtracted from the \emph{registerY}, assuming binary 64 floating-point numbers. The result is rounded to the binary 64 floating-point format according to the rounding mode and stored into the \emph{registerW}. This instruction may raise exception bits in the CS register.


\paragraph{Modifier(s)}
\noindent\begin{itemize}
\item floatmode (cf. floatmode in section \ref{par:modifier-floatmode} on page \pageref{par:modifier-floatmode})
\end{itemize}

\paragraph{Execution} {Reservation table \textsf{ALU\_LITE} (Section~\ref{sec:reservation-ALU-LITE})}
~
{\begin{lstlisting}
stage ID:
  new floatmode = _ZX_3(floatmode);
stage RR:
  new argument3 = GPR[registerY];
  new argument2 = GPR[registerZ];
  new RM = decode_riscv_float_rounding_mode(floatmode);
  new result1 = f64_sub(RM, argument3, argument2);
stage E4:
  GPR[registerW] = result1;
  commit_float_exception_flags();
\end{lstlisting}}
\newpage
\subsubsection[{\small FSBFH: Floating-Point Subtract From Half Word}]{FSBFH\hfill Floating-Point Subtract From Half Word}
\label{instr:FSBFH}\index{fsbfh}
 
\paragraph{Encoding} Format \textsf{ALU\_FHADD} (Section~\ref{sec:format-ALU-FHADD}), \verb|alucode3 = 001|\\

\bitfields{ 1/P, 2/11, 1/1, 1/0, 3/floatmode, 6/registerW, 2/01, 3/001, 1/1, 6/registerY, 6/registerZ }


\paragraph{Syntax} \texttt{fsbfh}\emph{floatmode} \emph{registerW} = \emph{registerZ}, \emph{registerY}

\paragraph{Description} The \emph{registerZ} is subtracted from the \emph{registerY}, assuming binary 16 floating-point numbers. The result is rounded to the binary 16 floating-point format according to the rounding mode and stored into the \emph{registerW}. This instruction may raise exception bits in the CS register.


\paragraph{Modifier(s)}
\noindent\begin{itemize}
\item floatmode (cf. floatmode in section \ref{par:modifier-floatmode} on page \pageref{par:modifier-floatmode})
\end{itemize}

\paragraph{Execution} {Reservation table \textsf{ALU\_LITE} (Section~\ref{sec:reservation-ALU-LITE})}
~
{\begin{lstlisting}
stage ID:
  new floatmode = _ZX_3(floatmode);
stage RR:
  new argument3 = GPR[registerY];
  new argument2 = GPR[registerZ];
  new RM = decode_riscv_float_rounding_mode(floatmode);
  new result1 = f16_sub(RM, argument3, argument2);
stage E3:
  GPR[registerW] = result1;
  commit_float_exception_flags();
\end{lstlisting}}
\newpage
\subsubsection[{\small FSBFHQ: Floating-Point Subtract From Half Word Quadruple}]{FSBFHQ\hfill Floating-Point Subtract From Half Word Quadruple}
\label{instr:FSBFHQ}\index{fsbfhq}
 
\paragraph{Encoding} Format \textsf{ALU\_FHQADDE} (Section~\ref{sec:format-ALU-FHQADDE}), \verb|alucode3 = 001|\\

\bitfields{ 1/P, 2/11, 1/1, 1/0, 3/floatmode, 5/registerWe, 1/0, 2/11, 3/001, 1/1, 5/registerYe, 1/--, 5/registerZe, 1/1 }


\paragraph{Syntax} \texttt{fsbfhq}\emph{floatmode} \emph{registerWe} = \emph{registerZe}, \emph{registerYe}

\paragraph{Description} The \emph{registerZe} and the \emph{registerYe} are interpreted as four binary 16 floating-point numbers packed into 64 bits. The \emph{registerZe} half words are subtracted from the \emph{registerYe} half words. The four resulting half words are packed and stored into the \emph{registerWe}.


\paragraph{Modifier(s)}
\noindent\begin{itemize}
\item floatmode (cf. floatmode in section \ref{par:modifier-floatmode} on page \pageref{par:modifier-floatmode})
\end{itemize}

\paragraph{Execution} {Reservation table \textsf{ALU\_LITE} (Section~\ref{sec:reservation-ALU-LITE})}
~
{\begin{lstlisting}
stage ID:
  new floatmode = _ZX_3(floatmode);
stage RR:
  new argument3 = GPR[registerYe<<1];
  new argument2 = GPR[registerZe<<1];
  new RM = decode_riscv_float_rounding_mode(floatmode);
  for (I = 0; I < 4; I += 1) {
    result1.16[I] = f16_sub(RM, argument3.16[I], argument2.16[I])
  };
stage E4:
  GPR[registerWe<<1] = result1;
  commit_float_exception_flags();
\end{lstlisting}}
\newpage

\paragraph{Encoding} Format \textsf{ALU\_FHQADDO} (Section~\ref{sec:format-ALU-FHQADDO}), \verb|alucode3 = 001|\\

\bitfields{ 1/P, 2/11, 1/1, 1/0, 3/floatmode, 5/registerWo, 1/1, 2/11, 3/001, 1/1, 5/registerYo, 1/--, 5/registerZo, 1/1 }


\paragraph{Syntax} \texttt{fsbfhq}\emph{floatmode} \emph{registerWo} = \emph{registerZo}, \emph{registerYo}

\paragraph{Description} The \emph{registerZo} and the \emph{registerYo} are interpreted as four binary 16 floating-point numbers packed into 64 bits. The \emph{registerZo} half words are subtracted from the \emph{registerYo} half words. The four resulting half words are packed and stored into the \emph{registerWo}.


\paragraph{Modifier(s)}
\noindent\begin{itemize}
\item floatmode (cf. floatmode in section \ref{par:modifier-floatmode} on page \pageref{par:modifier-floatmode})
\end{itemize}

\paragraph{Execution} {Reservation table \textsf{ALU\_LITE} (Section~\ref{sec:reservation-ALU-LITE})}
~
{\begin{lstlisting}
stage ID:
  new floatmode = _ZX_3(floatmode);
stage RR:
  new argument3 = GPR[(registerYo<<1)+1];
  new argument2 = GPR[(registerZo<<1)+1];
  new RM = decode_riscv_float_rounding_mode(floatmode);
  for (I = 0; I < 4; I += 1) {
    result1.16[I] = f16_sub(RM, argument3.16[I], argument2.16[I])
  };
stage E4:
  GPR[(registerWo<<1)+1] = result1;
  commit_float_exception_flags();
\end{lstlisting}}
\newpage
\subsubsection[{\small FSBFW: Floating-Point Subtract From Word}]{FSBFW\hfill Floating-Point Subtract From Word}
\label{instr:FSBFW}\index{fsbfw}
 
\paragraph{Encoding} Format \textsf{ALU\_FWADD} (Section~\ref{sec:format-ALU-FWADD}), \verb|alucode3 = 001|\\

\bitfields{ 1/P, 2/11, 1/1, 1/0, 3/floatmode, 6/registerW, 2/01, 3/001, 1/0, 6/registerY, 6/registerZ }


\paragraph{Syntax} \texttt{fsbfw}\emph{floatmode} \emph{registerW} = \emph{registerZ}, \emph{registerY}

\paragraph{Description} The \emph{registerZ} is subtracted from the \emph{registerY}, assuming binary 32 floating-point numbers. The result is rounded to the binary 32 floating-point format according to the rounding mode and stored into the \emph{registerW}. This instruction may raise exception bits in the CS register.


\paragraph{Modifier(s)}
\noindent\begin{itemize}
\item floatmode (cf. floatmode in section \ref{par:modifier-floatmode} on page \pageref{par:modifier-floatmode})
\end{itemize}

\paragraph{Execution} {Reservation table \textsf{ALU\_LITE} (Section~\ref{sec:reservation-ALU-LITE})}
~
{\begin{lstlisting}
stage ID:
  new floatmode = _ZX_3(floatmode);
stage RR:
  new argument3 = GPR[registerY];
  new argument2 = GPR[registerZ];
  new RM = decode_riscv_float_rounding_mode(floatmode);
  new result1 = f32_sub(RM, argument3, argument2);
stage E4:
  GPR[registerW] = result1;
  commit_float_exception_flags();
\end{lstlisting}}
\newpage
\subsubsection[{\small FSBFWC: Floating-Point Subtract from Word Complex}]{FSBFWC\hfill Floating-Point Subtract from Word Complex}
\label{instr:FSBFWC}\index{fsbfwc}
 
\paragraph{Encoding} Format \textsf{ALU\_FWCOP} (Section~\ref{sec:format-ALU-FWCOP}), \verb|alucode6 = 01|\\

\bitfields{ 1/P, 2/11, 1/1, 1/imultiply, 3/floatmode, 6/registerW, 2/00, 2/01, 1/conjugate, 1/1, 6/registerY, 6/registerZ }


\paragraph{Syntax} \texttt{fsbfwc}\emph{conjugate}\emph{imultiply}\emph{floatmode} \emph{registerW} = \emph{registerZ}, \emph{registerY}

\paragraph{Description} The \emph{registerY} and the \emph{registerZ} are interpreted as binary 32 floating-point complex numbers. If \emph{conjugate} is set, the complex conjugate of the \emph{registerZ} is used in the operation. The real and imaginary parts are respectively subtracted from and rounded according to the rounding mode. If \emph{imultiply} is set, the result is multiplied by the imaginary unit. The resulting binary 32 floating-point complex number is stored back into the \emph{registerW}. This instruction may raise exception bits in the CS register.


\paragraph{Modifier(s)}
\noindent\begin{itemize}
\item floatmode (cf. floatmode in section \ref{par:modifier-floatmode} on page \pageref{par:modifier-floatmode})
\item imultiply (cf. imultiply in section \ref{par:modifier-imultiply} on page \pageref{par:modifier-imultiply})
\item conjugate (cf. conjugate in section \ref{par:modifier-conjugate} on page \pageref{par:modifier-conjugate})
\end{itemize}

\paragraph{Execution} {Reservation table \textsf{ALU\_LITE} (Section~\ref{sec:reservation-ALU-LITE})}
~
{\begin{lstlisting}
stage ID:
  new conjugate = _ZX_1(conjugate);
  new imultiply = _ZX_1(imultiply);
  new floatmode = _ZX_3(floatmode);
stage RR:
  new argument3 = GPR[registerY];
  new argument2 = GPR[registerZ];
  new argument1 = GPR[registerW];
  new RM = decode_riscv_float_rounding_mode(floatmode);
  new fconj = conjugate ? 0x80000000 : 0;
  result1.32[0] = f32_sub(RM, argument3.32[0], argument2.32[0]);
  result1.32[1] = f32_sub(RM, argument3.32[1], argument2.32[1] ^ fconj);
  if (imultiply) {
    new result1 = _SWAP_32(result1);
    result1.32[0] = result1.32[0] ^ 0x80000000;
  }
stage E4:
  GPR[registerW] = result1;
  commit_float_exception_flags();
\end{lstlisting}}
\newpage
\subsubsection[{\small FSBFWP: Floating-Point Subtract From Word Pair}]{FSBFWP\hfill Floating-Point Subtract From Word Pair}
\label{instr:FSBFWP}\index{fsbfwp}
 
\paragraph{Encoding} Format \textsf{ALU\_FWPADDE} (Section~\ref{sec:format-ALU-FWPADDE}), \verb|alucode3 = 001|\\

\bitfields{ 1/P, 2/11, 1/1, 1/0, 3/floatmode, 5/registerWe, 1/0, 2/11, 3/001, 1/1, 5/registerYe, 1/--, 5/registerZe, 1/0 }


\paragraph{Syntax} \texttt{fsbfwp}\emph{floatmode} \emph{registerWe} = \emph{registerZe}, \emph{registerYe}

\paragraph{Description} The \emph{registerZe} and the \emph{registerYe} are interpreted as two binary 32 floating-point numbers packed into 64 bits. The \emph{registerZe} words are subtracted from to the \emph{registerYe} words. The two resulting words are packed and stored into the \emph{registerWe}.


\paragraph{Modifier(s)}
\noindent\begin{itemize}
\item floatmode (cf. floatmode in section \ref{par:modifier-floatmode} on page \pageref{par:modifier-floatmode})
\end{itemize}

\paragraph{Execution} {Reservation table \textsf{ALU\_LITE} (Section~\ref{sec:reservation-ALU-LITE})}
~
{\begin{lstlisting}
stage ID:
  new floatmode = _ZX_3(floatmode);
stage RR:
  new argument3 = GPR[registerYe<<1];
  new argument2 = GPR[registerZe<<1];
  new RM = decode_riscv_float_rounding_mode(floatmode);
  for (I = 0; I < 2; I += 1) {
    result1.32[I] = f32_sub(RM, argument3.32[I], argument2.32[I])
  };
stage E4:
  GPR[registerWe<<1] = result1;
  commit_float_exception_flags();
\end{lstlisting}}
\newpage

\paragraph{Encoding} Format \textsf{ALU\_FWPADDO} (Section~\ref{sec:format-ALU-FWPADDO}), \verb|alucode3 = 001|\\

\bitfields{ 1/P, 2/11, 1/1, 1/0, 3/floatmode, 5/registerWo, 1/1, 2/11, 3/001, 1/1, 5/registerYo, 1/--, 5/registerZo, 1/0 }


\paragraph{Syntax} \texttt{fsbfwp}\emph{floatmode} \emph{registerWo} = \emph{registerZo}, \emph{registerYo}

\paragraph{Description} The \emph{registerZo} and the \emph{registerYo} are interpreted as two binary 32 floating-point numbers packed into 64 bits. The \emph{registerZo} words are subtracted from to the \emph{registerYo} words. The two resulting words are packed and stored into the \emph{registerWo}.


\paragraph{Modifier(s)}
\noindent\begin{itemize}
\item floatmode (cf. floatmode in section \ref{par:modifier-floatmode} on page \pageref{par:modifier-floatmode})
\end{itemize}

\paragraph{Execution} {Reservation table \textsf{ALU\_LITE} (Section~\ref{sec:reservation-ALU-LITE})}
~
{\begin{lstlisting}
stage ID:
  new floatmode = _ZX_3(floatmode);
stage RR:
  new argument3 = GPR[(registerYo<<1)+1];
  new argument2 = GPR[(registerZo<<1)+1];
  new RM = decode_riscv_float_rounding_mode(floatmode);
  for (I = 0; I < 2; I += 1) {
    result1.32[I] = f32_sub(RM, argument3.32[I], argument2.32[I])
  };
stage E4:
  GPR[(registerWo<<1)+1] = result1;
  commit_float_exception_flags();
\end{lstlisting}}
\newpage
\subsubsection[{\small FSIGND: Floating-Point Sign Transfer Double Word}]{FSIGND\hfill Floating-Point Sign Transfer Double Word}
\label{instr:FSIGND}\index{fsignd}
 
\paragraph{Encoding} Format \textsf{ALU\_FDMMWR} (Section~\ref{sec:format-ALU-FDMMWR}), \verb|floatmode = 100|\\

\bitfields{ 1/P, 2/11, 1/1, 1/1, 3/100, 6/registerW, 2/00, 3/000, 1/0, 6/registerY, 6/registerZ }


\paragraph{Syntax} \texttt{fsignd} \emph{registerW} = \emph{registerZ}, \emph{registerY}

\paragraph{Description} The \emph{registerZ} and the \emph{registerY} are interpreted as binary 64 floating-point numbers. The non-sign bits of the \emph{registerZ} are concatenated with the sign bit of the \emph{registerY}. The result is stored into the \emph{registerW}.


\paragraph{Execution} {Reservation table \textsf{ALU\_LITE} (Section~\ref{sec:reservation-ALU-LITE})}
~
{\begin{lstlisting}
stage RR:
  new argument3 = GPR[registerY];
  new argument2 = GPR[registerZ];
  new result1 = (argument2 & 0x3FFFFFFFFFFFFFFF) | (argument3 & 0x8000000000000000);
stage E1:
  GPR[registerW] = result1;
\end{lstlisting}}
\newpage
\subsubsection[{\small FSIGNH: Floating-Point Sign Transfer Half Word}]{FSIGNH\hfill Floating-Point Sign Transfer Half Word}
\label{instr:FSIGNH}\index{fsignh}
 
\paragraph{Encoding} Format \textsf{ALU\_FHMMWR} (Section~\ref{sec:format-ALU-FHMMWR}), \verb|floatmode = 100|\\

\bitfields{ 1/P, 2/11, 1/1, 1/1, 3/100, 6/registerW, 2/01, 3/011, 1/1, 6/registerY, 6/registerZ }


\paragraph{Syntax} \texttt{fsignh} \emph{registerW} = \emph{registerZ}, \emph{registerY}

\paragraph{Description} The \emph{registerZ} and the \emph{registerY} are interpreted as binary 16 floating-point numbers. The non-sign bits of the \emph{registerZ} are concatenated with the sign bit of the \emph{registerY}. The result is stored into the \emph{registerW}.


\paragraph{Execution} {Reservation table \textsf{ALU\_LITE} (Section~\ref{sec:reservation-ALU-LITE})}
~
{\begin{lstlisting}
stage RR:
  new argument3 = GPR[registerY];
  new argument2 = GPR[registerZ];
  new result1 = (argument2 & 0x3FFF) | (argument3 & 0x8000);
stage E1:
  GPR[registerW] = result1;
\end{lstlisting}}
\newpage
\subsubsection[{\small FSIGNHQ: Floating-Point Sign Transfer Half Word Quadruple}]{FSIGNHQ\hfill Floating-Point Sign Transfer Half Word Quadruple}
\label{instr:FSIGNHQ}\index{fsignhq}
 
\paragraph{Encoding} Format \textsf{ALU\_FHQMMWRE} (Section~\ref{sec:format-ALU-FHQMMWRE}), \verb|floatmode = 100|\\

\bitfields{ 1/P, 2/11, 1/1, 1/1, 3/100, 5/registerWe, 1/0, 2/11, 3/000, 1/1, 5/registerYe, 1/0, 5/registerZe, 1/1 }


\paragraph{Syntax} \texttt{fsignhq} \emph{registerWe} = \emph{registerZe}, \emph{registerYe}

\paragraph{Description} The \emph{registerZe} and the \emph{registerYe} are interpreted as four binary 16 floating-point numbers packed into 64 bits. The non-sign bits of the \emph{registerZe} half words are concatenated with the sign bit of the \emph{registerYe} half words. The four resulting half words are packed and stored into the \emph{registerWe}.


\paragraph{Execution} {Reservation table \textsf{ALU\_LITE} (Section~\ref{sec:reservation-ALU-LITE})}
~
{\begin{lstlisting}
stage RR:
  new argument3 = GPR[registerYe<<1];
  new argument2 = GPR[registerZe<<1];
  for (I = 0; I < 4; I += 1) {
    result1.16[I] = (argument2.16[I] & 0x3FFF) | (argument3.16[I] & 0x8000)
  };
stage E1:
  GPR[registerWe<<1] = result1;
\end{lstlisting}}
\newpage

\paragraph{Encoding} Format \textsf{ALU\_FHQMMWRO} (Section~\ref{sec:format-ALU-FHQMMWRO}), \verb|floatmode = 100|\\

\bitfields{ 1/P, 2/11, 1/1, 1/1, 3/100, 5/registerWo, 1/1, 2/11, 3/000, 1/1, 5/registerYo, 1/0, 5/registerZo, 1/1 }


\paragraph{Syntax} \texttt{fsignhq} \emph{registerWo} = \emph{registerZo}, \emph{registerYo}

\paragraph{Description} The \emph{registerZo} and the \emph{registerYo} are interpreted as four binary 16 floating-point numbers packed into 64 bits. The non-sign bits of the \emph{registerZo} half words are concatenated with the sign bit of the \emph{registerYo} half words. The four resulting half words are packed and stored into the \emph{registerWo}.


\paragraph{Execution} {Reservation table \textsf{ALU\_LITE} (Section~\ref{sec:reservation-ALU-LITE})}
~
{\begin{lstlisting}
stage RR:
  new argument3 = GPR[(registerYo<<1)+1];
  new argument2 = GPR[(registerZo<<1)+1];
  for (I = 0; I < 4; I += 1) {
    result1.16[I] = (argument2.16[I] & 0x3FFF) | (argument3.16[I] & 0x8000)
  };
stage E1:
  GPR[(registerWo<<1)+1] = result1;
\end{lstlisting}}
\newpage
\subsubsection[{\small FSIGNMD: Floating-Point Sign Combine Transfer Double Word}]{FSIGNMD\hfill Floating-Point Sign Combine Transfer Double Word}
\label{instr:FSIGNMD}\index{fsignmd}
 
\paragraph{Encoding} Format \textsf{ALU\_FDMMWR} (Section~\ref{sec:format-ALU-FDMMWR}), \verb|floatmode = 110|\\

\bitfields{ 1/P, 2/11, 1/1, 1/1, 3/110, 6/registerW, 2/00, 3/000, 1/0, 6/registerY, 6/registerZ }


\paragraph{Syntax} \texttt{fsignmd} \emph{registerW} = \emph{registerZ}, \emph{registerY}

\paragraph{Description} The \emph{registerZ} and the \emph{registerY} are interpreted as binary 64 floating-point numbers. The sign bits of the \emph{registerZ} is combined with the sign bit of the \emph{registerY} as if multiplying. The result is stored into the \emph{registerW}.


\paragraph{Execution} {Reservation table \textsf{ALU\_LITE} (Section~\ref{sec:reservation-ALU-LITE})}
~
{\begin{lstlisting}
stage RR:
  new argument3 = GPR[registerY];
  new argument2 = GPR[registerZ];
  new result1 = argument2 ^ (~argument3 & 0x8000000000000000);
stage E1:
  GPR[registerW] = result1;
\end{lstlisting}}
\newpage
\subsubsection[{\small FSIGNMH: Floating-Point Sign Combine Transfer Half Word}]{FSIGNMH\hfill Floating-Point Sign Combine Transfer Half Word}
\label{instr:FSIGNMH}\index{fsignmh}
 
\paragraph{Encoding} Format \textsf{ALU\_FHMMWR} (Section~\ref{sec:format-ALU-FHMMWR}), \verb|floatmode = 110|\\

\bitfields{ 1/P, 2/11, 1/1, 1/1, 3/110, 6/registerW, 2/01, 3/011, 1/1, 6/registerY, 6/registerZ }


\paragraph{Syntax} \texttt{fsignmh} \emph{registerW} = \emph{registerZ}, \emph{registerY}

\paragraph{Description} The \emph{registerZ} and the \emph{registerY} are interpreted as binary 16 floating-point numbers. The sign bits of the \emph{registerZ} is combined with the sign bit of the \emph{registerY} as if multiplying. The result is stored into the \emph{registerW}.


\paragraph{Execution} {Reservation table \textsf{ALU\_LITE} (Section~\ref{sec:reservation-ALU-LITE})}
~
{\begin{lstlisting}
stage RR:
  new argument3 = GPR[registerY];
  new argument2 = GPR[registerZ];
  new result1 = argument2 ^ (~argument3 & 0x8000);
stage E1:
  GPR[registerW] = result1;
\end{lstlisting}}
\newpage
\subsubsection[{\small FSIGNMHQ: Floating-Point Sign Combine Transfer Half Word Quadruple}]{FSIGNMHQ\hfill Floating-Point Sign Combine Transfer Half Word Quadruple}
\label{instr:FSIGNMHQ}\index{fsignmhq}
 
\paragraph{Encoding} Format \textsf{ALU\_FHQMMWRE} (Section~\ref{sec:format-ALU-FHQMMWRE}), \verb|floatmode = 110|\\

\bitfields{ 1/P, 2/11, 1/1, 1/1, 3/110, 5/registerWe, 1/0, 2/11, 3/000, 1/1, 5/registerYe, 1/0, 5/registerZe, 1/1 }


\paragraph{Syntax} \texttt{fsignmhq} \emph{registerWe} = \emph{registerZe}, \emph{registerYe}

\paragraph{Description} The \emph{registerZe} and the \emph{registerYe} are interpreted as four binary 16 floating-point numbers packed into 64 bits. The sign bits of the \emph{registerZe} half words are combined with the sign bit of the \emph{registerYe} half words as if multiplying. The four resulting half words are packed and stored into the \emph{registerWe}.


\paragraph{Execution} {Reservation table \textsf{ALU\_LITE} (Section~\ref{sec:reservation-ALU-LITE})}
~
{\begin{lstlisting}
stage RR:
  new argument3 = GPR[registerYe<<1];
  new argument2 = GPR[registerZe<<1];
  for (I = 0; I < 4; I += 1) {
    result1.16[I] = argument2.16[I] ^ (~argument3.16[I] & 0x8000000000000000)
  };
stage E1:
  GPR[registerWe<<1] = result1;
\end{lstlisting}}
\newpage

\paragraph{Encoding} Format \textsf{ALU\_FHQMMWRO} (Section~\ref{sec:format-ALU-FHQMMWRO}), \verb|floatmode = 110|\\

\bitfields{ 1/P, 2/11, 1/1, 1/1, 3/110, 5/registerWo, 1/1, 2/11, 3/000, 1/1, 5/registerYo, 1/0, 5/registerZo, 1/1 }


\paragraph{Syntax} \texttt{fsignmhq} \emph{registerWo} = \emph{registerZo}, \emph{registerYo}

\paragraph{Description} The \emph{registerZo} and the \emph{registerYo} are interpreted as four binary 16 floating-point numbers packed into 64 bits. The sign bits of the \emph{registerZo} half words are combined with the sign bit of the \emph{registerYo} half words as if multiplying. The four resulting half words are packed and stored into the \emph{registerWo}.


\paragraph{Execution} {Reservation table \textsf{ALU\_LITE} (Section~\ref{sec:reservation-ALU-LITE})}
~
{\begin{lstlisting}
stage RR:
  new argument3 = GPR[(registerYo<<1)+1];
  new argument2 = GPR[(registerZo<<1)+1];
  for (I = 0; I < 4; I += 1) {
    result1.16[I] = argument2.16[I] ^ (~argument3.16[I] & 0x8000000000000000)
  };
stage E1:
  GPR[(registerWo<<1)+1] = result1;
\end{lstlisting}}
\newpage
\subsubsection[{\small FSIGNMW: Floating-Point Sign Combine Transfer Word}]{FSIGNMW\hfill Floating-Point Sign Combine Transfer Word}
\label{instr:FSIGNMW}\index{fsignmw}
 
\paragraph{Encoding} Format \textsf{ALU\_FWMMWR} (Section~\ref{sec:format-ALU-FWMMWR}), \verb|floatmode = 110|\\

\bitfields{ 1/P, 2/11, 1/1, 1/1, 3/110, 6/registerW, 2/01, 3/000, 1/0, 6/registerY, 6/registerZ }


\paragraph{Syntax} \texttt{fsignmw} \emph{registerW} = \emph{registerZ}, \emph{registerY}

\paragraph{Description} The \emph{registerZ} and the \emph{registerY} are interpreted as binary 32 floating-point numbers. The sign bits of the \emph{registerZ} is combined with the sign bit of the \emph{registerY} as if multiplying. The result is stored into the \emph{registerW}.


\paragraph{Execution} {Reservation table \textsf{ALU\_LITE} (Section~\ref{sec:reservation-ALU-LITE})}
~
{\begin{lstlisting}
stage RR:
  new argument3 = GPR[registerY];
  new argument2 = GPR[registerZ];
  new result1 = argument2 ^ (~argument3 & 0x80000000);
stage E1:
  GPR[registerW] = result1;
\end{lstlisting}}
\newpage
\subsubsection[{\small FSIGNMWP: Floating-Point Sign Combine Transfer Double Word Pair}]{FSIGNMWP\hfill Floating-Point Sign Combine Transfer Double Word Pair}
\label{instr:FSIGNMWP}\index{fsignmwp}
 
\paragraph{Encoding} Format \textsf{ALU\_FWPMMWRE} (Section~\ref{sec:format-ALU-FWPMMWRE}), \verb|floatmode = 110|\\

\bitfields{ 1/P, 2/11, 1/1, 1/1, 3/110, 5/registerWe, 1/0, 2/11, 3/000, 1/1, 5/registerYe, 1/0, 5/registerZe, 1/0 }


\paragraph{Syntax} \texttt{fsignmwp} \emph{registerWe} = \emph{registerZe}, \emph{registerYe}

\paragraph{Description} The \emph{registerZe} and the \emph{registerYe} are interpreted as two binary 32 floating-point numbers packed into 128 bits. The sign bits of the \emph{registerZe} words are combined with the sign bit of the \emph{registerYe} words as if multiplying. The two resulting words are packed and stored into the \emph{registerWe}.


\paragraph{Execution} {Reservation table \textsf{ALU\_LITE} (Section~\ref{sec:reservation-ALU-LITE})}
~
{\begin{lstlisting}
stage RR:
  new argument3 = GPR[registerYe<<1];
  new argument2 = GPR[registerZe<<1];
  for (I = 0; I < 2; I += 1) {
    result1.32[I] = argument2.32[I] ^ (~argument3.32[I] & 0x80000000)
  };
stage E1:
  GPR[registerWe<<1] = result1;
\end{lstlisting}}
\newpage

\paragraph{Encoding} Format \textsf{ALU\_FWPMMWRO} (Section~\ref{sec:format-ALU-FWPMMWRO}), \verb|floatmode = 110|\\

\bitfields{ 1/P, 2/11, 1/1, 1/1, 3/110, 5/registerWo, 1/1, 2/11, 3/000, 1/1, 5/registerYo, 1/0, 5/registerZo, 1/0 }


\paragraph{Syntax} \texttt{fsignmwp} \emph{registerWo} = \emph{registerZo}, \emph{registerYo}

\paragraph{Description} The \emph{registerZo} and the \emph{registerYo} are interpreted as two binary 32 floating-point numbers packed into 128 bits. The sign bits of the \emph{registerZo} words are combined with the sign bit of the \emph{registerYo} words as if multiplying. The two resulting words are packed and stored into the \emph{registerWo}.


\paragraph{Execution} {Reservation table \textsf{ALU\_LITE} (Section~\ref{sec:reservation-ALU-LITE})}
~
{\begin{lstlisting}
stage RR:
  new argument3 = GPR[(registerYo<<1)+1];
  new argument2 = GPR[(registerZo<<1)+1];
  for (I = 0; I < 2; I += 1) {
    result1.32[I] = argument2.32[I] ^ (~argument3.32[I] & 0x80000000)
  };
stage E1:
  GPR[(registerWo<<1)+1] = result1;
\end{lstlisting}}
\newpage
\subsubsection[{\small FSIGNND: Floating-Point Sign Negate Transfer Double Word}]{FSIGNND\hfill Floating-Point Sign Negate Transfer Double Word}
\label{instr:FSIGNND}\index{fsignnd}
 
\paragraph{Encoding} Format \textsf{ALU\_FDMMWR} (Section~\ref{sec:format-ALU-FDMMWR}), \verb|floatmode = 101|\\

\bitfields{ 1/P, 2/11, 1/1, 1/1, 3/101, 6/registerW, 2/00, 3/000, 1/0, 6/registerY, 6/registerZ }


\paragraph{Syntax} \texttt{fsignnd} \emph{registerW} = \emph{registerZ}, \emph{registerY}

\paragraph{Description} The \emph{registerZ} and the \emph{registerY} are interpreted as binary 64 floating-point numbers. The non-sign bits of the \emph{registerZ} are concatenated with the negated sign bit of the \emph{registerY}. The result is stored into the \emph{registerW}.


\paragraph{Execution} {Reservation table \textsf{ALU\_LITE} (Section~\ref{sec:reservation-ALU-LITE})}
~
{\begin{lstlisting}
stage RR:
  new argument3 = GPR[registerY];
  new argument2 = GPR[registerZ];
  new result1 = (argument2 & 0x3FFFFFFFFFFFFFFF) | (~argument3 & 0x8000000000000000);
stage E1:
  GPR[registerW] = result1;
\end{lstlisting}}
\newpage
\subsubsection[{\small FSIGNNH: Floating-Point Sign Negate Transfer Half Word}]{FSIGNNH\hfill Floating-Point Sign Negate Transfer Half Word}
\label{instr:FSIGNNH}\index{fsignnh}
 
\paragraph{Encoding} Format \textsf{ALU\_FHMMWR} (Section~\ref{sec:format-ALU-FHMMWR}), \verb|floatmode = 101|\\

\bitfields{ 1/P, 2/11, 1/1, 1/1, 3/101, 6/registerW, 2/01, 3/011, 1/1, 6/registerY, 6/registerZ }


\paragraph{Syntax} \texttt{fsignnh} \emph{registerW} = \emph{registerZ}, \emph{registerY}

\paragraph{Description} The \emph{registerZ} and the \emph{registerY} are interpreted as binary 16 floating-point numbers. The non-sign bits of the \emph{registerZ} are concatenated with the negated sign bit of the \emph{registerY}. The result is stored into the \emph{registerW}.


\paragraph{Execution} {Reservation table \textsf{ALU\_LITE} (Section~\ref{sec:reservation-ALU-LITE})}
~
{\begin{lstlisting}
stage RR:
  new argument3 = GPR[registerY];
  new argument2 = GPR[registerZ];
  new result1 = (argument2 & 0x3FFF) | (~argument3 & 0x8000);
stage E1:
  GPR[registerW] = result1;
\end{lstlisting}}
\newpage
\subsubsection[{\small FSIGNNHQ: Floating-Point Sign Negate Transfer Half Word Quadruple}]{FSIGNNHQ\hfill Floating-Point Sign Negate Transfer Half Word Quadruple}
\label{instr:FSIGNNHQ}\index{fsignnhq}
 
\paragraph{Encoding} Format \textsf{ALU\_FHQMMWRE} (Section~\ref{sec:format-ALU-FHQMMWRE}), \verb|floatmode = 101|\\

\bitfields{ 1/P, 2/11, 1/1, 1/1, 3/101, 5/registerWe, 1/0, 2/11, 3/000, 1/1, 5/registerYe, 1/0, 5/registerZe, 1/1 }


\paragraph{Syntax} \texttt{fsignnhq} \emph{registerWe} = \emph{registerZe}, \emph{registerYe}

\paragraph{Description} The \emph{registerZe} and the \emph{registerYe} are interpreted as four binary 16 floating-point numbers packed into 64 bits. The non-sign bits of the \emph{registerZe} half words are concatenated with the negated sign bit of the \emph{registerYe} half words. The four resulting half words are packed and stored into the \emph{registerWe}.


\paragraph{Execution} {Reservation table \textsf{ALU\_LITE} (Section~\ref{sec:reservation-ALU-LITE})}
~
{\begin{lstlisting}
stage RR:
  new argument3 = GPR[registerYe<<1];
  new argument2 = GPR[registerZe<<1];
  for (I = 0; I < 4; I += 1) {
    result1.16[I] = (argument2.16[I] & 0x3FFF) | (~argument3.16[I] & 0x8000)
  };
stage E1:
  GPR[registerWe<<1] = result1;
\end{lstlisting}}
\newpage

\paragraph{Encoding} Format \textsf{ALU\_FHQMMWRO} (Section~\ref{sec:format-ALU-FHQMMWRO}), \verb|floatmode = 101|\\

\bitfields{ 1/P, 2/11, 1/1, 1/1, 3/101, 5/registerWo, 1/1, 2/11, 3/000, 1/1, 5/registerYo, 1/0, 5/registerZo, 1/1 }


\paragraph{Syntax} \texttt{fsignnhq} \emph{registerWo} = \emph{registerZo}, \emph{registerYo}

\paragraph{Description} The \emph{registerZo} and the \emph{registerYo} are interpreted as four binary 16 floating-point numbers packed into 64 bits. The non-sign bits of the \emph{registerZo} half words are concatenated with the negated sign bit of the \emph{registerYo} half words. The four resulting half words are packed and stored into the \emph{registerWo}.


\paragraph{Execution} {Reservation table \textsf{ALU\_LITE} (Section~\ref{sec:reservation-ALU-LITE})}
~
{\begin{lstlisting}
stage RR:
  new argument3 = GPR[(registerYo<<1)+1];
  new argument2 = GPR[(registerZo<<1)+1];
  for (I = 0; I < 4; I += 1) {
    result1.16[I] = (argument2.16[I] & 0x3FFF) | (~argument3.16[I] & 0x8000)
  };
stage E1:
  GPR[(registerWo<<1)+1] = result1;
\end{lstlisting}}
\newpage
\subsubsection[{\small FSIGNNW: Floating-Point Sign Negate Transfer Word}]{FSIGNNW\hfill Floating-Point Sign Negate Transfer Word}
\label{instr:FSIGNNW}\index{fsignnw}
 
\paragraph{Encoding} Format \textsf{ALU\_FWMMWR} (Section~\ref{sec:format-ALU-FWMMWR}), \verb|floatmode = 101|\\

\bitfields{ 1/P, 2/11, 1/1, 1/1, 3/101, 6/registerW, 2/01, 3/000, 1/0, 6/registerY, 6/registerZ }


\paragraph{Syntax} \texttt{fsignnw} \emph{registerW} = \emph{registerZ}, \emph{registerY}

\paragraph{Description} The \emph{registerZ} and the \emph{registerY} are interpreted as binary 32 floating-point numbers. The non-sign bits of the \emph{registerZ} are concatenated with the negated sign bit of the \emph{registerY}. The result is stored into the \emph{registerW}.


\paragraph{Execution} {Reservation table \textsf{ALU\_LITE} (Section~\ref{sec:reservation-ALU-LITE})}
~
{\begin{lstlisting}
stage RR:
  new argument3 = GPR[registerY];
  new argument2 = GPR[registerZ];
  new result1 = (argument2 & 0x3FFFFFFF) | (~argument3 & 0x80000000);
stage E1:
  GPR[registerW] = result1;
\end{lstlisting}}
\newpage
\subsubsection[{\small FSIGNNWP: Floating-Point Sign Negate Transfer Word Pair}]{FSIGNNWP\hfill Floating-Point Sign Negate Transfer Word Pair}
\label{instr:FSIGNNWP}\index{fsignnwp}
 
\paragraph{Encoding} Format \textsf{ALU\_FWPMMWRE} (Section~\ref{sec:format-ALU-FWPMMWRE}), \verb|floatmode = 101|\\

\bitfields{ 1/P, 2/11, 1/1, 1/1, 3/101, 5/registerWe, 1/0, 2/11, 3/000, 1/1, 5/registerYe, 1/0, 5/registerZe, 1/0 }


\paragraph{Syntax} \texttt{fsignnwp} \emph{registerWe} = \emph{registerZe}, \emph{registerYe}

\paragraph{Description} The \emph{registerZe} and the \emph{registerYe} are interpreted as two binary 32 floating-point numbers packed into 64 bits. The non-sign bits of the \emph{registerZe} words are concatenated with the negated sign bit of the \emph{registerYe} words. The two resulting words are packed and stored into the \emph{registerWe}.


\paragraph{Execution} {Reservation table \textsf{ALU\_LITE} (Section~\ref{sec:reservation-ALU-LITE})}
~
{\begin{lstlisting}
stage RR:
  new argument3 = GPR[registerYe<<1];
  new argument2 = GPR[registerZe<<1];
  for (I = 0; I < 2; I += 1) {
    result1.32[I] = (argument2.32[I] & 0x3FFFFFFF) | (~argument3.32[I] & 0x80000000)
  };
stage E1:
  GPR[registerWe<<1] = result1;
\end{lstlisting}}
\newpage

\paragraph{Encoding} Format \textsf{ALU\_FWPMMWRO} (Section~\ref{sec:format-ALU-FWPMMWRO}), \verb|floatmode = 101|\\

\bitfields{ 1/P, 2/11, 1/1, 1/1, 3/101, 5/registerWo, 1/1, 2/11, 3/000, 1/1, 5/registerYo, 1/0, 5/registerZo, 1/0 }


\paragraph{Syntax} \texttt{fsignnwp} \emph{registerWo} = \emph{registerZo}, \emph{registerYo}

\paragraph{Description} The \emph{registerZo} and the \emph{registerYo} are interpreted as two binary 32 floating-point numbers packed into 64 bits. The non-sign bits of the \emph{registerZo} words are concatenated with the negated sign bit of the \emph{registerYo} words. The two resulting words are packed and stored into the \emph{registerWo}.


\paragraph{Execution} {Reservation table \textsf{ALU\_LITE} (Section~\ref{sec:reservation-ALU-LITE})}
~
{\begin{lstlisting}
stage RR:
  new argument3 = GPR[(registerYo<<1)+1];
  new argument2 = GPR[(registerZo<<1)+1];
  for (I = 0; I < 2; I += 1) {
    result1.32[I] = (argument2.32[I] & 0x3FFFFFFF) | (~argument3.32[I] & 0x80000000)
  };
stage E1:
  GPR[(registerWo<<1)+1] = result1;
\end{lstlisting}}
\newpage
\subsubsection[{\small FSIGNW: Floating-Point Sign Transfer Word}]{FSIGNW\hfill Floating-Point Sign Transfer Word}
\label{instr:FSIGNW}\index{fsignw}
 
\paragraph{Encoding} Format \textsf{ALU\_FWMMWR} (Section~\ref{sec:format-ALU-FWMMWR}), \verb|floatmode = 100|\\

\bitfields{ 1/P, 2/11, 1/1, 1/1, 3/100, 6/registerW, 2/01, 3/000, 1/0, 6/registerY, 6/registerZ }


\paragraph{Syntax} \texttt{fsignw} \emph{registerW} = \emph{registerZ}, \emph{registerY}

\paragraph{Description} The \emph{registerZ} and the \emph{registerY} are interpreted as binary 32 floating-point numbers. The non-sign bits of the \emph{registerZ} are concatenated with the sign bit of the \emph{registerY}. The result is stored into the \emph{registerW}.


\paragraph{Execution} {Reservation table \textsf{ALU\_LITE} (Section~\ref{sec:reservation-ALU-LITE})}
~
{\begin{lstlisting}
stage RR:
  new argument3 = GPR[registerY];
  new argument2 = GPR[registerZ];
  new result1 = (argument2 & 0x3FFFFFFF) | (argument3 & 0x80000000);
stage E1:
  GPR[registerW] = result1;
\end{lstlisting}}
\newpage
\subsubsection[{\small FSIGNWP: Floating-Point Sign Transfer Word Pair}]{FSIGNWP\hfill Floating-Point Sign Transfer Word Pair}
\label{instr:FSIGNWP}\index{fsignwp}
 
\paragraph{Encoding} Format \textsf{ALU\_FWPMMWRE} (Section~\ref{sec:format-ALU-FWPMMWRE}), \verb|floatmode = 100|\\

\bitfields{ 1/P, 2/11, 1/1, 1/1, 3/100, 5/registerWe, 1/0, 2/11, 3/000, 1/1, 5/registerYe, 1/0, 5/registerZe, 1/0 }


\paragraph{Syntax} \texttt{fsignwp} \emph{registerWe} = \emph{registerZe}, \emph{registerYe}

\paragraph{Description} The \emph{registerZe} and the \emph{registerYe} are interpreted as two binary 32 floating-point numbers packed into 64 bits. The non-sign bits of the \emph{registerZe} words are concatenated with the sign bit of the \emph{registerYe} words. The two resulting words are packed and stored into the \emph{registerWe}.


\paragraph{Execution} {Reservation table \textsf{ALU\_LITE} (Section~\ref{sec:reservation-ALU-LITE})}
~
{\begin{lstlisting}
stage RR:
  new argument3 = GPR[registerYe<<1];
  new argument2 = GPR[registerZe<<1];
  for (I = 0; I < 2; I += 1) {
    result1.32[I] = (argument2.32[I] & 0x3FFFFFFF) | (argument3.32[I] & 0x80000000)
  };
stage E1:
  GPR[registerWe<<1] = result1;
\end{lstlisting}}
\newpage

\paragraph{Encoding} Format \textsf{ALU\_FWPMMWRO} (Section~\ref{sec:format-ALU-FWPMMWRO}), \verb|floatmode = 100|\\

\bitfields{ 1/P, 2/11, 1/1, 1/1, 3/100, 5/registerWo, 1/1, 2/11, 3/000, 1/1, 5/registerYo, 1/0, 5/registerZo, 1/0 }


\paragraph{Syntax} \texttt{fsignwp} \emph{registerWo} = \emph{registerZo}, \emph{registerYo}

\paragraph{Description} The \emph{registerZo} and the \emph{registerYo} are interpreted as two binary 32 floating-point numbers packed into 64 bits. The non-sign bits of the \emph{registerZo} words are concatenated with the sign bit of the \emph{registerYo} words. The two resulting words are packed and stored into the \emph{registerWo}.


\paragraph{Execution} {Reservation table \textsf{ALU\_LITE} (Section~\ref{sec:reservation-ALU-LITE})}
~
{\begin{lstlisting}
stage RR:
  new argument3 = GPR[(registerYo<<1)+1];
  new argument2 = GPR[(registerZo<<1)+1];
  for (I = 0; I < 2; I += 1) {
    result1.32[I] = (argument2.32[I] & 0x3FFFFFFF) | (argument3.32[I] & 0x80000000)
  };
stage E1:
  GPR[(registerWo<<1)+1] = result1;
\end{lstlisting}}
\newpage
\subsubsection[{\small FSQRTD: Floating Point Square Root Double Word}]{FSQRTD\hfill Floating Point Square Root Double Word}
\label{instr:FSQRTD}\index{fsqrtd}
 
\paragraph{Encoding} Format \textsf{ALU\_FDMO0} (Section~\ref{sec:format-ALU-FDMO0}), \verb|fpucode6 = 1000|\\

\bitfields{ 1/P, 2/11, 1/1, 1/1, 3/floatmode, 6/registerW, 2/00, 3/001, 1/0, 4/1000, 1/--, 1/0, 6/registerZ }


\paragraph{Syntax} \texttt{fsqrtd}\emph{floatmode} \emph{registerW} = \emph{registerZ}

\paragraph{Description} The square root of the \emph{registerZ} interpreted as a 64-bit floating-point number is computed according to the rounding mode and stored into the \emph{registerW}. This instruction may raise exception bits in the CS register.


\paragraph{Modifier(s)}
\noindent\begin{itemize}
\item floatmode (cf. floatmode in section \ref{par:modifier-floatmode} on page \pageref{par:modifier-floatmode})
\end{itemize}

\paragraph{Execution} {Reservation table \textsf{ALU\_FULL} (Section~\ref{sec:reservation-ALU-FULL})}
~
{\begin{lstlisting}
stage ID:
  new floatmode = _ZX_3(floatmode);
stage RR:
  new argument2 = GPR[registerZ];
  new RM = decode_riscv_float_rounding_mode(floatmode);
  new result1 = f64_sqrt(RM, argument2);
stage E4:
  GPR[registerW] = result1;
  commit_float_exception_flags();
\end{lstlisting}}
\newpage
\subsubsection[{\small FSQRTH: Floating Point Square Root Half Word}]{FSQRTH\hfill Floating Point Square Root Half Word}
\label{instr:FSQRTH}\index{fsqrth}
 
\paragraph{Encoding} Format \textsf{ALU\_FHMO0} (Section~\ref{sec:format-ALU-FHMO0}), \verb|fpucode6 = 1000|\\

\bitfields{ 1/P, 2/11, 1/1, 1/1, 3/floatmode, 6/registerW, 2/01, 3/001, 1/1, 4/1000, 1/--, 1/0, 6/registerZ }


\paragraph{Syntax} \texttt{fsqrth}\emph{floatmode} \emph{registerW} = \emph{registerZ}

\paragraph{Description} The square root of the \emph{registerZ} interpreted as a 16-bit floating-point number is computed according to the rounding mode and stored into the \emph{registerW}. This instruction may raise exception bits in the CS register.


\paragraph{Modifier(s)}
\noindent\begin{itemize}
\item floatmode (cf. floatmode in section \ref{par:modifier-floatmode} on page \pageref{par:modifier-floatmode})
\end{itemize}

\paragraph{Execution} {Reservation table \textsf{ALU\_FULL} (Section~\ref{sec:reservation-ALU-FULL})}
~
{\begin{lstlisting}
stage ID:
  new floatmode = _ZX_3(floatmode);
stage RR:
  new argument2 = GPR[registerZ];
  new RM = decode_riscv_float_rounding_mode(floatmode);
  new result1 = f16_sqrt(RM, argument2);
stage SF:
  GPR[registerW] = result1;
  commit_float_exception_flags();
\end{lstlisting}}
\newpage
\subsubsection[{\small FSQRTW: Floating Point Square Root Word}]{FSQRTW\hfill Floating Point Square Root Word}
\label{instr:FSQRTW}\index{fsqrtw}
 
\paragraph{Encoding} Format \textsf{ALU\_FWMO0} (Section~\ref{sec:format-ALU-FWMO0}), \verb|fpucode6 = 1000|\\

\bitfields{ 1/P, 2/11, 1/1, 1/1, 3/floatmode, 6/registerW, 2/01, 3/001, 1/0, 4/1000, 1/--, 1/0, 6/registerZ }


\paragraph{Syntax} \texttt{fsqrtw}\emph{floatmode} \emph{registerW} = \emph{registerZ}

\paragraph{Description} The square root of the \emph{registerZ} interpreted as a 32-bit floating-point number is computed according to the rounding mode and stored into the \emph{registerW}. This instruction may raise exception bits in the CS register.


\paragraph{Modifier(s)}
\noindent\begin{itemize}
\item floatmode (cf. floatmode in section \ref{par:modifier-floatmode} on page \pageref{par:modifier-floatmode})
\end{itemize}

\paragraph{Execution} {Reservation table \textsf{ALU\_FULL} (Section~\ref{sec:reservation-ALU-FULL})}
~
{\begin{lstlisting}
stage ID:
  new floatmode = _ZX_3(floatmode);
stage RR:
  new argument2 = GPR[registerZ];
  new RM = decode_riscv_float_rounding_mode(floatmode);
  new result1 = f32_sqrt(RM, argument2);
stage SF:
  GPR[registerW] = result1;
  commit_float_exception_flags();
\end{lstlisting}}
\newpage
\subsubsection[{\small FSRECD: Floating-Point Seed for Reciprocal Double Word}]{FSRECD\hfill Floating-Point Seed for Reciprocal Double Word}
\label{instr:FSRECD}\index{fsrecd}
 
\paragraph{Encoding} Format \textsf{ALU\_FDMO1} (Section~\ref{sec:format-ALU-FDMO1}), \verb|fpucode6 = 0100|\\

\bitfields{ 1/P, 2/11, 1/1, 1/1, 3/111, 6/registerW, 2/00, 3/000, 1/0, 4/0100, 1/--, 1/0, 6/registerZ }


\paragraph{Syntax} \texttt{fsrecd} \emph{registerW} = \emph{registerZ}

\paragraph{Description} An approximation of the reciprocal of the \emph{registerZ}, assuming binary 64 floating-point numbers, is stored into the \emph{registerW}. This instruction may raise exception bits in the CS register.


\paragraph{Execution} {Reservation table \textsf{ALU\_LITE} (Section~\ref{sec:reservation-ALU-LITE})}
~
{\begin{lstlisting}
stage RR:
  new argument2 = GPR[registerZ];
  new result1 = f64_fast_rec(argument2);
stage E1:
  GPR[registerW] = result1;
\end{lstlisting}}
\newpage
\subsubsection[{\small FSRECW: Floating-Point Seed for Reciprocal Word}]{FSRECW\hfill Floating-Point Seed for Reciprocal Word}
\label{instr:FSRECW}\index{fsrecw}
 
\paragraph{Encoding} Format \textsf{ALU\_FWMO1} (Section~\ref{sec:format-ALU-FWMO1}), \verb|fpucode6 = 0100|\\

\bitfields{ 1/P, 2/11, 1/1, 1/1, 3/111, 6/registerW, 2/01, 3/000, 1/0, 4/0100, 1/--, 1/0, 6/registerZ }


\paragraph{Syntax} \texttt{fsrecw} \emph{registerW} = \emph{registerZ}

\paragraph{Description} An approximation of the reciprocal of the \emph{registerZ}, assuming binary 32 floating-point numbers, is stored into the \emph{registerW}. This instruction may raise exception bits in the CS register.


\paragraph{Execution} {Reservation table \textsf{ALU\_LITE} (Section~\ref{sec:reservation-ALU-LITE})}
~
{\begin{lstlisting}
stage RR:
  new argument2 = GPR[registerZ];
  new result1 = f32_fast_rec(argument2);
stage E1:
  GPR[registerW] = result1;
\end{lstlisting}}
\newpage
\subsubsection[{\small FSRECWP: Floating-Point Seed for Reciprocal Word Pair}]{FSRECWP\hfill Floating-Point Seed for Reciprocal Word Pair}
\label{instr:FSRECWP}\index{fsrecwp}
 
\paragraph{Encoding} Format \textsf{ALU\_FWPMO1} (Section~\ref{sec:format-ALU-FWPMO1}), \verb|fpucode6 = 0100|\\

\bitfields{ 1/P, 2/11, 1/1, 1/1, 3/111, 6/registerW, 2/00, 3/000, 1/0, 4/0100, 1/--, 1/1, 6/registerZ }


\paragraph{Syntax} \texttt{fsrecwp} \emph{registerW} = \emph{registerZ}

\paragraph{Description} An approximation of the reciprocal of the \emph{registerZ}, assuming a pair of binary 32 floating-point numbers, is stored into the \emph{registerW}. This instruction may raise exception bits in the CS register.


\paragraph{Execution} {Reservation table \textsf{ALU\_LITE} (Section~\ref{sec:reservation-ALU-LITE})}
~
{\begin{lstlisting}
stage RR:
  new argument2 = GPR[registerZ];
  for (I = 0; I < 2; I += 1) {
    result1.32[I] = f32_fast_rec(argument2.32[I])
  };
stage E1:
  GPR[registerW] = result1;
\end{lstlisting}}
\newpage
\subsubsection[{\small FSRSRD: Floating-Point Seed for Reciprocal Square Root Double Word}]{FSRSRD\hfill Floating-Point Seed for Reciprocal Square Root Double Word}
\label{instr:FSRSRD}\index{fsrsrd}
 
\paragraph{Encoding} Format \textsf{ALU\_FDMO1} (Section~\ref{sec:format-ALU-FDMO1}), \verb|fpucode6 = 0101|\\

\bitfields{ 1/P, 2/11, 1/1, 1/1, 3/111, 6/registerW, 2/00, 3/000, 1/0, 4/0101, 1/--, 1/0, 6/registerZ }


\paragraph{Syntax} \texttt{fsrsrd} \emph{registerW} = \emph{registerZ}

\paragraph{Description} An approximation of the reciprocal of square root of the \emph{registerZ}, assuming binary 64 floating-point numbers, is stored into the \emph{registerW}.


\paragraph{Execution} {Reservation table \textsf{ALU\_LITE} (Section~\ref{sec:reservation-ALU-LITE})}
~
{\begin{lstlisting}
stage RR:
  new argument2 = GPR[registerZ];
  new result1 = f64_fast_rsqrt(argument2);
stage E1:
  GPR[registerW] = result1;
\end{lstlisting}}
\newpage
\subsubsection[{\small FSRSRW: Floating-Point Seed for Reciprocal Square Root Word}]{FSRSRW\hfill Floating-Point Seed for Reciprocal Square Root Word}
\label{instr:FSRSRW}\index{fsrsrw}
 
\paragraph{Encoding} Format \textsf{ALU\_FWMO1} (Section~\ref{sec:format-ALU-FWMO1}), \verb|fpucode6 = 0101|\\

\bitfields{ 1/P, 2/11, 1/1, 1/1, 3/111, 6/registerW, 2/01, 3/000, 1/0, 4/0101, 1/--, 1/0, 6/registerZ }


\paragraph{Syntax} \texttt{fsrsrw} \emph{registerW} = \emph{registerZ}

\paragraph{Description} An approximation of the reciprocal of square root of the \emph{registerZ}, assuming binary 32 floating-point numbers, is stored into the \emph{registerW}.


\paragraph{Execution} {Reservation table \textsf{ALU\_LITE} (Section~\ref{sec:reservation-ALU-LITE})}
~
{\begin{lstlisting}
stage RR:
  new argument2 = GPR[registerZ];
  new result1 = f32_fast_rsqrt(argument2);
stage E1:
  GPR[registerW] = result1;
\end{lstlisting}}
\newpage
\subsubsection[{\small FSRSRWP: Floating-Point Seed for Reciprocal Square Root Word Pair}]{FSRSRWP\hfill Floating-Point Seed for Reciprocal Square Root Word Pair}
\label{instr:FSRSRWP}\index{fsrsrwp}
 
\paragraph{Encoding} Format \textsf{ALU\_FWPMO1} (Section~\ref{sec:format-ALU-FWPMO1}), \verb|fpucode6 = 0101|\\

\bitfields{ 1/P, 2/11, 1/1, 1/1, 3/111, 6/registerW, 2/00, 3/000, 1/0, 4/0101, 1/--, 1/1, 6/registerZ }


\paragraph{Syntax} \texttt{fsrsrwp} \emph{registerW} = \emph{registerZ}

\paragraph{Description} An approximation of the reciprocal of square root of the \emph{registerZ}, assuming a pair of binary 32 floating-point numbers, is stored into the \emph{registerW}.


\paragraph{Execution} {Reservation table \textsf{ALU\_LITE} (Section~\ref{sec:reservation-ALU-LITE})}
~
{\begin{lstlisting}
stage RR:
  new argument2 = GPR[registerZ];
  for (I = 0; I < 2; I += 1) {
    result1.32[I] = f32_fast_rsqrt(argument2.32[I])
  };
stage E1:
  GPR[registerW] = result1;
\end{lstlisting}}
\newpage
\subsubsection[{\small FWIDENHW: Floating Point Widen Half Word to Word}]{FWIDENHW\hfill Floating Point Widen Half Word to Word}
\label{instr:FWIDENHW}\index{fwidenhw}
 
\paragraph{Encoding} Format \textsf{ALU\_FWWIDEN} (Section~\ref{sec:format-ALU-FWWIDEN}), \verb|fpucode6 = 1000|\\

\bitfields{ 1/P, 2/11, 1/1, 1/1, 3/111, 6/registerW, 2/01, 3/000, 1/0, 4/1000, 1/mostsig, 1/0, 6/registerZ }


\paragraph{Syntax} \texttt{fwidenhw}\emph{mostsig} \emph{registerW} = \emph{registerZ}

\paragraph{Description} The \emph{registerZ} is interpreted as a binary 16 floating-point number. The least or most significant half word of the \emph{registerZ} word is selected, depending on \emph{mostsig}. The number is converted to a binary 32 floating-point number and stored into the \emph{registerW}.


\paragraph{Modifier(s)}
\noindent\begin{itemize}
\item mostsig (cf. mostsig in section \ref{par:modifier-mostsig} on page \pageref{par:modifier-mostsig})
\end{itemize}

\paragraph{Execution} {Reservation table \textsf{ALU\_LITE} (Section~\ref{sec:reservation-ALU-LITE})}
~
{\begin{lstlisting}
stage ID:
  new mostsig = _ZX_1(mostsig);
stage RR:
  new argument2 = GPR[registerZ];
  if (mostsig) {
    argument2 = argument2 >> 16;
  }
  new result1 = f16_to_f32(argument2.16[0]);
stage E1:
  GPR[registerW] = result1;
\end{lstlisting}}
\newpage
\subsubsection[{\small FWIDENHWQ: Floating Point Widen Half Word to Word Quadruple}]{FWIDENHWQ\hfill Floating Point Widen Half Word to Word Quadruple}
\label{instr:FWIDENHWQ}\index{fwidenhwq}
 
\paragraph{Encoding} Format \textsf{ALU\_FWQWIDEWR} (Section~\ref{sec:format-ALU-FWQWIDEWR}), \verb|fpucode6 = 1000|\\

\bitfields{ 1/P, 2/11, 1/1, 1/1, 3/111, 5/registerM, 1/--, 2/11, 3/000, 1/1, 4/1000, 1/mostsig, 1/0, 5/registerP, 1/0 }


\paragraph{Syntax} \texttt{fwidenhwq}\emph{mostsig} \emph{registerM} = \emph{registerP}

\paragraph{Description} The \emph{registerP} interpreted as eight binary 16 floating-point numbers packed into 128 bits. The least or most significant half words of the \emph{registerP} are selected, depending on \emph{mostsig}. The four selected half words are converted to binary 32 floating-point numbers, packed and stored into the \emph{registerM}. This instruction may raise invalid exception bit in the CS register.


\paragraph{Modifier(s)}
\noindent\begin{itemize}
\item mostsig (cf. mostsig in section \ref{par:modifier-mostsig} on page \pageref{par:modifier-mostsig})
\end{itemize}

\paragraph{Execution} {Reservation table \textsf{ALU\_LITE} (Section~\ref{sec:reservation-ALU-LITE})}
~
{\begin{lstlisting}
stage ID:
  new mostsig = _ZX_1(mostsig);
stage RR:
  new argument2 = PGR[registerP];
  if (mostsig) {
    argument2 = argument2 >> 64;
  }
  for (I = 0; I < 4; I += 1) {
    result1.32[I] = f16_to_f32(_ZX_16(argument2.16[I]))
  };
stage E1:
  PGR[registerM] = result1;
\end{lstlisting}}
\newpage
\subsubsection[{\small FWIDENWD: Floating Point Widen Word to Double Word}]{FWIDENWD\hfill Floating Point Widen Word to Double Word}
\label{instr:FWIDENWD}\index{fwidenwd}
 
\paragraph{Encoding} Format \textsf{ALU\_FDWIDEN} (Section~\ref{sec:format-ALU-FDWIDEN}), \verb|fpucode6 = 1000|\\

\bitfields{ 1/P, 2/11, 1/1, 1/1, 3/111, 6/registerW, 2/00, 3/000, 1/0, 4/1000, 1/mostsig, 1/0, 6/registerZ }


\paragraph{Syntax} \texttt{fwidenwd}\emph{mostsig} \emph{registerW} = \emph{registerZ}

\paragraph{Description} The \emph{registerZ} is interpreted as a binary 32 floating-point number. The least or most significant word of the \emph{registerZ} double word is selected, depending on \emph{mostsig}. The number is converted to a binary 64 floating-point number and stored into the \emph{registerW}.


\paragraph{Modifier(s)}
\noindent\begin{itemize}
\item mostsig (cf. mostsig in section \ref{par:modifier-mostsig} on page \pageref{par:modifier-mostsig})
\end{itemize}

\paragraph{Execution} {Reservation table \textsf{ALU\_LITE} (Section~\ref{sec:reservation-ALU-LITE})}
~
{\begin{lstlisting}
stage ID:
  new mostsig = _ZX_1(mostsig);
stage RR:
  new argument2 = GPR[registerZ];
  if (mostsig) {
    argument2 = argument2 >> 32;
  }
  new result1 = f32_to_f64(argument2.32[0]);
stage E1:
  GPR[registerW] = result1;
\end{lstlisting}}
\newpage
\subsection{EXT Instructions}
\label{sec:EXT-instructions}

\subsubsection[{\small XACCESSO: Access from Extension into Octuple Word}]{XACCESSO\hfill Access from Extension into Octuple Word}
\label{instr:XACCESSO}\index{xaccesso}
 
\paragraph{Encoding} Format \textsf{EXT\_XACCESSO2} (Section~\ref{sec:format-EXT-XACCESSO2}), \verb|log2size = 000|\\

\bitfields{ 1/P, 2/10, 1/0, 1/0, 3/000, 4/registerN, 2/11, 2/00, 1/1, 3/000, 5/registerCg, 1/0, 6/registerZ }


\paragraph{Syntax} \texttt{xaccesso} \emph{registerN} = \emph{registerCg}, \emph{registerZ}

\paragraph{Description} The 5 least significant bits of the \emph{registerZ} are interpreted as a byte count. The most significant bits of the \emph{registerZ} are interpreted as a vector register index into the buffer specified by the \emph{registerCg}. The vector register indexed and its successor modulo the buffer size are concatenated, and shifted right bt the byte count. The resulting 32 bytes are written to the \emph{registerN}.


\paragraph{Execution} {Reservation table \textsf{EXT\_MISC\_AUXW} (Section~\ref{sec:reservation-EXT-MISC-AUXW})}
~
{\begin{lstlisting}
stage RR:
  new mask = 1;
  new buffer = registerCg << 1;
  new address = GPR[registerZ];
  new shift = (address & 31) * 8;
  new index_0 = address >> 5;
  new index_1 = index_0 + 1;
  new where_0 = buffer + (index_0 & mask);
  new where_1 = buffer + (index_1 & mask);
  new result1_0 = XVR[where_0];
  new result1_1 = XVR[where_1];
  new result1 = result1_0;
  if (shift) {
    result1 = (_ZX_256(result1_0) >> shift) | (_ZX_256(result1_1) << (256 - shift));
  }
stage E1:
  CS.XMF = 1;
stage E3:
  QGR[registerN] = result1;
\end{lstlisting}}
\newpage

\paragraph{Encoding} Format \textsf{EXT\_XACCESSO4} (Section~\ref{sec:format-EXT-XACCESSO4}), \verb|log2size = 000|\\

\bitfields{ 1/P, 2/10, 1/0, 1/0, 3/000, 4/registerN, 2/11, 2/00, 1/1, 3/000, 4/registerCh, 2/01, 6/registerZ }


\paragraph{Syntax} \texttt{xaccesso} \emph{registerN} = \emph{registerCh}, \emph{registerZ}

\paragraph{Description} The 5 least significant bits of the \emph{registerZ} are interpreted as a byte count. The most significant bits of the \emph{registerZ} are interpreted as a vector register index into the buffer specified by the \emph{registerCh}. The vector register indexed and its successor modulo the buffer size are concatenated, and shifted right bt the byte count. The resulting 32 bytes are written to the \emph{registerN}.


\paragraph{Execution} {Reservation table \textsf{EXT\_MISC\_AUXW} (Section~\ref{sec:reservation-EXT-MISC-AUXW})}
~
{\begin{lstlisting}
stage RR:
  new mask = 3;
  new buffer = registerCh << 2;
  new address = GPR[registerZ];
  new shift = (address & 31) * 8;
  new index_0 = address >> 5;
  new index_1 = index_0 + 1;
  new where_0 = buffer + (index_0 & mask);
  new where_1 = buffer + (index_1 & mask);
  new result1_0 = XVR[where_0];
  new result1_1 = XVR[where_1];
  new result1 = result1_0;
  if (shift) {
    result1 = (_ZX_256(result1_0) >> shift) | (_ZX_256(result1_1) << (256 - shift));
  }
stage E1:
  CS.XMF = 1;
stage E3:
  QGR[registerN] = result1;
\end{lstlisting}}
\newpage

\paragraph{Encoding} Format \textsf{EXT\_XACCESSO8} (Section~\ref{sec:format-EXT-XACCESSO8}), \verb|log2size = 000|\\

\bitfields{ 1/P, 2/10, 1/0, 1/0, 3/000, 4/registerN, 2/11, 2/00, 1/1, 3/000, 3/registerCi, 3/011, 6/registerZ }


\paragraph{Syntax} \texttt{xaccesso} \emph{registerN} = \emph{registerCi}, \emph{registerZ}

\paragraph{Description} The 5 least significant bits of the \emph{registerZ} are interpreted as a byte count. The most significant bits of the \emph{registerZ} are interpreted as a vector register index into the buffer specified by the \emph{registerCi}. The vector register indexed and its successor modulo the buffer size are concatenated, and shifted right bt the byte count. The resulting 32 bytes are written to the \emph{registerN}.


\paragraph{Execution} {Reservation table \textsf{EXT\_MISC\_AUXW} (Section~\ref{sec:reservation-EXT-MISC-AUXW})}
~
{\begin{lstlisting}
stage RR:
  new mask = 7;
  new buffer = registerCi << 3;
  new address = GPR[registerZ];
  new shift = (address & 31) * 8;
  new index_0 = address >> 5;
  new index_1 = index_0 + 1;
  new where_0 = buffer + (index_0 & mask);
  new where_1 = buffer + (index_1 & mask);
  new result1_0 = XVR[where_0];
  new result1_1 = XVR[where_1];
  new result1 = result1_0;
  if (shift) {
    result1 = (_ZX_256(result1_0) >> shift) | (_ZX_256(result1_1) << (256 - shift));
  }
stage E1:
  CS.XMF = 1;
stage E3:
  QGR[registerN] = result1;
\end{lstlisting}}
\newpage

\paragraph{Encoding} Format \textsf{EXT\_XACCESSO16} (Section~\ref{sec:format-EXT-XACCESSO16}), \verb|log2size = 000|\\

\bitfields{ 1/P, 2/10, 1/0, 1/0, 3/000, 4/registerN, 2/11, 2/00, 1/1, 3/000, 2/registerCj, 4/0111, 6/registerZ }


\paragraph{Syntax} \texttt{xaccesso} \emph{registerN} = \emph{registerCj}, \emph{registerZ}

\paragraph{Description} The 5 least significant bits of the \emph{registerZ} are interpreted as a byte count. The most significant bits of the \emph{registerZ} are interpreted as a vector register index into the buffer specified by the \emph{registerCj}. The vector register indexed and its successor modulo the buffer size are concatenated, and shifted right bt the byte count. The resulting 32 bytes are written to the \emph{registerN}.


\paragraph{Execution} {Reservation table \textsf{EXT\_MISC\_AUXW} (Section~\ref{sec:reservation-EXT-MISC-AUXW})}
~
{\begin{lstlisting}
stage RR:
  new mask = 15;
  new buffer = registerCj << 4;
  new address = GPR[registerZ];
  new shift = (address & 31) * 8;
  new index_0 = address >> 5;
  new index_1 = index_0 + 1;
  new where_0 = buffer + (index_0 & mask);
  new where_1 = buffer + (index_1 & mask);
  new result1_0 = XVR[where_0];
  new result1_1 = XVR[where_1];
  new result1 = result1_0;
  if (shift) {
    result1 = (_ZX_256(result1_0) >> shift) | (_ZX_256(result1_1) << (256 - shift));
  }
stage E1:
  CS.XMF = 1;
stage E3:
  QGR[registerN] = result1;
\end{lstlisting}}
\newpage

\paragraph{Encoding} Format \textsf{EXT\_XACCESSO32} (Section~\ref{sec:format-EXT-XACCESSO32}), \verb|log2size = 000|\\

\bitfields{ 1/P, 2/10, 1/0, 1/0, 3/000, 4/registerN, 2/11, 2/00, 1/1, 3/000, 1/registerCk, 5/01111, 6/registerZ }


\paragraph{Syntax} \texttt{xaccesso} \emph{registerN} = \emph{registerCk}, \emph{registerZ}

\paragraph{Description} The 5 least significant bits of the \emph{registerZ} are interpreted as a byte count. The most significant bits of the \emph{registerZ} are interpreted as a vector register index into the buffer specified by the \emph{registerCk}. The vector register indexed and its successor modulo the buffer size are concatenated, and shifted right bt the byte count. The resulting 32 bytes are written to the \emph{registerN}.


\paragraph{Execution} {Reservation table \textsf{EXT\_MISC\_AUXW} (Section~\ref{sec:reservation-EXT-MISC-AUXW})}
~
{\begin{lstlisting}
stage RR:
  new mask = 31;
  new buffer = registerCk << 5;
  new address = GPR[registerZ];
  new shift = (address & 31) * 8;
  new index_0 = address >> 5;
  new index_1 = index_0 + 1;
  new where_0 = buffer + (index_0 & mask);
  new where_1 = buffer + (index_1 & mask);
  new result1_0 = XVR[where_0];
  new result1_1 = XVR[where_1];
  new result1 = result1_0;
  if (shift) {
    result1 = (_ZX_256(result1_0) >> shift) | (_ZX_256(result1_1) << (256 - shift));
  }
stage E1:
  CS.XMF = 1;
stage E3:
  QGR[registerN] = result1;
\end{lstlisting}}
\newpage

\paragraph{Encoding} Format \textsf{EXT\_XACCESSO64} (Section~\ref{sec:format-EXT-XACCESSO64}), \verb|log2size = 000|\\

\bitfields{ 1/P, 2/10, 1/0, 1/0, 3/000, 4/registerN, 2/11, 2/00, 1/1, 3/000, 1/registerCl, 5/11111, 6/registerZ }


\paragraph{Syntax} \texttt{xaccesso} \emph{registerN} = \emph{registerCl}, \emph{registerZ}

\paragraph{Description} The 5 least significant bits of the \emph{registerZ} are interpreted as a byte count. The most significant bits of the \emph{registerZ} are interpreted as a vector register index into the buffer specified by the \emph{registerCl}. The vector register indexed and its successor modulo the buffer size are concatenated, and shifted right bt the byte count. The resulting 32 bytes are written to the \emph{registerN}.


\paragraph{Execution} {Reservation table \textsf{EXT\_MISC\_AUXW} (Section~\ref{sec:reservation-EXT-MISC-AUXW})}
~
{\begin{lstlisting}
stage RR:
  new mask = 63;
  new buffer = registerCl;
  new address = GPR[registerZ];
  new shift = (address & 31) * 8;
  new index_0 = address >> 5;
  new index_1 = index_0 + 1;
  new where_0 = buffer + (index_0 & mask);
  new where_1 = buffer + (index_1 & mask);
  new result1_0 = XVR[where_0];
  new result1_1 = XVR[where_1];
  new result1 = result1_0;
  if (shift) {
    result1 = (_ZX_256(result1_0) >> shift) | (_ZX_256(result1_1) << (256 - shift));
  }
stage E1:
  CS.XMF = 1;
stage E3:
  QGR[registerN] = result1;
\end{lstlisting}}
\newpage
\subsubsection[{\small XALIGNO: Align Extension into Octuple Word}]{XALIGNO\hfill Align Extension into Octuple Word}
\label{instr:XALIGNO}\index{xaligno}
 
\paragraph{Encoding} Format \textsf{EXT\_XALIGNO2} (Section~\ref{sec:format-EXT-XALIGNO2}), \verb|log2size = 000|\\

\bitfields{ 1/P, 2/10, 1/0, 1/0, 3/001, 6/registerA, 2/00, 1/1, 3/000, 5/registerCg, 1/0, 6/registerZ }


\paragraph{Syntax} \texttt{xaligno} \emph{registerA} = \emph{registerCg}, \emph{registerZ}

\paragraph{Description} The 5 least significant bits of the \emph{registerZ} are interpreted as a byte count. The most significant bits of the \emph{registerZ} are interpreted as a vector register index into the buffer specified by the \emph{registerCg}. The vector register indexed and its successor modulo the buffer size are concatenated, and shifted right bt the byte count. The resulting 32 bytes are written to the \emph{registerA}.


\paragraph{Execution} {Reservation table \textsf{EXT\_MISC\_AUXW} (Section~\ref{sec:reservation-EXT-MISC-AUXW})}
~
{\begin{lstlisting}
stage RR:
  new mask = 1;
  new buffer = registerCg << 1;
  new address = GPR[registerZ];
  new shift = (address & 31) * 8;
  new index_0 = address >> 5;
  new index_1 = index_0 + 1;
  new where_0 = buffer + (index_0 & mask);
  new where_1 = buffer + (index_1 & mask);
  new result1_0 = XVR[where_0];
  new result1_1 = XVR[where_1];
  new result1 = result1_0;
  if (shift) {
    result1 = (_ZX_256(result1_0) >> shift) | (_ZX_256(result1_1) << (256 - shift));
  }
stage E1:
  CS.XMF = 1;
stage E3:
  XVR[registerA] = result1;
\end{lstlisting}}
\newpage

\paragraph{Encoding} Format \textsf{EXT\_XALIGNO4} (Section~\ref{sec:format-EXT-XALIGNO4}), \verb|log2size = 000|\\

\bitfields{ 1/P, 2/10, 1/0, 1/0, 3/001, 6/registerA, 2/00, 1/1, 3/000, 4/registerCh, 2/01, 6/registerZ }


\paragraph{Syntax} \texttt{xaligno} \emph{registerA} = \emph{registerCh}, \emph{registerZ}

\paragraph{Description} The 5 least significant bits of the \emph{registerZ} are interpreted as a byte count. The most significant bits of the \emph{registerZ} are interpreted as a vector register index into the buffer specified by the \emph{registerCh}. The vector register indexed and its successor modulo the buffer size are concatenated, and shifted right bt the byte count. The resulting 32 bytes are written to the \emph{registerA}.


\paragraph{Execution} {Reservation table \textsf{EXT\_MISC\_AUXW} (Section~\ref{sec:reservation-EXT-MISC-AUXW})}
~
{\begin{lstlisting}
stage RR:
  new mask = 3;
  new buffer = registerCh << 2;
  new address = GPR[registerZ];
  new shift = (address & 31) * 8;
  new index_0 = address >> 5;
  new index_1 = index_0 + 1;
  new where_0 = buffer + (index_0 & mask);
  new where_1 = buffer + (index_1 & mask);
  new result1_0 = XVR[where_0];
  new result1_1 = XVR[where_1];
  new result1 = result1_0;
  if (shift) {
    result1 = (_ZX_256(result1_0) >> shift) | (_ZX_256(result1_1) << (256 - shift));
  }
stage E1:
  CS.XMF = 1;
stage E3:
  XVR[registerA] = result1;
\end{lstlisting}}
\newpage

\paragraph{Encoding} Format \textsf{EXT\_XALIGNO8} (Section~\ref{sec:format-EXT-XALIGNO8}), \verb|log2size = 000|\\

\bitfields{ 1/P, 2/10, 1/0, 1/0, 3/001, 6/registerA, 2/00, 1/1, 3/000, 3/registerCi, 3/011, 6/registerZ }


\paragraph{Syntax} \texttt{xaligno} \emph{registerA} = \emph{registerCi}, \emph{registerZ}

\paragraph{Description} The 5 least significant bits of the \emph{registerZ} are interpreted as a byte count. The most significant bits of the \emph{registerZ} are interpreted as a vector register index into the buffer specified by the \emph{registerCi}. The vector register indexed and its successor modulo the buffer size are concatenated, and shifted right bt the byte count. The resulting 32 bytes are written to the \emph{registerA}.


\paragraph{Execution} {Reservation table \textsf{EXT\_MISC\_AUXW} (Section~\ref{sec:reservation-EXT-MISC-AUXW})}
~
{\begin{lstlisting}
stage RR:
  new mask = 7;
  new buffer = registerCi << 3;
  new address = GPR[registerZ];
  new shift = (address & 31) * 8;
  new index_0 = address >> 5;
  new index_1 = index_0 + 1;
  new where_0 = buffer + (index_0 & mask);
  new where_1 = buffer + (index_1 & mask);
  new result1_0 = XVR[where_0];
  new result1_1 = XVR[where_1];
  new result1 = result1_0;
  if (shift) {
    result1 = (_ZX_256(result1_0) >> shift) | (_ZX_256(result1_1) << (256 - shift));
  }
stage E1:
  CS.XMF = 1;
stage E3:
  XVR[registerA] = result1;
\end{lstlisting}}
\newpage

\paragraph{Encoding} Format \textsf{EXT\_XALIGNO16} (Section~\ref{sec:format-EXT-XALIGNO16}), \verb|log2size = 000|\\

\bitfields{ 1/P, 2/10, 1/0, 1/0, 3/001, 6/registerA, 2/00, 1/1, 3/000, 2/registerCj, 4/0111, 6/registerZ }


\paragraph{Syntax} \texttt{xaligno} \emph{registerA} = \emph{registerCj}, \emph{registerZ}

\paragraph{Description} The 5 least significant bits of the \emph{registerZ} are interpreted as a byte count. The most significant bits of the \emph{registerZ} are interpreted as a vector register index into the buffer specified by the \emph{registerCj}. The vector register indexed and its successor modulo the buffer size are concatenated, and shifted right bt the byte count. The resulting 32 bytes are written to the \emph{registerA}.


\paragraph{Execution} {Reservation table \textsf{EXT\_MISC\_AUXW} (Section~\ref{sec:reservation-EXT-MISC-AUXW})}
~
{\begin{lstlisting}
stage RR:
  new mask = 15;
  new buffer = registerCj << 4;
  new address = GPR[registerZ];
  new shift = (address & 31) * 8;
  new index_0 = address >> 5;
  new index_1 = index_0 + 1;
  new where_0 = buffer + (index_0 & mask);
  new where_1 = buffer + (index_1 & mask);
  new result1_0 = XVR[where_0];
  new result1_1 = XVR[where_1];
  new result1 = result1_0;
  if (shift) {
    result1 = (_ZX_256(result1_0) >> shift) | (_ZX_256(result1_1) << (256 - shift));
  }
stage E1:
  CS.XMF = 1;
stage E3:
  XVR[registerA] = result1;
\end{lstlisting}}
\newpage

\paragraph{Encoding} Format \textsf{EXT\_XALIGNO32} (Section~\ref{sec:format-EXT-XALIGNO32}), \verb|log2size = 000|\\

\bitfields{ 1/P, 2/10, 1/0, 1/0, 3/001, 6/registerA, 2/00, 1/1, 3/000, 1/registerCk, 5/01111, 6/registerZ }


\paragraph{Syntax} \texttt{xaligno} \emph{registerA} = \emph{registerCk}, \emph{registerZ}

\paragraph{Description} The 5 least significant bits of the \emph{registerZ} are interpreted as a byte count. The most significant bits of the \emph{registerZ} are interpreted as a vector register index into the buffer specified by the \emph{registerCk}. The vector register indexed and its successor modulo the buffer size are concatenated, and shifted right bt the byte count. The resulting 32 bytes are written to the \emph{registerA}.


\paragraph{Execution} {Reservation table \textsf{EXT\_MISC\_AUXW} (Section~\ref{sec:reservation-EXT-MISC-AUXW})}
~
{\begin{lstlisting}
stage RR:
  new mask = 31;
  new buffer = registerCk << 5;
  new address = GPR[registerZ];
  new shift = (address & 31) * 8;
  new index_0 = address >> 5;
  new index_1 = index_0 + 1;
  new where_0 = buffer + (index_0 & mask);
  new where_1 = buffer + (index_1 & mask);
  new result1_0 = XVR[where_0];
  new result1_1 = XVR[where_1];
  new result1 = result1_0;
  if (shift) {
    result1 = (_ZX_256(result1_0) >> shift) | (_ZX_256(result1_1) << (256 - shift));
  }
stage E1:
  CS.XMF = 1;
stage E3:
  XVR[registerA] = result1;
\end{lstlisting}}
\newpage

\paragraph{Encoding} Format \textsf{EXT\_XALIGNO64} (Section~\ref{sec:format-EXT-XALIGNO64}), \verb|log2size = 000|\\

\bitfields{ 1/P, 2/10, 1/0, 1/0, 3/001, 6/registerA, 2/00, 1/1, 3/000, 1/registerCl, 5/11111, 6/registerZ }


\paragraph{Syntax} \texttt{xaligno} \emph{registerA} = \emph{registerCl}, \emph{registerZ}

\paragraph{Description} The 5 least significant bits of the \emph{registerZ} are interpreted as a byte count. The most significant bits of the \emph{registerZ} are interpreted as a vector register index into the buffer specified by the \emph{registerCl}. The vector register indexed and its successor modulo the buffer size are concatenated, and shifted right bt the byte count. The resulting 32 bytes are written to the \emph{registerA}.


\paragraph{Execution} {Reservation table \textsf{EXT\_MISC\_AUXW} (Section~\ref{sec:reservation-EXT-MISC-AUXW})}
~
{\begin{lstlisting}
stage RR:
  new mask = 63;
  new buffer = registerCl;
  new address = GPR[registerZ];
  new shift = (address & 31) * 8;
  new index_0 = address >> 5;
  new index_1 = index_0 + 1;
  new where_0 = buffer + (index_0 & mask);
  new where_1 = buffer + (index_1 & mask);
  new result1_0 = XVR[where_0];
  new result1_1 = XVR[where_1];
  new result1 = result1_0;
  if (shift) {
    result1 = (_ZX_256(result1_0) >> shift) | (_ZX_256(result1_1) << (256 - shift));
  }
stage E1:
  CS.XMF = 1;
stage E3:
  XVR[registerA] = result1;
\end{lstlisting}}
\newpage
\subsubsection[{\small XCOPYO: Copy Extension Octuple Word}]{XCOPYO\hfill Copy Extension Octuple Word}
\label{instr:XCOPYO}\index{xcopyo}
 
\paragraph{Encoding} Format \textsf{EXT\_XCOPYO} (Section~\ref{sec:format-EXT-XCOPYO}), \verb|log2size = 000|\\

\bitfields{ 1/P, 2/10, 1/0, 1/0, 3/001, 6/registerA, 2/00, 1/0, 3/000, 6/registerC, 6/-- -- -- -- -- -- }


\paragraph{Syntax} \texttt{xcopyo} \emph{registerA} = \emph{registerC}

\paragraph{Description} The \emph{registerA} operand receives the value of the \emph{registerC} operand.


\paragraph{Execution} {Reservation table \textsf{EXT\_MISC\_AUXW} (Section~\ref{sec:reservation-EXT-MISC-AUXW})}
~
{\begin{lstlisting}
stage RR:
  new argument2 = XVR[registerC];
  new result1 = argument2;
stage E1:
  CS.XMF = 1;
stage E3:
  XVR[registerA] = result1;
\end{lstlisting}}
\newpage
\subsubsection[{\small XMOVEFD: Move From Extension Double Word}]{XMOVEFD\hfill Move From Extension Double Word}
\label{instr:XMOVEFD}\index{xmovefd}
 
\paragraph{Encoding} Format \textsf{EXT\_XMOVEFD} (Section~\ref{sec:format-EXT-XMOVEFD}), \verb|tcacode1 = 000|\\

\bitfields{ 1/P, 2/10, 1/0, 1/0, 3/000, 6/registerW, 2/00, 1/1, 1/1, 2/qselectC, 6/registerCc, 6/-- -- -- -- -- -- }


\paragraph{Syntax} \texttt{xmovefd} \emph{registerW} = \emph{registerCc\_\discretionary{}{}{}qselectC}

\paragraph{Description} The \emph{registerW} operand receives the contents of the \emph{registerCc\_\discretionary{}{}{}qselectC}.


\paragraph{Execution} {Reservation table \textsf{EXT\_MISC\_AUXW} (Section~\ref{sec:reservation-EXT-MISC-AUXW})}
~
{\begin{lstlisting}
stage RR:
  new argument2 = XCR[registerCc_qselectC];
  new result1 = _ZX_64(argument2);
stage E1:
  CS.XMF = 1;
stage E3:
  GPR[registerW] = result1;
\end{lstlisting}}
\newpage
\subsubsection[{\small XMOVEFO: Move From Extension Octuple Word}]{XMOVEFO\hfill Move From Extension Octuple Word}
\label{instr:XMOVEFO}\index{xmovefo}
 
\paragraph{Encoding} Format \textsf{EXT\_XMOVEFO} (Section~\ref{sec:format-EXT-XMOVEFO}), \verb|tcacode1 = 000|\\

\bitfields{ 1/P, 2/10, 1/0, 1/0, 3/000, 4/registerN, 2/01, 2/00, 1/1, 3/000, 6/registerC, 6/-- -- -- -- -- -- }


\paragraph{Syntax} \texttt{xmovefo} \emph{registerN} = \emph{registerC}

\paragraph{Description} The \emph{registerN} operand receives the contents of the \emph{registerC}.


\paragraph{Execution} {Reservation table \textsf{EXT\_MISC\_AUXW} (Section~\ref{sec:reservation-EXT-MISC-AUXW})}
~
{\begin{lstlisting}
stage RR:
  new argument2 = XVR[registerC];
  new result1 = _ZX_256(argument2);
stage E1:
  CS.XMF = 1;
stage E3:
  QGR[registerN] = result1;
\end{lstlisting}}
\newpage
\subsubsection[{\small XMOVEFQ: Move From Extension Quadruple Word}]{XMOVEFQ\hfill Move From Extension Quadruple Word}
\label{instr:XMOVEFQ}\index{xmovefq}
 
\paragraph{Encoding} Format \textsf{EXT\_XMOVEFQ} (Section~\ref{sec:format-EXT-XMOVEFQ}), \verb|tcacode1 = 000|\\

\bitfields{ 1/P, 2/10, 1/0, 1/0, 3/000, 5/registerM, 1/1, 2/00, 1/0, 1/0, 1/0, 1/hselectC, 6/registerCb, 6/-- -- -- -- -- -- }


\paragraph{Syntax} \texttt{xmovefq} \emph{registerM} = \emph{registerCb\_\discretionary{}{}{}hselectC}

\paragraph{Description} The \emph{registerM} operand receives the contents of the \emph{registerCb\_\discretionary{}{}{}hselectC}.


\paragraph{Execution} {Reservation table \textsf{EXT\_MISC\_AUXW} (Section~\ref{sec:reservation-EXT-MISC-AUXW})}
~
{\begin{lstlisting}
stage RR:
  new argument2 = XBR[registerCb_hselectC];
  new result1 = _ZX_128(argument2);
stage E1:
  CS.XMF = 1;
stage E3:
  PGR[registerM] = result1;
\end{lstlisting}}
\newpage
\subsection{RV32B Instructions}
\label{sec:RV32B-instructions}

\subsubsection[{\small ADD.UW: Integer Add Unsigned Word}]{ADD.UW\hfill Integer Add Unsigned Word}

\paragraph{Encoding} Format \textsf{RV32B\_OPREG} (Section~\ref{sec:format-RV32B-OPREG}), \verb|funct7 = 0000100|, \verb|funct3 = 000|\\

\bitfields{ 7/0000100, 5/rs2, 5/rs1, 3/000, 5/rd, 7/0111011 }


\paragraph{Syntax} \texttt{add.uw} \emph{rd}, \emph{rs1}, \emph{rs2}

\paragraph{Description} The zero-extended \emph{rs1} word is added to the \emph{rs2}. The 64-bit result is stored into the \emph{rd}.


\paragraph{Execution} {Reservation table \textsf{ALU\_TINY} (Section~\ref{sec:reservation-ALU-TINY})}
~
{\begin{lstlisting}
stage RR:
  new rs2 = BIR[rs2];
  new rs1 = BIR[rs1];
  new rd = _ZX_32(rs1) + rs2;
stage E1:
  if (rd != 0) {
    BIR[rd] = rd;
  }
\end{lstlisting}}
\newpage
\subsubsection[{\small CLZW: Count Leading Zero Bits in Word}]{CLZW\hfill Count Leading Zero Bits in Word}

\paragraph{Encoding} Format \textsf{RV32B\_OPCNT\_0x30} (Section~\ref{sec:format-RV32B-OPCNT-0x30}), \verb|i_4_0 = 00000|\\

\bitfields{ 7/0110000, 5/00000, 5/rs1, 3/001, 5/rd, 7/0011011 }


\paragraph{Syntax} \texttt{clzw} \emph{rd}, \emph{rs1}

\paragraph{Description} Counts the number of 0s before the first 1 of the least significant \emph{rs1} word, starting at the most significant bit. The result is stored into the \emph{rd}.


\paragraph{Execution} {Reservation table \textsf{ALU\_LITE} (Section~\ref{sec:reservation-ALU-LITE})}
~
{\begin{lstlisting}
stage RR:
  new rs1 = BIR[rs1];
  new rd = _CLZ_32(rs1);
stage E1:
  if (rd != 0) {
    BIR[rd] = rd;
  }
\end{lstlisting}}
\newpage
\subsubsection[{\small CPOPW: Count Set Bits in Word}]{CPOPW\hfill Count Set Bits in Word}

\paragraph{Encoding} Format \textsf{RV32B\_OPCNT\_0x30} (Section~\ref{sec:format-RV32B-OPCNT-0x30}), \verb|i_4_0 = 00010|\\

\bitfields{ 7/0110000, 5/00010, 5/rs1, 3/001, 5/rd, 7/0011011 }


\paragraph{Syntax} \texttt{cpopw} \emph{rd}, \emph{rs1}

\paragraph{Description} Counts the number of 1s of the least significant \emph{rs1} word. The result is stored into the \emph{rd}.


\paragraph{Execution} {Reservation table \textsf{ALU\_LITE} (Section~\ref{sec:reservation-ALU-LITE})}
~
{\begin{lstlisting}
stage RR:
  new rs1 = BIR[rs1];
  new rd = _CBS_32(rs1);
stage E1:
  if (rd != 0) {
    BIR[rd] = rd;
  }
\end{lstlisting}}
\newpage
\subsubsection[{\small CTZW: Count Trailing Zero Bits in Word}]{CTZW\hfill Count Trailing Zero Bits in Word}

\paragraph{Encoding} Format \textsf{RV32B\_OPCNT\_0x30} (Section~\ref{sec:format-RV32B-OPCNT-0x30}), \verb|i_4_0 = 00001|\\

\bitfields{ 7/0110000, 5/00001, 5/rs1, 3/001, 5/rd, 7/0011011 }


\paragraph{Syntax} \texttt{ctzw} \emph{rd}, \emph{rs1}

\paragraph{Description} Counts the number of 0s before the first 1 of the least significant \emph{rs1} word, starting at the least significant bit. The result is stored into the \emph{rd}.


\paragraph{Execution} {Reservation table \textsf{ALU\_LITE} (Section~\ref{sec:reservation-ALU-LITE})}
~
{\begin{lstlisting}
stage RR:
  new rs1 = BIR[rs1];
  new rd = _CTZ_32(rs1);
stage E1:
  if (rd != 0) {
    BIR[rd] = rd;
  }
\end{lstlisting}}
\newpage
\subsubsection[{\small ROLW: Rotate Left Word}]{ROLW\hfill Rotate Left Word}

\paragraph{Encoding} Format \textsf{RV32B\_OPREG} (Section~\ref{sec:format-RV32B-OPREG}), \verb|funct7 = 0110000|, \verb|funct3 = 001|\\

\bitfields{ 7/0110000, 5/rs2, 5/rs1, 3/001, 5/rd, 7/0111011 }


\paragraph{Syntax} \texttt{rolw} \emph{rd}, \emph{rs1}, \emph{rs2}

\paragraph{Description} The \emph{rs1} word is rotated left by the \emph{rs2} modulo 32. The result is sign-extended and stored into the \emph{rd}.


\paragraph{Execution} {Reservation table \textsf{ALU\_LITE} (Section~\ref{sec:reservation-ALU-LITE})}
~
{\begin{lstlisting}
stage RR:
  new rs2 = BIR[rs2];
  new rs1 = BIR[rs1];
  new shift = _ZX_5(rs2);
  new rd = _SX_32(_ROL_32(rs1, shift));
stage E1:
  if (rd != 0) {
    BIR[rd] = rd;
  }
\end{lstlisting}}
\newpage
\subsubsection[{\small RORIW: Rotate Right Immediate Word}]{RORIW\hfill Rotate Right Immediate Word}

\paragraph{Encoding} Format \textsf{RV32B\_OPIMM\_0x30} (Section~\ref{sec:format-RV32B-OPIMM-0x30}), \verb|funct3 = 101|\\

\bitfields{ 7/0110000, 5/i-- 4-- 0, 5/rs1, 3/101, 5/rd, 7/0011011 }


\paragraph{Syntax} \texttt{roriw} \emph{rd}, \emph{rs1}, \emph{i\_\discretionary{}{}{}4\_\discretionary{}{}{}0}

\paragraph{Description} The \emph{rs1} word is rotated right by the \emph{i\_\discretionary{}{}{}4\_\discretionary{}{}{}0}. The result is sign-extended and stored into the \emph{rd}.


\paragraph{Execution} {Reservation table \textsf{ALU\_LITE} (Section~\ref{sec:reservation-ALU-LITE})}
~
{\begin{lstlisting}
stage RR:
  new unsigned5 = _ZX_5(i_4_0);
  new rs1 = BIR[rs1];
  new rd = _SX_32(_ROR_32(rs1, unsigned5));
stage E1:
  if (rd != 0) {
    BIR[rd] = rd;
  }
\end{lstlisting}}
\newpage
\subsubsection[{\small RORW: Rotate Right Word}]{RORW\hfill Rotate Right Word}

\paragraph{Encoding} Format \textsf{RV32B\_OPREG} (Section~\ref{sec:format-RV32B-OPREG}), \verb|funct7 = 0110000|, \verb|funct3 = 101|\\

\bitfields{ 7/0110000, 5/rs2, 5/rs1, 3/101, 5/rd, 7/0111011 }


\paragraph{Syntax} \texttt{rorw} \emph{rd}, \emph{rs1}, \emph{rs2}

\paragraph{Description} The \emph{rs1} word is rotated right by the \emph{rs2} modulo 32. The result is sign-extended and stored into the \emph{rd}.


\paragraph{Execution} {Reservation table \textsf{ALU\_LITE} (Section~\ref{sec:reservation-ALU-LITE})}
~
{\begin{lstlisting}
stage RR:
  new rs2 = BIR[rs2];
  new rs1 = BIR[rs1];
  new shift = _ZX_5(rs2);
  new rd = _SX_32(_ROR_32(rs1, shift));
stage E1:
  if (rd != 0) {
    BIR[rd] = rd;
  }
\end{lstlisting}}
\newpage
\subsubsection[{\small SH1ADD.UW: Shift unsigned word left by 1 and ADD}]{SH1ADD.UW\hfill Shift unsigned word left by 1 and ADD}

\paragraph{Encoding} Format \textsf{RV32B\_OPREG} (Section~\ref{sec:format-RV32B-OPREG}), \verb|funct7 = 0010000|, \verb|funct3 = 010|\\

\bitfields{ 7/0010000, 5/rs2, 5/rs1, 3/010, 5/rd, 7/0111011 }


\paragraph{Syntax} \texttt{sh1add.uw} \emph{rd}, \emph{rs1}, \emph{rs2}

\paragraph{Description} The \emph{rs1} word is zero-extended, shifted left by one bit and added to the \emph{rs2}. The result is stored into the \emph{rd}.


\paragraph{Execution} {Reservation table \textsf{ALU\_TINY} (Section~\ref{sec:reservation-ALU-TINY})}
~
{\begin{lstlisting}
stage RR:
  new rs2 = BIR[rs2];
  new rs1 = BIR[rs1];
  new rd = rs2 + (_ZX_32(rs1) << 1);
stage E1:
  if (rd != 0) {
    BIR[rd] = rd;
  }
\end{lstlisting}}
\newpage
\subsubsection[{\small SH2ADD.UW: Shift unsigned word left by 2 and ADD}]{SH2ADD.UW\hfill Shift unsigned word left by 2 and ADD}

\paragraph{Encoding} Format \textsf{RV32B\_OPREG} (Section~\ref{sec:format-RV32B-OPREG}), \verb|funct7 = 0010000|, \verb|funct3 = 100|\\

\bitfields{ 7/0010000, 5/rs2, 5/rs1, 3/100, 5/rd, 7/0111011 }


\paragraph{Syntax} \texttt{sh2add.uw} \emph{rd}, \emph{rs1}, \emph{rs2}

\paragraph{Description} The \emph{rs1} word is zero-extended, shifted left by two bits and added to the \emph{rs2}. The result is stored into the \emph{rd}.


\paragraph{Execution} {Reservation table \textsf{ALU\_TINY} (Section~\ref{sec:reservation-ALU-TINY})}
~
{\begin{lstlisting}
stage RR:
  new rs2 = BIR[rs2];
  new rs1 = BIR[rs1];
  new rd = rs2 + (_ZX_32(rs1) << 2);
stage E1:
  if (rd != 0) {
    BIR[rd] = rd;
  }
\end{lstlisting}}
\newpage
\subsubsection[{\small SH3ADD.UW: Shift unsigned word left by 3 and ADD}]{SH3ADD.UW\hfill Shift unsigned word left by 3 and ADD}

\paragraph{Encoding} Format \textsf{RV32B\_OPREG} (Section~\ref{sec:format-RV32B-OPREG}), \verb|funct7 = 0010000|, \verb|funct3 = 110|\\

\bitfields{ 7/0010000, 5/rs2, 5/rs1, 3/110, 5/rd, 7/0111011 }


\paragraph{Syntax} \texttt{sh3add.uw} \emph{rd}, \emph{rs1}, \emph{rs2}

\paragraph{Description} The \emph{rs1} word is zero-extended, shifted left by three bits and added to the \emph{rs2}. The result is stored into the \emph{rd}.


\paragraph{Execution} {Reservation table \textsf{ALU\_TINY} (Section~\ref{sec:reservation-ALU-TINY})}
~
{\begin{lstlisting}
stage RR:
  new rs2 = BIR[rs2];
  new rs1 = BIR[rs1];
  new rd = rs2 + (_ZX_32(rs1) << 3);
stage E1:
  if (rd != 0) {
    BIR[rd] = rd;
  }
\end{lstlisting}}
\newpage
\subsubsection[{\small SLLI.UW: Shift Left Logical Immediate Unsigned Word}]{SLLI.UW\hfill Shift Left Logical Immediate Unsigned Word}

\paragraph{Encoding} Format \textsf{RV32B\_OPIMM\_0x20} (Section~\ref{sec:format-RV32B-OPIMM-0x20}), \verb|funct3 = 001|\\

\bitfields{ 6/000010, 6/i-- 5-- 0, 5/rs1, 3/001, 5/rd, 7/0011011 }


\paragraph{Syntax} \texttt{slli.uw} \emph{rd}, \emph{rs1}, \emph{i\_\discretionary{}{}{}5\_\discretionary{}{}{}0}

\paragraph{Description} The \emph{rs1} is 32-bit zero-extended and shifted left by the \emph{i\_\discretionary{}{}{}5\_\discretionary{}{}{}0}. The result is stored into the \emph{rd}.


\paragraph{Execution} {Reservation table \textsf{ALU\_TINY} (Section~\ref{sec:reservation-ALU-TINY})}
~
{\begin{lstlisting}
stage RR:
  new unsigned6 = _ZX_6(i_5_0);
  new rs1 = BIR[rs1];
  new rd = _ZX_32(rs1) << unsigned6;
stage E1:
  if (rd != 0) {
    BIR[rd] = rd;
  }
\end{lstlisting}}
\newpage
\subsubsection[{\small ZEXT.H: Zero-Extend Half Word}]{ZEXT.H\hfill Zero-Extend Half Word}

\paragraph{Encoding} Format \textsf{RV32B\_OPEXT\_0x04} (Section~\ref{sec:format-RV32B-OPEXT-0x04}), \verb|funct3 = 100|\\

\bitfields{ 7/0000100, 5/00000, 5/rs1, 3/100, 5/rd, 7/0111011 }


\paragraph{Syntax} \texttt{zext.h} \emph{rd}, \emph{rs1}

\paragraph{Description} The \emph{rs1} half word is zero-extended. The result is stored into the \emph{rd}.


\paragraph{Execution} {Reservation table \textsf{ALU\_TINY} (Section~\ref{sec:reservation-ALU-TINY})}
~
{\begin{lstlisting}
stage RR:
  new rs1 = BIR[rs1];
  new rd = _ZX_16(rs1);
stage E1:
  if (rd != 0) {
    BIR[rd] = rd;
  }
\end{lstlisting}}
\newpage
\subsection{RV32D Instructions}
\label{sec:RV32D-instructions}

\subsubsection[{\small FADD.D: Floating-Point Add Double}]{FADD.D\hfill Floating-Point Add Double}

\paragraph{Encoding} Format \textsf{RV32D\_R3TypeRM} (Section~\ref{sec:format-RV32D-R3TypeRM}), \verb|funct5 = 00000|\\

\bitfields{ 5/00000, 2/01, 5/frs2, 5/frs1, 3/frm, 5/frd, 7/1010011 }


\paragraph{Syntax} \texttt{fadd.d} \emph{frd}, \emph{frs1}, \emph{frs2}, \emph{frm}

\paragraph{Description} 

\paragraph{Modifier(s)}
\noindent\begin{itemize}
\item frm (cf. froundmode in section \ref{par:modifier-froundmode} on page \pageref{par:modifier-froundmode})
\end{itemize}

\paragraph{Execution} {Reservation table \textsf{ALU\_LITE} (Section~\ref{sec:reservation-ALU-LITE})}
~
{\begin{lstlisting}
stage ID:
  new frm = _ZX_3(frm);
stage RR:
  new frs2 = FPR[frs2];
  new frs1 = FPR[frs1];
  new frm = decode_riscv_float_rounding_mode(frm);
  if (check_float_rounding_mode_trap(frm) == 0) {
    _THROW(OPCODE);
  } else {
    new frd = f64_add(frm, frs1, frs2);
    FPR[frd] = frd;
    commit_float_exception_flags();
  }
\end{lstlisting}}
\newpage
\subsubsection[{\small FCLASS.D: Floating-Point Classify Double}]{FCLASS.D\hfill Floating-Point Classify Double}

\paragraph{Encoding} Format \textsf{RV32D\_RTypeFD2I} (Section~\ref{sec:format-RV32D-RTypeFD2I}), \verb|funct3 = 001|\\

\bitfields{ 5/11100, 2/01, 5/00000, 5/frs1, 3/001, 5/rd, 7/1010011 }


\paragraph{Syntax} \texttt{fclass.d} \emph{rd}, \emph{frs1}

\paragraph{Description} 

\paragraph{Execution} {Reservation table \textsf{ALU\_LITE} (Section~\ref{sec:reservation-ALU-LITE})}
~
{\begin{lstlisting}
stage RR:
  new frs1 = FPR[frs1];
  new rd = f64_classify(frs1);
stage E1:
  if (rd != 0) {
    BIR[rd] = rd;
  }
\end{lstlisting}}
\newpage
\subsubsection[{\small FCVT.D.S: Floating-Point Convert Double from Single}]{FCVT.D.S\hfill Floating-Point Convert Double from Single}

\paragraph{Encoding} Format \textsf{RV32D\_RTypeFCF2D} (Section~\ref{sec:format-RV32D-RTypeFCF2D}), \verb|rs2 = 00000|\\

\bitfields{ 5/01000, 2/01, 5/00000, 5/frs1, 3/frm, 5/frd, 7/1010011 }


\paragraph{Syntax} \texttt{fcvt.d.s} \emph{frd}, \emph{frs1}, \emph{frm}

\paragraph{Description} 

\paragraph{Modifier(s)}
\noindent\begin{itemize}
\item frm (cf. froundmode in section \ref{par:modifier-froundmode} on page \pageref{par:modifier-froundmode})
\end{itemize}

\paragraph{Execution} {Reservation table \textsf{ALU\_LITE} (Section~\ref{sec:reservation-ALU-LITE})}
~
{\begin{lstlisting}
stage ID:
  new frm = _ZX_3(frm);
stage RR:
  new frs1 = FPR[frs1];
  new frm = decode_riscv_float_rounding_mode(frm);
  if (check_float_rounding_mode_trap(frm) == 0) {
    _THROW(OPCODE);
  } else {
    new frd = f32_to_f64(frs1);
    FPR[frd] = frd;
    commit_float_exception_flags();
  }
\end{lstlisting}}
\newpage
\subsubsection[{\small FCVT.D.W: Floating-Point Convert Double from Word}]{FCVT.D.W\hfill Floating-Point Convert Double from Word}

\paragraph{Encoding} Format \textsf{RV32D\_RTypeFCI2D} (Section~\ref{sec:format-RV32D-RTypeFCI2D}), \verb|rs2 = 00000|\\

\bitfields{ 5/11010, 2/01, 5/00000, 5/rs1, 3/frm, 5/frd, 7/1010011 }


\paragraph{Syntax} \texttt{fcvt.d.w} \emph{frd}, \emph{rs1}, \emph{frm}

\paragraph{Description} 

\paragraph{Modifier(s)}
\noindent\begin{itemize}
\item frm (cf. froundmode in section \ref{par:modifier-froundmode} on page \pageref{par:modifier-froundmode})
\end{itemize}

\paragraph{Execution} {Reservation table \textsf{ALU\_LITE} (Section~\ref{sec:reservation-ALU-LITE})}
~
{\begin{lstlisting}
stage ID:
  new frm = _ZX_3(frm);
stage RR:
  new rs1 = BIR[rs1];
  new frm = decode_riscv_float_rounding_mode(frm);
  if (check_float_rounding_mode_trap(frm) == 0) {
    _THROW(OPCODE);
  } else {
    new frd = i32_to_f64(frm, rs1);
    FPR[frd] = frd;
    commit_float_exception_flags();
  }
\end{lstlisting}}
\newpage
\subsubsection[{\small FCVT.D.WU: Floating-Point Convert Double from Word Unsigned}]{FCVT.D.WU\hfill Floating-Point Convert Double from Word Unsigned}

\paragraph{Encoding} Format \textsf{RV32D\_RTypeFCI2D} (Section~\ref{sec:format-RV32D-RTypeFCI2D}), \verb|rs2 = 00001|\\

\bitfields{ 5/11010, 2/01, 5/00001, 5/rs1, 3/frm, 5/frd, 7/1010011 }


\paragraph{Syntax} \texttt{fcvt.d.wu} \emph{frd}, \emph{rs1}, \emph{frm}

\paragraph{Description} 

\paragraph{Modifier(s)}
\noindent\begin{itemize}
\item frm (cf. froundmode in section \ref{par:modifier-froundmode} on page \pageref{par:modifier-froundmode})
\end{itemize}

\paragraph{Execution} {Reservation table \textsf{ALU\_LITE} (Section~\ref{sec:reservation-ALU-LITE})}
~
{\begin{lstlisting}
stage ID:
  new frm = _ZX_3(frm);
stage RR:
  new rs1 = BIR[rs1];
  new frm = decode_riscv_float_rounding_mode(frm);
  if (check_float_rounding_mode_trap(frm) == 0) {
    _THROW(OPCODE);
  } else {
    new frd = ui32_to_f64(frm, rs1);
    FPR[frd] = frd;
    commit_float_exception_flags();
  }
\end{lstlisting}}
\newpage
\subsubsection[{\small FCVT.S.D: Floating-Point Convert Single from Double}]{FCVT.S.D\hfill Floating-Point Convert Single from Double}

\paragraph{Encoding} Format \textsf{RV32D\_RTypeFCD2F} (Section~\ref{sec:format-RV32D-RTypeFCD2F}), \verb|rs2 = 00001|\\

\bitfields{ 5/01000, 2/00, 5/00001, 5/frs1, 3/frm, 5/frd, 7/1010011 }


\paragraph{Syntax} \texttt{fcvt.s.d} \emph{frd}, \emph{frs1}, \emph{frm}

\paragraph{Description} 

\paragraph{Modifier(s)}
\noindent\begin{itemize}
\item frm (cf. froundmode in section \ref{par:modifier-froundmode} on page \pageref{par:modifier-froundmode})
\end{itemize}

\paragraph{Execution} {Reservation table \textsf{ALU\_LITE} (Section~\ref{sec:reservation-ALU-LITE})}
~
{\begin{lstlisting}
stage ID:
  new frm = _ZX_3(frm);
stage RR:
  new frs1 = FPR[frs1];
  new frm = decode_riscv_float_rounding_mode(frm);
  if (check_float_rounding_mode_trap(frm) == 0) {
    _THROW(OPCODE);
  } else {
    new frd = f64_to_f32(frm, frs1) | 0xFFFFFFFF00000000;
    FPR[frd] = frd;
    commit_float_exception_flags();
  }
\end{lstlisting}}
\newpage
\subsubsection[{\small FCVT.W.D: Floating-Point Convert Word from Double}]{FCVT.W.D\hfill Floating-Point Convert Word from Double}

\paragraph{Encoding} Format \textsf{RV32D\_RTypeFCD2I} (Section~\ref{sec:format-RV32D-RTypeFCD2I}), \verb|rs2 = 00000|\\

\bitfields{ 5/11000, 2/01, 5/00000, 5/frs1, 3/frm, 5/rd, 7/1010011 }


\paragraph{Syntax} \texttt{fcvt.w.d} \emph{rd}, \emph{frs1}, \emph{frm}

\paragraph{Description} 

\paragraph{Modifier(s)}
\noindent\begin{itemize}
\item frm (cf. froundmode in section \ref{par:modifier-froundmode} on page \pageref{par:modifier-froundmode})
\end{itemize}

\paragraph{Execution} {Reservation table \textsf{ALU\_LITE} (Section~\ref{sec:reservation-ALU-LITE})}
~
{\begin{lstlisting}
stage ID:
  new frm = _ZX_3(frm);
stage RR:
  new frs1 = FPR[frs1];
  new frm = decode_riscv_float_rounding_mode(frm);
  if (check_float_rounding_mode_trap(frm) == 0) {
    _THROW(OPCODE);
  } else {
    new rd = _SX_32(f64_to_i32(frm, frs1));
    if (rd != 0) {
      BIR[rd] = rd;
    };
    commit_float_exception_flags();
  }
\end{lstlisting}}
\newpage
\subsubsection[{\small FCVT.WU.D: Floating-Point Convert Word Unsigned from Double}]{FCVT.WU.D\hfill Floating-Point Convert Word Unsigned from Double}

\paragraph{Encoding} Format \textsf{RV32D\_RTypeFCD2I} (Section~\ref{sec:format-RV32D-RTypeFCD2I}), \verb|rs2 = 00001|\\

\bitfields{ 5/11000, 2/01, 5/00001, 5/frs1, 3/frm, 5/rd, 7/1010011 }


\paragraph{Syntax} \texttt{fcvt.wu.d} \emph{rd}, \emph{frs1}, \emph{frm}

\paragraph{Description} 

\paragraph{Modifier(s)}
\noindent\begin{itemize}
\item frm (cf. froundmode in section \ref{par:modifier-froundmode} on page \pageref{par:modifier-froundmode})
\end{itemize}

\paragraph{Execution} {Reservation table \textsf{ALU\_LITE} (Section~\ref{sec:reservation-ALU-LITE})}
~
{\begin{lstlisting}
stage ID:
  new frm = _ZX_3(frm);
stage RR:
  new frs1 = FPR[frs1];
  new frm = decode_riscv_float_rounding_mode(frm);
  if (check_float_rounding_mode_trap(frm) == 0) {
    _THROW(OPCODE);
  } else {
    new rd = _SX_32(f64_to_ui32(frm, frs1));
    if (rd != 0) {
      BIR[rd] = rd;
    };
    commit_float_exception_flags();
  }
\end{lstlisting}}
\newpage
\subsubsection[{\small FDIV.D: Floating-Point Divide Double}]{FDIV.D\hfill Floating-Point Divide Double}

\paragraph{Encoding} Format \textsf{RV32D\_R3TypeRM} (Section~\ref{sec:format-RV32D-R3TypeRM}), \verb|funct5 = 00011|\\

\bitfields{ 5/00011, 2/01, 5/frs2, 5/frs1, 3/frm, 5/frd, 7/1010011 }


\paragraph{Syntax} \texttt{fdiv.d} \emph{frd}, \emph{frs1}, \emph{frs2}, \emph{frm}

\paragraph{Description} 

\paragraph{Modifier(s)}
\noindent\begin{itemize}
\item frm (cf. froundmode in section \ref{par:modifier-froundmode} on page \pageref{par:modifier-froundmode})
\end{itemize}

\paragraph{Execution} {Reservation table \textsf{ALU\_FULL} (Section~\ref{sec:reservation-ALU-FULL})}
~
{\begin{lstlisting}
stage ID:
  new frm = _ZX_3(frm);
stage RR:
  new frs2 = FPR[frs2];
  new frs1 = FPR[frs1];
  new frm = decode_riscv_float_rounding_mode(frm);
  if (check_float_rounding_mode_trap(frm) == 0) {
    _THROW(OPCODE);
  } else {
    new frd = f64_div(frm, frs1, frs2);
    FPR[frd] = frd;
    commit_float_exception_flags();
  }
\end{lstlisting}}
\newpage
\subsubsection[{\small FEQ.D: Floating-Point Compare Equal Double}]{FEQ.D\hfill Floating-Point Compare Equal Double}

\paragraph{Encoding} Format \textsf{RV32D\_RTypeFCMP} (Section~\ref{sec:format-RV32D-RTypeFCMP}), \verb|funct3 = 010|\\

\bitfields{ 5/10100, 2/01, 5/frs2, 5/frs1, 3/010, 5/rd, 7/1010011 }


\paragraph{Syntax} \texttt{feq.d} \emph{rd}, \emph{frs1}, \emph{frs2}

\paragraph{Description} 

\paragraph{Execution} {Reservation table \textsf{ALU\_LITE} (Section~\ref{sec:reservation-ALU-LITE})}
~
{\begin{lstlisting}
stage RR:
  new frs2 = FPR[frs2];
  new frs1 = FPR[frs1];
  new rd = f64_eq(frs1, frs2);
stage E1:
  if (rd != 0) {
    BIR[rd] = rd;
  }
stage E4:
  commit_float_exception_flags();
\end{lstlisting}}
\newpage
\subsubsection[{\small FLD: Floating-Point Load Double Word}]{FLD\hfill Floating-Point Load Double Word}

\paragraph{Encoding} Format \textsf{RV32D\_ILoad} (Section~\ref{sec:format-RV32D-ILoad}), \verb|funct3 = 011|\\

\bitfields{ 12/i-- 11-- 0, 5/rs1, 3/011, 5/frd, 7/0000111 }


\paragraph{Syntax} \texttt{fld} \emph{frd}, \emph{i\_\discretionary{}{}{}11\_\discretionary{}{}{}0}[\emph{rs1}]

\paragraph{Description} The \emph{frd} receives the memory double word at address \emph{rs1} + \emph{i\_\discretionary{}{}{}11\_\discretionary{}{}{}0}.


\paragraph{Execution} {Reservation table \textsf{LSU\_AUXW} (Section~\ref{sec:reservation-LSU-AUXW})}
~
{\begin{lstlisting}
stage RR:
  new signed12 = _SX_12(i_11_0);
  new rs1 = BIR[rs1];
  new address = rs1 + signed12;
  new frd = _ZX_64(MEM.load(address, 0xFF, 0, frd));
stage E1:
  FPR[frd] = frd;
\end{lstlisting}}
\newpage
\subsubsection[{\small FLE.D: Floating-Point Compare Less Than or Equal Double}]{FLE.D\hfill Floating-Point Compare Less Than or Equal Double}

\paragraph{Encoding} Format \textsf{RV32D\_RTypeFCMP} (Section~\ref{sec:format-RV32D-RTypeFCMP}), \verb|funct3 = 000|\\

\bitfields{ 5/10100, 2/01, 5/frs2, 5/frs1, 3/000, 5/rd, 7/1010011 }


\paragraph{Syntax} \texttt{fle.d} \emph{rd}, \emph{frs1}, \emph{frs2}

\paragraph{Description} 

\paragraph{Execution} {Reservation table \textsf{ALU\_LITE} (Section~\ref{sec:reservation-ALU-LITE})}
~
{\begin{lstlisting}
stage RR:
  new frs2 = FPR[frs2];
  new frs1 = FPR[frs1];
  new rd = f64_le(frs1, frs2);
stage E1:
  if (rd != 0) {
    BIR[rd] = rd;
  }
stage E4:
  commit_float_exception_flags();
\end{lstlisting}}
\newpage
\subsubsection[{\small FLT.D: Floating-Point Compare Less Than Double}]{FLT.D\hfill Floating-Point Compare Less Than Double}

\paragraph{Encoding} Format \textsf{RV32D\_RTypeFCMP} (Section~\ref{sec:format-RV32D-RTypeFCMP}), \verb|funct3 = 001|\\

\bitfields{ 5/10100, 2/01, 5/frs2, 5/frs1, 3/001, 5/rd, 7/1010011 }


\paragraph{Syntax} \texttt{flt.d} \emph{rd}, \emph{frs1}, \emph{frs2}

\paragraph{Description} 

\paragraph{Execution} {Reservation table \textsf{ALU\_LITE} (Section~\ref{sec:reservation-ALU-LITE})}
~
{\begin{lstlisting}
stage RR:
  new frs2 = FPR[frs2];
  new frs1 = FPR[frs1];
  new rd = f64_lt(frs1, frs2);
stage E1:
  if (rd != 0) {
    BIR[rd] = rd;
  }
stage E4:
  commit_float_exception_flags();
\end{lstlisting}}
\newpage
\subsubsection[{\small FMADD.D: Floating-Point Fused Multiply Add Double}]{FMADD.D\hfill Floating-Point Fused Multiply Add Double}

\paragraph{Encoding} Format \textsf{RV32D\_R4TypeRM} (Section~\ref{sec:format-RV32D-R4TypeRM}), \verb|opcode = 1000011|\\

\bitfields{ 5/frs3, 2/01, 5/frs2, 5/frs1, 3/frm, 5/frd, 7/1000011 }


\paragraph{Syntax} \texttt{fmadd.d} \emph{frd}, \emph{frs1}, \emph{frs2}, \emph{frs3}, \emph{frm}

\paragraph{Description} 

\paragraph{Modifier(s)}
\noindent\begin{itemize}
\item frm (cf. froundmode in section \ref{par:modifier-froundmode} on page \pageref{par:modifier-froundmode})
\end{itemize}

\paragraph{Execution} {Reservation table \textsf{ALU\_LITE} (Section~\ref{sec:reservation-ALU-LITE})}
~
{\begin{lstlisting}
stage ID:
  new frm = _ZX_3(frm);
stage RR:
  new frs3 = FPR[frs3];
  new frs2 = FPR[frs2];
  new frs1 = FPR[frs1];
  new frm = decode_riscv_float_rounding_mode(frm);
  if (check_float_rounding_mode_trap(frm) == 0) {
    _THROW(OPCODE);
  } else {
    new frd = f64_mulAdd(frm, frs1, frs2, frs3);
    FPR[frd] = frd;
    commit_float_exception_flags();
  }
\end{lstlisting}}
\newpage
\subsubsection[{\small FMAX.D: Floating-Point Maximum Double}]{FMAX.D\hfill Floating-Point Maximum Double}

\paragraph{Encoding} Format \textsf{RV32D\_RTypeFMIN} (Section~\ref{sec:format-RV32D-RTypeFMIN}), \verb|funct3 = 001|\\

\bitfields{ 5/00101, 2/01, 5/frs2, 5/frs1, 3/001, 5/frd, 7/1010011 }


\paragraph{Syntax} \texttt{fmax.d} \emph{frd}, \emph{frs1}, \emph{frs2}

\paragraph{Description} 

\paragraph{Execution} {Reservation table \textsf{ALU\_LITE} (Section~\ref{sec:reservation-ALU-LITE})}
~
{\begin{lstlisting}
stage RR:
  new frs2 = FPR[frs2];
  new frs1 = FPR[frs1];
  new frd = f64_maxNum(frs1, frs2);
stage E1:
  FPR[frd] = frd;
  commit_float_exception_flags();
\end{lstlisting}}
\newpage
\subsubsection[{\small FMIN.D: Floating-Point Minimum Double}]{FMIN.D\hfill Floating-Point Minimum Double}

\paragraph{Encoding} Format \textsf{RV32D\_RTypeFMIN} (Section~\ref{sec:format-RV32D-RTypeFMIN}), \verb|funct3 = 000|\\

\bitfields{ 5/00101, 2/01, 5/frs2, 5/frs1, 3/000, 5/frd, 7/1010011 }


\paragraph{Syntax} \texttt{fmin.d} \emph{frd}, \emph{frs1}, \emph{frs2}

\paragraph{Description} 

\paragraph{Execution} {Reservation table \textsf{ALU\_LITE} (Section~\ref{sec:reservation-ALU-LITE})}
~
{\begin{lstlisting}
stage RR:
  new frs2 = FPR[frs2];
  new frs1 = FPR[frs1];
  new frd = f64_minNum(frs1, frs2);
stage E1:
  FPR[frd] = frd;
  commit_float_exception_flags();
\end{lstlisting}}
\newpage
\subsubsection[{\small FMSUB.D: Floating-Point Fused Multiply Subtract Double}]{FMSUB.D\hfill Floating-Point Fused Multiply Subtract Double}

\paragraph{Encoding} Format \textsf{RV32D\_R4TypeRM} (Section~\ref{sec:format-RV32D-R4TypeRM}), \verb|opcode = 1000111|\\

\bitfields{ 5/frs3, 2/01, 5/frs2, 5/frs1, 3/frm, 5/frd, 7/1000111 }


\paragraph{Syntax} \texttt{fmsub.d} \emph{frd}, \emph{frs1}, \emph{frs2}, \emph{frs3}, \emph{frm}

\paragraph{Description} 

\paragraph{Modifier(s)}
\noindent\begin{itemize}
\item frm (cf. froundmode in section \ref{par:modifier-froundmode} on page \pageref{par:modifier-froundmode})
\end{itemize}

\paragraph{Execution} {Reservation table \textsf{ALU\_LITE} (Section~\ref{sec:reservation-ALU-LITE})}
~
{\begin{lstlisting}
stage ID:
  new frm = _ZX_3(frm);
stage RR:
  new frs3 = FPR[frs3];
  new frs2 = FPR[frs2];
  new frs1 = FPR[frs1];
  new frm = decode_riscv_float_rounding_mode(frm);
  if (check_float_rounding_mode_trap(frm) == 0) {
    _THROW(OPCODE);
  } else {
    new frd = f64_mulAdd(frm, frs1, frs2, frs3 ^ 0x8000000000000000);
    FPR[frd] = frd;
    commit_float_exception_flags();
  }
\end{lstlisting}}
\newpage
\subsubsection[{\small FMUL.D: Floating-Point Multiply Double}]{FMUL.D\hfill Floating-Point Multiply Double}

\paragraph{Encoding} Format \textsf{RV32D\_R3TypeRM} (Section~\ref{sec:format-RV32D-R3TypeRM}), \verb|funct5 = 00010|\\

\bitfields{ 5/00010, 2/01, 5/frs2, 5/frs1, 3/frm, 5/frd, 7/1010011 }


\paragraph{Syntax} \texttt{fmul.d} \emph{frd}, \emph{frs1}, \emph{frs2}, \emph{frm}

\paragraph{Description} 

\paragraph{Modifier(s)}
\noindent\begin{itemize}
\item frm (cf. froundmode in section \ref{par:modifier-froundmode} on page \pageref{par:modifier-froundmode})
\end{itemize}

\paragraph{Execution} {Reservation table \textsf{ALU\_LITE} (Section~\ref{sec:reservation-ALU-LITE})}
~
{\begin{lstlisting}
stage ID:
  new frm = _ZX_3(frm);
stage RR:
  new frs2 = FPR[frs2];
  new frs1 = FPR[frs1];
  new frm = decode_riscv_float_rounding_mode(frm);
  if (check_float_rounding_mode_trap(frm) == 0) {
    _THROW(OPCODE);
  } else {
    new frd = f64_mul(frm, frs1, frs2);
    FPR[frd] = frd;
    commit_float_exception_flags();
  }
\end{lstlisting}}
\newpage
\subsubsection[{\small FNMADD.D: Floating-Point Fused Negate Multiply Add Double}]{FNMADD.D\hfill Floating-Point Fused Negate Multiply Add Double}

\paragraph{Encoding} Format \textsf{RV32D\_R4TypeRM} (Section~\ref{sec:format-RV32D-R4TypeRM}), \verb|opcode = 1001111|\\

\bitfields{ 5/frs3, 2/01, 5/frs2, 5/frs1, 3/frm, 5/frd, 7/1001111 }


\paragraph{Syntax} \texttt{fnmadd.d} \emph{frd}, \emph{frs1}, \emph{frs2}, \emph{frs3}, \emph{frm}

\paragraph{Description} 

\paragraph{Modifier(s)}
\noindent\begin{itemize}
\item frm (cf. froundmode in section \ref{par:modifier-froundmode} on page \pageref{par:modifier-froundmode})
\end{itemize}

\paragraph{Execution} {Reservation table \textsf{ALU\_LITE} (Section~\ref{sec:reservation-ALU-LITE})}
~
{\begin{lstlisting}
stage ID:
  new frm = _ZX_3(frm);
stage RR:
  new frs3 = FPR[frs3];
  new frs2 = FPR[frs2];
  new frs1 = FPR[frs1];
  new frm = decode_riscv_float_rounding_mode(frm);
  if (check_float_rounding_mode_trap(frm) == 0) {
    _THROW(OPCODE);
  } else {
    new frd = f64_mulAdd(frm, frs1 ^ 0x8000000000000000, frs2, frs3 ^ 0x8000000000000000);
    FPR[frd] = frd;
    commit_float_exception_flags();
  }
\end{lstlisting}}
\newpage
\subsubsection[{\small FNMSUB.D: Floating-Point Fused Negate Multiply Subtract Double}]{FNMSUB.D\hfill Floating-Point Fused Negate Multiply Subtract Double}

\paragraph{Encoding} Format \textsf{RV32D\_R4TypeRM} (Section~\ref{sec:format-RV32D-R4TypeRM}), \verb|opcode = 1001011|\\

\bitfields{ 5/frs3, 2/01, 5/frs2, 5/frs1, 3/frm, 5/frd, 7/1001011 }


\paragraph{Syntax} \texttt{fnmsub.d} \emph{frd}, \emph{frs1}, \emph{frs2}, \emph{frs3}, \emph{frm}

\paragraph{Description} 

\paragraph{Modifier(s)}
\noindent\begin{itemize}
\item frm (cf. froundmode in section \ref{par:modifier-froundmode} on page \pageref{par:modifier-froundmode})
\end{itemize}

\paragraph{Execution} {Reservation table \textsf{ALU\_LITE} (Section~\ref{sec:reservation-ALU-LITE})}
~
{\begin{lstlisting}
stage ID:
  new frm = _ZX_3(frm);
stage RR:
  new frs3 = FPR[frs3];
  new frs2 = FPR[frs2];
  new frs1 = FPR[frs1];
  new frm = decode_riscv_float_rounding_mode(frm);
  if (check_float_rounding_mode_trap(frm) == 0) {
    _THROW(OPCODE);
  } else {
    new frd = f64_mulAdd(frm, frs1 ^ 0x8000000000000000, frs2, frs3);
    FPR[frd] = frd;
    commit_float_exception_flags();
  }
\end{lstlisting}}
\newpage
\subsubsection[{\small FSD: Floating-Point Store Double Word}]{FSD\hfill Floating-Point Store Double Word}

\paragraph{Encoding} Format \textsf{RV32D\_SType} (Section~\ref{sec:format-RV32D-SType}), \verb|funct3 = 011|\\

\bitfields{ 7/i-- 11-- 5, 5/frs2, 5/rs1, 3/011, 5/i-- 4-- 0s, 7/0100111 }


\paragraph{Syntax} \texttt{fsd} \emph{frs2}, \emph{i\_\discretionary{}{}{}11\_\discretionary{}{}{}5\_\discretionary{}{}{}i\_\discretionary{}{}{}4\_\discretionary{}{}{}0s}[\emph{rs1}]

\paragraph{Description} The \emph{frs2} is stored into the memory double word at address \emph{rs1} + \emph{i\_\discretionary{}{}{}11\_\discretionary{}{}{}5\_\discretionary{}{}{}i\_\discretionary{}{}{}4\_\discretionary{}{}{}0s}.


\paragraph{Execution} {Reservation table \textsf{LSU\_MEMW\_AUXR} (Section~\ref{sec:reservation-LSU-MEMW-AUXR})}
~
{\begin{lstlisting}
stage RR:
  new signed12 = _SX_12(i_11_5_i_4_0s);
  new rs1 = BIR[rs1];
  new address = rs1 + signed12;
stage E1:
  new frs2 = FPR[frs2];
stage E3:
  MEM.store(address, 0xFF, 0, frs2, frs2);
\end{lstlisting}}
\newpage
\subsubsection[{\small FSGNJ.D: Floating-Point Sign Inject Double}]{FSGNJ.D\hfill Floating-Point Sign Inject Double}

\paragraph{Encoding} Format \textsf{RV32D\_RTypeFSGN} (Section~\ref{sec:format-RV32D-RTypeFSGN}), \verb|funct3 = 000|\\

\bitfields{ 5/00100, 2/01, 5/frs2, 5/frs1, 3/000, 5/frd, 7/1010011 }


\paragraph{Syntax} \texttt{fsgnj.d} \emph{frd}, \emph{frs1}, \emph{frs2}

\paragraph{Description} 

\paragraph{Execution} {Reservation table \textsf{ALU\_LITE} (Section~\ref{sec:reservation-ALU-LITE})}
~
{\begin{lstlisting}
stage RR:
  new frs2 = FPR[frs2];
  new frs1 = FPR[frs1];
  new frd = (frs1 & 0x7FFFFFFFFFFFFFFF) | (frs2 & 0x8000000000000000);
stage E1:
  FPR[frd] = frd;
\end{lstlisting}}
\newpage
\subsubsection[{\small FSGNJN.D: Floating-Point Sign Inject Negate Double}]{FSGNJN.D\hfill Floating-Point Sign Inject Negate Double}

\paragraph{Encoding} Format \textsf{RV32D\_RTypeFSGN} (Section~\ref{sec:format-RV32D-RTypeFSGN}), \verb|funct3 = 001|\\

\bitfields{ 5/00100, 2/01, 5/frs2, 5/frs1, 3/001, 5/frd, 7/1010011 }


\paragraph{Syntax} \texttt{fsgnjn.d} \emph{frd}, \emph{frs1}, \emph{frs2}

\paragraph{Description} 

\paragraph{Execution} {Reservation table \textsf{ALU\_LITE} (Section~\ref{sec:reservation-ALU-LITE})}
~
{\begin{lstlisting}
stage RR:
  new frs2 = FPR[frs2];
  new frs1 = FPR[frs1];
  new frd = (frs1 & 0x7FFFFFFFFFFFFFFF) | (~frs2 & 0x8000000000000000);
stage E1:
  FPR[frd] = frd;
\end{lstlisting}}
\newpage
\subsubsection[{\small FSGNJX.D: Floating-Point Sign Inject XOR Double}]{FSGNJX.D\hfill Floating-Point Sign Inject XOR Double}

\paragraph{Encoding} Format \textsf{RV32D\_RTypeFSGN} (Section~\ref{sec:format-RV32D-RTypeFSGN}), \verb|funct3 = 010|\\

\bitfields{ 5/00100, 2/01, 5/frs2, 5/frs1, 3/010, 5/frd, 7/1010011 }


\paragraph{Syntax} \texttt{fsgnjx.d} \emph{frd}, \emph{frs1}, \emph{frs2}

\paragraph{Description} 

\paragraph{Execution} {Reservation table \textsf{ALU\_LITE} (Section~\ref{sec:reservation-ALU-LITE})}
~
{\begin{lstlisting}
stage RR:
  new frs2 = FPR[frs2];
  new frs1 = FPR[frs1];
  new frd = frs1 ^ (frs2 & 0x8000000000000000);
stage E1:
  FPR[frd] = frd;
\end{lstlisting}}
\newpage
\subsubsection[{\small FSQRT.D: Floating-Point Square Root Double}]{FSQRT.D\hfill Floating-Point Square Root Double}

\paragraph{Encoding} Format \textsf{RV32D\_R2TypeRM} (Section~\ref{sec:format-RV32D-R2TypeRM}), \verb|funct5 = 01011|\\

\bitfields{ 5/01011, 2/01, 5/00000, 5/frs1, 3/frm, 5/frd, 7/1010011 }


\paragraph{Syntax} \texttt{fsqrt.d} \emph{frd}, \emph{frs1}, \emph{frm}

\paragraph{Description} 

\paragraph{Modifier(s)}
\noindent\begin{itemize}
\item frm (cf. froundmode in section \ref{par:modifier-froundmode} on page \pageref{par:modifier-froundmode})
\end{itemize}

\paragraph{Execution} {Reservation table \textsf{ALU\_FULL} (Section~\ref{sec:reservation-ALU-FULL})}
~
{\begin{lstlisting}
stage ID:
  new frm = _ZX_3(frm);
stage RR:
  new frs1 = FPR[frs1];
  new frm = decode_riscv_float_rounding_mode(frm);
  if (check_float_rounding_mode_trap(frm) == 0) {
    _THROW(OPCODE);
  } else {
    new frd = f64_sqrt(frm, frs1);
    FPR[frd] = frd;
    commit_float_exception_flags();
  }
\end{lstlisting}}
\newpage
\subsubsection[{\small FSUB.D: Floating-Point Subtract Double}]{FSUB.D\hfill Floating-Point Subtract Double}

\paragraph{Encoding} Format \textsf{RV32D\_R3TypeRM} (Section~\ref{sec:format-RV32D-R3TypeRM}), \verb|funct5 = 00001|\\

\bitfields{ 5/00001, 2/01, 5/frs2, 5/frs1, 3/frm, 5/frd, 7/1010011 }


\paragraph{Syntax} \texttt{fsub.d} \emph{frd}, \emph{frs1}, \emph{frs2}, \emph{frm}

\paragraph{Description} 

\paragraph{Modifier(s)}
\noindent\begin{itemize}
\item frm (cf. froundmode in section \ref{par:modifier-froundmode} on page \pageref{par:modifier-froundmode})
\end{itemize}

\paragraph{Execution} {Reservation table \textsf{ALU\_LITE} (Section~\ref{sec:reservation-ALU-LITE})}
~
{\begin{lstlisting}
stage ID:
  new frm = _ZX_3(frm);
stage RR:
  new frs2 = FPR[frs2];
  new frs1 = FPR[frs1];
  new frm = decode_riscv_float_rounding_mode(frm);
  if (check_float_rounding_mode_trap(frm) == 0) {
    _THROW(OPCODE);
  } else {
    new frd = f64_sub(frm, frs1, frs2);
    FPR[frd] = frd;
    commit_float_exception_flags();
  }
\end{lstlisting}}
\newpage
\subsection{RV32F Instructions}
\label{sec:RV32F-instructions}

\subsubsection[{\small FADD.S: Floating-Point Add Single}]{FADD.S\hfill Floating-Point Add Single}

\paragraph{Encoding} Format \textsf{RV32F\_R3TypeRM} (Section~\ref{sec:format-RV32F-R3TypeRM}), \verb|funct5 = 00000|\\

\bitfields{ 5/00000, 2/00, 5/frs2, 5/frs1, 3/frm, 5/frd, 7/1010011 }


\paragraph{Syntax} \texttt{fadd.s} \emph{frd}, \emph{frs1}, \emph{frs2}, \emph{frm}

\paragraph{Description} 

\paragraph{Modifier(s)}
\noindent\begin{itemize}
\item frm (cf. froundmode in section \ref{par:modifier-froundmode} on page \pageref{par:modifier-froundmode})
\end{itemize}

\paragraph{Execution} {Reservation table \textsf{ALU\_LITE} (Section~\ref{sec:reservation-ALU-LITE})}
~
{\begin{lstlisting}
stage ID:
  new frm = _ZX_3(frm);
stage RR:
  new frs2 = FPR[frs2];
  new frs1 = FPR[frs1];
  new frm = decode_riscv_float_rounding_mode(frm);
  if (check_float_rounding_mode_trap(frm) == 0) {
    _THROW(OPCODE);
  } else {
    new frd = f32_add(frm, frs1, frs2) | 0xFFFFFFFF00000000;
    FPR[frd] = frd;
    commit_float_exception_flags();
  }
\end{lstlisting}}
\newpage
\subsubsection[{\small FCLASS.S: Floating-Point Classify Single}]{FCLASS.S\hfill Floating-Point Classify Single}

\paragraph{Encoding} Format \textsf{RV32F\_RTypeFF2I} (Section~\ref{sec:format-RV32F-RTypeFF2I}), \verb|funct3 = 001|\\

\bitfields{ 5/11100, 2/00, 5/00000, 5/frs1, 3/001, 5/rd, 7/1010011 }


\paragraph{Syntax} \texttt{fclass.s} \emph{rd}, \emph{frs1}

\paragraph{Description} 

\paragraph{Execution} {Reservation table \textsf{ALU\_LITE} (Section~\ref{sec:reservation-ALU-LITE})}
~
{\begin{lstlisting}
stage RR:
  new frs1 = FPR[frs1];
  new rd = f32_classify(frs1);
stage E1:
  if (rd != 0) {
    BIR[rd] = rd;
  }
\end{lstlisting}}
\newpage
\subsubsection[{\small FCVT.S.W: Floating-Point Convert Single from Word}]{FCVT.S.W\hfill Floating-Point Convert Single from Word}

\paragraph{Encoding} Format \textsf{RV32F\_RTypeFCI2F} (Section~\ref{sec:format-RV32F-RTypeFCI2F}), \verb|rs2 = 00000|\\

\bitfields{ 5/11010, 2/00, 5/00000, 5/rs1, 3/frm, 5/frd, 7/1010011 }


\paragraph{Syntax} \texttt{fcvt.s.w} \emph{frd}, \emph{rs1}, \emph{frm}

\paragraph{Description} 

\paragraph{Modifier(s)}
\noindent\begin{itemize}
\item frm (cf. froundmode in section \ref{par:modifier-froundmode} on page \pageref{par:modifier-froundmode})
\end{itemize}

\paragraph{Execution} {Reservation table \textsf{ALU\_LITE} (Section~\ref{sec:reservation-ALU-LITE})}
~
{\begin{lstlisting}
stage ID:
  new frm = _ZX_3(frm);
stage RR:
  new rs1 = BIR[rs1];
  new frm = decode_riscv_float_rounding_mode(frm);
  if (check_float_rounding_mode_trap(frm) == 0) {
    _THROW(OPCODE);
  } else {
    new frd = i32_to_f32(frm, rs1) | 0xFFFFFFFF00000000;
    FPR[frd] = frd;
    commit_float_exception_flags();
  }
\end{lstlisting}}
\newpage
\subsubsection[{\small FCVT.S.WU: Floating-Point Convert Single from Word Unsigned}]{FCVT.S.WU\hfill Floating-Point Convert Single from Word Unsigned}

\paragraph{Encoding} Format \textsf{RV32F\_RTypeFCI2F} (Section~\ref{sec:format-RV32F-RTypeFCI2F}), \verb|rs2 = 00001|\\

\bitfields{ 5/11010, 2/00, 5/00001, 5/rs1, 3/frm, 5/frd, 7/1010011 }


\paragraph{Syntax} \texttt{fcvt.s.wu} \emph{frd}, \emph{rs1}, \emph{frm}

\paragraph{Description} 

\paragraph{Modifier(s)}
\noindent\begin{itemize}
\item frm (cf. froundmode in section \ref{par:modifier-froundmode} on page \pageref{par:modifier-froundmode})
\end{itemize}

\paragraph{Execution} {Reservation table \textsf{ALU\_LITE} (Section~\ref{sec:reservation-ALU-LITE})}
~
{\begin{lstlisting}
stage ID:
  new frm = _ZX_3(frm);
stage RR:
  new rs1 = BIR[rs1];
  new frm = decode_riscv_float_rounding_mode(frm);
  if (check_float_rounding_mode_trap(frm) == 0) {
    _THROW(OPCODE);
  } else {
    new frd = ui32_to_f32(frm, rs1) | 0xFFFFFFFF00000000;
    FPR[frd] = frd;
    commit_float_exception_flags();
  }
\end{lstlisting}}
\newpage
\subsubsection[{\small FCVT.W.S: Floating-Point Convert Word from Single}]{FCVT.W.S\hfill Floating-Point Convert Word from Single}

\paragraph{Encoding} Format \textsf{RV32F\_RTypeFCF2I} (Section~\ref{sec:format-RV32F-RTypeFCF2I}), \verb|rs2 = 00000|\\

\bitfields{ 5/11000, 2/00, 5/00000, 5/frs1, 3/frm, 5/rd, 7/1010011 }


\paragraph{Syntax} \texttt{fcvt.w.s} \emph{rd}, \emph{frs1}, \emph{frm}

\paragraph{Description} 

\paragraph{Modifier(s)}
\noindent\begin{itemize}
\item frm (cf. froundmode in section \ref{par:modifier-froundmode} on page \pageref{par:modifier-froundmode})
\end{itemize}

\paragraph{Execution} {Reservation table \textsf{ALU\_LITE} (Section~\ref{sec:reservation-ALU-LITE})}
~
{\begin{lstlisting}
stage ID:
  new frm = _ZX_3(frm);
stage RR:
  new frs1 = FPR[frs1];
  new frm = decode_riscv_float_rounding_mode(frm);
  if (check_float_rounding_mode_trap(frm) == 0) {
    _THROW(OPCODE);
  } else {
    new rd = _SX_32(f32_to_i32(frm, frs1));
    if (rd != 0) {
      BIR[rd] = rd;
    };
    commit_float_exception_flags();
  }
\end{lstlisting}}
\newpage
\subsubsection[{\small FCVT.WU.S: Floating-Point Convert Word Unsigned from Single}]{FCVT.WU.S\hfill Floating-Point Convert Word Unsigned from Single}

\paragraph{Encoding} Format \textsf{RV32F\_RTypeFCF2I} (Section~\ref{sec:format-RV32F-RTypeFCF2I}), \verb|rs2 = 00001|\\

\bitfields{ 5/11000, 2/00, 5/00001, 5/frs1, 3/frm, 5/rd, 7/1010011 }


\paragraph{Syntax} \texttt{fcvt.wu.s} \emph{rd}, \emph{frs1}, \emph{frm}

\paragraph{Description} 

\paragraph{Modifier(s)}
\noindent\begin{itemize}
\item frm (cf. froundmode in section \ref{par:modifier-froundmode} on page \pageref{par:modifier-froundmode})
\end{itemize}

\paragraph{Execution} {Reservation table \textsf{ALU\_LITE} (Section~\ref{sec:reservation-ALU-LITE})}
~
{\begin{lstlisting}
stage ID:
  new frm = _ZX_3(frm);
stage RR:
  new frs1 = FPR[frs1];
  new frm = decode_riscv_float_rounding_mode(frm);
  if (check_float_rounding_mode_trap(frm) == 0) {
    _THROW(OPCODE);
  } else {
    new rd = _SX_32(f32_to_ui32(frm, frs1));
    if (rd != 0) {
      BIR[rd] = rd;
    };
    commit_float_exception_flags();
  }
\end{lstlisting}}
\newpage
\subsubsection[{\small FDIV.S: Floating-Point Divide Single}]{FDIV.S\hfill Floating-Point Divide Single}

\paragraph{Encoding} Format \textsf{RV32F\_R3TypeRM} (Section~\ref{sec:format-RV32F-R3TypeRM}), \verb|funct5 = 00011|\\

\bitfields{ 5/00011, 2/00, 5/frs2, 5/frs1, 3/frm, 5/frd, 7/1010011 }


\paragraph{Syntax} \texttt{fdiv.s} \emph{frd}, \emph{frs1}, \emph{frs2}, \emph{frm}

\paragraph{Description} 

\paragraph{Modifier(s)}
\noindent\begin{itemize}
\item frm (cf. froundmode in section \ref{par:modifier-froundmode} on page \pageref{par:modifier-froundmode})
\end{itemize}

\paragraph{Execution} {Reservation table \textsf{ALU\_FULL} (Section~\ref{sec:reservation-ALU-FULL})}
~
{\begin{lstlisting}
stage ID:
  new frm = _ZX_3(frm);
stage RR:
  new frs2 = FPR[frs2];
  new frs1 = FPR[frs1];
  new frm = decode_riscv_float_rounding_mode(frm);
  if (check_float_rounding_mode_trap(frm) == 0) {
    _THROW(OPCODE);
  } else {
    new frd = f32_div(frm, frs1, frs2) | 0xFFFFFFFF00000000;
    FPR[frd] = frd;
    commit_float_exception_flags();
  }
\end{lstlisting}}
\newpage
\subsubsection[{\small FEQ.S: Floating-Point Compare Equal Single}]{FEQ.S\hfill Floating-Point Compare Equal Single}

\paragraph{Encoding} Format \textsf{RV32F\_RTypeFCMP} (Section~\ref{sec:format-RV32F-RTypeFCMP}), \verb|funct3 = 010|\\

\bitfields{ 5/10100, 2/00, 5/frs2, 5/frs1, 3/010, 5/rd, 7/1010011 }


\paragraph{Syntax} \texttt{feq.s} \emph{rd}, \emph{frs1}, \emph{frs2}

\paragraph{Description} 

\paragraph{Execution} {Reservation table \textsf{ALU\_LITE} (Section~\ref{sec:reservation-ALU-LITE})}
~
{\begin{lstlisting}
stage RR:
  new frs2 = FPR[frs2];
  new frs1 = FPR[frs1];
  new rd = f32_eq(frs1, frs2);
stage E1:
  if (rd != 0) {
    BIR[rd] = rd;
  }
stage E4:
  commit_float_exception_flags();
\end{lstlisting}}
\newpage
\subsubsection[{\small FLE.S: Floating-Point Compare Less Than or Equal Single}]{FLE.S\hfill Floating-Point Compare Less Than or Equal Single}

\paragraph{Encoding} Format \textsf{RV32F\_RTypeFCMP} (Section~\ref{sec:format-RV32F-RTypeFCMP}), \verb|funct3 = 000|\\

\bitfields{ 5/10100, 2/00, 5/frs2, 5/frs1, 3/000, 5/rd, 7/1010011 }


\paragraph{Syntax} \texttt{fle.s} \emph{rd}, \emph{frs1}, \emph{frs2}

\paragraph{Description} 

\paragraph{Execution} {Reservation table \textsf{ALU\_LITE} (Section~\ref{sec:reservation-ALU-LITE})}
~
{\begin{lstlisting}
stage RR:
  new frs2 = FPR[frs2];
  new frs1 = FPR[frs1];
  new rd = f32_le(frs1, frs2);
stage E1:
  if (rd != 0) {
    BIR[rd] = rd;
  }
stage E4:
  commit_float_exception_flags();
\end{lstlisting}}
\newpage
\subsubsection[{\small FLT.S: Floating-Point Compare Less Than Single}]{FLT.S\hfill Floating-Point Compare Less Than Single}

\paragraph{Encoding} Format \textsf{RV32F\_RTypeFCMP} (Section~\ref{sec:format-RV32F-RTypeFCMP}), \verb|funct3 = 001|\\

\bitfields{ 5/10100, 2/00, 5/frs2, 5/frs1, 3/001, 5/rd, 7/1010011 }


\paragraph{Syntax} \texttt{flt.s} \emph{rd}, \emph{frs1}, \emph{frs2}

\paragraph{Description} 

\paragraph{Execution} {Reservation table \textsf{ALU\_LITE} (Section~\ref{sec:reservation-ALU-LITE})}
~
{\begin{lstlisting}
stage RR:
  new frs2 = FPR[frs2];
  new frs1 = FPR[frs1];
  new rd = f32_lt(frs1, frs2);
stage E1:
  if (rd != 0) {
    BIR[rd] = rd;
  }
stage E4:
  commit_float_exception_flags();
\end{lstlisting}}
\newpage
\subsubsection[{\small FLW: Floating-Point Load Word}]{FLW\hfill Floating-Point Load Word}

\paragraph{Encoding} Format \textsf{RV32F\_ILoad} (Section~\ref{sec:format-RV32F-ILoad}), \verb|funct3 = 010|\\

\bitfields{ 12/i-- 11-- 0, 5/rs1, 3/010, 5/frd, 7/0000111 }


\paragraph{Syntax} \texttt{flw} \emph{frd}, \emph{i\_\discretionary{}{}{}11\_\discretionary{}{}{}0}[\emph{rs1}]

\paragraph{Description} The \emph{frd} receives the memory word at address \emph{rs1} + \emph{i\_\discretionary{}{}{}11\_\discretionary{}{}{}0}.


\paragraph{Execution} {Reservation table \textsf{LSU\_AUXW} (Section~\ref{sec:reservation-LSU-AUXW})}
~
{\begin{lstlisting}
stage RR:
  new signed12 = _SX_12(i_11_0);
  new rs1 = BIR[rs1];
  new address = rs1 + signed12;
  new frd = _ZX_32(MEM.load(address, 0xF, 0, frd)) | 0xFFFFFFFF00000000;
stage E1:
  FPR[frd] = frd;
\end{lstlisting}}
\newpage
\subsubsection[{\small FMADD.S: Floating-Point Fused Multiply Add Single}]{FMADD.S\hfill Floating-Point Fused Multiply Add Single}

\paragraph{Encoding} Format \textsf{RV32F\_R4TypeRM} (Section~\ref{sec:format-RV32F-R4TypeRM}), \verb|opcode = 1000011|\\

\bitfields{ 5/frs3, 2/00, 5/frs2, 5/frs1, 3/frm, 5/frd, 7/1000011 }


\paragraph{Syntax} \texttt{fmadd.s} \emph{frd}, \emph{frs1}, \emph{frs2}, \emph{frs3}, \emph{frm}

\paragraph{Description} 

\paragraph{Modifier(s)}
\noindent\begin{itemize}
\item frm (cf. froundmode in section \ref{par:modifier-froundmode} on page \pageref{par:modifier-froundmode})
\end{itemize}

\paragraph{Execution} {Reservation table \textsf{ALU\_LITE} (Section~\ref{sec:reservation-ALU-LITE})}
~
{\begin{lstlisting}
stage ID:
  new frm = _ZX_3(frm);
stage RR:
  new frs3 = FPR[frs3];
  new frs2 = FPR[frs2];
  new frs1 = FPR[frs1];
  new frm = decode_riscv_float_rounding_mode(frm);
  if (check_float_rounding_mode_trap(frm) == 0) {
    _THROW(OPCODE);
  } else {
    new frd = f32_mulAdd(frm, frs1, frs2, frs3) | 0xFFFFFFFF00000000;
    FPR[frd] = frd;
    commit_float_exception_flags();
  }
\end{lstlisting}}
\newpage
\subsubsection[{\small FMAX.S: Floating-Point Maximum Single}]{FMAX.S\hfill Floating-Point Maximum Single}

\paragraph{Encoding} Format \textsf{RV32F\_RTypeFMIN} (Section~\ref{sec:format-RV32F-RTypeFMIN}), \verb|funct3 = 001|\\

\bitfields{ 5/00101, 2/00, 5/frs2, 5/frs1, 3/001, 5/frd, 7/1010011 }


\paragraph{Syntax} \texttt{fmax.s} \emph{frd}, \emph{frs1}, \emph{frs2}

\paragraph{Description} 

\paragraph{Execution} {Reservation table \textsf{ALU\_LITE} (Section~\ref{sec:reservation-ALU-LITE})}
~
{\begin{lstlisting}
stage RR:
  new frs2 = FPR[frs2];
  new frs1 = FPR[frs1];
  new frd = f32_maxNum(frs1, frs2) | 0xFFFFFFFF00000000;
stage E1:
  FPR[frd] = frd;
  commit_float_exception_flags();
\end{lstlisting}}
\newpage
\subsubsection[{\small FMIN.S: Floating-Point Minimum Single}]{FMIN.S\hfill Floating-Point Minimum Single}

\paragraph{Encoding} Format \textsf{RV32F\_RTypeFMIN} (Section~\ref{sec:format-RV32F-RTypeFMIN}), \verb|funct3 = 000|\\

\bitfields{ 5/00101, 2/00, 5/frs2, 5/frs1, 3/000, 5/frd, 7/1010011 }


\paragraph{Syntax} \texttt{fmin.s} \emph{frd}, \emph{frs1}, \emph{frs2}

\paragraph{Description} 

\paragraph{Execution} {Reservation table \textsf{ALU\_LITE} (Section~\ref{sec:reservation-ALU-LITE})}
~
{\begin{lstlisting}
stage RR:
  new frs2 = FPR[frs2];
  new frs1 = FPR[frs1];
  new frd = f32_minNum(frs1, frs2) | 0xFFFFFFFF00000000;
stage E1:
  FPR[frd] = frd;
  commit_float_exception_flags();
\end{lstlisting}}
\newpage
\subsubsection[{\small FMSUB.S: Floating-Point Fused Multiply Subtract Single}]{FMSUB.S\hfill Floating-Point Fused Multiply Subtract Single}

\paragraph{Encoding} Format \textsf{RV32F\_R4TypeRM} (Section~\ref{sec:format-RV32F-R4TypeRM}), \verb|opcode = 1000111|\\

\bitfields{ 5/frs3, 2/00, 5/frs2, 5/frs1, 3/frm, 5/frd, 7/1000111 }


\paragraph{Syntax} \texttt{fmsub.s} \emph{frd}, \emph{frs1}, \emph{frs2}, \emph{frs3}, \emph{frm}

\paragraph{Description} 

\paragraph{Modifier(s)}
\noindent\begin{itemize}
\item frm (cf. froundmode in section \ref{par:modifier-froundmode} on page \pageref{par:modifier-froundmode})
\end{itemize}

\paragraph{Execution} {Reservation table \textsf{ALU\_LITE} (Section~\ref{sec:reservation-ALU-LITE})}
~
{\begin{lstlisting}
stage ID:
  new frm = _ZX_3(frm);
stage RR:
  new frs3 = FPR[frs3];
  new frs2 = FPR[frs2];
  new frs1 = FPR[frs1];
  new frm = decode_riscv_float_rounding_mode(frm);
  if (check_float_rounding_mode_trap(frm) == 0) {
    _THROW(OPCODE);
  } else {
    new frd = f32_mulAdd(frm, frs1, frs2, frs3 ^ 0x80000000) | 0xFFFFFFFF00000000;
    FPR[frd] = frd;
    commit_float_exception_flags();
  }
\end{lstlisting}}
\newpage
\subsubsection[{\small FMUL.S: Floating-Point Multiply Single}]{FMUL.S\hfill Floating-Point Multiply Single}

\paragraph{Encoding} Format \textsf{RV32F\_R3TypeRM} (Section~\ref{sec:format-RV32F-R3TypeRM}), \verb|funct5 = 00010|\\

\bitfields{ 5/00010, 2/00, 5/frs2, 5/frs1, 3/frm, 5/frd, 7/1010011 }


\paragraph{Syntax} \texttt{fmul.s} \emph{frd}, \emph{frs1}, \emph{frs2}, \emph{frm}

\paragraph{Description} 

\paragraph{Modifier(s)}
\noindent\begin{itemize}
\item frm (cf. froundmode in section \ref{par:modifier-froundmode} on page \pageref{par:modifier-froundmode})
\end{itemize}

\paragraph{Execution} {Reservation table \textsf{ALU\_LITE} (Section~\ref{sec:reservation-ALU-LITE})}
~
{\begin{lstlisting}
stage ID:
  new frm = _ZX_3(frm);
stage RR:
  new frs2 = FPR[frs2];
  new frs1 = FPR[frs1];
  new frm = decode_riscv_float_rounding_mode(frm);
  if (check_float_rounding_mode_trap(frm) == 0) {
    _THROW(OPCODE);
  } else {
    new frd = f32_mul(frm, frs1, frs2) | 0xFFFFFFFF00000000;
    FPR[frd] = frd;
    commit_float_exception_flags();
  }
\end{lstlisting}}
\newpage
\subsubsection[{\small FMV.W.X: Floating-Point Move Single from Integer}]{FMV.W.X\hfill Floating-Point Move Single from Integer}

\paragraph{Encoding} Format \textsf{RV32F\_RTypeFMI2F} (Section~\ref{sec:format-RV32F-RTypeFMI2F}), \verb|funct3 = 000|\\

\bitfields{ 5/11110, 2/00, 5/00000, 5/rs1, 3/000, 5/frd, 7/1010011 }


\paragraph{Syntax} \texttt{fmv.w.x} \emph{frd}, \emph{rs1}

\paragraph{Description} 

\paragraph{Execution} {Reservation table \textsf{ALU\_TINY} (Section~\ref{sec:reservation-ALU-TINY})}
~
{\begin{lstlisting}
stage RR:
  new rs1 = BIR[rs1];
  new frd = rs1 | 0xFFFFFFFF00000000;
stage E1:
  FPR[frd] = frd;
\end{lstlisting}}
\newpage
\subsubsection[{\small FMV.X.W: Floating-Point Move Integer from Single}]{FMV.X.W\hfill Floating-Point Move Integer from Single}

\paragraph{Encoding} Format \textsf{RV32F\_RTypeFMF2I} (Section~\ref{sec:format-RV32F-RTypeFMF2I}), \verb|funct3 = 000|\\

\bitfields{ 5/11100, 2/00, 5/00000, 5/frs1, 3/000, 5/rd, 7/1010011 }


\paragraph{Syntax} \texttt{fmv.x.w} \emph{rd}, \emph{frs1}

\paragraph{Description} 

\paragraph{Execution} {Reservation table \textsf{ALU\_TINY} (Section~\ref{sec:reservation-ALU-TINY})}
~
{\begin{lstlisting}
stage RR:
  new frs1 = FPR[frs1];
  new rd = _SX_32(frs1);
stage E1:
  if (rd != 0) {
    BIR[rd] = rd;
  }
\end{lstlisting}}
\newpage
\subsubsection[{\small FNMADD.S: Floating-Point Fused Negate Multiply Add Single}]{FNMADD.S\hfill Floating-Point Fused Negate Multiply Add Single}

\paragraph{Encoding} Format \textsf{RV32F\_R4TypeRM} (Section~\ref{sec:format-RV32F-R4TypeRM}), \verb|opcode = 1001111|\\

\bitfields{ 5/frs3, 2/00, 5/frs2, 5/frs1, 3/frm, 5/frd, 7/1001111 }


\paragraph{Syntax} \texttt{fnmadd.s} \emph{frd}, \emph{frs1}, \emph{frs2}, \emph{frs3}, \emph{frm}

\paragraph{Description} 

\paragraph{Modifier(s)}
\noindent\begin{itemize}
\item frm (cf. froundmode in section \ref{par:modifier-froundmode} on page \pageref{par:modifier-froundmode})
\end{itemize}

\paragraph{Execution} {Reservation table \textsf{ALU\_LITE} (Section~\ref{sec:reservation-ALU-LITE})}
~
{\begin{lstlisting}
stage ID:
  new frm = _ZX_3(frm);
stage RR:
  new frs3 = FPR[frs3];
  new frs2 = FPR[frs2];
  new frs1 = FPR[frs1];
  new frm = decode_riscv_float_rounding_mode(frm);
  if (check_float_rounding_mode_trap(frm) == 0) {
    _THROW(OPCODE);
  } else {
    new frd = f32_mulAdd(frm, frs1 ^ 0x80000000, frs2, frs3 ^ 0x80000000) | 0xFFFFFFFF00000000;
    FPR[frd] = frd;
    commit_float_exception_flags();
  }
\end{lstlisting}}
\newpage
\subsubsection[{\small FNMSUB.S: Floating-Point Fused Negate Multiply Subtract Single}]{FNMSUB.S\hfill Floating-Point Fused Negate Multiply Subtract Single}

\paragraph{Encoding} Format \textsf{RV32F\_R4TypeRM} (Section~\ref{sec:format-RV32F-R4TypeRM}), \verb|opcode = 1001011|\\

\bitfields{ 5/frs3, 2/00, 5/frs2, 5/frs1, 3/frm, 5/frd, 7/1001011 }


\paragraph{Syntax} \texttt{fnmsub.s} \emph{frd}, \emph{frs1}, \emph{frs2}, \emph{frs3}, \emph{frm}

\paragraph{Description} 

\paragraph{Modifier(s)}
\noindent\begin{itemize}
\item frm (cf. froundmode in section \ref{par:modifier-froundmode} on page \pageref{par:modifier-froundmode})
\end{itemize}

\paragraph{Execution} {Reservation table \textsf{ALU\_LITE} (Section~\ref{sec:reservation-ALU-LITE})}
~
{\begin{lstlisting}
stage ID:
  new frm = _ZX_3(frm);
stage RR:
  new frs3 = FPR[frs3];
  new frs2 = FPR[frs2];
  new frs1 = FPR[frs1];
  new frm = decode_riscv_float_rounding_mode(frm);
  if (check_float_rounding_mode_trap(frm) == 0) {
    _THROW(OPCODE);
  } else {
    new frd = f32_mulAdd(frm, frs1 ^ 0x80000000, frs2, frs3) | 0xFFFFFFFF00000000;
    FPR[frd] = frd;
    commit_float_exception_flags();
  }
\end{lstlisting}}
\newpage
\subsubsection[{\small FSGNJ.S: Floating-Point Sign Inject Single}]{FSGNJ.S\hfill Floating-Point Sign Inject Single}

\paragraph{Encoding} Format \textsf{RV32F\_RTypeFSGN} (Section~\ref{sec:format-RV32F-RTypeFSGN}), \verb|funct3 = 000|\\

\bitfields{ 5/00100, 2/00, 5/frs2, 5/frs1, 3/000, 5/frd, 7/1010011 }


\paragraph{Syntax} \texttt{fsgnj.s} \emph{frd}, \emph{frs1}, \emph{frs2}

\paragraph{Description} 

\paragraph{Execution} {Reservation table \textsf{ALU\_TINY} (Section~\ref{sec:reservation-ALU-TINY})}
~
{\begin{lstlisting}
stage RR:
  new frs2 = FPR[frs2];
  new frs1 = FPR[frs1];
  new frd = ((frs1 & 0x7FFFFFFF) | (frs2 & 0x80000000)) | 0xFFFFFFFF00000000;
stage E1:
  FPR[frd] = frd;
\end{lstlisting}}
\newpage
\subsubsection[{\small FSGNJN.S: Floating-Point Sign Inject Negate Single}]{FSGNJN.S\hfill Floating-Point Sign Inject Negate Single}

\paragraph{Encoding} Format \textsf{RV32F\_RTypeFSGN} (Section~\ref{sec:format-RV32F-RTypeFSGN}), \verb|funct3 = 001|\\

\bitfields{ 5/00100, 2/00, 5/frs2, 5/frs1, 3/001, 5/frd, 7/1010011 }


\paragraph{Syntax} \texttt{fsgnjn.s} \emph{frd}, \emph{frs1}, \emph{frs2}

\paragraph{Description} 

\paragraph{Execution} {Reservation table \textsf{ALU\_TINY} (Section~\ref{sec:reservation-ALU-TINY})}
~
{\begin{lstlisting}
stage RR:
  new frs2 = FPR[frs2];
  new frs1 = FPR[frs1];
  new frd = ((frs1 & 0x7FFFFFFF) | (~frs2 & 0x80000000)) | 0xFFFFFFFF00000000;
stage E1:
  FPR[frd] = frd;
\end{lstlisting}}
\newpage
\subsubsection[{\small FSGNJX.S: Floating-Point Sign Inject XOR Single}]{FSGNJX.S\hfill Floating-Point Sign Inject XOR Single}

\paragraph{Encoding} Format \textsf{RV32F\_RTypeFSGN} (Section~\ref{sec:format-RV32F-RTypeFSGN}), \verb|funct3 = 010|\\

\bitfields{ 5/00100, 2/00, 5/frs2, 5/frs1, 3/010, 5/frd, 7/1010011 }


\paragraph{Syntax} \texttt{fsgnjx.s} \emph{frd}, \emph{frs1}, \emph{frs2}

\paragraph{Description} 

\paragraph{Execution} {Reservation table \textsf{ALU\_TINY} (Section~\ref{sec:reservation-ALU-TINY})}
~
{\begin{lstlisting}
stage RR:
  new frs2 = FPR[frs2];
  new frs1 = FPR[frs1];
  new frd = (frs1 ^ (frs2 & 0x80000000)) | 0xFFFFFFFF00000000;
stage E1:
  FPR[frd] = frd;
\end{lstlisting}}
\newpage
\subsubsection[{\small FSQRT.S: Floating-Point Square Root Single}]{FSQRT.S\hfill Floating-Point Square Root Single}

\paragraph{Encoding} Format \textsf{RV32F\_R2TypeRM} (Section~\ref{sec:format-RV32F-R2TypeRM}), \verb|funct5 = 01011|\\

\bitfields{ 5/01011, 2/00, 5/00000, 5/frs1, 3/frm, 5/frd, 7/1010011 }


\paragraph{Syntax} \texttt{fsqrt.s} \emph{frd}, \emph{frs1}, \emph{frm}

\paragraph{Description} 

\paragraph{Modifier(s)}
\noindent\begin{itemize}
\item frm (cf. froundmode in section \ref{par:modifier-froundmode} on page \pageref{par:modifier-froundmode})
\end{itemize}

\paragraph{Execution} {Reservation table \textsf{ALU\_FULL} (Section~\ref{sec:reservation-ALU-FULL})}
~
{\begin{lstlisting}
stage ID:
  new frm = _ZX_3(frm);
stage RR:
  new frs1 = FPR[frs1];
  new frm = decode_riscv_float_rounding_mode(frm);
  if (check_float_rounding_mode_trap(frm) == 0) {
    _THROW(OPCODE);
  } else {
    new frd = f32_sqrt(frm, frs1) | 0xFFFFFFFF00000000;
    FPR[frd] = frd;
    commit_float_exception_flags();
  }
\end{lstlisting}}
\newpage
\subsubsection[{\small FSUB.S: Floating-Point Subtract Single}]{FSUB.S\hfill Floating-Point Subtract Single}

\paragraph{Encoding} Format \textsf{RV32F\_R3TypeRM} (Section~\ref{sec:format-RV32F-R3TypeRM}), \verb|funct5 = 00001|\\

\bitfields{ 5/00001, 2/00, 5/frs2, 5/frs1, 3/frm, 5/frd, 7/1010011 }


\paragraph{Syntax} \texttt{fsub.s} \emph{frd}, \emph{frs1}, \emph{frs2}, \emph{frm}

\paragraph{Description} 

\paragraph{Modifier(s)}
\noindent\begin{itemize}
\item frm (cf. froundmode in section \ref{par:modifier-froundmode} on page \pageref{par:modifier-froundmode})
\end{itemize}

\paragraph{Execution} {Reservation table \textsf{ALU\_LITE} (Section~\ref{sec:reservation-ALU-LITE})}
~
{\begin{lstlisting}
stage ID:
  new frm = _ZX_3(frm);
stage RR:
  new frs2 = FPR[frs2];
  new frs1 = FPR[frs1];
  new frm = decode_riscv_float_rounding_mode(frm);
  if (check_float_rounding_mode_trap(frm) == 0) {
    _THROW(OPCODE);
  } else {
    new frd = f32_sub(frm, frs1, frs2) | 0xFFFFFFFF00000000;
    FPR[frd] = frd;
    commit_float_exception_flags();
  }
\end{lstlisting}}
\newpage
\subsubsection[{\small FSW: Floating-Point Store Word}]{FSW\hfill Floating-Point Store Word}

\paragraph{Encoding} Format \textsf{RV32F\_SType} (Section~\ref{sec:format-RV32F-SType}), \verb|funct3 = 010|\\

\bitfields{ 7/i-- 11-- 5, 5/frs2, 5/rs1, 3/010, 5/i-- 4-- 0s, 7/0100111 }


\paragraph{Syntax} \texttt{fsw} \emph{frs2}, \emph{i\_\discretionary{}{}{}11\_\discretionary{}{}{}5\_\discretionary{}{}{}i\_\discretionary{}{}{}4\_\discretionary{}{}{}0s}[\emph{rs1}]

\paragraph{Description} The \emph{frs2} is stored into the memory word at address \emph{rs1} + \emph{i\_\discretionary{}{}{}11\_\discretionary{}{}{}5\_\discretionary{}{}{}i\_\discretionary{}{}{}4\_\discretionary{}{}{}0s}.


\paragraph{Execution} {Reservation table \textsf{LSU\_MEMW\_AUXR} (Section~\ref{sec:reservation-LSU-MEMW-AUXR})}
~
{\begin{lstlisting}
stage RR:
  new signed12 = _SX_12(i_11_5_i_4_0s);
  new rs1 = BIR[rs1];
  new address = rs1 + signed12;
stage E1:
  new frs2 = FPR[frs2];
stage E3:
  MEM.store(address, 0xF, 0, frs2, frs2);
\end{lstlisting}}
\newpage
\subsection{RV32I Instructions}
\label{sec:RV32I-instructions}

\subsubsection[{\small ADD: Integer Add}]{ADD\hfill Integer Add}

\paragraph{Encoding} Format \textsf{RV32I\_RF00} (Section~\ref{sec:format-RV32I-RF00}), \verb|funct3 = 000|\\

\bitfields{ 7/0000000, 5/rs2, 5/rs1, 3/000, 5/rd, 7/0110011 }


\paragraph{Syntax} \texttt{add} \emph{rd}, \emph{rs1}, \emph{rs2}

\paragraph{Description} The \emph{rs1} is added to the \emph{rs2}. The result is stored into the \emph{rd}.


\paragraph{Execution} {Reservation table \textsf{ALU\_TINY} (Section~\ref{sec:reservation-ALU-TINY})}
~
{\begin{lstlisting}
stage RR:
  new rs2 = BIR[rs2];
  new rs1 = BIR[rs1];
  new rd = rs1 + rs2;
stage E1:
  if (rd != 0) {
    BIR[rd] = rd;
  }
\end{lstlisting}}
\newpage
\subsubsection[{\small ADDI: Integer ADD Immediate}]{ADDI\hfill Integer ADD Immediate}

\paragraph{Encoding} Format \textsf{RV32I\_I\_ALU} (Section~\ref{sec:format-RV32I-I-ALU}), \verb|funct3 = 000|\\

\bitfields{ 12/i-- 11-- 0, 5/rs1, 3/000, 5/rd, 7/0010011 }


\paragraph{Syntax} \texttt{addi} \emph{rd}, \emph{rs1}, \emph{i\_\discretionary{}{}{}11\_\discretionary{}{}{}0}

\paragraph{Description} The \emph{rs1} is added to the \emph{i\_\discretionary{}{}{}11\_\discretionary{}{}{}0}. The result is stored into the \emph{rd}.


\paragraph{Execution} {Reservation table \textsf{ALU\_TINY} (Section~\ref{sec:reservation-ALU-TINY})}
~
{\begin{lstlisting}
stage RR:
  new signed12 = _SX_12(i_11_0);
  new rs1 = BIR[rs1];
  new rd = rs1 + signed12;
stage E1:
  if (rd != 0) {
    BIR[rd] = rd;
  }
\end{lstlisting}}
\newpage
\subsubsection[{\small AND: Bitwise AND}]{AND\hfill Bitwise AND}

\paragraph{Encoding} Format \textsf{RV32I\_RF00} (Section~\ref{sec:format-RV32I-RF00}), \verb|funct3 = 111|\\

\bitfields{ 7/0000000, 5/rs2, 5/rs1, 3/111, 5/rd, 7/0110011 }


\paragraph{Syntax} \texttt{and} \emph{rd}, \emph{rs1}, \emph{rs2}

\paragraph{Description} The \emph{rs1} is ANDed with the \emph{rs2}. The result is stored into the \emph{rd}.


\paragraph{Execution} {Reservation table \textsf{ALU\_TINY} (Section~\ref{sec:reservation-ALU-TINY})}
~
{\begin{lstlisting}
stage RR:
  new rs2 = BIR[rs2];
  new rs1 = BIR[rs1];
  new rd = rs1 & rs2;
stage E1:
  if (rd != 0) {
    BIR[rd] = rd;
  }
\end{lstlisting}}
\newpage
\subsubsection[{\small ANDI: Bitwise AND Immediate}]{ANDI\hfill Bitwise AND Immediate}

\paragraph{Encoding} Format \textsf{RV32I\_I\_ALU} (Section~\ref{sec:format-RV32I-I-ALU}), \verb|funct3 = 111|\\

\bitfields{ 12/i-- 11-- 0, 5/rs1, 3/111, 5/rd, 7/0010011 }


\paragraph{Syntax} \texttt{andi} \emph{rd}, \emph{rs1}, \emph{i\_\discretionary{}{}{}11\_\discretionary{}{}{}0}

\paragraph{Description} The \emph{rs1} is ANDed with the \emph{i\_\discretionary{}{}{}11\_\discretionary{}{}{}0}. The result is stored into the \emph{rd}.


\paragraph{Execution} {Reservation table \textsf{ALU\_TINY} (Section~\ref{sec:reservation-ALU-TINY})}
~
{\begin{lstlisting}
stage RR:
  new signed12 = _SX_12(i_11_0);
  new rs1 = BIR[rs1];
  new rd = rs1 & signed12;
stage E1:
  if (rd != 0) {
    BIR[rd] = rd;
  }
\end{lstlisting}}
\newpage
\subsubsection[{\small AUIPC: Add Upper Immediate to PC}]{AUIPC\hfill Add Upper Immediate to PC}

\paragraph{Encoding} Format \textsf{RV32I\_UType} (Section~\ref{sec:format-RV32I-UType}), \verb|opcode = 0010111|\\

\bitfields{ 20/i-- 31-- 12, 5/rd, 7/0010111 }


\paragraph{Syntax} \texttt{auipc} \emph{rd}, \emph{i\_\discretionary{}{}{}31\_\discretionary{}{}{}12}

\paragraph{Description} The \emph{i\_\discretionary{}{}{}31\_\discretionary{}{}{}12} is shifted left 12 bits and added to the PC. The result is sign-extended from 32 bits and stored into the \emph{rd}.


\paragraph{Execution} {Reservation table \textsf{ALU\_TINY} (Section~\ref{sec:reservation-ALU-TINY})}
~
{\begin{lstlisting}
stage ID:
  new signed20 = _SX_20(i_31_12);
stage RR:
  new rd = (signed20 << 12) + PC;
stage E1:
  if (rd != 0) {
    BIR[rd] = rd;
  }
\end{lstlisting}}
\newpage
\subsubsection[{\small BEQ: Branch if Equal}]{BEQ\hfill Branch if Equal}

\paragraph{Encoding} Format \textsf{RV32I\_BType} (Section~\ref{sec:format-RV32I-BType}), \verb|funct3 = 000|\\

\bitfields{ 1/i-- 12, 6/i-- 10-- 5, 5/rs2, 5/rs1, 3/000, 4/i-- 4-- 1, 1/i-- 11b, 7/1100011 }


\paragraph{Syntax} \texttt{beq} \emph{rs1}, \emph{rs2}, \emph{i\_\discretionary{}{}{}12\_\discretionary{}{}{}i\_\discretionary{}{}{}11b\_\discretionary{}{}{}i\_\discretionary{}{}{}10\_\discretionary{}{}{}5\_\discretionary{}{}{}i\_\discretionary{}{}{}4\_\discretionary{}{}{}1}

\paragraph{Description} Conditional branch to PC+\emph{i\_\discretionary{}{}{}12\_\discretionary{}{}{}i\_\discretionary{}{}{}11b\_\discretionary{}{}{}i\_\discretionary{}{}{}10\_\discretionary{}{}{}5\_\discretionary{}{}{}i\_\discretionary{}{}{}4\_\discretionary{}{}{}1} if \emph{rs1} $==$ \emph{rs2}.


\paragraph{Execution} {Reservation table \textsf{BCU\_BRRP2} (Section~\ref{sec:reservation-BCU-BRRP2})}
~
{\begin{lstlisting}
stage ID:
  new pcrel12s1 = (_SX_12(i_12_i_11b_i_10_5_i_4_1)<<1);
  new rs2 = BIR[rs2];
  new rs1 = BIR[rs1];
stage RR:
  if (comp64_eq(rs1, rs2)) {
    new address = PC + pcrel12s1;
    NPC = address;
    branch_info(1, address);
    srhpc_update();
  } else {
    branch_info(0, 0);
  }
\end{lstlisting}}
\newpage
\subsubsection[{\small BEQZ: Branch if Equal Zero}]{BEQZ\hfill Branch if Equal Zero}

\paragraph{Encoding} Format \textsf{RV32I\_BTypeZ2} (Section~\ref{sec:format-RV32I-BTypeZ2}), \verb|funct3 = 000|\\

\bitfields{ 1/i-- 12, 6/i-- 10-- 5, 5/00000, 5/rs1, 3/000, 4/i-- 4-- 1, 1/i-- 11b, 7/1100011 }


\paragraph{Syntax} \texttt{beqz} \emph{rs1}, \emph{i\_\discretionary{}{}{}12\_\discretionary{}{}{}i\_\discretionary{}{}{}11b\_\discretionary{}{}{}i\_\discretionary{}{}{}10\_\discretionary{}{}{}5\_\discretionary{}{}{}i\_\discretionary{}{}{}4\_\discretionary{}{}{}1}

\paragraph{Description} Conditional branch to PC+\emph{i\_\discretionary{}{}{}12\_\discretionary{}{}{}i\_\discretionary{}{}{}11b\_\discretionary{}{}{}i\_\discretionary{}{}{}10\_\discretionary{}{}{}5\_\discretionary{}{}{}i\_\discretionary{}{}{}4\_\discretionary{}{}{}1} if \emph{rs1} $==$ 0.


\paragraph{Execution} {Reservation table \textsf{BCU\_BRRP} (Section~\ref{sec:reservation-BCU-BRRP})}
~
{\begin{lstlisting}
stage ID:
  new pcrel12s1 = (_SX_12(i_12_i_11b_i_10_5_i_4_1)<<1);
  new rs1 = BIR[rs1];
stage RR:
  if (comp64_eq(rs1, 0)) {
    new address = PC + pcrel12s1;
    NPC = address;
    branch_info(1, address);
    srhpc_update();
  } else {
    branch_info(0, 0);
  }
\end{lstlisting}}
\newpage
\subsubsection[{\small BGE: Branch if Greater Than or Equal}]{BGE\hfill Branch if Greater Than or Equal}

\paragraph{Encoding} Format \textsf{RV32I\_BType} (Section~\ref{sec:format-RV32I-BType}), \verb|funct3 = 101|\\

\bitfields{ 1/i-- 12, 6/i-- 10-- 5, 5/rs2, 5/rs1, 3/101, 4/i-- 4-- 1, 1/i-- 11b, 7/1100011 }


\paragraph{Syntax} \texttt{bge} \emph{rs1}, \emph{rs2}, \emph{i\_\discretionary{}{}{}12\_\discretionary{}{}{}i\_\discretionary{}{}{}11b\_\discretionary{}{}{}i\_\discretionary{}{}{}10\_\discretionary{}{}{}5\_\discretionary{}{}{}i\_\discretionary{}{}{}4\_\discretionary{}{}{}1}

\paragraph{Description} Conditional branch to PC+\emph{i\_\discretionary{}{}{}12\_\discretionary{}{}{}i\_\discretionary{}{}{}11b\_\discretionary{}{}{}i\_\discretionary{}{}{}10\_\discretionary{}{}{}5\_\discretionary{}{}{}i\_\discretionary{}{}{}4\_\discretionary{}{}{}1} if \emph{rs1} $\ge$ \emph{rs2} signed.


\paragraph{Execution} {Reservation table \textsf{BCU\_BRRP2} (Section~\ref{sec:reservation-BCU-BRRP2})}
~
{\begin{lstlisting}
stage ID:
  new pcrel12s1 = (_SX_12(i_12_i_11b_i_10_5_i_4_1)<<1);
  new rs2 = BIR[rs2];
  new rs1 = BIR[rs1];
stage RR:
  if (comp64_ge(rs1, rs2)) {
    new address = PC + pcrel12s1;
    NPC = address;
    branch_info(1, address);
    srhpc_update();
  } else {
    branch_info(0, 0);
  }
\end{lstlisting}}
\newpage
\subsubsection[{\small BGEU: Branch if Greater Than or Equal Unsigned}]{BGEU\hfill Branch if Greater Than or Equal Unsigned}

\paragraph{Encoding} Format \textsf{RV32I\_BType} (Section~\ref{sec:format-RV32I-BType}), \verb|funct3 = 111|\\

\bitfields{ 1/i-- 12, 6/i-- 10-- 5, 5/rs2, 5/rs1, 3/111, 4/i-- 4-- 1, 1/i-- 11b, 7/1100011 }


\paragraph{Syntax} \texttt{bgeu} \emph{rs1}, \emph{rs2}, \emph{i\_\discretionary{}{}{}12\_\discretionary{}{}{}i\_\discretionary{}{}{}11b\_\discretionary{}{}{}i\_\discretionary{}{}{}10\_\discretionary{}{}{}5\_\discretionary{}{}{}i\_\discretionary{}{}{}4\_\discretionary{}{}{}1}

\paragraph{Description} Conditional branch to PC+\emph{i\_\discretionary{}{}{}12\_\discretionary{}{}{}i\_\discretionary{}{}{}11b\_\discretionary{}{}{}i\_\discretionary{}{}{}10\_\discretionary{}{}{}5\_\discretionary{}{}{}i\_\discretionary{}{}{}4\_\discretionary{}{}{}1} if \emph{rs1} $\ge$ \emph{rs2} unsigned.


\paragraph{Execution} {Reservation table \textsf{BCU\_BRRP2} (Section~\ref{sec:reservation-BCU-BRRP2})}
~
{\begin{lstlisting}
stage ID:
  new pcrel12s1 = (_SX_12(i_12_i_11b_i_10_5_i_4_1)<<1);
  new rs2 = BIR[rs2];
  new rs1 = BIR[rs1];
stage RR:
  if (comp64_geu(rs1, rs2)) {
    new address = PC + pcrel12s1;
    NPC = address;
    branch_info(1, address);
    srhpc_update();
  } else {
    branch_info(0, 0);
  }
\end{lstlisting}}
\newpage
\subsubsection[{\small BGEZ: Branch if Greater Than or Equal Zero}]{BGEZ\hfill Branch if Greater Than or Equal Zero}

\paragraph{Encoding} Format \textsf{RV32I\_BTypeZ2} (Section~\ref{sec:format-RV32I-BTypeZ2}), \verb|funct3 = 101|\\

\bitfields{ 1/i-- 12, 6/i-- 10-- 5, 5/00000, 5/rs1, 3/101, 4/i-- 4-- 1, 1/i-- 11b, 7/1100011 }


\paragraph{Syntax} \texttt{bgez} \emph{rs1}, \emph{i\_\discretionary{}{}{}12\_\discretionary{}{}{}i\_\discretionary{}{}{}11b\_\discretionary{}{}{}i\_\discretionary{}{}{}10\_\discretionary{}{}{}5\_\discretionary{}{}{}i\_\discretionary{}{}{}4\_\discretionary{}{}{}1}

\paragraph{Description} Conditional branch to PC+\emph{i\_\discretionary{}{}{}12\_\discretionary{}{}{}i\_\discretionary{}{}{}11b\_\discretionary{}{}{}i\_\discretionary{}{}{}10\_\discretionary{}{}{}5\_\discretionary{}{}{}i\_\discretionary{}{}{}4\_\discretionary{}{}{}1} if \emph{rs1} $\ge$ 0 signed.


\paragraph{Execution} {Reservation table \textsf{BCU\_BRRP} (Section~\ref{sec:reservation-BCU-BRRP})}
~
{\begin{lstlisting}
stage ID:
  new pcrel12s1 = (_SX_12(i_12_i_11b_i_10_5_i_4_1)<<1);
  new rs1 = BIR[rs1];
stage RR:
  if (comp64_ge(rs1, 0)) {
    new address = PC + pcrel12s1;
    NPC = address;
    branch_info(1, address);
    srhpc_update();
  } else {
    branch_info(0, 0);
  }
\end{lstlisting}}
\newpage
\subsubsection[{\small BGTZ: Branch if Greater Than Zero}]{BGTZ\hfill Branch if Greater Than Zero}

\paragraph{Encoding} Format \textsf{RV32I\_BTypeZ1} (Section~\ref{sec:format-RV32I-BTypeZ1}), \verb|funct3 = 100|\\

\bitfields{ 1/i-- 12, 6/i-- 10-- 5, 5/rs2, 5/00000, 3/100, 4/i-- 4-- 1, 1/i-- 11b, 7/1100011 }


\paragraph{Syntax} \texttt{bgtz} \emph{rs2}, \emph{i\_\discretionary{}{}{}12\_\discretionary{}{}{}i\_\discretionary{}{}{}11b\_\discretionary{}{}{}i\_\discretionary{}{}{}10\_\discretionary{}{}{}5\_\discretionary{}{}{}i\_\discretionary{}{}{}4\_\discretionary{}{}{}1}

\paragraph{Description} Conditional branch to PC+\emph{i\_\discretionary{}{}{}12\_\discretionary{}{}{}i\_\discretionary{}{}{}11b\_\discretionary{}{}{}i\_\discretionary{}{}{}10\_\discretionary{}{}{}5\_\discretionary{}{}{}i\_\discretionary{}{}{}4\_\discretionary{}{}{}1} if \emph{rs2} $>$ 0 signed.


\paragraph{Execution} {Reservation table \textsf{BCU\_BRRP} (Section~\ref{sec:reservation-BCU-BRRP})}
~
{\begin{lstlisting}
stage ID:
  new pcrel12s1 = (_SX_12(i_12_i_11b_i_10_5_i_4_1)<<1);
  new rs2 = BIR[rs2];
stage RR:
  if (comp64_lt(0, rs2)) {
    new address = PC + pcrel12s1;
    NPC = address;
    branch_info(1, address);
    srhpc_update();
  } else {
    branch_info(0, 0);
  }
\end{lstlisting}}
\newpage
\subsubsection[{\small BLEZ: Branch if Less Than or Equal Zero}]{BLEZ\hfill Branch if Less Than or Equal Zero}

\paragraph{Encoding} Format \textsf{RV32I\_BTypeZ1} (Section~\ref{sec:format-RV32I-BTypeZ1}), \verb|funct3 = 101|\\

\bitfields{ 1/i-- 12, 6/i-- 10-- 5, 5/rs2, 5/00000, 3/101, 4/i-- 4-- 1, 1/i-- 11b, 7/1100011 }


\paragraph{Syntax} \texttt{blez} \emph{rs2}, \emph{i\_\discretionary{}{}{}12\_\discretionary{}{}{}i\_\discretionary{}{}{}11b\_\discretionary{}{}{}i\_\discretionary{}{}{}10\_\discretionary{}{}{}5\_\discretionary{}{}{}i\_\discretionary{}{}{}4\_\discretionary{}{}{}1}

\paragraph{Description} Conditional branch to PC+\emph{i\_\discretionary{}{}{}12\_\discretionary{}{}{}i\_\discretionary{}{}{}11b\_\discretionary{}{}{}i\_\discretionary{}{}{}10\_\discretionary{}{}{}5\_\discretionary{}{}{}i\_\discretionary{}{}{}4\_\discretionary{}{}{}1} if \emph{rs2} $\le$ 0 signed.


\paragraph{Execution} {Reservation table \textsf{BCU\_BRRP} (Section~\ref{sec:reservation-BCU-BRRP})}
~
{\begin{lstlisting}
stage ID:
  new pcrel12s1 = (_SX_12(i_12_i_11b_i_10_5_i_4_1)<<1);
  new rs2 = BIR[rs2];
stage RR:
  if (comp64_ge(0, rs2)) {
    new address = PC + pcrel12s1;
    NPC = address;
    branch_info(1, address);
    srhpc_update();
  } else {
    branch_info(0, 0);
  }
\end{lstlisting}}
\newpage
\subsubsection[{\small BLT: Branch if Less Than}]{BLT\hfill Branch if Less Than}

\paragraph{Encoding} Format \textsf{RV32I\_BType} (Section~\ref{sec:format-RV32I-BType}), \verb|funct3 = 100|\\

\bitfields{ 1/i-- 12, 6/i-- 10-- 5, 5/rs2, 5/rs1, 3/100, 4/i-- 4-- 1, 1/i-- 11b, 7/1100011 }


\paragraph{Syntax} \texttt{blt} \emph{rs1}, \emph{rs2}, \emph{i\_\discretionary{}{}{}12\_\discretionary{}{}{}i\_\discretionary{}{}{}11b\_\discretionary{}{}{}i\_\discretionary{}{}{}10\_\discretionary{}{}{}5\_\discretionary{}{}{}i\_\discretionary{}{}{}4\_\discretionary{}{}{}1}

\paragraph{Description} Conditional branch to PC+\emph{i\_\discretionary{}{}{}12\_\discretionary{}{}{}i\_\discretionary{}{}{}11b\_\discretionary{}{}{}i\_\discretionary{}{}{}10\_\discretionary{}{}{}5\_\discretionary{}{}{}i\_\discretionary{}{}{}4\_\discretionary{}{}{}1} if \emph{rs1} $<$ \emph{rs2} signed.


\paragraph{Execution} {Reservation table \textsf{BCU\_BRRP2} (Section~\ref{sec:reservation-BCU-BRRP2})}
~
{\begin{lstlisting}
stage ID:
  new pcrel12s1 = (_SX_12(i_12_i_11b_i_10_5_i_4_1)<<1);
  new rs2 = BIR[rs2];
  new rs1 = BIR[rs1];
stage RR:
  if (comp64_lt(rs1, rs2)) {
    new address = PC + pcrel12s1;
    NPC = address;
    branch_info(1, address);
    srhpc_update();
  } else {
    branch_info(0, 0);
  }
\end{lstlisting}}
\newpage
\subsubsection[{\small BLTU: Branch if Less Than Unsigned}]{BLTU\hfill Branch if Less Than Unsigned}

\paragraph{Encoding} Format \textsf{RV32I\_BType} (Section~\ref{sec:format-RV32I-BType}), \verb|funct3 = 110|\\

\bitfields{ 1/i-- 12, 6/i-- 10-- 5, 5/rs2, 5/rs1, 3/110, 4/i-- 4-- 1, 1/i-- 11b, 7/1100011 }


\paragraph{Syntax} \texttt{bltu} \emph{rs1}, \emph{rs2}, \emph{i\_\discretionary{}{}{}12\_\discretionary{}{}{}i\_\discretionary{}{}{}11b\_\discretionary{}{}{}i\_\discretionary{}{}{}10\_\discretionary{}{}{}5\_\discretionary{}{}{}i\_\discretionary{}{}{}4\_\discretionary{}{}{}1}

\paragraph{Description} Conditional branch to PC+\emph{i\_\discretionary{}{}{}12\_\discretionary{}{}{}i\_\discretionary{}{}{}11b\_\discretionary{}{}{}i\_\discretionary{}{}{}10\_\discretionary{}{}{}5\_\discretionary{}{}{}i\_\discretionary{}{}{}4\_\discretionary{}{}{}1} if \emph{rs1} $<$ \emph{rs2} unsigned.


\paragraph{Execution} {Reservation table \textsf{BCU\_BRRP2} (Section~\ref{sec:reservation-BCU-BRRP2})}
~
{\begin{lstlisting}
stage ID:
  new pcrel12s1 = (_SX_12(i_12_i_11b_i_10_5_i_4_1)<<1);
  new rs2 = BIR[rs2];
  new rs1 = BIR[rs1];
stage RR:
  if (comp64_ltu(rs1, rs2)) {
    new address = PC + pcrel12s1;
    NPC = address;
    branch_info(1, address);
    srhpc_update();
  } else {
    branch_info(0, 0);
  }
\end{lstlisting}}
\newpage
\subsubsection[{\small BLTZ: Branch if Less Than Zero}]{BLTZ\hfill Branch if Less Than Zero}

\paragraph{Encoding} Format \textsf{RV32I\_BTypeZ2} (Section~\ref{sec:format-RV32I-BTypeZ2}), \verb|funct3 = 100|\\

\bitfields{ 1/i-- 12, 6/i-- 10-- 5, 5/00000, 5/rs1, 3/100, 4/i-- 4-- 1, 1/i-- 11b, 7/1100011 }


\paragraph{Syntax} \texttt{bltz} \emph{rs1}, \emph{i\_\discretionary{}{}{}12\_\discretionary{}{}{}i\_\discretionary{}{}{}11b\_\discretionary{}{}{}i\_\discretionary{}{}{}10\_\discretionary{}{}{}5\_\discretionary{}{}{}i\_\discretionary{}{}{}4\_\discretionary{}{}{}1}

\paragraph{Description} Conditional branch to PC+\emph{i\_\discretionary{}{}{}12\_\discretionary{}{}{}i\_\discretionary{}{}{}11b\_\discretionary{}{}{}i\_\discretionary{}{}{}10\_\discretionary{}{}{}5\_\discretionary{}{}{}i\_\discretionary{}{}{}4\_\discretionary{}{}{}1} if \emph{rs1} $<$ 0 signed.


\paragraph{Execution} {Reservation table \textsf{BCU\_BRRP} (Section~\ref{sec:reservation-BCU-BRRP})}
~
{\begin{lstlisting}
stage ID:
  new pcrel12s1 = (_SX_12(i_12_i_11b_i_10_5_i_4_1)<<1);
  new rs1 = BIR[rs1];
stage RR:
  if (comp64_lt(rs1, 0)) {
    new address = PC + pcrel12s1;
    NPC = address;
    branch_info(1, address);
    srhpc_update();
  } else {
    branch_info(0, 0);
  }
\end{lstlisting}}
\newpage
\subsubsection[{\small BNE: Branch if Not Equal}]{BNE\hfill Branch if Not Equal}

\paragraph{Encoding} Format \textsf{RV32I\_BType} (Section~\ref{sec:format-RV32I-BType}), \verb|funct3 = 001|\\

\bitfields{ 1/i-- 12, 6/i-- 10-- 5, 5/rs2, 5/rs1, 3/001, 4/i-- 4-- 1, 1/i-- 11b, 7/1100011 }


\paragraph{Syntax} \texttt{bne} \emph{rs1}, \emph{rs2}, \emph{i\_\discretionary{}{}{}12\_\discretionary{}{}{}i\_\discretionary{}{}{}11b\_\discretionary{}{}{}i\_\discretionary{}{}{}10\_\discretionary{}{}{}5\_\discretionary{}{}{}i\_\discretionary{}{}{}4\_\discretionary{}{}{}1}

\paragraph{Description} Conditional branch to PC+\emph{i\_\discretionary{}{}{}12\_\discretionary{}{}{}i\_\discretionary{}{}{}11b\_\discretionary{}{}{}i\_\discretionary{}{}{}10\_\discretionary{}{}{}5\_\discretionary{}{}{}i\_\discretionary{}{}{}4\_\discretionary{}{}{}1} if \emph{rs1} $\ne$ \emph{rs2}.


\paragraph{Execution} {Reservation table \textsf{BCU\_BRRP2} (Section~\ref{sec:reservation-BCU-BRRP2})}
~
{\begin{lstlisting}
stage ID:
  new pcrel12s1 = (_SX_12(i_12_i_11b_i_10_5_i_4_1)<<1);
  new rs2 = BIR[rs2];
  new rs1 = BIR[rs1];
stage RR:
  if (comp64_ne(rs1, rs2)) {
    new address = PC + pcrel12s1;
    NPC = address;
    branch_info(1, address);
    srhpc_update();
  } else {
    branch_info(0, 0);
  }
\end{lstlisting}}
\newpage
\subsubsection[{\small BNEZ: Branch if Not Equal Zero}]{BNEZ\hfill Branch if Not Equal Zero}

\paragraph{Encoding} Format \textsf{RV32I\_BTypeZ2} (Section~\ref{sec:format-RV32I-BTypeZ2}), \verb|funct3 = 001|\\

\bitfields{ 1/i-- 12, 6/i-- 10-- 5, 5/00000, 5/rs1, 3/001, 4/i-- 4-- 1, 1/i-- 11b, 7/1100011 }


\paragraph{Syntax} \texttt{bnez} \emph{rs1}, \emph{i\_\discretionary{}{}{}12\_\discretionary{}{}{}i\_\discretionary{}{}{}11b\_\discretionary{}{}{}i\_\discretionary{}{}{}10\_\discretionary{}{}{}5\_\discretionary{}{}{}i\_\discretionary{}{}{}4\_\discretionary{}{}{}1}

\paragraph{Description} Conditional branch to PC+\emph{i\_\discretionary{}{}{}12\_\discretionary{}{}{}i\_\discretionary{}{}{}11b\_\discretionary{}{}{}i\_\discretionary{}{}{}10\_\discretionary{}{}{}5\_\discretionary{}{}{}i\_\discretionary{}{}{}4\_\discretionary{}{}{}1} if \emph{rs1} $\ne$ 0.


\paragraph{Execution} {Reservation table \textsf{BCU\_BRRP} (Section~\ref{sec:reservation-BCU-BRRP})}
~
{\begin{lstlisting}
stage ID:
  new pcrel12s1 = (_SX_12(i_12_i_11b_i_10_5_i_4_1)<<1);
  new rs1 = BIR[rs1];
stage RR:
  if (comp64_ne(rs1, 0)) {
    new address = PC + pcrel12s1;
    NPC = address;
    branch_info(1, address);
    srhpc_update();
  } else {
    branch_info(0, 0);
  }
\end{lstlisting}}
\newpage
\subsubsection[{\small CBO.CLEAN: Perform a clean operation on the cache block containing the effective address}]{CBO.CLEAN\hfill Perform a clean operation on the cache block containing the effective address}

\paragraph{Encoding} Format \textsf{RV32I\_ZICBOM} (Section~\ref{sec:format-RV32I-ZICBOM}), \verb|i_11_0 = 000000000001|\\

\bitfields{ 12/000000000001, 5/rs1, 3/010, 5/00000, 7/0001111 }


\paragraph{Syntax} \texttt{cbo.clean} [\emph{rs1}]

\paragraph{Description} A CBO.CLEAN instruction performs a clean operation on the cache block whose effective address is the base address specified in \emph{rs1}. The instruction operates on the set of coherent caches accessed by the agent executing the instruction.


\paragraph{Execution} {Reservation table \textsf{LSU} (Section~\ref{sec:reservation-LSU})}
~
{\begin{lstlisting}
stage RR:
  new address = BIR[rs1];
stage E3:
  MEM.dflushl(address);
\end{lstlisting}}
\newpage
\subsubsection[{\small CBO.FLUSH: Perform a flush operation on the cache block containing the effective address}]{CBO.FLUSH\hfill Perform a flush operation on the cache block containing the effective address}

\paragraph{Encoding} Format \textsf{RV32I\_ZICBOM} (Section~\ref{sec:format-RV32I-ZICBOM}), \verb|i_11_0 = 000000000010|\\

\bitfields{ 12/000000000010, 5/rs1, 3/010, 5/00000, 7/0001111 }


\paragraph{Syntax} \texttt{cbo.flush} [\emph{rs1}]

\paragraph{Description} A CBO.FLUSH instruction preforms a flush operation on the cache block whose effective address is the base address specified in \emph{rs1}. The instruction operates on the set of coherent caches accessed by the agent executing the instruction.


\paragraph{Execution} {Reservation table \textsf{LSU} (Section~\ref{sec:reservation-LSU})}
~
{\begin{lstlisting}
stage RR:
  new address = BIR[rs1];
stage E3:
  MEM.dflushl(address);
\end{lstlisting}}
\newpage
\subsubsection[{\small CBO.INVAL: Perform an invalidate operation on the cache block containing the effective address}]{CBO.INVAL\hfill Perform an invalidate operation on the cache block containing the effective address}

\paragraph{Encoding} Format \textsf{RV32I\_ZICBOM} (Section~\ref{sec:format-RV32I-ZICBOM}), \verb|i_11_0 = 000000000000|\\

\bitfields{ 12/000000000000, 5/rs1, 3/010, 5/00000, 7/0001111 }


\paragraph{Syntax} \texttt{cbo.inval} [\emph{rs1}]

\paragraph{Description} A CBO.INVAL instruction performs an invalidate operation on the cache block whose effective address is the base address specified in \emph{rs1}. The instruction operates on the set of coherent caches accessed by the agent executing the instruction. Depending on CSR programming, the instruction may perfom a flush operation instead of an invalidate operation.


\paragraph{Execution} {Reservation table \textsf{LSU} (Section~\ref{sec:reservation-LSU})}
~
{\begin{lstlisting}
stage RR:
  new address = BIR[rs1];
stage E3:
  MEM.dpurgel(address);
\end{lstlisting}}
\newpage
\subsubsection[{\small CBO.ZERO: Store zeros to the set of the bytes corresponding to the cache block containing the effective address}]{CBO.ZERO\hfill Store zeros to the set of the bytes corresponding to the cache block containing the effective address}

\paragraph{Encoding} Format \textsf{RV32I\_ZICBOM} (Section~\ref{sec:format-RV32I-ZICBOM}), \verb|i_11_0 = 000000000100|\\

\bitfields{ 12/000000000100, 5/rs1, 3/010, 5/00000, 7/0001111 }


\paragraph{Syntax} \texttt{cbo.zero} [\emph{rs1}]

\paragraph{Description} A CBO.ZERO instruction performs stores to the set of bytes corresponding to a cache block.


\paragraph{Execution} {Reservation table \textsf{LSU} (Section~\ref{sec:reservation-LSU})}
~
{\begin{lstlisting}
stage RR:
  new address = BIR[rs1];
stage E3:
  MEM.dzerol(address, 1);
\end{lstlisting}}
\newpage
\subsubsection[{\small EBREAK: Debugger Breakpoint}]{EBREAK\hfill Debugger Breakpoint}

\paragraph{Encoding} Format \textsf{RV32I\_IEnv} (Section~\ref{sec:format-RV32I-IEnv}), \verb|i_11_0 = 000000000001|\\

\bitfields{ 12/000000000001, 5/00000, 3/000, 5/00000, 7/1110011 }


\paragraph{Syntax} \texttt{ebreak}

\paragraph{Description} Breakpoint for the debugger.


\paragraph{Execution} {Reservation table \textsf{ALL} (Section~\ref{sec:reservation-ALL})}
~
{\begin{lstlisting}
stage RR:
  break(PS.PL);
\end{lstlisting}}
\newpage
\subsubsection[{\small ECALL: Environment Call}]{ECALL\hfill Environment Call}

\paragraph{Encoding} Format \textsf{RV32I\_IEnv} (Section~\ref{sec:format-RV32I-IEnv}), \verb|i_11_0 = 000000000000|\\

\bitfields{ 12/000000000000, 5/00000, 3/000, 5/00000, 7/1110011 }


\paragraph{Syntax} \texttt{ecall}

\paragraph{Description} Environment call.


\paragraph{Execution} {Reservation table \textsf{ALL} (Section~\ref{sec:reservation-ALL})}
~
{\begin{lstlisting}
stage ID:
  syscall(GPR[RVC[RVC_ECRID:5]]);
  branch_info(1, 0);
\end{lstlisting}}
\newpage
\subsubsection[{\small FENCE.MEM: Memory Access Fence}]{FENCE.MEM\hfill Memory Access Fence}

\paragraph{Encoding} Format \textsf{RV32I\_IFence} (Section~\ref{sec:format-RV32I-IFence}), \verb|fm = 0000|\\

\bitfields{ 4/0000, 4/pred, 4/succ, 5/00000, 3/000, 5/00000, 7/0001111 }


\paragraph{Syntax} \texttt{fence.mem}\emph{pred}\emph{succ}

\paragraph{Description} Ensures that all issued memory operations are committed to memory.


\paragraph{Modifier(s)}
\noindent\begin{itemize}
\item pred (cf. ordering in section \ref{par:modifier-ordering} on page \pageref{par:modifier-ordering})
\item succ (cf. ordering in section \ref{par:modifier-ordering} on page \pageref{par:modifier-ordering})
\end{itemize}

\paragraph{Execution} {Reservation table \textsf{LSU2\_MEMW} (Section~\ref{sec:reservation-LSU2-MEMW})}
~
{\begin{lstlisting}
stage RR:
  new pred = _ZX_4(pred);
  new succ = _ZX_4(succ);
stage E3:
  MEM.rv_fence(pred, succ);
\end{lstlisting}}
\newpage
\subsubsection[{\small J: Jump}]{J\hfill Jump}

\paragraph{Encoding} Format \textsf{RV32I\_JTypeZD} (Section~\ref{sec:format-RV32I-JTypeZD}), \verb|opcode = 1101111|\\

\bitfields{ 1/i-- 20, 10/i-- 10-- 1, 1/i-- 11j, 8/i-- 19-- 12, 5/00000, 7/1101111 }


\paragraph{Syntax} \texttt{j} \emph{i\_\discretionary{}{}{}20\_\discretionary{}{}{}i\_\discretionary{}{}{}19\_\discretionary{}{}{}12\_\discretionary{}{}{}i\_\discretionary{}{}{}11j\_\discretionary{}{}{}i\_\discretionary{}{}{}10\_\discretionary{}{}{}1}

\paragraph{Description} The next PC is computed by adding the \emph{} to the PC.


\paragraph{Execution} {Reservation table \textsf{BCU\_XFER} (Section~\ref{sec:reservation-BCU-XFER})}
~
{\begin{lstlisting}
stage ID:
  new pcrel20s1 = (_SX_20(i_20_i_19_12_i_11j_i_10_1)<<1);
  NPC = PC + pcrel20s1;
\end{lstlisting}}
\newpage
\subsubsection[{\small JAL: Jump And Link}]{JAL\hfill Jump And Link}

\paragraph{Encoding} Format \textsf{RV32I\_JType} (Section~\ref{sec:format-RV32I-JType}), \verb|opcode = 1101111|\\

\bitfields{ 1/i-- 20, 10/i-- 10-- 1, 1/i-- 11j, 8/i-- 19-- 12, 5/rd, 7/1101111 }


\paragraph{Syntax} \texttt{jal} \emph{rd}, \emph{i\_\discretionary{}{}{}20\_\discretionary{}{}{}i\_\discretionary{}{}{}19\_\discretionary{}{}{}12\_\discretionary{}{}{}i\_\discretionary{}{}{}11j\_\discretionary{}{}{}i\_\discretionary{}{}{}10\_\discretionary{}{}{}1}

\paragraph{Description} The PC of the next instruction is written into the \emph{rd}. The next PC is computed by adding the \emph{i\_\discretionary{}{}{}20\_\discretionary{}{}{}i\_\discretionary{}{}{}19\_\discretionary{}{}{}12\_\discretionary{}{}{}i\_\discretionary{}{}{}11j\_\discretionary{}{}{}i\_\discretionary{}{}{}10\_\discretionary{}{}{}1} to the PC.


\paragraph{Execution} {Reservation table \textsf{BCU2\_TINY\_LSU} (Section~\ref{sec:reservation-BCU2-TINY-LSU})}
~
{\begin{lstlisting}
stage ID:
  new pcrel20s1 = (_SX_20(i_20_i_19_12_i_11j_i_10_1)<<1);
  NPC = PC + pcrel20s1;
stage RR:
  new rd = PC + 4;
stage E1:
  if (rd != 0) {
    BIR[rd] = rd;
  }
\end{lstlisting}}
\newpage
\subsubsection[{\small JALR: Jump And Link Register}]{JALR\hfill Jump And Link Register}

\paragraph{Encoding} Format \textsf{RV32I\_JALR} (Section~\ref{sec:format-RV32I-JALR}), \verb|opcode = 1100111|\\

\bitfields{ 12/i-- 11-- 0, 5/rs1, 3/000, 5/rd, 7/1100111 }


\paragraph{Syntax} \texttt{jalr} \emph{rd}, \emph{rs1}, \emph{i\_\discretionary{}{}{}11\_\discretionary{}{}{}0}

\paragraph{Description} The PC of the next instruction is written into the \emph{rd}. The \emph{i\_\discretionary{}{}{}11\_\discretionary{}{}{}0} is sign-extended from 12 bits and added to the \emph{rs1}. The next PC is set to the result with the least-significant bit cleared.


\paragraph{Execution} {Reservation table \textsf{BCU2\_TINY\_LSU} (Section~\ref{sec:reservation-BCU2-TINY-LSU})}
~
{\begin{lstlisting}
stage ID:
  NPC = (rs1 + signed12) & 0xFFFFFFFFFFFFFFFE;
stage RR:
  new signed12 = _SX_12(i_11_0);
  new rs1 = BIR[rs1];
  new rd = PC + 4;
stage E1:
  if (rd != 0) {
    BIR[rd] = rd;
  }
\end{lstlisting}}
\newpage
\subsubsection[{\small JR: Jump Register}]{JR\hfill Jump Register}

\paragraph{Encoding} Format \textsf{RV32I\_JALRZDI} (Section~\ref{sec:format-RV32I-JALRZDI}), \verb|opcode = 1100111|\\

\bitfields{ 12/000000000000, 5/rs1, 3/000, 5/00000, 7/1100111 }


\paragraph{Syntax} \texttt{jr} \emph{rs1}

\paragraph{Description} The next PC is set to the \emph{rs1} with the least-significant bit cleared.


\paragraph{Execution} {Reservation table \textsf{BCU\_XFER\_BRRP} (Section~\ref{sec:reservation-BCU-XFER-BRRP})}
~
{\begin{lstlisting}
stage ID:
  NPC = rs1 & 0xFFFFFFFFFFFFFFFE;
stage RR:
  new rs1 = BIR[rs1];
\end{lstlisting}}
\newpage
\subsubsection[{\small LB: Load Byte}]{LB\hfill Load Byte}

\paragraph{Encoding} Format \textsf{RV32I\_ILoad} (Section~\ref{sec:format-RV32I-ILoad}), \verb|funct3 = 000|\\

\bitfields{ 12/i-- 11-- 0, 5/rs1, 3/000, 5/rd, 7/0000011 }


\paragraph{Syntax} \texttt{lb} \emph{rd}, \emph{i\_\discretionary{}{}{}11\_\discretionary{}{}{}0}[\emph{rs1}]

\paragraph{Description} The \emph{rd} receives the sign-extended memory byte at address \emph{rs1} + \emph{i\_\discretionary{}{}{}11\_\discretionary{}{}{}0}.


\paragraph{Execution} {Reservation table \textsf{LSU\_AUXW} (Section~\ref{sec:reservation-LSU-AUXW})}
~
{\begin{lstlisting}
stage RR:
  new signed12 = _SX_12(i_11_0);
  new rs1 = BIR[rs1];
  new address = rs1 + signed12;
  new rd = _SX_8(MEM.load(address, 0x1, 0, rd));
stage E1:
  if (rd != 0) {
    BIR[rd] = rd;
  }
\end{lstlisting}}
\newpage
\subsubsection[{\small LBU: Load Byte Unsigned}]{LBU\hfill Load Byte Unsigned}

\paragraph{Encoding} Format \textsf{RV32I\_ILoad} (Section~\ref{sec:format-RV32I-ILoad}), \verb|funct3 = 100|\\

\bitfields{ 12/i-- 11-- 0, 5/rs1, 3/100, 5/rd, 7/0000011 }


\paragraph{Syntax} \texttt{lbu} \emph{rd}, \emph{i\_\discretionary{}{}{}11\_\discretionary{}{}{}0}[\emph{rs1}]

\paragraph{Description} The \emph{rd} receives the zero-extended memory byte at address \emph{rs1} + \emph{i\_\discretionary{}{}{}11\_\discretionary{}{}{}0}.


\paragraph{Execution} {Reservation table \textsf{LSU\_AUXW} (Section~\ref{sec:reservation-LSU-AUXW})}
~
{\begin{lstlisting}
stage RR:
  new signed12 = _SX_12(i_11_0);
  new rs1 = BIR[rs1];
  new address = rs1 + signed12;
  new rd = _ZX_8(MEM.load(address, 0x1, 0, rd));
stage E1:
  if (rd != 0) {
    BIR[rd] = rd;
  }
\end{lstlisting}}
\newpage
\subsubsection[{\small LH: Load Half Word}]{LH\hfill Load Half Word}

\paragraph{Encoding} Format \textsf{RV32I\_ILoad} (Section~\ref{sec:format-RV32I-ILoad}), \verb|funct3 = 001|\\

\bitfields{ 12/i-- 11-- 0, 5/rs1, 3/001, 5/rd, 7/0000011 }


\paragraph{Syntax} \texttt{lh} \emph{rd}, \emph{i\_\discretionary{}{}{}11\_\discretionary{}{}{}0}[\emph{rs1}]

\paragraph{Description} The \emph{rd} receives the sign-extended memory half word at address \emph{rs1} + \emph{i\_\discretionary{}{}{}11\_\discretionary{}{}{}0}.


\paragraph{Execution} {Reservation table \textsf{LSU\_AUXW} (Section~\ref{sec:reservation-LSU-AUXW})}
~
{\begin{lstlisting}
stage RR:
  new signed12 = _SX_12(i_11_0);
  new rs1 = BIR[rs1];
  new address = rs1 + signed12;
  new rd = _SX_16(MEM.load(address, 0x3, 0, rd));
stage E1:
  if (rd != 0) {
    BIR[rd] = rd;
  }
\end{lstlisting}}
\newpage
\subsubsection[{\small LHU: Load Half Word Unsigned}]{LHU\hfill Load Half Word Unsigned}

\paragraph{Encoding} Format \textsf{RV32I\_ILoad} (Section~\ref{sec:format-RV32I-ILoad}), \verb|funct3 = 101|\\

\bitfields{ 12/i-- 11-- 0, 5/rs1, 3/101, 5/rd, 7/0000011 }


\paragraph{Syntax} \texttt{lhu} \emph{rd}, \emph{i\_\discretionary{}{}{}11\_\discretionary{}{}{}0}[\emph{rs1}]

\paragraph{Description} The \emph{rd} receives the zero-extended memory half word at address \emph{rs1} + \emph{i\_\discretionary{}{}{}11\_\discretionary{}{}{}0}.


\paragraph{Execution} {Reservation table \textsf{LSU\_AUXW} (Section~\ref{sec:reservation-LSU-AUXW})}
~
{\begin{lstlisting}
stage RR:
  new signed12 = _SX_12(i_11_0);
  new rs1 = BIR[rs1];
  new address = rs1 + signed12;
  new rd = _ZX_16(MEM.load(address, 0x3, 0, rd));
stage E1:
  if (rd != 0) {
    BIR[rd] = rd;
  }
\end{lstlisting}}
\newpage
\subsubsection[{\small LUI: Load Upper Immediate}]{LUI\hfill Load Upper Immediate}

\paragraph{Encoding} Format \textsf{RV32I\_UType} (Section~\ref{sec:format-RV32I-UType}), \verb|opcode = 0110111|\\

\bitfields{ 20/i-- 31-- 12, 5/rd, 7/0110111 }


\paragraph{Syntax} \texttt{lui} \emph{rd}, \emph{i\_\discretionary{}{}{}31\_\discretionary{}{}{}12}

\paragraph{Description} The \emph{i\_\discretionary{}{}{}31\_\discretionary{}{}{}12} is shifted left 12 bits. The result is sign-extended from 32 bits and stored into the \emph{rd}.


\paragraph{Execution} {Reservation table \textsf{ALU\_TINY} (Section~\ref{sec:reservation-ALU-TINY})}
~
{\begin{lstlisting}
stage ID:
  new signed20 = _SX_20(i_31_12);
stage RR:
  new rd = signed20 << 12;
stage E1:
  if (rd != 0) {
    BIR[rd] = rd;
  }
\end{lstlisting}}
\newpage
\subsubsection[{\small LW: Load Word}]{LW\hfill Load Word}

\paragraph{Encoding} Format \textsf{RV32I\_ILoad} (Section~\ref{sec:format-RV32I-ILoad}), \verb|funct3 = 010|\\

\bitfields{ 12/i-- 11-- 0, 5/rs1, 3/010, 5/rd, 7/0000011 }


\paragraph{Syntax} \texttt{lw} \emph{rd}, \emph{i\_\discretionary{}{}{}11\_\discretionary{}{}{}0}[\emph{rs1}]

\paragraph{Description} The \emph{rd} receives the sign-extended memory word at address \emph{rs1} + \emph{i\_\discretionary{}{}{}11\_\discretionary{}{}{}0}.


\paragraph{Execution} {Reservation table \textsf{LSU\_AUXW} (Section~\ref{sec:reservation-LSU-AUXW})}
~
{\begin{lstlisting}
stage RR:
  new signed12 = _SX_12(i_11_0);
  new rs1 = BIR[rs1];
  new address = rs1 + signed12;
  new rd = _SX_32(MEM.load(address, 0xF, 0, rd));
stage E1:
  if (rd != 0) {
    BIR[rd] = rd;
  }
\end{lstlisting}}
\newpage
\subsubsection[{\small OR: Inclusive OR}]{OR\hfill Inclusive OR}

\paragraph{Encoding} Format \textsf{RV32I\_RF00} (Section~\ref{sec:format-RV32I-RF00}), \verb|funct3 = 110|\\

\bitfields{ 7/0000000, 5/rs2, 5/rs1, 3/110, 5/rd, 7/0110011 }


\paragraph{Syntax} \texttt{or} \emph{rd}, \emph{rs1}, \emph{rs2}

\paragraph{Description} The \emph{rs1} is ORed with the \emph{rs2}. The result is stored into the \emph{rd}.


\paragraph{Execution} {Reservation table \textsf{ALU\_TINY} (Section~\ref{sec:reservation-ALU-TINY})}
~
{\begin{lstlisting}
stage RR:
  new rs2 = BIR[rs2];
  new rs1 = BIR[rs1];
  new rd = rs1 | rs2;
stage E1:
  if (rd != 0) {
    BIR[rd] = rd;
  }
\end{lstlisting}}
\newpage
\subsubsection[{\small ORI: Inclusive OR Immediate}]{ORI\hfill Inclusive OR Immediate}

\paragraph{Encoding} Format \textsf{RV32I\_I\_ALU} (Section~\ref{sec:format-RV32I-I-ALU}), \verb|funct3 = 110|\\

\bitfields{ 12/i-- 11-- 0, 5/rs1, 3/110, 5/rd, 7/0010011 }


\paragraph{Syntax} \texttt{ori} \emph{rd}, \emph{rs1}, \emph{i\_\discretionary{}{}{}11\_\discretionary{}{}{}0}

\paragraph{Description} The \emph{rs1} is ORed with the \emph{i\_\discretionary{}{}{}11\_\discretionary{}{}{}0}. The result is stored into the \emph{rd}.


\paragraph{Execution} {Reservation table \textsf{ALU\_TINY} (Section~\ref{sec:reservation-ALU-TINY})}
~
{\begin{lstlisting}
stage RR:
  new signed12 = _SX_12(i_11_0);
  new rs1 = BIR[rs1];
  new rd = rs1 | signed12;
stage E1:
  if (rd != 0) {
    BIR[rd] = rd;
  }
\end{lstlisting}}
\newpage
\subsubsection[{\small PREFETCH.I: Provide a hint to hardware that the cache block containing the effective address is likely to be accessed by an instruction fetch in the near future
}]{PREFETCH.I\hfill Provide a hint to hardware that the cache block containing the effective address is likely to be accessed by an instruction fetch in the near future
}

\paragraph{Encoding} Format \textsf{RV32I\_STypeZICBOP} (Section~\ref{sec:format-RV32I-STypeZICBOP}), \verb|rs2 = 00000|\\

\bitfields{ 7/i-- 11-- 5, 5/00000, 5/rs1, 3/110, 5/00000, 7/0010011 }


\paragraph{Syntax} \texttt{prefetch.i} \emph{i\_\discretionary{}{}{}11\_\discretionary{}{}{}5}[\emph{rs1}]

\paragraph{Description} A prefetch.i instruction indicates to hardware that the cache block containing the effective address is likely to be accessed by an instruction fetch in the near future. An implementation may opt to cache a copy of the cache block in a cache accessed by an instruction fetch in order to improve memory access latency, but this behavior is not required.


\paragraph{Execution} {Reservation table \textsf{LSU} (Section~\ref{sec:reservation-LSU})}
~
{\begin{lstlisting}
stage RR:
  new offset = i_11_5 << 5;
  new base = BIR[rs1];
  new address = base + offset;
\end{lstlisting}}
\newpage
\subsubsection[{\small PREFETCH.R: Provide a hint to hardware that the cache block containing the effective address is likely to be accessed by a data read in the near future
}]{PREFETCH.R\hfill Provide a hint to hardware that the cache block containing the effective address is likely to be accessed by a data read in the near future
}

\paragraph{Encoding} Format \textsf{RV32I\_STypeZICBOP} (Section~\ref{sec:format-RV32I-STypeZICBOP}), \verb|rs2 = 00001|\\

\bitfields{ 7/i-- 11-- 5, 5/00001, 5/rs1, 3/110, 5/00000, 7/0010011 }


\paragraph{Syntax} \texttt{prefetch.r} \emph{i\_\discretionary{}{}{}11\_\discretionary{}{}{}5}[\emph{rs1}]

\paragraph{Description} A prefetch.r instruction indicates to hardware that the cache block containing the effective address is likely to be accessed by a data read (i.e. load) in the near future. An implementation may opt to cache a copy of the cache block in a cache accessed by a data read in order to improve memory access latency, but this behavior is not required.


\paragraph{Execution} {Reservation table \textsf{LSU} (Section~\ref{sec:reservation-LSU})}
~
{\begin{lstlisting}
stage RR:
  new offset = i_11_5 << 5;
  new base = BIR[rs1];
  new address = base + offset;
stage E3:
  MEM.dtouchl(address);
\end{lstlisting}}
\newpage
\subsubsection[{\small PREFETCH.W: Provide a hint to hardware that the cache block containing the effective address is likely to be accessed by a data write in the near future
}]{PREFETCH.W\hfill Provide a hint to hardware that the cache block containing the effective address is likely to be accessed by a data write in the near future
}

\paragraph{Encoding} Format \textsf{RV32I\_STypeZICBOP} (Section~\ref{sec:format-RV32I-STypeZICBOP}), \verb|rs2 = 00011|\\

\bitfields{ 7/i-- 11-- 5, 5/00011, 5/rs1, 3/110, 5/00000, 7/0010011 }


\paragraph{Syntax} \texttt{prefetch.w} \emph{i\_\discretionary{}{}{}11\_\discretionary{}{}{}5}[\emph{rs1}]

\paragraph{Description} A prefetch.w instruction indicates to hardware that the cache block containing the effective address is likely to be accessed by a data write (i.e. store) in the near future. An implementation may opt to cache a copy of the cache block in a cache accessed by a data write in order to improve memory access latency, but this behavior is not required


\paragraph{Execution} {Reservation table \textsf{LSU} (Section~\ref{sec:reservation-LSU})}
~
{\begin{lstlisting}
stage RR:
  new offset = i_11_5 << 5;
  new base = BIR[rs1];
  new address = base + offset;
stage E3:
  MEM.prefetch_w(address);
\end{lstlisting}}
\newpage
\subsubsection[{\small SB: Store Byte}]{SB\hfill Store Byte}

\paragraph{Encoding} Format \textsf{RV32I\_SType} (Section~\ref{sec:format-RV32I-SType}), \verb|funct3 = 000|\\

\bitfields{ 7/i-- 11-- 5, 5/rs2, 5/rs1, 3/000, 5/i-- 4-- 0s, 7/0100011 }


\paragraph{Syntax} \texttt{sb} \emph{rs2}, \emph{i\_\discretionary{}{}{}11\_\discretionary{}{}{}5\_\discretionary{}{}{}i\_\discretionary{}{}{}4\_\discretionary{}{}{}0s}[\emph{rs1}]

\paragraph{Description} The \emph{rs2} is stored into the memory byte at address \emph{rs1} + \emph{i\_\discretionary{}{}{}11\_\discretionary{}{}{}5\_\discretionary{}{}{}i\_\discretionary{}{}{}4\_\discretionary{}{}{}0s}.


\paragraph{Execution} {Reservation table \textsf{LSU\_MEMW\_AUXR} (Section~\ref{sec:reservation-LSU-MEMW-AUXR})}
~
{\begin{lstlisting}
stage RR:
  new signed12 = _SX_12(i_11_5_i_4_0s);
  new rs1 = BIR[rs1];
  new address = rs1 + signed12;
stage E1:
  new rs2 = BIR[rs2];
stage E3:
  MEM.store(address, 0x1, 0, rs2, rs2);
\end{lstlisting}}
\newpage
\subsubsection[{\small SH: Store Half Word}]{SH\hfill Store Half Word}

\paragraph{Encoding} Format \textsf{RV32I\_SType} (Section~\ref{sec:format-RV32I-SType}), \verb|funct3 = 001|\\

\bitfields{ 7/i-- 11-- 5, 5/rs2, 5/rs1, 3/001, 5/i-- 4-- 0s, 7/0100011 }


\paragraph{Syntax} \texttt{sh} \emph{rs2}, \emph{i\_\discretionary{}{}{}11\_\discretionary{}{}{}5\_\discretionary{}{}{}i\_\discretionary{}{}{}4\_\discretionary{}{}{}0s}[\emph{rs1}]

\paragraph{Description} The \emph{rs2} is stored into the memory half word at address \emph{rs1} + \emph{i\_\discretionary{}{}{}11\_\discretionary{}{}{}5\_\discretionary{}{}{}i\_\discretionary{}{}{}4\_\discretionary{}{}{}0s}.


\paragraph{Execution} {Reservation table \textsf{LSU\_MEMW\_AUXR} (Section~\ref{sec:reservation-LSU-MEMW-AUXR})}
~
{\begin{lstlisting}
stage RR:
  new signed12 = _SX_12(i_11_5_i_4_0s);
  new rs1 = BIR[rs1];
  new address = rs1 + signed12;
stage E1:
  new rs2 = BIR[rs2];
stage E3:
  MEM.store(address, 0x3, 0, rs2, rs2);
\end{lstlisting}}
\newpage
\subsubsection[{\small SLL: Shift Left Logical}]{SLL\hfill Shift Left Logical}

\paragraph{Encoding} Format \textsf{RV32I\_RF00} (Section~\ref{sec:format-RV32I-RF00}), \verb|funct3 = 001|\\

\bitfields{ 7/0000000, 5/rs2, 5/rs1, 3/001, 5/rd, 7/0110011 }


\paragraph{Syntax} \texttt{sll} \emph{rd}, \emph{rs1}, \emph{rs2}

\paragraph{Description} The \emph{rs1} is shifted left by the \emph{rs2} modulo 64. The result is stored into the \emph{rd}.


\paragraph{Execution} {Reservation table \textsf{ALU\_TINY} (Section~\ref{sec:reservation-ALU-TINY})}
~
{\begin{lstlisting}
stage RR:
  new rs2 = BIR[rs2];
  new rs1 = BIR[rs1];
  new rd = rs1 << (rs2 & 0x3F);
stage E1:
  if (rd != 0) {
    BIR[rd] = rd;
  }
\end{lstlisting}}
\newpage
\subsubsection[{\small SLT: Set Less Than}]{SLT\hfill Set Less Than}

\paragraph{Encoding} Format \textsf{RV32I\_RF00} (Section~\ref{sec:format-RV32I-RF00}), \verb|funct3 = 010|\\

\bitfields{ 7/0000000, 5/rs2, 5/rs1, 3/010, 5/rd, 7/0110011 }


\paragraph{Syntax} \texttt{slt} \emph{rd}, \emph{rs1}, \emph{rs2}

\paragraph{Description} The \emph{rs1} is compared for less than signed the \emph{rs2}. The result is stored into the \emph{rd}.


\paragraph{Execution} {Reservation table \textsf{ALU\_TINY} (Section~\ref{sec:reservation-ALU-TINY})}
~
{\begin{lstlisting}
stage RR:
  new rs2 = BIR[rs2];
  new rs1 = BIR[rs1];
  new rd = _SX_64(rs1) < _SX_64(rs2);
stage E1:
  if (rd != 0) {
    BIR[rd] = rd;
  }
\end{lstlisting}}
\newpage
\subsubsection[{\small SLTI: Set Less Than Immediate}]{SLTI\hfill Set Less Than Immediate}

\paragraph{Encoding} Format \textsf{RV32I\_I\_ALU} (Section~\ref{sec:format-RV32I-I-ALU}), \verb|funct3 = 010|\\

\bitfields{ 12/i-- 11-- 0, 5/rs1, 3/010, 5/rd, 7/0010011 }


\paragraph{Syntax} \texttt{slti} \emph{rd}, \emph{rs1}, \emph{i\_\discretionary{}{}{}11\_\discretionary{}{}{}0}

\paragraph{Description} The \emph{rs1} is compared for less than signed the \emph{i\_\discretionary{}{}{}11\_\discretionary{}{}{}0}. The result is stored into the \emph{rd}.


\paragraph{Execution} {Reservation table \textsf{ALU\_TINY} (Section~\ref{sec:reservation-ALU-TINY})}
~
{\begin{lstlisting}
stage RR:
  new signed12 = _SX_12(i_11_0);
  new rs1 = BIR[rs1];
  new rd = _SX_64(rs1) < _SX_64(signed12);
stage E1:
  if (rd != 0) {
    BIR[rd] = rd;
  }
\end{lstlisting}}
\newpage
\subsubsection[{\small SLTIU: Set Less Than Immediate Unsigned}]{SLTIU\hfill Set Less Than Immediate Unsigned}

\paragraph{Encoding} Format \textsf{RV32I\_I\_ALU} (Section~\ref{sec:format-RV32I-I-ALU}), \verb|funct3 = 011|\\

\bitfields{ 12/i-- 11-- 0, 5/rs1, 3/011, 5/rd, 7/0010011 }


\paragraph{Syntax} \texttt{sltiu} \emph{rd}, \emph{rs1}, \emph{i\_\discretionary{}{}{}11\_\discretionary{}{}{}0}

\paragraph{Description} The \emph{rs1} is compared for less than signed the \emph{i\_\discretionary{}{}{}11\_\discretionary{}{}{}0}. The result is stored into the \emph{rd}.


\paragraph{Execution} {Reservation table \textsf{ALU\_TINY} (Section~\ref{sec:reservation-ALU-TINY})}
~
{\begin{lstlisting}
stage RR:
  new signed12 = _SX_12(i_11_0);
  new rs1 = BIR[rs1];
  new rd = _ZX_64(rs1) < _ZX_64(signed12);
stage E1:
  if (rd != 0) {
    BIR[rd] = rd;
  }
\end{lstlisting}}
\newpage
\subsubsection[{\small SLTU: Set Less Than Unsigned}]{SLTU\hfill Set Less Than Unsigned}

\paragraph{Encoding} Format \textsf{RV32I\_RF00} (Section~\ref{sec:format-RV32I-RF00}), \verb|funct3 = 011|\\

\bitfields{ 7/0000000, 5/rs2, 5/rs1, 3/011, 5/rd, 7/0110011 }


\paragraph{Syntax} \texttt{sltu} \emph{rd}, \emph{rs1}, \emph{rs2}

\paragraph{Description} The \emph{rs1} is compared for less than signed the \emph{rs2}. The result is stored into the \emph{rd}.


\paragraph{Execution} {Reservation table \textsf{ALU\_TINY} (Section~\ref{sec:reservation-ALU-TINY})}
~
{\begin{lstlisting}
stage RR:
  new rs2 = BIR[rs2];
  new rs1 = BIR[rs1];
  new rd = _ZX_64(rs1) < _ZX_64(rs2);
stage E1:
  if (rd != 0) {
    BIR[rd] = rd;
  }
\end{lstlisting}}
\newpage
\subsubsection[{\small SRA: Shift Right Arithmetic}]{SRA\hfill Shift Right Arithmetic}

\paragraph{Encoding} Format \textsf{RV32I\_RF20} (Section~\ref{sec:format-RV32I-RF20}), \verb|funct3 = 101|\\

\bitfields{ 7/0100000, 5/rs2, 5/rs1, 3/101, 5/rd, 7/0110011 }


\paragraph{Syntax} \texttt{sra} \emph{rd}, \emph{rs1}, \emph{rs2}

\paragraph{Description} The \emph{rs1} assuming signed integer is shifted right by the \emph{rs2} modulo 64. The result is stored into the \emph{rd}.


\paragraph{Execution} {Reservation table \textsf{ALU\_TINY} (Section~\ref{sec:reservation-ALU-TINY})}
~
{\begin{lstlisting}
stage RR:
  new rs2 = BIR[rs2];
  new rs1 = BIR[rs1];
  new rd = _SX_64(rs1) >> (rs2 & 0x3F);
stage E1:
  if (rd != 0) {
    BIR[rd] = rd;
  }
\end{lstlisting}}
\newpage
\subsubsection[{\small SRL: Shift Right Logical}]{SRL\hfill Shift Right Logical}

\paragraph{Encoding} Format \textsf{RV32I\_RF00} (Section~\ref{sec:format-RV32I-RF00}), \verb|funct3 = 101|\\

\bitfields{ 7/0000000, 5/rs2, 5/rs1, 3/101, 5/rd, 7/0110011 }


\paragraph{Syntax} \texttt{srl} \emph{rd}, \emph{rs1}, \emph{rs2}

\paragraph{Description} The \emph{rs1} assuming unsigned integer is shifted right by the \emph{rs2} modulo 64. The result is stored into the \emph{rd}.


\paragraph{Execution} {Reservation table \textsf{ALU\_TINY} (Section~\ref{sec:reservation-ALU-TINY})}
~
{\begin{lstlisting}
stage RR:
  new rs2 = BIR[rs2];
  new rs1 = BIR[rs1];
  new rd = _ZX_64(rs1) >> (rs2 & 0x3F);
stage E1:
  if (rd != 0) {
    BIR[rd] = rd;
  }
\end{lstlisting}}
\newpage
\subsubsection[{\small SUB: Integer Subtract}]{SUB\hfill Integer Subtract}

\paragraph{Encoding} Format \textsf{RV32I\_RF20} (Section~\ref{sec:format-RV32I-RF20}), \verb|funct3 = 000|\\

\bitfields{ 7/0100000, 5/rs2, 5/rs1, 3/000, 5/rd, 7/0110011 }


\paragraph{Syntax} \texttt{sub} \emph{rd}, \emph{rs1}, \emph{rs2}

\paragraph{Description} The \emph{rs1} is added to the \emph{rs2}. The result is stored into the \emph{rd}.


\paragraph{Execution} {Reservation table \textsf{ALU\_TINY} (Section~\ref{sec:reservation-ALU-TINY})}
~
{\begin{lstlisting}
stage RR:
  new rs2 = BIR[rs2];
  new rs1 = BIR[rs1];
  new rd = rs1 - rs2;
stage E1:
  if (rd != 0) {
    BIR[rd] = rd;
  }
\end{lstlisting}}
\newpage
\subsubsection[{\small SW: Store Word}]{SW\hfill Store Word}

\paragraph{Encoding} Format \textsf{RV32I\_SType} (Section~\ref{sec:format-RV32I-SType}), \verb|funct3 = 010|\\

\bitfields{ 7/i-- 11-- 5, 5/rs2, 5/rs1, 3/010, 5/i-- 4-- 0s, 7/0100011 }


\paragraph{Syntax} \texttt{sw} \emph{rs2}, \emph{i\_\discretionary{}{}{}11\_\discretionary{}{}{}5\_\discretionary{}{}{}i\_\discretionary{}{}{}4\_\discretionary{}{}{}0s}[\emph{rs1}]

\paragraph{Description} The \emph{rs2} is stored into the memory word at address \emph{rs1} + \emph{i\_\discretionary{}{}{}11\_\discretionary{}{}{}5\_\discretionary{}{}{}i\_\discretionary{}{}{}4\_\discretionary{}{}{}0s}.


\paragraph{Execution} {Reservation table \textsf{LSU\_MEMW\_AUXR} (Section~\ref{sec:reservation-LSU-MEMW-AUXR})}
~
{\begin{lstlisting}
stage RR:
  new signed12 = _SX_12(i_11_5_i_4_0s);
  new rs1 = BIR[rs1];
  new address = rs1 + signed12;
stage E1:
  new rs2 = BIR[rs2];
stage E3:
  MEM.store(address, 0xF, 0, rs2, rs2);
\end{lstlisting}}
\newpage
\subsubsection[{\small XOR: Exclusive OR}]{XOR\hfill Exclusive OR}

\paragraph{Encoding} Format \textsf{RV32I\_RF00} (Section~\ref{sec:format-RV32I-RF00}), \verb|funct3 = 100|\\

\bitfields{ 7/0000000, 5/rs2, 5/rs1, 3/100, 5/rd, 7/0110011 }


\paragraph{Syntax} \texttt{xor} \emph{rd}, \emph{rs1}, \emph{rs2}

\paragraph{Description} The \emph{rs1} is XORed with the \emph{rs2}. The result is stored into the \emph{rd}.


\paragraph{Execution} {Reservation table \textsf{ALU\_TINY} (Section~\ref{sec:reservation-ALU-TINY})}
~
{\begin{lstlisting}
stage RR:
  new rs2 = BIR[rs2];
  new rs1 = BIR[rs1];
  new rd = rs1 ^ rs2;
stage E1:
  if (rd != 0) {
    BIR[rd] = rd;
  }
\end{lstlisting}}
\newpage
\subsubsection[{\small XORI: Exclusive OR Immediate}]{XORI\hfill Exclusive OR Immediate}

\paragraph{Encoding} Format \textsf{RV32I\_I\_ALU} (Section~\ref{sec:format-RV32I-I-ALU}), \verb|funct3 = 100|\\

\bitfields{ 12/i-- 11-- 0, 5/rs1, 3/100, 5/rd, 7/0010011 }


\paragraph{Syntax} \texttt{xori} \emph{rd}, \emph{rs1}, \emph{i\_\discretionary{}{}{}11\_\discretionary{}{}{}0}

\paragraph{Description} The \emph{rs1} is XORed with the \emph{i\_\discretionary{}{}{}11\_\discretionary{}{}{}0}. The result is stored into the \emph{rd}.


\paragraph{Execution} {Reservation table \textsf{ALU\_TINY} (Section~\ref{sec:reservation-ALU-TINY})}
~
{\begin{lstlisting}
stage RR:
  new signed12 = _SX_12(i_11_0);
  new rs1 = BIR[rs1];
  new rd = rs1 ^ signed12;
stage E1:
  if (rd != 0) {
    BIR[rd] = rd;
  }
\end{lstlisting}}
\newpage
\subsection{RV32M Instructions}
\label{sec:RV32M-instructions}

\subsubsection[{\small DIV: Integer Division}]{DIV\hfill Integer Division}

\paragraph{Encoding} Format \textsf{RV32M\_RType} (Section~\ref{sec:format-RV32M-RType}), \verb|funct3 = 100|\\

\bitfields{ 7/0000001, 5/rs2, 5/rs1, 3/100, 5/rd, 7/0110011 }


\paragraph{Syntax} \texttt{div} \emph{rd}, \emph{rs1}, \emph{rs2}

\paragraph{Description} The \emph{rs1} assuming signed integer is divided by the \emph{rs2} assuming signed integer. The quotient is stored into the \emph{rd}.


\paragraph{Execution} {Reservation table \textsf{ALU\_FULL} (Section~\ref{sec:reservation-ALU-FULL})}
~
{\begin{lstlisting}
stage RR:
  new rs2 = BIR[rs2];
  new rs1 = BIR[rs1];
  if (_ZX_64(rs2) == 0) {
    new rd = 0xFFFFFFFFFFFFFFFF;
  } else if ((_ZX_64(rs1) == 0x8000000000000000) && (_ZX_64(rs2) == 0xFFFFFFFFFFFFFFFF)) {
    rd = rs1;
  } else {
    rd = _SX_64(rs1) / _SX_64(rs2);
  }
stage E1:
  if (rd != 0) {
    BIR[rd] = rd;
  }
\end{lstlisting}}
\newpage
\subsubsection[{\small DIVU: Integer Division Unsigned}]{DIVU\hfill Integer Division Unsigned}

\paragraph{Encoding} Format \textsf{RV32M\_RType} (Section~\ref{sec:format-RV32M-RType}), \verb|funct3 = 101|\\

\bitfields{ 7/0000001, 5/rs2, 5/rs1, 3/101, 5/rd, 7/0110011 }


\paragraph{Syntax} \texttt{divu} \emph{rd}, \emph{rs1}, \emph{rs2}

\paragraph{Description} The \emph{rs1} assuming unsigned integer is divided by the \emph{rs2} assuming unsigned integer. The quotient is stored into the \emph{rd}.


\paragraph{Execution} {Reservation table \textsf{ALU\_FULL} (Section~\ref{sec:reservation-ALU-FULL})}
~
{\begin{lstlisting}
stage RR:
  new rs2 = BIR[rs2];
  new rs1 = BIR[rs1];
  if (_ZX_64(rs2) == 0) {
    new rd = 0xFFFFFFFFFFFFFFFF;
  } else {
    rd = _ZX_64(rs1) / _ZX_64(rs2);
  }
stage E1:
  if (rd != 0) {
    BIR[rd] = rd;
  }
\end{lstlisting}}
\newpage
\subsubsection[{\small MUL: Integer Multiply}]{MUL\hfill Integer Multiply}

\paragraph{Encoding} Format \textsf{RV32M\_RType} (Section~\ref{sec:format-RV32M-RType}), \verb|funct3 = 000|\\

\bitfields{ 7/0000001, 5/rs2, 5/rs1, 3/000, 5/rd, 7/0110011 }


\paragraph{Syntax} \texttt{mul} \emph{rd}, \emph{rs1}, \emph{rs2}

\paragraph{Description} The \emph{rs1} is multiplied with the \emph{rs2}. The result is stored into the \emph{rd}.


\paragraph{Execution} {Reservation table \textsf{ALU\_LITE} (Section~\ref{sec:reservation-ALU-LITE})}
~
{\begin{lstlisting}
stage RR:
  new rs2 = BIR[rs2];
  new rs1 = BIR[rs1];
  new rd = rs1 * rs2;
stage E1:
  if (rd != 0) {
    BIR[rd] = rd;
  }
\end{lstlisting}}
\newpage
\subsubsection[{\small MULH: Integer Multiply High}]{MULH\hfill Integer Multiply High}

\paragraph{Encoding} Format \textsf{RV32M\_RType} (Section~\ref{sec:format-RV32M-RType}), \verb|funct3 = 001|\\

\bitfields{ 7/0000001, 5/rs2, 5/rs1, 3/001, 5/rd, 7/0110011 }


\paragraph{Syntax} \texttt{mulh} \emph{rd}, \emph{rs1}, \emph{rs2}

\paragraph{Description} The \emph{rs1} assuming signed integer is multiplied with the \emph{rs2} assuming signed integer. The result is shifted right by 64 bits and stored into the \emph{rd}.


\paragraph{Execution} {Reservation table \textsf{ALU\_LITE} (Section~\ref{sec:reservation-ALU-LITE})}
~
{\begin{lstlisting}
stage RR:
  new rs2 = BIR[rs2];
  new rs1 = BIR[rs1];
  new rd = (_SX_64(rs1) * _SX_64(rs2)) >> 64;
stage E1:
  if (rd != 0) {
    BIR[rd] = rd;
  }
\end{lstlisting}}
\newpage
\subsubsection[{\small MULHSU: Integer Multiply High Signed by Unsigned}]{MULHSU\hfill Integer Multiply High Signed by Unsigned}

\paragraph{Encoding} Format \textsf{RV32M\_RType} (Section~\ref{sec:format-RV32M-RType}), \verb|funct3 = 010|\\

\bitfields{ 7/0000001, 5/rs2, 5/rs1, 3/010, 5/rd, 7/0110011 }


\paragraph{Syntax} \texttt{mulhsu} \emph{rd}, \emph{rs1}, \emph{rs2}

\paragraph{Description} The \emph{rs1} assuming signed integer is multiplied with the \emph{rs2} assuming unsigned integer. The result is shifted right by 64 bits and stored into the \emph{rd}.


\paragraph{Execution} {Reservation table \textsf{ALU\_LITE} (Section~\ref{sec:reservation-ALU-LITE})}
~
{\begin{lstlisting}
stage RR:
  new rs2 = BIR[rs2];
  new rs1 = BIR[rs1];
  new rd = (_SX_64(rs1) * _ZX_64(rs2)) >> 64;
stage E1:
  if (rd != 0) {
    BIR[rd] = rd;
  }
\end{lstlisting}}
\newpage
\subsubsection[{\small MULHU: Integer Multiply High Signed by Unsigned}]{MULHU\hfill Integer Multiply High Signed by Unsigned}

\paragraph{Encoding} Format \textsf{RV32M\_RType} (Section~\ref{sec:format-RV32M-RType}), \verb|funct3 = 011|\\

\bitfields{ 7/0000001, 5/rs2, 5/rs1, 3/011, 5/rd, 7/0110011 }


\paragraph{Syntax} \texttt{mulhu} \emph{rd}, \emph{rs1}, \emph{rs2}

\paragraph{Description} The \emph{rs1} assuming unsigned integer is multiplied with the \emph{rs2} assuming unsigned integer. The result is shifted right by 64 bits and stored into the \emph{rd}.


\paragraph{Execution} {Reservation table \textsf{ALU\_LITE} (Section~\ref{sec:reservation-ALU-LITE})}
~
{\begin{lstlisting}
stage RR:
  new rs2 = BIR[rs2];
  new rs1 = BIR[rs1];
  new rd = (_ZX_64(rs1) * _ZX_64(rs2)) >> 64;
stage E1:
  if (rd != 0) {
    BIR[rd] = rd;
  }
\end{lstlisting}}
\newpage
\subsubsection[{\small REM: Integer Remainder}]{REM\hfill Integer Remainder}

\paragraph{Encoding} Format \textsf{RV32M\_RType} (Section~\ref{sec:format-RV32M-RType}), \verb|funct3 = 110|\\

\bitfields{ 7/0000001, 5/rs2, 5/rs1, 3/110, 5/rd, 7/0110011 }


\paragraph{Syntax} \texttt{rem} \emph{rd}, \emph{rs1}, \emph{rs2}

\paragraph{Description} The \emph{rs1} assuming signed integer is divided by the \emph{rs2} assuming signed integer. The remainder is stored into the \emph{rd}.


\paragraph{Execution} {Reservation table \textsf{ALU\_FULL} (Section~\ref{sec:reservation-ALU-FULL})}
~
{\begin{lstlisting}
stage RR:
  new rs2 = BIR[rs2];
  new rs1 = BIR[rs1];
  if (_ZX_64(rs2) == 0) {
    new rd = rs1;
  } else if ((_ZX_64(rs1) == 0x8000000000000000) && (_ZX_64(rs2) == 0xFFFFFFFFFFFFFFFF)) {
    rd = 0;
  } else {
    rd = _SX_64(rs1) % _SX_64(rs2);
  }
stage E1:
  if (rd != 0) {
    BIR[rd] = rd;
  }
\end{lstlisting}}
\newpage
\subsubsection[{\small REMU: Integer Remainder Unsigned}]{REMU\hfill Integer Remainder Unsigned}

\paragraph{Encoding} Format \textsf{RV32M\_RType} (Section~\ref{sec:format-RV32M-RType}), \verb|funct3 = 111|\\

\bitfields{ 7/0000001, 5/rs2, 5/rs1, 3/111, 5/rd, 7/0110011 }


\paragraph{Syntax} \texttt{remu} \emph{rd}, \emph{rs1}, \emph{rs2}

\paragraph{Description} The \emph{rs1} assuming unsigned integer is divided by the \emph{rs2} assuming unsigned integer. The remainder is stored into the \emph{rd}.


\paragraph{Execution} {Reservation table \textsf{ALU\_FULL} (Section~\ref{sec:reservation-ALU-FULL})}
~
{\begin{lstlisting}
stage RR:
  new rs2 = BIR[rs2];
  new rs1 = BIR[rs1];
  if (_ZX_64(rs2) == 0) {
    new rd = rs1;
  } else {
    rd = _ZX_64(rs1) % _ZX_64(rs2);
  }
stage E1:
  if (rd != 0) {
    BIR[rd] = rd;
  }
\end{lstlisting}}
\newpage
\subsection{RV64A Instructions}
\label{sec:RV64A-instructions}

\subsubsection[{\small AMOADD.B: Atomic Memory Operation ADD byte}]{AMOADD.B\hfill Atomic Memory Operation ADD byte}

\paragraph{Encoding} Format \textsf{RV64A\_RMW} (Section~\ref{sec:format-RV64A-RMW}), \verb|funct5 = 00000|, \verb|funct3 = 000|\\

\bitfields{ 5/00000, 2/acqrel, 5/rs2, 5/rs1, 3/000, 5/rd, 7/0101111 }


\paragraph{Syntax} \texttt{amoadd.b}\emph{acqrel} \emph{rd}, \emph{rs2}, [\emph{rs1}]

\paragraph{Description} This instruction implements atomic load-ADD-store on byte. The effective address is in \emph{rs1}.


\paragraph{Modifier(s)}
\noindent\begin{itemize}
\item acqrel (cf. acqrel in section \ref{par:modifier-acqrel} on page \pageref{par:modifier-acqrel})
\end{itemize}

\paragraph{Execution} {Reservation table \textsf{LSU2\_MEMW\_AUXR\_AUXW} (Section~\ref{sec:reservation-LSU2-MEMW-AUXR-AUXW})}
~
{\begin{lstlisting}
stage ID:
  new acqrel = _ZX_1(acqrel);
stage RR:
  new acq = 1 & (acqrel >> 1);
  new rel = 1 & acqrel;
  new rs2 = BIR[rs2];
  new rs1 = BIR[rs1];
  new rd = _SX_8(MEM.atomic_add(rs1, 0x1, acq | (rel << 16), rs2.32[0], rs1));
stage E1:
  if (rd != 0) {
    BIR[rd] = rd;
  }
\end{lstlisting}}
\newpage
\subsubsection[{\small AMOADD.D: Atomic Memory Operation ADD Double Word}]{AMOADD.D\hfill Atomic Memory Operation ADD Double Word}

\paragraph{Encoding} Format \textsf{RV64A\_RMW} (Section~\ref{sec:format-RV64A-RMW}), \verb|funct5 = 00000|, \verb|funct3 = 011|\\

\bitfields{ 5/00000, 2/acqrel, 5/rs2, 5/rs1, 3/011, 5/rd, 7/0101111 }


\paragraph{Syntax} \texttt{amoadd.d}\emph{acqrel} \emph{rd}, \emph{rs2}, [\emph{rs1}]

\paragraph{Description} This instruction implements atomic load-ADD-store on double word. The effective address in \emph{rs1} must be double word aligned (multiple of 8).


\paragraph{Modifier(s)}
\noindent\begin{itemize}
\item acqrel (cf. acqrel in section \ref{par:modifier-acqrel} on page \pageref{par:modifier-acqrel})
\end{itemize}

\paragraph{Execution} {Reservation table \textsf{LSU2\_MEMW\_AUXR\_AUXW} (Section~\ref{sec:reservation-LSU2-MEMW-AUXR-AUXW})}
~
{\begin{lstlisting}
stage ID:
  new acqrel = _ZX_1(acqrel);
stage RR:
  new acq = 1 & (acqrel >> 1);
  new rel = 1 & acqrel;
  new rs2 = BIR[rs2];
  new rs1 = BIR[rs1];
  new rd = MEM.atomic_add(rs1, 0xFF, acq | (rel << 16), rs2.64[0], rs1);
stage E1:
  if (rd != 0) {
    BIR[rd] = rd;
  }
\end{lstlisting}}
\newpage
\subsubsection[{\small AMOADD.H: Atomic Memory Operation ADD half Word}]{AMOADD.H\hfill Atomic Memory Operation ADD half Word}

\paragraph{Encoding} Format \textsf{RV64A\_RMW} (Section~\ref{sec:format-RV64A-RMW}), \verb|funct5 = 00000|, \verb|funct3 = 001|\\

\bitfields{ 5/00000, 2/acqrel, 5/rs2, 5/rs1, 3/001, 5/rd, 7/0101111 }


\paragraph{Syntax} \texttt{amoadd.h}\emph{acqrel} \emph{rd}, \emph{rs2}, [\emph{rs1}]

\paragraph{Description} This instruction implements atomic load-ADD-store on half word. The effective address in \emph{rs1} must be word aligned (multiple of 2).


\paragraph{Modifier(s)}
\noindent\begin{itemize}
\item acqrel (cf. acqrel in section \ref{par:modifier-acqrel} on page \pageref{par:modifier-acqrel})
\end{itemize}

\paragraph{Execution} {Reservation table \textsf{LSU2\_MEMW\_AUXR\_AUXW} (Section~\ref{sec:reservation-LSU2-MEMW-AUXR-AUXW})}
~
{\begin{lstlisting}
stage ID:
  new acqrel = _ZX_1(acqrel);
stage RR:
  new acq = 1 & (acqrel >> 1);
  new rel = 1 & acqrel;
  new rs2 = BIR[rs2];
  new rs1 = BIR[rs1];
  new rd = _SX_16(MEM.atomic_add(rs1, 0x3, acq | (rel << 16), rs2.32[0], rs1));
stage E1:
  if (rd != 0) {
    BIR[rd] = rd;
  }
\end{lstlisting}}
\newpage
\subsubsection[{\small AMOADD.W: Atomic Memory Operation ADD Word}]{AMOADD.W\hfill Atomic Memory Operation ADD Word}

\paragraph{Encoding} Format \textsf{RV64A\_RMW} (Section~\ref{sec:format-RV64A-RMW}), \verb|funct5 = 00000|, \verb|funct3 = 010|\\

\bitfields{ 5/00000, 2/acqrel, 5/rs2, 5/rs1, 3/010, 5/rd, 7/0101111 }


\paragraph{Syntax} \texttt{amoadd.w}\emph{acqrel} \emph{rd}, \emph{rs2}, [\emph{rs1}]

\paragraph{Description} This instruction implements atomic load-ADD-store on word. The effective address in \emph{rs1} must be word aligned (multiple of 4).


\paragraph{Modifier(s)}
\noindent\begin{itemize}
\item acqrel (cf. acqrel in section \ref{par:modifier-acqrel} on page \pageref{par:modifier-acqrel})
\end{itemize}

\paragraph{Execution} {Reservation table \textsf{LSU2\_MEMW\_AUXR\_AUXW} (Section~\ref{sec:reservation-LSU2-MEMW-AUXR-AUXW})}
~
{\begin{lstlisting}
stage ID:
  new acqrel = _ZX_1(acqrel);
stage RR:
  new acq = 1 & (acqrel >> 1);
  new rel = 1 & acqrel;
  new rs2 = BIR[rs2];
  new rs1 = BIR[rs1];
  new rd = _SX_32(MEM.atomic_add(rs1, 0xF, acq | (rel << 16), rs2.32[0], rs1));
stage E1:
  if (rd != 0) {
    BIR[rd] = rd;
  }
\end{lstlisting}}
\newpage
\subsubsection[{\small AMOAND.B: Atomic Memory Operation AND byte}]{AMOAND.B\hfill Atomic Memory Operation AND byte}

\paragraph{Encoding} Format \textsf{RV64A\_RMW} (Section~\ref{sec:format-RV64A-RMW}), \verb|funct5 = 01100|, \verb|funct3 = 000|\\

\bitfields{ 5/01100, 2/acqrel, 5/rs2, 5/rs1, 3/000, 5/rd, 7/0101111 }


\paragraph{Syntax} \texttt{amoand.b}\emph{acqrel} \emph{rd}, \emph{rs2}, [\emph{rs1}]

\paragraph{Description} This instruction implements atomic load-AND-store on byte. The effective address is in \emph{rs1}.


\paragraph{Modifier(s)}
\noindent\begin{itemize}
\item acqrel (cf. acqrel in section \ref{par:modifier-acqrel} on page \pageref{par:modifier-acqrel})
\end{itemize}

\paragraph{Execution} {Reservation table \textsf{LSU2\_MEMW\_AUXR\_AUXW} (Section~\ref{sec:reservation-LSU2-MEMW-AUXR-AUXW})}
~
{\begin{lstlisting}
stage ID:
  new acqrel = _ZX_1(acqrel);
stage RR:
  new acq = 1 & (acqrel >> 1);
  new rel = 1 & acqrel;
  new rs2 = BIR[rs2];
  new rs1 = BIR[rs1];
  new rd = _SX_8(MEM.atomic_and(rs1, 0x1, acq | (rel << 16), rs2.32[0], rs1));
stage E1:
  if (rd != 0) {
    BIR[rd] = rd;
  }
\end{lstlisting}}
\newpage
\subsubsection[{\small AMOAND.D: Atomic Memory Operation AND Double Word}]{AMOAND.D\hfill Atomic Memory Operation AND Double Word}

\paragraph{Encoding} Format \textsf{RV64A\_RMW} (Section~\ref{sec:format-RV64A-RMW}), \verb|funct5 = 01100|, \verb|funct3 = 011|\\

\bitfields{ 5/01100, 2/acqrel, 5/rs2, 5/rs1, 3/011, 5/rd, 7/0101111 }


\paragraph{Syntax} \texttt{amoand.d}\emph{acqrel} \emph{rd}, \emph{rs2}, [\emph{rs1}]

\paragraph{Description} This instruction implements atomic load-AND-store on double word. The effective address in \emph{rs1} must be double word aligned (multiple of 8).


\paragraph{Modifier(s)}
\noindent\begin{itemize}
\item acqrel (cf. acqrel in section \ref{par:modifier-acqrel} on page \pageref{par:modifier-acqrel})
\end{itemize}

\paragraph{Execution} {Reservation table \textsf{LSU2\_MEMW\_AUXR\_AUXW} (Section~\ref{sec:reservation-LSU2-MEMW-AUXR-AUXW})}
~
{\begin{lstlisting}
stage ID:
  new acqrel = _ZX_1(acqrel);
stage RR:
  new acq = 1 & (acqrel >> 1);
  new rel = 1 & acqrel;
  new rs2 = BIR[rs2];
  new rs1 = BIR[rs1];
  new rd = MEM.atomic_and(rs1, 0xFF, acq | (rel << 16), rs2.64[0], rs1);
stage E1:
  if (rd != 0) {
    BIR[rd] = rd;
  }
\end{lstlisting}}
\newpage
\subsubsection[{\small AMOAND.H: Atomic Memory Operation AND half Word}]{AMOAND.H\hfill Atomic Memory Operation AND half Word}

\paragraph{Encoding} Format \textsf{RV64A\_RMW} (Section~\ref{sec:format-RV64A-RMW}), \verb|funct5 = 01100|, \verb|funct3 = 001|\\

\bitfields{ 5/01100, 2/acqrel, 5/rs2, 5/rs1, 3/001, 5/rd, 7/0101111 }


\paragraph{Syntax} \texttt{amoand.h}\emph{acqrel} \emph{rd}, \emph{rs2}, [\emph{rs1}]

\paragraph{Description} This instruction implements atomic load-AND-store on half word. The effective address in \emph{rs1} must be word aligned (multiple of 2).


\paragraph{Modifier(s)}
\noindent\begin{itemize}
\item acqrel (cf. acqrel in section \ref{par:modifier-acqrel} on page \pageref{par:modifier-acqrel})
\end{itemize}

\paragraph{Execution} {Reservation table \textsf{LSU2\_MEMW\_AUXR\_AUXW} (Section~\ref{sec:reservation-LSU2-MEMW-AUXR-AUXW})}
~
{\begin{lstlisting}
stage ID:
  new acqrel = _ZX_1(acqrel);
stage RR:
  new acq = 1 & (acqrel >> 1);
  new rel = 1 & acqrel;
  new rs2 = BIR[rs2];
  new rs1 = BIR[rs1];
  new rd = _SX_16(MEM.atomic_and(rs1, 0x3, acq | (rel << 16), rs2.32[0], rs1));
stage E1:
  if (rd != 0) {
    BIR[rd] = rd;
  }
\end{lstlisting}}
\newpage
\subsubsection[{\small AMOAND.W: Atomic Memory Operation AND Word}]{AMOAND.W\hfill Atomic Memory Operation AND Word}

\paragraph{Encoding} Format \textsf{RV64A\_RMW} (Section~\ref{sec:format-RV64A-RMW}), \verb|funct5 = 01100|, \verb|funct3 = 010|\\

\bitfields{ 5/01100, 2/acqrel, 5/rs2, 5/rs1, 3/010, 5/rd, 7/0101111 }


\paragraph{Syntax} \texttt{amoand.w}\emph{acqrel} \emph{rd}, \emph{rs2}, [\emph{rs1}]

\paragraph{Description} This instruction implements atomic load-AND-store on word. The effective address in \emph{rs1} must be word aligned (multiple of 4).


\paragraph{Modifier(s)}
\noindent\begin{itemize}
\item acqrel (cf. acqrel in section \ref{par:modifier-acqrel} on page \pageref{par:modifier-acqrel})
\end{itemize}

\paragraph{Execution} {Reservation table \textsf{LSU2\_MEMW\_AUXR\_AUXW} (Section~\ref{sec:reservation-LSU2-MEMW-AUXR-AUXW})}
~
{\begin{lstlisting}
stage ID:
  new acqrel = _ZX_1(acqrel);
stage RR:
  new acq = 1 & (acqrel >> 1);
  new rel = 1 & acqrel;
  new rs2 = BIR[rs2];
  new rs1 = BIR[rs1];
  new rd = _SX_32(MEM.atomic_and(rs1, 0xF, acq | (rel << 16), rs2.32[0], rs1));
stage E1:
  if (rd != 0) {
    BIR[rd] = rd;
  }
\end{lstlisting}}
\newpage
\subsubsection[{\small AMOCAS.B: Atomic Memory Operation Compare-And-Swap byte}]{AMOCAS.B\hfill Atomic Memory Operation Compare-And-Swap byte}

\paragraph{Encoding} Format \textsf{RV64A\_CAS} (Section~\ref{sec:format-RV64A-CAS}), \verb|funct3 = 000|\\

\bitfields{ 5/00101, 2/acqrel, 5/rs2, 5/rs1, 3/000, 5/rd, 7/0101111 }


\paragraph{Syntax} \texttt{amocas.b}\emph{acqrel} \emph{rd}, \emph{rs2}, [\emph{rs1}]

\paragraph{Description} This instruction implements atomic compare-and-swap on byte. The effective address is in \emph{rs1}.


\paragraph{Modifier(s)}
\noindent\begin{itemize}
\item acqrel (cf. acqrel in section \ref{par:modifier-acqrel} on page \pageref{par:modifier-acqrel})
\end{itemize}

\paragraph{Execution} {Reservation table \textsf{LSU2\_MEMW\_AUXR\_AUXW} (Section~\ref{sec:reservation-LSU2-MEMW-AUXR-AUXW})}
~
{\begin{lstlisting}
stage ID:
  new acqrel = _ZX_1(acqrel);
stage RR:
  new acq = 1 & (acqrel >> 1);
  new rel = 1 & acqrel;
  new rs2 = BIR[rs2];
  new rs1 = BIR[rs1];
  new rd = BIR[rd];
  rd = _SX_8(MEM.atomic_cas(rs1, 0x1, (acq | (rel << 16)) ^ (acq | (rel << 16)), rs2.32[0], rd.32[0], rs1));
stage E1:
  if (rd != 0) {
    BIR[rd] = rd;
  }
\end{lstlisting}}
\newpage
\subsubsection[{\small AMOCAS.D: Atomic Memory Operation Compare-And-Swap Double Word}]{AMOCAS.D\hfill Atomic Memory Operation Compare-And-Swap Double Word}

\paragraph{Encoding} Format \textsf{RV64A\_CAS} (Section~\ref{sec:format-RV64A-CAS}), \verb|funct3 = 011|\\

\bitfields{ 5/00101, 2/acqrel, 5/rs2, 5/rs1, 3/011, 5/rd, 7/0101111 }


\paragraph{Syntax} \texttt{amocas.d}\emph{acqrel} \emph{rd}, \emph{rs2}, [\emph{rs1}]

\paragraph{Description} This instruction implements atomic compare-and-swap on double word. The effective address in \emph{rs1} must be double word aligned (multiple of 8).


\paragraph{Modifier(s)}
\noindent\begin{itemize}
\item acqrel (cf. acqrel in section \ref{par:modifier-acqrel} on page \pageref{par:modifier-acqrel})
\end{itemize}

\paragraph{Execution} {Reservation table \textsf{LSU2\_MEMW\_AUXR\_AUXW} (Section~\ref{sec:reservation-LSU2-MEMW-AUXR-AUXW})}
~
{\begin{lstlisting}
stage ID:
  new acqrel = _ZX_1(acqrel);
stage RR:
  new acq = 1 & (acqrel >> 1);
  new rel = 1 & acqrel;
  new rs2 = BIR[rs2];
  new rs1 = BIR[rs1];
  new rd = BIR[rd];
  rd = MEM.atomic_cas(rs1, 0xFF, (acq | (rel << 16)) ^ (acq | (rel << 16)), rs2.64[0], rd.64[0], rs1);
stage E1:
  if (rd != 0) {
    BIR[rd] = rd;
  }
\end{lstlisting}}
\newpage
\subsubsection[{\small AMOCAS.H: Atomic Memory Operation Compare-And-Swap half word}]{AMOCAS.H\hfill Atomic Memory Operation Compare-And-Swap half word}

\paragraph{Encoding} Format \textsf{RV64A\_CAS} (Section~\ref{sec:format-RV64A-CAS}), \verb|funct3 = 001|\\

\bitfields{ 5/00101, 2/acqrel, 5/rs2, 5/rs1, 3/001, 5/rd, 7/0101111 }


\paragraph{Syntax} \texttt{amocas.h}\emph{acqrel} \emph{rd}, \emph{rs2}, [\emph{rs1}]

\paragraph{Description} This instruction implements atomic compare-and-swap on half word. The effective address in \emph{rs1} must be word aligned (multiple of 2).


\paragraph{Modifier(s)}
\noindent\begin{itemize}
\item acqrel (cf. acqrel in section \ref{par:modifier-acqrel} on page \pageref{par:modifier-acqrel})
\end{itemize}

\paragraph{Execution} {Reservation table \textsf{LSU2\_MEMW\_AUXR\_AUXW} (Section~\ref{sec:reservation-LSU2-MEMW-AUXR-AUXW})}
~
{\begin{lstlisting}
stage ID:
  new acqrel = _ZX_1(acqrel);
stage RR:
  new acq = 1 & (acqrel >> 1);
  new rel = 1 & acqrel;
  new rs2 = BIR[rs2];
  new rs1 = BIR[rs1];
  new rd = BIR[rd];
  rd = _SX_16(MEM.atomic_cas(rs1, 0x3, (acq | (rel << 16)) ^ (acq | (rel << 16)), rs2.32[0], rd.32[0], rs1));
stage E1:
  if (rd != 0) {
    BIR[rd] = rd;
  }
\end{lstlisting}}
\newpage
\subsubsection[{\small AMOCAS.Q: Atomic Memory Operation Compare-And-Swap Quadruple Word}]{AMOCAS.Q\hfill Atomic Memory Operation Compare-And-Swap Quadruple Word}

\paragraph{Encoding} Format \textsf{RV64A\_CAS\_PAIR} (Section~\ref{sec:format-RV64A-CAS-PAIR}), \verb|funct3 = 100|\\

\bitfields{ 5/00101, 2/acqrel, 4/rs2-- pair, 1/0, 5/rs1, 3/100, 4/rd-- pair, 1/0, 7/0101111 }


\paragraph{Syntax} \texttt{amocas.q}\emph{acqrel} \emph{rd\_\discretionary{}{}{}pair}, \emph{rs2\_\discretionary{}{}{}pair}, [\emph{rs1}]

\paragraph{Description} This instruction implements atomic compare-and-swap on quadruple word. The effective address in \emph{rs1} must be quadruple word aligned (multiple of 16).


\paragraph{Modifier(s)}
\noindent\begin{itemize}
\item acqrel (cf. acqrel in section \ref{par:modifier-acqrel} on page \pageref{par:modifier-acqrel})
\end{itemize}

\paragraph{Execution} {Reservation table \textsf{LSU2\_MEMW\_AUXR\_AUXW} (Section~\ref{sec:reservation-LSU2-MEMW-AUXR-AUXW})}
~
{\begin{lstlisting}
stage ID:
  new acqrel = _ZX_1(acqrel);
stage RR:
  new acq = 1 & (acqrel >> 1);
  new rel = 1 & acqrel;
  new rs2 = BIRP[rs2_pair];
  new rs1 = BIR[rs1];
  new rd = BIRP[rd_pair];
  rd = MEM.atomic_cas(rs1, 0xFFFF, (acq | (rel << 16)) ^ (acq | (rel << 16)), rs2.128[0], rd.128[0], rs1);
stage E1:
  if (rd_pair != 0) {
    BIRP[rd_pair] = rd;
  }
\end{lstlisting}}
\newpage
\subsubsection[{\small AMOCAS.W: Atomic Memory Operation Compare-And-Swap Word}]{AMOCAS.W\hfill Atomic Memory Operation Compare-And-Swap Word}

\paragraph{Encoding} Format \textsf{RV64A\_CAS} (Section~\ref{sec:format-RV64A-CAS}), \verb|funct3 = 010|\\

\bitfields{ 5/00101, 2/acqrel, 5/rs2, 5/rs1, 3/010, 5/rd, 7/0101111 }


\paragraph{Syntax} \texttt{amocas.w}\emph{acqrel} \emph{rd}, \emph{rs2}, [\emph{rs1}]

\paragraph{Description} This instruction implements atomic compare-and-swap on word. The effective address in \emph{rs1} must be word aligned (multiple of 4).


\paragraph{Modifier(s)}
\noindent\begin{itemize}
\item acqrel (cf. acqrel in section \ref{par:modifier-acqrel} on page \pageref{par:modifier-acqrel})
\end{itemize}

\paragraph{Execution} {Reservation table \textsf{LSU2\_MEMW\_AUXR\_AUXW} (Section~\ref{sec:reservation-LSU2-MEMW-AUXR-AUXW})}
~
{\begin{lstlisting}
stage ID:
  new acqrel = _ZX_1(acqrel);
stage RR:
  new acq = 1 & (acqrel >> 1);
  new rel = 1 & acqrel;
  new rs2 = BIR[rs2];
  new rs1 = BIR[rs1];
  new rd = BIR[rd];
  rd = _SX_32(MEM.atomic_cas(rs1, 0xF, (acq | (rel << 16)) ^ (acq | (rel << 16)), rs2.32[0], rd.32[0], rs1));
stage E1:
  if (rd != 0) {
    BIR[rd] = rd;
  }
\end{lstlisting}}
\newpage
\subsubsection[{\small AMOMAX.B: Atomic Memory Operation MAX byte}]{AMOMAX.B\hfill Atomic Memory Operation MAX byte}

\paragraph{Encoding} Format \textsf{RV64A\_RMW} (Section~\ref{sec:format-RV64A-RMW}), \verb|funct5 = 10100|, \verb|funct3 = 000|\\

\bitfields{ 5/10100, 2/acqrel, 5/rs2, 5/rs1, 3/000, 5/rd, 7/0101111 }


\paragraph{Syntax} \texttt{amomax.b}\emph{acqrel} \emph{rd}, \emph{rs2}, [\emph{rs1}]

\paragraph{Description} This instruction implements atomic load-MAX-store on byte. The effective address is in \emph{rs1}.


\paragraph{Modifier(s)}
\noindent\begin{itemize}
\item acqrel (cf. acqrel in section \ref{par:modifier-acqrel} on page \pageref{par:modifier-acqrel})
\end{itemize}

\paragraph{Execution} {Reservation table \textsf{LSU2\_MEMW\_AUXR\_AUXW} (Section~\ref{sec:reservation-LSU2-MEMW-AUXR-AUXW})}
~
{\begin{lstlisting}
stage ID:
  new acqrel = _ZX_1(acqrel);
stage RR:
  new acq = 1 & (acqrel >> 1);
  new rel = 1 & acqrel;
  new rs2 = BIR[rs2];
  new rs1 = BIR[rs1];
  new rd = _SX_8(MEM.atomic_max(rs1, 0x1, acq | (rel << 16), rs2.32[0], rs1));
stage E1:
  if (rd != 0) {
    BIR[rd] = rd;
  }
\end{lstlisting}}
\newpage
\subsubsection[{\small AMOMAX.D: Atomic Memory Operation MAX Double Word}]{AMOMAX.D\hfill Atomic Memory Operation MAX Double Word}

\paragraph{Encoding} Format \textsf{RV64A\_RMW} (Section~\ref{sec:format-RV64A-RMW}), \verb|funct5 = 10100|, \verb|funct3 = 011|\\

\bitfields{ 5/10100, 2/acqrel, 5/rs2, 5/rs1, 3/011, 5/rd, 7/0101111 }


\paragraph{Syntax} \texttt{amomax.d}\emph{acqrel} \emph{rd}, \emph{rs2}, [\emph{rs1}]

\paragraph{Description} This instruction implements atomic load-MAX-store on double word. The effective address in \emph{rs1} must be double word aligned (multiple of 8).


\paragraph{Modifier(s)}
\noindent\begin{itemize}
\item acqrel (cf. acqrel in section \ref{par:modifier-acqrel} on page \pageref{par:modifier-acqrel})
\end{itemize}

\paragraph{Execution} {Reservation table \textsf{LSU2\_MEMW\_AUXR\_AUXW} (Section~\ref{sec:reservation-LSU2-MEMW-AUXR-AUXW})}
~
{\begin{lstlisting}
stage ID:
  new acqrel = _ZX_1(acqrel);
stage RR:
  new acq = 1 & (acqrel >> 1);
  new rel = 1 & acqrel;
  new rs2 = BIR[rs2];
  new rs1 = BIR[rs1];
  new rd = MEM.atomic_max(rs1, 0xFF, acq | (rel << 16), rs2.64[0], rs1);
stage E1:
  if (rd != 0) {
    BIR[rd] = rd;
  }
\end{lstlisting}}
\newpage
\subsubsection[{\small AMOMAX.H: Atomic Memory Operation MAX half Word}]{AMOMAX.H\hfill Atomic Memory Operation MAX half Word}

\paragraph{Encoding} Format \textsf{RV64A\_RMW} (Section~\ref{sec:format-RV64A-RMW}), \verb|funct5 = 10100|, \verb|funct3 = 001|\\

\bitfields{ 5/10100, 2/acqrel, 5/rs2, 5/rs1, 3/001, 5/rd, 7/0101111 }


\paragraph{Syntax} \texttt{amomax.h}\emph{acqrel} \emph{rd}, \emph{rs2}, [\emph{rs1}]

\paragraph{Description} This instruction implements atomic load-MAX-store on half word. The effective address in \emph{rs1} must be word aligned (multiple of 2).


\paragraph{Modifier(s)}
\noindent\begin{itemize}
\item acqrel (cf. acqrel in section \ref{par:modifier-acqrel} on page \pageref{par:modifier-acqrel})
\end{itemize}

\paragraph{Execution} {Reservation table \textsf{LSU2\_MEMW\_AUXR\_AUXW} (Section~\ref{sec:reservation-LSU2-MEMW-AUXR-AUXW})}
~
{\begin{lstlisting}
stage ID:
  new acqrel = _ZX_1(acqrel);
stage RR:
  new acq = 1 & (acqrel >> 1);
  new rel = 1 & acqrel;
  new rs2 = BIR[rs2];
  new rs1 = BIR[rs1];
  new rd = _SX_16(MEM.atomic_max(rs1, 0x3, acq | (rel << 16), rs2.32[0], rs1));
stage E1:
  if (rd != 0) {
    BIR[rd] = rd;
  }
\end{lstlisting}}
\newpage
\subsubsection[{\small AMOMAX.W: Atomic Memory Operation MAX Word}]{AMOMAX.W\hfill Atomic Memory Operation MAX Word}

\paragraph{Encoding} Format \textsf{RV64A\_RMW} (Section~\ref{sec:format-RV64A-RMW}), \verb|funct5 = 10100|, \verb|funct3 = 010|\\

\bitfields{ 5/10100, 2/acqrel, 5/rs2, 5/rs1, 3/010, 5/rd, 7/0101111 }


\paragraph{Syntax} \texttt{amomax.w}\emph{acqrel} \emph{rd}, \emph{rs2}, [\emph{rs1}]

\paragraph{Description} This instruction implements atomic load-MAX-store on word. The effective address in \emph{rs1} must be word aligned (multiple of 4).


\paragraph{Modifier(s)}
\noindent\begin{itemize}
\item acqrel (cf. acqrel in section \ref{par:modifier-acqrel} on page \pageref{par:modifier-acqrel})
\end{itemize}

\paragraph{Execution} {Reservation table \textsf{LSU2\_MEMW\_AUXR\_AUXW} (Section~\ref{sec:reservation-LSU2-MEMW-AUXR-AUXW})}
~
{\begin{lstlisting}
stage ID:
  new acqrel = _ZX_1(acqrel);
stage RR:
  new acq = 1 & (acqrel >> 1);
  new rel = 1 & acqrel;
  new rs2 = BIR[rs2];
  new rs1 = BIR[rs1];
  new rd = _SX_32(MEM.atomic_max(rs1, 0xF, acq | (rel << 16), rs2.32[0], rs1));
stage E1:
  if (rd != 0) {
    BIR[rd] = rd;
  }
\end{lstlisting}}
\newpage
\subsubsection[{\small AMOMAXU.B: Atomic Memory Operation MAX Unsigned byte}]{AMOMAXU.B\hfill Atomic Memory Operation MAX Unsigned byte}

\paragraph{Encoding} Format \textsf{RV64A\_RMW} (Section~\ref{sec:format-RV64A-RMW}), \verb|funct5 = 11100|, \verb|funct3 = 000|\\

\bitfields{ 5/11100, 2/acqrel, 5/rs2, 5/rs1, 3/000, 5/rd, 7/0101111 }


\paragraph{Syntax} \texttt{amomaxu.b}\emph{acqrel} \emph{rd}, \emph{rs2}, [\emph{rs1}]

\paragraph{Description} This instruction implements atomic load-MAXU-store on byte. The effective address is in \emph{rs1}.


\paragraph{Modifier(s)}
\noindent\begin{itemize}
\item acqrel (cf. acqrel in section \ref{par:modifier-acqrel} on page \pageref{par:modifier-acqrel})
\end{itemize}

\paragraph{Execution} {Reservation table \textsf{LSU2\_MEMW\_AUXR\_AUXW} (Section~\ref{sec:reservation-LSU2-MEMW-AUXR-AUXW})}
~
{\begin{lstlisting}
stage ID:
  new acqrel = _ZX_1(acqrel);
stage RR:
  new acq = 1 & (acqrel >> 1);
  new rel = 1 & acqrel;
  new rs2 = BIR[rs2];
  new rs1 = BIR[rs1];
  new rd = _SX_8(MEM.atomic_maxu(rs1, 0x1, acq | (rel << 16), rs2.32[0], rs1));
stage E1:
  if (rd != 0) {
    BIR[rd] = rd;
  }
\end{lstlisting}}
\newpage
\subsubsection[{\small AMOMAXU.D: Atomic Memory Operation MAX Unsigned Double Word}]{AMOMAXU.D\hfill Atomic Memory Operation MAX Unsigned Double Word}

\paragraph{Encoding} Format \textsf{RV64A\_RMW} (Section~\ref{sec:format-RV64A-RMW}), \verb|funct5 = 11100|, \verb|funct3 = 011|\\

\bitfields{ 5/11100, 2/acqrel, 5/rs2, 5/rs1, 3/011, 5/rd, 7/0101111 }


\paragraph{Syntax} \texttt{amomaxu.d}\emph{acqrel} \emph{rd}, \emph{rs2}, [\emph{rs1}]

\paragraph{Description} This instruction implements atomic load-MAXU-store on double word. The effective address in \emph{rs1} must be double word aligned (multiple of 8).


\paragraph{Modifier(s)}
\noindent\begin{itemize}
\item acqrel (cf. acqrel in section \ref{par:modifier-acqrel} on page \pageref{par:modifier-acqrel})
\end{itemize}

\paragraph{Execution} {Reservation table \textsf{LSU2\_MEMW\_AUXR\_AUXW} (Section~\ref{sec:reservation-LSU2-MEMW-AUXR-AUXW})}
~
{\begin{lstlisting}
stage ID:
  new acqrel = _ZX_1(acqrel);
stage RR:
  new acq = 1 & (acqrel >> 1);
  new rel = 1 & acqrel;
  new rs2 = BIR[rs2];
  new rs1 = BIR[rs1];
  new rd = MEM.atomic_maxu(rs1, 0xFF, acq | (rel << 16), rs2.64[0], rs1);
stage E1:
  if (rd != 0) {
    BIR[rd] = rd;
  }
\end{lstlisting}}
\newpage
\subsubsection[{\small AMOMAXU.H: Atomic Memory Operation MAX Unsigned half Word}]{AMOMAXU.H\hfill Atomic Memory Operation MAX Unsigned half Word}

\paragraph{Encoding} Format \textsf{RV64A\_RMW} (Section~\ref{sec:format-RV64A-RMW}), \verb|funct5 = 11100|, \verb|funct3 = 001|\\

\bitfields{ 5/11100, 2/acqrel, 5/rs2, 5/rs1, 3/001, 5/rd, 7/0101111 }


\paragraph{Syntax} \texttt{amomaxu.h}\emph{acqrel} \emph{rd}, \emph{rs2}, [\emph{rs1}]

\paragraph{Description} This instruction implements atomic load-MAXU-store on half word. The effective address in \emph{rs1} must be word aligned (multiple of 2).


\paragraph{Modifier(s)}
\noindent\begin{itemize}
\item acqrel (cf. acqrel in section \ref{par:modifier-acqrel} on page \pageref{par:modifier-acqrel})
\end{itemize}

\paragraph{Execution} {Reservation table \textsf{LSU2\_MEMW\_AUXR\_AUXW} (Section~\ref{sec:reservation-LSU2-MEMW-AUXR-AUXW})}
~
{\begin{lstlisting}
stage ID:
  new acqrel = _ZX_1(acqrel);
stage RR:
  new acq = 1 & (acqrel >> 1);
  new rel = 1 & acqrel;
  new rs2 = BIR[rs2];
  new rs1 = BIR[rs1];
  new rd = _SX_16(MEM.atomic_maxu(rs1, 0x3, acq | (rel << 16), rs2.32[0], rs1));
stage E1:
  if (rd != 0) {
    BIR[rd] = rd;
  }
\end{lstlisting}}
\newpage
\subsubsection[{\small AMOMAXU.W: Atomic Memory Operation MAX Unsigned Word}]{AMOMAXU.W\hfill Atomic Memory Operation MAX Unsigned Word}

\paragraph{Encoding} Format \textsf{RV64A\_RMW} (Section~\ref{sec:format-RV64A-RMW}), \verb|funct5 = 11100|, \verb|funct3 = 010|\\

\bitfields{ 5/11100, 2/acqrel, 5/rs2, 5/rs1, 3/010, 5/rd, 7/0101111 }


\paragraph{Syntax} \texttt{amomaxu.w}\emph{acqrel} \emph{rd}, \emph{rs2}, [\emph{rs1}]

\paragraph{Description} This instruction implements atomic load-MAXU-store on word. The effective address in \emph{rs1} must be word aligned (multiple of 4).


\paragraph{Modifier(s)}
\noindent\begin{itemize}
\item acqrel (cf. acqrel in section \ref{par:modifier-acqrel} on page \pageref{par:modifier-acqrel})
\end{itemize}

\paragraph{Execution} {Reservation table \textsf{LSU2\_MEMW\_AUXR\_AUXW} (Section~\ref{sec:reservation-LSU2-MEMW-AUXR-AUXW})}
~
{\begin{lstlisting}
stage ID:
  new acqrel = _ZX_1(acqrel);
stage RR:
  new acq = 1 & (acqrel >> 1);
  new rel = 1 & acqrel;
  new rs2 = BIR[rs2];
  new rs1 = BIR[rs1];
  new rd = _SX_32(MEM.atomic_maxu(rs1, 0xF, acq | (rel << 16), rs2.32[0], rs1));
stage E1:
  if (rd != 0) {
    BIR[rd] = rd;
  }
\end{lstlisting}}
\newpage
\subsubsection[{\small AMOMIN.B: Atomic Memory Operation MIN byte}]{AMOMIN.B\hfill Atomic Memory Operation MIN byte}

\paragraph{Encoding} Format \textsf{RV64A\_RMW} (Section~\ref{sec:format-RV64A-RMW}), \verb|funct5 = 10000|, \verb|funct3 = 000|\\

\bitfields{ 5/10000, 2/acqrel, 5/rs2, 5/rs1, 3/000, 5/rd, 7/0101111 }


\paragraph{Syntax} \texttt{amomin.b}\emph{acqrel} \emph{rd}, \emph{rs2}, [\emph{rs1}]

\paragraph{Description} This instruction implements atomic load-MIN-store on byte. The effective address is in \emph{rs1}.


\paragraph{Modifier(s)}
\noindent\begin{itemize}
\item acqrel (cf. acqrel in section \ref{par:modifier-acqrel} on page \pageref{par:modifier-acqrel})
\end{itemize}

\paragraph{Execution} {Reservation table \textsf{LSU2\_MEMW\_AUXR\_AUXW} (Section~\ref{sec:reservation-LSU2-MEMW-AUXR-AUXW})}
~
{\begin{lstlisting}
stage ID:
  new acqrel = _ZX_1(acqrel);
stage RR:
  new acq = 1 & (acqrel >> 1);
  new rel = 1 & acqrel;
  new rs2 = BIR[rs2];
  new rs1 = BIR[rs1];
  new rd = _SX_8(MEM.atomic_min(rs1, 0x1, acq | (rel << 16), rs2.32[0], rs1));
stage E1:
  if (rd != 0) {
    BIR[rd] = rd;
  }
\end{lstlisting}}
\newpage
\subsubsection[{\small AMOMIN.D: Atomic Memory Operation MIN Double Word}]{AMOMIN.D\hfill Atomic Memory Operation MIN Double Word}

\paragraph{Encoding} Format \textsf{RV64A\_RMW} (Section~\ref{sec:format-RV64A-RMW}), \verb|funct5 = 10000|, \verb|funct3 = 011|\\

\bitfields{ 5/10000, 2/acqrel, 5/rs2, 5/rs1, 3/011, 5/rd, 7/0101111 }


\paragraph{Syntax} \texttt{amomin.d}\emph{acqrel} \emph{rd}, \emph{rs2}, [\emph{rs1}]

\paragraph{Description} This instruction implements atomic load-MIN-store on double word. The effective address in \emph{rs1} must be double word aligned (multiple of 8).


\paragraph{Modifier(s)}
\noindent\begin{itemize}
\item acqrel (cf. acqrel in section \ref{par:modifier-acqrel} on page \pageref{par:modifier-acqrel})
\end{itemize}

\paragraph{Execution} {Reservation table \textsf{LSU2\_MEMW\_AUXR\_AUXW} (Section~\ref{sec:reservation-LSU2-MEMW-AUXR-AUXW})}
~
{\begin{lstlisting}
stage ID:
  new acqrel = _ZX_1(acqrel);
stage RR:
  new acq = 1 & (acqrel >> 1);
  new rel = 1 & acqrel;
  new rs2 = BIR[rs2];
  new rs1 = BIR[rs1];
  new rd = MEM.atomic_min(rs1, 0xFF, acq | (rel << 16), rs2.64[0], rs1);
stage E1:
  if (rd != 0) {
    BIR[rd] = rd;
  }
\end{lstlisting}}
\newpage
\subsubsection[{\small AMOMIN.H: Atomic Memory Operation MIN half Word}]{AMOMIN.H\hfill Atomic Memory Operation MIN half Word}

\paragraph{Encoding} Format \textsf{RV64A\_RMW} (Section~\ref{sec:format-RV64A-RMW}), \verb|funct5 = 10000|, \verb|funct3 = 001|\\

\bitfields{ 5/10000, 2/acqrel, 5/rs2, 5/rs1, 3/001, 5/rd, 7/0101111 }


\paragraph{Syntax} \texttt{amomin.h}\emph{acqrel} \emph{rd}, \emph{rs2}, [\emph{rs1}]

\paragraph{Description} This instruction implements atomic load-MIN-store on half word. The effective address in \emph{rs1} must be word aligned (multiple of 2).


\paragraph{Modifier(s)}
\noindent\begin{itemize}
\item acqrel (cf. acqrel in section \ref{par:modifier-acqrel} on page \pageref{par:modifier-acqrel})
\end{itemize}

\paragraph{Execution} {Reservation table \textsf{LSU2\_MEMW\_AUXR\_AUXW} (Section~\ref{sec:reservation-LSU2-MEMW-AUXR-AUXW})}
~
{\begin{lstlisting}
stage ID:
  new acqrel = _ZX_1(acqrel);
stage RR:
  new acq = 1 & (acqrel >> 1);
  new rel = 1 & acqrel;
  new rs2 = BIR[rs2];
  new rs1 = BIR[rs1];
  new rd = _SX_16(MEM.atomic_min(rs1, 0x3, acq | (rel << 16), rs2.32[0], rs1));
stage E1:
  if (rd != 0) {
    BIR[rd] = rd;
  }
\end{lstlisting}}
\newpage
\subsubsection[{\small AMOMIN.W: Atomic Memory Operation MIN Word}]{AMOMIN.W\hfill Atomic Memory Operation MIN Word}

\paragraph{Encoding} Format \textsf{RV64A\_RMW} (Section~\ref{sec:format-RV64A-RMW}), \verb|funct5 = 10000|, \verb|funct3 = 010|\\

\bitfields{ 5/10000, 2/acqrel, 5/rs2, 5/rs1, 3/010, 5/rd, 7/0101111 }


\paragraph{Syntax} \texttt{amomin.w}\emph{acqrel} \emph{rd}, \emph{rs2}, [\emph{rs1}]

\paragraph{Description} This instruction implements atomic load-MIN-store on word. The effective address in \emph{rs1} must be word aligned (multiple of 4).


\paragraph{Modifier(s)}
\noindent\begin{itemize}
\item acqrel (cf. acqrel in section \ref{par:modifier-acqrel} on page \pageref{par:modifier-acqrel})
\end{itemize}

\paragraph{Execution} {Reservation table \textsf{LSU2\_MEMW\_AUXR\_AUXW} (Section~\ref{sec:reservation-LSU2-MEMW-AUXR-AUXW})}
~
{\begin{lstlisting}
stage ID:
  new acqrel = _ZX_1(acqrel);
stage RR:
  new acq = 1 & (acqrel >> 1);
  new rel = 1 & acqrel;
  new rs2 = BIR[rs2];
  new rs1 = BIR[rs1];
  new rd = _SX_32(MEM.atomic_min(rs1, 0xF, acq | (rel << 16), rs2.32[0], rs1));
stage E1:
  if (rd != 0) {
    BIR[rd] = rd;
  }
\end{lstlisting}}
\newpage
\subsubsection[{\small AMOMINU.B: Atomic Memory Operation MIN Unsigned byte}]{AMOMINU.B\hfill Atomic Memory Operation MIN Unsigned byte}

\paragraph{Encoding} Format \textsf{RV64A\_RMW} (Section~\ref{sec:format-RV64A-RMW}), \verb|funct5 = 11000|, \verb|funct3 = 000|\\

\bitfields{ 5/11000, 2/acqrel, 5/rs2, 5/rs1, 3/000, 5/rd, 7/0101111 }


\paragraph{Syntax} \texttt{amominu.b}\emph{acqrel} \emph{rd}, \emph{rs2}, [\emph{rs1}]

\paragraph{Description} This instruction implements atomic load-MINU-store on byte. The effective address is in \emph{rs1}.


\paragraph{Modifier(s)}
\noindent\begin{itemize}
\item acqrel (cf. acqrel in section \ref{par:modifier-acqrel} on page \pageref{par:modifier-acqrel})
\end{itemize}

\paragraph{Execution} {Reservation table \textsf{LSU2\_MEMW\_AUXR\_AUXW} (Section~\ref{sec:reservation-LSU2-MEMW-AUXR-AUXW})}
~
{\begin{lstlisting}
stage ID:
  new acqrel = _ZX_1(acqrel);
stage RR:
  new acq = 1 & (acqrel >> 1);
  new rel = 1 & acqrel;
  new rs2 = BIR[rs2];
  new rs1 = BIR[rs1];
  new rd = _SX_8(MEM.atomic_minu(rs1, 0x1, acq | (rel << 16), rs2.32[0], rs1));
stage E1:
  if (rd != 0) {
    BIR[rd] = rd;
  }
\end{lstlisting}}
\newpage
\subsubsection[{\small AMOMINU.D: Atomic Memory Operation MIN Unsigned Double Word}]{AMOMINU.D\hfill Atomic Memory Operation MIN Unsigned Double Word}

\paragraph{Encoding} Format \textsf{RV64A\_RMW} (Section~\ref{sec:format-RV64A-RMW}), \verb|funct5 = 11000|, \verb|funct3 = 011|\\

\bitfields{ 5/11000, 2/acqrel, 5/rs2, 5/rs1, 3/011, 5/rd, 7/0101111 }


\paragraph{Syntax} \texttt{amominu.d}\emph{acqrel} \emph{rd}, \emph{rs2}, [\emph{rs1}]

\paragraph{Description} This instruction implements atomic load-MINU-store on double word. The effective address in \emph{rs1} must be double word aligned (multiple of 8).


\paragraph{Modifier(s)}
\noindent\begin{itemize}
\item acqrel (cf. acqrel in section \ref{par:modifier-acqrel} on page \pageref{par:modifier-acqrel})
\end{itemize}

\paragraph{Execution} {Reservation table \textsf{LSU2\_MEMW\_AUXR\_AUXW} (Section~\ref{sec:reservation-LSU2-MEMW-AUXR-AUXW})}
~
{\begin{lstlisting}
stage ID:
  new acqrel = _ZX_1(acqrel);
stage RR:
  new acq = 1 & (acqrel >> 1);
  new rel = 1 & acqrel;
  new rs2 = BIR[rs2];
  new rs1 = BIR[rs1];
  new rd = MEM.atomic_minu(rs1, 0xFF, acq | (rel << 16), rs2.64[0], rs1);
stage E1:
  if (rd != 0) {
    BIR[rd] = rd;
  }
\end{lstlisting}}
\newpage
\subsubsection[{\small AMOMINU.H: Atomic Memory Operation MIN Unsigned half Word}]{AMOMINU.H\hfill Atomic Memory Operation MIN Unsigned half Word}

\paragraph{Encoding} Format \textsf{RV64A\_RMW} (Section~\ref{sec:format-RV64A-RMW}), \verb|funct5 = 11000|, \verb|funct3 = 001|\\

\bitfields{ 5/11000, 2/acqrel, 5/rs2, 5/rs1, 3/001, 5/rd, 7/0101111 }


\paragraph{Syntax} \texttt{amominu.h}\emph{acqrel} \emph{rd}, \emph{rs2}, [\emph{rs1}]

\paragraph{Description} This instruction implements atomic load-MINU-store on half word. The effective address in \emph{rs1} must be word aligned (multiple of 2).


\paragraph{Modifier(s)}
\noindent\begin{itemize}
\item acqrel (cf. acqrel in section \ref{par:modifier-acqrel} on page \pageref{par:modifier-acqrel})
\end{itemize}

\paragraph{Execution} {Reservation table \textsf{LSU2\_MEMW\_AUXR\_AUXW} (Section~\ref{sec:reservation-LSU2-MEMW-AUXR-AUXW})}
~
{\begin{lstlisting}
stage ID:
  new acqrel = _ZX_1(acqrel);
stage RR:
  new acq = 1 & (acqrel >> 1);
  new rel = 1 & acqrel;
  new rs2 = BIR[rs2];
  new rs1 = BIR[rs1];
  new rd = _SX_16(MEM.atomic_minu(rs1, 0x3, acq | (rel << 16), rs2.32[0], rs1));
stage E1:
  if (rd != 0) {
    BIR[rd] = rd;
  }
\end{lstlisting}}
\newpage
\subsubsection[{\small AMOMINU.W: Atomic Memory Operation MIN Unsigned Word}]{AMOMINU.W\hfill Atomic Memory Operation MIN Unsigned Word}

\paragraph{Encoding} Format \textsf{RV64A\_RMW} (Section~\ref{sec:format-RV64A-RMW}), \verb|funct5 = 11000|, \verb|funct3 = 010|\\

\bitfields{ 5/11000, 2/acqrel, 5/rs2, 5/rs1, 3/010, 5/rd, 7/0101111 }


\paragraph{Syntax} \texttt{amominu.w}\emph{acqrel} \emph{rd}, \emph{rs2}, [\emph{rs1}]

\paragraph{Description} This instruction implements atomic load-MINU-store on word. The effective address in \emph{rs1} must be word aligned (multiple of 4).


\paragraph{Modifier(s)}
\noindent\begin{itemize}
\item acqrel (cf. acqrel in section \ref{par:modifier-acqrel} on page \pageref{par:modifier-acqrel})
\end{itemize}

\paragraph{Execution} {Reservation table \textsf{LSU2\_MEMW\_AUXR\_AUXW} (Section~\ref{sec:reservation-LSU2-MEMW-AUXR-AUXW})}
~
{\begin{lstlisting}
stage ID:
  new acqrel = _ZX_1(acqrel);
stage RR:
  new acq = 1 & (acqrel >> 1);
  new rel = 1 & acqrel;
  new rs2 = BIR[rs2];
  new rs1 = BIR[rs1];
  new rd = _SX_32(MEM.atomic_minu(rs1, 0xF, acq | (rel << 16), rs2.32[0], rs1));
stage E1:
  if (rd != 0) {
    BIR[rd] = rd;
  }
\end{lstlisting}}
\newpage
\subsubsection[{\small AMOOR.B: Atomic Memory Operation OR byte}]{AMOOR.B\hfill Atomic Memory Operation OR byte}

\paragraph{Encoding} Format \textsf{RV64A\_RMW} (Section~\ref{sec:format-RV64A-RMW}), \verb|funct5 = 01000|, \verb|funct3 = 000|\\

\bitfields{ 5/01000, 2/acqrel, 5/rs2, 5/rs1, 3/000, 5/rd, 7/0101111 }


\paragraph{Syntax} \texttt{amoor.b}\emph{acqrel} \emph{rd}, \emph{rs2}, [\emph{rs1}]

\paragraph{Description} This instruction implements atomic load-OR-store on byte. The effective address is in \emph{rs1}.


\paragraph{Modifier(s)}
\noindent\begin{itemize}
\item acqrel (cf. acqrel in section \ref{par:modifier-acqrel} on page \pageref{par:modifier-acqrel})
\end{itemize}

\paragraph{Execution} {Reservation table \textsf{LSU2\_MEMW\_AUXR\_AUXW} (Section~\ref{sec:reservation-LSU2-MEMW-AUXR-AUXW})}
~
{\begin{lstlisting}
stage ID:
  new acqrel = _ZX_1(acqrel);
stage RR:
  new acq = 1 & (acqrel >> 1);
  new rel = 1 & acqrel;
  new rs2 = BIR[rs2];
  new rs1 = BIR[rs1];
  new rd = _SX_8(MEM.atomic_ior(rs1, 0x1, acq | (rel << 16), rs2.32[0], rs1));
stage E1:
  if (rd != 0) {
    BIR[rd] = rd;
  }
\end{lstlisting}}
\newpage
\subsubsection[{\small AMOOR.D: Atomic Memory Operation OR Double Word}]{AMOOR.D\hfill Atomic Memory Operation OR Double Word}

\paragraph{Encoding} Format \textsf{RV64A\_RMW} (Section~\ref{sec:format-RV64A-RMW}), \verb|funct5 = 01000|, \verb|funct3 = 011|\\

\bitfields{ 5/01000, 2/acqrel, 5/rs2, 5/rs1, 3/011, 5/rd, 7/0101111 }


\paragraph{Syntax} \texttt{amoor.d}\emph{acqrel} \emph{rd}, \emph{rs2}, [\emph{rs1}]

\paragraph{Description} This instruction implements atomic load-OR-store on double word. The effective address in \emph{rs1} must be double word aligned (multiple of 8).


\paragraph{Modifier(s)}
\noindent\begin{itemize}
\item acqrel (cf. acqrel in section \ref{par:modifier-acqrel} on page \pageref{par:modifier-acqrel})
\end{itemize}

\paragraph{Execution} {Reservation table \textsf{LSU2\_MEMW\_AUXR\_AUXW} (Section~\ref{sec:reservation-LSU2-MEMW-AUXR-AUXW})}
~
{\begin{lstlisting}
stage ID:
  new acqrel = _ZX_1(acqrel);
stage RR:
  new acq = 1 & (acqrel >> 1);
  new rel = 1 & acqrel;
  new rs2 = BIR[rs2];
  new rs1 = BIR[rs1];
  new rd = MEM.atomic_ior(rs1, 0xFF, acq | (rel << 16), rs2.64[0], rs1);
stage E1:
  if (rd != 0) {
    BIR[rd] = rd;
  }
\end{lstlisting}}
\newpage
\subsubsection[{\small AMOOR.H: Atomic Memory Operation OR half Word}]{AMOOR.H\hfill Atomic Memory Operation OR half Word}

\paragraph{Encoding} Format \textsf{RV64A\_RMW} (Section~\ref{sec:format-RV64A-RMW}), \verb|funct5 = 01000|, \verb|funct3 = 001|\\

\bitfields{ 5/01000, 2/acqrel, 5/rs2, 5/rs1, 3/001, 5/rd, 7/0101111 }


\paragraph{Syntax} \texttt{amoor.h}\emph{acqrel} \emph{rd}, \emph{rs2}, [\emph{rs1}]

\paragraph{Description} This instruction implements atomic load-OR-store on half word. The effective address in \emph{rs1} must be word aligned (multiple of 2).


\paragraph{Modifier(s)}
\noindent\begin{itemize}
\item acqrel (cf. acqrel in section \ref{par:modifier-acqrel} on page \pageref{par:modifier-acqrel})
\end{itemize}

\paragraph{Execution} {Reservation table \textsf{LSU2\_MEMW\_AUXR\_AUXW} (Section~\ref{sec:reservation-LSU2-MEMW-AUXR-AUXW})}
~
{\begin{lstlisting}
stage ID:
  new acqrel = _ZX_1(acqrel);
stage RR:
  new acq = 1 & (acqrel >> 1);
  new rel = 1 & acqrel;
  new rs2 = BIR[rs2];
  new rs1 = BIR[rs1];
  new rd = _SX_16(MEM.atomic_ior(rs1, 0x3, acq | (rel << 16), rs2.32[0], rs1));
stage E1:
  if (rd != 0) {
    BIR[rd] = rd;
  }
\end{lstlisting}}
\newpage
\subsubsection[{\small AMOOR.W: Atomic Memory Operation OR Word}]{AMOOR.W\hfill Atomic Memory Operation OR Word}

\paragraph{Encoding} Format \textsf{RV64A\_RMW} (Section~\ref{sec:format-RV64A-RMW}), \verb|funct5 = 01000|, \verb|funct3 = 010|\\

\bitfields{ 5/01000, 2/acqrel, 5/rs2, 5/rs1, 3/010, 5/rd, 7/0101111 }


\paragraph{Syntax} \texttt{amoor.w}\emph{acqrel} \emph{rd}, \emph{rs2}, [\emph{rs1}]

\paragraph{Description} This instruction implements atomic load-OR-store on word. The effective address in \emph{rs1} must be word aligned (multiple of 4).


\paragraph{Modifier(s)}
\noindent\begin{itemize}
\item acqrel (cf. acqrel in section \ref{par:modifier-acqrel} on page \pageref{par:modifier-acqrel})
\end{itemize}

\paragraph{Execution} {Reservation table \textsf{LSU2\_MEMW\_AUXR\_AUXW} (Section~\ref{sec:reservation-LSU2-MEMW-AUXR-AUXW})}
~
{\begin{lstlisting}
stage ID:
  new acqrel = _ZX_1(acqrel);
stage RR:
  new acq = 1 & (acqrel >> 1);
  new rel = 1 & acqrel;
  new rs2 = BIR[rs2];
  new rs1 = BIR[rs1];
  new rd = _SX_32(MEM.atomic_ior(rs1, 0xF, acq | (rel << 16), rs2.32[0], rs1));
stage E1:
  if (rd != 0) {
    BIR[rd] = rd;
  }
\end{lstlisting}}
\newpage
\subsubsection[{\small AMOSWAP.B: Atomic Memory Operation Swap byte}]{AMOSWAP.B\hfill Atomic Memory Operation Swap byte}

\paragraph{Encoding} Format \textsf{RV64A\_RMW} (Section~\ref{sec:format-RV64A-RMW}), \verb|funct5 = 00001|, \verb|funct3 = 000|\\

\bitfields{ 5/00001, 2/acqrel, 5/rs2, 5/rs1, 3/000, 5/rd, 7/0101111 }


\paragraph{Syntax} \texttt{amoswap.b}\emph{acqrel} \emph{rd}, \emph{rs2}, [\emph{rs1}]

\paragraph{Description} This instruction implements atomic swap on byte. The effective address is in \emph{rs1}.


\paragraph{Modifier(s)}
\noindent\begin{itemize}
\item acqrel (cf. acqrel in section \ref{par:modifier-acqrel} on page \pageref{par:modifier-acqrel})
\end{itemize}

\paragraph{Execution} {Reservation table \textsf{LSU2\_MEMW\_AUXR\_AUXW} (Section~\ref{sec:reservation-LSU2-MEMW-AUXR-AUXW})}
~
{\begin{lstlisting}
stage ID:
  new acqrel = _ZX_1(acqrel);
stage RR:
  new acq = 1 & (acqrel >> 1);
  new rel = 1 & acqrel;
  new rs2 = BIR[rs2];
  new rs1 = BIR[rs1];
  new rd = _SX_8(MEM.atomic_swap(rs1, 0x1, acq | (rel << 16), rs2.32[0], rs1));
stage E1:
  if (rd != 0) {
    BIR[rd] = rd;
  }
\end{lstlisting}}
\newpage
\subsubsection[{\small AMOSWAP.D: Atomic Memory Operation Swap Double Word}]{AMOSWAP.D\hfill Atomic Memory Operation Swap Double Word}

\paragraph{Encoding} Format \textsf{RV64A\_RMW} (Section~\ref{sec:format-RV64A-RMW}), \verb|funct5 = 00001|, \verb|funct3 = 011|\\

\bitfields{ 5/00001, 2/acqrel, 5/rs2, 5/rs1, 3/011, 5/rd, 7/0101111 }


\paragraph{Syntax} \texttt{amoswap.d}\emph{acqrel} \emph{rd}, \emph{rs2}, [\emph{rs1}]

\paragraph{Description} This instruction implements atomic swap on double word. The effective address in \emph{rs1} must be double word aligned (multiple of 8).


\paragraph{Modifier(s)}
\noindent\begin{itemize}
\item acqrel (cf. acqrel in section \ref{par:modifier-acqrel} on page \pageref{par:modifier-acqrel})
\end{itemize}

\paragraph{Execution} {Reservation table \textsf{LSU2\_MEMW\_AUXR\_AUXW} (Section~\ref{sec:reservation-LSU2-MEMW-AUXR-AUXW})}
~
{\begin{lstlisting}
stage ID:
  new acqrel = _ZX_1(acqrel);
stage RR:
  new acq = 1 & (acqrel >> 1);
  new rel = 1 & acqrel;
  new rs2 = BIR[rs2];
  new rs1 = BIR[rs1];
  new rd = MEM.atomic_swap(rs1, 0xFF, acq | (rel << 16), rs2.64[0], rs1);
stage E1:
  if (rd != 0) {
    BIR[rd] = rd;
  }
\end{lstlisting}}
\newpage
\subsubsection[{\small AMOSWAP.H: Atomic Memory Operation Swap half Word}]{AMOSWAP.H\hfill Atomic Memory Operation Swap half Word}

\paragraph{Encoding} Format \textsf{RV64A\_RMW} (Section~\ref{sec:format-RV64A-RMW}), \verb|funct5 = 00001|, \verb|funct3 = 001|\\

\bitfields{ 5/00001, 2/acqrel, 5/rs2, 5/rs1, 3/001, 5/rd, 7/0101111 }


\paragraph{Syntax} \texttt{amoswap.h}\emph{acqrel} \emph{rd}, \emph{rs2}, [\emph{rs1}]

\paragraph{Description} This instruction implements atomic swap on half word. The effective address in \emph{rs1} must be word aligned (multiple of 2).


\paragraph{Modifier(s)}
\noindent\begin{itemize}
\item acqrel (cf. acqrel in section \ref{par:modifier-acqrel} on page \pageref{par:modifier-acqrel})
\end{itemize}

\paragraph{Execution} {Reservation table \textsf{LSU2\_MEMW\_AUXR\_AUXW} (Section~\ref{sec:reservation-LSU2-MEMW-AUXR-AUXW})}
~
{\begin{lstlisting}
stage ID:
  new acqrel = _ZX_1(acqrel);
stage RR:
  new acq = 1 & (acqrel >> 1);
  new rel = 1 & acqrel;
  new rs2 = BIR[rs2];
  new rs1 = BIR[rs1];
  new rd = _SX_16(MEM.atomic_swap(rs1, 0x3, acq | (rel << 16), rs2.32[0], rs1));
stage E1:
  if (rd != 0) {
    BIR[rd] = rd;
  }
\end{lstlisting}}
\newpage
\subsubsection[{\small AMOSWAP.W: Atomic Memory Operation Swap Word}]{AMOSWAP.W\hfill Atomic Memory Operation Swap Word}

\paragraph{Encoding} Format \textsf{RV64A\_RMW} (Section~\ref{sec:format-RV64A-RMW}), \verb|funct5 = 00001|, \verb|funct3 = 010|\\

\bitfields{ 5/00001, 2/acqrel, 5/rs2, 5/rs1, 3/010, 5/rd, 7/0101111 }


\paragraph{Syntax} \texttt{amoswap.w}\emph{acqrel} \emph{rd}, \emph{rs2}, [\emph{rs1}]

\paragraph{Description} This instruction implements atomic swap on word. The effective address in \emph{rs1} must be word aligned (multiple of 4).


\paragraph{Modifier(s)}
\noindent\begin{itemize}
\item acqrel (cf. acqrel in section \ref{par:modifier-acqrel} on page \pageref{par:modifier-acqrel})
\end{itemize}

\paragraph{Execution} {Reservation table \textsf{LSU2\_MEMW\_AUXR\_AUXW} (Section~\ref{sec:reservation-LSU2-MEMW-AUXR-AUXW})}
~
{\begin{lstlisting}
stage ID:
  new acqrel = _ZX_1(acqrel);
stage RR:
  new acq = 1 & (acqrel >> 1);
  new rel = 1 & acqrel;
  new rs2 = BIR[rs2];
  new rs1 = BIR[rs1];
  new rd = _SX_32(MEM.atomic_swap(rs1, 0xF, acq | (rel << 16), rs2.32[0], rs1));
stage E1:
  if (rd != 0) {
    BIR[rd] = rd;
  }
\end{lstlisting}}
\newpage
\subsubsection[{\small AMOXOR.B: Atomic Memory Operation XOR byte}]{AMOXOR.B\hfill Atomic Memory Operation XOR byte}

\paragraph{Encoding} Format \textsf{RV64A\_RMW} (Section~\ref{sec:format-RV64A-RMW}), \verb|funct5 = 00100|, \verb|funct3 = 000|\\

\bitfields{ 5/00100, 2/acqrel, 5/rs2, 5/rs1, 3/000, 5/rd, 7/0101111 }


\paragraph{Syntax} \texttt{amoxor.b}\emph{acqrel} \emph{rd}, \emph{rs2}, [\emph{rs1}]

\paragraph{Description} This instruction implements atomic load-XOR-store on byte. The effective address is in \emph{rs1}.


\paragraph{Modifier(s)}
\noindent\begin{itemize}
\item acqrel (cf. acqrel in section \ref{par:modifier-acqrel} on page \pageref{par:modifier-acqrel})
\end{itemize}

\paragraph{Execution} {Reservation table \textsf{LSU2\_MEMW\_AUXR\_AUXW} (Section~\ref{sec:reservation-LSU2-MEMW-AUXR-AUXW})}
~
{\begin{lstlisting}
stage ID:
  new acqrel = _ZX_1(acqrel);
stage RR:
  new acq = 1 & (acqrel >> 1);
  new rel = 1 & acqrel;
  new rs2 = BIR[rs2];
  new rs1 = BIR[rs1];
  new rd = _SX_8(MEM.atomic_eor(rs1, 0x1, acq | (rel << 16), rs2.32[0], rs1));
stage E1:
  if (rd != 0) {
    BIR[rd] = rd;
  }
\end{lstlisting}}
\newpage
\subsubsection[{\small AMOXOR.D: Atomic Memory Operation XOR Double Word}]{AMOXOR.D\hfill Atomic Memory Operation XOR Double Word}

\paragraph{Encoding} Format \textsf{RV64A\_RMW} (Section~\ref{sec:format-RV64A-RMW}), \verb|funct5 = 00100|, \verb|funct3 = 011|\\

\bitfields{ 5/00100, 2/acqrel, 5/rs2, 5/rs1, 3/011, 5/rd, 7/0101111 }


\paragraph{Syntax} \texttt{amoxor.d}\emph{acqrel} \emph{rd}, \emph{rs2}, [\emph{rs1}]

\paragraph{Description} This instruction implements atomic load-XOR-store on double word. The effective address in \emph{rs1} must be double word aligned (multiple of 8).


\paragraph{Modifier(s)}
\noindent\begin{itemize}
\item acqrel (cf. acqrel in section \ref{par:modifier-acqrel} on page \pageref{par:modifier-acqrel})
\end{itemize}

\paragraph{Execution} {Reservation table \textsf{LSU2\_MEMW\_AUXR\_AUXW} (Section~\ref{sec:reservation-LSU2-MEMW-AUXR-AUXW})}
~
{\begin{lstlisting}
stage ID:
  new acqrel = _ZX_1(acqrel);
stage RR:
  new acq = 1 & (acqrel >> 1);
  new rel = 1 & acqrel;
  new rs2 = BIR[rs2];
  new rs1 = BIR[rs1];
  new rd = MEM.atomic_eor(rs1, 0xFF, acq | (rel << 16), rs2.64[0], rs1);
stage E1:
  if (rd != 0) {
    BIR[rd] = rd;
  }
\end{lstlisting}}
\newpage
\subsubsection[{\small AMOXOR.H: Atomic Memory Operation XOR half Word}]{AMOXOR.H\hfill Atomic Memory Operation XOR half Word}

\paragraph{Encoding} Format \textsf{RV64A\_RMW} (Section~\ref{sec:format-RV64A-RMW}), \verb|funct5 = 00100|, \verb|funct3 = 001|\\

\bitfields{ 5/00100, 2/acqrel, 5/rs2, 5/rs1, 3/001, 5/rd, 7/0101111 }


\paragraph{Syntax} \texttt{amoxor.h}\emph{acqrel} \emph{rd}, \emph{rs2}, [\emph{rs1}]

\paragraph{Description} This instruction implements atomic load-XOR-store on half word. The effective address in \emph{rs1} must be word aligned (multiple of 2).


\paragraph{Modifier(s)}
\noindent\begin{itemize}
\item acqrel (cf. acqrel in section \ref{par:modifier-acqrel} on page \pageref{par:modifier-acqrel})
\end{itemize}

\paragraph{Execution} {Reservation table \textsf{LSU2\_MEMW\_AUXR\_AUXW} (Section~\ref{sec:reservation-LSU2-MEMW-AUXR-AUXW})}
~
{\begin{lstlisting}
stage ID:
  new acqrel = _ZX_1(acqrel);
stage RR:
  new acq = 1 & (acqrel >> 1);
  new rel = 1 & acqrel;
  new rs2 = BIR[rs2];
  new rs1 = BIR[rs1];
  new rd = _SX_16(MEM.atomic_eor(rs1, 0x3, acq | (rel << 16), rs2.32[0], rs1));
stage E1:
  if (rd != 0) {
    BIR[rd] = rd;
  }
\end{lstlisting}}
\newpage
\subsubsection[{\small AMOXOR.W: Atomic Memory Operation XOR Word}]{AMOXOR.W\hfill Atomic Memory Operation XOR Word}

\paragraph{Encoding} Format \textsf{RV64A\_RMW} (Section~\ref{sec:format-RV64A-RMW}), \verb|funct5 = 00100|, \verb|funct3 = 010|\\

\bitfields{ 5/00100, 2/acqrel, 5/rs2, 5/rs1, 3/010, 5/rd, 7/0101111 }


\paragraph{Syntax} \texttt{amoxor.w}\emph{acqrel} \emph{rd}, \emph{rs2}, [\emph{rs1}]

\paragraph{Description} This instruction implements atomic load-XOR-store on word. The effective address in \emph{rs1} must be word aligned (multiple of 4).


\paragraph{Modifier(s)}
\noindent\begin{itemize}
\item acqrel (cf. acqrel in section \ref{par:modifier-acqrel} on page \pageref{par:modifier-acqrel})
\end{itemize}

\paragraph{Execution} {Reservation table \textsf{LSU2\_MEMW\_AUXR\_AUXW} (Section~\ref{sec:reservation-LSU2-MEMW-AUXR-AUXW})}
~
{\begin{lstlisting}
stage ID:
  new acqrel = _ZX_1(acqrel);
stage RR:
  new acq = 1 & (acqrel >> 1);
  new rel = 1 & acqrel;
  new rs2 = BIR[rs2];
  new rs1 = BIR[rs1];
  new rd = _SX_32(MEM.atomic_eor(rs1, 0xF, acq | (rel << 16), rs2.32[0], rs1));
stage E1:
  if (rd != 0) {
    BIR[rd] = rd;
  }
\end{lstlisting}}
\newpage
\subsubsection[{\small LR.D: Load-Reserved Double Word}]{LR.D\hfill Load-Reserved Double Word}

\paragraph{Encoding} Format \textsf{RV64A\_LR} (Section~\ref{sec:format-RV64A-LR}), \verb|funct3 = 011|\\

\bitfields{ 5/00010, 2/acqrel, 5/00000, 5/rs1, 3/011, 5/rd, 7/0101111 }


\paragraph{Syntax} \texttt{lr.d}\emph{acqrel} \emph{rd}, [\emph{rs1}]

\paragraph{Description} The \emph{rd} receives the sign-extended memory double word at address \emph{rs1}. A reservation is registered on the memory address. The effective address in \emph{rs1} must be double word aligned (multiple of 8).


\paragraph{Modifier(s)}
\noindent\begin{itemize}
\item acqrel (cf. acqrel in section \ref{par:modifier-acqrel} on page \pageref{par:modifier-acqrel})
\end{itemize}

\paragraph{Execution} {Reservation table \textsf{LSU2\_MEMW\_AUXR\_AUXW} (Section~\ref{sec:reservation-LSU2-MEMW-AUXR-AUXW})}
~
{\begin{lstlisting}
stage ID:
  new acqrel = _ZX_1(acqrel);
stage RR:
  new acq = 1 & (acqrel >> 1);
  new rel = 1 & acqrel;
  new rs1 = BIR[rs1];
  new rd = BIR[rd];
  rd = _SX_64(MEM.load_reserve(rs1, 0xFF, acq | (rel << 16), rs1));
stage E1:
  if (rd != 0) {
    BIR[rd] = rd;
  }
\end{lstlisting}}
\newpage
\subsubsection[{\small LR.W: Load-Reserved Word}]{LR.W\hfill Load-Reserved Word}

\paragraph{Encoding} Format \textsf{RV64A\_LR} (Section~\ref{sec:format-RV64A-LR}), \verb|funct3 = 010|\\

\bitfields{ 5/00010, 2/acqrel, 5/00000, 5/rs1, 3/010, 5/rd, 7/0101111 }


\paragraph{Syntax} \texttt{lr.w}\emph{acqrel} \emph{rd}, [\emph{rs1}]

\paragraph{Description} The \emph{rd} receives the sign-extended memory word at address \emph{rs1}. A reservation is registered on the memory address. The effective address in \emph{rs1} must be word aligned (multiple of 4).


\paragraph{Modifier(s)}
\noindent\begin{itemize}
\item acqrel (cf. acqrel in section \ref{par:modifier-acqrel} on page \pageref{par:modifier-acqrel})
\end{itemize}

\paragraph{Execution} {Reservation table \textsf{LSU2\_MEMW\_AUXR\_AUXW} (Section~\ref{sec:reservation-LSU2-MEMW-AUXR-AUXW})}
~
{\begin{lstlisting}
stage ID:
  new acqrel = _ZX_1(acqrel);
stage RR:
  new acq = 1 & (acqrel >> 1);
  new rel = 1 & acqrel;
  new rs1 = BIR[rs1];
  new rd = BIR[rd];
  rd = _SX_32(MEM.load_reserve(rs1, 0xF, acq | (rel << 16), rs1));
stage E1:
  if (rd != 0) {
    BIR[rd] = rd;
  }
\end{lstlisting}}
\newpage
\subsubsection[{\small SC.D: Store-Conditional Double Word}]{SC.D\hfill Store-Conditional Double Word}

\paragraph{Encoding} Format \textsf{RV64A\_SC} (Section~\ref{sec:format-RV64A-SC}), \verb|funct3 = 011|\\

\bitfields{ 5/00011, 2/acqrel, 5/rs2, 5/rs1, 3/011, 5/rd, 7/0101111 }


\paragraph{Syntax} \texttt{sc.d}\emph{acqrel} \emph{rd}, \emph{rs2}, [\emph{rs1}]

\paragraph{Description} The \emph{rs2} is written to the memory double word at address \emph{rs1} if the reservation is still valid. If the store succeeds, \emph{rd} is set to 0, else to 1. The effective address in \emph{rs1} must be double word aligned (multiple of 8).


\paragraph{Modifier(s)}
\noindent\begin{itemize}
\item acqrel (cf. acqrel in section \ref{par:modifier-acqrel} on page \pageref{par:modifier-acqrel})
\end{itemize}

\paragraph{Execution} {Reservation table \textsf{LSU2\_MEMW\_AUXR\_AUXW} (Section~\ref{sec:reservation-LSU2-MEMW-AUXR-AUXW})}
~
{\begin{lstlisting}
stage ID:
  new acqrel = _ZX_1(acqrel);
stage RR:
  new acq = 1 & (acqrel >> 1);
  new rel = 1 & acqrel;
  new rs1 = BIR[rs1];
  new rs2 = BIR[rs2];
  new rd = BIR[rd];
  rd = _SX_64(MEM.store_conditional(rs1, 0xFF, acq | (rel << 16), rs2, rs1));
stage E1:
  if (rd != 0) {
    BIR[rd] = rd;
  }
\end{lstlisting}}
\newpage
\subsubsection[{\small SC.W: Store-Conditional Word}]{SC.W\hfill Store-Conditional Word}

\paragraph{Encoding} Format \textsf{RV64A\_SC} (Section~\ref{sec:format-RV64A-SC}), \verb|funct3 = 010|\\

\bitfields{ 5/00011, 2/acqrel, 5/rs2, 5/rs1, 3/010, 5/rd, 7/0101111 }


\paragraph{Syntax} \texttt{sc.w}\emph{acqrel} \emph{rd}, \emph{rs2}, [\emph{rs1}]

\paragraph{Description} The \emph{rs2} is written to the memory word at address \emph{rs1} if the reservation is still valid. If the store succeeds, \emph{rd} is set to 0, else to 1. The effective address in \emph{rs2} must be word aligned (multiple of 4).


\paragraph{Modifier(s)}
\noindent\begin{itemize}
\item acqrel (cf. acqrel in section \ref{par:modifier-acqrel} on page \pageref{par:modifier-acqrel})
\end{itemize}

\paragraph{Execution} {Reservation table \textsf{LSU2\_MEMW\_AUXR\_AUXW} (Section~\ref{sec:reservation-LSU2-MEMW-AUXR-AUXW})}
~
{\begin{lstlisting}
stage ID:
  new acqrel = _ZX_1(acqrel);
stage RR:
  new acq = 1 & (acqrel >> 1);
  new rel = 1 & acqrel;
  new rs1 = BIR[rs1];
  new rs2 = BIR[rs2];
  new rd = BIR[rd];
  rd = _SX_32(MEM.store_conditional(rs1, 0xF, acq | (rel << 16), rs2, rs1));
stage E1:
  if (rd != 0) {
    BIR[rd] = rd;
  }
\end{lstlisting}}
\newpage
\subsection{RV64B Instructions}
\label{sec:RV64B-instructions}

\subsubsection[{\small ANDN: AND with Not operand}]{ANDN\hfill AND with Not operand}

\paragraph{Encoding} Format \textsf{RV64B\_OPREG\_0x20} (Section~\ref{sec:format-RV64B-OPREG-0x20}), \verb|funct3 = 111|\\

\bitfields{ 7/0100000, 5/rs2, 5/rs1, 3/111, 5/rd, 7/0110011 }


\paragraph{Syntax} \texttt{andn} \emph{rd}, \emph{rs1}, \emph{rs2}

\paragraph{Description} The \emph{rs1} is AND-ed with the bitwise NOT of the \emph{rs2}. The result is stored into the \emph{rd}.


\paragraph{Execution} {Reservation table \textsf{ALU\_TINY} (Section~\ref{sec:reservation-ALU-TINY})}
~
{\begin{lstlisting}
stage RR:
  new rs2 = BIR[rs2];
  new rs1 = BIR[rs1];
  new rd = rs1 & ~rs2;
stage E1:
  if (rd != 0) {
    BIR[rd] = rd;
  }
\end{lstlisting}}
\newpage
\subsubsection[{\small BCLR: Single Bit Clear}]{BCLR\hfill Single Bit Clear}

\paragraph{Encoding} Format \textsf{RV64B\_OPREG\_0x24} (Section~\ref{sec:format-RV64B-OPREG-0x24}), \verb|funct3 = 001|\\

\bitfields{ 7/0100100, 5/rs2, 5/rs1, 3/001, 5/rd, 7/0110011 }


\paragraph{Syntax} \texttt{bclr} \emph{rd}, \emph{rs1}, \emph{rs2}

\paragraph{Description} The single bit of the \emph{rs1} specified by the \emph{rs2} is cleared. The result is stored into the \emph{rd}.


\paragraph{Execution} {Reservation table \textsf{ALU\_TINY} (Section~\ref{sec:reservation-ALU-TINY})}
~
{\begin{lstlisting}
stage RR:
  new rs2 = BIR[rs2];
  new rs1 = BIR[rs1];
  new index = rs2 & 63;
  new rd = rs1 & ~(0x1 << index);
stage E1:
  if (rd != 0) {
    BIR[rd] = rd;
  }
\end{lstlisting}}
\newpage
\subsubsection[{\small BCLRI: Single Bit Clear Immediate}]{BCLRI\hfill Single Bit Clear Immediate}

\paragraph{Encoding} Format \textsf{RV64B\_OPIMM\_0x12} (Section~\ref{sec:format-RV64B-OPIMM-0x12}), \verb|funct3 = 001|\\

\bitfields{ 6/010010, 6/i-- 5-- 0, 5/rs1, 3/001, 5/rd, 7/0010011 }


\paragraph{Syntax} \texttt{bclri} \emph{rd}, \emph{rs1}, \emph{i\_\discretionary{}{}{}5\_\discretionary{}{}{}0}

\paragraph{Description} The single bit of the \emph{rs1} specified by the \emph{i\_\discretionary{}{}{}5\_\discretionary{}{}{}0} is cleared. The result is stored into the \emph{rd}.


\paragraph{Execution} {Reservation table \textsf{ALU\_TINY} (Section~\ref{sec:reservation-ALU-TINY})}
~
{\begin{lstlisting}
stage RR:
  new unsigned6 = _ZX_6(i_5_0);
  new rs1 = BIR[rs1];
  new index = unsigned6 & 63;
  new rd = rs1 & ~(0x1 << index);
stage E1:
  if (rd != 0) {
    BIR[rd] = rd;
  }
\end{lstlisting}}
\newpage
\subsubsection[{\small BEXT: Single Bit Extract}]{BEXT\hfill Single Bit Extract}

\paragraph{Encoding} Format \textsf{RV64B\_OPREG\_0x24} (Section~\ref{sec:format-RV64B-OPREG-0x24}), \verb|funct3 = 101|\\

\bitfields{ 7/0100100, 5/rs2, 5/rs1, 3/101, 5/rd, 7/0110011 }


\paragraph{Syntax} \texttt{bext} \emph{rd}, \emph{rs1}, \emph{rs2}

\paragraph{Description} The single bit of the \emph{rs1} specified by the \emph{rs2} is extracted. The result is stored into the \emph{rd}.


\paragraph{Execution} {Reservation table \textsf{ALU\_TINY} (Section~\ref{sec:reservation-ALU-TINY})}
~
{\begin{lstlisting}
stage RR:
  new rs2 = BIR[rs2];
  new rs1 = BIR[rs1];
  new index = rs2 & 63;
  new rd = (rs1 >> index) & 0x1;
stage E1:
  if (rd != 0) {
    BIR[rd] = rd;
  }
\end{lstlisting}}
\newpage
\subsubsection[{\small BEXTI: Single Bit Extract Immediate}]{BEXTI\hfill Single Bit Extract Immediate}

\paragraph{Encoding} Format \textsf{RV64B\_OPIMM\_0x12} (Section~\ref{sec:format-RV64B-OPIMM-0x12}), \verb|funct3 = 101|\\

\bitfields{ 6/010010, 6/i-- 5-- 0, 5/rs1, 3/101, 5/rd, 7/0010011 }


\paragraph{Syntax} \texttt{bexti} \emph{rd}, \emph{rs1}, \emph{i\_\discretionary{}{}{}5\_\discretionary{}{}{}0}

\paragraph{Description} The single bit of the \emph{rs1} specified by the \emph{i\_\discretionary{}{}{}5\_\discretionary{}{}{}0} is extracted. The result is stored into the \emph{rd}.


\paragraph{Execution} {Reservation table \textsf{ALU\_TINY} (Section~\ref{sec:reservation-ALU-TINY})}
~
{\begin{lstlisting}
stage RR:
  new unsigned6 = _ZX_6(i_5_0);
  new rs1 = BIR[rs1];
  new index = unsigned6 & 63;
  new rd = (rs1 >> index) & 0x1;
stage E1:
  if (rd != 0) {
    BIR[rd] = rd;
  }
\end{lstlisting}}
\newpage
\subsubsection[{\small BINV: Single Bit Invert}]{BINV\hfill Single Bit Invert}

\paragraph{Encoding} Format \textsf{RV64B\_OPREG\_0x34} (Section~\ref{sec:format-RV64B-OPREG-0x34}), \verb|funct3 = 001|\\

\bitfields{ 7/0110100, 5/rs2, 5/rs1, 3/001, 5/rd, 7/0110011 }


\paragraph{Syntax} \texttt{binv} \emph{rd}, \emph{rs1}, \emph{rs2}

\paragraph{Description} The single bit of the \emph{rs1} specified by the \emph{rs2} is inverted. The result is stored into the \emph{rd}.


\paragraph{Execution} {Reservation table \textsf{ALU\_TINY} (Section~\ref{sec:reservation-ALU-TINY})}
~
{\begin{lstlisting}
stage RR:
  new rs2 = BIR[rs2];
  new rs1 = BIR[rs1];
  new index = rs2 & 63;
  new rd = rs1 ^ (0x1 << index);
stage E1:
  if (rd != 0) {
    BIR[rd] = rd;
  }
\end{lstlisting}}
\newpage
\subsubsection[{\small BINVI: Single Bit Invert Immediate}]{BINVI\hfill Single Bit Invert Immediate}

\paragraph{Encoding} Format \textsf{RV64B\_OPIMM\_0x1a} (Section~\ref{sec:format-RV64B-OPIMM-0x1a}), \verb|funct3 = 001|\\

\bitfields{ 6/011010, 6/i-- 5-- 0, 5/rs1, 3/001, 5/rd, 7/0010011 }


\paragraph{Syntax} \texttt{binvi} \emph{rd}, \emph{rs1}, \emph{i\_\discretionary{}{}{}5\_\discretionary{}{}{}0}

\paragraph{Description} The single bit of the \emph{rs1} specified by the \emph{i\_\discretionary{}{}{}5\_\discretionary{}{}{}0} is inverted. The result is stored into the \emph{rd}.


\paragraph{Execution} {Reservation table \textsf{ALU\_TINY} (Section~\ref{sec:reservation-ALU-TINY})}
~
{\begin{lstlisting}
stage RR:
  new unsigned6 = _ZX_6(i_5_0);
  new rs1 = BIR[rs1];
  new index = unsigned6 & 63;
  new rd = rs1 ^ (0x1 << index);
stage E1:
  if (rd != 0) {
    BIR[rd] = rd;
  }
\end{lstlisting}}
\newpage
\subsubsection[{\small BSET: Single Bit Set}]{BSET\hfill Single Bit Set}

\paragraph{Encoding} Format \textsf{RV64B\_OPREG\_0x14} (Section~\ref{sec:format-RV64B-OPREG-0x14}), \verb|funct3 = 001|\\

\bitfields{ 7/0010100, 5/rs2, 5/rs1, 3/001, 5/rd, 7/0110011 }


\paragraph{Syntax} \texttt{bset} \emph{rd}, \emph{rs1}, \emph{rs2}

\paragraph{Description} The single bit of the \emph{rs1} specified by the \emph{rs2} is set. The result is stored into the \emph{rd}.


\paragraph{Execution} {Reservation table \textsf{ALU\_TINY} (Section~\ref{sec:reservation-ALU-TINY})}
~
{\begin{lstlisting}
stage RR:
  new rs2 = BIR[rs2];
  new rs1 = BIR[rs1];
  new index = rs2 & 63;
  new rd = rs1 | (0x1 << index);
stage E1:
  if (rd != 0) {
    BIR[rd] = rd;
  }
\end{lstlisting}}
\newpage
\subsubsection[{\small BSETI: Single Bit Set Immediate}]{BSETI\hfill Single Bit Set Immediate}

\paragraph{Encoding} Format \textsf{RV64B\_OPIMM\_0x0a} (Section~\ref{sec:format-RV64B-OPIMM-0x0a}), \verb|funct3 = 001|\\

\bitfields{ 6/001010, 6/i-- 5-- 0, 5/rs1, 3/001, 5/rd, 7/0010011 }


\paragraph{Syntax} \texttt{bseti} \emph{rd}, \emph{rs1}, \emph{i\_\discretionary{}{}{}5\_\discretionary{}{}{}0}

\paragraph{Description} The single bit of the \emph{rs1} specified by the \emph{i\_\discretionary{}{}{}5\_\discretionary{}{}{}0} is set. The result is stored into the \emph{rd}.


\paragraph{Execution} {Reservation table \textsf{ALU\_TINY} (Section~\ref{sec:reservation-ALU-TINY})}
~
{\begin{lstlisting}
stage RR:
  new unsigned6 = _ZX_6(i_5_0);
  new rs1 = BIR[rs1];
  new index = unsigned6 & 63;
  new rd = rs1 | (0x1 << index);
stage E1:
  if (rd != 0) {
    BIR[rd] = rd;
  }
\end{lstlisting}}
\newpage
\subsubsection[{\small CLMUL: Carryless Multiply}]{CLMUL\hfill Carryless Multiply}

\paragraph{Encoding} Format \textsf{RV64B\_OPREG\_0x05} (Section~\ref{sec:format-RV64B-OPREG-0x05}), \verb|funct3 = 001|\\

\bitfields{ 7/0000101, 5/rs2, 5/rs1, 3/001, 5/rd, 7/0110011 }


\paragraph{Syntax} \texttt{clmul} \emph{rd}, \emph{rs1}, \emph{rs2}

\paragraph{Description} The carryless multiply of the \emph{rs1} and the \emph{rs2} is computed. The lower 64 bits of the result are stored into the \emph{rd}.


\paragraph{Execution} {Reservation table \textsf{ALU\_LITE} (Section~\ref{sec:reservation-ALU-LITE})}
~
{\begin{lstlisting}
stage RR:
  new rs2 = BIR[rs2];
  new rs1 = BIR[rs1];
  new rd = clm_64_128(rs1, rs2);
stage E1:
  if (rd != 0) {
    BIR[rd] = rd;
  }
\end{lstlisting}}
\newpage
\subsubsection[{\small CLMULH: Carryless Multiply High}]{CLMULH\hfill Carryless Multiply High}

\paragraph{Encoding} Format \textsf{RV64B\_OPREG\_0x05} (Section~\ref{sec:format-RV64B-OPREG-0x05}), \verb|funct3 = 011|\\

\bitfields{ 7/0000101, 5/rs2, 5/rs1, 3/011, 5/rd, 7/0110011 }


\paragraph{Syntax} \texttt{clmulh} \emph{rd}, \emph{rs1}, \emph{rs2}

\paragraph{Description} The carryless multiply of the \emph{rs1} and the \emph{rs2} is computed. The upper 64 bits of the result are stored into the \emph{rd}.


\paragraph{Execution} {Reservation table \textsf{ALU\_LITE} (Section~\ref{sec:reservation-ALU-LITE})}
~
{\begin{lstlisting}
stage RR:
  new rs2 = BIR[rs2];
  new rs1 = BIR[rs1];
  new rd = clm_64_128(rs1, rs2) >> 64;
stage E1:
  if (rd != 0) {
    BIR[rd] = rd;
  }
\end{lstlisting}}
\newpage
\subsubsection[{\small CLMULR: Carryless Multiply Reversed}]{CLMULR\hfill Carryless Multiply Reversed}

\paragraph{Encoding} Format \textsf{RV64B\_OPREG\_0x05} (Section~\ref{sec:format-RV64B-OPREG-0x05}), \verb|funct3 = 010|\\

\bitfields{ 7/0000101, 5/rs2, 5/rs1, 3/010, 5/rd, 7/0110011 }


\paragraph{Syntax} \texttt{clmulr} \emph{rd}, \emph{rs1}, \emph{rs2}

\paragraph{Description} The bit-reversed carryless multiply of the \emph{rs1} and the \emph{rs2} is computed. The lower 64 bits of the result are stored into the \emph{rd}.


\paragraph{Execution} {Reservation table \textsf{ALU\_LITE} (Section~\ref{sec:reservation-ALU-LITE})}
~
{\begin{lstlisting}
stage RR:
  new rs2 = BIR[rs2];
  new rs1 = BIR[rs1];
  new rd = brev_64(clm_64_128(brev_64(rs1), brev_64(rs2)));
stage E1:
  if (rd != 0) {
    BIR[rd] = rd;
  }
\end{lstlisting}}
\newpage
\subsubsection[{\small CLZ: Count Leading Zero Bits}]{CLZ\hfill Count Leading Zero Bits}

\paragraph{Encoding} Format \textsf{RV64B\_OPCNT\_0x18} (Section~\ref{sec:format-RV64B-OPCNT-0x18}), \verb|i_5_0 = 000000|\\

\bitfields{ 6/011000, 6/000000, 5/rs1, 3/001, 5/rd, 7/0010011 }


\paragraph{Syntax} \texttt{clz} \emph{rd}, \emph{rs1}

\paragraph{Description} Counts the number of 0s before the first 1 of the \emph{rs1}, starting at the most significant bit. The result is stored into the \emph{rd}.


\paragraph{Execution} {Reservation table \textsf{ALU\_LITE} (Section~\ref{sec:reservation-ALU-LITE})}
~
{\begin{lstlisting}
stage RR:
  new rs1 = BIR[rs1];
  new rd = _CLZ_64(rs1);
stage E1:
  if (rd != 0) {
    BIR[rd] = rd;
  }
\end{lstlisting}}
\newpage
\subsubsection[{\small CPOP: Count Set Bits}]{CPOP\hfill Count Set Bits}

\paragraph{Encoding} Format \textsf{RV64B\_OPCNT\_0x18} (Section~\ref{sec:format-RV64B-OPCNT-0x18}), \verb|i_5_0 = 000010|\\

\bitfields{ 6/011000, 6/000010, 5/rs1, 3/001, 5/rd, 7/0010011 }


\paragraph{Syntax} \texttt{cpop} \emph{rd}, \emph{rs1}

\paragraph{Description} Counts the number of 1s of the \emph{rs1}. The result is stored into the \emph{rd}.


\paragraph{Execution} {Reservation table \textsf{ALU\_LITE} (Section~\ref{sec:reservation-ALU-LITE})}
~
{\begin{lstlisting}
stage RR:
  new rs1 = BIR[rs1];
  new rd = _CBS_64(rs1);
stage E1:
  if (rd != 0) {
    BIR[rd] = rd;
  }
\end{lstlisting}}
\newpage
\subsubsection[{\small CTZ: Count Trailing Zero Bits}]{CTZ\hfill Count Trailing Zero Bits}

\paragraph{Encoding} Format \textsf{RV64B\_OPCNT\_0x18} (Section~\ref{sec:format-RV64B-OPCNT-0x18}), \verb|i_5_0 = 000001|\\

\bitfields{ 6/011000, 6/000001, 5/rs1, 3/001, 5/rd, 7/0010011 }


\paragraph{Syntax} \texttt{ctz} \emph{rd}, \emph{rs1}

\paragraph{Description} Counts the number of 0s before the first 1 of the \emph{rs1}, starting at the least significant bit. The result is stored into the \emph{rd}.


\paragraph{Execution} {Reservation table \textsf{ALU\_LITE} (Section~\ref{sec:reservation-ALU-LITE})}
~
{\begin{lstlisting}
stage RR:
  new rs1 = BIR[rs1];
  new rd = _CTZ_64(rs1);
stage E1:
  if (rd != 0) {
    BIR[rd] = rd;
  }
\end{lstlisting}}
\newpage
\subsubsection[{\small MAX: Maximum of Double Words}]{MAX\hfill Maximum of Double Words}

\paragraph{Encoding} Format \textsf{RV64B\_OPREG\_0x05} (Section~\ref{sec:format-RV64B-OPREG-0x05}), \verb|funct3 = 110|\\

\bitfields{ 7/0000101, 5/rs2, 5/rs1, 3/110, 5/rd, 7/0110011 }


\paragraph{Syntax} \texttt{max} \emph{rd}, \emph{rs1}, \emph{rs2}

\paragraph{Description} The maximum of the \emph{rs1} and the \emph{rs2} is computed. The result is stored into the \emph{rd}.


\paragraph{Execution} {Reservation table \textsf{ALU\_TINY} (Section~\ref{sec:reservation-ALU-TINY})}
~
{\begin{lstlisting}
stage RR:
  new rs2 = BIR[rs2];
  new rs1 = BIR[rs1];
  new rd = _MAX(_SX_64(rs1), _SX_64(rs2));
stage E1:
  if (rd != 0) {
    BIR[rd] = rd;
  }
\end{lstlisting}}
\newpage
\subsubsection[{\small MAXU: Maximum of Double Words Unsigned}]{MAXU\hfill Maximum of Double Words Unsigned}

\paragraph{Encoding} Format \textsf{RV64B\_OPREG\_0x05} (Section~\ref{sec:format-RV64B-OPREG-0x05}), \verb|funct3 = 111|\\

\bitfields{ 7/0000101, 5/rs2, 5/rs1, 3/111, 5/rd, 7/0110011 }


\paragraph{Syntax} \texttt{maxu} \emph{rd}, \emph{rs1}, \emph{rs2}

\paragraph{Description} The unsigned maximum of the \emph{rs1} and the \emph{rs2} is computed. The result is stored into the \emph{rd}.


\paragraph{Execution} {Reservation table \textsf{ALU\_TINY} (Section~\ref{sec:reservation-ALU-TINY})}
~
{\begin{lstlisting}
stage RR:
  new rs2 = BIR[rs2];
  new rs1 = BIR[rs1];
  new rd = _MAX(_ZX_64(rs1), _ZX_64(rs2));
stage E1:
  if (rd != 0) {
    BIR[rd] = rd;
  }
\end{lstlisting}}
\newpage
\subsubsection[{\small MIN: Minimum of Double Words}]{MIN\hfill Minimum of Double Words}

\paragraph{Encoding} Format \textsf{RV64B\_OPREG\_0x05} (Section~\ref{sec:format-RV64B-OPREG-0x05}), \verb|funct3 = 100|\\

\bitfields{ 7/0000101, 5/rs2, 5/rs1, 3/100, 5/rd, 7/0110011 }


\paragraph{Syntax} \texttt{min} \emph{rd}, \emph{rs1}, \emph{rs2}

\paragraph{Description} The minimum of the \emph{rs1} and the \emph{rs2} is computed. The result is stored into the \emph{rd}.


\paragraph{Execution} {Reservation table \textsf{ALU\_TINY} (Section~\ref{sec:reservation-ALU-TINY})}
~
{\begin{lstlisting}
stage RR:
  new rs2 = BIR[rs2];
  new rs1 = BIR[rs1];
  new rd = _MIN(_SX_64(rs1), _SX_64(rs2));
stage E1:
  if (rd != 0) {
    BIR[rd] = rd;
  }
\end{lstlisting}}
\newpage
\subsubsection[{\small MINU: Minimum of Double Words Unsigned}]{MINU\hfill Minimum of Double Words Unsigned}

\paragraph{Encoding} Format \textsf{RV64B\_OPREG\_0x05} (Section~\ref{sec:format-RV64B-OPREG-0x05}), \verb|funct3 = 101|\\

\bitfields{ 7/0000101, 5/rs2, 5/rs1, 3/101, 5/rd, 7/0110011 }


\paragraph{Syntax} \texttt{minu} \emph{rd}, \emph{rs1}, \emph{rs2}

\paragraph{Description} The unsigned minimum of the \emph{rs1} and the \emph{rs2} is computed. The result is stored into the \emph{rd}.


\paragraph{Execution} {Reservation table \textsf{ALU\_TINY} (Section~\ref{sec:reservation-ALU-TINY})}
~
{\begin{lstlisting}
stage RR:
  new rs2 = BIR[rs2];
  new rs1 = BIR[rs1];
  new rd = _MIN(_ZX_64(rs1), _ZX_64(rs2));
stage E1:
  if (rd != 0) {
    BIR[rd] = rd;
  }
\end{lstlisting}}
\newpage
\subsubsection[{\small ORC.B: Bitwise OR-Combine, Byte Granule}]{ORC.B\hfill Bitwise OR-Combine, Byte Granule}

\paragraph{Encoding} Format \textsf{RV64B\_OPIMM\_0x287} (Section~\ref{sec:format-RV64B-OPIMM-0x287}), \verb|funct3 = 101|\\

\bitfields{ 12/001010000111, 5/rs1, 3/101, 5/rd, 7/0010011 }


\paragraph{Syntax} \texttt{orc.b} \emph{rd}, \emph{rs1}

\paragraph{Description} Combines the bits within each byte using bitwise logical OR.

\paragraph{Execution} {Reservation table \textsf{ALU\_TINY} (Section~\ref{sec:reservation-ALU-TINY})}
~
{\begin{lstlisting}
stage RR:
  new rs1 = BIR[rs1];
  for (I = 0; I < 8; I += 1) {
    rd.8[I] = -rs1.8[I]
  };
stage E1:
  if (rd != 0) {
    BIR[rd] = rd;
  }
\end{lstlisting}}
\newpage
\subsubsection[{\small ORN: OR with Not operand}]{ORN\hfill OR with Not operand}

\paragraph{Encoding} Format \textsf{RV64B\_OPREG\_0x20} (Section~\ref{sec:format-RV64B-OPREG-0x20}), \verb|funct3 = 110|\\

\bitfields{ 7/0100000, 5/rs2, 5/rs1, 3/110, 5/rd, 7/0110011 }


\paragraph{Syntax} \texttt{orn} \emph{rd}, \emph{rs1}, \emph{rs2}

\paragraph{Description} The \emph{rs1} is OR-ed with the bitwise NOT of the \emph{rs2}. The result is stored into the \emph{rd}.


\paragraph{Execution} {Reservation table \textsf{ALU\_TINY} (Section~\ref{sec:reservation-ALU-TINY})}
~
{\begin{lstlisting}
stage RR:
  new rs2 = BIR[rs2];
  new rs1 = BIR[rs1];
  new rd = rs1 | ~rs2;
stage E1:
  if (rd != 0) {
    BIR[rd] = rd;
  }
\end{lstlisting}}
\newpage
\subsubsection[{\small REV8: Reverse Bytes}]{REV8\hfill Reverse Bytes}

\paragraph{Encoding} Format \textsf{RV64B\_OPIMM\_0x6b8} (Section~\ref{sec:format-RV64B-OPIMM-0x6b8}), \verb|funct3 = 101|\\

\bitfields{ 12/011010111000, 5/rs1, 3/101, 5/rd, 7/0010011 }


\paragraph{Syntax} \texttt{rev8} \emph{rd}, \emph{rs1}

\paragraph{Description} Reverse the order of bytes.


\paragraph{Execution} {Reservation table \textsf{ALU\_TINY} (Section~\ref{sec:reservation-ALU-TINY})}
~
{\begin{lstlisting}
stage RR:
  new rs1 = BIR[rs1];
  new rd = _SWAP_32(_SWAP_16(_SWAP_8(rs1)));
stage E1:
  if (rd != 0) {
    BIR[rd] = rd;
  }
\end{lstlisting}}
\newpage
\subsubsection[{\small ROL: Rotate Left}]{ROL\hfill Rotate Left}

\paragraph{Encoding} Format \textsf{RV64B\_OPREG\_0x30} (Section~\ref{sec:format-RV64B-OPREG-0x30}), \verb|funct3 = 001|\\

\bitfields{ 7/0110000, 5/rs2, 5/rs1, 3/001, 5/rd, 7/0110011 }


\paragraph{Syntax} \texttt{rol} \emph{rd}, \emph{rs1}, \emph{rs2}

\paragraph{Description} The \emph{rs1} is rotated left by the \emph{rs2} modulo 64. The result is stored into the \emph{rd}.


\paragraph{Execution} {Reservation table \textsf{ALU\_LITE} (Section~\ref{sec:reservation-ALU-LITE})}
~
{\begin{lstlisting}
stage RR:
  new rs2 = BIR[rs2];
  new rs1 = BIR[rs1];
  new shift = _ZX_6(rs2);
  new rd = _ROL_64(rs1, shift);
stage E1:
  if (rd != 0) {
    BIR[rd] = rd;
  }
\end{lstlisting}}
\newpage
\subsubsection[{\small ROR: Rotate Right}]{ROR\hfill Rotate Right}

\paragraph{Encoding} Format \textsf{RV64B\_OPREG\_0x30} (Section~\ref{sec:format-RV64B-OPREG-0x30}), \verb|funct3 = 101|\\

\bitfields{ 7/0110000, 5/rs2, 5/rs1, 3/101, 5/rd, 7/0110011 }


\paragraph{Syntax} \texttt{ror} \emph{rd}, \emph{rs1}, \emph{rs2}

\paragraph{Description} The \emph{rs1} is rotated right by the \emph{rs2} modulo 64. The result is stored into the \emph{rd}.


\paragraph{Execution} {Reservation table \textsf{ALU\_LITE} (Section~\ref{sec:reservation-ALU-LITE})}
~
{\begin{lstlisting}
stage RR:
  new rs2 = BIR[rs2];
  new rs1 = BIR[rs1];
  new shift = _ZX_6(rs2);
  new rd = _ROR_64(rs1, shift);
stage E1:
  if (rd != 0) {
    BIR[rd] = rd;
  }
\end{lstlisting}}
\newpage
\subsubsection[{\small RORI: Rotate Right Immediate}]{RORI\hfill Rotate Right Immediate}

\paragraph{Encoding} Format \textsf{RV64B\_OPIMM\_0x18} (Section~\ref{sec:format-RV64B-OPIMM-0x18}), \verb|funct3 = 101|\\

\bitfields{ 6/011000, 6/i-- 5-- 0, 5/rs1, 3/101, 5/rd, 7/0010011 }


\paragraph{Syntax} \texttt{rori} \emph{rd}, \emph{rs1}, \emph{i\_\discretionary{}{}{}5\_\discretionary{}{}{}0}

\paragraph{Description} The \emph{rs1} is rotated right by the \emph{i\_\discretionary{}{}{}5\_\discretionary{}{}{}0}. The result is stored into the \emph{rd}.


\paragraph{Execution} {Reservation table \textsf{ALU\_LITE} (Section~\ref{sec:reservation-ALU-LITE})}
~
{\begin{lstlisting}
stage RR:
  new unsigned6 = _ZX_6(i_5_0);
  new rs1 = BIR[rs1];
  new rd = _ROR_64(rs1, unsigned6);
stage E1:
  if (rd != 0) {
    BIR[rd] = rd;
  }
\end{lstlisting}}
\newpage
\subsubsection[{\small SEXT.B: Sign-Extend Byte}]{SEXT.B\hfill Sign-Extend Byte}

\paragraph{Encoding} Format \textsf{RV64B\_OPEXT\_0x30} (Section~\ref{sec:format-RV64B-OPEXT-0x30}), \verb|i_4_0 = 00100|\\

\bitfields{ 7/0110000, 5/00100, 5/rs1, 3/001, 5/rd, 7/0010011 }


\paragraph{Syntax} \texttt{sext.b} \emph{rd}, \emph{rs1}

\paragraph{Description} The \emph{rs1} byte is sign-extended. The result is stored into the \emph{rd}.


\paragraph{Execution} {Reservation table \textsf{ALU\_TINY} (Section~\ref{sec:reservation-ALU-TINY})}
~
{\begin{lstlisting}
stage RR:
  new rs1 = BIR[rs1];
  new rd = _SX_8(rs1);
stage E1:
  if (rd != 0) {
    BIR[rd] = rd;
  }
\end{lstlisting}}
\newpage
\subsubsection[{\small SEXT.H: Sign-Extend Half Word}]{SEXT.H\hfill Sign-Extend Half Word}

\paragraph{Encoding} Format \textsf{RV64B\_OPEXT\_0x30} (Section~\ref{sec:format-RV64B-OPEXT-0x30}), \verb|i_4_0 = 00101|\\

\bitfields{ 7/0110000, 5/00101, 5/rs1, 3/001, 5/rd, 7/0010011 }


\paragraph{Syntax} \texttt{sext.h} \emph{rd}, \emph{rs1}

\paragraph{Description} The \emph{rs1} half word is sign-extended. The result is stored into the \emph{rd}.


\paragraph{Execution} {Reservation table \textsf{ALU\_TINY} (Section~\ref{sec:reservation-ALU-TINY})}
~
{\begin{lstlisting}
stage RR:
  new rs1 = BIR[rs1];
  new rd = _SX_16(rs1);
stage E1:
  if (rd != 0) {
    BIR[rd] = rd;
  }
\end{lstlisting}}
\newpage
\subsubsection[{\small SH1ADD: Shift left by 1 and ADD}]{SH1ADD\hfill Shift left by 1 and ADD}

\paragraph{Encoding} Format \textsf{RV64B\_OPREG\_0x10} (Section~\ref{sec:format-RV64B-OPREG-0x10}), \verb|funct3 = 010|\\

\bitfields{ 7/0010000, 5/rs2, 5/rs1, 3/010, 5/rd, 7/0110011 }


\paragraph{Syntax} \texttt{sh1add} \emph{rd}, \emph{rs1}, \emph{rs2}

\paragraph{Description} The \emph{rs1} is shifted left by one bit and added to the \emph{rs2}. The result is stored into the \emph{rd}.


\paragraph{Execution} {Reservation table \textsf{ALU\_TINY} (Section~\ref{sec:reservation-ALU-TINY})}
~
{\begin{lstlisting}
stage RR:
  new rs2 = BIR[rs2];
  new rs1 = BIR[rs1];
  new rd = rs2 + (rs1 << 1);
stage E1:
  if (rd != 0) {
    BIR[rd] = rd;
  }
\end{lstlisting}}
\newpage
\subsubsection[{\small SH2ADD: Shift left by 2 and ADD}]{SH2ADD\hfill Shift left by 2 and ADD}

\paragraph{Encoding} Format \textsf{RV64B\_OPREG\_0x10} (Section~\ref{sec:format-RV64B-OPREG-0x10}), \verb|funct3 = 100|\\

\bitfields{ 7/0010000, 5/rs2, 5/rs1, 3/100, 5/rd, 7/0110011 }


\paragraph{Syntax} \texttt{sh2add} \emph{rd}, \emph{rs1}, \emph{rs2}

\paragraph{Description} The \emph{rs1} is shifted left by two bits and added to the \emph{rs2}. The result is stored into the \emph{rd}.


\paragraph{Execution} {Reservation table \textsf{ALU\_TINY} (Section~\ref{sec:reservation-ALU-TINY})}
~
{\begin{lstlisting}
stage RR:
  new rs2 = BIR[rs2];
  new rs1 = BIR[rs1];
  new rd = rs2 + (rs1 << 2);
stage E1:
  if (rd != 0) {
    BIR[rd] = rd;
  }
\end{lstlisting}}
\newpage
\subsubsection[{\small SH3ADD: Shift left by 3 and ADD}]{SH3ADD\hfill Shift left by 3 and ADD}

\paragraph{Encoding} Format \textsf{RV64B\_OPREG\_0x10} (Section~\ref{sec:format-RV64B-OPREG-0x10}), \verb|funct3 = 110|\\

\bitfields{ 7/0010000, 5/rs2, 5/rs1, 3/110, 5/rd, 7/0110011 }


\paragraph{Syntax} \texttt{sh3add} \emph{rd}, \emph{rs1}, \emph{rs2}

\paragraph{Description} The \emph{rs1} is shifted left by three bits and added to the \emph{rs2}. The result is stored into the \emph{rd}.


\paragraph{Execution} {Reservation table \textsf{ALU\_TINY} (Section~\ref{sec:reservation-ALU-TINY})}
~
{\begin{lstlisting}
stage RR:
  new rs2 = BIR[rs2];
  new rs1 = BIR[rs1];
  new rd = rs2 + (rs1 << 3);
stage E1:
  if (rd != 0) {
    BIR[rd] = rd;
  }
\end{lstlisting}}
\newpage
\subsubsection[{\small XNOR: Exclusive NOR}]{XNOR\hfill Exclusive NOR}

\paragraph{Encoding} Format \textsf{RV64B\_OPREG\_0x20} (Section~\ref{sec:format-RV64B-OPREG-0x20}), \verb|funct3 = 100|\\

\bitfields{ 7/0100000, 5/rs2, 5/rs1, 3/100, 5/rd, 7/0110011 }


\paragraph{Syntax} \texttt{xnor} \emph{rd}, \emph{rs1}, \emph{rs2}

\paragraph{Description} The \emph{rs1} is XOR-ed with the bitwise NOT of the \emph{rs2}. The result is stored into the \emph{rd}.


\paragraph{Execution} {Reservation table \textsf{ALU\_TINY} (Section~\ref{sec:reservation-ALU-TINY})}
~
{\begin{lstlisting}
stage RR:
  new rs2 = BIR[rs2];
  new rs1 = BIR[rs1];
  new rd = rs1 ^ ~rs2;
stage E1:
  if (rd != 0) {
    BIR[rd] = rd;
  }
\end{lstlisting}}
\newpage
\subsection{RV64D Instructions}
\label{sec:RV64D-instructions}

\subsubsection[{\small FCVT.D.L: Floating-Point Convert Double from Long}]{FCVT.D.L\hfill Floating-Point Convert Double from Long}

\paragraph{Encoding} Format \textsf{RV64D\_RTypeFCL2D} (Section~\ref{sec:format-RV64D-RTypeFCL2D}), \verb|rs2 = 00010|\\

\bitfields{ 5/11010, 2/01, 5/00010, 5/rs1, 3/frm, 5/frd, 7/1010011 }


\paragraph{Syntax} \texttt{fcvt.d.l} \emph{frd}, \emph{rs1}, \emph{frm}

\paragraph{Description} 

\paragraph{Modifier(s)}
\noindent\begin{itemize}
\item frm (cf. froundmode in section \ref{par:modifier-froundmode} on page \pageref{par:modifier-froundmode})
\end{itemize}

\paragraph{Execution} {Reservation table \textsf{ALU\_LITE} (Section~\ref{sec:reservation-ALU-LITE})}
~
{\begin{lstlisting}
stage ID:
  new frm = _ZX_3(frm);
stage RR:
  new rs1 = BIR[rs1];
  new frm = decode_riscv_float_rounding_mode(frm);
  if (check_float_rounding_mode_trap(frm) == 0) {
    _THROW(OPCODE);
  } else {
    new frd = i64_to_f64(frm, rs1);
    FPR[frd] = frd;
    commit_float_exception_flags();
  }
\end{lstlisting}}
\newpage
\subsubsection[{\small FCVT.D.LU: Floating-Point Convert Double from Long Unsigned}]{FCVT.D.LU\hfill Floating-Point Convert Double from Long Unsigned}

\paragraph{Encoding} Format \textsf{RV64D\_RTypeFCL2D} (Section~\ref{sec:format-RV64D-RTypeFCL2D}), \verb|rs2 = 00011|\\

\bitfields{ 5/11010, 2/01, 5/00011, 5/rs1, 3/frm, 5/frd, 7/1010011 }


\paragraph{Syntax} \texttt{fcvt.d.lu} \emph{frd}, \emph{rs1}, \emph{frm}

\paragraph{Description} 

\paragraph{Modifier(s)}
\noindent\begin{itemize}
\item frm (cf. froundmode in section \ref{par:modifier-froundmode} on page \pageref{par:modifier-froundmode})
\end{itemize}

\paragraph{Execution} {Reservation table \textsf{ALU\_LITE} (Section~\ref{sec:reservation-ALU-LITE})}
~
{\begin{lstlisting}
stage ID:
  new frm = _ZX_3(frm);
stage RR:
  new rs1 = BIR[rs1];
  new frm = decode_riscv_float_rounding_mode(frm);
  if (check_float_rounding_mode_trap(frm) == 0) {
    _THROW(OPCODE);
  } else {
    new frd = ui64_to_f64(frm, rs1);
    FPR[frd] = frd;
    commit_float_exception_flags();
  }
\end{lstlisting}}
\newpage
\subsubsection[{\small FCVT.L.D: Floating-Point Convert Long from Double}]{FCVT.L.D\hfill Floating-Point Convert Long from Double}

\paragraph{Encoding} Format \textsf{RV64D\_RTypeFCD2L} (Section~\ref{sec:format-RV64D-RTypeFCD2L}), \verb|rs2 = 00010|\\

\bitfields{ 5/11000, 2/01, 5/00010, 5/frs1, 3/frm, 5/rd, 7/1010011 }


\paragraph{Syntax} \texttt{fcvt.l.d} \emph{rd}, \emph{frs1}, \emph{frm}

\paragraph{Description} 

\paragraph{Modifier(s)}
\noindent\begin{itemize}
\item frm (cf. froundmode in section \ref{par:modifier-froundmode} on page \pageref{par:modifier-froundmode})
\end{itemize}

\paragraph{Execution} {Reservation table \textsf{ALU\_LITE} (Section~\ref{sec:reservation-ALU-LITE})}
~
{\begin{lstlisting}
stage ID:
  new frm = _ZX_3(frm);
stage RR:
  new frs1 = FPR[frs1];
  new frm = decode_riscv_float_rounding_mode(frm);
  if (check_float_rounding_mode_trap(frm) == 0) {
    _THROW(OPCODE);
  } else {
    new rd = f64_to_i64(frm, frs1);
    if (rd != 0) {
      BIR[rd] = rd;
    };
    commit_float_exception_flags();
  }
\end{lstlisting}}
\newpage
\subsubsection[{\small FCVT.LU.D: Floating-Point Convert Long Unsigned from Double}]{FCVT.LU.D\hfill Floating-Point Convert Long Unsigned from Double}

\paragraph{Encoding} Format \textsf{RV64D\_RTypeFCD2L} (Section~\ref{sec:format-RV64D-RTypeFCD2L}), \verb|rs2 = 00011|\\

\bitfields{ 5/11000, 2/01, 5/00011, 5/frs1, 3/frm, 5/rd, 7/1010011 }


\paragraph{Syntax} \texttt{fcvt.lu.d} \emph{rd}, \emph{frs1}, \emph{frm}

\paragraph{Description} 

\paragraph{Modifier(s)}
\noindent\begin{itemize}
\item frm (cf. froundmode in section \ref{par:modifier-froundmode} on page \pageref{par:modifier-froundmode})
\end{itemize}

\paragraph{Execution} {Reservation table \textsf{ALU\_LITE} (Section~\ref{sec:reservation-ALU-LITE})}
~
{\begin{lstlisting}
stage ID:
  new frm = _ZX_3(frm);
stage RR:
  new frs1 = FPR[frs1];
  new frm = decode_riscv_float_rounding_mode(frm);
  if (check_float_rounding_mode_trap(frm) == 0) {
    _THROW(OPCODE);
  } else {
    new rd = f64_to_ui64(frm, frs1);
    if (rd != 0) {
      BIR[rd] = rd;
    };
    commit_float_exception_flags();
  }
\end{lstlisting}}
\newpage
\subsubsection[{\small FMV.D.X: Floating-Point Move Double from Integer}]{FMV.D.X\hfill Floating-Point Move Double from Integer}

\paragraph{Encoding} Format \textsf{RV64D\_RTypeFMI2F} (Section~\ref{sec:format-RV64D-RTypeFMI2F}), \verb|funct3 = 000|\\

\bitfields{ 5/11110, 2/01, 5/00000, 5/rs1, 3/000, 5/frd, 7/1010011 }


\paragraph{Syntax} \texttt{fmv.d.x} \emph{frd}, \emph{rs1}

\paragraph{Description} 

\paragraph{Execution} {Reservation table \textsf{ALU\_TINY} (Section~\ref{sec:reservation-ALU-TINY})}
~
{\begin{lstlisting}
stage RR:
  new rs1 = BIR[rs1];
  new frd = rs1;
stage E1:
  FPR[frd] = frd;
\end{lstlisting}}
\newpage
\subsubsection[{\small FMV.X.D: Floating-Point Move Integer from Double}]{FMV.X.D\hfill Floating-Point Move Integer from Double}

\paragraph{Encoding} Format \textsf{RV64D\_RTypeFMF2I} (Section~\ref{sec:format-RV64D-RTypeFMF2I}), \verb|funct3 = 000|\\

\bitfields{ 5/11100, 2/01, 5/00000, 5/frs1, 3/000, 5/rd, 7/1010011 }


\paragraph{Syntax} \texttt{fmv.x.d} \emph{rd}, \emph{frs1}

\paragraph{Description} 

\paragraph{Execution} {Reservation table \textsf{ALU\_TINY} (Section~\ref{sec:reservation-ALU-TINY})}
~
{\begin{lstlisting}
stage RR:
  new frs1 = FPR[frs1];
  new rd = frs1;
stage E1:
  if (rd != 0) {
    BIR[rd] = rd;
  }
\end{lstlisting}}
\newpage
\subsection{RV64F Instructions}
\label{sec:RV64F-instructions}

\subsubsection[{\small FCVT.L.S: Floating-Point Convert Long from Single}]{FCVT.L.S\hfill Floating-Point Convert Long from Single}

\paragraph{Encoding} Format \textsf{RV64F\_RTypeFCF2L} (Section~\ref{sec:format-RV64F-RTypeFCF2L}), \verb|rs2 = 00010|\\

\bitfields{ 5/11000, 2/00, 5/00010, 5/frs1, 3/frm, 5/rd, 7/1010011 }


\paragraph{Syntax} \texttt{fcvt.l.s} \emph{rd}, \emph{frs1}, \emph{frm}

\paragraph{Description} 

\paragraph{Modifier(s)}
\noindent\begin{itemize}
\item frm (cf. froundmode in section \ref{par:modifier-froundmode} on page \pageref{par:modifier-froundmode})
\end{itemize}

\paragraph{Execution} {Reservation table \textsf{ALU\_LITE} (Section~\ref{sec:reservation-ALU-LITE})}
~
{\begin{lstlisting}
stage ID:
  new frm = _ZX_3(frm);
stage RR:
  new frs1 = FPR[frs1];
  new frm = decode_riscv_float_rounding_mode(frm);
  if (check_float_rounding_mode_trap(frm) == 0) {
    _THROW(OPCODE);
  } else {
    new rd = f32_to_i64(frm, frs1);
    if (rd != 0) {
      BIR[rd] = rd;
    };
    commit_float_exception_flags();
  }
\end{lstlisting}}
\newpage
\subsubsection[{\small FCVT.LU.S: Floating-Point Convert Long Unsigned from Single}]{FCVT.LU.S\hfill Floating-Point Convert Long Unsigned from Single}

\paragraph{Encoding} Format \textsf{RV64F\_RTypeFCF2L} (Section~\ref{sec:format-RV64F-RTypeFCF2L}), \verb|rs2 = 00011|\\

\bitfields{ 5/11000, 2/00, 5/00011, 5/frs1, 3/frm, 5/rd, 7/1010011 }


\paragraph{Syntax} \texttt{fcvt.lu.s} \emph{rd}, \emph{frs1}, \emph{frm}

\paragraph{Description} 

\paragraph{Modifier(s)}
\noindent\begin{itemize}
\item frm (cf. froundmode in section \ref{par:modifier-froundmode} on page \pageref{par:modifier-froundmode})
\end{itemize}

\paragraph{Execution} {Reservation table \textsf{ALU\_LITE} (Section~\ref{sec:reservation-ALU-LITE})}
~
{\begin{lstlisting}
stage ID:
  new frm = _ZX_3(frm);
stage RR:
  new frs1 = FPR[frs1];
  new frm = decode_riscv_float_rounding_mode(frm);
  if (check_float_rounding_mode_trap(frm) == 0) {
    _THROW(OPCODE);
  } else {
    new rd = f32_to_ui64(frm, frs1);
    if (rd != 0) {
      BIR[rd] = rd;
    };
    commit_float_exception_flags();
  }
\end{lstlisting}}
\newpage
\subsubsection[{\small FCVT.S.L: Floating-Point Convert Single from Long}]{FCVT.S.L\hfill Floating-Point Convert Single from Long}

\paragraph{Encoding} Format \textsf{RV64F\_RTypeFCL2F} (Section~\ref{sec:format-RV64F-RTypeFCL2F}), \verb|rs2 = 00010|\\

\bitfields{ 5/11010, 2/00, 5/00010, 5/rs1, 3/frm, 5/frd, 7/1010011 }


\paragraph{Syntax} \texttt{fcvt.s.l} \emph{frd}, \emph{rs1}, \emph{frm}

\paragraph{Description} 

\paragraph{Modifier(s)}
\noindent\begin{itemize}
\item frm (cf. froundmode in section \ref{par:modifier-froundmode} on page \pageref{par:modifier-froundmode})
\end{itemize}

\paragraph{Execution} {Reservation table \textsf{ALU\_LITE} (Section~\ref{sec:reservation-ALU-LITE})}
~
{\begin{lstlisting}
stage ID:
  new frm = _ZX_3(frm);
stage RR:
  new rs1 = BIR[rs1];
  new frm = decode_riscv_float_rounding_mode(frm);
  if (check_float_rounding_mode_trap(frm) == 0) {
    _THROW(OPCODE);
  } else {
    new frd = i64_to_f32(frm, rs1) | 0xFFFFFFFF00000000;
    FPR[frd] = frd;
    commit_float_exception_flags();
  }
\end{lstlisting}}
\newpage
\subsubsection[{\small FCVT.S.LU: Floating-Point Convert Single from Long Unsigned}]{FCVT.S.LU\hfill Floating-Point Convert Single from Long Unsigned}

\paragraph{Encoding} Format \textsf{RV64F\_RTypeFCL2F} (Section~\ref{sec:format-RV64F-RTypeFCL2F}), \verb|rs2 = 00011|\\

\bitfields{ 5/11010, 2/00, 5/00011, 5/rs1, 3/frm, 5/frd, 7/1010011 }


\paragraph{Syntax} \texttt{fcvt.s.lu} \emph{frd}, \emph{rs1}, \emph{frm}

\paragraph{Description} 

\paragraph{Modifier(s)}
\noindent\begin{itemize}
\item frm (cf. froundmode in section \ref{par:modifier-froundmode} on page \pageref{par:modifier-froundmode})
\end{itemize}

\paragraph{Execution} {Reservation table \textsf{ALU\_LITE} (Section~\ref{sec:reservation-ALU-LITE})}
~
{\begin{lstlisting}
stage ID:
  new frm = _ZX_3(frm);
stage RR:
  new rs1 = BIR[rs1];
  new frm = decode_riscv_float_rounding_mode(frm);
  if (check_float_rounding_mode_trap(frm) == 0) {
    _THROW(OPCODE);
  } else {
    new frd = ui64_to_f32(frm, rs1) | 0xFFFFFFFF00000000;
    FPR[frd] = frd;
    commit_float_exception_flags();
  }
\end{lstlisting}}
\newpage
\subsection{RV64I Instructions}
\label{sec:RV64I-instructions}

\subsubsection[{\small ADDIW: Integer ADD Immediate Word}]{ADDIW\hfill Integer ADD Immediate Word}

\paragraph{Encoding} Format \textsf{RV64I\_I\_ALU} (Section~\ref{sec:format-RV64I-I-ALU}), \verb|funct3 = 000|\\

\bitfields{ 12/i-- 11-- 0, 5/rs1, 3/000, 5/rd, 7/0011011 }


\paragraph{Syntax} \texttt{addiw} \emph{rd}, \emph{rs1}, \emph{i\_\discretionary{}{}{}11\_\discretionary{}{}{}0}

\paragraph{Description} The \emph{rs1} is added to the \emph{i\_\discretionary{}{}{}11\_\discretionary{}{}{}0}. The 32-bit sign-extended result is stored into the \emph{rd}.


\paragraph{Execution} {Reservation table \textsf{ALU\_TINY} (Section~\ref{sec:reservation-ALU-TINY})}
~
{\begin{lstlisting}
stage RR:
  new signed12 = _SX_12(i_11_0);
  new rs1 = BIR[rs1];
  new rd = _SX_32(rs1 + signed12);
stage E1:
  if (rd != 0) {
    BIR[rd] = rd;
  }
\end{lstlisting}}
\newpage
\subsubsection[{\small ADDW: Integer Add Word}]{ADDW\hfill Integer Add Word}

\paragraph{Encoding} Format \textsf{RV64I\_RF00} (Section~\ref{sec:format-RV64I-RF00}), \verb|funct3 = 000|\\

\bitfields{ 7/0000000, 5/rs2, 5/rs1, 3/000, 5/rd, 7/0111011 }


\paragraph{Syntax} \texttt{addw} \emph{rd}, \emph{rs1}, \emph{rs2}

\paragraph{Description} The \emph{rs1} is added to the \emph{rs2}. The 32-bit sign-extended result is stored into the \emph{rd}.


\paragraph{Execution} {Reservation table \textsf{ALU\_TINY} (Section~\ref{sec:reservation-ALU-TINY})}
~
{\begin{lstlisting}
stage RR:
  new rs2 = BIR[rs2];
  new rs1 = BIR[rs1];
  new rd = _SX_32(rs1 + rs2);
stage E1:
  if (rd != 0) {
    BIR[rd] = rd;
  }
\end{lstlisting}}
\newpage
\subsubsection[{\small LD: Load Double Word}]{LD\hfill Load Double Word}

\paragraph{Encoding} Format \textsf{RV64I\_ILoad} (Section~\ref{sec:format-RV64I-ILoad}), \verb|funct3 = 011|\\

\bitfields{ 12/i-- 11-- 0, 5/rs1, 3/011, 5/rd, 7/0000011 }


\paragraph{Syntax} \texttt{ld} \emph{rd}, \emph{i\_\discretionary{}{}{}11\_\discretionary{}{}{}0}[\emph{rs1}]

\paragraph{Description} The \emph{rd} receives the memory double word at address \emph{rs1} + \emph{i\_\discretionary{}{}{}11\_\discretionary{}{}{}0}.


\paragraph{Execution} {Reservation table \textsf{LSU\_AUXW} (Section~\ref{sec:reservation-LSU-AUXW})}
~
{\begin{lstlisting}
stage RR:
  new signed12 = _SX_12(i_11_0);
  new rs1 = BIR[rs1];
  new address = rs1 + signed12;
  new rd = _SX_64(MEM.load(address, 0xFF, 0, rd));
stage E1:
  if (rd != 0) {
    BIR[rd] = rd;
  }
\end{lstlisting}}
\newpage
\subsubsection[{\small LWU: Load Word Unsigned}]{LWU\hfill Load Word Unsigned}

\paragraph{Encoding} Format \textsf{RV64I\_ILoad} (Section~\ref{sec:format-RV64I-ILoad}), \verb|funct3 = 110|\\

\bitfields{ 12/i-- 11-- 0, 5/rs1, 3/110, 5/rd, 7/0000011 }


\paragraph{Syntax} \texttt{lwu} \emph{rd}, \emph{i\_\discretionary{}{}{}11\_\discretionary{}{}{}0}[\emph{rs1}]

\paragraph{Description} The \emph{rd} receives the zero-extended memory word at address \emph{rs1} + \emph{i\_\discretionary{}{}{}11\_\discretionary{}{}{}0}.


\paragraph{Execution} {Reservation table \textsf{LSU\_AUXW} (Section~\ref{sec:reservation-LSU-AUXW})}
~
{\begin{lstlisting}
stage RR:
  new signed12 = _SX_12(i_11_0);
  new rs1 = BIR[rs1];
  new address = rs1 + signed12;
  new rd = _ZX_32(MEM.load(address, 0xF, 0, rd));
stage E1:
  if (rd != 0) {
    BIR[rd] = rd;
  }
\end{lstlisting}}
\newpage
\subsubsection[{\small SD: Store Double Word}]{SD\hfill Store Double Word}

\paragraph{Encoding} Format \textsf{RV64I\_SType} (Section~\ref{sec:format-RV64I-SType}), \verb|funct3 = 011|\\

\bitfields{ 7/i-- 11-- 5, 5/rs2, 5/rs1, 3/011, 5/i-- 4-- 0s, 7/0100011 }


\paragraph{Syntax} \texttt{sd} \emph{rs2}, \emph{i\_\discretionary{}{}{}11\_\discretionary{}{}{}5\_\discretionary{}{}{}i\_\discretionary{}{}{}4\_\discretionary{}{}{}0s}[\emph{rs1}]

\paragraph{Description} The \emph{rs2} is stored into the memory double word at address \emph{rs1} + \emph{i\_\discretionary{}{}{}11\_\discretionary{}{}{}5\_\discretionary{}{}{}i\_\discretionary{}{}{}4\_\discretionary{}{}{}0s}.


\paragraph{Execution} {Reservation table \textsf{LSU\_MEMW\_AUXR} (Section~\ref{sec:reservation-LSU-MEMW-AUXR})}
~
{\begin{lstlisting}
stage RR:
  new signed12 = _SX_12(i_11_5_i_4_0s);
  new rs1 = BIR[rs1];
  new address = rs1 + signed12;
stage E1:
  new rs2 = BIR[rs2];
stage E3:
  MEM.store(address, 0xFF, 0, rs2, rs2);
\end{lstlisting}}
\newpage
\subsubsection[{\small SLLI: Shift Left Logical Immediate}]{SLLI\hfill Shift Left Logical Immediate}

\paragraph{Encoding} Format \textsf{RV64I\_XS00} (Section~\ref{sec:format-RV64I-XS00}), \verb|funct3 = 001|\\

\bitfields{ 6/000000, 6/i-- 5-- 0, 5/rs1, 3/001, 5/rd, 7/0010011 }


\paragraph{Syntax} \texttt{slli} \emph{rd}, \emph{rs1}, \emph{i\_\discretionary{}{}{}5\_\discretionary{}{}{}0}

\paragraph{Description} The \emph{rs1} is shifted left by the \emph{i\_\discretionary{}{}{}5\_\discretionary{}{}{}0}. The result is stored into the \emph{rd}.


\paragraph{Execution} {Reservation table \textsf{ALU\_TINY} (Section~\ref{sec:reservation-ALU-TINY})}
~
{\begin{lstlisting}
stage RR:
  new unsigned6 = _ZX_6(i_5_0);
  new rs1 = BIR[rs1];
  new rd = rs1 << unsigned6;
stage E1:
  if (rd != 0) {
    BIR[rd] = rd;
  }
\end{lstlisting}}
\newpage
\subsubsection[{\small SLLIW: Shift Left Logical Immediate Word}]{SLLIW\hfill Shift Left Logical Immediate Word}

\paragraph{Encoding} Format \textsf{RV64I\_IS00} (Section~\ref{sec:format-RV64I-IS00}), \verb|funct3 = 001|\\

\bitfields{ 7/0000000, 5/i-- 4-- 0, 5/rs1, 3/001, 5/rd, 7/0011011 }


\paragraph{Syntax} \texttt{slliw} \emph{rd}, \emph{rs1}, \emph{i\_\discretionary{}{}{}4\_\discretionary{}{}{}0}

\paragraph{Description} The \emph{rs1} is shifted left by the \emph{i\_\discretionary{}{}{}4\_\discretionary{}{}{}0}. The 32-bit sign-extended result is stored into the \emph{rd}.


\paragraph{Execution} {Reservation table \textsf{ALU\_TINY} (Section~\ref{sec:reservation-ALU-TINY})}
~
{\begin{lstlisting}
stage RR:
  new unsigned5 = _ZX_5(i_4_0);
  new rs1 = BIR[rs1];
  new rd = _SX_32(rs1 << unsigned5);
stage E1:
  if (rd != 0) {
    BIR[rd] = rd;
  }
\end{lstlisting}}
\newpage
\subsubsection[{\small SLLW: Shift Left Logical Word}]{SLLW\hfill Shift Left Logical Word}

\paragraph{Encoding} Format \textsf{RV64I\_RF00} (Section~\ref{sec:format-RV64I-RF00}), \verb|funct3 = 001|\\

\bitfields{ 7/0000000, 5/rs2, 5/rs1, 3/001, 5/rd, 7/0111011 }


\paragraph{Syntax} \texttt{sllw} \emph{rd}, \emph{rs1}, \emph{rs2}

\paragraph{Description} The \emph{rs1} is shifted left by the \emph{rs2} modulo 32. The 32-bit sign-extended result is stored into the \emph{rd}.


\paragraph{Execution} {Reservation table \textsf{ALU\_TINY} (Section~\ref{sec:reservation-ALU-TINY})}
~
{\begin{lstlisting}
stage RR:
  new rs2 = BIR[rs2];
  new rs1 = BIR[rs1];
  new rd = _SX_32(rs1 << (rs2 & 0x1F));
stage E1:
  if (rd != 0) {
    BIR[rd] = rd;
  }
\end{lstlisting}}
\newpage
\subsubsection[{\small SRAI: Shift Right Arithmetic Immediate}]{SRAI\hfill Shift Right Arithmetic Immediate}

\paragraph{Encoding} Format \textsf{RV64I\_XS10} (Section~\ref{sec:format-RV64I-XS10}), \verb|funct3 = 101|\\

\bitfields{ 6/010000, 6/i-- 5-- 0, 5/rs1, 3/101, 5/rd, 7/0010011 }


\paragraph{Syntax} \texttt{srai} \emph{rd}, \emph{rs1}, \emph{i\_\discretionary{}{}{}5\_\discretionary{}{}{}0}

\paragraph{Description} The \emph{rs1} assuming signed integer is shifted right by the \emph{i\_\discretionary{}{}{}5\_\discretionary{}{}{}0}. The result is stored into the \emph{rd}.


\paragraph{Execution} {Reservation table \textsf{ALU\_TINY} (Section~\ref{sec:reservation-ALU-TINY})}
~
{\begin{lstlisting}
stage RR:
  new unsigned6 = _ZX_6(i_5_0);
  new rs1 = BIR[rs1];
  new rd = _SX_64(rs1) >> unsigned6;
stage E1:
  if (rd != 0) {
    BIR[rd] = rd;
  }
\end{lstlisting}}
\newpage
\subsubsection[{\small SRAIW: Shift Right Arithmetic Immediate Word}]{SRAIW\hfill Shift Right Arithmetic Immediate Word}

\paragraph{Encoding} Format \textsf{RV64I\_IS20} (Section~\ref{sec:format-RV64I-IS20}), \verb|funct3 = 101|\\

\bitfields{ 7/0100000, 5/i-- 4-- 0, 5/rs1, 3/101, 5/rd, 7/0011011 }


\paragraph{Syntax} \texttt{sraiw} \emph{rd}, \emph{rs1}, \emph{i\_\discretionary{}{}{}4\_\discretionary{}{}{}0}

\paragraph{Description} The \emph{rs1} assuming signed integer is shifted right by the \emph{i\_\discretionary{}{}{}4\_\discretionary{}{}{}0}. The 32-bit sign-extended result is stored into the \emph{rd}.


\paragraph{Execution} {Reservation table \textsf{ALU\_TINY} (Section~\ref{sec:reservation-ALU-TINY})}
~
{\begin{lstlisting}
stage RR:
  new unsigned5 = _ZX_5(i_4_0);
  new rs1 = BIR[rs1];
  new rd = _SX_32(_SX_32(rs1) >> unsigned5);
stage E1:
  if (rd != 0) {
    BIR[rd] = rd;
  }
\end{lstlisting}}
\newpage
\subsubsection[{\small SRAW: Shift Right Arithmetic Word}]{SRAW\hfill Shift Right Arithmetic Word}

\paragraph{Encoding} Format \textsf{RV64I\_RF20} (Section~\ref{sec:format-RV64I-RF20}), \verb|funct3 = 101|\\

\bitfields{ 7/0100000, 5/rs2, 5/rs1, 3/101, 5/rd, 7/0111011 }


\paragraph{Syntax} \texttt{sraw} \emph{rd}, \emph{rs1}, \emph{rs2}

\paragraph{Description} The \emph{rs1} assuming signed integer is shifted right by the \emph{rs2} modulo 32. The 32-bit sign-extended result is stored into the \emph{rd}.


\paragraph{Execution} {Reservation table \textsf{ALU\_TINY} (Section~\ref{sec:reservation-ALU-TINY})}
~
{\begin{lstlisting}
stage RR:
  new rs2 = BIR[rs2];
  new rs1 = BIR[rs1];
  new rd = _SX_32(_SX_32(rs1) >> (rs2 & 0x1F));
stage E1:
  if (rd != 0) {
    BIR[rd] = rd;
  }
\end{lstlisting}}
\newpage
\subsubsection[{\small SRLI: Shift Right Logical Immediate}]{SRLI\hfill Shift Right Logical Immediate}

\paragraph{Encoding} Format \textsf{RV64I\_XS00} (Section~\ref{sec:format-RV64I-XS00}), \verb|funct3 = 101|\\

\bitfields{ 6/000000, 6/i-- 5-- 0, 5/rs1, 3/101, 5/rd, 7/0010011 }


\paragraph{Syntax} \texttt{srli} \emph{rd}, \emph{rs1}, \emph{i\_\discretionary{}{}{}5\_\discretionary{}{}{}0}

\paragraph{Description} The \emph{rs1} assuming unsigned integer is shifted right by the \emph{i\_\discretionary{}{}{}5\_\discretionary{}{}{}0}. The result is stored into the \emph{rd}.


\paragraph{Execution} {Reservation table \textsf{ALU\_TINY} (Section~\ref{sec:reservation-ALU-TINY})}
~
{\begin{lstlisting}
stage RR:
  new unsigned6 = _ZX_6(i_5_0);
  new rs1 = BIR[rs1];
  new rd = _ZX_64(rs1) >> unsigned6;
stage E1:
  if (rd != 0) {
    BIR[rd] = rd;
  }
\end{lstlisting}}
\newpage
\subsubsection[{\small SRLIW: Shift Right Logical Immediate Word}]{SRLIW\hfill Shift Right Logical Immediate Word}

\paragraph{Encoding} Format \textsf{RV64I\_IS00} (Section~\ref{sec:format-RV64I-IS00}), \verb|funct3 = 101|\\

\bitfields{ 7/0000000, 5/i-- 4-- 0, 5/rs1, 3/101, 5/rd, 7/0011011 }


\paragraph{Syntax} \texttt{srliw} \emph{rd}, \emph{rs1}, \emph{i\_\discretionary{}{}{}4\_\discretionary{}{}{}0}

\paragraph{Description} The \emph{rs1} assuming unsigned integer is shifted right by the \emph{i\_\discretionary{}{}{}4\_\discretionary{}{}{}0}. The 32-bit sign-extended result is stored into the \emph{rd}.


\paragraph{Execution} {Reservation table \textsf{ALU\_TINY} (Section~\ref{sec:reservation-ALU-TINY})}
~
{\begin{lstlisting}
stage RR:
  new unsigned5 = _ZX_5(i_4_0);
  new rs1 = BIR[rs1];
  new rd = _SX_32(_ZX_32(rs1) >> unsigned5);
stage E1:
  if (rd != 0) {
    BIR[rd] = rd;
  }
\end{lstlisting}}
\newpage
\subsubsection[{\small SRLW: Shift Right Logical Word}]{SRLW\hfill Shift Right Logical Word}

\paragraph{Encoding} Format \textsf{RV64I\_RF00} (Section~\ref{sec:format-RV64I-RF00}), \verb|funct3 = 101|\\

\bitfields{ 7/0000000, 5/rs2, 5/rs1, 3/101, 5/rd, 7/0111011 }


\paragraph{Syntax} \texttt{srlw} \emph{rd}, \emph{rs1}, \emph{rs2}

\paragraph{Description} The \emph{rs1} assuming unsigned integer is shifted right by the \emph{rs2} modulo 32. The 32-bit sign-extended result is stored into the \emph{rd}.


\paragraph{Execution} {Reservation table \textsf{ALU\_TINY} (Section~\ref{sec:reservation-ALU-TINY})}
~
{\begin{lstlisting}
stage RR:
  new rs2 = BIR[rs2];
  new rs1 = BIR[rs1];
  new rd = _SX_32(_ZX_32(rs1) >> (rs2 & 0x1F));
stage E1:
  if (rd != 0) {
    BIR[rd] = rd;
  }
\end{lstlisting}}
\newpage
\subsubsection[{\small SUBW: Integer Subtract Word}]{SUBW\hfill Integer Subtract Word}

\paragraph{Encoding} Format \textsf{RV64I\_RF20} (Section~\ref{sec:format-RV64I-RF20}), \verb|funct3 = 000|\\

\bitfields{ 7/0100000, 5/rs2, 5/rs1, 3/000, 5/rd, 7/0111011 }


\paragraph{Syntax} \texttt{subw} \emph{rd}, \emph{rs1}, \emph{rs2}

\paragraph{Description} The \emph{rs1} is added to the \emph{rs2}. The result is stored into the \emph{rd}.


\paragraph{Execution} {Reservation table \textsf{ALU\_TINY} (Section~\ref{sec:reservation-ALU-TINY})}
~
{\begin{lstlisting}
stage RR:
  new rs2 = BIR[rs2];
  new rs1 = BIR[rs1];
  new rd = _SX_32(rs1 - rs2);
stage E1:
  if (rd != 0) {
    BIR[rd] = rd;
  }
\end{lstlisting}}
\newpage
\subsection{RV64M Instructions}
\label{sec:RV64M-instructions}

\subsubsection[{\small DIVUW: Integer Division Unsigned Word}]{DIVUW\hfill Integer Division Unsigned Word}

\paragraph{Encoding} Format \textsf{RV64M\_RType} (Section~\ref{sec:format-RV64M-RType}), \verb|funct3 = 101|\\

\bitfields{ 7/0000001, 5/rs2, 5/rs1, 3/101, 5/rd, 7/0111011 }


\paragraph{Syntax} \texttt{divuw} \emph{rd}, \emph{rs1}, \emph{rs2}

\paragraph{Description} The \emph{rs1} assuming unsigned integer is divided by the \emph{rs2} assuming unsigned integer. The 32-bit sign-extended quotient is stored into the \emph{rd}.


\paragraph{Execution} {Reservation table \textsf{ALU\_FULL} (Section~\ref{sec:reservation-ALU-FULL})}
~
{\begin{lstlisting}
stage RR:
  new rs2 = BIR[rs2];
  new rs1 = BIR[rs1];
  if (_ZX_32(rs2) == 0) {
    new rd = 0xFFFFFFFFFFFFFFFF;
  } else {
    rd = _SX_32(_ZX_32(rs1) / _ZX_32(rs2));
  }
stage E1:
  if (rd != 0) {
    BIR[rd] = rd;
  }
\end{lstlisting}}
\newpage
\subsubsection[{\small DIVW: Integer Division Word}]{DIVW\hfill Integer Division Word}

\paragraph{Encoding} Format \textsf{RV64M\_RType} (Section~\ref{sec:format-RV64M-RType}), \verb|funct3 = 100|\\

\bitfields{ 7/0000001, 5/rs2, 5/rs1, 3/100, 5/rd, 7/0111011 }


\paragraph{Syntax} \texttt{divw} \emph{rd}, \emph{rs1}, \emph{rs2}

\paragraph{Description} The \emph{rs1} assuming signed integer is divided by the \emph{rs2} assuming signed integer. The 32-bit sign-extended quotient is stored into the \emph{rd}.


\paragraph{Execution} {Reservation table \textsf{ALU\_FULL} (Section~\ref{sec:reservation-ALU-FULL})}
~
{\begin{lstlisting}
stage RR:
  new rs2 = BIR[rs2];
  new rs1 = BIR[rs1];
  if (_ZX_32(rs2) == 0) {
    new rd = 0xFFFFFFFFFFFFFFFF;
  } else if ((_ZX_32(rs1) == 0x80000000) && (_ZX_32(rs2) == 0xFFFFFFFF)) {
    rd = _SX_32(rs1);
  } else {
    rd = _SX_32(_SX_32(rs1) / _SX_32(rs2));
  }
stage E1:
  if (rd != 0) {
    BIR[rd] = rd;
  }
\end{lstlisting}}
\newpage
\subsubsection[{\small MULW: Integer Multiply Word}]{MULW\hfill Integer Multiply Word}

\paragraph{Encoding} Format \textsf{RV64M\_RType} (Section~\ref{sec:format-RV64M-RType}), \verb|funct3 = 000|\\

\bitfields{ 7/0000001, 5/rs2, 5/rs1, 3/000, 5/rd, 7/0111011 }


\paragraph{Syntax} \texttt{mulw} \emph{rd}, \emph{rs1}, \emph{rs2}

\paragraph{Description} The \emph{rs1} is multiplied with the \emph{rs2}. The 32-bit sign-extended result is stored into the \emph{rd}.


\paragraph{Execution} {Reservation table \textsf{ALU\_LITE} (Section~\ref{sec:reservation-ALU-LITE})}
~
{\begin{lstlisting}
stage RR:
  new rs2 = BIR[rs2];
  new rs1 = BIR[rs1];
  new rd = _SX_32(rs1 * rs2);
stage E1:
  if (rd != 0) {
    BIR[rd] = rd;
  }
\end{lstlisting}}
\newpage
\subsubsection[{\small REMUW: Integer Remainder Unsigned Word}]{REMUW\hfill Integer Remainder Unsigned Word}

\paragraph{Encoding} Format \textsf{RV64M\_RType} (Section~\ref{sec:format-RV64M-RType}), \verb|funct3 = 111|\\

\bitfields{ 7/0000001, 5/rs2, 5/rs1, 3/111, 5/rd, 7/0111011 }


\paragraph{Syntax} \texttt{remuw} \emph{rd}, \emph{rs1}, \emph{rs2}

\paragraph{Description} The \emph{rs1} assuming unsigned integer is divided by the \emph{rs2} assuming unsigned integer. The 32-bit sign-extended remainder is stored into the \emph{rd}.


\paragraph{Execution} {Reservation table \textsf{ALU\_FULL} (Section~\ref{sec:reservation-ALU-FULL})}
~
{\begin{lstlisting}
stage RR:
  new rs2 = BIR[rs2];
  new rs1 = BIR[rs1];
  if (_ZX_32(rs2) == 0) {
    new rd = _SX_32(rs1);
  } else {
    rd = _SX_32(_ZX_32(rs1) % _ZX_32(rs2));
  }
stage E1:
  if (rd != 0) {
    BIR[rd] = rd;
  }
\end{lstlisting}}
\newpage
\subsubsection[{\small REMW: Integer Remainder Word}]{REMW\hfill Integer Remainder Word}

\paragraph{Encoding} Format \textsf{RV64M\_RType} (Section~\ref{sec:format-RV64M-RType}), \verb|funct3 = 110|\\

\bitfields{ 7/0000001, 5/rs2, 5/rs1, 3/110, 5/rd, 7/0111011 }


\paragraph{Syntax} \texttt{remw} \emph{rd}, \emph{rs1}, \emph{rs2}

\paragraph{Description} The \emph{rs1} assuming signed integer is divided by the \emph{rs2} assuming signed integer. The 32-bit sign-extended remainder is stored into the \emph{rd}.


\paragraph{Execution} {Reservation table \textsf{ALU\_FULL} (Section~\ref{sec:reservation-ALU-FULL})}
~
{\begin{lstlisting}
stage RR:
  new rs2 = BIR[rs2];
  new rs1 = BIR[rs1];
  if (_ZX_32(rs2) == 0) {
    new rd = _SX_32(rs1);
  } else if ((_ZX_32(rs1) == 0x80000000) && (_ZX_32(rs2) == 0xFFFFFFFF)) {
    rd = 0;
  } else {
    rd = _SX_32(_SX_32(rs1) % _SX_32(rs2));
  }
stage E1:
  if (rd != 0) {
    BIR[rd] = rd;
  }
\end{lstlisting}}
\newpage
\subsection{RVZI Instructions}
\label{sec:RVZI-instructions}

\subsubsection[{\small CSRRC: CSR Atomic Read and Clear}]{CSRRC\hfill CSR Atomic Read and Clear}

\paragraph{Encoding} Format \textsf{RVZI\_CSRR} (Section~\ref{sec:format-RVZI-CSRR}), \verb|funct3 = 011|\\

\bitfields{ 12/csr, 5/rs1, 3/011, 5/rd, 7/1110011 }


\paragraph{Syntax} \texttt{csrrc} \emph{rd}, \emph{csr}, \emph{rs1}

\paragraph{Description} 

\paragraph{Execution} {Reservation table \textsf{BCU2\_TINY\_LSU} (Section~\ref{sec:reservation-BCU2-TINY-LSU})}
~
{\begin{lstlisting}
stage RR:
  new rs1 = BIR[rs1];
  new csr = CSR[csr];
stage E1:
  if (rd != 0) {
    BIR[rd] = csr;
  }
stage E3:
  if (rs1 != 0) {
    CSR[csr] = csr & ~rs1;
  }
\end{lstlisting}}
\newpage
\subsubsection[{\small CSRRCI: CSR Atomic Read and Clear Immediate}]{CSRRCI\hfill CSR Atomic Read and Clear Immediate}

\paragraph{Encoding} Format \textsf{RVZI\_CSRI} (Section~\ref{sec:format-RVZI-CSRI}), \verb|funct3 = 111|\\

\bitfields{ 12/csr, 5/uimm, 3/111, 5/rd, 7/1110011 }


\paragraph{Syntax} \texttt{csrrci} \emph{rd}, \emph{csr}, \emph{uimm}

\paragraph{Description} 

\paragraph{Execution} {Reservation table \textsf{BCU2\_TINY\_LSU} (Section~\ref{sec:reservation-BCU2-TINY-LSU})}
~
{\begin{lstlisting}
stage RR:
  new unsigned5 = _ZX_5(uimm);
  new csr = CSR[csr];
stage E1:
  if (rd != 0) {
    BIR[rd] = csr;
  }
stage E3:
  if (unsigned5 != 0) {
    CSR[csr] = csr & ~unsigned5;
  }
\end{lstlisting}}
\newpage
\subsubsection[{\small CSRRS: CSR Atomic Read and Set}]{CSRRS\hfill CSR Atomic Read and Set}

\paragraph{Encoding} Format \textsf{RVZI\_CSRR} (Section~\ref{sec:format-RVZI-CSRR}), \verb|funct3 = 010|\\

\bitfields{ 12/csr, 5/rs1, 3/010, 5/rd, 7/1110011 }


\paragraph{Syntax} \texttt{csrrs} \emph{rd}, \emph{csr}, \emph{rs1}

\paragraph{Description} 

\paragraph{Execution} {Reservation table \textsf{BCU2\_TINY\_LSU} (Section~\ref{sec:reservation-BCU2-TINY-LSU})}
~
{\begin{lstlisting}
stage RR:
  new rs1 = BIR[rs1];
  new csr = CSR[csr];
stage E1:
  if (rd != 0) {
    BIR[rd] = csr;
  }
stage E3:
  if (rs1 != 0) {
    CSR[csr] = csr | rs1;
  }
\end{lstlisting}}
\newpage
\subsubsection[{\small CSRRSI: CSR Atomic Read and Set Immediate}]{CSRRSI\hfill CSR Atomic Read and Set Immediate}

\paragraph{Encoding} Format \textsf{RVZI\_CSRI} (Section~\ref{sec:format-RVZI-CSRI}), \verb|funct3 = 110|\\

\bitfields{ 12/csr, 5/uimm, 3/110, 5/rd, 7/1110011 }


\paragraph{Syntax} \texttt{csrrsi} \emph{rd}, \emph{csr}, \emph{uimm}

\paragraph{Description} 

\paragraph{Execution} {Reservation table \textsf{BCU2\_TINY\_LSU} (Section~\ref{sec:reservation-BCU2-TINY-LSU})}
~
{\begin{lstlisting}
stage RR:
  new unsigned5 = _ZX_5(uimm);
  new csr = CSR[csr];
stage E1:
  if (rd != 0) {
    BIR[rd] = csr;
  }
stage E3:
  if (unsigned5 != 0) {
    CSR[csr] = csr | unsigned5;
  }
\end{lstlisting}}
\newpage
\subsubsection[{\small CSRRW: CSR Atomic Read Write}]{CSRRW\hfill CSR Atomic Read Write}

\paragraph{Encoding} Format \textsf{RVZI\_CSRR} (Section~\ref{sec:format-RVZI-CSRR}), \verb|funct3 = 001|\\

\bitfields{ 12/csr, 5/rs1, 3/001, 5/rd, 7/1110011 }


\paragraph{Syntax} \texttt{csrrw} \emph{rd}, \emph{csr}, \emph{rs1}

\paragraph{Description} 

\paragraph{Execution} {Reservation table \textsf{BCU2\_TINY\_LSU} (Section~\ref{sec:reservation-BCU2-TINY-LSU})}
~
{\begin{lstlisting}
stage RR:
  new rs1 = BIR[rs1];
  new csr = CSR[csr];
stage E1:
  if (rd != 0) {
    BIR[rd] = csr;
  }
stage E3:
  CSR[csr] = rs1;
\end{lstlisting}}
\newpage
\subsubsection[{\small CSRRWI: CSR Atomic Read Write Immediate}]{CSRRWI\hfill CSR Atomic Read Write Immediate}

\paragraph{Encoding} Format \textsf{RVZI\_CSRI} (Section~\ref{sec:format-RVZI-CSRI}), \verb|funct3 = 101|\\

\bitfields{ 12/csr, 5/uimm, 3/101, 5/rd, 7/1110011 }


\paragraph{Syntax} \texttt{csrrwi} \emph{rd}, \emph{csr}, \emph{uimm}

\paragraph{Description} 

\paragraph{Execution} {Reservation table \textsf{BCU2\_TINY\_LSU} (Section~\ref{sec:reservation-BCU2-TINY-LSU})}
~
{\begin{lstlisting}
stage RR:
  new unsigned5 = _ZX_5(uimm);
  new csr = CSR[csr];
stage E1:
  if (rd != 0) {
    BIR[rd] = csr;
  }
stage E3:
  CSR[csr] = unsigned5;
\end{lstlisting}}
\newpage
\subsubsection[{\small CZERO.EQZ: Conditional Copy Register}]{CZERO.EQZ\hfill Conditional Copy Register}

\paragraph{Encoding} Format \textsf{RVZI\_COND} (Section~\ref{sec:format-RVZI-COND}), \verb|funct3 = 101|\\

\bitfields{ 7/0000111, 5/rs2, 5/rs1, 3/101, 5/rd, 7/0110011 }


\paragraph{Syntax} \texttt{czero.eqz} \emph{rd}, \emph{rs1}, \emph{rs2}

\paragraph{Description} If \emph{rs2} is 0, 0 is copied into \emph{rd}, else \emph{rs1} is copied into \emph{rd}.


\paragraph{Execution} {Reservation table \textsf{ALU\_TINY} (Section~\ref{sec:reservation-ALU-TINY})}
~
{\begin{lstlisting}
stage RR:
  new rs2 = BIR[rs2];
  new rs1 = BIR[rs1];
  if (_ZX_64(rs2) == 0) {
    new rd = 0;
  } else {
    rd = _ZX_64(rs1);
  }
stage E1:
  if (rd != 0) {
    BIR[rd] = rd;
  }
\end{lstlisting}}
\newpage
\subsubsection[{\small CZERO.NEZ: Conditional Copy Register}]{CZERO.NEZ\hfill Conditional Copy Register}

\paragraph{Encoding} Format \textsf{RVZI\_COND} (Section~\ref{sec:format-RVZI-COND}), \verb|funct3 = 111|\\

\bitfields{ 7/0000111, 5/rs2, 5/rs1, 3/111, 5/rd, 7/0110011 }


\paragraph{Syntax} \texttt{czero.nez} \emph{rd}, \emph{rs1}, \emph{rs2}

\paragraph{Description} If \emph{rs2} is not 0, 0 is copied into \emph{rd}, else \emph{rs1} is copied into \emph{rd}.


\paragraph{Execution} {Reservation table \textsf{ALU\_TINY} (Section~\ref{sec:reservation-ALU-TINY})}
~
{\begin{lstlisting}
stage RR:
  new rs2 = BIR[rs2];
  new rs1 = BIR[rs1];
  if (_ZX_64(rs2) != 0) {
    new rd = 0;
  } else {
    rd = _ZX_64(rs1);
  }
stage E1:
  if (rd != 0) {
    BIR[rd] = rd;
  }
\end{lstlisting}}
\newpage
\subsubsection[{\small FENCE.I: Synchronize instruction and data.}]{FENCE.I\hfill Synchronize instruction and data.}

\paragraph{Encoding} Format \textsf{RVZI\_FenceI} (Section~\ref{sec:format-RVZI-FenceI}), \verb|funct3 = 001|\\

\bitfields{ 12/000000000000, 5/00000, 3/001, 5/00000, 7/0001111 }


\paragraph{Syntax} \texttt{fence.i}

\paragraph{Description} Ensures that a subsequent instruction fetch will observe the local hart’s stores.


\paragraph{Execution} {Reservation table \textsf{ALL} (Section~\ref{sec:reservation-ALL})}
~
{\begin{lstlisting}
stage RR:
  barrier();
stage E3:
  MEM.i1inval();
\end{lstlisting}}
\newpage
\subsection{RV64KV Instructions}
\label{sec:RV64KV-instructions}

\subsubsection[{\small KV.LQ: Load Quadruple Word}]{KV.LQ\hfill Load Quadruple Word}

\paragraph{Encoding} Format \textsf{RV64KV\_ILoad} (Section~\ref{sec:format-RV64KV-ILoad}), \verb|funct3 = 111|\\

\bitfields{ 12/i-- 11-- 0, 5/rs1, 3/111, 4/rd-- pair, 1/0, 7/0000011 }


\paragraph{Syntax} \texttt{kv.lq} \emph{rd\_\discretionary{}{}{}pair}, \emph{i\_\discretionary{}{}{}11\_\discretionary{}{}{}0}[\emph{rs1}]

\paragraph{Description} The \emph{rd\_\discretionary{}{}{}pair} receives the memory double word at address \emph{rs1} + \emph{i\_\discretionary{}{}{}11\_\discretionary{}{}{}0}.


\paragraph{Execution} {Reservation table \textsf{LSU\_AUXW} (Section~\ref{sec:reservation-LSU-AUXW})}
~
{\begin{lstlisting}
stage RR:
  new signed12 = _SX_12(i_11_0);
  new rs1 = BIR[rs1];
  new address = rs1 + signed12;
  new rd = _SX_128(MEM.load(address, 0xFFFF, 0, rd_pair << 1));
stage E1:
  if (rd_pair != 0) {
    BIRP[rd_pair] = rd;
  }
\end{lstlisting}}
\newpage
\subsubsection[{\small KV.SQ: Store Quadruple Word}]{KV.SQ\hfill Store Quadruple Word}

\paragraph{Encoding} Format \textsf{RV64KV\_SType} (Section~\ref{sec:format-RV64KV-SType}), \verb|funct3 = 100|\\

\bitfields{ 7/i-- 11-- 5, 4/rs2-- pair, 1/0, 5/rs1, 3/100, 5/i-- 4-- 0s, 7/0100011 }


\paragraph{Syntax} \texttt{kv.sq} \emph{rs2\_\discretionary{}{}{}pair}, \emph{i\_\discretionary{}{}{}11\_\discretionary{}{}{}5\_\discretionary{}{}{}i\_\discretionary{}{}{}4\_\discretionary{}{}{}0s}[\emph{rs1}]

\paragraph{Description} The \emph{rs2\_\discretionary{}{}{}pair} is stored into the memory double word at address \emph{rs1} + \emph{i\_\discretionary{}{}{}11\_\discretionary{}{}{}5\_\discretionary{}{}{}i\_\discretionary{}{}{}4\_\discretionary{}{}{}0s}.


\paragraph{Execution} {Reservation table \textsf{LSU\_MEMW\_AUXR} (Section~\ref{sec:reservation-LSU-MEMW-AUXR})}
~
{\begin{lstlisting}
stage RR:
  new signed12 = _SX_12(i_11_5_i_4_0s);
  new rs1 = BIR[rs1];
  new address = rs1 + signed12;
stage E1:
  new rs2 = BIRP[rs2_pair];
stage E3:
  MEM.store(address, 0xFFFF, 0, rs2, rs2_pair << 1);
\end{lstlisting}}
\newpage
\subsection{Instruction operands}
\label{sec:Instruction-operands}

\subsubsection{Immediate classes}
\label{opnd:Immediate-classes}
{\scriptsize
\begin{tabular}{|l|l|l|} \hline
\bf{Immediate classes} & \bf{Operands} & \bf{What} \\ \hline
brknumber & brknumber & \verb| Break Number |\\
pcrel11s2 & pcrel11s2 & \verb| CCB Offset |\\
pcrel12s1 & i\_\discretionary{}{}{}12\_\discretionary{}{}{}i\_\discretionary{}{}{}11b\_\discretionary{}{}{}i\_\discretionary{}{}{}10\_\discretionary{}{}{}5\_\discretionary{}{}{}i\_\discretionary{}{}{}4\_\discretionary{}{}{}1 & \verb| PC-Relative Immediate 12 bits |\\
pcrel17s2 & pcrel17s2 & \verb| CB Offset |\\
pcrel20s1 & i\_\discretionary{}{}{}20\_\discretionary{}{}{}i\_\discretionary{}{}{}19\_\discretionary{}{}{}12\_\discretionary{}{}{}i\_\discretionary{}{}{}11j\_\discretionary{}{}{}i\_\discretionary{}{}{}10\_\discretionary{}{}{}1 & \verb| PC-Relative Immediate 20 bits |\\
pcrel27s2 & pcrel27s2 & \verb| GOTO / CALL Offset |\\
pcrel38s2 & upper27\_\discretionary{}{}{}lower11 & \verb| CCB Offset Extended |\\
pcrel44s2 & upper27\_\discretionary{}{}{}lower17 & \verb| CB Offset Extended |\\
pcrel54s2 & upper27\_\discretionary{}{}{}lower27 & \verb| GOTO Offset Extended |\\
signed10 & signed10 & \verb| Sign Extended Immediate 10 bits |\\
signed12 & i\_\discretionary{}{}{}11\_\discretionary{}{}{}0, i\_\discretionary{}{}{}11\_\discretionary{}{}{}5\_\discretionary{}{}{}i\_\discretionary{}{}{}4\_\discretionary{}{}{}0s & \verb| Sign Extended Immediate 12 bits |\\
signed16 & signed16 & \verb| Sign Extended Immediate 16 bits |\\
signed20 & i\_\discretionary{}{}{}31\_\discretionary{}{}{}12 & \verb| Sign Extended Immediate 20 bits |\\
signed27 & offset27 & \verb| Sign Extended Immediate 27 bits |\\
signed37 & upper27\_\discretionary{}{}{}lower10 & \verb| Sign Extended Immediate 37 bits |\\
signed43 & extend6\_\discretionary{}{}{}upper27\_\discretionary{}{}{}lower10 & \verb| Sign Extended Immediate 43 bits |\\
signed54 & extend27\_\discretionary{}{}{}offset27 & \verb| Sign Extended Immediate 54 bits |\\
signed6 & i\_\discretionary{}{}{}11\_\discretionary{}{}{}5 & \verb| Sign Extended Immediate 6 bits |\\
sysnumber & sysnumber & \verb| System Number |\\
unsigned5 & i\_\discretionary{}{}{}4\_\discretionary{}{}{}0, uimm & \verb| Zero Extended Immediate 5 bits |\\
unsigned6 & activate, i\_\discretionary{}{}{}5\_\discretionary{}{}{}0, startbit, stopbit2\_\discretionary{}{}{}stopbit4, unsigned6 & \verb| Zero Extended Immediate 6 bits |\\
wrapped32 & upper27\_\discretionary{}{}{}lower5 & \verb| Wrapped Immediate 32 bits |\\
wrapped64 & extend27\_\discretionary{}{}{}upper27\_\discretionary{}{}{}lower10 & \verb| Wrapped Immediate 64 bits |\\
\hline \end{tabular}
}
\newpage
{\scriptsize
\begin{tabular}{|l|l|l|l|l|} \hline
\bf{Immediate classes} & \bf{Extend} & \bf{Shift} & \bf{Width} & \bf{Values} \\ \hline
brknumber & Unsigned & 0 & 2 & 0x0 to 0x3 \\
pcrel11s2 & Signed & 2 & 11 & 0x1000 to 0xffc \\
pcrel12s1 & Signed & 1 & 12 & 0x1000 to 0xffe \\
pcrel17s2 & Signed & 2 & 17 & 0x40000 to 0x3fffc \\
pcrel20s1 & Signed & 1 & 20 & 0x100000 to 0xffffe \\
pcrel27s2 & Signed & 2 & 27 & 0x10000000 to 0xffffffc \\
pcrel38s2 & Signed & 2 & 38 & 0x8000000000 to 0x7ffffffffc \\
pcrel44s2 & Signed & 2 & 44 & 0x200000000000 to 0x1ffffffffffc \\
pcrel54s2 & Signed & 2 & 54 & 0x80000000000000 to 0x7ffffffffffffc \\
signed10 & Signed & 0 & 10 & 0x200 to 0x1ff \\
signed12 & Signed & 0 & 12 & 0x800 to 0x7ff \\
signed16 & Signed & 0 & 16 & 0x8000 to 0x7fff \\
signed20 & Signed & 0 & 20 & 0x80000 to 0x7ffff \\
signed27 & Signed & 0 & 27 & 0x4000000 to 0x3ffffff \\
signed37 & Signed & 0 & 37 & 0x1000000000 to 0xfffffffff \\
signed43 & Signed & 0 & 43 & 0x40000000000 to 0x3ffffffffff \\
signed54 & Signed & 0 & 54 & 0x20000000000000 to 0x1fffffffffffff \\
signed6 & Signed & 0 & 6 & 0x20 to 0x1f \\
sysnumber & Unsigned & 0 & 12 & 0x0 to 0xfff \\
unsigned5 & Unsigned & 0 & 5 & 0x0 to 0x1f \\
unsigned6 & Unsigned & 0 & 6 & 0x0 to 0x3f \\
wrapped32 & Wrap & 0 & 32 & 0x80000000 to 0xffffffff \\
wrapped64 & Wrap & 0 & 64 & 0x8000000000000000 to 0x7fffffffffffffff \\
\hline \end{tabular}
}
\newpage
\subsubsection{Operand Modifiers}
\label{opnd:Operand-Modifiers}
{\scriptsize
\begin{tabular}{|l|l|l|} \hline
\bf{Operand Modifiers} & \bf{Instances} \\ \hline
accesses &  \\
acqrel &  \\
bcucond & , ,  \\
boolcas &  \\
cachelev &  \\
ccbcomp &  \\
coherency &  \\
conjugate &  \\
floatcomp &  \\
floatmode &  \\
fnegate &  \\
froundmode &  \\
highmult &  \\
imultiply &  \\
intcomp &  \\
lanecond &  \\
lanesize &  \\
lanetodo &  \\
mostsig &  \\
oddlanes &  \\
ordering & ,  \\
qindex &  \\
signextw &  \\
splat32 &  \\
variant &  \\
widemult &  \\
ziplanes &  \\
\hline \end{tabular}
}
\newpage

\paragraph{Modifier} \textsf{accesses}:\\
\phantomsection
\label{par:modifier-accesses}
\noindent \begin{tabular}{|l|l|r|} \hline
\bf{Name} & \bf{Case} & \bf{Value} \\ \hline
. & All Accesses & 0 \\
.W & Write Accesses & 1 \\
.R & Read Accesses & 2 \\
.WA & Write Accesses (asynchronous) & 3 \\
\hline \end{tabular}
\newpage

\paragraph{Modifier} \textsf{acqrel}:\\
\phantomsection
\label{par:modifier-acqrel}
\noindent \begin{tabular}{|l|l|r|} \hline
\bf{Name} & \bf{Case} & \bf{Value} \\ \hline
. & None & 0 \\
.RL & Release & 1 \\
.AQ & Acquire & 2 \\
.AQRL & Acquire/Release & 3 \\
\hline \end{tabular}
\newpage

\paragraph{Modifier} \textsf{bcucond}:\\
\phantomsection
\label{par:modifier-bcucond}
\noindent \begin{tabular}{|l|l|r|} \hline
\bf{Name} & \bf{Case} & \bf{Value} \\ \hline
.DLTZ & Double Less Than Zero & 0 \\
.DGEZ & Double Greater Than or Equal to Zero & 1 \\
.DLEZ & Double Less Than or Equal to Zero & 2 \\
.DGTZ & Double Greater Than Zero & 3 \\
.DEQZ & Double Equal to Zero & 4 \\
.DNEZ & Double Not Equal to Zero & 5 \\
.ODD & Odd (LSB Set) & 6 \\
.EVEN & Even (LSB Clear) & 7 \\
.WLTZ & Word Less Than Zero & 8 \\
.WGEZ & Word Greater Than or Equal to Zero & 9 \\
.WLEZ & Word Less Than or Equal to Zero & 10 \\
.WGTZ & Word Greater Than Zero & 11 \\
.WEQZ & Word Equal to Zero & 12 \\
.WNEZ & Word Not Equal to Zero & 13 \\
\hline \end{tabular}
\newpage

\paragraph{Modifier} \textsf{boolcas}:\\
\phantomsection
\label{par:modifier-boolcas}
\noindent \begin{tabular}{|l|l|r|} \hline
\bf{Name} & \bf{Case} & \bf{Value} \\ \hline
.V & Valued CAS & 0 \\
. & Boolean CAS & 1 \\
\hline \end{tabular}
\newpage

\paragraph{Modifier} \textsf{cachelev}:\\
\phantomsection
\label{par:modifier-cachelev}
\noindent \begin{tabular}{|l|l|r|} \hline
\bf{Name} & \bf{Case} & \bf{Value} \\ \hline
.L1 & L1 Cache & 0 \\
.L2 & L2 Cache & 1 \\
\hline \end{tabular}
\newpage

\paragraph{Modifier} \textsf{ccbcomp}:\\
\phantomsection
\label{par:modifier-ccbcomp}
\noindent \begin{tabular}{|l|l|r|} \hline
\bf{Name} & \bf{Case} & \bf{Value} \\ \hline
.DLT & Double Word Less Than & 0 \\
.DGE & Double Word Greater Than or Equal & 1 \\
.DLTU & Double Word Less Than Unsigned & 2 \\
.DGEU & Double Word Greater Than or Equal Unsigned & 3 \\
.DEQ & Double Word Equal & 4 \\
.DNE & Double Word Not Equal & 5 \\
.DANY & Double Any bit set after AND-ing & 6 \\
.DNONE & Double No bits set after AND-ing & 7 \\
.WLT & Word Less Than & 8 \\
.WGE & Word Greater Than or Equal & 9 \\
.WLTU & Word Less Than Unsigned & 10 \\
.WGEU & Word Greater Than or Equal Unsigned & 11 \\
.WEQ & Word Equal & 12 \\
.WNE & Word Not Equal & 13 \\
.WANY & Word Any bit set after AND-ing & 14 \\
.WNONE & Word No bits set after AND-ing & 15 \\
\hline \end{tabular}
\newpage

\paragraph{Modifier} \textsf{coherency}:\\
\phantomsection
\label{par:modifier-coherency}
\noindent \begin{tabular}{|l|l|r|} \hline
\bf{Name} & \bf{Case} & \bf{Value} \\ \hline
. & Local & 0 \\
.G & Global & 1 \\
.S & Reserved & 2 \\
\hline \end{tabular}
\newpage

\paragraph{Modifier} \textsf{conjugate}:\\
\phantomsection
\label{par:modifier-conjugate}
\noindent \begin{tabular}{|l|l|r|} \hline
\bf{Name} & \bf{Case} & \bf{Value} \\ \hline
. & Default & 0 \\
.C & Conjugate & 1 \\
\hline \end{tabular}
\newpage

\paragraph{Modifier} \textsf{floatcomp}:\\
\phantomsection
\label{par:modifier-floatcomp}
\noindent \begin{tabular}{|l|l|r|} \hline
\bf{Name} & \bf{Case} & \bf{Value} \\ \hline
.ONE & Ordered and Not Equal & 0 \\
.UEQ & Unordered or Equal & 1 \\
.OEQ & Ordered and Equal & 2 \\
.UNE & Unordered or Not Equal & 3 \\
.OLT & Ordered and Less Than & 4 \\
.UGE & Unordered or Greater Than or Equal & 5 \\
.OGE & Ordered and Greater Than or Equal & 6 \\
.ULT & Unordered or Less Than & 7 \\
\hline \end{tabular}
\newpage

\paragraph{Modifier} \textsf{floatmode}:\\
\phantomsection
\label{par:modifier-floatmode}
\noindent \begin{tabular}{|l|l|r|} \hline
\bf{Name} & \bf{Case} & \bf{Value} \\ \hline
.RN & Round to Nearest, ties to Even & 0 \\
.RZ & Round toward Zero & 1 \\
.RD & Round Downward & 2 \\
.RU & Round Upward & 3 \\
.RM & Round to Nearest, ties to Max Magnitude & 4 \\
.R5 & Reserved 5 & 5 \\
.RO & Round to Odd & 6 \\
. & Use CS rounding & 7 \\
\hline \end{tabular}
\newpage

\paragraph{Modifier} \textsf{fnegate}:\\
\phantomsection
\label{par:modifier-fnegate}
\noindent \begin{tabular}{|l|l|r|} \hline
\bf{Name} & \bf{Case} & \bf{Value} \\ \hline
. & Default Result & 0 \\
.N & Negate Result & 1 \\
\hline \end{tabular}
\newpage

\paragraph{Modifier} \textsf{froundmode}:\\
\phantomsection
\label{par:modifier-froundmode}
\noindent \begin{tabular}{|l|l|r|} \hline
\bf{Name} & \bf{Case} & \bf{Value} \\ \hline
RNE & Round to Nearest, ties to Even & 0 \\
RTZ & Round Towards Zero & 1 \\
RDN & Round Down (towards -Inf) & 2 \\
RUP & Round Up (towards +Inf) & 3 \\
RMM & Round to Nearest, ties to Max Magnitude & 4 \\
\_RES5 & Reserved value & 5 \\
\_RES6 & Reserved value & 6 \\
DYN & Dynamic rounding mode & 7 \\
\hline \end{tabular}
\newpage

\paragraph{Modifier} \textsf{highmult}:\\
\phantomsection
\label{par:modifier-highmult}
\noindent \begin{tabular}{|l|l|r|} \hline
\bf{Name} & \bf{Case} & \bf{Value} \\ \hline
. & Low & 0 \\
.H & High & 1 \\
.HU & High Unsigned & 2 \\
.HSU & High Signed Unsigned & 3 \\
\hline \end{tabular}
\newpage

\paragraph{Modifier} \textsf{imultiply}:\\
\phantomsection
\label{par:modifier-imultiply}
\noindent \begin{tabular}{|l|l|r|} \hline
\bf{Name} & \bf{Case} & \bf{Value} \\ \hline
. & Default & 0 \\
.MI & Multiply by I & 1 \\
\hline \end{tabular}
\newpage

\paragraph{Modifier} \textsf{intcomp}:\\
\phantomsection
\label{par:modifier-intcomp}
\noindent \begin{tabular}{|l|l|r|} \hline
\bf{Name} & \bf{Case} & \bf{Value} \\ \hline
.LT & Less Than & 0 \\
.GE & Greater Than or Equal & 1 \\
.LTU & Less Than Unsigned & 2 \\
.GEU & Greater Than or Equal Unsigned & 3 \\
.EQ & Equal & 4 \\
.NE & Not Equal & 5 \\
.ANY & Any bits set in after AND-ing & 6 \\
.NONE & No bits set after AND-ing & 7 \\
.LE & Less Than or Equal & 8 \\
.GT & Greater Than & 9 \\
.LEU & Less Than or Equal Unsigned & 10 \\
.GTU & Greater Than Unsigned & 11 \\
\hline \end{tabular}
\newpage

\paragraph{Modifier} \textsf{lanecond}:\\
\phantomsection
\label{par:modifier-lanecond}
\noindent \begin{tabular}{|l|l|r|} \hline
\bf{Name} & \bf{Case} & \bf{Value} \\ \hline
.LTZ & Less Than Zero & 0 \\
.GEZ & Greater Than or Equal to Zero & 1 \\
.LEZ & Less Than or Equal to Zero & 2 \\
.GTZ & Greater Than Zero & 3 \\
.EQZ & Equal to Zero & 4 \\
.NEZ & Not Equal to Zero & 5 \\
.ODD & Odd (LSB Set) & 6 \\
.EVEN & Even (LSB Clear) & 7 \\
\hline \end{tabular}
\newpage

\paragraph{Modifier} \textsf{lanesize}:\\
\phantomsection
\label{par:modifier-lanesize}
\noindent \begin{tabular}{|l|l|r|} \hline
\bf{Name} & \bf{Case} & \bf{Value} \\ \hline
. & Byte Size & 0 \\
.H & Half Word Size & 1 \\
.W & Word Size & 2 \\
.D & Double Word Size & 3 \\
\hline \end{tabular}
\newpage

\paragraph{Modifier} \textsf{lanetodo}:\\
\phantomsection
\label{par:modifier-lanetodo}
\noindent \begin{tabular}{|l|l|r|} \hline
\bf{Name} & \bf{Case} & \bf{Value} \\ \hline
.MT & Mask True & 0 \\
.MF & Mask False & 1 \\
.MTC & Mask True or Clear & 2 \\
.MFC & Mask False or Clear & 3 \\
\hline \end{tabular}
\newpage

\paragraph{Modifier} \textsf{mostsig}:\\
\phantomsection
\label{par:modifier-mostsig}
\noindent \begin{tabular}{|l|l|r|} \hline
\bf{Name} & \bf{Case} & \bf{Value} \\ \hline
. & Least Significant Bits & 0 \\
.M & Most Significant Bits & 1 \\
\hline \end{tabular}
\newpage

\paragraph{Modifier} \textsf{oddlanes}:\\
\phantomsection
\label{par:modifier-oddlanes}
\noindent \begin{tabular}{|l|l|r|} \hline
\bf{Name} & \bf{Case} & \bf{Value} \\ \hline
. & Even Lanes & 0 \\
.O & Odd Lanes & 1 \\
\hline \end{tabular}
\newpage

\paragraph{Modifier} \textsf{ordering}:\\
\phantomsection
\label{par:modifier-ordering}
\noindent \begin{tabular}{|l|l|r|} \hline
\bf{Name} & \bf{Case} & \bf{Value} \\ \hline
.NONE & None & 0 \\
.W & Write & 1 \\
.R & Read & 2 \\
.RW & Read Write & 3 \\
.O & Out & 4 \\
.OW & Out Write & 5 \\
.OR & Out Read & 6 \\
.ORW & Out Read Write & 7 \\
.I & In & 8 \\
.IW & In Write & 9 \\
.IR & In Read & 10 \\
.IRW & In Read Write & 11 \\
.IO & In Out & 12 \\
.IOW & In Out Write & 13 \\
.IOR & In Out Read & 14 \\
.IORW & In Out Read Write & 15 \\
\hline \end{tabular}
\newpage

\paragraph{Modifier} \textsf{qindex}:\\
\phantomsection
\label{par:modifier-qindex}
\noindent \begin{tabular}{|l|l|r|} \hline
\bf{Name} & \bf{Case} & \bf{Value} \\ \hline
.Q0 & Quarter 0 & 0 \\
.Q1 & Quarter 1 & 1 \\
.Q2 & Quarter 2 & 2 \\
.Q3 & Quarter 3 & 3 \\
\hline \end{tabular}
\newpage

\paragraph{Modifier} \textsf{signextw}:\\
\phantomsection
\label{par:modifier-signextw}
\noindent \begin{tabular}{|l|l|r|} \hline
\bf{Name} & \bf{Case} & \bf{Value} \\ \hline
. & Zero Extend & 0 \\
.SX & Sign Extend & 1 \\
\hline \end{tabular}
\newpage

\paragraph{Modifier} \textsf{splat32}:\\
\phantomsection
\label{par:modifier-splat32}
\noindent \begin{tabular}{|l|l|r|} \hline
\bf{Name} & \bf{Case} & \bf{Value} \\ \hline
. & No splat & 0 \\
.@ & Splat x2 & 1 \\
\hline \end{tabular}
\newpage

\paragraph{Modifier} \textsf{variant}:\\
\phantomsection
\label{par:modifier-variant}
\noindent \begin{tabular}{|l|l|r|} \hline
\bf{Name} & \bf{Case} & \bf{Value} \\ \hline
. & Cached & 0 \\
.S & Speculative & 1 \\
.U & Uncached & 2 \\
.US & Uncached Speculative & 3 \\
\hline \end{tabular}
\newpage

\paragraph{Modifier} \textsf{widemult}:\\
\phantomsection
\label{par:modifier-widemult}
\noindent \begin{tabular}{|l|l|r|} \hline
\bf{Name} & \bf{Case} & \bf{Value} \\ \hline
. & Signed & 0 \\
.U & Unsigned & 1 \\
.SU & Signed by Unsigned & 2 \\
\hline \end{tabular}
\newpage

\paragraph{Modifier} \textsf{ziplanes}:\\
\phantomsection
\label{par:modifier-ziplanes}
\noindent \begin{tabular}{|l|l|r|} \hline
\bf{Name} & \bf{Case} & \bf{Value} \\ \hline
. & No zip & 0 \\
.Z & Zip lanes & 1 \\
\hline \end{tabular}
\newpage
\subsubsection{Register Classes}
\label{opnd:Register-Classes}
{\scriptsize
\begin{tabular}{|l|l|l|} \hline
\bf{Register Classes} & \bf{Instances} & \bf{Width} \\ \hline
aloneReg & systemAlone & 64 \\
buffer16Reg & registerCj, registerGj & 0 \\
buffer2Reg & registerCg, registerGg & 0 \\
buffer32Reg & registerCk, registerGk & 0 \\
buffer4Reg & registerCh, registerGh & 0 \\
buffer64Reg & registerCl, registerGl & 0 \\
buffer8Reg & registerCi, registerGi & 0 \\
csReg & csr & 64 \\
floatReg & frd, frs1, frs2, frs3 & 64 \\
mainReg & rd, rs1, rs2 & 64 \\
mainRegPair & rd\_\discretionary{}{}{}pair, rs2\_\discretionary{}{}{}pair & 128 \\
onlyfxReg & systemT2 & 64 \\
onlygetReg & systemS2 & 64 \\
onlyraReg & systemRA & 64 \\
onlysetReg & systemT3 & 64 \\
onlyswapReg & systemS4 & 64 \\
pairedReg & registerM, registerO, registerP, registerU & 128 \\
quadReg & registerN, registerQ, registerR, registerV & 256 \\
singleReg & registerT, registerW, registerY, registerZ & 64 \\
worddRegE & registerWe, registerYe, registerZe & 64 \\
worddRegO & registerWo, registerYo, registerZo & 64 \\
xworddReg & registerCc\_\discretionary{}{}{}qselectC & 64 \\
xworddReg0M4 & registerAx & 64 \\
xworddReg1M4 & registerAy & 64 \\
xworddReg2M4 & registerAz & 64 \\
xworddReg3M4 & registerAt & 64 \\
xwordoReg & registerA, registerC, registerE, registerG & 256 \\
xwordqReg & registerCb\_\discretionary{}{}{}hselectC & 128 \\
xwordqRegE & registerAE & 128 \\
xwordqRegO & registerAO & 128 \\
xwordvReg & registerEq, registerGq & 1024 \\
\hline \end{tabular}
}
\newpage
{\scriptsize
\begin{tabular}{|l|l|l|} \hline
\bf{Register Class} & \bf{Register File} & \bf{Registers} \\ \hline
 aloneReg & SFR & \texttt{ } \\
 & & \texttt{ } \\
 & & \texttt{ } \\
 & & \texttt{ } \\
 & & \texttt{ } \\
 & & \texttt{ } \\
 & & \texttt{ } \\
 & & \texttt{ } \\
 & & \texttt{ } \\
 & & \texttt{1 } \\
 & & \texttt{1 } \\
 & & \texttt{1 } \\
 & & \texttt{1 } \\
\hline \end{tabular}
}
\newpage
{\scriptsize
\begin{tabular}{|l|l|l|} \hline
\bf{Register Class} & \bf{Register File} & \bf{Registers} \\ \hline
 buffer16Reg & X16R & \texttt{} \\
 & & \texttt{} \\
 & & \texttt{} \\
 & & \texttt{} \\
\hline \end{tabular}
}
\newpage
{\scriptsize
\begin{tabular}{|l|l|l|} \hline
\bf{Register Class} & \bf{Register File} & \bf{Registers} \\ \hline
 buffer2Reg & X2R & \texttt{} \\
 & & \texttt{} \\
 & & \texttt{} \\
 & & \texttt{} \\
 & & \texttt{} \\
 & & \texttt{} \\
 & & \texttt{} \\
 & & \texttt{} \\
 & & \texttt{} \\
 & & \texttt{} \\
 & & \texttt{} \\
 & & \texttt{} \\
 & & \texttt{} \\
 & & \texttt{} \\
 & & \texttt{} \\
 & & \texttt{} \\
 & & \texttt{} \\
 & & \texttt{} \\
 & & \texttt{} \\
 & & \texttt{} \\
 & & \texttt{} \\
 & & \texttt{} \\
 & & \texttt{} \\
 & & \texttt{} \\
 & & \texttt{} \\
 & & \texttt{} \\
 & & \texttt{} \\
 & & \texttt{} \\
 & & \texttt{} \\
 & & \texttt{} \\
 & & \texttt{} \\
 & & \texttt{} \\
\hline \end{tabular}
}
\newpage
{\scriptsize
\begin{tabular}{|l|l|l|} \hline
\bf{Register Class} & \bf{Register File} & \bf{Registers} \\ \hline
 buffer32Reg & X32R & \texttt{} \\
 & & \texttt{} \\
\hline \end{tabular}
}
\newpage
{\scriptsize
\begin{tabular}{|l|l|l|} \hline
\bf{Register Class} & \bf{Register File} & \bf{Registers} \\ \hline
 buffer4Reg & X4R & \texttt{} \\
 & & \texttt{} \\
 & & \texttt{} \\
 & & \texttt{} \\
 & & \texttt{} \\
 & & \texttt{} \\
 & & \texttt{} \\
 & & \texttt{} \\
 & & \texttt{} \\
 & & \texttt{} \\
 & & \texttt{} \\
 & & \texttt{} \\
 & & \texttt{} \\
 & & \texttt{} \\
 & & \texttt{} \\
 & & \texttt{} \\
\hline \end{tabular}
}
\newpage
{\scriptsize
\begin{tabular}{|l|l|l|} \hline
\bf{Register Class} & \bf{Register File} & \bf{Registers} \\ \hline
 buffer64Reg & X64R & \texttt{} \\
\hline \end{tabular}
}
\newpage
{\scriptsize
\begin{tabular}{|l|l|l|} \hline
\bf{Register Class} & \bf{Register File} & \bf{Registers} \\ \hline
 buffer8Reg & X8R & \texttt{} \\
 & & \texttt{} \\
 & & \texttt{} \\
 & & \texttt{} \\
 & & \texttt{} \\
 & & \texttt{} \\
 & & \texttt{} \\
 & & \texttt{} \\
\hline \end{tabular}
}
\newpage
{\scriptsize
\begin{tabular}{|l|l|l|} \hline
\bf{Register Class} & \bf{Register File} & \bf{Registers} \\ \hline
 csReg & CSR & \texttt{1} \\
 & & \texttt{1} \\
 & & \texttt{1} \\
 & & \texttt{1} \\
 & & \texttt{1} \\
 & & \texttt{1} \\
 & & \texttt{1} \\
 & & \texttt{1} \\
 & & \texttt{1} \\
 & & \texttt{1} \\
 & & \texttt{1} \\
 & & \texttt{1} \\
 & & \texttt{1} \\
 & & \texttt{1} \\
 & & \texttt{1} \\
 & & \texttt{1} \\
 & & \texttt{1} \\
 & & \texttt{1} \\
 & & \texttt{1} \\
 & & \texttt{1} \\
 & & \texttt{1} \\
 & & \texttt{1} \\
 & & \texttt{1} \\
 & & \texttt{1} \\
 & & \texttt{1} \\
 & & \texttt{1} \\
 & & \texttt{1} \\
 & & \texttt{1} \\
 & & \texttt{1} \\
 & & \texttt{1} \\
 & & \texttt{1} \\
 & & \texttt{1} \\
 & & \texttt{1} \\
 & & \texttt{1} \\
 & & \texttt{1} \\
 & & \texttt{1} \\
 & & \texttt{1} \\
 & & \texttt{1} \\
 & & \texttt{1} \\
 & & \texttt{1} \\
 & & \texttt{1} \\
 & & \texttt{1} \\
 & & \texttt{1} \\
 & & \texttt{1} \\
 & & \texttt{1} \\
 & & \texttt{1} \\
 & & \texttt{1} \\
 & & \texttt{1} \\
 & & \texttt{1} \\
 & & \texttt{1} \\
 & & \texttt{1} \\
 & & \texttt{1} \\
 & & \texttt{1} \\
 & & \texttt{1} \\
 & & \texttt{1} \\
 & & \texttt{1} \\
 & & \texttt{1} \\
 & & \texttt{1} \\
 & & \texttt{1} \\
 & & \texttt{1} \\
 & & \texttt{1} \\
 & & \texttt{1} \\
 & & \texttt{1} \\
 & & \texttt{1} \\
 & & \texttt{1} \\
 & & \texttt{1} \\
 & & \texttt{1} \\
 & & \texttt{1} \\
 & & \texttt{1} \\
 & & \texttt{1} \\
 & & \texttt{1} \\
 & & \texttt{1} \\
 & & \texttt{1} \\
 & & \texttt{1} \\
 & & \texttt{1} \\
 & & \texttt{1} \\
 & & \texttt{1} \\
 & & \texttt{1} \\
 & & \texttt{1} \\
 & & \texttt{1} \\
 & & \texttt{1} \\
 & & \texttt{1} \\
 & & \texttt{1} \\
 & & \texttt{1} \\
 & & \texttt{1} \\
 & & \texttt{1} \\
 & & \texttt{1} \\
 & & \texttt{1} \\
 & & \texttt{1} \\
 & & \texttt{1} \\
 & & \texttt{1} \\
 & & \texttt{1} \\
 & & \texttt{1} \\
 & & \texttt{1} \\
 & & \texttt{1} \\
 & & \texttt{1} \\
 & & \texttt{1} \\
 & & \texttt{1} \\
 & & \texttt{1} \\
 & & \texttt{1} \\
 & & \texttt{1} \\
 & & \texttt{1} \\
 & & \texttt{1} \\
 & & \texttt{1} \\
 & & \texttt{1} \\
 & & \texttt{1} \\
 & & \texttt{1} \\
 & & \texttt{1} \\
 & & \texttt{1} \\
 & & \texttt{1} \\
 & & \texttt{1} \\
 & & \texttt{1} \\
 & & \texttt{1} \\
 & & \texttt{1} \\
 & & \texttt{1} \\
 & & \texttt{1} \\
 & & \texttt{1} \\
 & & \texttt{1} \\
 & & \texttt{1} \\
 & & \texttt{1} \\
 & & \texttt{1} \\
 & & \texttt{1} \\
 & & \texttt{1} \\
 & & \texttt{1} \\
 & & \texttt{1} \\
 & & \texttt{1} \\
 & & \texttt{1} \\
 & & \texttt{1} \\
 & & \texttt{1} \\
 & & \texttt{1} \\
 & & \texttt{1} \\
 & & \texttt{1} \\
 & & \texttt{1} \\
 & & \texttt{1} \\
 & & \texttt{1} \\
 & & \texttt{1} \\
 & & \texttt{1} \\
 & & \texttt{1} \\
 & & \texttt{1} \\
 & & \texttt{1} \\
 & & \texttt{1} \\
 & & \texttt{1} \\
 & & \texttt{1} \\
 & & \texttt{1} \\
 & & \texttt{1} \\
 & & \texttt{1} \\
 & & \texttt{1} \\
 & & \texttt{1} \\
 & & \texttt{1} \\
 & & \texttt{1} \\
 & & \texttt{1} \\
 & & \texttt{1} \\
 & & \texttt{1} \\
 & & \texttt{1} \\
 & & \texttt{1} \\
 & & \texttt{1} \\
 & & \texttt{1} \\
 & & \texttt{1} \\
 & & \texttt{1} \\
 & & \texttt{1} \\
 & & \texttt{1} \\
 & & \texttt{1} \\
 & & \texttt{1} \\
 & & \texttt{1} \\
 & & \texttt{1} \\
 & & \texttt{1} \\
 & & \texttt{1} \\
 & & \texttt{1} \\
 & & \texttt{1} \\
 & & \texttt{1} \\
 & & \texttt{1} \\
 & & \texttt{1} \\
 & & \texttt{1} \\
 & & \texttt{1} \\
 & & \texttt{1} \\
 & & \texttt{1} \\
 & & \texttt{1} \\
 & & \texttt{1} \\
 & & \texttt{1} \\
 & & \texttt{1} \\
 & & \texttt{1} \\
 & & \texttt{1} \\
 & & \texttt{1} \\
 & & \texttt{1} \\
 & & \texttt{1} \\
 & & \texttt{1} \\
 & & \texttt{1} \\
 & & \texttt{1} \\
 & & \texttt{1} \\
 & & \texttt{1} \\
 & & \texttt{1} \\
 & & \texttt{1} \\
 & & \texttt{1} \\
 & & \texttt{1} \\
 & & \texttt{1} \\
 & & \texttt{1} \\
 & & \texttt{1} \\
 & & \texttt{1} \\
 & & \texttt{1} \\
 & & \texttt{1} \\
 & & \texttt{1} \\
 & & \texttt{1} \\
 & & \texttt{1} \\
 & & \texttt{1} \\
 & & \texttt{1} \\
 & & \texttt{1} \\
 & & \texttt{1} \\
 & & \texttt{1} \\
 & & \texttt{1} \\
 & & \texttt{1} \\
 & & \texttt{1} \\
 & & \texttt{1} \\
 & & \texttt{1} \\
 & & \texttt{1} \\
 & & \texttt{1} \\
 & & \texttt{1} \\
 & & \texttt{1} \\
 & & \texttt{1} \\
 & & \texttt{1} \\
 & & \texttt{1} \\
 & & \texttt{1} \\
 & & \texttt{1} \\
 & & \texttt{1} \\
 & & \texttt{1} \\
 & & \texttt{1} \\
 & & \texttt{1} \\
 & & \texttt{1} \\
 & & \texttt{1} \\
 & & \texttt{1} \\
 & & \texttt{1} \\
 & & \texttt{1} \\
 & & \texttt{1} \\
 & & \texttt{1} \\
 & & \texttt{1} \\
 & & \texttt{1} \\
 & & \texttt{1} \\
 & & \texttt{1} \\
 & & \texttt{1} \\
 & & \texttt{1} \\
 & & \texttt{1} \\
 & & \texttt{1} \\
 & & \texttt{1} \\
 & & \texttt{1} \\
 & & \texttt{1} \\
 & & \texttt{1} \\
 & & \texttt{1} \\
 & & \texttt{1} \\
 & & \texttt{1} \\
 & & \texttt{1} \\
 & & \texttt{1} \\
 & & \texttt{1} \\
 & & \texttt{1} \\
 & & \texttt{1} \\
 & & \texttt{1} \\
 & & \texttt{1} \\
 & & \texttt{1} \\
 & & \texttt{1} \\
 & & \texttt{1} \\
 & & \texttt{1} \\
 & & \texttt{1} \\
 & & \texttt{1} \\
 & & \texttt{1} \\
 & & \texttt{1} \\
 & & \texttt{1} \\
 & & \texttt{1} \\
 & & \texttt{1} \\
 & & \texttt{1} \\
 & & \texttt{1} \\
 & & \texttt{1} \\
 & & \texttt{1} \\
 & & \texttt{1} \\
 & & \texttt{1} \\
 & & \texttt{1} \\
 & & \texttt{1} \\
 & & \texttt{1} \\
 & & \texttt{1} \\
 & & \texttt{1} \\
 & & \texttt{1} \\
 & & \texttt{1} \\
 & & \texttt{1} \\
 & & \texttt{1} \\
 & & \texttt{1} \\
 & & \texttt{1} \\
 & & \texttt{1} \\
 & & \texttt{1} \\
 & & \texttt{1} \\
 & & \texttt{1} \\
 & & \texttt{1} \\
 & & \texttt{1} \\
 & & \texttt{1} \\
 & & \texttt{1} \\
 & & \texttt{1} \\
 & & \texttt{1} \\
 & & \texttt{1} \\
 & & \texttt{1} \\
 & & \texttt{1} \\
 & & \texttt{1} \\
 & & \texttt{1} \\
 & & \texttt{1} \\
 & & \texttt{1} \\
 & & \texttt{1} \\
 & & \texttt{1} \\
 & & \texttt{1} \\
 & & \texttt{1} \\
 & & \texttt{1} \\
 & & \texttt{1} \\
 & & \texttt{1} \\
 & & \texttt{1} \\
 & & \texttt{1} \\
 & & \texttt{1} \\
 & & \texttt{1} \\
 & & \texttt{1} \\
 & & \texttt{1} \\
 & & \texttt{1} \\
 & & \texttt{1} \\
 & & \texttt{1} \\
 & & \texttt{1} \\
 & & \texttt{1} \\
 & & \texttt{1} \\
 & & \texttt{1} \\
 & & \texttt{1} \\
 & & \texttt{1} \\
 & & \texttt{1} \\
 & & \texttt{1} \\
 & & \texttt{1} \\
 & & \texttt{1} \\
 & & \texttt{1} \\
 & & \texttt{1} \\
 & & \texttt{1} \\
 & & \texttt{1} \\
 & & \texttt{1} \\
 & & \texttt{1} \\
 & & \texttt{1} \\
 & & \texttt{1} \\
 & & \texttt{1} \\
 & & \texttt{1} \\
 & & \texttt{1} \\
 & & \texttt{1} \\
 & & \texttt{1} \\
 & & \texttt{1} \\
 & & \texttt{1} \\
 & & \texttt{1} \\
 & & \texttt{1} \\
 & & \texttt{1} \\
 & & \texttt{1} \\
 & & \texttt{1} \\
 & & \texttt{1} \\
 & & \texttt{1} \\
 & & \texttt{1} \\
 & & \texttt{1} \\
 & & \texttt{1} \\
 & & \texttt{1} \\
 & & \texttt{1} \\
 & & \texttt{1} \\
 & & \texttt{1} \\
 & & \texttt{1} \\
 & & \texttt{1} \\
 & & \texttt{1} \\
 & & \texttt{1} \\
 & & \texttt{1} \\
 & & \texttt{1} \\
 & & \texttt{1} \\
 & & \texttt{1} \\
 & & \texttt{1} \\
 & & \texttt{1} \\
 & & \texttt{1} \\
 & & \texttt{1} \\
 & & \texttt{1} \\
 & & \texttt{1} \\
 & & \texttt{1} \\
 & & \texttt{1} \\
 & & \texttt{1} \\
 & & \texttt{1} \\
 & & \texttt{1} \\
 & & \texttt{1} \\
 & & \texttt{1} \\
 & & \texttt{1} \\
 & & \texttt{1} \\
 & & \texttt{1} \\
 & & \texttt{1} \\
 & & \texttt{1} \\
 & & \texttt{1} \\
 & & \texttt{1} \\
 & & \texttt{1} \\
 & & \texttt{1} \\
 & & \texttt{1} \\
 & & \texttt{1} \\
 & & \texttt{1} \\
 & & \texttt{1} \\
 & & \texttt{1} \\
 & & \texttt{1} \\
 & & \texttt{1} \\
 & & \texttt{1} \\
 & & \texttt{1} \\
 & & \texttt{1} \\
 & & \texttt{1} \\
 & & \texttt{1} \\
 & & \texttt{1} \\
 & & \texttt{1} \\
 & & \texttt{1} \\
 & & \texttt{1} \\
 & & \texttt{1} \\
 & & \texttt{1} \\
 & & \texttt{1} \\
 & & \texttt{1} \\
 & & \texttt{1} \\
 & & \texttt{1} \\
 & & \texttt{1} \\
 & & \texttt{1} \\
 & & \texttt{1} \\
 & & \texttt{1} \\
 & & \texttt{1} \\
 & & \texttt{1} \\
 & & \texttt{1} \\
 & & \texttt{1} \\
 & & \texttt{1} \\
 & & \texttt{1} \\
 & & \texttt{1} \\
 & & \texttt{1} \\
 & & \texttt{1} \\
 & & \texttt{1} \\
 & & \texttt{1} \\
 & & \texttt{1} \\
 & & \texttt{1} \\
 & & \texttt{1} \\
 & & \texttt{1} \\
 & & \texttt{1} \\
 & & \texttt{1} \\
 & & \texttt{1} \\
 & & \texttt{1} \\
 & & \texttt{1} \\
 & & \texttt{1} \\
 & & \texttt{1} \\
 & & \texttt{1} \\
 & & \texttt{1} \\
 & & \texttt{1} \\
 & & \texttt{1} \\
 & & \texttt{1} \\
 & & \texttt{1} \\
 & & \texttt{1} \\
 & & \texttt{1} \\
 & & \texttt{1} \\
 & & \texttt{1} \\
 & & \texttt{1} \\
 & & \texttt{1} \\
 & & \texttt{1} \\
 & & \texttt{1} \\
 & & \texttt{1} \\
 & & \texttt{1} \\
 & & \texttt{1} \\
 & & \texttt{1} \\
 & & \texttt{1} \\
 & & \texttt{1} \\
 & & \texttt{1} \\
 & & \texttt{1} \\
 & & \texttt{1} \\
 & & \texttt{1} \\
 & & \texttt{1} \\
 & & \texttt{1} \\
 & & \texttt{1} \\
 & & \texttt{1} \\
 & & \texttt{1} \\
 & & \texttt{1} \\
 & & \texttt{1} \\
 & & \texttt{1} \\
 & & \texttt{1} \\
 & & \texttt{1} \\
 & & \texttt{1} \\
 & & \texttt{1} \\
 & & \texttt{1} \\
 & & \texttt{1} \\
 & & \texttt{1} \\
 & & \texttt{1} \\
 & & \texttt{1} \\
 & & \texttt{1} \\
 & & \texttt{1} \\
 & & \texttt{1} \\
 & & \texttt{1} \\
 & & \texttt{1} \\
 & & \texttt{1} \\
 & & \texttt{1} \\
 & & \texttt{1} \\
 & & \texttt{1} \\
 & & \texttt{1} \\
 & & \texttt{1} \\
 & & \texttt{1} \\
 & & \texttt{1} \\
 & & \texttt{1} \\
 & & \texttt{1} \\
 & & \texttt{1} \\
 & & \texttt{1} \\
 & & \texttt{1} \\
 & & \texttt{1} \\
 & & \texttt{1} \\
 & & \texttt{1} \\
 & & \texttt{1} \\
 & & \texttt{1} \\
 & & \texttt{1} \\
 & & \texttt{1} \\
 & & \texttt{1} \\
 & & \texttt{1} \\
 & & \texttt{1} \\
 & & \texttt{1} \\
 & & \texttt{1} \\
 & & \texttt{1} \\
 & & \texttt{1} \\
 & & \texttt{1} \\
 & & \texttt{1} \\
 & & \texttt{1} \\
 & & \texttt{1} \\
 & & \texttt{1} \\
 & & \texttt{1} \\
 & & \texttt{1} \\
 & & \texttt{1} \\
 & & \texttt{1} \\
 & & \texttt{1} \\
 & & \texttt{1} \\
 & & \texttt{1} \\
 & & \texttt{1} \\
 & & \texttt{1} \\
 & & \texttt{1} \\
 & & \texttt{1} \\
 & & \texttt{1} \\
 & & \texttt{1} \\
 & & \texttt{1} \\
 & & \texttt{1} \\
 & & \texttt{1} \\
 & & \texttt{1} \\
 & & \texttt{1} \\
 & & \texttt{1} \\
 & & \texttt{1} \\
 & & \texttt{1} \\
 & & \texttt{1} \\
 & & \texttt{1} \\
 & & \texttt{1} \\
 & & \texttt{1} \\
 & & \texttt{1} \\
 & & \texttt{1} \\
 & & \texttt{1} \\
 & & \texttt{1} \\
 & & \texttt{1} \\
 & & \texttt{1} \\
 & & \texttt{1} \\
 & & \texttt{1} \\
 & & \texttt{1} \\
 & & \texttt{1} \\
 & & \texttt{1} \\
 & & \texttt{1} \\
 & & \texttt{1} \\
 & & \texttt{1} \\
 & & \texttt{1} \\
 & & \texttt{1} \\
 & & \texttt{1} \\
 & & \texttt{1} \\
 & & \texttt{1} \\
 & & \texttt{1} \\
 & & \texttt{1} \\
 & & \texttt{1} \\
 & & \texttt{1} \\
 & & \texttt{1} \\
 & & \texttt{1} \\
 & & \texttt{1} \\
 & & \texttt{1} \\
 & & \texttt{1} \\
 & & \texttt{1} \\
 & & \texttt{1} \\
 & & \texttt{1} \\
 & & \texttt{1} \\
 & & \texttt{1} \\
 & & \texttt{1} \\
 & & \texttt{1} \\
 & & \texttt{1} \\
 & & \texttt{1} \\
 & & \texttt{1} \\
 & & \texttt{1} \\
 & & \texttt{1} \\
 & & \texttt{1} \\
 & & \texttt{1} \\
 & & \texttt{1} \\
 & & \texttt{1} \\
 & & \texttt{1} \\
 & & \texttt{1} \\
 & & \texttt{1} \\
 & & \texttt{1} \\
 & & \texttt{1} \\
 & & \texttt{1} \\
 & & \texttt{1} \\
 & & \texttt{1} \\
 & & \texttt{1} \\
 & & \texttt{1} \\
 & & \texttt{1} \\
 & & \texttt{1} \\
 & & \texttt{1} \\
 & & \texttt{1} \\
 & & \texttt{1} \\
 & & \texttt{1} \\
 & & \texttt{1} \\
 & & \texttt{1} \\
 & & \texttt{1} \\
 & & \texttt{1} \\
 & & \texttt{1} \\
 & & \texttt{1} \\
 & & \texttt{1} \\
 & & \texttt{1} \\
 & & \texttt{1} \\
 & & \texttt{1} \\
 & & \texttt{1} \\
 & & \texttt{1} \\
 & & \texttt{1} \\
 & & \texttt{1} \\
 & & \texttt{1} \\
 & & \texttt{1} \\
 & & \texttt{1} \\
 & & \texttt{1} \\
 & & \texttt{1} \\
 & & \texttt{1} \\
 & & \texttt{1} \\
 & & \texttt{1} \\
 & & \texttt{1} \\
 & & \texttt{1} \\
 & & \texttt{1} \\
 & & \texttt{1} \\
 & & \texttt{1} \\
 & & \texttt{1} \\
 & & \texttt{1} \\
 & & \texttt{1} \\
 & & \texttt{1} \\
 & & \texttt{1} \\
 & & \texttt{1} \\
 & & \texttt{1} \\
 & & \texttt{1} \\
 & & \texttt{1} \\
 & & \texttt{1} \\
 & & \texttt{1} \\
 & & \texttt{1} \\
 & & \texttt{1} \\
 & & \texttt{1} \\
 & & \texttt{1} \\
 & & \texttt{1} \\
 & & \texttt{1} \\
 & & \texttt{1} \\
 & & \texttt{1} \\
 & & \texttt{1} \\
 & & \texttt{1} \\
 & & \texttt{1} \\
 & & \texttt{1} \\
 & & \texttt{1} \\
 & & \texttt{1} \\
 & & \texttt{1} \\
 & & \texttt{1} \\
 & & \texttt{1} \\
 & & \texttt{1} \\
 & & \texttt{1} \\
 & & \texttt{1} \\
 & & \texttt{1} \\
 & & \texttt{1} \\
 & & \texttt{1} \\
 & & \texttt{1} \\
 & & \texttt{1} \\
 & & \texttt{1} \\
 & & \texttt{1} \\
 & & \texttt{1} \\
 & & \texttt{1} \\
 & & \texttt{1} \\
 & & \texttt{1} \\
 & & \texttt{1} \\
 & & \texttt{1} \\
 & & \texttt{1} \\
 & & \texttt{1} \\
 & & \texttt{1} \\
 & & \texttt{1} \\
 & & \texttt{1} \\
 & & \texttt{1} \\
 & & \texttt{1} \\
 & & \texttt{1} \\
 & & \texttt{1} \\
 & & \texttt{1} \\
 & & \texttt{1} \\
 & & \texttt{1} \\
 & & \texttt{1} \\
 & & \texttt{1} \\
 & & \texttt{1} \\
 & & \texttt{1} \\
 & & \texttt{1} \\
 & & \texttt{1} \\
 & & \texttt{1} \\
 & & \texttt{1} \\
 & & \texttt{1} \\
 & & \texttt{1} \\
 & & \texttt{1} \\
 & & \texttt{1} \\
 & & \texttt{1} \\
 & & \texttt{1} \\
 & & \texttt{1} \\
 & & \texttt{1} \\
 & & \texttt{1} \\
 & & \texttt{1} \\
 & & \texttt{1} \\
 & & \texttt{1} \\
 & & \texttt{1} \\
 & & \texttt{1} \\
 & & \texttt{1} \\
 & & \texttt{1} \\
 & & \texttt{1} \\
 & & \texttt{1} \\
 & & \texttt{1} \\
 & & \texttt{1} \\
 & & \texttt{1} \\
 & & \texttt{1} \\
 & & \texttt{1} \\
 & & \texttt{1} \\
 & & \texttt{1} \\
 & & \texttt{1} \\
 & & \texttt{1} \\
 & & \texttt{1} \\
 & & \texttt{1} \\
 & & \texttt{1} \\
 & & \texttt{1} \\
 & & \texttt{1} \\
 & & \texttt{1} \\
 & & \texttt{1} \\
 & & \texttt{1} \\
 & & \texttt{1} \\
 & & \texttt{1} \\
 & & \texttt{1} \\
 & & \texttt{1} \\
 & & \texttt{1} \\
 & & \texttt{1} \\
 & & \texttt{1} \\
 & & \texttt{1} \\
 & & \texttt{1} \\
 & & \texttt{1} \\
 & & \texttt{1} \\
 & & \texttt{1} \\
 & & \texttt{1} \\
 & & \texttt{1} \\
 & & \texttt{1} \\
 & & \texttt{1} \\
 & & \texttt{1} \\
 & & \texttt{1} \\
 & & \texttt{1} \\
 & & \texttt{1} \\
 & & \texttt{1} \\
 & & \texttt{1} \\
 & & \texttt{1} \\
 & & \texttt{1} \\
 & & \texttt{1} \\
 & & \texttt{1} \\
 & & \texttt{1} \\
 & & \texttt{1} \\
 & & \texttt{1} \\
 & & \texttt{1} \\
 & & \texttt{1} \\
 & & \texttt{1} \\
 & & \texttt{1} \\
 & & \texttt{1} \\
 & & \texttt{1} \\
 & & \texttt{1} \\
 & & \texttt{1} \\
 & & \texttt{1} \\
 & & \texttt{1} \\
 & & \texttt{1} \\
 & & \texttt{1} \\
 & & \texttt{1} \\
 & & \texttt{1} \\
 & & \texttt{1} \\
 & & \texttt{1} \\
 & & \texttt{1} \\
 & & \texttt{1} \\
 & & \texttt{1} \\
 & & \texttt{1} \\
 & & \texttt{1} \\
 & & \texttt{1} \\
 & & \texttt{1} \\
 & & \texttt{1} \\
 & & \texttt{1} \\
 & & \texttt{1} \\
 & & \texttt{1} \\
 & & \texttt{1} \\
 & & \texttt{1} \\
 & & \texttt{1} \\
 & & \texttt{1} \\
 & & \texttt{1} \\
 & & \texttt{1} \\
 & & \texttt{1} \\
 & & \texttt{1} \\
 & & \texttt{1} \\
 & & \texttt{1} \\
 & & \texttt{1} \\
 & & \texttt{1} \\
 & & \texttt{1} \\
 & & \texttt{1} \\
 & & \texttt{1} \\
 & & \texttt{1} \\
 & & \texttt{1} \\
 & & \texttt{1} \\
 & & \texttt{1} \\
 & & \texttt{1} \\
 & & \texttt{1} \\
 & & \texttt{1} \\
 & & \texttt{1} \\
 & & \texttt{1} \\
 & & \texttt{1} \\
 & & \texttt{1} \\
 & & \texttt{1} \\
 & & \texttt{1} \\
 & & \texttt{1} \\
 & & \texttt{1} \\
 & & \texttt{1} \\
 & & \texttt{1} \\
 & & \texttt{1} \\
 & & \texttt{1} \\
 & & \texttt{1} \\
 & & \texttt{1} \\
 & & \texttt{1} \\
 & & \texttt{1} \\
 & & \texttt{1} \\
 & & \texttt{1} \\
 & & \texttt{1} \\
 & & \texttt{1} \\
 & & \texttt{1} \\
 & & \texttt{1} \\
 & & \texttt{1} \\
 & & \texttt{1} \\
 & & \texttt{1} \\
 & & \texttt{1} \\
 & & \texttt{1} \\
 & & \texttt{1} \\
 & & \texttt{1} \\
 & & \texttt{1} \\
 & & \texttt{1} \\
 & & \texttt{1} \\
 & & \texttt{1} \\
 & & \texttt{1} \\
 & & \texttt{1} \\
 & & \texttt{1} \\
 & & \texttt{1} \\
 & & \texttt{1} \\
 & & \texttt{1} \\
 & & \texttt{1} \\
 & & \texttt{1} \\
 & & \texttt{1} \\
 & & \texttt{1} \\
 & & \texttt{1} \\
 & & \texttt{1} \\
 & & \texttt{1} \\
 & & \texttt{1} \\
 & & \texttt{1} \\
 & & \texttt{1} \\
 & & \texttt{1} \\
 & & \texttt{1} \\
 & & \texttt{1} \\
 & & \texttt{1} \\
 & & \texttt{1} \\
 & & \texttt{1} \\
 & & \texttt{1} \\
 & & \texttt{1} \\
 & & \texttt{1} \\
 & & \texttt{1} \\
 & & \texttt{1} \\
 & & \texttt{1} \\
 & & \texttt{1} \\
 & & \texttt{1} \\
 & & \texttt{1} \\
 & & \texttt{1} \\
 & & \texttt{1} \\
 & & \texttt{1} \\
 & & \texttt{1} \\
 & & \texttt{1} \\
 & & \texttt{1} \\
 & & \texttt{1} \\
 & & \texttt{1} \\
 & & \texttt{1} \\
 & & \texttt{1} \\
 & & \texttt{1} \\
 & & \texttt{1} \\
 & & \texttt{1} \\
 & & \texttt{1} \\
 & & \texttt{1} \\
 & & \texttt{1} \\
 & & \texttt{1} \\
 & & \texttt{1} \\
 & & \texttt{1} \\
 & & \texttt{1} \\
 & & \texttt{1} \\
 & & \texttt{1} \\
 & & \texttt{1} \\
 & & \texttt{1} \\
 & & \texttt{1} \\
 & & \texttt{1} \\
 & & \texttt{1} \\
 & & \texttt{1} \\
 & & \texttt{1} \\
 & & \texttt{1} \\
 & & \texttt{1} \\
 & & \texttt{1} \\
 & & \texttt{1} \\
 & & \texttt{1} \\
 & & \texttt{1} \\
 & & \texttt{1} \\
 & & \texttt{1} \\
 & & \texttt{1} \\
 & & \texttt{1} \\
 & & \texttt{1} \\
 & & \texttt{1} \\
 & & \texttt{1} \\
 & & \texttt{1} \\
 & & \texttt{1} \\
 & & \texttt{1} \\
 & & \texttt{1} \\
 & & \texttt{1} \\
 & & \texttt{1} \\
 & & \texttt{1} \\
 & & \texttt{1} \\
 & & \texttt{1} \\
 & & \texttt{1} \\
 & & \texttt{1} \\
 & & \texttt{1} \\
 & & \texttt{1} \\
 & & \texttt{1} \\
 & & \texttt{1} \\
 & & \texttt{1} \\
 & & \texttt{1} \\
 & & \texttt{1} \\
 & & \texttt{1} \\
 & & \texttt{1} \\
 & & \texttt{1} \\
 & & \texttt{1} \\
 & & \texttt{1} \\
 & & \texttt{1} \\
 & & \texttt{1} \\
 & & \texttt{1} \\
 & & \texttt{1} \\
 & & \texttt{1} \\
 & & \texttt{1} \\
 & & \texttt{1} \\
 & & \texttt{1} \\
 & & \texttt{1} \\
 & & \texttt{1} \\
 & & \texttt{1} \\
 & & \texttt{1} \\
 & & \texttt{1} \\
 & & \texttt{1} \\
 & & \texttt{1} \\
 & & \texttt{1} \\
 & & \texttt{1} \\
 & & \texttt{1} \\
 & & \texttt{1} \\
 & & \texttt{1} \\
 & & \texttt{1} \\
 & & \texttt{1} \\
 & & \texttt{1} \\
 & & \texttt{1} \\
 & & \texttt{1} \\
 & & \texttt{1} \\
 & & \texttt{1} \\
 & & \texttt{1} \\
 & & \texttt{1} \\
 & & \texttt{1} \\
 & & \texttt{1} \\
 & & \texttt{1} \\
 & & \texttt{1} \\
 & & \texttt{1} \\
 & & \texttt{1} \\
 & & \texttt{1} \\
 & & \texttt{1} \\
 & & \texttt{1} \\
 & & \texttt{1} \\
 & & \texttt{1} \\
 & & \texttt{1} \\
 & & \texttt{1} \\
 & & \texttt{1} \\
 & & \texttt{1} \\
 & & \texttt{1} \\
 & & \texttt{1} \\
 & & \texttt{1} \\
 & & \texttt{1} \\
 & & \texttt{1} \\
 & & \texttt{1} \\
 & & \texttt{1} \\
 & & \texttt{1} \\
 & & \texttt{1} \\
 & & \texttt{1} \\
 & & \texttt{1} \\
 & & \texttt{1} \\
 & & \texttt{1} \\
 & & \texttt{1} \\
 & & \texttt{1} \\
 & & \texttt{1} \\
 & & \texttt{1} \\
 & & \texttt{1} \\
 & & \texttt{1} \\
 & & \texttt{1} \\
 & & \texttt{1} \\
 & & \texttt{1} \\
 & & \texttt{1} \\
 & & \texttt{1} \\
 & & \texttt{1} \\
 & & \texttt{1} \\
 & & \texttt{1} \\
 & & \texttt{1} \\
 & & \texttt{1} \\
 & & \texttt{1} \\
 & & \texttt{1} \\
 & & \texttt{1} \\
 & & \texttt{1} \\
 & & \texttt{1} \\
 & & \texttt{1} \\
 & & \texttt{1} \\
 & & \texttt{1} \\
 & & \texttt{1} \\
 & & \texttt{1} \\
 & & \texttt{1} \\
 & & \texttt{1} \\
 & & \texttt{1} \\
 & & \texttt{1} \\
 & & \texttt{1} \\
 & & \texttt{1} \\
 & & \texttt{1} \\
 & & \texttt{1} \\
 & & \texttt{1} \\
 & & \texttt{1} \\
 & & \texttt{1} \\
 & & \texttt{1} \\
 & & \texttt{1} \\
 & & \texttt{1} \\
 & & \texttt{1} \\
 & & \texttt{1} \\
 & & \texttt{1} \\
 & & \texttt{1} \\
 & & \texttt{1} \\
 & & \texttt{1} \\
 & & \texttt{1} \\
 & & \texttt{1} \\
 & & \texttt{1} \\
 & & \texttt{1} \\
 & & \texttt{1} \\
 & & \texttt{1} \\
 & & \texttt{1} \\
 & & \texttt{1} \\
 & & \texttt{1} \\
 & & \texttt{1} \\
 & & \texttt{1} \\
 & & \texttt{1} \\
 & & \texttt{1} \\
 & & \texttt{1} \\
 & & \texttt{1} \\
 & & \texttt{1} \\
 & & \texttt{1} \\
 & & \texttt{1} \\
 & & \texttt{1} \\
 & & \texttt{1} \\
 & & \texttt{1} \\
 & & \texttt{1} \\
 & & \texttt{1} \\
 & & \texttt{1} \\
 & & \texttt{1} \\
 & & \texttt{1} \\
 & & \texttt{1} \\
 & & \texttt{1} \\
 & & \texttt{1} \\
 & & \texttt{1} \\
 & & \texttt{1} \\
 & & \texttt{1} \\
 & & \texttt{1} \\
 & & \texttt{1} \\
 & & \texttt{1} \\
 & & \texttt{1} \\
 & & \texttt{1} \\
 & & \texttt{1} \\
 & & \texttt{1} \\
 & & \texttt{1} \\
 & & \texttt{1} \\
 & & \texttt{1} \\
 & & \texttt{1} \\
 & & \texttt{1} \\
 & & \texttt{1} \\
 & & \texttt{1} \\
 & & \texttt{1} \\
 & & \texttt{1} \\
 & & \texttt{1} \\
 & & \texttt{1} \\
 & & \texttt{1} \\
 & & \texttt{1} \\
 & & \texttt{1} \\
 & & \texttt{1} \\
 & & \texttt{1} \\
 & & \texttt{1} \\
 & & \texttt{1} \\
 & & \texttt{1} \\
 & & \texttt{1} \\
 & & \texttt{1} \\
 & & \texttt{1} \\
 & & \texttt{1} \\
 & & \texttt{1} \\
 & & \texttt{1} \\
 & & \texttt{1} \\
 & & \texttt{1} \\
 & & \texttt{1} \\
 & & \texttt{1} \\
 & & \texttt{1} \\
 & & \texttt{1} \\
 & & \texttt{1} \\
 & & \texttt{1} \\
 & & \texttt{1} \\
 & & \texttt{1} \\
 & & \texttt{1} \\
 & & \texttt{1} \\
 & & \texttt{1} \\
 & & \texttt{1} \\
 & & \texttt{1} \\
 & & \texttt{1} \\
 & & \texttt{1} \\
 & & \texttt{1} \\
 & & \texttt{1} \\
 & & \texttt{1} \\
 & & \texttt{1} \\
 & & \texttt{1} \\
 & & \texttt{1} \\
 & & \texttt{1} \\
 & & \texttt{1} \\
 & & \texttt{1} \\
 & & \texttt{1} \\
 & & \texttt{1} \\
 & & \texttt{1} \\
 & & \texttt{1} \\
 & & \texttt{1} \\
 & & \texttt{1} \\
 & & \texttt{1} \\
 & & \texttt{1} \\
 & & \texttt{1} \\
 & & \texttt{1} \\
 & & \texttt{1} \\
 & & \texttt{1} \\
 & & \texttt{1} \\
 & & \texttt{1} \\
 & & \texttt{1} \\
 & & \texttt{1} \\
 & & \texttt{1} \\
 & & \texttt{1} \\
 & & \texttt{1} \\
 & & \texttt{1} \\
 & & \texttt{1} \\
 & & \texttt{1} \\
 & & \texttt{1} \\
 & & \texttt{1} \\
 & & \texttt{1} \\
 & & \texttt{1} \\
 & & \texttt{1} \\
 & & \texttt{1} \\
 & & \texttt{1} \\
 & & \texttt{1} \\
 & & \texttt{1} \\
 & & \texttt{1} \\
 & & \texttt{1} \\
 & & \texttt{1} \\
 & & \texttt{1} \\
 & & \texttt{1} \\
 & & \texttt{1} \\
 & & \texttt{1} \\
 & & \texttt{1} \\
 & & \texttt{1} \\
 & & \texttt{1} \\
 & & \texttt{1} \\
 & & \texttt{1} \\
 & & \texttt{1} \\
 & & \texttt{1} \\
 & & \texttt{1} \\
 & & \texttt{1} \\
 & & \texttt{1} \\
 & & \texttt{1} \\
 & & \texttt{1} \\
 & & \texttt{1} \\
 & & \texttt{1} \\
 & & \texttt{1} \\
 & & \texttt{1} \\
 & & \texttt{1} \\
 & & \texttt{1} \\
 & & \texttt{1} \\
 & & \texttt{1} \\
 & & \texttt{1} \\
 & & \texttt{1} \\
 & & \texttt{1} \\
 & & \texttt{1} \\
 & & \texttt{1} \\
 & & \texttt{1} \\
 & & \texttt{1} \\
 & & \texttt{1} \\
 & & \texttt{1} \\
 & & \texttt{1} \\
 & & \texttt{1} \\
 & & \texttt{1} \\
 & & \texttt{1} \\
 & & \texttt{1} \\
 & & \texttt{1} \\
 & & \texttt{1} \\
 & & \texttt{1} \\
 & & \texttt{1} \\
 & & \texttt{1} \\
 & & \texttt{1} \\
 & & \texttt{1} \\
 & & \texttt{1} \\
 & & \texttt{1} \\
 & & \texttt{1} \\
 & & \texttt{1} \\
 & & \texttt{1} \\
 & & \texttt{1} \\
 & & \texttt{1} \\
 & & \texttt{1} \\
 & & \texttt{1} \\
 & & \texttt{1} \\
 & & \texttt{1} \\
 & & \texttt{1} \\
 & & \texttt{1} \\
 & & \texttt{1} \\
 & & \texttt{1} \\
 & & \texttt{1} \\
 & & \texttt{1} \\
 & & \texttt{1} \\
 & & \texttt{1} \\
 & & \texttt{1} \\
 & & \texttt{1} \\
 & & \texttt{1} \\
 & & \texttt{1} \\
 & & \texttt{1} \\
 & & \texttt{1} \\
 & & \texttt{1} \\
 & & \texttt{1} \\
 & & \texttt{1} \\
 & & \texttt{1} \\
 & & \texttt{1} \\
 & & \texttt{1} \\
 & & \texttt{1} \\
 & & \texttt{1} \\
 & & \texttt{1} \\
 & & \texttt{1} \\
 & & \texttt{1} \\
 & & \texttt{1} \\
 & & \texttt{1} \\
 & & \texttt{1} \\
 & & \texttt{1} \\
 & & \texttt{1} \\
 & & \texttt{1} \\
 & & \texttt{1} \\
 & & \texttt{1} \\
 & & \texttt{1} \\
 & & \texttt{1} \\
 & & \texttt{1} \\
 & & \texttt{1} \\
 & & \texttt{1} \\
 & & \texttt{1} \\
 & & \texttt{1} \\
 & & \texttt{1} \\
 & & \texttt{1} \\
 & & \texttt{1} \\
 & & \texttt{1} \\
 & & \texttt{1} \\
 & & \texttt{1} \\
 & & \texttt{1} \\
 & & \texttt{1} \\
 & & \texttt{1} \\
 & & \texttt{1} \\
 & & \texttt{1} \\
 & & \texttt{1} \\
 & & \texttt{1} \\
 & & \texttt{1} \\
 & & \texttt{1} \\
 & & \texttt{1} \\
 & & \texttt{1} \\
 & & \texttt{1} \\
 & & \texttt{1} \\
 & & \texttt{1} \\
 & & \texttt{1} \\
 & & \texttt{1} \\
 & & \texttt{1} \\
 & & \texttt{1} \\
 & & \texttt{1} \\
 & & \texttt{1} \\
 & & \texttt{1} \\
 & & \texttt{1} \\
 & & \texttt{1} \\
 & & \texttt{1} \\
 & & \texttt{1} \\
 & & \texttt{1} \\
 & & \texttt{1} \\
 & & \texttt{1} \\
 & & \texttt{1} \\
 & & \texttt{1} \\
 & & \texttt{1} \\
 & & \texttt{1} \\
 & & \texttt{1} \\
 & & \texttt{1} \\
 & & \texttt{1} \\
 & & \texttt{1} \\
 & & \texttt{1} \\
 & & \texttt{1} \\
 & & \texttt{1} \\
 & & \texttt{1} \\
 & & \texttt{1} \\
 & & \texttt{1} \\
 & & \texttt{1} \\
 & & \texttt{1} \\
 & & \texttt{1} \\
 & & \texttt{1} \\
 & & \texttt{1} \\
 & & \texttt{1} \\
 & & \texttt{1} \\
 & & \texttt{1} \\
 & & \texttt{1} \\
 & & \texttt{1} \\
 & & \texttt{1} \\
 & & \texttt{1} \\
 & & \texttt{1} \\
 & & \texttt{1} \\
 & & \texttt{1} \\
 & & \texttt{1} \\
 & & \texttt{1} \\
 & & \texttt{1} \\
 & & \texttt{1} \\
 & & \texttt{1} \\
 & & \texttt{1} \\
 & & \texttt{1} \\
 & & \texttt{1} \\
 & & \texttt{1} \\
 & & \texttt{1} \\
 & & \texttt{1} \\
 & & \texttt{1} \\
 & & \texttt{1} \\
 & & \texttt{1} \\
 & & \texttt{1} \\
 & & \texttt{1} \\
 & & \texttt{1} \\
 & & \texttt{1} \\
 & & \texttt{1} \\
 & & \texttt{1} \\
 & & \texttt{1} \\
 & & \texttt{1} \\
 & & \texttt{1} \\
 & & \texttt{1} \\
 & & \texttt{1} \\
 & & \texttt{1} \\
 & & \texttt{1} \\
 & & \texttt{1} \\
 & & \texttt{1} \\
 & & \texttt{1} \\
 & & \texttt{1} \\
 & & \texttt{1} \\
 & & \texttt{1} \\
 & & \texttt{1} \\
 & & \texttt{1} \\
 & & \texttt{1} \\
 & & \texttt{1} \\
 & & \texttt{1} \\
 & & \texttt{1} \\
 & & \texttt{1} \\
 & & \texttt{1} \\
 & & \texttt{1} \\
 & & \texttt{1} \\
 & & \texttt{1} \\
 & & \texttt{1} \\
 & & \texttt{1} \\
 & & \texttt{1} \\
 & & \texttt{1} \\
 & & \texttt{1} \\
 & & \texttt{1} \\
 & & \texttt{1} \\
 & & \texttt{1} \\
 & & \texttt{1} \\
 & & \texttt{1} \\
 & & \texttt{1} \\
 & & \texttt{1} \\
 & & \texttt{1} \\
 & & \texttt{1} \\
 & & \texttt{1} \\
 & & \texttt{1} \\
 & & \texttt{1} \\
 & & \texttt{1} \\
 & & \texttt{1} \\
 & & \texttt{1} \\
 & & \texttt{1} \\
 & & \texttt{1} \\
 & & \texttt{1} \\
 & & \texttt{1} \\
 & & \texttt{1} \\
 & & \texttt{1} \\
 & & \texttt{1} \\
 & & \texttt{1} \\
 & & \texttt{1} \\
 & & \texttt{1} \\
 & & \texttt{1} \\
 & & \texttt{1} \\
 & & \texttt{1} \\
 & & \texttt{1} \\
 & & \texttt{1} \\
 & & \texttt{1} \\
 & & \texttt{1} \\
 & & \texttt{1} \\
 & & \texttt{1} \\
 & & \texttt{1} \\
 & & \texttt{1} \\
 & & \texttt{1} \\
 & & \texttt{1} \\
 & & \texttt{1} \\
 & & \texttt{1} \\
 & & \texttt{1} \\
 & & \texttt{1} \\
 & & \texttt{1} \\
 & & \texttt{1} \\
 & & \texttt{1} \\
 & & \texttt{1} \\
 & & \texttt{1} \\
 & & \texttt{1} \\
 & & \texttt{1} \\
 & & \texttt{1} \\
 & & \texttt{1} \\
 & & \texttt{1} \\
 & & \texttt{1} \\
 & & \texttt{1} \\
 & & \texttt{1} \\
 & & \texttt{1} \\
 & & \texttt{1} \\
 & & \texttt{1} \\
 & & \texttt{1} \\
 & & \texttt{1} \\
 & & \texttt{1} \\
 & & \texttt{1} \\
 & & \texttt{1} \\
 & & \texttt{1} \\
 & & \texttt{1} \\
 & & \texttt{1} \\
 & & \texttt{1} \\
 & & \texttt{1} \\
 & & \texttt{1} \\
 & & \texttt{1} \\
 & & \texttt{1} \\
 & & \texttt{1} \\
 & & \texttt{1} \\
 & & \texttt{1} \\
 & & \texttt{1} \\
 & & \texttt{1} \\
 & & \texttt{1} \\
 & & \texttt{1} \\
 & & \texttt{1} \\
 & & \texttt{1} \\
 & & \texttt{1} \\
 & & \texttt{1} \\
 & & \texttt{1} \\
 & & \texttt{1} \\
 & & \texttt{1} \\
 & & \texttt{1} \\
 & & \texttt{1} \\
 & & \texttt{1} \\
 & & \texttt{1} \\
 & & \texttt{1} \\
 & & \texttt{1} \\
 & & \texttt{1} \\
 & & \texttt{1} \\
 & & \texttt{1} \\
 & & \texttt{1} \\
 & & \texttt{1} \\
 & & \texttt{1} \\
 & & \texttt{1} \\
 & & \texttt{1} \\
 & & \texttt{1} \\
 & & \texttt{1} \\
 & & \texttt{1} \\
 & & \texttt{1} \\
 & & \texttt{1} \\
 & & \texttt{1} \\
 & & \texttt{1} \\
 & & \texttt{1} \\
 & & \texttt{1} \\
 & & \texttt{1} \\
 & & \texttt{1} \\
 & & \texttt{1} \\
 & & \texttt{1} \\
 & & \texttt{1} \\
 & & \texttt{1} \\
 & & \texttt{1} \\
 & & \texttt{1} \\
 & & \texttt{1} \\
 & & \texttt{1} \\
 & & \texttt{1} \\
 & & \texttt{1} \\
 & & \texttt{1} \\
 & & \texttt{1} \\
 & & \texttt{1} \\
 & & \texttt{1} \\
 & & \texttt{1} \\
 & & \texttt{1} \\
 & & \texttt{1} \\
 & & \texttt{1} \\
 & & \texttt{1} \\
 & & \texttt{1} \\
 & & \texttt{1} \\
 & & \texttt{1} \\
 & & \texttt{1} \\
 & & \texttt{1} \\
 & & \texttt{1} \\
 & & \texttt{1} \\
 & & \texttt{1} \\
 & & \texttt{1} \\
 & & \texttt{1} \\
 & & \texttt{1} \\
 & & \texttt{1} \\
 & & \texttt{1} \\
 & & \texttt{1} \\
 & & \texttt{1} \\
 & & \texttt{1} \\
 & & \texttt{1} \\
 & & \texttt{1} \\
 & & \texttt{1} \\
 & & \texttt{1} \\
 & & \texttt{1} \\
 & & \texttt{1} \\
 & & \texttt{1} \\
 & & \texttt{1} \\
 & & \texttt{1} \\
 & & \texttt{1} \\
 & & \texttt{1} \\
 & & \texttt{1} \\
 & & \texttt{1} \\
 & & \texttt{1} \\
 & & \texttt{1} \\
 & & \texttt{1} \\
 & & \texttt{1} \\
 & & \texttt{1} \\
 & & \texttt{1} \\
 & & \texttt{1} \\
 & & \texttt{1} \\
 & & \texttt{1} \\
 & & \texttt{1} \\
 & & \texttt{1} \\
 & & \texttt{1} \\
 & & \texttt{1} \\
 & & \texttt{1} \\
 & & \texttt{1} \\
 & & \texttt{1} \\
 & & \texttt{1} \\
 & & \texttt{1} \\
 & & \texttt{1} \\
 & & \texttt{1} \\
 & & \texttt{1} \\
 & & \texttt{1} \\
 & & \texttt{1} \\
 & & \texttt{1} \\
 & & \texttt{1} \\
 & & \texttt{1} \\
 & & \texttt{1} \\
 & & \texttt{1} \\
 & & \texttt{1} \\
 & & \texttt{1} \\
 & & \texttt{1} \\
 & & \texttt{1} \\
 & & \texttt{1} \\
 & & \texttt{1} \\
 & & \texttt{1} \\
 & & \texttt{1} \\
 & & \texttt{1} \\
 & & \texttt{1} \\
 & & \texttt{1} \\
 & & \texttt{1} \\
 & & \texttt{1} \\
 & & \texttt{1} \\
 & & \texttt{1} \\
 & & \texttt{1} \\
 & & \texttt{1} \\
 & & \texttt{1} \\
 & & \texttt{1} \\
 & & \texttt{1} \\
 & & \texttt{1} \\
 & & \texttt{1} \\
 & & \texttt{1} \\
 & & \texttt{1} \\
 & & \texttt{1} \\
 & & \texttt{1} \\
 & & \texttt{1} \\
 & & \texttt{1} \\
 & & \texttt{1} \\
 & & \texttt{1} \\
 & & \texttt{1} \\
 & & \texttt{1} \\
 & & \texttt{1} \\
 & & \texttt{1} \\
 & & \texttt{1} \\
 & & \texttt{1} \\
 & & \texttt{1} \\
 & & \texttt{1} \\
 & & \texttt{1} \\
 & & \texttt{1} \\
 & & \texttt{1} \\
 & & \texttt{1} \\
 & & \texttt{1} \\
 & & \texttt{1} \\
 & & \texttt{1} \\
 & & \texttt{1} \\
 & & \texttt{1} \\
 & & \texttt{1} \\
 & & \texttt{1} \\
 & & \texttt{1} \\
 & & \texttt{1} \\
 & & \texttt{1} \\
 & & \texttt{1} \\
 & & \texttt{1} \\
 & & \texttt{1} \\
 & & \texttt{1} \\
 & & \texttt{1} \\
 & & \texttt{1} \\
 & & \texttt{1} \\
 & & \texttt{1} \\
 & & \texttt{1} \\
 & & \texttt{1} \\
 & & \texttt{1} \\
 & & \texttt{1} \\
 & & \texttt{1} \\
 & & \texttt{1} \\
 & & \texttt{1} \\
 & & \texttt{1} \\
 & & \texttt{1} \\
 & & \texttt{1} \\
 & & \texttt{1} \\
 & & \texttt{1} \\
 & & \texttt{1} \\
 & & \texttt{1} \\
 & & \texttt{1} \\
 & & \texttt{1} \\
 & & \texttt{1} \\
 & & \texttt{1} \\
 & & \texttt{1} \\
 & & \texttt{1} \\
 & & \texttt{1} \\
 & & \texttt{1} \\
 & & \texttt{1} \\
 & & \texttt{1} \\
 & & \texttt{1} \\
 & & \texttt{1} \\
 & & \texttt{1} \\
 & & \texttt{1} \\
 & & \texttt{1} \\
 & & \texttt{1} \\
 & & \texttt{1} \\
 & & \texttt{1} \\
 & & \texttt{1} \\
 & & \texttt{1} \\
 & & \texttt{1} \\
 & & \texttt{1} \\
 & & \texttt{1} \\
 & & \texttt{1} \\
 & & \texttt{1} \\
 & & \texttt{1} \\
 & & \texttt{1} \\
 & & \texttt{1} \\
 & & \texttt{1} \\
 & & \texttt{1} \\
 & & \texttt{1} \\
 & & \texttt{1} \\
 & & \texttt{1} \\
 & & \texttt{1} \\
 & & \texttt{1} \\
 & & \texttt{1} \\
 & & \texttt{1} \\
 & & \texttt{1} \\
 & & \texttt{1} \\
 & & \texttt{1} \\
 & & \texttt{1} \\
 & & \texttt{1} \\
 & & \texttt{1} \\
 & & \texttt{1} \\
 & & \texttt{1} \\
 & & \texttt{1} \\
 & & \texttt{1} \\
 & & \texttt{1} \\
 & & \texttt{1} \\
 & & \texttt{1} \\
 & & \texttt{1} \\
 & & \texttt{1} \\
 & & \texttt{1} \\
 & & \texttt{1} \\
 & & \texttt{1} \\
 & & \texttt{1} \\
 & & \texttt{1} \\
 & & \texttt{1} \\
 & & \texttt{1} \\
 & & \texttt{1} \\
 & & \texttt{1} \\
 & & \texttt{1} \\
 & & \texttt{1} \\
 & & \texttt{1} \\
 & & \texttt{1} \\
 & & \texttt{1} \\
 & & \texttt{1} \\
 & & \texttt{1} \\
 & & \texttt{1} \\
 & & \texttt{1} \\
 & & \texttt{1} \\
 & & \texttt{1} \\
 & & \texttt{1} \\
 & & \texttt{1} \\
 & & \texttt{1} \\
 & & \texttt{1} \\
 & & \texttt{1} \\
 & & \texttt{1} \\
 & & \texttt{1} \\
 & & \texttt{1} \\
 & & \texttt{1} \\
 & & \texttt{1} \\
 & & \texttt{1} \\
 & & \texttt{1} \\
 & & \texttt{1} \\
 & & \texttt{1} \\
 & & \texttt{1} \\
 & & \texttt{1} \\
 & & \texttt{1} \\
 & & \texttt{1} \\
 & & \texttt{1} \\
 & & \texttt{1} \\
 & & \texttt{1} \\
 & & \texttt{1} \\
 & & \texttt{1} \\
 & & \texttt{1} \\
 & & \texttt{1} \\
 & & \texttt{1} \\
 & & \texttt{1} \\
 & & \texttt{1} \\
 & & \texttt{1} \\
 & & \texttt{1} \\
 & & \texttt{1} \\
 & & \texttt{1} \\
 & & \texttt{1} \\
 & & \texttt{1} \\
 & & \texttt{1} \\
 & & \texttt{1} \\
 & & \texttt{1} \\
 & & \texttt{1} \\
 & & \texttt{1} \\
 & & \texttt{1} \\
 & & \texttt{1} \\
 & & \texttt{1} \\
 & & \texttt{1} \\
 & & \texttt{1} \\
 & & \texttt{1} \\
 & & \texttt{1} \\
 & & \texttt{1} \\
 & & \texttt{1} \\
 & & \texttt{1} \\
 & & \texttt{1} \\
 & & \texttt{1} \\
 & & \texttt{1} \\
 & & \texttt{1} \\
 & & \texttt{1} \\
 & & \texttt{1} \\
 & & \texttt{1} \\
 & & \texttt{1} \\
 & & \texttt{1} \\
 & & \texttt{1} \\
 & & \texttt{1} \\
 & & \texttt{1} \\
 & & \texttt{1} \\
 & & \texttt{1} \\
 & & \texttt{1} \\
 & & \texttt{1} \\
 & & \texttt{1} \\
 & & \texttt{1} \\
 & & \texttt{1} \\
 & & \texttt{1} \\
 & & \texttt{1} \\
 & & \texttt{1} \\
 & & \texttt{1} \\
 & & \texttt{1} \\
 & & \texttt{1} \\
 & & \texttt{1} \\
 & & \texttt{1} \\
 & & \texttt{1} \\
 & & \texttt{1} \\
 & & \texttt{1} \\
 & & \texttt{1} \\
 & & \texttt{1} \\
 & & \texttt{1} \\
 & & \texttt{1} \\
 & & \texttt{1} \\
 & & \texttt{1} \\
 & & \texttt{1} \\
 & & \texttt{1} \\
 & & \texttt{1} \\
 & & \texttt{1} \\
 & & \texttt{1} \\
 & & \texttt{1} \\
 & & \texttt{1} \\
 & & \texttt{1} \\
 & & \texttt{1} \\
 & & \texttt{1} \\
 & & \texttt{1} \\
 & & \texttt{1} \\
 & & \texttt{1} \\
 & & \texttt{1} \\
 & & \texttt{1} \\
 & & \texttt{1} \\
 & & \texttt{1} \\
 & & \texttt{1} \\
 & & \texttt{1} \\
 & & \texttt{1} \\
 & & \texttt{1} \\
 & & \texttt{1} \\
 & & \texttt{1} \\
 & & \texttt{1} \\
 & & \texttt{1} \\
 & & \texttt{1} \\
 & & \texttt{1} \\
 & & \texttt{1} \\
 & & \texttt{1} \\
 & & \texttt{1} \\
 & & \texttt{1} \\
 & & \texttt{1} \\
 & & \texttt{1} \\
 & & \texttt{1} \\
 & & \texttt{1} \\
 & & \texttt{1} \\
 & & \texttt{1} \\
 & & \texttt{1} \\
 & & \texttt{1} \\
 & & \texttt{1} \\
 & & \texttt{1} \\
 & & \texttt{1} \\
 & & \texttt{1} \\
 & & \texttt{1} \\
 & & \texttt{1} \\
 & & \texttt{1} \\
 & & \texttt{1} \\
 & & \texttt{1} \\
 & & \texttt{1} \\
 & & \texttt{1} \\
 & & \texttt{1} \\
 & & \texttt{1} \\
 & & \texttt{1} \\
 & & \texttt{1} \\
 & & \texttt{1} \\
 & & \texttt{1} \\
 & & \texttt{1} \\
 & & \texttt{1} \\
 & & \texttt{1} \\
 & & \texttt{1} \\
 & & \texttt{1} \\
 & & \texttt{1} \\
 & & \texttt{1} \\
 & & \texttt{1} \\
 & & \texttt{1} \\
 & & \texttt{1} \\
 & & \texttt{1} \\
 & & \texttt{1} \\
 & & \texttt{1} \\
 & & \texttt{1} \\
 & & \texttt{1} \\
 & & \texttt{1} \\
 & & \texttt{1} \\
 & & \texttt{1} \\
 & & \texttt{1} \\
 & & \texttt{1} \\
 & & \texttt{1} \\
 & & \texttt{1} \\
 & & \texttt{1} \\
 & & \texttt{1} \\
 & & \texttt{1} \\
 & & \texttt{1} \\
 & & \texttt{1} \\
 & & \texttt{1} \\
 & & \texttt{1} \\
 & & \texttt{1} \\
 & & \texttt{1} \\
 & & \texttt{1} \\
 & & \texttt{1} \\
 & & \texttt{1} \\
 & & \texttt{1} \\
 & & \texttt{1} \\
 & & \texttt{1} \\
 & & \texttt{1} \\
 & & \texttt{1} \\
 & & \texttt{1} \\
 & & \texttt{1} \\
 & & \texttt{1} \\
 & & \texttt{1} \\
 & & \texttt{1} \\
 & & \texttt{1} \\
 & & \texttt{1} \\
 & & \texttt{1} \\
 & & \texttt{1} \\
 & & \texttt{1} \\
 & & \texttt{1} \\
 & & \texttt{1} \\
 & & \texttt{1} \\
 & & \texttt{1} \\
 & & \texttt{1} \\
 & & \texttt{1} \\
 & & \texttt{1} \\
 & & \texttt{1} \\
 & & \texttt{1} \\
 & & \texttt{1} \\
 & & \texttt{1} \\
 & & \texttt{1} \\
 & & \texttt{1} \\
 & & \texttt{1} \\
 & & \texttt{1} \\
 & & \texttt{1} \\
 & & \texttt{1} \\
 & & \texttt{1} \\
 & & \texttt{1} \\
 & & \texttt{1} \\
 & & \texttt{1} \\
 & & \texttt{1} \\
 & & \texttt{1} \\
 & & \texttt{1} \\
 & & \texttt{1} \\
 & & \texttt{1} \\
 & & \texttt{1} \\
 & & \texttt{1} \\
 & & \texttt{1} \\
 & & \texttt{1} \\
 & & \texttt{1} \\
 & & \texttt{1} \\
 & & \texttt{1} \\
 & & \texttt{1} \\
 & & \texttt{1} \\
 & & \texttt{1} \\
 & & \texttt{1} \\
 & & \texttt{1} \\
 & & \texttt{1} \\
 & & \texttt{1} \\
 & & \texttt{1} \\
 & & \texttt{1} \\
 & & \texttt{1} \\
 & & \texttt{1} \\
 & & \texttt{1} \\
 & & \texttt{1} \\
 & & \texttt{1} \\
 & & \texttt{1} \\
 & & \texttt{1} \\
 & & \texttt{1} \\
 & & \texttt{1} \\
 & & \texttt{1} \\
 & & \texttt{1} \\
 & & \texttt{1} \\
 & & \texttt{1} \\
 & & \texttt{1} \\
 & & \texttt{1} \\
 & & \texttt{1} \\
 & & \texttt{1} \\
 & & \texttt{1} \\
 & & \texttt{1} \\
 & & \texttt{1} \\
 & & \texttt{1} \\
 & & \texttt{1} \\
 & & \texttt{1} \\
 & & \texttt{1} \\
 & & \texttt{1} \\
 & & \texttt{1} \\
 & & \texttt{1} \\
 & & \texttt{1} \\
 & & \texttt{1} \\
 & & \texttt{1} \\
 & & \texttt{1} \\
 & & \texttt{1} \\
 & & \texttt{1} \\
 & & \texttt{1} \\
 & & \texttt{1} \\
 & & \texttt{1} \\
 & & \texttt{1} \\
 & & \texttt{1} \\
 & & \texttt{1} \\
 & & \texttt{1} \\
 & & \texttt{1} \\
 & & \texttt{1} \\
 & & \texttt{1} \\
 & & \texttt{1} \\
 & & \texttt{1} \\
 & & \texttt{1} \\
 & & \texttt{1} \\
 & & \texttt{1} \\
 & & \texttt{1} \\
 & & \texttt{1} \\
 & & \texttt{1} \\
 & & \texttt{1} \\
 & & \texttt{1} \\
 & & \texttt{1} \\
 & & \texttt{1} \\
 & & \texttt{1} \\
 & & \texttt{1} \\
 & & \texttt{1} \\
 & & \texttt{1} \\
 & & \texttt{1} \\
 & & \texttt{1} \\
 & & \texttt{1} \\
 & & \texttt{1} \\
 & & \texttt{1} \\
 & & \texttt{1} \\
 & & \texttt{1} \\
 & & \texttt{1} \\
 & & \texttt{1} \\
 & & \texttt{1} \\
 & & \texttt{1} \\
 & & \texttt{1} \\
 & & \texttt{1} \\
 & & \texttt{1} \\
 & & \texttt{1} \\
 & & \texttt{1} \\
 & & \texttt{1} \\
 & & \texttt{1} \\
 & & \texttt{1} \\
 & & \texttt{1} \\
 & & \texttt{1} \\
 & & \texttt{1} \\
 & & \texttt{1} \\
 & & \texttt{1} \\
 & & \texttt{1} \\
 & & \texttt{1} \\
 & & \texttt{1} \\
 & & \texttt{1} \\
 & & \texttt{1} \\
 & & \texttt{1} \\
 & & \texttt{1} \\
 & & \texttt{1} \\
 & & \texttt{1} \\
 & & \texttt{1} \\
 & & \texttt{1} \\
 & & \texttt{1} \\
 & & \texttt{1} \\
 & & \texttt{1} \\
 & & \texttt{1} \\
 & & \texttt{1} \\
 & & \texttt{1} \\
 & & \texttt{1} \\
 & & \texttt{1} \\
 & & \texttt{1} \\
 & & \texttt{1} \\
 & & \texttt{1} \\
 & & \texttt{1} \\
 & & \texttt{1} \\
 & & \texttt{1} \\
 & & \texttt{1} \\
 & & \texttt{1} \\
 & & \texttt{1} \\
 & & \texttt{1} \\
 & & \texttt{1} \\
 & & \texttt{1} \\
 & & \texttt{1} \\
 & & \texttt{1} \\
 & & \texttt{1} \\
 & & \texttt{1} \\
 & & \texttt{1} \\
 & & \texttt{1} \\
 & & \texttt{1} \\
 & & \texttt{1} \\
 & & \texttt{1} \\
 & & \texttt{1} \\
 & & \texttt{1} \\
 & & \texttt{1} \\
 & & \texttt{1} \\
 & & \texttt{1} \\
 & & \texttt{1} \\
 & & \texttt{1} \\
 & & \texttt{1} \\
 & & \texttt{1} \\
 & & \texttt{1} \\
 & & \texttt{1} \\
 & & \texttt{1} \\
 & & \texttt{1} \\
 & & \texttt{1} \\
 & & \texttt{1} \\
 & & \texttt{1} \\
 & & \texttt{1} \\
 & & \texttt{1} \\
 & & \texttt{1} \\
 & & \texttt{1} \\
 & & \texttt{1} \\
 & & \texttt{1} \\
 & & \texttt{1} \\
 & & \texttt{1} \\
 & & \texttt{1} \\
 & & \texttt{1} \\
 & & \texttt{1} \\
 & & \texttt{1} \\
 & & \texttt{1} \\
 & & \texttt{1} \\
 & & \texttt{1} \\
 & & \texttt{1} \\
 & & \texttt{1} \\
 & & \texttt{1} \\
 & & \texttt{1} \\
 & & \texttt{1} \\
 & & \texttt{1} \\
 & & \texttt{1} \\
 & & \texttt{1} \\
 & & \texttt{1} \\
 & & \texttt{1} \\
 & & \texttt{1} \\
 & & \texttt{1} \\
 & & \texttt{1} \\
 & & \texttt{1} \\
 & & \texttt{1} \\
 & & \texttt{1} \\
 & & \texttt{1} \\
 & & \texttt{1} \\
 & & \texttt{1} \\
 & & \texttt{1} \\
 & & \texttt{1} \\
 & & \texttt{1} \\
 & & \texttt{1} \\
 & & \texttt{1} \\
 & & \texttt{1} \\
 & & \texttt{1} \\
 & & \texttt{1} \\
 & & \texttt{1} \\
 & & \texttt{1} \\
 & & \texttt{1} \\
 & & \texttt{1} \\
 & & \texttt{1} \\
 & & \texttt{1} \\
 & & \texttt{1} \\
 & & \texttt{1} \\
 & & \texttt{1} \\
 & & \texttt{1} \\
 & & \texttt{1} \\
 & & \texttt{1} \\
 & & \texttt{1} \\
 & & \texttt{1} \\
 & & \texttt{1} \\
 & & \texttt{1} \\
 & & \texttt{1} \\
 & & \texttt{1} \\
 & & \texttt{1} \\
 & & \texttt{1} \\
 & & \texttt{1} \\
 & & \texttt{1} \\
 & & \texttt{1} \\
 & & \texttt{1} \\
 & & \texttt{1} \\
 & & \texttt{1} \\
 & & \texttt{1} \\
 & & \texttt{1} \\
 & & \texttt{1} \\
 & & \texttt{1} \\
 & & \texttt{1} \\
 & & \texttt{1} \\
 & & \texttt{1} \\
 & & \texttt{1} \\
 & & \texttt{1} \\
 & & \texttt{1} \\
 & & \texttt{1} \\
 & & \texttt{1} \\
 & & \texttt{1} \\
 & & \texttt{1} \\
 & & \texttt{1} \\
 & & \texttt{1} \\
 & & \texttt{1} \\
 & & \texttt{1} \\
 & & \texttt{1} \\
 & & \texttt{1} \\
 & & \texttt{1} \\
 & & \texttt{1} \\
 & & \texttt{1} \\
 & & \texttt{1} \\
 & & \texttt{1} \\
 & & \texttt{1} \\
 & & \texttt{1} \\
 & & \texttt{1} \\
 & & \texttt{1} \\
 & & \texttt{1} \\
 & & \texttt{1} \\
 & & \texttt{1} \\
 & & \texttt{1} \\
 & & \texttt{1} \\
 & & \texttt{1} \\
 & & \texttt{1} \\
 & & \texttt{1} \\
 & & \texttt{1} \\
 & & \texttt{1} \\
 & & \texttt{1} \\
 & & \texttt{1} \\
 & & \texttt{1} \\
 & & \texttt{1} \\
 & & \texttt{1} \\
 & & \texttt{1} \\
 & & \texttt{1} \\
 & & \texttt{1} \\
 & & \texttt{1} \\
 & & \texttt{1} \\
 & & \texttt{1} \\
 & & \texttt{1} \\
 & & \texttt{1} \\
 & & \texttt{1} \\
 & & \texttt{1} \\
 & & \texttt{1} \\
 & & \texttt{1} \\
 & & \texttt{1} \\
 & & \texttt{1} \\
 & & \texttt{1} \\
 & & \texttt{1} \\
 & & \texttt{1} \\
 & & \texttt{1} \\
 & & \texttt{1} \\
 & & \texttt{1} \\
 & & \texttt{1} \\
 & & \texttt{1} \\
 & & \texttt{1} \\
 & & \texttt{1} \\
 & & \texttt{1} \\
 & & \texttt{1} \\
 & & \texttt{1} \\
 & & \texttt{1} \\
 & & \texttt{1} \\
 & & \texttt{1} \\
 & & \texttt{1} \\
 & & \texttt{1} \\
 & & \texttt{1} \\
 & & \texttt{1} \\
 & & \texttt{1} \\
 & & \texttt{1} \\
 & & \texttt{1} \\
 & & \texttt{1} \\
 & & \texttt{1} \\
 & & \texttt{1} \\
 & & \texttt{1} \\
 & & \texttt{1} \\
 & & \texttt{1} \\
 & & \texttt{1} \\
 & & \texttt{1} \\
 & & \texttt{1} \\
 & & \texttt{1} \\
 & & \texttt{1} \\
 & & \texttt{1} \\
 & & \texttt{1} \\
 & & \texttt{1} \\
 & & \texttt{1} \\
 & & \texttt{1} \\
 & & \texttt{1} \\
 & & \texttt{1} \\
 & & \texttt{1} \\
 & & \texttt{1} \\
 & & \texttt{1} \\
 & & \texttt{1} \\
 & & \texttt{1} \\
 & & \texttt{1} \\
 & & \texttt{1} \\
 & & \texttt{1} \\
 & & \texttt{1} \\
 & & \texttt{1} \\
 & & \texttt{1} \\
 & & \texttt{1} \\
 & & \texttt{1} \\
 & & \texttt{1} \\
 & & \texttt{1} \\
 & & \texttt{1} \\
 & & \texttt{1} \\
 & & \texttt{1} \\
 & & \texttt{1} \\
 & & \texttt{1} \\
 & & \texttt{1} \\
 & & \texttt{1} \\
 & & \texttt{1} \\
 & & \texttt{1} \\
 & & \texttt{1} \\
 & & \texttt{1} \\
 & & \texttt{1} \\
 & & \texttt{1} \\
 & & \texttt{1} \\
 & & \texttt{1} \\
 & & \texttt{1} \\
 & & \texttt{1} \\
 & & \texttt{1} \\
 & & \texttt{1} \\
 & & \texttt{1} \\
 & & \texttt{1} \\
 & & \texttt{1} \\
 & & \texttt{1} \\
 & & \texttt{1} \\
 & & \texttt{1} \\
 & & \texttt{1} \\
 & & \texttt{1} \\
 & & \texttt{1} \\
 & & \texttt{1} \\
 & & \texttt{1} \\
 & & \texttt{1} \\
 & & \texttt{1} \\
 & & \texttt{1} \\
 & & \texttt{1} \\
 & & \texttt{1} \\
 & & \texttt{1} \\
 & & \texttt{1} \\
 & & \texttt{1} \\
 & & \texttt{1} \\
 & & \texttt{1} \\
 & & \texttt{1} \\
 & & \texttt{1} \\
 & & \texttt{1} \\
 & & \texttt{1} \\
 & & \texttt{1} \\
 & & \texttt{1} \\
 & & \texttt{1} \\
 & & \texttt{1} \\
 & & \texttt{1} \\
 & & \texttt{1} \\
 & & \texttt{1} \\
 & & \texttt{1} \\
 & & \texttt{1} \\
 & & \texttt{1} \\
 & & \texttt{1} \\
 & & \texttt{1} \\
 & & \texttt{1} \\
 & & \texttt{1} \\
 & & \texttt{1} \\
 & & \texttt{1} \\
 & & \texttt{1} \\
 & & \texttt{1} \\
 & & \texttt{1} \\
 & & \texttt{1} \\
 & & \texttt{1} \\
 & & \texttt{1} \\
 & & \texttt{1} \\
 & & \texttt{1} \\
 & & \texttt{1} \\
 & & \texttt{1} \\
 & & \texttt{1} \\
 & & \texttt{1} \\
 & & \texttt{1} \\
 & & \texttt{1} \\
 & & \texttt{1} \\
 & & \texttt{1} \\
 & & \texttt{1} \\
 & & \texttt{1} \\
 & & \texttt{1} \\
 & & \texttt{1} \\
 & & \texttt{1} \\
 & & \texttt{1} \\
 & & \texttt{1} \\
 & & \texttt{1} \\
 & & \texttt{1} \\
 & & \texttt{1} \\
 & & \texttt{1} \\
 & & \texttt{1} \\
 & & \texttt{1} \\
 & & \texttt{1} \\
 & & \texttt{1} \\
 & & \texttt{1} \\
 & & \texttt{1} \\
 & & \texttt{1} \\
 & & \texttt{1} \\
 & & \texttt{1} \\
 & & \texttt{1} \\
 & & \texttt{1} \\
 & & \texttt{1} \\
 & & \texttt{1} \\
 & & \texttt{1} \\
 & & \texttt{1} \\
 & & \texttt{1} \\
 & & \texttt{1} \\
 & & \texttt{1} \\
 & & \texttt{1} \\
 & & \texttt{1} \\
 & & \texttt{1} \\
 & & \texttt{1} \\
 & & \texttt{1} \\
 & & \texttt{1} \\
 & & \texttt{1} \\
 & & \texttt{1} \\
 & & \texttt{1} \\
 & & \texttt{1} \\
 & & \texttt{1} \\
 & & \texttt{1} \\
 & & \texttt{1} \\
 & & \texttt{1} \\
 & & \texttt{1} \\
 & & \texttt{1} \\
 & & \texttt{1} \\
 & & \texttt{1} \\
 & & \texttt{1} \\
 & & \texttt{1} \\
 & & \texttt{1} \\
 & & \texttt{1} \\
 & & \texttt{1} \\
 & & \texttt{1} \\
 & & \texttt{1} \\
 & & \texttt{1} \\
 & & \texttt{1} \\
 & & \texttt{1} \\
 & & \texttt{1} \\
 & & \texttt{1} \\
 & & \texttt{1} \\
 & & \texttt{1} \\
 & & \texttt{1} \\
 & & \texttt{1} \\
 & & \texttt{1} \\
 & & \texttt{1} \\
 & & \texttt{1} \\
 & & \texttt{1} \\
 & & \texttt{1} \\
 & & \texttt{1} \\
 & & \texttt{1} \\
 & & \texttt{1} \\
 & & \texttt{1} \\
 & & \texttt{1} \\
 & & \texttt{1} \\
 & & \texttt{1} \\
 & & \texttt{1} \\
 & & \texttt{1} \\
 & & \texttt{1} \\
 & & \texttt{1} \\
 & & \texttt{1} \\
 & & \texttt{1} \\
 & & \texttt{1} \\
 & & \texttt{1} \\
 & & \texttt{1} \\
 & & \texttt{1} \\
 & & \texttt{1} \\
 & & \texttt{1} \\
 & & \texttt{1} \\
 & & \texttt{1} \\
 & & \texttt{1} \\
 & & \texttt{1} \\
 & & \texttt{1} \\
 & & \texttt{1} \\
 & & \texttt{1} \\
 & & \texttt{1} \\
 & & \texttt{1} \\
 & & \texttt{1} \\
 & & \texttt{1} \\
 & & \texttt{1} \\
 & & \texttt{1} \\
 & & \texttt{1} \\
 & & \texttt{1} \\
 & & \texttt{1} \\
 & & \texttt{1} \\
 & & \texttt{1} \\
 & & \texttt{1} \\
 & & \texttt{1} \\
 & & \texttt{1} \\
 & & \texttt{1} \\
 & & \texttt{1} \\
 & & \texttt{1} \\
 & & \texttt{1} \\
 & & \texttt{1} \\
 & & \texttt{1} \\
 & & \texttt{1} \\
 & & \texttt{1} \\
 & & \texttt{1} \\
 & & \texttt{1} \\
 & & \texttt{1} \\
 & & \texttt{1} \\
 & & \texttt{1} \\
 & & \texttt{1} \\
 & & \texttt{1} \\
 & & \texttt{1} \\
 & & \texttt{1} \\
 & & \texttt{1} \\
 & & \texttt{1} \\
 & & \texttt{1} \\
 & & \texttt{1} \\
 & & \texttt{1} \\
 & & \texttt{1} \\
 & & \texttt{1} \\
 & & \texttt{1} \\
 & & \texttt{1} \\
 & & \texttt{1} \\
 & & \texttt{1} \\
 & & \texttt{1} \\
 & & \texttt{1} \\
 & & \texttt{1} \\
 & & \texttt{1} \\
 & & \texttt{1} \\
 & & \texttt{1} \\
 & & \texttt{1} \\
 & & \texttt{1} \\
 & & \texttt{1} \\
 & & \texttt{1} \\
 & & \texttt{1} \\
 & & \texttt{1} \\
 & & \texttt{1} \\
 & & \texttt{1} \\
 & & \texttt{1} \\
 & & \texttt{1} \\
 & & \texttt{1} \\
 & & \texttt{1} \\
 & & \texttt{1} \\
 & & \texttt{1} \\
 & & \texttt{1} \\
 & & \texttt{1} \\
 & & \texttt{1} \\
 & & \texttt{1} \\
 & & \texttt{1} \\
 & & \texttt{1} \\
 & & \texttt{1} \\
 & & \texttt{1} \\
 & & \texttt{1} \\
 & & \texttt{1} \\
 & & \texttt{1} \\
 & & \texttt{1} \\
 & & \texttt{1} \\
 & & \texttt{1} \\
 & & \texttt{1} \\
 & & \texttt{1} \\
 & & \texttt{1} \\
 & & \texttt{1} \\
 & & \texttt{1} \\
 & & \texttt{1} \\
 & & \texttt{1} \\
 & & \texttt{1} \\
 & & \texttt{1} \\
 & & \texttt{1} \\
 & & \texttt{1} \\
 & & \texttt{1} \\
 & & \texttt{1} \\
 & & \texttt{1} \\
 & & \texttt{1} \\
 & & \texttt{1} \\
 & & \texttt{1} \\
 & & \texttt{1} \\
 & & \texttt{1} \\
 & & \texttt{1} \\
 & & \texttt{1} \\
 & & \texttt{1} \\
 & & \texttt{1} \\
 & & \texttt{1} \\
 & & \texttt{1} \\
 & & \texttt{1} \\
 & & \texttt{1} \\
 & & \texttt{1} \\
 & & \texttt{1} \\
 & & \texttt{1} \\
 & & \texttt{1} \\
 & & \texttt{1} \\
 & & \texttt{1} \\
 & & \texttt{1} \\
 & & \texttt{1} \\
 & & \texttt{1} \\
 & & \texttt{1} \\
 & & \texttt{1} \\
 & & \texttt{1} \\
 & & \texttt{1} \\
 & & \texttt{1} \\
 & & \texttt{1} \\
 & & \texttt{1} \\
 & & \texttt{1} \\
 & & \texttt{1} \\
 & & \texttt{1} \\
 & & \texttt{1} \\
 & & \texttt{1} \\
 & & \texttt{1} \\
 & & \texttt{1} \\
 & & \texttt{1} \\
 & & \texttt{1} \\
 & & \texttt{1} \\
 & & \texttt{1} \\
 & & \texttt{1} \\
 & & \texttt{1} \\
 & & \texttt{1} \\
 & & \texttt{1} \\
 & & \texttt{1} \\
 & & \texttt{1} \\
 & & \texttt{1} \\
 & & \texttt{1} \\
 & & \texttt{1} \\
 & & \texttt{1} \\
 & & \texttt{1} \\
 & & \texttt{1} \\
 & & \texttt{1} \\
 & & \texttt{1} \\
 & & \texttt{1} \\
 & & \texttt{1} \\
 & & \texttt{1} \\
 & & \texttt{1} \\
 & & \texttt{1} \\
 & & \texttt{1} \\
 & & \texttt{1} \\
 & & \texttt{1} \\
 & & \texttt{1} \\
 & & \texttt{1} \\
 & & \texttt{1} \\
 & & \texttt{1} \\
 & & \texttt{1} \\
 & & \texttt{1} \\
 & & \texttt{1} \\
 & & \texttt{1} \\
 & & \texttt{1} \\
 & & \texttt{1} \\
 & & \texttt{1} \\
 & & \texttt{1} \\
 & & \texttt{1} \\
 & & \texttt{1} \\
 & & \texttt{1} \\
 & & \texttt{1} \\
 & & \texttt{1} \\
 & & \texttt{1} \\
 & & \texttt{1} \\
 & & \texttt{1} \\
 & & \texttt{1} \\
 & & \texttt{1} \\
 & & \texttt{1} \\
 & & \texttt{1} \\
 & & \texttt{1} \\
 & & \texttt{1} \\
 & & \texttt{1} \\
 & & \texttt{1} \\
 & & \texttt{1} \\
 & & \texttt{1} \\
 & & \texttt{1} \\
 & & \texttt{1} \\
 & & \texttt{1} \\
 & & \texttt{1} \\
 & & \texttt{1} \\
 & & \texttt{1} \\
 & & \texttt{1} \\
 & & \texttt{1} \\
 & & \texttt{1} \\
 & & \texttt{1} \\
 & & \texttt{1} \\
 & & \texttt{1} \\
 & & \texttt{1} \\
 & & \texttt{1} \\
 & & \texttt{1} \\
 & & \texttt{1} \\
 & & \texttt{1} \\
 & & \texttt{1} \\
 & & \texttt{1} \\
 & & \texttt{1} \\
 & & \texttt{1} \\
 & & \texttt{1} \\
 & & \texttt{1} \\
 & & \texttt{1} \\
 & & \texttt{1} \\
 & & \texttt{1} \\
 & & \texttt{1} \\
 & & \texttt{1} \\
 & & \texttt{1} \\
 & & \texttt{1} \\
 & & \texttt{1} \\
 & & \texttt{1} \\
 & & \texttt{1} \\
 & & \texttt{1} \\
 & & \texttt{1} \\
 & & \texttt{1} \\
 & & \texttt{1} \\
 & & \texttt{1} \\
 & & \texttt{1} \\
 & & \texttt{1} \\
 & & \texttt{1} \\
 & & \texttt{1} \\
 & & \texttt{1} \\
 & & \texttt{1} \\
 & & \texttt{1} \\
 & & \texttt{1} \\
 & & \texttt{1} \\
 & & \texttt{1} \\
 & & \texttt{1} \\
 & & \texttt{1} \\
 & & \texttt{1} \\
 & & \texttt{1} \\
 & & \texttt{1} \\
 & & \texttt{1} \\
 & & \texttt{1} \\
 & & \texttt{1} \\
 & & \texttt{1} \\
 & & \texttt{1} \\
 & & \texttt{1} \\
 & & \texttt{1} \\
 & & \texttt{1} \\
 & & \texttt{1} \\
 & & \texttt{1} \\
 & & \texttt{1} \\
 & & \texttt{1} \\
 & & \texttt{1} \\
 & & \texttt{1} \\
 & & \texttt{1} \\
 & & \texttt{1} \\
 & & \texttt{1} \\
 & & \texttt{1} \\
 & & \texttt{1} \\
 & & \texttt{1} \\
 & & \texttt{1} \\
 & & \texttt{1} \\
 & & \texttt{1} \\
 & & \texttt{1} \\
 & & \texttt{1} \\
 & & \texttt{1} \\
 & & \texttt{1} \\
 & & \texttt{1} \\
 & & \texttt{1} \\
 & & \texttt{1} \\
 & & \texttt{1} \\
 & & \texttt{1} \\
 & & \texttt{1} \\
 & & \texttt{1} \\
 & & \texttt{1} \\
 & & \texttt{1} \\
 & & \texttt{1} \\
 & & \texttt{1} \\
 & & \texttt{1} \\
 & & \texttt{1} \\
 & & \texttt{1} \\
 & & \texttt{1} \\
 & & \texttt{1} \\
 & & \texttt{1} \\
 & & \texttt{1} \\
 & & \texttt{1} \\
 & & \texttt{1} \\
 & & \texttt{1} \\
 & & \texttt{1} \\
 & & \texttt{1} \\
 & & \texttt{1} \\
 & & \texttt{1} \\
 & & \texttt{1} \\
 & & \texttt{1} \\
 & & \texttt{1} \\
 & & \texttt{1} \\
 & & \texttt{1} \\
 & & \texttt{1} \\
 & & \texttt{1} \\
 & & \texttt{1} \\
 & & \texttt{1} \\
 & & \texttt{1} \\
 & & \texttt{1} \\
 & & \texttt{1} \\
 & & \texttt{1} \\
 & & \texttt{1} \\
 & & \texttt{1} \\
 & & \texttt{1} \\
 & & \texttt{1} \\
 & & \texttt{1} \\
 & & \texttt{1} \\
 & & \texttt{1} \\
 & & \texttt{1} \\
 & & \texttt{1} \\
 & & \texttt{1} \\
 & & \texttt{1} \\
 & & \texttt{1} \\
 & & \texttt{1} \\
 & & \texttt{1} \\
 & & \texttt{1} \\
 & & \texttt{1} \\
 & & \texttt{1} \\
 & & \texttt{1} \\
 & & \texttt{1} \\
 & & \texttt{1} \\
 & & \texttt{1} \\
 & & \texttt{1} \\
 & & \texttt{1} \\
 & & \texttt{1} \\
 & & \texttt{1} \\
 & & \texttt{1} \\
 & & \texttt{1} \\
 & & \texttt{1} \\
 & & \texttt{1} \\
 & & \texttt{1} \\
 & & \texttt{1} \\
 & & \texttt{1} \\
 & & \texttt{1} \\
 & & \texttt{1} \\
 & & \texttt{1} \\
 & & \texttt{1} \\
 & & \texttt{1} \\
 & & \texttt{1} \\
 & & \texttt{1} \\
 & & \texttt{1} \\
 & & \texttt{1} \\
 & & \texttt{1} \\
 & & \texttt{1} \\
 & & \texttt{1} \\
 & & \texttt{1} \\
 & & \texttt{1} \\
 & & \texttt{1} \\
 & & \texttt{1} \\
 & & \texttt{1} \\
 & & \texttt{1} \\
 & & \texttt{1} \\
 & & \texttt{1} \\
 & & \texttt{1} \\
 & & \texttt{1} \\
 & & \texttt{1} \\
 & & \texttt{1} \\
 & & \texttt{1} \\
 & & \texttt{1} \\
 & & \texttt{1} \\
 & & \texttt{1} \\
 & & \texttt{1} \\
 & & \texttt{1} \\
 & & \texttt{1} \\
 & & \texttt{1} \\
 & & \texttt{1} \\
 & & \texttt{1} \\
 & & \texttt{1} \\
 & & \texttt{1} \\
 & & \texttt{1} \\
 & & \texttt{1} \\
 & & \texttt{1} \\
 & & \texttt{1} \\
 & & \texttt{1} \\
 & & \texttt{1} \\
 & & \texttt{1} \\
 & & \texttt{1} \\
 & & \texttt{1} \\
 & & \texttt{1} \\
 & & \texttt{1} \\
 & & \texttt{1} \\
 & & \texttt{1} \\
 & & \texttt{1} \\
 & & \texttt{1} \\
 & & \texttt{1} \\
 & & \texttt{1} \\
 & & \texttt{1} \\
 & & \texttt{1} \\
 & & \texttt{1} \\
 & & \texttt{1} \\
 & & \texttt{1} \\
 & & \texttt{1} \\
 & & \texttt{1} \\
 & & \texttt{1} \\
 & & \texttt{1} \\
 & & \texttt{1} \\
 & & \texttt{1} \\
 & & \texttt{1} \\
 & & \texttt{1} \\
 & & \texttt{1} \\
 & & \texttt{1} \\
 & & \texttt{1} \\
 & & \texttt{1} \\
 & & \texttt{1} \\
 & & \texttt{1} \\
 & & \texttt{1} \\
 & & \texttt{1} \\
 & & \texttt{1} \\
 & & \texttt{1} \\
 & & \texttt{1} \\
 & & \texttt{1} \\
 & & \texttt{1} \\
 & & \texttt{1} \\
 & & \texttt{1} \\
 & & \texttt{1} \\
 & & \texttt{1} \\
 & & \texttt{1} \\
 & & \texttt{1} \\
 & & \texttt{1} \\
 & & \texttt{1} \\
 & & \texttt{1} \\
 & & \texttt{1} \\
 & & \texttt{1} \\
 & & \texttt{1} \\
 & & \texttt{1} \\
 & & \texttt{1} \\
 & & \texttt{1} \\
 & & \texttt{1} \\
 & & \texttt{1} \\
 & & \texttt{1} \\
 & & \texttt{1} \\
 & & \texttt{1} \\
 & & \texttt{1} \\
 & & \texttt{1} \\
 & & \texttt{1} \\
 & & \texttt{1} \\
 & & \texttt{1} \\
 & & \texttt{1} \\
 & & \texttt{1} \\
 & & \texttt{1} \\
 & & \texttt{1} \\
 & & \texttt{1} \\
 & & \texttt{1} \\
 & & \texttt{1} \\
 & & \texttt{1} \\
 & & \texttt{1} \\
 & & \texttt{1} \\
 & & \texttt{1} \\
 & & \texttt{1} \\
 & & \texttt{1} \\
 & & \texttt{1} \\
 & & \texttt{1} \\
 & & \texttt{1} \\
 & & \texttt{1} \\
 & & \texttt{1} \\
 & & \texttt{1} \\
 & & \texttt{1} \\
 & & \texttt{1} \\
 & & \texttt{1} \\
 & & \texttt{1} \\
 & & \texttt{1} \\
 & & \texttt{1} \\
 & & \texttt{1} \\
 & & \texttt{1} \\
 & & \texttt{1} \\
 & & \texttt{1} \\
 & & \texttt{1} \\
 & & \texttt{1} \\
 & & \texttt{1} \\
 & & \texttt{1} \\
 & & \texttt{1} \\
 & & \texttt{1} \\
 & & \texttt{1} \\
 & & \texttt{1} \\
 & & \texttt{1} \\
 & & \texttt{1} \\
 & & \texttt{1} \\
 & & \texttt{1} \\
 & & \texttt{1} \\
 & & \texttt{1} \\
 & & \texttt{1} \\
 & & \texttt{1} \\
 & & \texttt{1} \\
 & & \texttt{1} \\
 & & \texttt{1} \\
 & & \texttt{1} \\
 & & \texttt{1} \\
 & & \texttt{1} \\
 & & \texttt{1} \\
 & & \texttt{1} \\
 & & \texttt{1} \\
 & & \texttt{1} \\
 & & \texttt{1} \\
 & & \texttt{1} \\
 & & \texttt{1} \\
 & & \texttt{1} \\
 & & \texttt{1} \\
 & & \texttt{1} \\
 & & \texttt{1} \\
 & & \texttt{1} \\
 & & \texttt{1} \\
 & & \texttt{1} \\
 & & \texttt{1} \\
 & & \texttt{1} \\
 & & \texttt{1} \\
 & & \texttt{1} \\
 & & \texttt{1} \\
 & & \texttt{1} \\
 & & \texttt{1} \\
 & & \texttt{1} \\
 & & \texttt{1} \\
 & & \texttt{1} \\
 & & \texttt{1} \\
 & & \texttt{1} \\
 & & \texttt{1} \\
 & & \texttt{1} \\
 & & \texttt{1} \\
 & & \texttt{1} \\
 & & \texttt{1} \\
 & & \texttt{1} \\
 & & \texttt{1} \\
 & & \texttt{1} \\
 & & \texttt{1} \\
 & & \texttt{1} \\
 & & \texttt{1} \\
 & & \texttt{1} \\
 & & \texttt{1} \\
 & & \texttt{1} \\
 & & \texttt{1} \\
 & & \texttt{1} \\
 & & \texttt{1} \\
 & & \texttt{1} \\
 & & \texttt{1} \\
 & & \texttt{1} \\
 & & \texttt{1} \\
 & & \texttt{1} \\
 & & \texttt{1} \\
 & & \texttt{1} \\
 & & \texttt{1} \\
 & & \texttt{1} \\
 & & \texttt{1} \\
 & & \texttt{1} \\
 & & \texttt{1} \\
 & & \texttt{1} \\
 & & \texttt{1} \\
 & & \texttt{1} \\
 & & \texttt{1} \\
 & & \texttt{1} \\
 & & \texttt{1} \\
 & & \texttt{1} \\
 & & \texttt{1} \\
 & & \texttt{1} \\
 & & \texttt{1} \\
 & & \texttt{1} \\
 & & \texttt{1} \\
 & & \texttt{1} \\
 & & \texttt{1} \\
 & & \texttt{1} \\
 & & \texttt{1} \\
 & & \texttt{1} \\
 & & \texttt{1} \\
 & & \texttt{1} \\
 & & \texttt{1} \\
 & & \texttt{1} \\
 & & \texttt{1} \\
 & & \texttt{1} \\
 & & \texttt{1} \\
 & & \texttt{1} \\
 & & \texttt{1} \\
 & & \texttt{1} \\
 & & \texttt{1} \\
 & & \texttt{1} \\
 & & \texttt{1} \\
 & & \texttt{1} \\
 & & \texttt{1} \\
 & & \texttt{1} \\
 & & \texttt{1} \\
 & & \texttt{1} \\
 & & \texttt{1} \\
 & & \texttt{1} \\
 & & \texttt{1} \\
 & & \texttt{1} \\
 & & \texttt{1} \\
 & & \texttt{1} \\
 & & \texttt{1} \\
 & & \texttt{1} \\
 & & \texttt{1} \\
 & & \texttt{1} \\
 & & \texttt{1} \\
 & & \texttt{1} \\
 & & \texttt{1} \\
 & & \texttt{1} \\
 & & \texttt{1} \\
 & & \texttt{1} \\
 & & \texttt{1} \\
 & & \texttt{1} \\
 & & \texttt{1} \\
 & & \texttt{1} \\
 & & \texttt{1} \\
 & & \texttt{1} \\
 & & \texttt{1} \\
 & & \texttt{1} \\
 & & \texttt{1} \\
 & & \texttt{1} \\
 & & \texttt{1} \\
 & & \texttt{1} \\
 & & \texttt{1} \\
 & & \texttt{1} \\
 & & \texttt{1} \\
 & & \texttt{1} \\
 & & \texttt{1} \\
 & & \texttt{1} \\
 & & \texttt{1} \\
 & & \texttt{1} \\
 & & \texttt{1} \\
 & & \texttt{1} \\
 & & \texttt{1} \\
 & & \texttt{1} \\
 & & \texttt{1} \\
 & & \texttt{1} \\
 & & \texttt{1} \\
 & & \texttt{1} \\
 & & \texttt{1} \\
 & & \texttt{1} \\
 & & \texttt{1} \\
 & & \texttt{1} \\
 & & \texttt{1} \\
 & & \texttt{1} \\
 & & \texttt{1} \\
 & & \texttt{1} \\
 & & \texttt{1} \\
 & & \texttt{1} \\
 & & \texttt{1} \\
 & & \texttt{1} \\
 & & \texttt{1} \\
 & & \texttt{1} \\
 & & \texttt{1} \\
 & & \texttt{1} \\
 & & \texttt{1} \\
 & & \texttt{1} \\
 & & \texttt{1} \\
 & & \texttt{1} \\
 & & \texttt{1} \\
 & & \texttt{1} \\
 & & \texttt{1} \\
 & & \texttt{1} \\
 & & \texttt{1} \\
 & & \texttt{1} \\
 & & \texttt{1} \\
 & & \texttt{1} \\
 & & \texttt{1} \\
 & & \texttt{1} \\
 & & \texttt{1} \\
 & & \texttt{1} \\
 & & \texttt{1} \\
 & & \texttt{1} \\
 & & \texttt{1} \\
 & & \texttt{1} \\
 & & \texttt{1} \\
 & & \texttt{1} \\
 & & \texttt{1} \\
 & & \texttt{1} \\
 & & \texttt{1} \\
 & & \texttt{1} \\
 & & \texttt{1} \\
 & & \texttt{1} \\
 & & \texttt{1} \\
 & & \texttt{1} \\
 & & \texttt{1} \\
 & & \texttt{1} \\
 & & \texttt{1} \\
 & & \texttt{1} \\
 & & \texttt{1} \\
 & & \texttt{1} \\
 & & \texttt{1} \\
 & & \texttt{1} \\
 & & \texttt{1} \\
 & & \texttt{1} \\
 & & \texttt{1} \\
 & & \texttt{1} \\
 & & \texttt{1} \\
 & & \texttt{1} \\
 & & \texttt{1} \\
 & & \texttt{1} \\
 & & \texttt{1} \\
 & & \texttt{1} \\
 & & \texttt{1} \\
 & & \texttt{1} \\
 & & \texttt{1} \\
 & & \texttt{1} \\
 & & \texttt{1} \\
 & & \texttt{1} \\
 & & \texttt{1} \\
 & & \texttt{1} \\
 & & \texttt{1} \\
 & & \texttt{1} \\
 & & \texttt{1} \\
 & & \texttt{1} \\
 & & \texttt{1} \\
 & & \texttt{1} \\
 & & \texttt{1} \\
 & & \texttt{1} \\
 & & \texttt{1} \\
 & & \texttt{1} \\
 & & \texttt{1} \\
 & & \texttt{1} \\
 & & \texttt{1} \\
 & & \texttt{1} \\
 & & \texttt{1} \\
 & & \texttt{1} \\
 & & \texttt{1} \\
 & & \texttt{1} \\
 & & \texttt{1} \\
 & & \texttt{1} \\
 & & \texttt{1} \\
 & & \texttt{1} \\
 & & \texttt{1} \\
 & & \texttt{1} \\
 & & \texttt{1} \\
 & & \texttt{1} \\
 & & \texttt{1} \\
 & & \texttt{1} \\
 & & \texttt{1} \\
 & & \texttt{1} \\
 & & \texttt{1} \\
 & & \texttt{1} \\
 & & \texttt{1} \\
 & & \texttt{1} \\
 & & \texttt{1} \\
 & & \texttt{1} \\
 & & \texttt{1} \\
 & & \texttt{1} \\
 & & \texttt{1} \\
 & & \texttt{1} \\
 & & \texttt{1} \\
 & & \texttt{1} \\
 & & \texttt{1} \\
 & & \texttt{1} \\
 & & \texttt{1} \\
 & & \texttt{1} \\
 & & \texttt{1} \\
 & & \texttt{1} \\
 & & \texttt{1} \\
 & & \texttt{1} \\
 & & \texttt{1} \\
 & & \texttt{1} \\
 & & \texttt{1} \\
 & & \texttt{1} \\
 & & \texttt{1} \\
 & & \texttt{1} \\
 & & \texttt{1} \\
 & & \texttt{1} \\
 & & \texttt{1} \\
 & & \texttt{1} \\
 & & \texttt{1} \\
 & & \texttt{1} \\
 & & \texttt{1} \\
 & & \texttt{1} \\
 & & \texttt{1} \\
 & & \texttt{1} \\
 & & \texttt{1} \\
 & & \texttt{1} \\
 & & \texttt{1} \\
 & & \texttt{1} \\
 & & \texttt{1} \\
 & & \texttt{1} \\
 & & \texttt{1} \\
 & & \texttt{1} \\
 & & \texttt{1} \\
 & & \texttt{1} \\
 & & \texttt{1} \\
 & & \texttt{1} \\
 & & \texttt{1} \\
 & & \texttt{1} \\
 & & \texttt{1} \\
 & & \texttt{1} \\
 & & \texttt{1} \\
 & & \texttt{1} \\
 & & \texttt{1} \\
 & & \texttt{1} \\
 & & \texttt{1} \\
 & & \texttt{1} \\
 & & \texttt{1} \\
 & & \texttt{1} \\
 & & \texttt{1} \\
 & & \texttt{1} \\
 & & \texttt{1} \\
 & & \texttt{1} \\
 & & \texttt{1} \\
 & & \texttt{1} \\
 & & \texttt{1} \\
 & & \texttt{1} \\
 & & \texttt{1} \\
 & & \texttt{1} \\
 & & \texttt{1} \\
 & & \texttt{1} \\
 & & \texttt{1} \\
 & & \texttt{1} \\
 & & \texttt{1} \\
 & & \texttt{1} \\
 & & \texttt{1} \\
 & & \texttt{1} \\
 & & \texttt{1} \\
 & & \texttt{1} \\
 & & \texttt{1} \\
 & & \texttt{1} \\
 & & \texttt{1} \\
 & & \texttt{1} \\
 & & \texttt{1} \\
 & & \texttt{1} \\
 & & \texttt{1} \\
 & & \texttt{1} \\
 & & \texttt{1} \\
 & & \texttt{1} \\
 & & \texttt{1} \\
 & & \texttt{1} \\
 & & \texttt{1} \\
 & & \texttt{1} \\
 & & \texttt{1} \\
 & & \texttt{1} \\
 & & \texttt{1} \\
 & & \texttt{1} \\
 & & \texttt{1} \\
 & & \texttt{1} \\
 & & \texttt{1} \\
 & & \texttt{1} \\
 & & \texttt{1} \\
 & & \texttt{1} \\
 & & \texttt{1} \\
 & & \texttt{1} \\
 & & \texttt{1} \\
 & & \texttt{1} \\
 & & \texttt{1} \\
 & & \texttt{1} \\
 & & \texttt{1} \\
 & & \texttt{1} \\
 & & \texttt{1} \\
 & & \texttt{1} \\
 & & \texttt{1} \\
 & & \texttt{1} \\
 & & \texttt{1} \\
 & & \texttt{1} \\
 & & \texttt{1} \\
 & & \texttt{1} \\
 & & \texttt{1} \\
 & & \texttt{1} \\
 & & \texttt{1} \\
 & & \texttt{1} \\
 & & \texttt{1} \\
 & & \texttt{1} \\
 & & \texttt{1} \\
 & & \texttt{1} \\
 & & \texttt{1} \\
 & & \texttt{1} \\
 & & \texttt{1} \\
 & & \texttt{1} \\
 & & \texttt{1} \\
 & & \texttt{1} \\
 & & \texttt{1} \\
 & & \texttt{1} \\
 & & \texttt{1} \\
 & & \texttt{1} \\
 & & \texttt{1} \\
 & & \texttt{1} \\
 & & \texttt{1} \\
 & & \texttt{1} \\
 & & \texttt{1} \\
 & & \texttt{1} \\
 & & \texttt{1} \\
 & & \texttt{1} \\
 & & \texttt{1} \\
 & & \texttt{1} \\
 & & \texttt{1} \\
 & & \texttt{1} \\
 & & \texttt{1} \\
 & & \texttt{1} \\
 & & \texttt{1} \\
 & & \texttt{1} \\
 & & \texttt{1} \\
 & & \texttt{1} \\
 & & \texttt{1} \\
 & & \texttt{1} \\
 & & \texttt{1} \\
 & & \texttt{1} \\
 & & \texttt{1} \\
 & & \texttt{1} \\
 & & \texttt{1} \\
 & & \texttt{1} \\
 & & \texttt{1} \\
 & & \texttt{1} \\
 & & \texttt{1} \\
 & & \texttt{1} \\
 & & \texttt{1} \\
 & & \texttt{1} \\
 & & \texttt{1} \\
 & & \texttt{1} \\
 & & \texttt{1} \\
 & & \texttt{1} \\
 & & \texttt{1} \\
 & & \texttt{1} \\
 & & \texttt{1} \\
 & & \texttt{1} \\
 & & \texttt{1} \\
 & & \texttt{1} \\
 & & \texttt{1} \\
 & & \texttt{1} \\
 & & \texttt{1} \\
 & & \texttt{1} \\
 & & \texttt{1} \\
 & & \texttt{1} \\
 & & \texttt{1} \\
 & & \texttt{1} \\
 & & \texttt{1} \\
 & & \texttt{1} \\
 & & \texttt{1} \\
 & & \texttt{1} \\
 & & \texttt{1} \\
 & & \texttt{1} \\
 & & \texttt{1} \\
 & & \texttt{1} \\
 & & \texttt{1} \\
 & & \texttt{1} \\
 & & \texttt{1} \\
 & & \texttt{1} \\
 & & \texttt{1} \\
 & & \texttt{1} \\
 & & \texttt{1} \\
 & & \texttt{1} \\
 & & \texttt{1} \\
 & & \texttt{1} \\
 & & \texttt{1} \\
 & & \texttt{1} \\
 & & \texttt{1} \\
 & & \texttt{1} \\
 & & \texttt{1} \\
 & & \texttt{1} \\
 & & \texttt{1} \\
 & & \texttt{1} \\
 & & \texttt{1} \\
 & & \texttt{1} \\
 & & \texttt{1} \\
 & & \texttt{1} \\
 & & \texttt{1} \\
 & & \texttt{1} \\
 & & \texttt{1} \\
 & & \texttt{1} \\
 & & \texttt{1} \\
 & & \texttt{1} \\
 & & \texttt{1} \\
 & & \texttt{1} \\
 & & \texttt{1} \\
 & & \texttt{1} \\
 & & \texttt{1} \\
 & & \texttt{1} \\
 & & \texttt{1} \\
 & & \texttt{1} \\
 & & \texttt{1} \\
 & & \texttt{1} \\
 & & \texttt{1} \\
 & & \texttt{1} \\
 & & \texttt{1} \\
 & & \texttt{1} \\
 & & \texttt{1} \\
 & & \texttt{1} \\
 & & \texttt{1} \\
 & & \texttt{1} \\
 & & \texttt{1} \\
 & & \texttt{1} \\
 & & \texttt{1} \\
 & & \texttt{1} \\
 & & \texttt{1} \\
 & & \texttt{1} \\
 & & \texttt{1} \\
 & & \texttt{1} \\
 & & \texttt{1} \\
 & & \texttt{1} \\
 & & \texttt{1} \\
 & & \texttt{1} \\
 & & \texttt{1} \\
 & & \texttt{1} \\
 & & \texttt{1} \\
 & & \texttt{1} \\
 & & \texttt{1} \\
 & & \texttt{1} \\
 & & \texttt{1} \\
 & & \texttt{1} \\
 & & \texttt{1} \\
 & & \texttt{1} \\
 & & \texttt{1} \\
 & & \texttt{1} \\
 & & \texttt{1} \\
 & & \texttt{1} \\
 & & \texttt{1} \\
 & & \texttt{1} \\
 & & \texttt{1} \\
 & & \texttt{1} \\
 & & \texttt{1} \\
 & & \texttt{1} \\
 & & \texttt{1} \\
 & & \texttt{1} \\
 & & \texttt{1} \\
 & & \texttt{1} \\
 & & \texttt{1} \\
 & & \texttt{1} \\
 & & \texttt{1} \\
 & & \texttt{1} \\
 & & \texttt{1} \\
 & & \texttt{1} \\
 & & \texttt{1} \\
 & & \texttt{1} \\
 & & \texttt{1} \\
 & & \texttt{1} \\
 & & \texttt{1} \\
 & & \texttt{1} \\
 & & \texttt{1} \\
 & & \texttt{1} \\
 & & \texttt{1} \\
 & & \texttt{1} \\
 & & \texttt{1} \\
 & & \texttt{1} \\
 & & \texttt{1} \\
 & & \texttt{1} \\
 & & \texttt{1} \\
 & & \texttt{1} \\
 & & \texttt{1} \\
 & & \texttt{1} \\
 & & \texttt{1} \\
 & & \texttt{1} \\
 & & \texttt{1} \\
 & & \texttt{1} \\
 & & \texttt{1} \\
 & & \texttt{1} \\
 & & \texttt{1} \\
 & & \texttt{1} \\
 & & \texttt{1} \\
 & & \texttt{1} \\
 & & \texttt{1} \\
 & & \texttt{1} \\
 & & \texttt{1} \\
 & & \texttt{1} \\
 & & \texttt{1} \\
 & & \texttt{1} \\
 & & \texttt{1} \\
 & & \texttt{1} \\
 & & \texttt{1} \\
 & & \texttt{1} \\
 & & \texttt{1} \\
 & & \texttt{1} \\
 & & \texttt{1} \\
 & & \texttt{1} \\
 & & \texttt{1} \\
 & & \texttt{1} \\
 & & \texttt{1} \\
 & & \texttt{1} \\
 & & \texttt{1} \\
 & & \texttt{1} \\
 & & \texttt{1} \\
 & & \texttt{1} \\
 & & \texttt{1} \\
 & & \texttt{1} \\
 & & \texttt{1} \\
 & & \texttt{1} \\
 & & \texttt{1} \\
 & & \texttt{1} \\
 & & \texttt{1} \\
 & & \texttt{1} \\
 & & \texttt{1} \\
 & & \texttt{1} \\
 & & \texttt{1} \\
 & & \texttt{1} \\
 & & \texttt{1} \\
 & & \texttt{1} \\
 & & \texttt{1} \\
 & & \texttt{1} \\
 & & \texttt{1} \\
 & & \texttt{1} \\
 & & \texttt{1} \\
 & & \texttt{1} \\
 & & \texttt{1} \\
 & & \texttt{1} \\
 & & \texttt{1} \\
 & & \texttt{1} \\
 & & \texttt{1} \\
 & & \texttt{1} \\
 & & \texttt{1} \\
 & & \texttt{1} \\
 & & \texttt{1} \\
 & & \texttt{1} \\
 & & \texttt{1} \\
 & & \texttt{1} \\
 & & \texttt{1} \\
 & & \texttt{1} \\
 & & \texttt{1} \\
 & & \texttt{1} \\
 & & \texttt{1} \\
 & & \texttt{1} \\
 & & \texttt{1} \\
 & & \texttt{1} \\
 & & \texttt{1} \\
 & & \texttt{1} \\
 & & \texttt{1} \\
 & & \texttt{1} \\
 & & \texttt{1} \\
 & & \texttt{1} \\
 & & \texttt{1} \\
 & & \texttt{1} \\
 & & \texttt{1} \\
 & & \texttt{1} \\
 & & \texttt{1} \\
 & & \texttt{1} \\
 & & \texttt{1} \\
 & & \texttt{1} \\
 & & \texttt{1} \\
 & & \texttt{1} \\
 & & \texttt{1} \\
 & & \texttt{1} \\
 & & \texttt{1} \\
 & & \texttt{1} \\
 & & \texttt{1} \\
 & & \texttt{1} \\
 & & \texttt{1} \\
 & & \texttt{1} \\
 & & \texttt{1} \\
 & & \texttt{1} \\
 & & \texttt{1} \\
 & & \texttt{1} \\
 & & \texttt{1} \\
 & & \texttt{1} \\
 & & \texttt{1} \\
 & & \texttt{1} \\
 & & \texttt{1} \\
 & & \texttt{1} \\
 & & \texttt{1} \\
 & & \texttt{1} \\
 & & \texttt{1} \\
 & & \texttt{1} \\
 & & \texttt{1} \\
 & & \texttt{1} \\
 & & \texttt{1} \\
 & & \texttt{1} \\
 & & \texttt{1} \\
 & & \texttt{1} \\
 & & \texttt{1} \\
 & & \texttt{1} \\
 & & \texttt{1} \\
 & & \texttt{1} \\
 & & \texttt{1} \\
 & & \texttt{1} \\
 & & \texttt{1} \\
 & & \texttt{1} \\
 & & \texttt{1} \\
 & & \texttt{1} \\
 & & \texttt{1} \\
 & & \texttt{1} \\
 & & \texttt{1} \\
 & & \texttt{1} \\
 & & \texttt{1} \\
 & & \texttt{1} \\
 & & \texttt{1} \\
 & & \texttt{1} \\
 & & \texttt{1} \\
 & & \texttt{1} \\
 & & \texttt{1} \\
 & & \texttt{1} \\
 & & \texttt{1} \\
 & & \texttt{1} \\
 & & \texttt{1} \\
 & & \texttt{1} \\
 & & \texttt{1} \\
 & & \texttt{1} \\
 & & \texttt{1} \\
 & & \texttt{1} \\
 & & \texttt{1} \\
 & & \texttt{1} \\
 & & \texttt{1} \\
 & & \texttt{1} \\
 & & \texttt{1} \\
 & & \texttt{1} \\
 & & \texttt{1} \\
 & & \texttt{1} \\
 & & \texttt{1} \\
 & & \texttt{1} \\
 & & \texttt{1} \\
 & & \texttt{1} \\
 & & \texttt{1} \\
 & & \texttt{1} \\
 & & \texttt{1} \\
 & & \texttt{1} \\
 & & \texttt{1} \\
 & & \texttt{1} \\
 & & \texttt{1} \\
 & & \texttt{1} \\
 & & \texttt{1} \\
 & & \texttt{1} \\
 & & \texttt{1} \\
 & & \texttt{1} \\
 & & \texttt{1} \\
 & & \texttt{1} \\
 & & \texttt{1} \\
 & & \texttt{1} \\
 & & \texttt{1} \\
 & & \texttt{1} \\
 & & \texttt{1} \\
 & & \texttt{1} \\
 & & \texttt{1} \\
 & & \texttt{1} \\
 & & \texttt{1} \\
 & & \texttt{1} \\
 & & \texttt{1} \\
 & & \texttt{1} \\
 & & \texttt{1} \\
 & & \texttt{1} \\
 & & \texttt{1} \\
 & & \texttt{1} \\
 & & \texttt{1} \\
 & & \texttt{1} \\
 & & \texttt{1} \\
 & & \texttt{1} \\
 & & \texttt{1} \\
 & & \texttt{1} \\
 & & \texttt{1} \\
 & & \texttt{1} \\
 & & \texttt{1} \\
 & & \texttt{1} \\
 & & \texttt{1} \\
 & & \texttt{1} \\
 & & \texttt{1} \\
 & & \texttt{1} \\
 & & \texttt{1} \\
 & & \texttt{1} \\
 & & \texttt{1} \\
 & & \texttt{1} \\
 & & \texttt{1} \\
 & & \texttt{1} \\
 & & \texttt{1} \\
 & & \texttt{1} \\
 & & \texttt{1} \\
 & & \texttt{1} \\
 & & \texttt{1} \\
 & & \texttt{1} \\
 & & \texttt{1} \\
 & & \texttt{1} \\
 & & \texttt{1} \\
 & & \texttt{1} \\
 & & \texttt{1} \\
 & & \texttt{1} \\
 & & \texttt{1} \\
 & & \texttt{1} \\
 & & \texttt{1} \\
 & & \texttt{1} \\
 & & \texttt{1} \\
 & & \texttt{1} \\
 & & \texttt{1} \\
 & & \texttt{1} \\
 & & \texttt{1} \\
 & & \texttt{1} \\
 & & \texttt{1} \\
 & & \texttt{1} \\
 & & \texttt{1} \\
 & & \texttt{1} \\
 & & \texttt{1} \\
 & & \texttt{1} \\
 & & \texttt{1} \\
 & & \texttt{1} \\
 & & \texttt{1} \\
 & & \texttt{1} \\
 & & \texttt{1} \\
 & & \texttt{1} \\
 & & \texttt{1} \\
 & & \texttt{1} \\
 & & \texttt{1} \\
 & & \texttt{1} \\
 & & \texttt{1} \\
 & & \texttt{1} \\
 & & \texttt{1} \\
 & & \texttt{1} \\
 & & \texttt{1} \\
 & & \texttt{1} \\
 & & \texttt{1} \\
 & & \texttt{1} \\
 & & \texttt{1} \\
 & & \texttt{1} \\
 & & \texttt{1} \\
 & & \texttt{1} \\
 & & \texttt{1} \\
 & & \texttt{1} \\
 & & \texttt{1} \\
 & & \texttt{1} \\
 & & \texttt{1} \\
 & & \texttt{1} \\
 & & \texttt{1} \\
 & & \texttt{1} \\
 & & \texttt{1} \\
 & & \texttt{1} \\
 & & \texttt{1} \\
 & & \texttt{1} \\
 & & \texttt{1} \\
 & & \texttt{1} \\
 & & \texttt{1} \\
 & & \texttt{1} \\
 & & \texttt{1} \\
 & & \texttt{1} \\
 & & \texttt{1} \\
 & & \texttt{1} \\
 & & \texttt{1} \\
 & & \texttt{1} \\
 & & \texttt{1} \\
 & & \texttt{1} \\
 & & \texttt{1} \\
 & & \texttt{1} \\
 & & \texttt{1} \\
 & & \texttt{1} \\
 & & \texttt{1} \\
 & & \texttt{1} \\
 & & \texttt{1} \\
 & & \texttt{1} \\
 & & \texttt{1} \\
 & & \texttt{1} \\
 & & \texttt{1} \\
 & & \texttt{1} \\
 & & \texttt{1} \\
 & & \texttt{1} \\
 & & \texttt{1} \\
 & & \texttt{1} \\
 & & \texttt{1} \\
 & & \texttt{1} \\
 & & \texttt{1} \\
 & & \texttt{1} \\
 & & \texttt{1} \\
 & & \texttt{1} \\
 & & \texttt{1} \\
 & & \texttt{1} \\
 & & \texttt{1} \\
 & & \texttt{1} \\
 & & \texttt{1} \\
 & & \texttt{1} \\
 & & \texttt{1} \\
 & & \texttt{1} \\
 & & \texttt{1} \\
 & & \texttt{1} \\
 & & \texttt{1} \\
 & & \texttt{1} \\
 & & \texttt{1} \\
 & & \texttt{1} \\
 & & \texttt{1} \\
 & & \texttt{1} \\
 & & \texttt{1} \\
 & & \texttt{1} \\
 & & \texttt{1} \\
 & & \texttt{1} \\
 & & \texttt{1} \\
 & & \texttt{1} \\
 & & \texttt{1} \\
 & & \texttt{1} \\
 & & \texttt{1} \\
 & & \texttt{1} \\
 & & \texttt{1} \\
 & & \texttt{1} \\
 & & \texttt{1} \\
 & & \texttt{1} \\
 & & \texttt{1} \\
 & & \texttt{1} \\
 & & \texttt{1} \\
 & & \texttt{1} \\
 & & \texttt{1} \\
 & & \texttt{1} \\
 & & \texttt{1} \\
 & & \texttt{1} \\
 & & \texttt{1} \\
 & & \texttt{1} \\
 & & \texttt{1} \\
 & & \texttt{1} \\
 & & \texttt{1} \\
 & & \texttt{1} \\
 & & \texttt{1} \\
 & & \texttt{1} \\
 & & \texttt{1} \\
 & & \texttt{1} \\
 & & \texttt{1} \\
 & & \texttt{1} \\
 & & \texttt{1} \\
 & & \texttt{1} \\
 & & \texttt{1} \\
 & & \texttt{1} \\
 & & \texttt{1} \\
 & & \texttt{1} \\
 & & \texttt{1} \\
 & & \texttt{1} \\
 & & \texttt{1} \\
 & & \texttt{1} \\
 & & \texttt{1} \\
 & & \texttt{1} \\
 & & \texttt{1} \\
 & & \texttt{1} \\
 & & \texttt{1} \\
 & & \texttt{1} \\
 & & \texttt{1} \\
 & & \texttt{1} \\
 & & \texttt{1} \\
 & & \texttt{1} \\
 & & \texttt{1} \\
 & & \texttt{1} \\
 & & \texttt{1} \\
 & & \texttt{1} \\
 & & \texttt{1} \\
 & & \texttt{1} \\
 & & \texttt{1} \\
 & & \texttt{1} \\
 & & \texttt{1} \\
 & & \texttt{1} \\
 & & \texttt{1} \\
 & & \texttt{1} \\
 & & \texttt{1} \\
 & & \texttt{1} \\
 & & \texttt{1} \\
 & & \texttt{1} \\
 & & \texttt{1} \\
 & & \texttt{1} \\
 & & \texttt{1} \\
 & & \texttt{1} \\
 & & \texttt{1} \\
 & & \texttt{1} \\
 & & \texttt{1} \\
 & & \texttt{1} \\
 & & \texttt{1} \\
 & & \texttt{1} \\
 & & \texttt{1} \\
 & & \texttt{1} \\
 & & \texttt{1} \\
 & & \texttt{1} \\
 & & \texttt{1} \\
 & & \texttt{1} \\
 & & \texttt{1} \\
 & & \texttt{1} \\
 & & \texttt{1} \\
 & & \texttt{1} \\
 & & \texttt{1} \\
 & & \texttt{1} \\
 & & \texttt{1} \\
 & & \texttt{1} \\
 & & \texttt{1} \\
 & & \texttt{1} \\
 & & \texttt{1} \\
 & & \texttt{1} \\
 & & \texttt{1} \\
 & & \texttt{1} \\
 & & \texttt{1} \\
 & & \texttt{1} \\
 & & \texttt{1} \\
 & & \texttt{1} \\
 & & \texttt{1} \\
 & & \texttt{1} \\
 & & \texttt{1} \\
 & & \texttt{1} \\
 & & \texttt{1} \\
 & & \texttt{1} \\
 & & \texttt{1} \\
 & & \texttt{1} \\
 & & \texttt{1} \\
 & & \texttt{1} \\
 & & \texttt{1} \\
 & & \texttt{1} \\
 & & \texttt{1} \\
 & & \texttt{1} \\
 & & \texttt{1} \\
 & & \texttt{1} \\
 & & \texttt{1} \\
 & & \texttt{1} \\
 & & \texttt{1} \\
 & & \texttt{1} \\
 & & \texttt{1} \\
 & & \texttt{1} \\
 & & \texttt{1} \\
 & & \texttt{1} \\
 & & \texttt{1} \\
 & & \texttt{1} \\
 & & \texttt{1} \\
 & & \texttt{1} \\
 & & \texttt{1} \\
 & & \texttt{1} \\
 & & \texttt{1} \\
 & & \texttt{1} \\
 & & \texttt{1} \\
 & & \texttt{1} \\
 & & \texttt{1} \\
 & & \texttt{1} \\
 & & \texttt{1} \\
 & & \texttt{1} \\
 & & \texttt{1} \\
 & & \texttt{1} \\
 & & \texttt{1} \\
 & & \texttt{1} \\
 & & \texttt{1} \\
 & & \texttt{1} \\
 & & \texttt{1} \\
 & & \texttt{1} \\
 & & \texttt{1} \\
 & & \texttt{1} \\
 & & \texttt{1} \\
 & & \texttt{1} \\
 & & \texttt{1} \\
 & & \texttt{1} \\
 & & \texttt{1} \\
 & & \texttt{1} \\
 & & \texttt{1} \\
 & & \texttt{1} \\
 & & \texttt{1} \\
 & & \texttt{1} \\
 & & \texttt{1} \\
 & & \texttt{1} \\
 & & \texttt{1} \\
 & & \texttt{1} \\
 & & \texttt{1} \\
 & & \texttt{1} \\
 & & \texttt{1} \\
 & & \texttt{1} \\
 & & \texttt{1} \\
 & & \texttt{1} \\
 & & \texttt{1} \\
 & & \texttt{1} \\
 & & \texttt{1} \\
 & & \texttt{1} \\
 & & \texttt{1} \\
 & & \texttt{1} \\
 & & \texttt{1} \\
 & & \texttt{1} \\
 & & \texttt{1} \\
 & & \texttt{1} \\
 & & \texttt{1} \\
 & & \texttt{1} \\
 & & \texttt{1} \\
 & & \texttt{1} \\
 & & \texttt{1} \\
 & & \texttt{1} \\
 & & \texttt{1} \\
 & & \texttt{1} \\
 & & \texttt{1} \\
 & & \texttt{1} \\
 & & \texttt{1} \\
 & & \texttt{1} \\
 & & \texttt{1} \\
 & & \texttt{1} \\
 & & \texttt{1} \\
 & & \texttt{1} \\
 & & \texttt{1} \\
 & & \texttt{1} \\
 & & \texttt{1} \\
 & & \texttt{1} \\
 & & \texttt{1} \\
 & & \texttt{1} \\
 & & \texttt{1} \\
 & & \texttt{1} \\
 & & \texttt{1} \\
 & & \texttt{1} \\
 & & \texttt{1} \\
 & & \texttt{1} \\
 & & \texttt{1} \\
 & & \texttt{1} \\
 & & \texttt{1} \\
 & & \texttt{1} \\
 & & \texttt{1} \\
 & & \texttt{1} \\
 & & \texttt{1} \\
 & & \texttt{1} \\
 & & \texttt{1} \\
 & & \texttt{1} \\
 & & \texttt{1} \\
 & & \texttt{1} \\
 & & \texttt{1} \\
 & & \texttt{1} \\
 & & \texttt{1} \\
 & & \texttt{1} \\
 & & \texttt{1} \\
 & & \texttt{1} \\
 & & \texttt{1} \\
 & & \texttt{1} \\
 & & \texttt{1} \\
 & & \texttt{1} \\
 & & \texttt{1} \\
 & & \texttt{1} \\
 & & \texttt{1} \\
 & & \texttt{1} \\
 & & \texttt{1} \\
 & & \texttt{1} \\
 & & \texttt{1} \\
 & & \texttt{1} \\
 & & \texttt{1} \\
 & & \texttt{1} \\
 & & \texttt{1} \\
 & & \texttt{1} \\
 & & \texttt{1} \\
 & & \texttt{1} \\
 & & \texttt{1} \\
 & & \texttt{1} \\
 & & \texttt{1} \\
 & & \texttt{1} \\
 & & \texttt{1} \\
 & & \texttt{1} \\
 & & \texttt{1} \\
 & & \texttt{1} \\
 & & \texttt{1} \\
 & & \texttt{1} \\
 & & \texttt{1} \\
 & & \texttt{1} \\
 & & \texttt{1} \\
 & & \texttt{1} \\
 & & \texttt{1} \\
 & & \texttt{1} \\
 & & \texttt{1} \\
 & & \texttt{1} \\
 & & \texttt{1} \\
 & & \texttt{1} \\
 & & \texttt{1} \\
 & & \texttt{1} \\
 & & \texttt{1} \\
 & & \texttt{1} \\
 & & \texttt{1} \\
 & & \texttt{1} \\
 & & \texttt{1} \\
 & & \texttt{1} \\
 & & \texttt{1} \\
 & & \texttt{1} \\
 & & \texttt{1} \\
 & & \texttt{1} \\
 & & \texttt{1} \\
 & & \texttt{1} \\
 & & \texttt{1} \\
 & & \texttt{1} \\
 & & \texttt{1} \\
 & & \texttt{1} \\
 & & \texttt{1} \\
 & & \texttt{1} \\
 & & \texttt{1} \\
 & & \texttt{1} \\
 & & \texttt{1} \\
 & & \texttt{1} \\
 & & \texttt{1} \\
 & & \texttt{1} \\
 & & \texttt{1} \\
 & & \texttt{1} \\
 & & \texttt{1} \\
 & & \texttt{1} \\
 & & \texttt{1} \\
 & & \texttt{1} \\
 & & \texttt{1} \\
 & & \texttt{1} \\
 & & \texttt{1} \\
 & & \texttt{1} \\
 & & \texttt{1} \\
 & & \texttt{1} \\
 & & \texttt{1} \\
 & & \texttt{1} \\
 & & \texttt{1} \\
 & & \texttt{1} \\
 & & \texttt{1} \\
 & & \texttt{1} \\
 & & \texttt{1} \\
 & & \texttt{1} \\
 & & \texttt{1} \\
 & & \texttt{1} \\
 & & \texttt{1} \\
 & & \texttt{1} \\
 & & \texttt{1} \\
 & & \texttt{1} \\
 & & \texttt{1} \\
 & & \texttt{1} \\
 & & \texttt{1} \\
 & & \texttt{1} \\
 & & \texttt{1} \\
 & & \texttt{1} \\
 & & \texttt{1} \\
 & & \texttt{1} \\
 & & \texttt{1} \\
 & & \texttt{1} \\
 & & \texttt{1} \\
 & & \texttt{1} \\
 & & \texttt{1} \\
 & & \texttt{1} \\
\hline \end{tabular}
}
\newpage
{\scriptsize
\begin{tabular}{|l|l|l|} \hline
\bf{Register Class} & \bf{Register File} & \bf{Registers} \\ \hline
 floatReg & RV\_FPR & \texttt{1 1} \\
 & & \texttt{1 1} \\
 & & \texttt{1 1} \\
 & & \texttt{1 1} \\
 & & \texttt{1 1} \\
 & & \texttt{1 1} \\
 & & \texttt{1 1} \\
 & & \texttt{1 1} \\
 & & \texttt{1 1} \\
 & & \texttt{1 1} \\
 & & \texttt{1 1} \\
 & & \texttt{1 1} \\
 & & \texttt{1 1} \\
 & & \texttt{1 1} \\
 & & \texttt{1 1} \\
 & & \texttt{1 1} \\
 & & \texttt{1 1} \\
 & & \texttt{1 1} \\
 & & \texttt{1 1} \\
 & & \texttt{1 1} \\
 & & \texttt{1 1} \\
 & & \texttt{1 1} \\
 & & \texttt{1 1} \\
 & & \texttt{1 1} \\
 & & \texttt{1 1} \\
 & & \texttt{1 1} \\
 & & \texttt{1 1} \\
 & & \texttt{1 1} \\
 & & \texttt{1 1} \\
 & & \texttt{1 1} \\
 & & \texttt{1 1} \\
 & & \texttt{1 1} \\
\hline \end{tabular}
}
\newpage
{\scriptsize
\begin{tabular}{|l|l|l|} \hline
\bf{Register Class} & \bf{Register File} & \bf{Registers} \\ \hline
 mainReg & RV\_BIR & \texttt{1 1} \\
 & & \texttt{1 1} \\
 & & \texttt{1 1} \\
 & & \texttt{1 1} \\
 & & \texttt{1 1} \\
 & & \texttt{1 1} \\
 & & \texttt{1 1} \\
 & & \texttt{1 1} \\
 & & \texttt{1 1} \\
 & & \texttt{1 1} \\
 & & \texttt{1 1} \\
 & & \texttt{1 1} \\
 & & \texttt{1 1} \\
 & & \texttt{1 1} \\
 & & \texttt{1 1} \\
 & & \texttt{1 1} \\
 & & \texttt{1 1} \\
 & & \texttt{1 1} \\
 & & \texttt{1 1} \\
 & & \texttt{1 1} \\
 & & \texttt{1 1} \\
 & & \texttt{1 1} \\
 & & \texttt{1 1} \\
 & & \texttt{1 1} \\
 & & \texttt{1 1} \\
 & & \texttt{1 1} \\
 & & \texttt{1 1} \\
 & & \texttt{1 1} \\
 & & \texttt{1 1} \\
 & & \texttt{1 1} \\
 & & \texttt{1 1} \\
 & & \texttt{1 1} \\
\hline \end{tabular}
}
\newpage
{\scriptsize
\begin{tabular}{|l|l|l|} \hline
\bf{Register Class} & \bf{Register File} & \bf{Registers} \\ \hline
 mainRegPair & RV\_BIRP & \texttt{1} \\
 & & \texttt{1} \\
 & & \texttt{1} \\
 & & \texttt{1} \\
 & & \texttt{1} \\
 & & \texttt{1} \\
 & & \texttt{1} \\
 & & \texttt{1} \\
 & & \texttt{1} \\
 & & \texttt{1} \\
 & & \texttt{1} \\
 & & \texttt{1} \\
 & & \texttt{1} \\
 & & \texttt{1} \\
 & & \texttt{1} \\
 & & \texttt{1} \\
\hline \end{tabular}
}
\newpage
{\scriptsize
\begin{tabular}{|l|l|l|} \hline
\bf{Register Class} & \bf{Register File} & \bf{Registers} \\ \hline
 onlyfxReg & SFR & \texttt{ } \\
 & & \texttt{ } \\
 & & \texttt{ } \\
 & & \texttt{ } \\
 & & \texttt{ } \\
 & & \texttt{ } \\
 & & \texttt{ } \\
 & & \texttt{ } \\
 & & \texttt{ } \\
 & & \texttt{ } \\
 & & \texttt{ } \\
 & & \texttt{ } \\
 & & \texttt{ } \\
 & & \texttt{ } \\
 & & \texttt{ } \\
 & & \texttt{ } \\
 & & \texttt{ } \\
 & & \texttt{ } \\
 & & \texttt{ } \\
 & & \texttt{ } \\
 & & \texttt{ } \\
 & & \texttt{ } \\
 & & \texttt{1 } \\
 & & \texttt{1 } \\
 & & \texttt{1 } \\
 & & \texttt{1 } \\
 & & \texttt{1 } \\
 & & \texttt{1 } \\
 & & \texttt{1 } \\
 & & \texttt{1 } \\
 & & \texttt{ } \\
 & & \texttt{ } \\
 & & \texttt{1 } \\
 & & \texttt{1 } \\
 & & \texttt{1 } \\
 & & \texttt{1 } \\
 & & \texttt{ } \\
\hline \end{tabular}
}
\newpage
{\scriptsize
\begin{tabular}{|l|l|l|} \hline
\bf{Register Class} & \bf{Register File} & \bf{Registers} \\ \hline
 onlygetReg & SFR & \texttt{ } \\
 & & \texttt{ } \\
 & & \texttt{ } \\
 & & \texttt{ } \\
 & & \texttt{ } \\
 & & \texttt{ } \\
 & & \texttt{ } \\
 & & \texttt{ } \\
 & & \texttt{ } \\
 & & \texttt{ } \\
 & & \texttt{ } \\
 & & \texttt{ } \\
 & & \texttt{ } \\
 & & \texttt{ } \\
 & & \texttt{ } \\
 & & \texttt{ } \\
 & & \texttt{ } \\
 & & \texttt{ } \\
 & & \texttt{ } \\
 & & \texttt{ } \\
 & & \texttt{ } \\
 & & \texttt{ } \\
 & & \texttt{ } \\
 & & \texttt{ } \\
 & & \texttt{ } \\
 & & \texttt{ } \\
 & & \texttt{ } \\
 & & \texttt{ } \\
 & & \texttt{ } \\
 & & \texttt{ } \\
 & & \texttt{ } \\
 & & \texttt{ } \\
 & & \texttt{ } \\
 & & \texttt{ } \\
 & & \texttt{ } \\
 & & \texttt{ } \\
 & & \texttt{ } \\
 & & \texttt{ } \\
 & & \texttt{ } \\
 & & \texttt{ } \\
 & & \texttt{ } \\
 & & \texttt{ } \\
 & & \texttt{ } \\
 & & \texttt{ } \\
 & & \texttt{ } \\
 & & \texttt{ } \\
 & & \texttt{ } \\
 & & \texttt{ } \\
 & & \texttt{ } \\
 & & \texttt{ } \\
 & & \texttt{ } \\
 & & \texttt{ } \\
 & & \texttt{ } \\
 & & \texttt{ } \\
 & & \texttt{ } \\
 & & \texttt{ } \\
 & & \texttt{ } \\
 & & \texttt{ } \\
 & & \texttt{ } \\
 & & \texttt{ } \\
 & & \texttt{ } \\
 & & \texttt{ } \\
 & & \texttt{ } \\
 & & \texttt{ } \\
 & & \texttt{1 } \\
 & & \texttt{1 } \\
 & & \texttt{1 } \\
 & & \texttt{1 } \\
 & & \texttt{1 } \\
 & & \texttt{1 } \\
 & & \texttt{1 } \\
 & & \texttt{1 } \\
 & & \texttt{1 } \\
 & & \texttt{1 } \\
 & & \texttt{1 } \\
 & & \texttt{1 } \\
 & & \texttt{1 } \\
 & & \texttt{1 } \\
 & & \texttt{1 } \\
 & & \texttt{1 } \\
 & & \texttt{1 } \\
 & & \texttt{1 } \\
 & & \texttt{1 } \\
 & & \texttt{1 } \\
 & & \texttt{1 } \\
 & & \texttt{1 } \\
 & & \texttt{1 } \\
 & & \texttt{1 } \\
 & & \texttt{1 } \\
 & & \texttt{1 } \\
 & & \texttt{1 } \\
 & & \texttt{1 } \\
 & & \texttt{1 } \\
 & & \texttt{1 } \\
 & & \texttt{1 } \\
 & & \texttt{1 } \\
 & & \texttt{1 } \\
 & & \texttt{1 } \\
 & & \texttt{1 } \\
 & & \texttt{1 } \\
 & & \texttt{ } \\
 & & \texttt{ } \\
 & & \texttt{ } \\
 & & \texttt{ } \\
 & & \texttt{ } \\
 & & \texttt{ } \\
 & & \texttt{ } \\
 & & \texttt{ } \\
 & & \texttt{ } \\
 & & \texttt{ } \\
 & & \texttt{ } \\
 & & \texttt{ } \\
 & & \texttt{ } \\
 & & \texttt{ } \\
 & & \texttt{ } \\
 & & \texttt{ } \\
 & & \texttt{ } \\
 & & \texttt{ } \\
 & & \texttt{ } \\
 & & \texttt{ } \\
 & & \texttt{ } \\
 & & \texttt{ } \\
 & & \texttt{ } \\
 & & \texttt{ } \\
 & & \texttt{ } \\
 & & \texttt{ } \\
 & & \texttt{ } \\
 & & \texttt{ } \\
 & & \texttt{ } \\
 & & \texttt{ } \\
 & & \texttt{ } \\
 & & \texttt{ } \\
 & & \texttt{ } \\
 & & \texttt{ } \\
 & & \texttt{ } \\
 & & \texttt{ } \\
 & & \texttt{ } \\
 & & \texttt{ } \\
 & & \texttt{ } \\
 & & \texttt{ } \\
 & & \texttt{ } \\
 & & \texttt{ } \\
 & & \texttt{ } \\
 & & \texttt{ } \\
 & & \texttt{ } \\
 & & \texttt{ } \\
 & & \texttt{ } \\
 & & \texttt{ } \\
 & & \texttt{ } \\
 & & \texttt{ } \\
 & & \texttt{ } \\
 & & \texttt{ } \\
 & & \texttt{ } \\
 & & \texttt{ } \\
 & & \texttt{ } \\
 & & \texttt{ } \\
 & & \texttt{ } \\
 & & \texttt{ } \\
 & & \texttt{ } \\
 & & \texttt{ } \\
 & & \texttt{ } \\
 & & \texttt{ } \\
 & & \texttt{ } \\
 & & \texttt{ } \\
 & & \texttt{ } \\
 & & \texttt{ } \\
 & & \texttt{ } \\
 & & \texttt{ } \\
 & & \texttt{ } \\
 & & \texttt{ } \\
 & & \texttt{ } \\
 & & \texttt{ } \\
 & & \texttt{ } \\
 & & \texttt{ } \\
 & & \texttt{ } \\
 & & \texttt{ } \\
 & & \texttt{ } \\
 & & \texttt{ } \\
 & & \texttt{ } \\
 & & \texttt{ } \\
 & & \texttt{ } \\
 & & \texttt{ } \\
 & & \texttt{ } \\
 & & \texttt{ } \\
 & & \texttt{ } \\
 & & \texttt{ } \\
 & & \texttt{ } \\
 & & \texttt{ } \\
 & & \texttt{ } \\
 & & \texttt{ } \\
 & & \texttt{ } \\
 & & \texttt{ } \\
 & & \texttt{ } \\
 & & \texttt{ } \\
 & & \texttt{ } \\
 & & \texttt{ } \\
 & & \texttt{ } \\
 & & \texttt{ } \\
 & & \texttt{ } \\
 & & \texttt{ } \\
 & & \texttt{ } \\
 & & \texttt{ } \\
 & & \texttt{ } \\
 & & \texttt{ } \\
 & & \texttt{ } \\
 & & \texttt{ } \\
 & & \texttt{ } \\
 & & \texttt{ } \\
 & & \texttt{ } \\
 & & \texttt{ } \\
 & & \texttt{ } \\
 & & \texttt{ } \\
 & & \texttt{ } \\
 & & \texttt{ } \\
 & & \texttt{ } \\
 & & \texttt{ } \\
 & & \texttt{ } \\
 & & \texttt{ } \\
 & & \texttt{ } \\
 & & \texttt{ } \\
 & & \texttt{ } \\
 & & \texttt{ } \\
 & & \texttt{ } \\
 & & \texttt{ } \\
 & & \texttt{ } \\
 & & \texttt{ } \\
 & & \texttt{ } \\
 & & \texttt{ } \\
 & & \texttt{ } \\
 & & \texttt{ } \\
 & & \texttt{ } \\
 & & \texttt{ } \\
 & & \texttt{ } \\
 & & \texttt{ } \\
 & & \texttt{ } \\
 & & \texttt{ } \\
 & & \texttt{ } \\
 & & \texttt{ } \\
 & & \texttt{ } \\
 & & \texttt{ } \\
 & & \texttt{ } \\
 & & \texttt{ } \\
 & & \texttt{ } \\
 & & \texttt{ } \\
 & & \texttt{ } \\
 & & \texttt{ } \\
 & & \texttt{ } \\
 & & \texttt{ } \\
 & & \texttt{ } \\
 & & \texttt{ } \\
 & & \texttt{ } \\
 & & \texttt{ } \\
 & & \texttt{ } \\
 & & \texttt{ } \\
 & & \texttt{ } \\
 & & \texttt{ } \\
 & & \texttt{ } \\
 & & \texttt{ } \\
 & & \texttt{ } \\
 & & \texttt{ } \\
 & & \texttt{ } \\
 & & \texttt{ } \\
 & & \texttt{ } \\
 & & \texttt{ } \\
 & & \texttt{ } \\
 & & \texttt{ } \\
 & & \texttt{ } \\
 & & \texttt{ } \\
 & & \texttt{ } \\
 & & \texttt{ } \\
 & & \texttt{ } \\
 & & \texttt{ } \\
 & & \texttt{ } \\
 & & \texttt{ } \\
 & & \texttt{ } \\
 & & \texttt{ } \\
 & & \texttt{ } \\
 & & \texttt{ } \\
 & & \texttt{ } \\
 & & \texttt{ } \\
 & & \texttt{ } \\
 & & \texttt{ } \\
 & & \texttt{ } \\
 & & \texttt{ } \\
 & & \texttt{ } \\
 & & \texttt{ } \\
 & & \texttt{ } \\
 & & \texttt{ } \\
 & & \texttt{ } \\
 & & \texttt{ } \\
 & & \texttt{ } \\
 & & \texttt{ } \\
 & & \texttt{ } \\
 & & \texttt{ } \\
 & & \texttt{ } \\
 & & \texttt{ } \\
 & & \texttt{ } \\
 & & \texttt{ } \\
 & & \texttt{ } \\
 & & \texttt{ } \\
 & & \texttt{ } \\
 & & \texttt{ } \\
 & & \texttt{ } \\
 & & \texttt{ } \\
 & & \texttt{ } \\
 & & \texttt{ } \\
 & & \texttt{ } \\
 & & \texttt{ } \\
 & & \texttt{ } \\
 & & \texttt{ } \\
 & & \texttt{ } \\
 & & \texttt{ } \\
 & & \texttt{ } \\
 & & \texttt{ } \\
 & & \texttt{ } \\
 & & \texttt{ } \\
 & & \texttt{ } \\
 & & \texttt{ } \\
 & & \texttt{ } \\
 & & \texttt{ } \\
 & & \texttt{ } \\
 & & \texttt{ } \\
 & & \texttt{ } \\
 & & \texttt{ } \\
 & & \texttt{ } \\
 & & \texttt{ } \\
 & & \texttt{ } \\
 & & \texttt{ } \\
 & & \texttt{ } \\
 & & \texttt{ } \\
 & & \texttt{ } \\
 & & \texttt{ } \\
 & & \texttt{ } \\
 & & \texttt{ } \\
 & & \texttt{ } \\
 & & \texttt{ } \\
 & & \texttt{ } \\
 & & \texttt{ } \\
 & & \texttt{ } \\
 & & \texttt{ } \\
 & & \texttt{ } \\
 & & \texttt{ } \\
 & & \texttt{ } \\
 & & \texttt{ } \\
 & & \texttt{ } \\
 & & \texttt{ } \\
 & & \texttt{ } \\
 & & \texttt{ } \\
 & & \texttt{ } \\
 & & \texttt{ } \\
 & & \texttt{ } \\
 & & \texttt{ } \\
 & & \texttt{ } \\
 & & \texttt{ } \\
 & & \texttt{ } \\
 & & \texttt{ } \\
 & & \texttt{ } \\
 & & \texttt{ } \\
 & & \texttt{ } \\
 & & \texttt{ } \\
 & & \texttt{ } \\
 & & \texttt{ } \\
 & & \texttt{ } \\
 & & \texttt{ } \\
 & & \texttt{ } \\
\hline \end{tabular}
}
\newpage
{\scriptsize
\begin{tabular}{|l|l|l|} \hline
\bf{Register Class} & \bf{Register File} & \bf{Registers} \\ \hline
 onlyraReg & SFR & \texttt{ } \\
\hline \end{tabular}
}
\newpage
{\scriptsize
\begin{tabular}{|l|l|l|} \hline
\bf{Register Class} & \bf{Register File} & \bf{Registers} \\ \hline
 onlysetReg & SFR & \texttt{ } \\
 & & \texttt{ } \\
 & & \texttt{ } \\
 & & \texttt{ } \\
 & & \texttt{ } \\
 & & \texttt{ } \\
 & & \texttt{ } \\
 & & \texttt{ } \\
 & & \texttt{ } \\
 & & \texttt{ } \\
 & & \texttt{ } \\
 & & \texttt{ } \\
 & & \texttt{ } \\
 & & \texttt{ } \\
 & & \texttt{ } \\
 & & \texttt{ } \\
 & & \texttt{ } \\
 & & \texttt{ } \\
 & & \texttt{ } \\
 & & \texttt{ } \\
 & & \texttt{ } \\
 & & \texttt{ } \\
 & & \texttt{ } \\
 & & \texttt{ } \\
 & & \texttt{ } \\
 & & \texttt{ } \\
 & & \texttt{ } \\
 & & \texttt{ } \\
 & & \texttt{ } \\
 & & \texttt{ } \\
 & & \texttt{ } \\
 & & \texttt{ } \\
 & & \texttt{ } \\
 & & \texttt{ } \\
 & & \texttt{ } \\
 & & \texttt{ } \\
 & & \texttt{ } \\
 & & \texttt{ } \\
 & & \texttt{ } \\
 & & \texttt{ } \\
 & & \texttt{ } \\
 & & \texttt{ } \\
 & & \texttt{ } \\
 & & \texttt{ } \\
 & & \texttt{ } \\
 & & \texttt{ } \\
 & & \texttt{ } \\
 & & \texttt{ } \\
 & & \texttt{ } \\
 & & \texttt{ } \\
 & & \texttt{ } \\
 & & \texttt{ } \\
 & & \texttt{ } \\
 & & \texttt{ } \\
 & & \texttt{ } \\
 & & \texttt{ } \\
 & & \texttt{ } \\
 & & \texttt{1 } \\
 & & \texttt{1 } \\
 & & \texttt{1 } \\
 & & \texttt{1 } \\
 & & \texttt{1 } \\
 & & \texttt{1 } \\
 & & \texttt{1 } \\
 & & \texttt{1 } \\
 & & \texttt{1 } \\
 & & \texttt{1 } \\
 & & \texttt{1 } \\
 & & \texttt{1 } \\
 & & \texttt{1 } \\
 & & \texttt{1 } \\
 & & \texttt{1 } \\
 & & \texttt{1 } \\
 & & \texttt{1 } \\
 & & \texttt{1 } \\
 & & \texttt{1 } \\
 & & \texttt{1 } \\
 & & \texttt{1 } \\
 & & \texttt{1 } \\
 & & \texttt{1 } \\
 & & \texttt{1 } \\
 & & \texttt{1 } \\
 & & \texttt{1 } \\
 & & \texttt{1 } \\
 & & \texttt{1 } \\
 & & \texttt{1 } \\
 & & \texttt{1 } \\
 & & \texttt{1 } \\
 & & \texttt{1 } \\
 & & \texttt{1 } \\
 & & \texttt{1 } \\
 & & \texttt{1 } \\
 & & \texttt{1 } \\
 & & \texttt{ } \\
 & & \texttt{ } \\
 & & \texttt{ } \\
 & & \texttt{ } \\
 & & \texttt{ } \\
 & & \texttt{ } \\
 & & \texttt{ } \\
 & & \texttt{ } \\
 & & \texttt{ } \\
 & & \texttt{ } \\
 & & \texttt{ } \\
 & & \texttt{ } \\
 & & \texttt{ } \\
 & & \texttt{ } \\
 & & \texttt{ } \\
 & & \texttt{ } \\
 & & \texttt{ } \\
 & & \texttt{ } \\
 & & \texttt{ } \\
 & & \texttt{ } \\
 & & \texttt{ } \\
 & & \texttt{ } \\
 & & \texttt{ } \\
 & & \texttt{ } \\
 & & \texttt{ } \\
 & & \texttt{ } \\
 & & \texttt{ } \\
 & & \texttt{ } \\
 & & \texttt{ } \\
 & & \texttt{ } \\
 & & \texttt{ } \\
 & & \texttt{ } \\
 & & \texttt{ } \\
 & & \texttt{ } \\
 & & \texttt{ } \\
 & & \texttt{ } \\
 & & \texttt{ } \\
 & & \texttt{ } \\
 & & \texttt{ } \\
 & & \texttt{ } \\
 & & \texttt{ } \\
 & & \texttt{ } \\
 & & \texttt{ } \\
 & & \texttt{ } \\
 & & \texttt{ } \\
 & & \texttt{ } \\
 & & \texttt{ } \\
 & & \texttt{ } \\
 & & \texttt{ } \\
 & & \texttt{ } \\
 & & \texttt{ } \\
 & & \texttt{ } \\
 & & \texttt{ } \\
 & & \texttt{ } \\
 & & \texttt{ } \\
 & & \texttt{ } \\
 & & \texttt{ } \\
 & & \texttt{ } \\
 & & \texttt{ } \\
 & & \texttt{ } \\
 & & \texttt{ } \\
 & & \texttt{ } \\
 & & \texttt{ } \\
 & & \texttt{ } \\
 & & \texttt{ } \\
 & & \texttt{ } \\
 & & \texttt{ } \\
 & & \texttt{ } \\
 & & \texttt{ } \\
 & & \texttt{ } \\
 & & \texttt{ } \\
 & & \texttt{ } \\
 & & \texttt{ } \\
 & & \texttt{ } \\
 & & \texttt{ } \\
 & & \texttt{ } \\
 & & \texttt{ } \\
 & & \texttt{ } \\
 & & \texttt{ } \\
 & & \texttt{ } \\
 & & \texttt{ } \\
 & & \texttt{ } \\
 & & \texttt{ } \\
 & & \texttt{ } \\
 & & \texttt{ } \\
 & & \texttt{ } \\
 & & \texttt{ } \\
 & & \texttt{ } \\
 & & \texttt{ } \\
 & & \texttt{ } \\
 & & \texttt{ } \\
 & & \texttt{ } \\
 & & \texttt{ } \\
 & & \texttt{ } \\
 & & \texttt{ } \\
 & & \texttt{ } \\
 & & \texttt{ } \\
 & & \texttt{ } \\
 & & \texttt{ } \\
 & & \texttt{ } \\
 & & \texttt{ } \\
 & & \texttt{ } \\
 & & \texttt{ } \\
 & & \texttt{ } \\
 & & \texttt{ } \\
 & & \texttt{ } \\
 & & \texttt{ } \\
 & & \texttt{ } \\
 & & \texttt{ } \\
 & & \texttt{ } \\
 & & \texttt{ } \\
 & & \texttt{ } \\
 & & \texttt{ } \\
 & & \texttt{ } \\
 & & \texttt{ } \\
 & & \texttt{ } \\
 & & \texttt{ } \\
 & & \texttt{ } \\
 & & \texttt{ } \\
 & & \texttt{ } \\
 & & \texttt{ } \\
 & & \texttt{ } \\
 & & \texttt{ } \\
 & & \texttt{ } \\
 & & \texttt{ } \\
 & & \texttt{ } \\
 & & \texttt{ } \\
 & & \texttt{ } \\
 & & \texttt{ } \\
 & & \texttt{ } \\
 & & \texttt{ } \\
 & & \texttt{ } \\
 & & \texttt{ } \\
 & & \texttt{ } \\
 & & \texttt{ } \\
 & & \texttt{ } \\
 & & \texttt{ } \\
 & & \texttt{ } \\
 & & \texttt{ } \\
 & & \texttt{ } \\
 & & \texttt{ } \\
 & & \texttt{ } \\
 & & \texttt{ } \\
 & & \texttt{ } \\
 & & \texttt{ } \\
 & & \texttt{ } \\
 & & \texttt{ } \\
 & & \texttt{ } \\
 & & \texttt{ } \\
 & & \texttt{ } \\
 & & \texttt{ } \\
 & & \texttt{ } \\
 & & \texttt{ } \\
 & & \texttt{ } \\
 & & \texttt{ } \\
 & & \texttt{ } \\
 & & \texttt{ } \\
 & & \texttt{ } \\
 & & \texttt{ } \\
 & & \texttt{ } \\
 & & \texttt{ } \\
 & & \texttt{ } \\
 & & \texttt{ } \\
 & & \texttt{ } \\
 & & \texttt{ } \\
 & & \texttt{ } \\
 & & \texttt{ } \\
 & & \texttt{ } \\
 & & \texttt{ } \\
 & & \texttt{ } \\
 & & \texttt{ } \\
 & & \texttt{ } \\
 & & \texttt{ } \\
 & & \texttt{ } \\
 & & \texttt{ } \\
 & & \texttt{ } \\
 & & \texttt{ } \\
 & & \texttt{ } \\
 & & \texttt{ } \\
 & & \texttt{ } \\
 & & \texttt{ } \\
 & & \texttt{ } \\
 & & \texttt{ } \\
 & & \texttt{ } \\
 & & \texttt{ } \\
 & & \texttt{ } \\
 & & \texttt{ } \\
 & & \texttt{ } \\
 & & \texttt{ } \\
 & & \texttt{ } \\
 & & \texttt{ } \\
 & & \texttt{ } \\
 & & \texttt{ } \\
 & & \texttt{ } \\
 & & \texttt{ } \\
 & & \texttt{ } \\
 & & \texttt{ } \\
 & & \texttt{ } \\
 & & \texttt{ } \\
 & & \texttt{ } \\
 & & \texttt{ } \\
 & & \texttt{ } \\
 & & \texttt{ } \\
 & & \texttt{ } \\
 & & \texttt{ } \\
 & & \texttt{ } \\
 & & \texttt{ } \\
 & & \texttt{ } \\
 & & \texttt{ } \\
 & & \texttt{ } \\
 & & \texttt{ } \\
 & & \texttt{ } \\
 & & \texttt{ } \\
 & & \texttt{ } \\
 & & \texttt{ } \\
 & & \texttt{ } \\
 & & \texttt{ } \\
 & & \texttt{ } \\
 & & \texttt{ } \\
 & & \texttt{ } \\
 & & \texttt{ } \\
 & & \texttt{ } \\
 & & \texttt{ } \\
 & & \texttt{ } \\
 & & \texttt{ } \\
 & & \texttt{ } \\
 & & \texttt{ } \\
 & & \texttt{ } \\
 & & \texttt{ } \\
 & & \texttt{ } \\
 & & \texttt{ } \\
 & & \texttt{ } \\
 & & \texttt{ } \\
 & & \texttt{ } \\
 & & \texttt{ } \\
 & & \texttt{ } \\
 & & \texttt{ } \\
 & & \texttt{ } \\
 & & \texttt{ } \\
 & & \texttt{ } \\
 & & \texttt{ } \\
 & & \texttt{ } \\
 & & \texttt{ } \\
 & & \texttt{ } \\
 & & \texttt{ } \\
 & & \texttt{ } \\
 & & \texttt{ } \\
 & & \texttt{ } \\
 & & \texttt{ } \\
 & & \texttt{ } \\
 & & \texttt{ } \\
 & & \texttt{ } \\
 & & \texttt{ } \\
 & & \texttt{ } \\
 & & \texttt{ } \\
 & & \texttt{ } \\
 & & \texttt{ } \\
 & & \texttt{ } \\
 & & \texttt{ } \\
 & & \texttt{ } \\
 & & \texttt{ } \\
 & & \texttt{ } \\
 & & \texttt{ } \\
 & & \texttt{ } \\
 & & \texttt{ } \\
\hline \end{tabular}
}
\newpage
{\scriptsize
\begin{tabular}{|l|l|l|} \hline
\bf{Register Class} & \bf{Register File} & \bf{Registers} \\ \hline
 onlyswapReg & SFR & \texttt{ } \\
 & & \texttt{ } \\
 & & \texttt{ } \\
 & & \texttt{ } \\
 & & \texttt{ } \\
 & & \texttt{ } \\
 & & \texttt{ } \\
 & & \texttt{ } \\
 & & \texttt{ } \\
 & & \texttt{ } \\
 & & \texttt{ } \\
 & & \texttt{ } \\
 & & \texttt{ } \\
 & & \texttt{ } \\
 & & \texttt{ } \\
 & & \texttt{ } \\
 & & \texttt{ } \\
 & & \texttt{ } \\
 & & \texttt{ } \\
 & & \texttt{ } \\
 & & \texttt{ } \\
 & & \texttt{ } \\
 & & \texttt{ } \\
 & & \texttt{ } \\
 & & \texttt{ } \\
 & & \texttt{ } \\
 & & \texttt{ } \\
 & & \texttt{ } \\
 & & \texttt{ } \\
 & & \texttt{ } \\
 & & \texttt{ } \\
 & & \texttt{ } \\
 & & \texttt{ } \\
 & & \texttt{ } \\
 & & \texttt{ } \\
 & & \texttt{ } \\
 & & \texttt{ } \\
 & & \texttt{ } \\
 & & \texttt{ } \\
 & & \texttt{ } \\
 & & \texttt{ } \\
 & & \texttt{ } \\
 & & \texttt{ } \\
 & & \texttt{ } \\
 & & \texttt{ } \\
 & & \texttt{ } \\
 & & \texttt{ } \\
 & & \texttt{ } \\
 & & \texttt{ } \\
 & & \texttt{ } \\
 & & \texttt{ } \\
 & & \texttt{ } \\
 & & \texttt{ } \\
 & & \texttt{ } \\
 & & \texttt{ } \\
 & & \texttt{ } \\
 & & \texttt{ } \\
 & & \texttt{ } \\
 & & \texttt{ } \\
 & & \texttt{ } \\
 & & \texttt{ } \\
 & & \texttt{ } \\
 & & \texttt{ } \\
 & & \texttt{1 } \\
 & & \texttt{1 } \\
 & & \texttt{1 } \\
 & & \texttt{1 } \\
 & & \texttt{1 } \\
 & & \texttt{1 } \\
 & & \texttt{1 } \\
 & & \texttt{1 } \\
 & & \texttt{1 } \\
 & & \texttt{1 } \\
 & & \texttt{1 } \\
 & & \texttt{1 } \\
 & & \texttt{1 } \\
 & & \texttt{1 } \\
 & & \texttt{1 } \\
 & & \texttt{1 } \\
 & & \texttt{1 } \\
 & & \texttt{1 } \\
 & & \texttt{1 } \\
 & & \texttt{1 } \\
 & & \texttt{1 } \\
 & & \texttt{1 } \\
 & & \texttt{1 } \\
 & & \texttt{1 } \\
 & & \texttt{1 } \\
 & & \texttt{1 } \\
 & & \texttt{1 } \\
 & & \texttt{1 } \\
 & & \texttt{1 } \\
 & & \texttt{1 } \\
 & & \texttt{1 } \\
 & & \texttt{1 } \\
 & & \texttt{ } \\
 & & \texttt{ } \\
 & & \texttt{ } \\
 & & \texttt{ } \\
 & & \texttt{ } \\
 & & \texttt{ } \\
 & & \texttt{ } \\
 & & \texttt{ } \\
 & & \texttt{1 } \\
 & & \texttt{1 } \\
 & & \texttt{1 } \\
 & & \texttt{1 } \\
 & & \texttt{ } \\
 & & \texttt{ } \\
 & & \texttt{ } \\
 & & \texttt{ } \\
 & & \texttt{ } \\
 & & \texttt{ } \\
 & & \texttt{ } \\
 & & \texttt{ } \\
 & & \texttt{ } \\
 & & \texttt{ } \\
 & & \texttt{ } \\
 & & \texttt{ } \\
 & & \texttt{ } \\
 & & \texttt{ } \\
 & & \texttt{ } \\
 & & \texttt{ } \\
 & & \texttt{ } \\
 & & \texttt{ } \\
 & & \texttt{ } \\
 & & \texttt{ } \\
 & & \texttt{ } \\
 & & \texttt{ } \\
 & & \texttt{ } \\
 & & \texttt{ } \\
 & & \texttt{ } \\
 & & \texttt{ } \\
 & & \texttt{ } \\
 & & \texttt{ } \\
 & & \texttt{ } \\
 & & \texttt{ } \\
 & & \texttt{ } \\
 & & \texttt{ } \\
 & & \texttt{ } \\
 & & \texttt{ } \\
 & & \texttt{ } \\
 & & \texttt{ } \\
 & & \texttt{ } \\
 & & \texttt{ } \\
 & & \texttt{ } \\
 & & \texttt{ } \\
 & & \texttt{ } \\
 & & \texttt{ } \\
 & & \texttt{ } \\
 & & \texttt{ } \\
 & & \texttt{ } \\
 & & \texttt{ } \\
 & & \texttt{ } \\
 & & \texttt{ } \\
 & & \texttt{ } \\
 & & \texttt{ } \\
 & & \texttt{ } \\
 & & \texttt{ } \\
 & & \texttt{ } \\
 & & \texttt{ } \\
 & & \texttt{ } \\
 & & \texttt{ } \\
 & & \texttt{ } \\
 & & \texttt{ } \\
 & & \texttt{ } \\
 & & \texttt{ } \\
 & & \texttt{ } \\
 & & \texttt{ } \\
 & & \texttt{ } \\
 & & \texttt{ } \\
 & & \texttt{ } \\
 & & \texttt{ } \\
 & & \texttt{ } \\
 & & \texttt{ } \\
 & & \texttt{ } \\
 & & \texttt{ } \\
 & & \texttt{ } \\
 & & \texttt{ } \\
 & & \texttt{ } \\
 & & \texttt{ } \\
 & & \texttt{ } \\
 & & \texttt{ } \\
 & & \texttt{ } \\
 & & \texttt{ } \\
 & & \texttt{ } \\
 & & \texttt{ } \\
 & & \texttt{ } \\
 & & \texttt{ } \\
 & & \texttt{ } \\
 & & \texttt{ } \\
 & & \texttt{ } \\
 & & \texttt{ } \\
 & & \texttt{ } \\
 & & \texttt{ } \\
 & & \texttt{ } \\
 & & \texttt{ } \\
 & & \texttt{ } \\
 & & \texttt{ } \\
 & & \texttt{ } \\
 & & \texttt{ } \\
 & & \texttt{ } \\
 & & \texttt{ } \\
 & & \texttt{ } \\
 & & \texttt{ } \\
 & & \texttt{ } \\
 & & \texttt{ } \\
 & & \texttt{ } \\
 & & \texttt{ } \\
 & & \texttt{ } \\
 & & \texttt{ } \\
 & & \texttt{ } \\
 & & \texttt{ } \\
 & & \texttt{ } \\
 & & \texttt{ } \\
 & & \texttt{ } \\
 & & \texttt{ } \\
 & & \texttt{ } \\
 & & \texttt{ } \\
 & & \texttt{ } \\
 & & \texttt{ } \\
 & & \texttt{ } \\
 & & \texttt{ } \\
 & & \texttt{ } \\
 & & \texttt{ } \\
 & & \texttt{ } \\
 & & \texttt{ } \\
 & & \texttt{ } \\
 & & \texttt{ } \\
 & & \texttt{ } \\
 & & \texttt{ } \\
 & & \texttt{ } \\
 & & \texttt{ } \\
 & & \texttt{ } \\
 & & \texttt{ } \\
 & & \texttt{ } \\
 & & \texttt{ } \\
 & & \texttt{ } \\
 & & \texttt{ } \\
 & & \texttt{ } \\
 & & \texttt{ } \\
 & & \texttt{ } \\
 & & \texttt{ } \\
 & & \texttt{ } \\
 & & \texttt{ } \\
 & & \texttt{ } \\
 & & \texttt{ } \\
 & & \texttt{ } \\
 & & \texttt{ } \\
 & & \texttt{ } \\
 & & \texttt{ } \\
 & & \texttt{ } \\
 & & \texttt{ } \\
 & & \texttt{ } \\
 & & \texttt{ } \\
 & & \texttt{ } \\
 & & \texttt{ } \\
 & & \texttt{ } \\
 & & \texttt{ } \\
 & & \texttt{ } \\
 & & \texttt{ } \\
 & & \texttt{ } \\
 & & \texttt{ } \\
 & & \texttt{ } \\
 & & \texttt{ } \\
 & & \texttt{ } \\
 & & \texttt{ } \\
 & & \texttt{ } \\
 & & \texttt{ } \\
 & & \texttt{ } \\
 & & \texttt{ } \\
 & & \texttt{ } \\
 & & \texttt{ } \\
 & & \texttt{ } \\
 & & \texttt{ } \\
 & & \texttt{ } \\
 & & \texttt{ } \\
 & & \texttt{ } \\
 & & \texttt{ } \\
 & & \texttt{ } \\
 & & \texttt{ } \\
 & & \texttt{ } \\
 & & \texttt{ } \\
 & & \texttt{ } \\
 & & \texttt{ } \\
 & & \texttt{ } \\
 & & \texttt{ } \\
 & & \texttt{ } \\
 & & \texttt{ } \\
 & & \texttt{ } \\
 & & \texttt{ } \\
 & & \texttt{ } \\
 & & \texttt{ } \\
 & & \texttt{ } \\
 & & \texttt{ } \\
 & & \texttt{ } \\
 & & \texttt{ } \\
 & & \texttt{ } \\
 & & \texttt{ } \\
 & & \texttt{ } \\
 & & \texttt{ } \\
 & & \texttt{ } \\
 & & \texttt{ } \\
 & & \texttt{ } \\
 & & \texttt{ } \\
 & & \texttt{ } \\
 & & \texttt{ } \\
 & & \texttt{ } \\
 & & \texttt{ } \\
 & & \texttt{ } \\
 & & \texttt{ } \\
 & & \texttt{ } \\
 & & \texttt{ } \\
 & & \texttt{ } \\
 & & \texttt{ } \\
 & & \texttt{ } \\
 & & \texttt{ } \\
 & & \texttt{ } \\
 & & \texttt{ } \\
 & & \texttt{ } \\
 & & \texttt{ } \\
 & & \texttt{ } \\
 & & \texttt{ } \\
 & & \texttt{ } \\
 & & \texttt{ } \\
 & & \texttt{ } \\
 & & \texttt{ } \\
 & & \texttt{ } \\
 & & \texttt{ } \\
 & & \texttt{ } \\
 & & \texttt{ } \\
 & & \texttt{ } \\
 & & \texttt{ } \\
 & & \texttt{ } \\
 & & \texttt{ } \\
 & & \texttt{ } \\
 & & \texttt{ } \\
 & & \texttt{ } \\
 & & \texttt{ } \\
 & & \texttt{ } \\
 & & \texttt{ } \\
 & & \texttt{ } \\
 & & \texttt{ } \\
 & & \texttt{ } \\
 & & \texttt{ } \\
 & & \texttt{ } \\
 & & \texttt{ } \\
 & & \texttt{ } \\
 & & \texttt{ } \\
 & & \texttt{ } \\
 & & \texttt{ } \\
 & & \texttt{ } \\
 & & \texttt{ } \\
 & & \texttt{ } \\
 & & \texttt{ } \\
 & & \texttt{ } \\
 & & \texttt{ } \\
 & & \texttt{ } \\
 & & \texttt{ } \\
 & & \texttt{ } \\
 & & \texttt{ } \\
 & & \texttt{ } \\
 & & \texttt{ } \\
 & & \texttt{ } \\
 & & \texttt{ } \\
 & & \texttt{ } \\
 & & \texttt{ } \\
 & & \texttt{ } \\
 & & \texttt{ } \\
 & & \texttt{ } \\
 & & \texttt{ } \\
 & & \texttt{ } \\
 & & \texttt{ } \\
 & & \texttt{ } \\
 & & \texttt{ } \\
 & & \texttt{ } \\
 & & \texttt{ } \\
 & & \texttt{ } \\
 & & \texttt{ } \\
 & & \texttt{ } \\
 & & \texttt{ } \\
 & & \texttt{ } \\
 & & \texttt{ } \\
 & & \texttt{ } \\
 & & \texttt{ } \\
 & & \texttt{ } \\
 & & \texttt{ } \\
 & & \texttt{ } \\
 & & \texttt{ } \\
 & & \texttt{ } \\
 & & \texttt{ } \\
 & & \texttt{ } \\
 & & \texttt{ } \\
 & & \texttt{ } \\
 & & \texttt{ } \\
 & & \texttt{ } \\
 & & \texttt{ } \\
 & & \texttt{ } \\
 & & \texttt{ } \\
 & & \texttt{ } \\
 & & \texttt{ } \\
 & & \texttt{ } \\
 & & \texttt{ } \\
 & & \texttt{ } \\
 & & \texttt{ } \\
 & & \texttt{ } \\
 & & \texttt{ } \\
 & & \texttt{ } \\
 & & \texttt{ } \\
 & & \texttt{ } \\
 & & \texttt{ } \\
 & & \texttt{ } \\
 & & \texttt{ } \\
 & & \texttt{ } \\
 & & \texttt{ } \\
 & & \texttt{ } \\
 & & \texttt{ } \\
 & & \texttt{ } \\
 & & \texttt{ } \\
 & & \texttt{ } \\
 & & \texttt{ } \\
 & & \texttt{ } \\
 & & \texttt{ } \\
 & & \texttt{ } \\
 & & \texttt{ } \\
 & & \texttt{ } \\
 & & \texttt{ } \\
 & & \texttt{ } \\
 & & \texttt{ } \\
 & & \texttt{ } \\
 & & \texttt{ } \\
 & & \texttt{ } \\
 & & \texttt{ } \\
 & & \texttt{ } \\
 & & \texttt{ } \\
 & & \texttt{ } \\
 & & \texttt{ } \\
 & & \texttt{ } \\
 & & \texttt{ } \\
 & & \texttt{ } \\
 & & \texttt{ } \\
 & & \texttt{ } \\
 & & \texttt{ } \\
 & & \texttt{ } \\
 & & \texttt{ } \\
 & & \texttt{ } \\
 & & \texttt{ } \\
 & & \texttt{ } \\
 & & \texttt{ } \\
 & & \texttt{ } \\
 & & \texttt{ } \\
 & & \texttt{ } \\
 & & \texttt{ } \\
 & & \texttt{ } \\
 & & \texttt{ } \\
 & & \texttt{ } \\
 & & \texttt{ } \\
 & & \texttt{ } \\
 & & \texttt{ } \\
 & & \texttt{ } \\
 & & \texttt{ } \\
 & & \texttt{ } \\
 & & \texttt{ } \\
 & & \texttt{ } \\
 & & \texttt{ } \\
 & & \texttt{ } \\
 & & \texttt{ } \\
 & & \texttt{ } \\
 & & \texttt{ } \\
 & & \texttt{ } \\
 & & \texttt{ } \\
 & & \texttt{ } \\
 & & \texttt{ } \\
 & & \texttt{ } \\
 & & \texttt{ } \\
 & & \texttt{ } \\
 & & \texttt{ } \\
 & & \texttt{ } \\
 & & \texttt{ } \\
 & & \texttt{ } \\
 & & \texttt{ } \\
 & & \texttt{ } \\
 & & \texttt{ } \\
 & & \texttt{ } \\
 & & \texttt{ } \\
 & & \texttt{ } \\
 & & \texttt{ } \\
 & & \texttt{ } \\
 & & \texttt{ } \\
 & & \texttt{ } \\
 & & \texttt{ } \\
 & & \texttt{ } \\
 & & \texttt{ } \\
 & & \texttt{ } \\
 & & \texttt{ } \\
 & & \texttt{ } \\
 & & \texttt{ } \\
 & & \texttt{ } \\
 & & \texttt{ } \\
 & & \texttt{ } \\
 & & \texttt{ } \\
 & & \texttt{ } \\
 & & \texttt{ } \\
 & & \texttt{ } \\
 & & \texttt{ } \\
 & & \texttt{ } \\
 & & \texttt{ } \\
 & & \texttt{ } \\
 & & \texttt{ } \\
 & & \texttt{ } \\
 & & \texttt{ } \\
\hline \end{tabular}
}
\newpage
{\scriptsize
\begin{tabular}{|l|l|l|} \hline
\bf{Register Class} & \bf{Register File} & \bf{Registers} \\ \hline
 pairedReg & PGR & \texttt{ } \\
 & & \texttt{ } \\
 & & \texttt{ } \\
 & & \texttt{ } \\
 & & \texttt{ } \\
 & & \texttt{ } \\
 & & \texttt{ } \\
 & & \texttt{ } \\
 & & \texttt{ } \\
 & & \texttt{ } \\
 & & \texttt{ } \\
 & & \texttt{ } \\
 & & \texttt{ } \\
 & & \texttt{ } \\
 & & \texttt{ } \\
 & & \texttt{ } \\
 & & \texttt{ } \\
 & & \texttt{ } \\
 & & \texttt{ } \\
 & & \texttt{ } \\
 & & \texttt{ } \\
 & & \texttt{ } \\
 & & \texttt{ } \\
 & & \texttt{ } \\
 & & \texttt{ } \\
 & & \texttt{ } \\
 & & \texttt{ } \\
 & & \texttt{ } \\
 & & \texttt{ } \\
 & & \texttt{ } \\
 & & \texttt{ } \\
 & & \texttt{ } \\
\hline \end{tabular}
}
\newpage
{\scriptsize
\begin{tabular}{|l|l|l|} \hline
\bf{Register Class} & \bf{Register File} & \bf{Registers} \\ \hline
 quadReg & QGR & \texttt{} \\
 & & \texttt{} \\
 & & \texttt{} \\
 & & \texttt{} \\
 & & \texttt{} \\
 & & \texttt{} \\
 & & \texttt{} \\
 & & \texttt{} \\
 & & \texttt{} \\
 & & \texttt{} \\
 & & \texttt{} \\
 & & \texttt{} \\
 & & \texttt{} \\
 & & \texttt{} \\
 & & \texttt{} \\
 & & \texttt{} \\
\hline \end{tabular}
}
\newpage
{\scriptsize
\begin{tabular}{|l|l|l|} \hline
\bf{Register Class} & \bf{Register File} & \bf{Registers} \\ \hline
 singleReg & GPR & \texttt{  } \\
 & & \texttt{  } \\
 & & \texttt{  } \\
 & & \texttt{  } \\
 & & \texttt{  } \\
 & & \texttt{  } \\
 & & \texttt{  } \\
 & & \texttt{  } \\
 & & \texttt{  } \\
 & & \texttt{  } \\
 & & \texttt{  } \\
 & & \texttt{  } \\
 & & \texttt{ } \\
 & & \texttt{ } \\
 & & \texttt{  } \\
 & & \texttt{  } \\
 & & \texttt{  } \\
 & & \texttt{  } \\
 & & \texttt{  } \\
 & & \texttt{  } \\
 & & \texttt{  } \\
 & & \texttt{  } \\
 & & \texttt{  } \\
 & & \texttt{  } \\
 & & \texttt{  } \\
 & & \texttt{  } \\
 & & \texttt{  } \\
 & & \texttt{  } \\
 & & \texttt{  } \\
 & & \texttt{  } \\
 & & \texttt{  } \\
 & & \texttt{  } \\
 & & \texttt{  } \\
 & & \texttt{  } \\
 & & \texttt{  } \\
 & & \texttt{  } \\
 & & \texttt{  } \\
 & & \texttt{  } \\
 & & \texttt{  } \\
 & & \texttt{  } \\
 & & \texttt{  } \\
 & & \texttt{  } \\
 & & \texttt{  } \\
 & & \texttt{  } \\
 & & \texttt{  } \\
 & & \texttt{  } \\
 & & \texttt{  } \\
 & & \texttt{  } \\
 & & \texttt{  } \\
 & & \texttt{  } \\
 & & \texttt{  } \\
 & & \texttt{  } \\
 & & \texttt{  } \\
 & & \texttt{  } \\
 & & \texttt{  } \\
 & & \texttt{  } \\
 & & \texttt{  } \\
 & & \texttt{  } \\
 & & \texttt{  } \\
 & & \texttt{  } \\
 & & \texttt{  } \\
 & & \texttt{  } \\
 & & \texttt{  } \\
 & & \texttt{  } \\
\hline \end{tabular}
}
\newpage
{\scriptsize
\begin{tabular}{|l|l|l|} \hline
\bf{Register Class} & \bf{Register File} & \bf{Registers} \\ \hline
 worddRegE & GPR & \texttt{  } \\
 & & \texttt{  } \\
 & & \texttt{  } \\
 & & \texttt{  } \\
 & & \texttt{  } \\
 & & \texttt{  } \\
 & & \texttt{ } \\
 & & \texttt{  } \\
 & & \texttt{  } \\
 & & \texttt{  } \\
 & & \texttt{  } \\
 & & \texttt{  } \\
 & & \texttt{  } \\
 & & \texttt{  } \\
 & & \texttt{  } \\
 & & \texttt{  } \\
 & & \texttt{  } \\
 & & \texttt{  } \\
 & & \texttt{  } \\
 & & \texttt{  } \\
 & & \texttt{  } \\
 & & \texttt{  } \\
 & & \texttt{  } \\
 & & \texttt{  } \\
 & & \texttt{  } \\
 & & \texttt{  } \\
 & & \texttt{  } \\
 & & \texttt{  } \\
 & & \texttt{  } \\
 & & \texttt{  } \\
 & & \texttt{  } \\
 & & \texttt{  } \\
\hline \end{tabular}
}
\newpage
{\scriptsize
\begin{tabular}{|l|l|l|} \hline
\bf{Register Class} & \bf{Register File} & \bf{Registers} \\ \hline
 worddRegO & GPR & \texttt{  } \\
 & & \texttt{  } \\
 & & \texttt{  } \\
 & & \texttt{  } \\
 & & \texttt{  } \\
 & & \texttt{  } \\
 & & \texttt{ } \\
 & & \texttt{  } \\
 & & \texttt{  } \\
 & & \texttt{  } \\
 & & \texttt{  } \\
 & & \texttt{  } \\
 & & \texttt{  } \\
 & & \texttt{  } \\
 & & \texttt{  } \\
 & & \texttt{  } \\
 & & \texttt{  } \\
 & & \texttt{  } \\
 & & \texttt{  } \\
 & & \texttt{  } \\
 & & \texttt{  } \\
 & & \texttt{  } \\
 & & \texttt{  } \\
 & & \texttt{  } \\
 & & \texttt{  } \\
 & & \texttt{  } \\
 & & \texttt{  } \\
 & & \texttt{  } \\
 & & \texttt{  } \\
 & & \texttt{  } \\
 & & \texttt{  } \\
 & & \texttt{  } \\
\hline \end{tabular}
}
\newpage
{\scriptsize
\begin{tabular}{|l|l|l|} \hline
\bf{Register Class} & \bf{Register File} & \bf{Registers} \\ \hline
 xworddReg & XCR & \texttt{1 } \\
 & & \texttt{1 } \\
 & & \texttt{1 } \\
 & & \texttt{1 } \\
 & & \texttt{1 } \\
 & & \texttt{1 } \\
 & & \texttt{1 } \\
 & & \texttt{1 } \\
 & & \texttt{1 } \\
 & & \texttt{1 } \\
 & & \texttt{1 } \\
 & & \texttt{1 } \\
 & & \texttt{1 } \\
 & & \texttt{1 } \\
 & & \texttt{1 } \\
 & & \texttt{1 } \\
 & & \texttt{1 } \\
 & & \texttt{1 } \\
 & & \texttt{1 } \\
 & & \texttt{1 } \\
 & & \texttt{1 } \\
 & & \texttt{1 } \\
 & & \texttt{1 } \\
 & & \texttt{1 } \\
 & & \texttt{1 } \\
 & & \texttt{1 } \\
 & & \texttt{1 } \\
 & & \texttt{1 } \\
 & & \texttt{1 } \\
 & & \texttt{1 } \\
 & & \texttt{1 } \\
 & & \texttt{1 } \\
 & & \texttt{1 } \\
 & & \texttt{1 } \\
 & & \texttt{1 } \\
 & & \texttt{1 } \\
 & & \texttt{1 } \\
 & & \texttt{1 } \\
 & & \texttt{1 } \\
 & & \texttt{1 } \\
 & & \texttt{1 } \\
 & & \texttt{1 } \\
 & & \texttt{1 } \\
 & & \texttt{1 } \\
 & & \texttt{1 } \\
 & & \texttt{1 } \\
 & & \texttt{1 } \\
 & & \texttt{1 } \\
 & & \texttt{1 } \\
 & & \texttt{1 } \\
 & & \texttt{1 } \\
 & & \texttt{1 } \\
 & & \texttt{1 } \\
 & & \texttt{1 } \\
 & & \texttt{1 } \\
 & & \texttt{1 } \\
 & & \texttt{1 } \\
 & & \texttt{1 } \\
 & & \texttt{1 } \\
 & & \texttt{1 } \\
 & & \texttt{1 } \\
 & & \texttt{1 } \\
 & & \texttt{1 } \\
 & & \texttt{1 } \\
 & & \texttt{1 } \\
 & & \texttt{1 } \\
 & & \texttt{1 } \\
 & & \texttt{1 } \\
 & & \texttt{1 } \\
 & & \texttt{1 } \\
 & & \texttt{1 } \\
 & & \texttt{1 } \\
 & & \texttt{1 } \\
 & & \texttt{1 } \\
 & & \texttt{1 } \\
 & & \texttt{1 } \\
 & & \texttt{1 } \\
 & & \texttt{1 } \\
 & & \texttt{1 } \\
 & & \texttt{1 } \\
 & & \texttt{1 } \\
 & & \texttt{1 } \\
 & & \texttt{1 } \\
 & & \texttt{1 } \\
 & & \texttt{1 } \\
 & & \texttt{1 } \\
 & & \texttt{1 } \\
 & & \texttt{1 } \\
 & & \texttt{1 } \\
 & & \texttt{1 } \\
 & & \texttt{1 } \\
 & & \texttt{1 } \\
 & & \texttt{1 } \\
 & & \texttt{1 } \\
 & & \texttt{1 } \\
 & & \texttt{1 } \\
 & & \texttt{1 } \\
 & & \texttt{1 } \\
 & & \texttt{1 } \\
 & & \texttt{1 } \\
 & & \texttt{1 } \\
 & & \texttt{1 } \\
 & & \texttt{1 } \\
 & & \texttt{1 } \\
 & & \texttt{1 } \\
 & & \texttt{1 } \\
 & & \texttt{1 } \\
 & & \texttt{1 } \\
 & & \texttt{1 } \\
 & & \texttt{1 } \\
 & & \texttt{1 } \\
 & & \texttt{1 } \\
 & & \texttt{1 } \\
 & & \texttt{1 } \\
 & & \texttt{1 } \\
 & & \texttt{1 } \\
 & & \texttt{1 } \\
 & & \texttt{1 } \\
 & & \texttt{1 } \\
 & & \texttt{1 } \\
 & & \texttt{1 } \\
 & & \texttt{1 } \\
 & & \texttt{1 } \\
 & & \texttt{1 } \\
 & & \texttt{1 } \\
 & & \texttt{1 } \\
 & & \texttt{1 } \\
 & & \texttt{1 } \\
 & & \texttt{1 } \\
 & & \texttt{1 } \\
 & & \texttt{1 } \\
 & & \texttt{1 } \\
 & & \texttt{1 } \\
 & & \texttt{1 } \\
 & & \texttt{1 } \\
 & & \texttt{1 } \\
 & & \texttt{1 } \\
 & & \texttt{1 } \\
 & & \texttt{1 } \\
 & & \texttt{1 } \\
 & & \texttt{1 } \\
 & & \texttt{1 } \\
 & & \texttt{1 } \\
 & & \texttt{1 } \\
 & & \texttt{1 } \\
 & & \texttt{1 } \\
 & & \texttt{1 } \\
 & & \texttt{1 } \\
 & & \texttt{1 } \\
 & & \texttt{1 } \\
 & & \texttt{1 } \\
 & & \texttt{1 } \\
 & & \texttt{1 } \\
 & & \texttt{1 } \\
 & & \texttt{1 } \\
 & & \texttt{1 } \\
 & & \texttt{1 } \\
 & & \texttt{1 } \\
 & & \texttt{1 } \\
 & & \texttt{1 } \\
 & & \texttt{1 } \\
 & & \texttt{1 } \\
 & & \texttt{1 } \\
 & & \texttt{1 } \\
 & & \texttt{1 } \\
 & & \texttt{1 } \\
 & & \texttt{1 } \\
 & & \texttt{1 } \\
 & & \texttt{1 } \\
 & & \texttt{1 } \\
 & & \texttt{1 } \\
 & & \texttt{1 } \\
 & & \texttt{1 } \\
 & & \texttt{1 } \\
 & & \texttt{1 } \\
 & & \texttt{1 } \\
 & & \texttt{1 } \\
 & & \texttt{1 } \\
 & & \texttt{1 } \\
 & & \texttt{1 } \\
 & & \texttt{1 } \\
 & & \texttt{1 } \\
 & & \texttt{1 } \\
 & & \texttt{1 } \\
 & & \texttt{1 } \\
 & & \texttt{1 } \\
 & & \texttt{1 } \\
 & & \texttt{1 } \\
 & & \texttt{1 } \\
 & & \texttt{1 } \\
 & & \texttt{1 } \\
 & & \texttt{1 } \\
 & & \texttt{1 } \\
 & & \texttt{1 } \\
 & & \texttt{1 } \\
 & & \texttt{1 } \\
 & & \texttt{1 } \\
 & & \texttt{1 } \\
 & & \texttt{1 } \\
 & & \texttt{1 } \\
 & & \texttt{1 } \\
 & & \texttt{1 } \\
 & & \texttt{1 } \\
 & & \texttt{1 } \\
 & & \texttt{1 } \\
 & & \texttt{1 } \\
 & & \texttt{1 } \\
 & & \texttt{1 } \\
 & & \texttt{1 } \\
 & & \texttt{1 } \\
 & & \texttt{1 } \\
 & & \texttt{1 } \\
 & & \texttt{1 } \\
 & & \texttt{1 } \\
 & & \texttt{1 } \\
 & & \texttt{1 } \\
 & & \texttt{1 } \\
 & & \texttt{1 } \\
 & & \texttt{1 } \\
 & & \texttt{1 } \\
 & & \texttt{1 } \\
 & & \texttt{1 } \\
 & & \texttt{1 } \\
 & & \texttt{1 } \\
 & & \texttt{1 } \\
 & & \texttt{1 } \\
 & & \texttt{1 } \\
 & & \texttt{1 } \\
 & & \texttt{1 } \\
 & & \texttt{1 } \\
 & & \texttt{1 } \\
 & & \texttt{1 } \\
 & & \texttt{1 } \\
 & & \texttt{1 } \\
 & & \texttt{1 } \\
 & & \texttt{1 } \\
 & & \texttt{1 } \\
 & & \texttt{1 } \\
 & & \texttt{1 } \\
 & & \texttt{1 } \\
 & & \texttt{1 } \\
 & & \texttt{1 } \\
 & & \texttt{1 } \\
 & & \texttt{1 } \\
 & & \texttt{1 } \\
 & & \texttt{1 } \\
 & & \texttt{1 } \\
 & & \texttt{1 } \\
 & & \texttt{1 } \\
 & & \texttt{1 } \\
 & & \texttt{1 } \\
 & & \texttt{1 } \\
 & & \texttt{1 } \\
 & & \texttt{1 } \\
 & & \texttt{1 } \\
 & & \texttt{1 } \\
\hline \end{tabular}
}
\newpage
{\scriptsize
\begin{tabular}{|l|l|l|} \hline
\bf{Register Class} & \bf{Register File} & \bf{Registers} \\ \hline
 xworddReg0M4 & XCR & \texttt{1 } \\
 & & \texttt{1 } \\
 & & \texttt{1 } \\
 & & \texttt{1 } \\
 & & \texttt{1 } \\
 & & \texttt{1 } \\
 & & \texttt{1 } \\
 & & \texttt{1 } \\
 & & \texttt{1 } \\
 & & \texttt{1 } \\
 & & \texttt{1 } \\
 & & \texttt{1 } \\
 & & \texttt{1 } \\
 & & \texttt{1 } \\
 & & \texttt{1 } \\
 & & \texttt{1 } \\
 & & \texttt{1 } \\
 & & \texttt{1 } \\
 & & \texttt{1 } \\
 & & \texttt{1 } \\
 & & \texttt{1 } \\
 & & \texttt{1 } \\
 & & \texttt{1 } \\
 & & \texttt{1 } \\
 & & \texttt{1 } \\
 & & \texttt{1 } \\
 & & \texttt{1 } \\
 & & \texttt{1 } \\
 & & \texttt{1 } \\
 & & \texttt{1 } \\
 & & \texttt{1 } \\
 & & \texttt{1 } \\
 & & \texttt{1 } \\
 & & \texttt{1 } \\
 & & \texttt{1 } \\
 & & \texttt{1 } \\
 & & \texttt{1 } \\
 & & \texttt{1 } \\
 & & \texttt{1 } \\
 & & \texttt{1 } \\
 & & \texttt{1 } \\
 & & \texttt{1 } \\
 & & \texttt{1 } \\
 & & \texttt{1 } \\
 & & \texttt{1 } \\
 & & \texttt{1 } \\
 & & \texttt{1 } \\
 & & \texttt{1 } \\
 & & \texttt{1 } \\
 & & \texttt{1 } \\
 & & \texttt{1 } \\
 & & \texttt{1 } \\
 & & \texttt{1 } \\
 & & \texttt{1 } \\
 & & \texttt{1 } \\
 & & \texttt{1 } \\
 & & \texttt{1 } \\
 & & \texttt{1 } \\
 & & \texttt{1 } \\
 & & \texttt{1 } \\
 & & \texttt{1 } \\
 & & \texttt{1 } \\
 & & \texttt{1 } \\
 & & \texttt{1 } \\
\hline \end{tabular}
}
\newpage
{\scriptsize
\begin{tabular}{|l|l|l|} \hline
\bf{Register Class} & \bf{Register File} & \bf{Registers} \\ \hline
 xworddReg1M4 & XCR & \texttt{1 } \\
 & & \texttt{1 } \\
 & & \texttt{1 } \\
 & & \texttt{1 } \\
 & & \texttt{1 } \\
 & & \texttt{1 } \\
 & & \texttt{1 } \\
 & & \texttt{1 } \\
 & & \texttt{1 } \\
 & & \texttt{1 } \\
 & & \texttt{1 } \\
 & & \texttt{1 } \\
 & & \texttt{1 } \\
 & & \texttt{1 } \\
 & & \texttt{1 } \\
 & & \texttt{1 } \\
 & & \texttt{1 } \\
 & & \texttt{1 } \\
 & & \texttt{1 } \\
 & & \texttt{1 } \\
 & & \texttt{1 } \\
 & & \texttt{1 } \\
 & & \texttt{1 } \\
 & & \texttt{1 } \\
 & & \texttt{1 } \\
 & & \texttt{1 } \\
 & & \texttt{1 } \\
 & & \texttt{1 } \\
 & & \texttt{1 } \\
 & & \texttt{1 } \\
 & & \texttt{1 } \\
 & & \texttt{1 } \\
 & & \texttt{1 } \\
 & & \texttt{1 } \\
 & & \texttt{1 } \\
 & & \texttt{1 } \\
 & & \texttt{1 } \\
 & & \texttt{1 } \\
 & & \texttt{1 } \\
 & & \texttt{1 } \\
 & & \texttt{1 } \\
 & & \texttt{1 } \\
 & & \texttt{1 } \\
 & & \texttt{1 } \\
 & & \texttt{1 } \\
 & & \texttt{1 } \\
 & & \texttt{1 } \\
 & & \texttt{1 } \\
 & & \texttt{1 } \\
 & & \texttt{1 } \\
 & & \texttt{1 } \\
 & & \texttt{1 } \\
 & & \texttt{1 } \\
 & & \texttt{1 } \\
 & & \texttt{1 } \\
 & & \texttt{1 } \\
 & & \texttt{1 } \\
 & & \texttt{1 } \\
 & & \texttt{1 } \\
 & & \texttt{1 } \\
 & & \texttt{1 } \\
 & & \texttt{1 } \\
 & & \texttt{1 } \\
 & & \texttt{1 } \\
\hline \end{tabular}
}
\newpage
{\scriptsize
\begin{tabular}{|l|l|l|} \hline
\bf{Register Class} & \bf{Register File} & \bf{Registers} \\ \hline
 xworddReg2M4 & XCR & \texttt{1 } \\
 & & \texttt{1 } \\
 & & \texttt{1 } \\
 & & \texttt{1 } \\
 & & \texttt{1 } \\
 & & \texttt{1 } \\
 & & \texttt{1 } \\
 & & \texttt{1 } \\
 & & \texttt{1 } \\
 & & \texttt{1 } \\
 & & \texttt{1 } \\
 & & \texttt{1 } \\
 & & \texttt{1 } \\
 & & \texttt{1 } \\
 & & \texttt{1 } \\
 & & \texttt{1 } \\
 & & \texttt{1 } \\
 & & \texttt{1 } \\
 & & \texttt{1 } \\
 & & \texttt{1 } \\
 & & \texttt{1 } \\
 & & \texttt{1 } \\
 & & \texttt{1 } \\
 & & \texttt{1 } \\
 & & \texttt{1 } \\
 & & \texttt{1 } \\
 & & \texttt{1 } \\
 & & \texttt{1 } \\
 & & \texttt{1 } \\
 & & \texttt{1 } \\
 & & \texttt{1 } \\
 & & \texttt{1 } \\
 & & \texttt{1 } \\
 & & \texttt{1 } \\
 & & \texttt{1 } \\
 & & \texttt{1 } \\
 & & \texttt{1 } \\
 & & \texttt{1 } \\
 & & \texttt{1 } \\
 & & \texttt{1 } \\
 & & \texttt{1 } \\
 & & \texttt{1 } \\
 & & \texttt{1 } \\
 & & \texttt{1 } \\
 & & \texttt{1 } \\
 & & \texttt{1 } \\
 & & \texttt{1 } \\
 & & \texttt{1 } \\
 & & \texttt{1 } \\
 & & \texttt{1 } \\
 & & \texttt{1 } \\
 & & \texttt{1 } \\
 & & \texttt{1 } \\
 & & \texttt{1 } \\
 & & \texttt{1 } \\
 & & \texttt{1 } \\
 & & \texttt{1 } \\
 & & \texttt{1 } \\
 & & \texttt{1 } \\
 & & \texttt{1 } \\
 & & \texttt{1 } \\
 & & \texttt{1 } \\
 & & \texttt{1 } \\
 & & \texttt{1 } \\
\hline \end{tabular}
}
\newpage
{\scriptsize
\begin{tabular}{|l|l|l|} \hline
\bf{Register Class} & \bf{Register File} & \bf{Registers} \\ \hline
 xworddReg3M4 & XCR & \texttt{1 } \\
 & & \texttt{1 } \\
 & & \texttt{1 } \\
 & & \texttt{1 } \\
 & & \texttt{1 } \\
 & & \texttt{1 } \\
 & & \texttt{1 } \\
 & & \texttt{1 } \\
 & & \texttt{1 } \\
 & & \texttt{1 } \\
 & & \texttt{1 } \\
 & & \texttt{1 } \\
 & & \texttt{1 } \\
 & & \texttt{1 } \\
 & & \texttt{1 } \\
 & & \texttt{1 } \\
 & & \texttt{1 } \\
 & & \texttt{1 } \\
 & & \texttt{1 } \\
 & & \texttt{1 } \\
 & & \texttt{1 } \\
 & & \texttt{1 } \\
 & & \texttt{1 } \\
 & & \texttt{1 } \\
 & & \texttt{1 } \\
 & & \texttt{1 } \\
 & & \texttt{1 } \\
 & & \texttt{1 } \\
 & & \texttt{1 } \\
 & & \texttt{1 } \\
 & & \texttt{1 } \\
 & & \texttt{1 } \\
 & & \texttt{1 } \\
 & & \texttt{1 } \\
 & & \texttt{1 } \\
 & & \texttt{1 } \\
 & & \texttt{1 } \\
 & & \texttt{1 } \\
 & & \texttt{1 } \\
 & & \texttt{1 } \\
 & & \texttt{1 } \\
 & & \texttt{1 } \\
 & & \texttt{1 } \\
 & & \texttt{1 } \\
 & & \texttt{1 } \\
 & & \texttt{1 } \\
 & & \texttt{1 } \\
 & & \texttt{1 } \\
 & & \texttt{1 } \\
 & & \texttt{1 } \\
 & & \texttt{1 } \\
 & & \texttt{1 } \\
 & & \texttt{1 } \\
 & & \texttt{1 } \\
 & & \texttt{1 } \\
 & & \texttt{1 } \\
 & & \texttt{1 } \\
 & & \texttt{1 } \\
 & & \texttt{1 } \\
 & & \texttt{1 } \\
 & & \texttt{1 } \\
 & & \texttt{1 } \\
 & & \texttt{1 } \\
 & & \texttt{1 } \\
\hline \end{tabular}
}
\newpage
{\scriptsize
\begin{tabular}{|l|l|l|} \hline
\bf{Register Class} & \bf{Register File} & \bf{Registers} \\ \hline
 xwordoReg & XVR & \texttt{  } \\
 & & \texttt{  } \\
 & & \texttt{  } \\
 & & \texttt{  } \\
 & & \texttt{  } \\
 & & \texttt{  } \\
 & & \texttt{  } \\
 & & \texttt{  } \\
 & & \texttt{  } \\
 & & \texttt{  } \\
 & & \texttt{  } \\
 & & \texttt{  } \\
 & & \texttt{  } \\
 & & \texttt{  } \\
 & & \texttt{  } \\
 & & \texttt{  } \\
 & & \texttt{  } \\
 & & \texttt{  } \\
 & & \texttt{  } \\
 & & \texttt{  } \\
 & & \texttt{  } \\
 & & \texttt{  } \\
 & & \texttt{  } \\
 & & \texttt{  } \\
 & & \texttt{  } \\
 & & \texttt{  } \\
 & & \texttt{  } \\
 & & \texttt{  } \\
 & & \texttt{  } \\
 & & \texttt{  } \\
 & & \texttt{  } \\
 & & \texttt{  } \\
 & & \texttt{  } \\
 & & \texttt{  } \\
 & & \texttt{  } \\
 & & \texttt{  } \\
 & & \texttt{  } \\
 & & \texttt{  } \\
 & & \texttt{  } \\
 & & \texttt{  } \\
 & & \texttt{  } \\
 & & \texttt{  } \\
 & & \texttt{  } \\
 & & \texttt{  } \\
 & & \texttt{  } \\
 & & \texttt{  } \\
 & & \texttt{  } \\
 & & \texttt{  } \\
 & & \texttt{  } \\
 & & \texttt{  } \\
 & & \texttt{  } \\
 & & \texttt{  } \\
 & & \texttt{  } \\
 & & \texttt{  } \\
 & & \texttt{  } \\
 & & \texttt{  } \\
 & & \texttt{  } \\
 & & \texttt{  } \\
 & & \texttt{  } \\
 & & \texttt{  } \\
 & & \texttt{  } \\
 & & \texttt{  } \\
 & & \texttt{  } \\
 & & \texttt{  } \\
\hline \end{tabular}
}
\newpage
{\scriptsize
\begin{tabular}{|l|l|l|} \hline
\bf{Register Class} & \bf{Register File} & \bf{Registers} \\ \hline
 xwordqReg & XBR & \texttt{1 } \\
 & & \texttt{1 } \\
 & & \texttt{1 } \\
 & & \texttt{1 } \\
 & & \texttt{1 } \\
 & & \texttt{1 } \\
 & & \texttt{1 } \\
 & & \texttt{1 } \\
 & & \texttt{1 } \\
 & & \texttt{1 } \\
 & & \texttt{1 } \\
 & & \texttt{1 } \\
 & & \texttt{1 } \\
 & & \texttt{1 } \\
 & & \texttt{1 } \\
 & & \texttt{1 } \\
 & & \texttt{1 } \\
 & & \texttt{1 } \\
 & & \texttt{1 } \\
 & & \texttt{1 } \\
 & & \texttt{1 } \\
 & & \texttt{1 } \\
 & & \texttt{1 } \\
 & & \texttt{1 } \\
 & & \texttt{1 } \\
 & & \texttt{1 } \\
 & & \texttt{1 } \\
 & & \texttt{1 } \\
 & & \texttt{1 } \\
 & & \texttt{1 } \\
 & & \texttt{1 } \\
 & & \texttt{1 } \\
 & & \texttt{1 } \\
 & & \texttt{1 } \\
 & & \texttt{1 } \\
 & & \texttt{1 } \\
 & & \texttt{1 } \\
 & & \texttt{1 } \\
 & & \texttt{1 } \\
 & & \texttt{1 } \\
 & & \texttt{1 } \\
 & & \texttt{1 } \\
 & & \texttt{1 } \\
 & & \texttt{1 } \\
 & & \texttt{1 } \\
 & & \texttt{1 } \\
 & & \texttt{1 } \\
 & & \texttt{1 } \\
 & & \texttt{1 } \\
 & & \texttt{1 } \\
 & & \texttt{1 } \\
 & & \texttt{1 } \\
 & & \texttt{1 } \\
 & & \texttt{1 } \\
 & & \texttt{1 } \\
 & & \texttt{1 } \\
 & & \texttt{1 } \\
 & & \texttt{1 } \\
 & & \texttt{1 } \\
 & & \texttt{1 } \\
 & & \texttt{1 } \\
 & & \texttt{1 } \\
 & & \texttt{1 } \\
 & & \texttt{1 } \\
 & & \texttt{1 } \\
 & & \texttt{1 } \\
 & & \texttt{1 } \\
 & & \texttt{1 } \\
 & & \texttt{1 } \\
 & & \texttt{1 } \\
 & & \texttt{1 } \\
 & & \texttt{1 } \\
 & & \texttt{1 } \\
 & & \texttt{1 } \\
 & & \texttt{1 } \\
 & & \texttt{1 } \\
 & & \texttt{1 } \\
 & & \texttt{1 } \\
 & & \texttt{1 } \\
 & & \texttt{1 } \\
 & & \texttt{1 } \\
 & & \texttt{1 } \\
 & & \texttt{1 } \\
 & & \texttt{1 } \\
 & & \texttt{1 } \\
 & & \texttt{1 } \\
 & & \texttt{1 } \\
 & & \texttt{1 } \\
 & & \texttt{1 } \\
 & & \texttt{1 } \\
 & & \texttt{1 } \\
 & & \texttt{1 } \\
 & & \texttt{1 } \\
 & & \texttt{1 } \\
 & & \texttt{1 } \\
 & & \texttt{1 } \\
 & & \texttt{1 } \\
 & & \texttt{1 } \\
 & & \texttt{1 } \\
 & & \texttt{1 } \\
 & & \texttt{1 } \\
 & & \texttt{1 } \\
 & & \texttt{1 } \\
 & & \texttt{1 } \\
 & & \texttt{1 } \\
 & & \texttt{1 } \\
 & & \texttt{1 } \\
 & & \texttt{1 } \\
 & & \texttt{1 } \\
 & & \texttt{1 } \\
 & & \texttt{1 } \\
 & & \texttt{1 } \\
 & & \texttt{1 } \\
 & & \texttt{1 } \\
 & & \texttt{1 } \\
 & & \texttt{1 } \\
 & & \texttt{1 } \\
 & & \texttt{1 } \\
 & & \texttt{1 } \\
 & & \texttt{1 } \\
 & & \texttt{1 } \\
 & & \texttt{1 } \\
 & & \texttt{1 } \\
 & & \texttt{1 } \\
 & & \texttt{1 } \\
 & & \texttt{1 } \\
 & & \texttt{1 } \\
 & & \texttt{1 } \\
\hline \end{tabular}
}
\newpage
{\scriptsize
\begin{tabular}{|l|l|l|} \hline
\bf{Register Class} & \bf{Register File} & \bf{Registers} \\ \hline
 xwordqRegE & XBR & \texttt{1 } \\
 & & \texttt{1 } \\
 & & \texttt{1 } \\
 & & \texttt{1 } \\
 & & \texttt{1 } \\
 & & \texttt{1 } \\
 & & \texttt{1 } \\
 & & \texttt{1 } \\
 & & \texttt{1 } \\
 & & \texttt{1 } \\
 & & \texttt{1 } \\
 & & \texttt{1 } \\
 & & \texttt{1 } \\
 & & \texttt{1 } \\
 & & \texttt{1 } \\
 & & \texttt{1 } \\
 & & \texttt{1 } \\
 & & \texttt{1 } \\
 & & \texttt{1 } \\
 & & \texttt{1 } \\
 & & \texttt{1 } \\
 & & \texttt{1 } \\
 & & \texttt{1 } \\
 & & \texttt{1 } \\
 & & \texttt{1 } \\
 & & \texttt{1 } \\
 & & \texttt{1 } \\
 & & \texttt{1 } \\
 & & \texttt{1 } \\
 & & \texttt{1 } \\
 & & \texttt{1 } \\
 & & \texttt{1 } \\
 & & \texttt{1 } \\
 & & \texttt{1 } \\
 & & \texttt{1 } \\
 & & \texttt{1 } \\
 & & \texttt{1 } \\
 & & \texttt{1 } \\
 & & \texttt{1 } \\
 & & \texttt{1 } \\
 & & \texttt{1 } \\
 & & \texttt{1 } \\
 & & \texttt{1 } \\
 & & \texttt{1 } \\
 & & \texttt{1 } \\
 & & \texttt{1 } \\
 & & \texttt{1 } \\
 & & \texttt{1 } \\
 & & \texttt{1 } \\
 & & \texttt{1 } \\
 & & \texttt{1 } \\
 & & \texttt{1 } \\
 & & \texttt{1 } \\
 & & \texttt{1 } \\
 & & \texttt{1 } \\
 & & \texttt{1 } \\
 & & \texttt{1 } \\
 & & \texttt{1 } \\
 & & \texttt{1 } \\
 & & \texttt{1 } \\
 & & \texttt{1 } \\
 & & \texttt{1 } \\
 & & \texttt{1 } \\
 & & \texttt{1 } \\
\hline \end{tabular}
}
\newpage
{\scriptsize
\begin{tabular}{|l|l|l|} \hline
\bf{Register Class} & \bf{Register File} & \bf{Registers} \\ \hline
 xwordqRegO & XBR & \texttt{1 } \\
 & & \texttt{1 } \\
 & & \texttt{1 } \\
 & & \texttt{1 } \\
 & & \texttt{1 } \\
 & & \texttt{1 } \\
 & & \texttt{1 } \\
 & & \texttt{1 } \\
 & & \texttt{1 } \\
 & & \texttt{1 } \\
 & & \texttt{1 } \\
 & & \texttt{1 } \\
 & & \texttt{1 } \\
 & & \texttt{1 } \\
 & & \texttt{1 } \\
 & & \texttt{1 } \\
 & & \texttt{1 } \\
 & & \texttt{1 } \\
 & & \texttt{1 } \\
 & & \texttt{1 } \\
 & & \texttt{1 } \\
 & & \texttt{1 } \\
 & & \texttt{1 } \\
 & & \texttt{1 } \\
 & & \texttt{1 } \\
 & & \texttt{1 } \\
 & & \texttt{1 } \\
 & & \texttt{1 } \\
 & & \texttt{1 } \\
 & & \texttt{1 } \\
 & & \texttt{1 } \\
 & & \texttt{1 } \\
 & & \texttt{1 } \\
 & & \texttt{1 } \\
 & & \texttt{1 } \\
 & & \texttt{1 } \\
 & & \texttt{1 } \\
 & & \texttt{1 } \\
 & & \texttt{1 } \\
 & & \texttt{1 } \\
 & & \texttt{1 } \\
 & & \texttt{1 } \\
 & & \texttt{1 } \\
 & & \texttt{1 } \\
 & & \texttt{1 } \\
 & & \texttt{1 } \\
 & & \texttt{1 } \\
 & & \texttt{1 } \\
 & & \texttt{1 } \\
 & & \texttt{1 } \\
 & & \texttt{1 } \\
 & & \texttt{1 } \\
 & & \texttt{1 } \\
 & & \texttt{1 } \\
 & & \texttt{1 } \\
 & & \texttt{1 } \\
 & & \texttt{1 } \\
 & & \texttt{1 } \\
 & & \texttt{1 } \\
 & & \texttt{1 } \\
 & & \texttt{1 } \\
 & & \texttt{1 } \\
 & & \texttt{1 } \\
 & & \texttt{1 } \\
\hline \end{tabular}
}
\newpage
{\scriptsize
\begin{tabular}{|l|l|l|} \hline
\bf{Register Class} & \bf{Register File} & \bf{Registers} \\ \hline
 xwordvReg & XMR & \texttt{} \\
 & & \texttt{} \\
 & & \texttt{} \\
 & & \texttt{} \\
 & & \texttt{} \\
 & & \texttt{} \\
 & & \texttt{} \\
 & & \texttt{} \\
 & & \texttt{} \\
 & & \texttt{} \\
 & & \texttt{} \\
 & & \texttt{} \\
 & & \texttt{} \\
 & & \texttt{} \\
 & & \texttt{} \\
 & & \texttt{} \\
\hline \end{tabular}
}
\newpage
